\raggedbottom
\frontmatter
\title{Lean Formalization of the Geometrization Conjecture}
\author{Project blueprint}
\date{Revision 204: Actual history interfaces and conditional descent, September 2026}
\subjclass[2020]{Primary 53E20; Secondary 57K30}
\keywords{geometrization conjecture, Ricci flow, surgery, three-manifolds, Lean}
\maketitle
\begingroup
\small
\tableofcontents
\endgroup

\chapter*{Purpose and status}

This is revision 204 of the blueprint for a kernel-checked Lean
formalization of the geometrization conjecture using Ricci flow with
surgery. It specifies theorem contracts, dependencies, source roles,
milestones and unresolved obligations. Written mathematical arguments
support the specifications; checked formal proofs remain separate.

% BEGIN REVISION204 FRONT STATUS
\paragraph{Actual history interfaces and conditional descent.}
Revision204 connects the endpoint to the actual PC observed histories and their
initial markings. An explicit oriented refinement of each event reconstruction,
certificates on every discarded component and on each required sphere-product
factor, and final-stage component certificates imply the initial endpoint by
backward induction. The terminal branch has no late-geometry premise; the
zero-event sphere example is proved. General production of those event inputs
and late certificates remains open. Protected ports, the prescribed cycle
ledger, a flattened decorated history and one controlled global profile are
not inferred from this bounded composition.
% END REVISION204 FRONT STATUS

% BEGIN REVISION203 FRONT STATUS
\paragraph{Current scope: reviewed Lean interfaces and proof consumers.}
The user has reauthorized coding and prioritized a bounded interface-building
phase around the existing DAG. This supersedes the planning-only instruction
recorded historically below. The public geometrization proposition now has
concrete Lean definitions for prime factors, torus cuts, whole-interior geometry
and oriented reconstruction; the universal theorem is not yet proved.
Revision203 supplies actual smooth homogeneous complete metrics for the five
coordinate models, and conditional finite-family assembly into that endpoint.
The next interfaces are actual late-slice production, marked finite-history
reconstruction, the terminal alternative and one controlled global flow profile.
Parent contracts and the complete premise census remain unaccepted.
% END REVISION203 FRONT STATUS

% BEGIN REVISION202 FRONT STATUS
\paragraph{Current scope: blueprint planning and verification only.}
The user has stopped Lean development. Revision202 consolidates the plan and
its evidence boundary; it introduces no new Lean code or mathematical theorem.
Revision201 is the frozen implementation baseline, pushed to the new branch
\code{codex/geometrization-baseline-201} of
\code{qinz1yang/differential-geometry-dev}. The exact remote commit and unchanged
source hashes are recorded in the planning handoff. Earlier implementation
agendas and status notices below are historical, not current work instructions.
The Geometrization endpoint and complete semantic dependency census remain open.
% END REVISION202 FRONT STATUS

% BEGIN REVISION201 FRONT STATUS
\paragraph{One coherent general surgery tower.}
Revision201 concatenates actual finite retained histories while preserving the
old prefix and initial marking. Recursive extension produces one compatible
tower for arbitrary compact oriented initial data, with locally finite surgery
times and actual event geometry. Its late-nonempty or absorbing-empty alternative
is proved. A single globally controlled cutoff profile, the required long-time
estimates, and full marked reconstruction remain separate obligations.
% END REVISION201 FRONT STATUS

% BEGIN REVISION200 FRONT STATUS
\paragraph{Actual continuation compatibility.}
Revision200 derives agreement with the existing smooth tail from compact Ricci
flow uniqueness, excludes singular events before the current horizon, and
preserves the given initial marking. Time translation transports the actual
terminal metric and retained-core surgery event. Finite-history concatenation
and one controlled global tower remain the next tasks. Earlier notices retain
their historical scope.
% END REVISION200 FRONT STATUS

% BEGIN REVISION199 FRONT STATUS
\paragraph{The general initial bound and finite-horizon analytic producers.}
Revision199 proves the uniform initial finite-free-factor order bound from the
actual finite decompositions. The bound now applies to every retained surgery
history in the unchanged cutoff class. Actual consumers produce general
small-scale noncollapsing, strong canonical neighborhoods and finite geometric
histories to every prescribed positive horizon, without simple connectivity or
supplied completion. Coherence into one global flow and marked reconstruction
remain open. Earlier notices record their historical revision scopes.
% END REVISION199 FRONT STATUS

% BEGIN REVISION198 FRONT STATUS
\paragraph{Actual finite decompositions of the initial groups.}
Revision198 binds a pinned Grushko rank-additivity proof to the accepted Lean
foundation. Strong induction on the minimal number of generators now produces
an actual finite indexed free product of nontrivial, finitely generated,
freely indecomposable groups for every initial component group. The trivial
group uses the empty family. The further comparison needed to bound all finite
free factors remains open; decomposition existence alone is not that bound.
% END REVISION198 FRONT STATUS

% BEGIN REVISION197 FRONT STATUS
\paragraph{A checked Kurosh supplier and its actual finite-subgroup consumer.}
Revision197 binds a pinned public Kurosh formalization to the accepted foundation,
with a recorded compatibility patch. Every finite nontrivial subgroup of an
indexed free product is proved to lie in an actual conjugate of an original
factor. An injective homomorphism into that factor and its conjugation equation
are produced. The general initial-group decomposition and order bound remain
open; neither follows merely by naming Kurosh or Grushko.
% END REVISION197 FRONT STATUS

% BEGIN REVISION196 FRONT STATUS
\paragraph{Finite generation of actual initial groups.}
Revision196 proves finite generation for every fundamental group of an actual
initial stage, using compactness and its real-model charts. A finite cover,
its actual connector identities and ordinary subgroup generation are all
produced. Arbitrary carrier universes and disconnected stages are included.
Every initial free-factor witness also has a finitely generated complement.
The remaining initial-order-bound question still requires the Grushko comparison;
finite generation alone is not counted as that bound.
% END REVISION196 FRONT STATUS

% BEGIN REVISION195 FRONT STATUS
\paragraph{Actual finite cuts and history ancestry.}
Revision195 proves the finite-collar induction and applies it to the actual
surgery tube systems, child cores and parent components. Every retained history
now carries actual fundamental groups as free factors of an identified initial
group, without a capping-completion assumption. The exact cutoff-class degree
and strong canonical consumers retain only the initial finite-free-factor bound.
That bound is produced when all initial component groups are finite. The general
initial bound, coherent flow and marked reconstruction remain open.
% END REVISION195 FRONT STATUS

% BEGIN REVISION194 FRONT STATUS
\paragraph{Every component after one collared cut.}
Revision194 supplies the separating case and combines it with revision193.
For one compact simply connected two-sided collared hypersurface in a connected,
Hausdorff, locally path-connected ambient space, the group of every actual
complement component is a free factor of the ambient group. Both cases are
proved internally from the two component labels; ambient simple connectivity
is not assumed. Actual finite surgery families and their history transport
remain to be produced before the all-history completion gate can be removed.
% END REVISION194 FRONT STATUS

% BEGIN REVISION193 FRONT STATUS
\paragraph{An actual free complement for disconnected overlap.}
Revision193 proves the open-cover free-product formula with one free generator
per overlap path component other than the base component. Both inverse maps
and agreement with the actual retained-region inclusion are checked. The
actual two-sided collar geometry supplies the hypotheses when its simply
connected compact hypersurface has path-connected complement. Separating cuts,
finite surgery families and general history ancestry remain to be assembled;
this result does not remove the all-history completion premise by itself.
% END REVISION193 FRONT STATUS

% BEGIN REVISION192 FRONT STATUS
\paragraph{Finite free-factor prerequisites and a genuine initial-bound case.}
Revision192 proves the general reduced-word infinitude needed for finite-factor
indecomposability, and proves that a free complement of a finitely generated
group is finitely generated. It also produces the actual uniform initial bound
when every initial component group is finite, then applies the unchanged
history-degree and strong-canonical consumers. Completion for every cutoff
history remains supplied. The general initial-group case still requires the
Grushko/free-factor argument; neither that case nor the whole program is closed.
% END REVISION192 FRONT STATUS

% BEGIN REVISION191 FRONT STATUS
\paragraph{Actual cap attachments preserve the component group.}
Revision191 proves that the actual child component, its core, and the punctured
component selected by the transition's core retraction have isomorphic
fundamental groups. All cap-cover hypotheses are produced from the actual
transition, without smooth completion or simple connectivity of the parent.
The parent-side free-factor theorem and the initial finite-free-factor bound
remain open; this partial direct route does not yet discharge revision190's
all-history completion premise.
% END REVISION191 FRONT STATUS

% BEGIN REVISION190 FRONT STATUS
\paragraph{Actual free factors now connect reconstruction to history degree.}
Revision190 extracts actual free-product complements from the inherited finite
connected-sum theorem and revision177's cut/cap reconstruction. These give
retained-component ancestry through completed histories, with exact component
group transport into revision189's analytic degree premise. The remaining
inputs are a single initial finite-free-factor order bound and completion for
every history in the exact cutoff class. Neither is silently assumed proved.
See Section~\ref{sec:actual-free-factor-history-revision190}.
% END REVISION190 FRONT STATUS

% BEGIN REVISION189 FRONT STATUS
\paragraph{The exact history consumer now accepts a degree bound.}
Revision188's degree-volume chain now feeds the entire inherited small-scale
noncollapsing contract, including event slabs and every required closed prefix
of terminal continuation. Actual checked consumers then produce general strong
canonical neighborhoods and finite geometric horizons. The supplied premise
is one positive integer bounding finite component-group orders uniformly over
the exact cutoff history class. Producing that integer remains the topological
frontier; no coherent global tower or final endpoint is accepted.
See Section~\ref{sec:history-degree-small-scale-revision189}.
% END REVISION189 FRONT STATUS

% BEGIN REVISION188 FRONT STATUS
\paragraph{Actual round-component degree and volume.}
A new checked chain produces the finite spherical covering of a general
constant-curvature model, identifies its degree with the actual target's
fundamental-group order, and supplies volume to the spatial canonical and
ambient-ball consumers. A fixed bound $N>0$ on these orders gives a volume
reserve independent of model noncollapsing. Producing one such $N$ uniformly
for the actual surgery history remains open. A possible initial-topology
free-factor argument is source-compared, not accepted as its producer.
See Section~\ref{sec:round-degree-volume-revision188}.
% END REVISION188 FRONT STATUS

% BEGIN REVISION187 FRONT STATUS
\paragraph{Current general analytic dependency.}
Six general analytic suppliers from the accepted PC release now feed the actual
strong canonical-neighborhood assembly and revision186's finite geometric
histories. The remaining supplied analytic contract is precisely small-scale
noncollapsing on the same initial metric. Its volume constant must precede the
later canonical thresholds. The source's simply-connected specialization uses
that hypothesis for high-curvature spatial-witness volume. This is a narrowed,
checked dependency, not a general small-scale producer or a coherent infinite
flow. See Section~\ref{sec:general-canonical-small-scale-revision187}.
% END REVISION187 FRONT STATUS

% BEGIN REVISION186 FRONT STATUS
\paragraph{Current finite-horizon iteration.}
The actual metric surgery event now supplies its marked discarded-component
geometry. Three proof modules construct a finite history to every prescribed
positive horizon, with the actual initial marking, one common cutoff-record
parameter function and geometric outputs at all actual events. The general
strong canonical-neighborhood package remains an explicit input; the public
uniform-debit step is applied without a simply-connected assumption. Histories
at different horizons are not yet a coherent tower. See
Section~\ref{sec:finite-geometric-horizons-revision186}.
% END REVISION186 FRONT STATUS

% BEGIN REVISION185 FRONT STATUS
\paragraph{Current discarded-component iteration.}
Three proof modules connect the actual surgery library's discarded-component
classification to complete fixed-model metrics on nonempty oriented factor lists.
The empty connected-sum list is explicitly represented by its sphere unit, with
the original smooth map retained. Actual event, component and factor indices feed
one labelled reconstruction, including a nonvacuous empty-retained consumer.
The tower, analytic classification, prime recognition, prescribed ports and whole
history reconstruction remain separate obligations. See
Section~\ref{sec:discarded-geometry-revision185}. Earlier notices are historical.
% END REVISION185 FRONT STATUS

% BEGIN REVISION184 FRONT STATUS
\paragraph{Current geometric-witness iteration.}
Four proof modules and thirteen axiom queries supply fixed elementary model
metrics, actual complete spherical quotient and sphere-times-circle atlases,
and a geometry witness on every literal standard-factor slot of one oriented
presentation. The model metric is fixed before chart choices. This is a bounded
part of M1d and AT25, not acceptance of the full eight-model endpoint, prime
recognition, prescribed surgery ports or the general-flow producer. See
Section~\ref{sec:fixed-model-atlases-revision184}. Earlier notices are historical.
% END REVISION184 FRONT STATUS

% BEGIN REVISION183 FRONT STATUS
\paragraph{Current composition iteration.}
Five actual applications now resolve the declaration-pair findings from
revision182; the Euclidean pair is correctly retained as two parallel
specializations. The new consumers preserve literal reconstruction slots,
the exact stored initial identification in both tower outcomes, and direct
bounds for one profile selected before its actual cutoff records. Three proof
modules and seven axiom queries pass; see
Section~\ref{sec:actual-composition-revision183}. General-flow production,
full marked reconstruction and the endpoint premise census remain open.
% END REVISION183 FRONT STATUS

% BEGIN REVISION182 FRONT STATUS
\paragraph{Current declaration audit.}
The accepted proof pilots now have an actual Lean declaration graph with complete
stored type expressions, separate proof/type references and axiom evidence.
Among twenty-eight registered producer/consumer pairs, twenty-two have a
syntactic proof-value path; six are related checked results whose advertised
composition remains unestablished. This repairs the evidence interpretation,
not the mathematics of the individual proved statements. See
Section~\ref{sec:actual-declaration-graph-revision182}. Earlier notices are historical.
% END REVISION182 FRONT STATUS

% BEGIN REVISION181 FRONT STATUS
\paragraph{Current implementation iteration.}
Round windowed models now produce the actual target canonical packet with its
quantitative volume reserve. A checked continuation argument retains that reserve,
and an explicit constant calculation rejects feeding the new bound into the same
positive noncollapsing constant. Four modules and eleven axiom queries pass.
History-wide analytic production and its interval quantifier order remain open;
see Section~\ref{sec:reserved-canonical-revision181}. Earlier notices are historical.
% END REVISION181 FRONT STATUS

% BEGIN REVISION180 FRONT STATUS
\paragraph{Current implementation iteration.}
Actual ancient-model noncollapsing now supplies a normalized-volume reserve,
including round models. Checked consumers transport it to the actual target
component and then to its ambient balls without simple connectivity. Four
modules and seven axiom queries pass. Uniform model data across the general
history and the coherent analytic induction remain open; earlier notices below
retain their historical scope. See Section~\ref{sec:model-volume-revision180}.
% END REVISION180 FRONT STATUS

% BEGIN REVISION179 FRONT STATUS
\paragraph{Current implementation iteration.}
Revision179 proves a component-volume-to-ball-volume estimate without simple
connectivity and an explicit common positive nonincreasing profile from
supplied positive interval budgets. Actual canonical-witness and cutoff-record
consumers compile; fourteen transitive axiom queries are standard. The uniform
volume and analytic budget producers, coherent general flow and full GC remain
open. Public PC endpoint acceptance from revision178 remains in force.
See Section~\ref{sec:general-flow-interfaces-revision179}.
% END REVISION179 FRONT STATUS

% BEGIN REVISION178 FRONT STATUS
\paragraph{Public foundation acceptance.}
The pinned public PC root target now builds successfully. Exact topological
and smooth endpoint consumers compile, and seven transitive axiom queries
use only the three standard classical axioms. All 17,874 tracked source
hashes were revalidated. Revision178 binds this scoped acceptance into the
current task graph and rejects unsupported promotion to general GC flow.
Historical preparation records remain unchanged. The full Geometrization
program, remaining producers and premise census continue; see
Section~\ref{sec:public-pc-acceptance-revision178}. The earlier iteration
notices below are historical: their pending-build statements describe
their own freeze dates, not the accepted status in this revision.
% END REVISION178 FRONT STATUS

% BEGIN REVISION177 FRONT STATUS
\paragraph{Current implementation iteration.}
Revision177 binds an existing smooth finite-sum reconstruction from the
actual cut/cap transition. Checked consumers retain every capped-component
label, identify its actual retained or discarded target, and prove
bijection for literal factor-list positions. The empty retained case
cannot lose a nonempty source into an empty discard list. Eighteen axiom
queries pass. The final reconstruction's orientation, marked sphere/cap
ports, cycle-ledger comparison and relative history substitution remain
separate; full AT21/AT24 and GC are not accepted. The public endpoint
audit remains pending. See Section~\ref{sec:event-reconstruction-revision177}.
% END REVISION177 FRONT STATUS

% BEGIN REVISION176 FRONT STATUS
\paragraph{Current implementation iteration.}
Revision176 binds actual spherical quotient metrics, completeness under
whole-carrier pullback, positive-space-form recognition without simple
connectivity, and orientation correction preserving the original factor
slots. Four bounded atoms and fourteen axiom queries are checked on the
pinned public PC source. An actual antipodal quotient tests the
nonsimply connected case. Actual discard production, prime witnesses,
marked surgery ports and general flow/profile production remain open.
See Section~\ref{sec:spherical-consumers-revision176}. The public PC
endpoint build is still pending at this iteration's freeze.
% END REVISION176 FRONT STATUS

% BEGIN REVISION175 FRONT STATUS
\paragraph{Current implementation iteration.}
Revision175 compiles actual observation-tower consumers: exact finite
prefixes, later regular smooth slabs, inherited initial identifications,
slice-metric coherence and the nonempty-or-absorbing-empty alternative.
The tower is an explicit input, not a constructed general GC flow.
A checked counterexample rules out exchanging separate finite-horizon
profile choices for one common positive profile. The source comparison
now points to KL's successive-interval construction for the missing
producer interface. The public PC endpoint build remains pending at
this iteration's freeze; the full user program remains active.
% END REVISION175 FRONT STATUS

% BEGIN REVISION174 FRONT STATUS
\paragraph{Current implementation iteration.}
Revision174 adds checked based-map and finite-side consumers on actual
PC fundamental groups and tori, and a concrete Euclidean smooth-metric
consumer. The public release comparison now isolates the round-component
restriction in the inspected noncollapsing producer and records the
existing explicit spherical-factor models. These are bounded bindings
and source comparisons; general flow/profile, marked event reconstruction
and the full Geometrization endpoint remain unaccepted.
% END REVISION174 FRONT STATUS

% BEGIN REVISION173 FRONT STATUS
Revision173 consolidates the later argument records into an individually
addressable current graph and audits the main integration paths through
both outcomes. It separates proof producers from supplied definitions,
keeps refuted and optional routes inactive, and records open leaves
explicitly. A further Lean consumer verifies metric-scale normalization
using the public PC library's actual smooth metrics and closed balls.
The graph is not a proof-closure certificate; the complete premise census,
general-flow/profile binding, event reconstruction and exceptional-factor
producers remain work. The public PC root build is still pending at this
iteration's freeze. See Section~\ref{sec:active-graph-revision173}.
% END REVISION173 FRONT STATUS

% BEGIN REVISION172 FRONT STATUS
Revision172 begins execution of the autoformalization program. It pins a
separate checkout of the public PC release, repairs task provenance and
current FC28 status, and supplies checked Lean implementations of the
MC06 radial estimate and MC07--MC09. Their actual downstream consumer
constructs MC24's finite-stage comparison map; the remaining limit
argument is not claimed. All ten queried pilot declarations report only
\code{propext}, \code{Classical.choice} and \code{Quot.sound}.
The public PC endpoint build remains in progress at this revision's
freeze; its release and GC consumers are not yet accepted. The whole
Geometrization implementation and the remaining work program are open.
Section~\ref{sec:implementation-revision172} records the exact scope.
% END REVISION172 FRONT STATUS

% BEGIN REVISION171 FRONT STATUS
Revision171 completes the requested written relative graph refinement and
geometric realization application under its explicit retained classical
foundations. The ONE final mathematical audit is complete in
Section~\ref{sec:graph-realization-final-audit}; GAU03 binds AT13/AT16 and
the graph part of GA06 to the actual late carriers. Complete metrics and
all cap, boundary and reconstruction markings now belong to the same
chosen pieces. Accepted-release Lean bindings and final history assembly
remain separate. Earlier revision-specific open-status text is historical.
% END REVISION171 FRONT STATUS

% BEGIN REVISION170 FRONT STATUS
Revision170 supplies complete metrics on all actual terminal blocks of the
relative graph construction. RG01--RG06 and GM01--GM06 give the written
AT13/AT16 and graph-region GA06 chain, retaining essential extra tori under
E1. The requested ONE final mathematical audit is next; it is not yet
recorded as completed. No accepted-release or Lean acceptance is asserted.
% END REVISION170 FRONT STATUS

% BEGIN REVISION169 FRONT STATUS
Revision169 supplies the actual smooth relative graph refinement RG01--RG06,
including AT13/AT16 for the chosen factors. It retains all external markings
and supplies essential good-block cuts by the existing finite pants phase.
Model realization is the next iteration; the requested final mathematical
audit has not yet begun. Historical formal statuses remain unaccepted in Lean.
% END REVISION169 FRONT STATUS

% BEGIN REVISION168 FRONT STATUS
Revision168 completes the requested written ambient cusp-incompressibility
application under its named classical foundations. The ONE final audit in
Section~\ref{sec:ambient-final-mathematical-audit} makes compact-ball and
surgery-profile hypotheses explicit and verifies the common carrier for
all nearby minimizing disks. The actual continuous area function and
its surgery-time upper barriers now give injection of the hyperbolic--thin
interfaces into the late flow manifold. Meeks--Yau and accepted-foundation
Lean bindings remain separate; no implementation acceptance is asserted.
Earlier revision-specific open-status paragraphs below are historical.
% END REVISION168 FRONT STATUS

% BEGIN REVISION167 FRONT STATUS
Revision167 develops ambient cusp incompressibility from the actual
persistent family. IMS01--IMS10 use a prescribed meridian, the
fixed-boundary Meeks--Yau theorem, a written local stability radius
estimate, unchanged compact carriers through surgery, and the actual
continuous least-area function with smooth upper barriers. This supplies
the ambient interface injections under explicit classical foundations.
The independent Meeks--Yau and accepted-PC Lean bindings remain open;
the historical free-boundary contracts are replaced for this endpoint.
The ONE final mathematical audit of this new goal is still pending.
% END REVISION167 FRONT STATUS

% BEGIN REVISION166 FRONT STATUS
Revision166 completes the requested long-time analytic goal as written
mathematics under its explicit analytic, surgery, geometric and topological
foundations. The ONE final mathematical audit in
Section~\ref{sec:collapse-longtime-final-analytic-audit} repairs the
finite-family starting threshold and clarifies the normalized estimate's
component hypothesis. The actual persistent hyperbolic regions and actual
thin graph carriers are now supplied. Accepted-PC/Lean verification,
ambient incompressibility and final reconstruction remain separate.
The revision-specific status paragraphs below preserve the development
history; their former open-frontier statements are superseded here.
% END REVISION166 FRONT STATUS

% BEGIN REVISION165 FRONT STATUS
Revision165 supplies TCF01--TCF07: actual slowly truncated thin carriers,
long cusp collars, equality of interior intrinsic and ambient curvature
scales and balls, boundary-near derivatives, large-radius exclusion,
and one derivative function chosen before the static collapse tolerance.
The actual late thin complement now has a labeled graph presentation
under the retained foundations. All requested analytic producers are
written; the ONE final mathematical audit is next. This does not assert
ambient incompressibility, final reconstruction, or Lean verification.
% END REVISION165 FRONT STATUS

% BEGIN REVISION164 FRONT STATUS
Revision164 supplies HPI01--HPI08, the actual HG08 persistent family.
Minimal-cusp corrections have vanishing jumps; thin cusp barriers
preserve the remaining regions; the volume budget gives a finite
disjoint family; and a separate cusp estimate proves thick coverage
on the matching radius-$1/\alpha(t)$ balls. The maps have the
controlled physical time dependence from HPS04. The intrinsic thin
complement and its static-collapse application remain in development;
the final analytic audit has not begun.
% END REVISION164 FRONT STATUS

% BEGIN REVISION163 FRONT STATUS
Revision163 supplies LTF01--LTF05, the actual HG07 flow producer:
normalized volume deficit forces hyperbolic limits; thick points have
macroscopic noncollapsed seeds; large spatial neighborhoods become
hyperbolic and continue through unscathed dyadic windows with all
finite smooth controls. The persistent minimal-cusp family and actual
intrinsic thin-carrier application remain in development. The final
analytic audit has not begun.
% END REVISION163 FRONT STATUS

% BEGIN REVISION162 FRONT STATUS
Revision162 proves WBD01--WBD04, the actual whole-ambient-ball
curvature and finite derivative estimates on late flow slices. The proof
retains the enlarged parabolic region above the surgery scale and
normalizes the canonical threshold below it. It uses the source's
center estimate without changing its conclusion. The intrinsic carrier,
large-radius and persistent-coverage adapters remain in development;
the final analytic audit has not begun.
% END REVISION162 FRONT STATUS

% BEGIN REVISION161 FRONT STATUS
Revision161 supplies the controlled time-smoothing construction HPS01--HPS05:
close maps with controlled pullback metrics have small smooth transitions,
which extend to compactly supported isotopies and interpolate the discrete
maps inside their old images. Embeddings, disjointness and physical speed
are proved under the explicit discrete-window hypotheses. The actual
minimal-cusp induction, flow estimates and thin-complement application
remain in development; the final analytic audit has not begun.

% END REVISION161 FRONT STATUS
% BEGIN REVISION160 FRONT STATUS
Revision160 begins the long-time analytic development from the completed
static blueprint. QRG04--QRG09 supply buffered injectivity comparison,
actual radial transversality, all target boundary components and cusp
counting, hence QRG00 and the existing HG06 quantitative consumer.
The persistent time-dependent family and actual thin-complement estimates
remain in development. This is written mathematics under retained
foundations, not Lean verification; the final audit of this new goal
has not begun.

% END REVISION160 FRONT STATUS
Revision159 completes the written Chapter14 static goal in
Section~\ref{sec:fibration-final-static-goal-audit}. The ONE final
mathematical audit checks both closed and nearly cuspidal boundary
constructions and repairs the explicit edge-exclusion parameter bound.
Actual clouds, whole fibers, compact pieces and compatible faces now
supply both static endpoints under their retained foundations.
Accepted-foundation Lean verification remains a separate obligation.

The project requires a completed, audited Poincar\'e release before production
GC implementation uses its foundation. An accepted completed PC release has not yet been bound for GC;
no historical snapshot or placeholder count establishes its present status. Meanwhile,
literature comparisons, mathematical contracts and proof decompositions can
proceed independently of its changing declaration names. The eventual
released foundation is imported through adapters and is not duplicated here.

The front matter retains Roman page numbers. The main body enters
\code{mainmatter} immediately before Chapter 1, so Arabic numbering
begins there at page 1. This revision preserves the LaTeX repairs and pauses
further PC-specific adapter expansion. The integrated Margulis and
Mostow developments now give a written proof of qualitative finite-volume
Mostow--Prasad rigidity relative to the specified inherited interfaces.
MPR61 supplies marked controlled representatives, MPR68 their boundary
maps, and MPR69--MPR73 regularity. Revision117 adds continuous
quasi-isometry orbit extraction and returning zooms in
Section~\ref{sec:mostow-qi-rigidity}: MPR78 supplies marked boundary
rigidity and MPR79 assembles existence and uniqueness in the original
homotopy class. Earlier statements and proofs are retained except for the explicit
homeomorphism hypotheses clarified in revision120. Accepted
inherited bindings, Lean verification, quantitative flow consumers and
ambient-flow peripheral injection remain pending. Revision118 adds the
end-to-end dependency audit in Section~\ref{sec:mostow-dependency-audit},
with twelve checked transfers, twelve inherited-interface families and
two corrections separating proof prerequisites from contract/application
references. Revision119 writes the quantitative compactness step in
Section~\ref{sec:hyp-quantitative-consumer}, conditional on the source's
global topological comparison. MPR79 applies directly to its homotopy
equivalence; the older product-collar lemma is optional. Revision120
completes the requested final working audit in
Section~\ref{sec:mostow-final-audit}, with two explicit homeomorphism
hypothesis clarifications and no further gap found relative to the
stated inherited interfaces. No completed formal HG04 theorem is asserted.
Revision104's whole-book audit and quantitative refinements are retained.
Hamilton--Ivey and smooth Schoenflies remain inherited PC candidates. Existing PC-dependent Lean
drafts stay provisional and unrun; the completed-release gate is pending.
The current implementation guide is
Section~\ref{sec:autoformalization-protocol}; its gates distinguish
written mathematics, source-qualified inputs and future Lean evidence.
The 15 atomic packets and five refinements are views of existing parent
contracts, not accepted declarations or new PC absence claims.
\paragraph{Revision126: local collapse and finite regularity.}
Sections~\ref{sec:collapse-finite-consumer-bridge}
and~\ref{sec:collapse-rank-two-construction} write controlled distance
smoothing with the metric unchanged, finite-regularity buffered isotopy,
smooth fiber-type recovery including boundary, and actual rank-two
adapted coordinates. The circle packet follows under the retained
tested residual-factor hypotheses. Finite-order compactness, actual
finite model embeddings, soul/core and classification producers, and
edge/slim geometric packets remain explicit obligations. This is the
first iteration of the current local-model task, not its final audit or
a Lean completion claim. The Chapter14 developments through revision125
are preserved and consume these local results only when their stated
hypotheses have been supplied.

\paragraph{Revision127: actual finite model maps and exact products.}
Section~\ref{sec:collapse-finite-embeddings-products} propagates basepoint
noncollapse to uniform local injectivity using retained PC CGT, proves
that the actual pointed comparison maps embed whole buffers and cover
smaller source balls, and constructs a $C^{K+1}$ product map and $C^K$
factor metric from an exact metric splitting. Its derivative budget
closes those clauses of the LC81 category bridge without smoothing
the metric. The finite-compactness chart/map producer and the geometric
edge/slim core and interpolation estimates remain; the full local-model
goal and its final audit are not complete. All earlier developments
and the historical parent contracts are retained.

\paragraph{Revision128: constructing the finite models.}
Section~\ref{sec:collapse-finite-compactness-construction} gives a
normal-coordinate construction under the actual smooth-source,
finite-curvature hypotheses. Its output is a $C^{K-1}$ metric
with actual pointed $C^K$ embeddings and complete nonnegative
factors tracking the supplied approximate splitting. The one
derivative loss is explicit and sufficient for the named local
fiber consumers at $K\geq10$. KL3.5's stronger $C^K$ metric
statement remains distinct. Geometric model cores, classification
and edge/slim interpolation estimates still remain; this is not
the final audit or completion of the local-model goal.

\paragraph{Revision129: the slim surface packet.}
Section~\ref{sec:collapse-finite-slim-construction} now constructs
LC85's compact finite surface factor, smooth sphere-or-torus fiber,
whole-fiber enclosure, transverse compact interpolation, proper smooth
bundle and cutoff. Curvature-derivative bounds are only eventual on
fixed balls. Actual minimizing directions supply the needed derivative
estimate; a value approximation is not substituted for it. Auxiliary
smooth metrics are used only to pass the Gauss--Bonnet integral, with
no curvature-sign preservation claim. The edge disk and zero-model
core/classification work remains; this is not the full goal's final audit.

\paragraph{Revision130: the finite edge model core.}
Section~\ref{sec:collapse-finite-edge-core} supplies the finite surface
topology adapter and an actual smooth disk core with a controlled
collar on the unchanged finite metric. It expands the connected-level
and flat-torus exclusions in the endpoint argument. A constant core
modification makes the model function smooth throughout its domain;
collar-only smoothing would not suffice. The source edge-set comparison,
variable-scale derivative and entire source-fiber transfer remain next.
This is not the final audit of the full local-model goal.

\paragraph{Revision131: transfer of the finite edge disk.}
Section~\ref{sec:collapse-finite-edge-transfer} constructs the source
disk bundle from an explicit buffered coarse-border hypothesis.
Triangle comparison controls every nearest-set direction before
noncollapse enters; the same finite model then transfers the collar.
The actual source sublevel is preserved by a smooth constant-core
modification, and its ENTIRE fiber is enclosed before interpolation.
The coarse-border producer from edge predicates and the full adapted
collar remain explicit, separate obligations. This is not the final
full-goal audit.

\paragraph{Revision132: producing the closed edge border.}
Section~\ref{sec:collapse-edge-border-producer} obtains the coarse-border
input of LFR25--LFR28 from actual strong/weak edge maps. A finite
four-point obstruction, buffered lifts, scale recentering and endpoint
extension replace an unexpanded set-limit assertion. The closure step
uses coarse estimates, not continuity of the approximation. One
physical-metric smoothing supplies a common normalized distance to any
finite verified edge family. Full adapted collars and strong-edge
coverage remain separate. The source has $C\geq200\Delta$; the
retained packet's strict inequality has an explicit adapter. This is
not the final full-goal audit.

\paragraph{Revision143: local consumer audit.}
Section~\ref{sec:collapse-finite-final-consumer-audit} completes the written
finite-category review. Foundation bindings and global assembly remain separate.

\paragraph{Retained development history.}
Section~\ref{sec:independent-collapse-contracts} records the earlier task and
records the version distinction discovered in the Morgan--Tian archive.
Section~\ref{sec:graph-refinement-contracts} tracks the graph-refinement
work. Revision 47 resolves the endpoint contract mismatch A25 by adopting
complete model geometry on every piece and requiring finite volume on every
hyperbolic-tagged piece; see Section~\ref{sec:endpoint-compatibility-gate}.
The torus-only cut architecture is retained. This specification decision
is not a proof of the geometric conversion.
The relative prime, filling and graph-rewrite targets are developed in
Sections~\ref{sec:relative-prime-refinement}--\ref{sec:matveev-graph-rewrites},
using the local archive and separately identified author texts. The
extended-pants and prime-terminal arguments are in
Sections~\ref{sec:extended-pants-splitting}--\ref{sec:graph-prime-normal-form}.
Revision 53 writes finite gluing injection and sphere reconstruction in
Section~\ref{sec:gluing-injection-transport}. Revision 54 constructs
simultaneous cap avoidance and the relative torus-family output in
Section~\ref{sec:cap-avoidance-output}, preserving the separate prime ledger
and the outer ambient-injection obligation.
Revision 55 audits the producer, assembly and carrier boundaries in
Section~\ref{sec:revision55-audit}, with source checks and fresh isolated
sandbox runs. Revision 56 writes the conditional late-slice ambient
injection argument and starts the independent infrastructure program in
Sections~\ref{sec:ambient-history} and~\ref{sec:independent-infrastructure}.
Revision 57 assigns mathematical topics to individual chapters and moves
implementation/release records to appendices. Environment pinning and
new Lean coding remain deferred. Revisions 58--59 expand Chapter 3 and
incorporate the BBI author errata. Revision 60 audits the metric chapter
and its source conversions. Revision 61 writes eventual pointed
extraction, the proper-source finite-prefix argument, and the conditional
packing reduction. Revision 62 develops Chapter 4's local radial/chart
covering estimates and inherited smooth specialization; the general
local-structure producer and limit geometry remain open. Revision 63
adds a quantitative distance-coordinate proof and records the new AKP,
KLP and Petrunin sources. Revisions 64--65 add complete buffers,
almost-shortest-path corrections and local openness. Revision 66
proves nearby paired distance charts with dimension-dependent
distortion before local compactness. Revision 67 transports finite
packing by almost-shortest paths, supplying AC16 and both clauses
of AC07 in written mathematics. Revision 68 proves midpoint transfer,
complete dyadic length realization, geodesicity and proper pointed
limit uniqueness in Chapter 3. Revision 69 passes lower comparison
to the pointed limit with controlled source domains and an explicit
zero-curvature angle limit. Revision 70 proves the Hausdorff upper
bound from quantitative nets and assembles the written compactness
bridge, retaining source-qualified local Toponogov. Revision 71
adds exact splitting data, a quantitative factor-dimension proof and
pointed low-dimensional product models, with splitting and local
one-dimensional recognition explicitly source-qualified. Revision 72
writes convergence of actual approximate product maps, preservation of
Euclidean coordinates and the prescribed-coordinate transfer. The
long-strainer passage was retained as a named source input. Revision 73
writes the opposite-endpoint defect, almost-shortest prefix extraction,
limiting lines and a quantitative Busemann-coordinate estimate.
Revision 74 proves cross-pair orthogonality by a numerical shortening
followed by buffered comparison, closing the written AC54 geometric
argument with the existing comparison and splitting qualifications.
Revision 75 writes the reverse implication with exact-radius endpoints
and an explicit approximation tolerance. Revision 76 writes compatibility
of the original approximate product maps under a fixed higher-rank
exclusion, conditional on the existing comparison and splitting inputs.
Revision 77 writes exact target-basepoint repair, quantitative recentering,
compatibility transport and convergence with a moving marked point.
Revision 78 writes radial cone data, the two-apex line, closure under
proper pointed limits and the conditional overlapping-cone theorem.
Revision 79's audit of Chapters 3--4 and their dependencies is in
Section~\ref{sec:revision79-dependency-audit}. Revisions 80--89 develop
Chapter 13's cone scales, radial smoothing, normal-flow cores and
uniform collar transfer. The PC comparison retains its scoped Lean
checks. Revision 96 completes and audits Chapter 13's blueprint coverage;
its source-qualified producers and inherited formal bindings remain
explicit. See Section~\ref{sec:collapse-completion-audit}.
Revision 101 completes Chapter 14's blueprint coverage and its
full dependency/source audit; see
Section~\ref{sec:fibration-complete-audit}. Revision 102 completes
Chapter 15 coverage and audits its entire reconstruction route; see
Section~\ref{sec:transport-complete-audit}. Source-qualified producers
and inherited release bindings remain open. Chapters 3, 4, 13 and 14
are preserved. Revision 103 makes finite-volume Mostow--Prasad rigidity,
hyperbolic persistence and ambient cusp injection explicit in Chapter 8,
binds the actual late-or-terminal producers in Chapter 12, and completes
Chapter 16 topic coverage and its full audit; see
Section~\ref{sec:assembly-complete-audit}. Qualified producers and
accepted-release bindings remain open. Lean formalization and environment
pinning remain deferred.

\section*{Prerequisite gate}

Production GC Lean implementation begins only after the Poincar\'e repository has a pinned,
buildable release containing its completed public theorem and no project-local
\code{sorry}, \code{admit}, or unreviewed \code{axiom}. The release gate requires an exact
import path, a stable theorem statement in the chosen category, a successful
root build, and a \code{#print axioms} audit showing only the approved
foundational axioms.

\mainmatter

\chapter{Target theorem, conventions and chapter guide}
\label{ch:topic-guide}

% BEGIN migrated B001
\label{ch:target}

\begin{roadmap}{Freeze the mathematical target and all conventions before
building the analytic machinery. These definitions are shared interfaces and
must not drift between workstreams.}
\end{roadmap}

\section{Public theorem}

The public theorem has one stable interface:
\begin{equation}
\begin{aligned}
 &\texttt{geometrization} : \\
 &\quad \texttt{ClosedConnectedOrientableSmoothThreeManifold M} \\
 &\quad\longrightarrow \texttt{Geometrizes M}.
\end{aligned}
\end{equation}

The project defines both predicates. \code{Geometrizes M} contains explicit finite
decomposition data: prime pieces, embedded incompressible tori, boundary and
gluing maps, and geometric structures on the pieces. The eight Thurston model
geometries are explicit objects. No essential convention is hidden in an
informal phrase such as ``has a geometric decomposition.''

The public theorem also needs a construction of its analytic input.
The planned \code{GlobalRicciFlowSurgery} contract must supply
\code{Nonempty (RicciFlowSurgeryPackage M)} for the public manifold,
using an inherited smooth initial metric, a proved normalization and
admissible surgery parameters. The exact completed-release hypotheses
remain to be bound. A conditional theorem from a supplied flow package
alone does not establish this existence. The final composition uses that
construction, the separate \code{long_time_outcome} theorem, both outcome
branches and the audited endpoint adapter; no such construction is
claimed checked here.

\section{Proposed endpoint certificate}
\label{sec:endpoint-certificate}

The first target will use an existential certificate rather than requiring a
canonical JSJ decomposition as part of the public theorem. This keeps the
formalization focused on the geometrization conclusion. The certificate has
the following mathematical fields:
\begin{enumerate}[label=\textbf{C\arabic*:},leftmargin=*]
\item a finite prime decomposition of $M$;
\item for each prime factor, a finite family of pairwise disjoint embedded
  incompressible tori;
\item the compact pieces obtained by cutting along those tori, with their
  boundary identifications;
\item for every resulting piece, a complete metric on its interior locally
  modeled on one of the eight Thurston geometries, with finite volume
  required whenever the selected model is hyperbolic.
\end{enumerate}
Thus \code{Geometrizes M} is initially the proposition that such a finite
certificate exists. Canonicality and uniqueness of the JSJ system are useful
later theorems, but they are not prerequisites for the first public
geometrization theorem. Revision 47 adopts the torus-only convention explained in
Section~\ref{sec:endpoint-compatibility-gate}. Complete infinite-volume
flat product and twisted-bundle interiors are allowed; hyperbolic pieces
still require finite volume. Other branches may provide separate volume
proofs for their chosen metrics. This corrects C4 consistently across the
interfaces without adding one-sided surface cuts.

The corresponding Lean-facing structures should be named only after the
existing Poincar\'e project and mathlib APIs have been audited. The theorem
inventory records the proposed names \code{GeometrizationCertificate},
\code{PrimeDecomposition}, \code{TorusCutSystem}, \code{CutPiece}, and
\code{HasThurstonGeometry} as provisional identifiers.

\section{Top-down contract freeze}
\label{sec:top-down-contract}

The next design task is to replace the provisional endpoint description by a
precise data contract. We work from the public theorem downward in the
following order:

\begin{enumerate}[label=\textbf{T\arabic*:},leftmargin=*]
\item \textbf{Piece geometry.} Define a \code{GeometricPieceCertificate} for a
  compact oriented 3-manifold with boundary. It records the model among
  spherical, Euclidean, hyperbolic, $\Sph^2\times\R$, $\mathbb{H}^2\times\R$,
  $\widetilde{\mathrm{SL}}_2(\R)$, Nil, and Sol, together with the complete
  locally homogeneous metric on the interior, finite volume when the tag is
  hyperbolic, and the appropriate cusp or boundary data.
\item \textbf{Cut piece.} Define \code{CutPiece} with its compact carrier,
  oriented boundary components, inclusion into the cut manifold, and the
  identification of each boundary torus with a side of the torus-cut system.
\item \textbf{Torus system.} Define \code{TorusCutSystem} with a finite index
  type, pairwise disjoint two-sided embedded tori, incompressibility proofs,
  and a cut-space construction whose boundary maps are explicit.
\item \textbf{Prime layer.} Define \code{PrimeDecomposition} and its connected-sum
  reconstruction. The torus system is attached to each prime factor, so the
  assembly theorem can distinguish sphere cuts from torus cuts.
\item \textbf{Certificate and predicate.} Package these fields as
  \code{GeometrizationCertificate M} and set \code{Geometrizes M} to its nonemptiness.
  No canonicality or uniqueness field is included in the first public target.
\end{enumerate}

The structural assembly theorem is then fixed before any Ricci-flow proof:
\begin{equation}
\begin{aligned}
 &\texttt{assemble\_geometrization} : \\
 &\quad \texttt{PoincareAdapter} \\
 &\quad\longrightarrow \texttt{ThickThinCollapsePackage M} \\
 &\quad\longrightarrow \texttt{GeometrizationCertificate M}.
\end{aligned}
\end{equation}
The analytic producer is deliberately downstream:
\begin{equation}
\begin{aligned}
 &\texttt{PoincareAdapter} \\
 &\quad\longrightarrow \texttt{RicciFlowSurgeryPackage M} \\
 &\quad\longrightarrow \texttt{LongTimeOutcome M} \\
 &\quad\longrightarrow \texttt{GeometrizationCertificate M}.
\end{aligned}
\end{equation}
This dependency order prevents a convenient Ricci-flow representation from
silently defining the theorem we are trying to prove.

\section{Top-down acceptance criteria}

The endpoint contract is frozen when the following questions have explicit
answers in the theorem inventory:
\begin{itemize}[leftmargin=*]
\item What is the carrier of each compact cut piece, and how is its interior
  represented?
\item What exactly proves that a torus is embedded, two-sided, pairwise
  disjoint, and incompressible?
\item Which boundary components are paired by gluing maps, and how does the
  reconstruction recover the original oriented manifold?
\item What data identifies a complete locally homogeneous metric, its
  conditional hyperbolic volume proof, and its cusp or boundary collars?
\item Which fields come from topology alone, which are produced by the
  thick/thin package, and where does the completed PC adapter enter?
\end{itemize}

Until these answers are fixed, further analytic formalization is premature.
Once they are fixed, topology, hyperbolic geometry, collapse, and the PC
adapter can proceed as independent workstreams against the same endpoint.
% END migrated B001

\section{How to use the topic chapters}
\label{sec:chapter-framework}

Revision 57 organizes the existing blueprint by mathematical topic.
The mathematical architecture and theorem contracts are retained.
Mathlib environment pinning and new Lean implementation are deferred at
the user's request; the active work is the detailed chapter blueprints.
All chapters and the bibliography remain in this one TeX file.

Each topic chapter will contain its purpose, definitions and conventions,
precise theorem statements, required inputs, downstream outputs,
lemma-level proof plan, exact source/errata record, and open obligations.
A chapter outline is not a proof or a claim that its infrastructure is
missing from the inherited foundation. Early placement of a static topic
is an exposition choice, not a change to the M3 analytic priority.
The explicit chapter IDs below describe this book. C0--C13 remain the
legacy implementation-scaffold IDs in Appendix~\ref{ch:implementation-crosswalk}.
Milestone numbers, endpoint fields and theorem IDs keep their meanings.

\setlength{\tablebodywidth}{\dimexpr\textwidth-4\tabcolsep\relax}
\begin{longtable}{@{}P{0.08\tablebodywidth}P{0.50\tablebodywidth}P{0.42\tablebodywidth}@{}}
\textbf{Ch.} & \textbf{Topic} & \textbf{Current development}\\\hline
\endhead
\ref{ch:topic-guide} & Target theorem, conventions and chapter guide & Endpoint and current implementation guide.\\
\ref{ch:inherited-topic} & Inherited foundation and the PC reuse boundary & Historical discovery; accepted-release gate pending.\\
\ref{ch:metric-geometry} & Metric geometry and pointed limits & 24 contracts; written estimates and extraction; smooth bridge open.\\
\ref{ch:alexandrov-geometry} & Alexandrov geometry and local models & 87 contracts; written refinements; comparison and splitting qualified.\\
\ref{ch:three-manifold-topology} & Three-manifold topology, prime factors and torus cuts & Actual cut/prime contracts; relative producers open.\\
\ref{ch:seifert-graph} & Seifert fibrations and graph manifolds & Filling and graph progress; essential-cut producers open.\\
\ref{ch:model-realization} & Thurston geometries and model realization & E1 models; orientation correction; retained bundle cuts.\\
\ref{ch:hyperbolic-geometry} & Hyperbolic geometry, cusps and thick limits & 18 contracts; top-down rigidity proof begun; qualified producers open.\\
\ref{ch:ricci-flow-pde} & Ricci-flow PDE interfaces and surviving gaps & Surviving PDE interfaces; completed-release comparison pending.\\
\ref{ch:perelman-estimates} & Perelman estimates and geometric compactness & Inherited estimate interfaces; exact consumer bindings pending.\\
\ref{ch:surgery-topic} & Ricci flow with surgery & Flow/event contracts; exact inherited bindings pending.\\
\ref{ch:long-time-topic} & Long-time behavior and thick/thin decomposition & 12 contracts; actual common-slice or terminal producers.\\
\ref{ch:local-collapse-topic} & Local collapsing theory & 90 contracts; coverage audited; geometric/category gates open.\\
\ref{ch:compatible-fibrations} & Compatible fibrations and collapse assembly & 46 contracts; coverage audited; constraint register begun.\\
\ref{ch:ambient-reconstruction} & Ambient transport and reconstruction & 31 contracts; coverage audited; qualified transport inputs.\\
\ref{ch:global-assembly-topic} & Global assembly and Geometrization & 24 contracts; coverage audited; conditional packaging.\\
\end{longtable}

\code{chapter_framework.csv} records chapter ownership and section plans.
\code{chapter_assignments.csv} assigns every theorem/interface inventory
row a primary chapter without changing its statement, status or proof
dependencies. Source rows also carry topic-chapter navigation. Existing
labels remain valid across the move; historical revision/page locators
still refer to their original files.

The current pass refines concrete contracts and removes implementation
bottlenecks; see Section~\ref{sec:autoformalization-protocol}. Chapters 3--4
have expanded proofs and qualified structural inputs, and Chapters 13--16
have completed blueprint coverage with open producers. MG01's later
environment audit remains separate from mathematical statement refinement.

\chapter{Inherited foundation and the PC reuse boundary}
\label{ch:inherited-topic}

\paragraph{Scope.} The geometric and analytic foundation to inherit, with completed-release scope left pending.

\paragraph{Planned sections.}
\begin{enumerate}[leftmargin=*]
\item Reusable objects.
\item Convention ledger.
\item Adapter boundary.
\item Release prerequisites.
\end{enumerate}

\paragraph{Existing material and present task.} W0 audit and PC reuse policy; detailed M2A records in Appendix C. Clarify chapter inputs at the mathematical level; do not expand snapshot probes.

\paragraph{Current reading boundary (revision104).}
The following W0 and initial classification material describes earlier
discovery snapshots. It is not a census of the completed PC release and
must not be mined for current absence claims or instructions to repair
that changing checkout. Preserve its evidence; bind current GC consumers
only to an accepted release or an explicitly stabilized interface. The
current task packets keep accepted declarations and imports unbound.

% BEGIN migrated B002
\section{W0 audit of the existing Lean base}
\label{ch:w0-audit}

The current development checkout already contains a substantial reusable
Ricci-flow and geometric-analysis library. The audit below is a source-level
inventory, not a claim that these results automatically solve their GC
interfaces. Their exact signatures and transitive axiom closures must still be
checked before import.

\setlength{\tablebodywidth}{\dimexpr\textwidth-2\tabcolsep\relax}
\begin{longtable}{@{}P{0.348315\tablebodywidth}P{0.651685\tablebodywidth}@{}}
\textbf{Existing component} & \textbf{Evidence and GC relevance}\\\hline
\endhead
Short-time Ricci flow &
\code{Geometry/Flow/RicciFlow/ShortTime/Existence.lean} contains
\code{ricci_flow_short_time_existence}; the DeTurck construction and
gauge recovery are already developed.\\
Reduced-volume monotonicity &
\code{Geometry/Flow/RicciFlow/Perelman/LGeometry/ReducedVolume/Basic.lean}
contains \code{redVolume_anti}, together with the reduced exponential,
Jacobian, and cut-locus infrastructure.\\
Three-dimensional pinching &
\code{Geometry/Flow/RicciFlow/DimensionThree/HamiltonIvey/MaximumPrinciple.lean}
contains \code{hamilton_ivey_pinching} and the asymptotic pinching
estimate.\\
Flow compactness &
The \code{Compactness} tree contains pointed-flow limits, Shi estimates,
curvature convergence, and the definition
\code{HasSmoothCheegerGromovLimit}.\\
Surgery-local geometry &
The \code{Geometry/Flow/RicciFlow/Surgery/StandardCap} tree contains
standard-cap metrics, collars, truncation, and collapse-map constructions.
The present pre-release checkout exposes this mainly as local surgery
infrastructure. Under the project assumption, the completed PC release is
expected to export the global flow-with-surgery object and its canonical-
neighborhood/event/continuation interfaces; GC must verify the public names
and adapt them rather than rebuild them.\\
Three-manifold cut/cap topology &
The \code{Topology/ThreeManifold} tree contains connected-sum laws,
cut/cap reconstruction, spherical space-form reductions, and
\code{isPoincareStandard} interfaces. These are reusable inputs for the
spherical and endpoint-geometry branches.\\
Boundary and smooth-gluing infrastructure &
The source-level audit of the current checkout finds extensive
\code{Topology/Manifold/Boundary}, \code{BoundaryCollar},
\code{Diffeomorph}, and connected-sum smooth-structure modules, including
cut/cap quotient and reconstruction files. These support the M1 wrapper, but
no single public GC certificate type yet packages arbitrary compact smooth
3-manifolds with torus boundary.\\
Missing dedicated GC layer &
The initial inventory found no dedicated public modules for prime
decomposition, incompressible tori, JSJ systems, Seifert pieces, graph
manifolds, or the eight geometric structures. These become the first genuinely
new topology contracts.\\
\end{longtable}

\subsection{Consequences for the work plan}

The project should not restart the analytic foundations. W2/M2 first audits
and adapts the existing declarations. W1 remains the endpoint topology
workstream. After M2, W3/M3 is the first substantial new mathematical
workstream: formalize the Ricci-flow PDE layer and record precisely which
evolution and regularity results are genuinely absent. W4/M4 then handles
Perelman estimates and compactness, W5/M5 integrates surgery, and W6/M6
supplies the long-time output consumed by W1.

The current checkout contains substantial Poincar\'e-standard and spherical
space-form infrastructure. Under the project assumption, the finished
Poincar\'e release becomes a hard prerequisite rather than an active GC
workstream. The adapter is written once against that pinned public theorem;
the GC proof never reopens or duplicates the Poincar\'e development.

\subsection{Audit acceptance}

Before a component is marked \code{existing-reuse} in the theorem inventory, record
its exact import, declaration signature, build target, and \code{#print axioms}
output. A README summary is useful navigation but is not proof evidence.

Before the GC project is declared ready to start, additionally record the
Poincar\'e release tag or commit, its public theorem name, the adapter theorem,
and the completed axiom audit. This is a gate for project initialization, not a
temporary hypothesis inside the GC proof.
% END migrated B002

% BEGIN migrated B003
\section{PC reuse boundary and the minimal GC path}
\label{ch:pc-reuse}

The completed Poincar\'e project is a major input to GC. We should reuse its
proved analytic and local-surgery results instead of rebuilding them. The
word ``reuse'' means that the result is available through a stable public
declaration in the pinned release, with its hypotheses and axiom closure
audited. A file name, a large import tree, or a theorem whose conclusion is
conditional on an input structure does not by itself satisfy this condition.

\subsection{Three reuse classes}

\setlength{\tablebodywidth}{\dimexpr\textwidth-4\tabcolsep\relax}
\begin{longtable}{@{}P{0.233333\tablebodywidth}P{0.388889\tablebodywidth}P{0.377778\tablebodywidth}@{}}
\textbf{Class} & \textbf{Likely PC contribution} & \textbf{GC action}\\\hline
\endhead
Direct reuse & Riemannian and tensor APIs; Ricci-flow equation; short-time
existence and uniqueness on smooth intervals; evolution identities; maximum
principles; Hamilton--Ivey pinching; reduced-volume machinery; Shi estimates,
flow compactness, any exported noncollapsing theorem, and the completed
canonical-neighborhood and surgery-invariant theorems. & Import through the
public modules, check exact signatures and \code{#print axioms}, and state GC
lemmas in the notation of the imported API. Do not duplicate proofs.\\
Adapter reuse & Standard-cap metrics, collars, truncation, collapse maps,
finite surgery histories, global surgery-event records, cut/cap
reconstruction, spherical-space-form interfaces, and the completed
Poincar\'e conclusion. & Build a thin GC adapter that translates the PC
event/history types into the GC long-time interface. Preserve the imported
invariants; do not expose PC implementation details in the public GC theorem.\\
GC-new & Long-time thick/thin output;
hyperbolic limits; application of local collapse; graph-manifold to geometric
decomposition; and prime/JSJ endpoint topology. & Formalize the missing
theorem contracts and proofs. These are not made reusable merely by importing
a module that contains local cap or extinction vocabulary.\\
\end{longtable}

This classification is intentionally conservative. In particular, the PC
checkout contains a substantial \code{Extinction} tree, but its public
\code{PoincareControlledExtinction} structure records a finite surgery history,
control data, and an extinction witness for the Poincar\'e endgame. We do not
make that machinery a GC dependency merely because it is present. The
completed-release audit records it as reusable support only if a later GC
theorem explicitly needs it.

\subsection{Finite extinction is outside the default GC route}
\label{sec:finite-extinction-role}

Finite-time extinction is not part of the statement of GC, and it is not part
of the default GC work plan. The central analytic output is the long-time
behavior of Ricci flow with surgery: thick regions with hyperbolic limits and
collapsed regions with graph-manifold structure. Morgan--Tian's completion
identifies that long-time and collapse stage as the completion of the
geometrization proof.

The completed PC theorem lets GC use the simply connected spherical result
without replaying the Poincar\'e argument. It does not logically imply the
broader finite-extinction theorem for every manifold whose prime factors have
no aspherical pieces. That distinction is harmless for the project plan: W1
will give the spherical-space-form and $\Sph^2\times\Sph^1$ cases the explicit
geometric structures required by the endpoint certificate, while W5--W6 prove
the global surgery and long-time decomposition that GC actually uses.

Perelman's extinction paper remains a bibliography and a possible diagnostic
source. We activate it only if a later, precisely stated GC theorem cannot be
proved by the direct endpoint and long-time route. Until such a dependency is
found, no \code{FiniteExtinctionAdapter}, finite-extinction work package, or
curve-shortening formalization belongs in the GC critical path.

\subsection{Revised dependency consequences}

The dependency graph is therefore not a second copy of the PC project:

\begin{itemize}[leftmargin=*]
\item W2--W4 begin with an import-and-audit pass. New work is limited to
  notation adapters, missing hypotheses, and genuinely GC-specific estimates.
\item W5 imports the completed PC global flow-with-surgery, canonical-
  neighborhood, surgery-event, continuation, and invariant-preservation
  interfaces. It builds the GC adapter around those declarations and adds only
  fields needed for long-time geometric output.
\item W6 is the main GC analytic frontier: long-time thick pieces, hyperbolic
  limits, local collapse, and graph-manifold output.
\item W1 formalizes the prime/torus/geometry endpoint and converts W6 output
  into the existential certificate; it imports the PC adapter only where the
  spherical branch is used.
\end{itemize}

The first audit after the completed PC release is a matrix with one row per
candidate theorem and columns for exact declaration, public import path,
hypothesis compatibility, axiom closure, and GC status (\code{direct}, \code{adapter},
or \code{new}). This matrix is the evidence for shrinking the GC work plan; until
then, the blueprint should not say that finite extinction is the only missing
result.
% END migrated B003


\paragraph{Current PC audit.} Revision 87 audits the current private PC branch at
\code{54ad4d8ec040}; see Section~\ref{sec:collapse-pc-reuse-binding}.
Its scoped module builds and axiom checks are separate from the preserved,
unrun M2A preparation probes. The full root build and completed-release gate
remain unsuccessful. Historical snapshot-discovery restrictions do not
cancel this user-requested audit; production GC environment pinning remains
deferred.

\chapter{Metric geometry and pointed limits}
\label{ch:metric-geometry}

This chapter supplies the static metric language needed to extract local
collapse models and to interpret the metric consequences of smooth
hyperbolic limits. Its main outputs are finite-net compactness, pointed
approximation calculus, uniqueness of pointed limits, and a conditional
smooth-to-metric bridge. Curvature supplies hypotheses for these results
in later chapters; metric convergence by itself supplies neither a smooth
structure nor a fibration nor the topology of a collapsed manifold.

The source backbone is BBI, Chapters 1--2 and 7--8
\cite{BuragoBuragoIvanov}, compared with Kleiner--Lott's local-collapse
Sections 2--3 and Lemma 6.1 \cite{KleinerLottCollapse2014}.
All statements in this chapter are mathematical specifications or written
arguments. None is newly checked in Lean. The July 6, 2024 author errata
\cite{BuragoBuragoIvanovErrata2024} have now been read and the relevant
corrections incorporated; the remaining proof and reuse obligations are
separate from that completed source check.

\section{Metric, length and compactness conventions}
\label{sec:metric-basic}

\begin{definition}[Metric objects and balls; MC01]
All distances in this chapter are finite real numbers. A metric on $X$
is a symmetric function $d:X\times X\to[0,\infty)$ satisfying the triangle
inequality and $d(a,b)=0$ exactly when $a=b$. A pseudometric retains
$d(a,a)=0$ but allows distinct points to have zero distance. A pointed space $(X,p)$ includes a
chosen $p\in X$, hence is nonempty. Set
\[
 B_X(p,R)=\{x:d(p,x)<R\},\qquad
 \overline B_X(p,R)=\{x:d(p,x)\leq R\}.
\]
The second notation means a closed metric ball, not the closure of the
first set; those sets need not agree in a general metric space. Subsets
carry restricted ambient distances unless an intrinsic metric is
explicitly named. An isometric embedding preserves distances; an
isometry is additionally onto. A pointed isometry sends the chosen
basepoint to the chosen basepoint.
\end{definition}

BBI Definitions 1.1.1--1.1.5 allow infinite distances; we take their
finite-distance specialization and their zero-distance quotient.
For a disconnected Riemannian slice we apply the metric construction
componentwise. We do not turn an infinite path distance between components
into a finite metric by convention.

\begin{definition}[Lengths, completeness and properness]
For a continuous curve $\gamma:[a,b]\to X$, define
\[
 L_d(\gamma)=\sup_{a=t_0<\cdots<t_m=b}
       \sum_{j=1}^{m} d(\gamma(t_{j-1}),\gamma(t_j)),
\]
with value in $[0,\infty]$. The space is a \emph{length space} if each
distance is the infimum of the lengths of curves joining its endpoints.
It is a \emph{geodesic space} if every pair is joined by a distance
minimizer, parameterizable as an isometric map from an interval of length
$d(x,y)$. This does not assert that every locally minimizing geodesic is
globally minimizing. Completeness means convergence of every Cauchy
sequence. Properness means compactness of every closed bounded subset,
equivalently of all closed balls. Local compactness means that every
point has a neighborhood with compact closure.
\end{definition}

\begin{lemma}[Lower semicontinuity of length; MC01]
\label{lem:metric-length-lsc}
Let $\gamma_i,\gamma:[a,b]\to X$ be continuous curves in one metric
space, with $\gamma_i(t)\to\gamma(t)$ for every $t$. Then, in the
extended nonnegative reals,
\[
 L_d(\gamma)\leq\liminf_{i\to\infty}L_d(\gamma_i).
\]
No finite-length hypothesis on the limit curve is required.
\end{lemma}
\begin{proof}
For each fixed finite partition $P=(t_0,\ldots,t_m)$, continuity of the
distance function gives
\[
 \sum_{j=1}^m d(\gamma(t_{j-1}),\gamma(t_j))
 =\lim_{i\to\infty}\sum_{j=1}^m
       d(\gamma_i(t_{j-1}),\gamma_i(t_j))
 \leq\liminf_{i\to\infty}L_d(\gamma_i).
\]
Taking the supremum over $P$ proves the claim, including infinite length.
The partition is held fixed before taking the limit. This implements the
repair to BBI Proposition 2.3.4(iv), printed page 35/PDF page 50, in
errata page 2. The original proof's number of partition points depends on
its error parameter, so its unadjusted accumulated error need not vanish.
\end{proof}

\begin{lemma}[Compactness of uniformly Lipschitz curves; MC01]
\label{lem:metric-curve-compactness}
Let $X$ be a compact metric space and $C\geq0$. A sequence of
$C$-Lipschitz maps $\gamma_i:[0,1]\to X$ has a subsequence converging
uniformly to a $C$-Lipschitz map $\gamma:[0,1]\to X$.
\end{lemma}
\begin{proof}
Choose the rational parameters in $[0,1]$, including the endpoints. Compactness and
a diagonal subsequence give convergence at every point of this set.
The limiting values satisfy the same Lipschitz inequality, so completeness
of $X$ extends them uniquely to a $C$-Lipschitz map on $[0,1]$.
If $C=0$, the curves are constant and convergence at one point is uniform.
If $C>0$, given $\varepsilon>0$, choose an integer $N>4C/\varepsilon$
and the rational grid $0,1/N,\ldots,1$. Along the subsequence,
the errors at its finitely many vertices are eventually less than
$\varepsilon/2$. For each $t$ choose a grid vertex $u$ within $1/N$.
Then
\[
 d(\gamma_i(t),\gamma(t))
 \leq C|t-u|+d(\gamma_i(u),\gamma(u))+C|u-t|
 <2C/N+\varepsilon/2<\varepsilon.
\]
Thus convergence is uniform. This is the common-Lipschitz specialization
of BBI Theorem 2.5.14, printed pages 47--48/PDF pages 62--63, with
the mesh constants corrected on errata page 3.
\end{proof}

\begin{theorem}[Metric Hopf--Rinow; source-qualified]
\label{thm:metric-hopf-rinow}
A complete locally compact length space with finite distances is proper
and geodesic. In particular a proper length space is geodesic.
\end{theorem}
\begin{proof}[Proof plan from BBI Proposition 2.5.22 and Theorem 2.5.23]
Fix a center. Local compactness gives a compact ball of positive radius.
If the supremum of compact-ball radii were finite, near-minimizing curves
would show that the ball at that supremum has finite nets at every scale.
Completeness makes it compact. A finite cover by relatively compact
neighborhoods then gives a slightly larger compact ball, a contradiction.
For two endpoints, restrict a minimizing sequence of curves to a fixed
compact ball, parameterize with uniformly bounded speed, extract a
uniform limit, and use lower semicontinuity of length. This proves
existence of a minimizer. Lemmas~\ref{lem:metric-curve-compactness}
and~\ref{lem:metric-length-lsc} now supply written arguments for the
curve compactness and length passage. The arclength reparameterization
input is BBI Proposition 2.5.9, printed page 46/PDF page 61; its body and
proof were read. The compact-ball argument, reparameterization and these
lemmas still require formal proof or reuse bindings.
\end{proof}

The length hypothesis is essential: an infinite discrete space with all
nonzero distances equal to one is complete and locally compact, but its
closed unit ball is not compact. The inherited Riemannian-distance and
Hopf--Rinow interfaces must establish these metric hypotheses for the
complete connected manifolds used later. This is an adapter boundary,
not a request to reconstruct PC's differential geometry.

\section{Finite nets, packing and elementary estimates}
\label{sec:metric-nets}

\begin{definition}[Witnessed finite nets; MC02]
For $\varepsilon>0$, a finite $\varepsilon$-net in a subset $A\subset X$ is a finite
$S\subset A$ such that for every $a\in A$ there is $s\in S$ with
$d(a,s)\leq\varepsilon$. The empty set has the empty finite net; a net
in a nonempty set is necessarily nonempty. A set is $\rho$-separated if distinct points
have distance at least $\rho>0$. Total boundedness means that finite
$\varepsilon$-nets exist for every $\varepsilon>0$.
\end{definition}

BBI Definition 1.6.1 uses $d(a,S)\leq\varepsilon$. For a nonempty finite
$S$ the minimum is attained, giving precisely our witnessed convention.
For arbitrary nonclosed $S$, an infimum bound alone need not supply a
witness with the same non-strict bound. Open covering statements below
instead require an actual distance strictly less than the error.

\begin{lemma}[Nets and packing]
If $A$ has a finite $\varepsilon$-net of size $N$ and
$\rho>2\varepsilon$, every $\rho$-separated subset of $A$ has size at most $N$.
If every finite $\rho$-separated subset of $A$ has size at most $N$, then
$A$ has a finite net with at most $N$ points and covering distances
strictly less than $\rho$. A metric space is compact exactly when it is
complete and totally bounded. Every proper metric space is complete.
\end{lemma}
\begin{proof}
Assign each separated point a net point within $\varepsilon$. Two points
assigned the same center have distance at most $2\varepsilon<\rho$,
so the assignment is injective. For the converse, successively add a
point whose distance from all existing centers is at least $\rho$.
The finite packing bound forces this procedure to stop by $N$ points;
at stopping all distances to some center are strictly less than $\rho$.
For compactness, finite nets at scales tending to zero give a Cauchy
subsequence of any sequence by repeated selection of an infinite subset
in one covering ball; completeness gives its limit. Conversely compact
metric spaces are complete and finite subcovers give finite nets. A
Cauchy sequence in a proper space lies in one compact closed ball and
therefore converges. See BBI Exercise 1.6.4 and Theorem 1.6.5.
\end{proof}

A map $f:X\to Y$ is $C$-Lipschitz, for $C\geq0$, if
$d_Y(f(x),f(x'))\leq C d_X(x,x')$ for every $x,x'\in X$.

\begin{lemma}[Distance functions]
For a nonempty $A\subset X$, $d_A(x)=\inf_{a\in A}d(x,a)$ is finite and
$1$-Lipschitz. A $C$-Lipschitz map satisfies
$L(f\circ\gamma)\leq C L(\gamma)$ for every rectifiable curve.
\end{lemma}
\begin{proof}
The triangle inequality gives $d_A(x)\leq d(x,y)+d_A(y)$; interchange
$x,y$. For lengths, apply the Lipschitz estimate to every partition sum
and take the supremum. No differentiability is needed.
\end{proof}

These facts turn geometric packing bounds into the exact finite-net
hypotheses of compactness. Chapter~\ref{ch:alexandrov-geometry} must
supply the curvature-dependent uniform packing bound. A separate
finite net in each member, with no uniform cardinality, is insufficient.

% BEGIN migrated B058
\section{MG02: a positive Lipschitz envelope}
\label{sec:metric-envelope}

Kleiner--Lott's local-collapse Lemma 6.1 (printed page 41, PDF page 36)
gives an elementary metric lemma that can be checked independently of
PC. Let $(X,d)$ be a metric space, $C>0$, and $\ell,u:X\to(0,\infty)$.
There exists a $C$-Lipschitz function $r:X\to(0,\infty)$ with
$\ell\leq r\leq u$ if and only if
\begin{equation}
 \ell(p)-C d(p,q)\leq u(q)\qquad\text{for all }p,q\in X.
 \label{eq:metric-envelope-condition}
\end{equation}
For necessity, $\ell(p)\leq r(p)\leq r(q)+Cd(p,q)\leq
u(q)+Cd(p,q)$. For sufficiency define, pointwise,
\[
 r(q)=\sup_{p\in X}\{\ell(p)-Cd(p,q)\}.
\]
For each $q$, the set is nonempty by $p=q$, and it is bounded above
by $u(q)$ by \eqref{eq:metric-envelope-condition}. Its real supremum
therefore exists as a finite number and satisfies
$0<\ell(q)\leq r(q)\leq u(q)$. The triangle inequality gives
\[
 \ell(p)-Cd(p,q)
 \leq \ell(p)-Cd(p,q')+Cd(q,q')
 \leq r(q')+Cd(q,q').
\]
Taking the supremum gives $r(q)\leq r(q')+Cd(q,q')$; exchanging $q,q'$
gives $|r(q)-r(q')|\leq Cd(q,q')$. For $X$ empty the assertion is
vacuous; no global nonemptiness premise is needed if the construction
and its boundedness proofs are pointwise. A Lean implementation must
supply the nonempty and bounded-above side conditions to each real
supremum rule, with the positivity coercions recorded.

This proves a Lipschitz envelope, not a smooth function or a volume
comparison. The corrected Corollary 6.5 is a separate geometric target
requiring scale bounds, slack and smoothing. The author correction
changes its interval to $(0,1+2\Lambda]$; no smoothness follows from
the preceding supremum construction. The full smoothing proof remains
a later source/proof-expansion task. MG02 should first be compared with
existing library extension/envelope results before implementing it.
% END migrated B058

\paragraph{Consumer of the envelope (MC19).}
In the collapse route, lower and upper scale functions must first satisfy
\eqref{eq:metric-envelope-condition} with one uniform $C$. Only then does
the envelope provide a positive scale varying in a controlled way between
nearby centers. The volume-scale comparison and smooth replacement are
separate tasks. In particular the author-corrected Corollary 6.5 is not
proved by the displayed supremum argument.

\section{Compact Gromov--Hausdorff distance}
\label{sec:metric-gh}

\begin{definition}[Hausdorff and Gromov--Hausdorff distances; MC03]
For nonempty compact subsets $A,B$ of a metric space $Z$, set
\[
 d_H^Z(A,B)=\max\left\{\sup_{a\in A}d(a,B),\sup_{b\in B}d(b,A)\right\}.
\]
For nonempty compact metric spaces $X,Y$, define
\[
 d_{GH}(X,Y)=\inf_{Z,\iota,\jmath}d_H^Z(\iota(X),\jmath(Y)),
\]
where $\iota:X\to Z$ and $\jmath:Y\to Z$ are isometric embeddings.
Compact GH convergence means $d_{GH}(X_i,X)\to0$.
\end{definition}

The ambient embeddings must preserve the given distances, not merely
curve lengths. For example the usual round sphere's inclusion in
Euclidean space does not preserve its intrinsic Riemannian distance.
The definitions and this distinction are BBI 7.3.1 and 7.3.10, printed
pages 252 and 254 (PDF pages 267 and 269).

\begin{definition}[Correspondences and distortion; MC04]
A correspondence $\mathcal R\subset X\times Y$ projects onto both $X$
and $Y$. Its distortion is
\[
 \operatorname{dis}(\mathcal R)=
 \sup_{(x,y),(x',y')\in\mathcal R}|d_X(x,x')-d_Y(y,y')|.
\]
For a map $f$, its distortion is the same supremum over its graph,
without requiring surjectivity.
\end{definition}

\begin{theorem}[Correspondence formula; source-qualified]
\label{thm:metric-correspondence}
For nonempty compact spaces,
\[
 d_{GH}(X,Y)=\tfrac12\inf_{\mathcal R}\operatorname{dis}(\mathcal R).
\]
Moreover $d_{GH}$ is a metric on compact spaces up to isometry. If
$S\subset X$ is a finite $\varepsilon$-net with its restricted metric,
then $d_{GH}(X,S)\leq\varepsilon$.
\end{theorem}
\begin{proof}[Proof decomposition]
An embedding with Hausdorff distance less than $a$ pairs points at
ambient distance less than $a$ and gives distortion at most $2a$.
Conversely, if a correspondence has distortion $D$, choose $c>D/2$ and
put, on the disjoint union,
\[
 d(x,y)=\inf_{(u,v)\in\mathcal R}
             \bigl(d_X(x,u)+c+d_Y(v,y)\bigr)
       \quad(x\in X,\ y\in Y).
\]
Keep the original distances within each side. The distortion bound
implies that crossing to the other side and back cannot shorten an
original distance; the triangle inequalities follow. Corresponding
points are at distance at most $c$, so $d_{GH}\leq c$; let $c$ decrease
to $D/2$. Composing correspondences adds their distortion bounds, proving
the triangle inequality. For zero distance, diagonal limits of maps on
a countable dense set and compactness give an onto distance-preserving
map; the surjectivity step uses the covering property. Finally, embed
$S$ and $X$ into $X$. These are BBI Theorems 7.3.25 and 7.3.30 and
Example 7.3.11; the diagonal extension is an implementation obligation.
\end{proof}

\begin{lemma}[Maps and explicit inverse estimates; MC05]
\label{lem:metric-global-approx}
Let $X,Y$ be nonempty compact metric spaces, $\varepsilon>0$, and
$f:X\to Y$ have distortion at most $\varepsilon$, with for each $y$ an
$x$ satisfying $d_Y(y,f(x))\leq\varepsilon$. Then
\[
 d_{GH}(X,Y)\leq\tfrac32\varepsilon.
\]
Choosing such an $x$ as $g(y)$ gives distortion at most $3\varepsilon$
for $g$ and
\[
 d_Y(fg(y),y)\leq\varepsilon,\qquad
 d_X(gf(x),x)\leq2\varepsilon.
\]
Conversely $d_{GH}(X,Y)<a$ gives a map of distortion at most $2a$ whose
image is a $2a$-net. These maps need not be continuous.
\end{lemma}
\begin{proof}
The relation $d_Y(y,f(x))\leq\varepsilon$ is a correspondence of
distortion at most $\varepsilon+2\varepsilon$. For $g$, insert $fg(y)$
and $fg(y')$ in the triangle inequality and then use the distortion of
$f$. Applying that distortion to $gf(x),x$ proves the second inverse
bound. For the converse choose one partner for each $x$ in a
correspondence of distortion less than $2a$. Every $y$ has a partner
$x$; its distance from the selected $f(x)$ is at most that distortion.
This is the estimate in BBI Corollary 7.3.28, with its constants exposed.
\end{proof}

% BEGIN migrated B059
\section{MG03--MG04: quantitative pointed approximations}

Our closed-ball variant of Burago--Burago--Ivanov's Definition 8.1.1
(printed page 272, PDF page 287), for pointed spaces $(X,x)$ and $(Y,y)$,
consists, for $0<\varepsilon<R$, of a map
\[
 f:\overline B_X(x,R)\longrightarrow Y,\qquad f(x)=y,
\]
with pairwise distance errors less than $\varepsilon$ and the open
$\varepsilon$-neighborhood of its image containing
$\overline B_Y(y,R-\varepsilon)$. The map need not be continuous.
Record the closed source ball, the smaller target ball and the exact
basepoint convention. This is approximation data, not a homeomorphism. BBI itself uses open
balls: the conversion, with slack, is proved below rather than treated as
an identical fixed-parameter definition.

MG04 derives restriction, composition, quasi-inverse and rescaling
lemmas with explicit losses in radius and error. Do not assert symmetry
with the same constants or restrict a map to a smaller ball while
silently retaining its old coverage witnesses. For composition, check
that the intermediate images lie in the next map's domain. Rescaling
$d$ by $\lambda$ scales distances by $\lambda$; rescaling a Riemannian
metric by $c$ scales its induced distance by $\sqrt c$, through a
separate inherited adapter. The source's approximation exercises are
proof obligations, not already supplied proofs.
% END migrated B059

\begin{definition}[Closed-ball approximation; MC06]
\label{def:metric-approx}
Write $\mathcal A_{R,\varepsilon}(f:(X,p)\to(Y,q))$ when
$0<\varepsilon<R$, $f:\overline B_X(p,R)\to Y$, $f(p)=q$, and
\begin{align*}
 |d_Y(f(x),f(x'))-d_X(x,x')|&<\varepsilon
       &&(x,x'\in\overline B_X(p,R)),\\
 \forall y\in\overline B_Y(q,R-\varepsilon)\quad
 \exists x\in\overline B_X(p,R)&:\ d_Y(y,f(x))<\varepsilon.
\end{align*}
The first requirement is a pairwise strict bound; its supremum need
only be at most $\varepsilon$. Whenever a source requires strictly
smaller supremum distortion, enlarge the error. We always retain the
actual map and coverage witnesses, not merely a truth value.
\end{definition}

In particular $|d_Y(q,f(x))-d_X(p,x)|<\varepsilon$. We can increase
the error while keeping it below $R$, since the target ball then shrinks.
The following constants are deliberately generous and proved here;
they are not numerical assertions attributed to a source exercise.

\begin{lemma}[Restriction with slack; MC07]
\label{lem:metric-restrict}
If $f$ satisfies $\mathcal A_{R,\varepsilon}$ and
$2\varepsilon<s\leq R$, its restriction to the closed $s$-ball satisfies
$\mathcal A_{s,2\varepsilon}$.
\end{lemma}
\begin{proof}
For $d_Y(q,y)\leq s-2\varepsilon$, original coverage applies. Its witness
$x$ obeys
\[
 d_X(p,x)<d_Y(q,f(x))+\varepsilon
       <d_Y(q,y)+2\varepsilon\leq s.
\]
It therefore survives the restriction. The distortion and coverage
errors are still less than $\varepsilon$, hence less than
$2\varepsilon$. Naive restriction without the radius loss has no such
witness guarantee.
\end{proof}

\begin{lemma}[Composition on a controlled domain; MC08]
\label{lem:metric-compose}
Suppose $f:(X,p)\to(Y,q)$ satisfies $\mathcal A_{R,\varepsilon}$ and
$g:(Y,q)\to(Z,z)$ satisfies $\mathcal A_{S,\eta}$. If
\[
 2(\varepsilon+\eta)<s\leq R,\qquad s+\varepsilon\leq S,
\]
then $g\circ f$ on $\overline B_X(p,s)$ satisfies
$\mathcal A_{s,2(\varepsilon+\eta)}$.
\end{lemma}
\begin{proof}
The radial estimate puts $f(\overline B_X(p,s))$ inside the domain of
$g$. Distortions add to less than $\varepsilon+\eta$. For a target
point $w$ of radius at most $s-2\varepsilon-2\eta$, first choose a
$g$-witness $y$; it has radius strictly less than $s-2\varepsilon$.
Then choose an $f$-witness $x$, which has radius less than $s$ by the
preceding argument. The total coverage error is less than
$\eta+(\varepsilon+\eta)=\varepsilon+2\eta$, using the distortion
of $g$ between $f(x)$ and $y$. Both points are in its domain. This is
less than the asserted $2(\varepsilon+\eta)$.
\end{proof}

\begin{lemma}[A local quasi-inverse; MC09]
\label{lem:metric-inverse}
If $\mathcal A_{R,\varepsilon}(f)$, $4\varepsilon<s$ and
$s+\varepsilon\leq R$, there is
$g:\overline B_Y(q,s)\to X$, $g(q)=p$, satisfying
$\mathcal A_{s,4\varepsilon}$ in the reverse direction. Its distortion
is less than $3\varepsilon$, and
\[
 d_Y(fg(y),y)<\varepsilon
 \quad(y\in\overline B_Y(q,s)),\qquad
 d_X(gf(x),x)<2\varepsilon
 \quad(x\in\overline B_X(p,s-4\varepsilon)).
\]
\end{lemma}
\begin{proof}
Coverage selects $g(y)$ in the original $R$-ball; take $g(q)=p$.
Two coverage errors and the original distortion give $3\varepsilon$.
For $x$ in the displayed smaller ball, $f(x)$ has radius less than
$s-3\varepsilon$, so $g(f(x))$ is defined. The original distortion
and $d_Y(fgf(x),f(x))<\varepsilon$ give the $2\varepsilon$ bound.
Thus every such $x$ is within $4\varepsilon$ of the image of $g$.
This proves the required covering condition as well as the inverse
estimates. It does not give symmetry at the original error.
\end{proof}

\begin{lemma}[Scaling and recentering; MC10]
\label{lem:metric-scale-recenter}
For $\lambda>0$, the same map satisfying $\mathcal A_{R,\varepsilon}$
satisfies $\mathcal A_{\lambda R,\lambda\varepsilon}$ after both
distances are multiplied by $\lambda$. If $a\in X$,
$D=d_X(p,a)$, $3\varepsilon<s$ and $D+s\leq R$, restriction of $f$ to
$\overline B_X(a,s)$ satisfies $\mathcal A_{s,3\varepsilon}$ with
basepoints $a,f(a)$.
\end{lemma}
\begin{proof}
Scaling multiplies every distance, radius and error by $\lambda$.
For recentering, a target point within $s-3\varepsilon$ of $f(a)$
has distance less than $D+s-2\varepsilon\leq R-\varepsilon$ from $q$.
An original coverage witness $x$ satisfies
$d_X(a,x)<d_Y(f(a),y)+2\varepsilon<s$, so it lies in the new source
ball. Distortion is unchanged. The basepoint is exactly $f(a)$; replacing
it by another nearby point needs a further estimate.
\end{proof}

\section{Pointed convergence and the source conventions}
\label{sec:metric-conventions}

\begin{definition}[Pointed GH convergence]
Write $(X_i,p_i)\longrightarrow(X,p)$ if for every $R>0$ and every
$0<\varepsilon<R$, all sufficiently large $i$ admit a map satisfying
$\mathcal A_{R,\varepsilon}((X_i,p_i)\to(X,p))$.
The index threshold may depend on $R,\varepsilon$ but not on a point
of the ball. We require a complete limit, and use proper limits in the
applications here.
\end{definition}

\begin{proposition}[Open balls, closed balls and KL parameters; MC11]
\label{prop:metric-conventions}
On complete limits, this convergence agrees with BBI Definition 8.1.1.
It is equivalent to the existence of an index $i_0$ and, for every
$i\geq i_0$, a whole-space pointed Kleiner--Lott
$\delta_i$-approximation with $0<\delta_i<1$ and $\delta_i\to0$.
The comparison is of eventual convergence notions; fixed-error predicates
are not identified.
\end{proposition}
\begin{proof}
Let $\mathcal O_{R,\varepsilon}$ denote the corresponding open-ball
version: source $B(p,R)$, coverage of $B(q,R-\varepsilon)$ and the same
pointwise estimates. Either an open or a closed approximation on radius
$R$ restricts to either convention on radius $s<R$, with error
$2\varepsilon$, provided $2\varepsilon<s$. The witness calculation in
Lemma~\ref{lem:metric-restrict} puts every witness strictly inside $s$,
even for a target on its closed boundary. For prescribed final radius,
start with a larger source radius. If a strict supremum bound is needed,
use error $3\varepsilon$ instead of $2\varepsilon$.

For $0<\delta<1$, KL Definition 3.2 uses a whole-space pointed map,
\emph{non-strict} distortion bounds $\leq\delta$ on $B(p,\delta^{-1})$,
and distance-to-image bounds $\leq\delta$ on
$B(q,\delta^{-1}-\delta)$. The image need not be closed; an infimum
bound does not give a witness at distance at most $\delta$.
It does give a witness at distance strictly less than $2\delta$.
If $3\delta<R<\delta^{-1}$, restrict the source map to the closed
$R$-ball. For a target $y$ of radius at most $R-3\delta$, such a witness
$x$ has
\[
 d_X(p,x)\leq d_Y(q,f(x))+\delta
       <d_Y(q,y)+3\delta\leq R.
\]
Thus it survives the restriction. Distortion is at most $\delta$ and
coverage error is less than $2\delta$, proving
$\mathcal A_{R,3\delta}$. This slack handles both non-strict inequalities
and a possibly unattained infimum.

Conversely start with
$\mathcal A_{\delta^{-1}+\delta,\delta/4}$. Restrict to radius
$\delta^{-1}$: a witness for a target of radius at most
$\delta^{-1}-\delta$ has radius less than
$\delta^{-1}-\delta/2$, and its coverage error is less than $\delta$.
Extend outside this source ball by the target basepoint. The resulting
whole-space pointed map satisfies the non-strict KL tests.

For the forward convergence implication, for each integer $j\geq2$
choose a tail admitting $\mathcal A_{j+1/j,1/(4j)}$. Choose increasing
thresholds $N_j\geq j$ for these tails. For $i\geq N_2$ let
$j(i)=\max\{j:N_j\leq i\}$ and set $\delta_i=1/j(i)$.
The preceding conversion supplies the KL maps and $\delta_i\to0$.
Conversely, for any fixed $R,\varepsilon$, eventually
$3\delta_i<\min\{R,\varepsilon\}$ and $R<\delta_i^{-1}$.
Apply the KL restriction estimate and enlarge the error to
$\varepsilon$. This proves the all-radii definition.

\end{proof}

The finite prefix matters to the literal statement. Let the first source
be $[-1,1]$ pointed at zero and all later sources and the limit be a
singleton. The sequence converges, but the first source has no KL
approximation to the singleton with $0<\delta<1$: its tested ball
contains both endpoints, whose distance $2$ cannot have distortion
at most $\delta$. We use the source's convergence language in its
eventual topological sense, with this qualification explicit.

The result does not choose one map valid simultaneously for all indices
and all accuracies. Nor does it assert that a small KL parameter is a
symmetric numerical distance; KL explicitly warns otherwise. For
subsequences the restriction, composition and quasi-inverse lemmas keep
the necessary data visible.

\begin{proposition}[From pointed to compact convergence; MC11]
\label{prop:metric-pointed-to-compact}
Let $(X_i,p_i)$ and $(X,p)$ be nonempty compact metric spaces, with
pointed convergence to $(X,p)$. Assume one finite $D\geq0$ satisfies
$\operatorname{diam}X_i\leq D$ for every $i$. Then
$d_{GH}(X_i,X)\to0$.
\end{proposition}
\begin{proof}
For $\varepsilon>0$, choose
$R>\max\{D,\operatorname{diam}X\}+\varepsilon$.
Eventually the pointed approximation on radius $R$ has domain all of
$X_i$ and covers all of $X$ within $\varepsilon$. Its distortion is
at most $\varepsilon$, so Lemma~\ref{lem:metric-global-approx} gives
$d_{GH}(X_i,X)\leq3\varepsilon/2$. Let $\varepsilon$ tend to zero.
\end{proof}

The uniform diameter hypothesis repairs BBI Exercise 8.1.2(1), printed
page 272/PDF page 287 (errata page 9). It cannot be omitted: take a
two-point space $X_i=\{p_i,q_i\}$ with $d(p_i,q_i)=i$, for integers
$i\geq1$. It converges pointedly to a singleton, since every fixed
basepoint ball is eventually a singleton, but its GH distance from that
singleton is $i/2$ by the correspondence formula. Compactness of each
member alone does not prevent points escaping every observed radius.

% BEGIN migrated B060
\section{MG05--MG07: compactness and uniqueness}

Definition 7.4.13 and Theorem 7.4.15 of BBI (printed pages 263--264,
PDF pages 278--279) provide the compact precompactness criterion.
For a family of nonempty compact metric spaces, require a uniform
diameter bound and, for each $\varepsilon>0$, a finite bound
$N(\varepsilon)$ on the size of an $\varepsilon$-net in every member.
The order is ``for each scale, one bound for the family,'' not a bound
chosen after its member. The proof targets are nested finite nets,
bounded distance matrices, a diagonal limit, the zero-distance quotient,
completion, and quantitative nets in the limit that prove compactness
and convergence. The archived proof was read; it is not Lean evidence.

For pointed spaces, MG06 uses complete proper spaces and a bound
$N(R,\varepsilon)$ for every $R>0$ and $\varepsilon>0$, independent
of the family member, on nets in each basepoint $R$-ball. The output
includes a complete proper limit and coherent approximation data on
all radii. The diagonal extraction must handle both radius and accuracy.
Arbitrarily chosen compact limits of individual balls do not themselves
supply a compatible pointed limit. Theorem 8.1.10 (statement on printed page 274, PDF page 289;
proof outline on printed page 275, PDF page 290) gives an outline by double diagonalization; expanding
that outline is a real implementation task.

MG07 separates uniqueness up to a pointed isometry for complete proper
limits (Theorem 8.1.7, printed pages 273--274, PDF pages 288--289) from
preservation of the length-space property (Theorem 8.1.9, printed page
274, PDF page 289). Retain the hypotheses and the all-radii approximation
data. Statements about convergence of the actual closed radius-$R$
balls have additional length-space qualifications in the surrounding
exercises; they are not free consequences for arbitrary pointed spaces.
Several details in these sources are outlines, references or exercises
and must become proved lemmas before promotion.
% END migrated B060

\subsection{Compact extraction as a distance-matrix construction}

\begin{theorem}[Compact precompactness; MC12, source-qualified]
\label{thm:metric-compactness}
Let $(X_i)$ be nonempty compact metric spaces. Suppose there is $D<\infty$
with $\operatorname{diam}X_i\leq D$ for every $i$, and
\begin{gather*}
 \forall\varepsilon>0\ \exists N(\varepsilon)\in\mathbb N\
 \forall i\ \exists S_i\subset X_i:\\
 |S_i|\leq N(\varepsilon),\qquad S_i\text{ an }\varepsilon\text{-net}.
\end{gather*}
There are a strictly increasing subsequence $i_j$ and a nonempty compact
metric space $X_\infty$ with $d_{GH}(X_{i_j},X_\infty)\to0$.
\end{theorem}
\begin{proof}[Expanded construction from BBI Theorem 7.4.15]
For each $i$ enumerate the union of finite $1/k$-nets in nested blocks.
By repeating centers, arrange that the first
$M_k=\sum_{h=1}^k N(1/h)$ entries contain a $1/k$-net, with the same
$M_k$ for all $i$. All pairwise distances lie in $[0,D]$. A diagonal
subsequence makes every matrix entry converge. Relabel this subsequence
by $i$. The resulting function
on the countable index set is a pseudometric: symmetry and the triangle
inequality pass to the limit, but separation can fail.

Take the zero-distance quotient and then its metric completion. For
any fixed index $a$, the covering center among the first $M_k$ entries
has a constant choice along an infinite subsequence. Passing to the
limit shows that the quotient point $a$ lies within $1/k$ of one of
those centers. This remains true in the completion because the union
of these finitely many closed $1/k$-balls is closed. Thus the completion
is totally bounded and complete, hence compact.

Let $F_{i,k}$ and $F_{\infty,k}$ denote the sets of the first $M_k$
centers. The relation pairing equal indices is a correspondence even
when entries repeat or merge in the limit. Its distortion tends to
zero for fixed $k$, since it is a maximum over a finite matrix. Hence
\begin{align*}
 d_{GH}(X_i,X_\infty)&\leq
 \frac{2}{k}+d_{GH}(F_{i,k},F_{\infty,k}),\\
 \limsup_{i\to\infty}d_{GH}(X_i,X_\infty)&\leq\frac{2}{k}.
\end{align*}
First choose $k$ for the desired error, then a sufficiently large
subsequence index. This proves convergence and records the order of
choices needed in a formal proof.
\end{proof}

\subsection{Pointed extraction and what coherence means}

\begin{theorem}[Pointed precompactness; MC13 source specialization]
\label{thm:metric-pointed-compactness}
Let $(X_i,p_i)$ be complete proper pointed metric spaces. Suppose
\begin{gather*}
 \forall R>0\ \forall\varepsilon>0\ \exists N(R,\varepsilon)\in\mathbb N\
 \forall i:\\
 \overline B_{X_i}(p_i,R)
 \text{ has an }\varepsilon\text{-net of size at most }N(R,\varepsilon).
\end{gather*}
A subsequence converges pointedly to a complete proper pointed metric
space $(X_\infty,p_\infty)$. No lower volume bound and no preservation
of dimension are part of this metric theorem.
\end{theorem}
\begin{proof}
Apply Theorem~\ref{thm:metric-eventual-compactness} below with initial
index $I=1$ for every radius and accuracy. The more general construction
also supplies the finite-map details required here.
\end{proof}

\begin{theorem}[Eventual local nets suffice; MC13 refinement]
\label{thm:metric-eventual-compactness}
Let $(X_i,p_i)$ be pointed metric spaces. Assume the following
\emph{eventual net condition}, with centers inside the indicated ball:
\begin{gather*}
 \forall R>0\ \forall\eta>0\ \exists N\in\mathbb N\ \exists I\in\mathbb N\
 \forall i\geq I\ \exists F_i\subseteq\overline B_{X_i}(p_i,R):\\
 |F_i|\leq N,\qquad
 \forall x\in\overline B_{X_i}(p_i,R)\ \exists a\in F_i:
 d_i(x,a)\leq\eta.
\end{gather*}
There are a strictly increasing sequence $i_j$, a complete proper pointed
metric space $(Y,q)$, and, for every $0<\varepsilon<R$, maps witnessing
$\mathcal A_{R,\varepsilon}$ from $(X_{i_j},p_{i_j})$ to $(Y,q)$ for
all sufficiently large $j$. Source completeness, properness and length
structure are not assumed. The convergence assertion uses
Definition~\ref{def:metric-approx}; it does not put a nonproper source
inside the metric space of proper pointed isometry classes.
\end{theorem}
\begin{proof}[Countable construction and quantitative maps]
\emph{1. Fixed labels despite exceptional prefixes.}
For each integer $m\geq1$, set $\eta_m=2^{-m}$ and choose $N_m,I_m$
and finite $\eta_m$-nets of the closed $m$-balls for $i\geq I_m$.
These balls contain $p_i$, so $N_m\geq1$. Repeat centers to get exactly
$N_m$ labels $(m,1),\ldots,(m,N_m)$, and introduce a separate label $0$
for $p_i$. For $i<I_m$, set every level-$m$ point equal to $p_i$;
no net assertion is made at those indices. Write $x_a^i$ for the point
with label $a$. The radius of a level-$m$ point is at most $m$ for
every $i$. Hence the distances between fixed level-$m$ and level-$l$
labels lie in $[0,m+l]$.

Enumerate the countably many pairs of labels. Successively take infinite
subsequences on which each distance converges, and then choose a diagonal
subsequence $i_j$ tending strictly to infinity, with
$i_j\geq\max(j,I_1,\ldots,I_j)$. Each fixed matrix entry converges on
this single subsequence. Put
\[
 d_\infty(a,b)=\lim_{j\to\infty}d_{i_j}(x_a^{i_j},x_b^{i_j}).
\]
Symmetry, nonnegativity, zero diagonal and the triangle inequality pass
to the limit. Quotient the label set by zero distance and complete it.
Call the result $Y$, write $z_a$ for a label's image, and put $q=z_0$.
The $z_a$ are dense in $Y$. Different labels may represent the same point.

\emph{2. Nets in the completion and properness.}
Let $E_m$ consist of label $0$ and all level-$m$ labels, and put
$Z_m=\{z_a:a\in E_m\}$. Fix a label $b$ with $d_Y(q,z_b)<m$.
For all sufficiently large $j$, its representative lies in the source
$m$-ball and the level-$m$ net is available. Therefore
\[
 \min_{a\in E_m}d_{i_j}(x_b^{i_j},x_a^{i_j})\leq\eta_m.
\]
There are finitely many entries in this minimum. Passing to their limits
gives $\min_{a\in E_m}d_Y(z_b,z_a)\leq\eta_m$.
For any $y\in Y$ with $d_Y(q,y)<m$, approximate $y$ by labels whose
limiting radii are also less than $m$. The distance to the finite set
$Z_m$ is continuous, so $d_Y(y,Z_m)\leq\eta_m$, with a minimizing
center in $Z_m$.

Consequently, for any fixed $R<m$, $Z_m$ is an ambient
$\eta_m$-net for $\overline B_Y(q,R)$. Its centers need not lie in
that ball. For each center whose closed $\eta_m$-ball meets
$\overline B_Y(q,R)$, choose one point of that intersection. These
finitely many chosen points form a $2\eta_m$-net inside the ball.
Given any positive accuracy, choose $m>R$ with $2\eta_m$ smaller
than that accuracy. This proves total boundedness. The ball is closed
in the complete space $Y$, hence complete and compact. Thus $Y$ is
proper. Notice that we used nets at all sufficiently fine levels;
one covering estimate at one scale would not suffice.

\emph{3. Actual approximations, including coverage.}
Fix $0<\varepsilon<R$. Choose $m>R+1$ with
$\eta=\eta_m<\varepsilon/6$. For all sufficiently large $j$, the
level-$m$ net is available and its finite matrix, including label $0$,
satisfies
\[
 \bigl|d_{i_j}(x_a^{i_j},x_b^{i_j})-d_Y(z_a,z_b)\bigr|<\eta
 \qquad(a,b\in E_m).
\]
For $u\in\overline B_{X_{i_j}}(p_{i_j},R)$ choose a label
$a(u)\in E_m$ with $d_{i_j}(u,x_{a(u)}^{i_j})\leq\eta$, choosing
label $0$ when $u=p_{i_j}$. Define $f_j(u)=z_{a(u)}$.
It fixes the basepoint exactly, and the triangle inequality gives
\[
 \bigl|d_Y(f_j(u),f_j(v))-d_{i_j}(u,v)\bigr|
 <\eta+2\eta=3\eta<\varepsilon.
\]
To prove coverage, take $y\in\overline B_Y(q,R-\varepsilon)$.
Step~2 gives $a\in E_m$ with $d_Y(y,z_a)\leq\eta$. The source
representative $u=x_a^{i_j}$ satisfies
\[
 d_{i_j}(p_{i_j},u)<d_Y(q,z_a)+\eta
 \leq R-\varepsilon+2\eta<R,
\]
so $u$ lies in the domain of $f_j$. Its chosen label $b=a(u)$ may
differ from $a$, since labels can repeat. Nevertheless
$d_{i_j}(x_a^{i_j},x_b^{i_j})\leq\eta$, whence
$d_Y(z_a,z_b)<2\eta$ by the finite matrix estimate. It follows that
$d_Y(y,f_j(u))<3\eta<\varepsilon$. This proves the required
coverage of the smaller closed target ball, including the case $u=p_{i_j}$.
All choices of $m$ depend only on the requested $R,\varepsilon$;
the previously extracted subsequence and completed space stay fixed.
\end{proof}

This is a written adaptation of the distance-matrix construction in
BBI Theorem 7.4.15 (printed 264/PDF 279) and the pointed outline in
Theorem 8.1.10 (printed 274--275/PDF 289--290). The eventual quantifiers,
absence of source properness, completed-ball proof and $3\eta$ map
estimates are established above, not quoted as a stronger BBI statement.
The union of chosen source nets need not be dense in any whole $X_i$;
density is required only of the labels in their constructed completion.
The source's proper-space setting is recovered by the original MC13
specialization. Formal reuse, indexed choices and kernel verification
remain deferred.

\begin{corollary}[Absorbing a finite prefix when sources are proper]
\label{cor:metric-finite-prefix}
If every $X_i$ is proper and the eventual net condition holds, the
all-index net hypothesis of Theorem~\ref{thm:metric-pointed-compactness}
holds as well.
\end{corollary}
\begin{proof}
Fix $R,\eta>0$ and choose the eventual bound $N$ and threshold $I$.
For each of the finitely many $i<I$, compactness of the closed $R$-ball
gives a finite $\eta$-net with $N_i$ centers inside that ball. Take
\[
 N^*(R,\eta)=\max\bigl(\{N\}\cup\{N_i:1\leq i<I\}\bigr).
\]
This also covers the empty-prefix case. Proper spaces are complete by
the earlier metric toolkit, so the remaining MC13 hypothesis holds.
The threshold may depend on $R,\eta$; no one universal tail for all
scales is claimed. For complete connected Riemannian sources, properness
is an inherited Hopf--Rinow consequence whose release binding remains
pending, not a new foundational theorem to implement here.
\end{proof}

\subsection{Uniqueness, length structure and actual balls}

\subsubsection{Approximate midpoints and length closure.}
The following proofs expand BBI Theorems 2.4.16, 7.5.1 and 8.1.9
\cite{BuragoBuragoIvanov}. The author errata remove local compactness
from the discussion preceding 2.4.16; only completeness is needed
for the approximate-midpoint criterion. All distances here are finite.
We call $z$ an $h$-midpoint of $a,b$ when
\[
 \max\{d(a,z),d(b,z)\}\leq\tfrac12d(a,b)+h.
\]
The triangle inequality then also gives each lower bound
$d(a,z),d(b,z)\geq d(a,b)/2-h$. Thus this convention implies
the convention in BBI Lemma 2.4.10 with error $2h$; the error parameters are
not identified literally.

\begin{lemma}[Midpoint transfer on a padded domain; MC21]
\label{lem:metric-midpoint-transfer}
Let $(Z,o)$ be a pointed length space, $(Y,q)$ a pointed metric
space, and $a,b\in Y$. Put
\[
 M=\max\{1,d_Y(q,a),d_Y(q,b)\},\qquad L=d_Y(a,b),
 \qquad R=4M+4.
\]
If $0<\varepsilon<1/4$ and
$\mathcal A_{R,\varepsilon}(f:(Z,o)\to(Y,q))$, there is $z\in Y$
such that
\[
 d_Y(a,z),d_Y(b,z)<L/2+4\varepsilon.
\]
No completeness or properness is required for either space here,
and $f$ need not be continuous.
\end{lemma}
\begin{proof}
For $a=b$ take $z=a$. Otherwise $a,b$ belong to the covered
closed $(R-\varepsilon)$-ball. Choose witnesses $a',b'\in Z$
with $d_Y(f(a'),a),d_Y(f(b'),b)<\varepsilon$. The basepoint and
distortion estimates give
\[
 d_Z(o,a'),d_Z(o,b')<M+2\varepsilon,
 \qquad d_Z(a',b')<L+3\varepsilon.
\]
Choose a rectifiable path $\gamma$ from $a'$ to $b'$ of length
less than $d_Z(a',b')+\varepsilon<L+4\varepsilon$.
Every point of $\gamma$ has radius less than
\[
 M+2\varepsilon+L+4\varepsilon
 \leq3M+6\varepsilon<R.
\]
Thus all needed points lie in the actual domain of $f$, even
though a restricted ball is not assumed to be a length space.
The continuous function
$d_Z(a',\gamma(t))-d_Z(b',\gamma(t))$ changes sign between the
endpoints. At a zero, call it $z'$, the two distances are equal.
Each is at most half the path length: their sum is bounded by the
sum of the lengths of the two subpaths. This also holds for a
constant path if $a'=b'$. Set $z=f(z')$. For either endpoint,
coverage and distortion contribute less than $2\varepsilon$;
hence
\[
 d_Y(a,z)<\varepsilon+\varepsilon+\tfrac12(L+4\varepsilon)
          =L/2+4\varepsilon,
\]
and likewise for $b$. No curve is mapped continuously by $f$;
only its selected points are evaluated.
\end{proof}

\begin{theorem}[Complete approximate-midpoint criterion; MC22]
\label{thm:metric-approximate-midpoints}
Let $Y$ be a complete metric space with finite distances. Suppose
for every $a,b\in Y$ and $h>0$ there is an $h$-midpoint.
For every $a,b$ and $\tau>0$ there is a path
$\gamma:[0,1]\to Y$ from $a$ to $b$ with
\[
 \operatorname{Lip}(\gamma)\leq d_Y(a,b)+\tau/2,
 \qquad L(\gamma)<d_Y(a,b)+\tau.
\]
In particular $Y$ is a length space. Local compactness is not an
assumption.
\end{theorem}
\begin{proof}
Write $D=d_Y(a,b)$. For $a=b$ the constant path suffices.
Assign $\gamma(0)=a$, $\gamma(1)=b$. Inductively keep all values
at dyadic level $n$ and choose the values at the new midpoints
using error
\[
 h_n=\tau\,2^{-2n-3}\qquad(n\geq0).
\]
Let $B_n$ be $2^n$ times the largest adjacent distance at level
$n$. There are finitely many edges, so this maximum is defined.
The new midpoint inequalities imply
\[
 B_{n+1}\leq B_n+2^{n+1}h_n
          =B_n+\tau\,2^{-n-2},\qquad B_0=D.
\]
Summing this geometric series yields
\[
 B_n\leq D+\frac{\tau}{2}(1-2^{-n})\leq D+\tau/2=:C.
\]
Any two dyadic parameters occur at a common level. Summing the
adjacent distances between them gives
$d_Y(\gamma(s),\gamma(t))\leq C|s-t|$.
For each $t\in[0,1]$, choose dyadic $t_j\to t$. Their images
are Cauchy, so completeness supplies a limit. The same Lipschitz
bound shows that this limit is independent of the approximating
sequence. It agrees with every previously assigned value and is
$C$-Lipschitz by passage to the limit. This is the dense-domain
extension argument of BBI Proposition 1.5.9, supplied here explicitly.
For every partition of $[0,1]$, the sum of successive image distances
is at most $C$ times the sum of parameter increments, hence at most
$C$. Taking the supremum proves $L(\gamma)\leq C<D+\tau$.
The lower bound $L(\gamma)\geq D$ and arbitrary $\tau$ identify
the given metric with its induced length metric.
\end{proof}

\begin{corollary}[Proper midpoint limits are geodesic; MC23]
\label{cor:metric-proper-midpoint-geodesic}
A proper metric space with an $h$-midpoint for every pair and every
$h>0$ is geodesic. In particular the proper limits obtained below
admit a path $\gamma$ for every $a,b$ with
\[
 \gamma(0)=a,\quad\gamma(1)=b,\qquad
 d(\gamma(s),\gamma(t))=|s-t|d(a,b)\quad(s,t\in[0,1]).
\]
\end{corollary}
\begin{proof}
For fixed $a,b$, put $D=d(a,b)$ and choose $1/j$-midpoints
$z_j$, $j\geq1$. They lie in the compact closed ball
$\overline B(a,D/2+1)$. A convergent subsequence has limit $z$
with $d(a,z),d(b,z)\leq D/2$. The triangle inequality forces
both distances to equal $D/2$, so every pair has an exact midpoint.
A proper metric space is complete: any Cauchy sequence is bounded,
has a subsequence converging in a compact closed ball, and then
converges itself.
Repeat MC22's dyadic construction with exact midpoints and zero
error. At level $n$, each adjacent distance equals $D/2^n$, and
the extension is $D$-Lipschitz. For $0\leq s\leq t\leq1$,
\[
 D\leq d(a,\gamma(s))+d(\gamma(s),\gamma(t))+d(\gamma(t),b)
 \leq Ds+d(\gamma(s),\gamma(t))+D(1-t).
\]
Together with the Lipschitz upper bound this gives the displayed
equality. If $D=0$, use the constant path. This proof does not
assume that a closed ball with its restricted metric is a length space.
\end{proof}

\subsubsection{Uniqueness with actual coverage witnesses.}

\begin{theorem}[Proper pointed limits are isometric; MC24]
\label{thm:metric-proper-limit-uniqueness}
If one sequence $(Z_i,o_i)$ converges in the convention of MC06
to both complete proper pointed spaces $(X,p)$ and $(Y,q)$, there
is a surjective distance-preserving map $F:X\to Y$ with $F(p)=q$.
The source spaces need not be proper, complete or length spaces.
The conclusion is existence of a pointed isometry, not uniqueness
of that isometry as a function.
\end{theorem}
\begin{proof}
For $j\geq1$, set
\[
 e_j=\frac1{100(j+1)},\qquad R_j=4j+4,\qquad
 \delta_j=10e_j.
\]
Choose increasing source indices $i_j$ large enough that the
same source member admits both approximations
$\mathcal A_{R_j,e_j}(Z_{i_j}\to X)$ and
$\mathcal A_{R_j,e_j}(Z_{i_j}\to Y)$.
MC09 inverts the first on radius $2j+2$ with error $4e_j$.
Compose with the second by MC08, using output radius $j$.
The domain conditions are
$4e_j<2j+2$, $2j+2+e_j\leq R_j$,
$10e_j<j$ and $j+4e_j\leq R_j$.
We obtain maps
\[
 F_j:\overline B_X(p,j)\to Y,\qquad
 \mathcal A_{j,\delta_j}(F_j),\qquad \delta_j\longrightarrow0.
\]

Finite $1/k$-nets in the compact closed $k$-balls of $X$, together
with $p$, give a countable dense set $A$. Repetitions are allowed
when enumerating $A$. For each fixed $a\in A$ and all sufficiently
large $j$, $F_j(a)$ is defined and belongs to the compact ball
$\overline B_Y(q,d_X(p,a)+1)$. Successive subsequences and a
diagonal choice give $j_\ell\to\infty$ such that every
$F_{j_\ell}(a)$ converges. The limit map on $A$ preserves distances
because the distortion tends to zero, and sends $p$ to $q$.
By the complete dense-domain extension argument in MC22, it extends
to a distance-preserving map $F:X\to Y$.

It remains to prove surjectivity; distance preservation alone
would not do this. Fix $y\in Y$. For all sufficiently large
$j=j_\ell$, coverage gives $x_j\in\overline B_X(p,j)$ such that
$d_Y(F_j(x_j),y)<\delta_j$. Radial distortion gives
\[
 d_X(p,x_j)<d_Y(q,y)+2\delta_j\leq d_Y(q,y)+1.
\]
Properness of $X$ supplies a further subsequence with $x_j\to x$.
For any fixed $a\in A$, both $a$ and $x_j$ lie in the domain of
$F_j$ for large $j$, and
\[
\begin{split}
 d_Y(F_j(x_j),F(x))
 &\leq d_X(x_j,a)+\delta_j\\
 &\quad+d_Y(F_j(a),F(a))+d_X(a,x).
\end{split}
\]
Its upper limit is at most $2d_X(a,x)$. Let $a$ approach $x$
through $A$; then $F_j(x_j)\to F(x)$. Since these same images
approach $y$, we have $F(x)=y$. This proves onto-ness without
assuming continuity of any approximation map or convergence of
images at moving points without an estimate. Properness of $Y$
is used for the first diagonal selection and properness of $X$
for the coverage preimages. This proves the present two-proper-limit
contract; the stronger one-proper-limit form in BBI 8.1.7 is not
asserted here.
\end{proof}

\begin{theorem}[Limit identification and length closure; MC14]
\label{thm:metric-limit-structure}
If the same sequence converges to two complete proper pointed spaces,
there is a pointed isometry between them. If each source is a length
space and the limit is complete, the limit is a length space. In the
proper case it is therefore geodesic.
\end{theorem}
\begin{proof}
MC24 proves the pointed isometry assertion. For length closure, fix
$a,b$ in the limit and $h>0$. Use MC21's fixed radius $R=4M+4$
and choose $0<\varepsilon<\min\{1/4,h/4\}$. Convergence supplies
such an approximation from a sufficiently late length source, so
MC21 gives an $h$-midpoint. MC22 then constructs paths in the
complete limit with lengths arbitrarily close to the endpoint
distance. If the limit is proper, MC23 gives minimizing geodesics.
Neither source properness nor continuity of the approximation maps
was used for length closure. These are written metric proofs;
Lean declarations and library reuse remain deferred.
\end{proof}

\subsubsection{Closed-ball consequence.}

\begin{proposition}[Closed balls in proper length limits; MC15]
\label{prop:metric-closed-balls}
If proper pointed spaces $(X_i,p_i)$ converge to a proper length space
$(X,p)$, then for every $r>0$ their closed $r$-balls converge in compact
GH distance to $\overline B_X(p,r)$, with restricted ambient distances.
The source balls need not themselves be length spaces.
\end{proposition}
\begin{proof}
Use an approximation on a radius larger than $r+1$, with error
$0<\varepsilon<r/4$. For each point of the source closed $r$-ball,
its image has radius less than $r+\varepsilon$. Move that image toward
$p$ along a minimizing geodesic, if needed, into the closed $r$-ball,
a distance at most $\varepsilon$. The resulting map has distortion
less than $3\varepsilon$.

For $y\in\overline B_X(p,r)$, first move radially to $y'$ of radius
at most $r-3\varepsilon$, a distance at most $3\varepsilon$. A coverage
witness for $y'$ has source radius less than $r-\varepsilon$, so belongs
to the source closed $r$-ball. Its adjusted image is within
$2\varepsilon$ of $y'$, hence within $5\varepsilon$ of $y$.
Lemma~\ref{lem:metric-global-approx} gives GH distance at most
$15\varepsilon/2$; let $\varepsilon\to0$. Properness makes both
closed balls compact, and MC23 gives the minimizing radial geodesics,
using approximate midpoints from the length property. This supplies the relevant qualified
ball-convergence argument behind BBI Exercises 8.1.3--8.1.4.
\end{proof}

The length qualification cannot be deleted: take two-point spaces with
basepoint one point and other point at distance $1+1/i$. They converge
pointedly to distance $1$. Their closed unit balls are singletons,
whereas the limit's closed unit ball has two points. The theorem also
does not turn unpointed isometries of each radius ball into a single
basepoint-preserving global isometry.

\section{A collapse example and the local model extraction input}
\label{sec:metric-collapse-example}

\begin{example}[Shrinking a compact factor; MC16]
Let $Y,F$ be nonempty compact metric spaces. Give $Y\times F$ the metric
\[
 d_t((y,a),(y',a'))=
       \sqrt{d_Y(y,y')^2+t^2d_F(a,a')^2},\qquad t>0.
\]
Projection onto $Y$ is onto and its distortion is at most
$t\operatorname{diam}F$, since
$0\leq\sqrt{u^2+v^2}-u\leq v$. Its graph is a correspondence, so
\[
 d_{GH}((Y\times F,d_t),Y)\leq\tfrac12 t\operatorname{diam}F.
\]
For flat tori, $T^2\times (tS^1)$ therefore converges to $T^2$ as
$t\to0$. The product metric is the Pythagorean convention in KL
Section 2.2. The estimate is proved here. Even these smooth complete
three-dimensional examples have a two-dimensional limit; metric
convergence cannot certify a diffeomorphism or a Seifert fibration.
\end{example}

\begin{theorem}[Curvature-to-model extraction input; MC18]
\label{thm:metric-kl-extraction}
Fix a positive integer $n$. Let $(X_i,p_i)$ be complete pointed length
spaces, and let $c_i,r_i>0$ with $c_i\to0$, $r_i\to\infty$. Suppose
the open region $B_{X_i}(p_i,r_i)$ has Alexandrov curvature at least
$-c_i$ and Hausdorff dimension at most $n$. A subsequence converges pointedly to
a nonnegatively curved Alexandrov space of dimension at most $n$.
\end{theorem}

This is KL Lemma 3.10, printed page 24/PDF page 19, in its standing
finite-dimensional Alexandrov setting. Its written conditional proof is assembled in
Corollary~\ref{cor:alexandrov-assembled-extraction}. Uniform local
nets, lower comparison and dimension require separate arguments; Chapter~\ref{ch:alexandrov-geometry} owns those steps and
the local completeness safeguards in KL Section 3.3. When applied to
rescaled Riemannian manifolds, sectional curvature, the domain on which
it is bounded, and dimension must first be translated to this input.
A sequence of unrelated small neighborhoods with no radii tending to
infinity does not meet this statement.

\paragraph{Local-to-global compactness bridge (audit A40).}
The hypotheses do not assert that each whole $X_i$ is proper. Nor do
local curvature bounds on growing balls give uniform nets for every
radius in every member from the first index. Before applying
Theorem~\ref{thm:metric-pointed-compactness}, the proof must obtain
\[
 \forall R,\varepsilon>0\ \exists N,I\ \forall i\geq I:\quad
 \overline B_{X_i}(p_i,R)\text{ has an }\varepsilon\text{-net of size }\leq N.
\]
Once these bounds have been produced,
Theorem~\ref{thm:metric-eventual-compactness} supplies extraction without
global source properness. Alternatively,
Corollary~\ref{cor:metric-finite-prefix} proves the finite-prefix route
when properness has been independently established. These two metric
arguments are now written; neither manufactures the net hypothesis
from local curvature. Closed restricted balls are not silently declared
to be length spaces or Alexandrov spaces.

For the distinction, attach countably many unit segments at the point
$2i$ of a half-line, with its path metric and basepoint zero. This is a
complete length space; its open radius-$i$ region is an interval, of
curvature at least $-1/i$ and dimension one. But its closed radius-$(2i+1)$
ball contains infinitely many endpoints at pairwise distance $2$ and is
not compact. Thus the local hypotheses do not give global source
properness. Chapter~\ref{ch:alexandrov-geometry} now supplies local
structure AC07 through AC29 and AC33/AC17. Its covering argument
uses AC02's source-qualified comparison input and then feeds MC13.
MC14 proves length closure and geodesicity. Chapter~\ref{ch:alexandrov-geometry} supplies lower-comparison
stability through AC34--AC37, conditional on the recorded local
Toponogov input. AC38--AC40 prove dimension closure from the
quantitative covering rate, and AC41 assembles this written bridge.
BBI Theorem 10.7.2 (printed 376--377/PDF 391--392) uses the
Section 10.8 dimension-closure results; our proof keeps the
growing-region domains and derives the Hausdorff bound directly. Splitting and fibration compatibility belong
to later selected consumers.

\section{From controlled smooth data to metric approximations}
\label{sec:metric-smooth-bridge}

\begin{proposition}[A sufficient metric bridge; MC17]
\label{prop:metric-smooth-bridge}
Let $R>0$, $0<\varepsilon<R$, $0\leq a<1$, and suppose
$f:\overline B_X(p,R)\to Y$ fixes basepoints and obeys the
\emph{ambient-distance} inequalities
\[
 (1-a)d_X(x,x')\leq d_Y(f(x),f(x'))\leq(1+a)d_X(x,x').
\]
Assume $2aR<\varepsilon$ and that every point of
$\overline B_Y(q,R-\varepsilon)$ is within a distance strictly less
than $\varepsilon$ of the image. Then $f$ satisfies
$\mathcal A_{R,\varepsilon}$.
\end{proposition}
\begin{proof}
Two source points have distance at most $2R$, so the distortion is at
most $2aR<\varepsilon$. The other conditions are the assumed coverage
and exact basepoint. This isolates the two outputs a smooth-domain
adapter must actually establish.
\end{proof}

For a diffeomorphism $F:U\to V$ between Riemannian domains with
$0\leq\theta<1$, a tensor estimate
\[
 (1-\theta)g\leq F^*h\leq(1+\theta)g
\]
implies the curve-length bounds with factors
$\sqrt{1-\theta}$ and $\sqrt{1+\theta}$. Taking infima over curves
\emph{remaining in the domains} gives the corresponding intrinsic
$U,V$ distance bounds. To substitute ambient manifold distances, prove
that the relevant minimizing curves, or sufficiently short competitors,
stay in those domains. A padded exhaustion, an exit-distance estimate,
or an appropriate convexity statement can supply this proof, but it
must be specified. Separately show that the image contains or densely
covers the required target ball. Convergence in coordinate charts alone
does not state either assertion.

Thus a sufficient smooth-to-pointed-limit contract supplies, for every
fixed $R$, pointed maps on the ambient closed $R$-balls for all large
$i$, numbers $a_i(R)\to0$ satisfying the ambient inequalities above,
and coverage of each smaller target ball with errors tending to zero.
Proposition~\ref{prop:metric-smooth-bridge} then proves pointed metric
convergence. The production adapter must be bound to the accepted
smooth convergence and exhaustion convention, using inherited geometry.
It does not reconstruct Hamilton's compactness theorem.

\paragraph{Normalization.}
Multiplying a distance by $\lambda>0$ multiplies lengths, radii and
approximation errors by $\lambda$. Multiplying a Riemannian tensor by
$c>0$ multiplies curve speeds and distances by $\sqrt c$. Consequently
the tensor $r^{-2}g$ has distance $r^{-1}d_g$; the tensor $t^{-1}g(t)$
has distance $t^{-1/2}d_{g(t)}$. Curvature and volume scaling belong to
the inherited normalization bridge, not to the definition of GH
convergence. The positive factor and the particular time slice must
remain explicit for the long-time consumers.

\paragraph{Comparison with smooth precompactness (MC20).}
KL Definition 3.4 builds a GH approximation into its definition of a
$\delta$--$C^K$ approximation; it cannot be used to infer that GH
condition from a tensor estimate alone. KL Lemma 3.5, printed page 22/
PDF page 17, assumes fixed $v,r>0$, dimension $n$, and one function
$A:(0,\infty)\to(0,\infty)$ for the family. For complete pointed
$C^{K+2}$ Riemannian manifolds, with the source's standing integer
$K\geq10$, its hypotheses are
\[
 \operatorname{vol}B(p,r)\geq v,\qquad
 |\nabla^k\mathrm{Rm}|\leq A(R)\text{ on }B(p,R)
 \quad(0\leq k\leq K,\ R>0).
\]
Its conclusion is pointed $C^K$ precompactness. This is a
source-qualified smooth compactness statement to compare with inherited
interfaces, not a new theorem proved in this chapter. Unlike metric
precompactness, it retains a noncollapse hypothesis and derivative
bounds. Hyperbolic thick-limit identification uses such smooth data
plus the flow and curvature arguments of its own chapter. The thin
branch can have the dimension loss of the preceding example.

% BEGIN migrated B061
\section{MG08--MG09: comparison and foundation boundaries}

Kleiner--Lott's Sections 3.1--3.3 (printed pages 21--24, PDF pages
16--19) use their own pointed $\delta$-approximation conventions.
MG08 now has the written convergence translation in
Proposition~\ref{prop:metric-conventions}; its formalization remains open. Equivalent notions of convergence do not
make two fixed-$\delta$ predicates identical. MG09 is the separate
bridge from controlled smooth metric convergence to metric
approximations on smaller balls. Chart convergence alone does not
supply ambient ball containment; record the domain and distance
comparison hypotheses. Audit the completed geometric foundation for
this bridge instead of rebuilding flow compactness here.

The July 6, 2024 author errata have been retrieved, read in full and
retained with a SHA-256 record; the earlier retrieval check is closed.
Lemmas~\ref{lem:metric-length-lsc} and~\ref{lem:metric-curve-compactness}
repair the length and curve compactness arguments.
Proposition~\ref{prop:metric-pointed-to-compact} retains the missing
uniform diameter hypothesis. The finite-net convention already includes
the finiteness required by the correction to Exercise 1.6.2, printed
page 13/PDF page 28. The Hausdorff-completeness proof on printed
page 253/PDF page 268 quantifies over every point; the finite-subset
construction on printed page 263/PDF page 278 converges to $S$, not $X$.

\code{bbi_errata_review.csv} and \code{reference_checks_revision59.md}
record exact correction locators, current consumers and a later-chapter
watchlist for quotient, covering, comparison, strainer and hyperbolic
arguments. Their book proofs require review when selected. The log also
records an apparent further typo in the Wronskian calculation on errata
page 14; that replacement proof is not imported here. Compactness and
smooth-domain proof obligations remain open. Earlier source checks in
\code{reference_checks_revision56.md} and
\code{reference_checks_revision58.md} remain historical evidence.

The next mathematical task is to supply Chapter~\ref{ch:alexandrov-geometry}'s
comparison-to-packing and local-model arguments, using this chapter's
explicit metric inputs. Expand the remaining diagonal and smooth-domain
proof obligations as those consumers require. MG01 environment pinning
remains deferred, and no proposed name below is a verified Lean declaration.
% END migrated B061

\paragraph{Current refinement of the preserved roadmap.}
The general pointed extraction construction is now written in MC13.
References in the preserved implementation roadmap to remaining diagonal
work concern formal binding. Revision 68 supplies MC14 identification
and length closure; Chapter 4 supplies AC07 through AC29 and AC33/AC17.
AC34--AC41 supply comparison, dimension and compactness assembly; smooth-domain producers retain their obligations.

\section{Results, consumers and remaining proof work}
\label{sec:metric-consumers}

The records MC01--MC24 in \code{metric_geometry_contracts.csv}
track statements, sources, dependencies and evidence. MC13 supplies
pointed extraction; MC14 and MC21--MC24 supply identification,
length closure and geodesics. These refine MG02--MG09 without
adding verified production declarations. Chapter 4 supplies local
structure for packing; formal compactness binding and smooth-domain
arguments retain their obligations.

\begingroup
\small
\begin{longtable}{P{0.22\textwidth}P{0.34\textwidth}P{0.34\textwidth}}
Output & Consumer & Additional input still required\\\hline
\endfirsthead
\hline Output & Consumer & Additional input still required\\\hline
\endhead
Nets and pointed extraction (MC02, MC12--MC14) &
Alexandrov and local-collapse chapters; local-model subsequences &
Chapter 4 local covering uses its stated comparison input;
AC37 writes lower-comparison stability; its local Toponogov input
remains source-qualified.\\
Pointed maps and source translation (MC06--MC11, MC15) &
KL local models and compatible-fibration chapter &
Chapter 4 writes metric strainer and compatibility arguments;
Chapter 14 still needs smooth coordinates and fibration compatibility.\\
Positive Lipschitz scale (MC19) &
Collapse hypotheses and rescaled neighborhoods &
Lower/upper-scale compatibility, corrected smoothing and volume bounds.\\
Smooth-to-metric and scaling bridge (MC17, MC20) &
Hyperbolic thick output and long-time normalization &
Inherited metric API, actual domain coverage, smooth compactness and hyperbolicity.\\
Collapse example and curvature input (MC16, MC18) &
Separation of thin-limit geometry from graph-manifold topology &
Local topology, fibration construction and global reconstruction.\\\hline
\end{longtable}
\endgroup

No result here supplies the exceptional terminal branch or changes the
late-thick/thin-or-terminal outcome. The independent metric material
can be elaborated mathematically while the PC release is pending;
new implementation remains subject to the deferred reuse audit.

% BEGIN migrated B057
\section{MG01: deferred implementation preparation}

Environment pinning is deferred by the user. When implementation is
resumed, a generic metric lemma can live in an isolated project with an exact
Lean toolchain, mathlib commit, imports and axiom evidence. This neither
accepts M2A/K1 nor imports the unfinished PC snapshot. Before creating a
new declaration, inspect the pinned metric, supremum, Lipschitz,
compactness and Gromov--Hausdorff APIs; record direct reuse, a small
adapter, or a genuine leaf gap. Proposed module names are not evidence
that an API exists or is absent. Revision 56 did not locate a pinned mathlib project under the GC
directory; no new environment search, setup or compilation is requested
for the present mathematical chapter expansion.

Detailed leaf records are in \code{metric_geometry_targets.csv}. The
compactness results in this chapter are specifications and proof decompositions;
MG02 and several quantitative metric lemmas now have written proofs. Smooth manifolds, tensors, connections,
curvature and volume are still inherited-foundation candidates, not
new projects inside these work streams.
% END migrated B057

\chapter{Alexandrov geometry and local models}
\label{ch:alexandrov-geometry}

\paragraph{Scope.} The selected comparison and local-model facts consumed by collapse.

\paragraph{Planned sections.}
\begin{enumerate}[leftmargin=*]
\item Comparison conventions.
\item Dimension.
\item Strainers.
\item Splitting.
\item Local models.
\item Consumer-by-consumer source map.
\end{enumerate}

\paragraph{Existing material and present task.} I02 and L36--L37; MC18 selects the first consumer. AC01--AC08 below develop local compactness, comparison domains and covering bounds. AC30--AC33 supply local structure. AC34--AC37 pass lower comparison to the pointed limit. AC38--AC41 prove the Hausdorff upper bound and assemble the written compactness bridge, conditional on local Toponogov. AC42--AC48 develop splitting and low-dimensional product models for later collapse consumers. AC49--AC56 supply approximate-product limits and coordinate control, with the long-strainer geometric input explicit. AC57--AC70 write both strainer directions; AC71--AC81 write metric compatibility and basepoint transport; AC82--AC87 write cone limits and overlapping-cone splitting. The audit in Section~\ref{sec:revision79-dependency-audit} records the remaining source qualifications.

\section{Comparison and local-model worksheet}
For each proposed comparison, strainer or splitting lemma, record the
metric-space class, curvature normalization, dimension hypothesis,
completeness/properness, radius and error constants, and the exact local
collapse consumer. The present I02 plan does not assert that all familiar
Alexandrov theorems are needed. First fill this worksheet from the
selected Kleiner--Lott route, then expand only the necessary proofs.

\section{The compactness producer consumed by Chapter 3}
\label{sec:alexandrov-packing-producer}

The first selected consumer is MC18, KL Lemma 3.10. Chapter~\ref{ch:metric-geometry}
now proves extraction from eventual local nets. This chapter must supply
those nets from the curvature and dimension hypotheses. The following
elementary reduction makes a possible producer precise without assuming
that a restricted ball is itself an Alexandrov length space.

\begin{lemma}[A compact comparison target supplies finite nets]
\label{lem:alexandrov-model-packing}
Let $A$ be a nonempty metric space and $K$ a compact metric space.
Suppose $\Phi:A\to K$ satisfies
\[
 d_K(\Phi(x),\Phi(y))\geq d_A(x,y)\qquad(x,y\in A).
\]
If $K$ has an $\eta/3$-net of $N$ points, where $\eta>0$, then $A$
has an $\eta$-net of at most $N$ points, with centers in $A$.
Continuity and surjectivity of $\Phi$ are not required.
\end{lemma}
\begin{proof}
Any set in $A$ with pairwise distances greater than $\eta$ maps to
such a set in $K$. Two of its images cannot lie within $\eta/3$ of
the same net center, since their distance would then be at most
$2\eta/3<\eta$. Thus every finite such set has at most $N$ elements.
Starting with one point of $A$, add a point at distance greater than
$\eta$ from all previous choices whenever those choices fail to be
an $\eta$-net. A failure after $N$ choices would give $N+1$ separated
points, a contradiction. The process therefore stops with the required
net. This is the finite packing argument behind BBI Exercise 1.6.4
(printed 14/PDF 29); it uses no completeness or compactness of $A$.
\end{proof}

For each $R>0$, a fixed compact target $K_R$, an index $I_R$, and maps
\[
 \Phi_{i,R}:\overline B_{X_i}(p_i,R)\longrightarrow K_R
 \quad(i\geq I_R)
\]
with the displayed distance inequality would therefore give MC13's
eventual net condition: choose a finite $\eta/3$-net of $K_R$ and use
its cardinality for all $i\geq I_R$. The distances on the source ball
here are the ambient restricted distances. The target must be uniform
in $i$; separate compact targets with unbounded covering numbers do
not suffice. Such maps are one sufficient route, not an extra
assumption inserted into the statement of MC18.

\section{Local compactness and the comparison domain}
\label{sec:alexandrov-local-domain}

The route selected here follows the packing proof in BBI Lemma 10.9.2
(printed 390--391/PDF 405--406): contract radially into a small
bilipschitz chart, then count separated points in Euclidean space.
This avoids using the space of directions to obtain the first packing
bound. Indeed, the proof of BBI Theorem 10.8.6 about directions passes
through Section 10.9; Proposition 10.9.1 there already uses Lemma
10.9.2. The logarithm-map route of Proposition 10.6.10 is still useful
later, after those prerequisites are supplied. It is not the dependency
chosen for the first local packing estimate.

All balls in $X$ below use the ambient distance $d_X$. A closed ball is
not thereby a length space with its restricted distance. The auxiliary
intrinsic metric is introduced only on an \emph{open} region, and its
relation to $d_X$ is proved where needed. AC01--AC08 in
\code{alexandrov_geometry_contracts.csv} index the following written
arguments and source-qualified inputs; they do not introduce a second
mathematical document or claim Lean acceptance.

\begin{lemma}[Compact balls inside a locally compact region; AC01]
\label{lem:alexandrov-local-balls}
Let $X$ be a complete length space, let $p\in X$, and let $L>0$.
Suppose $B_X(p,L)$ is locally compact. Then
$\overline B_X(p,r)$ is compact whenever $0<r<L$.
If $2r<L$, any two points of $\overline B_X(p,r)$ admit an
ambient minimizing path, lying in $\overline B_X(p,2r)$.
\end{lemma}
\begin{proof}
This is a localization of the proof of BBI Theorem 2.5.22
(printed 50--51/PDF 65--66). Local compactness at $p$ gives a positive
compact-ball radius. Compact radii are downward closed. If their
supremum $\rho$ were less than $L$, the ball of radius $\rho$ would
be totally bounded: for small $\epsilon>0$, trim an almost minimizing
path from $p$ to a point in that ball at radius $\rho-\epsilon/3$.
When trimming is needed, the remaining path can have length less than
$\epsilon/2$. A finite $\epsilon/3$-net in the compact smaller ball
then covers the larger ball with error less than $\epsilon$.
Completeness makes the closed $\rho$-ball compact.

Cover that compact ball by finitely many open neighborhoods whose
closures are compact and contained in $B_X(p,L)$. Their union contains
some positive metric neighborhood of the ball. The length property,
using the same trimming argument, puts a slightly larger closed ball
in that neighborhood. It is a closed subset of the finite union of
compact closures, hence compact, a contradiction. This proves the
first assertion without asserting compactness at radius $L$.

For $x,y\in\overline B_X(p,r)$ choose paths with lengths decreasing
to $d_X(x,y)\leq2r$. Each point on such a path has distance from $p$
at most $r$ plus half its total length, by using the nearer endpoint.
Thus these paths lie in one compact ball of radius $2r+\alpha<L$,
for a fixed small $\alpha>0$, after discarding finitely many paths.
Arclength parametrization on a common interval, compactness of uniformly
Lipschitz curves and lower semicontinuity of length give a minimizing
limit path. Those curve facts are the inherited MC01 inputs. The same
endpoint estimate puts the minimizing path in the closed $2r$-ball.
\end{proof}

\begin{lemma}[A sufficient controlled radius for radial comparison; AC02]
\label{lem:alexandrov-eight-radius}
Let $R>0$ and let $X$ be a complete length space. Suppose
$U=B_X(p,8R)$ is locally compact and is finite-dimensional with local
Alexandrov curvature at least $-1$, interpreted in its intrinsic length
metric $d_U$. For $q\in B_X(p,R/2)$ and
$x,y\in\overline B_X(p,R)$ choose ambient minimizing paths from $q$
to $x$ and $y$. Radial subtriangles on these paths satisfy the lower
comparison inequality in curvature $-1$ supplied by KL Section 3.3.
All their distances agree whether computed using $d_X$ or $d_U$.
\end{lemma}
\begin{proof}[Domain verification for the cited local Toponogov input]
First $U$ is path connected: a sufficiently short almost minimizing
path from $p$ to any point of $U$ stays in $U$. Its intrinsic metric
is finite and satisfies $d_U\geq d_X$. Near each point, sufficiently
short ambient almost minimizing paths remain in $U$; hence the two
metrics agree on a smaller neighborhood and induce the same local
length and curvature conventions.

For any $v\in B_X(p,2R)$ the intrinsic closed $6R$-ball at $v$ is
complete. An intrinsic Cauchy sequence has an ambient limit, and
\[
 d_X(p,v)+6R<8R
\]
puts that limit in $U$. Local equality of the metrics gives intrinsic
convergence; closedness retains the intrinsic radius bound. Thus the
complete balls required here are established for $d_U$, not presumed
from the completeness of an open subset.

Lemma~\ref{lem:alexandrov-local-balls} supplies the chosen radial paths.
Their lengths are less than $3R/2$; every point of either lies in
$B_X(p,2R)$. Any two such radial points have ambient distance less than
$3R$, using the path through $q$. Their ambient minimizing paths also
exist: choose $r<2R$ containing the two endpoints and apply that lemma.
These paths lie in $B_X(p,4R)\subset U$, so their intrinsic and
ambient endpoint distances coincide. For every radial subtriangle
we may therefore take
\[
 D=3R,\qquad 2D=6R.
\]
Its sides have length strictly less than $D$ and its intrinsic closed
$2D$-balls at the vertices are complete, as proved above. This verifies
the safeguard in KL Section 3.3 (printed 23/PDF 18) for local
Toponogov and its monotonicity of comparison angles. The cited
comparison theorem remains a source-qualified input; its full
Alexandrov proof is not reproduced here. The factor $8$ is our
sufficient domain margin, not a constant quoted from KL.
\end{proof}

\section{Radial contraction and a uniform covering bound}
\label{sec:alexandrov-radial-covering}

\begin{lemma}[Quantitative radial contraction at curvature minus one; AC03]
\label{lem:alexandrov-radial-contraction}
Under Lemma~\ref{lem:alexandrov-eight-radius}, choose $0<\delta<R$,
put $t=\delta/(2R)$, and choose one minimizing path from $q$ to each
$x\in A=\overline B_X(p,R)$. Write $h_t(x)$ for the point at fraction
$t$ of its length, taking $h_t(q)=q$ if necessary. Then
\[
 h_t(A)\subset B_X(q,\delta),\qquad
 d_X(h_t(x),h_t(y))\geq
 \frac{\delta}{\sinh(2R)}\,d_X(x,y).
\]
No continuity of the chosen map is asserted or needed.
\end{lemma}
\begin{proof}
All radial lengths are less than $3R/2<2R$, so the image inclusion
holds. The following calculation makes the coefficient explicit in
the normalization $k=-1$ used here; compare BBI Lemma 10.6.2
(printed 370/PDF 385). In a hyperbolic comparison triangle let the
radial lengths be $a,b\leq S=2R$, the included angle be $\theta$,
and the opposite side be $c$. Let $c_t$ be the opposite distance at
radial lengths $ta,tb$ with the same angle. The hyperbolic cosine law gives
\[
 \sinh^2(c/2)=\sinh^2((a-b)/2)
                 +\sinh(a)\sinh(b)\sin^2(\theta/2).
\]
Set $\lambda=tS/\sinh S$. For $0\leq u\leq S$, monotonicity of
$\sinh(u)/u$ (with its value at zero defined by continuity) gives
$\sinh(tu)\geq tu\geq\lambda\sinh(u)$. Apply this to
$|a-b|/2,a,b$ in the displayed identity to obtain
$\sinh(c_t/2)\geq\lambda\sinh(c/2)$. Since $0<\lambda\leq1$,
concavity of $\operatorname{arsinh}$ on $[0,\infty)$ gives
$c_t\geq\lambda c$. Lower comparison from
Lemma~\ref{lem:alexandrov-eight-radius} puts the actual radial
subpoint distance at least this model distance. Finally
$\lambda=\delta/\sinh(2R)$. Degenerate triangles follow from the
same formula by continuity; if a radial endpoint is $q$, the assertion
also follows directly from $\lambda\leq t$.
\end{proof}

\begin{theorem}[A small chart supplies uniform nets in a large ball; AC04]
\label{thm:alexandrov-chart-covering}
Assume Lemma~\ref{lem:alexandrov-eight-radius}. Suppose $n\geq1$,
$L\geq1$, and, for the point $q\in B_X(p,R/2)$, there are a number
$\delta_0>0$ and a map
$\phi:B_X(q,\delta_0)\to\mathbb R^n$ with $\phi(q)=0$ such that
\[
 L^{-1}d_X(u,v)\leq |\phi(u)-\phi(v)|\leq Ld_X(u,v).
\]
For every $\epsilon>0$, the ball $A=\overline B_X(p,R)$ has an
internal $\epsilon$-net with at most
\[
 N(n,L,R,\epsilon)=
 \left(1+\left\lceil
   \frac{4L^2\sqrt n\,\sinh(2R)}{\epsilon}
 \right\rceil\right)^n
\]
points. In particular, the bound does not depend on $\delta_0$.
\end{theorem}
\begin{proof}
Choose $0<\delta<\min\{R,\delta_0\}$ and use
Lemma~\ref{lem:alexandrov-radial-contraction}. If a finite subset of
$A$ has pairwise distances greater than $\epsilon$, its images under
$\phi\circ h_t$ have pairwise Euclidean distances greater than
\[
 s=\frac{\delta\epsilon}{L\sinh(2R)}.
\]
They lie in the cube $[-L\delta,L\delta]^n$. Divide each coordinate
interval into bins of length at most $s/(2\sqrt n)$, assigning boundary
points consistently to one bin. At most
$1+\lceil4L^2\sqrt n\sinh(2R)/\epsilon\rceil$ bins per coordinate
suffice. Each resulting box has diameter at most $s/2$, so it contains
at most one image point. This bounds every finite separated subset by
$N(n,L,R,\epsilon)$. The bounded greedy argument of
Lemma~\ref{lem:alexandrov-model-packing} yields an internal
$\epsilon$-net with the same bound. The factor $\delta$ cancels
between the size of the target cube and the separation of its points.
\end{proof}

\lean{The mathematical implementation order is AC01, AC02, AC03, AC04.
Only the distortion $L$ must be uniform across a family; the chart
radius may tend to zero. A chart into $\mathbb R^m$ with $1\leq m\leq n$
may be followed by the standard isometric inclusion into $\mathbb R^n$.
Singleton spaces have a one-point net and are handled separately.
No lower volume bound, injectivity-radius estimate or noncollapse
assumption is used. Formal declarations and imports remain deferred.}

\section{The inherited smooth specialization}
\label{sec:alexandrov-smooth-covering}

\begin{lemma}[Small smooth charts control ambient distance; AC05]
\label{lem:alexandrov-smooth-chart}
In a smooth Riemannian manifold of positive dimension $n$, every point
$q$ has a radius $\delta_0>0$ and a chart on $B_M(q,\delta_0)$
satisfying the preceding theorem with $L=2$ and $\phi(q)=0$.
Distances are those of the whole manifold.
\end{lemma}
\begin{proof}[Local adapter using inherited smooth geometry]
Normalize coordinates at $q$ so the metric at zero is Euclidean.
Continuity of the metric gives an $a>0$ such that the closed coordinate
ball of radius $4a$ lies inside the chart and the tangent norm satisfies
\[
 \tfrac12|v|\leq |v|_g\leq2|v|
\]
throughout it. For two points with coordinates in $B(0,a)$, the
straight coordinate segment gives $d_M(u,v)\leq2|\phi(u)-\phi(v)|$.
Any connecting curve staying in the coordinate $4a$-ball has length
at least $\tfrac12|\phi(u)-\phi(v)|$. If it exits that ball, its
initial segment has length at least $(4a-a)/2=3a/2$, also at least
$\tfrac12|\phi(u)-\phi(v)|$ since the coordinate separation is at
most $2a$. Taking the infimum over all connecting curves therefore gives
\[
 \tfrac12|\phi(u)-\phi(v)|\leq d_M(u,v)
       \leq2|\phi(u)-\phi(v)|.
\]
A sufficiently small ambient ball at $q$ is contained in the coordinate
$a$-ball: any curve from zero to its exterior must first travel to the
coordinate sphere of radius $a$, at cost at least $a/2$.
Take $\delta_0<a/2$. Rearranging the displayed inequalities proves
the stated bilipschitz bound for the ambient metric.

This uses inherited coordinates, continuity, curve length and the
Riemannian distance construction; it does not rebuild PC geometry.
BBI Lemmas 5.1.13 and 5.1.16 (printed 146--147/PDF
161--162) motivate the near-Euclidean input. BBI leaves Lemma 5.1.13
as an exercise; the exit-path argument above supplies the particular
ambient-distance detail needed here.
\end{proof}

\begin{corollary}[Eventual nets for smooth sources; AC06]
\label{cor:alexandrov-smooth-eventual-nets}
Let $(M_i,p_i)$ be complete connected smooth Riemannian $n$-manifolds,
with fixed $n\geq1$. Suppose $c_i\geq0$, $c_i\to0$, $r_i\to\infty$,
and sectional curvature is at least $-c_i$ on $B_{M_i}(p_i,r_i)$.
For every $R,\epsilon>0$, all sufficiently late closed $R$-balls
have internal $\epsilon$-nets of at most $N(n,2,R,\epsilon)$ points.
Consequently Theorem~\ref{thm:metric-eventual-compactness} produces
one complete proper pointed metric limit along a subsequence.
\end{corollary}
\begin{proof}[Conditional on the inherited Riemannian comparison input]
Fix $R,\epsilon>0$. Choose $I$ such that
\[
 c_i\leq1\quad\hbox{and}\quad r_i>8R\qquad(i\geq I).
\]
The Riemannian metric is a length metric and is locally compact;
completeness is an assumption. The inherited sectional-curvature to
local Alexandrov-comparison theorem supplies curvature at least $-1$
on the open $8R$-ball in the convention of AC02. Choose $q=p_i$ and
apply Lemma~\ref{lem:alexandrov-smooth-chart} with $L=2$. Its chart
radius may depend on $i$ and need have no positive uniform lower bound.
Theorem~\ref{thm:alexandrov-chart-covering} gives the stated uniform
net bound. MC13 now applies without any finite-prefix adjustment.
\end{proof}

The smooth comparison and metric interfaces just used are inherited
reuse candidates. Their eventual declaration/import binding is subject
to the completed PC release audit; no claim about the unfinished
snapshot's exact exports is made. The conclusion here is metric
compactness. Length closure, comparison of the limit and dimension
closure are separate conclusions, not consequences of the word
``proper.''

\section{The general Alexandrov producer and remaining limit geometry}
\label{sec:alexandrov-local-structure-gate}

\begin{definition}[Local structure input to be bound; AC07]
\label{def:alexandrov-local-chart-input}
For complete length sources with a growing controlled open region
having curvature at least $-1$ and Hausdorff dimension at most $n$, the required
local structure producer must establish the following facts in each
sufficiently interior $8R$-ball:
\begin{enumerate}[leftmargin=*]
\item Local compactness at all points of that ball, in the ambient
metric (equivalently in its locally identical intrinsic metric).
\item Unless the source is a singleton, a point $q\in B_X(p,R/2)$
and some $\delta_0>0$ admitting the ambient-distance chart of
Theorem~\ref{thm:alexandrov-chart-covering}, with a distortion
$L(n)\geq1$ depending only on the dimension bound. The chart dimension
may be any positive integer at most $n$.
\end{enumerate}
The uniform chart \emph{distortion}, not its radius, is the required
quantifier. Both clauses now have written producers below: AC33 and
AC17 supply local compactness, and
Corollary~\ref{cor:alexandrov-ambient-paired-chart} supplies the chart.
This does not accept the limit-geometry conclusions of KL Lemma 3.10.
\end{definition}

The source route is BBI Proposition 10.8.15, Corollary 10.8.16, Theorem 10.8.18,
Corollaries 10.8.20 and 10.8.23 (printed 385--389/PDF 400--404),
then the radial packing proof in Lemma 10.9.2. The relevant proofs and
Lemmas 10.8.13--10.8.14 were read with the author corrections to
printed 383--387. BBI explicitly simplifies these local proofs to
curvature zero and leaves the small negative-curvature corrections to
be supplied (printed 380/PDF 395). The original route requires local-domain, completeness-buffer and
negative-curvature adaptations. AC24--AC29 below now supply the nearby
chart by a finite hierarchy of paired configurations; AC30--AC33
then supply compactness at all points through AC17. In particular, the
printed chart constant $500m$ is not accepted here at curvature $-1$
without that adaptation. Local compactness in AC01--AC04 is an explicit
input, now supplied by AC33 and AC17; those lemmas are not used to
prove their own compactness input.

KL Definition 3.6 fixes equal strainer radii and uses one quality
parameter; BBI Definition 10.8.9 has different cross-angle tolerances.
No equality of the two parameterized predicates is assumed. The
chart-output interface avoids choosing a premature strainer-parameter
adapter. The dimension-zero alternative is supplied separately by
Lemma~\ref{lem:alexandrov-zero-dimensional-singleton} below.

\begin{corollary}[General extraction conditional on local structure; AC08]
\label{cor:alexandrov-conditional-eventual-nets}
Take the complete length sources of MC18, with fixed finite dimension
bound $n\geq1$, curvature at least $-c_i$ on growing $r_i$-regions,
$c_i\to0$ and $r_i\to\infty$. Suppose the local structure input AC07
has been supplied with one $L(n)$ for all sufficiently late controlled
balls. Then their closed $R$-balls have eventual internal
$\epsilon$-nets of at most $N(n,L(n),R,\epsilon)$ points, for every
$R,\epsilon>0$. They admit a subsequence converging in MC13's
approximation sense to one complete proper pointed metric space.
\end{corollary}
\begin{proof}
For fixed $R,\epsilon$, pass to a tail with $c_i\leq1$ and $r_i>8R$.
Apply AC07 and Theorem~\ref{thm:alexandrov-chart-covering} on this
tail, using a one-point net for any singleton source. The bound is
independent of both the source index and its chart scale. Apply
Theorem~\ref{thm:metric-eventual-compactness}.
\end{proof}

\paragraph{Intrinsic dimension in the AC08 application.}
The ambient dimension hypothesis also supplies AC02's intrinsic one;
this is a domain translation, not an extra assumption. After AC07 gives
local compactness on the larger controlled region, AC01 makes each
closed ball of radius less than $r_i$ compact. Thus the open $8R$-ball
is separable and has a countable cover by neighborhoods on which its
intrinsic metric equals the ambient metric, as in AC02/AC13.
Each neighborhood has Hausdorff dimension at most $n$; countable
stability gives the same bound for the intrinsic open ball. Local
curvature uses that same local metric equality. This order does not
use intrinsic dimension to obtain the initial local compactness.

\paragraph{What extraction alone does not prove.}
AC29 and AC33/AC17 below supply both clauses of AC07, and MC14
supplies length closure and geodesicity. AC34--AC37 separately verify
the comparison domains and limit passage. Generic finite nets alone
do not imply dimension closure: AC38--AC40 use the explicit uniform
power in AC04's bound to construct Hausdorff covers. AC41 assembles
these written arguments. AC06 remains the smooth specialization using
inherited inputs. The source-qualified local Toponogov dependency and
future formal bindings remain; no Lean verification is claimed.

\section{Quantitative distance coordinates for the local producer}
\label{sec:alexandrov-distance-coordinates}

The new source comparison uses Alexander--Kapovitch--Petrunin (AKP),
\emph{Alexandrov geometry: foundations}, Theorem 15.10 and Lemma 15.12,
printed/PDF 233--235 in the archived 303-page version, and
Kapovitch--Lebedeva--Petrunin (KLP), Theorem 6.11, Exercise 6.12 and
its semisolution, printed 66,137/PDF 68,139
\cite{AlexanderKapovitchPetruninFoundations,KapovitchLebedevaPetrunin2026}.
AKP also provides TeX sources at
\url{https://github.com/anton-petrunin/book}.
The source snapshot used here is pinned to a commit in
\code{reference_checks_revision63.md}; it is not asserted to be
identical to the archived PDF or to every later online version.

AKP uses $m+1$ pairwise obtuse \emph{struts}; BBI and KL use paired
\emph{strainers}. These are different predicates. AKP's dimension in
this chapter is initially \emph{linear dimension}; KL's MC18 hypothesis
is Hausdorff dimension on a controlled open region. Moreover, the cited
AKP results concern complete length spaces with a global lower bound,
and KLP's standing Alexandrov spaces are complete and geodesic.
Neither theorem may be applied directly to an incomplete open region
or to a closed ball with the restricted metric. The following argument
extracts the quantitative chart estimate while retaining those structural
translations as explicit inputs.

\subsection{A short hyperbolic hinge gives a definite distance drop}

\begin{lemma}[An explicit negative-curvature estimate; AC09]
\label{lem:alexandrov-short-hinge}
Let $0<\theta<\pi/2$, $a_0>0$, and put $\sigma=\sin\theta$.
Suppose an ambient minimizing hinge $[xy],[xz]$ has lengths
$a=d(x,z)\geq a_0$ and $\ell=d(x,y)>0$, angle at most
$\pi/2-\theta$, and satisfies the lower-curvature $-1$ hinge comparison.
If
\[
 \ell\leq\tau(a_0,\theta):=
 \min\{1,a_0/2,(\sigma/2)\tanh a_0\},
\]
then
\[
 d(x,z)-d(y,z)\geq(\sigma/2)\,d(x,y).
\]
The comparison must hold for this actual hinge in its controlled domain.
\end{lemma}
\begin{proof}
Let $c$ be the opposite distance in the hyperbolic model with the same
radial lengths and included angle. By comparison $d(y,z)\leq c$;
by the hyperbolic cosine law,
\[
 \cosh c\leq\cosh a\cosh\ell-\sigma\sinh a\sinh\ell.
\]
For $0\leq\ell\leq1$, the power series gives
$\cosh\ell-1\leq\ell^2$, and $\sinh\ell\geq\ell$. Hence
\begin{align*}
 \cosh c
 &\leq\cosh a+\cosh a\,\ell^2-\sigma\sinh a\,\ell\\
 &\leq\cosh a-(\sigma/2)\sinh a\,\ell\\
 &\leq\cosh\bigl(a-(\sigma/2)\ell\bigr).
\end{align*}
The middle inequality uses
$\ell\leq(\sigma/2)\tanh a_0\leq(\sigma/2)\tanh a$;
the last uses convexity of $\cosh$. The argument
$a-(\sigma/2)\ell$ is positive because $\ell\leq a_0/2$.
Monotonicity of $\cosh$ on $[0,\infty)$ therefore gives
$c\leq a-(\sigma/2)\ell$, proving the claim. No comparison between
a tensor and an intrinsic domain distance is hidden in this estimate.
\end{proof}

\subsection{The angular input and the distance-coordinate estimate}

\begin{definition}[Angular obstruction required at the chart dimension]
\label{def:alexandrov-angular-obstruction}
At a point $x$ with an angular metric on its space of directions,
write $H(m,\theta;x)$ for the following assertion: there do not exist
unit directions $\xi,\zeta^0,\ldots,\zeta^m$ such that
\[
 \angle(\zeta^j,\zeta^k)>\pi/2+\theta\quad(j\ne k),
 \qquad
 \angle(\xi,\zeta^j)>\pi/2-\theta\quad(0\leq j\leq m).
\]
Here $m$ is the actual chart dimension. An upper bound $n\geq m$
does not by itself supply $H(m,\theta;x)$.
\end{definition}

AKP Lemma 15.12 supplies this assertion in its complete global
$m$-dimensional setting. Its proof first moves a finite configuration
to a nearby straight point, uses Corollary 13.40 to put its directions
in a Euclidean subcone, and then proves a hemispherical packing bound.
Corollary 13.40 depends on the linear-subcone theorem, whose supplied
proof uses the anti-sum construction and splitting. Those are genuine
inputs, not elementary consequences of a Hausdorff-dimension inequality.
Their local versions and the dimension translation remain unbound here.

\begin{theorem}[Quantitative distance-coordinate embedding; AC10]
\label{thm:alexandrov-distance-embedding}
Let $V=B_X(q,\rho)$, $m\geq1$, $0<\theta<\pi/2$, and $a_0>0$.
Choose anchors $z_0=b,z_1,\ldots,z_m$ outside $V$. Assume:
\begin{enumerate}[leftmargin=*]
\item $d(x,z_j)\geq a_0$ for every $x\in V$ and every $j$, and
$2\rho\leq\tau(a_0,\theta)$.
\item For every $x\in V$ there are chosen ambient minimizing paths to
all anchors. Their initial directions have pairwise angles greater than
$\pi/2+\theta$, and $H(m,\theta;x)$ holds.
\item For every ordered pair of distinct points $x,y\in V$, an ambient
minimizing path $[xy]$ is available, and the curvature $-1$ hinge
comparison holds between that path and each chosen anchor path.
\end{enumerate}
Set $\sigma=\sin\theta$ and
\[
 F(x)=\bigl(d(x,z_1)-d(q,z_1),\ldots,
                 d(x,z_m)-d(q,z_m)\bigr).
\]
Then $F(q)=0$ and for all $x,y\in V$,
\[
 (\sigma/2)d(x,y)\leq |F(x)-F(y)|
                         \leq\sqrt m\,d(x,y).
\]
Thus $F$ is a bilipschitz embedding for the ambient distances.
No openness or surjectivity onto a Euclidean neighborhood is asserted.
\end{theorem}
\begin{proof}
For distinct $x,y$, apply the angular obstruction at $x$ to the direction
of $[xy]$ and the $m+1$ anchor directions. At least one anchor makes
angle at most $\pi/2-\theta$ with $[xy]$. Apply the same argument
at $y$ to $[yx]$. Lemma~\ref{lem:alexandrov-short-hinge} applies
because $d(x,y)<2\rho\leq\tau(a_0,\theta)$.
If both witnesses were $b$, it would give simultaneously
\[
 d(x,b)-d(y,b)\geq(\sigma/2)d(x,y),\qquad
 d(y,b)-d(x,b)\geq(\sigma/2)d(x,y),
\]
which is impossible. Hence at least one witness is $z_j$ with $j\geq1$,
and its coordinate difference has absolute value at least
$(\sigma/2)d(x,y)$. This proves the lower Euclidean norm bound.
Every distance coordinate is $1$-Lipschitz, so summing the squares of
its $m$ coordinate differences proves the upper bound. Equal points
are immediate.
\end{proof}

This spells out the endpoint reversal in AKP's proof and KLP's
semisolution, with an explicit hyperbolic scale instead of the phrase
``sufficiently small.'' The right-inverse theorem is used by AKP for
openness of the chart. AC04 only needs a bilipschitz map on a small
ambient ball, so its use of AC10 does not add that right-inverse
conclusion. This does not remove the deeper dependencies of the
angular obstruction itself.

If $X$ is complete, AC10 also proves compactness of each closed ball
$\overline B_X(q,r)$ with $r<\rho$: its image under $F$ is bounded,
and is complete in Euclidean distance by the lower Lipschitz bound.
The image is therefore closed and compact in $\mathbb R^m$; the
continuous inverse proves compactness of the original ball. This is
compactness near this chart point only. It supplies neither geodesics
used earlier in AC10 nor local compactness at every other point.

\subsection{Uniform angle quality from a Euclidean tangent}

\begin{lemma}[A dimension-controlled strut margin; AC11]
\label{lem:alexandrov-simplex-struts}
Let $1\leq m\leq n$ and put $\theta_n=1/(8n)$. Suppose the tangent
cone at $q$ is cone-isometric to $\mathbb R^m$, its unit directions are
the completion of geodesic directions, and the curvature $-1$
comparison angles of shortened representative geodesics tend to their
angular distances. Then there are anchors $z_0,\ldots,z_m$ at a common
positive distance $s$ from $q$, as small as prescribed, with
\[
 \widetilde\angle_{-1}(z_j,q,z_k)>\pi/2+5\theta_n
                       \qquad(j\ne k).
\]
After shrinking an ambient neighborhood $B_X(q,\rho)$, the same
comparison angles at every $x$ in that neighborhood exceed
$\pi/2+4\theta_n$, and $d(x,z_j)>s/2$.
\end{lemma}
\begin{proof}
In the hyperplane of $\mathbb R^{m+1}$ perpendicular to
$u=(1,\ldots,1)$, take the regular-simplex unit vectors
\[
 v_j=\sqrt{\frac{m+1}{m}}
       \left(e_j-\frac{u}{m+1}\right),\qquad 0\leq j\leq m.
\]
Their off-diagonal inner products are $-1/m$. Thus their angular
separations are
$\pi/2+\arcsin(1/m)\geq\pi/2+1/n=\pi/2+8\theta_n$.
Use the cone isometry and approximate these finitely many directions
by geodesic directions within angular distance strictly less than
$\theta_n$ each. Their mutual angles then exceed
$\pi/2+6\theta_n$. Shorten the corresponding geodesics to one common
length $s$; by the assumed comparison-angle limit, choose $s$ small
enough that all comparison angles differ from their angular limits
by less than $\theta_n$. This proves the first assertion, including
$m=1$. The endpoints are distinct, and their radial distances are
positive. Continuity of the hyperbolic comparison angle in its three
side lengths, including collinear limiting configurations, now supplies
a neighborhood with the stated strict margin. Decrease $\rho$ below
$s/2$ to get the distance bound by the triangle inequality.
\end{proof}

\begin{corollary}[Conditional uniform chart part of AC07; AC12]
\label{cor:alexandrov-uniform-chart-part}
Fix $n\geq1$, $R>0$, and a complete length source with a locally compact
curvature-controlled region $B_X(p,8R)$ as in AC02. Suppose there is
$q\in B_X(p,R/2)$ and an integer $1\leq m\leq n$ satisfying AC11's
tangent and angle hypotheses. Suppose also that
$H(m,\theta_n;x)$ holds throughout some neighborhood of $q$.
Then the chart part of AC07 is supplied with the uniform distortion
\[
 L(n)=\max\{\sqrt n,\,2/\sin(1/(8n))\}.
\]
The chart radius may depend on the source and on $q$.
\end{corollary}
\begin{proof}
Choose AC11's anchors with $s<R/8$. Take $\rho<s/8$ small enough
for its comparison-angle margins, the angular obstruction, and
$2\rho\leq\tau(s/2,\theta_n)$. All chart points and anchors lie in
$B_X(p,R)$. AC01 supplies ambient minimizing paths among them.
The same intrinsic-domain verification as AC02 applies: these paths
lie in the controlled region, their endpoint distances agree, and the
closed intrinsic $6R$-balls about their vertices are complete. KL's
local comparison input with $D=3R$ therefore supplies the required
hinge comparisons, including those between two anchor paths at a
chart point. Actual anchor angles are at least their comparison angles,
so AC10 applies with $a_0=s/2$ and $\theta=\theta_n$.
Its upper constant is at most $\sqrt n$ and its inverse constant is
$2/\sin\theta_n$. Translate the distance coordinates by their values
at $q$ and, if necessary, embed $\mathbb R^m$ isometrically in
$\mathbb R^n$. The resulting map satisfies AC04 with the displayed
$L(n)$.
\end{proof}

\paragraph{Remaining structural work.}
AC09--AC12 settle the explicit hyperbolic distance estimate and the
uniform-distortion assembly \emph{conditional on} their stated inputs.
They do not close AC07. Their own route still needs local compactness
on the entire controlled region, a nearby Euclidean tangent of
some dimension $m\leq n$, and the localized angular obstruction
at that same $m$. The translation from the given Hausdorff-dimension
bound to those tangent and local-dimension statements must be proved.
AC29 below supplies the nearby chart by a separate paired-configuration
argument that does not require those tangent inputs.
AKP Theorem 15.13 and Corollaries 15.14--15.16 identify the global
source route; their proof dependencies include straight points,
linear subcones, and the right-inverse/dimension arguments. The
present source reading does not discharge their local versions.
The singleton/dimension-zero branch is handled separately below. In particular,
AC12 uses local compactness as an input, so it cannot be used to prove
that input by circularly invoking its own chart construction.

\section{Local completeness and the finite-dimensional structure route}
\label{sec:alexandrov-local-completeness}

The newly archived Burago--Gromov--Perelman paper (BGP)
\cite{BuragoGromovPerelman1992} supplies a local source route for the
first part of AC07. Its Definition 2.3 assumes local completeness;
Theorem 5.4 constructs distance-coordinate maps; Lemmas 5.9, 6.3
and 6.4 connect their rank to local rough and Hausdorff dimension.
The relevant printed pages are 5,17--22 (PDF 6,18--23 in the
59-page archived English copy). The global precompactness statement,
Corollary 6.8, assumes completeness and is not directly a theorem
about our open controlled region.

The source has two explicit qualifications in Remark 5.1.1:
its displayed constants are calculated at curvature zero, and its
proofs in Sections 5--6 assume available minimizing geodesics,
leaving the general modifications unstated. We retain both tasks at
curvature $-1$. The following metric arguments do not require local
compactness or minimizing paths. They make the completeness part of
that adaptation precise and isolate the remaining geometric input.

\subsection{Complete inner neighborhoods before local compactness}

\begin{lemma}[Intrinsic local completeness without compactness; AC13]
\label{lem:alexandrov-intrinsic-local-completeness}
Let $(X,d_X)$ be a complete length space, $p\in X$, $L>0$, and
$U=B_X(p,L)$. Give $U$ the intrinsic length distance $d_U$, obtained
by taking lengths of curves in $U$ in the ambient metric.
Then $d_U$ is finite, is a length metric and induces the ambient
subspace topology. For $x\in U$ and $h>0$ with
\[
 d_X(p,x)+4h<L,
\]
the metrics agree on $\overline B_X(x,h)$, and
\[
 \overline B_{(U,d_U)}(x,h)=\overline B_X(x,h).
\]
This common closed ball is complete in either metric. Local
compactness and existence of minimizing geodesics are not assumptions.
\end{lemma}
\begin{proof}
For $y\in U$, choose an almost-shortest path from $p$ to $y$ of
length less than $L$. Every point on it is in $U$. Concatenating
such paths shows $d_U(y,z)<\infty$ for all $y,z\in U$; also
$d_U\geq d_X$. The induced length construction is a length metric:
for a rectifiable path in $U$, its $d_U$-length equals its ambient
length, since each subpath bounds the intrinsic distance of its
endpoints, and $d_U\geq d_X$ gives the reverse inequality.

For $y,z\in\overline B_X(x,h)$, choose an ambient path of length
less than $d_X(y,z)+\eta$, where $0<\eta<h$. Its points have
$d_X$-distance from $x$ less than $3h+\eta<4h$, so the path is in
$U$. Taking $\eta\downarrow0$ proves $d_U(y,z)=d_X(y,z)$.
This also proves the equality of the closed balls, using
$d_U\geq d_X$ in the other direction. The ambient closed ball is
complete as a closed subset of $X$, hence is complete for $d_U$.
Varying $x,h$ proves equality of the topologies and local completeness.
\end{proof}

In particular, local four-point comparison and Hausdorff-dimension
bounds transfer on these common neighborhoods. We do not need an
unproved countable cover to infer a global dimension bound for
$(U,d_U)$: AC16 below is a local assertion. Nor do we declare these
closed balls to be length spaces. This separates a complete inner
neighborhood from a compact inner neighborhood; AC01 requires the
latter, whereas AC13 does not supply it.

\subsection{A buffered limit for successive corrections}

\begin{lemma}[Successive corrections stay in a complete buffer; AC14]
\label{lem:alexandrov-buffered-correction}
Let $C=\overline B_Y(x_0,r)$ be a complete closed ball in a metric
space $Y$, with $r>0$. Let $F:O\to E$ be continuous, where
$C\subset O\subset Y$ and $E$ is a normed space. Fix $v\in E$,
$\mu>0$ and $0\leq\lambda<1$, and set $e(x)=\|F(x)-v\|$.
Assume $e(x_0)<\mu r$. For every $x\in C$ with $e(x)>0$, suppose
there exists $x'\in O$ satisfying
\[
 e(x')\leq\lambda e(x),\qquad
 \mu d_Y(x,x')\leq e(x)-e(x').
\]
Then there is $x_\infty\in C$ with
\[
 F(x_\infty)=v,\qquad
 d_Y(x_0,x_\infty)\leq e(x_0)/\mu<r.
\]
The correction point need not be assumed to lie in $C$ in advance;
the cumulative displacement estimate proves this for the iteration.
\end{lemma}
\begin{proof}
If $e(x_0)=0$, take $x_\infty=x_0$. Otherwise choose successive
corrections and keep a zero-residual point fixed if one is reached.
Write $e_j=e(x_j)$. By induction,
\[
 e_j\leq\lambda^j e_0,\qquad
 \sum_{j=0}^{k-1}d_Y(x_j,x_{j+1})
       \leq (e_0-e_k)/\mu\leq e_0/\mu<r.
\]
Thus every new point is in $C$, so the next correction is legitimate.
For $k>j$, the same telescoping estimate gives
$d_Y(x_j,x_k)\leq e_j/\mu\leq\lambda^j e_0/\mu$.
The sequence is Cauchy in the complete ball $C$ and converges there.
Continuity of $F$ and $e_j\to0$ give $F(x_\infty)=v$;
passing to the limit in the displacement bound proves the assertion.
\end{proof}

If the correction hypotheses hold for every target with
$\|F(x_0)-v\|<\mu r$, the image contains that open target ball.
For the distance coordinates in BGP's proof of Theorem 5.4
(printed 18--19/PDF 19--20), the residual is an $\ell^1$ norm.
Its Euclidean calculation aims at
$\mu=1-2m\delta$ and $\lambda=1-\mu/m$.
AC14 explains why the displacement bound uses $1/\mu$ and why a
closed completeness buffer is needed. It does not assert those
particular correction constants at curvature $-1$, or construct the
correction moves. Approximate paths and their error budget must be
supplied before using this lemma without known minimizing geodesics.

\begin{lemma}[An open distance image bounds coordinate rank; AC15]
\label{lem:alexandrov-open-image-rank}
Suppose $A$ is a subset of a metric region of Hausdorff dimension at
most $n$, and $F:A\to\mathbb R^m$ is Lipschitz. If $F(A)$ contains
a nonempty Euclidean open ball, then $m\leq n$.
\end{lemma}
\begin{proof}
A Lipschitz map does not increase Hausdorff dimension: applying it to
arbitrarily fine covers multiplies every $s$-dimensional diameter sum
by at most the $s$th power of its Lipschitz constant. A nonempty open
ball in $\mathbb R^m$ has Hausdorff dimension $m$. Monotonicity under
inclusion therefore gives
\[
 m\leq\dim_H F(A)\leq\dim_H A\leq n.
\]
\end{proof}

The Hausdorff-cover interfaces and the dimension of Euclidean balls
are standard metric-measure inputs to be audited before Lean binding.
This is the rank-ceiling step in BGP Lemma 6.4. An injective map alone
would not suffice; it is the open image of a Lipschitz map that gives
the needed direction of the dimension inequality.

\subsection{The local producer contract and its compactness consequence}

For a bounded nonempty set $A$, let $P_A(\epsilon)$ be the supremum
of the cardinalities of its finite subsets with pairwise distances
at least $\epsilon$, allowing $+\infty$. Following BGP Definition 6.2,
put
\[
 V_a^{\rm rough}(A)=
   \limsup_{\epsilon\downarrow0}\epsilon^a P_A(\epsilon),
 \qquad
 \dim_{\rm rough}A=
   \inf\{a\geq0:V_a^{\rm rough}(A)=0\}.
\]
The infimum of the empty set is $+\infty$. Only finiteness of this
dimension will be used, not a value of rough volume at the critical
exponent. The non-strict separation convention here differs from the
strict finite-packing convention in AC04; shrinking a positive scale
translates the finite-bound statements needed below.

\begin{definition}[Local rough-dimension producer; AC16]
\label{def:alexandrov-local-rough-dimension}
The geometric producer has the following contract, supplied by AC33 below.
Let $Y$ be a locally complete length space with local curvature
comparison at least $-1$. Suppose each point has a neighborhood of
Hausdorff dimension at most a fixed finite $n$.
For every $x\in Y$, produce a bounded neighborhood $W$ of $x$ such that
\[
 \dim_{\rm rough}W<\infty.
\]
No local compactness or available minimizing geodesics may be
inserted as extra assumptions. Theorem~\ref{thm:alexandrov-local-rough-dimension-produced}
now supplies this implication in written mathematics, not Lean.
\end{definition}

BGP Lemma 6.4 (printed 20--21/PDF 21--22) states the stronger local
identification of burst index, rough dimension and Hausdorff dimension.
Its proof supplies a local route through openness, a bilipschitz patch
and rough-dimension transport. The current progress is:
\begin{enumerate}[leftmargin=*]
\item AC19--AC23 supply the negative-curvature correction moves and
local openness for a paired packet in one comparison neighborhood,
using almost-shortest paths and a complete closed buffer.
\item AC24--AC28 below prove quantitative next-rank exclusion and
produce a nearby bilipschitz patch under the finite Hausdorff bound.
A finite hierarchy of qualities replaces arbitrary burst improvement
for this consumer. All chosen anchors remain in the same comparison
neighborhood; no distant-anchor localization is assumed.
\item AC30--AC33 below supply the finite-packing transport suggested
by BGP Lemma 6.3 and BBI Corollary 10.8.20. Four-point comparison
and finite distance-array compactness control almost-shortest-path
errors at curvature $-1$, transferring the patch bound to a
neighborhood of every prescribed point.
\end{enumerate}
AC29 supplies the ambient nearby-chart clause of AC07 separately.
The paired configurations are not identified with AKP's obtuse struts;
AC09--AC12 remain the distinct conditional tangent/strut route.

\begin{proposition}[Local compactness from AC16; AC17]
\label{prop:alexandrov-local-compactness-rough}
Let $X$ be a complete length space and $U=B_X(p,L)$ a controlled
region with local lower curvature comparison at least $-1$ and
ambient Hausdorff dimension at most $n<\infty$. If AC16 is supplied,
then every point of $U$ has a compact ambient closed-ball neighborhood
contained in $U$. Thus the local compactness part of AC07 follows.
No family-uniform covering estimate is asserted by this proposition.
\end{proposition}
\begin{proof}
First, if a bounded set $W$ has finite rough dimension, there is a
finite $a\geq0$ for which $V_a^{\rm rough}(W)=0$. Hence
$P_W(\epsilon)<\infty$ for every sufficiently small positive
$\epsilon$. For any $\eta>0$, take such an $\epsilon<\eta$.
A greedy selection of points at distance greater than $\eta$ from
those already selected must stop: its finite subsets are
$\epsilon$-separated and their cardinalities have a fixed finite
bound. On stopping it gives an internal $\eta$-net. The same argument
works for each nonempty subset of $W$. Thus $W$ is totally bounded.
This argument neither assumes nor constructs a maximal infinite
separated set.

Apply AC13 to $(U,d_U)$. Its metric agrees locally with the ambient
metric, so the local curvature and dimension hypotheses of AC16 hold.
For $x\in U$, let $W$ be the resulting intrinsic neighborhood of
finite rough dimension. Choose $r>0$ small enough that the common
ambient/intrinsic closed $r$-ball from AC13 is contained in $W$.
It is complete by AC13 and totally bounded by the preceding argument,
so is compact. The identical metrics give ambient compactness and
prove the assertion at every $x\in U$.
\end{proof}

This proof uses only local neighborhoods. It does not apply BGP's
complete-space Corollary 6.8 to an incomplete open ball. Its output
now supplies AC01's local compactness input through AC33;
AC01, AC02 and AC12 are not used to establish AC16 in advance.
The chart part of AC07 is supplied below by AC29; the limit's
comparison and dimension remain separate tasks.

\subsection{The zero-dimensional branch}

\begin{lemma}[A nontrivial length space has positive local dimension; AC18]
\label{lem:alexandrov-zero-dimensional-singleton}
Let $X$ be a length space, $p\in X$ and $L>0$. If
$\dim_H B_X(p,L)<1$, then $X=\{p\}$. In particular this handles
the dimension-zero branch of AC07 without curvature, completeness,
strainers or tangent cones.
\end{lemma}
\begin{proof}
If $q\ne p$, take a continuous rectifiable path from $p$ to $q$ and
choose $0<a<\min\{L,d_X(p,q)\}$. Stop the path the first time its
distance from $p$ reaches $a$. This first time exists by continuity
on the compact parameter interval; the initial path lies in
$\overline B_X(p,a)\subset B_X(p,L)$. The $1$-Lipschitz distance
function $x\mapsto d_X(p,x)$ maps that initial path onto $[0,a]$.
Since the interval has Hausdorff dimension one, Lipschitz monotonicity
forces $\dim_H B_X(p,L)\geq1$, a contradiction.
\end{proof}

The zero-dimensional alternative is now a written metric proof.
Positive-dimensional compactness is supplied by AC33 and AC17 below. The separate
AC12 route retains its regular-tangent/angular inputs. The following section supplies
the curvature-$-1$ correction move and buffered limit for a specified
paired packet; its production and the later dimension steps remain.

\section[Almost-shortest paths and local openness]{Almost-shortest paths and local openness at curvature minus one}
\label{sec:alexandrov-approximate-corrections}

This section supplies the correction moves left open in AC14 for a
specified paired configuration. The strategy follows BGP Theorem 5.4
and its proof in 5.8 (printed 17--19/PDF 18--20), compared with BBI
Proposition 10.8.15 (printed 385/PDF 400)
\cite{BuragoGromovPerelman1992,BuragoBuragoIvanov}. The estimates and constants
below are derived here for curvature $-1$ and almost-shortest paths.
They do not import the curvature-zero or assumed-geodesics
simplifications of BGP Remark 5.1.1. In particular, the selected
coordinate need not reach its target in one move.

Write $\theta_z(u,v)=\widetilde\angle_{-1}(u,z,v)$ for the
hyperbolic comparison angle computed from the three distances.
Degenerate triangles are included by continuity. A \emph{common
comparison neighborhood} $\Omega$ means an open set on which
\[
 \theta_z(u,v)+\theta_z(v,w)+\theta_z(w,u)\leq2\pi
\]
for every four distinct points $z,u,v,w\in\Omega$, using the metric
of the length space under consideration. Every later application
names endpoints in this one neighborhood. No globalization or
comparison along an unavailable minimizing geodesic is used.

\subsection{Model estimates and almost-radial points}

\begin{lemma}[Uniform hyperbolic finite difference; AC19]
\label{lem:alexandrov-hyperbolic-finite-difference}
Fix $0<a\leq A$ and define
\[
 K=K(a,A):=\frac{4\cosh(A+1)}{\sinh a}.
\]
For any metric triangle $x,y,z$ with $r=d(x,z)\in[a,A]$ and
$0<t=d(x,y)\leq\min\{1,a/2\}$, put $c=d(y,z)$ and
$\theta=\theta_x(z,y)$. Then
\[
 |c-r+t\cos\theta|\leq Kt^2.
\]
This is a model identity estimate; no curvature hypothesis is needed.
\end{lemma}
\begin{proof}
The triangle inequality gives $|c-r|\leq t$. The hyperbolic cosine
law and Taylor's formula give
\begin{align*}
 \cosh c&=\cosh r\cosh t-\sinh r\sinh t\cos\theta,\\
 |\cosh c-\cosh r-(c-r)\sinh r|
     &\leq\tfrac12\cosh(A+1)t^2.
\end{align*}
For $0\leq t\leq1$, $\cosh t-1\leq t^2$ and
$0\leq\sinh t-t\leq t^2$. Thus
\[
 \sinh r\,|c-r+t\cos\theta|
 \leq\bigl(\cosh r+\sinh r+\tfrac12\cosh(A+1)\bigr)t^2
 \leq4\cosh(A+1)t^2.
\]
Divide by $\sinh r\geq\sinh a>0$. The same formulas cover
degenerate model triangles.
\end{proof}

\begin{lemma}[An almost-short path supplies an almost-radial point; AC20]
\label{lem:alexandrov-almost-radial-point}
Let $X$ be a length space, $0<\delta\leq1/100$, and $0<a\leq A$.
Set
\[
 C=\frac{\sinh(A+1)}{\sinh(a/2)},\qquad
 \nu=\min\{\delta^2,(1-\cos\delta)/C,1\}>0.
\]
If $r=d(x,u)\in[a,A]$ and $0<t\leq\min\{1,a/2\}$,
there is $y$ with $d(x,y)=t$ such that, writing $s=d(y,u)$,
\[
 r-t\leq s<r-t+\nu t,\qquad
 \theta_y(x,u)\geq\pi-\delta,\qquad
 \theta_x(y,u)\leq\delta.
\]
In particular $(1-\delta^2)t\leq r-s\leq t$.
The initial path from $x$ to $y$ can be kept in $\overline B(x,t)$.
\end{lemma}
\begin{proof}
Choose a continuous rectifiable path from $x$ to $u$ of length
less than $r+\eta$, where $\eta=\nu t$. Since $t<r$, its first
point $y$ at distance $t$ from $x$ exists. The initial segment stays
in $\overline B(x,t)$ and has length at least $t$; additivity of
length bounds the remaining segment by $r+\eta-t$. This proves
the bounds on $s$, including $s\geq a/2$ and $s+t<r+\eta\leq A+1$.
The model cosine law at $y$ and at $x$ gives, respectively,
\begin{align*}
 1+\cos\theta_y(x,u)
   &=\frac{\cosh(s+t)-\cosh r}{\sinh s\sinh t}
     \leq C\eta/t\leq1-\cos\delta,\\
 1-\cos\theta_x(y,u)
   &=\frac{\cosh s-\cosh(r-t)}{\sinh r\sinh t}
     \leq C\eta/t\leq1-\cos\delta.
\end{align*}
Here the mean value theorem bounds each numerator by
$\sinh(A+1)\eta$, and $\sinh t\geq t$.
Cosine is decreasing on $[0,\pi]$, giving the two angles.
No part of this construction assumes a minimizing path. Only the
initial segment is required to stay near $x$; the rest of the chosen
path need not lie in a comparison neighborhood.
\end{proof}

\subsection{A paired configuration gives a contracting correction}

For the next two lemmas fix $m\geq1$ and
$0<\delta\leq1/(100m)$. A \emph{paired packet on $V$} consists
of $2m$ distinct anchors $a_1,b_1,\ldots,a_m,b_m$ in a common
comparison neighborhood $\Omega$, an open set $V\subset\Omega$,
and bounds $0<a\leq A$, such that for every $z\in V$ all anchor
distances belong to $[a,A]$ and
\begin{align*}
 \theta_z(a_i,b_i)&>\pi-\delta,\\
 \theta_z(c_i,c_j)&>\pi/2-\delta
 \quad(i\ne j,\ c_i\in\{a_i,b_i\},\ c_j\in\{a_j,b_j\}).
\end{align*}
The positive distance bound excludes anchors from $V$.
This is a stated comparison-angle packet, with no equal-radius
requirement. It is not identified with AKP's obtuse-strut predicate
or with KL's equal-radius convention. Set
\[
 t_0=\min\{1,a/2,\delta^2/K(a,A)\}.
\]

\begin{lemma}[One geometric correction move; AC21]
\label{lem:alexandrov-paired-correction-move}
Let $X$ be a length space with the paired packet just specified.
Choose $x\in V$, $0<t\leq t_0$, and
$\overline B(x,t)\subset V$. For either ordered choice
$(u,v)=(a_i,b_i)$ or $(b_i,a_i)$, there is $y\in V$ with
$d(x,y)=t$ and
\begin{align*}
 (1-\delta^2)t&\leq d(x,u)-d(y,u)\leq t,\\
 (1-3\delta^2)t&\leq d(y,v)-d(x,v)\leq t,\\
 |d(y,a_j)-d(x,a_j)|&\leq6\delta t\qquad(j\ne i).
\end{align*}
\end{lemma}
\begin{proof}
Choose $y$ by AC20 toward $u$. Its first estimate gives the drop
in distance to $u$. At $y$, the four-point inequality for $x,u,v$
and the inequalities $\theta_y(x,u)\geq\pi-\delta$ and
$\theta_y(u,v)>\pi-\delta$ give $\theta_y(x,v)\leq2\delta$.
AC19 at $y$ then yields
\[
 d(y,v)-d(x,v)\geq\cos(2\delta)t-Kt^2
                         \geq(1-3\delta^2)t.
\]
The upper bound is the triangle inequality. Apply AC19 at $x$
to the same anchor $v$. It follows that
\[
 1+\cos\theta_x(v,y)\leq4\delta^2,
 \qquad \theta_x(v,y)\geq\pi-4\delta.
\]
For the second implication use
$1-\cos(4\delta)\geq4\delta^2$, which follows from
$1-\cos w\geq w^2/4$ for $0\leq w\leq1$.

Fix $j\ne i$ and write $c=a_j$. The four-point inequality at $x$
for $v,c,y$ gives $\theta_x(c,y)\leq\pi/2+5\delta$.
At $y$, the inequality for $u,c,x$ gives
$\theta_y(c,x)\leq\pi/2+2\delta$. AC19 at the two endpoints
therefore gives
\begin{align*}
 d(y,c)-d(x,c)&\leq(\sin(5\delta)+\delta^2)t\leq6\delta t,\\
 d(x,c)-d(y,c)&\leq(\sin(2\delta)+\delta^2)t\leq3\delta t.
\end{align*}
All four-point comparisons use actual endpoints in $\Omega$;
anchor bounds at both endpoints justify both uses of AC19.
There is no use of tangent angles, local compactness or Toponogov.
\end{proof}

\begin{lemma}[Largest-residual correction with a displacement budget; AC22]
\label{lem:alexandrov-residual-correction}
Use the packet of AC21 and the distance map
$F(z)=(d(z,a_1),\ldots,d(z,a_m))$. For a target $w\in\mathbb R^m$
write $e(z)=\lVert F(z)-w\rVert_1$, and put
\[
 \mu=1-3\delta^2-6(m-1)\delta,\qquad
 \lambda=1-\mu/m.
\]
Then $9/10\leq\mu<1$ and $0\leq\lambda<1$.
If $e(x)>0$, set $h=\max_j|F_j(x)-w_j|$.
Provided $h\leq t_0$ and $\overline B(x,h)\subset V$, there is
$y\in V$ such that
\[
 d(x,y)=h,\qquad e(y)\leq\lambda e(x),\qquad
 \mu d(x,y)\leq e(x)-e(y).
\]
\end{lemma}
\begin{proof}
Choose an index $i$ realizing $h$. Take the AC21 move of length
$t=h$ toward $a_i$ if $F_i(x)>w_i$, and toward $b_i$ otherwise.
In either case the $i$th coordinate moves in the required direction
by an amount between $(1-3\delta^2)h$ and $h$.
The selected coordinate does not overshoot its target. Its residual
decreases by at least $(1-3\delta^2)h$, whereas each other residual
increases by at most $6\delta h$. Hence
\[
 e(y)\leq e(x)-\mu h\leq(1-\mu/m)e(x),
\]
since $h\geq e(x)/m$. This also proves the displacement budget.
Finally $3\delta^2\leq3/10000$ and
$6(m-1)\delta\leq6/100$ imply $\mu>9/10$;
$\delta>0$ implies $\mu<1$. These bounds include $m=1$.
\end{proof}

\subsection{The complete-buffer limit and its precise scope}

\begin{theorem}[Local openness from a common paired packet; AC23]
\label{thm:alexandrov-local-distance-openness}
Let $X$ be a locally complete length space and let $q$ and $2m$
distinct anchors lie in one common comparison neighborhood
$\Omega$, with $m\geq1$. Fix $0<\delta\leq1/(100m)$.
Assume the paired angle inequalities at $q$ hold with quality
$\delta/2$ in place of $\delta$. Then there is an open neighborhood
$V$ of $q$ with a paired packet of quality $\delta$ as above, on
which the distance map $F:V\longrightarrow\mathbb R^m$ is open.
It is $m$-Lipschitz for the $\ell^1$ norm and $\sqrt m$-Lipschitz
for the Euclidean norm. If $\dim_H V\leq n<\infty$, then $m\leq n$.
No injectivity or bilipschitz lower bound is asserted.
\end{theorem}
\begin{proof}
Continuity of finitely many model comparison angles and anchor
distances gives $V\subset\Omega$ and common bounds $0<a\leq A$
with the quality-$\delta$ packet throughout $V$. The initial angles
in particular require $q$ to differ from every anchor. Shrink $V$
to exclude all anchors. Choose $r>0$ such that
$C_q=\overline B(q,r)$ is complete and $B(q,2r)\subset V$.
Local completeness gives a complete closed ball at $q$; any smaller
closed ball is its closed subset. Thus this choice requires no
compactness. Set
\[
 \varepsilon_0=\min\{t_0,\mu r/2\}>0.
\]
Fix $w$ with $e_0=\lVert F(q)-w\rVert_1<\varepsilon_0$.
If $e_0=0$, use $q$. Otherwise, at any $x\in C_q$ with
$0<e(x)\leq e_0$, the largest residual satisfies
$h\leq e_0<t_0$ and $h<r$, so
$\overline B(x,h)\subset B(q,2r)\subset V$.
AC22 supplies its correction.

The correction hypothesis here is needed only on the residual
sublevel $e(x)\leq e_0$, not at all points of $C_q$.
We spell out the corresponding variant of AC14's induction.
Start at $x_0=q$ and stop if a residual vanishes. Otherwise choose
the AC22 correction successively. For every finite initial segment,
\[
 e(x_k)\leq\lambda^k e_0,\qquad
 \sum_{j<k}d(x_j,x_{j+1})
       \leq\frac{e_0-e(x_k)}{\mu}<r/2.
\]
Thus every iterate, including the newly chosen correction, stays
in $C_q$ and in the residual sublevel. The same telescoping bound
on a tail is at most $e(x_k)/\mu$, so the iterates are Cauchy.
Completeness of $C_q$ supplies $x_\infty$, and continuity gives
\[
 F(x_\infty)=w,\qquad d(q,x_\infty)\leq e_0/\mu<r/2.
\]
Stopping at a zero residual gives the same conclusion. Consequently
the $\ell^1$ ball of radius $\varepsilon_0$ about $F(q)$ lies in
$F(V)$. The radius may depend on the packet and the center.

To prove openness, repeat the same argument at any $z\in V$ and
inside any open set $W\subset V$ containing $z$. Local completeness
allows a complete closed $r$-ball with $B(z,2r)\subset W$;
the existing quality-$\delta$ packet on $V$ is sufficient for
AC21 and AC22. This gives a target neighborhood inside $F(W)$.
Each distance coordinate is $1$-Lipschitz, giving both upper bounds.
AC15 applied to the nonempty open image gives the rank ceiling.
\end{proof}

The constants $K,C,\nu,t_0,\mu,\lambda$ are explicit project
estimates, not the printed constants of BGP or BBI. The path excess
$\eta=\nu t$ tends to zero with each correction scale; a fixed
positive path error would not support this iteration. The proof
derives the geometric moves and the complete-buffer limit before
minimizing geodesics or local compactness have been obtained.

AC23 handles a supplied packet whose anchors and comparison
endpoints share one neighborhood. The next section constructs such
packets and proves injectivity using a finite hierarchy; AC30--AC33
then supply rough-dimension transport for AC16. No packet with distant anchors is silently localized by
shrinking its center neighborhood. AC23 alone does not close AC16,
AC07, or the limit-geometry parts of MC18.

\section{Paired charts before local compactness}
\label{sec:alexandrov-paired-charts}

BGP Lemma 5.6, Corollary 5.7 and the injectivity part of Theorem 5.4
suggest adjoining a nearly straight pair at an almost midpoint of
points with nearly equal distance coordinates (printed 18--19/PDF
19--20). BBI Lemma 10.8.14 and Theorem 10.8.18 give the corresponding
strainer argument (printed 383--384,387--388/PDF 398--399,402--403).
The author correction reverses the almost-midpoint angle inequality
in BBI Proposition 10.8.17: it must be close to $\pi$ from below.
The following estimates prove the needed version directly at
curvature $-1$, with no minimizing paths or local compactness.

For the present finite Hausdorff-dimension consumer, arbitrary burst
improvement is not needed first. We choose a finite hierarchy of
rank-dependent qualities and maximize the rank within one comparison
neighborhood. BGP Remark 6.9 (printed 22/PDF 23) motivates this
hierarchy; its additional approximate-comparison and stratification
claims are not imported here. Our exact-comparison proof, constants
and local scope are given below.

\subsection{Endpoint reversal and a balanced almost midpoint}

\begin{lemma}[Comparison-angle reversal and complements; AC24]
\label{lem:alexandrov-angle-reversal}
Let $x,z,u$ be points with $s=d(x,z)>0$,
$d(x,u),d(z,u)\in[a,A]$, and
$s\leq\min\{1,a/2\}$. Put $K=K(a,A)$ as in AC19 and
$\omega=\pi\sqrt{Ks}$. Then
\[
 |\theta_z(u,x)+\theta_x(u,z)-\pi|\leq\omega.
\]
If $u,v,x,z$ lie in one common comparison neighborhood, both anchors
have these bounds at $x,z$, and
$\theta_x(u,v),\theta_z(u,v)>\pi-\delta$, then
\[
 \pi-\delta-2\omega
 \leq\theta_z(u,x)+\theta_z(v,x)\leq\pi+\delta.
\]
\end{lemma}
\begin{proof}
For $\alpha,\beta\in[0,\pi]$ one has
\[
 |\cos\alpha-\cos\beta|
 \geq1-\cos|\alpha-\beta|
 \geq2|\alpha-\beta|^2/\pi^2.
\]
The first inequality follows from the sine product formula, the
second from $\sin t\geq2t/\pi$ on $[0,\pi/2]$.
Apply AC19 to the anchor $u$ at each endpoint. Adding the two
remainders gives
\[
 |\cos\theta_z(u,x)+\cos\theta_x(u,z)|\leq2Ks.
\]
Compare the second angle with $\pi-\theta_z(u,x)$ using the
preceding cosine estimate. This proves reversal. The four-point
inequality at $x$ gives
$\theta_x(u,z)+\theta_x(v,z)\leq\pi+\delta$.
Reverse both angles to obtain the lower complement bound at $z$.
The upper bound is the four-point inequality at $z$ itself.
\end{proof}

\begin{lemma}[Balanced almost midpoint without a minimizer; AC25]
\label{lem:alexandrov-balanced-midpoint}
Let $X$ be a length space, $0<L=d(x,y)\leq1$ and $0<\tau\leq1$.
Define
\[
 \nu_\tau=\min\{1/2,(1-\cos\tau)/(8\cosh2)\}>0.
\]
There is $z$ with
\[
 d(x,z)=d(y,z)=s,\qquad L/2\leq s<(1+\nu_\tau)L/2\leq3L/4,
 \qquad \theta_z(x,y)>\pi-\tau.
\]
\end{lemma}
\begin{proof}
Choose a continuous rectifiable path from $x$ to $y$ of length
less than $(1+\nu_\tau)L$. Along it the continuous function
$d(x,\cdot)-d(y,\cdot)$ changes sign, so vanishes at some $z$.
The triangle inequality gives $2s\geq L$; the sum of the two
subpath lengths gives $2s<(1+\nu_\tau)L$. In particular $z$ is
neither endpoint. The model cosine law gives
\[
 1+\cos\theta_z(x,y)
 =\frac{\cosh(2s)-\cosh L}{\sinh^2s}
 <8\cosh2\,\nu_\tau\leq1-\cos\tau.
\]
Indeed $2s-L<\nu_\tau L$, $\sinh^2s\geq L^2/4$, and
$\sinh((1+\nu_\tau)L)\leq2L\cosh2$ since $L\leq1$.
The angle conclusion follows by monotonicity of cosine. The
construction uses neither comparison in $X$ nor completeness.
\end{proof}

\subsection{Almost equal coordinates force an extra pair}

Use the paired-packet convention of Section~\ref{sec:alexandrov-approximate-corrections}:
all anchors and working points lie in one common comparison
neighborhood $\Omega$; throughout the open working set $V$,
anchor distances are in $[a,A]$, opposite angles exceed
$\pi-\delta$, and cross-pair angles exceed $\pi/2-\delta$.

\begin{lemma}[Quantitative rank increase at an almost midpoint; AC26]
\label{lem:alexandrov-midpoint-rank-increase}
Let $m\geq1$, $0<\beta\leq1/(200(m+1))$ and
$0<\delta\leq\beta/100$. Suppose the rank-$m$ packet just
specified is given in a length space. Set
\[
 \epsilon_\beta=\left(\frac{\beta}{100\pi}\right)^2,
 \qquad \ell_\beta=\min\{1,a/2,\epsilon_\beta/K(a,A)\}.
\]
If $x,y\in V$, $0<L=d(x,y)\leq\ell_\beta$,
$\overline B(x,L)\subset V$, and
\[
 |d(x,a_i)-d(y,a_i)|\leq\epsilon_\beta L
                         \qquad(1\leq i\leq m),
\]
then there is $z\in V$ at which the old pairs together with $(x,y)$
form a rank-$(m+1)$ packet of quality $\beta$ at that point, with
all anchors in $\Omega$. No uniform neighborhood is asserted yet.
\end{lemma}
\begin{proof}
Apply AC25 with $\tau=\beta/100$ and put
$s=d(x,z)=d(y,z)$. Then $s<3L/4$, so $z\in V$.
The endpoints and $z$ are distinct from every old anchor, because
$V$ has positive anchor-distance bounds. Fix an old pair $(u,v)=(a_i,b_i)$
and abbreviate
\[
 \alpha=\theta_z(u,x),\quad \alpha'=\theta_z(u,y),\quad
 \gamma=\theta_z(v,x),\quad \gamma'=\theta_z(v,y).
\]
AC19 at $z$ applied to $x$ and $y$, using the equal lengths $s$,
gives
\[
 |\cos\alpha-\cos\alpha'|
 \leq\epsilon_\beta L/s+2Ks\leq2\epsilon_\beta+2Ks.
\]
Consequently, by AC24's elementary cosine estimate,
\[
 |\alpha-\alpha'|\leq\sigma:=\pi\sqrt{\epsilon_\beta+Ks}.
\]
The four-point inequality at $z$ for $u,x,y$ implies
$\alpha+\alpha'\leq\pi+\tau$, so both are at most
$\pi/2+H$, where $H=(\tau+\sigma)/2$.
Put $\omega=\pi\sqrt{Ks}$ and $D=\delta+2\omega$.
AC24 for the pair $(u,v)$ and each endpoint gives
\[
 \alpha+\gamma\geq\pi-D,\qquad
 \alpha'+\gamma'\geq\pi-D.
\]
Thus $\gamma,\gamma'\geq\pi/2-D-H$.
The four-point inequality for $v,x,y$ gives
$\gamma+\gamma'\leq\pi+\tau$, hence both are at most
$\pi/2+\tau+D+H$. Substituting back gives
$\alpha,\alpha'\geq\pi/2-\tau-2D-H$.
All four angles therefore differ from $\pi/2$ by at most
\[
 E=2\delta+\tfrac32\tau+4\omega+\tfrac12\sigma<\beta/10.
\]
For the last inequality, use $Ks\leq\epsilon_\beta$,
$\omega\leq\beta/100$, $\sigma\leq\sqrt2\,\beta/100<\beta/50$,
and $\delta,\tau\leq\beta/100$.
The new opposite angle exceeds $\pi-\tau>\pi-\beta$;
all old inequalities have quality $\delta<\beta$ at $z$.
The four cross angles just estimated complete the required packet.
All comparisons were made at actual endpoints within $\Omega$.
\end{proof}

\begin{theorem}[A chart from exclusion of the next rank; AC27]
\label{thm:alexandrov-rank-exclusion-chart}
Let $X$ be a locally complete length space with the packet and
parameters of AC26. Suppose there is no rank-$(m+1)$ quality-$\beta$
packet centered in $V$ with all anchors in $\Omega$.
Choose $q\in V$ and $r>0$ with
\[
 B(q,3r)\subset V,\qquad 2r\leq\ell_\beta.
\]
Then the map $F=(d(\cdot,a_1),\ldots,d(\cdot,a_m))$ on
$B(q,r)$ is a homeomorphism onto an open subset of $\mathbb R^m$,
and for all $x,y$ in that ball
\[
 \epsilon_\beta d(x,y)\leq|F(x)-F(y)|
                              \leq\sqrt m\,d(x,y).
\]
\end{theorem}
\begin{proof}
If $x\ne y$, then $L=d(x,y)<2r\leq\ell_\beta$ and
$\overline B(x,L)\subset B(q,3r)\subset V$.
If every coordinate difference were at most $\epsilon_\beta L$,
AC26 would give the excluded next-rank packet. Therefore some
coordinate difference is greater than $\epsilon_\beta L$,
which proves the Euclidean lower bound and injectivity. The upper
bound follows from the $1$-Lipschitz distance coordinates.

At each point of $B(q,r)$, AC23 applies with its quality parameter
$2\delta$: our strict quality-$\delta$ inequalities supply its
initial quality, and $2\delta\leq1/(100m)$.
Shrink the resulting neighborhood inside $B(q,r)$. These local open
maps show that the restriction to the whole ball is open. Combined
with injectivity and continuity, this gives the homeomorphism and
the Lipschitz inverse on its image. This does not assume that the
restricted metric on the ball is itself a length metric.
\end{proof}

\subsection{A finite hierarchy supplies the nearby chart}

\begin{theorem}[A chart in every comparison neighborhood; AC28]
\label{thm:alexandrov-finite-rank-chart}
Let $X$ be a non-singleton locally complete length space, let
$\Omega$ be a common comparison neighborhood with
$\dim_H\Omega\leq n$, where $n\geq1$ is an integer, and let
$O\subset\Omega$ be nonempty and open. Define
\[
 \tau_k=\frac{200^{\,k-(n+1)}}{400(n+1)}
          \quad(1\leq k\leq n+1),\qquad
 L_n=\max\left\{\sqrt n,\left(\frac{100\pi}{\tau_1}\right)^2\right\}.
\]
There are $q\in O$, $r>0$ and $1\leq m\leq n$ with
$B_X(q,r)\subset O$ and a distance-coordinate chart $\phi$ on
that ball, centered by $\phi(q)=0$, such that
\[
 L_n^{-1}d_X(x,y)\leq|\phi(x)-\phi(y)|\leq L_n d_X(x,y).
\]
Its image is open in $\mathbb R^m$. The distortion depends only
on $n$; no uniform lower bound on the chart radius is claimed.
\end{theorem}
\begin{proof}
For $1\leq k\leq n+1$, call $k$ admissible if there is a
rank-$k$ quality-$\tau_k$ packet centered in $O$ with all anchors
in $\Omega$. These are pointwise strict inequalities.

The rank $1$ is admissible. To see this, choose $p\in O$ and a
ball $B(p,\rho)\subset O$. A non-singleton length space has points
arbitrarily close to each of its points: along a continuous path to
another point, the distance from its starting point takes every
sufficiently small positive value. Choose $y\ne p$ with
$d(p,y)<\min\{1,\rho/2\}$, and apply AC25 with $x=p$ and
$\tau=\tau_1$. Its center $z$ and both anchors lie in $O$ and its
opposite angle is greater than $\pi-\tau_1$.

Every admissible $k$ satisfies $k\leq n$. Indeed AC23 with quality
parameter $2\tau_k$ applies, since
$2\tau_k\leq1/(100k)$. Its neighborhood may be shrunk inside $O$;
the Lipschitz open image and AC15 then give $k\leq\dim_H\Omega\leq n$.
Choose the largest admissible rank $m$, and one witnessing center
$q$ and anchors. In particular no rank-$(m+1)$ quality-$\tau_{m+1}$
packet with center in $O$ and anchors in $\Omega$ exists.

By continuity shrink an open neighborhood $V\subset O$ of $q$
so that the chosen anchors form a uniform packet of quality
$\delta=2\tau_m$, with distances in some $[a,A]$.
Take $\beta=\tau_{m+1}=200\tau_m$. Then
$\delta=\beta/100$ and $\beta\leq1/(200(m+1))$.
Choose $r>0$ with $B(q,3r)\subset V$ and $2r\leq\ell_\beta$.
AC27 applies. Its upper constant is at most $\sqrt n$ and its
inverse constant is
$1/\epsilon_\beta=(100\pi/\beta)^2\leq(100\pi/\tau_1)^2$.
Subtracting $F(q)$ gives the stated chart.
\end{proof}

A bounded chart ball in AC28 has finite rough dimension: its
$\epsilon$-separated subsets map to bounded Euclidean sets separated
by at least $\epsilon/L_n$. The finite grid argument of MC02 bounds
their cardinalities by $C(1+1/\epsilon)^m$, for a constant depending
on this chart. Thus its rough volume vanishes at every exponent
strictly greater than $m$. This is a patch near some $q\in O$,
not yet a neighborhood of every prescribed point of $O$.
Local completeness also gives a complete closed inner ball at $q$;
the bilipschitz image of that ball is complete and bounded in
Euclidean space, hence compact. Compactness at all points still
requires transport from this patch.

\begin{corollary}[The nearby ambient chart part of AC07; AC29]
\label{cor:alexandrov-ambient-paired-chart}
Let $X$ be a non-singleton complete length space, $p\in X$ and
$R>0$. Suppose $U=B_X(p,8R)$ has local lower comparison at least
$-1$ and ambient Hausdorff dimension at most an integer $n\geq1$,
in the conventions of AC07. Then there are $q\in B_X(p,R/2)$,
$\rho>0$ and $1\leq m\leq n$ with an ambient-distance chart
$\phi:B_X(q,\rho)\to\mathbb R^m$, $\phi(q)=0$, satisfying
\[
 L_n^{-1}d_X(x,y)\leq|\phi(x)-\phi(y)|\leq L_n d_X(x,y).
\]
No local compactness, regular tangent or angular obstruction is an
additional hypothesis of this corollary.
\end{corollary}
\begin{proof}
By AC13, $Y=(U,d_U)$ is a locally complete length space, with
ambient and intrinsic metrics agreeing on sufficiently small closed
balls. It is non-singleton because $X$ is a non-singleton length
space. Choose a common comparison neighborhood $\Omega$ of $p$
in $Y$, small enough to lie in $B_X(p,R/2)$ and in one such common
metric ball. Thus all pairwise distances on $\Omega$ agree with
$d_X$, and $\dim_H\Omega\leq n$. Apply AC28 with $O=\Omega$.
Its chart uses anchors in $\Omega$, so the coordinate distances
also agree with ambient distances. Shrink the chart radius at $q$
if necessary so that its intrinsic ball equals the ambient ball,
as provided by AC13, and remains in $\Omega$.
The bilipschitz bounds and centered value are unchanged.
\end{proof}

This proves the chart clause of AC07 in written mathematics, using
one dimension-dependent constant $L_n$ for all sources and allowing
arbitrarily small chart radii. Padding $\mathbb R^m$ into $\mathbb R^n$
gives AC04's target convention. The zero-dimensional/singleton
alternative remains AC18. The stronger tangent/strut route AC09--AC12
is preserved as a separate conditional result; its unbound regularity
inputs are no longer prerequisites for this particular chart producer.
AC01, AC02 and local compactness are not inputs to AC24--AC29.

The remaining geometric task is to transfer finite packing or rough
dimension from the bounded chart patch to a neighborhood of every
prescribed point, with curvature-$-1$ and almost-shortest-path errors
controlled. BGP Lemma 6.3 and BBI Corollary 10.8.20 indicate this
step, but their displayed radial arguments still require adaptation.
AC33 below transports that bound to every prescribed point;
AC17 then gives local compactness. Neither
AC07 as a whole nor the full MC18 limit statement is closed by the
new chart proof. Arbitrary burst improvement, regular tangents and
precise dimension identification remain separate future results.

\section{Finite packing transport before minimizing geodesics}
\label{sec:alexandrov-finite-packing-transport}

BGP Lemma 6.3 (printed 20/PDF 21) transfers rough dimension by
radial contraction. BBI Corollary 10.8.20 (printed 388--389/PDF
403--404) explicitly calls for an almost-shortest-path version on
each finite separated set \cite{BuragoGromovPerelman1992,BuragoBuragoIvanov}.
We give that version at curvature $-1$. Compactness below is first
used only for ten real distances, never for a subset of the source
space. The comparison domain is always one common neighborhood
$\Omega$ in the sense fixed above.

\subsection{Exact finite configurations and their stability}

\begin{lemma}[Finite radial comparison; AC30]
\label{lem:alexandrov-finite-radial-comparison}
In a metric space whose relevant finite configurations satisfy the
curvature-$-1$ four-point inequality, let $q,x,y,u,v$ obey
\[
\begin{gathered}
0<r=d(q,x)\leq D,\quad 0<s=d(q,y)\leq D,\quad 0<t<1,\\
d(q,u)=tr,\quad d(u,x)=(1-t)r,\\
d(q,v)=ts,\quad d(v,y)=(1-t)s.
\end{gathered}
\]
Then, with $\lambda=tD/\sinh D>0$,
\[
 d(u,v)\geq\lambda d(x,y).
\]
Only these five points are required; the exact metric equalities do
not assert the existence of paths realizing them.
\end{lemma}
\begin{proof}
First prove a one-sided assertion. Suppose $d(q,u)=a>0$,
$d(u,x)=b>0$ and $d(q,x)=a+b$, and let $z\ne q$.
If $q,u,x,z$ are distinct, write $l=d(u,z)>0$,
$S=d(q,z)$ and $C=d(x,z)$. The comparison angle
$\theta_u(q,x)$ is $\pi$, so four-point comparison gives
$\theta_u(q,z)+\theta_u(x,z)\leq\pi$.
For two angles in $[0,\pi]$, a sum at most $\pi$ is equivalent
to their cosines having nonnegative sum. The hyperbolic cosine law
therefore gives
\[
 \sinh(a+b)\cosh l\geq
       \sinh b\cosh S+\sinh a\cosh C.
\]
In the model triangle with sides $a+b,S,C$, the distance from
$\bar z$ to the point $\bar u$ at distance $a$ along
$[\bar q,\bar x]$ has hyperbolic cosine equal to the right side divided by
$\sinh(a+b)$, by the same cosine law and the addition formulas.
Thus $d(u,z)\geq d(\bar u,\bar z)$, and comparison at $q$ gives
$\theta_q(u,z)\geq\theta_q(x,z)$.
If $z=u$ or $z=x$, both angles in this last inequality are zero;
these are the only excluded coincidences. Here an angle with equal
noncentral endpoints is defined to be zero.

Apply this assertion first to $q,u,x$ with $z=y$, and then to
$q,v,y$ with $z=u$. It follows that
\[
 \theta_q(u,v)\geq\theta_q(u,y)\geq\theta_q(x,y).
\]
If $x=y$, the desired conclusion is automatic. Otherwise let
$c=d(x,y)$ and let $c_t$ be the model distance between the points
at radii $tr,ts$ with included angle $\theta=\theta_q(x,y)$.
Monotonicity of the cosine law gives $d(u,v)\geq c_t$.
For completeness the purely model estimate is
\[
 \sinh^2(c/2)=\sinh^2((r-s)/2)
       +\sinh r\sinh s\sin^2(\theta/2).
\]
For $0\leq h\leq D$, monotonicity of $\sinh h/h$ gives
$\sinh(th)\geq th\geq\lambda\sinh h$.
Apply this to $r,s,|r-s|/2$ in the displayed identity to obtain
$\sinh(c_t/2)\geq\lambda\sinh(c/2)$.
Since $0<\lambda\leq1$ and $\operatorname{arsinh}$ is concave and
vanishes at zero, $c_t\geq\lambda c$.
This rederives the model part of AC03 without importing AC03's
minimizing-path or compactness hypotheses.
\end{proof}

\begin{lemma}[Uniform finite-distance stability; AC31]
\label{lem:alexandrov-radial-defect-stability}
Fix $0<a\leq D$, $0<t<1$ and $\epsilon>0$, and put
$\lambda=tD/\sinh D$. There is $\eta>0$, depending on these
four parameters, such that every five-point metric configuration
satisfying four-point comparison and
\[
\begin{gathered}
r=d(q,x),\ s=d(q,y),\quad r,s\in[a,D],\qquad d(x,y)\geq\epsilon,\\
d(q,u)=tr,\quad (1-t)r\leq d(u,x)\leq(1-t)r+\eta,\\
d(q,v)=ts,\quad (1-t)s\leq d(v,y)\leq(1-t)s+\eta
\end{gathered}
\]
satisfies $d(u,v)\geq(\lambda/2)d(x,y)$.
The same $\eta$ works for all such configurations; no effective
formula for it is asserted.
\end{lemma}
\begin{proof}
If no such $\eta$ exists, choose a violating configuration with
$\eta=1/j$ for each positive integer $j$. Every point is at distance
at most $D$ from $q$, so its ten pairwise distances lie in
$[0,2D]^{10}$. A subsequence of these real arrays converges.
Its limit is a pseudometric on five labels; quotient labels at zero
distance. Triangle inequalities and all displayed radial equalities
pass to this finite metric quotient. In particular,
$d(q,x),d(q,y)\geq a$, $d(q,u),d(q,v)\geq ta$,
$d(u,x),d(v,y)\geq(1-t)a$, and $d(x,y)\geq\epsilon$.

For any four distinct points of the quotient, choose their label
representatives. Their distances are positive in the limit, so the
corresponding original points are distinct for all sufficiently large
$j$. The comparison angles depend continuously on their distances,
including degenerate triangles when both central sides are positive.
Their four-point inequality therefore passes to the limit.
AC30 applies to the quotient and gives $d(u,v)\geq\lambda d(x,y)$.
The violating inequalities give
$d(u,v)\leq(\lambda/2)d(x,y)$ in the limit, a contradiction since
$d(x,y)\geq\epsilon>0$. Possible collisions between other labels,
including $u=v$, do not invalidate this quotient argument.
\end{proof}

\subsection{Finite sets suffice for rough dimension}

\begin{proposition}[Packing transport through almost-short paths; AC32]
\label{prop:alexandrov-finite-packing-transport}
Let $Y$ be a length space and $\Omega$ a common comparison
neighborhood for its metric. Suppose
\[
 q\in\Omega,\qquad B_Y(q,\rho)\subset\Omega,\qquad
 W\subset\Omega\cap\{x:a\leq d(q,x)\leq D\},
\]
where $\rho>0$ and $0<a\leq D$. Fix $0<t<1$ with $tD<\rho$,
and set $\lambda=tD/\sinh D$.
For every $\epsilon>0$ and every finite $\epsilon$-separated set
$S\subset W$, there is a map $f_S:S\to B_Y(q,\rho)$ such that
\[
 d(f_S(x),f_S(y))\geq(\lambda/2)d(x,y)
       \quad(x,y\in S).
\]
Consequently, with the non-strict packing convention of AC16,
\[
 P_W(\epsilon)\leq P_{B_Y(q,\rho)}(\lambda\epsilon/2).
\]
No local completeness, local compactness, or minimizing path is
assumed. The map may depend on the finite set and on $\epsilon$.
\end{proposition}
\begin{proof}
Choose $\eta$ from AC31 for $a,D,t,\epsilon$. For each $x\in S$,
choose a rectifiable path from $q$ to $x$ of length less than
$r_x+\eta$, where $r_x=d(q,x)$. Let $f_S(x)$ be its first point
at distance $tr_x$ from $q$. Continuity on the compact parameter
interval supplies this point, and the initial path remains in the
closed $tr_x$-ball. Triangle inequality and additivity of path length
give
\[
 (1-t)r_x\leq d(f_S(x),x)<(1-t)r_x+\eta.
\]
Also $f_S(x)\in B_Y(q,\rho)$. For distinct $x,y\in S$, all five
points $q,x,y,f_S(x),f_S(y)$ lie in $\Omega$, so AC31 applies.
There is no need for the rest of either path to stay in $\Omega$:
only these actual endpoints enter the comparison inequalities.
The resulting separated image has the same cardinality as $S$.
Taking suprema over finite $S$ proves the packing bound, including
its extended-value interpretation.
\end{proof}

\begin{theorem}[Local rough dimension is produced; AC33]
\label{thm:alexandrov-local-rough-dimension-produced}
The producer AC16 has a written proof: in a locally complete length
space with local lower comparison $-1$ and a finite local Hausdorff
bound, every prescribed point has a bounded neighborhood of finite
rough dimension. No local compactness or minimizing-geodesic
hypothesis is added.
\end{theorem}
\begin{proof}
A singleton has packing number one and rough dimension zero. Otherwise
fix the prescribed point $p$. Choose a common comparison neighborhood
$\Omega$ of $p$ within a neighborhood of finite Hausdorff dimension,
and an integer $n\geq1$ bounding that dimension. AC28, with centers
in $\Omega$ and all anchors in $\Omega$, supplies $q\in\Omega$,
$\rho>0$ and a $\Lambda$-bilipschitz map
$F:B_Y(q,\rho)\to\mathbb R^m$, where $1\leq m\leq n$.
Shrink $\rho$ so that this ball lies in $\Omega$, and translate
$F(q)=0$. If $q=p$, a smaller chart ball is already the desired
neighborhood, by the Euclidean packing estimate below.

If $q\ne p$, choose $0<h<d(q,p)/2$ with $W=B_Y(p,h)\subset\Omega$.
Put $a=d(q,p)-h>0$, $D=d(q,p)+h$, and choose
$t=\min\{1/2,\rho/(2D)\}$. Then $W$ lies in AC32's annulus
and $tD<\rho$. For a finite $\sigma$-separated subset of the chart
ball, its $F$-image lies in $[-\Lambda\rho,\Lambda\rho]^m$
and is $(\sigma/\Lambda)$-separated. Partition each coordinate
interval into at most $2+4\Lambda^2\rho\sqrt m/\sigma$ half-open
intervals of length at most $\sigma/(2\Lambda\sqrt m)$, assigning
the outer endpoints to their adjacent cells. Each resulting box
contains at most one image point. AC32 therefore gives
\[
 P_W(\epsilon)\leq
 \left(2+\frac{8\Lambda^2\rho\sqrt m}{\lambda\epsilon}\right)^m.
\]
For every $b>m$, multiplication by $\epsilon^b$ tends to zero as
$\epsilon\downarrow0$. Thus $\dim_{\rm rough}W\leq m<\infty$.
The same grid estimate directly treats $q=p$.
The constants and neighborhood may depend on $p$ and the source;
no family-uniform radius or packing constant is inferred here.
\end{proof}

\paragraph{Local structure and the extraction consumer.}
AC33 supplies AC16, so AC17 gives a compact ambient closed-ball
neighborhood at every point of a controlled region. Together with
AC29's nearby chart and AC18's singleton branch, this supplies both
clauses of AC07 in written mathematics. The dependency order is
AC28, then AC30--AC33, then AC16/AC17, and only then the
compact-domain comparison and covering in AC01--AC04.
There is no appeal to AC01 or AC02 in AC30--AC33.
AC08 can now use this producer with the dimension-only distortion
from AC29, and MC13 supplies metric extraction. AC02 still records
its source-qualified local Toponogov input and later formal binding;
this revision does not prove that source theorem anew.

Local compactness and finite covering do not yet identify the limit
as a space with the required curvature and dimension. MC14 now
supplies length closure; AC34--AC37 below pass comparison to the
limit. AC38--AC40 below supply dimension closure, and AC41 assembles
the compactness statement with the local comparison qualification retained.
Nor does the present one-way packing transport prove equality of
local dimensions, full burst improvement, or regular-tangent results.
No new Lean declaration or acceptance of the full KL Lemma 3.10 is
claimed.

\section{Passing lower comparison to the pointed limit}
\label{sec:alexandrov-limit-comparison}

BBI Proposition 10.7.1 (printed 376/PDF 391) passes the four-point
condition to a metric limit. Our sources have curvature bounds only
on growing regions, and the bounds vary. We therefore separate finite
point lifting, model continuity at zero curvature, and a controlled
domain for the source inequalities. Throughout this section distances
are finite. For $k\leq0$ and $a,b>0$, write
$\Theta_k(a,b,h)\in[0,\pi]$ for the model angle between sides of
lengths $a,b$, opposite the side of length $h$, where
$|a-b|\leq h\leq a+b$. Degenerate triangles are included. In a metric
space put
\[
 \theta^k_x(y,z)=\Theta_k(d(x,y),d(x,z),d(y,z))
       \qquad(y,z\ne x).
\]
The four-point inequality is the sum of the three such angles at
one vertex being at most $2\pi$. If two outer points coincide, that
inequality holds automatically: the sum is twice an angle in
$[0,\pi]$. Thus only four distinct points need a curvature hypothesis.

\begin{lemma}[Lifting a finite configuration; AC34]
\label{lem:alexandrov-finite-lifting}
Let $f:\overline B_Z(o,R)\to Y$ satisfy
$\mathcal A_{R,\epsilon}$, where $0<\epsilon<1/4$.
Let $y_0,\ldots,y_m\in Y$, with distinguished point $q\in Y$, and put
\[
 M=\max\{1,d(q,y_0),\ldots,d(q,y_m)\},\qquad R=M+2.
\]
Here $o,q$ are the basepoints in $\mathcal A_{R,\epsilon}$.
There are $z_0,\ldots,z_m$ in the actual domain of $f$ such that
\[
 \begin{split}
 d_Y(f(z_a),y_a)&<\epsilon,\qquad
 d_Z(o,z_a)<d_Y(q,y_a)+2\epsilon<R,\\
 |d_Z(z_a,z_b)-d_Y(y_a,y_b)|&<3\epsilon.
 \end{split}
\]
If the $y_a$ are pairwise distinct with minimum pairwise distance
$\delta>0$, then the $z_a$ are pairwise distinct whenever
$3\epsilon<\delta$. No continuity of $f$ is required.
\end{lemma}
\begin{proof}
All $y_a$ lie in $\overline B_Y(q,R-\epsilon)$, so coverage gives
actual witnesses $z_a$. MC06 fixes $f(o)=q$ exactly. Thus
\[
 d_Z(o,z_a)<d_Y(q,f(z_a))+\epsilon
            <d_Y(q,y_a)+2\epsilon<R.
\]
For a pair, the distortion contributes less than $\epsilon$ and the
two coverage errors contribute less than $2\epsilon$. This gives the
three-error bound. If $z_a=z_b$ for distinct target labels, that bound
would imply $d_Y(y_a,y_b)<3\epsilon<\delta$, a contradiction.
\end{proof}

\begin{lemma}[Comparison angles as negative curvature tends to zero; AC35]
\label{lem:alexandrov-angle-zero-limit}
Suppose $\kappa_j\geq0$ tends to zero and
$(a_j,b_j,h_j)\to(a,b,h)$, where $a,b>0$ and each triple has
positive first two entries and satisfies the triangle inequalities.
Then
\[
 \Theta_{-\kappa_j}(a_j,b_j,h_j)
       \longrightarrow\Theta_0(a,b,h).
\]
This includes $h=|a-b|$ and $h=a+b$, and does not require the
angles to stay away from zero or $\pi$.
\end{lemma}
\begin{proof}
Write $s=\sqrt\kappa$. For $s>0$ the hyperbolic cosine law gives
\[
 \cos\Theta_{-\kappa}(a,b,h)=
 \frac{\cosh(sa)\cosh(sb)-\cosh(sh)}{\sinh(sa)\sinh(sb)}.
\]
The length scale is $\sqrt\kappa$, not $\kappa$. This normalization
is also the pinned AKP \mbox{\code{model.tex}} cosine law, with curvature
parameter $-\kappa$. Define continuous functions on $[0,\infty)$ by
\[
 S(u)=\begin{cases}\sinh(u)/u&u>0,\\1&u=0,\end{cases}
 \qquad
 Q(u)=\begin{cases}(\cosh(u)-1)/u^2&u>0,\\1/2&u=0.\end{cases}
\]
Their values at zero follow, for example, from the convergent power
series for $\sinh$ and $\cosh$. Dividing numerator and denominator
by $s^2$ rewrites the cosine as
\[
 \frac{a^2Q(sa)+b^2Q(sb)-h^2Q(sh)
              +s^2a^2b^2Q(sa)Q(sb)}{abS(sa)S(sb)}.
\]
At $s=0$ this is $(a^2+b^2-h^2)/(2ab)$, the Euclidean formula.
The denominator tends to $ab>0$ when the side lengths also vary.
Thus the cosines converge, including indices with $\kappa_j=0$.
Every cosine lies in $[-1,1]$ by the triangle inequalities.
Continuity of $\arccos$ on that closed interval proves the result;
no derivative estimate for $\arccos$ is used.
\end{proof}

\begin{lemma}[Four-point comparison on an interior ball; AC36]
\label{lem:alexandrov-buffered-four-point}
Let $R>0$, $\kappa\geq0$, and let $(Z,o)$ be a complete length
space. Suppose $U=B_Z(o,16R)$ is locally compact, and its intrinsic
length metric $d_U$ has finite Hausdorff dimension and local
Alexandrov curvature at least $-\kappa$. Using the local Toponogov
input recorded in AC02, every four distinct points
$x,y_1,y_2,y_3\in B_Z(o,R)$ satisfy, for ambient distances,
\[
 \theta^{-\kappa}_x(y_1,y_2)+
 \theta^{-\kappa}_x(y_2,y_3)+
 \theta^{-\kappa}_x(y_3,y_1)\leq2\pi.
\]
The whole space $Z$ is not assumed proper, and $U$ is not assumed
complete.
\end{lemma}
\begin{proof}[Domain proof and reduction to the source comparison input]
As in AC02, almost-short paths from $o$ show that $U$ is path
connected. Its intrinsic metric is finite, dominates $d_Z$, and
agrees with $d_Z$ on a sufficiently small neighborhood of every
point. In particular it has the same local topology.
AC01 gives ambient minimizing paths from $x$ to each $y_a$.
Their lengths are less than $2R$. The nearer-endpoint estimate puts
all their points in $B_Z(o,2R)$. Any two points $v,w$ on these
paths admit an ambient minimizing path in $B_Z(o,4R)$: choose
$r<2R$ containing both and apply AC01. Hence
\[
 d_U(v,w)=d_Z(v,w)<4R.
\]
The same applies to the endpoints. These paths are minimizing for
$d_U$ too, since $d_U\geq d_Z$ and their lengths agree locally.

For every $v\in B_Z(o,2R)$, the intrinsic closed $10R$-ball at $v$
is complete. Indeed an intrinsic Cauchy sequence is ambient Cauchy,
so has an ambient limit $w$. Its radii give
\[
 d_Z(o,w)\leq d_Z(o,v)+10R<12R<16R.
\]
Thus $w\in U$; local equality of metrics gives intrinsic convergence,
and closedness retains the $10R$-bound. We may therefore use
\[
 D=5R,\qquad 2D=10R
\]
for every triangle with vertices on the chosen radial paths.
Every side is strictly shorter than $D$, and the required intrinsic
closed $2D$-balls at all vertices are complete. This is precisely
the safeguard of KL Section 3.3 (printed 23/PDF 18). The local
Toponogov theorem supplies comparison-angle monotonicity along
each of the two radial sides of every such triangle. There is no
positive-curvature perimeter restriction because $-\kappa\leq0$.

Choose a positive common radial distance smaller than all three
path lengths and so small that the three radial points $u_a$ and
$x$ lie in one local four-point neighborhood at $x$. Monotonicity
and the metric agreement just proved give
\[
 \theta^{-\kappa}_x(y_a,y_b)
       \leq\theta^{-\kappa}_x(u_a,u_b).
\]
Summing and using the local four-point inequality gives the claim.
If radial points coincide, the repeated-outer-point observation at
the start of this section gives the same sum bound. No assertion
that the geodesics separate near $x$ is needed.
\end{proof}

The factor $16$ and the constants $5,10,12$ are sufficient project
margins, not constants quoted from KL. The local globalization theorem
itself remains source-qualified. BBI 10.3.1 is its complete-space
comparison; its printed proof uses local compactness. The retained
BBI errata correct the limiting neighborhood and the failing-angle
choice on printed 361, and the figure references on 362. None of
those corrections turns an incomplete open region into a complete
space. Our argument verifies the separate KL domain safeguard.

\begin{theorem}[The extracted limit has nonnegative comparison; AC37]
\label{thm:alexandrov-limit-nonnegative-comparison}
Under the source hypotheses of MC18, let $(Y,q)$ be a complete
proper pointed metric limit of a subsequence, as produced by
MC13 and AC08. Then $Y$ is a geodesic length space and satisfies
the global curvature-zero four-point inequality. For every geodesic
$\gamma:[0,1]\to Y$ from $a$ to $b$, parametrized proportionally
to length, every $v\in Y$ and $0\leq t\leq1$, it also satisfies
\[
 d(v,\gamma(t))^2\geq
 (1-t)d(v,a)^2+t\,d(v,b)^2-t(1-t)d(a,b)^2.
\]
This is the Euclidean lower triangle-comparison inequality. The
Hausdorff bound $\dim_HY\leq n$ is not a conclusion of this theorem;
it is still required for MC18 and KL's finite-dimensional convention.
\end{theorem}
\begin{proof}
MC14 supplies the length property and geodesics. Fix four distinct
points $y_0,y_1,y_2,y_3\in Y$, set
$M=\max\{1,d(q,y_0),\ldots,d(q,y_3)\}$ and $R=M+2$.
Work along the already convergent subsequence; no new limit is chosen.
For each $j$, eventual pointed convergence supplies an increasing
source index $i_j$ and an actual map
\[
 \mathcal A_{R,\epsilon_j}
   (f_j:(X_{i_j},p_{i_j})\to(Y,q)),\qquad
 \epsilon_j=1/(100(j+1)),
\]
with $c_{i_j}\leq1$ and $r_{i_j}>16R$. All these requirements
hold on a tail for each $j$, so the indices can be chosen recursively.
AC34 lifts the four points to $z_{j,a}\in B_{X_{i_j}}(p_{i_j},R)$.
Their six distances differ from the target distances by less than
$3\epsilon_j$; the lifts are distinct for all sufficiently large $j$.

Here is the domain translation needed before applying AC36.
AC13 supplies complete inner neighborhoods and equality of the
ambient and intrinsic metrics locally. The source curvature bound
$-c_{i_j}\geq-1$ and the finite local dimension bound allow AC07
to supply local compactness throughout the controlled region.
AC01 then makes the ambient closed $16R$-ball compact, because
$16R<r_{i_j}$. In particular $U_j=B_{X_{i_j}}(p_{i_j},16R)$ is
separable. Its intrinsic metric has the same topology and is locally
equal to the ambient metric. A countable subcover of these equality
neighborhoods shows that its Hausdorff dimension is at most $n$:
on each neighborhood the dimensions agree, and Hausdorff dimension
of a countable union is the supremum. This uses only the source
dimension bound, not dimension closure for $Y$.
The local curvature interpretation likewise agrees. Thus AC36 applies
to $U_j$ with its actual intrinsic metric.

At the source center $z_{j,0}$, AC36 gives
\[
 \sum_{1\leq a<b\leq3}
 \Theta_{-c_{i_j}}
   \bigl(d(z_{j,0},z_{j,a}),d(z_{j,0},z_{j,b}),
                 d(z_{j,a},z_{j,b})\bigr)\leq2\pi.
\]
Every central distance has a positive limit. AC35 passes each of
the three angles to its Euclidean comparison angle, including
collinear triples. The resulting four-point inequality holds for
the chosen quadruple, hence for every quadruple in $Y$. No source
comparison-neighborhood radius is required to be uniform in $j$:
AC36 is applied separately on each complete interior buffer.

For completeness, the triangle consequence has a direct algebraic
proof. Put $z=\gamma(t)$, $L=d(a,b)$, $A=d(v,a)$, $B=d(v,b)$
and $h=d(v,z)$. Cases $L=0$, $t=0,1$, or $v\in\{a,b,z\}$ give
equality directly. Otherwise let $\alpha=d(a,z)=tL>0$ and
$\beta=d(z,b)=(1-t)L>0$. At $z$, the comparison angle between
$a$ and $b$ is $\pi$, so the four-point inequality implies
\[
 \theta^0_z(a,v)+\theta^0_z(v,b)\leq\pi.
\]
For two angles in $[0,\pi]$ this implies that the sum of their
cosines is nonnegative. Substituting the Euclidean cosine formulas,
\[
 \frac{\alpha^2+h^2-A^2}{2\alpha h}
 +\frac{\beta^2+h^2-B^2}{2\beta h}\geq0.
\]
Multiplication by the positive denominator and rearrangement give
\[
 h^2\geq\frac{\beta A^2+\alpha B^2}{\alpha+\beta}
                 -\alpha\beta,
\]
which is the stated formula. The right side is the squared Euclidean
distance from the comparison vertex to the corresponding point of
the opposite side. This includes degenerate comparison triangles.
\end{proof}

\paragraph{What is now supplied to MC18.}
AC34--AC37 write the lower-comparison passage for the existing proper
pointed limit, using the explicitly retained local Toponogov input.
They use neither global source properness nor continuity of the
approximation maps, and never equip a closed restricted ball with
an unproved length-space structure. Length closure and identification
remain in MC14. AC38--AC40 below prove dimension closure separately.
In particular, calling
$Y$ finite-dimensional at the start of the limit proof would be
circular. The stronger tangent and exact-dimension results remain
separate from this comparison argument.

\section{Hausdorff dimension and the assembled compactness bridge}
\label{sec:alexandrov-dimension-closure}

AC04 supplies more than finiteness of nets: its bound has a uniform
power $\epsilon^{-n}$. The next argument preserves that exponent under
pointed approximation and uses the resulting covers directly.
BBI Corollary 10.8.25 (printed 390/PDF 405) instead uses the equality
of Hausdorff dimension and strainer number, together with continuity
of comparison angles. That stronger identification is unnecessary for
this upper bound. The source definitions are BBI 1.7.7--1.7.8 and
1.7.17--1.7.19 (printed 19--23/PDF 34--38); the quantitative packing
comparison is BBI Lemma 10.9.2 (printed 390--391/PDF 405--406).
Exact reading, corrections and scope are recorded in
\code{reference_checks_revision70.md}.

\begin{lemma}[Transferring internal nets to a limit ball; AC38]
\label{lem:alexandrov-net-transfer}
Let $(Z,o),(Y,q)$ be pointed metric spaces, $R>0$, $\delta>0$, and
$0<\varepsilon<\min(1/4,\delta/8)$. Suppose $f$ satisfies
$\mathcal A_{R+2,\varepsilon}((Z,o),(Y,q))$ in MC06's convention,
and $F\subset\overline B_Z(o,R+1)$ is an internal
$\delta/4$-net with at most $N$ points. Then
$\overline B_Y(q,R)$ has an internal strict $\delta$-net of at most
$N$ points. No continuity of $f$ is required.
\end{lemma}
\begin{proof}
For $y\in\overline B_Y(q,R)$, coverage supplies $z$ in the domain
of $f$ with $d_Y(f(z),y)<\varepsilon$. Exact preservation of the
basepoint and distortion give
\[
 d_Z(o,z)<d_Y(q,y)+2\varepsilon<R+1.
\]
Choose $a\in F$ with $d_Z(z,a)\leq\delta/4$. Both points are in
$f$'s domain, so
\[
 d_Y(y,f(a))<\delta/4+2\varepsilon<\delta/2.
\]
The centers $f(a)$ need not belong to the target $R$-ball.
For each nonempty set
\[
 V_a=B_Y(f(a),\delta/2)\cap\overline B_Y(q,R),
\]
choose one point $y_a\in V_a$ and discard empty sets. The $V_a$
cover the target ball, and $y,y_a\in V_a$ implies
$d_Y(y,y_a)<\delta$. Thus the chosen points give the required
internal net. There are at most $|F|\leq N$ of them. Only finitely
many choices are involved; no compactness or length property is used.
\end{proof}

\begin{lemma}[Polynomial nets give a Hausdorff upper bound; AC39]
\label{lem:alexandrov-polynomial-hausdorff}
Let $A$ be a metric space, $n\geq0$, and $0<C<\infty$.
Suppose that for every $0<\delta\leq1$ there is an internal
$\delta$-net of size at most $C\delta^{-n}$. Then
$\mathcal H^s(A)=0$ for every $s>n$, and $\dim_H A\leq n$.
The same vanishing conclusion holds for a countable union of subsets
each having zero $s$-dimensional Hausdorff measure in the restricted
metric.
\end{lemma}
\begin{proof}
For $s>0$ and $\rho>0$, use the unnormalized Hausdorff content
\[
 \begin{gathered}
 \mathcal H^s_\rho(A)=
 \inf\left\{\sum_j (\operatorname{diam}U_j)^s:
 A\subseteq\bigcup_j U_j,\quad
 \operatorname{diam}U_j<\rho\right\},\\
 \mathcal H^s(A)=\lim_{\rho\downarrow0}\mathcal H^s_\rho(A).
 \end{gathered}
\]
Covers are finite or countable; the empty set has the empty cover.
A positive normalization depending only on $s$ does not affect
vanishing or dimension. All exponents used below are positive, so
no convention about $0^0$ is needed.

Fix $s>n$, a cutoff $\rho>0$ and a budget $b>0$. Since $s-n>0$,
choose $0<\delta\leq1$ such that $2\delta<\rho$ and
$2^s C\delta^{s-n}<b$. The intersections with $A$ of closed
$\delta$-balls at the net centers cover $A$ and have diameter at
most $2\delta$. Their total weight is at most
\[
 (C\delta^{-n})(2\delta)^s
       =2^s C\delta^{s-n}<b.
\]
Hence $\mathcal H^s_\rho(A)=0$ for every $\rho>0$, and
$\mathcal H^s(A)=0$. With
$\dim_H A=\inf\{s>0:\mathcal H^s(A)=0\}$ this gives the claimed
upper bound, including $n=0$.

For $A=\bigcup_{j\geq1}A_j$ with $\mathcal H^s(A_j)=0$,
fix the same $\rho,b$. Cover $A_j$ by sets of diameter $<\rho$
and total weight $<b2^{-j}$. Such covers exist because every
$\mathcal H^s_\rho(A_j)$ is zero. Restricted and ambient diameters
agree, and the combined countable cover has total weight $<b$.
Taking arbitrary $b$ and then arbitrary $\rho$ proves the assertion.
No uniform constant across the subsets $A_j$ is required.
\end{proof}

\begin{theorem}[Dimension of the extracted pointed limit; AC40]
\label{thm:alexandrov-limit-dimension}
Use MC18's source hypotheses and an existing pointed limit $(Y,q)$
in MC06's approximation sense along a subsequence. Suppose the
AC08 nets are supplied with the same $L=L(n)$ for all sufficiently
late controlled source balls. Then, for every $R>0$ and
$0<\delta\leq1$, the closed $R$-ball in $Y$ has an internal
$\delta$-net with at most $C_R\delta^{-n}$ points, where
\[
 A_R=16L(n)^2\sqrt n\,\sinh(2(R+1)),
 \qquad C_R=(2+A_R)^n.
\]
Consequently $\mathcal H^s(Y)=0$ for all $s>n$, and
$\dim_H Y\leq n$. No curvature or dimension hypothesis on the
limit is used in this proof.
\end{theorem}
\begin{proof}
Fix $R,\delta$ and then choose
$0<\varepsilon<\min(1/4,\delta/8)$. On a sufficiently late index
of the given convergent subsequence there is an actual
$\mathcal A_{R+2,\varepsilon}$ map, and AC08 applies at radius
$R+1$ and mesh $\delta/4$. In particular, require
$c_i\leq1$ and $r_i>8(R+1)$ as well as the approximation condition.
The source has an internal net of size at most
\[
 N(n,L(n),R+1,\delta/4)
       =\bigl(1+\lceil A_R/\delta\rceil\bigr)^n.
\]
AC38 transfers this cardinality bound to the target ball. Since
$\lceil u\rceil\leq u+1$ and $\delta\leq1$,
\[
 \bigl(1+\lceil A_R/\delta\rceil\bigr)^n
 \leq (2+A_R/\delta)^n
 \leq (2+A_R)^n\delta^{-n}=C_R\delta^{-n}.
\]
The constant $C_R$ is finite and depends only on $n,R$ and the
fixed $L(n)$, not on the source index or $\delta$. The selected
index may depend on $R,\delta,\varepsilon$; no single tail works
for all scales by assumption. No lower bound on the chart radii
is required. AC39 now gives
$\mathcal H^s(\overline B_Y(q,R))=0$ for every $s>n$.
Since distances are finite,
\[
 Y=\bigcup_{m=1}^{\infty}\overline B_Y(q,m).
\]
The countable-union conclusion of AC39 proves global vanishing
and the dimension bound. This step does not need AC37's comparison
conclusion, exact dimension identification or noncollapse.
\end{proof}

\begin{corollary}[Assembled growing-region compactness bridge; AC41]
\label{cor:alexandrov-assembled-extraction}
Under MC18's hypotheses, the written producers above yield a
subsequence converging to a complete proper geodesic pointed space
with nonnegative Alexandrov comparison and Hausdorff dimension
at most $n$. This is the conclusion of KL Lemma 3.10, conditional
on the retained source-qualified local Toponogov theorem in
AC02/AC36 and the stated metric and analytic inputs.
\end{corollary}
\begin{proof}
AC07 is supplied by AC29 and AC33/AC17. AC08 therefore gives
uniform eventual internal nets and uses MC13 to extract one complete
proper pointed metric limit. MC14 gives length closure and geodesicity
for that limit. AC37 supplies its zero-curvature four-point and lower
triangle comparison, with the explicit source completeness buffers.
AC40 supplies its Hausdorff-dimension bound independently of that
comparison conclusion. These properties give a nonnegatively curved
Alexandrov space in KL Section 3.3's finite-dimensional convention.
No global properness of the source spaces is added, and no second
subsequence limit is substituted. This establishes the conclusion
stated in Theorem~\ref{thm:metric-kl-extraction}.
\end{proof}

The written compactness bridge is now assembled. This does not accept
a Lean declaration or remove the local Toponogov source dependency.
The dependency graph points from MC18 to AC41; AC41 uses the earlier
lemmas and MC18's \emph{hypotheses}, not MC18 as an input theorem.
The upper dimension bound allows dimension loss, as in MC16; it
asserts neither positive $n$-dimensional measure nor equality with $n$.
Splitting, classification of low-dimensional limits, fibration and
topological reconstruction still require their own contracts.

\clearpage
\section{Splitting and low-dimensional product models}
\label{sec:alexandrov-splitting-models}

The next consumers are KL Lemmas 8.1 and 9.1, printed 47 and 49/PDF
42 and 44. They use exact splittings of a limit whose Hausdorff
dimension is at most two. The existence and dimension of that collapsed
limit require additional hypotheses: KL Lemma 6.18 is a separate
collapse input, not a consequence of MC18 with $n=3$.
This section supplies the static product geometry. Chapter 13 owns the
collapse application; Chapter 14 owns adapted fibrations and compatibility.
A general classification of two-dimensional Alexandrov spaces is not
needed for the particular products below.

\subsection{Exact splitting data and the global line theorem}

\begin{definition}[Pointed Euclidean splitting; AC42]
\label{def:alexandrov-exact-splitting}
A pointed $k$-splitting, $k\in\mathbb N$, of a nonempty pointed metric
space $(X,p)$ consists of a nonempty pointed metric space $(Y,y_0)$ and
an onto isometry
\[
 \alpha:X\longrightarrow\mathbb R^k\times Y,
 \qquad \alpha(p)=(0,y_0),
\]
where the product distance is
\[
 d((u,y),(v,z))=\sqrt{\|u-v\|_2^2+d_Y(y,z)^2}.
\]
For a specified unit-speed line $\gamma:\mathbb R\to X$, meaning
$d_X(\gamma(s),\gamma(t))=|s-t|$ for all real $s,t$, a $1$-splitting
is aligned with $\gamma$ if $p=\gamma(0)$ and
$\alpha(\gamma(t))=(t,y_0)$. A ray only has domain $[0,\infty)$
and is not a line. Both coordinate projections are $1$-Lipschitz;
this follows directly from the product formula.
\end{definition}

KL Section 2.2 and Definitions 4.1--4.2 use this Pythagorean product
metric. Normalizing the Euclidean basepoint to zero is a translation,
not a restriction on the transverse basepoint $y_0$. Equivalence of
splittings uses isometries of the two factors; it does not identify
arbitrary coordinate maps. Compatibility and approximate splittings
will require separate contracts. In particular, a distance-preserving
map into a product is insufficient without surjectivity.

\begin{theorem}[Global line splitting input; AC43]
\label{thm:alexandrov-line-splitting-input}
Let $X$ be a nonempty complete proper geodesic Alexandrov space of
global curvature at least zero and finite Hausdorff dimension.
Given a unit-speed line $\gamma:\mathbb R\to X$, there exists an
aligned pointed $1$-splitting in the sense of AC42. Its transverse
factor may be chosen to be a complete proper geodesic nonnegatively
curved Alexandrov space.
\end{theorem}

This is a \emph{source-qualified theorem}, not a new proof of splitting.
BBI Theorem 10.5.1 and its proof, printed 366--369/PDF 381--384,
construct parallel lines through every point, prove the projection
relation symmetric and transitive, and identify the transverse slice.
BBI's Chapter 10 convention includes completeness; properness here
supplies the local compactness required by its theorem. The chosen
original line is the Euclidean factor; its unit-speed parameter fixes
orientation and scale after translating the origin. AC44 below
verifies the factor qualifications for any such onto product.
The retained author errata correct the parallel-line definition on
366, the comparison-distance chain on 367 and the line through $x$
on 369. All three enter this source route.

AKP Theorem 16.22, archived printed/PDF 251--253, gives a comparison
proof through affine Busemann functions and their gradient flows.
Its direct-sum definition on 251 explicitly requires the projections
to give an onto map. This matches the published-edition correction at
page 228, but the two page number systems are distinct. We retain
BBI's parallel-line route as the selected implementation plan; the
AKP gradient-flow route is an alternative with additional unbound
inputs. No proof of those inputs is inferred here.
A curvature bound on a bounded open region does not authorize AC43.
Even the complete nonnegatively curved ray $[0,\infty)$ has a ray
but admits no line and hence cannot supply a line splitting.

\begin{lemma}[Geometry inherited by an exact factor; AC44]
\label{lem:alexandrov-factor-geometry}
Suppose $X$ is complete, proper and geodesic, with global nonnegative
Alexandrov comparison, and $\alpha$ is an onto pointed $k$-splitting.
Then $Y$ is complete, proper and geodesic and has the same comparison
lower bound. Every minimizing geodesic in $X$ between points in
$\alpha^{-1}(\{0\}\times Y)$ stays in that slice. Moreover,
$\dim_H Y\leq\dim_H X$, and $k\leq\dim_H X$.
\end{lemma}
\begin{proof}
The inverse image of $\{0\}\times Y$ is closed and isometric to $Y$.
Completeness and properness follow by taking closed subsets and
closed bounded balls. For endpoints $(0,a),(0,b)$, write
$D=d_Y(a,b)$ and let $(u,y)$ lie on a minimizing product geodesic.
The triangle equality gives
\[
 D=\sqrt{\|u\|_2^2+d_Y(a,y)^2}
       +\sqrt{\|u\|_2^2+d_Y(y,b)^2}
 \geq d_Y(a,y)+d_Y(y,b)\geq D.
\]
If $u\ne0$ the first inequality is strict, a contradiction. Thus the
whole geodesic stays in the zero slice, including the constant case
$D=0$. Its projection is a minimizing geodesic in $Y$. Distances and
triangles in the slice are unchanged, so comparison restricts to it.
Finally, $Y$ and $\mathbb R^k\times\{y_0\}$ embed isometrically in
$X$. Hausdorff-dimension monotonicity and the dimension of Euclidean
space give the two inequalities, as in BBI Proposition 1.7.19.
These inequalities alone do not yet subtract the Euclidean dimension.
\end{proof}

\subsection{Removing the Euclidean exponent from covering bounds}

\begin{lemma}[Quantitative dimension bound for an exact factor; AC45]
\label{lem:alexandrov-factor-dimension}
Suppose a pointed metric space $(X,p)$ has an onto $k$-splitting,
where $0\leq k\leq n$ are integers. Assume that for each $S>0$ there
is a finite $C_S>0$ such that, for all $0<\varepsilon\leq1$,
$\overline B_X(p,S)$ has an internal $\varepsilon$-net with at most
$C_S\varepsilon^{-n}$ points. Then, for every $R>0$ and
$0<\delta\leq1$, the ball $\overline B_Y(y_0,R)$ has an internal
$\delta$-net of size at most
\[
 3^n C_{R+\sqrt{k}+1}\,\delta^{-(n-k)}.
\]
In particular $\dim_H Y\leq n-k$.
For a complete proper geodesic nonnegatively curved Alexandrov $X$
with $\dim_H X\leq n$, $n\geq1$, the required net hypothesis follows
from AC07 and AC04. No general product-dimension equality is assumed.
\end{lemma}
\begin{proof}
Fix $R$, put $S=R+\sqrt{k}+1$, and first take
$0<\varepsilon\leq1/3$. In $[0,1]^k$ use the grid
\[
 G_\varepsilon=
 \{0,3\varepsilon,\ldots,
       3\varepsilon\lfloor(3\varepsilon)^{-1}\rfloor\}^k,
 \qquad
 |G_\varepsilon|\geq(3\varepsilon)^{-k}.
\]
For $k=0$ this is the singleton zero vector. Let $F$ be any finite
subset of $\overline B_Y(y_0,R)$ with pairwise distances strictly
greater than $3\varepsilon$. The set $G_\varepsilon\times F$ lies
in the closed $S$-ball of the product and has pairwise distances
strictly greater than $2\varepsilon$: distinct grid coordinates are
at least $3\varepsilon$ apart, and otherwise the factor separates
the points. Assign each such point to a center in an ambient
$\varepsilon$-net. Two assigned to the same center would be at
distance at most $2\varepsilon$, so the assignment is injective.
Surjectivity of $\alpha$ permits this argument in the source ball.
Consequently
\[
 |F|(3\varepsilon)^{-k}
 \leq |F|\,|G_\varepsilon|
 \leq C_S\varepsilon^{-n},
 \qquad
 |F|\leq3^k C_S\varepsilon^{-(n-k)}.
\]
The finite greedy packing argument of MC02 therefore gives an
internal $3\varepsilon$-net with that bound: keep adding a point
at distance $>3\varepsilon$ from all earlier points until the finite
cardinality bound forces termination. Set $\varepsilon=\delta/3$.
The resulting bound is $3^n C_S\delta^{-(n-k)}$. AC39 gives
Hausdorff dimension at most $n-k$ on each factor ball and then on
$Y$ by its countable integer-radius exhaustion.

For the Alexandrov specialization, AC44 first gives $k\leq n$.
Apply AC07/AC04 to the constant source space $X$ on any required
controlled ball, or equivalently to the constant sequence with
$c_i=1/i$ and $r_i=i$. Global curvature zero meets all growing-region
hypotheses. With the dimension-only constant $L(n)$, AC04 gives
\[
 C_S=\bigl(2+4L(n)^2\sqrt n\,\sinh(2S)\bigr)^n,
\]
since $1+\lceil u\rceil\leq2+u$ and $\varepsilon\leq1$.
A singleton source is harmless. This reuses the written local
structure and the retained comparison input; it does not assume
uniform chart radii or an exact dimension formula.
\end{proof}

\subsection{One-dimensional recognition and the pointed model list}

\begin{proposition}[One-dimensional local recognition input; AC46]
\label{prop:alexandrov-one-dimensional-recognition}
A nonempty complete proper geodesic nonnegatively curved Alexandrov
space $Y$ with $\dim_H Y\leq1$ is either a singleton or has the
following local structure: for every $y\in Y$ there is $r>0$ such
that $B_Y(y,r)$, with its restricted metric and basepoint $y$, is
isometric to either $((-r,r),0)$ or $([0,r),0)$.
\end{proposition}

The nontrivial local-recognition step remains source-qualified.
If $Y$ has two distinct points, a minimizing segment and BBI 1.7.19
give $\dim_H Y\geq1$, hence equality. AKP Corollary 15.16
(archived printed/PDF 236--237) translates Hausdorff dimension to
its linear-dimension convention. AKP Theorem 15.18 and its proof
(238) then give exactly these metric interval charts. The proof
excludes a third direction off an interior segment by Proposition
15.7; the endpoint case is separate. Thus the rank/dimension theorem
and local recognition are explicit source inputs. They are not
consequences of AC45 or of a bare Hausdorff upper bound for arbitrary
metric spaces. BGP Section 7.19, printed 30/PDF 31, uses the
one-dimensional manifold boundary as the initial case of its
Alexandrov-boundary definition; this is consistent with retaining
half-interval charts, rather than deleting endpoints.

\begin{proposition}[Complete pointed one-dimensional models; AC47]
\label{prop:alexandrov-one-dimensional-models}
Conditional on AC46's local-recognition input, every space in AC46
is isometric to exactly one of the following metric types, with the
indicated positive parameters:
\[
 \{*\},\qquad \mathbb R,\qquad [0,\infty),\qquad
 [0,\ell]\ (\ell>0),\qquad
 S^1_\ell=\mathbb R/\ell\mathbb Z\ (\ell>0).
\]
Intervals have their usual distance. The circle has its length metric
\[
 d_{S^1_\ell}([s],[t])=\min\{|s-t+j\ell|:j\in\mathbb Z\}.
\]
The pointed types retain $a\geq0$ on the ray and $a\in[0,\ell]$
on the segment. On the segment, reflection changes $a$ to $\ell-a$;
it does not move every basepoint to an endpoint. Translations allow
the line and circle basepoints to be represented by zero.
\end{proposition}
\begin{proof}
Treat the singleton separately. Start in the interior of an interval
chart and continue its unit-speed parametrization in both directions.
In an overlapping interval chart there is a unique continuation with
the chosen direction: an isometry between real intervals has slope
$1$ or $-1$, and matching the incoming direction fixes the choice.
At a finite end of the parameter interval the curve is Cauchy,
since its speed is one. Completeness supplies the endpoint. An
interior chart extends it further, so a finite maximal end must be
a boundary point with a half-interval chart and is included.

If the parametrization returns to a point, its outgoing direction
must agree with the earlier outgoing direction. The opposite choice
would, by local uniqueness continued back to the midpoint of the
return segment, make the parametrization reverse there, contrary to
its local unit-speed interval form. Hence a return produces a period;
repeat that segment to continue the parametrization on all of
$\mathbb R$. Its periods form a closed subgroup of $\mathbb R$; local injectivity
excludes arbitrarily small positive periods. There is therefore a
least positive period $\ell$, and the parametrization descends
injectively to the circle of length $\ell$. It has no finite
boundary end. In the absence of a return, the parametrization is
injective on an interval whose finite maximal ends are included.
Up to translation or reflection, this interval is $\mathbb R$,
$[0,\infty)$ or $[0,\ell]$ with $\ell>0$.

The image is all of $Y$: any point can be joined to the starting
point by a minimizing geodesic. In interval charts that geodesic
must follow one of the two continued directions and cannot leave
the parametrized curve or continue past a boundary endpoint.
Its finite length and the included finite endpoints ensure that
its last point is also in the image. Finally, paths in interval
charts measure ordinary arclength. On an interval the distance is
the difference of parameters; on the periodic curve it is the
shorter of the two arc lengths, giving the displayed quotient metric.
This proves an isometry, not merely a homeomorphism. The basepoint
statements follow from the isometries of these explicit models.
\end{proof}

\begin{corollary}[Exact products in dimension at most two; AC48]
\label{cor:alexandrov-low-dimensional-products}
Let $X$ be a nonempty complete proper geodesic nonnegatively curved
Alexandrov space with $\dim_H X\leq2$.
If an exact pointed $2$-splitting is supplied, its transverse factor
is a singleton and $X$ is the Euclidean plane.
If an exact pointed $1$-splitting is supplied, then, conditional on
AC46, $X$ is one of the Pythagorean products
\[
 \mathbb R,\quad \mathbb R^2,\quad
 \mathbb R\times[0,\infty),\quad
 \mathbb R\times[0,\ell],\quad
 \mathbb R\times S^1_\ell\qquad(\ell>0).
\]
For a specified line in $X$, AC43 supplies the required $1$-splitting
as a separate source-qualified input. The half-plane and strip
basepoints retain their transverse coordinate $a$ from AC47.
\end{corollary}
\begin{proof}
AC44 supplies the factor geometry. AC45 with $n=2$ gives factor
dimension at most $2-k$. If $k=2$, a nonsingleton geodesic factor
would contain a nondegenerate interval of Hausdorff dimension one,
a contradiction (also the singleton criterion in AC18). For $k=1$,
apply AC46 and AC47 and insert the five possibilities into the product.
Compose the actual factor isometry with the supplied splitting.
The line case uses AC43 first. All products are global metric
isometries; no smooth structure on $X$ is asserted.
\end{proof}

A maximal $1$-splitting, meaning that no $2$-splitting exists, excludes
the plane alternative. The mere existence of a $1$-splitting does not
exclude it. Compact or curved two-dimensional limits without a line
are outside this product list. Likewise, a basepoint at positive
height in a half-plane cannot be moved to its boundary by a pointed
isometry. KL's edge-point condition is an additional approximation
condition at the basepoint, not part of exact splitting.

\subsection{The remaining passage to collapse consumers}

KL Lemmas 4.15--4.16 (printed 31--32/PDF 26--27) connect long good
strainers, approximate splittings and exact splittings of pointed
limits. AC49--AC56 below expand these contracts. A metric limit from
AC41 alone does not contain a line or a prescribed Euclidean factor.
One must produce the limiting opposite rays, their line and
orthogonality relations, and the convergence of the selected
coordinate maps. Exact splitting of a limit does not itself prove
quantitative compatibility of approximating maps.

KL Lemma 8.1 uses the $2$-splitting branch after a separate
collapsed-limit dimension bound, while Lemma 9.1 uses the transverse
dimension bound for a $1$-splitting. The latter is
$\dim_H Y\leq1$, not a strict inequality. Passing these exact
conclusions back to approximate factors is still an obligation;
in particular, small bounded balls in a factor do not by themselves
bound its global diameter. Smooth submersions and circle/disk/torus
fibers require the later adapted-coordinate and proper-fibration
arguments. Those conclusions are not outputs of AC42--AC48.

\clearpage
\section{Approximate splittings and their pointed limits}
\label{sec:alexandrov-approximate-splitting-limits}

KL Definitions 3.2 and 4.7 (printed 21 and 30/PDF 16 and 25)
retain an actual approximation to a product. Lemma 4.15(2)
(printed 31--32/PDF 26--27) prescribes its Euclidean coordinates;
Lemma 4.16 (printed 32/PDF 27) concerns a previously specified
pointed limit. We separate the metric passage from the geometric
long-strainer input. In particular, AC50--AC51 below do not use
curvature or the splitting theorem.

\subsection{Actual product maps and compactness of the factors}

\begin{definition}[Normalized approximate splitting; AC49]
\label{def:alexandrov-approximate-splitting}
For an integer $k\geq0$ and $0<\delta<1$, a normalized
$(k,\delta)$-splitting of $(X,p)$ consists of a pointed metric
space $(Y,y)$ and a whole-space pointed map
\[
 \phi=(u,v):(X,p)\longrightarrow
       (\mathbb R^k\times Y,(0,y)).
\]
Write $P=\mathbb R^k\times Y$ and let $\pi_E$ be its Euclidean
projection. The map is a KL $\delta$-approximation: distortion on
$B_X(p,\delta^{-1})$ is at most $\delta$, and every point of
$B_P((0,y),\delta^{-1}-\delta)$ has distance
at most $\delta$ from the image of that source ball. Products use
the Pythagorean metric. This is an infimum bound, not an assertion
that the infimum is attained. The factor $Y$ need not be complete,
proper or a length space. The map need not be continuous.
\end{definition}

Translation in KL's Euclidean factor gives this normalization without
changing the transverse basepoint. The strict closed-ball convention
MC06 is a different predicate: MC11 gives
$\mathcal A_{R,3\delta}(\phi)$ whenever
$3\delta<R<\delta^{-1}$, using the same map. Thus a sequence of
approximate splittings with $\delta_i\to0$ supplies actual product
maps on every fixed radius with errors tending to zero. This conversion
preserves $u_i$, and uses slack for open balls and unattained infima.
No fixed-error symmetry or exact surjectivity is inferred.

\begin{lemma}[Eventual nets in the approximate factors; AC50]
\label{lem:alexandrov-approximate-factor-nets}
Suppose $(X_i,p_i)\to(X,p)$, where $X$ is complete and proper,
and each $X_i$ has a supplied normalized $(k,\delta_i)$-splitting
$\phi_i:X_i\to\mathbb R^k\times Y_i$, with fixed $k$ and
$\delta_i\to0$. Then the $Y_i$ satisfy the eventual internal-net
condition of MC13. A subsequence converges to a complete proper
pointed space $(Y,y)$. No properness of the individual factors
or source spaces is required.
\end{lemma}
\begin{proof}
Fix $R,\eta>0$, and choose
$0<e<\min\{1/10,\eta/6\}$. On a sufficiently late tail, there
are maps $f_i:X_i\to X$ and the given product maps $\phi_i$,
restricted as necessary, both satisfying $\mathcal A_{R+4,e}$.
Choose a finite internal $e$-net in the compact ball
$\overline B_X(p,R+1)$; let its size be $N$.

For each $z\in\overline B_{Y_i}(y_i,R)$, product coverage selects
$x_z\in X_i$ with
\[
 d(\phi_i(x_z),(0,z))<e,\qquad d(p_i,x_z)<R+2e.
\]
Thus $f_i(x_z)$ lies in $\overline B_X(p,R+1)$. Assign $z$ to
one of its net centers within $e$. If $z,z'$ have the same assigned
center, then
\[
 \begin{split}
 d_X(f_i(x_z),f_i(x_{z'}))&\leq2e,\\
 d_{X_i}(x_z,x_{z'})&<3e,\\
 d(\phi_i(x_z),\phi_i(x_{z'}))&<4e,\\
 d_{Y_i}(z,z')&<6e<\eta.
 \end{split}
\]
All evaluated source points are inside the common $R+4$ domain.
Choose one factor point from every nonempty assignment class.
These at most $N$ representatives lie inside the factor $R$-ball
and form an $\eta$-net there. The bound and the tail depend on
$R,\eta$, not on a point or on the member of the tail. Apply the
\emph{eventual} version of MC13, which does not assume proper sources,
to all radii and accuracies. One subsequence and one completed factor
limit suffice; independent limits of separate balls are not used.
\end{proof}

\subsection{A product on the specified limit, with coordinate control}

\begin{theorem}[Limit of actual approximate product maps; AC51]
\label{thm:alexandrov-approximate-product-limit}
Under AC50's hypotheses, after passing to a subsequence there is
an onto pointed isometry
\[
 \alpha:(X,p)\longrightarrow(\mathbb R^k\times Y,(0,y)).
\]
The factor and subsequence can be chosen so that, for suitable
pointed approximations $f_i:X_i\to X$ on growing balls and every
fixed $R>0$,
\[
 \sup_{x\in\overline B_{X_i}(p_i,R)}
 \|u_i(x)-\pi_E\alpha(f_i(x))\|_2
          \longrightarrow0.
\]
The assertion is along the selected subsequence. It makes no
uniqueness claim for $\alpha$ or the limiting coordinate functions.
\end{theorem}
\begin{proof}
AC50 supplies $Y_i\to Y$ along a subsequence. First verify that
products converge while keeping their Euclidean coordinates.
Put $P_i=\mathbb R^k\times Y_i$. If $g_i:\overline B_{Y_i}(y_i,S)\to Y$ satisfies
$\mathcal A_{S,e}$ and $3e<S$, set
\[
 H_i(u,z)=(u,g_i(z))
 \quad\hbox{on }\overline B_{P_i}((0,y_i),S).
\]
The map $(a,b)\mapsto\sqrt{a^2+b^2}$ is $1$-Lipschitz in either
nonnegative argument, so $H_i$ has distortion less than $e$.
For $(v,w)$ of product radius at most $S-3e$, coverage for $g_i$
chooses $z$ with $d(g_i(z),w)<e$ and
$d(y_i,z)<d(y,w)+2e$. Consequently
\[
 \sqrt{\|v\|_2^2+d(y_i,z)^2}
 <\sqrt{\|v\|_2^2+d(y,w)^2}+2e<S.
\]
The witness $(v,z)$ lies in the actual product domain. Hence
$\mathcal A_{S,3e}(H_i)$. This argument does not assume that
closed balls in $Y_i$ are length spaces.

Compose $H_i$ with $\phi_i$ on padded domains using MC08. A diagonal
choice of radii and errors gives maps $q_i:X_i\to P$, where
$P=\mathbb R^k\times Y$, on growing balls, whose approximation
errors tend to zero and whose first coordinate is exactly $u_i$.
The product $P$ is complete and proper: a closed product ball is
a closed subset of the product of two compact coordinate balls.
Both $X$ and $P$ are therefore complete proper limits of the same
subsequence. MC24 gives a pointed onto isometry. We retain its
construction to obtain the asserted control of the actual maps.

Choose a common subsequence and approximation domains with radii
much larger than $j$, and errors tending to zero. MC09 supplies
local quasi-inverses $t_i$ to $f_i$; MC08 forms
$F_i=q_i\circ t_i$ on the closed $j$-ball of $X$. The domain margins
can be chosen exactly as in MC24, enlarging the original radii
before composition. Its dense-set diagonal proof gives
$F_i\to\alpha$ on a countable dense subset. It also gives uniform
convergence on each fixed ball: for a finite $a$-net in a compact
ball and its nearest center $b$,
\[
 d_P(F_i(x),\alpha(x))
 \leq 2a+\operatorname{dis}(F_i)
       +\max_b d_P(F_i(b),\alpha(b)).
\]
First take $i\to\infty$, then $a\downarrow0$. Here distortion
means the supremum bound on the chosen domain. The coverage
preimages used in MC24 remain in a fixed compact ball; their
convergent subsequences prove onto-ness, not just an embedding.

For $x\in\overline B_{X_i}(p_i,R)$, the quasi-inverse estimate
gives $d_{X_i}(t_i(f_i(x)),x)\to0$, uniformly on that ball.
Its image under $f_i$ lies in the fixed compact $(R+1)$-ball for
large $i$. Therefore
\[
 \begin{split}
 d_P(q_i(x),\alpha(f_i(x)))
 &\leq d_{X_i}(x,t_i(f_i(x)))+\operatorname{dis}(q_i)\\
 &\quad+d_P(F_i(f_i(x)),\alpha(f_i(x)))\longrightarrow0.
 \end{split}
\]
Projection to the Euclidean factor proves the stated estimate.
All maps are evaluated only on their controlled domains; no continuity
of an approximation map has been assumed.
\end{proof}

Even for $X_i=X=\mathbb R$, alternating the maps $x\mapsto x$ and
$x\mapsto-x$ prevents convergence of the entire sequence of coordinates.
A subsequence or a separately specified alignment is necessary.
AC51 neither proves KL 4.8 compatibility between two supplied splittings
nor bounds the global diameter of an approximate factor.

\begin{lemma}[Restoring prescribed Euclidean coordinates; AC52]
\label{lem:alexandrov-prescribed-coordinate-transfer}
Let $\alpha:(X,p)\to(\mathbb R^k\times Y,(0,y))$ be an exact
splitting, let $\mathcal A_{S,e}(f:(Z,o)\to(X,p))$, and let
$h:Z\to\mathbb R^k$ satisfy $h(o)=0$. Suppose on
$\overline B_Z(o,S)$ that
\[
 \|h-\pi_E\alpha f\|_2\leq\rho,
 \qquad E=e+2\rho<S.
\]
Then on that ball the map
\[
 \psi=(h,\pi_Y\alpha f)
\]
satisfies $\mathcal A_{S,E}$. In particular, if $(Z_i,o_i)\to(X,p)$
and the displayed coordinate discrepancy tends to zero on every
fixed ball, then for every fixed $0<\delta<1$, all sufficiently
large $i$ have a normalized $(k,\delta)$-splitting with Euclidean
coordinate exactly $h_i$.
\end{lemma}
\begin{proof}
The distance from $\psi(x)$ to $\alpha f(x)$ is at most $\rho$.
Two such changes and the distortion of $f$ give distortion less
than $e+2\rho=E$. A target of radius at most $S-E$ is in the
$S-e$ coverage region for $\alpha f$. The original witness is
within $e+\rho\leq E$ after the change, with strict inequality
because the original error is strict. The basepoint is exact.

For the whole-space KL statement, choose
$S=\delta^{-1}+\delta$ and take $e,\rho<\delta/8$ on this
fixed domain. Then $E<3\delta/8$. Extend the map outside the
$S$-ball by $(h_i,y)$. A target of radius less than
$\delta^{-1}-\delta$ has a coverage witness $x$ satisfying
\[
 d(o_i,x)<d((0,y),\psi_i(x))+E
 <\delta^{-1}-\delta+2E<\delta^{-1}.
\]
Thus the witness lies in the required open source ball, and
both KL bounds hold with error at most $\delta$. The extension
preserves $h_i$ everywhere, including outside the controlled region.
No smoothness or submersion property follows from this construction.
\end{proof}

\subsection{The geometric input behind the long-strainer implication}

\begin{lemma}[Oriented orthogonal lines give the coordinate splitting; AC53]
\label{lem:alexandrov-orthogonal-line-coordinates}
Let $(X,p)$ be a complete proper finite-dimensional geodesic
nonnegatively curved Alexandrov space. Suppose unit-speed lines
$\gamma_1,\ldots,\gamma_k:\mathbb R\to X$ pass through $p$ at
parameter zero, and their positive rays have pairwise Alexandrov
angle $\pi/2$ at $p$. Conditional on AC43, there is an exact
pointed splitting $\alpha:X\to\mathbb R^k\times Y$ with
\[
 \alpha(\gamma_j(t))=(t e_j,y),\qquad
 \pi_j\alpha(x)=-b_j(x),\qquad
 b_j(x)=\lim_{t\to\infty}\bigl(d(x,\gamma_j(t))-t\bigr).
\]
The sign is fixed by the chosen positive ray. For $k=0$ use the
identity splitting with factor $X$.
\end{lemma}
\begin{proof}
For any unit-speed ray, the expression defining $b_j$ is
nonincreasing in $t$ by the triangle inequality and is bounded
below by $-d(x,p)$, so its finite limit exists.
Apply AC43 to $\gamma_1$, aligned so it becomes $(t,y)$ in
$\mathbb R\times Y$. Each other line has a linear real coordinate
$at$ and a constant-speed factor geodesic, of speed $b$, where
$a^2+b^2=1$. To see this product-geodesic fact, equality in the
product triangle inequality forces equality in both coordinate
triangle inequalities and proportional coordinate-distance pairs,
by equality in the Euclidean Cauchy--Schwarz inequality. Applying
this to every three ordered parameters gives constant speeds;
the real coordinate is affine, and its value at zero is zero.
This is KL Sublemma 4.5(1), with its equality argument exposed.

The squared distance from $(s,y)$ to that line at parameter $t>0$
is $s^2+t^2-2ast$. Hence the angle with the positive first axis
is $\arccos(a)$, so orthogonality forces $a=0$. All remaining
lines lie in the zero slice. Their mutual angles and distances
are unchanged there, and AC44 supplies the factor geometry.
Repeat on that factor. This proves the asserted alignment of all
$k$ oriented axes without an unspecified final rotation.

In the resulting product, write $\alpha(x)=(u,z)$. Then
\[
 d(x,\gamma_j(t))^2=(t-u_j)^2+
          \sum_{h\ne j}u_h^2+d_Y(z,y)^2.
\]
Subtract $t$ and rationalize, or divide numerator and denominator
by $t$, to obtain $b_j(x)=-u_j$. Thus the prescribed negative
Busemann functions are exactly the Euclidean coordinates.
\end{proof}

KL's proof of 4.15(2) compresses the following geometric passage into
one paragraph. AC54 records its precise contract; AC57--AC66 below
expand its proof while retaining the stated geometric source inputs.

\begin{proposition}[Long-strainer ray and coordinate input; AC54]
\label{prop:alexandrov-long-strainer-input}
Let $1\leq k\leq n$, $\sigma_i>0$ with $\sigma_i\to0$, and
let $(X_i,p_i)$ be complete pointed length spaces with curvature
at least $-\sigma_i$ and Hausdorff dimension at most $n$ on
$B(p_i,\sigma_i^{-1})$, in the local convention of AC41.
Suppose $p_i$ has a KL $k$-strainer
$\{a_{i,j}^+,a_{i,j}^-\}_{j=1}^k$ of quality $\sigma_i$ and
scale $\sigma_i^{-1}$. In particular, both endpoint distances
are exactly $\sigma_i^{-1}$; the comparison angles use curvature
$-\sigma_i$, with opposite angles greater than $\pi-\sigma_i$
and cross-pair angles for distinct indices greater than
$\pi/2-\sigma_i$.
Let $X$ be a complete proper nonnegatively curved limit supplied
by AC41. After a further subsequence there are opposite rays
$\gamma_j^\pm$ from $p$, forming unit-speed lines $\gamma_j$
whose positive rays have pairwise angle $\pi/2$, and approximations
$f_i$ realizing that same limit such that, on each fixed source ball,
\[
 \sup_{x\in\overline B_{X_i}(p_i,R)}
 \left|h_{i,j}(x)+b_j(f_i(x))\right|\longrightarrow0,
 \quad
 h_{i,j}(x)=d(p_i,a_{i,j}^+)-d(x,a_{i,j}^+),
\]
where $b_j$ is the Busemann function of the positive limiting ray.
\emph{AC66 below supplies the written proof, conditional on the retained comparison and splitting inputs.}
\end{proposition}

The source proof uses the sequence $\sigma_i=1/i$. The general positive
sequence formulation now has written line and coordinate arguments in
AC57--AC61 below. They construct controlled almost-shortest prefixes,
extract globally minimizing lines and identify negative Busemann limits
with an explicit two-sided estimate. AC62--AC66 subsequently supply cross-pair orthogonality using
buffered comparison. Finite model-angle continuity AC35 alone does not
handle endpoints escaping to infinity.

Each $h_{i,j}$ is $1$-Lipschitz and vanishes at $p_i$, and $b_j$
is $1$-Lipschitz. Once convergence at lifted points of a dense set
has been proved geometrically, finite nets and the distortion bound
upgrade it to the displayed uniform-on-balls statement. Concretely,
choose a finite $a$-net in the compact $(R+1)$-ball of $X$ and lift
its centers by $f_i$; comparison with a nearest center gives an error
at most $2a+o(1)$ plus the finite maximum of the center errors.
This elementary upgrade does not prove the required pointwise
Busemann identification. AC60 supplies the previously missing pointwise identification and a
direct uniform estimate; AC65 supplies orthogonality. Local comparison
and global line splitting retain their separate source qualifications.

The scale is $\sigma_i^{-1}$ and the curvature is $-\sigma_i$;
it is not a claim at curvature $-\sigma_i^2$. All complete-buffer
and intrinsic/ambient qualifications in AC41 remain in force.
No minimum over far endpoints, tangent-cone identification, global
factor-diameter bound, or differentiability of $h_{i,j}$ is asserted.

\begin{theorem}[Long strainers produce the prescribed coordinates; AC55]
\label{thm:alexandrov-long-strainer-splitting}
Conditional on AC54 and AC43, for every integer $n\geq1$ and
$0<\delta<1$ there is $0<\sigma=\sigma(n,\delta)<1$ with the
following property, uniformly for $1\leq k\leq n$.
If a complete pointed length space $(Z,o)$ has curvature at least
$-\sigma$ and Hausdorff dimension at most $n$ on
$B(o,\sigma^{-1})$, and a KL $k$-strainer
$\{a_j^+,a_j^-\}$ of quality $\sigma$ at scale $\sigma^{-1}$,
then it has a normalized $(k,\delta)$-splitting $\psi$ whose
$j$th Euclidean coordinate is exactly
\[
 h_j(z)=d(o,a_j^+)-d(z,a_j^+).
\]
The local geometric convention is that of AC41. This is the
source-qualified geometric content of KL 4.15(2), with the final
metric contradiction written below. No explicit modulus or
monotonicity in the parameter $\sigma$ is claimed.
\end{theorem}
\begin{proof}[Conditional compactness reduction]
If no such $\sigma$ exists, choose counterexamples at
$\sigma_i=1/i$ for integers $i\geq2$. Their strainer ranks belong
to the finite set $\{1,\ldots,n\}$, so restrict to an infinite
subsequence with a fixed $k$. The original indices still tend
to infinity. AC41 gives one complete proper nonnegative limit
of dimension at most $n$. AC54 supplies its orthogonal lines
and the locally uniform limits $h_{i,j}\to-b_j$.
AC53 then gives an onto exact splitting whose Euclidean
coordinates are exactly $(-b_1,\ldots,-b_k)$.
The scalar uniform estimates imply the vector estimate, with
at most a factor $\sqrt{k}$. AC52 produces, for all sufficiently
late members of this subsequence, a $(k,\delta)$-splitting with
Euclidean coordinate exactly $h_i$. This contradicts the chosen
counterexamples. This proof closes only the reduction from the
stated geometric inputs, not AC54 itself.
\end{proof}

The reverse implication KL 4.15(1) assumes a complete
\emph{globally nonnegatively curved} Alexandrov space and produces
a strainer at the prescribed scale $1/\delta$. It is not silently
applied to the locally almost-nonnegative source spaces above.
Its exact-scale endpoint adjustment is written in AC67--AC70
below; the following limit theorem does not need it.

\subsection{The approximate-splitting closure used by collapse}

\begin{corollary}[Splitting of the given Alexandrov limit; AC56]
\label{cor:alexandrov-given-limit-splitting}
Let $n\geq1$, $0\leq k\leq n$ be integers. Let $(X_i,p_i)$
be complete pointed length spaces with positive $\delta_i\to0$,
curvature at least $-\delta_i$ and Hausdorff dimension at most
$n$ on $B(p_i,\delta_i^{-1})$, in the convention of AC41.
Suppose they have supplied normalized $(k,\delta_i)$-splittings
on a tail and converge to a specified complete proper pointed
space $(X,p)$. Then this same $X$ is geodesic and nonnegatively
curved, has $\dim_H X\leq n$, and admits an exact pointed
$k$-splitting. Its factor is complete proper geodesic and
nonnegatively curved, with Hausdorff dimension at most $n-k$.
The original Euclidean coordinates converge to those of a
chosen exact splitting along a subsequence as in AC51.
\end{corollary}
\begin{proof}
AC41 applied to any suitable subsequence supplies a complete proper
geodesic nonnegative limit with the dimension bound. MC24 identifies
it by a pointed isometry with the already specified $X$, since the
same source subsequence converges to both. These intrinsic properties
therefore hold on $X$. AC51 gives the exact product on $X$ and the
coordinate convergence. AC44 and AC45 give the factor properties
and dimension bound. For $k=0$ the identity splitting also suffices.
This route uses no long-strainer producer: AC54 and AC55 are not
dependencies. Its geometric qualification is precisely the local
comparison input already retained by AC41.
\end{proof}

This writes a metric route to KL 4.16's conclusion; KL itself says
it follows from 4.15. The specified limit is complete and proper
in our selected contract, consistent with the proper-space pointed
category recalled in KL Section 3.1. No new classification of
arbitrary two-dimensional spaces follows. KL 6.18 must still supply
the collapse-specific bound $\dim_H X\leq2$ before AC48 can be
used for its intended low-dimensional applications. AC71--AC80 below
write metric compatibility and its basepoint transport in this chapter.
Smooth adapted coordinates and compatibility of actual fibrations retain
their separate Chapter 14 obligations.

\section{Opposite endpoints, limiting lines and distance coordinates}
\label{sec:alexandrov-opposite-endpoint-limits}

The next five contracts expand the line and coordinate parts of AC54.
They avoid requiring minimizing segments all the way to the distant
strainer endpoints, which lie on the boundary of the controlled open
ball. Almost-shortest paths suffice. The only new comparison-model
calculation is an estimate for the opposite-endpoint distance defect;
it is a calculation with the three given distances, not an application
of Toponogov to a long source triangle. Orthogonality of different limiting lines is proved separately in
AC62--AC66; it does not follow from the opposite-pair data alone.

\subsection{An absolute defect estimate and controlled prefixes}

\begin{lemma}[Opposite comparison angles give vanishing excess; AC57]
\label{lem:alexandrov-opposite-endpoint-excess}
Let $0<\sigma\leq1$ and let $p,a^+,a^-$ be points in a metric
space with $d(p,a^+)=d(p,a^-)=L=\sigma^{-1}$. If the comparison
angle at $p$ in the curvature $-\sigma$ model is at least
$\pi-\sigma$, then, writing $c=d(a^+,a^-)$,
\[
 0\leq E:=2L-c\leq
 \omega(\sigma):=-\frac{2}{\sqrt\sigma}\log\cos(\sigma/2)
 \leq\frac{\sigma^{3/2}}2.
\]
No curvature or geodesic assumption on the source space is needed
beyond the stated comparison-angle condition on its distances.
\end{lemma}
\begin{proof}
Set $\lambda=\sqrt\sigma$, $A=\lambda L$, and
$q=\cos(\sigma/2)>0$. For the comparison angle $\theta$, the
hyperbolic cosine law and the half-angle identity give
\[
 \sinh(\lambda c/2)=\sinh A\sin(\theta/2)\geq q\sinh A.
\]
The identity
\[
 \sinh(A+\log q)
 =\tfrac12(qe^A-q^{-1}e^{-A})
 \leq q\sinh A
\]
and strict monotonicity of $\sinh$ imply
$\lambda c/2\geq A+\log q$. Rearranging proves the middle
inequality, while $E\geq0$ is the triangle inequality.
For $0\leq t\leq1/2$, $\cos t\geq1-t^2/2\geq7/8$, so
$(\tan t)'=1/\cos^2t\leq64/49<2$. Thus $\tan t\leq2t$, and
\[
 -\log\cos(\sigma/2)
 =\int_0^{\sigma/2}\tan t\,dt\leq\sigma^2/4.
\]
This proves the final bound. The degenerate angle $\pi$ gives
$c=2L$ and causes no exception. Notice the absolute error
$E\to0$, not merely the relative statement $E/L\to0$.
\end{proof}

\begin{lemma}[Almost-shortest prefixes and their calibration; AC58]
\label{lem:alexandrov-short-prefix-calibration}
Let $(Z,o)$ be a length space with finite-valued metric,
$d(o,a^+)=d(o,a^-)=L>0$, and $E=2L-d(a^+,a^-)\geq0$.
For every $\eta>0$ choose paths from $o$ to $a^\pm$ of lengths
$\ell_\pm<L+\eta$, parametrized by arc length, and let
$q^\pm(t)$ be their points at time $0\leq t\leq L$.
These paths exist without completeness or local compactness.
For $0\leq s\leq t\leq L$ and $0\leq u\leq L$,
\[
 \begin{aligned}
 t-\eta&\leq d(o,q^\pm(t))\leq t,\\
 t-s-\eta&\leq d(q^\pm(s),q^\pm(t))\leq t-s,\\
 t+u-E-2\eta&\leq d(q^+(t),q^-(u))\leq t+u.
 \end{aligned}
\]
For $h(z)=L-d(z,a^+)$ one also has
\[
 t-\eta\leq h(q^+(t))\leq t,\qquad
 -t\leq h(q^-(t))\leq -t+E+\eta.
\]
The whole prefix up to time $T\leq L$ is contained in
$\overline B(o,T)$; the full path need not remain in $B(o,L)$.
\end{lemma}
\begin{proof}
The length-space characterization supplies a rectifiable path of
length less than $L+\eta$ for each endpoint. Arc-length
reparametrization supplies $q^\pm$ on $[0,\ell_\pm]$;
$\ell_\pm\geq L$ ensures the stated prefixes exist. Their
subpaths have length equal to the parameter difference, and
\[
 d(q^\pm(t),a^\pm)\leq\ell_\pm-t<L+\eta-t.
\]
The radial and same-branch lower bounds follow by comparing a
shortcut through the indicated points with $d(o,a^\pm)=L$.
For different branches, the triangle inequality gives
\[
 2L-E\leq (L+\eta-t)
       +d(q^+(t),q^-(u))+(L+\eta-u).
\]
All upper bounds follow from the two prefix paths through $o$.
For calibration on the positive branch, use the displayed tail
estimate and $L\leq d(o,q^+(t))+d(q^+(t),a^+)$. On the
negative branch, $d(q^-(t),a^+)\leq t+L$, whereas
\[
 d(q^-(t),a^+)\geq d(a^+,a^-)-d(q^-(t),a^-)
 \geq L+t-E-\eta.
\]
Subtracting from $L$ proves both calibration inequalities.
The containment follows from the radial upper bound at every
prefix time. No existence of a globally minimizing source segment
has been used.
\end{proof}

\subsection{One diagonal subsequence and an explicit two-sided squeeze}

\begin{proposition}[A limiting line from almost opposite endpoints; AC59]
\label{prop:alexandrov-endpoint-limit-line}
Let pointed length spaces $(X_i,p_i)$ with finite-valued metrics
converge to a complete proper pointed metric space $(X,p)$.
Suppose $d(p_i,a_i^+)=d(p_i,a_i^-)=L_i\to\infty$ and
$E_i=2L_i-d(a_i^+,a_i^-)\to0$. Choose $\eta_i>0$ tending to
zero and the prefixes of AC58. Choose actual approximations
$f_i$ to this same $X$ with domains tending to infinity and
errors tending to zero, as in MC06. After a subsequence there is
an isometric line $\gamma:\mathbb R\to X$, $\gamma(0)=p$, with
\[
 f_i(q_i^+(t))\longrightarrow\gamma(t),\qquad
 f_i(q_i^-(t))\longrightarrow\gamma(-t)
 \quad\text{for every fixed }t\geq0.
\]
Convergence is uniform for $t$ in each bounded parameter interval.
The same subsequence can be chosen for any fixed finite number
of endpoint pairs satisfying these hypotheses. Individual $X_i$
need not be proper, complete, geodesic or Alexandrov.
\end{proposition}
\begin{proof}
For each nonnegative rational $t$, the prefix is eventually defined
and its image lies in the compact ball $\overline B_X(p,t+1)$.
A countable diagonal selection makes all these images converge,
for both signs and all finitely many pairs. Applying the distortion
bound to AC58 shows that the limiting distances on one branch
are $|t-s|$, and between branches are $t+s$. At zero both
images equal $p$. Thus the signed rational parameters give an
isometric copy of the rational line in $X$.

Completeness extends this map uniquely to an isometric map
$\gamma:\mathbb R\to X$: for a real parameter use any rational
sequence converging to it, whose images are Cauchy, and pass the
exact distance identity to the limit. In particular, every segment
of $\gamma$ is globally minimizing. This is stronger than merely
having an angle $\pi$ between two short rays.

For an arbitrary real $t$ in a fixed interval $[0,T]$, select a
rational $s$ in that interval close to $t$. The prefix maps are
$1$-Lipschitz before applying $f_i$, so, with approximation error
$e_i$ on a ball containing the prefixes,
\[
 d_X(f_i(q_i^\pm(t)),\gamma(\pm t))
 \leq 2|t-s|+e_i+
       d_X(f_i(q_i^\pm(s)),\gamma(\pm s)).
\]
A finite rational net and convergence at its finitely many points
prove uniform convergence on $[0,T]$. No continuity of $f_i$ is
required. Taking all rational times at once ensures one whole
line in the specified limit, not independent limits of balls.
\end{proof}

Absolute excess cannot be replaced here by relative excess. For example,
join an interval of length $a>0$ at one end to two infinite rays, and
let $p$ be its other end. Points at distance $L_i\to\infty$ from
$p$ on the two rays have $E_i=2a$ and $E_i/L_i\to0$. The constant
pointed tree is a proper limit, but has no line through its terminal
point $p$; the two prefixes share the initial interval.

\begin{lemma}[Distance coordinates converge along a split line; AC60]
\label{lem:alexandrov-two-sided-busemann-squeeze}
Use the data, prefixes, subsequence and line of AC59. Suppose an
onto pointed isometry $\alpha:X\to\mathbb R\times Y$ is supplied
with $\alpha(\gamma(t))=(t,y)$, using the Pythagorean metric.
Write $\alpha(x)=(u(x),z(x))$, and set
$h_i(x)=L_i-d(x,a_i^+)$. Then for every fixed $R>0$,
\[
 \sup_{x\in\overline B_{X_i}(p_i,R)}
 |h_i(x)-u(f_i(x))|\longrightarrow0.
\]
Moreover $u=-b_+$, where $b_+(x)=\lim_{T\to\infty}
(d(x,\gamma(T))-T)$ is the Busemann function of the positive ray.
No orthogonality between different lines is needed for this
single-coordinate assertion.
\end{lemma}
\begin{proof}
Put $B=R+1$ and fix $T>B$. For all sufficiently large $i$,
$L_i>T$, and let $e_i$ bound the distortion of $f_i$ on a ball
containing both the source $R$-ball and the two time-$T$ prefixes.
Also require its basepoint error to be zero and $e_i<1$, so
$d_X(p,f_i(x))\leq B$ on that source ball. Let
\[
 \zeta_i(T)=\max_{\epsilon\in\{+,-\}}
 d_X(f_i(q_i^\epsilon(T)),\gamma(\epsilon T)),\qquad
 D_i(T)=E_i+\eta_i+e_i+\zeta_i(T).
\]
Here $\epsilon T$ means $T$ or $-T$, respectively, and
$D_i(T)\to0$ for this fixed $T$. The scalar functions $h_i$
are $1$-Lipschitz. Calibration on the positive prefix gives a
lower bound, and calibration on the negative prefix gives an
upper bound. Thus, writing $w=f_i(x)$,
\[
 T-d_X(w,\gamma(T))-D_i(T)
 \leq h_i(x)\leq
 d_X(w,\gamma(-T))-T+D_i(T).
\]
In the product, put $u=u(w)$ and $r=d_Y(z(w),y)$, so
$u^2+r^2\leq B^2$. The Pythagorean formula and rationalization
show, for either sign,
\[
 0\leq\sqrt{(T\pm u)^2+r^2}-(T\pm u)
 =\frac{r^2}{\sqrt{(T\pm u)^2+r^2}+T\pm u}
 \leq\frac{B^2}{2(T-B)}.
\]
Combining the last two displays gives the explicit uniform estimate
\[
 \sup_{x\in\overline B_{X_i}(p_i,R)}
 |h_i(x)-u(f_i(x))|
 \leq D_i(T)+\frac{B^2}{2(T-B)}.
\]
Given an accuracy, first fix $T$ large enough for the second term,
then take $i$ large enough for $D_i(T)$. There is no exchange of
uncontrolled limits. For any fixed $w\in X$, the same product
formula yields $d_X(w,\gamma(T))-T\to-u(w)$, proving the
Busemann identity and its sign. The positive endpoint therefore
produces the negative of the positive-ray Busemann function.
\end{proof}

\subsection{Separating cross-pair orthogonality from the individual lines}

\begin{corollary}[Long-strainer lines and their individual coordinates; AC61]
\label{cor:alexandrov-long-strainer-individual-lines}
In the setting of AC54, with the complete proper nonnegative limit
from AC41, the following part of its conclusion has a written proof,
conditional on AC43. After one further subsequence there are
unit-speed lines $\gamma_1,\ldots,\gamma_k$ through $p$, with
positive and negative rays obtained from the respective endpoint
prefixes, and, on every fixed source ball,
\[
 \sup_x|h_{i,j}(x)+b_j(f_i(x))|\longrightarrow0
 \qquad (1\leq j\leq k).
\]
This conclusion does not yet assert that distinct lines are
orthogonal, distinct as unoriented sets, or axes of one common
$k$-splitting.
\end{corollary}
\begin{proof}
Discard finitely many indices so that $0<\sigma_i\leq1$.
AC57 gives $0\leq E_{i,j}\leq\sigma_i^{3/2}/2\to0$ for every
opposite pair, while $L_i=\sigma_i^{-1}\to\infty$. Choose,
for example, $\eta_i=\sigma_i$ in AC58. For any fixed $T$, the
resulting prefixes eventually lie in
$\overline B(p_i,T)\subset B(p_i,\sigma_i^{-1})$.
Only their existence and metric estimates are used here; no
comparison is applied to the full paths or the distant endpoints.
AC59 gives all $k$ lines and the same actual approximations to
$X$ on one common subsequence.

AC41 supplies the complete proper finite-dimensional geodesic
nonnegative geometry of $X$. Apply source-qualified AC43 to each
line separately, aligned with its positive direction. Each such
one-dimensional splitting satisfies AC60, so the asserted scalar
convergence follows for all $j$. These splittings may be different;
no simultaneous alignment is inferred. This proves the line and
coordinate portions of AC54 for arbitrary positive
$\sigma_i\to0$. It uses no cross-pair comparison-angle inequality.
\end{proof}

The additional requirement in AC54 is precise: use the cross-pair strainer
inequalities and the local lower-curvature hypotheses to show
$\angle(\gamma_j^+,\gamma_h^+)=\pi/2$ for $j\ne h$, for a
suitable common prefix extraction. AC57--AC61 alone cannot do
this: in $\mathbb R^2$, endpoint pairs at $\pm L_i v_j$ have
zero opposite excess for any chosen unit vectors $v_j$, including
coincident or nonorthogonal ones. The cross-pair hypotheses are
therefore essential to the full producer. Finite-scale comparison
must justify its source domains; AC35 cannot be applied directly
to the endpoints escaping to infinity. For $k=1$ the orthogonality condition is vacuous, so AC61 supplies
the full AC54 conclusion conditional on AC41 and AC43. For higher rank, AC65 below establishes orthogonality; AC53 and AC52
supply the already written simultaneous
coordinate splitting and the closing step of AC55.

\section{Buffered shortening and orthogonality of the limiting lines}
\label{sec:alexandrov-buffered-orthogonality}

The far endpoints in AC54 lie at the boundary of the controlled open
ball. Applying AC02 to that whole configuration would not satisfy its
radius requirement. Instead we first shorten every radial segment to
one hundredth of its length. A numerical estimate in the negative
curvature model preserves the limiting cross-pair angles under this
first shortening; it uses only the source triangle inequality.
The shortened configuration then fits AC02, which handles a second
shortening to each fixed radius. This supplies the missing
orthogonality argument without a new global comparison assumption.

\subsection{Radial endpoint geodesics and a numerical shortening estimate}

\begin{lemma}[Geodesics to the boundary of a locally compact ball; AC62]
\label{lem:alexandrov-boundary-radial-geodesic}
Let $X$ be a complete length space, let $L>0$, and suppose
$B_X(p,L)$ is locally compact. Every point $a\in X$ with
$d(p,a)=L$ is joined to $p$ by an isometric segment
$\gamma:[0,L]\to X$, with $\gamma(0)=p$ and $\gamma(L)=a$.
No local compactness at $a$ or outside $B_X(p,L)$ is assumed.
\end{lemma}
\begin{proof}
Choose $\eta_m>0$ tending to zero and arc-length paths from $p$
to $a$ of lengths less than $L+\eta_m$. Restrict them to $[0,L]$.
By the single-path estimates in AC58, for $0\leq s\leq t\leq L$,
\[
 t-s-\eta_m\leq d(q_m(s),q_m(t))\leq t-s,
 \qquad d(q_m(t),a)\leq L+\eta_m-t.
\]
For each rational $t\in[0,L)$ the points $q_m(t)$ lie in a
compact ball: choose $t<r<L$ and use AC01 for
$\overline B(p,r)$. A countable diagonal makes all these
rational-time points converge. Their limiting distances are exactly
$|s-t|$. Completeness extends the limiting map isometrically to
$[0,L]$, just as in AC59. For rational $t<L$, the second estimate
passes to $d(\gamma(t),a)\leq L-t$; letting $t$ increase to $L$
identifies the endpoint with $a$. The original paths need not stay
inside the open $L$-ball, but all fixed interior-time prefixes do.
\end{proof}

\begin{lemma}[Shortening equal radial legs in the hyperbolic model; AC63]
\label{lem:alexandrov-model-angle-shortening}
Let $\lambda>0$, $L\geq r>0$, and let two isometric segments
$[p,a]$ and $[p,b]$ have length $L$. Write $a_r,b_r$ for their
points at distance $r$ from $p$. Suppose the comparison angle of
$(p;a,b)$ in curvature $-\lambda^2$ is at least
$\theta_0\in(0,\pi]$. Put $q=\sin(\theta_0/2)$ and
$B=\lambda r$. If $\theta_r$ is the comparison angle of
$(p;a_r,b_r)$ in that same model, then
\[
 \sin(\theta_r/2)\geq
 q-\frac{q^{-1}-q}{e^{2B}-1}.
\]
This assertion uses no curvature hypothesis on the source space.
In particular, if $\theta_{0,i}\to\pi/2$ and
$\lambda_i r_i\to\infty$, then
$\liminf_i\theta_{r_i}\geq\pi/2$.
\end{lemma}
\begin{proof}
Set $c=d(a,b)$ and $A=\lambda L$. The model cosine law gives
$\sinh(\lambda c/2)\geq q\sinh A$.
As in AC57, $\sinh(A+\log q)\leq q\sinh A$, so
\[
 c\geq2L+\frac{2}{\lambda}\log q.
\]
The source triangle inequality, using the two known radial tails,
gives
\[
 c_r:=d(a_r,b_r)\geq c-2(L-r)
       \geq2r+\frac{2}{\lambda}\log q.
\]
Apply the equal-leg model identity a second time and the
monotonicity of $\sinh$:
\[
 \begin{aligned}
 \sin(\theta_r/2)
 &=\frac{\sinh(\lambda c_r/2)}{\sinh B}\\
 &\geq\frac{\sinh(B+\log q)}{\sinh B}
 =\frac{q-q^{-1}e^{-2B}}{1-e^{-2B}}
 =q-\frac{q^{-1}-q}{e^{2B}-1}.
 \end{aligned}
\]
If the lower bound is negative it is simply uninformative, not
invalid; no positivity of $B+\log q$ was used. In the sequential
case $q_i\to1/\sqrt2$, so the error term tends to zero. The
continuity and monotonicity of $\arcsin$ on $[0,1]$, after
replacing a negative lower bound by zero, give the liminf assertion.
Degenerate comparison triangles are included by the same identity.
\end{proof}

For the specific cross-pair data of AC54, discard a finite prefix so
$0<\sigma_i\leq1$, and take
\[
 L_i=\sigma_i^{-1},\qquad r_i=L_i/100,\qquad
 \lambda_i=\sqrt{\sigma_i},\qquad
 q_i=\sin(\pi/4-\sigma_i/2).
\]
Then $8r_i<L_i$ and $\lambda_i r_i=1/(100\sqrt{\sigma_i})\to\infty$.
The explicit common lower angle for all cross pairs is
\[
 \beta_i=2\arcsin\!\left(
 \max\left\{0,q_i-\frac{q_i^{-1}-q_i}
                         {e^{2\lambda_i r_i}-1}\right\}\right)
 \longrightarrow\pi/2.
\]
The first shortening does not assert exact angle monotonicity;
AC63 proves precisely this asymptotic replacement. No comparison
on a triangle reaching the original endpoint sphere is used.

\subsection{Comparison inside the buffer and passage to the limit}

\begin{lemma}[Comparison-angle monotonicity on buffered radial segments; AC64]
\label{lem:alexandrov-buffered-angle-monotonicity}
Let $X$ be a complete length space, $p\in X$, $\kappa>0$, and
$0<8r<L$. Suppose $B_X(p,L)$ is locally compact and has local
Alexandrov curvature at least $-\kappa$ and Hausdorff dimension
at most $n$, in the ambient/intrinsic convention of AC41.
For any two isometric radial segments $\alpha,\beta:[0,r]\to X$
starting at $p$, and $0<s,t\leq r$, one has, conditional on
the local comparison input retained in AC02,
\[
 \Theta_{-\kappa}\bigl(s,t,d(\alpha(s),\beta(t))\bigr)
 \geq
 \Theta_{-\kappa}\bigl(r,r,d(\alpha(r),\beta(r))\bigr).
\]
All endpoint distances here are ambient distances. The comparison
is used only inside the open $8r$-ball, not at radius $L$.
\end{lemma}
\begin{proof}
Put $U=B_X(p,8r)$. Its intrinsic metric is finite and locally agrees
with the ambient metric, by the length-space argument in AC02.
AC01 makes the ambient closed $8r$-ball compact, since $8r<L$.
Thus $U$ is separable and second countable in the common local
topology. A countable cover by metric-equality neighborhoods shows
that its intrinsic Hausdorff dimension is at most $n$. The local
curvature condition transfers on the same neighborhoods.

Rescale all distances by $\lambda=\sqrt\kappa$. The local
curvature lower bound becomes $-1$, and apply AC02 with
$R=\lambda r$ and center $p$. Its open $8R$-ball is exactly
$U$ with the rescaled intrinsic metric. Its proof supplies complete
intrinsic $6R$-balls at all relevant vertices, and ambient/intrinsic
distance equality for the radial triangles. In original units,
$D=3r$ and the required complete intrinsic radii are $2D=6r$;
all evaluated vertices lie inside $B(p,2r)$, with
$d(p,v)+6r<8r$. These are the KL Section 3.3 safeguards.

Apply point-on-side lower comparison to the triangle
$(p,\alpha(r),\beta(r))$, replacing $\alpha(r)$ by
$\alpha(s)$. Its distance to $\beta(r)$ is at least the model
distance, hence its comparison angle at $p$ cannot decrease,
by the model cosine law. Apply the same argument to
$(p,\alpha(s),\beta(r))$, replacing $\beta(r)$ by $\beta(t)$.
Both triangles and their radial subtriangles are covered by AC02.
This gives the displayed inequality. Endpoints of a side give
equality, and degenerate triangles follow by the model formulas.
This expands the monotonicity consequence, not the source-qualified
local Toponogov theorem itself.
\end{proof}

\begin{proposition}[The long-strainer lines are mutually orthogonal; AC65]
\label{prop:alexandrov-long-strainer-orthogonality}
Use the hypotheses and the specified proper limit of AC54, without
assuming its conclusion. Choose the source radial segments to each
$a_{i,j}^\pm$ by AC62. On a common subsequence their fixed-time
prefixes converge to lines $\gamma_j$ as in AC59. For every pair
of distinct indices $j,h$, their positive rays have Alexandrov
angle $\pi/2$ at $p$. This assertion is conditional on AC41's
comparison input and the same local Toponogov input in AC02;
it requires no line-splitting theorem.
\end{proposition}
\begin{proof}
Work on the tail $0<\sigma_i\leq1$. The local structure producer
AC07 supplies local compactness throughout $B(p_i,L_i)$, so AC62
applies even to endpoints on its boundary. Each chosen radial
segment is an admissible AC58 path for any positive error
$\eta_i\to0$. AC57 gives the vanishing absolute opposite
excess. AC59 then extracts all $k$ lines using these same segments
and the actual maps to the specified AC41 limit.

Shorten these segments to $r_i=L_i/100$ and use AC63 with the
parameters above. Every cross pair of signs has comparison angle
at least $\beta_i\to\pi/2$ in curvature $-\sigma_i$.
Fix $T>0$. Eventually $T<r_i$, and AC64 gives the same lower
bound for the comparison angle between the time-$T$ prefixes.
Their distances to $p_i$ are exactly $T$, and their mutual
distances converge by AC59 and the distortion bound. AC35 now
applies to finite side triples, so for each choice of signs
\[
 \Theta_0\bigl(T,T,
 d_X(\gamma_j(\epsilon T),\gamma_h(\epsilon' T))\bigr)
 \geq\pi/2,\qquad \epsilon,\epsilon'\in\{-1,1\}.
\]
In particular each such cross distance is at least $\sqrt2\,T$.
The points $p$, $\gamma_j(T)$, $\gamma_h(T)$ and $\gamma_h(-T)$ are
therefore distinct. The line identity gives comparison angle
$\pi$ between the last two at $p$. Apply the global nonnegative
four-point comparison supplied by AC41:
\[
 \pi+
 \Theta_0\bigl(T,T,d(\gamma_j(T),\gamma_h(T))\bigr)+
 \Theta_0\bigl(T,T,d(\gamma_j(T),\gamma_h(-T))\bigr)
 \leq2\pi.
\]
Each of the last two angles is at least $\pi/2$, so both equal
$\pi/2$. This holds for every $T>0$ on the one chosen subsequence.
Letting $T$ decrease to zero gives the asserted Alexandrov angle.
No interchange with an escaping-endpoint limit occurs: the first
shortening was numerical, the second had a verified buffer, and
AC35 was used only after fixing $T$.
\end{proof}

\subsection{Closing the geometric input without changing its hypotheses}

\begin{theorem}[Written assembly of the long-strainer coordinate input; AC66]
\label{thm:alexandrov-long-strainer-input-assembly}
The conclusion of AC54 has a written proof conditional on the
local comparison input in AC02/AC41 and the aligned line-splitting
input AC43. In its specified limit $X$ there is an onto pointed
splitting $\alpha:X\to\mathbb R^k\times Y$ with
\[
 \alpha(\gamma_j(t))=(te_j,y),\qquad
 \pi_j\alpha=-b_j,\qquad
 \sup_{x\in\overline B_{X_i}(p_i,R)}
 |h_{i,j}(x)-\pi_j\alpha(f_i(x))|\longrightarrow0
\]
for each fixed $R>0$ and $j$. The same original coordinates
$h_{i,j}=d(p_i,a_{i,j}^+)-d(\,\cdot\,,a_{i,j}^+)$ are retained.
No stronger source curvature, larger controlled ball or changed
strainer scale is required.
\end{theorem}
\begin{proof}
AC65 supplies mutually orthogonal oriented lines from the chosen
source segments. Apply AC43 to each line individually and then
AC60 to identify $h_{i,j}\to-b_j$ locally uniformly, as in AC61.
The lines are the same lines used in the orthogonality argument.
AC53 now gives one onto splitting aligning all their oriented
axes; its coordinate formula is $\pi_j\alpha=-b_j$. Substitution
gives the displayed convergence. For $k=1$ orthogonality is vacuous;
for all finite $k$ the extraction and coordinate arguments retain
one common subsequence. This proves the entire AC54 conclusion
under the existing source-qualified comparison and splitting inputs.
\end{proof}

Consequently AC55's fixed-rank compactness contradiction now has a
written geometric middle: AC66 followed by AC52 produces exactly
the prescribed Euclidean coordinates. This does not reprove local
Toponogov or the global splitting theorem, nor supply an explicit
parameter modulus. AC56's separate metric route remains independent
of AC54/AC55. AC67--AC70 next expand the exact-scale endpoint
adjustment in reverse KL 4.15(1) under globally nonnegative
curvature; AC71--AC76 below write the separate compatibility argument
for KL 4.17.

\section{Exact-scale strainers from an approximate splitting}
\label{sec:alexandrov-reverse-strainer}

KL Lemma 4.15(1) (printed 31/PDF 26) goes in the reverse
direction: a sufficiently accurate approximate $k$-splitting of a
complete nonnegatively curved Alexandrov space gives a good
$k$-strainer at the exact prescribed scale. Its supplied proof is one
sentence. The following estimates expand the endpoint choice and
radius adjustment. They use the actual map of AC49, including its
non-strict infimum bound, and do not assume a continuous inverse.

\subsection{Lift farther out, then shorten to the prescribed sphere}

\begin{lemma}[Exact-radius lifting of the coordinate axes; AC67]
\label{lem:alexandrov-exact-radius-lifting}
Let $X$ be a geodesic metric space, $k\geq1$, $r>0$, and
$0<\varepsilon<1$. Suppose $\phi:(X,p)\to
(\mathbb R^k\times Y,(0,y))$ is a normalized
$(k,\varepsilon)$-splitting as in AC49, and
\[
 r+5\varepsilon<\varepsilon^{-1}.
\]
For every signed coordinate vector $v\in\{e_1,-e_1,\ldots,e_k,-e_k\}$
there is $a_v\in X$ such that
\[
 \begin{aligned}
 d(p,a_v)&=r,\\
 d_P\bigl(\phi(a_v),(rv,y)\bigr)&<14\varepsilon,\\
 \bigl|d(a_v,a_w)-r\|v-w\|_2\bigr|&<27\varepsilon
 \qquad\hbox{for all signed axes }v,w.
 \end{aligned}
\]
The factor $Y$ need not be proper, complete or geodesic. No
continuity of $\phi$ or extension of a geodesic beyond its endpoint
is required.
\end{lemma}
\begin{proof}
Put $s=r+4\varepsilon$ and $z_v=(sv,y)$. The strict domain
inequality gives $s<\varepsilon^{-1}-\varepsilon$, so each $z_v$
is in the target coverage ball. The distance-to-image bound is only
an infimum bound at most $\varepsilon$. Choose actual points
$x_v\in B_X(p,\varepsilon^{-1})$ with
\[
 d_P(\phi(x_v),z_v)<2\varepsilon.
\]
Pointedness and distortion give, with $\ell_v=d(p,x_v)$,
\[
 |\ell_v-s|<3\varepsilon,\qquad
 r+\varepsilon<\ell_v<r+7\varepsilon.
\]
Choose a minimizing unit-speed segment from $p$ to $x_v$ and let
$a_v$ be its point at parameter $r$. Then $d(p,a_v)=r$ and
$d(a_v,x_v)=\ell_v-r<7\varepsilon$. Both $a_v$ and $x_v$ are
in the source distortion ball. Therefore
\[
 \begin{aligned}
 d_P(\phi(a_v),(rv,y))
 &\leq d_P(\phi(a_v),\phi(x_v))
       +d_P(\phi(x_v),z_v)+d_P(z_v,(rv,y))\\
 &<7\varepsilon+\varepsilon+2\varepsilon+4\varepsilon
 =14\varepsilon.
 \end{aligned}
\]
For two axes, distortion and the two lifting errors give
\[
 \bigl|d(x_v,x_w)-s\|v-w\|_2\bigr|<5\varepsilon.
\]
The two radial tails change the source distance by less than
$14\varepsilon$. Replacing $s$ by $r$ in the model distance
changes it by at most $8\varepsilon$, since $\|v-w\|_2\leq2$.
Their sum is less than $27\varepsilon$, proving the last bound.
Only finitely many choices were made, with constants independent of
$k$. In particular we never choose an attained nearest point from
the infimum bound, or extend a lifted point outward along a geodesic.
\end{proof}

\subsection{Convert chord errors into strict comparison-angle bounds}

\begin{lemma}[A positive margin for equal-radius strainer angles; AC68]
\label{lem:alexandrov-strainer-angle-margin}
For $r>0$ and $0<\alpha\leq1$, define
\[
 \begin{aligned}
 \mu(r,\alpha)=\min\bigl\{
 &2r\bigl(1-\cos(\alpha/2)\bigr),\\
 &\sqrt2\,r-2r\sin(\pi/4-\alpha/2)\bigr\}>0.
 \end{aligned}
\]
Suppose points $a_v$ in a metric space satisfy $d(p,a_v)=r$ for
all signed coordinate axes, and, for some $0\leq D<\mu(r,\alpha)$,
\[
 \bigl|d(a_v,a_w)-r\|v-w\|_2\bigr|\leq D.
\]
Then their Euclidean comparison angles at $p$ satisfy
\[
 \begin{aligned}
 \Theta_0(p;a_{e_j},a_{-e_j})&>\pi-\alpha,\\
 \Theta_0(p;a_{\epsilon e_j},a_{\epsilon' e_h})
 &>\pi/2-\alpha
 \quad(j\ne h,\ \epsilon,\epsilon'\in\{-1,1\}).
 \end{aligned}
\]
All the labeled endpoints are distinct. These are comparison
angles in curvature zero; the assertion requires no curvature
hypothesis on the metric space itself.
\end{lemma}
\begin{proof}
The first entry in the minimum is positive because
$0<\alpha/2\leq1/2$. The second is positive by strict monotonicity
of sine on $[0,\pi/2]$, since
$0<\pi/4-\alpha/2<\pi/4$. For an equal-leg comparison triangle
with legs $r$ and opposite side $c\in[0,2r]$, the Euclidean cosine
law, including the degenerate endpoints, gives
\[
 \Theta_0(r,r,c)=2\arcsin\bigl(c/(2r)\bigr).
\]
For an opposite pair the assumed error gives
$c\geq2r-D>2r\cos(\alpha/2)$, hence its angle is greater than
$\pi-\alpha$. For different axes it gives
$c\geq\sqrt2\,r-D>2r\sin(\pi/4-\alpha/2)$, hence the angle
is greater than $\pi/2-\alpha$. The latter lower threshold and
the opposite-pair threshold are positive. Thus different signed
axes give different endpoints, also distinct from $p$ because
$r>0$. No angle comparison between the source and a model triangle
was used: these angles are calculated from the three distances.
\end{proof}

\begin{proposition}[A quantitative metric reverse implication; AC69]
\label{prop:alexandrov-metric-reverse-strainer}
For any $r>0$ and $0<\alpha\leq1$, put
\[
 \varepsilon_0(r,\alpha)=
 \min\left\{\frac18,\frac1{4(r+1)},
                   \frac{\mu(r,\alpha)}{60}\right\}>0.
\]
For every finite $k\geq1$, every geodesic pointed metric space
with a normalized $(k,\varepsilon)$-splitting,
$0<\varepsilon\leq\varepsilon_0(r,\alpha)$, has the equal-radius
configuration of AC68 at radius $r$, with its strict angle bounds.
Its endpoints additionally obey the $14\varepsilon$ map-coordinate
bound of AC67. The same $\varepsilon_0$ works for every rank $k$.
\end{proposition}
\begin{proof}
Since $\varepsilon\leq1/8$ and
$\varepsilon^{-1}\geq4(r+1)$, one has
\[
 r+5\varepsilon\leq r+5/8<4(r+1)\leq\varepsilon^{-1}.
\]
Thus AC67 supplies the endpoints with pairwise chord error less
than $27\varepsilon$. Take $D=27\varepsilon$; then
\[
 D\leq\frac{27}{60}\mu(r,\alpha)<\mu(r,\alpha).
\]
AC68 gives all strict angle inequalities. Positivity of the minimum
and $\varepsilon_0<1$ verify the approximation-parameter convention.
No dimension estimate, compactness limit, Toponogov theorem or
line-splitting theorem enters this metric estimate.
\end{proof}

\subsection{Specialization to the exact KL hypothesis and convention}

\begin{corollary}[The exact-scale reverse KL strainer implication; AC70]
\label{cor:alexandrov-kl-reverse-strainer}
Let $\delta>0$, $r=\delta^{-1}$, and
$\alpha=\min\{\delta,1\}$. Choose
$\delta'=\varepsilon_0(r,\alpha)$. For every $k\geq1$, if
$(X,p)$ is a complete pointed globally nonnegatively curved
finite-dimensional Alexandrov space with a normalized
$(k,\delta')$-splitting, then $p$ has a KL $k$-strainer of
quality $\delta$ at scale exactly $\delta^{-1}$, using
curvature-zero comparison triangles. The same conclusion holds
for any supplied approximation parameter
$0<\varepsilon\leq\delta'$. The chosen strainer obeys AC67's
map-coordinate estimate for that supplied map.
\end{corollary}
\begin{proof}
In KL's standing convention the source is a complete length space
of finite Hausdorff dimension. AC07 supplies local compactness
at every point, viewing the nonnegative lower bound as a local
lower bound of $-1$ for that producer. This does not change the
curvature-zero comparison model in the conclusion. For each
endpoint $x\ne p$, use AC62 with $L=d(p,x)$ to obtain a minimizing
segment from $p$ to $x$. The same argument at every center proves
that $X$ is geodesic. This makes explicit the shortest-path input
also supplied by BBI Theorem 2.5.23.

Apply AC69 at radius $r=\delta^{-1}$ and quality
$\alpha=\min\{\delta,1\}$. Since $\alpha\leq\delta$, its strict
angle inequalities imply the KL inequalities with quality $\delta$.
Every endpoint is at distance exactly $r$; for $\delta>1$ only the
auxiliary quality was reduced, not the requested radius. The same
argument applies to every smaller positive approximation parameter.
The constants are independent of $k$ and of a dimension bound.
This is a quantitative expansion of KL 4.15(1), conditional only
on the existing local-structure and geodesic inputs when obtaining
the latter from the Alexandrov hypotheses.
\end{proof}

The endpoint convention matters here. KL Definition 3.6 uses a
single exact radius and the cross-angle tolerance $\delta$.
BBI Definition 10.8.9 (printed 381/PDF 396) has no prescribed
common radius and uses cross-angle tolerance $10\delta$.
AC70 follows the KL convention directly. The constants $4$, $14$,
$27$ and $60$ above are this blueprint's explicit estimates, not
constants quoted from either source.

The metric core AC67--AC69 deliberately separates the Euclidean
comparison-angle calculation from curvature of the source.
It does not assert that an arbitrary geodesic metric space is
Alexandrov, or provide strainers measured in a negative-curvature
model without another estimate. AC70 keeps KL's global
nonnegative-curvature hypothesis. Neither the local almost-nonnegative
producer AC55 nor its compactness/splitting proof is a dependency
of this reverse argument. In particular, the two implications of
KL 4.15 are not used to prove each other. Compatibility of two
different approximate splittings is addressed separately below.

\section{Compatibility of the original approximate splittings}
\label{sec:alexandrov-splitting-compatibility}

KL Definitions 4.1--4.2 and 4.8 distinguish equivalence of exact
products, closeness of approximate products, and compatibility of
different ranks. Lemma 4.17 (printed 32--33/PDF 27--28) makes
compatibility uniform when one excludes an extra Euclidean factor.
We retain its direction of comparison and both original maps.

\subsection{Directed compatibility and cancellation of a common coordinate}

\begin{definition}[Normalized directed approximate compatibility; AC71]
\label{def:alexandrov-directed-compatibility}
Let $0\leq j\leq k$, $0<\tau<1$, and let
\[
 \phi=(u,z):(X,p)\to(\mathbb R^j\times A,(0,a)),\qquad
 \psi=(v,w):(X,p)\to(\mathbb R^k\times B,(0,b))
\]
be normalized approximate splittings. Say that $\phi$ is
$\tau$-compatible with $\psi$ if there exist an orthogonal
isomorphism $Q=(Q_1,Q_2):\mathbb R^k\to
\mathbb R^j\times\mathbb R^{k-j}$ and whole-space pointed
KL $\tau$-approximations
\[
 E:\mathbb R^j\to\mathbb R^j,\qquad
 F:(\mathbb R^{k-j}\times B,(0,b))\to(A,a)
\]
such that, writing $P=\mathbb R^j\times A$, for every
$x\in B_X(p,\tau^{-1})$,
\[
 d_P
 \bigl((E(Q_1v(x)),F(Q_2v(x),w(x))),\phi(x)\bigr)\leq\tau.
\]
All products are Pythagorean. The approximation tests for $E$ and
$F$ include distortion and coverage of their own factor balls;
control just on the image of $X$ is insufficient.
\end{definition}

This is KL 4.8 after translating Euclidean basepoints to zero:
the $j$-splitting induced from $\psi$ is close \emph{to} $\phi$,
so the factor maps run from the induced factors to those of $\phi$.
A pointed Euclidean isometry is orthogonal, by preservation of
inner products and linearity on an orthonormal basis.
We do not assert symmetry at the same error. When $j=k$, the
zero-dimensional Euclidean factor is a singleton; the transverse
map $F$ is still required.

\begin{lemma}[Cancellation with a shared Euclidean coordinate; AC72]
\label{lem:alexandrov-shared-coordinate-cancellation}
Suppose $a:X\to\mathbb R^j\times A$ and
$c:X\to\mathbb R^j\times C$ are onto pointed isometries and
$\pi_E a=\pi_E c$. There is a unique pointed onto isometry
$H:C\to A$ such that
\[
 a=(\operatorname{Id},H)\circ c.
\]
No curvature, geodesicity or completeness is needed.
\end{lemma}
\begin{proof}
For $q\in C$, define $H(q)$ by
$a(c^{-1}(0,q))=(0,H(q))$. Distances on the zero slices show
that $H$ is an isometry. For any $z\in A$, the point
$a^{-1}(0,z)$ has zero first coordinate under $c$, so $H$ is onto.
It preserves the basepoint. If $c(x)=(t,q)$ and $a(x)=(t,z)$,
compare $x$ with $c^{-1}(0,q)$ in both products:
\[
 \|t\|_2^2=\|t\|_2^2+d_A(z,H(q))^2.
\]
Thus $z=H(q)$, proving the formula everywhere. Its restriction
to the zero slice also proves uniqueness. No ambient embedding of
one transverse factor into the other was assumed.
\end{proof}

\subsection{Exact compatibility when the larger splitting is maximal}

\begin{proposition}[Maximal exact splittings are compatible; AC73]
\label{prop:alexandrov-maximal-exact-compatibility}
Let $X$ be a complete, proper, geodesic, finite-dimensional
nonnegatively curved Alexandrov space. Suppose there are exact pointed splittings
\[
 a:X\to\mathbb R^j\times A,\qquad
 b:X\to\mathbb R^k\times B,\qquad 1\leq j\leq k,
\]
with factor basepoints $a_0,b_0$, and suppose $X$ has no exact
$(k+1)$-splitting. Conditional on AC43,
there are an orthogonal $Q=(Q_1,Q_2)$ and a pointed onto isometry
$H:\mathbb R^{k-j}\times B\to A$ with
\[
 a(x)=\bigl(Q_1\pi_E b(x),
             H(Q_2\pi_E b(x),\pi_B b(x))\bigr)
 \quad(x\in X).
\]
This is the CBB-zero case needed below, not a claim to have
reproved KL 4.4 for every metric space.
\end{proposition}
\begin{proof}
For each $1\leq\ell\leq j$ use the oriented axis line
$\gamma_\ell(t)=a^{-1}(te_\ell,a_0)$ through the basepoint.
In the $b$-product it has the form
\[
 b(\gamma_\ell(t))=(t\xi_\ell,\eta_\ell(t)),\qquad
 d_B(\eta_\ell(s),\eta_\ell(t))=c_\ell|s-t|,
 \qquad \|\xi_\ell\|_2^2+c_\ell^2=1.
\]
Indeed, equality in the product triangle inequality gives equality
in each coordinate and proportional coordinate-distance pairs.
Applying this to every three ordered parameters gives constant-speed
coordinate geodesics, as in AC53 and KL 4.5(1). The Euclidean
coordinate is linear since it vanishes at zero.

AC44 supplies the complete proper finite-dimensional nonnegative
geometry of $B$. If $c_\ell>0$, reparametrizing $\eta_\ell$ gives
an actual line in $B$. AC43 then splits $B$ as $\mathbb R\times D$,
and composition with $b$ gives an exact $(k+1)$-splitting of $X$,
a contradiction. Hence $c_\ell=0$ for all $\ell$.
The vectors $\xi_\ell$ are unit and mutually orthogonal: the
pairwise distances between the axis points at time one are
$\sqrt2$ in the $a$-product and are unchanged in the $b$-product.

For $a(x)=(u,z)$ and $b(x)=(v,w)$, the positive-ray distance
limit computed in either product gives
\[
 \lim_{t\to\infty}\bigl(d(x,\gamma_\ell(t))-t\bigr)
 =-u_\ell=-\langle v,\xi_\ell\rangle.
\]
These identities follow directly by rationalizing the Pythagorean
formulas; they do not assume compatibility in advance. Extend
$\xi_1,\ldots,\xi_j$ to an orthonormal basis of $\mathbb R^k$.
Let $Q$ be the orthogonal coordinate map in that basis, so its first
$j$ coordinates are the displayed inner products. The induced exact
splitting $c=(Q_1\pi_E b,(Q_2\pi_E b,\pi_B b))$ has the same
first coordinate as $a$. Apply AC72 to obtain $H$ and the global
identity. This proof uses a line in the transverse factor to invoke
AC43; it imports no general parallel-line foliation theorem.
\end{proof}

\subsection{Keep both full product maps on the same limiting space}

\begin{lemma}[Simultaneous full product convergence; AC74]
\label{lem:alexandrov-two-product-convergence}
Suppose $(X_i,p_i)\to(X,p)$ with $X$ complete and proper, and
there are supplied normalized splittings
\[
 \phi_i=(u_i,z_i):X_i\to\mathbb R^j\times A_i,\qquad
 \psi_i=(v_i,w_i):X_i\to\mathbb R^k\times B_i
\]
of fixed ranks whose two positive approximation parameters tend
to zero. On a common subsequence there are complete proper factor
limits $A,B$, exact pointed splittings $a:X\to\mathbb R^j\times A$
and $b:X\to\mathbb R^k\times B$, and pointed approximations
$f_i:X_i\to X$, $g_i:A_i\to A$, $h_i:B_i\to B$ on growing
balls, such that on every fixed source ball
\[
 \begin{aligned}
 d\bigl((u_i(x),g_i(z_i(x))),a(f_i(x))\bigr)&\longrightarrow0,\\
 d\bigl((v_i(x),h_i(w_i(x))),b(f_i(x))\bigr)&\longrightarrow0
 \end{aligned}
\]
uniformly in $x$. Write $a_i,b_i$ for the source factor basepoints
and $a_0,b_0$ for those of $A,B$. The same $f_i$ and subsequence
work for both maps.
Neither the $X_i$ nor the individual factors need be proper.
\end{lemma}
\begin{proof}
Fix actual maps $f_i$ realizing the given convergence on growing
balls. Apply AC50 and the proof of AC51 first to $\phi_i$.
That proof constructs $g_i$, forms $q_i=(\operatorname{Id},g_i)\phi_i$,
and proves the \emph{full-product} estimate
$d(q_i(x),a(f_i(x)))\to0$ before projecting onto Euclidean
coordinates. Retain this stronger estimate and these factor maps.
On the resulting subsequence run the same construction for $\psi_i$,
using the already chosen $f_i$. This gives $h_i$ and $b$ on a
further subsequence; the first uniform convergence persists there.
The factor limits are complete and proper by AC50. All factors of
images of bounded source points have bounded radius by pointedness
and distortion. Thus the compositions are defined on padded domains,
exactly as in AC51. No independent replacement of the specified
limit or of either original Euclidean coordinate is made.
\end{proof}

\begin{proposition}[Transfer compatible limits to the original maps; AC75]
\label{prop:alexandrov-compatibility-transfer}
Use the full convergence data in AC74 and suppose there are
$Q,H$ with the exact identity in AC73. Then for every fixed
$0<\tau<1$, all sufficiently large $i$ have $\phi_i$
$\tau$-compatible with $\psi_i$ in AC71. One can take $E$ to be
the identity and the same $Q$ as in the limit. The required maps
$F_i:\mathbb R^{k-j}\times B_i\to A_i$ are constructed explicitly
from factor approximations and local quasi-inverses.
This implication is metric; it uses no curvature or splitting theorem.
\end{proposition}
\begin{proof}
Fix $\tau$ and put $R=\tau^{-1}+2$. On a sufficiently late tail,
choose the AC74 factor maps with
$\mathcal A_{8R,e_i}(g_i)$ and $\mathcal A_{8R,e_i}(h_i)$,
where $e_i\to0$ and $e_i<\tau/100$.
Let $\rho_i\to0$ bound both full-product discrepancies in AC74
on the source $R$-ball. The original approximation parameters
also tend to zero; require them to be less than one.

Write $C_i=\mathbb R^{k-j}\times B_i$ and
$C=\mathbb R^{k-j}\times B$. The product estimate in AC51 gives
$K_i=(\operatorname{Id},h_i):C_i\to C$ with
$\mathcal A_{8R,3e_i}$.
MC09 gives a pointed $t_i:\overline B_A(a_0,4R)\to A_i$ with
$\mathcal A_{4R,4e_i}$ and
\[
 d_A(g_i(t_i(q)),q)<e_i\qquad(d_A(a_0,q)\leq4R).
\]
All its values lie in the original $8R$-ball of $A_i$.
Set $J_i=H K_i$ and $c_i=(0,b_i)$. For
$q\in\overline B_{C_i}(c_i,2R)$,
its image $J_i(q)$ has radius less than $2R+3e_i<4R$, so define
\[
 F_i(q)=t_i(J_i(q)).
\]
The MC08 conditions $14e_i<2R$ and $2R+3e_i\leq4R$ hold.
Consequently $\mathcal A_{2R,14e_i}(F_i)$.
Extend $F_i$ outside this ball by $a_i$. A target point of radius
less than $\tau^{-1}-\tau$ has a coverage witness of radius
\[
 <\tau^{-1}-\tau+28e_i<\tau^{-1}.
\]
Thus the extension is a whole-space pointed KL $\tau$-approximation
on the transverse factors. This proves its own factor coverage,
not merely closeness on source images.

For $x\in B_{X_i}(p_i,\tau^{-1})$, put
$q_i(x)=(Q_2v_i(x),w_i(x))$. Both $q_i(x)$ and $z_i(x)$ have
radius less than $R$, since the original maps are pointed and their
distortions are less than one. The limit identity and full-product
bounds give
\[
 \begin{aligned}
 \|u_i(x)-Q_1v_i(x)\|_2&\leq2\rho_i,\\
 d_A(g_i(z_i(x)),J_i(q_i(x)))&\leq2\rho_i.
 \end{aligned}
\]
The inverse estimate above and distortion of $g_i$ then imply
\[
 d_{A_i}(z_i(x),F_i(q_i(x)))<2e_i+2\rho_i.
\]
Both points are in its $8R$ domain. The product error is therefore
less than $2e_i+4\rho_i$, uniformly on the required source ball.
Take $i$ large enough that this is less than $\tau$. Together with
the identity Euclidean map and the factor approximation just proved,
this is precisely the directed compatibility in AC71. No continuity
or properness of the approximating factors was used.
\end{proof}

\subsection{The uniform exclusion theorem and its source correction}

\begin{theorem}[Uniform compatibility without an extra factor; AC76]
\label{thm:alexandrov-uniform-compatibility}
Fix integers $1\leq j\leq k\leq n$ and $0<\tau,\nu<1$.
There exists $0<\varepsilon<1$ such that the following holds,
conditional on AC41's local comparison input and AC43.
Let $(Z,o)$ be a complete pointed length space with local
curvature at least $-\varepsilon$ and Hausdorff dimension at most
$n$ on $B(o,\varepsilon^{-1})$, in the convention of AC41.
Suppose $(Z,o)$ has no normalized $(k+1,\nu)$-splitting.
Then \emph{every} supplied normalized $(j,\varepsilon)$-splitting
is $\tau$-compatible with \emph{every} supplied normalized
$(k,\varepsilon)$-splitting.
The number $\varepsilon$ depends only on $j,k,n,\tau,\nu$.
\end{theorem}
\begin{proof}
If no such parameter existed, for each integer $i\geq1$ choose a
counterexample at $\varepsilon_i=1/(i+2)$, including its two
incompatible maps. The controlled balls have radii $i+2\to\infty$
and curvature bounds $-1/(i+2)\to0$. AC41 gives one complete
proper geodesic nonnegative limit $X$ of dimension at most $n$.
AC74 supplies the two exact limiting splittings on this same $X$,
with the full-product convergence of both original maps.

There is no exact $(k+1)$-splitting of $X$. Otherwise compose it
with the convergence maps from the sources. AC52, using its exact
Euclidean coordinate and discrepancy zero, produces a normalized
$(k+1,\nu)$-splitting of every sufficiently late source. This
contradicts the fixed exclusion hypothesis. This argument transfers
an existing exact splitting back to the sources; it does not assume
that the absence of a splitting automatically passes to a limit.

Apply AC73 to the two limiting splittings to obtain $Q,H$.
AC75 makes the original two source maps $\tau$-compatible for all
large $i$, contradicting their selection. The counterexample choice
retains the universal quantifier over both maps. There is no claimed
explicit modulus for this compactness argument.
\end{proof}

This is KL 4.17 in the normalized parameter range. Its fixed
parameters are $\tau=\beta'_k$ and $\nu=\beta_{k+1}$; its three
chosen parameters may be taken equal,
$\delta=\beta_j=\beta_k=\varepsilon$. The radii are part of the
hypothesis, so no unstated monotonicity of a curvature/scale packet
is used. The proof does not require either strainer direction
AC54--AC55 or AC67--AC70.

\paragraph{Source discrepancy.}
The archived KL proof on printed 32/PDF 27 writes
$B(\star_{X_i},i^{-1})$ in items (4)--(5). The lemma's hypothesis,
with $\delta=i^{-1}$, instead requires $B(\star_{X_i},i)$;
shrinking balls would not supply Lemma 3.10's growing-region input.
This is an inferred correction verified against the rendered page,
not an entry on the inspected May 15, 2015 author correction sheet.
AC76 chooses its growing balls directly from the lemma's statement.
Its convergence errors are allowed to tend to zero along the
extracted subsequence; no prescribed $1/i$ convergence rate is assumed.

AC71--AC76 write the factor-map and maximality steps compressed in
KL's proof. The general metric-space common-refinement theorem KL 4.4
is broader than the CBB-zero exact argument used here and is not
claimed as newly proved. Basepoint transport is written below; smooth
adapted coordinates and fibration compatibility remain separate obligations.

\section{Basepoint transport and compatibility on overlapping balls}
\label{sec:alexandrov-basepoint-transport}

KL Lemma 4.10 (printed 30--31/PDF 25--26) recenters an approximate
splitting. Lemma 4.14 (printed 31/PDF 26) asserts that compatibility
survives, but omits the details. Two operations must be distinguished:
using the actual image as the new target basepoint, and replacing that
image by a nearby prescribed basepoint. MC10 supplies the first;
the following estimate supplies the second. All statements here are
metric and use no curvature, dimension or PC input.

\newpage
\subsection{Repair the target basepoint with a controlled error}

\begin{lemma}[One-point target repair; AC77]
\label{lem:alexandrov-target-basepoint-repair}
Let $f:\overline B_X(a,R)\to Y$ satisfy $\mathcal A_{R,e}$ with
$f(a)=q$. Suppose $d_Y(q,b)\leq\eta$, where $\eta\geq0$ and
$E=e+2\eta<R$. Define
\[
 \widehat f(x)=
 \begin{cases} b,&x=a,\\ f(x),&x\ne a.\end{cases}
\]
Then $\widehat f$, based at $a,b$, satisfies $\mathcal A_{R,E}$,
and $d_Y(\widehat f(x),f(x))\leq\eta$ everywhere on its domain.
No continuity of the repaired map is asserted.
\end{lemma}
\begin{proof}
For any pair of source points, the two image changes cost at most
$2\eta$. Adding the original strict distortion gives error less than
$e+2\eta=E$. If $d_Y(b,y)\leq R-E$, then
\[
 d_Y(q,y)\leq R-E+\eta=R-e-\eta\leq R-e.
\]
Original coverage gives $x\in\overline B_X(a,R)$ with
$d_Y(f(x),y)<e$. Its repaired image has distance less than
$e+\eta\leq E$ from $y$. This also covers $\eta=0$, because the
original error is strict. The new basepoint is exact by definition.
\end{proof}

\begin{proposition}[Uniform recentering with a nearby target; AC78]
\label{prop:alexandrov-kl-recentering}
For $0<\delta<1$ and $C\geq0$, put
\[
 \vartheta(\delta,C)=
 \min\left\{\frac{\delta}{100},
             \frac{1}{100(C+\delta^{-1}+2)}\right\}.
\]
Let $f:(X,p)\to(Y,q)$ be a whole-space pointed KL
$\varepsilon$-approximation with
$0<\varepsilon\leq\vartheta(\delta,C)$. If
$d_X(p,a)\leq C$ and $d_Y(f(a),b)\leq\varepsilon$, change
only the value at $a$ to $b$. The resulting whole-space map
$\widehat f:(X,a)\to(Y,b)$ is a KL $\delta$-approximation.
No completeness, properness or length-space hypothesis is needed.
\end{proposition}
\begin{proof}
Set $s=\delta^{-1}+1$ and $R_0=C+s+1$. The chosen tolerance gives
\[
 R_0<\varepsilon^{-1},\qquad 3\varepsilon<R_0,
 \qquad 9\varepsilon<s,\qquad 22\varepsilon<\delta.
\]
MC11 converts the given open-ball, non-strict KL approximation to
$\mathcal A_{R_0,3\varepsilon}$ on the closed $R_0$-ball.
This uses actual coverage witnesses within $2\varepsilon$,
without assuming attainment of the original infimum.
Since $d(p,a)+s\leq C+s<R_0$, MC10 recenters its restriction to
$\overline B_X(a,s)$ with error $9\varepsilon$ and target basepoint
$f(a)$. AC77 changes that target to $b$, giving error at most
$11\varepsilon$; enlarge the strict approximation error to
$11\varepsilon$ if necessary.

The entire open source $\delta^{-1}$-ball is inside this domain,
so its distortion is less than $11\varepsilon<\delta$.
For a target $y$ with $d_Y(b,y)<\delta^{-1}-\delta$, coverage of
$\mathcal A_{s,11\varepsilon}$ applies. Its witness $x$ satisfies
\[
 \begin{aligned}
 d_X(a,x)&<d_Y(b,\widehat f(x))+11\varepsilon\\
         &<d_Y(b,y)+22\varepsilon\\
         &<\delta^{-1}-\delta+22\varepsilon<\delta^{-1}.
 \end{aligned}
\]
Thus the witness lies in the required open source ball and has
image error less than $11\varepsilon<\delta$.
Outside the closed $s$-ball the original whole-space map remains
unchanged; no new values there enter the tests. This proves the KL
predicate with its exact basepoint and own target coverage.
\end{proof}

\subsection{Transport the original splittings and their factor maps}

\begin{corollary}[Normalized splitting at a new basepoint; AC79]
\label{cor:alexandrov-splitting-recenter}
Let $k\geq0$, $0<\delta<1$ and $C\geq0$. Suppose
$\phi=(u,z):(X,p)\to(\mathbb R^k\times A,(0,a_0))$ is a
normalized $(k,\varepsilon)$-splitting and
$0<\varepsilon\leq\vartheta(\delta,C)$. For every $c\in X$
with $d_X(p,c)\leq C$, the actual map
\[
 \phi_c(x)=(u(x)-u(c),z(x))
\]
is a normalized $(k,\delta)$-splitting of $(X,c)$ with residual
basepoint $z(c)$. The tolerance is independent of $k$ and $X$.
\end{corollary}
\begin{proof}
Apply AC78 with new target basepoint exactly $\phi(c)$, so no value
is changed. Translate the Euclidean target by $-u(c)$, an isometry.
This preserves every approximation test and places the Euclidean
basepoint at zero. The transverse metric and the actual residual
coordinate $z$ are unchanged. This proves KL 4.10 in our normalized
range, and also shows that completeness in its stated hypothesis is
not needed for this metric argument.
\end{proof}

\begin{theorem}[Compatibility after changing basepoint; AC80]
\label{thm:alexandrov-compatibility-recenter}
Fix $0<\delta<1$, $C\geq0$ and integers $0\leq j\leq k$.
Suppose the normalized $(j,\varepsilon)$-splitting
$\phi=(u,z):X\to\mathbb R^j\times A$ is
$\varepsilon$-compatible with the normalized $(k,\varepsilon)$-
splitting $\psi=(v,w):X\to\mathbb R^k\times B$ at $p$, and
\[
 0<\varepsilon\leq\vartheta(\delta,C+1).
\]
For every $c$ with $d_X(p,c)\leq C$, the original recentered maps
$\phi_c=(u-u(c),z)$ and $\psi_c=(v-v(c),w)$ are normalized
splittings of quality $\delta$, and $\phi_c$ is
$\delta$-compatible with $\psi_c$. One may retain the same
orthogonal map $Q$ as in the original compatibility witness.
This is uniform over $c$, both maps and both ranks.
\end{theorem}
\begin{proof}
The tolerance $\vartheta(\delta,C+1)$ is at most
$\vartheta(\delta,C)$, so AC79 gives the two recentered splittings.
Let $Q=(Q_1,Q_2),E,F$ witness the original compatibility in AC71.
Write
\[
 \begin{aligned}
 t_c&=Q_1v(c),& r_c&=(Q_2v(c),w(c)),\\
 \eta_E&=\|E(t_c)-u(c)\|_2,&
 \eta_F&=d_A(F(r_c),z(c)).
 \end{aligned}
\]
The original source point $c$ lies in $B(p,\varepsilon^{-1})$.
Pointedness and distortion of $\psi$ put $t_c$ and $r_c$ at distance
at most $C+\varepsilon<C+1$ from their respective old factor
basepoints. Compatibility at $c$ gives the stronger joint bound
\[
 \eta_E^2+\eta_F^2\leq\varepsilon^2.
\]
Apply AC78 separately to the whole-space KL $\varepsilon$-
approximations $E$ and $F$, with displacement bound $C+1$.
Repair $E(t_c)$ to $u(c)$ and $F(r_c)$ to $z(c)$, leaving their
other values unchanged; call the maps $\widehat E,\widehat F$.
They are KL $\delta$-approximations with these new factor basepoints.
Define the normalized witnesses
\[
 \begin{aligned}
 E_c(t)&=\widehat E(t+t_c)-u(c),\\
 F_c(s,b)&=\widehat F(s+Q_2v(c),b).
 \end{aligned}
\]
The second map is from
$(\mathbb R^{k-j}\times B,(0,w(c)))$ to $(A,z(c))$.
These source translations and the Euclidean target translation are
pointed isometries for the new basepoints, so $E_c,F_c$ retain
\emph{their own} KL $\delta$-approximation tests.

If $x\in B_X(c,\delta^{-1})$, then
\[
 d_X(p,x)<C+\delta^{-1}<\varepsilon^{-1},
\]
so the original compatibility estimate applies at $x$.
The new factor composition, after translating back the first target,
is exactly
\[
 \bigl(\widehat E(Q_1v(x)),
       \widehat F(Q_2v(x),w(x))\bigr).
\]
Its distance from the old factor composition is at most
$(\eta_E^2+\eta_F^2)^{1/2}\leq\varepsilon$.
This remains true when a factor argument equals its new basepoint
for many different source points: the corresponding map value changes
by at most its specified $\eta$ each time it is evaluated.
The original product error is at most $\varepsilon$, hence the new
error relative to $\phi_c$ is at most $2\varepsilon<\delta$.
Together with the factor tests, this is AC71 at the new basepoint.
The maps $\phi_c,\psi_c$ themselves were never repaired or replaced.
\end{proof}

This expands the omitted proof of KL 4.14, including the target
basepoint correction absent from a bare application of MC10.
The repairs need not be continuous, as allowed by the approximation
predicate. They establish neither smooth coordinate compatibility
nor a cocycle identity among maps chosen at different centers.
No CBB condition, dimension bound, higher-rank exclusion or use of
AC76 enters this transport theorem. It includes $j=k$ and $j=0$.

\subsection{A moving marked point in the same pointed limit}

\begin{proposition}[Moving-basepoint convergence; AC81]
\label{prop:alexandrov-moving-basepoint-limit}
Suppose $X$ is complete and there are actual pointed maps
$f_i:\overline B_{X_i}(p_i,R_i)\to(X,p)$ satisfying
$\mathcal A_{R_i,e_i}$, with $R_i\to\infty$ and $e_i\to0$.
Let $d_{X_i}(p_i,c_i)\leq C$ for a fixed $C\geq0$.
If $f_i(c_i)\to c\in X$, then $(X_i,c_i)\to(X,c)$ in MC06's
pointed convergence convention.
If $X$ is also proper, a subsequence with such a limit $c$ always
exists. If moreover $d_{X_i}(p_i,c_i)\to d$, then $d_X(p,c)=d$.
In particular, source separation one gives distinct limit points
at separation one, not an unidentified second copy of the limit.
\end{proposition}
\begin{proof}
Ignore the finite prefix on which $f_i(c_i)$ is not yet defined.
Fix a radius $s>0$ and an error $0<\alpha<s$.
For all sufficiently large $i$, $C+s\leq R_i$ and $3e_i<s$.
MC10 recenters the restriction to the closed $s$-ball about $c_i$,
with error $3e_i$ and target basepoint $f_i(c_i)$.
Put $\eta_i=d_X(f_i(c_i),c)\to0$ and use AC77 to repair the value
at $c_i$ to $c$. For large $i$ this has error
\[
 3e_i+2\eta_i<\alpha<s.
\]
Increasing the error to $\alpha$ proves the required all-radii,
all-errors convergence to the same underlying $X$.

For the subsequence assertion, the radial distortion bound gives
\[
 \bigl|d_X(p,f_i(c_i))-d_{X_i}(p_i,c_i)\bigr|<e_i.
\]
Eventually the images lie in the compact closed $(C+1)$-ball of $X$;
choose a convergent subsequence there. If the source radii tend to $d$,
the displayed bound and continuity of distance give $d_X(p,c)=d$.
Only the limit is required to be proper for this extraction. No
source properness, length structure or curvature is used.
\end{proof}

AC81 supplies the marked-point step compressed in KL 4.20's proof
(printed 33--34/PDF 28--29). It does not yet produce the first
limit, prove that either limiting basepoint is a cone apex, or
prove KL 4.19's two-apex line. AC82--AC87 below supply these
separate steps, with the existing comparison and splitting qualifications.

\section{Metric cones, two apices and closure under pointed limits}
\label{sec:alexandrov-overlapping-cones}

KL 4.19--4.20 (printed 33--34/PDF 28--29) require more than
invariance of a pointed space under rescaling. We retain actual radial
maps satisfying the Euclidean two-ray distance law. This makes the
line argument and the closure of cone structures explicit. No new
comparison or splitting theorem is assumed in the metric steps below.

\subsection{Radial data and the corrected cone convention}

\begin{definition}[Radial metric-cone data; AC82]
\label{def:alexandrov-radial-cone}
For a pointed metric space $(X,p)$ put
\[
 r_x=d(p,x),\qquad
 B_p(x,y)=\frac{r_x^2+r_y^2-d(x,y)^2}{2}.
\]
A radial metric-cone structure at $p$ consists of maps
$H_t:X\to X$ for every real $t\geq0$, with $H_0x=p$, $H_1x=x$,
and the identity
\[
 d(H_sx,H_ty)^2=s^2r_x^2+t^2r_y^2-2stB_p(x,y)
 \qquad(s,t\geq0,\ x,y\in X).
\]
The function $B_p$ is a metric kernel, not an asserted bilinear form.
The one-point space is allowed, with all maps constant. Neither
curvature, completeness, properness nor the length property is part
of this definition.
\end{definition}

For the BBI/AKP Euclidean cone over an angular metric space $\Sigma$,
the corrected distance law is
\[
 d_C((r,a),(s,b))^2
 =r^2+s^2-2rs\cos\bigl(\min\{\pi,d_\Sigma(a,b)\}\bigr).
\]
Here the zero-radius representatives are identified. BBI 3.6.16,
printed 93/PDF 108, misprints the cutoff as $r+s\leq\pi$;
the July 6, 2024 author errata, PDF 5, corrects it to
$d_\Sigma(a,b)\leq\pi$. AKP 6E, archived printed/PDF 67,
uses the clipped angular distance. Radial scaling
$H_t(r,a)=(tr,a)$ satisfies AC82 by this formula. We use this
comparison on the existing cone metric; we do not claim to have
formalized the angular quotient or its triangle inequality.

\begin{lemma}[Radial identities and the two-ray convention; AC83]
\label{lem:alexandrov-radial-cone-identities}
AC82 implies
\[
 \begin{aligned}
 H_t p&=p,& r_{H_t x}&=t r_x,\\
 d(H_sx,H_tx)&=|s-t|r_x,&d(H_tx,H_ty)&=t\,d(x,y),\\
 B_p(H_tx,y)&=tB_p(x,y),&H_s(H_tx)&=H_{st}x.
 \end{aligned}
\]
For $t>0$, $H_t$ is onto and has inverse $H_{1/t}$.
Every non-apex point lies on a unique unit-speed ray from $p$,
and its minimizing segment from $p$ is unique and radial.
For a nontrivial space, AC82 is equivalent to KL's two-ray
metric-cone convention: the space is covered by rays from its
apex, and each pair has the metric of two Euclidean rays from
the origin. The one-point case is the explicit degenerate extension.
\end{lemma}
\begin{proof}
Use $s=0$ in AC82 for the radial norm, then use $x=y$ and
$s=t$ for the two distance identities. The norm identity with
$x=p$ gives $H_tp=p$. Substituting $H_1y=y$ in the defining
identity and expanding $B_p(H_tx,y)$ gives the kernel identity.
In particular $B_p(H_tx,x)=t r_x^2$. Applying AC82 to the two
points $H_tx,x$, at parameters $s,st$, shows
$d(H_s(H_tx),H_{st}x)^2=0$. This also proves the inverse assertion.

For $x\ne p$, the map $u\mapsto H_{u/r_x}x$, $u\geq0$,
is a unit-speed ray through $x$. If a point $z$ on any minimizing
segment from $p$ to $x$ has $r_z=t r_x$ with $0\leq t\leq1$,
then $d(z,x)=(1-t)r_x$ and $B_p(z,x)=t r_x^2$. AC82 gives
$d(z,H_tx)^2=0$. For a point $z$ farther along any ray through
$x$, this uniqueness applied to the segment ending at $z$ gives
$x=H_{r_x/r_z}z$; invert that map to identify $z$ as well.

In the other direction, start with KL's family of rays.
Two such rays meeting at a positive radius have angle zero in
their Euclidean realization and therefore coincide at every
radius. Thus scaling along the ray through a non-apex point is
well defined. The Euclidean cosine law for each pair gives AC82,
including zero parameters. Conversely, in AC82 the triangle
inequality gives
\[
 -r_xr_y\leq B_p(x,y)\leq r_xr_y.
\]
For non-apex $x,y$, choose the angle
$\arccos(B_p(x,y)/(r_xr_y))\in[0,\pi]$.
The kernel identity realizes the distances between all points
of their two radial rays as distances on the corresponding
Euclidean rays. If the angle is zero the two rays coincide by
the same identity, so there is no failure of injectivity.
These rays cover $X$ by $H_1=\mathrm{id}$.
\end{proof}

This verifies the convention in KL Section 2.2, printed 19/PDF 14,
and Section 4.5. It does not identify arbitrary rescaling isometries
with cone data; BBI 8.2.1, printed 275/PDF 290, warns of that
distinction. KL's later point-cone case (printed 62/PDF 57) motivates
retaining the degenerate case. No angular-link curvature or
dimension theorem is needed here.

\subsection{Two different apices determine a global line}

\begin{theorem}[Two cone apices give a line; AC84]
\label{thm:alexandrov-two-cone-apices-line}
Let the same metric space $X$ carry AC82 structures at distinct
points $p,q$, and put $L=d(p,q)>0$. Then there is an isometric
map $\gamma:\mathbb R\to X$ with $\gamma(0)=p$ and
$\gamma(L)=q$. No curvature, completeness or length hypothesis
is needed for this metric assertion.
\end{theorem}
\begin{proof}
Write $H^p,H^q$ for the two radial structures and define, for
$s,u\geq0$,
\[
 \alpha(s)=H^p_{s/L}q,\qquad
 \beta(u)=H^q_{u/L}p,\qquad z_t=\beta(L+t)\quad(t\geq0).
\]
Both $\alpha$ and $\beta$ are unit-speed rays and $\beta(L)=p$.
Consequently $d(p,z_t)=t$ and $d(q,z_t)=L+t$, so
\[
 B_p(q,z_t)=\frac{L^2+t^2-(L+t)^2}{2}=-Lt.
\]
The cone identity at $p$ now proves the global cross-distance law
\[
 d(\alpha(s),z_t)^2=s^2+t^2+2st=(s+t)^2.
\]
Define $\gamma(s)=\alpha(s)$ for $s\geq0$ and
$\gamma(-t)=z_t$ for $t\geq0$. These prescriptions agree at zero.
Distances on either branch are parameter differences; the displayed
identity handles opposite signs. Thus every pair of real parameters
satisfies $d(\gamma(a),\gamma(b))=|a-b|$, as required.
\end{proof}

If $X$ is also complete proper geodesic, globally nonnegatively
curved and finite dimensional, the existing source-qualified
AC43 applies to this line. Its aligned isometry sends
$p$ to $(0,y)$ and $q$ to $(L,y)$ in $\mathbb R\times Y$.
Translating the real coordinate by $-L$ gives a pointed splitting
at $q$. This is the written reduction of KL 4.19 to AC43;
a locally minimizing extension alone would not suffice. We do
not additionally assert that the residual factor is a metric cone.
AKP 16.24, archived printed/PDF 252, has a stronger cone-factor
conclusion which is not needed here.

\subsection{Passing the actual radial structure to one proper limit}

\begin{proposition}[Proper pointed limits retain cone data; AC85]
\label{prop:alexandrov-cone-limit-closure}
Suppose $(X_i,p_i)$ carry AC82 structures and converge in MC06's
pointed convention to a complete proper pointed metric space
$(X,p)$. Then $(X,p)$ has an AC82 structure. The source spaces
need not be proper, length spaces or Alexandrov spaces.
\end{proposition}
\begin{proof}
Choose actual pointed maps $f_i$ satisfying
$\mathcal A_{R_i,e_i}$ with $R_i\to\infty$, $e_i\to0$,
by a diagonal choice in MC06. Properness supplies finite nets
in each closed integer ball of $X$. Their union, together with
$p$, gives a countable dense set $D$.
For each $a\in D$, coverage supplies, for all sufficiently large
$i$, a point $x_i(a)$ with
\[
 d(f_i(x_i(a)),a)<e_i,\qquad
 \bigl|d(p_i,x_i(a))-d(p,a)\bigr|<2e_i.
\]
Take $x_i(p)=p_i$ and assign $p_i$ on any undefined finite prefix.
For fixed $a,b\in D$, distortion and the two coverage errors
show $d(x_i(a),x_i(b))\to d(a,b)$. In particular the radial
norms and their metric kernels converge.

For each pair $(a,t)\in D\times\mathbb Q_{\geq0}$, AC83 bounds
the source radius of $H^i_t x_i(a)$ by
$t(d(p,a)+2e_i)$. Thus its image under $f_i$ is eventually defined
and stays in a fixed compact ball of $X$. One countable diagonal
subsequence makes all these images converge; write the limits as
$G_t(a)$. In particular $G_0(a)=p$ and $G_1(a)=a$.
Passing the source cone identity through the distortion bound gives
\[
 d(G_sa,G_tb)^2=s^2r_a^2+t^2r_b^2-2stB_p(a,b)
 \quad(a,b\in D,\ s,t\in\mathbb Q_{\geq0}).
\]
All parameters and lifts in each use of this identity are fixed and
bounded; no assertion about convergence on an unbounded ball is used.

For each rational $t\geq0$, the identity with $s=t$ shows that
$G_t$ is $t$-Lipschitz on $D$. Extend it uniquely to $X$ by
completeness. Continuity extends the displayed identity to every
$x,y\in X$ for rational $s,t$. In particular
\[
 d(G_sx,G_tx)=|s-t|r_x.
\]
For real $t\geq0$, choose nonnegative rationals tending to $t$.
This last identity makes their images Cauchy and makes the limit
independent of that choice; define it as $H_tx$. Passing to limits
in both rational parameters gives the full AC82 identity, with
$H_0x=p$ and $H_1x=x$. AC83 then supplies the semigroup law and
surjectivity. These are actual cone maps on the specified $X$,
not independent identifications of rescaled spaces.
\end{proof}

\begin{corollary}[Compactness of the nonnegative cone models; AC86]
\label{cor:alexandrov-cone-model-compactness}
Fix an integer $n\geq1$. Every sequence of complete pointed length
spaces with global curvature at least zero, Hausdorff dimension
at most $n$, and AC82 cone data has a subsequence converging to a
complete proper geodesic nonnegatively curved space of dimension
at most $n$, carrying AC82 data at its limiting basepoint.
This is conditional on the local comparison input retained by AC41.
\end{corollary}
\begin{proof}
Apply AC41 with curvature allowance $c_i=1/(i+1)$ and radii
$r_i=i+1$. Global curvature zero can be weakened to the required
local lower bound $-c_i$, and the global dimension bound implies
the local one. AC41 supplies one complete proper geodesic limit
with the asserted curvature and dimension. AC85 supplies its cone
data. Dimension may drop and the limit may be a point.
\end{proof}

\subsection{Uniform splitting from overlapping cone approximations}

\begin{theorem}[Normalized overlapping-cone splitting; AC87]
\label{thm:alexandrov-overlapping-cone-splitting}
Conditional on AC41's retained local comparison and AC43,
for every integer $n\geq1$ and $0<\delta<1$ there is
$0<\varepsilon=\varepsilon(n,\delta)<1$ with the following
property. Let $Z$ be a complete length space and let $p,q\in Z$
satisfy $d_Z(p,q)=1$. Suppose there are actual pointed KL
$\varepsilon$-approximations
\[
 (Z,p)\longrightarrow(C,o),\qquad
 (Z,q)\longrightarrow(D,v),
\]
where $C,D$ are complete nonnegatively curved length spaces of
Hausdorff dimension at most $n$, with AC82 structures at $o,v$.
Then $(Z,q)$ has a normalized $(1,\delta)$-splitting.
No curvature or dimension hypothesis is imposed on $Z$.
The tolerance is uniform over $Z,p,q$, the two models and their
approximation maps; no explicit modulus is asserted.
\end{theorem}
\begin{proof}
If this fails for fixed $n,\delta$, choose counterexamples with
$\varepsilon_i=1/(i+2)$. Apply AC86 to the first cone models and
then to the second along a further common subsequence. Denote
the complete proper cone limits by $(X,p_\infty)$ and $(Y,v_\infty)$.

Composition gives $(Z_i,p_i)\to(X,p_\infty)$ and
$(Z_i,q_i)\to(Y,v_\infty)$. Here is the domain check: for each
fixed $R>0$ choose $S=2R+4$. MC11 converts the first KL maps,
for large $i$, to $\mathcal A_{S,3\varepsilon_i}$ maps. Choose
$\mathcal A_{S+1,\eta_i}$ convergence maps from the cone models
to their limits, with $\eta_i\to0$. MC08 composes and restricts
to radius $R$, with error $2(3\varepsilon_i+\eta_i)\to0$;
the required margins $R+3\varepsilon_i\leq S+1$ and
$2(3\varepsilon_i+\eta_i)<R$ eventually hold. The same reasoning
applies to the second models. Taking growing radii diagonally
retains actual maps for the first convergence.

Use AC81 and properness of $X$ to pass to a further subsequence
on which the images of $q_i$ converge to $q_\infty\in X$.
It gives $d_X(p_\infty,q_\infty)=1$ and
$(Z_i,q_i)\to(X,q_\infty)$ on this same underlying $X$.
MC24 identifies the two complete proper pointed limits by an
onto isometry $J:(Y,v_\infty)\to(X,q_\infty)$.
Conjugating the radial maps on $Y$ by $J$ gives a second AC82
structure on $X$, now at $q_\infty$.

AC84 supplies a line through the two distinct apices of $X$.
AC43 applies to its complete proper finite-dimensional
nonnegative geometry and gives an exact splitting aligned with
that line. Translate its real coordinate so that $q_\infty$
is the product basepoint. Compose this fixed pointed isometry
with actual growing-ball approximations from $(Z_i,q_i)$.
Extend the Euclidean coordinate by zero outside each approximation
domain; on every fixed ball this extension is eventually irrelevant.
AC52, with prescribed Euclidean coordinate equal to this very
composition and hence coordinate discrepancy zero, gives normalized
$(1,\delta)$-splittings for all sufficiently large $i$.
This contradicts the selected counterexamples.
\end{proof}

This is KL 4.20 expressed using actual approximation witnesses,
with the normalized range $0<\delta<1$ used throughout the blueprint.
KL's wording in terms of pointed GH closeness does not select an
identical numerical metric on the collection of pointed spaces;
MC06/MC11 supply the convention translation, not an asserted
identity of tolerance constants. Neither strainer implication nor
AC76 compatibility is used. The two cone structures, separation
one, proper-limit uniqueness and the transfer back to the original
spaces are all explicit. Local comparison and AC43 still require
their future formal bindings.

\section{Audit of Chapters 3--4 and dependency scope}
\label{sec:revision79-dependency-audit}

Revision 79 reviews all 24 metric and 87 Alexandrov contracts, their
written arguments and the declared dependencies. No additional false
statement or invalid proof step was identified in this manual review.
This is not kernel verification: source-qualified inputs and formal
bindings remain. \code{AUDIT_REVISION78.md} records the findings;
\code{reference_checks_revision79.md} records the targeted source and
errata checks. Other chapters were reviewed for interfaces and global
dependency boundaries, not subjected to a new line-by-line proof audit.

\subsection{The order is between results, not whole chapters}

The generic metric toolkit supplies approximation calculus, MC13
extraction, and MC21--MC24 length closure and identification.
Chapter 4 first obtains local charts and compactness without using
AC01/AC02 to prove their own inputs. AC04/AC08 then supply uniform
nets for MC13. AC37 supplies limit comparison, AC40 supplies the
dimension bound, and AC41 assembles these with MC14. Finally MC18
exports that result to Chapter 3's consumer interface. There is no
reverse dependency from MC13 or AC41 to MC18.

The coarse implementation worksheet previously included MC18 inside
MG06 while scheduling MG07 after MG06. That could be read as requiring
the completed Alexandrov application before the length-closure results
it consumes. MG06 now owns only generic extraction MC13; MG07 follows
the approximation calculus MG03/MG04 and need not wait for extraction.
MC18 remains the downstream application, with its sole direct theorem
dependency AC41. I02's ``after metric'' priority means after the required
generic leaves, not after every application printed in Chapter 3.
This repairs the schedule without changing a mathematical statement or
the detailed theorem graph. The worksheet remains a planning aid,
not an exact import graph or a reason to pin an environment now.

The 111-node MC/AC graph and the 136-row production graph are acyclic.
The latter still has 21 unresolved external interface aliases.
Consumer columns, chapter assignments and source citations are distinct
from theorem dependencies. A node with no listed predecessor may still
use ordinary real analysis, metric topology, Hausdorff dimension facts
or an explicitly cited theorem; it is not an assertion of no axioms.

\subsection{Inputs that the audit does not discharge}

\begin{itemize}[leftmargin=*]
\item MC01 retains the stated arclength and Hopf--Rinow bindings;
MC04's compact GH correspondence and separation proof retains its
source qualifications. The subsequent written arguments do not turn
these into checked declarations.
\item MC17 requires actual ambient-distance and coverage producers.
MC20's smooth precompactness retains its regularity, derivative bounds,
noncollapse and inherited-foundation obligations.
\item AC02/AC36 verify buffers for the cited local Toponogov theorem;
they do not prove that theorem. AC41 and MC18 retain this qualification.
The intrinsic/ambient dimension translation in the AC08 application
uses local compactness first and a countable local-isometry cover.
\item AC43 retains the source-qualified nonnegative splitting theorem;
AC46 retains local one-dimensional recognition. Their hypotheses must
be supplied before the product classification AC48 is applied.
The conditional angular/tangent route AC09--AC12 is not needed by
the selected AC41 compactness route.
\item AC56's product-limit closure does not use AC54/AC55. AC76's
metric compatibility does not use either strainer direction.
AC87 uses compactness, actual cone data, proper-limit identification
and AC43, without AC76 or either strainer direction.
\end{itemize}

These source inputs are mathematical dependencies, not new PC gap
classifications. All new formal binding remains deferred. Existing
finite K0 evidence verifies its own sandbox only. No Chapter 3 or 4
theorem is promoted to Lean-checked status by this audit.

\subsection{Boundary before Chapter 13}

AC41 gives $\dim_H X\leq n$, not the collapse-specific conclusion
$\dim_H X\leq2$ when $n=3$. Chapter 13 now supplies the
volume-scale and local positive-volume arguments for KL 6.18, retains
the author correction to 6.5, and excludes the three-stratum using
AC41 and AC51. Section~\ref{sec:collapse-cone-scales} adds cone scales
and outward annulus points with explicit smooth/Tits inputs.
Neither metric convergence nor approximate splitting alone provides
smooth adapted coordinates or compatible fibrations; Chapter 14 owns
those constructions. The production graph still separates raw flow
input from the long-time outcome, raw graph data from essential cuts,
and structural assembly from analytic producers. The late-thick/thin
or terminal-exceptional alternative, H1--H4 gates and PC release
boundary remain unchanged.

\section{Source scope and next implementation targets}
See \mbox{\code{reference_checks_revision78.md}} for KL 4.19--4.20,
the corrected BBI cone law, AKP's angular convention, author corrections
and the manual proof review. AC82--AC85 are metric; AC86 uses the
existing AC41 comparison qualification; AC87 also uses AC43 splitting.
These are written mathematical arguments, not Lean verification.

Chapter 13 now expands KL 6.18 through LC01--LC09, with corrected
scale selection, volume estimates and a quantified model argument.
Its local full-dimensional volume input LC06 is supplied by LC15 in
Section~\ref{sec:collapse-local-volume-proof}, with inherited bindings retained;
see Section~\ref{sec:collapse-volume-models}. The remaining good-annulus topology
in KL Section 11 is subsequent work. Chapter 14 retains smooth
adapted coordinates and fibrations. Mathlib pinning and PC-dependent
probes remain deferred.

\section{A written global line-splitting proof}
\label{sec:alexandrov-written-line-splitting}

The following subordinate arguments ALS01--ALS05 discharge AC43's
splitting producer from its GLOBAL nonnegative four-point hypothesis.
They are metric arguments on the actual space. They neither produce
global comparison from local comparison nor smooth any metric.
The finite Riemannian models have this global input by LFR55; the
collapsed limits have it by AC37, whose upstream local-comparison
qualification remains explicit.

\paragraph{Sources and proof decomposition.}
BBI10.5.1--10.5.6 and the full three-step proof, printed366--369/
PDF381--384, supply the selected parallel-line route. The retained
July6,2024 author errata, PDF13, require PARALLEL plane lines in
Definition10.5.3, correct the distance chain and its constant on367,
and use the line through $x$, not through $q$, in the slice map on369.
The author errata endpoint was retried but returned502; its current
bytes were not verified. We use the retained corrected source with
that limitation. BBI10.1.1, printed353/PDF368, supplies the four-point
comparison convention; its forward implication is the algebra already
expanded at the end of AC37. The errata's correction to the reverse
implication on353 is recorded but not needed by the argument below.

We expand the source's flat-triangle and parallel-line steps through
quadratic squared distances. A signed Busemann coordinate is then
obtained as a distance limit, and all its needed properties are
proved here. This is not an invocation of a Busemann gradient flow
or the alternative AKP splitting proof. It avoids presuming convergence
of actual angles when finite geodesic segments converge to rays.
The pairwise distance formula below also supplies the source's
compressed uniqueness, projection-transitivity and final onto-product
steps. Exact reading scopes and source hashes are recorded in
\code{reference_checks_revision138.md}.

\begin{lemma}[Global four-point consequences; ALS01]
\label{lem:alexandrov-global-four-point-side}
Let $X$ be a geodesic metric space satisfying the Euclidean four-point
inequality at every center, with the repeated-outer-point convention
of AC34--AC37. If $z=\sigma(u)$, $0\leq u\leq1$, on a constant-speed
minimizing segment from $a$ to $b$, then for every $v\in X$,
\[
 d(v,z)^2\geq(1-u)d(v,a)^2+u\,d(v,b)^2-u(1-u)d(a,b)^2.
\]
For two chosen minimizing segments from a common center, their
Euclidean comparison angle is nonincreasing in each positive arm
length. This uses the GLOBAL four-point condition directly, without
a local-to-global theorem or preexisting angles between geodesics.
\end{lemma}
\begin{proof}
This is the algebraic last step of AC37, independent of its
compactness hypotheses. Set $L=d(a,b)$ and suppose first that
$L>0$, $0<u<1$, and $v\notin\{a,b,z\}$. At $z$ the comparison
angle between $a$ and $b$ is $\pi$. Thus the sum of the other two
comparison angles is at most $\pi$, and the sum of their cosines
is nonnegative. With $h=d(v,z)$, $A=d(v,a)$ and $B=d(v,b)$,
the cosine laws give
\[
 \frac{u^2L^2+h^2-A^2}{2uLh}
 +\frac{(1-u)^2L^2+h^2-B^2}{2(1-u)Lh}\geq0.
\]
Multiplication by the positive denominator gives the claimed
inequality. The excluded cases give equality directly.

For the angle statement, let one endpoint $x$ have distance $A>0$
from the center $p$, and let the other endpoint $y$ have distance
$B>0$, on its chosen minimizing segment. At the point $y_u$ of
that segment at distance $uB$, the first inequality gives
\[
 d(x,y_u)^2\geq (1-u)A^2+u\,d(x,y)^2-u(1-u)B^2.
\]
Consequently
\[
 \frac{A^2+u^2B^2-d(x,y_u)^2}{2AuB}
 \leq\frac{A^2+B^2-d(x,y)^2}{2AB}.
\]
Since $\arccos$ is decreasing, shortening the second arm can only
increase its comparison angle. Interchanging the arms proves the
other assertion. Degenerate triangles are allowed by continuity of
$\arccos$ on its closed domain. No endpoint is identified with the
center in these positive-arm formulas.
\end{proof}

\begin{lemma}[Exact squared distance to any line; ALS02]
\label{lem:alexandrov-line-distance-quadratic}
In the space of ALS01, let $\gamma:\mathbb R\to X$ be a unit-speed
line. For each $x\in X$ there are uniquely determined numbers
$b(x)\in\mathbb R$ and $h(x)\geq0$ such that
\[
 d(x,\gamma(t))^2=(t-b(x))^2+h(x)^2\qquad(t\in\mathbb R).
\]
The point $\gamma(b(x))$ is the unique closest point of the line
to $x$. This assertion holds for ANY chosen line in $X$.
\end{lemma}
\begin{proof}
If $x=\gamma(0)$ the formula is immediate. Otherwise put
$a=d(x,\gamma(0))>0$. Let $\alpha_T$ and $\beta_S$ be the
comparison angles at $\gamma(0)$ between $x$ and $\gamma(T)$,
and between $x$ and $\gamma(-S)$, for $T,S>0$. ALS01 implies
that both decrease as their respective parameters increase, so they
have limits $\alpha_\infty,\beta_\infty\in[0,\pi]$.
Four-point comparison at $\gamma(0)$ gives
\[
 \alpha_T+\beta_S\leq\pi
 \qquad\text{for EVERY independent pair }T,S>0.
\]
If an outer point repeats, its prescribed automatic inequality
still gives this statement.

The triangle inequality gives
$d(x,\gamma(T))+d(x,\gamma(-T))\geq2T$.
Rationalizing the cosine laws yields
\[
 \begin{split}
 d(x,\gamma(T))-T
 &=\frac{a^2-2Ta\cos\alpha_T}{d(x,\gamma(T))+T}
   \longrightarrow-a\cos\alpha_\infty,\\
 d(x,\gamma(-T))-T&\longrightarrow-a\cos\beta_\infty.
 \end{split}
\]
Indeed each denominator divided by $T$ tends to two by
$T-a\leq d(x,\gamma(\pm T))\leq T+a$.
It follows that $\cos\alpha_\infty+\cos\beta_\infty\leq0$.
Two angles in $[0,\pi]$ with sum strictly less than $\pi$ have a
strictly positive cosine sum; hence their limiting sum is at least
$\pi$. It is also at most $\pi$ by the four-point inequality.
Thus their sum is exactly $\pi$. Monotonicity and the inequality
for every independent $T,S$ now force
$\alpha_T=\alpha_\infty$ and $\beta_S=\beta_\infty$ for ALL
positive parameters. There is no interchange of an unproved angle
limit with a limit of geodesics.

Set $b(x)=a\cos\alpha_\infty$ and
$h(x)^2=a^2-b(x)^2\geq0$. The cosine laws on the two half-lines,
using $\cos\beta_\infty=-\cos\alpha_\infty$, give the asserted
formula, also at $t=0$. The coefficients of this polynomial determine
$b$ and $h^2$ uniquely. Its unique minimum occurs at $t=b(x)$,
which is an actual point of the given line. No projection or flat
strip theorem has been assumed.
\end{proof}

\begin{lemma}[An affine coordinate and lines through every point; ALS03]
\label{lem:alexandrov-affine-line-coordinate}
Let $X$ additionally be complete and proper, and fix the line
$\gamma$ of ALS02. Its coordinate $b$ is $1$-Lipschitz,
$b(\gamma(t))=t$, and is affine on EVERY minimizing segment.
For each $x\in X$ there is a unit-speed line
$\ell_x:\mathbb R\to X$ with
\[
 \ell_x(0)=x,\qquad b(\ell_x(s))=b(x)+s\quad(s\in\mathbb R).
\]
\end{lemma}
\begin{proof}
Rationalizing ALS02's formula shows
\[
 b(x)=\lim_{T\to\infty}\bigl(T-d(x,\gamma(T))\bigr).
\]
The functions in this limit are $1$-Lipschitz, so $b$ is
$1$-Lipschitz. On the original line the limit is its unit-speed
parameter. If $z$ is the point at fraction $u$ of a minimizing
segment from $x$ to $y$, apply ALS01 with vertex $\gamma(T)$.
ALS02 expresses the three squared distances as
$T^2-2Tb(w)+c(w)$, where $c(w)=d(w,\gamma(0))^2$.
Cancellation gives
\[
 -2T\bigl(b(z)-(1-u)b(x)-ub(y)\bigr)+C\geq0
 \qquad(T\in\mathbb R),
\]
where $C$ is independent of $T$. Since BOTH signs and arbitrarily
large absolute values of $T$ are allowed, its linear coefficient
vanishes. This proves affineness on every minimizing choice.

Choose $T_i\to\infty$ and minimizing unit-speed segments from $x$
to $\gamma(T_i)$, of lengths $L_i$. ALS02 gives $L_i\to\infty$
and $(T_i-b(x))/L_i\to1$. On each fixed parameter interval these
segments lie in one compact ball about $x$. The inherited metric
curve compactness and a diagonal subsequence give a minimizing ray
$r_x^+:[0,\infty)\to X$: distances between its fixed parameter
values are limits of the same unit-speed distances on the segments.
Affineness on each approximating segment gives, for $0\leq s\leq L_i$,
\[
 b(\sigma_i(s))=b(x)+s\frac{T_i-b(x)}{L_i}.
\]
Continuity of $b$ shows $b(r_x^+(s))=b(x)+s$. The identical
construction towards $\gamma(-T_i)$ yields a ray $r_x^-$ with
$b(r_x^-(s))=b(x)-s$. Separate subsequences suffice; their common
point is exactly $x$.

Join these rays by $\ell_x(s)=r_x^+(s)$ for $s\geq0$ and
$\ell_x(s)=r_x^-(-s)$ for $s\leq0$. On each half this is already
minimizing. For points on opposite halves the $1$-Lipschitz function
$b$ gives the lower distance bound $|s-t|$, while the path through
$x$ gives the upper bound $|s|+|t|=|s-t|$. Hence the joined curve
is a line with the stated orientation and parameter. Only compact
metric curve extraction was used; no convergence of actual angles
or gradient-flow construction is required.
\end{proof}

\begin{lemma}[Unique parallel lines and the full cross-distance law; ALS04]
\label{lem:alexandrov-parallel-line-distance-law}
The lines of ALS03 are unique among unit-speed lines through $x$
on which $b$ increases with slope one. For all $x,y\in X$ and
$s,t\in\mathbb R$ they satisfy
\[
 d(\ell_x(s),\ell_y(t))^2
 =\bigl(b(x)+s-b(y)-t\bigr)^2+d(x,y)^2-(b(x)-b(y))^2.
\]
Thus translation $F_t(x)=\ell_x(t)$ is an isometry,
$F_0=\mathrm{id}$, $F_sF_t=F_{s+t}$, and $F_{-t}$ is its inverse.
\end{lemma}
\begin{proof}
For ANY oriented line $\ell_y$ with $b(\ell_y(t))=b(y)+t$, ALS02
applied to this line and an arbitrary point $x$ writes
\[
 d(x,\ell_y(t))^2=t^2-2ct+d(x,y)^2
\]
for a coefficient $c$. Since $b$ is $1$-Lipschitz, the right side
is at least $(t-(b(x)-b(y)))^2$ for every real $t$. Subtracting
these quadratics gives an affine function nonnegative on ALL of
$\mathbb R$, so its linear coefficient must vanish. Therefore
\[
 d(x,\ell_y(t))^2
 =t^2-2t(b(x)-b(y))+d(x,y)^2.                 \tag{ALS04a}
\]
Suppose two such lines through $x$ are $\ell$ and $\ell'$. Apply
this formula to the point $\ell(s)$ and line $\ell'$ at parameter
$t=s$. Since $b(\ell(s))-b(x)=s$ and $d(\ell(s),x)=|s|$, the
resulting squared distance is zero. Thus the lines agree at every
parameter, with no angle-zero or nonbranching theorem imported.
The reparametrized line $u\mapsto\ell_x(t+u)$ has the same slope
and passes through $F_t(x)$, so uniqueness gives the group law.

Apply (ALS04a) to $\ell_x(s)$ and $\ell_y$, and then apply it to
$y$ and $\ell_x$. Substitution yields
\[
 \begin{split}
 d(\ell_x(s),\ell_y(t))^2
 &=t^2-2t(b(x)+s-b(y))
   +s^2-2s(b(y)-b(x))+d(x,y)^2,
 \end{split}
\]
which is the displayed cross-distance identity after completing the
square. At $s=t$ it gives $d(F_t x,F_t y)=d(x,y)$; the group law
supplies surjectivity and the inverse. In particular the squared
transverse separation $d(x,y)^2-(b(x)-b(y))^2$ is nonnegative.
The formula identifies any two members, after matching the $b$
coordinate, with parallel Euclidean lines at that separation (with
the coincident case allowed). This is the corrected BBI parallel-line
convention, including the entire union's restricted metric.
\end{proof}

\begin{theorem}[The actual aligned onto splitting; ALS05]
\label{thm:alexandrov-written-aligned-splitting}
Let $X$ be nonempty, complete, proper and geodesic, with global
nonnegative four-point comparison, and let $\gamma$ be a specified
unit-speed line. Put $p=\gamma(0)$. There is an actual onto pointed
isometry
\[
 \alpha:(X,p)\longrightarrow
 (\mathbb R\times Y,(0,p)),\qquad
 \alpha(\gamma(t))=(t,p),
\]
where $Y$ is a complete proper geodesic space with the same global
four-point comparison. If $\dim_H X$ is finite, then
$\dim_HY\leq\dim_HX$. This is AC43's required aligned splitting.
\end{theorem}
\begin{proof}
Use ALS02--ALS04 and let $Y=b^{-1}(0)$ with its restricted metric.
It is nonempty because it contains $p$, closed because $b$ is
continuous, and convex along EVERY minimizing segment because $b$
is affine. Thus closed-subspace completeness/properness and ambient
minimizing segments give completeness, properness and geodesicity
for $Y$. Its four-point inequalities use exactly the same distances
as in $X$. The Hausdorff dimension bound follows from its isometric
inclusion; no product-dimension equality is assumed.

Define the actual maps
\[
 \alpha(x)=\bigl(b(x),F_{-b(x)}x\bigr),\qquad
 \Psi(t,y)=F_t y.
\]
The second component of $\alpha$ belongs to $Y$ by the slope formula.
The group law gives $\Psi\alpha=\mathrm{id}_X$ and
$\alpha\Psi=\mathrm{id}_{\mathbb R\times Y}$. For $y,z\in Y$,
ALS04 is precisely
\[
 d(\Psi(t,y),\Psi(s,z))^2=(t-s)^2+d_Y(y,z)^2.
\]
Thus these are inverse onto isometries for the Pythagorean product,
not merely distance-preserving maps into a product. The original
$\gamma$ has $b(\gamma(t))=t$ and so is the unique slope-one line
through $p$. Hence $F_t(p)=\gamma(t)$, proving alignment and exact
basepoint preservation. The slice transport always follows the line
through the point being transported, as required by the correction
to BBI369. The same formula shows projection between any two
parallel lines preserves the $b$ coordinate, so the projection
relation used in BBI's proof is symmetric and transitive as well.

All comparison used here was the supplied GLOBAL four-point
inequality and ALS01's direct algebra. Properness was used to
extract the two rays. No local comparison radius, local-to-global
producer, general tangent-space theory, finite-dimensional
stratification, soul theorem or metric regularization is invoked.
Finite Hausdorff dimension is retained to match AC43 and its
consumers, although it is not needed for the metric product construction.
\end{proof}

\paragraph{Binding the existing consumers.}
ALS05 now supplies AC43 in AC48, AC53, AC60, AC73 and AC87 once
their written hypotheses give a complete proper geodesic globally
nonnegative space and an actual line. Factor applications use AC44;
finite-model applications use LFR55. AC37 supplies global four-point
comparison for the extracted collapsed limit before splitting is
applied, so ALS01 does not call that compactness theorem as a
producer: it reuses only its standalone algebraic last paragraph.
This prevents a dependency cycle through compactness, strainers or
local packets. No ray is substituted for a line, and no local
curvature bound is substituted for global comparison.

The historical source-qualified AC43 statement and audits are
preserved; their splitting-producer qualification is superseded by
ALS01--ALS05. Their accepted-foundation and Lean bindings are still
pending. AC02/AC36's local-to-global comparison and AC46's local
one-dimensional recognition remain independent mathematical inputs.
The final simultaneous local parameter/export binding is also still
required. This completes the written splitting argument, not the
full local-collapse goal or a Lean formalization.


\section{One-dimensional recognition from the written paired charts}
\label{sec:alexandrov-written-one-dimensional-recognition}

This section supplies AC46 for its stated nonnegative-curvature
consumers. The metric space has GLOBAL zero-curvature four-point
comparison. No local-to-global theorem is inferred here.
The existing finite paired-chart proof supplies injectivity at each
interior point of a minimizing segment; its actual distance coordinate
then gives an ISOMETRIC interval, not merely a topological chart.
An exit-point argument treats the remaining endpoints.

\paragraph{Source comparison.}
AKP Theorem15.18, archived printed/PDF238, separates interior points
and endpoints in this way \cite{AlexanderKapovitchPetruninFoundations}.
Its interior argument invokes Proposition15.7 (230), which uses
Corollary13.40 and tangent-space packing; Corollary15.16 (236--237)
also translates its linear dimension to Hausdorff dimension.
Those additional structural theorems are NOT imported here.
Instead, AC23, AC15 and AC27 supply the needed rank exclusion and
injectivity directly in the given Hausdorff convention. Their
BGP/BBI source comparisons and written proofs are retained unchanged;
see Sections~\ref{sec:alexandrov-approximate-corrections} and
\ref{sec:alexandrov-paired-charts}. The simple curvature weakening
needed to reuse their minus-one normalization is proved below.
The pinned AKP \code{model.tex} cosine laws and \code{dim.tex},
label \code{thm:dim=1.CBB}, remain distinct from the archived PDF
and the published edition. The author errata body dated July12,2026
was checked externally; its published-page238 correction concerns
SectionJ, not this archived-page238 theorem. Exact identities,
locators and limits are in \code{reference_checks_revision139.md}.

\begin{lemma}[Curvature weakening and exclusion of rank two; ALR01]
\label{lem:alexandrov-recognition-rank-ceiling}
For any metric triangle with positive adjacent sides, its
minus-one comparison angle is at most its zero-curvature comparison
angle, including degenerate triangles. Consequently, if a locally
complete length space $Y$ has GLOBAL zero-curvature four-point
comparison and $\dim_HY\leq1$, there is no rank-two paired packet
of quality $\beta=1/1000$ anywhere in $Y$, with anchors in $Y$.
The paired-packet angles here are the minus-one angles of AC23.
\end{lemma}
\begin{proof}
Write the adjacent sides as $a,b>0$, the opposite side as $c$, and
put $v=(a^2+b^2-c^2)/(2ab)\in[-1,1]$. The function
\[
 H(u)=\cosh\sqrt u=\sum_{j=0}^{\infty}\frac{u^j}{(2j)!}
 \qquad(u\geq0)
\]
is convex: each summand is convex and the finite convexity inequalities
pass to the convergent sum. The number $c^2$ is the convex combination
of $(a-b)^2$ and $(a+b)^2$ with weights $(1+v)/2$ and $(1-v)/2$.
The chord bound therefore gives
\[
 \cosh c\leq\frac{1+v}{2}\cosh(a-b)
                  +\frac{1-v}{2}\cosh(a+b)
             =\cosh a\cosh b-v\sinh a\sinh b.
\]
Divide by the positive $\sinh a\sinh b$ and use both cosine laws:
$\cos\widetilde\angle_{-1}\geq v
=\cos\widetilde\angle_0$. Monotonicity of cosine on $[0,\pi]$
proves the assertion, including $v=\pm1$.
Summing the three angle inequalities gives the global minus-one
four-point condition. Thus $\Omega=Y$ is a common comparison
neighborhood for the existing paired-chart arguments.

If a rank-two quality-$\beta$ packet existed at $q$, apply AC23
with $m=2$ and quality parameter $2\beta\leq1/200$. Its initial
quality is exactly $\beta$. The resulting Lipschitz distance map
has a nonempty open image in $\mathbb R^2$. AC15 then gives
$2\leq\dim_HY\leq1$, a contradiction. This uses the written
complete-buffer correction proof, not AKP's linear-dimension theorem.
\end{proof}

\begin{lemma}[Exact small intervals at every segment interior; ALR02]
\label{lem:alexandrov-recognition-interior-chart}
Let $Y$ have the hypotheses of ALR01 and be geodesic. Suppose
$\sigma:[-A,B]\to Y$ is an isometric minimizing segment,
$A,B>0$, and $p=\sigma(0)$. There is $\rho>0$ such that
\[
 B_Y(p,\rho)=\sigma((-\rho,\rho)),\qquad \rho<\min\{A,B\}.
\]
The parametrization is an onto pointed isometry for the RESTRICTED
metric on this ambient ball.
\end{lemma}
\begin{proof}
Take the anchors $a_1=\sigma(-A)$ and $b_1=\sigma(B)$.
Their minus-one comparison angle at $p$ equals $\pi$, because the
triangle is degenerate with opposite length $A+B$.
Set $\beta=1/1000$ and $\delta=\beta/100$. Continuity supplies
an open $V\ni p$ on which this rank-one packet has quality
$\delta$ and positive finite anchor-distance bounds $[a_*,A_*]$.
All anchors lie in the common comparison neighborhood $\Omega=Y$.
ALR01 excludes rank-two quality-$\beta$ packets centered in $V$.
Choose $\rho>0$ with $B(p,3\rho)\subset V$,
$2\rho\leq\ell_\beta$ of AC26, and
$\rho<\min\{A,B\}$. The inequalities
$\beta\leq1/400$ and $\delta\leq\beta/100$ allow AC27 with $m=1$.
In particular $F=d_Y(a_1,\cdot)$ is injective on $B(p,\rho)$.

For any $z$ in that ball put $t=F(z)-F(p)$. The distance function
is $1$-Lipschitz, so $|t|\leq d(p,z)<\rho$. The point $\sigma(t)$
belongs to the same ball and, since the ORIGINAL segment is
isometric, $F(\sigma(t))=A+t=F(z)$. Injectivity forces
$z=\sigma(t)$. The reverse inclusion follows from
$d(p,\sigma(t))=|t|$. Restricting the original isometric segment
proves the exact metric conclusion; no length-space structure on
the restricted ball or uniform lower bound on $\rho$ is assumed.
\end{proof}

\begin{lemma}[No exit before a segment endpoint; ALR03]
\label{lem:alexandrov-recognition-no-exit}
Under ALR02's hypotheses, for EVERY $0<r<\min\{A,B\}$,
\[
 B_Y(p,r)=\sigma((-r,r)).
\]
In particular the entire interior of any minimizing segment is
open in $Y$.
\end{lemma}
\begin{proof}
ALR02 applied at each interior point already shows that the
segment interior is open in $Y$. Suppose $z\in B(p,r)$ is outside
the segment. Choose a minimizing segment $\eta$ from $p$ to $z$,
parametrized by arclength on $[0,d(p,z)]$. Its intersection with
$\sigma([-A,B])$ has a nonempty compact parameter set: the latter
segment is a compact, hence closed, subset of the metric space.
Let $t_*$ be the largest such parameter and $q=\eta(t_*)$.
Since $z$ is outside, $t_*<d(p,z)$. Also
$d(p,q)=t_*<r<\min\{A,B\}$, so $q$ is an INTERIOR point of
the original segment. Its open neighborhood from ALR02 contains
$\eta(t)$ for $t>t_*$ sufficiently close to $t_*$. Those points
are still on the original segment, contradicting maximality.
Thus every such $z$ is on the segment, and its unique original
parameter has absolute value $d(p,z)<r$. The other inclusion
and all pairwise distances follow from the original isometry.
No assertion about distinct tangent directions or nonbranching
geodesics has been inserted into the exit argument.
\end{proof}

\begin{lemma}[Exact half intervals at the remaining points; ALR04]
\label{lem:alexandrov-recognition-endpoint-chart}
Let $Y$ have the hypotheses of ALR02 and let $p\in Y$ lie in
the interior of NO minimizing segment with distinct endpoints.
If $y\ne p$ and $\sigma:[0,D]\to Y$ is any minimizing segment
from $p$ to $y$, then for every $0<r<D$,
\[
 B_Y(p,r)=\sigma([0,r)).
\]
This is an onto pointed isometry for the restricted ambient metric.
\end{lemma}
\begin{proof}
Suppose $z\in B(p,r)\setminus\sigma([0,D])$. Choose
$0<a<(D-r)/3$ and put $w=\sigma(a)$; then $0<a<D$.
A minimizing segment $\eta$ from $w$ to $z$ has length
\[
 d(w,z)\leq a+d(p,z)<a+r<D-a=d(w,y).
\]
Thus it cannot pass through $y$. It cannot pass through $p$
either: both its endpoints differ from $p$, so such a passage
would put $p$ in the interior of a minimizing segment, contrary
to the hypothesis. The parameter set where $\eta$ meets the
original segment is compact and nonempty, since it contains $w$.
Take its last point $q$. It is different from $p,y$ by the two
exclusions just proved and different from $z$ by the choice of $z$.
ALR02 therefore supplies an open neighborhood of $q$ contained
in the interior of the original segment. Continuity of $\eta$
contradicts the last-intersection property exactly as in ALR03.
Hence the whole ball lies on the original segment. The original
distance parameter from $p$ proves the equality and the onto
isometry. This argument retains the endpoint; it does not extend
the segment through $p$ or assume a boundary theorem.
\end{proof}

\begin{theorem}[Written one-dimensional recognition and consumer binding; ALR05]
\label{thm:alexandrov-written-one-dimensional-recognition}
AC46's local-recognition conclusion holds under its existing
hypotheses: a nonempty complete proper geodesic space $Y$ with
GLOBAL nonnegative four-point comparison and $\dim_HY\leq1$
is a singleton or has the exact pointed open-interval/half-interval
charts stated there. Consequently AC47's written classification
and its existing local-collapse consumers no longer require a
separate source-qualified recognition premise.
\end{theorem}
\begin{proof}
If $Y$ is a singleton there is nothing to prove. Otherwise choose
any $p\in Y$. If it is interior to some minimizing segment,
ALR02 gives the pointed chart $((-r,r),0)$ for some $r>0$.
If it is interior to none, choose $y\ne p$ and a minimizing
segment from $p$ to $y$. ALR04 gives $([0,r),0)$ for, say,
$r=d(p,y)/2$. Completeness provides the local completeness used
by AC23; geodesicity supplies the chosen segments. Properness
is retained from AC46's contract, although the segment-intersection
arguments themselves only need compactness of each segment.
The Hausdorff bound is used directly in ALR01, with no translation
to linear or topological dimension. Every chart uses the ORIGINAL
ambient distance and its ACTUAL ball.

Insert this conclusion into the existing proof of AC47. Its
line, ray, compact interval and length-circle parameters, including
arbitrary pointed positions on a ray or interval, are unchanged.
AC48's one-dimensional factors and LFR43--LFR44's endpoint and
edge-coverage arguments therefore consume the written ALR05 in
place of AC46's previously supplied local-recognition assertion.
Those consumers retain their separate comparison, product, scale
and parameter hypotheses. No extra rank-two splitting or uniform
recognition radius is claimed.
\end{proof}

The dependency order is AC19--AC23/AC15, then AC24--AC27, then
ALR01--ALR05 and the existing AC47/AC48 and local-packet consumers.
There is no dependency on AC43, a regular-tangent theorem,
AKP15.7/15.16, gradient flow, or AC02/AC36's unproved globalization.
The global four-point input is still explicit. General local-to-global
comparison, final simultaneous local parameter/export binding, and
accepted foundation/PC/Lean bindings remain pending. This revision
is written mathematics and document verification, not a Lean proof
or the final full-goal audit.

\section{Written globalization with an explicit interior buffer}
\label{sec:alexandrov-written-buffered-globalization}

The remaining comparison input must distinguish a complete space from
an incomplete open region. BBI Theorem 10.3.1, printed 360--363/PDF
375--378, proves the former; its first reduction can move a test
triangle through a region much larger than that triangle. KL
Section 3.3, printed 23/PDF 18, states the sharper interior criterion
using complete closed $2D$-balls for triangles with sides less than
$D$. That precise criterion is still source-qualified here.
Instead the following proof gives an explicit, coarser buffer and
checks its actual consumers. This suffices for growing-region
compactness without strengthening its hypotheses.

The proof follows the cat's-cradle and almost-minimum arguments of
AKP, archived Theorem 8.31 and Lemmas 8.35--8.36, printed/PDF
94,96--100. The retained \code{vol1} source commit and exact labels
are recorded in \code{reference_checks_revision140.md}.
The local angle argument, quantitative nondegeneracy in the cradle,
and complete-ball localization are written below. Only curvature
$k\leq0$ is needed. No positive-curvature globalization, ultrapower,
nonbranching theorem or smooth approximation is used.

Write $\theta^k_x(a,b)$ for the comparison angle determined by the
three distances, and $S_k(a,b,\theta)$ for the opposite model side
with arms $a,b$ and angle $\theta\in[0,\pi]$.
For positive arms these are inverse monotone functions of the
opposite side and angle. Both zero and negative curvature have
continuous model formulas, including degenerate triangles;
$S_k(a,b,\pi)=a+b$ and $S_k(a,b,0)=|a-b|$.
A local four-point condition means that each point has an open
neighborhood on which the curvature-$k$ four-point inequality holds
for the RESTRICTED metric of the space under discussion.
Repeated outer points are handled by angle zero and the bound
$\theta\leq\pi$, as in AC36; distinctness is not a hidden condition.

\begin{lemma}[Model gluing and local germ angles; ALG01]
\label{lem:alexandrov-local-germ-angles}
Let a metric space satisfy the local curvature-$k$ four-point
condition, with $k\leq0$. For any two specified minimizing geodesic
germs starting at $x$, the comparison angles have a joint limit
$\alpha_x\in[0,\pi]$ as the two positive arm lengths independently
tend to zero. Shortening a germ does not change this limit. Three
germs have pairwise limits whose sum is at most $2\pi$; if two are
opposite parts of one minimizing segment, the two adjacent limits
have sum at most $\pi$. For hinges contained in one sufficiently
small common comparison neighborhood, $\alpha_x$ is at least the
comparison angle of their endpoints. These small-hinge radii are
bounded below by a positive number near each point.
\end{lemma}
\begin{proof}
First let $z$ be interior to a minimizing segment $[a,b]$, put
$l=d(a,z)$, $m=d(z,b)$, $h=d(v,z)$, $u=d(v,a)$ and $w=d(v,b)$,
and suppose
\[
 \theta^k_z(a,v)+\theta^k_z(v,b)\leq\pi.
\]
For $h>0$ the cosine of the first angle plus the cosine of the
second is nonnegative. The cosine laws give, respectively,
\[
 h^2\geq\frac{mu^2+lw^2}{l+m}-lm\quad(k=0),
\]
and, for $k=-s^2<0$,
\[
 \sinh(s(l+m))\cosh(sh)
 \geq\sinh(sm)\cosh(su)+\sinh(sl)\cosh(sw).
\]
The right sides are the corresponding interpolation formulas for the
model point on the side of length $l+m$. Thus $h$ is at least its
model distance to $v$. The cosine law at $a$ now gives
$\theta^k_a(v,z)\geq\theta^k_a(v,b)$: shortening one radial arm
cannot decrease its comparison angle. The case $v=z$ is direct,
since then $u=l,w=m$ and the original triangle is degenerate.
Coincident endpoints and zero shortened arms are treated by their
limiting formulas, not division by a zero length.

If $a,b,z,v$ lie in one common four-point neighborhood, its
inequality at $z$, together with $\theta^k_z(a,b)=\pi$, gives the
assumed adjacent-angle bound. Apply the resulting monotonicity to
each arm of the two germs inside a small common neighborhood of
$x$. The bounded angle function is monotone in each variable, so
its supremum is its joint independent-parameter limit as both
lengths decrease to zero. Taking common sufficiently small positive
lengths in the local four-point inequality gives the three-germ
sum bound. Opposite germs have comparison angle exactly $\pi$,
which gives the adjacent bound. This proves the required properties
without constructing a space of directions or asserting equality
of adjacent angles.

Finally choose a ball $B(x,4a)$ inside a comparison neighborhood.
If an endpoint lies in $B(x,a)$ and a hinge has sum of arms less
than $a$, its center, other endpoint and both specified geodesics
lie in $B(x,3a)$. The preceding local argument applies throughout
this common neighborhood. This supplies a positive lower bound
for endpoint-based small-hinge radii near $x$.
\end{proof}

\begin{lemma}[A quantitative cat's-cradle enlargement; ALG02]
\label{lem:alexandrov-cradle-enlargement}
Fix $\ell>0$ and endpoints $p,q$ in a locally curvature-$k$ space,
$k\leq0$. Suppose every hinge with one endpoint $p$ or $q$ and sum
of arms less than $2\ell/3$ satisfies $\alpha\geq\theta^k$.
Suppose also that the joins from each center $x$ with
$d(x,p)+d(x,q)<\ell$ to $p,q$ are available as minimizing segments
in a common region where the local four-point condition holds.
Then every specified hinge $(x;p,q)$ with sum of arms less than
$\ell$ satisfies $\alpha_x\geq\theta^k_x(p,q)$.
All new centers obey the same sum-of-arms bound. If the required
joins exist in $B(p,2\ell)$, the proof stays in that ball.
\end{lemma}
\begin{proof}
Put $S=S_k(d(x,p),d(x,q),\alpha_x)$; the desired assertion is
$S\geq d(p,q)$. Coincident endpoints or a zero arm give it directly.
A sum of arms below $2\ell/3$ is covered by the hypothesis.
Otherwise write $a$ for the shorter arm, $b$ for the longer,
and $u,v$ for their respective endpoints. Set
\[
 h=\frac{2\ell/3-a}{3},\qquad
 x'\in[x,v],\qquad d(x,x')=h.
\]
Since $2\ell/3\leq a+b<\ell$, one has
\[
 \ell/18<h\leq2\ell/9<b,\qquad
 a+h<5\ell/9,\qquad a+2h<11\ell/18<2\ell/3.
\]
Choose $[x',u]$. Both short hinges $(x;u,x')$ and $(x';u,x)$
have sums of arms less than $2\ell/3$. If $x'=u$, then $u$
lies on the original segment to $v$, and
$S\geq|b-a|=d(u,v)$; stop. Otherwise form the model triangle
$(\widetilde x,\widetilde{x'},\widetilde u)$ and extend its side
from $\widetilde x$ through $\widetilde{x'}$ to
$\widetilde v$ at distance $b$ from $\widetilde x$.
The first short-hinge comparison gives
\[
 S\geq d(\widetilde u,\widetilde v).
\]
At $x'$, ALG01 gives that the new angle toward $u,v$ is at most
$\pi$ minus the angle toward $u,x$. The second short-hinge
comparison bounds this above by $\pi$ minus its model angle,
which is the model angle toward $\widetilde u,\widetilde v$.
Consequently
\[
 d(\widetilde u,\widetilde v)\geq
 S':=S_k(d(x',u),d(x',v),\alpha_{x'}).
\]
These inequalities hold also for collinear model triangles by the
same cosine formulas. The triangle inequality gives
$d(x',p)+d(x',q)\leq a+b$.

Iterate, stopping if the small-hinge or endpoint case occurs.
Then the sequence of model opposite sides is nonincreasing, and
the stopping case proves the original assertion. In an infinite
iteration let $r_n=a_n+b_n$ and $\delta_n=r_n-r_{n+1}$.
The $r_n$ decrease and remain at least $2\ell/3$, so
$\delta_n\to0$. With $c_n=d(x_{n+1},u_n)$,
\[
 c_n=a_n+h_n-\delta_n,\qquad
 b_n-h_n\geq\ell/9.
\]
For all late $n$, $\delta_n<\ell/36$, hence
$c_n\geq\ell/36$. Thus BOTH arms at every sufficiently late
center are at least $\ell/36$; also $h_n>\ell/18$.
All these lengths are bounded above by $\ell$.
The model angle at $x_n$ of the triangle with adjacent sides
$a_n,h_n$ and opposite side $a_n+h_n-\delta_n$ therefore tends
to $\pi$, uniformly on this compact range of positive arms.
The actual germ limit $\alpha_{x_n}$ is at least that angle, so
it too tends to $\pi$. Continuity of $S_k$ on the compact range
now gives $r_n-S_n\to0$. Since $r_n\geq d(p,q)$ and $S_n$
is nonincreasing, $S_0\geq\lim S_n=\lim r_n\geq d(p,q)$.
Every moving center has distance less than $\ell$ from each
endpoint. The only new segments are the specified joins to those
endpoints and subsegments of previous joins, proving the domain claim.
\end{proof}

\begin{lemma}[Almost minimum using one complete ball; ALG03]
\label{lem:alexandrov-buffered-almost-minimum}
Let $r:Y\to(0,\infty)$ be finite-valued and locally bounded below
by a positive number. For $o\in Y$ and $0<\epsilon<1$, assume
$\overline B(o,r(o)/\epsilon^2)$ is complete. There is a point
$p$ in this closed ball with $r(p)\leq r(o)$ such that
\[
 r(q)>(1-\epsilon)r(p)
 \quad\hbox{whenever }d(p,q)\leq r(p)/\epsilon.
\]
The quantifier over $q$ is in $Y$, not a truncated ball.
\end{lemma}
\begin{proof}
If no such $p$ exists, start with $x_0=o$ and successively choose
$x_{j+1}$ with
\[
 d(x_j,x_{j+1})\leq r(x_j)/\epsilon,\qquad
 r(x_{j+1})\leq(1-\epsilon)r(x_j).
\]
Induction gives $r(x_j)\leq(1-\epsilon)^jr(o)$ and
\[
 d(o,x_j)\leq\frac{r(o)}{\epsilon}
       \sum_{h=0}^{j-1}(1-\epsilon)^h<r(o)/\epsilon^2.
\]
Thus each chosen point, even though chosen in $Y$, remains inside
the complete ball. The geometric tail bound makes the sequence
Cauchy, with limit there. Local positivity at the limit contradicts
$r(x_j)\to0$. This proves the asserted full-radius quantifier.
\end{proof}

\begin{theorem}[Complete geodesic globalization for nonpositive bounds; ALG04]
\label{thm:alexandrov-written-complete-globalization}
A complete geodesic space with local curvature-$k$ four-point
comparison, $k\leq0$, has global hinge and four-point comparison.
It also has point-on-side lower comparison and independent radial
comparison-angle monotonicity for all specified minimizing segments.
\end{theorem}
\begin{proof}
For $p\in Y$ let $C(p)\in(0,\infty]$ be the supremum of all
$s$ for which every EXISTING hinge with endpoint $p$ and sum of
arms less than $s$ satisfies $\alpha\geq\theta^k$.
ALG01 makes these radii locally bounded below by a positive number.
If a bad hinge exists, some $C(o)=r_0<\infty$. Choose $M>r_0$
and apply ALG03 to $r(p)=\min\{C(p),M\}$ with $\epsilon=1/4$.
It gives $p_*$ with
\[
 r_*:=r(p_*)=C(p_*)\leq r_0<M,
 \qquad C(q)>3r_*/4\quad(d(p_*,q)\leq4r_*).
\]
Put $\ell=17r_*/16$. Then $2\ell/3=17r_*/24<3r_*/4$.
For any hinge with endpoint $p_*$ and sum of arms less than
$\ell$, its other endpoint $q$ lies within $\ell<4r_*$ of
$p_*$. The short-hinge hypotheses of ALG02 hold for both
endpoints; geodesicity supplies every new join. ALG02 proves
comparison for every such hinge. Hence $C(p_*)\geq\ell>r_*$,
a contradiction. Supremum radii are used only with STRICT smaller
arm sums, so no attainment of $C$ is assumed.

For four-point comparison choose any three geodesics from the
center. Each full comparison angle is at most its germ limit;
ALG01 bounds the sum of these limits by $2\pi$. For a point $z$
interior to a side $[a,b]$, compare the two adjacent hinges with
third vertex $v$. Their model angles have sum at most the sum of
the germ limits, hence at most $\pi$. The model-gluing calculation
in ALG01 gives point-on-side lower comparison and radial angle
monotonicity, one arm at a time. Degeneracies are as in ALG01.
This proves all conclusions without using a global angle theorem.
\end{proof}

\begin{theorem}[Comparison inside a complete interior buffer; ALG05]
\label{thm:alexandrov-written-coarse-local-globalization}
Let $Y$ be a locally compact length space with local curvature-$k$
four-point comparison, $k\leq0$. If $\overline B_Y(o,L)$ is
complete, every specified hinge with endpoint $o$ and sum of arms
less than $L/20$ satisfies $\alpha\geq\theta^k$.
The space $Y$ itself need not be complete or globally geodesic.
\end{theorem}
\begin{proof}
First, every closed ball about $o$ of radius strictly less than $L$
is compact. The compact-radius argument of AC01 works with exactly
this weaker completeness input: if the supremum of compact radii
were $\rho<L$, trimming almost minimizing paths gives total
boundedness of the closed $\rho$-ball. It is a closed subset of
the complete $L$-ball, hence compact. Finitely many relatively
compact neighborhoods of that ball, and the same trimming argument,
then give a larger compact radius still below $L$, a contradiction.
This uses the length property of $Y$, not an assertion that a closed
ball is a length space. The argument applies equally at any center
whose required complete ball is supplied.

Define $C$ as in ALG04 using existing hinges in $Y$. A bad hinge
of arm sum less than $L/20$ would give $0<r_0=C(o)<L/20$.
Apply ALG03 to $\min\{C,M\}$, $M>r_0$, using the complete
closed $16r_0$-ball at $o$. It gives $p_*,r_*$ as in ALG04 with
$d(o,p_*)\leq16r_0$ and $r_*\leq r_0$.
Put $\ell=17r_*/16$. The closed $3\ell$-ball at $p_*$ is
complete because it is a closed subset of the complete $L$-ball:
\[
 d(o,p_*)+3\ell
 \leq (16+51/16)r_0<20r_0<L.
\]
Every strictly smaller closed ball at $p_*$ is compact by the
first paragraph. Any two points of $B(p_*,\ell)$ can consequently
be joined by a minimizing geodesic: almost minimizing paths lie
in a fixed compact ball of radius less than $3\ell$, using the
nearer endpoint and half their length as in AC01. The minimizing
limit lies in $B(p_*,2\ell)$. This uses only the inherited
compact-curve and length facts already isolated in MC01.

For any hinge with endpoint $p_*$ and arm sum less than $\ell$,
all cradle centers and the other endpoint lie in $B(p_*,\ell)$.
The preceding paragraph supplies exactly the new joins needed by
ALG02, inside $B(p_*,2\ell)$. Its short-hinge assumptions hold
because ALG03 gives $C(q)>3r_*/4>2\ell/3$ there.
Thus $C(p_*)\geq\ell>r_*$, again a contradiction.
The constant $20$ is a sufficient project bound. This argument
does NOT establish KL's sharper $2D$ criterion.
\end{proof}

\begin{corollary}[Ambient comparison with a written larger margin; ALG06]
\label{cor:alexandrov-written-ambient-buffer}
Let $(X,p)$ be a complete length space and $R>0$. Suppose
$U=B_X(p,256R)$ is locally compact and its intrinsic metric has
local curvature-$k$ four-point comparison, $k\leq0$.
Then four-point comparison holds for the ambient distances on
$B_X(p,R)$. All radial triangles used in AC02--AC04 at radius
$R$ have point-on-side lower comparison and independent-arm
comparison-angle monotonicity, for every choice of their minimizing
segments. With their remaining hypotheses, the numerical conclusions
of AC03--AC04 are unchanged.
\end{corollary}
\begin{proof}
The length argument of AC02 makes $d_U$ finite and locally equal
to $d_X$. AC01 supplies all ambient joins between points of
$B(p,2R)$ in $B(p,4R)\subset U$, so their intrinsic and ambient
distances agree. This also holds for pairs of points in $B(p,4R)$,
by taking the joins in $B(p,8R)$.
For every $v\in B(p,4R)$ the intrinsic closed $200R$-ball is
complete: an intrinsic Cauchy sequence has an ambient limit $w$,
\[
 d_X(p,w)\leq d_X(p,v)+200R<204R<256R;
\]
local metric equality gives intrinsic convergence in $U$.
ALG05 therefore proves comparison for every hinge with an endpoint
in $B(p,4R)$ and sum of intrinsic arms less than $10R$.

For four points of $B(p,R)$, geodesics from one to the other three
lie in $B(p,2R)$ and have length less than $2R$. Their hinge sums
are less than $4R$, so their comparison angles are bounded by the
germ limits, whose sum is at most $2\pi$ by ALG01.
For AC02's radial triangles, all radial vertices lie in $B(p,2R)$,
each side is less than $3R$, and any joining side lies in
$B(p,4R)$. At a point interior to a side the adjacent hinge sums
are less than $8R$: its distance to a side endpoint is less than
$3R$ and its distance to the opposite vertex is less than $5R$,
using the nearer endpoint. Intrinsic distances agree with the
ambient ones by the joins already established. ALG05 applies to
both adjacent hinges, and ALG01's gluing proves point-on-side
comparison. Applying it successively to the radial sides proves
the same monotonicity as AC02. AC03's hyperbolic calculation and
AC04's finite packing then give their original constants.
The same argument applies to radial triangles based at $p$ with
both arms at most $R$, including endpoints on the closed $R$-ball.
\end{proof}

\begin{theorem}[Growing-region compactness without the sharp local input; ALG07]
\label{thm:alexandrov-growing-region-written-comparison}
Under the ORIGINAL hypotheses of MC18, the conclusion of AC41
holds without importing the source-qualified local Toponogov
statements AC02 or AC36. All its metric, length, comparison and
Hausdorff-dimension conclusions, and the net constants in AC08
and AC40, are retained. Accepted-foundation and Lean bindings are
still separate obligations.
\end{theorem}
\begin{proof}
For any fixed $R>0$, the original condition $r_i\to\infty$ permits
passing to a tail with $r_i>256R$ and $c_i\leq1$.
AC33/AC17 supply local compactness on this controlled region using
only LOCAL four-point neighborhoods; AC29 supplies AC07's nearby
chart. AC13's local equality of metrics transfers local curvature
to the intrinsic open ball. ALG06 replaces the use of AC02 in
AC03--AC04, with exactly the same radial inequality and covering
constant. Thus AC08 and MC13 give the same eventual nets and one
complete proper pointed limit; MC14 supplies its length property
and geodesics. No global properness of the sources is added.

For the four chosen limit points in AC37, keep its radius $R=M+2$
and its actual lifted points. Pass to the tail $r_i>256R$.
ALG06 with $k=-c_i$ supplies the ambient four-point inequality
previously obtained using AC36. AC35 passes these inequalities to
curvature zero. AC37's direct point-on-side algebra is unchanged.
For the AC40 dimension proof choose the tail
$r_i>256(R+1)$ and use the just-proved AC08 nets at $R+1$.
Their cardinalities have not changed, so AC38--AC40 give exactly
$C_R\delta^{-n}$ and $\dim_HY\leq n$ as before.
These choices are made for each fixed scale and accuracy; no one
tail for all scales is asserted. A larger finite window is available
under the SAME growing-region hypothesis, not supplied as a new
assumption. Singleton and dimension-zero cases remain those
already treated in the metric and local-chart arguments.
\end{proof}

\begin{proposition}[Specified comparison consumers and their parameter choices; ALG08]
\label{prop:alexandrov-written-comparison-consumers}
The written comparison above supplies the formerly source-qualified
comparison uses in AC65's long-strainer orthogonality, LC12's
lifted packet, and LC78's product-anchor directions, with their
original quantified conclusions. It also supplies the global
comparison input whenever the space is already complete and
geodesic and the local comparison hypothesis holds. The old fixed-margin statements AC02, AC36 and AC64
are not thereby proved at their original constants.
\end{proposition}
\begin{proof}
In AC65 use the same original radial segments but shorten first to
$r_i=L_i/1024$ instead of $L_i/100$. AC63 is valid for every
$0<r_i\leq L_i$, and now
\[
 256r_i<L_i,\qquad
 \sqrt{\sigma_i}\,r_i=1/(1024\sqrt{\sigma_i})\longrightarrow\infty.
\]
Its displayed angle lower bound, with this $r_i$, still tends to
$\pi/2$. Apply ALG06 at radius $r_i$ and curvature $-\sigma_i$
to shorten further to each fixed $T>0$. The common subsequence,
AC59 prefix convergence, AC35 limit passage and the final
orthogonality argument of AC65 are unchanged. This gives the same
conclusion of AC65 and hence the same comparison input to AC54;
no endpoint is moved in the statement of the long strainer.

In LC12 the fixed radius $R$ is chosen from one finite limiting
configuration. Replace its tail requirement $H_i>16R$ by
$H_i>256R$, possible because $H_i\to\infty$. The inherited
Riemannian-to-local-comparison input followed by ALG06 at $k=-1$
gives the required common ambient four-point region. All continuity
bounds, the common neighborhood radius and the remaining packet
properties in LC12 are then exactly as written.

In LC78 retain the order: first choose $s$ from $R$ and the desired
error, then choose $\nu_*$ small. Impose additionally
$\nu^{-1}>512(s+R+1)$. All relevant vertices lie in
$B(q,s+R+1)$, their geodesics lie in its doubled ball, and their
hinge arm sums are less than $4(s+R+1)$. Intrinsic closed
$200(s+R+1)$-balls there are complete inside the controlled
curvature region, exactly as in ALG06. ALG05 therefore applies
to EVERY chosen minimizing hinge. The germ limits are the usual
Riemannian angles by the inherited first-order distance expansion
for equal small arm lengths in a metric chart (the fixed-$k$ cosine formula has the same
zero-scale limit). Thus the all-direction inner-product estimates
and coordinate conclusions of LC78, hence its LC79 use, are
unchanged. The new bound is a permitted extra choice of the
existential $\nu_*$, not a stronger hypothesis on its inputs.

Finally ALG04 applies directly whenever completeness and actual
geodesics have already been supplied. It leaves the original metric
unchanged. These are SPECIFIED consumer bindings; other fixed-window
uses must still have their buffers checked, and none is closed merely
by mentioning globalization.
\end{proof}

\paragraph{Current disposition of comparison and the local-model goal.}
ALG01--ALG08 replace the missing comparison producer in the
MC18/AC41 growing-region route and in the three consumers just
checked. They also give the complete geodesic theorem. The
stronger sharp-buffer AC02/AC36/AC64 statements remain historically
preserved and source-qualified at their original margins. The next
step is the remaining local-consumer and simultaneous-parameter
export check for LC81/LC83--LC85, using these written margins.
Finite metric regularity, the smooth consumer carrier and the
original metric continue to be distinguished as in LFR01--LFR59.
No smooth-metric upgrade of LC56, accepted PC release binding,
Lean verification or completion of the full user goal is asserted.


\chapter{Three-manifold topology, prime factors and torus cuts}
\label{ch:three-manifold-topology}

\paragraph{Scope.} Carriers, collars, embedded tori, cuts and relative prime data.

\paragraph{Planned sections.}
\begin{enumerate}[leftmargin=*]
\item Smooth boundary wrappers.
\item Torus data.
\item Cut carriers.
\item Relative prime factors.
\item Incompressibility conventions.
\end{enumerate}

\paragraph{Existing material and present task.} M1a-M1c and R01-R07; inherited Schoenflies boundary retained. Develop the relative-prime factor comparison while preserving actual maps and boundary marks.

% BEGIN migrated B016
\section{M1 category policy and reconstruction boundary}
\label{sec:m1-category-policy}

Revision 17 left the word ``equivalence'' deliberately general. We now choose
the smooth category for the primary certificate and retain the underlying
topological maps as derived data. This matches the Ricci-flow input and avoids
making a smoothing theorem an untracked dependency at the end of the project.

\subsection{Chosen categories}

The M1 target uses the following policy:

\setlength{\tablebodywidth}{\dimexpr\textwidth-4\tabcolsep\relax}
\begin{longtable}{@{}P{0.280899\tablebodywidth}P{0.471910\tablebodywidth}P{0.247191\tablebodywidth}@{}}
\textbf{Object} & \textbf{M1 contract} & \textbf{Reason}\\\hline
\endhead
Input $M$ & A closed, connected, orientable smooth 3-manifold. & This is the
category in which the Ricci-flow producer and the completed PC result are
stated.\\
Public prime carrier & A closed, connected, oriented smooth prime
3-manifold. Relative region factors may retain external torus boundary,
but are separate intermediate data. & Public connected-sum factors and
relative sphere-cap reconstruction have distinct boundary contracts.\\
Embedded torus & A smooth two-sided embedding of the standard 2-torus, with
an explicit orientation and incompressibility proof. & Smooth embeddings give
the boundary and collar data used by the metric witnesses.\\
Cut piece & A compact smooth 3-manifold with boundary, an open smooth interior,
and an explicit smooth inclusion of the interior. & Complete metrics live on
the interior while gluing uses the compact carrier.\\
Boundary gluing & A smooth orientation-reversing equivalence between matched
torus sides, together with its induced topological map. & The orientation
convention is needed for oriented reconstruction.\\
Global reconstruction & An orientation-preserving diffeomorphism from the
assembled smooth carrier to $M$; its underlying homeomorphism is stored or
derived for topological consumers. & The endpoint certificate remains
compatible with smooth Ricci-flow data and geometric metrics.\\
\end{longtable}

The word ``diffeomorphism'' in this table means the project's chosen
manifold-with-boundary equivalence. If the underlying Lean API represents
these objects as manifolds with corners, the GC wrapper records the condition
that all boundary faces used here are smooth codimension-one tori. The wrapper
is an implementation boundary; it does not weaken the mathematical target.

\subsection{Public certificate fields under this policy}

The certificate has the following category-sensitive fields:

\begin{itemize}[leftmargin=*]
\item \code{primeData} consists of closed smooth prime carriers and an
  orientation-preserving connected-sum diffeomorphism to $M$;
\item \code{torusData} consists of smooth embedded tori and their two-sided collars,
  with incompressibility stated using the selected fundamental-group API;
\item \code{cutData} consists of compact smooth cut carriers, their open interiors,
  and smooth boundary-side labels;
\item \code{geometryData} consists of a model tag, a smooth Riemannian metric on
  each interior, completeness, a finite-volume proof whenever the tag is
  hyperbolic, and the local homogeneity witness; and
\item \code{reconstruction} is an orientation-preserving diffeomorphism after all
  boundary sides are glued by the supplied smooth matching maps.
\end{itemize}

Every field has an underlying topological projection. The topology workstream
may use those projections to state prime, incompressibility, and gluing lemmas,
but it may not replace the smooth reconstruction field with a homeomorphism
without an explicit M1 API decision.

\subsection{Implementation boundary for manifolds with boundary}

The project should not implement quotient manifolds inside the public
certificate structures. Instead, introduce a small wrapper layer with three
interfaces:

\begin{enumerate}[label=\textbf{K\arabic*:},leftmargin=*]
\item a compact smooth carrier with boundary components and an interior;
\item a smooth boundary equivalence and its orientation-reversal proof; and
\item an assembled-carrier constructor with a universal reconstruction map.
\end{enumerate}

The wrapper may initially be implemented using existing quotient or manifold-
with-corners constructions. Those implementation choices remain behind the
interfaces. This keeps the endpoint theorem stable if the best quotient API
changes during the Lean audit.

The required API audit records, for each wrapper, the exact declaration,
import path, dimension hypotheses, orientation conventions, and transitive
axiom closure. A missing boundary API becomes a named M1 gap; it is not
silently replaced by a set-theoretic quotient or an unproved smoothing axiom.

The current source checkout gives a useful starting point. The
\code{Topology/Manifold/Boundary} and \code{BoundaryCollar} trees contain
smooth collar, boundary-orientation, attachment, and gluing infrastructure;
the \code{Topology/Diffeomorph} tree contains the ambient diffeomorphism and
isotopy API; and \code{Topology/ThreeManifold/ConnectedSum} and
\code{CutCap} contain smooth connected-sum and cut/cap reconstruction
components. The audit has not yet found a single public wrapper satisfying all
K1--K3 fields for arbitrary torus cuts. M1 should compose and expose those
pieces through the new GC wrapper instead of treating the existing file tree
as an already completed endpoint API.

\subsection{Why smooth-first is the default}

The smooth-first choice has three practical advantages. The flow and surgery
producer already transports smooth metrics; the eight model geometries are
smooth metric objects; and the PC adapter can expose its spherical conclusion
in the same category. A topological-only endpoint would require a later
smoothing and metric-transfer theorem precisely where the final assembly is
most difficult.

The cost is a stronger M1 manifold-with-boundary API. We accept that cost and
isolate it in the wrapper modules. If the API audit shows that a particular
topological lemma is easier on underlying spaces, it may be proved there and
lifted through the stored smooth projections. The public certificate and its
reconstruction field remain smooth.

\subsection{Category acceptance tests}

M1a is complete only after the following tests use the same category policy as
the final certificate:

\begin{itemize}[leftmargin=*]
\item the standard spherical, flat, hyperbolic, and product examples carry
  smooth model metrics on their interiors;
\item a matched pair of torus boundaries has a smooth orientation-reversing
  gluing map and an orientation-preserving assembled reconstruction;
\item the underlying homeomorphism projections compose as expected with the
  smooth equivalences;
\item no assembly theorem has a hidden coercion from a homeomorphism to a
  diffeomorphism; and
\item any temporary manifold-with-corners hypothesis is discharged or listed
  as a named theorem obligation before M1f is accepted.
\end{itemize}

The next implementation action is therefore an API audit and wrapper skeleton,
not a Ricci-flow theorem. Once this category boundary is stable, the M1
topology modules can be written without revisiting the endpoint notion of
equivalence.
% END migrated B016

% BEGIN migrated B017
\section{M1a boundary-wrapper contract}
\label{sec:m1a-boundary-wrapper}

M1a now has a concrete target: a small wrapper layer that represents compact
smooth 3-manifolds with boundary and hides the implementation of cutting and
gluing. The wrapper is not a second manifold library. It records only the
fields needed by the GC certificate and delegates analytic, smooth, and
topological facts to audited existing declarations.

\subsection{The three wrapper objects}

The first pass defines three mathematical objects. Their names are provisional
Lean identifiers; their fields are the contract.

\setlength{\tablebodywidth}{\dimexpr\textwidth-4\tabcolsep\relax}
\begin{longtable}{@{}P{0.285714\tablebodywidth}P{0.472527\tablebodywidth}P{0.241758\tablebodywidth}@{}}
\textbf{Object} & \textbf{Mandatory fields} & \textbf{Forbidden shortcut}\\\hline
\endhead
\code{GCManifoldWithBoundary} & A smooth manifold-with-boundary carrier; dimension
3; compactness; connectedness; orientation; a finite type of boundary
components; the open interior and its smooth inclusion; induced boundary
orientation; and a smooth collar for every selected boundary component. & A
bare topological space, an unproved finite boundary list, or a quotient whose
smooth structure is only asserted.\\
\code{SmoothBoundaryGluing} & Two boundary-side indices; the corresponding
boundary carriers; a smooth equivalence between them; proof that it reverses
the induced boundary orientations; and the fixed-point-free involution on all
matched sides. & An untyped pairing of names, or an orientation convention
that is left to the downstream reconstruction theorem.\\
\code{AssembledCarrier} & A finite family of compact carriers; compatible smooth
boundary gluings; the resulting compact smooth carrier; inclusion maps from
each piece; and an orientation-preserving diffeomorphism from the assembled
carrier to the target closed manifold. & Defining the assembled carrier to be
the target manifold, which would make reconstruction tautological.\\
\end{longtable}

The wrapper permits empty boundary. Thus a closed prime factor, a closed
hyperbolic piece, and a closed model example are represented without adding a
dummy torus. A cut piece has nonempty selected boundary only when the torus
system supplies it.

\subsection{Field-level obligations}

The wrapper exposes the following lemmas before any prime or geometry code is
written:

\begin{enumerate}[label=\textbf{B\arabic*:},leftmargin=*]
\item \code{interior_is_open}: the interior carrier is open in the compact carrier
  and its inclusion is smooth;
\item \code{boundary_components_disjoint}: distinct boundary-component indices
  have disjoint carriers;
\item \code{collar_restricts_to_boundary}: each collar restricts to the stated
  boundary inclusion and induces the stored boundary orientation;
\item \code{gluing_has_no_unmatched_side}: the side involution accounts for every
  selected boundary component exactly once;
\item \code{gluing_is_smooth}: the assembled smooth structure agrees with each
  piece away from the identified boundary and with the supplied boundary
  equivalence on the seam; and
\item \code{assembled_reconstruction}: the stored reconstruction map is an
  orientation-preserving diffeomorphism and its underlying homeomorphism has
  the expected restriction on each piece.
\end{enumerate}

B1--B3 are local wrapper facts. B4--B6 are the handoff to the assembly
theorems. In particular, \code{TorusSystemFormation} may use B4 to enumerate the
cut tori, while \code{GluingReconstruction} consumes B5 and B6 without opening the
quotient implementation.

\subsection{Existing API audit and adapter obligations}

The current checkout supplies candidate components, but each one needs a
signature and axiom audit before being classified as direct reuse:

\setlength{\tablebodywidth}{\dimexpr\textwidth-4\tabcolsep\relax}
\begin{longtable}{@{}P{0.386364\tablebodywidth}P{0.386364\tablebodywidth}P{0.227273\tablebodywidth}@{}}
\textbf{Candidate source} & \textbf{Potential contribution} & \textbf{M1 action}\\\hline
\endhead
\code{Topology/Manifold/Boundary} & Defining functions, boundary smooth
maps, boundary orientation, and sublevel/superlevel constructions. & Check
compatibility with the chosen wrapper and dimension-3 conventions.\\
\code{Topology/Manifold/BoundaryCollar} & Smooth collars, collar
attachments, seam charts, and collar diffeomorphisms. & Select the smallest
imports proving B2--B3; avoid importing flow-specific files.\\
\code{Topology/Diffeomorph} & Composition, restriction, extension, and
isotopy operations for smooth equivalences. & Build the wrapper's composition
and orientation-preservation lemmas.\\
\code{Topology/ThreeManifold/ConnectedSum} & Smooth connected-sum carriers,
orientation transport, seam charts, and reconstruction. & Adapt to
\code{PrimeDecomposition}; audit which results require closed factors.\\
\code{Topology/ThreeManifold/CutCap} & Cut/cap carriers, frontier labels,
quotient gluing, and reconstruction reductions. & Reuse only declarations
whose boundary and orientation hypotheses match B4--B6.\\
\code{Topology/VanKampen} & Boundary-collar and sphere-boundary injection
lemmas useful for later incompressibility arguments. & Keep out of the M1a
wrapper unless a named B-lemma needs it.\\
\end{longtable}

The audit is an adapter task even when a source file contains a theorem with a
similar name. Record the exact import, declaration, hypotheses, target type,
and \code{#print axioms} result. A theorem about a closed connected sum does not
automatically provide the compact-boundary gluing theorem required here.

\subsection{M1a acceptance matrix}

The wrapper is accepted when these examples and failure tests are available:

\begin{itemize}[leftmargin=*]
\item closed $\Sph^3$, $\mathbb{T}^3$, and a closed hyperbolic model use the
  empty-boundary branch;
\item $\mathbb{T}^2\times[0,1]$ supplies two boundary components with the
  induced opposite-side orientations;
\item gluing two compatible torus sides produces an assembled carrier with a
  reconstruction map, while reversing only one orientation proof makes the
  gluing contract fail;
\item deleting a boundary side from the matching involution makes B4
  unprovable; and
\item the M1a imports contain no Ricci-flow, surgery, Perelman, or finite-time
  extinction module.
\end{itemize}

After K1, the first geometric wrapper target should compile B1--B3 and
closed-model examples before implementing a general quotient assembly.
The earlier K0 boundary check uses finite labels only. B4--B6 may initially
be stated as wrapper interfaces in a clearly marked sandbox, but the release
target must replace them with proofs or audited imported theorems.

\subsection{Handoff to M1b}

M1b begins only after the wrapper exposes finite boundary indices and the
orientation-reversing side involution. It then defines \code{EmbeddedTorusData},
\code{TorusCutSystem}, and incompressibility. The prime-decomposition workstream
may proceed in parallel using the closed-carrier branch, but it must consume
the same \code{GCManifoldWithBoundary} type rather than introduce a second piece
representation.
% END migrated B017

% BEGIN migrated B018
\section{M1b embedded-torus and incompressibility contract}
\label{sec:m1b-torus-contract}

M1b supplies the torus data consumed by every later cut and geometric-piece
construction. The goal is a finite, typed system of smooth embedded tori with
explicit collars and $\pi_1$-injectivity. A phrase such as ``choose disjoint
incompressible tori'' is not an acceptable Lean input.

\subsection{One torus datum}

\code{EmbeddedTorusData X} is a certificate for one torus in a compact smooth
3-manifold $X$. It has the following fields:

\setlength{\tablebodywidth}{\dimexpr\textwidth-4\tabcolsep\relax}
\begin{longtable}{@{}P{0.285714\tablebodywidth}P{0.472527\tablebodywidth}P{0.241758\tablebodywidth}@{}}
\textbf{Field} & \textbf{Required content} & \textbf{Use}\\\hline
\endhead
Torus carrier & A fixed standard smooth 2-torus $\mathbb{T}^2$ with its
orientation and basepoint convention. & Gives a common domain for all torus
embeddings and fundamental-group maps.\\
Smooth embedding & A smooth embedding of $\mathbb{T}^2$ into the interior of
$X$, with its image and induced orientation. & Defines the embedded surface
that will become a cut boundary.\\
Two-sided collar & A smooth collar of the image, modeled on
$\mathbb{T}^2\times(-\varepsilon,\varepsilon)$, with center slice equal to
the embedding and disjoint positive and negative sides. & Supplies the two
boundary-side labels after cutting.\\
Incompressibility & Injectivity of the induced fundamental-group map, with
basepoint transport stated explicitly. & Prevents a compressible torus from
being admitted as a JSJ\slash geometrization cut.\\
Ambient orientation & Compatibility between the torus orientation, collar
normal direction, and orientation of $X$. & Fixes the induced orientations of
the two boundary sides.\\
\end{longtable}

The collar definition is the primary definition of two-sidedness in M1b. A
normal-bundle formulation may be proved equivalent later, but it is not used
as an unproved synonym. The positive and negative collar sides are ordered and
retain their order under every transport map.

\subsection{The incompressibility convention}

The project chooses $\pi_1$-injectivity as the public incompressibility
predicate. Because fundamental groups depend on basepoints, the wrapper
stores the following data and their compatibility explicitly:

\begin{enumerate}[label=\textbf{T\arabic*:},leftmargin=*]
\item a basepoint on $\mathbb{T}^2$ and its image in $X$;
\item the induced homomorphism on the selected fundamental-group objects;
\item a proof that this homomorphism is injective; and
\item transport lemmas showing that changing the chosen basepoint or collar
  identification preserves the predicate.
\end{enumerate}

The first implementation may use a project wrapper around the existing
fundamental-group API. It must not replace injectivity by a prose label such
as \code{incompressible}, nor by an assumption that every boundary torus of a
graph-manifold region is automatically incompressible. The graph-manifold and
hyperbolic producers must supply the actual proof for each torus retained in
the global system.

\subsection{Finite torus systems}

\code{TorusCutSystem P} is indexed by a finite family over a prime carrier $P$.
For a prime decomposition with index type $J$, the global system uses a
finite dependent index of the form
\[
  \{(j,t) : j\in J,\ t\in T_j\},
\]
where each $T_j$ is finite. This preserves which prime factor contains a
torus and avoids flattening the data into an untracked list.

The system contains:

\begin{itemize}[leftmargin=*]
\item one \code{EmbeddedTorusData} value for every dependent index;
\item a pairwise-disjointness proof for distinct torus images, including
  tori belonging to the same prime factor;
\item the finite boundary-side type obtained by taking the positive and
  negative collar side of every torus;
\item the fixed-point-free side involution identifying the two sides of each
  original torus; and
\item a cut-carrier interface connecting the system to
  \code{GCManifoldWithBoundary} and exposing the boundary-side labels.
\end{itemize}

The last field is deliberately an interface, not a definition by quotient.
It states that cutting along the system produces compact carriers whose
boundary sides are exactly the listed sides. The theorem
\code{TorusCutReconstruction} then proves that gluing each matched pair recovers the
original prime carrier. The general assembly theorem later adds the
geometric-piece metrics and the global reconstruction to $M$.

\subsection{Torus-system lemmas}

The first M1b theorem stubs are:

\setlength{\tablebodywidth}{\dimexpr\textwidth-4\tabcolsep\relax}
\begin{longtable}{@{}P{0.318182\tablebodywidth}P{0.443182\tablebodywidth}P{0.238636\tablebodywidth}@{}}
\textbf{Lemma} & \textbf{Statement shape} & \textbf{Consumer}\\\hline
\endhead
\code{TwoSidedTorusCollar} & An \code{EmbeddedTorusData} value supplies disjoint positive
and negative collar sides and induced boundary orientations. & Cut-piece
boundary labels.\\
\code{TorusIncompressible} & The stored induced $\pi_1$ map is injective and is
stable under the selected basepoint/collar transports. & \code{TorusCutSystem} and
hyperbolic/graph producers.\\
\code{TorusPairwiseDisjoint} & Distinct finite indices have disjoint images in the
relevant prime carrier. & Cut construction.\\
\code{TorusBoundarySides} & Cutting exposes exactly two labelled boundary sides for
each torus and no additional selected side. & \code{CutPiece} and gluing.\\
\code{TorusCutReconstruction} & Gluing the cut carriers along the side involution
recovers the prime carrier by an orientation-preserving equivalence. & Prime
and global reconstruction.\\
\end{longtable}

The proof of \code{TorusCutReconstruction} may use the existing cut/cap and collar
infrastructure, but its target must be the M1 wrapper types. A theorem that
reconstructs only an underlying set or only a capped manifold is an adapter
input, not the final M1b theorem.

\subsection{Failure tests and handoff}

M1b is accepted only when the following invalid constructions fail for the
stated reason:

\begin{itemize}[leftmargin=*]
\item a torus with no collar cannot inhabit \code{EmbeddedTorusData};
\item a compressible torus cannot inhabit the incompressibility field;
\item two intersecting torus images cannot inhabit a \code{TorusCutSystem};
\item a side involution with a fixed point or an unmatched side cannot produce
  a cut/reconstruction interface; and
\item a cut reconstruction with reversed orientation cannot satisfy the
  orientation-preserving target field.
\end{itemize}

The handoff to M1c is a \code{TorusCutSystem} together with its typed compact cut
carriers and boundary-side labels. M1c then defines \code{CutPiece} and proves the
boundary/interior relationship without changing the meaning of
incompressibility.

The field and theorem audit for this slice is recorded in
\code{m1_torus_api.csv}. It is maintained separately from \code{m1_boundary_api.csv}
because the torus layer adds the fundamental-group and pairwise-disjointness
obligations.
% END migrated B018

% BEGIN migrated B019
\section{M1c cut-piece and boundary-label contract}
\label{sec:m1c-cut-piece-contract}

M1c turns a \code{TorusCutSystem} into finitely many compact carriers with
boundary. It does not assign a geometry to a piece and does not decide which
pieces are hyperbolic or graph manifolds. Its only purpose is to make the
topological output of cutting precise enough for those later workstreams.

\subsection{One cut piece}

\code{CutPiece} is indexed by a prime factor and a component of the cut carrier.
It contains the following data:

\setlength{\tablebodywidth}{\dimexpr\textwidth-4\tabcolsep\relax}
\begin{longtable}{@{}P{0.300000\tablebodywidth}P{0.466667\tablebodywidth}P{0.233333\tablebodywidth}@{}}
\textbf{Field} & \textbf{Required content} & \textbf{Consumer}\\\hline
\endhead
Compact carrier & A \code{GCManifoldWithBoundary} with orientation and connected
component data. & Geometry witnesses and gluing.\\
Open interior & The open smooth interior and its smooth inclusion into the
compact carrier. & Complete metric and model-specific volume statements.\\
Prime/component label & The prime-factor index and a finite component index
for the cut complement. & Prime reconstruction and coverage.\\
Boundary labels & A finite family of boundary components and an equivalence
with the subtype of global sides assigned to this piece. & Matching
and torus-system formation.\\
Side realization & For every labelled side, a smooth boundary equivalence to
the corresponding torus-side carrier, preserving the stored orientation
convention. & Boundary gluing and collars.\\
Component relation & A proof that this carrier is the designated component of
the cut carrier and that its labelled boundary sides are exactly the sides
assigned to it. & Cut decomposition and boundary incidence.\\
\end{longtable}

The compact carrier is not defined as a closure in an untyped ambient set.
The cut operation supplies it as a smooth wrapper object, with the
complement relation and boundary labels proved separately. This avoids making
the correctness of the smooth structure depend on a later point-set closure
calculation.

\subsection{The cut decomposition object}

For one prime carrier $P$ and a finite \code{TorusCutSystem P}, define a
\code{CutDecomposition P} containing:

\begin{enumerate}[label=\textbf{C\arabic*:},leftmargin=*]
\item a finite component index type;
\item one \code{CutPiece} for each component;
\item an incidence map assigning every selected boundary side to exactly one
  component and boundary component;
\item a proof that every boundary component arising from a selected torus is
  represented exactly once with its side label; and
\item a reconstruction theorem that gluing the cut pieces along the torus
  side involution recovers $P$ with orientation preserved.
\end{enumerate}

Use the cut-at-the-tori reconstruction map
$q:\bigsqcup_v C_v\to P$ for these abstract compact carriers.
Its restriction to the disjoint union of interiors is a diffeomorphism
onto $P\setminus\bigcup_t T_t$. Before gluing, the carriers are disjoint
summands; their images under $q$ can meet on matched boundary sides.
In a self-pairing, two boundary copies of one carrier can have the same
image. No embedding of every closed cut carrier is asserted.

A realization by deleting open tubular neighborhoods is also permitted.
Its interiors map instead to the complement of the corresponding closed
collars. Store the collar comparison diffeomorphism to the abstract cut
and include the collars in reconstruction; do not equate these two
complements. Hatcher's archived splitting convention, printed 3--4/PDF
4--5, uses this neighborhood-deletion realization. This comparison is an
inherited collar/gluing obligation, not a new topology library.

The global \code{CutDecomposition M} is the dependent family of these objects over
the finite prime decomposition. Connected-sum sphere interfaces remain in the
prime layer; they are not silently mixed with torus boundary labels. An empty
torus system produces one cut piece with empty selected boundary, so closed
hyperbolic and flat examples use the same contract.

\subsection{Cut-piece theorem contracts}

The M1c theorem stubs are:

\setlength{\tablebodywidth}{\dimexpr\textwidth-4\tabcolsep\relax}
\begin{longtable}{@{}P{0.329545\tablebodywidth}P{0.431818\tablebodywidth}P{0.238636\tablebodywidth}@{}}
\textbf{Lemma} & \textbf{Statement shape} & \textbf{Consumer}\\\hline
\endhead
\code{CutPieceInterior} & Every cut piece has an open smooth interior with the
stored inclusion and induced orientation. & Geometry layer.\\
\code{CutPieceBoundarySides} & The boundary components of a piece are exactly the
side labels supplied by the torus system. & Gluing and geometry collars.\\
\code{CutPiecesDisjoint} & Images of distinct interiors under reconstruction
are disjoint; compact images meet only on matched boundary sides. & Coverage.\\
\code{CutPiecesCoverComplement} & Images of the abstract piece interiors
under reconstruction cover exactly the complement of the selected tori. & Reconstruction.\\
\code{CutDecompositionReconstructs} & Gluing all cut pieces along the side
involution yields an orientation-preserving equivalence to the prime carrier.
& Prime and global assembly.\\
\end{longtable}

\code{CutPiecesDisjoint} distinguishes the disjoint cut carriers from their
images after boundary identifications. Pairwise disjoint source carriers
are consistent with quotient gluing. Coverage is a field of
\code{CutDecomposition}, not of an individual \code{CutPiece}.

\subsection{Compatibility with the torus layer}

M1c consumes the following M1b fields without changing them:

\begin{itemize}[leftmargin=*]
\item the finite dependent torus index and prime-factor label;
\item the positive and negative collar sides and their induced orientations;
\item pairwise disjointness of the torus images;
\item the fixed-point-free side involution; and
\item the explicit $\pi_1$-injectivity proofs.
\end{itemize}

The cut operation may shrink collars or reparametrize their boundary copies,
but it must supply a named transport equivalence for every such change. A
later hyperbolic or graph-manifold theorem may use its preferred collar
normalization only after applying this transport lemma.

\subsection{M1c acceptance tests}

The cut layer is accepted with the following examples and negative tests:

\begin{itemize}[leftmargin=*]
\item an empty torus system on a closed prime produces one closed cut piece;
\item cutting $\mathbb{T}^2\times\Sph^1$ along a product torus produces the
  expected compact carrier with two labelled torus boundary sides;
\item the two sides of a separating torus may belong to distinct components,
  while the two sides of a nonseparating torus may belong to one component;
\item deleting one side label makes \code{CutPiecesCoverComplement} fail; and
\item changing a boundary orientation without changing its gluing map makes
  \code{CutDecompositionReconstructs} fail.
\end{itemize}

The handoff to M1d is a \code{CutDecomposition} whose compact carriers, interiors,
boundary labels, collars, and reconstruction equivalences are all typed. M1d
may now attach the eight model-geometry witnesses without defining cutting
again.

The field and theorem audit for this layer is recorded in
\code{m1_cut_api.csv}.
% END migrated B019

% BEGIN migrated B046
\section{Relative prime refinement: concrete targets}
\label{sec:relative-prime-refinement}

Revision 48 applies the standing source policy in \code{AGENTS.md}.
The following targets refine G02/G04/G05 within the existing topology and
reconstruction modules; they introduce no second foundation. They are
recorded in \code{relative_prime_refinement.csv} and remain unproved.

\paragraph{Source evidence and its limits.}
The archived Hatcher notes work in the smooth category, with compact,
connected, orientable manifolds possibly having boundary (printed 1/PDF 2).
Section 1.1 defines splitting and sphere capping on printed 3--4/PDF 4--5;
Proposition 1.4 and Theorem 1.5, including the existence proof and orientation
remark, are on printed 4--8/PDF 5--9. Their finite prime decomposition does
not assert that graph presentations survive. The proof uses smooth
Schoenflies, finite triangulations, and homological bounds, rather than a
claimed decrease in the number of auxiliary graph tori. Its discussion of
independence of cap choices also uses extension of a sphere diffeomorphism
over the ball. These are distinct foundation interfaces to compare at the
accepted release; citing Schoenflies alone does not bind all of them.

For the torus predicate, Section 1.2, printed 10--11/PDF 11--12, and
Corollary 3.3 with its proof, printed 48/PDF 49, connect the disk definition
to injectivity of the inclusion-induced fundamental-group map for a
two-sided surface. The latter proof removes trivial intersection loops
from a nullhomotopy and invokes the Loop Theorem, Theorem 3.1, printed
45/PDF 46. This records a dependency, not a new proof of the Loop Theorem.
There is no such automatic equivalence for a one-sided surface.

The author's current download page identifies a 2023 revision with 75
pages. Its smooth-category convention, Proposition 1.4, Theorem 1.5, and
Corollary 3.3 were checked at printed 1, 6--7, and 59--60/PDF 2, 7--8,
and 60--61. These selected statements agree with the uses here. This is a
separate remote comparison, not a replacement of the 61-page archive or
a full audit of the revised book. No separate applicable errata sheet was
located on the checked author page. See \cite{Hatcher3ManifoldNotes,Hatcher2023Comparison}.

\paragraph{Input and boundary contract.}
Work componentwise with a compact connected oriented smooth region $W$
whose external boundary consists of finitely many labelled tori; the
boundary may be empty. Store the actual boundary parametrizations and
collars. A supplied KL graph presentation is additional data. The generic
prime theorem is used before any claim of graph-class closure.
The proposed relative witness chooses its sphere cuts in the interior,
away from a smaller retained boundary collar. Only the newly created
spherical boundary components are capped. External tori remain labelled
boundary components, and the resulting correspondence must commute with
the reconstruction maps. This relative statement is an adapter obligation
extracted from the cutting construction, not an extra clause quoted from
Hatcher's theorem. No boundary incompressibility is assumed at this stage.

\begin{description}[style=nextline]
\item[R01: \code{RetainExternalBoundary}]
Given a finite interior sphere-cut system with collars, construct the cut
components, distinguish external tori from new sphere sides, and prove an
exact external-boundary correspondence before and after capping. Record
opposite orientations on paired sphere sides and preserve the external
boundary parametrizations. A count of components alone is insufficient.

\item[R02: \code{PrimeFactorsWithReconstruction}]
Use the prime-decomposition interface after its completed-foundation audit,
together with R01, to produce finite
prime factors together with the actual punctured components and oriented
connected-sum reconstruction of $W$. The output retains all factors,
including $S^3$ identity factors if chosen; it requires neither uniqueness
nor a minimal decomposition. A nonseparating-sphere argument must retain
its $S^2\times S^1$ contribution. Do not count it as an empty capped region.

\item[R03: \code{PrimeIrreducibleOrSphereProduct}]
Route each orientable prime factor to an irreducible carrier or an
explicit $S^2\times S^1$ alternative using Proposition 1.4. The reducible
Seifert example $\mathbb{RP}^3\#\mathbb{RP}^3$ from L10 must already be
split at the prime stage. This target does not produce model metrics.

\item[R04: \code{CapPi1Equivalence}]
For each connected punctured component, prove that attaching the chosen
balls along its sphere sides induces an isomorphism on $\pi_1$. The
proposed proof uses the inherited van Kampen interface with collars and
explicit basepoint paths. Keep the sphere-side condition: attaching a
solid torus along a torus can kill a meridian, so there is no analogous
unrestricted filling rule. This is a proposed application of foundation
results, not a source-checked Lean declaration. Section~\ref{sec:gluing-injection-transport}
now supplies the written application with its collar hypotheses.

\item[R05: \code{TorusDiskPi1Bridge}]
For a specified smooth two-sided torus in a specified carrier, identify
absence of a disk with essential boundary with injectivity of its actual
inclusion map on $\pi_1$. Preserve the carrier, basepoints, and collar
transport. Corollary 3.3 supplies the noninjective-to-disk direction;
the converse uses the surface disk criterion. Neither direction proves
that a given auxiliary torus satisfies either predicate.

\item[R06: \code{TransportRetainedTorus}]
Move the selected torus system off the cap balls with a proved simultaneous
isotopy, then use the punctured-component reconstruction and the group-map
comparison below. Prove disjointness and the relative-boundary statement
through this move. Ambient injection and later surgery-history transport
remain additional G05 obligations. Do not posit an embedding of a whole
capped prime factor into the original region.
Section~\ref{sec:cap-avoidance-output} supplies the written simultaneous
avoidance and relative torus output, with smaller fixed boundary collars.

\item[R07: \code{GraphPrimeClosure}]
After R02/R03, supply graph presentations on the prime factors and justify
the compression/filling simplifications leading to G04. This remains the
unproved relative implementation, now guided by the inspected 2007
Proposition 2.4.3 and its prime-uniqueness dependency. It is not a
consequence of R01--R06 or generic prime existence alone. Retain both ball-contained and solid-torus-side
compression, the fiber-parallel filling branch, and a separate termination
argument. No final essential-torus or geometric certificate is an input.
\end{description}

\paragraph{The carrier maps for G05.}
Let $R_a$ be a punctured cut component and $P_a$ its capped prime factor.
The maps actually supplied by the construction are
\[
 c_a:R_a\longrightarrow P_a,\qquad
 j_a:R_a\longrightarrow W.
\]
Here $c_a$ attaches caps and $j_a$ is the actual reconstruction map on
the punctured component (or inclusion of a shrunken core). A capped factor need not come with a map to $W$.
After R04, define the group map
\[
 \kappa_a=(j_a)_*\circ((c_a)_*)^{-1}:
 \pi_1(P_a)\longrightarrow\pi_1(W).
\]
For a retained torus $i:T^2\hookrightarrow R_a$, the required square is
\[
 (j_a\circ i)_*=\kappa_a\circ(c_a\circ i)_*.
\]
These formulas use compatible basepoints, or the recorded change-of-basepoint
isomorphisms. Proving $\kappa_a$ injective requires the connected-sum group
comparison; it does not follow just from writing the formula. Once both
it and $(c_a\circ i)_*$ are injective, composition gives injectivity into
$W$. For an ambient inclusion $e:W\hookrightarrow M$, prove injectivity of
$e_*\circ(j_a\circ i)_*$, for example using a separately proved injection
of $e_*$. Injectivity into $W$ alone is insufficient. This is the proposed
proof decomposition for the existing G05, not a claim that these geometric
maps or their Lean proofs have been constructed. The written group
injection argument is now supplied in
Section~\ref{sec:gluing-injection-transport}, conditional on the relative
reconstruction and cap-avoidance witnesses.

\paragraph{Reference access and release boundary.}
Matveev's 2003 Proposition 2.4.7 remains uninspected. The publisher confirms
separate 2003 and 2007 editions; its accessible chapter page is a preview,
and the indexed book mirror returned HTTP 404. The publisher's 2007
supplement contains a contents file and a Turaev--Viro chapter, not a
correction of Section 2.4. No applicable correction was found in these
checked materials; this does not certify that none exists. Edition
metadata and search snippets did not close A27. That historical access
record is in \code{reference_checks_revision48.md}. Revision 50 closes
only the source-access task by inspecting the separately supplied 2007
edition; the 2003 comparison is still unverified, and A28 records the
conditional-cut/proof-application obligation. The archive is read-only.

Smooth Schoenflies and Hamilton--Ivey are user-confirmed inherited PC
candidates. No declaration search in the unfinished snapshot is required
to preserve that classification. Other topological dependencies above
await comparison with the completed foundation; none is classified as
missing or assigned a reproof here. The seven R-targets are specifications,
all geometric regression tests remain unrun, and PC adapter/probe expansion
stays paused. The next subsection develops the separately sourced G03
filling obligations. R07's graph-definition and relative-boundary
comparison and G04's global termination argument remain open.
% END migrated B046

\chapter{Seifert fibrations and graph manifolds}
\label{ch:seifert-graph}

\paragraph{Scope.} Graph presentations, filling, compression and relative refinement.

\paragraph{Planned sections.}
\begin{enumerate}[leftmargin=*]
\item Graph conventions.
\item Retained boundary.
\item Primitive slopes.
\item Seifert filling.
\item Guarded rewrites.
\item Prime-phase progress.
\end{enumerate}

\paragraph{Existing material and present task.} G01-G05, F01-F06, J01-J05 and the written pants/prime arguments. Separate raw auxiliary cuts from retained essential cuts and expand remaining relative proof obligations.

% BEGIN migrated B040
\section[Graph refinement and endpoint gate]{Graph refinement and the endpoint compatibility gate}
\label{sec:graph-refinement-contracts}

Revisions 46--48 track nine geometric proof targets G01--G09 and the
source-reviewed endpoint decision G10, recorded in \code{graph_refinement_contracts.csv}. These are
mathematical specifications, not Lean declarations or checked proofs.
They refine existing theorem IDs and leave the inherited foundation and
PC release gates unchanged. In particular, a KL14 graph presentation does
not itself supply the incompressible cuts or metric witnesses required by
\code{ThinRegionToGeometricPieces}.
% END migrated B040

% BEGIN migrated B041
\section{Source versions and exact comparison points}

The Hatcher archive is the 61-page notes with small revisions in 1999 and
2000 mentioned in its preface; its printed page number is the PDF page
number minus one. It is not the separately available 2014 revision.
Scott's archived article has 87 PDF pages, with printed page number equal
to PDF page number plus 400. The Morgan--Tian locators below refer to the
August 7, 2012 manuscript body, not a verified 2014-book crosswalk.
Morgan--Tian 2014 remains the global spine.

\begingroup\small
\setlength{\tablebodywidth}{\dimexpr\linewidth-4\tabcolsep\relax}
\begin{longtable}{P{0.35\tablebodywidth}P{0.65\tablebodywidth}}
\textbf{Source and locator} & \textbf{Role and restriction}\\\hline
\endfirsthead
\textbf{Source and locator} & \textbf{Role and restriction}\\\hline
\endhead
Hatcher, Section 1.2, facts (4)--(5); printed 11/PDF 12 &
Compression dichotomy in an irreducible manifold; cutting along
incompressible surfaces and irreducibility. The ball-contained alternative
must remain.\\
Hatcher, Theorem 1.9; printed 14/PDF 15 &
Incompressible torus decomposition into atoroidal or Seifert pieces.
This alone does not eliminate the atoroidal alternatives for a graph
presentation.\\
Hatcher, Propositions 1.12--1.13; printed 18--19/PDF 19--20 &
Seifert irreducibility exceptions and essential vertical/horizontal
surfaces; hypotheses and compressible exceptions are explicit.\\
Hatcher, Section 2.1 and Proposition 2.1; printed 24--25/PDF 25--26 &
Slope-sensitive Seifert filling; a meridian parallel to the fiber is
excluded from the ordinary extension.\\
MT12, Lemma 5.1; printed 54--55/PDF 58--59 &
Solid-torus, product and twisted interval-bundle cases; finite-volume
interior geometry is a separate alternative.\\
MT12, Definition 5.2, Proposition 5.4; printed 55--56/PDF 59--60 &
Prime/Seifert graph definition and global reconstruction. The proof uses
Klein-bottle cuts in an exceptional case.\\
Scott, Nil/Sol discussion; printed 469--472/PDF 69--72 &
Monodromy distinguishes Nil from Euclidean and Sol torus bundles.\\
Scott, Theorem 5.3; printed 477--478/PDF 77--78 &
Classification of \textbf{closed} nonhyperbolic geometric manifolds;
not a theorem about every open Seifert interior.\\
\end{longtable}
\endgroup

These comparisons use \cite{Hatcher3ManifoldNotes,Scott1983,MorganTian2012Archive}.
The source ledger now includes L08--L32 in addition to L01--L07.
A located theorem is evidence for a proof target, not an automatic conversion
between the sources' categories.
% END migrated B041

% BEGIN migrated B042
\section{G01--G05: presentation, compression and essential cuts}

\begin{description}[leftmargin=1.3cm,style=nextline]
\item[G01] Extract the finite embedded auxiliary torus family and the
circle-bundle structures on the completed cut pieces from KL14 Definition
1.2. Regard each circle bundle as a Seifert structure with no exceptional
fibers. Keep external boundary labels and gluing maps. Do not identify
this auxiliary torus index with the final incompressible torus index.
An orientable total space need not have an orientable base surface.

\item[G02] Pass to prime factors relative to the retained external
boundary. Supply sphere cuts, capped factors and reconstruction data;
prove that the graph presentation survives the required operations.
For an embedded two-sided torus in an \emph{irreducible} carrier, Hatcher's
compression dichotomy has two alternatives: it bounds a solid torus or
it lies in a ball. They are alternatives, not asserted to be disjoint.
The second cannot be discarded or treated as an already identified
solid-torus side. The graph-class closure proof remains open; KL Appendix I points to
Matveev Proposition 2.4.7. The second edition is now inspected: its
Proposition 2.4.3 supplies the connected-sum closure argument, whereas
Proposition 2.4.7 already assumes incompressible Seifert cuts. See
Section~\ref{sec:matveev-graph-rewrites}; neither citation completes
the relative formalization.

\item[G03] When a solid-torus side is absorbed into a Seifert piece,
retain its attaching meridian slope $m$ and the regular fiber slope $f$
in $H_1(T;\mathbb Z)$. They are primitive unoriented slopes. The quantity
\[
             \Delta(m,f)=|\det(m,f)|
\]
is independent of an oriented basis up to the displayed absolute value.
For $\Delta>0$, the standard filling extends the Seifert structure with
core multiplicity $\Delta$; multiplicity one is a regular fiber.
For $\Delta=0$, that extension theorem is unavailable. Record a separate
fiber-parallel filling case and prove its topology before proceeding.
Do not assign it a model tag by the transverse-filling rule. F01--F06
in Section~\ref{sec:seifert-filling-targets} specify the two branches,
the pants splitting move, surviving slopes and local filling progress.

\item[G04] Construct essential torus cuts on the appropriate irreducible
prime carriers, with a termination argument for the chosen simplifications.
Hatcher's Seifert irreducibility result leaves $S^2\times S^1$ and
$\mathbb{RP}^3\#\mathbb{RP}^3$ as orientable exceptions; split or handle
them through the prime/exceptional layer. His Proposition 1.13 requires
an irreducible carrier and two-sided surfaces, and excludes a vertical
torus bounding a model-fibered solid torus with at most one multiple
fiber. It does not say that every vertical torus is incompressible.
Theorem 1.9 supplies a general torus decomposition, but a proof is still
needed that the retained graph presentation yields the desired Seifert
alternatives. Record this obligation instead of calling generic JSJ
existence the complete graph-to-Seifert bridge.

\item[G05] Separate injectivity into a cut piece, into a graph region,
and into the ambient prime carrier. For each retained torus, give the
actual inclusion maps and the injectivity or gluing theorem that composes
them. Incompressibility relative to one carrier does not imply
incompressibility in an arbitrary larger manifold. Keep disjointness,
external-boundary avoidance, orientations and the maps identifying final
labels. Later surgery-history transport is an additional obligation;
none of these steps follows merely from local volume collapse.
\end{description}
% END migrated B042

% BEGIN migrated B047
\section{G03 implementation targets: Seifert filling with retained boundary}
\label{sec:seifert-filling-targets}

Revision 49 refines a local topological operation needed by the graph
simplification. It does not alter the analytic spine or open the PC release
gate. The six targets F01--F06 in \code{seifert_filling_contracts.csv}
are unproved specifications. They use oriented smooth compact carriers,
explicit boundary identifications and inherited topology to be bound at
release. Neither a fundamental-group calculation nor a factor count is
accepted in place of a diffeomorphism and reconstruction data.

\paragraph{Source scope and conventions.}
The author-hosted \cite{Martelli2025v4} identifies itself as Version 4 of
September 2025, with 496 PDF pages; printed page $p$ is PDF page $p+8$.
Definitions 10.1.1 and 10.1.4 and Lemma 10.1.2, printed 305--307/PDF
313--315, fix Dehn filling and lens-space conventions. Section 10.3.1,
printed 316--317/PDF 324--325, fixes the Seifert notation.
The proof of Proposition 10.3.41, printed 332/PDF 340, gives the local
pants move. Exercise 10.3.42 and Corollaries 10.3.43--10.3.44, printed
332--333/PDF 340--341, locate its general application. The exercise has
no supplied proof; the stated corollaries do not supply the missing
relative boundary argument. The reading and correction-search limits,
including the exact downloaded hash, are in
\code{reference_checks_revision49.md}.

Compare the archived Hatcher Section 2.1, printed 24--26/PDF 25--27,
and its lens-space discussion, printed 31--32/PDF 32--33. With section
and fiber classes $(s,f)$, write the attaching meridian as
\[
       m=p s+q f,\qquad \gcd(p,q)=1,\qquad
       \Delta(m,f)=|p|.
\]
Hatcher's fibered-boundary fraction is $q/p$, not $p/q$.
This uses a chosen basis on that boundary, not the meridian--longitude
basis of an arbitrary solid torus. Martelli's lens convention kills
$q\mu+p\lambda$ in $L(p,q)$; Hatcher's $L_{\alpha/\beta}$ kills
$\alpha\mu+\beta\lambda$. Thus the numerical correspondence is
$(p,q)=(\beta,\alpha)$ only after matching the actual bases and attaching
maps. These symbols alone are not an oriented identification of lens
spaces. In particular, a fiber class that becomes a solid-torus meridian
must not be silently reused as its new regular Seifert fiber.

\begin{description}[leftmargin=1.3cm,style=nextline]
\item[F01: \code{PrimitiveSlopeTransport}]
Store the unoriented primitive classes, the boundary lattice and the
induced map of each attaching diffeomorphism. Prove that changing a
lattice basis by a unimodular matrix or reversing a slope representative
preserves $\Delta$. Prove that $\Delta=0$ means the same unoriented slope;
primitivity is essential. The vector $(0,2)$ is not a valid primitive
meridian, and a determinant calculation cannot repair that invalid input.
An orientation-reversing boundary gluing and the induced lattice map
must agree with the oriented reconstruction record.

\item[F02: \code{TransverseSeifertFill}]
For $\Delta>0$, construct the extended Seifert structure and its core
multiplicity $\Delta$, treating $\Delta=1$ as regular. Transport every
unfilled boundary label and collar. This is the Hatcher/Martelli
transverse-filling branch. It has no output for $\Delta=0$ and produces
no complete model metric by itself.

\item[F03: \code{PantsFiberFill}]
For $P\times S^1$, choose an essential arc in the pair of pants $P$
with both ends on the boundary being filled. Its vertical annulus closes
with two meridian disks of the added solid torus. Following the inspected
proof of Proposition 10.3.41, the resulting separating sphere gives two
punctured solid tori; capping gives their connected-sum presentation.
The implementation must retain the two external tori and show that their
old fiber slopes are meridians of the new solid-torus factors.
Sphere caps, punctured pieces and reconstruction maps use the R01/R02
conventions. The source proof does not itself provide the proposed
labelled-collar or orientation-sensitive API.

\item[F04: \code{AssembleFiberParallelFactors}]
Extend F03 to one selected boundary of a connected Seifert carrier.
The base surface $S$ here has the exceptional-fiber disks restored;
its $b\geq1$ remaining boundary circles are the external boundaries.
For coprime pairs $(p_i,q_i)$ with $p_i>0$, the expected factors from
Corollary 10.3.43 are
\[
 \mathop{\#}_{i=1}^{h} L(p_i,q_i)
 \ \#\ \mathop{\#}^{k}(S^2\times S^1)
 \ \#\ \mathop{\#}^{b-1}(D^2\times S^1),
 \qquad k=2-b-\chi(S).
\]
If $S$ is orientable of genus $g$, then $k=2g$; if its nonorientable
genus is $g$, then $k=g$. An orientable total space over a nonorientable
base uses the orientable circle bundle, not an assumed product.
The missing general bundle proof in Exercise 10.3.42 must be supplied
with the relative data needed here. Prove the boundary correspondence,
transported meridians, oriented factor choices and reconstruction;
do not silently strengthen the displayed source statement to these fields.
Retain factors with $p_i=1$ until an explicit $S^3$-identity move removes
them. An empty factor list reconstructs $S^3$, not the empty manifold.
No claim that every factor is irreducible is made.

\item[F05: \code{ResumeFillingOnSummands}]
If more original boundary components were selected for filling, transport
their attaching slopes through the actual relative diffeomorphism from
F04. They now lie on the labelled solid-torus factors. Choose and record
meridian--longitude bases there before identifying any subsequent lens
space. A remaining attaching slope equal to the transported meridian
has a sphere-product outcome; it is not a transverse extension of the
old Seifert structure. The disk-base, no-exceptional-fiber case gives the
$S^3$ identity already noted. Prove compatibility with connected-sum
reconstruction; an isomorphic fundamental group alone is insufficient.

\item[F06: \code{LocalFillingProgress}]
For a fixed finite set of selected original external boundary tori, use
the number still unfilled as a natural-number measure. F02 removes one
label and preserves all other external labels; F04 replaces the selected
piece by factors with exactly its other $b-1$ external labels. F05
transports the remaining requests. Prove one-step reconstruction and
decrease of this count. This gives only a finite filling procedure.
A G04 compression may change the graph presentation or introduce new
operations; its global termination proof is a separate obligation.
\end{description}

\paragraph{Required geometric regression specifications.}
The future tests must include a pants base giving two labelled solid
tori, an annulus base giving one, a disk base with no exceptional data
giving $S^3$, a once-punctured torus giving two sphere-product summands,
and a M\"obius base with orientable total space giving one sphere-product
summand. A disk base with pairs $(2,1),(3,1)$ must retain both lens factors.
Also reject a nonprimitive meridian and a transverse-filling call with
$\Delta=0$. These examples are specification cases, not executed geometric
tests. They require actual maps and slopes, beyond arithmetic counts.

\paragraph{Graph closure and inherited dependencies remain separate.}
Martelli Definition 11.5.8 and Exercise 11.5.9, printed 369/PDF 377,
offer another graph-refinement route through connected sums and geometric
decomposition. The latter is an exercise with a hint using Corollary
10.3.44; it is not a supplied proof of R07. Moreover the adjacent geometric
decomposition permits one-sided Klein surfaces, unlike E1. Any use needs
a graph-definition comparison and the G08/G05 return to torus-only cuts.
The new Matveev comparison in Section~\ref{sec:matveev-graph-rewrites}
closes A27 only as a source-access task and leaves A28 open. Local filling progress
does not close R07 or G04, and no filling result produces G09's metrics.

For filling independence, Martelli Lemma 10.1.2 uses sphere-cap
independence (Proposition 9.2.1, printed 275/PDF 283), which in turn uses
sphere-diffeomorphism isotopy/extension (Theorem 6.4.5 and Corollary 6.4.6,
printed 184--185/PDF 192--193). These are explicit inherited-foundation
interfaces to compare at release, not consequences to infer merely from
a theorem named Schoenflies. The user-confirmed smooth Schoenflies and
Hamilton--Ivey candidates retain their status; no PC snapshot probe was
added. The general relative proof of Exercise 10.3.42 behind F04 remains
open. The bounded extended-pants branch below uses F03 directly and
does not wait for that larger bundle exercise.
% END migrated B047

% BEGIN migrated B048
\section{Matveev second edition: source scope and guarded graph rewrites}
\label{sec:matveev-graph-rewrites}

The newly supplied \code{MatveevBook2ndEd.pdf} identifies itself as the
2007 second edition of \cite{Matveev2007}. Its 498 PDF pages do not have
a single printed-to-PDF offset. Here printed 84--88 are PDF 96--100,
printed 123--125 are PDF 135--137, and printed 248--249 are PDF 258--259.
The exact source hash, reading scope and publisher/correction checks are
in \code{reference_checks_revision50.md}. The 2003 edition cited by KL
remains a separate bibliographic record. A27 is closed only for obtaining
and reading a usable primary source; no first-edition equivalence or
formal graph-refinement theorem is thereby established.

\paragraph{What the inspected statements supply.}
Definition 2.4.1, printed 84/PDF 96, builds graph manifolds from solid
tori and pants times a circle, using boundary homeomorphisms. Free
boundary tori are permitted. Proposition 2.4.2, printed 84--85/PDF 96--97,
reduces orientable Seifert manifolds to these blocks. Its nonorientable-base
step uses the alternative disk-base fibration of the orientable twisted
interval bundle over the Klein bottle; Example 6.4.14, printed
248--249/PDF 258--259, describes that alternative. This supports a
comparison with KL's circle-bundle presentation, but the smooth,
orientation-preserving, labelled-collar version is an adapter to prove.
It does not authorize a one-sided Klein cut in E1.

Proposition 2.4.3, printed 85/PDF 97, states closure under connected sum
in both directions. Its supplied argument treats extended pants blocks
with attached solid tori, consolidates a solid-torus block, and splits the
fiber-parallel case. It then uses uniqueness of prime factors to transfer
graph membership to specified connected-sum factors. Thus R02's existence
and reconstruction fields alone do not supply that particular proof.
Either bind the needed inherited prime-uniqueness comparison or construct
and compare the required factors directly. Endpoint uniqueness remains
unnecessary; a proof's auxiliary uniqueness dependency is a different matter.
The phrases describing repeated simplification do not provide all the
relative maps or a formal measure for our combined procedure.

Corollary 2.4.4 and Remark 2.4.5, printed 86--87/PDF 98--99, treat a
proper disk cut followed by capping spherical boundary components. The
three reconstruction cases retain two capped components, an additional
$S^2\times S^1$ summand, or an additional solid-torus summand. A plain
boundary connected sum is not a graph-closure rule. The cap operation is
on sphere sides; it is not arbitrary torus Dehn filling. Applying this
result to external labels still requires relative collars and maps.

\paragraph{The exact canonical-subsystem input.}
Definition 2.4.6 and Proposition 2.4.7, printed 87--88/PDF 99--100,
require an irreducible, boundary-irreducible graph carrier and an already
incompressible torus system whose cut pieces are Seifert. The conclusion
selects a subsystem with no compatible fiber choices across a retained
seam. This is a conditional coarsening theorem. It does not, as stated,
construct the first incompressible Seifert decomposition from auxiliary
KL cuts. Boundary irreducibility must be established or the exceptional
solid-torus case treated separately. The stronger canonical conclusion
is optional for E1, which asks for neither uniqueness nor minimality.

There is a proof-application issue to retain explicitly. The second move
in the supplied proof discusses a compressible torus and asserts a
solid-torus side using irreducibility. Such a torus is absent from the
literal input above. If that move is adapted to an earlier raw
presentation, Hatcher's printed 11/PDF 12 compression dichotomy also has
a ball-contained alternative. Do not discard it or infer a solid-torus
witness from irreducibility alone. We record A28 as an open application
obligation; this observation does not declare the proposition false or
silently correct its wording.

The relevant incompressibility convention is checked in Matveev
Definitions 3.3.1--3.3.4 and Lemma 3.3.5, printed 123--125/PDF 135--137.
For two-sided tori, disk incompressibility agrees with injectivity of the
specified inclusion on $\pi_1$; the proof invokes the Loop Theorem and
Dehn Lemma. The one-sided discussion is different and is not imported
into \code{TorusCutSystem}. R05's carrier and two-sidedness requirements
remain in force.

\paragraph{Concrete obligations J01--J05.}
The subordinate inventory \code{graph_rewrite_contracts.csv} records the
following unproved implementation targets. These are source-guided proof
interfaces; the assertions about geometric moves remain to be proved.

\begin{description}[leftmargin=1.3cm,style=nextline]
\item[J01: \code{ElementaryGraphPresentation}]
Translate G01's circle-bundle pieces into Matveev elementary blocks,
including nonorientable bases and the alternate twisted-bundle fibration.
Keep actual gluing maps, free boundary labels, orientations and the
smooth-category comparison. Membership up to unspecified homeomorphism
is not the complete relative output required by GC.

\item[J02: \code{GraphPrimeClosureImplementation}]
Use the inspected Proposition 2.4.3 argument with F03 and the local
transport proof in Section~\ref{sec:extended-pants-splitting} for the
extended-pants filling case, and R01--R03 for sphere cuts and reconstruction.
F04--F05 retain their role for larger Seifert blocks.
Retain $S^3$ identities, sphere-product contributions, and the cap/cut cases
of Corollary 2.4.4. Prove the innermost-circle reduction, progress of the
chosen splitting process, and either the prime-uniqueness comparison or
an explicit direct factor comparison. Source access is no longer missing;
these proof and relative-data obligations remain.

\item[J03: \code{CompatibleFiberCoarsening}]
On a fixed carrier with an already incompressible Seifert-cut system,
remove a seam only after exhibiting compatible fibrations and proving
that they glue through the recorded boundary map. Preserve all other
embedded tori and external data. When the two sides belong to the same
cut component, the compatibility witness must use one coherent fibration
on that component; two independently chosen fibrations do not suffice.
Retained incompressibility concerns unchanged embeddings in the same
carrier. Do not infer it just from a new abstract block name.

\item[J04: \code{CertifiedSolidTorusAbsorption}]
For use before incompressibility is established, require an actual
embedded solid-torus region, a nonempty cluster of internal seams inside
it, and proof that the region is a union of complete cut pieces with the
required retained boundary. Erase those internal seams and replace the
cluster by the certified region, preserving reconstruction. A ball-contained
compression does not supply these fields. Nor may a selected compressible
interface be declared resolved if it is merely retained as the new region's
boundary. Eliminating all compressible interfaces requires a separate
coverage/progress argument.

\item[J05: \code{FullGraphRefinementProgress}]
Compose the justified relative prime, filling and torus moves without
assuming the desired essential-cut certificate as input. Supply an overall
well-founded measure or a proved decomposition into terminating phases,
and show that every terminal state has the required incompressible
Seifert cuts or a separately handled exceptional carrier. J03's fixed-carrier
coarsening and J04's local erasures do not establish this claim. F06's
count of selected filling requests also remains a different measure.
\end{description}

\paragraph{A proved finite bound with a conditional geometric interface.}
There is a useful elementary termination calculation independent of the
unfinished geometric moves. Fix a finite set $E_0$ of internal seam labels.
During an erasure-only phase let $E_i\subseteq E_0$ be the active labels,
and require every allowed step to carry a strict inclusion
$E_{i+1}\subsetneq E_i$. No step introduces a new seam. Then
\[
              |E_i|+i\leq |E_0|.
\]
Indeed the claim holds at $i=0$, and strict inclusion of finite sets gives
$|E_{i+1}|+1\leq |E_i|$. Adding $i$ proves the next instance by induction.
Consequently such a phase has at most $|E_0|$ steps. This is a written
finite-set proof, not new Lean evidence. Counting cut components instead
would not cover a compatible self-gluing seam: deleting that seam can
leave the number of components unchanged.

The calculation proves termination only after the geometric relation has
been shown to satisfy its strict-subset contract. It does not construct
any admissible move, show that a stopped state is incompressible, or prove
progress for prime splitting and filling phases that may change the label
universe. At a normal form, canonical fiber incompatibility additionally
needs the converse move-availability statement: every compatible seam
admits a J03 move. E1 needs only the weaker geometric endpoint, not this
canonical enhancement. These distinctions are guarded by the static audit.

\paragraph{Next proof task and release boundary.}
The next subsection develops J02's extended-block splitting using the
explicit pants move, including boundary and meridian transport. The
prime-phase innermost-circle and terminal argument is developed in
Section~\ref{sec:graph-prime-normal-form}; full J05 essential-cut progress
remains distinct. F04's general bundle exercise
remains open and is available when a larger Seifert block is used. The
canonical-subsystem statement cannot replace this work. All topology
interfaces eventually bind to the completed PC foundation; smooth
Schoenflies and Hamilton--Ivey remain inherited candidates, with no new
snapshot probes, foundational reproofs or changed analytic spine.
% END migrated B048

% BEGIN migrated B049
\section{Extended pants blocks: relative splitting and a bounded phase}
\label{sec:extended-pants-splitting}

This subsection develops J02's fiber-parallel branch, rather than the
whole graph-prime theorem. The local construction follows Martelli
Proposition 10.3.41, printed 332/PDF 340, and Matveev Proposition 2.4.3,
printed 85/PDF 97. The relative maps and the phase bound below are our
written elaboration of those passages. They are not checked Lean
implementations. In particular, this two-boundary calculation uses F03
directly: F04's general bundle exercise is not a prerequisite for this
branch. F05 still has its general F04 dependency for larger Seifert blocks.
The exact reading and inherited interfaces are recorded in
\code{reference_checks_revision51.md}.

\paragraph{Input and explicit sphere.}
Let $P$ be an oriented pair of pants with labelled boundary circles
$\gamma_0,\gamma_1,\gamma_2$, put $T_i=\gamma_i\times S^1$, and attach
an oriented solid torus $U_0$ to $T_0$ using a recorded boundary map
$a_0:\partial U_0\to T_0$. Its unoriented meridian slope is the fiber
slope. All gluings reverse induced boundary orientations. Work with
smooth collars and the actual maps; a bare homeomorphism-class label
from Matveev's elementary presentation first requires J01's adapter.
The other two tori may remain exposed or carry specified solid-torus
fillings. Exposed tori may subsequently be paired with outside ports,
including a pairing with each other.

Choose an embedded essential arc $\alpha$ in $P$ with endpoints on
$\gamma_0$, separating $\gamma_1$ from $\gamma_2$. Its product annulus
$\alpha\times S^1$ has two parallel boundary fibers on $T_0$. Under
$a_0^{-1}$ these are parallel meridians; choose disjoint meridian disks
$D_-,D_+\subset U_0$ with precisely these boundaries. Existence with these
boundary representatives uses meridian isotopy and collar extension,
not just equality of homology classes. Rounding corners gives
\[
 S=(\alpha\times S^1)\cup D_-\cup D_+.
\]
This sphere lies away from $T_1,T_2$ and their chosen small collars.
Cutting the filled block along $S$ produces two connected pieces $C_i$,
each with boundary $T_i\sqcup S_i$. Cutting the pants along $\alpha$
gives two annuli; on each side the corresponding portion of $U_0$
attaches a two-handle to $T^2\times I$ along a primitive curve.
This is the punctured-solid-torus identification in Martelli's proof.
Cap only $S_i$ to obtain a solid torus $V_i$, recording the cap ball
$Q_i$, a diffeomorphism $C_i\cong V_i\setminus\operatorname{int}Q_i$,
and its boundary restriction $h_i:T_i\to\partial V_i$.

The surviving fiber on $T_i$ bounds a meridian disk: take a base arc
from $\gamma_i$ to $\gamma_0$, in the corresponding annular side and
away from $\alpha$, form its product with the fiber, and close the
$\gamma_0$ end by a meridian disk in the corresponding cut portion of
$U_0$. Its boundary is that fiber on $T_i$. Thus $h_i$ sends the old
fiber slope to the meridian slope $\mu_i$ of $V_i$. The boundary map is
recorded, not replaced by an unspecified identification. Choose the
output collars so that its extension is $h_i\times\mathrm{id}$ there;
no claim of identity relative to an arbitrary preassigned output
parametrization is needed. The caps and sphere gluing retain the
punctured reconstruction, including the case that a capped factor is $S^3$.

\paragraph{Further fillings commute with the cut.}
If a solid torus $U_i$ was attached at $T_i$ by $a_i$, attach it on the
new side by $h_i\circ a_i$. Otherwise retain that labelled boundary and
its transported collar. Denote the resulting capped factor by $F_i$;
it is $V_i$ when unfilled, and a union of two solid tori when filled.
The attachment locus is disjoint from the sphere collar and cap balls.
Both orders of construction therefore use the same pieces and the same
identifications on disjoint boundary components. This gives a relative
identification of the two quotients, extending the specified collar maps.
Consequently the filled extended block is the punctured reconstruction
of $F_1\#F_2$. This is the local commuting argument; it does not invoke
the general bundle formula or identify manifolds from their groups.

For an exposed pairing $\psi:T_i\to T_j$, the replacement map is
\[
             \psi'=h_j\circ\psi\circ h_i^{-1}.
\]
Use the identity map on an unchanged outside port. This formula also
applies when $T_1$ and $T_2$ are paired to one another. For an unpaired
external port, compose the recorded boundary marking with $h_i$.
These equalities, not just preservation of boundary counts, are the
required reconstruction data.

Choose a meridian--longitude basis $(\mu_i,\lambda_i)$ on $\partial V_i$
and express a transported attaching meridian as
\[
       m_i'=q_i\mu_i+p_i\lambda_i,\qquad \gcd(p_i,q_i)=1.
\]
After matching the orientation and basis conventions, the filled factor
is Martelli's $L(p_i,q_i)$ by Definition 10.1.4, printed 307/PDF 315.
Keep the actual oriented filling until that comparison is supplied.
The longitude is a choice, not the old section by fiat. If $p_i=0$,
primitivity gives $q_i=\pm1$ and the factor is $S^2\times S^1$.
If $|p_i|=1$, set $\lambda_i'=\lambda_i+p_iq_i\mu_i$; then
$m_i'=p_i\lambda_i'$. This basis change extends by a solid-torus disk
twist, so the factor is $S^3$ by Proposition 10.1.6. The explicit disk
twist and the noncanonical longitude are described immediately after
that proposition. These recognitions require diffeomorphisms with
orientation data, not only the cyclic fundamental group calculation.
In particular a second old fiber filling becomes a meridian filling;
it must not be processed as a transverse Seifert extension.

\paragraph{Separating locally does not mean separating globally.}
Let $M$ be the connected oriented carrier obtained by all remaining
boundary pairings. Perform the replacement with those transported maps
and let $N$ be the capped cut $\widehat{M\mid S}$. The two caps stay as
marked balls in $N$. Although $S$ separates the isolated block, the
outside gluings can connect its two sides. Determine this using the
incidence graph of the replacement: vertices are the connected pieces
$F_1,F_2$ and unchanged outside blocks, and edges are the retained torus
pairings. A free external torus is a port, not a connecting edge.
Because the pieces and the glued tori are connected, paths across the
pieces correspond exactly to paths in this finite incidence graph.
In particular this graph tests whether the two cap marks are in the
same component. No topological recognition oracle is being assumed.

If the cap marks are in different components $N_1,N_2$, then
$M\cong N_1\#N_2$. If they are in the same component, $N$ is connected
and $M\cong N\#(S^2\times S^1)$. To see the latter reconstruction,
join the two marked balls by an embedded arc in the interior of $N$.
A regular neighborhood of the balls and arc is a ball; removing the
two marked ball interiors and identifying their spheres replaces this
ball by a punctured $S^2\times S^1$. The complement is unchanged, so
all external boundary collars are retained. This is the local handle
construction used in Martelli Proposition 9.2.14, printed 279/PDF 287,
and the sphere-pairing discussion of Matveev Corollary 2.4.4, printed
86/PDF 98. We use that construction without assuming that $M$ is prime;
the prime hypothesis in Martelli's proposition is needed for its further
classification conclusion. Cap independence, sphere gluing and relative
isotopy remain inherited-topology interfaces to bind at release.

Two different sphere-product contributions must be tracked separately:
a meridian filling can create one among the $F_i$, and a nonseparating
ambient sphere creates one from the sphere pairing. For example, pair
$T_1$ directly to $T_2$ after filling $T_0$ along its fiber. The capped
replacement is two solid tori joined by the transported pairing, hence
a lens space; the original carrier is that lens space connected-summed
with $S^2\times S^1$. If both $T_i$ are instead external, the result is
two solid-torus factors with their distinct external labels. No case
permits dropping a marked cap or a sphere-product contribution.

\paragraph{A bound for the extended-block phase.}
Keep a state consisting of the active elementary presentations of all
current factors, the accumulated sphere-reconstruction data, and a
separate list of closed two-solid-torus and sphere-product factors.
Let $p$ be the total number of active pants blocks, summed over all
components. A fiber-parallel replacement removes its one pants block
and its adjacent filling tori and adds only the two solid-torus sides
or their closed fillings. Therefore $p'=p-1$. A witnessed consolidation
of an extended block to a solid torus also removes one pants block;
more generally a witnessed cluster containing $r\geq1$ pants blocks
has $p'=p-r$. Retained torus pairings use the transported maps above.
Store closed lens and sphere-product factors without expanding them
back into pants presentations during this phase. The reconstruction
ledger can grow; it is not the decreasing measure.

For any sequence consisting of these witnessed replacements, induction
on the number $k$ of steps gives
\[
                       p_k+k\leq p_0.
\]
The induction uses $p_{k+1}+1\leq p_k$ and adds $k$. Thus this phase has
at most $p_0$ steps, even when sphere cuts increase the number of
connected components or a nonseparating cut leaves that number fixed.
This is a written conditional bound. It neither assumes fixed seam
labels nor requires the original filling-request set of F06.
Re-expanding a Seifert block, inserting pants blocks to represent a
connected sum, or interleaving arbitrary J03 rewrites is outside the
proved phase relation. Such operations require a separate progress proof.

\paragraph{Remaining proof frontier and regression cases.}
The branch now has a written local proof and explicit map formulas;
its geometric Lean implementation remains open. The next subsection
develops the following outstanding prime-phase obligations:
control intersection curves, justify the innermost-disk reduction, and
handle spheres contained in a block and closed exceptional blocks.
A stopped state must then be proved prime or separately classified;
absence of a fiber-parallel attachment alone is insufficient. Only
then can prime uniqueness or direct factor comparison finish R07/J02.
J05 additionally needs the essential-cut terminal theorem, including
A28's compression alternatives. The two bounded phases do not supply it.

Required future geometric cases are: two external labels; one and two
additional fillings; a zero transported longitudinal coefficient;
a unit longitudinal coefficient with a nonzero meridional coefficient;
a direct pairing of the two exposed ports; and an outside path joining
the two cap marks. Reject nonprimitive slopes and any replacement that
loses an external marking, cap mark or pairing conjugation. These are
unrun specifications. The next subsection develops the innermost-circle
and prime-phase terminal argument; unfinished-PC probes remain paused.
The general F04 exercise and all release gates remain open.
% END migrated B049

% BEGIN migrated B050
\section{Innermost circles and the prime-phase terminal argument}
\label{sec:graph-prime-normal-form}

The next step is to justify when the extended-block phase can stop.
The following is a written proof from the inspected topology sources,
with smooth relative implementation still pending. It concerns prime
refinement, not yet the ambient injectivity of all retained tori.
The source record is \code{reference_checks_revision52.md}; L08 and L10
are refined, and L30 records the boundary-disk lemma used below.

\paragraph{Local irreducibility supplies the disk removal.}
Call a connected cut block \emph{good} when it is irreducible and every
boundary torus is disk-incompressible in that block. Consider a finite
smooth gluing of good compact oriented blocks along paired boundary
tori, with some boundary tori possibly left free. Each port is used at
most once; a pair of distinct ports may belong to the same block.
Then each connected glued carrier is irreducible.

Here is the innermost-circle proof, including its local hypotheses.
Make a sphere $\Sigma$ transverse to the finite seam system $\mathcal T$;
there are finitely many intersection circles. If there are any, choose
one innermost on $\Sigma$, bounding a disk $D$ whose interior misses
$\mathcal T$. In the cut carrier, $D$ lies in one block $B$ and its
boundary is on a particular copy $T$ of a seam. Incompressibility of
$T$ in $B$ gives a disk $E\subset T$ with $\partial E=\partial D$.
Irreducibility of $B$ supplies a boundary parallelism between $D$ and
$E$: push $D\cup E$ slightly inward to a sphere, use its ball, and
restore the boundary collar. This is the disk argument in Martelli
Proposition 9.2.30, printed 287--288/PDF 295--296.

Push the sphere across that ball and slightly beyond $E$, as in
Hatcher Section 1.2, fact (5), printed 11/PDF 12. The isotopy removes
$\partial D$ and any other intersection circles lying in $E$, and
creates no new seam intersections. The disk $E$ need not have interior
disjoint from $\Sigma$; one moves the sphere by the ball-supported
isotopy, rather than asserting that a replacement disk alone is
embedded while all other portions stay fixed. The ball is in the cut
block and meets its boundary only in the specified disk after restoring
the collar. Thus the support avoids all other boundary components and
can be chosen away from the external boundary collars.

The number of circles strictly decreases. Eventually the sphere is
contained in a single cut block, where it bounds a ball by that block's
irreducibility. Undo the ambient isotopies to obtain a ball bounded by
the original sphere. The zero-intersection case is exactly this last
step. Working with separate boundary copies also covers self-pairings.
There is no assumption that the ambient carrier is already irreducible
or that its seams are already incompressible. Hatcher's stated fact (5)
has ambient incompressibility as a hypothesis; the argument here uses
its local disk-removal construction with the stronger supplied
block-boundary hypotheses. This distinction prevents circular use of
that stated theorem.

\paragraph{Why the nonsolid open Seifert blocks are good.}
Hatcher Proposition 1.12, printed 18/PDF 19, says that a compact connected
Seifert manifold is irreducible apart from two sphere bundles and
$\mathbb{RP}^3\#\mathbb{RP}^3$. In the oriented category the twisted
sphere bundle is excluded. The remaining exceptions are closed, so a
Seifert block with nonempty boundary is irreducible. Its proof modifies
the horizontal/vertical argument of Proposition 1.11 using surgery on
an essential sphere; it does not assume ambient irreducibility in
order to establish this conclusion. Those proof passages were reopened.

For a connected compact orientable irreducible manifold with a torus
boundary component, an essential disk on that component forces the
\emph{whole manifold} to be a solid torus. This is the second part of
Martelli Proposition 9.2.30. A compression of that boundary torus gives
a sphere; the interior-side ball and the reversed one-handle recover
the solid torus. Consequently an open Seifert block which is not a
solid torus has incompressible boundary and is good. This shorter route
uses the general boundary-disk lemma; it need not import the full
classification of essential surfaces in Martelli Proposition 10.4.9.
Corollary 10.4.10, printed 337--338/PDF 345--346, independently matches
these Seifert exceptions, but its broader proof is not the chosen
dependency for this step.

A boundary torus of a cut block is different from an arbitrary internal
torus in an irreducible carrier. The boundary-disk lemma does not remove
Hatcher's ball-contained alternative for an internal compressible torus.
It therefore does not resolve A28 by a change of terminology.

\paragraph{Exhausting the extended-block cases.}
Start with a smooth relative elementary presentation supplied by J01.
Each solid-torus vertex has one boundary port. Attach every such vertex
adjacent to a pants vertex to that pants block; keep all the attaching
maps and any unpaired external ports. This gives the extended blocks.
Distinct pants vertices share none of these solid tori. A component
with no pants vertices is either one solid torus with free boundary or
two solid tori joined along their boundaries. Handle these separately.
A pairing of two remaining ports of one pants block is retained as a
seam, not misclassified as a solid-torus filling.

For each extended pants block $B$, use the following exhaustive cases.
These are classical existence cases with geometric witnesses; this is
not a promised computable solid-torus recognition algorithm.
\begin{enumerate}[label=(\arabic*)]
\item If an attached meridian is fiber-parallel, use the relative
splitting from Section~\ref{sec:extended-pants-splitting}, including its
ambient separating/nonseparating reconstruction. This case takes
priority even if the block has no exposed ports. No essential sphere
needs to have been chosen in advance.
\item Otherwise every filling is transverse and F02 extends the Seifert
fibration. If $B$ has nonempty boundary and is a solid torus, consolidate
it using an actual oriented diffeomorphism and its boundary restriction.
There is exactly one remaining boundary torus; transport its marking
or gluing map. This removes the pants vertex, and the new solid torus
can be regrouped with its neighboring pants vertex. If $B$ is not a
solid torus, the preceding argument makes it good.
\item If all fillings are transverse and $B$ has no exposed ports, it
is an entire closed connected component. The closed Seifert alternative
is: irreducible, $S^2\times S^1$, or
$\mathbb{RP}^3\#\mathbb{RP}^3$. Retain the first as an irreducible
Seifert component, retain the sphere product explicitly, and split the
last into two $\mathbb{RP}^3$ factors with its sphere reconstruction.
This is a sufficient general classification; it does not assert that
every alternative occurs for an orientable pants base. Retire these
components into the completed-factor list with their actual witnesses.
\end{enumerate}

The components without pants vertices are solid tori or lens spaces
by their two-solid-torus construction. A solid torus is irreducible;
for a lens space retain the meridian-filling sphere product and the
$S^3$ identity, and use irreducibility for the remaining lens cases.
Hatcher Proposition 1.6 and its lens-space application, printed
8--9/PDF 9--10, give the covering-space argument. Martelli Proposition
10.1.10, printed 308/PDF 316, identifies the nonzero-coefficient Dehn
filling with a cyclic quotient, using a mirror in the displayed proof.
Irreducibility is unchanged by that mirror; the stored oriented filling
and its boundary maps are not replaced by an unverified oriented identification. The sphere product
is prime; $\mathbb{RP}^3$ is irreducible. These are topological
classifications, not constructions of the endpoint metrics.

If none of the reducible or consolidation cases remains in a connected
component with pants vertices, every extended block is open and good.
The local gluing proof above makes that component irreducible. In
particular, a hypothetical essential sphere in such a component would
contradict the terminal criterion. This supplies a move-availability
argument by contraposition after the exceptional cases have been
separated. It does not require the first innermost disk of an arbitrary
sphere to have essential boundary on a seam. Spheres wholly inside a
closed block are handled by the closed Seifert case, not omitted.

\paragraph{Progress and the resulting prime-phase output.}
Extend the phase state of the preceding subsection by allowing a
completed closed Seifert factor with an irreducibility witness, or its
explicit exceptional splitting, in the completed-factor list. Let $p$
be the total number of pants vertices still active. Fiber-parallel
splitting and open solid-torus consolidation lower $p$ by one. Retiring
a closed extended pants block also lowers $p$ by one, even when its
exceptional splitting adds two completed factors. Regrouping adjacent
solid tori changes no pants count. No completed factor is re-expanded
into active pants blocks.

Thus induction again gives $p_k+k\leq p_0$ for the steps that remove
pants vertices. There are at most $p_0$ such steps. Processing components
with no pants vertices and recording already-good terminal components
are finite bookkeeping operations, not additional strictly-decreasing
pants moves. Storing an irreducible graph component can instead remove
all its remaining pants vertices at once. In either implementation the
finite choice of cases ends with irreducible graph components, explicit
sphere products and $S^3$ identities, with all sphere reconstruction and
external boundary markings retained. This is a written prime-refinement
existence argument from the supplied local interfaces. Its Lean proof
and relative smooth implementations are still open.

To transfer graph structure to a separately specified prime decomposition,
J02/R07 still require the inherited prime-uniqueness comparison or an
explicit comparison of the constructed factors. If that comparison is
only an abstract diffeomorphism, it does not by itself match external
labels and collars. Carry the relative markings through the comparison.
No uniqueness or minimality requirement is added to the GC endpoint.

\paragraph{What remains before essential torus cuts.}
Goodness gives injection into each cut block after R05; this subsection
has proved ambient irreducibility, not the inclusion maps from seams or
external tori into the reassembled component. The next subsection proves injection for the finite gluing, including
self-pairings, and connects it to R06/G05's sphere reconstruction. It
also states the conditional G04 consequence for this normal form. It is not yet permission to
invoke Matveev's canonical-subsystem theorem on the original raw cuts.
A28, J05's full essential-cut contract, G09's metrics and H1--H4 stay open.

Geometric regression specifications now include a sphere meeting a seam
in an inessential circle, additional circles inside the boundary disk,
a sphere missing every seam, a self-pairing of two ports of one good
block, an open solid-torus consolidation, a fully filled closed block,
and the two-summand $\mathbb{RP}^3$ exception. Reject an ambient-irreducibility
premise in the gluing lemma, a missing closed-block branch, a twisted
sphere-bundle output in the oriented category, and a claim of ambient
torus injection based only on blockwise injection. These are unrun
geometric specifications; source and inventory audits are not their proofs.
% END migrated B050

% BEGIN REVISION169 ACTUAL RELATIVE GRAPH REFINEMENT
\section{Actual relative refinement of a smooth raw graph presentation}
\label{sec:actual-relative-graph-refinement}
% BEGIN REVISION171 SHORT GRAPH RUNNING HEAD
\markright{\thesection. Actual relative graph refinement}
% END REVISION171 SHORT GRAPH RUNNING HEAD


This section supplies the existence bindings left in J01, F02 and
AT13/AT16, for the chosen (not previously prescribed) prime factors.
The proofs use the written pants move, terminal criterion, finite gluing
and cap avoidance of Sections~\ref{sec:extended-pants-splitting},
\ref{sec:graph-prime-normal-form}, \ref{sec:gluing-injection-transport}
and~\ref{sec:cap-avoidance-output}. They do not use Matveev's
Proposition 2.4.7 to produce the first essential cuts. Source identities
and the category comparison are in
\code{reference_checks_revision169.md}. All statements here are written
mathematics over the retained smooth surface, collar, isotopy, Loop
Theorem and elementary three-manifold foundations, awaiting release
binding and Lean implementation.

\begin{lemma}[Smooth bundle normalization with actual markings; RG01]
\label{lem:graph-smooth-bundle-normalization}
A smooth circle bundle with oriented total three-manifold over a compact
surface with nonempty boundary has, after cutting the base into disks
and bands, the standard product gluings on untwisted bands and
simultaneous base/fiber reversal on twisted bands. This gives a smooth
bundle isomorphism to the usual oriented circle bundle over that surface.
For an orientable base it is the product. The isomorphism includes its
actual boundary restrictions; it need not realize boundary sections
specified in advance.
\end{lemma}
\begin{proof}
Give the vertical tangent line a smooth metric and normalize the length
of each fiber to one. Oriented arclength reduces each fiber coordinate
change to a rotation, or a reflection followed by a rotation when its
orientation is reversed. The orientation of the total space forces the
fiber reversal character to equal the orientation character of the base.
A disk admits a smooth trivialization. Choose a handle decomposition of
the boundary surface with one disk and finitely many bands. On each
attaching interval a rotation has a smooth real lift. Extend that lift
across a collar of the band end and remove it by a rotation of the band
trivialization. The two ends have disjoint collars, so these corrections
can be made independently. What remains is identity or reflection,
according to the character just described. There are no two-handles:
thus there is no remaining extension obstruction. Equivalently, the
bundle over this surface retracting to a graph has no degree-two Euler
obstruction, including with the orientation local system.

For completeness, the same normalization can be performed before the
arclength reduction: an increasing lift $F:\R\to\R$ of a circle
map satisfies $F(x+1)=F(x)+1$. The interpolation
$(1-u)F(x)+u(x+F(0))$ has strictly positive derivative and descends to a
smooth isotopy to a rotation. Compose a decreasing map with reflection
first. Smooth dependence on an attaching interval follows from a lift
on that interval. Collar cutoffs make the changes stationary near their
ends. Store every resulting bundle map and its inverse. On an external
port with original marking $m$, the model marking is $h\circ m$,
and its reconstruction is $h^{-1}$ there. Hence the resulting map to
the original carrier preserves the original marking exactly. This last
identity is not a claim that arbitrary prescribed section slopes can
all be extended by a single section.
\end{proof}

\begin{lemma}[M\"obius blocks have an elementary smooth presentation; RG02]
\label{lem:graph-mobius-elementary}
The oriented circle bundle over a M\"obius band is a smooth union of a
pants product and two fibered solid tori. Its external torus and collar
have recorded actual maps to the original bundle.
\end{lemma}
\begin{proof}
RG01 identifies the bundle with
\[
 Q=(S^1_x\times S^1_y\times[-1,1])/
 ((x,y,t)\sim(x+\tfrac12,-y,-t)),
\]
where $x,y$ are taken modulo one. The original bundle has fibers in
the $y$ direction and base the M\"obius band with coordinates $(x,t)$.
Instead use the smooth circle action translating $x$. Its quotient is
$(S^1_y\times[-1,1])/((y,t)\sim(-y,-t))$, a disk with exactly two
order-two cone points, at $(0,0)$ and $(1/2,0)$. To see the disk,
cut the annulus along the two invariant radial intervals: the involution
identifies the paired halves, leaving one disk with the two indicated
branched points. Local coordinates at each fixed point are a disk modulo
$z\mapsto-z$. The action on the total space has no fixed point, because
the accompanying translation of $x$ is by $1/2$.

Remove small disjoint invariant disk neighborhoods of those two points.
Their inverse images are smooth solid tori with multiplicity two; the
remaining base is pants and its bundle is a product by RG01. These are
actual embedded regions of $Q$. The cutting tori, orientation, collars
and external marking are transported by the inverse of the RG01 map.
This constructs the relative version of Matveev Example 6.4.14; no
numerical sign convention for its two displayed Seifert coefficients
is substituted for the actual attaching maps.
\end{proof}

\begin{proposition}[Actual elementary presentation from raw cuts; RG03]
\label{prop:graph-actual-elementary-presentation}
A finite smooth raw graph presentation by circle-bundle pieces, with
actual torus gluings and labelled external collars, gives a finite
smooth elementary presentation by solid tori and pants products with
the same oriented collared reconstruction. The same assertion holds
for already supplied smooth Seifert pieces.
\end{proposition}
\begin{proof}
In a Seifert piece remove the finite exceptional-fiber neighborhoods;
in a circle-bundle piece this list is empty. If the remaining base is
closed, first remove a regular disk and retain its inverse image as a
solid torus. Cut off M\"obius bands from a nonorientable base, leaving
an orientable surface. Use a finite smooth surface decomposition into
disks, annuli and pants. An annulus may be replaced by pants together
with a regular disk filling, so only disk and pants pieces remain.
This also treats a sphere, torus, disk and annulus base; no negative
Euler characteristic assumption is being hidden. Apply RG01 to disk
and pants bundles and RG02 to each M\"obius piece. Exceptional-fiber
neighborhoods are already solid tori. Thus every piece has the asserted
finite presentation.

All new cuts are inverse images of embedded circles, or the explicit
cuts of RG02. Choose disjoint smaller collars, retaining a smaller collar
at each original boundary torus. Round only the auxiliary corners in
these local product coordinates. If $h_i$ identifies a piece with its
model, an old port pairing $a_{ij}$ becomes
$h_j a_{ij}h_i^{-1}$; the external marking becomes $h_i m_i$.
Their inverses assemble to a diffeomorphism of the glued models to the
original carrier, equal to its original marking after composition.
A self-pairing still pairs two distinct labelled ports. No seam is yet
claimed essential.

This proof starts from smooth bundles and smooth reconstruction maps,
as delivered by the static construction on the actual smooth carrier.
If raw data are provided only as a topological homeomorphism type, an
additional relative smoothing theorem is required; this proposition
neither assumes nor proves that stronger category conversion. The
actual application uses its smooth cuts and bundle charts, not a
homeomorphism-only statement of Matveev's definition.
\end{proof}

\begin{lemma}[Transverse filling with the actual attaching map; RG04]
\label{lem:graph-actual-transverse-filling}
Let $(s,f)$ be a section/fiber basis on a boundary torus of a smooth
circle bundle. If the primitive attaching meridian is
$m=p s+q f$, $p>0$, the filling extends the Seifert structure smoothly,
with core multiplicity $p$. All other ports and their markings remain
as supplied. The case $p=0$ is assigned to the existing pants splitting.
\end{lemma}
\begin{proof}
Choose integers $a,b$ with $pb-qa=1$ and set
$\mu=m$, $\lambda=a s+b f$. Then $f=-a\mu+p\lambda$.
On $D^2\times S^1$, with meridian/longitude $(\mu,\lambda)$, use
\[
 \theta\cdot(z,w)=(e^{-2\pi i a\theta}z,
                         e^{2\pi i p\theta}w),\qquad\theta\in\R/\mathbb{Z}.
\]
Coprimality makes the action free off the core and gives core stabilizer
of order $p$. Its boundary orbit is precisely the pulled-back fiber
slope. The normalized circle actions agree after choosing a boundary
fiber coordinate. The source attaching diffeomorphism can be matched
as follows: use its induced primitive integral basis, then a smooth
isotopy from its remaining identity mapping class to the identity on
the torus. Extend that isotopy through a collar. This is the retained
smooth torus mapping-class/isotopy interface; no extension of an
arbitrary torus diffeomorphism over a solid torus is asserted. Equivalently,
pull the model filling back by the collar diffeomorphism of the glued
carrier. Store that diffeomorphism, so the output reconstruction uses
the original attaching map. Its support avoids all other ports.
For $p=1$ the action is a regular solid-torus bundle. Notice that the
rotation weight is $-a$, not automatically the coefficient $q$;
identifying these numbers without the Bezout change of basis would
lose the marked filling information.
\end{proof}

\begin{theorem}[Actual prime and good-block data; RG05]
\label{thm:graph-actual-relative-good-blocks}
Let $W$ be a compact connected oriented smooth raw graph carrier with
its finite smooth presentation and external torus collars. There is
an actual finite relative sphere reconstruction whose capped factors
are irreducible graph carriers, sphere products and recorded $S^3$
identities. The irreducible factors with boundary are either solid tori
or finite gluings of good Seifert blocks along essential tori. Closed
factors are either such gluings or closed Seifert/lens exceptions.
All cutting maps, cap balls, port pairings, basepoint paths and external
markings are supplied. If every external torus injects into $W$, no
solid-torus factor with an external port occurs. In this case the
output supplies AT13 and AT16 on these chosen factors.
\end{theorem}
\begin{proof}
Apply RG03. Regroup every degree-one solid-torus vertex with its adjacent
pants vertex; retain self-pairings between two remaining ports as seams.
Use the exact phase of Section~\ref{sec:graph-prime-normal-form}.
A fiber-parallel filling has priority and is replaced by the written
relative pants sphere split, with its two new meridians and conjugated
remaining gluing maps. If every attached meridian is transverse, RG04
makes the extended block Seifert. An open block is irreducible by the
retained Seifert sphere theorem (Hatcher Proposition 1.12). If it has a
compressible boundary torus, the retained boundary-disk argument forces
it to be a solid torus. Classical case distinction supplies either an
actual smooth oriented solid-torus diffeomorphism or the negation of
such an existence statement. In the latter case the block is good;
in the former consolidate and record the diffeomorphism's boundary
restriction. This is existence with witnesses, not an algorithm for
recognizing a triangulation.

A transverse closed extended block is an entire component. Retire it
using the existing closed Seifert irreducibility/exceptional alternative,
including the explicit $\mathbb{RP}^3\#\mathbb{RP}^3$ sphere splitting
if that branch is used. Components without pants are a free solid torus
or two solid tori glued as a lens space; the meridional lens case is a
sphere product, and the $S^3$ identity is retained. These witnesses are
those of the established phase, with their actual oriented maps.
The number of active pants decreases at every split, consolidation or
closed-block retirement. Consequently at most the initial pants count
such operations occur. Recording zero-pants components is finite and
creates no active pants. Every remaining component consists of good
blocks and is irreducible by the written good-block gluing argument.
R05 and the finite gluing lemma give injection of every surviving seam
into this same capped component, including loop edges.

Here is the relative implementation of that induction. Before a later
sphere move in a capped carrier, use AT14 on the proposed splitting
sphere together with the current disjoint seam decoration to put them
off the previously recorded caps, fixing smaller external collars.
Transport all block structures and pairing maps by that one isotopy.
Then the new cut takes place in the old punctured component. Record
the new cap embeddings and the transported old caps in their actual
components. Iterating gives a finite nested sphere reconstruction.
AT15 commutes the disjoint operations; AT20's spanning-tree and cycle
construction flattens it into the chosen connected-sum data, including
one sphere-product contribution for each nonseparating cycle. It also
provides the actual punctured and capped carriers, rather than merely
an equality of lists of prime types. The already written cap-avoidance
construction now moves the final torus decoration off all attachment
caps, and transports its block metrics later by pullback.

For an external port in $R_a$, its actual map to $W$ factors through
$R_a$. Injectivity in $W$ implies injectivity in $R_a$ by cancellation,
and then in $P_a$ by the sphere-cap inclusion isomorphism. A solid torus
cannot have an injective torus boundary, since its meridian dies. Hence
such an external-port factor is excluded under the stated hypothesis.
The other factors have exactly the AT16 good blocks or separately
recorded closed exceptions. The maps into $W$ are
$(j_a)_*(c_a)_*^{-1}$, with the stored basepoint paths; a capped factor
is not asserted to embed in $W$. No comparison to a different prescribed
prime decomposition, and thus no relative prime-uniqueness theorem, is
needed for this existence endpoint.
\end{proof}

\begin{corollary}[A small exhaustive list of terminal Seifert blocks; RG06]
\label{cor:graph-terminal-seifert-list}
The preceding construction may choose all good open blocks to be extended
pants with zero, one or two transverse fillings. After deleting regular
cone points, these are: a pants product; a bundle over an annulus with
one cone point; or a bundle over a disk with two cone points of orders
at least two. The zero-cone annulus is a product. The disk case with
orders $(2,2)$ is the twisted interval bundle over the Klein bottle.
The remaining closed non-lens blocks have base $S^2$ with exactly three
cone points of orders at least two. This list is sufficient for E1;
optional coarsening of compatible fibers is not required.
\end{corollary}
\begin{proof}
Each active extended block has one pants vertex and at most three attached
solid tori. In the open case at most two ports are filled. RG04 fills
a base boundary circle by a disk, introducing one cone of order $p$;
$p=1$ is regular. An open disk base with at most one exceptional fiber
is a solid torus: remove its core fiber to obtain a product annulus
bundle and use the ordinary collar of the model fibered solid torus.
It would have been consolidated. A regular filling of an annular base
therefore gives the product case or a disk with two genuine cones.
For $(2,2)$ the coefficients have odd numerator; changing a section
by twists along arcs to the free boundary reduces them to the RG02
coefficients. These twists are actual diffeomorphisms and their effect
on the free boundary marking is recorded. Hatcher Proposition 2.1's
section-twist proof supplies precisely this boundary allowance.

If a closed three-filled pants block has a regular fiber among its
three cores, absorb it, leaving a fibration over $S^2$ with at most
two exceptional fibers. Split that base into two disks each containing
at most one exceptional point. Their inverse images are solid tori,
so this is an actual two-solid-torus lens presentation. Otherwise there
are exactly three orders at least two. The optional closed exceptional
split in RG05 can consequently be bypassed by this more specific
classification and the metrics below. No general classification of a
nonorientable closed base or of an uncut torus semibundle is required:
RG02 and the retained seams have already dealt with those possibilities.
\end{proof}
% END REVISION169 ACTUAL RELATIVE GRAPH REFINEMENT

\chapter{Thurston geometries and model realization}
\label{ch:model-realization}

\paragraph{Scope.} Model witnesses on the actual cut-piece interiors under E1.

\paragraph{Planned sections.}
\begin{enumerate}[leftmargin=*]
\item Eight tags.
\item Completeness and volume.
\item Seifert cases.
\item Product chains.
\item Nil and Sol.
\item Twisted bundles.
\item Realization.
\end{enumerate}

\paragraph{Existing material and present task.} M1d, G06-G10 and the adopted endpoint policy. Detail model-by-model metric existence and transport without promoting policy decisions to proofs.

% BEGIN migrated B020
\section{M1d model geometries and geometric-piece witness}
\label{sec:m1d-model-geometry}

M1d attaches a geometric witness to a \code{CutPiece}. It defines the target
language of the geometrization theorem but does not prove that Ricci flow
produces the witnesses. The analytic production of hyperbolic witnesses and
the conversion of collapsed regions to graph-manifold witnesses remain later
work packages.

\subsection{The eight model tags}

The project fixes one finite model type with the following eight constructors:

\setlength{\tablebodywidth}{\dimexpr\textwidth-4\tabcolsep\relax}
\begin{longtable}{@{}P{0.303371\tablebodywidth}P{0.471910\tablebodywidth}P{0.224719\tablebodywidth}@{}}
\textbf{Tag} & \textbf{Normalized model} & \textbf{Typical endpoint use}\\\hline
\endhead
\code{Spherical} & $\Sph^3$ with its standard constant-positive-curvature metric. &
Spherical space forms and spherical prime factors.\\
\code{Euclidean} & $\R^3$ with its flat metric. & Flat manifolds and torus
bundles with Euclidean structure.\\
\code{Hyperbolic} & $\mathbb{H}^3$ with constant sectional curvature $-1$. & Thick
regions and finite-volume hyperbolic pieces.\\
\code{SphericalProduct} & $\Sph^2\times\R$ with the product metric. &
$\Sph^2\times\Sph^1$ and related quotients.\\
\code{HyperbolicProduct} & $\mathbb{H}^2\times\R$ with the product metric. &
Seifert pieces with hyperbolic base orbifold.\\
\code{TildeSL2} & $\widetilde{\mathrm{SL}}_2(\R)$ with the standard left-invariant
geometry. & Seifert pieces with the corresponding non-product geometry.\\
\code{Nil} & The Heisenberg/Nil geometry with its normalized left-invariant metric.
& Nil-manifold and Nil-Seifert pieces.\\
\code{Sol} & Sol geometry with its normalized solvable-group metric. & Sol
manifolds and hyperbolic torus bundles.\\
\end{longtable}

The tags are mathematical data, not strings. Each tag has a model carrier, a
normalized metric, an orientation convention, and a local-isometry interface.
Alternative normalizations may be used internally only after a proved
isometry or constant-rescaling equivalence is supplied.

\subsection{Common model interface}

\code{ModelGeometry G} contains the common data for a tag $G$:

\begin{itemize}[leftmargin=*]
\item a connected oriented smooth 3-dimensional model carrier;
\item a normalized Riemannian metric and its completeness proof;
\item a local-isometry predicate for maps from an open piece into the model;
\item a proof that the model metric is locally homogeneous; and
\item the model-specific volume and orientation conventions.
\end{itemize}

The first formal definition of local homogeneity is chart-based: every point
has a neighborhood and a smooth local isometry carrying it
to a neighborhood of a fixed model basepoint. A later group-action API may
prove an equivalent formulation, but it is not required to state the public
certificate.

The common interface does not assert that every geometric piece is globally a
quotient by a discrete subgroup. That quotient theorem is model-specific and
would make the common endpoint depend on eight unrelated group-action APIs.
For the first target, a local-isometry atlas, completeness, the conditional
hyperbolic volume proof and appropriate end data are the authoritative
witness. Nonhyperbolic volume proofs are separate properties of that metric.

\subsection{Geometric-piece certificate}

\code{GeometricPieceCertificate C} for a \code{CutPiece C} has these fields:

\setlength{\tablebodywidth}{\dimexpr\textwidth-4\tabcolsep\relax}
\begin{longtable}{@{}P{0.318182\tablebodywidth}P{0.454545\tablebodywidth}P{0.227273\tablebodywidth}@{}}
\textbf{Field} & \textbf{Required content} & \textbf{Audit question}\\\hline
\endhead
Model tag & One of the eight finite model constructors and its
\code{ModelGeometry} data. & Is the normalization fixed and compatible with the
reference source?\\
Interior metric & A smooth Riemannian metric on the open interior of $C$. &
Does the metric live on the same interior object supplied by M1c?\\
Completeness & The interior metric space is complete; geodesic completeness
may be used via an applicable equivalence theorem. & Which completeness
API and equivalence hypotheses are imported?\\
Hyperbolic volume & If the model tag is hyperbolic, the volume of this
interior metric is finite. & Is the proof about this metric and carrier?\\
Additional volume & Nonhyperbolic finite volume may be proved separately;
it is not a universal field. & Are finiteness and infinitude claims kept
separate from completeness?\\
Local homogeneity & A local-isometry atlas from the interior to the selected
model. & Do the charts realize this metric and model, independently of orientation signs?\\
End/boundary data & For each boundary side, a collar or end certificate
compatible with the model. Hyperbolic pieces additionally carry cusp data. &
Does the end datum match the M1c boundary label without truncating the
complete metric?\\
\end{longtable}

The phrase ``complete metric on the interior'' is intentional. For a compact
core with torus boundary, the interior is noncompact and its complete
finite-volume hyperbolic metric contains the cusp ends. A finite truncation is
stored as topological/collar data, not substituted for the complete metric.
More precisely, for a full hyperbolic manifold $(H,g_H)$ and compact
cut carrier $C$, store an end-compatible orientation-preserving
$\psi:\operatorname{int}(C)\to H$ that is a diffeomorphism, and use
$g_C=\psi^*g_H$. This makes $\psi$ an isometry for the pulled-back metric,
so completeness and finite volume transfer to that same interior.
The core embedding into $H$ or a late slice is a different map.

In cusp coordinates $g_H=du^2+e^{-2u}g_{T^2}$, a sequence approaching a
finite truncation level from the retained side is Cauchy for the restricted
metric but has no limit in the open truncated core. This directly rejects
restriction as the complete metric constructor. The cusp formula and
proof are in Benedetti--Petronio Proposition D.3.12, printed 151--152/PDF
162--163; the interior compactification is Martelli Corollary 4.2.18,
printed 121--122/PDF 129--130. The formal pullback, end-label and volume
transports remain implementation obligations.

\subsection{Model-specific obligations}

M1d proves only the common interface and elementary standard-model examples.
The model-specific obligations are recorded for later modules:

\begin{enumerate}[label=\textbf{G\arabic*:},leftmargin=*]
\item standard metrics satisfy completeness, local homogeneity and the
  volume requirement for their selected model;
\item local metric charts compose with actual interior diffeomorphisms;
  orientation and boundary-collar compatibility remain separate
  reconstruction obligations;
\item quotient or lattice constructions, when used, produce the common
  witness without changing the model tag; and
\item model-specific classifications of Seifert, Sol, and hyperbolic pieces
  are kept separate from the common certificate constructor.
\end{enumerate}

The hyperbolic module will later add finite-volume cusp coordinates, Margulis
truncations, and rigidity. The graph-manifold module will later add the
Seifert, product, and Sol conversions. Neither module may redefine
\code{GeometricPieceCertificate}.

\subsection{M1d acceptance tests}

The common geometry layer is accepted when it proves:

\begin{itemize}[leftmargin=*]
\item a standard spherical certificate for $\Sph^3$;
\item a flat Euclidean certificate for $\mathbb{T}^3$;
\item a product certificate for $\Sph^2\times\Sph^1$;
\item a closed finite-volume hyperbolic certificate for a chosen closed
  hyperbolic example;
\item complete flat interior witnesses for $T^2\times I$ and the orientable
  twisted interval bundle over the Klein bottle, despite their infinite
  volumes;
\item rejection of a finite cusp-truncation restriction as a complete
  hyperbolic interior metric;
\item rejection of an incomplete metric for every tag; and
\item rejection of a hyperbolic candidate lacking the required finite-volume
  proof, with that proof tied to the stored metric.
\end{itemize}

The examples are API tests, not applications of the geometrization theorem.
They do not use the PC adapter, finite-time extinction, surgery, or a
Ricci-flow limit.

\subsection{Handoff to M1e}

M1e will define the common \code{ThickThinCollapsePackage} fields that point to
these geometric-piece witnesses. The hyperbolic branch may supply a
\code{GeometricPieceCertificate} directly after proving its limit-to-piece and cusp
theorems; the thin branch may supply one after graph-manifold conversion.
Both branches consume the same M1c cut-piece and M1d witness interfaces.

The field and theorem audit for this layer is recorded in
\code{m1_geometry_api.csv}.
% END migrated B020

\paragraph{Current hyperbolic status.}
The historical planning paragraph above is superseded for cusp coordinates
and Margulis truncations by the written development in
Section~\ref{margulis:section}. Rigidity remains under development;
accepted inherited bindings and Lean implementation remain pending.


% BEGIN migrated B043
\section{G06--G08: product chains, monodromy and twisted bundles}

G06 removes redundant parallel cuts only with a reconstruction witness.
For a product $T^2\times I$ between two pieces, compose the boundary
gluing maps, update the side involution and record the induced orientations.
A chain that closes into a cycle carries monodromy; it cannot be deleted as
an empty region. Exceptional Seifert presentations of products and twisted
interval bundles also prevent assuming that every gluing preserves a
chosen fiber slope.

G07 handles \emph{closed} torus bundles and semibundles separately from
open interiors of cut pieces. At a minimum the eventual mathematical
regressions must distinguish these orientation-preserving torus-bundle
monodromies:
\[
\begin{array}{c|c}
 A\in\mathrm{SL}_2(\mathbb Z)&\text{model geometry}\\\hline
 \begin{pmatrix}1&0\\0&1\end{pmatrix}&\mathbb E^3\\[3pt]
 \begin{pmatrix}1&1\\0&1\end{pmatrix}&\mathrm{Nil}\\[3pt]
 \begin{pmatrix}2&1\\1&1\end{pmatrix}&\mathrm{Sol}
\end{array}
\]
Scott's Nil discussion gives the nonzero unipotent example; the last
matrix has determinant one and trace three and lies in the hyperbolic
monodromy case of his Sol construction. These are source-guided test
specifications, not executed Lean regressions. A semibundle needs its
finite torus-bundle cover and the relevant geometric descent argument;
its name alone is not a Sol certificate.

The parenthetical flat-or-Sol description of closed unions of product and
twisted interval-bundle pieces in the archived MT12 Proposition 5.4 must
not be copied as an exhaustive classifier: identifying the ends of one
product by the middle matrix above gives a Nil example. Use Scott's
classification and retain Nil explicitly. This comparison does not change
the global choice of Morgan--Tian as assembly guide.

For closed Seifert pieces, Scott's Theorem 5.3 classifies the geometry using
the base orbifold Euler characteristic and the Seifert Euler number. Its
closed-manifold hypothesis is essential here. The boundary case is handled
separately by MT12 Lemma 5.1, which singles out the product and the
orientable twisted interval bundle over the Klein bottle.

G08 records that twisted bundle with its single torus boundary and its
one-sided Klein-bottle zero section. In the finite-volume reconstruction
of MT12 Proposition 5.4, a boundary torus is sometimes replaced by that
zero section. Cutting a one-sided surface has boundary incidence governed
by its orientation double cover; it is not an ordinary two-sided torus
cut. No such cut may be passed to the current \code{TorusCutSystem}.
% END migrated B043

% BEGIN migrated B044
\section{G10: torus-only endpoint convention adopted}
\label{sec:endpoint-compatibility-gate}

\textbf{Revision 47 decision E1.} Retain the prime decomposition and
incompressible torus cuts. Every cut-piece interior must carry a complete
metric locally isometric to its selected Thurston model. Finite volume is
required when that model tag is \code{Hyperbolic}; on other tags it is a
separate property of the same metric, supplied when proved. This corrects
the universal finite-volume requirement in C4 and the common geometry API.
It does not add Klein-bottle cuts or change the M1b/M1c reconstruction.

\paragraph{Source comparison.}
Scott, printed page 403/PDF page 3, defines geometric structure using
completeness and local homogeneity and states a sphere/torus decomposition
formulation. His historical statement is a target formulation, not a proof
of Geometrization. The archived KL notes' Appendix I, printed
2854--2855/PDF 268--269, specifies finite-volume hyperbolic complements of
graph regions and ambient incompressibility of their boundary tori.
It then identifies the prime/Seifert refinement result as Matveev
Proposition 2.4.7. The archived notes explicitly say that they are the
February 20, 2013 revision of the 2008 article; the locators here refer to
that copy. The analytic exposition role of KL2008 is unchanged.

Lott's Section 2, printed 490--491/PDF 6--7, gives a torus-only formulation
with compact geometric quotients, noncompact finite-volume hyperbolic or
$\mathbb H^2\times\R$ quotients, and a twisted real-line bundle over the
Klein bottle. Remark 2.2 explains the alternative Klein-bottle cuts
\cite{Lott2010DimensionalReduction}. We use this statement solely to compare
endpoint conventions, not as an analytic input or a proof of our theorem.
The book-level assembly spine remains Morgan--Tian 2014.

The revised existential target does not require minimal cuts or uniqueness.
It permits redundant product cuts with complete flat interiors. The stronger
source decomposition implies this target by forgetting minimality and
retaining its metrics and hyperbolic volume proofs. This is a
specification-level implication; its geometric formalization remains to be
implemented. No inference from completeness alone to finite volume is used.

\paragraph{Explicit product and twisted-bundle witnesses.}
The following constructions explain the exceptional metrics without invoking
an additional Ricci-flow theorem. Write $T^2=\R^2/\mathbb Z^2$ with its unit
flat metric. The product $T^2\times\R$ is complete and flat, with
infinite volume. It identifies with the interior of $T^2\times[-1,1]$
by the diffeomorphism $(x,s)\mapsto(x,\operatorname{arctanh}(s))$;
the metric on the compact interval itself is not the complete interior
metric.

For the twisted case define
\[
 \tau([x,y])=[x+\tfrac12,-y],\qquad
 \sigma([x,y],t)=(\tau([x,y]),-t).
\]
The map $\tau$ is a free orientation-reversing isometric involution of
the square flat torus, and its quotient is a Klein bottle. Thus $\sigma$
is free and isometric, has order two, and preserves the product orientation:
its derivative reverses both the $y$ and $t$ directions. The quotient
\[
 Q=(T^2\times\R)/\langle\sigma\rangle
\]
is an oriented complete flat three-manifold. Completeness follows by metric
descent through the finite isometric covering; local quotient charts give
the Euclidean model witness. The odd interval diffeomorphism above is
equivariant, so $Q$ is the interior of the compact orientable twisted
interval bundle
\[
 W=(T^2\times[-1,1])/\langle\sigma\rangle.
\]
Its two covering boundary tori are exchanged, leaving one torus boundary.
The zero section is a Klein bottle. Retraction onto that section identifies
the boundary inclusion with its orientation double cover, so the boundary
fundamental-group map is injective with index two image. This proves
injectivity into $W$ in the mathematical construction; injectivity into an
ambient prime still needs G05.

The product cover has infinite volume, and its degree is two, so the
quotient flat metric also has infinite volume. These witnesses therefore
satisfy the revised Euclidean piece requirement and fail a demand for a
finite-volume proof of that same metric. They do not assert that every
metric on the underlying manifolds has infinite volume. The construction
and covering arguments are written mathematics here, not kernel-checked
Lean. Their eventual implementation must use the inherited quotient,
metric, completeness and volume APIs after release audit.

\paragraph{Exact certificate and consumer changes.}
For a compact cut carrier $C$, a certificate records a tag $G$, an interior
metric $g$, completeness, compatible local model charts and boundary/end
data. Its additional volume field has the mathematical type
\[
 G=\texttt{Hyperbolic}\quad\Longrightarrow\quad
                 \Vol_g(C^\circ)<\infty.
\]
Here equality is equality of model constructors. The implication is
vacuous for another constructor; it does not infer finiteness for that
constructor. Any separately supplied volume proof must refer to this same
$g$ and this same open interior. No Boolean ``volume checked'' flag or
unrelated finite-volume metric is an acceptable witness.
\code{HasThurstonGeometry} means existence of a certificate with this
policy. It includes the hyperbolic volume obligation, rather than merely
the existence of some local metric.

The eight-row \code{endpoint_volume_policy.csv} fixes the requirement by
model tag. The common certificate still requires completeness on all eight
models. Thick conversion supplies its already required finite-volume
hyperbolic witness; thin conversion supplies the appropriate nonhyperbolic
model and its proofs; exceptional-prime conversion follows the same rule.
Transport moves the metric, tag, charts and conditional volume proof
together. Relabelling a hyperbolic metric as Euclidean cannot bypass the
volume field because the local-isometry model witness must also type-check.
Cusp/core data and all hyperbolic limit hypotheses remain unchanged.

The M1d regressions now distinguish an incomplete metric (always rejected),
an infinite-volume hyperbolic candidate (rejected), and the two complete
infinite-volume flat interiors above (permitted once their geometric and
boundary proofs are supplied). Nonhyperbolic finite volume is neither
required nor forbidden. The geometric tests are specifications behind K1;
no new geometric Lean test has been executed.

\paragraph{Scope of closure and remaining work.}
A25 is closed as a \emph{contract mismatch}: E1 is now reflected in C4,
\code{HasThurstonGeometry}, \code{GeometricPieceCertificate}, the M1d field
and negative-test tables, assembly consumers and the audit. G10 is a
source-reviewed specification decision, not a geometric theorem. G01--G09,
H1--H4, K1 and the PC release gate are not discharged by that decision.
The universal finite-volume/Klein-cut alternative is not selected.
% END migrated B044

% BEGIN migrated B045
\section{G09: geometric conversion remains unproved}

G09 consumes the topology, transport, product-chain, bundle-classification
and twisted-bundle proofs G04--G08, together with the adopted G10 policy.
It constructs \code{ThinRegionToGeometricPieces} with the final internal
tori, model witnesses and carrier/label identifications. The required
orientation-compatible gluing and ambient incompressibility remain
explicit. A raw graph presentation must not store this final certificate
as a premise.

KL Appendix I cites the 2003 edition of Matveev, Chapter 2.4,
Propositions 2.4.2, 2.4.3 and 2.4.7. Revision 50 independently reads those
numbered passages in the user-supplied 2007 second edition. The 2003 text
remains uninspected; the second edition has its own reference and manifest
entry. Its conditional canonical-subsystem statement must not be promoted
to an initial essential-cut existence theorem. Section~\ref{sec:matveev-graph-rewrites}
records the exact source scope and the remaining construction work.
Chapter~\ref{ch:three-manifold-topology} makes the independently verified prime and torus
interfaces concrete while isolating the unresolved graph-class conversion.
G02/G04 remain unproved; PC-specific declaration probes remain paused.
% END migrated B045

\section{Metric charts and oriented reconstruction are separate}
\label{sec:model-atlas-orientation-audit}

The local-model clause in \code{GeometricPieceCertificate} asks for a
smooth local isometry to the selected normalized model; it imposes no
orientation sign on that local isometry. The compact cut carrier, its
prime, its labelled boundary sides, and every reconstruction map retain
their prescribed orientations. A model may also retain its convenient
reference orientation. Metric charts need not identify those two choices.
This interpretation supersedes the earlier orientation-compatible atlas
wording; it changes neither the eight model tags nor E1.

This distinction is necessary for a uniform transport API. Scott's
archived article, printed462/PDF62 and printed465/PDF65, shows that every
isometry of the standard $\widetilde{\mathrm{SL}}_2$ geometry preserves
orientation. The analogous statement and argument for Nil are on
printed467--469/PDF67--69. Thus one must not repair an arbitrary chart's
orientation by assuming a model reflection. Scott's metric definition
on printed403/PDF3 does not impose the extra oriented-chart condition.
The known author errata URL was checked but was inaccessible in this
revision; this is a correction of the project's extra condition, not a
claim that Scott's text has no errata.

If $e:C^\circ\to B^\circ$ is any smooth diffeomorphism and $g=e^*h$, a
local $h$-isometry $\varphi$ gives the local $g$-isometry
$\varphi\circ e$. Its metric proof is independent of the sign of $e$.
Completeness and volume use the same actual metric transport as AT26.
When the compact comparison is used in oriented reconstruction, its
separate orientation and marking proofs are still mandatory. In
particular this correction does not permit an orientation-reversing
boundary gluing to be replaced by an orientation-preserving one.

An implementation regression should pull back the standard Nil metric
on $\mathbb R^3$ by an orientation-reversing smooth diffeomorphism. The
composed local metric charts must be accepted without a model-reflection
premise. This tests transport of the given metric; it does not assert a
new classification of all geometric metrics on the underlying manifold.

\subsection{Retaining elementary cuts avoids unnecessary monodromy work}
\label{sec:retained-bundle-cut-route}

E1 allows a complete infinite-volume Euclidean interior and does not
require minimal tori. Accordingly the default assembly may retain
certified product and twisted-bundle cuts. Classification of a closed
torus bundle by its monodromy is an optional alternative constructor,
not a prerequisite for this retained-cut constructor. The following
arguments start with actual smooth presentations on a supplied prime
carrier. They do not recognize an arbitrary graph manifold as such a
presentation and do not discharge G04's essential-cut existence.

\begin{proposition}[Supplied torus-bundle presentation]
\label{prop:retained-torus-bundle}
Let $P$ be a supplied closed connected oriented prime carrier, with an
actual oriented diffeomorphism to the mapping torus of an
orientation-preserving smooth diffeomorphism $a:T^2\to T^2$ and the
associated smooth cut/collar comparison. Then one fiber torus gives an
E1 decomposition with one compact cut piece $T^2\times[0,1]$, two
distinct paired sides, and a complete Euclidean metric on its interior.
\end{proposition}
\begin{proof}
The cover $T^2\times\mathbb R$ of the mapping torus is obtained from
the action generated by $(x,s)\mapsto(a(x),s+1)$, up to the chosen
mapping-torus sign convention. Projection to $\mathbb R$ proves proper
discontinuity; every nonidentity deck transformation changes that
coordinate. The fiber inclusion into this cover induces an isomorphism
on fundamental groups because projection to $T^2$ is a deformation
retraction, and the covering map induces an injection. Their composite
is the actual fiber inclusion, with the chosen basepoint comparisons.
The supplied cut comparison identifies the compact cut with
$T^2\times[0,1]$ and retains the two sides even though they belong to
the same piece. Use the complete flat $T^2\times\mathbb R$ metric
already constructed under E1, transported to the whole open interior.
The boundary pairing reconstructs the given mapping torus. It need not
be an isometry of this metric. Transport all fields to $P$ along the
given marked diffeomorphism.
\end{proof}
No decision about the trace, eigenvalues, periodicity or Jordan form of
the monodromy is used. In particular identity, nontrivial unipotent and hyperbolic monodromy
examples all pass this same retained-cut constructor. The geometry of the
uncut carrier is a separate classification question.

\begin{proposition}[Supplied torus-semibundle presentation]
\label{prop:retained-torus-semibundle}
Let a supplied oriented prime carrier $P$ be presented, with actual
smooth collars and an oriented reconstruction, by gluing two copies of
the orientable twisted interval bundle $W$ over the Klein bottle along
their torus boundaries. Its displayed separating torus and these two
cut carriers give an E1 decomposition with Euclidean interior metrics.
\end{proposition}
\begin{proof}
The explicit E1 model for $W$ identifies its boundary fundamental-group
map with the index-two inclusion coming from the orientation double
cover of the Klein bottle. Apply the existing finite gluing injection
lemma to those actual port maps. The displayed torus is therefore
injective in $P$. Retain the two compact pieces and use their already
constructed complete flat interior metrics. The supplied orientation-
reversing boundary pairing and product collars reconstruct $P$.
\end{proof}
The hypotheses contain a genuine presented carrier and prime witness,
not a bare semibundle label or an arbitrary finite torus-bundle cover.
No geometric descent theorem for a finite cover is used here. Neither
proposition removes the closed Seifert realization problems, including
Nil and $\widetilde{\mathrm{SL}}_2$ cases without a supplied essential
torus presentation.

These are project-level conditional arguments using the retained E1
metrics, actual cut interfaces and covering/gluing inputs. The source
comparison is Morgan--Tian's archived2012 Lemma5.1 and Proposition5.4,
printed54--56/PDF58--60, and Scott's torus-bundle discussion,
printed469--472/PDF69--72. The stronger closed-model classification is
retained in G07 for callers that request an uncut model; its numerical
case analysis is not an input to the two propositions above.


% BEGIN REVISION170 ACTUAL SEIFERT MODEL METRICS
\section{Complete metrics on the actual terminal carriers}
\label{sec:actual-terminal-model-metrics}

RG06 reduces the model problem to the carriers below. The construction
uses the actual smooth locally free circle action on each extended-pants
block. A metric on a finite cover is never descended by an unproved
invariance assertion. Scott's Theorem 5.3 and its supplied proof,
printed 477--480/PDF 77--80, are comparisons for the closed cases;
Martelli Version4 Sections 6.2 and 12.4--12.6 supply the polygon and
model formulas. The following connection construction provides the
missing witnesses on our carriers, with its own normalization checks.
\code{reference_checks_revision170.md} records the newly checked Lee
correction required by the open-base primitive argument.

\begin{lemma}[Metrics on the required base orbifolds; GM01]
\label{lem:graph-terminal-base-metrics}
Give a filled port order $p\ge2$ and an unfilled port order $\infty$.
Every negative-characteristic base in RG06 has a complete finite-area
hyperbolic orbifold metric on its entire interior, with each free port
a cusp. A closed triangle base has a constant-curvature metric of
curvature $4$, $0$ or $-1$ according as
\[
                 \chi=-1+\frac1{p_1}+\frac1{p_2}+\frac1{p_3}
\]
is positive, zero or negative. The identifications retain the labels of
cone points and ends.
\end{lemma}
\begin{proof}
Double a geodesic triangle of angles $\pi/p_i$, interpreting $\pi/\infty=0$.
For negative sum defect use the hyperbolic triangle; ideal vertices are
omitted. For zero defect and three finite orders use a Euclidean triangle.
For positive defect use a spherical triangle, then rescale its metric
from curvature one to four. The elementary constant-curvature triangle
existence facts are the retained two-dimensional model geometry, as in
Martelli Exercises 3.3.2--3 and the constructions in 6.2.10--12; the
exercises are not recorded as already elaborated Lean proofs.
Reflection in a side identifies the two metric germs smoothly. At a
finite vertex the doubled angle is $2\pi/p_i$, so a $p_i$-fold disk
chart has an ordinary smooth constant-curvature metric. At an ideal
vertex a horoball description gives $dr^2+e^{-2r}dx^2$ with $x$ periodic.
Its radial length is infinite and its area is finite. Thus doubling is
complete, including its finite cone points in the orbifold sense, and
has finite area. The closed doubles are compact. The polygon area
formula also gives $\kappa A=2\pi\chi$ when $\kappa\ne0$.

For an open extended pants block there are three vertices counting cone
points and punctures. Negative characteristic is exactly the strict
angle-sum condition. The omitted equality case is the disk with orders
$(2,2)$ and one free port; it is treated by the existing flat twisted
model, not by a nonexistent Euclidean triangle with an ideal vertex.
A regular filling is deleted first, as in RG06. Smooth surface
classification with labelled disks and collars identifies these doubles
with the actual bases. Near a cone use its rotation disk chart; extend
the boundary identifications over the punctured sphere. This is a smooth
orbifold identification, not an assumption that a cone metric is smooth
in the quotient surface coordinate. Pull back by it and keep the actual
port labels.
\end{proof}

\begin{lemma}[Prescribing connection curvature on these bases; GM02]
\label{lem:graph-actual-connection-curvature}
Let $\pi:B^\circ\to O$ be the actual oriented circle action fibration
of an extended-pants block with orientable base. Normalize its regular
fiber parameter to $\R/\mathbb Z$, with generator $X$.
There is an invariant smooth one-form $\eta$ with $\eta(X)=1$ and
$d\eta=\pi^*\beta$ for a smooth orbifold two-form $\beta$.
On an open base, $\eta$ can be chosen flat. On a closed triangle base,
put $J=\int_O\beta$. For any smooth orbifold area form $\omega$ with
$A=\int_O\omega$, it can be chosen so that
\[
                    d\eta=(J/A)\pi^*\omega.
\]
In the oriented pants section convention with meridians
$p_i s_i+q_i f$, $J=-\sum_i q_i/p_i$; changing conventions reverses the
corresponding signs consistently. In particular the zero/nonzero test
uses the complete filling data, including integral twists.
\end{lemma}
\begin{proof}
RG04 produces the actions on the filling tori matching the product action
on pants. Average a smooth metric over this circle action and define
$\eta_0=g(X,\cdot)/g(X,X)$. The vector field $X$ never vanishes, even
on exceptional fibers. Invariance and Cartan's formula give
$\iota_Xd\eta_0=0$. In a finite cyclic local covering
$D^2\times S^1\to(D^2\times S^1)/C_p$ this says that $d\eta_0$
is the pullback of a smooth rotation-invariant disk form. These forms
agree on overlaps and define the orbifold form $\beta_0$.

We need only the following elementary extension of ordinary de Rham
calculus. Near each of the finitely many cone points, the radial
Poincar\'e homotopy gives an invariant primitive of a given two-form;
choose an invariant cutoff and subtract its exterior derivative.
The remaining form vanishes near the cone points. In the underlying
smooth punctured sphere it is now an ordinary smooth form, extended
by zero across those disks. If the underlying surface is open it is
exact, since its ordinary top de Rham group vanishes. If it is the
closed sphere it is exact precisely when its integral vanishes.
Pull an ordinary primitive back into the cone charts and add back the
cutoff primitives. This proves the orbifold statements used here.
Orbifold integration divides each disk-cover integral by $p$; Stokes'
theorem makes all the subtracted exact integrals zero in the closed case.
The ordinary noncompact result uses Lee's corrected Theorem17.32 proof
(author corrections dated January3,2026, PDF9--10), not the archived
proof's potentially disconnected exhaustion bands.

On an open base choose $\alpha$ with $d\alpha=\beta_0$ and set
$\eta=\eta_0-\pi^*\alpha$. On a closed base the form
$\beta_0-(J/A)\omega$ has zero integral, so the same construction
gives the stated curvature. No fiber-period or local isotropy data
change when a basic one-form is subtracted.

To compute $J$, first make the connection flat in the exceptional
fiber neighborhoods by the same local construction. On the pants
section, write $\eta=d\theta+\alpha$. On a filling boundary the
meridian bounds a disk in that flat neighborhood, so
$p_i\int_{s_i}\eta+q_i=0$. Stokes on pants yields
$J=\sum_i\int_{s_i}\alpha=-\sum_iq_i/p_i$; the removed neighborhoods
have zero curvature integral. Here $s_i$ has the induced pants-boundary
orientation and the regular fiber has period one. This computes the
curvature class on the actual fibration. It does not identify Euler
numbers from different source sign conventions without a translation.
\end{proof}

\begin{proposition}[The normalized model comparison; GM03]
\label{prop:graph-normalized-connection-models}
Use GM01's base metric $h$, area form $\omega$ and GM02's connection.
The following positive definite smooth metrics on the actual total
space are locally isometric to the indicated fixed models:
\[
\begin{array}{c|c|c|c}
\text{base curvature}&\text{condition}&g=\pi^*h+\ell^2\eta^2
                                      &\text{model}\\ \hline
-1&d\eta=0&\ell>0&H^2\times\R\\
0&d\eta=0&\ell>0&\R^3\\
4&J\ne0&\ell=2A/|J|&S^3\text{ (curvature one)}\\
0&J\ne0&\ell=A/|J|&\mathrm{Nil}\\
-1&J\ne0&\ell=A/|J|&\widetilde{\mathrm{SL}}_2\R
\end{array}
\]
The flat-connection hyperbolic row also applies to all negative-characteristic
open blocks. These metrics are complete on their entire interiors.
\end{proposition}
\begin{proof}
On a simply connected regular base chart, write $\eta=d\theta+\alpha$,
and put $u=\ell\theta$. Two connection potentials with the same
exterior derivative differ locally by an exact form; changing $u$ by
its primitive identifies the two metrics. For zero curvature this
is the product metric, with no condition on global holonomy.
For nonzero curvature change the sign of the local fiber coordinate
if needed. Then $\ell d\alpha$ is the required model curvature form.
The model formulas are
\[
 \begin{split}
 g_{\mathrm{Nil}}&=dx^2+dy^2+(du-x\,dy)^2,\\
 g_{\widetilde{\mathrm{SL}}_2\R}
   &=\frac{dx^2+dy^2}{y^2}+(du+dx/y)^2,\qquad y>0.
 \end{split}
\]
Their connection curvatures are respectively $-dx\wedge dy$ and
$dx\wedge dy/y^2$. The first formula follows directly from the
Heisenberg multiplication $(x,y,u)(x',y',u')=(x+x',y+y',u+u'+xy')$;
the second is the lifted unit-tangent-bundle metric in Martelli
Lemma12.6.1. For the spherical row the standard round Hopf fibration
has base curvature four, unit vertical coordinate $u$ of period $2\pi$,
and connection curvature of magnitude $2\omega$. Indeed in coordinates
\[
 (z_1,z_2)=
 (e^{i(u+\phi/2)}\cos(\vartheta/2),
  e^{i(u-\phi/2)}\sin(\vartheta/2))
\]
the round metric is
\[
 \tfrac14(d\vartheta^2+\sin^2\vartheta\,d\phi^2)
       +(du+\tfrac12\cos\vartheta\,d\phi)^2.
\]
Its connection curvature is minus twice the oriented base area form.
This explains both the factor two and the base curvature four in the
table. The metric is not merely an arbitrary Berger deformation of
$S^3$. Changes of signs in local charts do not change the retained
orientation of the actual carrier; E1 does not require oriented model
charts.

At an exceptional fiber perform the same comparison on its finite local
cover $D^2\times S^1$. The lifted metric is invariant under its cyclic
action because $h$ and $\eta$ are. That action is free on the total
space, including the fiber over the disk center: RG04's coprimality
ensures the nontrivial stabilizer rotation has a nontrivial fiber
translation. Hence the quotient has a smooth metric and every total
space point has a neighborhood isometric to a model neighborhood.
There is no cone singularity in the three-manifold metric.

Closed carriers are compact and the metrics are smooth, hence complete.
For an open carrier the projection to the complete base orbifold is
length nonincreasing. A Cauchy sequence projects into a compact base
set, enlarged slightly if necessary. The inverse image of a compact
base set is compact: this follows from finitely many disk charts and
compact circle fibers, including their finite quotients. On that compact
set the smooth total-space metric is uniformly comparable to coordinate
metrics. The sequence therefore converges. This proves metric completeness;
Hopf--Rinow gives the equivalent geodesic assertion. It uses properness
of the actual circle fibration, not boundedness of a chosen connection
potential at a cusp. In the hyperbolic open row its volume is additionally
$\ell\operatorname{Area}(O)<\infty$, although E1 does not require this
for a nonhyperbolic tag.
\end{proof}

\begin{lemma}[The positive triangle case cannot have zero curvature class; GM04]
\label{lem:graph-positive-triangle-euler}
For a closed extended-pants block with three genuine exceptional fibers,
$\chi>0$ implies $J\ne0$. The GM03 rows therefore cover every closed
triangle block.
\end{lemma}
\begin{proof}
The possible ordered triples, up to permutation, are $(2,2,n)$,
$(2,3,3)$, $(2,3,4)$ and $(2,3,5)$. For $(2,2,n)$ the first two
terms $q_1/2+q_2/2$ have integer sum because both numerators are odd,
whereas $q_3/n$ is not an integer. For the other triples use common
denominators $6$, $12$ and $30$. The respective numerators are
$3q_1+2q_2+2q_3$, $6q_1+4q_2+3q_3$, and
$15q_1+10q_2+6q_3$, all odd by primitivity of the even-order coefficient.
They cannot vanish. Integral section twists do not affect these tests.
If $\chi\le0$, either $J=0$ or $J\ne0$ selects exactly one of the
remaining rows of GM03. Bad spherical orbifolds with one cone or two
unequal cones have already been assigned to the lens branch of RG06;
they are not uniformized by GM01.
\end{proof}

\begin{proposition}[Explicit exceptional and flat witnesses; GM05]
\label{prop:graph-lens-product-flat-witnesses}
All product, twisted-interval, lens and sphere-product outputs of RG05--RG06
have complete model metrics on their actual interiors, with their stored
oriented reconstruction data. A free solid-torus factor, when allowed
without an external-injection assumption, also has such a witness.
\end{proposition}
\begin{proof}
For $T^2\times I$ and the twisted interval bundle use the already written
complete flat metrics in this chapter. RG02 and RG06 supply the actual
diffeomorphisms for the latter; record their boundary restrictions.
The interior of a solid torus is diffeomorphic to $\R^2\times S^1$:
use a smooth radial diffeomorphism of the open disk with $\R^2$ and
the identity on the circle. Pull back the product flat metric. No
boundary-injection assertion accompanies this optional branch.

For two solid tori, fix the actual first meridian/longitude $(\mu,\lambda)$
and write the other meridian as $q\mu+p\lambda$. If $p=0$,
primitivity gives $|q|=1$ and the existing disk-double construction is
$S^2\times S^1$; use the quotient of the product $S^2\times\R$ by
an integral translation. If $p\ne0$, change the unoriented meridian
sign so $p>0$. Put $\omega=e^{2\pi i/p}$ and act on the round sphere
in $\mathbb C^2$ by
\[
                       (z,w)\longmapsto(\omega z,\omega^{-q}w).
\]
This action is free: at least one coordinate is nonzero, and
$\gcd(p,q)=1$ forces any stabilizing power to be a multiple of $p$.
It is isometric and the quotient is compact, hence complete.
To identify the carrier, split $S^3$ along $|z|=|w|$. On the quotient
torus its lattice has basis
$v=(1/p,-q/p)$ and $w_0=(0,1)$.
Take $(\mu,\lambda)=(w_0,v)$ on the half whose core is the $z$ circle.
The other half's meridian $(1,0)$ equals $q\mu+p\lambda$.
Both quotient halves are solid tori; for the first the coordinates
$z/|z|$ to the $p$th power and $w(z/|z|)^q$ give a direct trivialization.
The analogous Bezout trivialization works on the other half.
Thus the two actual attaching meridians match. The residual torus
mapping class fixing a meridian is realized by disk twists and the
simultaneous sign changes that extend over a solid torus; an isotopy
extends through a collar. Composing these maps gives the actual smooth
identification with the original filled carrier. Pull back the round
quotient metric along it. Assign the model carrier the orientation
transported by this map; no orientation-reversing self-map of a chiral
lens space is assumed. For $p=1$ this is $S^3$, and for $p=2$ it is the
antipodal quotient $\mathbb{RP}^3$.

This uses the lattice proof of Martelli Proposition10.1.10 while recording
the sign explicitly: his displayed $(\omega z,\omega^q w)$ action
produces the opposite meridian sign before a mirror. Our choice avoids
losing that mirror in the marked identification. Finite group order
alone is never used to recognize the carrier. An $S^3$ identity and
each sphere-product ledger contribution receive these same explicit
metrics.
\end{proof}

\begin{theorem}[Marked graph refinement with complete E1 witnesses; GM06]
\label{thm:graph-marked-geometric-endpoint}
Every actual smooth raw graph carrier has the relative prime/sphere
reconstruction of RG05, essential torus cuts on its chosen irreducible
factors, and complete E1 model metrics on all their resulting interiors.
If its external ports inject, its boundary factors have good-block
refinements supplying AT16, and their metrics supply the graph-region
part of GA06. The construction preserves the external labels, sphere
and torus pairings, orientations and cap assignments. It neither
requires canonical JSJ uniqueness nor changes E1's volume convention.
\end{theorem}
\begin{proof}
Choose the RG03 elementary presentation and run RG05, retaining all
its good-block seams. RG06 gives the exhaustive terminal list. GM01--GM04
put $H^2\times\R$, Euclidean, Nil, spherical or
$\widetilde{\mathrm{SL}}_2\R$ metrics on the listed Seifert interiors;
GM05 treats every product, twisted, lens and sphere-product case.
The action and connection are defined on the actual block, or the metric
is pulled back by the explicitly stored diffeomorphism. Thus the metric
and topological certificates refer to the same carrier, not merely to
isomorphic fundamental groups. Nil has not been omitted when $\chi=0$
and $J\ne0$.

Take one simultaneous cap-avoidance isotopy as in RG05 and transport
all the decorations, block fibrations and metrics by its actual maps.
Completeness and local model charts are preserved by pullback.
The capped Seifert pieces carry the metrics; their sphere-punctured
parts serve only reconstruction. Do not claim that puncturing a
complete geometric piece leaves a complete restriction metric, or that
sphere reconstruction is itself geometric. Boundary labels refer to
the original compact collars. In the cusp rows, a chosen collar can be
identified with the corresponding cusp end by its base collar and the
circle bundle trivialization; a collar compression extends this to the
whole interior. The exceptional flat rows use the displayed full-line
coordinate instead. This is an end comparison of the same marked
compactification, not an isometric matching of arbitrary boundary
metrics.

Essentiality follows from RG05's actual good-block gluing injections
and the existing cap square. Extra product cuts may remain. A closed
Sol bundle, or another closed torus bundle or semibundle encountered
inside a raw presentation, therefore needs no uncut homogeneous metric:
the essential retained cuts have the complete flat or Seifert witnesses
already constructed. This is the previously authorized retained-cut
E1 route, not a proof of G07's stronger uncut classification/descent.
Optional J03/J04 coarsening and F04's general fiber-parallel bundle
formula are likewise unnecessary here; the elementary pants branch
supplies the required result directly. Reattaching hyperbolic cores
uses AT17/AT18 and the separately proved GA05 complete finite-volume
hyperbolic metric. Final surgery-history reconstruction and arbitrary
non-graph exceptional producers remain separate goals.
\end{proof}
% END REVISION170 ACTUAL SEIFERT MODEL METRICS

% BEGIN REVISION171 FINAL GRAPH AUDIT
\section{Final mathematical audit of relative refinement and realization}
\label{sec:graph-realization-final-audit}
\markright{\thesection. Final graph realization audit}

The ONE final mathematical audit of the goal begun from revision168 is
complete here. It reviewed RG01--RG06 and GM01--GM06 against their actual
consumers and retained source inputs. The two expansions below make
nested ledger substitution and the exceptional-fiber covering coordinates
explicit; they belong to this same audit. The earlier completed static,
long-time analytic and ambient-incompressibility audits are not restarted.
The requirement matrix and source limits are recorded in
\code{AUDIT_GRAPH_REALIZATION_REVISION171.md} and
\code{graph_goal_audit_revision171.csv}.

\begin{lemma}[Flattening the actual relative phase ledger; GAU01]
\label{lem:graph-audit-marked-ledger}
The nested sphere operations in RG05 give one actual AT19 sphere ledger
and the required jointly disjoint final torus decoration. Its cycle
normalization counts each sphere-product contribution exactly once.
The comparison preserves the external labels and the based maps used
in AT16/AT17.
\end{lemma}
\begin{proof}
Initially the elementary presentation reconstructs $W$ by RG03's actual
maps. At a splitting step keep the local pants sphere and all surviving
seams as a jointly disjoint decoration of the current capped vertex.
AT14 puts that entire decoration off the incoming attachment balls.
Apply AT23, with tagged inner and outer ports, to substitute the new
sphere reconstruction into the current ledger. Every old ball lies in
one new punctured component; its pairing is assigned to that component,
and no new cap is confused with an old one. The quotient equivalence
is AT07/AT15, composed with the prior equivalence. Consolidations and
model identifications act by the same conjugation of actual port maps.
Induction on the finite RG05 phase yields a flat ledger, with all
orientations and external collar maps retained. The topology argument
uses this substitution before invoking a cycle count.

Now apply AT20 once to the resulting finite graph. A converted cycle
edge is removed and replaced by its sphere-product vertex and tree
edge. A sphere product already created by a meridional lens filling
is an existing labelled vertex and is not counted as a new cycle.
Transport every unused cap and pairing through the displayed maps.
Finally use the simultaneous cap-avoidance isotopy for the final torus
system; its lifted maps transport the Seifert structures and metrics.
With $c_a:R_a\to P_a$ and $j_a:R_a\to W$, the unchanged comparison
is $(j_a)_*(c_a)_*^{-1}$ with its recorded paths. External injection
therefore cancels through exactly these maps. No arbitrary prime-type
isomorphism replaces the reconstruction, and no embedded capped factor
in $W$ is required.
\end{proof}

\begin{lemma}[Exceptional-fiber and normalization check; GAU02]
\label{lem:graph-audit-exceptional-metric}
GM02--GM03's connection adjustment preserves the actual local Seifert
type at every exceptional fiber and gives a smooth complete metric on
the whole intended interior. It requires neither an ordinary section
through a cone point nor invariance of an independently chosen metric
on a finite cover.
\end{lemma}
\begin{proof}
In RG04 put $r=-a$ and use the explicit local covering
\[
 F:D^2\times(\R/\mathbb Z)\longrightarrow D^2\times S^1,
 \qquad F(z,\theta)=(e^{2\pi i r\theta}z,e^{2\pi i p\theta}).
\]
It is a smooth covering of degree $p$. Its deck generator is
$(z,\theta)\mapsto(e^{-2\pi i r/p}z,\theta+1/p)$, which acts freely;
$\gcd(r,p)=1$ makes its disk rotation have exact order $p$. The lifted
circle generator is $\partial_\theta$, so the lifted connection has
form $d\theta+\alpha$ with rotation-invariant curvature. Subtracting
a pulled-back orbifold one-form changes $\alpha$, not $F$, this deck
action, or the period-one normalization. Thus the locally compared
metric descends to the same smooth filling with the same meridian.
A small neighborhood of any point avoids its nontrivial deck translates,
giving an ordinary local model chart even over the cone point.

For the spherical metric the base curvature is four and the normalized
connection curvature has magnitude two; for Nil and the lifted unit
tangent metric it has magnitude one. GM03's positive values of $\ell$
produce exactly these magnitudes. In the flat-connection rows a closed
connection may have nontrivial holonomy, which is allowed in a locally
product metric. On open bases the de Rham statement is applied to the
boundaryless interior with cone points filled as orbifold points, not
to an unspecified compact boundary convention. Cone primitives are
constructed in their disk covers. Completeness uses the proper compact
circle projection to the complete base; it is not inferred from a
finite cusp truncation. The zero/nonzero $J$ alternatives, GM04's
positive-triangle exclusion and GM05's bad-orbifold lens cases exhaust
the RG06 list. No arbitrary finite-group recognition is used.
\end{proof}

\paragraph{Scope after the audit.}
The requested \emph{written relative graph-refinement and geometric
realization application} is complete under the explicit retained
classical foundations. RG05/GAU01 produce actual chosen prime carriers
and essential cuts; GM01--GM06/GAU02 produce complete metrics on the
same cut interiors. The actual late-flow binding is
Corollary~\ref{cor:graph-audited-late-binding}.
Ordinary smooth surface/bundle/collar/isotopy and torus mapping-class
interfaces, the Loop Theorem and the already retained Seifert
irreducibility/boundary-disk inputs, elementary two-dimensional model
geometry, corrected de Rham calculus and completeness remain exact
accepted-release binding obligations. This is not a claim that those
Lean declarations are absent, or that all declarations have been checked.

Optional canonical coarsening, comparison to a prescribed prime
factorization, the general F04 fiber-parallel formula and G07's uncut
Sol/semibundle realization are not needed by this E1 construction.
The current application status supersedes their historical necessity
for G09/GA06 without asserting their stronger signatures. Hyperbolic
finite volume is retained; no volume condition is imposed on every
flat open piece. Finite surgery-history reconstruction, arbitrary
exceptional-factor realization outside this graph list, and final GC
assembly remain separate. No new PC inspection or Lean execution was
performed. Consistency checks and PDF builds below certify records and
layout only, not these mathematical assertions.
% END REVISION171 FINAL GRAPH AUDIT

\chapter{Hyperbolic geometry, cusps and thick limits}
\label{ch:hyperbolic-geometry}

\paragraph{Scope.} Static hyperbolic infrastructure and the separately gated connection to thick flow limits.

\paragraph{Planned sections.}
\begin{enumerate}[leftmargin=*]
\item Hyperbolic models.
\item Quotients.
\item Cusps and truncation.
\item Peripheral groups.
\item Full-interior metrics.
\item Limit/stability bridge.
\item Cusp incompressibility route.
\end{enumerate}

\paragraph{Existing material and present task.} Existing source stack, I03/I08 and L33/L39; flow-slice injection is separate from model injection. Expand the static cusp/truncation/metric-transport outline with precise source-backed statements.

\section{Static geometry and the flow bridge}
Develop the static chain first: specify the hyperbolic model and quotient
conventions; define finite-volume cusp data and truncation; name the
actual peripheral maps; and describe transport of a complete metric to
the full interior of a compact core. Keep the local metric on a truncated
core distinct from the complete full-interior witness required by E1.

The separate flow chain has limit extraction, embedding/stability and
ambient cusp incompressibility obligations. L39 concerns peripheral
injection into the hyperbolic model; L33 concerns the actual image in a
late slice. Their distinction is retained in
Chapter~\ref{ch:ambient-reconstruction}. I08's variational support remains
a route-comparison task, not a decision to rebuild a minimal-surface
library. Hamilton 1999 and MSM206 retain their established guide roles.

% BEGIN migrated B035
\section{Hyperbolic branch: a source stack}
\label{sec:hyperbolic-source-stack}

The hyperbolic part should be developed as a stack of sources with distinct
jobs. We should not ask one reference to provide the Ricci-flow theorem, the
hyperbolic-geometry foundations, and the final topological reconstruction all
at once.

\begin{description}[leftmargin=3.5cm,style=nextline]
\item[Hamilton 1999] Hamilton's \emph{Non-singular solutions of the Ricci
  flow on three-manifolds} is the primary source for the nonsingular-flow
  classification. Its proof separates sequential collapse, positive and zero
  curvature limits, and negative limits; the negative case produces finite
  volume hyperbolic limits and, when the limit is noncompact, hyperbolic
  pieces. The GC formalization must expose the normalized-flow and bounded-
  curvature hypotheses instead of importing the conclusion as an unqualified
  theorem.
\item[MSM206] The local \code{MSM206.tex} copy is the detailed bridge for this
  branch. Chapters 31--34 and Appendix K should be mined for Lean-sized
  contracts: finite-volume hyperbolic ends, Margulis collars, incompressible
  boundary tori, convergence to hyperbolic limits, harmonic-map stability,
  and the passage from a limit to a piece embedded in the original manifold.
  Record the exact chapter, theorem label, and page in the theorem inventory;
  the chapter numbers are not a substitute for a locator.
\item[Perelman] Perelman's second paper supplies the Ricci-flow-with-surgery
  estimates, canonical-neighborhood and long-time inputs. Its hyperbolic
  discussion is a primary proof source, but its compressed arguments should
  be expanded using MSM206, Kleiner--Lott, and Morgan--Tian.
\item[Kleiner--Lott] The 2008 \emph{Notes} provide an alternative expansion
  of Perelman's analytic arguments. They should be compared with MSM206
  whenever a theorem concerns limits, stability, or surgery-compatible
  convergence. The 2014 Ast\'erisque paper is the selected local-collapse
  proof candidate, not a replacement for this hyperbolic-limit comparison.
\item[Morgan--Tian] Morgan--Tian supplies the full-GC proof spine and the
  assembly from hyperbolic and collapsed regions to geometrization. Use it to
  decide which Hamilton/Perelman intermediate statements are actually needed
  by the final endpoint, and as an independent check on the reconstruction.
\item[Benedetti--Petronio] \emph{Lectures on Hyperbolic Geometry} is the
  foundational geometry reference for the formal definitions and standard
  theorems about complete finite-volume hyperbolic 3-manifolds. It is not the
  source of the Ricci-flow convergence theorem. In particular, it should
  support the definitions of hyperbolic pieces, cusp truncations, and rigidity,
  while MSM206 and Morgan--Tian supply the Ricci-flow-to-piece interface.
\end{description}

The first hyperbolic vertical slice should therefore prove only a clean
interface theorem: a certified finite-volume hyperbolic limit together with
the MSM206 boundary and embedding hypotheses yields a hyperbolic-piece field
in \code{ThickThinCollapsePackage}. The full Hamilton/Perelman limit theorem and
the Kleiner--Lott/Morgan--Tian comparison remain upstream dependencies.
% END migrated B035

\section{Explicit hyperbolic objects and rigidity}
\label{sec:hyperbolic-rigidity-contracts}

The following HG contracts make the Chapter 8 dependencies explicit. They
refine the existing hyperbolic inventory, not the inherited foundation.
{\upshape\code{hyperbolic_geometry_contracts.csv}} distinguishes definitions, written
conditional arguments and source-qualified producers. Mostow--Prasad rigidity
is required by the stabilization route below, even though the public endpoint
does not ask for a unique decomposition.

\begin{lemma}[Positive rescaling; HG01]
\label{lem:hyp-rescaling}
For a positive constant $c$ and a three-dimensional Riemannian metric $g$,
the metric $cg$ multiplies lengths by $\sqrt c$, sectional curvature by
$c^{-1}$ and volume by $c^{3/2}$. Completeness and the underlying inclusion
maps on fundamental groups are unchanged.
\end{lemma}
\begin{proof}
The length integrand is multiplied by $\sqrt c$; taking infima over curves
gives the same distance scaling. Cauchy sequences are therefore unchanged.
The Levi-Civita connection is unchanged under constant scaling, and the
sectional-curvature numerator and denominator scale by $c$ and $c^2$.
The determinant formula for the volume density gives $c^{3/2}$. The maps
of underlying spaces have not changed. These computations use the inherited
metric, connection and measure APIs.
\end{proof}
KL Section 90 uses $t^{-1}g(t)$ and constant curvature $-1/4$.
Multiplication by $1/4$ gives curvature $-1$; distances and volumes then
scale by $1/2$ and $1/8$. Every radius, collar width, derivative norm and
volume inequality must be translated, not just the curvature label.

\begin{definition}[Supplied finite-volume model and cusp core; HG02]
\label{def:hyp-model-core-data}
Fix a connected complete oriented smooth hyperbolic three-manifold
$(H,h)$ of finite volume, with its curvature normalization. Cusp-core
data consist of finitely many disjoint parametrized torus ends, their
product collars, a compact smooth truncation $C$, an embedding
$\iota:C\to H$, and the end identifications recovering $H$ from $C$.
The empty end family is allowed when $H$ is closed. The boundary labels,
orientations and specific inclusion maps are part of the record.
\end{definition}
This is supplied data. Defining the record does not prove its existence.
A core embedding $\iota$ and a diffeomorphism $C^\circ\to H$ have different
jobs; they must not be identified.

% HG03 and Margulis05 integrated by the user in master112; audited in revision113.
\begin{theorem}[Cusp structure and model peripheral injection; HG03]
\label{found:hyp-cusp-structure}
Let $(H,h)$ be a connected complete oriented smooth three-manifold
without boundary, with $\sec_h=-\kappa^2$, $\kappa>0$, and finite volume.
Then the HG02 cusp-core data exist. The cusp charts, including inner
collars, have metric $dr^2+e^{-2\kappa r}q_i$ over compact flat tori.
Their lifted tails are precisely invariant horoballs with full
rank-two peripheral stabilizers. The actual cusp-slice and
core-boundary inclusions are injective on fundamental groups.
At every $0<s<t<\mu/\kappa$, the open thin sets defined by
$2\operatorname{inj}<s$ and $2\operatorname{inj}<t$ have the buffered
tube/cusp description of MGL01. Moreover
$\Vol_h(H)\geq v_*/\kappa^3>0$, with universal $\mu,v_*$ at curvature
$-1$. Compact truncations can contain any prescribed compact set
and discard arbitrarily little volume.
\end{theorem}
\begin{proof}
MGL00 supplies the geometric cusp/core witnesses, MGL01 the buffered
thin comparison, and MGL06 the actual peripheral maps. MGL08 gives
the universal volume bound, MGL09 changes normalization, and MGL31
gives volume-capturing truncations. Their supporting Margulis and
Zassenhaus proofs are included below.
\end{proof}
These are model-geometric statements. They do not assert injection
of a subsequently embedded torus into a flow slice or construct the
quantitative core comparison required by HG06. The source comparisons
and inherited interfaces are recorded in the Margulis development;
Lean verification and accepted-release bindings remain pending.
MGL32 supplies the compact-flat-section variant without orientability,
and MGL30 supplies countability and compact lift sets.

% Retained Margulis05 mathematical development; historical fragment unchanged.
% Current subordinate register: margulis_proof_obligations.csv.
% Integration/source audit: reference_checks_revision113.md.

\section{Margulis theory and the geometric cusp theorem}
\label{margulis:section}

The purpose of this section is to supply the geometric content of HG03.
The main output is a compact core with finitely many \emph{geometric}
torus cusps, including their lifts to horoballs.  These witnesses give
HG02 and the geometric hypotheses of MPR06--MPR07 in the finite-volume
Mostow--Prasad development.  The curvature is initially $-1$.

\paragraph{Written proof status.}
This section gives the complete natural-language proof development for
HG03: Zassenhaus, the algebraic Margulis lemma, buffered thin geometry,
and the finite-volume cusp/core theorem.  MGL00--MGL29 contain the
proof chain; MGL30--MGL31 make the compact-lift and volume-exhaustion
consumer interfaces explicit.  MGL32 adds the geometric cusp theorem
without orientability, with compact flat sections which may be Klein
bottles.  The main oriented blueprint interface retains torus sections.

All thirty-three contracts have written proofs or assemblies relative
to the specified inherited geometry, algebra, analysis, and topology.
No Margulis or Zassenhaus statement is left as an unproved foundation
block.  These are mathematical arguments and proposed formalization
interfaces, not checked Lean declarations.  Accepted inherited bindings
and Lean proofs remain implementation work. Revision113 integrates this
written status under HG03 without asserting formal acceptance.

\subsection{The theorem required by the blueprint}
\label{margulis:target}

\begin{theorem}[Geometric cusp and core theorem; MGL00]
\label{margulis:cusp-theorem}
Let $(H,h)$ be a connected,
complete, oriented smooth three-manifold without boundary, with
$\sec_h=-1$ and $\Vol_h(H)<\infty$.  There exist a nonnegative integer
$m$, compact flat tori $(T_i,q_i)$, positive numbers $\eta_i$, and
smooth embeddings
\[
 \Phi_i:T_i\times(-\eta_i,\infty)\longrightarrow H
 \qquad (1\leq i\leq m)
\]
with the following properties.
\begin{enumerate}[label=(\roman*)]
\item Their images are pairwise disjoint, and
\begin{equation}
 \Phi_i^*h=du^2+e^{-2u}q_i.
 \label{margulis:cusp-metric}
\end{equation}
Thus the warping function is $e^{-u}$; the coefficient of the flat
metric is its square $e^{-2u}$.
\item The set
\[
 C=H\setminus\bigcup_{i=1}^m\Phi_i(T_i\times(0,\infty))
\]
is a compact connected smooth submanifold with boundary,
$\partial C=\coprod_i\Phi_i(T_i\times\{0\})$.
Moreover $H$ is the union of $C$ and the closed cusp tails, meeting
along exactly these boundary tori.  The negative-$u$ parts give actual
inner collars in $C$.
\item For a hyperbolic universal covering
$p:\mathbb H^3\to H=\mathbb H^3/\Gamma$, every component of the
inverse image of an open cusp tail is an open horoball.  The components
over a given cusp form one $\Gamma$-orbit.  Each chosen horoball is
precisely invariant under its stabilizer $P_i<\Gamma$: a translate by
an element outside $P_i$ is disjoint from it.  The group $P_i$ is the
\emph{full} stabilizer in $\Gamma$ of its ideal center.  In upper
half-space coordinates it acts by
\[
 (x,z)\longmapsto(x+\lambda,z),\qquad
 \lambda\in\Lambda_i\subset\R^2,
\]
where $\Lambda_i$ is a rank-two lattice.  The quotient map of this
horoball is the given cusp embedding.
\item Every actual inclusion of a cusp slice into $H$, and every
boundary inclusion into $C$, is injective on fundamental groups.
With basepoints and connecting paths chosen, its image in $\Gamma$
is the corresponding conjugate of $P_i$.
\item The manifold $H$ strongly deformation retracts onto $C$, and
there is a diffeomorphism $C^\circ\to H$.  The latter is a different
map from the inclusion $C^\circ\hookrightarrow H$.
\end{enumerate}
There are exactly $m$ topological ends, represented by these cusp
tails.  The empty family is allowed and occurs exactly when $H$ is
compact.  Replacing the level $0$ in each cusp by any prescribed
$R_i\geq0$ gives another compact core with the same properties.
\end{theorem}

\begin{proof}[Top-down proof]
Apply MGL01 to the constant $\mu$ supplied by MGL10,
and choose $a=\mu/4$ and $b=\mu/2$.  Lemma~\ref{margulis:thick-compact}
makes the $a$-thick set compact.  The buffered comparison in
Lemma~\ref{margulis:finite-components} proves that there are only
finitely many $a$-thin components.  By MGL01 each is a tube or a
parabolic horoball quotient.  Lemma~\ref{margulis:rank-and-metric}
excludes rank-one parabolic components by finite volume and writes the
exact metric of every remaining cusp.  The larger $b$-thin component
provides an embedded collar beyond its boundary.

Remove the open cusp components and retain the tubes.  The construction
in Proposition~\ref{margulis:core-assembly} gives the compact connected
smooth core, its collars, the retraction, the interior diffeomorphism,
and the statement about ends.  The covering and full-stabilizer witnesses
come from MGL01; Lemma~\ref{margulis:peripheral} verifies the actual
fundamental-group maps.  Finite further truncations add only compact
product bands.  MGL10 and MGL11 are proved below, completing the
argument relative to the specified inherited interfaces.
\end{proof}

\paragraph{Scope.}
The orientation hypothesis is essential for the torus formulation:
without it, flat Klein-bottle sections can occur.  MGL32 proves the
corresponding compact-flat-section theorem.  Completeness and
absence of boundary are also hypotheses.  No asymptotic error is
present in \eqref{margulis:cusp-metric}; it is an exact metric identity.
No universal cutoff eliminating \emph{all} tubes is asserted.  The
compact core here includes the finitely many short-geodesic tubes.

\lean{Export the chosen core, cusp charts, flat lattices, collar widths,
covering, and actual inclusion maps as witnesses.  HG02 records part of
this data; the lifted-horoball enhancement is needed by MPR06.  Keep the
basepoint-change conjugacies.  This section supplies model geometry;
HG10--HG15 and AT01 still supply injection into an ambient flow slice.}

\subsection{The local input and the finite-volume reduction}
\label{margulis:local-to-global}

We use the inherited space-form presentation of a complete connected
hyperbolic manifold, the properly discontinuous free deck action, the
upper half-space metric, path lifting, and Hopf--Rinow.  In the oriented
case $\Gamma<\operatorname{Isom}^+(\mathbb H^3)$ is discrete and
torsion-free.  These are inherited-foundation bindings, not new PC
absence claims.  Write
\[
 \rho(q)=\operatorname{inj}_h(q),\qquad
 H_{<s}=\{q:2\rho(q)<s\},\qquad
 H_{\geq s}=\{q:2\rho(q)\geq s\}.
\]
For $p(x)=q$ the covering formula is
\begin{equation}
 2\rho(q)=\inf_{\gamma\in\Gamma\setminus\{1\}}d(x,\gamma x).
 \label{margulis:inj-displacement}
\end{equation}
The empty infimum is $+\infty$.  For nontrivial $\Gamma$, properness
and discreteness make this a positive attained minimum.  Each
displacement function is $2$-Lipschitz.  Lifting a path between two
points and taking the infimum gives that $\rho$ is $1$-Lipschitz.
The trivial-group case is $\mathbb H^3$, of infinite volume, and so
does not occur in MGL00.

\begin{theorem}[Buffered thin-component description; MGL01]
\label{margulis:thin-input}
Let $\mu>0$ be the constant in MGL10.  The following
properties hold for every complete connected oriented hyperbolic
three-manifold and every $0<s<\mu$.

Every connected component $U$ of $H_{<s}$ has a chosen lifted component
$W\subset\mathbb H^3$ and subgroup $P<\Gamma$ such that
\[
 W=\{x:\text{some }\gamma\in P\setminus\{1\}
                   \text{ satisfies }d(x,\gamma x)<s\},
 \qquad U=W/P.
\]
The set $W$ is precisely invariant under $P$, and its quotient embeds
smoothly in $H$ with the restricted Riemannian metric.  Exactly one of
the following alternatives occurs.
\begin{enumerate}[label=(\roman*)]
\item $P$ is generated by a loxodromic isometry with axis $L$, positive
translation length $\ell$, and possibly nonzero rotational part.
It is the full stabilizer of $L$ in $\Gamma$.  The set $W$ is an open
tube of a finite positive radius about $L$.  Its closed quotient is
compact and is a smooth solid torus.
\item $P$ is the full stabilizer of an ideal point $\xi$.  After
moving $\xi$ to $\infty$, it is a nonzero discrete group of translations
of $\R^2$, of rank one or two, and $W=\{(x,z):z>z_s\}$ for some
$z_s>0$.
\end{enumerate}
For $0<s<t<\mu$, inclusion of thin sets sends distinct $s$-thin
components into distinct $t$-thin components.  A nonempty $s$-component
and its containing $t$-component have the same maximal subgroup, after
compatible lifting, and
\[
 \overline{W_s}\subset W_t.
\]
In particular its closure and a collar of its boundary embed inside
the larger quotient.  These are statements about the ordinary open
thin set; a geodesic whose length is exactly $s$ contributes no
degenerate component at that threshold.
\end{theorem}
\begin{proof}[Proof from the algebraic input]
By MGL13, each nontrivial deck transformation determines a unique
maximal elementary subgroup, which is the full stabilizer of its
boundary fixed-point data and equals its own normalizer.  MGL14
computes the region where some element of such a subgroup moves a
point by less than $s$: it is an open round tube, an open horoball,
or the empty set.  The algebraic small-displacement assertion of
MGL10 shows that regions belonging to different subgroups cannot
overlap.  MGL15 therefore identifies the nonempty regions with the
components of the lifted thin set.

MGL16 proves precise invariance and covering descent, retaining the
full stabilizer rather than just the subgroup generated by short
elements at one point.  MGL17 proves non-merging under $s<t$ and the
identification of the embedded closures.  MGL18 gives compact solid
tori with their rotational holonomy and supplies smooth product
collars in the larger component.  The parabolic stabilizer and its
rank come from MGL13 and MGL23.  These establish all the stated
conclusions.  The supporting proofs appear below, before the lower
algebraic producer.  None uses finite volume or the global cusp theorem.
\end{proof}

The local geometric reduction and its algebraic inputs MGL10--MGL11
now all have written proofs, relative to the inherited model and
covering interfaces specified below.
The strict buffer $s<t$ avoids equality-level degeneracy and gives
room for the collars; it is retained by every subsequent consumer.

\begin{lemma}[Compact thick sets; MGL02]
\label{margulis:thick-compact}
If $(H,h)$ is connected, complete, hyperbolic, and has finite volume, then
$H_{\geq s}$ is compact for every $s>0$.
\end{lemma}
\begin{proof}
It is closed by continuity of $\rho$.  Completeness and Hopf--Rinow
make $H$ a proper metric space.  If this closed set were noncompact,
it would be unbounded.  Inductively choose points $q_j$ in it with
$d(q_j,q_k)\geq s$ for $j\ne k$.  The open balls $B(q_j,s/4)$ are
disjoint.  Each has the hyperbolic ball volume $v(s/4)>0$, since
$\rho(q_j)\geq s/2$ and the covering is injective on a ball of radius
$s/4$.  Finite disjoint unions would have arbitrarily large volume,
contrary to $\Vol_h(H)<\infty$.
\end{proof}

\begin{lemma}[Finitely many thin components; MGL03]
\label{margulis:finite-components}
Assume finite volume.  If $0<a<b<\mu$, then $H_{<a}$ has
finitely many components, their closures are disjoint smooth
submanifolds with boundary, and $H_{\geq a}$ is nonempty.
\end{lemma}
\begin{proof}
If $H_{<a}=H$, its full inverse image would be the single connected
component $\mathbb H^3$.  MGL01 identifies such a component with a
proper tube or horoball, a contradiction.  Thus $H_{\geq a}$ is
nonempty.  Each nonempty thin component $U$ is a proper open subset of
connected $H$, so it has nonempty boundary.  At a boundary point
$q\in\partial U$ continuity gives $\rho(q)=a/2$: a point with smaller
injectivity radius has a connected neighborhood in $H_{<a}$ and cannot
be a boundary point of a different component.

Put $\delta=(b-a)/2$.  If boundary points $q\in\partial U$ and
$q'\in\partial U'$ satisfy $d(q,q')<\delta$, a minimizing segment
from $q$ to $q'$ lies in $H_{<b}$, by the $1$-Lipschitz bound on $\rho$.
Small neighborhoods of its endpoints meet $U$ and $U'$.  Hence both
components lie in the same $b$-thin component.  MGL01 implies $U=U'$.
Choosing one boundary point from each component therefore gives a
$\delta$-separated set in the compact set $H_{\geq a}$.  Such a set
is finite, by a finite cover with balls of radius $\delta/3$.

The closure of each component embeds inside its containing $b$-thin
component.  There it is a round closed tube quotient or a closed
horoball quotient, with smooth boundary.  Distinct smaller components
have distinct containing components; their closures are disjoint.
\end{proof}
This argument supplies the finiteness step without inferring that an
arbitrary compact subset has only finitely many connected components.
The buffer also keeps all collar assertions away from the Margulis
threshold and from a degenerate tube.

\begin{lemma}[Rank and exact metric of a cusp; MGL04]
\label{margulis:rank-and-metric}
Let a parabolic thin component be represented by
$\{z>z_0\}/\Lambda$, as in MGL01.  If it embeds in a finite-volume
$H$, then $\Lambda$ has rank two.  Its closure, with the restricted
Riemannian tensor, has the form
\[
 (T\times[0,\infty),\ du^2+e^{-2u}q),
 \quad T=\R^2/\Lambda,\quad
 q=z_0^{-2}q_{\mathrm{Euc}},\quad z=z_0e^u.
\]
These coordinates extend to a negative interval in $u$ inside the
larger buffered cusp.
\end{lemma}
\begin{proof}
The upper half-space metric and volume density are
\[
 h=z^{-2}(dx_1^2+dx_2^2+dz^2),\qquad
 d\Vol_h=z^{-3}\,dx_1\,dx_2\,dz.
\]
If $\Lambda$ has rank one, choose orthogonal Euclidean coordinates in
which its generator is $(L,0)$, $L>0$.  The slab fundamental domain
$[0,L)\times\R\times(z_0,2z_0)$ has infinite hyperbolic volume.
For example, its restrictions to $|x_2|<A$ have volumes growing
linearly in $A$.  This contradicts the volume of the embedded component.
Thus $\Lambda$ has rank two, and a compact lattice parallelogram gives
the flat torus $T$.

Substitute $z=z_0e^u$ and $dz=z_0e^u du$ in $h$ to obtain the displayed
tensor identity.  The substitution commutes with the translation
action and therefore descends to the quotient.  If the larger cusp
starts at $z_b<z_0$, choose any
$0<\eta<\log(z_0/z_b)$.  The same formula on $u> -\eta$ gives the
required embedded extension.
\end{proof}

\begin{proposition}[Assembly of the compact core; MGL05]
\label{margulis:core-assembly}
The finite cusp family constructed in MGL03--MGL04 gives
all the core, collar, and end conclusions of MGL00.
\end{proposition}
\begin{proof}
Let $C$ be the complement of the open cusp components.  It is the
union of the compact thick set and the closures of the finitely many
tube components.  Each closed tube quotient is compact, even when its
loxodromic generator has a rotational part.  Hence $C$ is compact.
The buffered cusp charts show that it is a smooth submanifold with
boundary, whose boundary is exactly the finite union of level-zero
tori.  All inner collars lie in $C$ and are mutually disjoint.

On each closed cusp define, for $0\leq t\leq1$,
\[
 F_t(\Phi_i(x,u))=\Phi_i(x,(1-t)u),\qquad u\geq0,
\]
and let $F_t$ be the identity on $C$.  The formulas agree at $u=0$.
The core and the finitely many closed cusps form a finite closed cover,
so they glue continuously, also with the homotopy parameter.  This is
a strong deformation retraction onto $C$.  Since $H$ is connected,
$C=F_1(H)$ is connected.

For the interior diffeomorphism use an inner collar
$T_i\times(-\eta_i,0]$.  Choose a smooth nondecreasing function
$\chi_i:(-\eta_i,0)\to[0,1]$ which vanishes near $-\eta_i$ and equals
$1$ near $0$.  Then
\[
 f_i(t)=t-\frac{\chi_i(t)}{t}
\]
has positive derivative, is the identity near the inner end, and is a
diffeomorphism $(-\eta_i,0)\to(-\eta_i,\infty)$.  On the interior
collar use $(x,t)\mapsto\Phi_i(x,f_i(t))$; away from the collars use
the inclusion.  The formulas and their inverses glue smoothly to a
diffeomorphism $C^\circ\to H$.  This uses the inherited smooth cutoff
construction, not any rigidity theorem.

For $R_i\geq0$, the core obtained by removing only $u>R_i$ is the
union of $C$ and finitely many compact bands.  The cores with all
$R_i=n$ exhaust $H$ and are cofinal among its compact subsets: their
interiors cover $H$, so compactness bounds the cusp depth of any
compact subset.  The complement at every such level has exactly the
$m$ connected tails, with each tail contained in the same labeled
tail at the preceding level.  This proves the assertion about
topological ends.  If $m=0$, $H=C$ is compact; if $m>0$, the closed
cusp product is noncompact and closed in $H$, so $H$ is noncompact.
\end{proof}

\begin{lemma}[The actual peripheral maps; MGL06]
\label{margulis:peripheral}
For every cusp chart supplied above and every $R\geq0$, the map
$T_i\to H$, $x\mapsto\Phi_i(x,R)$, is injective on fundamental
groups.  The boundary inclusion $T_i\to C$ at $R=0$ is also injective.
\end{lemma}
\begin{proof}
Choose a basepoint and one lift on the horizontal horosphere
$\R^2\times\{z_0e^R\}$.  This simply connected horosphere covers
the torus with deck group $P_i\cong\Lambda_i$.  A torus loop lifts
to a path ending at $\gamma x$ for its associated $\gamma\in P_i$.
The same path is its lift to the universal covering of $H$, so its
image in $\pi_1(H)$ corresponds to the same element of
$P_i<\Gamma$.  If the image is trivial, $\gamma=1$ and the lifted
loop contracts in the horosphere.  This proves injection and identifies
the image.  Other basepoint choices conjugate this subgroup by the
specified connecting path.  Since $C\hookrightarrow H$ is a homotopy
equivalence, the boundary map into $C$ is injective as well.
\end{proof}
Use one \emph{connected component} of the inverse image of a cusp
slice in this argument; its entire inverse image is a disjoint union
of horospheres.  A subgroup isomorphic to $\mathbb Z^2$ need not equal
a full peripheral subgroup: proper finite-index subgroups of a lattice
already give counterexamples to that stronger statement.

\subsection{Quantitative consequences and normalization}
\label{margulis:consequences}

\begin{corollary}[Cusp decay, volume, and second fundamental form; MGL07]
\label{margulis:decay}
Let $A_i=\operatorname{Area}_{q_i}(T_i)$ and let $D_i$ be its intrinsic
diameter.  At height $R\geq0$ the induced metric is $e^{-2R}q_i$.
Lengths of fixed curves and intrinsic diameter scale by $e^{-R}$,
and
\[
 \operatorname{Area}(T_i\times\{R\})=e^{-2R}A_i,
 \qquad \operatorname{diam}(T_i\times\{R\})=e^{-R}D_i,
\]
\[
 \Vol(T_i\times[R,\infty))=\frac{A_i}{2}e^{-2R}.
\]
With $\nu=\partial_u$ pointing toward the cusp end and convention
$\mathrm{II}(X,Y)=h(\nabla_X\nu,Y)$, one has
$\mathrm{II}=-h|_{T_i\times\{R\}}$ and mean curvature $-2$.
\end{corollary}
\begin{proof}
The slice metric gives the length and diameter identities.  Its
two-dimensional area density is multiplied by $e^{-2R}$, and the
three-dimensional density is $e^{-2u}dA_{q_i}\,du$.  Integration from
$R$ to infinity gives the tail volume.  For slice-tangent fields
independent of $u$, the Koszul formula gives
\[
 2h(\nabla_X\partial_u,Y)=\partial_u h(X,Y)=-2h(X,Y).
\]
Taking the trace on the two-dimensional slice gives the stated mean
curvature.  Reversing the normal reverses both signs.
\end{proof}
The displayed diameter uses the intrinsic slice metric.  In particular
these formulas do not assert that the ambient injectivity radius is
exactly a constant times $e^{-R}$.

\begin{corollary}[A uniform positive volume bound; MGL08]
\label{margulis:volume-bound}
There is a universal $v_*>0$ such that every connected
complete oriented hyperbolic three-manifold of finite volume and
curvature $-1$ has volume at least $v_*$.  One admissible choice is
\[
 r_* = \mu/16,\qquad
 v_* = \Vol(B(r_*))
      =\pi\bigl(\sinh(2r_*)-2r_*\bigr)>0.
\]
\end{corollary}
Here $B(r_*)$ is a ball in $\mathbb H^3$.
\begin{proof}
MGL03 supplies a point in $H_{\geq\mu/4}$.  Its injectivity radius
is at least $\mu/8$, so the covering embeds the radius-$r_*$ ball
with its hyperbolic Riemannian volume.  This ball has the displayed
volume, obtained by integrating $4\pi\sinh^2 r$.  Its volume is a
lower bound for that of $H$.
\end{proof}
This is sufficient for the HG09 volume-budget argument; no optimal
Margulis constant or sharp minimal-volume theorem is needed.

\begin{lemma}[Changing curvature normalization; MGL09]
\label{margulis:normalization}
For curvature $-\kappa^2$, $\kappa>0$, the unit-speed cusp coordinate
$r$ has metric
\[
 dr^2+e^{-2\kappa r}q_0.
\]
Its tail volume above $r=R$ is
$\operatorname{Area}_{q_0}(T)e^{-2\kappa R}/(2\kappa)$.
The same normal convention gives $\mathrm{II}=-\kappa h|_T$.
The constants $\mu$ and $v_*$ at curvature $-1$ become
$\mu/\kappa$ and $v_*/\kappa^3$.
\end{lemma}
\begin{proof}
Multiply the curvature-$-\kappa^2$ metric by $\kappa^2$ to obtain
curvature $-1$.  In the normalized cusp formula set $u=\kappa r$ and
$q_0=\kappa^{-2}q$.  Dividing the metric by $\kappa^2$ gives the
claimed identity.  Lengths and volumes scale by $\kappa^{-1}$ and
$\kappa^{-3}$.  The density integral and the same Koszul calculation
give the other two assertions.
\end{proof}
For the KL normalization $\sec=-1/4$, $\kappa=1/2$ and the formula is
$dr^2+e^{-r}q_0$.  Equivalently $h_{-1}=\tfrac14h_{-1/4}$, with
length factor $1/2$ and volume factor $1/8$.  The coordinate must be
rescaled along with the tensor.  A complete metric on $C^\circ$ is
obtained by pulling back $h$ through the diffeomorphism in MGL05;
the restricted metric from a finite truncation is generally incomplete.

\subsection{The exact geometric inputs for the consumers}
\label{margulis:consumer-adapters}

\begin{lemma}[Countable deck orbits and compact lift sets; MGL30]
\label{margulis:compact-lifts}
Let $p:\mathbb H^3\to H=\mathbb H^3/\Gamma$ be the hyperbolic
universal covering of a connected complete hyperbolic manifold.
No orientability or finite-volume assumption is needed here.
Then $\Gamma$ is countable.  For every compact $K\subset H$ there
is a compact $D\subset\mathbb H^3$ such that
\[
 \forall x\in p^{-1}(K)\ \exists\gamma\in\Gamma:
                      \gamma x\in D.
\]
If a finite family of cusp neighborhoods has one deck orbit of lifted
horoballs over each cusp, the set of all their ideal centers is countable.
\end{lemma}
\begin{proof}
Fix $o\in\mathbb H^3$.  The map $\gamma\mapsto\gamma o$ is injective
by freeness, and its image is a discrete subspace: an evenly covered
neighborhood of $p(o)$ isolates each point of the fiber.  Every
discrete subspace of a second-countable space is countable.  Indeed,
choose a countable basis and assign to each point the least basis
index whose open set meets that subspace only at the chosen point.
Different points receive different indices.  Thus $\Gamma$ is countable.

For each $q\in K$, choose an evenly covered open neighborhood $V_q$,
and then an open $U_q$ containing $q$ with compact closure contained
in $V_q$.  Such a shrink exists in a manifold.  Choose finitely many
$U_1,\ldots,U_k$ covering $K$ and a single inverse sheet
$\sigma_j:V_j\to\mathbb H^3$ for each.  The set
\[
 D=\bigcup_{j=1}^k\sigma_j(\overline U_j)
\]
is compact.  If $p(x)\in U_j$, regularity of the universal covering
gives a deck transformation taking $x$ to $\sigma_j(p(x))\in D$.
Finally choose one center $\xi_i$ over each of the finitely many
cusps.  All centers lie in $\bigcup_i\Gamma\xi_i$, which is countable.
\end{proof}
Together with MGL00, this supplies the compact-core, horoball,
countability, and compact-lift hypotheses of MPR06--MPR07 in the
oriented case.  MGL32 below supplies the geometric data without
orientability.  Properness of isometry evaluation and continuity of
its action on the compactification remain the stated inherited model
interfaces.  No boundary regularity or controlled representative of
a prescribed homotopy equivalence is asserted here.

\begin{lemma}[Truncations capturing volume; MGL31]
\label{margulis:volume-truncations}
Suppose a complete finite-volume hyperbolic three-manifold of curvature
$-\kappa^2$, $\kappa>0$, has the geometric cusp/core data above, allowing
any compact flat cusp sections.  Write $A_i$ for the area of section
$i$ at unit-speed height $r=0$, and let $C_{(R_i)}$ be the
truncation at heights $R_i\geq0$.  Then
\[
 \Vol(H\setminus C_{(R_i)})
   =\sum_{i=1}^m\frac{A_i}{2\kappa}e^{-2\kappa R_i}.
\]
For every compact $B\subset H$ and every $\varepsilon>0$ one can
choose the $R_i$ so that $B\subset C_{(R_i)}$ and the displayed
volume is less than $\varepsilon$.  In particular, if $\Vol(H)\geq v>0$,
one can choose $\Vol(C_{(R_i)})>3v/4$.
\end{lemma}
\begin{proof}
The complement is a finite disjoint union of open tails.  The density
calculation of MGL07--MGL09 applies to compact flat sections using
Riemannian densities, whether or not they are orientable, and gives
the displayed formula.  The common-height compact cores are cofinal
among compact sets by MGL05's exhaustion argument.  Increase their
common height further until the finite sum is below $\varepsilon$;
every summand tends to zero.  If there are no cusps, the sum is zero
and the core is all of $H$.  Finite additivity gives
$\Vol(C_{(R_i)})=\Vol(H)-\Vol(H\setminus C_{(R_i)})$.
Take $\varepsilon=v/4$ for the last assertion.
\end{proof}
This provides a useful margin for HG09: an embedding of such a core
whose pulled-back volume density is at least $2/3$ of the model
density has image volume greater than $v/2$, by change of variables.
Supplying that embedding and its metric estimate, at a common late
time for each finite selection of models, remains the flow argument.
This lemma imposes no one time for an infinite selection.

\begin{theorem}[Geometric cusps without orientability; MGL32]
\label{margulis:unoriented-cusps}
Let $(H,h)$ be a connected complete smooth hyperbolic three-manifold
without boundary, of curvature $-1$ and finite volume, with no
orientability assumption.  The conclusions of MGL00 hold with compact
flat connected surfaces $(S_i,q_i)$ in place of the flat tori, and with
cusp groups $P_i$ acting freely by Euclidean isometries on $\R^2$.
Each $P_i$ contains its rank-two translation subgroup
$P_i^+=P_i\cap\operatorname{Isom}^+(\mathbb H^3)$ with index at most two.
The full ideal-center stabilizer and the actual peripheral injection
conclusions are retained.  Every $S_i$ is a torus or a Klein bottle;
in the oriented case all are tori.  There is a universal volume bound
$\Vol(H)\geq v_*/2$, where $v_*$ is the oriented constant of MGL08.
At curvature $-\kappa^2$ use $v_*/(2\kappa^3)$ and
$dr^2+e^{-2\kappa r}q_i$.
\end{theorem}
\begin{proof}
Present $H=X/\Gamma$, $X=\mathbb H^3$, using the inherited space-form
covering, and set
\[
 \Gamma^+=\Gamma\cap\operatorname{Isom}^+(X),\qquad
 H^+=X/\Gamma^+,\qquad d=[\Gamma:\Gamma^+]\in\{1,2\}.
\]
This is a normal subgroup and $H^+\to H$ is a regular degree-$d$
local-isometry covering.  It is complete by geodesic lifting and
Hopf--Rinow, connected, and oriented.  Its volume is
$d\,\Vol(H)$.  Here is a measure justification independent of
orientability.  Take a countable cover of $H$ by evenly covered
coordinate neighborhoods $V_j$, and partition it into the disjoint
Borel sets $E_j=V_j\setminus\bigcup_{l<j}V_l$.  Over $E_j$ there
are exactly $d$ inverse sheets; each preserves Riemannian density.
Change of variables on those sheets and countable additivity give
the asserted volume equality.  Thus MGL00 applies to $H^+$.

Use specifically the cusp family constructed from the
$a$-thin components for $\Gamma^+$, with $a=\mu/4$ and buffer
$b=\mu/2$.  The deck group $\Gamma/\Gamma^+$ acts isometrically on
$H^+$ and permutes these components, preserving the parabolic
type.  Equivalently, on $X$ normality gives
\[
 \gamma W_a(P^+)=W_a(\gamma P^+\gamma^{-1})
                    \quad(\gamma\in\Gamma).
\]
The parabolic regions are disjoint horoballs by MGL15, so every
translate of a selected horoball is either itself or disjoint.
There are finitely many orbits downstairs.  Removing their open
images leaves the quotient of the invariant compact core $C^+$,
which is compact.

We identify the full stabilizer, a step not implied just by taking
the finite quotient.  Fix a selected horoball $W$ with center $\xi$
and let $P^+$ be its stabilizer in $\Gamma^+$, a full rank-two
translation group.  If $\gamma\in\Gamma$ fixes $\xi$, then
$\gamma^2\in\Gamma^+$ fixes $\xi$, hence $\gamma^2\in P^+$ by
MGL00.  In coordinates with $\xi=\infty$, the inherited normal form
for an arbitrary isometry fixing infinity is
\[
 \gamma(x,z)=(c+\lambda Qx,\lambda z),
 \qquad \lambda>0,\quad Q\in O(2).
\]
For clarity, the normal form also follows from the model facts already
used: an isometry fixing infinity permutes horizontal horospheres and
vertical geodesics.  Its restriction between two horizontal levels is
a Euclidean similarity with height ratio $\lambda$; the vertical
geodesics force the same horizontal formula at every height, and
unit-speed vertical lengths, with direction toward infinity preserved,
give $z\mapsto\lambda z$.
Since $\gamma^2$ is a translation, $\lambda^2=1$, whence $\lambda=1$.
Thus $\gamma$ preserves every horoball centered at $\xi$, including
$W$.  Conversely, an isometry preserving $W$ fixes its unique ideal
center.  Therefore $P=\operatorname{Stab}_\Gamma(W)$ is exactly
$\operatorname{Stab}_\Gamma(\xi)$, and $P\cap\Gamma^+=P^+$ has index
at most two in $P$.

The action of $P$ on each horizontal plane is Euclidean and free:
a fixed horizontal point would yield a fixed point of $X$, contrary
to deck freeness.  It is properly discontinuous by restriction of
the deck action.  Hence $S=\R^2/P$ is a smooth connected flat
surface.  The torus $\R^2/P^+$ covers it with degree one or two,
so $S$ is compact.  Precise invariance and the fiber argument of
MGL16 make $W/P\to H$ an injective local isometry onto its open
image.  The full inverse image is exactly $\bigcup_{\gamma\in\Gamma}\gamma W$;
these disjoint connected open horoballs are its components and form
one deck orbit.  These quotient maps identify the cusp with
$S\times(0,\infty)$ in the coordinate $u=\log(z/z_a)$, with tensor
$du^2+e^{-2u}q$.

The inner collars also descend.  In MGL14 the ratio of the heights
of the two cusp boundaries is
\[
 \frac{z_a}{z_b}=\frac{\sinh(b/2)}{\sinh(a/2)}>1,
\]
independent of the lattice.  Choose one
$0<\eta<\log(\sinh(b/2)/\sinh(a/2))$ for every cusp.  An isometry
carrying one selected horoball to another preserves the signed
distance $u$ from its boundary: in half-space coordinates both
$z$ and the boundary height acquire the same factor.  Thus the
extensions on $u>-\eta$ are permuted as products, and distinct
cusp orbits have disjoint extended images, inside distinct buffered
components.  The quotient core is therefore a connected smooth
compact submanifold with boundary the level-zero flat surfaces.

For peripheral injection, a chosen horizontal plane is simply
connected and covers $S$ with deck group $P$.  A loop in $S$ and
its image in $H$ have the same lifted endpoint, so their associated
deck element is the same element of $P<\Gamma$.  This proves
injection and the full image assertion, with the usual based
conjugacies.  The retraction, collar-stretching diffeomorphism,
further truncations, and cofinal end exhaustion in MGL05 work
verbatim for these connected compact flat sections.  In particular
the boundary injection into the core follows through its homotopy
equivalence with $H$.

For the topological names of the sections, if $P=P^+$ the section
is the torus already constructed.  Otherwise $P$ contains an
orientation-reversing Euclidean isometry, so $S$ is nonorientable.
Its torus double cover gives $0=\chi(T^2)=2\chi(S)$: triangulate $S$
and lift its cells to verify this equality.  The inherited
classification of closed connected surfaces gives
$S\cong\#^g\mathbb{RP}^2$ with $\chi(S)=2-g$, hence $g=2$ and $S$
is a Klein bottle.  Only this naming step uses surface classification;
the geometric cusp and covering statements above do not.

Finally MGL08 gives
$d\,\Vol(H)=\Vol(H^+)\geq v_*$, and $d\leq2$ gives the stated
lower bound.  Constant rescaling as in MGL09 proves the last claim.
\end{proof}
The cusp neighborhoods in this proof descend from the thin set of
the orientation cover.  We have not identified them with the thin
components of the nonorientable quotient at the same threshold:
additional orientation-reversing short elements can change those
components.  That stronger assertion is not used by MPR06--MPR07.
HG02--HG03 keep their oriented torus scope; MGL32 is the additional
geometric input available for HG04's broader manifold scope.

\lean{MGL30 supplies the compact lift witnesses and countability;
MGL31 supplies truncation volumes with an explicit margin; MGL32
supplies orientation-cover descent with full stabilizers and exact
warped metrics.  The new inherited bindings are standard covering,
density, second-countability, and finite-quotient interfaces; naming
the nonorientable flat section also uses surface classification.
No rigidity theorem is used in their proofs.}


\subsection{Proof of the local thin-component theorem}
\label{margulis:thin-proof}

Fix a discrete freely acting subgroup
$\Gamma<\operatorname{Isom}^+(\mathbb H^3)$ and its hyperbolic quotient
$p:\mathbb H^3\to H$.  Use the small-displacement
theorem MGL10.  The elementary classification used in MGL13
is supplied independently by MGL23.  No finite-volume hypothesis is
imposed in this subsection.

\paragraph{Inherited model facts used here.}
A nonidentity element of this group has either one parabolic fixed
point or two loxodromic fixed points at infinity.  Conjugation carries
these fixed sets to their images.  In upper half-space coordinates
$(x,z)\in\R^2\times(0,\infty)$ use
\begin{equation}
 \cosh d((x,z),(y,w))
   =1+\frac{|x-y|^2+(z-w)^2}{2zw},
 \qquad h=z^{-2}(|dx|^2+dz^2).
 \label{margulis:halfspace-distance}
\end{equation}
After normalizing its axis, an oriented loxodromic generator has the
form $(x,z)\mapsto(e^\ell A_\theta x,e^\ell z)$, where $\ell>0$ and
$A_\theta$ is planar rotation through an angle $\theta\in\R$.
Its nonzero powers have the same two fixed points.  A nonzero
parabolic translation fixes precisely $\infty$.  These single-isometry
and model formulas are inherited geometric bindings, distinct from
MGL10's assertion about entire elementary groups.  The calculations
below supply the new subgroup regions and quotient comparisons.

\begin{lemma}[Full elementary stabilizers; MGL13]
\label{margulis:full-stabilizers}
For $g\in\Gamma\setminus\{1\}$ let $F(g)$ be its boundary fixed set
and define
\[
 E(g)=\{\gamma\in\Gamma:\gamma F(g)=F(g)\}.
\]
Then $E(g)$ is the unique maximal elementary subgroup of $\Gamma$
containing $g$.  For every $h\in E(g)\setminus\{1\}$,
\[
 F(h)=F(g),\qquad E(h)=E(g).
\]
It is cyclic loxodromic when $g$ is loxodromic, and is a rank-one or
rank-two parabolic translation group when $g$ is parabolic.  It is
the full axis stabilizer in the first case and the full ideal-point
stabilizer in the second.  Moreover,
\[
 E(\gamma g\gamma^{-1})=\gamma E(g)\gamma^{-1},
 \qquad N_\Gamma(E(g))=E(g).
\]
\end{lemma}
\begin{proof}
The finite nonempty set $F(g)$ is preserved by $E(g)$, so $E(g)$ is
elementary.  It contains $g$ and inherits discreteness and freeness.
Apply MGL23's elementary-group classification.  Since it contains
$g$, its type is the type of $g$; in either alternative all its
nonidentity elements have the same boundary fixed set, namely $F(g)$.
This proves the assertions about $h$ and the two possible group forms.
Any elementary subgroup containing $g$ has, by the same classification,
this same fixed-point data and is contained in $E(g)$.  This proves
maximality and uniqueness.  Stabilizing the two endpoints of an axis
as a set is equivalent to stabilizing the axis itself.  For a singleton
it is exactly the full ideal-point stabilizer.

Conjugation carries fixed sets to fixed sets, proving the first
displayed equivariance identity.  If $\gamma$ normalizes $E(g)$, then
$\gamma g\gamma^{-1}$ is a nonidentity element of $E(g)$; hence
$\gamma F(g)=F(\gamma g\gamma^{-1})=F(g)$ and $\gamma\in E(g)$.
The reverse normalizer inclusion holds for every subgroup.
\end{proof}

Write $\mathcal E$ for the set of distinct subgroups $E(g)$ and, for
$P\in\mathcal E$ and $s>0$, put
\begin{equation}
 W_s(P)=\{x\in\mathbb H^3:
       \exists g\in P\setminus\{1\},\ d(x,gx)<s\}.
 \label{margulis:region-definition}
\end{equation}
The use of the full subgroup $P$ matters: the small-displacement group
$\Gamma_s(x)$ at a particular point need not equal $P$.  For example,
the shortest translations at that point can generate a cyclic subgroup
of a rank-two cusp lattice.  The peripheral group exported to HG03 is
the full stabilizer, not that smaller group.

\begin{lemma}[Explicit displacement regions; MGL14]
\label{margulis:region-shapes}
For each subgroup of either form in MGL13, the following descriptions
hold for every $s>0$.
\begin{enumerate}[label=(\roman*)]
\item If $P$ is the nonzero translation group of a discrete subgroup
$\Lambda\subset\R^2$, let
$\lambda_0=\min\{|\lambda|:\lambda\in\Lambda\setminus\{0\}\}>0$.
Then
\[
 W_s(P)=\{(x,z):z>z_s\},\qquad
 z_s=\frac{\lambda_0}{2\sinh(s/2)}.
\]
\item Suppose $P=\langle g\rangle$ with translation length $\ell>0$
and rotation $\theta$, and let $L$ be its axis.  If $s\leq\ell$,
then $W_s(P)$ is empty.  If $s>\ell$, put
\[
 I_s=\{k\in\N:k\geq1,\ k\ell<s\},
\]
\begin{equation}
 R_k(s)=\operatorname{arsinh}
 \sqrt{\frac{\cosh s-\cosh(k\ell)}
 {\cosh(k\ell)-\cos(k\theta)}},\qquad
 R_s=\max_{k\in I_s}R_k(s).
 \label{margulis:tube-radius}
\end{equation}
The index set is finite and nonempty, and
$W_s(P)=\{x:d(x,L)<R_s\}$, with $0<R_s<\infty$.
\end{enumerate}
In either nonempty case $W_s(P)$ is connected and proper.  For
$0<s<t$ with $W_s(P)$ nonempty, its boundary has displacement $s$
under some element of $P$, and
$\overline{W_s(P)}\subset W_t(P)$.  In the tube case $R_t>R_s$;
in the cusp case $z_t<z_s$.
\end{lemma}
\begin{proof}
A nonzero discrete additive subgroup of $\R^2$ has a shortest
nonzero vector: choose any nonzero vector and minimize over the
finitely many subgroup elements in the closed ball of that radius.
Finiteness follows from discreteness and the uniform separation
obtained by subtracting two subgroup elements.  For a translation,
\eqref{margulis:halfspace-distance} gives
\[
 \cosh d((x,z),(x+\lambda,z))
       =1+\frac{|\lambda|^2}{2z^2}.
\]
The strict displacement inequality is therefore
$|\lambda|<2z\sinh(s/2)$.  Existence of such a nonzero vector is
equivalent to $z>z_s$, including the strict endpoint convention.

For the loxodromic case normalize $L$ to the vertical axis.  If
$r=d((x,z),L)$, minimization over $(0,w)\in L$ gives
\[
 \cosh r=\min_{w>0}\frac{|x|^2+z^2+w^2}{2zw}
 =\frac{\sqrt{|x|^2+z^2}}{z},\qquad
 \sinh r=\frac{|x|}{z}.
\]
The minimum is attained at $w=\sqrt{|x|^2+z^2}$, as follows by
completing the square or the arithmetic--geometric mean inequality.
For a nonzero integer $k$, substitution of
$g^k(x,z)=(e^{k\ell}A_{k\theta}x,e^{k\ell}z)$ gives
\begin{equation}
 \cosh d((x,z),g^k(x,z))
 =\cosh(k\ell)
  +\bigl(\cosh(k\ell)-\cos(k\theta)\bigr)\sinh^2 r.
 \label{margulis:lox-displacement}
\end{equation}
Indeed $|x-e^{k\ell}A_{k\theta}x|^2$
equals $(1+e^{2k\ell}-2e^{k\ell}\cos(k\theta))|x|^2$;
the remaining vertical term gives $\cosh(k\ell)$.
The coefficient of $\sinh^2r$ is strictly positive.  Consequently
$d((x,z),g^k(x,z))\geq |k|\ell$, with equality on the axis, and
the displacement is strictly increasing with $r\geq0$.
The formulas for $k$ and $-k$ coincide.

Only the finitely many positive integers in $I_s$ can contribute.
For each of them the strict sublevel set is exactly $r<R_k(s)$.
Their union is the open tube of radius $R_s$ in
\eqref{margulis:tube-radius}.  If $s\leq\ell$, no nonzero power
contributes, so even the axis is excluded.  These formulas prove
connectedness and properness in both cases.

For a tube choose $k$ attaining $R_s$.  On its boundary that power
has displacement exactly $s$, and on the whole closed tube it has
displacement at most $s<t$.  Also $k\in I_t$ and $R_k(t)>R_k(s)$,
so $R_t>R_s$.  For a cusp use a shortest vector and the strict decrease
of $z_s$ with $s$.  This proves all the buffer assertions.
\end{proof}
The maximum over powers in \eqref{margulis:tube-radius} is necessary:
with rotation present, it is not justified to compute the thin radius
using only the chosen generator.

\begin{proposition}[Components of the lifted thin set; MGL15]
\label{margulis:lifted-components}
Choose the constant $\mu$ of MGL10.  For $0<s<\mu$ the
nonempty sets $W_s(P)$, $P\in\mathcal E$, are exactly the connected
components of $p^{-1}(H_{<s})$.  More generally, if $0<s,t<\mu$, then
\[
 W_s(P)\cap W_t(Q)\ne\varnothing\quad\Longrightarrow\quad P=Q.
\]
\end{proposition}
\begin{proof}
At a point of the indicated intersection choose nonidentity elements
$g\in P$ and $h\in Q$ with respective displacements less than
$s$ and $t$.  Both belong to $\Gamma_\mu(x)$.  This group is
nontrivial and elementary by MGL10.  Its classification implies
$F(g)=F(h)$.  By MGL13, $P=E(g)=E(h)=Q$.

The injectivity-radius covering formula gives
\[
 p^{-1}(H_{<s})
 =\bigcup_{g\in\Gamma\setminus\{1\}}
       \{x:d(x,gx)<s\}
 =\bigcup_{P\in\mathcal E}W_s(P).
\]
Each region is open by continuity of displacement; each nonempty
region is connected by MGL14.  The intersection assertion makes these
regions pairwise disjoint.  A region is both open and closed relative
to their union, since its relative complement is the union of the
other open regions.  Hence a connected subset of the union cannot
meet both it and its complement.  The regions are precisely the
connected components.  No local-finiteness assumption is used here.
\end{proof}

\begin{lemma}[Precise invariance and covering descent; MGL16]
\label{margulis:descent}
For a nonempty $W_s(P)$ at $0<s<\mu$, its stabilizer in $\Gamma$
is exactly $P$, and it is precisely invariant under $P$.
The induced map
\[
 W_s(P)/P\longrightarrow H
\]
is an injective smooth local isometry and a diffeomorphism onto an
open component $U_s(P)$ of $H_{<s}$.  Components downstairs correspond
to conjugacy classes of the subgroups $P$ with $W_s(P)$ nonempty.
The full inverse image of $U_s(P)$ consists of one $\Gamma$-orbit of
lifted components.
\end{lemma}
\begin{proof}
Invariance of distance under isometries gives
\[
 \gamma W_s(P)=W_s(\gamma P\gamma^{-1}).
\]
If this translate meets $W_s(P)$, MGL15 gives
$\gamma P\gamma^{-1}=P$, and the normalizer conclusion of MGL13
gives $\gamma\in P$.  Every element of $P$ preserves the region.
This proves precise invariance and identifies its stabilizer.

If two points of $W_s(P)$ have the same image in $H$, a deck
transformation carries one to the other.  Precise invariance forces
it to belong to $P$.  The induced quotient map is therefore injective.
It is a local isometry by the inherited quotient-covering construction
and is open, so it is a diffeomorphism onto its open image.

If two such images intersect, lift an intersection point.  The
corresponding regions meet after one is translated by a deck
transformation.  MGL15 identifies their subgroups up to that
conjugation, and equivariance then makes their images equal.
Consequently the distinct images form a disjoint open cover of
$H_{<s}$ by connected sets.  As in MGL15, these are exactly its
components.  The equality
\[
 p^{-1}(U_s(P))=\bigcup_{\gamma\in\Gamma}\gamma W_s(P)
\]
also follows directly from transitivity of the deck action on each
fiber.  MGL15 identifies these translates as the lifted components,
proving the asserted orbit statement.
\end{proof}
The tensor on the quotient agrees with the restricted tensor from
$H$.  No assertion that ambient distances between arbitrary points
of this subset equal its intrinsic path distances is needed.

\begin{lemma}[Two thresholds, closures, and non-merging; MGL17]
\label{margulis:buffer-descent}
Let $0<s<t<\mu$ and $W_s(P)\ne\varnothing$.  Then
$U_s(P)\subset U_t(P)$; distinct $s$-thin components are contained in
distinct $t$-thin components.  Their closures satisfy
\[
 \overline{U_s(P)}=p(\overline{W_s(P)})\subset U_t(P),
\]
and the map $\overline{W_s(P)}/P\to H$ is a smooth embedding with
this image.  Distinct $s$-component closures are disjoint.
\end{lemma}
\begin{proof}
MGL14 gives $\overline{W_s(P)}\subset W_t(P)$.  On the quotient,
the containing $t$-component thus has the same conjugacy class of
full subgroup.  MGL16 classifies components by these conjugacy
classes at both thresholds; distinct $s$-components cannot merge.

It remains to justify closure in the whole manifold, since this does
not follow merely from an embedding of the open region.  If
$q\in\overline{U_s(P)}$, continuity of the injectivity radius gives
$2\rho(q)\leq s<t$.  Thus $q\in H_{<t}$.  The component $U_t(P)$
is closed relative to $H_{<t}$ and contains $U_s(P)$, so $q$ belongs
to $U_t(P)$.  Inside the latter, use the diffeomorphism
$W_t(P)/P\to U_t(P)$.  The $P$-invariant subset
$\overline{W_s(P)}$ is closed in $W_t(P)$, hence its quotient is
closed in $W_t(P)/P$.  Since the quotient map is open, its image is
exactly the relative closure of $W_s(P)/P$.  This proves the displayed
global closure equality.

The closed horoball or closed positive-radius tube is a smooth
submanifold with boundary inside $W_t(P)$.  Its invariant quotient
inherits that structure, and restriction of the open quotient
diffeomorphism proves the claimed embedding.  Finally different
$s$-components have different containing $t$-components; their
closures are contained in these disjoint open sets.
\end{proof}

\begin{lemma}[Tube compactness and embedded product collars; MGL18]
\label{margulis:tube-collars}
A closed tube of finite positive radius modulo its full loxodromic
stabilizer is a compact smooth solid torus, with the metric induced
by the actual generator and its rotational part.  Every nonempty
$s$-component closure at $s<t<\mu$ has a smooth product collar of
its boundary inside its containing $t$-component.  For a parabolic
component the collar has the exact warped cusp metric at negative
as well as positive height, even when its flat section is noncompact.
\end{lemma}
\begin{proof}
For a tube of radius $R$ normalize the generator as in MGL14 and use
coordinates $y=x/z$, $v=\log z$.  Since $\sinh r=|x|/z$, its closed
lift and the deck map become
\[
 \overline D_{\sinh R}\times\R,
 \qquad (y,v)\longmapsto(A_\theta y,v+\ell),
\]
where $\overline D_{\sinh R}$ is the closed Euclidean disk of that
radius.  Substitution in the half-space metric gives the invariant
tensor
\begin{equation}
 |dy+y\,dv|^2+dv^2.
 \label{margulis:twisted-tube-metric}
\end{equation}
The compact set $\overline D_{\sinh R}\times[0,\ell]$ maps onto
the quotient, proving compactness.  The formula
\[
 [(y,v)]\longmapsto
 \bigl(A_{-\theta v/\ell}y,\ [v]\bigr)
 \quad\hbox{in}\quad
 \overline D_{\sinh R}\times(\R/\ell\mathbb Z)
\]
is well-defined and is a smooth diffeomorphism, with inverse obtained
by choosing any representative of $[v]$ and applying
$A_{\theta v/\ell}$ to the disk coordinate.  Changing the
representative applies exactly a deck transformation.  This proves
the solid-torus assertion, including its smooth boundary.  The
diffeomorphism untwists the topology; its differential must still be
applied to \eqref{margulis:twisted-tube-metric}, so no untwisted
product-metric conclusion is drawn.

For nested tubes, $0<R_s<R_t$.  Vary the radial coordinate $r$ in
a small interval around $R_s$ contained in $(0,R_t)$, keeping the
angular coordinate of $y$ and $v$ fixed.  Rotation and translation
in the deck map preserve $r$, so this construction descends to a
smooth product collar of the boundary torus.  MGL17 embeds it in $H$.

For a cusp, $z_t<z_s$ and $u=\log(z/z_s)$.  Choosing
$0<\eta<\log(z_s/z_t)$ gives the chart on $u> -\eta$ and hence a
product collar of $u=0$ inside the larger component.  The translation
action keeps $u$ fixed.  Substituting $z=z_se^u$ in
\eqref{margulis:halfspace-distance} gives the tensor
$du^2+e^{-2u}z_s^{-2}q_{\mathrm{Euc}}$, which descends to the flat
quotient.  This calculation does not require the parabolic
cross-section to be compact; finite volume supplies that condition
in MGL04.
\end{proof}

\lean{The local output records a full subgroup $P$, a region $W_s(P)$,
its quotient identification, the covering map, and its inclusion into a
region at a larger threshold.  The proof uses explicit displacement
identities, finite maxima, open component partitions, and quotient
descent.  MGL17 includes a separate global-closure argument: embedding
an open region alone is insufficient.  These written reductions do
not substitute for a checked binding of MGL10 or of the inherited
model, subgroup, smooth quotient, and covering APIs.}

\subsection{The algebraic Margulis theorem}
\label{margulis:algebra}

\begin{theorem}[Small displacement and elementary groups; MGL10]
\label{margulis:algebra-target}
There exists $\mu>0$ such that, for every discrete subgroup
$\Gamma<\operatorname{Isom}^+(\mathbb H^3)$ and every $x\in\mathbb H^3$,
\[
 \Gamma_\mu(x)=
 \bigl\langle\gamma\in\Gamma:d(x,\gamma x)<\mu\bigr\rangle
\]
is virtually nilpotent.  If $\Gamma$ acts freely, this group is trivial
or elementary.  Every nontrivial discrete torsion-free elementary
subgroup of $\operatorname{Isom}^+(\mathbb H^3)$ is either cyclic
loxodromic or a rank-one or rank-two group of parabolic translations
with a common ideal fixed point.
\end{theorem}
Here \emph{elementary} means having a finite orbit in the closed-ball
compactification.  It is equivalent to preserving a nonempty finite
subset: an orbit in such a subset is finite.  The trivial group is
allowed by this definition, though the conclusions distinguish it.
MGL10 is independent of finite volume.
\begin{proof}
MGL11 is proved below; MGL12 then gives the virtually nilpotent conclusion, with a uniform bound
on the index as well as a uniform displacement constant.  If $\Gamma$
acts freely, its subgroup $\Gamma_\mu(x)$ is discrete and acts freely;
apply MGL25.  The classification assertion is MGL23 and does not
require MGL11.  These applications use the inherited proper evaluation
and transitivity interfaces stated in MGL12, and the individual-isometry
interfaces stated before MGL19.
\end{proof}

\subsection{Elementary groups and the virtually nilpotent reduction}
\label{margulis:elementary-proofs}

The arguments in this subsection use no Margulis constant, no Zassenhaus
neighborhood, and no finite-volume hypothesis.  The inherited geometric
inputs are the single-isometry elliptic/parabolic/loxodromic trichotomy,
conjugacy of boundary fixed sets, the upper half-space normal forms,
and compactness of an interior-point stabilizer.  In particular, an
orientation-preserving isometry fixing $\infty$ has the form
\[
 (x,z)\longmapsto(a+rAx,rz),\qquad
 a\in\R^2,\quad r>0,\quad A\in\operatorname{SO}(2).
\]
The parameterization is continuous in the isometry-group topology.
For a parabolic $r=1$ and the boundary Euclidean isometry has no fixed
point; for a loxodromic with fixed points $0,\infty$, one has $a=0$
and $r\ne1$.  These are individual-isometry inputs, not an assumption
about elementary subgroups.  They must eventually be bound to the
accepted geometric foundation.

\begin{lemma}[Discreteness, freeness, and periodic points; MGL19]
\label{margulis:discrete-powers}
A discrete subgroup of a metrizable topological group meets each
compact subset in finitely many points.  A discrete subgroup
$\Gamma<\operatorname{Isom}^+(\mathbb H^3)$ is torsion-free if and
only if it acts freely.  In that case, for $g\in\Gamma\setminus\{1\}$
and every integer $k\ne0$, $g^k\ne1$ and
\[
 \operatorname{Fix}_{\overline{\mathbb H^3}}(g^k)
   =\operatorname{Fix}_{\overline{\mathbb H^3}}(g)
   =F(g)\subset\partial\mathbb H^3,
\]
where $F(g)$ has one or two points.  Every finite nonempty
$g$-invariant subset of the closed-ball compactification is contained
in $F(g)$.
\end{lemma}
\begin{proof}
Choose an identity neighborhood $U$ meeting the subgroup only in $1$.
If subgroup elements $g_j$ converge in the ambient group, then
$g_j^{-1}g_l\in U$ for all sufficiently large $j,l$, by continuity of
$(a,b)\mapsto a^{-1}b$ at the limiting pair.  Thus the sequence is
eventually constant.  An infinite intersection with a compact set
would have a convergent subsequence of distinct elements, a
contradiction.  This also proves closedness of the subgroup.

If $g$ fixes an interior point, all its powers lie in that compact
point stabilizer.  They form a finite subgroup, so $g$ has finite
order.  Conversely, a nonelliptic isometry has infinite order and the
same boundary fixed set under every nonzero power.  Here is the
normal-form verification in dimension three.  A parabolic fixing
$\infty$ acts by $(x,z)\mapsto(a+Ax,z)$.  If $A\ne I$, then
\[
 \det(I-A)=2-2\cos\theta>0
\]
for its nontrivial rotation angle, so $(I-A)x=a$ has a solution and
the isometry fixes the vertical line through that solution.  This is
impossible for a parabolic.  Hence $A=I$, $a\ne0$, and its $k$th
power is translation by $ka$, with the same unique fixed point
$\infty$.  A loxodromic has the form $(x,z)\mapsto(rAx,rz)$ with
$r\ne1$.  Its $k$th power has scale $r^k\ne1$, no interior fixed
point, and exactly the boundary fixed points $0,\infty$: the equation
$r^kA^kx=x$ forces $x=0$ by taking norms.  Thus a finite-order
isometry is elliptic.  The two implications prove the claimed
equivalence of torsion-freeness and freeness, as well as the power
assertions.  Finally, on a finite invariant set $S$, $g$ acts as a
permutation.  Some positive power fixes every point of $S$; the power
assertion gives $S\subset F(g)$.
\end{proof}

\begin{lemma}[A shared endpoint of a loxodromic; MGL20]
\label{margulis:shared-endpoint}
Let $\Gamma<\operatorname{Isom}^+(\mathbb H^3)$ be discrete.
If $g\in\Gamma$ is loxodromic and $h\in\Gamma$ fixes one of its
ideal endpoints, then $h$ fixes both endpoints.
\end{lemma}
\begin{proof}
Put the common endpoint at $\infty$ and the other endpoint of $g$
at $0$.  Replacing $g$ by its inverse if necessary, write
\[
 g(x,z)=(sBx,sz),\quad 0<s<1,\qquad
 h(x,z)=(a+rAx,rz).
\]
The planar rotations $A,B$ commute.  Consequently
\[
 g^nhg^{-n}(x,z)=(s^nB^na+rAx,rz).
\]
If $a\ne0$, these are distinct, since their translation vectors have
distinct norms $s^n|a|$, but they converge to the isometry
$(x,z)\mapsto(rAx,rz)$.  This contradicts MGL19's discreteness
argument.  Hence $a=0$, proving that $h$ also fixes $0$.
\end{proof}
This proof uses $\operatorname{SO}(2)$ and is deliberately restricted
to oriented dimension three.  It does not assert that an arbitrary
discrete group has a common boundary fixed point.

\begin{lemma}[A discrete free axis stabilizer is cyclic; MGL21]
\label{margulis:axis-cyclic}
Let $P<\operatorname{Isom}^+(\mathbb H^3)$ be a nontrivial discrete
freely acting subgroup preserving a geodesic $L$ setwise.  Then $P$
is infinite cyclic, generated by a loxodromic with positive translation
length and possibly nonzero rotational part.  Every nonidentity
element has exactly the endpoints of $L$ as its fixed set.
\end{lemma}
\begin{proof}
An element interchanging the endpoints restricts to an orientation-
reversing isometry of the real line $L$ and therefore fixes a point
of $L$.  Freeness excludes this.  Put the endpoints at $0,\infty$;
every element then has the form $(x,z)\mapsto(e^tAx,e^tz)$.
The signed axial translation $\tau:P\to\R$, $\tau(g)=t$, is a
homomorphism.  Its kernel fixes the axis pointwise, so it is trivial.

It is necessary to prove that $\tau(P)$ is discrete; discreteness of
an image is not automatic.  For each $R>0$, all elements with
$|\tau(g)|\leq R$ lie in the compact image of
$[-R,R]\times\operatorname{SO}(2)$ in the isometry group.  MGL19
makes this a finite set.  Thus $\tau(P)$ has finitely many points in
each bounded interval.  Choose any positive value $t_0$ in the image
and let $\ell$ be the minimum of its finite nonempty subset
$(0,t_0]\cap\tau(P)$.  This is the least positive value in the whole
image.  For $t\in\tau(P)$ choose $n\in\mathbb Z$ with
$0\leq t-n\ell<\ell$.  The remainder lies in $\tau(P)$ and must
vanish.  Therefore $\tau(P)=\ell\mathbb Z$.  Choose $g_0$ with
$\tau(g_0)=\ell$; injectivity of $\tau$ gives $P=\langle g_0\rangle$.
The normal form and MGL19 give the final fixed-point assertion.
\end{proof}

\begin{lemma}[Discrete translation groups in the plane; MGL22]
\label{margulis:planar-lattice}
Every nonzero discrete additive subgroup $\Lambda<\R^2$ is either
$\mathbb Z v$ for some $v\ne0$, or
$\mathbb Z v\oplus\mathbb Z w$ for linearly independent $v,w$.  In the first
case $\R^2/\Lambda$ is a flat cylinder of infinite area; in the
second it is a compact flat torus.  This conclusion does not assume
in advance that $\Lambda$ is finitely generated or cocompact.
\end{lemma}
\begin{proof}
MGL19 applied to the additive group gives finiteness in every closed
ball.  Choose a shortest nonzero vector $v$ by minimizing among the
nonzero vectors no longer than any one chosen nonzero vector, and
write $a=|v|>0$.  Reducing a multiple of $v$ modulo integer
multiples shows
\[
 \Lambda\cap\R v=\mathbb Z v:
\]
otherwise a nonzero remainder of length less than $a$ would contradict
minimality.  Choose orthonormal coordinates with $v=(a,0)$ and let
$\pi$ be projection onto the second coordinate.

The group $\pi(\Lambda)<\R$ is discrete.  Indeed, for any projected
value $b$, subtract an integer multiple of $v$ from its lift to put
the first coordinate in $[-a/2,a/2]$.  If $0<|b|<a/2$, the resulting
nonzero vector has norm less than $a/\sqrt2<a$, a contradiction.
Thus $0$ is isolated in $\pi(\Lambda)$.  If this projected group is
zero, the rank-one conclusion follows.  Otherwise it has a least
positive element $b_0$: compact finiteness from MGL19 gives a minimum
in any nonempty bounded positive interval reaching one of its
elements.  Division with remainder, as in MGL21, gives
$\pi(\Lambda)=b_0\mathbb Z$.

Lift $b_0$ to $w\in\Lambda$.  For $u\in\Lambda$, subtracting the
appropriate integer multiple of $w$ leaves an element of
$\Lambda\cap\R v=\mathbb Z v$.  The two vectors are independent because
$\pi(w)=b_0\ne0$.  This proves the stated rank-two decomposition.
The linear isomorphism $(s,t)\mapsto sv+tw$ descends to a smooth
identification of the quotient with $(\R/\mathbb Z)^2$.  The closed
parallelogram spanned by $v,w$ maps onto the quotient, proving
compactness; its pulled-back Euclidean metric is a constant positive
definite metric and descends to the flat torus.  In rank one the
chosen orthonormal coordinates identify the metric quotient with
$(\R/a\mathbb Z)\times\R$.  Its disjoint unit-height bands each have
area $a$, proving infinite area.
\end{proof}

\begin{theorem}[Classification of the required elementary groups; MGL23]
\label{margulis:elementary-classification}
Every nontrivial discrete torsion-free elementary subgroup
$\Gamma<\operatorname{Isom}^+(\mathbb H^3)$ is either infinite
cyclic loxodromic or a rank-one or rank-two group of parabolic
translations fixing a common ideal point.  Conversely, every group
of either form is elementary.  In either case all nonidentity
elements have the same boundary fixed set.
\end{theorem}
\begin{proof}
By MGL19 the action is free.  Choose a finite nonempty orbit $S$ in
the closed-ball compactification and any $g\ne1$.  The final
assertion of MGL19 gives $S\subset F(g)$, so $S$ lies at infinity
and has one or two points.  If it has two, $\Gamma$ preserves their
joining geodesic, and MGL21 applies.

If $S=\{\xi\}$, all group elements fix $\xi$.  If the group contains
a loxodromic, MGL20 says that every element also fixes its other
endpoint; again MGL21 applies.  Otherwise every nonidentity element
is parabolic by the single-isometry trichotomy and freeness.  Put
$\xi=\infty$.  The dimension-three normal-form calculation in
MGL19 makes all these maps translations $(x,z)\mapsto(x+\lambda,z)$.
Their vectors form a nonzero additive subgroup $\Lambda<\R^2$.
It is discrete: the continuous translation parameterization pulls
an identity neighborhood isolating $1$ in $\Gamma$ back to a
neighborhood isolating $0$ in $\Lambda$.  Apply MGL22.  The normal
forms give the same fixed set for every nonidentity element.  The
converse follows because the indicated one-point or two-point set
is invariant and contains a finite orbit.
\end{proof}

\begin{lemma}[Finite-index transfer of finite orbits; MGL24]
\label{margulis:finite-orbit-transfer}
Let $D$ act on a set $X$ and let $A<D$ have finite index.  If the
$A$-orbit of $x\in X$ is finite, then its $D$-orbit is finite.
In particular, a finite-index extension of an elementary subgroup
of hyperbolic isometries is elementary.  No normality is required.
\end{lemma}
\begin{proof}
Choose left-coset representatives $d_1,\ldots,d_k$ so that
$D=\bigcup_{i=1}^k d_iA$.  Then
\[
 D x=\bigcup_{i=1}^k d_i(Ax),\qquad |Dx|\leq k\,|Ax|.
\]
The union is finite.  Apply this to the closed-ball compactification
for the second assertion.  This proof does not replace a finite
orbit by a global fixed point.
\end{proof}

\begin{proposition}[Virtually nilpotent implies elementary here; MGL25]
\label{margulis:virtually-elementary}
A discrete freely acting virtually nilpotent subgroup
$D<\operatorname{Isom}^+(\mathbb H^3)$ is either trivial or
elementary.  If nontrivial, it has one of the two forms in MGL23.
\end{proposition}
\begin{proof}
Choose a nilpotent subgroup $A<D$ of finite index; it need not be
normal or finitely generated.  If $A=\{1\}$, then $D$ is finite.
Every element of $D$ has finite order, and MGL19 and freeness imply
$D=\{1\}$.

Suppose $A\ne\{1\}$.  Define its lower central series by
$A_0=A$, $A_{j+1}=[A_j,A]$, where the bracket denotes the subgroup
generated by all commutators of elements of the two groups.
Nilpotence supplies a last nontrivial term $A_j$ with
$[A_j,A]=\{1\}$; hence $A_j$ lies in the center of $A$.
Choose $z\in A_j\setminus\{1\}$.  For every $a\in A$,
$aza^{-1}=z$, so conjugacy of fixed sets gives
\[
 aF(z)=F(aza^{-1})=F(z).
\]
MGL19 makes $F(z)$ a nonempty set of one or two ideal points.
Thus $A$ has a finite orbit in this set and is elementary.  MGL24
makes $D$ elementary.  Freeness implies torsion-freeness by MGL19,
so MGL23 gives its precise form.
\end{proof}

\lean{MGL19--MGL25 are written arguments relative to the specified
individual-isometry and topological-group interfaces.  They do not
depend on MGL10--MGL12 or on the thin-component theorem.  The key
outputs are the full group classification, the actual rank-one or
rank-two translation basis, and finite-orbit transfer without
normality.  No general Bieberbach theorem or finite-generation
theorem for nilpotent groups is needed.  These are proposed proof
decompositions, not checked Lean declarations.}

\subsection{The Zassenhaus theorem}
\label{margulis:zassenhaus-reduction}

\begin{theorem}[Zassenhaus neighborhood; MGL11]
\label{margulis:zassenhaus}
For $G=\operatorname{Isom}^+(\mathbb H^3)$ there is an open identity
neighborhood $U$ such that, for every discrete subgroup $D<G$, the
subgroup generated by $D\cap U$ is nilpotent.
\end{theorem}
\begin{proof}
Take the inverse image of the radius-$1/8$ operator-norm neighborhood
under the hyperboloid representation in MGL29.  That proposition
proves openness and transports the matrix Zassenhaus theorem MGL28.
MGL28 follows from the explicit estimate MGL27 and the group criterion
MGL26.  The neighborhood is independent of the discrete subgroup.
\end{proof}

\subsection{An explicit proof of the Zassenhaus input}
\label{margulis:zassenhaus-proof}

We use the commutator convention $[a,b]=aba^{-1}b^{-1}$.
For a group $L$ put $L_0=L$ and $L_{j+1}=[L_j,L]$, as in MGL25.
Thus $L_c=\{1\}$ means nilpotence of class at most $c$; the case
$c=0$ is the trivial group.  The group argument below applies to
arbitrary generating sets, without finiteness or symmetry assumptions.

\begin{lemma}[Uniform generator commutators imply nilpotence; MGL26]
\label{margulis:commutator-criterion}
Let $S$ generate a group $L$.  Define right-nested commutators by
\[
 C_0(s_0)=s_0,\qquad
 C_{j+1}(s_0,\ldots,s_{j+1})
   =[s_0,C_j(s_1,\ldots,s_{j+1})].
\]
If, for some integer $c\geq0$, every $C_c$ with entries in $S$
equals $1$, then $L_c=\{1\}$.
\end{lemma}
\begin{proof}
We induct on $c$, simultaneously for all groups and generating sets.
For $c=0$ the assumption is $S\subset\{1\}$, so $L=\{1\}$.
Suppose $c\geq1$ and fix a value
$w=C_{c-1}(s_1,\ldots,s_c)$.  For every $s\in S$ the assumption
gives $[s,w]=1$.  The centralizer of $w$ is a subgroup containing
$S$, and hence is all of $L$.  Explicitly, elements commuting with
$w$ are closed under products and inverses, which proves this
subgroup assertion without any closure assumption on $S$.
Thus every such $w$ lies in the center $Z(L)$.

Let $\pi:L\to\overline L=L/Z(L)$ be the quotient map.
The set $\pi(S)$ generates $\overline L$ and every $C_{c-1}$ in
these generators is trivial, since homomorphisms preserve
commutators.  The inductive hypothesis gives
$\overline L_{c-1}=\{1\}$.  For every $j$,
\[
 \pi(L_j)=\overline L_j.
\]
Indeed this holds for $j=0$ by surjectivity.  If it holds for $j$,
the image of the subgroup generated by $[u,v]$, $u\in L_j$, $v\in L$,
is precisely the subgroup generated by
$[\pi(u),\pi(v)]$; surjectivity onto each factor gives the next
equality.  Therefore $L_{c-1}\subset\ker\pi=Z(L)$, and
$L_c=[L_{c-1},L]=\{1\}$, as required.
\end{proof}
The hypothesis has one common depth $c$ for all tuples of generators.
This proof does not infer nilpotence from termination of separately
chosen commutator sequences.

\begin{lemma}[An operator-norm commutator estimate; MGL27]
\label{margulis:matrix-estimate}
Fix an integer $m\geq1$.  Use the Euclidean operator norm on real
$m\times m$ matrices and set $e(A)=\|A-I\|$ for
$A\in\operatorname{GL}_m(\R)$.  If $a=e(A)<1$ and $b=e(B)<1$,
then
\begin{equation}
 \|[A,B]-I\|
 \leq\frac{2ab}{(1-a)(1-b)}.
 \label{margulis:commutator-bound}
\end{equation}
In particular, for $a,b\leq r=1/8$,
\begin{equation}
 e([A,B])\leq\frac{16}{49}\min\{a,b\}
             \leq\frac12\min\{a,b\}.
 \label{margulis:commutator-contraction}
\end{equation}
These inequalities include the cases $A=I$ or $B=I$.
\end{lemma}
\begin{proof}
The operator norm satisfies $\|XY\|\leq\|X\|\|Y\|$ and
$\|I\|=1$.  For every vector $v$,
\[
 \|Av\|\geq\|v\|-\|(A-I)v\|\geq(1-a)\|v\|.
\]
Substitute $v=A^{-1}u$ and take the operator norm to get
$\|A^{-1}\|\leq(1-a)^{-1}$.  The same argument applies to $B$.
Write $X=A-I$, $Y=B-I$.  Direct multiplication gives
\[
 [A,B]-I=(AB-BA)A^{-1}B^{-1}
         =(XY-YX)A^{-1}B^{-1}.
\]
The triangle inequality and submultiplicativity prove
\eqref{margulis:commutator-bound}.  If $a,b\leq r$, then
$ab\leq r\min\{a,b\}$ and
\[
 \frac{2r}{(1-r)^2}=\frac{16}{49}<\frac12,
\]
proving \eqref{margulis:commutator-contraction}.  No inverse bound is
being inferred for a singular matrix: $A,B$ are already invertible.
\end{proof}

\begin{theorem}[A matrix Zassenhaus neighborhood; MGL28]
\label{margulis:matrix-zassenhaus}
For every $m\geq1$, the open identity neighborhood
\[
 U_m=\{A\in\operatorname{GL}_m(\R):\|A-I\|<1/8\}
\]
has the following property.  For every discrete subgroup
$D<\operatorname{GL}_m(\R)$, the subgroup
$L=\langle D\cap U_m\rangle$ is nilpotent.  More precisely, choose
$\delta_D>0$ such that
\[
 A\in D,\quad\|A-I\|<\delta_D\quad\Longrightarrow\quad A=I.
\]
For every integer $c\geq1$ satisfying $2^{-c}/8<\delta_D$,
the lower central term $L_c$ is trivial.
\end{theorem}
\begin{proof}
Continuity of the norm makes $U_m$ an open neighborhood.  Discreteness
isolates $I$ in $D$, giving $\delta_D$.  Put $S=D\cap U_m$ and
$r=1/8$.  Induction using MGL27 gives, for all $j\geq0$ and every
tuple $(s_0,\ldots,s_j)$ in $S$,
\begin{equation}
 C_j(s_0,\ldots,s_j)\in D,\qquad
 e\bigl(C_j(s_0,\ldots,s_j)\bigr)<2^{-j}r.
 \label{margulis:uniform-commutator-depth}
\end{equation}
For $j=0$ this is the definition of $S$.  In the inductive step both
the new outer generator and the inner commutator have defect at most
$r$.  Their commutator lies in $D$, and its defect is at most half
the defect of the inner commutator, giving the strict bound at the
next depth.  If the inner commutator is $I$, the new one is $I$ too,
and the same bound holds.

Choose $c\geq1$ with $2^{-c}r<\delta_D$, possible since
$2^{-c}r\to0$.  Equation~\eqref{margulis:uniform-commutator-depth}
and the defining property of $\delta_D$ make every $C_c$ trivial.
MGL26 gives $L_c=\{1\}$.  The neighborhood was fixed before $D$;
only $\delta_D$ and the nilpotence depth depend on $D$.  The proof
does not require $D$, $S$, or $L$ to be finitely generated.
\end{proof}
The chosen $U_m$ need not be symmetric under inversion; MGL26 does
not require a symmetric generating set.  The later MGL12 proof
chooses its own smaller neighborhood $V$ with $V^{-1}V\subset U$.
The number $1/8$ is a matrix-neighborhood radius, not a claimed
hyperbolic displacement Margulis constant.

\begin{proposition}[Transfer to hyperbolic isometries; MGL29]
\label{margulis:lorentz-transfer}
In the inherited hyperboloid model, the faithful linear representation
\[
 \iota:\operatorname{Isom}^+(\mathbb H^3)
       \longrightarrow\operatorname{GL}_4(\R)
\]
is a topological embedding.  Its inverse image $U=\iota^{-1}(U_4)$
is an open identity neighborhood satisfying MGL11.
\end{proposition}
\begin{proof}
Put $G=\operatorname{Isom}^+(\mathbb H^3)$ and spell out the
topological comparison.  The inherited algebraic
hyperboloid model identifies isometries with the linear Lorentz maps
preserving the upper sheet; restrict that representation to the
orientation-preserving subgroup.  Write the Lorentz form as
\[
 \langle x,y\rangle_L=x_1y_1+x_2y_2+x_3y_3-x_4y_4,
\]
and use the inherited distance identity
$\cosh d(x,y)=-\langle x,y\rangle_L$ on the upper unit hyperboloid.
The isometry group has the compact-open topology, equivalently
uniform convergence on compact sets with the hyperbolic metric.

Matrix convergence $A_j\to A$ implies convergence of the induced
isometries uniformly on each compact set $K$ in the hyperboloid.
Indeed, the Euclidean norms of $x\in K$ are bounded and
\[
 |\langle A_jx,Ax\rangle_L-\langle Ax,Ax\rangle_L|
 \leq\|A_j-A\|\,\|A\|\,\|x\|^2
\]
tends uniformly to zero.  The distance identity and continuity of
$\operatorname{arcosh}$ at $1$ imply uniform convergence of
$d(A_jx,Ax)$ to zero.  This estimate also proves continuity directly
in neighborhoods, without a sequential-topology assumption.

For the inverse, use the four hyperboloid points
\[
 v_0=e_4,\qquad v_i=e_i+\sqrt2e_4\quad(1\leq i\leq3).
\]
They span $\R^4$.  Evaluation at these points is continuous for the
compact-open topology, and the inherited model topology identifies
hyperbolic with Euclidean convergence on the hyperboloid.  The
matrix columns are recovered by
\[
 Ae_4=Av_0,\qquad Ae_i=Av_i-\sqrt2Av_0.
\]
These formulas make the inverse from isometries to their matrices
continuous in the finite-dimensional matrix topology, equivalently
in operator norm.  Thus $\iota$ is a topological embedding.

If $D$ is a discrete subgroup of the isometry group, $\iota(D)$
is a discrete subgroup of $\operatorname{GL}_4(\R)$.  To justify
this point, take an identity neighborhood in the isometry group
meeting $D$ only in $1$; the embedding expresses its image as a
relatively open neighborhood in $\iota(G)$, obtained by intersecting
an ambient open matrix set with $\iota(G)$.  That matrix set meets
$\iota(D)$ only in $I$.

Moreover
\[
 \iota(D\cap U)=\iota(D)\cap U_4,\qquad
 \iota\bigl(\langle D\cap U\rangle\bigr)
   =\langle\iota(D)\cap U_4\rangle.
\]
The second equality follows by mapping finite words in generators
and their inverses in both directions.  MGL28 makes the latter
group nilpotent.  An injective homomorphism identifies a group with
its image and preserves its lower central series, so
$\langle D\cap U\rangle$ is nilpotent.  The inverse image $U$ is
open and contains the identity, proving every assertion.
\end{proof}

\lean{The required matrix interface is the Euclidean operator norm
on finite-dimensional real linear maps, with submultiplicativity,
vector bounds, continuity, and finite-dimensional norm topology.
MGL26 proves the group-theoretic criterion; MGL27--MGL28 prove the
uniform contraction and subgroup-dependent depth; MGL29 checks
topological and generated-subgroup transport.  The algebraic
hyperboloid representation and its distance/model topology remain
inherited bindings.  No general linearity theorem for Lie groups,
projective-matrix lift, or new Margulis axiom is used.}

\subsection{The finite-index step and the completed written chain}
\label{margulis:finite-index-proof}

\begin{proposition}[From Zassenhaus to small displacement; MGL12]
\label{margulis:finite-index}
Use the inherited transitivity of $G$ on $\mathbb H^3$ and properness of
the evaluation map $G\to\mathbb H^3$, $g\mapsto g(o)$.
There are $\mu>0$ and an integer $N\geq1$ such that every
$\Gamma_\mu(x)$ in MGL10 contains a nilpotent subgroup of index at
most $N$.  Both constants are independent of $\Gamma$ and $x$.
\end{proposition}
\begin{proof}
Fix $o$.  Choose an open identity neighborhood $V$ with
$V^{-1}V\subset U$.  The stabilizer $K=G_o$ is compact, so finitely
many left translates $k_1V,\ldots,k_NV$ cover $K$.
There is $\delta>0$ such that
\[
 \{g:d(o,g(o))<\delta\}\subset\bigcup_{j=1}^N k_jV.
\]
Indeed, otherwise choose $g_n$ outside the open union with
$d(o,g_n(o))\to0$.  Eventually these lie in the compact inverse
image of $\overline B(o,1)$.  A convergent subsequence has limit in
$K$ and still outside the union, a contradiction.

Set $\mu=\delta/(N+1)$, $D=\Gamma_\mu(o)$, and
$A=\langle D\cap U\rangle$.  The subgroup $D$ is discrete, so $A$
is nilpotent by MGL11.  Let
\[
 S=\{\gamma\in\Gamma:d(o,\gamma o)<\mu\}.
\]
This set is symmetric and generates $D$.  Consider the graph of left
cosets $gA$ with edges $gA\longmapsto sgA$ for $s\in S$.
It is connected.  If it has more than $N$ vertices, the ball of graph
radius $N$ about $A$ contains at least $N+1$ vertices.  To see this,
each ball which is not the whole graph gains a vertex at the next
radius; if a ball is already the whole graph, it already has more
than $N$ vertices.  Choose representatives $g_0,\ldots,g_N$ for
$N+1$ distinct cosets, each represented by a word of length at most
$N$ in $S$.

For any word $s_1\cdots s_k$ the triangle inequality and invariance of
distance give
\[
 d(o,s_1\cdots s_k o)\leq\sum_{j=1}^k d(o,s_j o).
\]
Thus all $g_i$ lie in the displayed displacement neighborhood of
radius $\delta$.  Two lie in the same $k_jV$, say
$g_i=k_jv_i$ and $g_l=k_jv_l$.  Then
\[
 g_i^{-1}g_l=v_i^{-1}v_l\in V^{-1}V\subset U.
\]
This element also lies in $D$, hence in $A$.  Therefore $g_iA=g_lA$,
a contradiction.  We conclude $[D:A]\leq N$.  Finally, conjugating
an arbitrary point $x$ to $o$ transports the proof without changing
$\mu$ or $N$.
\end{proof}
The neighborhood shrink $V^{-1}V\subset U$ and the use of left cosets
are essential to this argument.  It does not assume $A$ is normal,
and it needs only the word-displacement inequality, not an equality
between products of displacement neighborhoods.

\paragraph{Completed written chain and implementation boundary.}
The chain MGL26--MGL29 $\longrightarrow$ MGL11
$\longrightarrow$ MGL12, together with MGL19--MGL25, proves MGL10.
MGL13--MGL18 then prove MGL01, and MGL02--MGL09 supply the
finite-volume cusp/core theorem and its consequences.  MGL30 gives
countability and compact lifts; MGL31 gives volume-capturing
truncations; MGL32 supplies the orientation-free geometric extension.
The end-to-end review retains strict thin thresholds, full stabilizers,
loxodromic rotation, actual based peripheral maps, uniform commutator
depth, and the exact curvature conversion.

There is no further unproved mathematical producer inside this
Margulis/cusp development beyond the explicit inherited interfaces.
Formalization still requires binding those interfaces to the accepted
foundation and checking the Lean implementations.  The general
variable-curvature and collapse-theoretic Margulis lemmas are separate
contracts.  The downstream rigidity, quantitative stability, and
ambient-flow arguments retain their own obligations.  The companion
register \code{margulis_proof_obligations.csv} and revision113's source
audit record the integrated parent dependencies and remaining bindings.

\subsection{Sources and blueprint integration}
\label{margulis:sources}

The following locators concern the actual archived versions.  PDF page
numbers count from the first physical PDF page.  The companion source
record contains hashes, proof-reading scope, and correction checks.
These citations locate mathematical evidence; they are not imported
Lean proofs.
\begin{enumerate}[label=(\arabic*)]
\item Benedetti--Petronio, \emph{Lectures on Hyperbolic Geometry}
(1992): Theorem D.1.1, printed 134/PDF 145, with proof through
printed 139/PDF 150; D.2.2 and D.2.6, printed 140--143/PDF 151--154;
D.3.3 and its proof, printed 145--150/PDF 156--161;
D.3.12--D.3.14, printed 151--156/PDF 162--167; D.3.18(1),
printed 157/PDF 168.  This is the principal comparison for the cusp
geometry.  We retain loxodromic rotational holonomy when interpreting
D.3.11/D.3.13, and use only the actual injection assertion of D.3.18(1).
\item Martelli, \emph{An Introduction to Geometric Topology},
version 4, September 2025: Definition 1.4.2 and Proposition 1.4.6,
printed 26--27/PDF 34--35, including Proposition 1.4.5;
Proposition 2.1.4, printed 47/PDF 55;
Proposition 2.4.6, printed 70/PDF 78;
Propositions 3.1.8 and 3.1.10, printed 75--76/PDF 83--84;
Section 3.4.8 and Proposition 3.4.12, printed 95--96/PDF 103--104;
Lemma 4.2.2, printed 115--116/PDF 123--124; Remark 4.1.6 and Proposition 4.1.7,
printed 112/PDF 120; Lemma 4.2.7, printed 116--117/PDF 124--125;
Lemma 4.2.8 and
Propositions 4.2.9--4.2.15, printed 116--120/PDF 124--128;
Proposition 4.2.17, Corollary 4.2.18, and Proposition 4.2.20,
printed 121--123/PDF 129--131.  This gives the shortest overall
route.  For MGL30--MGL32 also see Proposition 1.2.19, printed 22/PDF 30;
Proposition 1.2.20, printed 23/PDF 31; Proposition 1.7.3, printed 37/PDF 45;
Proposition 2.2.8, printed 63--64/PDF 71--72;
Exercise 3.1.7 and Section 3.1.4, printed 75--76/PDF 83--84;
and Propositions 6.2.6--6.2.7, printed 162/PDF 170.
The countability exercise is expanded in MGL30; the finite-cover volume
sketch is replaced by a Borel-partition argument in MGL32.
Surface classification is an explicit inherited input used only to name
the nonorientable section.  The invariant buffered-family descent is
a project argument, not an unquoted claim that the source proves this adapter.  MGL12 makes the neighborhood shrink and coset convention in
the finite-index step explicit; MGL03 supplies a separate finiteness
argument with two thresholds.
\item Ratcliffe, \emph{Foundations of Hyperbolic Manifolds},
\emph{first edition, 1994}: Theorem 4.6.1,
printed 136/PDF 147; Theorems 4.7.2--4.7.4,
printed 142--144/PDF 153--155; the matrix norm (5.1.4),
printed 149/PDF 160; Theorem 5.2.5, printed 156/PDF 167;
Section 5.3, Lemmas 1--3 and
Theorems 5.3.1--5.3.2, printed 161--164/PDF 172--175;
Section 5.5 and Theorems 5.5.3--5.5.8,
printed 180--185/PDF 191--196; Section 12.4, Lemma 7 and
Theorems 12.4.2--12.4.3, printed 625--627/PDF 636--638;
Theorems 12.5.1--12.5.2, printed 628--632/PDF 639--643;
Theorem 12.5.4 and Corollary 2, printed 634/PDF 645;
Theorem 12.6.6 and Corollary 4, printed 646--647/PDF 657--658.
The small-neighborhood matrix and coset proofs are particularly
useful for detailed formalization.  MGL20 specializes the contraction
argument in Theorem 5.5.4 to planar rotations; MGL21 expands the
compact-rotation step in the cyclic-axis discussion; MGL22 specializes
Theorem 5.3.2 with an explicit projection gap.  MGL27 uses the Euclidean operator norm and a smaller radius,
not the source's Euclidean norm on matrix entries.  MGL26 supplies
an explicit central-quotient proof of the commutator criterion.
The numbering is not that of the 2019 third edition.
\item Kapovich, \emph{Hyperbolic Manifolds and Discrete Groups},
archived author PDF: Sections 4.12--4.13, especially
Theorems 4.52--4.53 and Lemmas 4.54--4.55, printed 80--84/
PDF 101--105.  This is a detailed Zassenhaus/finite-index comparison.
Theorem 4.52 is compared with the uniform depth estimate in MGL28
and the explicit algebraic induction in MGL26.
The archive filename spells the author name
\code{KapovitchHyperbolicBook.pdf}; the author is Michael Kapovich.
\item Thurston, \emph{The Geometry and Topology of Three-Manifolds},
archived Chapter 5: 5.10.1, 5.10.2, and 5.11.1, printed
113--116/chapter-PDF 31--34.  This is a concise geometric comparison,
not the sole source for the lower algebraic interfaces.
\item Ballmann--Gromov--Schroeder, \emph{Manifolds of Nonpositive
Curvature} (1985): Sections 8--10, particularly 8.3--8.4,
printed 101--102/PDF 115--116; 9.3--9.5,
printed 105--109/PDF 119--123; 10.2--10.4,
printed 110--115/PDF 124--129.  This is the broader
variable-curvature comparison.  Its quantitative commutator,
generator, and parallel-transport estimates invoke Buser--Karcher;
the treatment of unbounded components also refers to Eberlein.
It is not a self-contained replacement for all those inputs.
\end{enumerate}

MSM206's manuscript interfaces are
\code{Nonsingular-Margulis-AlgebraicVersion},
\code{MargeLemmaGeoduckConsequences}, and
\code{theom Form of hyperbolic cusps}; its source notes explicitly
refer to the Benedetti--Petronio chain.  Its cusp formula and boundary
incompressibility discussion explain the hyperbolic-flow consumers.
For KL Section 90 use the curvature-$-1/4$ translation in MGL09;
the displayed (90.10) is not a lemma named ``Lemma 90.10.''

\paragraph{Placement and dependency integration.}
This section follows HG03 in Chapter 8 and precedes HG04.
It supplies HG02's geometric data and HG03's written proof.
The coarse dependency record includes HG03 $\longrightarrow$ HG04,
retaining the path through MPR06--MPR07.  MGL30 supplies countability
and compact lifts; MGL32 covers nonorientable models.  Isometry
properness, boundary-action continuity, and the other rigidity
arguments retain their stated bindings.

HG09 uses MGL08--MGL09 and MGL31's volume margin; HG06 retains its
quantitative core-comparison argument.  Chapter 7 consumes the
complete metric, Chapters 12 and 15 the truncations and collars, and
Chapter 16 the assembled certificate.  Ambient cusp incompressibility
remains in Chapters 8 and 15.  Chapter 13's collapse inputs are separate.
Inherited bindings remain in M2/M2A. Revision113 incorporates the parent
status and dependency edits from the retained integration note. The
33 subordinate rows refine HG03; they do not enlarge the 501-entry
parent catalog or change accepted-proof status. The integrated audit is
recorded in \code{reference_checks_revision113.md}.

% End of the integrated Margulis development; HG04 retains its full scope.


\begin{foundation}[Finite-volume Mostow--Prasad rigidity; HG04]
\label{found:hyp-mostow-prasad}
Let $H$ and $H'$ be connected complete finite-volume hyperbolic
three-manifolds with the same fixed curvature normalization. A homotopy
equivalence $u:H\to H'$ is homotopic to an isometry. In the marked
homotopy class the isometry is unique.
\end{foundation}
The group formulation requires basepoints and paths, or equality modulo
inner conjugacy. It must not claim that an unbased map induces a specified
based homomorphism without those choices. An orientation conclusion needs
the corresponding orientation-compatible comparison; it is not inferred
from an abstract group isomorphism alone.

BP Theorem C.0 proves the compact case; its Theorem C.5.4, printed 124/PDF
136, states the finite-volume extension without supplying its full proof.
MSM206's {\upshape\code{MosTowSection}} states the required finite-volume result.
Theorem~\ref{thm:mostow-full-marked-producer} now supplies its written
finite-volume proof relative to the explicit inherited interfaces.
Accepted-release reuse and exact bindings must still be checked before
Lean implementation. This is not a new axiom or a completed formal proof \cite{BenedettiPetronio1992,MSM206,KleinerLott2008}.

\section{Mostow--Prasad rigidity: top-down proof development}
\label{sec:mostow-top-down}

\paragraph{Target and present status.}
Foundation~\ref{found:hyp-mostow-prasad} is the target, with its full
finite-volume, homotopy-equivalence and uniqueness scope unchanged.
This section supplies its written proof relative to inherited interfaces,
culminating in Theorem~\ref{thm:mostow-full-marked-producer}.
All statements here are mathematical specifications or written arguments,
not elaborated Lean declarations. The accepted PC release must still be
searched before implementing any proposed reproof.

Work at sectional curvature $-1$ and write $X=\mathbb H^3$ and
$\partial X=S^2$. HG01 converts the existing flow normalization.
Use the full group $\operatorname{Isom}(X)$, including orientation-reversing
isometries. Choose universal local-isometric coverings
$p:X\to H$, $p':X\to H'$ and their deck groups $\Gamma,\Gamma'$.
The uniformization and covering identifications are inherited geometric
inputs to bind, not consequences of defining these symbols.
For a homotopy equivalence $u$, a lift $U$ determines an isomorphism
$\phi:\Gamma\to\Gamma'$ with
\[
 U\gamma=\phi(\gamma)U\qquad(\gamma\in\Gamma).
\]
Changing the lift conjugates $\phi$ in $\Gamma'$. Every construction below
retains this actual marking; the final statement is independent of that
choice. No orientation conclusion follows without an orientation hypothesis.

\subsection{The top-level argument and its open inputs}
\label{sec:mostow-assembly}

\begin{proposition}[Conditional HG04 assembly; MPR00]
\label{prop:mostow-conditional-assembly}
For the preceding $u,U,\phi$, suppose the following outputs have been
constructed:
\begin{enumerate}[leftmargin=*]
\item a continuous $\phi$-equivariant quasi-isometry $F:X\to X$;
\item a $\phi$-equivariant boundary homeomorphism $h:\partial X\to\partial X$
associated to $F$;
\item an isometry $A:X\to X$ with boundary map $h$;
\item the trivial-centralizer assertion of MPR04 below.
\end{enumerate}
Then $u$ is homotopic to a unique isometry $H\to H'$.
\end{proposition}
\begin{proof}
MPR10 below shows that boundary equivariance forces
$A\gamma=\phi(\gamma)A$, and that geodesic interpolation between $U$
and $A$ descends to a homotopy from $u$ to the quotient isometry.
MPR46 gives uniqueness using the now-written MPR04 producer. These steps use the same
$\phi$ throughout. In particular, an unmarked assertion that $H$ and $H'$
are isometric would not supply the third output in the required form.
\end{proof}

The four outputs are mathematical obligations, not fields whose mere
declaration closes HG04. MPR61, MPR68 and MPR45 now supply the first,
second and fourth, relative to inherited interfaces. MPR78 now supplies
the third, with the same actual marking. MPR79 assembles the full
written HG04 conclusion; the following target statements remain unchanged.

\begin{foundation}[Controlled representative; MPR01]
\label{found:mostow-controlled-representative}
Given the actual homotopy equivalence and marking above, construct a
continuous $\phi$-equivariant map $F:X\to X$, numbers $\lambda\geq1$,
$C\geq0$, and a quasi-inverse, such that for all $x,y\in X$
\[
 \lambda^{-1}d(x,y)-C\leq d(Fx,Fy)\leq\lambda d(x,y)+C,
 \qquad \forall z\in X\ \exists x\in X:\ d(z,Fx)\leq C.
\]
The inverse comparison must be equivariant for $\phi^{-1}$ and have
explicit coarse inverse bounds. Constants may depend on this comparison.
\end{foundation}
Proposition~\ref{prop:mostow-full-qi-producer} supplies this output
for every finite-volume input. It uses the compact and matched-cusp
construction of MPR54 and the actual marked cusp data supplied by MPR60. An arbitrary
lifted homotopy equivalence is not simply declared to satisfy these
estimates. MPR09 supplies a bilipschitz special case from an already
supplied collar-compatible core diffeomorphism; it does not produce that
diffeomorphism from an arbitrary homotopy equivalence. The new
quasi-isometry producer does not strengthen its output to bilipschitz.

\begin{foundation}[Boundary extension and normalization; MPR02]
\label{found:mostow-boundary-extension}
A quasi-isometry with the stated inverse bounds induces a boundary
homeomorphism, commuting with the specified deck actions. For a fixed
$K\geq1$, a sequence of $K$-bilipschitz homeomorphisms of $X$ whose
boundary maps send three fixed distinct points to three distinct limiting
points has bounded images of a fixed interior basepoint.
\end{foundation}
The bilipschitz extension is Schwartz's Theorem 1.2. His Lemma 1.3
states the three-point criterion for sequences derived from a fixed
equivariant pair. Lemma~\ref{lem:mostow-uniform-tameness} now proves
the uniform bilipschitz formulation above from the uniform Morse and
compact ideal-triangle inputs, following Section 2.1
\cite{SchwartzMostow2026}. MPR38 now constructs the compatible
bilipschitz boundary extension. MPR62--MPR68 in
Section~\ref{sec:mostow-qi-boundary} now supply the general quasi-isometry
boundary extension and uniform normalization; exact inherited bindings
and Lean implementations remain open.
Section~\ref{sec:mostow-boundary-compactness} now proves the needed
bilipschitz boundary compactness and disk control from the stated
uniform Morse and geometric inputs. Both boundary maps and inverses
are controlled. This avoids the compact-quotient orbit argument in
Schwartz's Lemma 1.4 and does not infer boundary convergence from
interior Arzela--Ascoli alone.

\begin{foundation}[Marked boundary rigidity; MPR03]
\label{found:mostow-boundary-rigidity}
For finite-volume deck groups as above, the equivariant boundary map
produced from MPR01--MPR02 is the boundary value of a hyperbolic isometry.
\end{foundation}
Theorem~\ref{thm:mostow-qi-boundary-rigidity} now proves this contract
for continuous controlled quasi-isometries in the full finite-volume
scope, including MPR61's representatives. MPR47 retains its supplied
bilipschitz special case. The selected Schwartz route uses regularity
and four returning zooms, with compact-core recentering in place of
cocompactness. No identification of quasi-isometries with bilipschitz
homeomorphisms is made. All accepted inherited bindings and Lean proofs
remain pending.

\begin{foundation}[Finite-volume centralizer input; MPR04]
\label{found:mostow-centralizer}
If $\Gamma<\operatorname{Isom}(X)$ is a torsion-free finite-covolume
deck group, an isometry commuting with every element of $\Gamma$ is
the identity.
\end{foundation}
Proposition~\ref{prop:mostow-trivial-centralizer} now proves this
assertion using finite-volume group arguments, including the
orientation-reversing case. Proposition~\ref{prop:mostow-homotopic-isometries}
gives the precise homotopy-lifting argument for uniqueness. Standard
inherited geometry, integration and covering bindings and the Lean proof
remain pending; density of loxodromic endpoints is not an extra input.

\subsection{The analytic and zoom obligations}
\label{sec:mostow-analytic-zoom}

The analytic output must leave enough choices of zoom center to exclude
the countable set of lifted cusp centers.

\begin{definition}[Boundary regularity point]
\label{def:mostow-asterisk}
Normalize a boundary homeomorphism by $h(\infty)=\infty$ and regard it
as a map $\mathbb R^2\to\mathbb R^2$. A point $z$ is an asterisk if
every directional derivative
\[
 D_vh(z)=\lim_{t\to0}\frac{h(z+tv)-h(z)}{t},
 \qquad v\in\mathbb Q^2,
\]
exists and $D_{(1,0)}h(z)\ne0$.
\end{definition}

\begin{foundation}[Regularity with room to avoid cusp centers; MPR05]
\label{found:mostow-regularity}
For a normalized boundary map of a $K$-bilipschitz homeomorphism of $X$,
prove that its asterisk set has positive planar outer measure.
The statement must not assume a compact quotient or equivariance.
\end{foundation}
Schwartz's Theorem 1.7 and Section 3, especially Lemma 3.1 and its
countable-direction argument, are the template. Sections 4--5 provide
the measure-theoretic details. The nonzero horizontal derivative occurs
on a set of positive outer measure; removing the countable union of
directional exceptional null sets preserves that positivity. This
stronger output, rather than existence of one unspecified point, permits
removal of the countable cusp-center set. The compactness and disk steps now have written conditional proofs
in MPR20--MPR24, and MPR25 supplies the finite packing inequality.
MPR26--MPR32 now prove the remaining project-specific analytic
implication from that disk control, using the stated inherited analysis
inputs; MPR33 feeds it into the bilipschitz boundary argument.
Use inherited Lebesgue measure, Fubini, absolute continuity and interval
calculus wherever their audited scope suffices. Public mathlib candidates
recorded in the research comparison are not yet project-checked imports.

\paragraph{Returning zooms; MPR08.}
\label{par:mostow-returning-zooms}
For a finite point $z$ and any real $s>0$, set
\[
 b_s(w,t)=\bigl(z+s^{-1}(w-z),s^{-1}t\bigr),\qquad
 a_s(w,t)=\bigl(h(z)+s(w-h(z)),st\bigr).
\]
These are hyperbolic isometries. With $o=(z,1)$ and $s=e^r$,
$b_s(o)=(z,e^{-r})$ follows the unit-speed ray to $z$.
If $z$ is not a lifted cusp center, MPR06 gives arbitrarily late $r$
for which $p(b_{e^r}(o))$ lies in a fixed compact core. At an asterisk,
the three boundary images of $z,z+1,z+2$ under $a_s h b_s$ tend to
$h(z),h(z)+D_1h(z),h(z)+2D_1h(z)$, which are distinct; here
$D_1=D_{(1,0)}$.
Thus MPR02 supplies tameness and MPR07 supplies an orbit-preserving
subsequence along these returning times. Integer zoom scales are not
required; rounding the return times is unnecessary.

The line-family and affine steps of MPR08 are proved conditionally in
Section~\ref{sec:mostow-zoom-proof}. MPR12--MPR17 provide the explicit
limit arguments; MPR18 assembles four returning zooms; MPR19 supplies
the three-point normalization adapter from its geometric inputs.
Revisions108--110 supplied the stated bilipschitz regularity,
finite-volume no-common-fixed-point and marked-uniqueness arguments.
Revision111 supplied compact and conditional matched-cusp quasi-isometry
representatives. MGL00/MGL30/MGL32 supply the individual cusp/core
and compact-lift geometry. Revision114 now supplies the arbitrary-marking
cusp matching and full controlled representative in
Section~\ref{sec:mostow-cusp-matching}. Revision115 supplies the general
quasi-isometry boundary and normalization producer MPR68 in
Section~\ref{sec:mostow-qi-boundary}. Revision116 supplies the corresponding
disk control and regularity in Section~\ref{sec:mostow-qi-regularity}.
Revision117 supplies the general quasi-isometry returning-zoom/orbit
adapter and marked assembly in Section~\ref{sec:mostow-qi-rigidity}.
These are explicit project extensions of Schwartz's short cusp remark.

\subsection{Return to the core and preservation of the isometry orbit}
\label{sec:mostow-cusp-return}

The next two lemmas are written project arguments. MGL00 and MGL32
now supply their finite-volume geometric cusp/core inputs, including
lifted horoballs, respectively with and without orientability. MGL30
supplies countability and compact lift sets. HG02's abstract record
alone does not prove these facts. Their accepted inherited bindings,
isometry-evaluation properness and compactification-action continuity
remain explicit implementation inputs. This bookkeeping update does not
change either lemma or produce matched cusp data for an arbitrary
homotopy equivalence.

\begin{lemma}[Non-cusp endpoints return to the core; MPR06]
\label{lem:mostow-core-return}
Let $p:X\to H$ be a regular hyperbolic covering with countable deck
group. Suppose $K\subset H$ is compact and $H\setminus K$ is a finite
disjoint union of open cusp neighborhoods, each with lifts that are
disjoint horoballs, forming one deck orbit for that cusp. Let
$P\subset\partial X$ be the set of centers of all these horoballs.
Then $P$ is countable. For every unit-speed geodesic ray
$r:[0,\infty)\to X$ with endpoint outside $P$,
\[
 \forall T\geq0\ \exists t\geq T:\quad p(r(t))\in K.
\]
Consequently such a ray has a sequence of return times tending to infinity.
\end{lemma}
\begin{proof}
Choose one lifted horoball for each of the finitely many cusps.
Every lifted center is the image of one of their centers under a deck
transformation. A finite union of images of a countable group is
countable, proving the first assertion.

If the displayed assertion fails, some tail of $p\circ r$ is contained
in $H\setminus K$. A connected tail lies in a single cusp because this
complement is a disjoint union of open components. Its lift, the actual
tail of $r$, lies in one component of the inverse image and hence in
one horoball. Move its center to $\infty$ in the upper half-space
model. A ray that remains in $\{t>a\}$ for all sufficiently large
parameters must be the upward vertical ray: a downward vertical ray
has height tending to zero, and a ray on an orthogonal semicircle also
has height tending to zero at its forward endpoint. These exhaust the
hyperbolic geodesics. Thus the endpoint is the center of that horoball,
contrary to the hypothesis. Finally choose successive return times
larger than both their index and the preceding time plus one.
\end{proof}
The empty cusp family is permitted: then $K=H$ and every ray returns.
This argument gives arbitrarily late returns, not a uniform return-time
bound or a proportion of time spent in the core. For a discrete deck
subgroup of the second-countable isometry group, countability follows
by assigning to each element a distinct member of a countable basis
isolating it. MGL30 now writes the covering/countability adapter; its inherited
interfaces must still be bound in Lean.

\begin{lemma}[Returning orbit extraction; MPR07]
\label{lem:mostow-returning-orbit}
Let $F:X\to X$ be a homeomorphism extending continuously to a boundary
homeomorphism $h$, equivariant for a fixed $\phi:\Gamma\to\Gamma'$. Let
$a_n,b_n\in\operatorname{Isom}(X)$ and $F_n=a_nFb_n$.
Assume $\{F_n(o)\}$ is bounded and $p(b_n(o))\in K$ for all $n$,
where $K$ is compact. Suppose a compact set $D\subset X$ contains a
deck translate of every point of $p^{-1}(K)$.
Then, on a subsequence, there are $a,b\in\operatorname{Isom}(X)$ such that
\[
 F_n\longrightarrow aFb,
 \qquad \partial a_n\,h\,\partial b_n\longrightarrow
 \partial a\,h\,\partial b.
\]
The first convergence is uniform on compact subsets of $X$; the second
is uniform on $S^2$.
Here use the standard properness of evaluation
$\operatorname{Isom}(X)\to X$, $g\mapsto g(o)$, and continuity of the
isometry action on the closed-ball compactification.
\end{lemma}
\begin{proof}
Choose $\gamma_n\in\Gamma$ with $\gamma_nb_n(o)\in D$ and put
\[
 B_n=\gamma_nb_n,\qquad A_n=a_n\phi(\gamma_n)^{-1}.
\]
Equivariance gives $F_n=A_nFB_n$, with the same identity at infinity.
The points $B_n(o)$ belong to the fixed compact set $D$. Properness of
evaluation gives a convergent subsequence of $B_n$, with limit an
isometry $b$. Since $F(D)$ is compact and
\[
 d(A_n(o),F_n(o))=d(o,F(B_n(o))),
\]
the points $A_n(o)$ are bounded. Properness of $X$ and of evaluation
therefore give a further subsequence $A_n\to a$ in the isometry group.
Composition with the fixed continuous $F$ gives uniform convergence on
each compact subset: the $B_n$ images of such a subset lie in one
compact set. On the sphere the isometries converge uniformly, and $h$
is uniformly continuous, giving the second convergence assertion.
Both limits are the displayed compositions; they are homeomorphisms.
\end{proof}

MGL30 supplies the compact $D$ required here. Its construction covers $K$ with finitely
many relatively compact evenly covered neighborhoods, lifting one
slightly larger neighborhood for each, and taking the union of the
lifted compact closures. Regularity of the covering supplies the deck
translates. No compactness of the whole quotient was used.
Properness of isometry evaluation is the usual hyperbolic-space fact
that a compact set of basepoint images and the compact orthogonal
frame stabilizer give a compact set of isometries. These covering,
frame and compactification facts are explicit inherited binding inputs.

The boundary limit $h'=\partial a\,h\,\partial b$ conjugates
$b^{-1}\Gamma b$ to $a\Gamma'a^{-1}$. Its lifted cusp centers are
$\partial b^{-1}(P)$, still countable. If $h'$ extends to an isometry,
then so does $h=\partial a^{-1}h'\partial b^{-1}$.
Convergence to an unspecified boundary homeomorphism would not retain
this information. This is the needed restricted substitute for
Schwartz's compact-quotient Lemma 1.4, not an assertion that every
normalized orbit sequence for a cusped lattice has an orbit limit.

\subsection{The returning-zoom proof with its remaining inputs}
\label{sec:mostow-zoom-proof}

We now discharge the line-limit and affine parts of MPR08 as written
arguments. The following proposition states the resulting boundary
theorem first; its proof uses the detailed lemmas below. MGL00/MGL32
and MGL30 supply the cusp/core and compact-lift data, and MPR41 supplies
the finite-volume no-fixed-boundary-point input. Isometry properness
and compactification-action continuity remain inherited bindings. MPR33 now
supplies its regularity hypothesis under the uniform disk/geometry
and inherited-analysis assumptions stated there.
MPR74--MPR79 later supply the broader quasi-isometry argument and
full HG04 assembly; the bilipschitz statement below is retained.

\begin{proposition}[Conditional bilipschitz boundary rigidity; MPR18]
\label{prop:mostow-bl-boundary-assembly}
Let $F:X\to X$ be a $K$-bilipschitz homeomorphism, equivariant for a
group isomorphism $\phi:\Gamma\to\Gamma'$, with boundary homeomorphism $h$ fixing
$\infty$. Suppose:
\begin{enumerate}[leftmargin=*]
\item the domain quotient has the geometric compact core, countable
cusp-center set and compact lift data of MPR06--MPR07, with the
isometry properness and boundary-action inputs stated there;
\item the uniform Morse and ideal-triangle inputs of MPR19 hold;
\item for every $a,b\in\operatorname{Isom}(X)$ for which
$\partial a\,h\,\partial b$ fixes $\infty$, this normalized boundary
map has an asterisk set of positive planar outer measure;
\item $\Gamma$ has no common fixed point on $S^2$, that is,
$\forall\xi\in S^2\ \exists\gamma\in\Gamma:\gamma\xi\ne\xi$.
\end{enumerate}
Then $h$ is the boundary map of an isometry of $X$.
\end{proposition}
\begin{proof}
Choose two distinct directions $D_1,D_2$ in the plane. The countable
set $P$ of lifted cusp centers is null, so the positive-outer-measure
asterisk set contains $z\notin P$. MPR13 gives a returning zoom limit
with every line through $z$ affine. MPR14, with a second center outside
the exceptional set for its actual conjugated domain lattice, gives a
second limit $h_2$ for which $D_1$ is good.

Both limits are of the form $\partial a\,h\,\partial b$, by MPR07,
and fix $\infty$. Thus the third hypothesis applies to $h_2$, not
merely to the initial $h$. Repeat MPR13--MPR14 with direction $D_2$.
MPR15 preserves $D_1$ at both further limits. The final homeomorphism
$h_4$ has two good directions and is real affine by MPR16.

Compose the four orbit identities to write
$h_4=\partial a\,h\,\partial b$. Its domain lattice is
$b^{-1}\Gamma b$. Apply the fourth hypothesis at $\partial b(\infty)$:
some $b^{-1}\gamma b$ moves $\infty$. MPR17 therefore extends $h_4$
to an isometry $A_4$. The isometry $a^{-1}A_4b^{-1}$ has boundary
value the original $h$. MPR10 subsequently gives descent in the
original marking whenever an actual comparison map supplies that marking.
\end{proof}

The geometric inputs transport under conjugation: $b$ induces an
isometry $X/(b^{-1}\Gamma b)\to X/\Gamma$, taking the new compact
core to the old one. Lifted horoballs and compact lift sets pull back
by $b^{-1}$; the new center set is $\partial b^{-1}(P)$.
Composition with isometries preserves the same bilipschitz constant.
All limits used here fix $\infty$ because each zoom does and the
boundary convergence is uniform. Thus their plane restrictions are
homeomorphisms and the hypotheses needed at the next step are preserved.

MPR32 now supplies the third hypothesis through MPR24, as recorded
in MPR33; its geometric and inherited-analysis inputs must still be
provided or bound. MPR41 supplies the fourth from finite covolume,
including for each conjugated domain group.
An element moving the original $\infty$ alone is insufficient: the
last lattice is conjugated. No uniqueness conclusion or controlled
representative for an arbitrary homotopy equivalence is proved by MPR18.

\paragraph{Affine lines and good directions.}
For $v\ne0$ a restriction to $x+\mathbb Rv$ is affine if it has the
form $t\mapsto p+tq$. For a plane homeomorphism $q\ne0$ automatically.
A direction $\mathbb Rv$ is good if every parallel line has such a
restriction. These definitions use the full real line and its affine
parameter, not merely a line segment or a set-theoretic line image.

\begin{lemma}[Affine restrictions on moving lines; MPR12]
\label{lem:mostow-moving-lines}
Suppose continuous maps $h_n:\mathbb R^2\to\mathbb R^2$ converge
uniformly on compact sets to a homeomorphism $h$. Let $x_n\to x$ and
$v_n\to v\ne0$, and suppose $t\mapsto h_n(x_n+tv_n)$ is affine
for all sufficiently large $n$. Then, for every $t\in\mathbb R$,
\[
 h(x+tv)=h(x)+t\bigl(h(x+v)-h(x)\bigr).
\]
Its slope is nonzero.
\end{lemma}
\begin{proof}
For each fixed $t$, the convergent sequence $x_n+tv_n$ and its limit
lie in a compact set. Local uniform convergence and continuity of $h$
give $h_n(x_n+tv_n)\to h(x+tv)$. The affine identity is
\[
 h_n(x_n+tv_n)=h_n(x_n)
       +t\bigl(h_n(x_n+v_n)-h_n(x_n)\bigr).
\]
Pass to the limit using $t,0,1$. The output slope is nonzero because
$x+v\ne x$ and $h$ is injective. A single locally uniformly convergent
subsequence works for every $x,t$ to which this argument is applied;
no new subsequence is chosen for a line or parameter.
\end{proof}

For application to boundary maps, uniform convergence on $S^2$ to a
homeomorphism fixing $\infty$ gives local uniform Euclidean convergence
on the plane. Indeed, the image of a fixed compact plane set under the
limit has positive spherical distance from $\infty$. Eventually all
its images lie in one compact subset of the same finite coordinate
chart. Spherical and Euclidean metrics are uniformly comparable there.
This chart argument is required; uniform spherical convergence alone
does not give a uniform Euclidean bound on the whole plane.

\begin{lemma}[A returning zoom produces an affine pencil; MPR13]
\label{lem:mostow-asterisk-pencil}
Let $(F,h)$ be a normalized equivariant bilipschitz pair with the
geometric inputs of MPR06--MPR07 and the tameness conclusion of MPR19.
If $z$ is an asterisk and is not in its domain cusp-center set, there
is a returning zoom limit $(F_1,h_1)$ in the isometry orbit of $(F,h)$
such that $h_1$ fixes $\infty$ and every line through $z$ is affine
for $h_1$.
\end{lemma}
\begin{proof}
Use MPR06 to choose return times $r_n\to\infty$ for the vertical ray
$r\mapsto(z,e^{-r})$ and set $s_n=e^{r_n}$. For the zooms already
defined in MPR08, the finite boundary formula is
\[
 h_n(w)=h(z)+s_n\bigl(h(z+(w-z)/s_n)-h(z)\bigr).
\]
The images of $z,z+1,z+2$ converge to three distinct points because
$D_1h(z)\ne0$. MPR19 gives tameness; MPR07 gives a subsequence with
$h_n\to h_1$ uniformly on $S^2$ and an actual orbit identity.
In particular $h_1$ is a homeomorphism fixing $\infty$.

Fix a rational vector $v\ne0$ and a real $t$. The directional
derivative, with parameter $t/s_n$, gives
\[
 h_1(z+tv)=h(z)+tD_vh(z).
\]
For $t=0$ the formula is exact before passage to the limit. For
$t\ne0$ it is the defining derivative limit, multiplied by $t$;
negative $t$ is allowed. Injectivity of $h_1$ makes its slope nonzero.
Every rational direction through $z$ is consequently affine.

For an arbitrary $v\ne0$ choose nonzero rational vectors $v_j\to v$.
Apply MPR12 with the constant sequence of maps $h_1$, constant center
$z$, and directions $v_j$. This proves affinity on $z+\mathbb Rv$.
No assertion of full differentiability of the original $h$ was used.
\end{proof}

\begin{lemma}[A pencil gives a prescribed good direction; MPR14]
\label{lem:mostow-pencil-parallel}
Let a normalized pair $(F,h)$ have the geometric and tameness inputs
of MPR13, and suppose every line through $z$ is affine for $h$.
For any direction $D$ there is a returning zoom limit in its isometry
orbit for which $D$ is good.
\end{lemma}
\begin{proof}
The line $z+D$ is uncountable, whereas the current domain cusp-center
set and $\{z\}$ are countable. Choose $w\in z+D$ outside that union
and put $v=w-z\ne0$. Zoom at $w$ through its returning scales
$s_n\to\infty$. On $z+D=w+\mathbb Rv$ the map $h$ is affine, so
the formula for $h_n$ equals $h$ on this entire line. The images of
$w,w+v,w+2v$ are fixed distinct points. MPR19 and MPR07 therefore
give a limiting homeomorphism $h_1$ in the same isometry orbit.

The map $h_n$ is affine on every line through
\[
 c_n=b_{s_n}^{-1}(z)=w+s_n(z-w)=w-s_nv.
\]
Fix an arbitrary $x\in\mathbb R^2$ and put
$v_n=v+(x-w)/s_n$. Then $v_n\to v\ne0$ and
$x-c_n=s_nv_n$. For all sufficiently large $n$, the line
$x+\mathbb Rv_n$ thus contains $c_n$, and $h_n$ is affine on it.
Local uniform convergence and MPR12 give affinity of $h_1$ on
$x+\mathbb Rv$. Since $x$ was arbitrary, $D$ is good for $h_1$.
The orbit subsequence was fixed before choosing $x$.
\end{proof}

\begin{lemma}[Good directions survive zooms and their limits; MPR15]
\label{lem:mostow-preserve-directions}
If $D$ is good for $h$, it is good for every plane zoom of $h$ by
positive homotheties, and for every locally uniform limit of these
zooms that is a plane homeomorphism.
\end{lemma}
\begin{proof}
A positive homothety sends a line of direction $D$ to a line of that
same direction and acts affinely on its parameter. Pre- and
post-composition therefore preserve the affine restriction. For a
fixed line in a limiting map apply MPR12 with constant $x_n=x$ and
$v_n=v\in D\setminus\{0\}$. Injectivity of the limit preserves the
nonzero slope. Repeat for each line, using the already fixed limit.
\end{proof}
This assertion concerns the actual homothety zoom sequences. Arbitrary
M\"obius changes of coordinates need not preserve a Euclidean direction.
The auxiliary isometries in MPR07 identify the limiting orbit; they do
not replace the zoom sequence in this good-direction argument.

\begin{lemma}[Two good directions force an affine map; MPR16]
\label{lem:mostow-two-directions-affine}
Let $h:\mathbb R^2\to\mathbb R^2$ be a homeomorphism with two
distinct good directions, spanned by linearly independent vectors
$v,w$. Put $p=h(v)-h(0)$ and $q=h(w)-h(0)$. Then $p,q$ are linearly
independent and
\[
 h(sv+tw)=h(0)+sp+tq\qquad(s,t\in\mathbb R).
\]
Thus $h$ is an invertible real affine map.
\end{lemma}
\begin{proof}
The two axes $\mathbb Rv,\mathbb Rw$ map onto entire lines, since
their affine slopes are nonzero. If their image directions were the
same, both image lines, meeting at $h(0)$, would coincide. A point
different from $h(0)$ would then have preimages on both axes,
contradicting injectivity. Hence $p,q$ are independent.

Distinct parallel domain lines are disjoint. Their full image lines
are disjoint by injectivity, and disjoint lines in the plane are
parallel. The line $sv+\mathbb Rw$ therefore maps to the line
through $h(sv)=h(0)+sp$ parallel to $q$; similarly
$tw+\mathbb Rv$ maps to the line through $h(0)+tq$ parallel to $p$.
Their unique intersection is $h(0)+sp+tq$ and contains $h(sv+tw)$.
The displayed formula follows for every pair $s,t$.
\end{proof}

Set-theoretic images of lines would not suffice: the homeomorphism
$(x,y)\mapsto(x,y^3)$ preserves both coordinate-parallel line families
but is not affine. MPR16 uses affine parametrizations on the full lines.

\begin{lemma}[Affine equivariance supplies an isometry; MPR17]
\label{lem:mostow-affine-isometry}
Suppose a real affine plane homeomorphism $h(x)=Lx+d$, extended by
$h(\infty)=\infty$, is equivariant for a fixed isomorphism
$\phi:\Gamma\to\Gamma'$ of hyperbolic isometry groups.
If some $\gamma\in\Gamma$ moves $\infty$, then
$L=cQ$ for $c>0$ and an orthogonal real linear map $Q$. Consequently
\[
 (x,t)\longmapsto(Lx+d,ct)
\]
is a hyperbolic isometry with boundary value $h$.
\end{lemma}
\begin{proof}
The point $\gamma^{-1}(\infty)$ is finite. Choose a Euclidean line
$\ell$ avoiding it. Since hyperbolic boundary isometries preserve
generalized circles, $\gamma(\ell\cup\{\infty\})$ is a finite
round circle $C$. Equivariance for this actual $\gamma$ gives
\[
 h(C)=\phi(\gamma)\bigl(h(\ell)\cup\{\infty\}\bigr).
\]
The right side is a generalized circle. The left side is a compact
subset of the plane, so it is a finite round circle. Thus $L$ sends
a round circle to a round circle, up to translation.

Let the input center and radius be $a,r$ and the output center and
radius be $a',r'$. Central symmetry of the input circle makes its
image centrally symmetric about $h(a)$. A bounded circle cannot have
two different symmetry centers: composing their central reflections
would give a nonzero translation preserving that bounded set, and
iteration would contradict boundedness. Hence $a'=h(a)$.
It follows that $\|Lu\|=r'/r=:c>0$ for every unit vector $u$.
Homogeneity gives $\|Lx\|=c\|x\|$ for every $x$; polarization
then shows $Q=c^{-1}L$ is orthogonal. In the upper half-space metric
$(\|dx\|^2+dt^2)/t^2$, the displayed extension multiplies numerator
and denominator by $c^2$. It is bijective with the analogous inverse,
so it is an isometry. Orthogonal $Q$ of either orientation is allowed.
\end{proof}

This expands Schwartz's Lemmas 1.5--1.6 with the lattice choice and
the affine-to-similarity calculation visible. MPR41 supplies the
finite-volume no-common-fixed-point input used by MPR18; choosing an
arbitrary ambient isometry in MPR17 would not justify equivariance.

\begin{lemma}[Uniform bilipschitz three-point tameness; MPR19]
\label{lem:mostow-uniform-tameness}
Let $F_n:X\to X$ be $K$-bilipschitz homeomorphisms with boundary
homeomorphisms $h_n$, with $K$ independent of $n$. Assume the uniform
Morse input that $F_n$ sends every complete geodesic into the closed
$M(K)$-neighborhood of the geodesic joining its boundary endpoints.
Assume also that the intersection of fixed-radius closed tubes about
the three sides of any ideal triangle is compact. If the images under
$h_n$ of three fixed distinct boundary points converge to three distinct
points, then $\{F_n(o)\}$ is bounded for every fixed $o\in X$.
\end{lemma}
\begin{proof}
There are boundary isometries $m_n$ tending to the identity and taking
the three image points to their three limits. One construction uses
the unique orientation-preserving M\"obius map between ordered triples:
in a finite chart containing the three limiting points, the coefficients of
$z\mapsto (z-a)(b-c)/((z-c)(b-a))$ vary continuously with the
distinct triple $(a,b,c)$. Compose its varying and limiting versions
to get $m_n$ and their hyperbolic extensions.

Set $G_n=m_nF_n$. The target ideal triangle is now fixed, and $G_n$
has the same bilipschitz constant and Morse bound. On each of the
three fixed domain geodesics choose one point $q_{ij}$. For fixed $o$,
the point $G_n(o)$ has distance at most
\[
 R=M(K)+K\max_{ij}d(o,q_{ij})
\]
from every corresponding target geodesic. All $G_n(o)$ therefore
lie in the compact intersection of the three closed $R$-tubes.
The maps $m_n^{-1}$ form a convergent sequence of isometries, so their
images of that compact set lie in a bounded set, by continuity of
the isometry action. This bounds $F_n(o)$. Finite initial indices, if
discarded when choosing the chart, do not affect boundedness.
\end{proof}
This is the scope adapter extracted from Schwartz's Lemma 1.3 and
Section 2.1 proof, with the distance from $o$ to the domain sides
made explicit. No equivariance or compact quotient is used. The
uniform Morse, boundary-extension and compact triple-tube inputs
are supplied for this bilipschitz class by MPR34--MPR39 below,
conditional on inherited model/calculus bindings;
MPR19 alone does not prove normalized boundary compactness or MPR05;
MPR22 below supplies the former from the additional inverse estimates.

\subsection{Boundary compactness and disk control without cocompactness}
\label{sec:mostow-boundary-compactness}

The next result supplies the geometric disk input to MPR05. Its proof
uses only uniformly bilipschitz maps of hyperbolic space, their supplied
boundary homeomorphisms, and the geometric inputs specified below. No
lattice or equivariance is needed. MPR34--MPR39 below supply boundary extension, uniform Morse and triple
tubes for this bilipschitz class, using inherited model/calculus inputs.
This subsection does not assert the general quasi-isometry version of MPR02.

\begin{proposition}[Uniform disk control; MPR24]
\label{prop:mostow-uniform-disks}
Fix $K\geq1$. Assume the standard curvature $-1$ ball/half-space model
identifications, the uniform Morse bound $M(K)$ for all $K$-bilipschitz
homeomorphisms and their inverses with the supplied boundary maps, and
the compact ideal-triangle tube input of MPR19. There is $C=C(K)>1$
such that, for every such boundary map $h$ fixing $\infty$ and every
Euclidean disk of center $c$ and radius $r>0$, there are concentric
closed disks $D_1,D_2$, centered at $h(c)$ and of positive radii, with
\[
 D_1\subset h(B(c,r)),\qquad
 h(\overline B(c,r))\subset D_2,\qquad
 \frac{\operatorname{diam}D_2}{\operatorname{diam}D_1}\leq C.
\]
The constant is chosen before the map, center and radius. Its dependence
on the fixed geometric inputs is suppressed in the notation $C(K)$.
\end{proposition}
\begin{proof}
Normalize each map and disk by boundary similarities:
\[
 g(z)=\frac{h(c+rz)-h(c)}{h(c+r)-h(c)}.
\]
The denominator is nonzero. Complex division here includes a rotation;
both similarities extend to hyperbolic isometries, so the corresponding
interior map still has constant $K$. The normalized map fixes
$0,1,\infty$ exactly. Consider the family of all these normalized maps
for the fixed $K$, allowing the original map to vary.

If no common ratio bound existed, choose a sequence $g_n$ for which no
pair of disks as above for the unit disk has ratio at most $n$. MPR19
gives a common bound for the image basepoints of this sequence. MPR22
below then gives, on one subsequence, uniform spherical convergence of
$g_n$ and $g_n^{-1}$ to mutually inverse homeomorphisms $g$ and $g^{-1}$.
Both fix $0,1,\infty$. MPR23 supplies fixed radii $0<\rho<R$ with
\[
 \overline B(0,\rho)\subset g_n(B(0,1)),\qquad
 g_n(\overline B(0,1))\subset\overline B(0,R)
\]
eventually on that subsequence. Its original indices tend to infinity,
contradicting the fixed ratio $R/\rho$. Thus a finite common bound exists;
increase it if necessary to make $C>1$. Undoing the target similarity
preserves the ratio and centers the disks at $h(c)$. The domain
similarity sends the open and closed unit disks to the indicated disks.
\end{proof}

This is a uniform version of Schwartz's Disk Theorem 3.6, printed/PDF
page 17 \cite{SchwartzMostow2026}. His displayed argument uses a fixed
equivariant pair and a boundary compactness step after tameness. The
proof above replaces that step by MPR20--MPR23. No limit in an isometry
orbit is asserted or needed here; the marked orbit statement MPR07 is
a different result used by the returning zooms.

\begin{lemma}[Visual distance and a complete geodesic; MPR20]
\label{lem:mostow-visual-distance}
Identify $X$ with the curvature $-1$ Poincare ball with basepoint $o=0$.
On its boundary put $\delta(\xi,\eta)=\tfrac12|\xi-\eta|$, so
$0\leq\delta\leq1$. For distinct boundary points, let $\ell_{\xi\eta}$
be the complete geodesic joining them and put $a=d(o,\ell_{\xi\eta})$.
Then
\[
 \delta(\xi,\eta)=\frac1{\cosh a},\qquad
 e^{-a}\leq\delta(\xi,\eta)\leq2e^{-a}.
\]
\end{lemma}
\begin{proof}
In the totally geodesic two-dimensional disk through the two endpoints
and $o$, rotate the endpoints to $e^{i\theta},e^{-i\theta}$ with
$0<\theta\leq\pi/2$. Their half chord length is $\sin\theta$.
If $\theta=\pi/2$, the geodesic is a diameter and $a=0$. Otherwise the
orthogonal circle has center $\sec\theta$ and radius $\tan\theta$.
Its point nearest the origin has Euclidean radius
$t=\cos\theta/(1+\sin\theta)$. Radial hyperbolic distance in the
curvature $-1$ ball is $a=\log((1+t)/(1-t))$, obtained by integrating
$2\,dt/(1-t^2)$. Hence
$\cosh a=(1+t^2)/(1-t^2)=1/\sin\theta$. Finally
$e^a/2\leq\cosh a\leq e^a$ for $a\geq0$.
\end{proof}

\begin{lemma}[Uniform boundary estimates in both directions; MPR21]
\label{lem:mostow-boundary-holder}
Let $F:X\to X$ be $K$-bilipschitz with supplied boundary homeomorphism
$h$, and let $d(o,F(o))\leq B$ with $B\geq0$. Assume that $F$ and
$F^{-1}$ send every complete geodesic into the closed $M$-tube about
the geodesic joining its corresponding boundary endpoints. In the
metric of MPR20, for all $\xi,\eta$,
\begin{align*}
 \delta(h\xi,h\eta)&\leq2e^{B+M/K}\delta(\xi,\eta)^{1/K},\\
 \delta(h^{-1}\xi,h^{-1}\eta)&\leq2e^{KB+M/K}\delta(\xi,\eta)^{1/K}.
\end{align*}
\end{lemma}
\begin{proof}
For distinct endpoints write $\ell=\ell_{\xi\eta}$,
$\ell'=\ell_{h\xi,h\eta}$, $a=d(o,\ell)$ and $a'=d(o,\ell')$.
For every $y\in\ell'$, the inverse Morse assumption gives
$d(F^{-1}y,\ell)\leq M$. The triangle inequality and bilipschitz bound
therefore give
\[
 d(o,y)\geq d(F(o),y)-B
 \geq K^{-1}d(o,F^{-1}y)-B
 \geq (a-M)/K-B.
\]
Taking the infimum over $y$ gives $a'\geq a/K-M/K-B$. By MPR20,
\[
 \delta(h\xi,h\eta)\leq2e^{-a'}
 \leq2e^{B+M/K}e^{-a/K}
 \leq2e^{B+M/K}\delta(\xi,\eta)^{1/K}.
\]
Apply the same argument to $F^{-1}$, whose basepoint displacement is
at most $KB$, to prove the second estimate. Equal endpoints give zero
on both sides. The common constant $2e^{KB+M/K}$ controls both maps.
\end{proof}

The inverse Morse hypothesis is used on the \emph{target} geodesic.
The one-sided inclusion $F(\ell)\subset N_M(\ell')$ alone would not
justify the lower bound for $d(o,\ell')$. Here the inverse is again
$K$-bilipschitz and its supplied boundary map is $h^{-1}$, so a Morse
theorem uniform over that class supplies both inclusions. These are
spherical estimates; they assert no global Euclidean Holder bound.

\begin{lemma}[Tame boundary compactness with inverses; MPR22]
\label{lem:mostow-boundary-compactness}
Let $F_n$ be $K$-bilipschitz homeomorphisms with supplied boundary
homeomorphisms $h_n$, with $K$ independent of $n$, uniform Morse
inputs for the maps and inverses, and $\sup_n d(o,F_n(o))<\infty$.
There is a subsequence on which $h_n$ and $h_n^{-1}$ converge uniformly
on $S^2$ to mutually inverse homeomorphisms $h$ and $g$.
In particular this conclusion holds under the three-point hypotheses
and geometric inputs of MPR19. If all $h_n$ fix $\infty$, then so do
both limits, and both convergences are locally uniform in the
Euclidean plane.
\end{lemma}
\begin{proof}
MPR21 gives a single modulus of continuity for both families.
The sphere is compact. Compact-domain Arzela--Ascoli first extracts a
uniformly convergent subsequence of $h_n$, and then a further one for
$h_n^{-1}$. Their limits $h,g$ are continuous. For each $x$,
\[
 \delta(h_n(h_n^{-1}x),h(gx))
 \leq\sup_y\delta(h_n y,h y)
       +\delta(h(h_n^{-1}x),h(gx))\longrightarrow0,
\]
uniformly in $x$, by uniform continuity of $h$. Thus $h\circ g$ is
the identity. Interchanging the two families proves $g\circ h$ is
the identity. This supplies injectivity and surjectivity of the limits,
not just continuity. MPR19 supplies the required basepoint bound when
the distinct three-point limit is given.

Fixed points at infinity persist under uniform convergence. For a
compact subset of the plane its image under either limit stays a
positive spherical distance from infinity. Eventually the corresponding
sequence of images stays in one bounded chart. Uniform equivalence of
the spherical and Euclidean metrics on that compact chart gives local
uniform Euclidean convergence, for each of the two families.
\end{proof}

This compactness theorem concerns boundary maps alone, which is exactly
what the disk argument consumes. It does not identify the limit as the
boundary map of an interior limit or place it in the original isometry
orbit. Neither conclusion is used in MPR24. The retained public
Arzela--Ascoli source is a candidate compact-domain interface, not an
accepted PC or Lean binding; no new foundational reproof is planned.

\begin{lemma}[Stable inner and outer disks; MPR23]
\label{lem:mostow-disk-stability}
Suppose plane homeomorphisms $g_n$ and their inverses converge locally
uniformly to mutually inverse plane homeomorphisms $g,g^{-1}$, all
fixing $0$. There are fixed $0<\rho<R$ such that, eventually,
\[
 \overline B(0,\rho)\subset g_n(B(0,1)),\qquad
 g_n(\overline B(0,1))\subset\overline B(0,R).
\]
\end{lemma}
\begin{proof}
By continuity of $g^{-1}$ at zero, choose $\rho>0$ so that
$g^{-1}(\overline B(0,\rho))\subset B(0,1/2)$.
Uniform convergence of $g_n^{-1}$ on this closed disk implies that
its images lie in $B(0,3/4)$ eventually. Applying $g_n$ proves the
inner inclusion. The image $g(\overline B(0,1))$ is compact, hence
bounded. Choose $R>\rho$ greater than its maximal norm plus $1$.
Uniform convergence of $g_n$ on the closed unit disk proves the outer
inclusion. Both hold after taking the larger of their thresholds.
\end{proof}

\begin{lemma}[Finite disk packing controls oscillation; MPR25]
\label{lem:mostow-disk-packing}
Let $h:\mathbb R^2\to\mathbb R^2$ be a homeomorphism satisfying
MPR24's open/closed disk conclusion with one finite constant $C>1$. Let
$I_1,\ldots,I_k$, $k\geq1$, be positive-length intervals on one
horizontal line with pairwise disjoint interiors, and let $\Delta_j$
be the open disk with diameter $I_j$. Suppose $h(\Delta_j)\subset S$.
For a linear map $P:\mathbb R^2\to\mathbb R$ of norm at most $1$,
write $a_j,b_j$ for the interval endpoints and
$q_j=|P(h(b_j))-P(h(a_j))|$. Define $\alpha(S)$ as the supremum
of sums of Euclidean areas of finite families of pairwise disjoint
closed positive-radius disks contained in $S$, allowing the empty
family and the value $+\infty$. Then
\[
 \alpha(S)\geq\frac{\pi}{4C^2}\sum_{j=1}^k q_j^2
 \geq\frac{\pi}{4C^2k}\left(\sum_{j=1}^k q_j\right)^2.
\]
\end{lemma}
\begin{proof}
The open disks $\Delta_j$ are pairwise disjoint: for collinear centers,
disjoint interval interiors give center distance at least the sum of
radii. The assumed disk property gives closed disks $D_{1,j}\subset h(\Delta_j)$ and
$h(\overline\Delta_j)\subset D_{2,j}$ with diameter ratio at most $C$.
Their inner disks are pairwise disjoint by injectivity of $h$, even if
interval endpoints coincide. Since the endpoints belong to $\overline\Delta_j$,
\[
 q_j\leq|h(b_j)-h(a_j)|\leq\operatorname{diam}D_{2,j}
 \leq C\operatorname{diam}D_{1,j}.
\]
The area formula gives
$\operatorname{area}(D_{1,j})\geq\pi q_j^2/(4C^2)$.
This finite family is admissible in $\alpha(S)$, giving the first
inequality. The second is finite Cauchy--Schwarz. If $\alpha(S)=+\infty$
the inequalities retain their extended-real meaning.
\end{proof}

This makes the packing step behind Schwartz's Theorem 3.7,
printed/PDF page 18, explicit. In a bounded image of a compact strip,
monotonicity and finite additivity of Euclidean area bound $\alpha$ by
the area of an enclosing disk, so it is finite there. For finitely many
pairwise disjoint subsets $S_j$ of that strip, union of finite disk
packings and approximation of each supremum give
$\sum_j\alpha(S_j)\leq\alpha(\bigcup_j S_j)$.
No measurability of an arbitrary $S$ is required for these packing
statements. The open/closed distinction is deliberate: inner closed
disks lie inside images of open diameter disks, while endpoints are
controlled through the outer bound for the closed disk.

The next subsection supplies the absolute-continuity and positive-measure
regularity arguments using this packing input. Its inherited measure
and one-dimensional analysis hypotheses remain explicit.

\subsection{Absolute continuity and positive-measure regularity}
\label{sec:mostow-acl-regularity}

We now prove the analytic implication needed by MPR05, starting from
the disk conclusion already isolated in MPR24. The inherited inputs
are ordinary Euclidean Lebesgue outer measure and area, countable
subadditivity, null-set completeness, Fubini--Tonelli for Borel sets,
and the following one-dimensional facts: an absolutely continuous
real function on a compact interval is differentiable almost
everywhere, and an absolutely continuous function with derivative zero
almost everywhere is constant. These standard analysis results are
reuse/binding obligations, not new project axioms or planned reproofs.
Absolute continuity below uses finite intervals with disjoint interiors.
All derivatives are two-sided at interior points; endpoints are null.

\begin{proposition}[Disk control implies positive-measure asterisks; MPR32]
\label{prop:mostow-disk-regularity}
Let $h:\mathbb R^2\to\mathbb R^2$ be a homeomorphism satisfying the
open/closed disk conclusion of MPR24 for a fixed finite $C>1$.
Then the asterisk set of Definition~\ref{def:mostow-asterisk} is Borel
and has positive planar measure, hence positive outer measure.
No equivariance, lattice or compact quotient is assumed.
\end{proposition}
\begin{proof}
MPR26--MPR29 below give absolute continuity of both coordinate
functions of $h$ on almost every horizontal unit segment. MPR30
supplies Borel derivative-existence and nonzero-derivative sets, and
MPR31 gives a null exceptional set for the horizontal vector derivative
on the plane, together with a positive-measure set $T$ in $(0,1)^2$
where that derivative is nonzero.

For each $v\in\mathbb Q^2\setminus\{0\}$ choose an orthogonal map
$R_v$ with $R_v e_1=v/|v|$. The map $h\circ R_v$ has the same disk
constant: $R_v$ sends open and closed disks to disks with the same
radii. Apply MPR31 on every integer translate of the unit square.
Their missing grid lines are null. Transporting by $R_v$, which
preserves null sets, shows that $D_{v/|v|}h$ exists away from a null
set. The substitution $s=|v|t$ gives existence of
$D_vh=|v|D_{v/|v|}h$ there. The derivative for $v=0$ is zero everywhere.

MPR30 shows that each full existence set $D_v$ is Borel, so its null
complement $N_v$ is Borel. The union $N=\bigcup_{v\in\mathbb Q^2}N_v$
is null. Consequently $T\setminus N$ has the same positive measure
as $T$ and consists of asterisks. In fact the entire asterisk set is
$T_{e_1}\cap\bigcap_{v\in\mathbb Q^2}D_v$, with the notation of
MPR30, so is Borel. The countable cusp-center set can subsequently be
removed without losing positive measure.
\end{proof}

An asterisk here means existence of every rational directional
derivative and a nonzero horizontal derivative, exactly as in the
existing definition. This proof does not assert Frechet
differentiability or nonsingularity of a full derivative. The
homeomorphic zoom limit in MPR13 supplies its own nonzero slopes.

\begin{corollary}[The bilipschitz analytic hypothesis is supplied; MPR33]
\label{cor:mostow-regularity-to-boundary}
Let $(F,h)$ be the normalized equivariant bilipschitz pair of MPR18.
Assume its geometric hypotheses (1), (2) and (4), and the uniform
class-wide Morse, compatible boundary and model inputs of MPR24.
Using the inherited analysis inputs above, $h$ is the boundary map
of an isometry of $X$.
\end{corollary}
\begin{proof}
For every $a,b\in\operatorname{Isom}(X)$ for which
$\partial a\,h\,\partial b$ fixes infinity, $aFb$ is still
$K$-bilipschitz with that boundary map. MPR24 and MPR32 apply to each
such pair and give its positive-measure asterisk set. This supplies
MPR18's universally quantified hypothesis (3), not just regularity of
the initial $h$. Apply MPR18.
\end{proof}

Thus MPR05 has a written analytic proof from disk control and inherited
analysis. MPR34--MPR40 below supply its bilipschitz geometric producers;
inherited bindings and Lean implementation remain open.
MPR33 does not supply the finite-volume no-common-boundary-fixed-point
input, an arbitrary-homotopy controlled representative, or uniqueness.
It does not close the general quasi-isometry branch or HG04. The broader
producer MPR72 now supplies this regularity also for controlled
quasi-isometries; MPR73 supplies cusp avoidance for each normalization.

\begin{lemma}[Interval selection without an attained maximum; MPR26]
\label{lem:mostow-fivefold-cover}
For any family $\mathcal B$ of positive-length closed intervals in
$[0,1]$, there is a finite or countable pairwise disjoint subfamily
$I_j$ such that $\bigcup\mathcal B\subset\bigcup_j5I_j$.
Here $5I$ is dilation about the interval's midpoint by factor five.
\end{lemma}
\begin{proof}
While intervals remain, choose one with length greater than half the
supremum of the remaining lengths, and remove all intervals intersecting
it. That supremum is positive and at most one, so such a choice exists;
no largest interval is required. If the process terminates the selected
family is finite. Otherwise countable dependent choice gives a sequence
of disjoint closed intervals. Every finite sum of their lengths is at
most one, hence their lengths tend to zero.

Each original interval $J$ is eventually selected or removed. Indeed,
if it remained forever, every selected length would be greater than
$|J|/2$, contradicting convergence to zero. When $J$ is removed it
intersects the selected $I_j$ and satisfies $|J|<2|I_j|$. The distance
of their centers is at most half the sum of their lengths. For any
$x\in J$, its distance from the center of $I_j$ is therefore at most
$|J|+|I_j|/2<5|I_j|/2$. Thus $J\subset5I_j$, proving the cover.
\end{proof}

This bounded-interval variant of Schwartz's Lemma 4.6, printed/PDF 22,
uses a factor five and approximate suprema. The source's instruction
to choose a largest available interval is not needed. A supplied
inherited covering theorem with this conclusion can replace MPR26.

\begin{lemma}[Outer-measure bound and stiff points; MPR27]
\label{lem:mostow-stiff-points}
Let $A$ assign finite nonnegative numbers to the closed intervals in
$[0,1]$. Suppose some finite $T\geq0$ bounds $\sum_j A(J_j)$ for
every finite family of pairwise disjoint positive-length intervals.
For $N>0$, let $S_N$ consist of points $y\in(0,1)$ admitting a
positive-length interval $J\subset[0,1]$ centered at $y$ with
$A(J)>N|J|$. Then
\[
 \mu_1^*(S_N)\leq\frac{5T}{N}.
\]
Outside a set of outer measure zero, there is an integer $N_y\geq1$
such that $A(J)\leq N_y|J|$ for every such centered interval $J$.
\end{lemma}
\begin{proof}
Apply MPR26 to all intervals witnessing membership in $S_N$. For
every finite portion of the selected family,
$N\sum_j|I_j|\leq\sum_j A(I_j)\leq T$. Passing to the nonnegative
series gives $\sum_j|I_j|\leq T/N$. The dilated intervals cover $S_N$,
so outer subadditivity gives the stated bound. This does not require
$S_N$ to be measurable. The non-stiff set is
$\bigcap_{m\geq1}S_m$, whose outer measure is bounded by $5T/m$ for
every $m$, hence is zero. Its complement has the asserted integer bound.
\end{proof}

\begin{lemma}[Comparable finite refinement; MPR28]
\label{lem:mostow-comparable-partition}
Given a nonempty finite family of positive-length intervals with
disjoint interiors, and any real-valued function $f$ on their union, there is a
finite refinement into $k$ positive-length intervals with disjoint
interiors such that, with $L$ the original total length and $\epsilon$
the largest refined length,
\[
 k\epsilon\leq2L,
 \qquad
 \sum_{\rm original}|f(b)-f(a)|
 \leq\sum_{\rm refined}|f(b)-f(a)|.
\]
The refined total length is $L$ and $\epsilon\leq L$.
\end{lemma}
\begin{proof}
Choose $0<\eta\leq\min_i|I_i|$. Subdivide $I_i$ into
$n_i=\lceil|I_i|/\eta\rceil$ equal pieces. Since
$|I_i|/\eta\leq n_i\leq2|I_i|/\eta$, each piece has length in
$[\eta/2,\eta]$. In particular each is at least $\epsilon/2$.
Summing gives $k\epsilon\leq2L$. Length is unchanged, and the
triangle inequality applied to each telescoping endpoint difference
gives the second assertion. Zero-length intervals can be discarded
beforehand because their endpoint contribution is zero.
\end{proof}

\begin{lemma}[Quantitative absolute continuity at stiff lines; MPR29]
\label{lem:mostow-acl-stiff}
Let $h:\mathbb R^2\to\mathbb R^2$ be a homeomorphism satisfying
MPR24's open/closed disk conclusion with one finite constant $C>1$, and put
$A(J)=\alpha(h([0,1]\times J))$ using MPR25's packing functional.
There is a null set of heights outside which each $y\in(0,1)$ has
an integer $N_y\geq1$ with the following property. For every linear
$P:\mathbb R^2\to\mathbb R$ of norm at most one and every finite
partial partition of $[0,1]$ of total length
$0<L<\min(y,1-y)$,
\[
 \sum_i|P(h(b_i,y))-P(h(a_i,y))|
 \leq4C\sqrt{N_y L/\pi}.
\]
Consequently $x\mapsto P(h(x,y))$ is absolutely continuous on $[0,1]$.
The same exceptional height set and $N_y$ work for all these $P$.
\end{lemma}
\begin{proof}
The image of $[0,1]^2$ is compact and lies in a finite-area disk.
Thus $A$ is finite and nonnegative. Disjoint closed intervals give
disjoint closed strips, and their images are disjoint by injectivity.
Finite union of inner-disk packings and supremum approximation give
$\sum_j A(J_j)\leq\alpha(h([0,1]^2))=:T<\infty$.
MPR27 supplies the null exceptional set and the bound $N_y$.

Refine a nonempty partial partition by MPR28. Its largest length
$\epsilon\leq L$ satisfies $[y-\epsilon,y+\epsilon]\subset[0,1]$.
All open diameter disks of the refined horizontal intervals lie in
$[0,1]\times[y-\epsilon,y+\epsilon]$. Stiffness bounds the packing
functional of its image by $2N_y\epsilon$. If $q_j$ are the refined
projection oscillations, MPR25 and $k\epsilon\leq2L$ give
\[
 \left(\sum_{j=1}^k q_j\right)^2
 \leq\frac{4C^2k}{\pi}\,2N_y\epsilon
 \leq\frac{16C^2N_y}{\pi}L.
\]
The original sum is no larger by MPR28. For a requested variation
$\tau>0$, take
\[
 \delta=\min\left\{y,1-y,\frac{\pi\tau^2}{16C^2N_y}\right\}>0.
\]
Total length less than $\delta$ gives total oscillation less than
$\tau$. Empty partitions and zero-length terms contribute zero.
This is precisely the finite-partition definition of absolute continuity.
\end{proof}

The estimate develops the subdivision and packing argument in
Schwartz's Theorem 3.7, printed/PDF 18. The stiffness constant depends
on the line; no uniform bound in $y$ is asserted. Disjoint interiors
are sufficient for the partial partitions and open diameter disks;
the covering step instead selected actually disjoint closed intervals.

\begin{lemma}[Borel directional derivative sets; MPR30]
\label{lem:mostow-borel-derivatives}
For continuous $h:\mathbb R^2\to\mathbb R^2$ and fixed
$v\in\mathbb R^2$, let $D_v$ be the set where the finite vector
directional derivative $D_vh$ exists, and let $T_v\subset D_v$ be the
set where it is nonzero. Both sets are Borel.
\end{lemma}
\begin{proof}
For $a\ne0$ put $Q_a(z)=(h(z+av)-h(z))/a$. Each $Q_a$ is continuous
in $z$. Completeness and the Cauchy criterion give
\[
 D_v=\bigcap_{k\geq1}\ \bigcup_{m\geq1}\quad
 \bigcap_{\substack{a,b\in\mathbb Q\\0<|a|,|b|<1/m}}
 \{z:|Q_a(z)-Q_b(z)|\leq1/k\}.
\]
To justify the rational criterion, it first makes $Q_{1/n}(z)$ a
Cauchy sequence with a finite limit. It controls all sufficiently
small nonzero rational parameters by the same limit. For fixed real
$t\ne0$, continuity of $a\mapsto Q_a(z)$ near $t$ allows approximation
by rationals still satisfying the chosen strict smallness bound.
This proves the full two-sided real limit. The converse follows
immediately from the ordinary Cauchy criterion. The displayed sets
are closed and all intersections and unions are countable, so $D_v$
is Borel. Moreover
\[
 T_v=D_v\cap\bigcup_{k\geq1}\ \bigcup_{m\geq1}\quad
              \bigcap_{n\geq m}\{z:|Q_{1/n}(z)|\geq1/k\}.
\]
On $D_v$ the sequence converges, and this formula says exactly that
the norm of its limit is positive. It is again a Borel formula.
\end{proof}

\begin{lemma}[Horizontal derivative and a positive nonzero set; MPR31]
\label{lem:mostow-horizontal-regularity}
Let $h$ be a plane homeomorphism with disk constant $C$. Using the
inherited analysis inputs stated above, $D_{e_1}h$ exists almost
everywhere on $\mathbb R^2$, and the Borel set
$T_{e_1}\cap(0,1)^2$ has positive measure.
\end{lemma}
\begin{proof}
Write $h=(h_1,h_2)$. For every stiff height from MPR29, both scalar
functions $x\mapsto h_i(x,y)$ are absolutely continuous. Their
derivatives exist almost everywhere on $(0,1)$, hence their vector
derivative exists there outside a null set. By MPR30 its failure set
in the open square is Borel. Fubini--Tonelli makes that failure set
planar null, since its sections are null for almost every height.
Translate this argument to all integer unit squares; the countable
union of their boundary lines is null, proving the whole-plane claim.

For each stiff height, injectivity gives $h(0,y)\ne h(1,y)$, so at
least one coordinate function has unequal endpoints. An absolutely
continuous scalar function with unequal endpoints has a positive-measure
set where its derivative exists and is nonzero: otherwise its derivative
is zero almost everywhere, and the inherited constancy theorem is a
contradiction. The one-dimensional scalar version of MPR30's rational
quotient criterion makes this derivative set Borel. Remove the null set
where the other coordinate derivative
fails. The remaining positive-measure subset lies in the horizontal
section of $T_{e_1}\cap(0,1)^2$.

This choice of coordinate may depend on the height; no measurable
choice of projections is used. The set $T_{e_1}$ itself is Borel by
MPR30. If its intersection with the open square were null, Fubini--Tonelli
would make its sections null for almost every height, contradicting
the preceding conclusion for every stiff height. Thus its measure is
positive.
\end{proof}

This closes the project-specific analytic reduction from disk control.
Schwartz's Lemma 3.1 and Theorems 3.3--3.5 supply the source template;
Sections 4.3 and 5.1--5.4 give the covering, Borel and standard-analysis
context. The retained public mathlib statements for almost-everywhere
differentiability of absolutely continuous real functions and constancy
from an almost-everywhere zero derivative are plausible inherited
interfaces, not accepted imports. Exact finite-partition/interval,
endpoint, real/vector and product-measure adapters must be bound against
the completed foundation before Lean implementation. No changing PC
snapshot has been inspected and no standard real-analysis reproof is
being added to the implementation plan.

The next subsection supplies the bilipschitz Morse, compatible boundary
extension and triple-tube inputs used by the disk argument. The analytic
lemmas above should not be replayed as unproved mathematical tasks merely
because their accepted library bindings and Lean proofs are pending.

\subsection{Uniform Morse tracking and compatible boundary extension}
\label{sec:mostow-morse-extension}

We next supply the bilipschitz geometric producers used above. Work in
curvature $-1$ hyperbolic three-space with its standard ball and
half-space models. Inherited inputs are the model distance and length
formulas, compactness of closed metric balls and of unit tangent vectors
over compact sets, continuous geodesic flow and endpoints, and ordinary
calculus and length estimates for Lipschitz curves. These are foundation
bindings, not additional rigidity hypotheses or proposed new axioms.
In particular a bilipschitz map is not assumed differentiable: the smooth
projection is composed with its Lipschitz image curves using the inherited
almost-everywhere chain rule and length formula.

\begin{proposition}[The bilipschitz geometric and analytic inputs are supplied; MPR40]
\label{prop:mostow-bl-geometric-inputs}
Let $F:X\to X$ be $K$-bilipschitz and equivariant for an isomorphism
$\phi:\Gamma\to\Gamma'$. Assume the geometric core/cusp, compact-lift,
isometry-properness and boundary-action inputs in MPR18(1), and that
$\Gamma$ has no common fixed point on $S^2$, as in MPR18(4).
With the inherited model and analysis inputs stated here and in
Section~\ref{sec:mostow-acl-regularity}, $F$ has a boundary map which
is the boundary map of a hyperbolic isometry.
\end{proposition}
\begin{proof}
MPR34--MPR38 below produce a compatible boundary homeomorphism $h$,
uniform Morse bounds for every map and its inverse, and compatibility
with isometry composition and equivariance. MPR39 supplies the compact
ideal-triangle tube input. Normalize $h$ by a target isometry so that
it fixes infinity. This preserves $K$ and the domain group hypotheses;
the target group and marking are conjugated by that actual isometry.
MPR24 then gives uniform disk control and MPR32 gives regularity
throughout the normalized isometry orbit. Equivalently apply MPR33.
Undoing the target normalization proves the claim.
\end{proof}

This removes the separate Morse, boundary-extension and triple-tube
producer assumptions in the bilipschitz branch, conditional on the
listed inherited model/calculus bindings. It does not prove the
finite-volume group hypothesis, produce a controlled representative
of an arbitrary homotopy equivalence, prove uniqueness, or settle the
general quasi-isometry branch. HG04 and HG06/HG08 retain their full scope.

\begin{lemma}[Exact contraction of projection to a line; MPR34]
\label{lem:mostow-exact-projection}
For a complete geodesic $L\subset X$, let $\pi_L$ be nearest-point
projection. If a Lipschitz curve $\beta$ stays at distance at least
$r\geq0$ from $L$, then
\[
 \operatorname{length}(\pi_L\circ\beta)
 \leq\frac{\operatorname{length}(\beta)}{\cosh r}.
\]
The curve may meet the boundary of the closed $r$-tube.
\end{lemma}
\begin{proof}
By an isometry take $L=\{0\}\times(0,\infty)$ in the half-space.
For $p=(z,t)$ put $\rho=(|z|^2+t^2)^{1/2}$. The model formula gives
\[
 \cosh d((z,t),(0,s))=\frac{|z|^2+t^2+s^2}{2ts}.
\]
Minimizing $s+\rho^2/s$ for $s>0$ gives the unique minimizer $s=\rho$.
Thus $\pi_L(p)=(0,\rho)$ and
$\cosh d(p,L)=\rho/t$. The arclength coordinate on $L$ is $\log\rho$.
For a Euclidean tangent vector $v$,
\[
 |d(\log\rho)_p(v)|=\frac{|\langle p,v\rangle|}{\rho^2}
 \leq\frac{|v|}{\rho}
 =\frac{t}{\rho}\,\|v\|_{\rm hyp}
 =\frac{\|v\|_{\rm hyp}}{\cosh d(p,L)}.
\]
Integrate along the Lipschitz curve using the inherited chain rule and
length formula. A compact parameter interval has image in a compact
subset of the half-space, so the Euclidean/hyperbolic local calculus
applies. This proves the bound even at equality $d(p,L)=r$.
\end{proof}

\begin{lemma}[Finite-arc tracking with an explicit constant; MPR35]
\label{lem:mostow-finite-arc-tracking}
Put $C_K=4K^3+2K$ for $K\geq1$. If $c:[a,b]\to X$ is a unit-speed
geodesic segment with $a<b$, and $F$ is $K$-bilipschitz, then
$F(c([a,b]))$ lies in the closed $C_K$-tube about the complete geodesic
through its two image endpoints. This conclusion concerns the complete
supporting line, not yet the segment between those endpoints.
\end{lemma}
\begin{proof}
Write $\beta=F\circ c$, and let $L$ be the supporting line. Every
subarc with endpoints $p,q$ has length $s\leq K^2d(p,q)$: its length
is at most $K$ times the original parameter length, which is at most
$K d(p,q)$. If some point of $\beta$ has distance greater than $C_K$
from $L$, continuity of distance gives an excursion subarc whose
endpoints have distance $2K$ and whose interior has distance greater
than $2K$. Both endpoints of the whole curve lie on $L$, so the
excursion has endpoints on both sides of the chosen parameter.
The distance-to-$L$ function is 1-Lipschitz, and hence the excursion
length satisfies $s>2(C_K-2K)=8K^3$.

MPR34 and the two end connectors give
\[
 d(p,q)\leq4K+\frac{s}{\cosh(2K)},\qquad
 s\leq K^2d(p,q).
\]
Since $\cosh(2K)\geq1+2K^2>2K^2$, these imply
$s<4K^3+s/2$, and therefore $s<8K^3$, a contradiction.
The same constant works for every map and every segment in the class.
\end{proof}

This develops the excursion argument of Schwartz's Lemma 2.4 (v5,
printed/PDF 13) using the exact projection factor. His weaker factor
$e^{-2K+1}$ also gives a valid contradiction: the function
$2K^2e^{-2K+1}$ is nonincreasing for $K\geq1$, since its logarithmic
derivative is $2/K-2\leq0$, and its value at one is $2/e<1$.
The exact factor above permits the direct bound
$\cosh(2K)>2K^2$ and retains the same finite-arc constant.
This is an alternative derivation, not a numerical correction to the source.

\begin{lemma}[Bounded distance and ideal endpoints; MPR37]
\label{lem:mostow-bounded-distance-endpoints}
If $x_n$ converges to $\xi\in S^2$ in the closed ball compactification
and $d(x_n,y_n)$ is bounded, then $y_n\to\xi$.
Unit-speed rays with the same ideal endpoint remain at bounded distance
when their nonnegative parameters are synchronized.
\end{lemma}
\begin{proof}
In the unit ball the distance formula is
\[
 \cosh d(x,y)=1+\frac{2|x-y|^2}{(1-|x|^2)(1-|y|^2)}.
\]
A distance bound $D$ gives
$|x_n-y_n|^2\leq(\cosh D-1)(1-|x_n|^2)/2\to0$.
For the ray assertion move their common endpoint to infinity. They are
$(z,a e^s)$ and $(w,b e^s)$ for fixed $a,b>0$, and
\[
 \cosh d((z,a e^s),(w,b e^s))
 =1+\frac{|z-w|^2e^{-2s}+(a-b)^2}{2ab}
 \leq1+\frac{|z-w|^2+(a-b)^2}{2ab}
\]
for $s\geq0$. This is a finite bound; no uniformity over the initial
rays is needed.
\end{proof}

\begin{lemma}[Complete-line tracking and both ends; MPR36]
\label{lem:mostow-complete-line-tracking}
For every $K$-bilipschitz $F$ and complete unit-speed geodesic
$c:\mathbb R\to X$, there is a unique complete geodesic $L$ at
finite Hausdorff distance from $F(c(\mathbb R))$. In fact
\[
 d_{\rm Haus}(F(c(\mathbb R)),L)\leq C_K.
\]
The two ends of $F\circ c$ converge to the two different ideal endpoints
of $L$. The same constant applies to $F^{-1}$.
\end{lemma}
\begin{proof}
Let $L_n$ support the endpoints of $F(c([-n,n]))$. MPR35 places
$F(c(0))$ within $C_K$ of $L_n$. Let $q_n$ be its nearest point on
$L_n$ and choose either unit tangent to $L_n$ at $q_n$. The points
$q_n$ lie in one compact metric ball; the unit tangent bundle over
that ball is compact. Pass to a subsequence of these pointed lines
converging to a pointed complete geodesic $L$.

Fix $t$. For all sufficiently large $n$, $F(c(t))$ is within $C_K$
of $L_n$. Its nearest points on $L_n$ stay in a fixed compact ball.
Their signed arclength coordinates relative to $q_n$ are bounded,
by the triangle inequality through $F(c(t))$ and $F(c(0))$.
Continuous dependence of the geodesics then gives
$d(F(c(t)),L)\leq C_K$. This holds for every $t$ on the same selected
subsequence; no diagonal choice over all real parameters is needed.

Give $L$ a signed unit-speed coordinate and put
$u(t)=\operatorname{coord}(\pi_L F(c(t)))$, with $u(0)=0$.
It is continuous, and the two projection connectors give
\[
 |u(t)-u(s)|\geq K^{-1}|t-s|-2C_K.
\]
Thus $|u(t)|\to\infty$ as $t\to\pm\infty$. Continuity makes the sign
constant on each sufficiently long tail. The two signs are different.
Otherwise, reversing the coordinate if necessary, both tails would tend
to $+\infty$. For every large $n$, the intermediate value theorem on
each tail would supply $s\leq-n$ and $t\geq n$ with $u(s)=u(t)$,
contradicting $|t-s|\leq2KC_K$.
Consequently $u$ is onto, so every point of $L$ is within $C_K$ of
some $F(c(t))$. MPR37 gives the asserted two ideal limits.

Any other line at finite Hausdorff distance has the same ideal endpoints,
by MPR37 and the two ends just proved, and hence is $L$.
Applying this argument to $F^{-1}$ gives the same bound since its
bilipschitz constant is also $K$.
\end{proof}

Taking pointed-line limits avoids assuming that two escaping image
endpoints have different ideal limits merely because their complete
supporting lines meet a fixed compact set. The signed projection and
coarse lower bound above supply that missing distinction explicitly.

\begin{proposition}[Canonical compatible boundary extension; MPR38]
\label{prop:mostow-canonical-boundary}
Every $K$-bilipschitz homeomorphism $F:X\to X$ extends uniquely to a
homeomorphism of the closed ball compactification. Its boundary map
$h$ sends the ideal endpoints of $L$ to those of the MPR36 line for
$F(L)$, and the inverse boundary map is the one constructed from
$F^{-1}$. Boundary extension respects composition and isometries;
if $F\gamma=\phi(\gamma)F$, then
$h\,\partial\gamma=\partial\phi(\gamma)\,h$.
The uniform Morse bound $M(K)=C_K$ is valid for all such maps and inverses.
\end{proposition}
\begin{proof}
For $\xi\in S^2$, take any unit-speed ray ending at $\xi$, complete
it to a line, and define $h(\xi)$ by the positive ideal limit in MPR36.
Two such rays stay at bounded distance by MPR37; their images do also
by the Lipschitz bound. MPR37 shows that the definition is independent
of the ray. The two ends of each complete line therefore map to
its two distinct tracked endpoints.

Construct $g$ in the same way from $F^{-1}$. If $c(t)\to\xi$ and
$L'$ tracks $F(c(\mathbb R))$, write
$F(c(t))$ within $C_K$ of $\ell'(u(t))$, where $\ell'$ is oriented
toward $h(\xi)$ and $u(t)\to+\infty$. Applying $F^{-1}$ gives
$d(c(t),F^{-1}(\ell'(u(t))))\leq KC_K$.
The second point tends to $g(h(\xi))$ by its defining ray limit.
MPR37 gives $gh(\xi)=\xi$; interchange $F,F^{-1}$ to get $hg=\mathrm{id}$.
Thus the endpoint map is a bijection before continuity is invoked.

For clarity, continuity does not assume the homeomorphism premise of
MPR21. Repeat its distance estimate now that endpoint compatibility
and the inverse endpoint map have been constructed. If $o$ is the ball
origin, $B=d(o,F(o))$, $a=d(o,L_{\xi\eta})$ and
$a'=d(o,L_{h\xi,h\eta})$, inverse tracking gives
\[
 a'\geq(a-C_K)/K-B.
\]
MPR20 then yields
$\delta(h\xi,h\eta)\leq2e^{B+C_K/K}\delta(\xi,\eta)^{1/K}$.
The inverse estimate uses $d(o,F^{-1}(o))\leq KB$.
Both endpoint maps are continuous, so they are homeomorphisms.

It remains to check continuity from the interior to the boundary.
If $x_n\to\xi\in S^2$, write $x_n=c_{\xi_n}(t_n)$ using unit-speed
rays from $o$. Then $t_n\to\infty$ and $\xi_n\to\xi$.
On each tracked line choose the foot of $F(o)$ as coordinate zero
and orient toward $h(\xi_n)$. Its signed projection $u_n(t)$ satisfies
$u_n(0)=0$ and $|u_n(t)|\geq t/K-2C_K$.
For $t>2KC_K$ the sign cannot change, and its positive-end limit
makes it positive. Thus $u_n(t_n)\to+\infty$.
The feet lie in the closed $C_K$-ball about $F(o)$. Compactness of
unit tangent vectors there implies that every subsequence admits a
further subsequence of oriented pointed lines with convergent frames.
Their positive endpoints tend to $h(\xi)$ by the continuity just
proved. The model geodesic formula, with the coordinate tending to
$+\infty$, makes their projected points tend to that endpoint.
MPR37 transfers this limit to $F(x_n)$. The subsequence criterion
proves $F(x_n)\to h(\xi)$. Repeat for $F^{-1}$.

We have obtained a homeomorphism of the compactification. It is unique
because the interior is dense and the target is Hausdorff. Uniqueness
also proves compatibility with composition, the standard extensions of
isometries and the displayed equivariance relation. No quotient
compactness or lattice assumption has been used.
\end{proof}

\begin{lemma}[Compact intersections of ideal-triangle tubes; MPR39]
\label{lem:mostow-triple-tube-compact}
For three distinct ideal points $\xi_1,\xi_2,\xi_3$ and $R\geq0$,
the intersection of the closed $R$-tubes about their three complete
joining geodesics is compact in $X$.
\end{lemma}
\begin{proof}
The intersection is closed in $X$. If it were unbounded, properness
and the compact ball compactification would give a sequence in it
converging to an ideal point $\xi$. Its nearest point on each side
is at distance at most $R$, hence tends to $\xi$ by MPR37. The closure
of each complete geodesic meets $S^2$ only in its two endpoints.
Consequently $\xi$ would belong to
\[
 \{\xi_1,\xi_2\}\cap\{\xi_2,\xi_3\}\cap\{\xi_3,\xi_1\}
 =\varnothing,
\]
a contradiction. The intersection is therefore closed and bounded
in the proper metric space $X$, hence compact; it may be empty.
\end{proof}

Schwartz's Theorem 1.2 and Section 2 (printed/PDF 11--14) provide
the geometric route. MPR34--MPR39 expand the projection, constant,
line-limit, inverse, interior-continuity and compactness steps. The
Gouezel--Shchur Theorem 1.1 is a distinct corrected general Morse
reference, not a correction to Schwartz or a new dependency here.
These written model arguments still require exact accepted-foundation
bindings. They settle neither arbitrary quasi-isometry boundary
regularity nor the controlled representative required by full HG04.
The next subsection supplies the finite-volume group facts and marked
uniqueness. Arbitrary marked representatives and the actual HG06
quantitative/core-map adapter remain separate.

\subsection{Finite-volume group facts and marked uniqueness}
\label{sec:mostow-finite-volume-groups}

We now supply the finite-volume boundary-action and centralizer inputs,
then combine them with the written bilipschitz argument. Throughout this
subsection $X=\mathbb H^3$ has curvature $-1$, $\Gamma$ acts freely and
properly discontinuously by isometries, and $H=X/\Gamma$ has finite
Riemannian volume. The volume is positive since a nonempty coordinate
ball has positive volume. Neither orientability nor compactness of $H$
is assumed. Volume always means the Riemannian density, so change of
variables applies also to orientation-reversing maps.

Inherited inputs are the standard hyperbolic models and their isometry
groups, covering and homotopy lifting, Riemannian integration and change
of variables, and finite-dimensional linear algebra. The elementary
volume and marking arguments below are project proof expansions using
those inputs. BP A.2.4, A.3.3, A.4.1--A.4.3, A.5.13 and B.1.6--B.1.7
supply the model and quotient conventions \cite{BenedettiPetronio1992}.
Martelli's Exercise 4.3.5 states the centralizer result without a proof;
Proposition 4.3.6 and Corollary 4.3.7 give its uniqueness application
\cite{Martelli2025v4}. We prove the required dimension-three case explicitly.

\begin{proposition}[Unique marked isometry from a supplied bilipschitz comparison; MPR47]
\label{prop:mostow-bl-marked-unique}
Let $u:H\to H'$ be the actual homotopy equivalence in HG04, with lift
$U:X\to X$ and marking $\phi:\Gamma\to\Gamma'$. Suppose there is a
$K$-bilipschitz homeomorphism $F:X\to X$ with
$F\gamma=\phi(\gamma)F$ for every $\gamma\in\Gamma$.
Assume MPR18(1)'s core/cusp, compact-lift, isometry-properness and
boundary-action inputs, and the inherited model/calculus/analysis and
covering inputs stated in this development. Then $u$ is homotopic to a
unique isometry $H\to H'$.
\end{proposition}
\begin{proof}
MPR41 below supplies the no-common-boundary-fixed-point input from
finite volume. MPR40 therefore constructs an isometry $A$ whose boundary
map equals the canonical boundary map of $F$. That boundary map is
$\phi$-equivariant by MPR38. MPR10 descends $A$ and interpolates between
$U$ and $A$, giving an isometry homotopic to the actual $u$. MPR46 below
gives uniqueness among all isometries in that homotopy class.
\end{proof}
The maps $U$ and $F$ need not agree or be a bounded distance apart;
they must carry the same actual marking. This proposition does not
produce $F$ from an arbitrary homotopy equivalence or remove the stated
core/orbit and inherited-interface bindings. In particular it does not
yet prove HG04. It does close the finite-volume group and uniqueness
steps in this supplied-comparison branch.

\begin{lemma}[No common boundary fixed point; MPR41]
\label{lem:mostow-no-boundary-fixed-point}
The finite-volume deck group $\Gamma$ has no common fixed point on
$\partial X$.
\end{lemma}
\begin{proof}
Suppose it fixes $\xi$. Conjugate by an isometry taking $\xi$ to infinity
in the upper half-space $\{(z,t):z\in\mathbb R^2,\ t>0\}$. Every
isometry fixing infinity has the form
\[
 (z,t)\longmapsto(\lambda Qz+b,\lambda t),
 \qquad \lambda>0,\quad Q\in O(2),\quad b\in\mathbb R^2.
\]
This includes the orientation-reversing case. For each $s\in\mathbb R$
the smooth diffeomorphism $\Psi_s(z,t)=(z,e^s t)$ commutes with every
such isometry. It consequently descends to a diffeomorphism
$\psi_s:H\to H$ with inverse $\psi_{-s}$. In these coordinates the
hyperbolic density is $t^{-3}\,|dz_1\,dz_2\,dt|$, and direct substitution
gives
\[
 \Psi_s^*d\operatorname{vol}_X=e^{-2s}d\operatorname{vol}_X,
 \qquad
 \psi_s^*d\operatorname{vol}_H=e^{-2s}d\operatorname{vol}_H.
\]
The second equality is local on the covering charts. Global change of
variables for the diffeomorphism $\psi_s$ now yields
\[
 0<\operatorname{vol}(H)
 =\int_H\psi_s^*d\operatorname{vol}_H
 =e^{-2s}\operatorname{vol}(H)<\infty.
\]
Taking $s=1$ contradicts positivity and finiteness. The map $\Psi_s$
is not asserted to be a hyperbolic isometry; its volume scaling is the
point of the argument.
\end{proof}

\begin{lemma}[Volume under a finite-index cover; MPR42]
\label{lem:mostow-finite-cover-volume}
If $\Gamma_0<\Gamma$ has index $m<\infty$, then
$X/\Gamma_0\to X/\Gamma$ is an $m$-sheeted local-isometry covering and
\[
 \operatorname{vol}(X/\Gamma_0)=m\operatorname{vol}(X/\Gamma).
\]
In particular $\Gamma_0$ also has positive finite covolume. Normality
of the subgroup is not required.
\end{lemma}
\begin{proof}
Freeness and proper discontinuity restrict to $\Gamma_0$. Over an evenly
covered neighborhood for $X\to H$, the sheets of the intermediate cover
are indexed by the $m$ cosets of $\Gamma_0$ in $\Gamma$; each is mapped
isometrically onto that neighborhood. Choose a countable cover of $H$
by such neighborhoods and replace it by the disjoint Borel partition
obtained by subtracting earlier neighborhoods. The preimage of each
partition member consists of $m$ Borel sets with equal volume, by local
change of variables. Countable additivity proves the formula. The
countable cover uses the inherited second-countability of the manifold.
\end{proof}

\begin{lemma}[No nonempty finite boundary orbit; MPR43]
\label{lem:mostow-no-finite-boundary-orbit}
There is no nonempty finite $\Gamma$-invariant subset of $\partial X$.
\end{lemma}
\begin{proof}
If $S$ were such a set, choose $\xi\in S$. Its stabilizer $\Gamma_\xi$
has finite index $|\Gamma\xi|\leq |S|$ by the orbit--coset bijection.
MPR42 gives positive finite covolume for this subgroup, whereas it fixes
$\xi$, contradicting MPR41. In particular $\Gamma$ cannot preserve an
unordered pair of boundary points.
\end{proof}

\begin{lemma}[No invariant totally geodesic plane; MPR44]
\label{lem:mostow-no-invariant-plane}
The group $\Gamma$ does not preserve a totally geodesic plane in $X$.
\end{lemma}
\begin{proof}
Suppose it preserves $P$. Use the hyperboloid model with bilinear form
\[
 \langle x,y\rangle=x_1y_1+x_2y_2+x_3y_3-x_4y_4,
 \qquad X=\{x:\langle x,x\rangle=-1,\ x_4>0\}.
\]
After an isometry, $P=X\cap\{x_3=0\}$ and its unit spacelike normal
is $e_3$. The map
\[
 \Phi:P\times\mathbb R\longrightarrow X,
 \qquad \Phi(p,r)=\cosh r\,p+\sinh r\,e_3
\]
is a diffeomorphism: the inverse has $r=\operatorname{arsinh}(x_3)$
and $p=(x-x_3e_3)/\sqrt{1+x_3^2}$. Differentiating and using
$\langle p,p\rangle=-1$, $\langle p,v\rangle=0$ for $v\in T_pP$,
and orthogonality to $e_3$, gives
\[
 \Phi^*g_X=dr^2+\cosh^2r\,g_P,
 \qquad
 \Phi^*d\operatorname{vol}_X=\cosh^2r\,d\operatorname{area}_P\,dr.
\]
The nearest point on $P$ to $\Phi(p,r)$ is uniquely $p$: for $q\in P$
the hyperboloid distance formula gives
\[
 \cosh d(\Phi(p,r),q)=\cosh r\,\cosh d_P(p,q),
\]
whose minimum occurs only at $q=p$. This nearest-point projection is
$\Gamma$-equivariant, even when an element interchanges the two sides
of $P$.

Choose a small open disk $D\subset P$ of positive finite area inside a
neighborhood $V\subset X$ with $\gamma V\cap V=\varnothing$ for
$\gamma\ne1$. For every $T>0$, the quotient projection is injective
on $\Phi(D\times(0,T))$. Indeed, if two such points differ by $\gamma$,
their nearest points lie in $D$ and also differ by $\gamma$; hence
$\gamma V\cap V\ne\varnothing$, forcing $\gamma=1$. This open set
therefore embeds by a local isometry into $H$. Its volume is
\[
 \operatorname{area}(D)\int_0^T\cosh^2 r\,dr
 \ \geq\ T\operatorname{area}(D).
\]
Finite volume of $H$ cannot dominate this expression for every $T$.
No global product decomposition of $H$, or orientation of its normal
bundle, was used.
\end{proof}

\begin{proposition}[Trivial centralizer in the full isometry group; MPR45]
\label{prop:mostow-trivial-centralizer}
Every $A\in\operatorname{Isom}(X)$ commuting with $\Gamma$ is the
identity. Thus the mathematical assertion of MPR04 is supplied.
\end{proposition}
\begin{proof}
First suppose $A$ preserves orientation and is not the identity. Its
boundary map is a nonidentity complex M\"obius transformation. Such a
map has a nonempty fixed set of cardinality at most two on the Riemann
sphere: its fixed lines are the eigenlines of a nonscalar invertible
complex $2\times2$ matrix. Existence of a complex eigenvalue gives at
least one, and three fixed lines would force the matrix to be scalar.
The boundary action is faithful, so the nonscalar case applies. Since
$A$ commutes with $\Gamma$, that finite set is $\Gamma$-invariant,
contradicting MPR43.

If $A$ reverses orientation, $A^2$ preserves orientation and commutes
with $\Gamma$, so the first case gives $A^2=1$. For any $x\in X$,
the midpoint of the segment from $x$ to $Ax$ is fixed by $A$, by
uniqueness of the segment. Conjugate that midpoint to the origin of
the ball model. The stabilizer there is $O(3)$, so $A$ is represented
by $Q\in O(3)$ with $Q^2=I$ and $\det Q=-1$. Its orthogonal
$+1$ and $-1$ eigenspaces decompose $\mathbb R^3$, and the negative
eigenspace has dimension one or three. In dimension three $Q=-I$,
so $A$ fixes exactly one point of $X$. Commutation makes that point
fixed by every deck transformation. Freeness forces $\Gamma=\{1\}$,
which contradicts MPR41 since the trivial group fixes all boundary
points. In dimension one $A$ fixes precisely the totally geodesic
plane given by its positive eigenspace in the ball. Commutation makes
this plane $\Gamma$-invariant, contradicting MPR44. Both cases are
impossible, proving the assertion in the full isometry group.
\end{proof}
This proof uses no density theorem for loxodromic endpoints and no
rigidity theorem as an input. The use of complex M\"obius fixed points
is specific to dimension three; no all-dimensional contract is asserted.

\begin{proposition}[Homotopic quotient isometries coincide; MPR46]
\label{prop:mostow-homotopic-isometries}
Let $H$ and $H'$ be complete connected hyperbolic three-manifolds,
with $H$ of finite volume. If isometries $a,b:H\to H'$ are homotopic,
then $a=b$.
\end{proposition}
\begin{proof}
Fix the hyperbolic universal covers and an isometric lift $A:X\to X$
of $a$, with actual equivariance $A\gamma=\phi(\gamma)A$.
An isometry lifts to an isometry: it lifts as a local isometry, and a
lift of its inverse chosen at matching basepoints is its inverse by
uniqueness of lifts. Lift a specified homotopy from $a$ to $b$ starting
at $A$, obtaining $\widetilde J:X\times[0,1]\to X$. For each $\gamma$,
the maps $\widetilde J(\gamma x,s)$ and
$\phi(\gamma)\widetilde J(x,s)$ are lifts of the same map to $H'$ and
agree at $s=0$. Uniqueness of lifts on the connected space
$X\times[0,1]$ makes them equal everywhere. In particular the endpoint
$B=\widetilde J(\cdot,1)$ is an isometric lift of $b$ with the same
$\phi$. Therefore
\[
 (A^{-1}B)\gamma=A^{-1}\phi(\gamma)B=\gamma(A^{-1}B).
\]
MPR45 gives $A^{-1}B=1$, hence $a=b$. This compares the lifts selected
by the homotopy; arbitrary independent choices of lifts need not be
equal. No boundedness or properness of the homotopy is required.
\end{proof}

The exact source locators and source-versus-project distinction are
recorded in \code{reference_checks_revision110.md}. The archived
Martelli version is byte-identical to the version checked in revision49;
its newly discovered local copy is now recorded in the source manifest.
Accepted inherited geometry/integration/covering bindings and all Lean
implementations remain pending.

\subsection{Controlled representatives with the original marking}
\label{sec:mostow-controlled-representatives}

The input to MPR47 is bilipschitz. An arbitrary homotopy equivalence
is not thereby a homeomorphism, and the construction below does not
make it one. We supply MPR01's weaker, quantitative quasi-isometry
output for compact manifolds and, with explicit peripheral data, for
cusped manifolds. This separates the input-map problem from the later
extension of the boundary/regularity argument to quasi-isometries.

Let $H=X/\Gamma$ and $H'=X/\Gamma'$ be complete connected hyperbolic
three-manifolds of curvature $-1$, with actual homotopy equivalence
$u$, continuous lift $U$, and isomorphism $\phi$ as above. Inherited
inputs are Whitney approximation for maps to manifolds without boundary,
covering and homotopy lifting, smooth isometry-equivariant geodesic
interpolation in $X$, smooth real cutoff functions, and Riemannian
path-length and differential-norm calculus. They require accepted
foundation bindings; no new foundational Lean development is asserted.

\begin{definition}[Matched geometric cusp data]
\label{def:mostow-matched-cusp-data}
The following data are supplied for the actual $\phi$.
\begin{enumerate}[leftmargin=*]
\item Finite families of cusp horoball representatives $B_j,B'_j$ whose
translates form the full lifted cusp families. Each has a chosen
logarithmic-height chart $(z,t)\in\mathbb R^2\times[0,\infty)$,
with metric $dt^2+e^{-2t}q_j$ or $dt^2+e^{-2t}q'_j$ for constant
Euclidean metrics. The quotient complements of the open tails $t>0$, including their
boundary sections, are compact cores.
The charts extend to $t>-\epsilon$ for some common $\epsilon>0$;
these enlarged cusp families are locally finite and pairwise disjoint.
\item The representatives are paired so that
$\phi(P_j)=P'_j$, where $P_j=\operatorname{Stab}_\Gamma(B_j)$ and
$P'_j=\operatorname{Stab}_{\Gamma'}(B'_j)$. The stabilizers act on
$z$ by Euclidean isometries and fix $t$. The pairing lists every cusp
orbit on both sides exactly once.
\item Invertible affine maps $a_j:\mathbb R^2\to\mathbb R^2$ satisfy
$a_j\rho_j(p)=\rho'_j(\phi(p))a_j$ for every $p\in P_j$, where
$\rho_j,\rho'_j$ are the actual horizontal actions in these charts.
\end{enumerate}
The family may be empty when both quotients are compact. The data
include the actual peripheral conjugacies after representatives have
been chosen; an unmarked bijection of cusp counts is insufficient.
\end{definition}

\begin{proposition}[Controlled representative, compact and matched-cusp cases; MPR54]
\label{prop:mostow-controlled-qi-producer}
Suppose either both $H,H'$ are compact, or matched geometric cusp data
are supplied. Then there are smooth maps $F,G:X\to X$, equivariant
for $\phi,\phi^{-1}$ respectively, and numbers $L\geq1$, $D,E\geq0$
such that both maps are $L$-Lipschitz and
\[
 d(GF(x),x)\leq D,\qquad d(FG(y),y)\leq E
 \quad\hbox{for all }x,y\in X.
\]
The quotient of $F$ is homotopic to the actual $u$. In the cusp case
$F$ and $G$ restrict to mutually inverse prescribed maps on the paired
horoballs. These outputs supply MPR01 with
$\lambda=L$ and $C=2\max\{D,E\}$, including an equivariant coarse
inverse and coarse surjectivity in both directions.
\end{proposition}
\begin{proof}
MPR49 supplies smooth starting maps $F_0,G_0$ with exactly the opposite
markings. In the compact case take these maps themselves. Their quotient
differential norms are bounded on the compact manifolds, so integration
along paths bounds both lifted maps by a common $L\geq1$. The
compositions commute with their respective deck groups; MPR50 gives
$D,E$ from compactness of the quotients.

In the cusp case MPR51 produces mutually inverse equivariant maps on
the two whole cusp families. Apply MPR52 independently to $F_0$ and
$G_0$, prescribing those inverse maps. The outputs $F,G$ are smooth,
retain their markings, and equal the prescribed maps on every horoball.
MPR53 gives global Lipschitz bounds. On the cusps $GF$ and $FG$ are
exact identities; on the complementary compact quotient cores MPR50
gives $D,E$. MPR48 proves the asserted common quasi-isometry constants.
Finally, geodesic interpolation between $U$ and $F$ is
$\phi$-equivariant and descends to a homotopy of their quotient maps,
as in MPR10. There is no need for $U$ itself to have any metric bound.
\end{proof}

This is an unconditional mathematical producer for the compact input
case, conditional only on the stated inherited standard interfaces.
For the full finite-volume case, MPR60 below supplies its matched cusp
data and MPR61 records the resulting full producer. The conditional
statement of MPR54 is retained as the reusable map-construction step.
MPR68 below supplies MPR02 for this quasi-isometry output; MPR78 later
supplies MPR03 in the same scope. This output cannot be substituted
into MPR47's bilipschitz hypothesis.

\begin{lemma}[Two Lipschitz coarse inverses; MPR48]
\label{lem:mostow-coarse-inverse-estimates}
Let $F:X\to Y$ and $G:Y\to X$ be $L$-Lipschitz, $L\geq1$, and
assume $d(GF(x),x)\leq D$ and $d(FG(y),y)\leq E$, with $D,E\geq0$.
Then
\[
 L^{-1}d(x_1,x_2)-2D/L\leq d(Fx_1,Fx_2)\leq Ld(x_1,x_2),
\]
and the analogous estimate for $G$ holds with $E$ in place of $D$.
Every $y$ is within $E$ of $F(X)$ and every $x$ within $D$ of $G(Y)$.
\end{lemma}
\begin{proof}
The triangle inequality and the Lipschitz bound for $G$ give
\[
 d(x_1,x_2)\leq 2D+d(GFx_1,GFx_2)
 \leq 2D+L\,d(Fx_1,Fx_2).
\]
Rearrange, and interchange $F,G$ for the other lower bound. The upper
bounds are the hypotheses. Choose $G(y)$ and $F(x)$ as the coarse
preimages. Since $L\geq1$, the common additive constant
$C=2\max\{D,E\}$ dominates all four errors. No injectivity, inverse
homeomorphism or small-scale lower Lipschitz bound is concluded.
\end{proof}

\begin{lemma}[Smooth lifts with exactly inverse markings; MPR49]
\label{lem:mostow-smooth-inverse-markings}
For the actual $u,U,\phi$ there are smooth maps $F_0,G_0:X\to X$
whose quotients are homotopic to $u$ and a homotopy inverse of $u$,
with $F_0\gamma=\phi(\gamma)F_0$ and
$G_0\gamma'=\phi^{-1}(\gamma')G_0$.
\end{lemma}
\begin{proof}
Choose a continuous homotopy inverse $v$ and a homotopy from the
identity of $H$ to $vu$. Lift that homotopy starting at the identity
of $X$, and call its endpoint $W$. Uniqueness of homotopy lifts gives
$W\gamma=\gamma W$. Choose a lift $V$ of $v$ so that
$V(U(x_0))=W(x_0)$ at one point. Then $VU=W$ by uniqueness of lifts.
If $\psi$ is the actual equivariance of $V$, comparison at $W(x_0)$
in
\[
 VU\gamma=\psi(\phi(\gamma))VU=\gamma VU
\]
and freeness of the deck action imply
$\psi(\phi(\gamma))=\gamma$. Thus $\psi=\phi^{-1}$, since $\phi$
is an isomorphism. This does not assume that arbitrarily chosen lifts
of homotopy inverses already have inverse markings.

Apply Whitney approximation to $u$ and $v$ on the quotient manifolds.
Lift the resulting homotopies starting at $U$ and $V$. Their smooth
endpoints $F_0,G_0$ retain the actual equivariances, again by uniqueness
of lifts on $X\times[0,1]$. Local covering charts imply smoothness of
the endpoints. No compactness or bound on the homotopy tracks is used.
\end{proof}

\begin{lemma}[Displacement over a compact quotient set; MPR50]
\label{lem:mostow-compact-displacement}
Let $p:X\to X/\Gamma$ be an isometric covering and let $P:X\to X$
be continuous and commute with $\Gamma$. For any compact
$K\subset X/\Gamma$ there is $D\geq0$ such that
$d(Px,x)\leq D$ whenever $p(x)\in K$.
\end{lemma}
\begin{proof}
The continuous function $x\mapsto d(Px,x)$ is $\Gamma$-invariant,
because the action is isometric and $P$ commutes with it. It descends
to a continuous nonnegative function on the quotient. Its restriction
to $K$ is bounded by compactness; take $D=0$ if $K$ is empty.
This bounds displacement upstairs, not merely the quotient distance
between $p(Px)$ and $p(x)$. Neither a compact fundamental domain nor
a rectifiable homotopy is needed for this step.
\end{proof}

\begin{lemma}[Equivariant mutually inverse cusp maps; MPR51]
\label{lem:mostow-equivariant-cusp-maps}
Matched geometric cusp data determine smooth inverse maps $A,B$ of the
full lifted cusp families, equivariant for $\phi,\phi^{-1}$.
They extend to the enlarged cusp collars in those data. On each paired
representative they are
\[
 A_j(z,t)=(a_j(z),t),\qquad B_j(z',t)=(a_j^{-1}(z'),t).
\]
\end{lemma}
\begin{proof}
The equality $\phi(P_j)=P'_j$ gives a bijection
$\Gamma/P_j\to\Gamma'/P'_j$, $\gamma P_j\mapsto\phi(\gamma)P'_j$.
Thus the orbit pairing $\gamma B_j\mapsto\phi(\gamma)B'_j$ is
well-defined and bijective. The same holds for the enlarged collars.
Define there
\[
 A(\gamma x)=\phi(\gamma)A_j(x).
\]
If a point has two such expressions, disjointness of the distinct
cusp pieces implies that both expressions use the same orbit and
the two representatives differ by an element of $P_j$. The affine
conjugacy identity makes their images equal. It also gives smoothness
in each chart and the claimed equivariance. Define $B$ by the inverse
formula. Checking on each representative, then translating, gives
$BA=1$ and $AB=1$ on the respective entire families. All this uses
actual subgroup equality, rather than equality only up to unspecified
conjugacy. Euclidean orientation reversals in the stabilizers are allowed.
\end{proof}

\begin{lemma}[Prescribing smooth cusp tails by a collar modification; MPR52]
\label{lem:mostow-prescribe-cusp-tails}
Let $F_0:X\to X$ be smooth and $\phi$-equivariant. Given $A$ from
MPR51, there is a smooth $\phi$-equivariant map $F$ equal to $A$ on
every cusp horoball, equal to $F_0$ away from the enlarged cusp family,
and equivariantly homotopic to $F_0$.
\end{lemma}
\begin{proof}
Choose a smooth function $\chi:\mathbb R\to[0,1]$ with
$\chi(t)=0$ for $t\leq-2\epsilon/3$ and $\chi(t)=1$ for
$t\geq-\epsilon/3$. Let $I(a,b,s)$ be the parameter-$s$ point of
the geodesic from $a$ to $b$ in $X$, smoothly defined also when $a=b$.
On each enlarged cusp chart set
\[
 F(x)=I(F_0(x),A(x),\chi(t(x))),
\]
and use $F_0$ outside the enlarged family. These definitions agree on
an open collar near its outer edge, where $\chi=0$, so they glue
smoothly. Local finiteness ensures that this is a global smooth map.
Height is invariant under each stabilizer and is transported to all
translates; isometry-equivariance of $I$ then gives
$F\gamma=\phi(\gamma)F$. The region $t\geq0$ has $\chi=1$, so
$F=A$ there. Replacing $\chi$ by $s\chi$, $0\leq s\leq1$, gives
the equivariant homotopy. Interpolation takes place in $X$; it need
not remain inside a target cusp or target core. Apply the identical
construction to $G_0$ and $B$ for the inverse marking.
\end{proof}

\begin{lemma}[Global Lipschitz bounds after cusp prescription; MPR53]
\label{lem:mostow-cusp-prescription-lipschitz}
The two maps constructed by MPR52 from matched geometric cusp data
have finite global Lipschitz constants. Their compositions are exactly
the identity on the corresponding cusp horoballs.
\end{lemma}
\begin{proof}
Write $a_j(z)=T_jz+b_j$. Choose $L_j\geq1$ bounding the operator
norms of $T_j$ and $T_j^{-1}$ between $q_j$ and $q'_j$. The product
form gives, for horizontal $v$ and vertical $a$,
\[
 \|dA_j(v,a)\|^2
 =a^2+e^{-2t}\|T_jv\|_{q'_j}^2
 \leq L_j^2\bigl(a^2+e^{-2t}\|v\|_{q_j}^2\bigr).
\]
The inverse has the analogous bound. Because the deck actions are
isometric, these bounds hold on every translate. There are finitely
many orbit types, so their maximum is finite. On the complementary
compact quotient cores, the newly constructed quotient maps are smooth;
their differential norms consequently have finite upper bounds. This
includes the transition collars, regardless of how far their two
interpolated endpoints are apart before modification. The bounds are
chosen after the actual smooth maps and data, not uniformly over all
homotopy equivalences.

Take a common bound $L\geq1$ for both maps on the core and cusps.
The covering local isometries transfer it to their lifts. Integrate
along piecewise smooth paths in $X$ and take the infimum of their
lengths, giving global $L$-Lipschitz estimates. Finally $F=A$ and
$G=B$ on the full paired horoballs, and $A,B$ carry those horoballs
onto one another. MPR51 therefore gives $GF=1$ and $FG=1$ there.
The maps may fail to be locally injective on the compact cores.
\end{proof}

\paragraph{Source comparison and next producer.}
BP C.1.4 (printed86--87/PDF98--99), Thurston Lemma5.9.1
(printed106/PDF24 of Chapter5), and Martelli Proposition5.2.5
(printed144/PDF152) give the compact pseudo-isometry construction
\cite{BenedettiPetronio1992,ThurstonNotes2002,Martelli2025v4}.
We expose the actual inverse marking, both composition errors and coarse
surjectivity. MPR50 replaces any appeal to the lengths of unspecified
continuous homotopy tracks by an invariant endpoint-displacement bound.
Lee's archived Whitney theorem, label
\code{thm:whitney-approx-manifold}, supplies the smooth approximation
interface \cite{JohnLeeISM}.

Prasad's Theorem C and Section4 (especially printed277--279 and283)
construct smooth controlled maps by prescribing the ends, interpolating
and smoothing, then estimating derivatives \cite{Prasad1973}.
His Lie-group, net-lattice and peripheral hypotheses are not imported
as an automatic adapter to HG04. MPR51--MPR53 are the explicit
hyperbolic-three-space construction from the data above. The next
subsection now produces those data: peripheral preservation, exact affine
conjugacies and the geometric collar/cover identifications. No Waldhausen
or rigidity theorem is used
to turn a homotopy equivalence into a core diffeomorphism.

The boundary development in MPR62--MPR68 below now covers these maps:
additive-error Morse bounds, canonical boundary inverse/equivariance
and uniform normalization. MPR69--MPR73 then supply disk control and
regularity for every normalized composition. MPR74--MPR79 now supply
the returning-zoom/orbit adapter and full marked assembly.
The existing bilipschitz proofs remain available, but their hypotheses
are not changed silently. Full HG04, HG06 and HG08 stay open. Exact
source versions and comparison limits are in
\code{reference_checks_revision111.md}.


\subsection{Cusp matching from the abstract marking}
\label{sec:mostow-cusp-matching}

We now supply the matched data used in MPR54. Work with the full deck
groups of complete connected finite-volume hyperbolic three-manifolds
without boundary, normalized to curvature $-1$. Orientability is not
assumed. The geometric input is the written MGL32 cusp/core theorem,
including its full ideal-center stabilizers and inner collars; MGL13
will be applied only inside the orientation-preserving subgroup.
No preservation of that subgroup by an abstract isomorphism is assumed.

\begin{proposition}[Controlled representative for every finite-volume input; MPR61]
\label{prop:mostow-full-qi-producer}
For every homotopy equivalence $u:H\to H'$ in HG04, fix its lift $U$
and actual marking $\phi:\Gamma\to\Gamma'$. There exist smooth
$\phi$- and $\phi^{-1}$-equivariant maps $F,G:X\to X$ and constants
$L\geq1$, $D,E\geq0$ such that both maps are $L$-Lipschitz,
\[
 d(GF(x),x)\leq D,\qquad d(FG(y),y)\leq E
 \quad(x,y\in X),
\]
and the quotient of $F$ is homotopic to $u$. With
$\lambda=L$ and $C=2\max\{D,E\}$ these maps supply MPR01, including
the quasi-isometry lower bounds and coarse surjectivity in both directions.
\end{proposition}
\begin{proof}
Apply MPR60 below to the actual $\phi$ to obtain matched geometric
cusp data. MPR54 gives the maps, constants and homotopy, using its
stated inherited smooth-approximation, covering, interpolation and
metric-calculus interfaces. MPR48 gives the displayed choice of common
quasi-isometry constants. Empty cusp data are allowed; MPR57 ensures
that the two families are simultaneously empty. For any common negative
constant curvature, first use HG01 to normalize both metrics, and
transport the resulting estimates back with their length scaling.
\end{proof}
This is a written producer relative to the specified inherited interfaces,
not a Lean proof. It neither makes $F$ injective nor supplies the
bilipschitz hypothesis of MPR47. MPR68 below now supplies MPR02's
general quasi-isometry boundary argument; MPR78 later supplies MPR03.

\begin{proposition}[Matched geometric cusp data exist; MPR60]
\label{prop:mostow-matched-cusp-producer}
Every isomorphism $\phi:\Gamma\to\Gamma'$ between these finite-volume
deck groups admits the matched geometric cusp data of
Definition~\ref{def:mostow-matched-cusp-data}, for that exact $\phi$.
\end{proposition}
\begin{proof}
Choose MGL32 cusp families and their compact cores on both quotients.
MPR57 pairs their orbits and chooses target lifts so that
$\phi(P_j)=P'_j$ as actual subgroups. In upper half-space coordinates
centered at each chosen ideal point, the logarithmic-height charts have
the required metrics $dt^2+e^{-2t}q_j$ and $dt^2+e^{-2t}q'_j$.
Their horizontal groups have full rank-two translation lattices of
finite index, by MGL32. Apply MPR58 to $\phi|_{P_j}$ to obtain the
invertible affine maps for the actual horizontal actions. Apply MPR59
to choose one positive collar width for both finite families and to
verify local finiteness of all lifted enlarged horoballs. These are
exactly the three requirements in the matched-data definition.
All choices are finite after orbit representatives are fixed. Neither
a prescribed core diffeomorphism nor a rigidity theorem is used.
\end{proof}

\begin{lemma}[Every parabolic center occurs in the cusp family; MPR55]
\label{lem:mostow-parabolic-center-cusp}
Choose any MGL32 geometric cusp family for $H=X/\Gamma$, with compact
complement $K$. If $a\in\Gamma^+=\Gamma\cap\operatorname{Isom}^+(X)$
is a nonidentity parabolic with fixed point $\xi$, then $\xi$ is the
center of exactly one horoball in the full lifted family. Its full
stabilizer is $\Gamma_\xi$. Conversely every center in that family
is fixed by a nonidentity parabolic in $\Gamma^+$.
\end{lemma}
\begin{proof}
Conjugate $\xi$ to infinity. The oriented parabolic normal form gives
$a(z,y)=(z+v,y)$ with $v\ne0$. On the vertical ray $x_t=(0,e^t)$,
the half-space distance formula gives
\[
 \cosh d(x_t,ax_t)=1+\frac{|v|^2}{2e^{2t}},
 \qquad \operatorname{inj}_H(p(x_t))\leq\tfrac12d(x_t,ax_t)\longrightarrow0.
\]
The positive continuous injectivity-radius function has a positive
minimum on the nonempty compact core $K$. Thus the projected ray is
outside $K$ for every sufficiently large $t$. If $K$ were empty, this
last assertion needs no minimum. Its connected tail lies in a single
open cusp component. Upstairs it therefore lies in one of the disjoint
horoballs covering that component. A geodesic ray eventually contained
in a horoball ends at its ideal center: send that center to infinity;
the downward vertical rays and the semicircular rays leave every fixed
upper horoball, while the upward vertical ray has endpoint infinity.
Consequently this center is $\xi$.

Two horoballs with the same ideal center are nested and intersect,
so disjointness of the lifted family gives uniqueness. MGL32 identifies
its full stabilizer with $\Gamma_\xi$. Conversely, the rank-two
translation subgroup supplied there contains a nonzero translation,
which is an oriented parabolic with that center. This argument also
shows that a compact quotient has no such parabolic.
\end{proof}

\begin{lemma}[Algebraic recognition of full cusp stabilizers; MPR56]
\label{lem:mostow-maximal-virtual-rank-two}
Call a group virtually $\mathbb Z^2$ if it contains a finite-index
subgroup isomorphic to $\mathbb Z^2$. Every virtually $\mathbb Z^2$
subgroup $V<\Gamma$ is contained in a unique full cusp stabilizer.
The full cusp stabilizers are precisely the subgroups maximal, under
inclusion in $\Gamma$, among virtually $\mathbb Z^2$ subgroups.
\end{lemma}
\begin{proof}
Choose $B<V$ of finite index with $B\cong\mathbb Z^2$. The kernel of
the action of $V$ on the finite set $V/B$ is a finite-index normal
subgroup contained in $B$. Intersect that kernel with $\Gamma^+$ to
obtain $A\triangleleft V$, still finite index in $V$ and in $B$.
Finite-index subgroups of $\mathbb Z^2$ are again isomorphic to
$\mathbb Z^2$: they contain $m\mathbb Z^2$ for some positive integer
$m$, and the subgroup-basis theorem over $\mathbb Z$ gives rank two.
This elementary integer-module interface is an inherited binding.

Choose $a\in A\setminus\{1\}$. Commutation makes every element of
$A$ preserve the fixed set $F(a)$ on the ideal boundary. Apply MGL13
in the discrete freely acting group $\Gamma^+$. If $a$ were loxodromic,
then $A$ would lie in its cyclic full elementary stabilizer, impossible
for $A\cong\mathbb Z^2$. Thus $a$ is parabolic, and MGL13 makes every
nonidentity element of $A$ parabolic with the same unique fixed point
$\xi$. For $v\in V$, normality gives $vav^{-1}\in A\setminus\{1\}$;
its fixed point is both $v\xi$ and $\xi$. Hence $V<\Gamma_\xi$.
MPR55 identifies $\Gamma_\xi$ as a full cusp stabilizer, and MGL32
makes it virtually $\mathbb Z^2$.

If $V$ lay in another cusp stabilizer $\Gamma_\eta$, the nonidentity
parabolic $a\in V$ would fix $\eta$, forcing $\eta=\xi$. This proves
uniqueness. In particular a virtually $\mathbb Z^2$ overgroup of
$\Gamma_\xi$ is contained in $\Gamma_\xi$, so the latter is maximal.
Conversely maximality of $V$ forces equality with its containing
stabilizer. There is no assertion that an arbitrary rank-two subgroup
is already the full stabilizer; a proper finite-index sublattice is
an immediate counterexample to that stronger assertion.
\end{proof}

\begin{lemma}[Cusp pairing in the exact marking; MPR57]
\label{lem:mostow-marked-peripheral-pairing}
An isomorphism $\phi:\Gamma\to\Gamma'$ bijects full cusp stabilizers
and their conjugacy classes. Given one source horoball for each cusp
orbit and any full target cusp family, target representatives may be
chosen so that $\phi(P_j)=P'_j$ exactly, with every orbit listed once.
In particular the cusp counts agree, and one quotient is compact if
and only if the other is compact.
\end{lemma}
\begin{proof}
Being virtually $\mathbb Z^2$ and being maximal among such subgroups
are invariant under group isomorphism. Apply MPR56 on both sides;
using $\phi^{-1}$ proves surjectivity as well as maximality of images.
Since $\phi(\gamma P\gamma^{-1})=\phi(\gamma)\phi(P)\phi(\gamma)^{-1}$,
this bijection descends to conjugacy classes, with inverse induced by
$\phi^{-1}$. Full stabilizers of distinct centers are distinct, by
their nonzero parabolic translations. Thus the conjugacy classes are
exactly the geometric cusp orbits.

If a chosen target representative has stabilizer $Q'_k$ and
$\phi(P_j)=\delta'_jQ'_k(\delta'_j)^{-1}$, replace that representative
by $\delta'_j B'_k$ and transport its chart by $\delta'_j$. Its
stabilizer is now exactly $\phi(P_j)$. This changes no full lifted
family or quotient core, and does not change the global marking $\phi$.
MGL32 says that absence of cusp orbits leaves the whole quotient as
its compact core; conversely a nonempty cusp has points of unbounded
height and is an end of the cofinal compact-core exhaustion. This
proves the compactness assertion.
\end{proof}

\begin{lemma}[Affine realization of a marked flat cusp group; MPR58]
\label{lem:mostow-affine-cusp-conjugacy}
Let $P,P'$ act faithfully by Euclidean isometries on two copies of
$\mathbb R^2$ with constant positive-definite metrics. Suppose their
translation subgroups $T,T'$ are full rank-two lattices of finite
index. Every group isomorphism $\psi:P\to P'$ is realized by an
invertible affine map $a$ satisfying
\[
 a\rho(p)=\rho'(\psi(p))a\qquad(p\in P).
\]
No orientation preservation is assumed, and $a$ need not be a similarity.
\end{lemma}
\begin{proof}
Write the actual actions as
$\rho(p)z=Q_pz+b_p$ and $\rho'(\psi(p))z'=R_pz'+b'_p$.
First recognize translations intrinsically. A translation $\tau_v$
has finitely many conjugates, since its conjugates have vectors $Q_pv$
and the linear parts factor through the finite group $P/T$.
If $Q_p\ne I$, choose a lattice vector $v$ with $(I-Q_p)v\ne0$;
one exists because the lattice spans $\mathbb R^2$. The conjugates
\[
 \tau_{nv}\rho(p)\tau_{-nv}(z)
     =Q_pz+b_p+n(I-Q_p)v\qquad(n\in\mathbb Z)
\]
are pairwise distinct. Therefore $T$ is exactly the set of elements
having finite conjugacy class. The same holds for $T'$, so
$\psi(T)=T'$. This characterization does not confuse the translation
lattice with a maximal cyclic subgroup generated by a glide reflection.

Choose a lattice basis $v_1,v_2$ for $T$. Write
$\psi(\tau_{v_i})=\tau_{w_i}$. The $w_i$ form a lattice basis for
$T'$, so the unique real linear map $L$ with $Lv_i=w_i$ is invertible.
For every lattice vector $v$, $\psi(\tau_v)=\tau_{Lv}$.
Apply $\psi$ to $p\tau_vp^{-1}=\tau_{Q_pv}$ to obtain
$LQ_pv=R_pLv$. Since the lattice spans the plane,
\[
 LQ_p=R_pL.
\]
Define $c_p=b'_p-Lb_p$. Multiplication of the two actual affine actions
and this intertwining identity give
\[
 c_{pq}=c_p+R_pc_q,\qquad c_t=0,\quad R_t=I\quad(t\in T).
\]
Thus $c_{pt}=c_p=c_{tp}$, and both $c_p$ and $R_p$ descend to the
finite quotient $F=P/T$. Put $m=|F|\geq1$ and
\[
 v_0=\frac1m\sum_{\bar q\in F}c_q.
\]
For fixed $p$, multiplication by $\bar p$ permutes $F$. Summing the
cocycle identity therefore gives
\[
 mv_0=mc_p+R_pmv_0,\qquad c_p=v_0-R_pv_0.
\]
For $a(z)=Lz+v_0$ it follows that
\[
 a(\rho(p)z)=LQ_pz+Lb_p+v_0
            =R_p(Lz+v_0)+b'_p=\rho'(\psi(p))a(z).
\]
Its inverse is $z'\mapsto L^{-1}(z'-v_0)$. This realizes the specified
isomorphism on every group element, not just on the translation subgroup.
For a torus cusp $F$ is trivial; for a Klein-bottle cusp it has order
two. The same finite-sum argument treats both.
\end{proof}

\begin{lemma}[Locally finite enlarged cusp families; MPR59]
\label{lem:mostow-cusp-collar-local-finiteness}
The MGL32 cusp data for two finite-volume quotients can be restricted
to one common width $\epsilon>0$ so that all their charts extend to
$t>-\epsilon$ and each full lifted enlarged family is pairwise
disjoint and locally finite in $X$. Their compact cores and exact
full stabilizers are unchanged.
\end{lemma}
\begin{proof}
MGL32 gives positive inner collar widths with pairwise disjoint
enlarged images. Take $\epsilon$ smaller than every width in the two
finite lists; for empty lists choose any positive number. Their lifts
are horoballs, and distinct lifted components are disjoint by the
covering description and disjointness on the quotient. Restricting
the negative-height collars changes neither the level-zero cores
nor the stabilizers of the horoballs and their centers.

Every pairwise disjoint family of open horoballs in $X$ is locally
finite. To see this, let $C\subset X$ be compact. For each horoball
meeting $C$, choose a point $x$ in the intersection and move distance
two along the ray toward its center to a point $y$. The open unit
ball about $y$ lies in that horoball, since signed distance from its
boundary has increased by two. These unit balls are disjoint and
their centers lie in the compact closed two-neighborhood of $C$.
Distinct centers have distance at least two: otherwise the midpoint
of their joining segment belongs to both unit balls. A compact metric
space contains only finitely many points separated by at least two,
as follows by a finite cover with balls of radius less than one.
Hence only finitely many horoballs meet $C$. Taking a compact closed
ball about each point proves local finiteness. Properness of $X$ and
its standard horoball-distance formula are inherited model interfaces.
\end{proof}

\paragraph{Source comparison and formalization boundary.}
Martelli V4, Lemmas4.2.1--4.2.2 and Corollary4.2.4,
printed115--116/PDF123--124, give the commuting-isometry and
rank-two parabolic comparison. BP D.3.18(2), printed157/PDF168,
must be read with care: a rank-two subgroup is contained in a peripheral
group, but need not equal it. MPR56 explicitly supplies the maximality
step. It also treats finite extensions needed for nonorientable cusps.
Ratcliffe's first edition (1994), Theorem7.4.4 and proof,
printed313--314/PDF324--325, is the marked affine-realization comparison.
Its finite-quotient averaging is made direct in MPR58, after proving the
translation subgroup characteristic by finite conjugacy classes; no
unqualified uniqueness claim about maximal free abelian subgroups is
used. The needed full translation lattices already come from MGL32.
No general Bieberbach classification or splitting-group construction
is an additional dependency. These comparisons and the precise source
qualification are recorded in \code{reference_checks_revision114.md}.

The new proofs use finite-index group operations, integer lattice bases,
finite sums, finite-dimensional real linear algebra, compactness and
the hyperbolic model/covering interfaces. Their eventual Lean bindings
must use the accepted foundation. MPR61 closes the written input-map
producer MPR01 in its full finite-volume scope. The next subsection
supplies the general quasi-isometry boundary extension, followed by
disk control and regularity in MPR69--MPR73. MPR74--MPR79 complete
the returning-zoom/orbit argument and qualitative marked assembly.
HG06/HG08's quantitative consumers and all Lean implementation remain open. In particular no bilipschitz upgrade follows from the cusp matching.


\subsection{Boundary extension with additive quasi-isometry errors}
\label{sec:mostow-qi-boundary}

We now supply MPR02 in its general quasi-isometry scope. Work in
curvature $-1$ with the inherited model, length, compactness and
geodesic interfaces of Section~\ref{sec:mostow-morse-extension}.
For this subsection a controlled pair consists of maps $F,G:X\to X$
and constants $\lambda\geq1$, $C,D,E\geq0$ such that both maps satisfy
\begin{equation}
\label{eq:mostow-qi-pair}
 \begin{aligned}
 \lambda^{-1}d(x,y)-C&\leq d(Jx,Jy)\leq\lambda d(x,y)+C
 \quad(J=F,G),\\
 d(GFx,x)&\leq D,\qquad d(FGy,y)\leq E.
 \end{aligned}
\end{equation}
Neither continuity nor injectivity is required for the boundary
construction. For continuous $F$, the resulting extension to the closed
ball is continuous, but need not be a homeomorphism on the interior.
The boundary maps will be mutually inverse homeomorphisms.

The symmetric convention loses no input from MPR01. If only $F$ has
the displayed quasi-isometry constants, its coarse inverse satisfies
\[
 \lambda^{-1}d(x,y)-(2E+C)/\lambda
 \leq d(Gx,Gy)\leq\lambda d(x,y)+\lambda(2E+C).
\]
Indeed apply the upper, respectively lower, bound for $F$ to $Gx,Gy$
and use the two errors at $FGx,FGy$. Enlarge $C$ to
$\max\{C,\lambda(2E+C)\}$ for a common convention. The smooth maps
of MPR61 already satisfy a common choice by MPR48.

\begin{proposition}[The full boundary and normalization producer; MPR68]
\label{prop:mostow-qi-boundary-producer}
Every controlled pair has canonical mutually inverse boundary
homeomorphisms $h,g:S^2\to S^2$. They respect composition, bounded
changes, and actual equivariance. For a family with common
$\lambda,C,D,E$, convergence of three fixed distinct boundary points
to three distinct limits gives bounded interior basepoint images and
a subsequence on which $h_n$ and $g_n$ converge uniformly to mutually
inverse homeomorphisms. Thus MPR02 is supplied, including for the
exactly marked maps of MPR61 and their isometry normalizations.
\end{proposition}
\begin{proof}
MPR62--MPR63 provide uniform complete-line tracking despite the
additive errors. MPR64 constructs the canonical boundary homeomorphisms,
using the estimates of MPR65 after first constructing an inverse pair
of endpoint bijections. MPR66 supplies three-point tameness and MPR67
supplies uniform boundary compactness with inverses. For actual deck
equivariance use MPR64; MPR61 supplies the input pair in every
finite-volume case. For isometries $a,b$, replace the pair by
$aFb,b^{-1}Ga^{-1}$. The same constants work and the boundary map is
$\partial a\,h\,\partial b$, by MPR64. The bilipschitz class is the
special case $\lambda=K,C=D=E=0$.
\end{proof}
This is a written producer relative to inherited interfaces. Boundary
compactness alone neither puts a limit in the original isometry orbit
nor proves marked boundary rigidity MPR03. Section~\ref{sec:mostow-qi-regularity}
supplies general quasi-isometry disk control and regularity;
Section~\ref{sec:mostow-qi-rigidity} then supplies returning zooms
with preservation of the marked isometry orbit.

\begin{lemma}[Additive-error tracking of Lipschitz curves; MPR62]
\label{lem:mostow-additive-curve-tracking}
Let $\beta:I\to X$ be $L$-Lipschitz, where $L,\lambda\geq1$, and suppose
\[
 d(\beta(s),\beta(t))\geq\lambda^{-1}|s-t|-A,
 \qquad A\geq0.
\]
Put $Q=L\lambda$, $r=2Q$, and $M=r+Q(2r+A)$.
For compact $I=[a,b]$ with distinct image endpoints, the image is in
the closed $M$-tube about their complete supporting line. For
$I=\mathbb R$, there is a unique complete line $\ell$ at finite
Hausdorff distance from the image, and that distance is at most $M$.
The two parameter ends tend to different ideal endpoints of $\ell$.
With the foot of $\beta(0)$ as coordinate zero and positive orientation
chosen by the positive end, its signed projection $v$ satisfies
\[
 \begin{aligned}
 |v(t)-v(s)|&\geq\lambda^{-1}|t-s|-A-2M,\\
 v(t)&\geq t/\lambda-A-2M
 \quad\text{if }t>\lambda(A+2M).
 \end{aligned}
\]
\end{lemma}
\begin{proof}
Every subarc of parameter length $t-s$ and image endpoints $p,q$
has length $b\leq L(t-s)\leq Q(d(p,q)+A)$.
For the finite-interval assertion, an excursion beyond distance $M$
from the supporting line has endpoints at distance $r$, remains
outside its open $r$-tube, and has length $b>2(M-r)=2Q(2r+A)$.
The endpoints of the entire curve are on the line, so both excursion
endpoints exist. MPR34 and the two connectors give
\[
 d(p,q)\leq2r+b/\cosh r,\qquad
 b\leq Q(2r+A)+Qb/\cosh r.
\]
Since $Q\geq1$ and $\cosh(2Q)\geq1+2Q^2>2Q$,
the last inequality implies $b<2Q(2r+A)$, a contradiction.
Repeated image points cause no problem: length is parameterized
length, and no inverse map of the curve was used.

For the whole line, take $n$ with $2n>\lambda A$; the endpoints of
$\beta([-n,n])$ are distinct. Their supporting lines meet the closed
$M$-ball about $\beta(0)$. Choose nearest points and unit tangents,
then extract convergent pointed frames over this compact ball.
For each fixed $t$, its projections on these lines remain in a compact
ball about $\beta(t)$. Their coordinates relative to the chosen frames
are bounded. Continuous geodesic flow gives $d(\beta(t),\ell)\leq M$
on the limiting line, on the same subsequence for every fixed $t$.

The two projection connectors yield the displayed lower bound for
$|v(t)-v(s)|$. Thus $|v(t)|\to\infty$ at both parameter ends;
continuity fixes the sign on each sufficiently long tail. If both
signs agreed, the intermediate value theorem on the two tails would
give arbitrarily separated $s<0<t$ with $v(s)=v(t)$, contradicting
$|t-s|\leq\lambda(A+2M)$. The signs are opposite, and $v$ is onto.
This proves the reverse Hausdorff inclusion. Orient toward the positive
end. The sign on $t>\lambda(A+2M)$ is positive, giving the last bound.
MPR37 transfers both projected ideal limits to $\beta$. Any other
line at finite Hausdorff distance has those same endpoints by MPR37,
and hence equals $\ell$.
\end{proof}

\begin{lemma}[Tracking a general quasi-isometry on lines; MPR63]
\label{lem:mostow-qi-line-tracking}
Suppose $F$ satisfies the two quasi-isometry inequalities with
$\lambda,C$, without any continuity assumption. Set
\[
 \begin{aligned}
 L&=\lambda+C,\quad R=2L,\quad A=C+2R,\\
 Q&=L\lambda,\quad M=2Q+Q(4Q+A),\quad N=M+R.
 \end{aligned}
\]
For each complete unit-speed geodesic $c$, the set $F(c(\mathbb R))$
has Hausdorff distance at most $N$ from a unique complete geodesic.
Each point of either set is within distance at most $N$ of an actual
point of the other. The limits as the parameter tends to positive and negative infinity
are its two distinct ideal endpoints. These constants depend only
on $\lambda,C$.
\end{lemma}
\begin{proof}
Join $F(c(k))$ to $F(c(k+1))$ by its constant-speed geodesic on
each $[k,k+1]$, $k\in\mathbb Z$, using a constant path when they
coincide. The resulting continuous curve $\beta$ is $L$-Lipschitz:
each piece is, and finite partitions at the intervening integers
give the same bound across pieces. For $k\leq t\leq k+1$,
\[
 d(\beta(t),F(c(t)))\leq
 d(\beta(t),F(c(k)))+d(F(c(k)),F(c(t)))\leq2L=R.
\]
It follows that
$d(\beta(s),\beta(t))\geq\lambda^{-1}|s-t|-C-2R$.
MPR62 applies with the displayed $A,M$. Combining its two Hausdorff
inclusions with the pointwise $R$-bound gives $N$. MPR37 transfers
both ideal limits from $\beta$ to $F\circ c$, and also proves
uniqueness of the tracked line. In particular the construction does
not assume that the possibly discontinuous image curve is rectifiable.
\end{proof}

\begin{proposition}[Canonical boundary maps of coarse inverse pairs; MPR64]
\label{prop:mostow-qi-canonical-boundary}
A controlled pair has unique boundary homeomorphisms $h,g$ with
$g=h^{-1}$ and the following sequential property:
\[
 x_n\to\xi\in S^2\ \Longrightarrow\ F(x_n)\to h(\xi),
 \qquad y_n\to\eta\in S^2\ \Longrightarrow\ G(y_n)\to g(\eta).
\]
They send line endpoints to the corresponding MPR63 endpoints.
The boundary map is unchanged by a bounded change of the interior map
and respects composition. If $F\gamma=\phi(\gamma)F$, then
$h\,\partial\gamma=\partial\phi(\gamma)\,h$ for that exact $\phi$.
When $F$ is continuous, $F\cup h$ is its unique continuous extension
to the closed ball; no interior homeomorphism is asserted.
\end{proposition}
\begin{proof}
For a ray ending at $\xi$, complete it to a line and define $h(\xi)$
by its positive image limit in MPR63. Rays with the same endpoint
stay a bounded synchronized distance apart by MPR37. Their images
do also, using the quasi-isometry upper bound including $C$.
MPR37 proves ray-independence. Distinct endpoints of a line have
distinct image limits. Construct $g$ in the same way from $G$.

For a complete line $c$, use the polygonal $\beta$ of MPR63 and
its tracked line $\ell'$, oriented toward $h(\xi)$ on the positive
end. Write $\ell'(v(t))$ for the projected point, so $v(t)\to+\infty$
and $d(F(c(t)),\ell'(v(t)))\leq N$. Then
\[
 d(c(t),G(\ell'(v(t))))\leq D+\lambda N+C.
\]
The second point tends to $g(h(\xi))$ by the definition of $g$;
MPR37 proves $gh=\mathrm{id}$. Interchanging the pair, with $E$
instead of $D$, proves $hg=\mathrm{id}$. This proves bijectivity
before any continuity is assumed. MPR65 below applies to these
endpoint bijections and the two-sided line tracking, proving that
both are continuous. They are therefore inverse homeomorphisms.

To check convergence from the interior, write $x_n=c_{\xi_n}(t_n)$
using unit-speed rays from the ball origin $o$. Then $t_n\to\infty$
and $\xi_n\to\xi$. Complete each ray to a line and use the integer
interpolation with parameter zero at $o$, so $\beta_n(0)=F(o)$.
The feet of $F(o)$ on the tracked lines lie in its closed $M$-ball.
Their oriented signed projection coordinates at $t_n$ satisfy
$v_n(t_n)\geq t_n/\lambda-A-2M\to\infty$ eventually, by MPR62.
Every subsequence has a further subsequence of convergent oriented
frames over this compact ball. Its positive endpoints converge to
$h(\xi)$, by boundary continuity. The model geodesic formula with
coordinate tending to positive infinity makes the projected points
tend to that endpoint. MPR37 and the uniform $N$-bound transfer this
limit to $F(x_n)$. The subsequence criterion proves the assertion;
repeat for $G$.

The ray limits make the boundary map unique and invariant under
bounded changes, by MPR37. Composing the sequential limits proves
compatibility with composition, whose quasi-isometry and coarse
inverse bounds follow by substitution. Isometries have their usual
boundary maps. Apply these statements to the exact equality
$F\gamma=\phi(\gamma)F$ to obtain the asserted equivariance.
If $F$ is continuous, the sequential assertion plus continuity of
$h$ proves continuity on the whole closed ball, also for sequences
mixing interior and boundary points. Density gives uniqueness.
For a discontinuous $F$, only its asserted boundary and sequential
properties are claimed.
\end{proof}

\begin{lemma}[Visual bounds with coarse errors; MPR65]
\label{lem:mostow-qi-visual-bounds}
Let a controlled pair have mutually inverse endpoint bijections $h,g$
which match the tracked lines of MPR63 in both directions. No continuity
of these bijections is assumed. If $B=d(o,F(o))$, then in the boundary
metric $\delta$ of MPR20,
\begin{align*}
 \delta(h\xi,h\eta)&\leq2e^{B+C+N}\delta(\xi,\eta)^{1/\lambda},\\
 \delta(g\xi,g\eta)&\leq2e^{\lambda B+2C+D+N}
                         \delta(\xi,\eta)^{1/\lambda}.
\end{align*}
In particular both endpoint maps are continuous.
\end{lemma}
\begin{proof}
For distinct $\xi,\eta$, put $a=d(o,\ell_{\xi\eta})$ and
$a'=d(o,\ell_{h\xi,h\eta})$. For every $y$ on the latter line,
the reverse Hausdorff inclusion supplies $x$ on the former with
$d(y,Fx)\leq N$. Hence
\[
 d(o,y)\geq d(Fo,Fx)-B-N\geq a/\lambda-C-B-N.
\]
Take the infimum and apply MPR20 to obtain
$\delta(h\xi,h\eta)\leq2e^{B+C+N}e^{-a/\lambda}
\leq2e^{B+C+N}\delta(\xi,\eta)^{1/\lambda}$.
For the inverse pair, the basepoint satisfies
\[
 d(o,G(o))\leq d(o,GF(o))+d(GF(o),G(o))
             \leq D+\lambda B+C.
\]
Apply the first calculation to $G$ and $g$, using the same $N$.
Equal endpoints give zero on both sides. These uniform power bounds
prove continuity without assuming it. The reverse Hausdorff inclusion,
not merely containment of $F(\ell)$ in a tube, is essential here.
\end{proof}

\begin{lemma}[Uniform quasi-isometry three-point tameness; MPR66]
\label{lem:mostow-qi-three-point-tameness}
Let $F_n$ occur in controlled pairs with common $\lambda,C$ and
boundary maps $h_n$. If $h_n(\xi_i)\to\eta_i$ for three fixed distinct
$\xi_i$ and three distinct $\eta_i$, then $F_n(o)$ is bounded.
All $F_n(x)$ are bounded for each fixed $x$. If also $D$ is common,
the coarse-inverse basepoint images $G_n(o)$ are bounded.
\end{lemma}
\begin{proof}
For all sufficiently large $n$, the ordered target triples are
distinct. The explicit three-point Mobius normalization supplies
boundary maps $m_n\to\mathrm{id}$ carrying $h_n(\xi_i)$ to $\eta_i$;
their hyperbolic isometries $a_n$ converge to the identity. This is
the inherited isometry/compactification interface also used in MPR19.
For each domain side $\ell_{ij}$ choose its nearest point $q_{ij}$
to $o$. By MPR63 and the upper quasi-isometry bound,
\[
 d(a_nF_n(o),\ell_{\eta_i\eta_j})
 \leq\lambda d(o,q_{ij})+C+N.
\]
The maximum of these three fixed radii is finite. MPR39 places
$a_nF_n(o)$ in a fixed compact triple-tube intersection. The inverses
$a_n^{-1}$ converge, so their action sends this compact set into
a bounded set. The finitely many omitted indices are harmless.
For fixed $x$, use $d(F_nx,F_no)\leq\lambda d(x,o)+C$.
The bound for $G_n(o)$ follows from MPR65's basepoint calculation.
\end{proof}

\begin{lemma}[Uniform quasi-isometry boundary compactness; MPR67]
\label{lem:mostow-qi-boundary-compactness}
Let controlled pairs have common $\lambda,C,D,E$ and bounded
$F_n(o)$. Some subsequence of their boundary maps and inverse boundary
maps converges uniformly on $S^2$ to mutually inverse homeomorphisms.
This applies in particular under MPR66's three-point hypotheses.
If every $h_n$ fixes infinity, both limits fix infinity and both
convergences are locally uniform in Euclidean coordinates on the plane.
\end{lemma}
\begin{proof}
MPR64 supplies the actual inverse boundary maps, even though $G_n$
need not equal an interior inverse. MPR65 gives a common modulus
for both boundary families. Compact-domain Arzela--Ascoli, applied
successively to the two families on the compact sphere, gives uniform
limits $h,g$. To retain inverse identities, for example, use
\[
 \delta(h_n(g_nx),h(gx))\leq
 \sup_y\delta(h_ny,hy)+\delta(h(g_nx),h(gx))\longrightarrow0
\]
uniformly in $x$ by uniform continuity of $h$. Thus $hg=\mathrm{id}$;
the other order gives $gh=\mathrm{id}$. MPR66 supplies the basepoint
bound under its hypotheses. Fixed infinity persists under uniform
convergence. The images of a compact plane set under either limit
stay a positive spherical distance from infinity; eventually the
corresponding images lie in one bounded chart. Uniform equivalence
of spherical and Euclidean metrics there proves both local uniform
convergences. No interior limit or membership in an isometry orbit
has been inferred.
\end{proof}

\paragraph{Source comparison and remaining work.}
Schwartz's v5 Section2, Lemmas2.1--2.4 and the extension/tameness proofs
(printed/PDF11--14; TeX670--955), provide the projection, line-limit
and normalization pattern. His theorem is stated for bilipschitz maps.
MPR62 retains the additive term in the parameterized length estimate;
MPR63 supplies unit-mesh interpolation; MPR64 uses coarse errors rather
than an interior inverse. These are explicit project extensions of
that pattern, not assertions that his BL statement already has this scope.
Gouezel--Shchur Theorem1.1 and conventions (printed/PDF1--2) concern
compact-interval quasi-geodesics and a sharper bound in general
Gromov-hyperbolic spaces. Its full proof and Isabelle certification
are not imported as Lean evidence. The elementary hyperbolic-model
argument here needs neither that theorem nor a hyperbolicity-constant
translation. Exact versions and correction checks are recorded in
\code{reference_checks_revision115.md}.

MPR68 supplies MPR02 after MPR61 supplies MPR01. MPR69--MPR73 below
supply general quasi-isometry disk control and regularity. MPR74--MPR79
then supply the returning-zoom/orbit adapter and marked conclusion.
The earlier BL arguments remain unchanged and retain their hypotheses.
All inherited bindings, Lean implementations, quantitative flow
consumers and ambient-flow peripheral injection remain open.


\subsection{General quasi-isometry disk control and regularity}
\label{sec:mostow-qi-regularity}

We continue with the controlled pairs and constants
$(\lambda,C,D,E)$ of Section~\ref{sec:mostow-qi-boundary}.
The following consumer supplies regularity for every normalized
isometry composition of the actual finite-volume representative.
Its proof is placed before the supporting disk estimates. The
returning-zoom and isometry-orbit argument remains a separate step.

\begin{proposition}[Regularity away from the current cusp centers; MPR73]
\label{prop:mostow-qi-orbit-regularity}
Let $F,G$ be the exactly marked controlled pair supplied by
Proposition~\ref{prop:mostow-full-qi-producer}, with source and target
full deck groups $\Gamma,\Gamma'$ and marking $\phi$.
Let $h=\partial F$, and let $P$ be the set of centers of all lifted
source cusp horoballs supplied by MGL32. For every pair of hyperbolic
isometries $a,b$ for which
\[
 h_{a,b}=\partial a\circ h\circ\partial b
 \quad\hbox{fixes }\infty,
\]
the set of asterisks of $h_{a,b}$ outside
$P_b=\partial b^{-1}(P)$ is Borel and has positive planar measure.
Here an asterisk has exactly the meaning of
Definition~\ref{def:mostow-asterisk}: every rational directional
derivative exists, and the horizontal derivative is nonzero.
The normalized pair has source group $b^{-1}\Gamma b$, target group
$a\Gamma'a^{-1}$ and exact marking
\[
 \phi_{a,b}(\beta)=a\phi(b\beta b^{-1})a^{-1}.
\]
For each selected asterisk outside $P_b$, every geodesic ray to that
point has arbitrarily late returns to the transported compact core.
There is no asserted common good set for all $a,b$, no uniform lower
bound for its measure, and no uniform return time.
\end{proposition}
\begin{proof}
MPR68 supplies the canonical boundary homeomorphism and shows that
$F_{a,b}=aFb$ and $G_{a,b}=b^{-1}Ga^{-1}$ have the same
$(\lambda,C,D,E)$. Direct substitution gives
$F_{a,b}\beta=\phi_{a,b}(\beta)F_{a,b}$; MPR64 transfers this
identity to the boundary. Proposition~\ref{prop:mostow-qi-regularity}
below gives a Borel asterisk set of positive measure for $h_{a,b}$.

The source quotient for $b^{-1}\Gamma b$ is isometric to the original
quotient by the map induced by $b$. Pull back the compact core and
each lifted horoball by $b^{-1}$. The new full cusp-center set is
exactly $P_b$. It is countable: MGL32 gives finitely many cusp orbits
and MGL30 gives countability of the full deck group. A countable
subset of the plane is Borel and null. Removing
$P_b\cap\R^2$ therefore retains a Borel set of positive measure.
The transported geometric data satisfy
Lemma~\ref{lem:mostow-core-return}, which gives the asserted late
returns for any chosen endpoint outside $P_b$. Empty cusp families
are allowed. This proves the statement separately for each $a,b$;
no intersection over the uncountable family of normalizations is taken.
\end{proof}

\begin{proposition}[General quasi-isometry boundary regularity; MPR72]
\label{prop:mostow-qi-regularity}
Let $(F,G)$ be any controlled pair on $\mathbb H^3$, and suppose its
canonical boundary map $h$ fixes $\infty$. Then the asterisk set of
$h|_{\R^2}$ is Borel and has positive planar measure. The same holds
for $h^{-1}|_{\R^2}$. No interior continuity or injectivity,
equivariance, lattice, or compact quotient is required.
\end{proposition}
\begin{proof}
MPR64 supplies inverse boundary homeomorphisms. The uniform disk
conclusion of Proposition~\ref{prop:mostow-qi-uniform-disks} below
is exactly the hypothesis of
Proposition~\ref{prop:mostow-disk-regularity}, with its fixed finite
constant $\kappa>1$. That proposition uses only a plane homeomorphism
and disk control; its packing, absolute-continuity, Borel and Fubini
arguments MPR25--MPR32 have no interior bilipschitz hypothesis.
It yields the required asterisk set. Apply the same argument to
$(G,F)$ for the inverse; this swaps $D,E$ and preserves the constant
chosen below. We retain the precise asterisk conclusion and make no
claim here of full Fr\'echet differentiability or a nonsingular
Fr\'echet derivative.
\end{proof}

\begin{proposition}[Uniform disks for controlled quasi-isometries; MPR71]
\label{prop:mostow-qi-uniform-disks}
For every $\lambda\geq1$ and $C,D,E\geq0$ there is
$\kappa=\kappa(\lambda,C,D,E)>1$ with the following property.
For every controlled pair with these constants whose canonical
boundary map $h$ fixes $\infty$, and for every $c\in\mathbb C$ and
$r>0$, there are closed disks $D_1,D_2$ of positive radii centered
at $h(c)$ such that
\[
 D_1\subset h(B(c,r)),\qquad
 h(\overline B(c,r))\subset D_2,\qquad
 \frac{\operatorname{diam}D_2}{\operatorname{diam}D_1}\leq\kappa.
\]
Here $B(c,r)$ is open. The constant is chosen before the map and disk.
Let $N=N(\lambda,C)$ be MPR63's tracking constant, and let
$B_*=B_*(\lambda,C)$ be the bound in
Lemma~\ref{lem:mostow-qi-fixed-triple-bound} below. One choice is
\[
 S=\max\{D,E\},\qquad
 H=2\exp(\lambda B_*+2C+S+N),\qquad
 \kappa=(2\sqrt2\,H^2)^\lambda.
\]
This is an explicit formula relative to the compactness bound $B_*$;
an effective or optimal numerical value of $B_*$ is not asserted.
\end{proposition}
\begin{proof}
Set $w=h(c+r)-h(c)$. The boundary homeomorphism gives $w\ne0$.
The complex similarities $z\mapsto c+rz$ and
$z\mapsto (z-h(c))/w$ extend to hyperbolic isometries $b,a$.
The normalized boundary map
\[
 q(z)=\frac{h(c+rz)-h(c)}{w}
\]
fixes $0,1,\infty$. By MPR68, it comes from the controlled pair
$(aFb,b^{-1}Ga^{-1})$ with the same constants.
Lemma~\ref{lem:mostow-qi-fixed-triple-bound} bounds $d(o,aFb(o))$
by $B_*$. MPR65 therefore gives, for the visual half-chord metric
$\delta$ based at $o=(0,1)$,
\[
 \delta(q^{-1}\xi,q^{-1}\eta)
 \leq H\,\delta(\xi,\eta)^{1/\lambda}.
\]
We used $D\leq S$; in particular $H\geq1$.
Lemma~\ref{lem:mostow-visual-disk-radii} below gives inner and outer
radii $\rho=(2H)^{-\lambda}$ and $R=(\sqrt2 H)^\lambda$ for the
unit disk. Undoing the target similarity gives $D_1,D_2$ with radii
$|w|\rho,|w|R$. The source similarity identifies the open and closed
unit disks with $B(c,r)$ and $\overline B(c,r)$, respectively.
The diameter ratio is $R/\rho=\kappa$, as required. No interior
homeomorphism or limit in an isometry orbit is used.
\end{proof}

\begin{lemma}[Fixed-triple basepoint bound for quasi-isometries; MPR69]
\label{lem:mostow-qi-fixed-triple-bound}
Fix $\lambda\geq1$ and $C\geq0$. There is a finite
$B_*=B_*(\lambda,C)\geq0$ such that every controlled pair with these
metric constants whose boundary map fixes $0,1,\infty$ satisfies
$d(o,F(o))\leq B_*$, where $o=(0,1)$ in the half-space model.
The bound does not depend on $D,E$ or on a quotient.
\end{lemma}
\begin{proof}
Write $\ell_{ij}$ for the three complete geodesics joining the
standard triple, and put
\[
 d_* = \max_{i<j}d(o,\ell_{ij}),\qquad
 R_* = \lambda d_*+C+N(\lambda,C).
\]
For a nearest point $x_{ij}\in\ell_{ij}$, MPR63 and MPR64 place
$F(x_{ij})$ within $N$ of the same geodesic $\ell_{ij}$, because
its endpoints are fixed. Consequently
\[
 d(F(o),\ell_{ij})
 \leq d(F(o),F(x_{ij}))+N
 \leq \lambda d(o,x_{ij})+C+N\leq R_*.
\]
By MPR39, the intersection $T$ of the three closed $R_*$-tubes is
compact. Choose $B_*\geq0$ bounding $d(o,z)$ on $T$ (take $B_*=0$
if $T$ is empty). This choice uses only the standard triple and
$\lambda,C$, and is made before $F$. Since $F(o)\in T$, it supplies
the claimed bound. No sequence of maps or cocompactness is needed.
\end{proof}

\begin{lemma}[Visual inverse bounds give explicit disk radii; MPR70]
\label{lem:mostow-visual-disk-radii}
Use the boundary chart $\mathbb C\cup\{\infty\}$ based at
$o=(0,1)$. Let $q$ be a boundary homeomorphism fixing $0,\infty$,
and suppose, for $H\geq1$ and $\lambda\geq1$, that
\[
 \delta(q^{-1}\xi,q^{-1}\eta)
 \leq H\,\delta(\xi,\eta)^{1/\lambda}
 \quad\hbox{for all }\xi,\eta.
\]
Then, with $\rho=(2H)^{-\lambda}$ and $R=(\sqrt2 H)^\lambda$,
\[
 \overline B(0,\rho)\subset q(B(0,1)),\qquad
 q(\overline B(0,1))\subset\overline B(0,R).
\]
\end{lemma}
\begin{proof}
Choose the isometry from the half-space model to the ball model
sending $o$ to its center. In boundary stereographic coordinates one
may take
\[
 \sigma(z)=\frac{(2\operatorname{Re}z,2\operatorname{Im}z,|z|^2-1)}
                   {1+|z|^2},\qquad
 \sigma(\infty)=(0,0,1).
\]
The visual metric in MPR20 is half the Euclidean chord length on
this unit sphere. Subtracting $\sigma(0)$ and $\sigma(\infty)$ gives
\[
 \delta(z,0)=\frac{|z|}{\sqrt{1+|z|^2}},\qquad
 \delta(z,\infty)=\frac1{\sqrt{1+|z|^2}}.
\]
If $|w|\leq\rho$, then $q^{-1}(w)$ is finite and
\[
 \delta(q^{-1}w,0)\leq H\,\delta(w,0)^{1/\lambda}
 \leq H|w|^{1/\lambda}\leq\tfrac12.
\]
Hence $|q^{-1}(w)|\leq1/\sqrt3<1$. This strict margin puts the
\emph{closed} inner disk inside the image of the \emph{open} unit disk.
For $|z|\leq1$, apply the same inverse estimate to $qz,\infty$:
\[
 \frac1{\sqrt2}\leq\delta(z,\infty)
 \leq H\,\delta(qz,\infty)^{1/\lambda}.
\]
Thus $\delta(qz,\infty)\geq(\sqrt2 H)^{-\lambda}=R^{-1}$,
so $|qz|\leq\sqrt{R^2-1}\leq R$. This includes the unit circle
and proves the second containment.
\end{proof}

\paragraph{Source comparison and formalization boundary.}
Schwartz's Lemma~2.2 and proof of Lemma~1.3 use the fixed ideal
triangle to control basepoints (archived TeX739--781, printed/PDF12).
His Theorem~3.6 and its proof, TeX1106--1141, printed/PDF17, give disk
control for the bilipschitz setting through normalization and compactness.
Here MPR69 makes the fixed-triple bound uniform over controlled pairs;
MPR70--MPR71 derive disk radii directly from MPR65's inverse visual
estimate. This general quasi-isometry extension and the displayed
constants are project arguments, not claims of Schwartz's theorem.
The analytic input is his Sections3.1--3.3, TeX960--1248,
printed/PDF15--18, as already expanded and qualified in MPR25--MPR32.
The current cusp family comes from the retained MGL32, with MPR06
handling late returns. Exact checks and inherited binding limits are
recorded in \code{reference_checks_revision116.md}.

MPR69--MPR73 now supply disk control and regularity in the full
controlled quasi-isometry class, including every normalized composition
needed for subsequent zooms. The original bilipschitz statements of
MPR05, MPR18, MPR24 and MPR33 are unchanged. The following subsection
supplies the general quasi-isometry returning-zoom/orbit adapter and
its marked assembly. Boundary regularity and late return times alone
do not assert an interior limit or an isometry-orbit conclusion.
All accepted inherited bindings and Lean proofs remain pending.

\subsection{Returning quasi-isometry zooms and the marked rigidity conclusion}
\label{sec:mostow-qi-rigidity}

The following argument completes the qualitative HG04 proof relative
to the inherited interfaces already stated. The final consumer is
presented first. Its new input is the returning-zoom argument for
continuous controlled quasi-isometries; the earlier bilipschitz
statements and proofs are retained.

\begin{theorem}[Finite-volume marked Mostow--Prasad producer; MPR79]
\label{thm:mostow-full-marked-producer}
Let $H,H'$ and $u:H\to H'$ satisfy HG04, with the same fixed negative
curvature normalization. Relative to the inherited model, covering,
smooth approximation, integration, compactness and analysis interfaces
specified in the Margulis and Mostow developments, $u$ is homotopic
to a unique isometry $H\to H'$. Neither compactness nor orientability
is required. This supplies HG04's written mathematical producer.
\end{theorem}
\begin{proof}
If the common sectional curvature is $-\kappa^2$, multiply both metrics
by $\kappa^2$ to obtain curvature $-1$. HG01 preserves completeness,
finite volume and the actual homotopy class, and an isometry for the
rescaled metrics is an isometry for the original metrics.
Choose universal local-isometric covers and a lift $U$ of the actual
$u$, with its induced marking $\phi$. MPR61 constructs smooth,
exactly $\phi$- and $\phi^{-1}$-equivariant controlled maps $F,G$ for
these inputs, including all cusped and nonorientable cases. MPR68
supplies their canonical boundary homeomorphism $h$, with exact
$\phi$-equivariance. Theorem~\ref{thm:mostow-qi-boundary-rigidity}
below gives an isometry $A$ with boundary value this same $h$.

Apply MPR10: faithfulness of the boundary isometry action gives
$A\gamma=\phi(\gamma)A$, so $A$ descends to an isometry $H\to H'$.
Geodesic interpolation between $U$ and $A$ descends to a homotopy
from the actual $u$. MPR46 proves uniqueness among isometries in that
homotopy class using MPR45's finite-volume centralizer result.
The marking throughout is the one determined by $U$; no equality of
arbitrarily chosen based homomorphisms is assumed. Rescale back if
necessary. This proves the stated conclusion.
\end{proof}

\begin{theorem}[Finite-volume quasi-isometry boundary rigidity; MPR78]
\label{thm:mostow-qi-boundary-rigidity}
Let $\Gamma,\Gamma'$ be full deck groups of complete connected
finite-volume hyperbolic three-manifolds of curvature $-1$.
Let $(F,G)$ be a controlled pair with $F$ continuous, equivariant
for an isomorphism $\phi:\Gamma\to\Gamma'$ and its inverse,
respectively. Under the inherited interfaces of this development,
the canonical boundary map $h=\partial F$ is the boundary value of
a hyperbolic isometry. This supplies MPR03 in its required scope.
\end{theorem}
\begin{proof}
First compose on the target with an isometry taking $h(\infty)$ to
$\infty$. Lemma~\ref{lem:mostow-qi-normalized-class} below retains
all controlled-pair and marked finite-volume hypotheses. We may
therefore work with $h(\infty)=\infty$ and undo this composition
at the end.

Choose two distinct plane directions $D_1,D_2$. By MPR75 below the
asterisk set has positive measure after deleting the current cusp
centers. Choose $z$ in this set. MPR76 gives a returning limit $h_1$
with an affine pencil through $z$. MPR77, applied to the current pair
and its own cusp-center set, gives a second returning limit $h_2$
for which $D_1$ is good. Both limits are actual isometry compositions
of the original pair, by MPR74, and fix infinity. MPR75 applies anew
to $h_2$, giving an asterisk outside its current cusp set. Repeat
MPR76 and MPR77 for $D_2$, obtaining $h_3,h_4$.

MPR15 preserves $D_1$ through these last two homothety zoom sequences
and their locally uniform plane limits. We keep those actual zoom
coordinates: the auxiliary isometries in the orbit identities are
not used as new plane coordinates. Thus $h_4$ has two distinct good
directions. MPR16 makes $h_4$ an invertible real affine map.

Compose the orbit identities, including the initial normalization,
to write
\[
 h_4=\partial a\,h\,\partial b,
 \qquad \Gamma_4=b^{-1}\Gamma b,
 \qquad \phi_4(\beta)=a\phi(b\beta b^{-1})a^{-1}.
\]
MPR41 applied to the original finite-volume $\Gamma$ at the point
$\partial b(\infty)$ gives $\gamma\in\Gamma$ moving that point.
Consequently $b^{-1}\gamma b\in\Gamma_4$ moves infinity.
MPR17, with this actual group element and $\phi_4$-equivariance,
extends $h_4$ to a hyperbolic isometry $A_4$, allowing either
orientation. The isometry $a^{-1}A_4b^{-1}$ has boundary value the
original $h$. No claim that the interior map $F$ itself is an isometry
or a homeomorphism was needed.
\end{proof}

\begin{lemma}[Returning orbit extraction for continuous controlled pairs; MPR74]
\label{lem:mostow-qi-returning-orbit}
Let $(F,G)$ be a controlled pair with $F$ continuous and with the
exact equivariance of MPR78. Write $h=\partial F$ and fix any $o\in X$.
For isometries $a_n,b_n$ set $F_n=a_nFb_n$. Suppose $F_n(o)$ is
bounded and $p(b_n(o))$ lies in a fixed compact set $K$ of the source
quotient. Suppose a compact set $D_0\subset X$ contains a deck
translate of every point of $p^{-1}(K)$.
Then a subsequence and isometries $a,b$ satisfy
\[
 \begin{gathered}
 F_n\longrightarrow aFb\quad\hbox{uniformly on compact subsets of }X,\\
 \partial a_n\,h\,\partial b_n\longrightarrow\partial a\,h\,\partial b
 \quad\hbox{uniformly on }S^2.
 \end{gathered}
\]
The limit pair $(aFb,b^{-1}Ga^{-1})$ has the same
$(\lambda,C,D,E)$ and the exact conjugated marking of MPR73.
The boundary limit is a homeomorphism. No interior injectivity or
continuity of $G$ is required, and no interior inverse is invoked.
\end{lemma}
\begin{proof}
Choose $\gamma_n\in\Gamma$ with $\gamma_nb_n(o)\in D_0$ and put
\[
 B_n=\gamma_nb_n,\qquad A_n=a_n\phi(\gamma_n)^{-1}.
\]
Exact equivariance, both inside and on the boundary by MPR64, gives
$F_n=A_nFB_n$ and $\partial a_n h\partial b_n=\partial A_n h\partial B_n$.
Properness of isometry evaluation at $o$ extracts $B_n\to b$.
Continuity of $F$ makes $F(D_0)$ compact, and
\[
 d(A_n(o),F_n(o))=d(o,F(B_n(o)))
\]
is bounded. Properness of $X$ and of evaluation extracts a further
subsequence $A_n\to a$.

For a compact $L\subset X$, the points of all $B_n(L)$ lie in one
closed bounded ball, hence in a compact set. Uniform continuity of
$F$ there and compact-uniform continuity of the isometry action give
$A_nFB_n\to aFb$ uniformly on $L$. On the compact boundary sphere,
$\partial A_n,\partial B_n$ converge uniformly and $h$ is uniformly
continuous. Their compositions therefore converge uniformly to
$\partial a h\partial b$, a homeomorphism. These are the stated limits,
not unidentified compactness limits.

Conjugation by isometries preserves the metric inequalities and both
composition errors: for example the new $GF$ is $b^{-1}GFb$.
Direct substitution gives the marking
$\beta\mapsto a\phi(b\beta b^{-1})a^{-1}$; MPR64 identifies the
canonical boundary of $aFb$ with the displayed limit. The argument
contains no use of $F^{-1}$. It does not assert convergence of the
normalized $G$ maps or a homeomorphism on the interior.
\end{proof}

\begin{lemma}[The normalized finite-volume controlled class; MPR75]
\label{lem:mostow-qi-normalized-class}
For every pair in MPR78, isometry compositions
$(aFb,b^{-1}Ga^{-1})$ retain its hypotheses and constants. The source
and target groups are $b^{-1}\Gamma b$ and $a\Gamma'a^{-1}$, and the
marking is $\beta\mapsto a\phi(b\beta b^{-1})a^{-1}$.
Each source quotient has a compact core, a countable full lifted
cusp-center set and compact lift data for that core. If its boundary
map fixes infinity, its asterisks outside its own cusp-center set have
positive planar measure. Every endpoint outside that set of centers
has arbitrarily late geodesic returns to the core. All these statements
also apply to the orbit representatives produced by MPR74.
\end{lemma}
\begin{proof}
MPR64 and MPR68 give the controlled-pair and boundary conclusions;
continuity of $F$ persists under composition with isometries. The
induced quotient isometries preserve completeness and finite volume.
MGL32 supplies the geometric compact core and full horoball family,
including nonorientable quotients and empty cusp families. MGL30
supplies countability of the deck group and compact lift sets.
For a chosen source family with centers $P$, pull back its horoballs
and compact lift sets by $b^{-1}$; its new centers are exactly
$\partial b^{-1}(P)$. Its core is transported by the induced quotient
isometry. MPR72 supplies positive-measure Borel asterisks for every
normalized boundary map; deleting its countable planar cusp set
preserves positivity. MPR06 applies to the transported data and
proves late returns. This is the same cusp-avoidance argument as
MPR73, now recorded for any pair in MPR78. Each orbit representative
from MPR74 has exactly this form and retains the hypotheses.
\end{proof}

\begin{lemma}[A returning quasi-isometry zoom gives an affine pencil; MPR76]
\label{lem:mostow-qi-asterisk-pencil}
Let a pair in MPR78 be normalized by $h(\infty)=\infty$, and let $z$
be an asterisk outside its current cusp-center set. There is a limit
of returning homothety zooms whose boundary map $h_1$ fixes infinity,
is in the marked isometry orbit of $h$, and is affine on every full
line through $z$. Its continuous interior representative and controlled
coarse inverse belong to the class of MPR75.
\end{lemma}
\begin{proof}
Use the fixed basepoint $o_z=(z,1)$ and the unit-speed ray
$t\mapsto(z,e^{-t})$. MPR75 and MPR06 give return times
$r_n\to\infty$. Put $s_n=e^{r_n}$ and define hyperbolic isometries
\[
 b_s(w,t)=(z+s^{-1}(w-z),s^{-1}t),\qquad
 a_s(w,t)=(h(z)+s(w-h(z)),st).
\]
In particular $b_{s_n}(o_z)=(z,e^{-r_n})$ lies over the compact core.
Their boundary zoom is
\[
 h_n(w)=h(z)+s_n\bigl(h(z+(w-z)/s_n)-h(z)\bigr).
\]
At $z,z+1,z+2$ the images converge to $h(z),h(z)+D_1h(z)$ and
$h(z)+2D_1h(z)$, respectively, which are distinct. MPR66 gives
bounded images of every fixed interior point, hence of $o_z$.
MPR74 extracts a single subsequence with the actual orbit limit
$h_1=\partial a h\partial b$. It is a sphere homeomorphism fixing
infinity. The chart argument following MPR12 gives local uniform
Euclidean convergence on the plane.

For each rational $v\ne0$ and each fixed real $t$, the derivative
at $z$, evaluated at parameter $t/s_n$, yields
\[
 h_1(z+tv)=h(z)+tD_vh(z).
\]
The case $t=0$ is exact; for $t\ne0$ multiply the defining two-sided
derivative limit by $t$. Injectivity of $h_1$ makes every such slope
nonzero. For any real $v\ne0$, choose nonzero rational vectors
$v_j\to v$ and apply MPR12 to the constant map sequence $h_1$, center
$z$ and these directions. This extends affinity to every full line
through $z$. The subsequence was chosen before $v,t$; no full
Fr\'echet derivative or interior injectivity is used.
\end{proof}

\begin{lemma}[A returning quasi-isometry pencil yields a good direction; MPR77]
\label{lem:mostow-qi-pencil-parallel}
Let a pair in MPR78 be normalized, and suppose every line through
$z$ is affine for its boundary map $h$. For every plane direction $D$
there is a returning homothety-zoom limit in its marked isometry orbit
for which $D$ is good. Its representative remains in the class of MPR75.
\end{lemma}
\begin{proof}
Choose $w\in z+D$ outside the current countable cusp-center set and
outside $\{z\}$. Such a point exists because a real line is uncountable.
Put $v=w-z\ne0$. Use MPR06 with basepoint $o_w=(w,1)$ and returning
scales $s_n=e^{r_n}\to\infty$. The zooms centered at $w$ have the
formula of MPR76 with $z$ replaced by $w$. On the affine line
$z+D=w+\mathbb Rv$ their boundary maps equal $h$ exactly.
Thus the images of $w,w+v,w+2v$ are fixed distinct points.
MPR66 bounds the interior images of $o_w$, and MPR74 gives a fixed
subsequence with a marked orbit limit $h_1$, again fixing infinity
and locally uniformly convergent on the plane.

The zoomed map $h_n$ is affine on every line through
$c_n=w+s_n(z-w)=w-s_nv$. Fix an arbitrary $x\in\mathbb R^2$ and set
\[
 v_n=v+(x-w)/s_n,\qquad x-c_n=s_nv_n.
\]
Then $v_n\to v\ne0$, so for all large $n$ the line
$x+\mathbb Rv_n$ passes through $c_n$. MPR12 gives affinity of
$h_1$ on $x+\mathbb Rv$. This holds for every $x$ with the same
subsequence, and $D$ is good. The center $w$ was selected from the
current pencil line, using countability; no assertion that the
asterisk set meets every prescribed line was needed.
\end{proof}

\paragraph{Source comparison, scope and next work.}
Schwartz's archived v5 TeX332--408, printed/PDF6, contains the
three-point and orbit lemmas, with compact quotient and an interior
inverse in Lemma~1.4's proof. MPR74 instead uses returns to a compact
core and the displayed bounded-basepoint identity, requiring only
continuity of $F$. TeX411--565, printed/PDF7--9, contains the four-zoom
argument, good directions and the affine-to-isometry step. MPR76--MPR77
supply the general-QI returning-scale adapters, with basepoints
$o_z=(z,1)$ and $o_w=(w,1)$ chosen so that the zoom images lie on the
actual returning rays. The real scales need not be integers.
The unchanged MPR12 and MPR15--MPR17 provide the line-limit and affine
arguments at their original, purely boundary-map scope.
Schwartz's Section6, TeX1892--1929, printed/PDF28, states the broader
scope briefly; it does not supply these finite-volume adapters.
Exact readings and correction-check limits are in
\code{reference_checks_revision117.md}.

MPR78 and MPR79 now supply the qualitative marked rigidity and unique
isometry conclusion in the full HG04 scope, relative to inherited
interfaces. The earlier BL theorem statements remain unchanged.
This is a written mathematical proof, not an accepted Lean theorem.
Section~\ref{sec:mostow-dependency-audit} now records the end-to-end
dependency and inherited-interface audit. Accepted-foundation binding
and independent mathematical review remain pending. The quantitative HG06/HG08
comparison and ambient-flow peripheral injection remain separate;
neither follows from qualitative rigidity alone.

\subsection{End-to-end dependency and inherited-interface audit}
\label{sec:mostow-dependency-audit}

The conclusion being audited is MPR79, with the original HG04
hypotheses. Revision118 checks its producer-to-consumer transfers
against the written arguments and the pinned source comparisons.
All preceding theorem statements and proof bodies are retained.
The review found two register edges that confused a supplied hypothesis
or an application with a proof prerequisite. Correcting these edges
removes an apparent bilipschitz dependency without changing the proof.
This is a mathematical reading and dependency audit by the working
assistant, not independent certification, an accepted foundation binding,
or Lean verification.

\paragraph{Proof edges versus references to contracts.}
MPR25 assumes the open/closed disk \emph{conclusion} specified in
MPR24. Its proof's phrase ``MPR24 gives'' uses those assumed disks;
it does not apply the bilipschitz producer to an arbitrary plane map.
In the general quasi-isometry branch, MPR71 supplies precisely those
disks before MPR72 applies MPR32 and its analytic lemmas.
Accordingly the register removes MPR24 from MPR25's proof dependencies
and records it as the locator of the assumed contract. The metric
hypotheses of the bilipschitz theorem do not enter MPR25--MPR32.

Likewise MPR68 proves its boundary and normalization assertions for
\emph{every} controlled pair, from MPR64, MPR66 and MPR67. The reference
to MPR61 instantiates that theorem for a later finite-volume application;
it is not needed to prove the general theorem. Its register edge is
removed and the application relation is retained explicitly. MPR79
still requires MPR61 to produce its actual maps. No theorem is being
made available without its stated input pair.

A transitive traversal of the corrected MPR79 register no longer
includes the bilipschitz disk producer MPR24 or its bilipschitz boundary
predecessors. The register remains conservative: MPR68 includes the
boundary compactness assertion MPR67, whose proof uses Arzela--Ascoli,
even though the actual returning-orbit extraction MPR74 does not use
that compactness theorem. No minimal-import claim is made. In particular
an acyclic graph is only a consistency check; the following transfers
supply the mathematical reasons for the applications.

\paragraph{Checked producer-to-consumer transfers.}
\begin{enumerate}[leftmargin=*,label=\textbf{T\arabic*.}]
\item \textbf{Original scope and normalization.}
HG01 and MPR79 multiply both metrics of curvature $-\kappa^2$ by
$\kappa^2$. This retains the actual map and homotopy class, and rescales
both sides equally when the final isometry is returned. The full deck
groups act freely; no orbifold or orientability hypothesis is inserted.
HG03's oriented torus statement is not the nonorientable cusp producer:
that producer is MGL32.
\item \textbf{Full marked cusp data.}
MGL32 supplies full ideal-center stabilizers, rank-two translation
lattices of index at most two, compact cores and inner collars.
MPR55--MPR57 recognize the full stabilizers as maximal virtually
$\mathbb Z^2$ subgroups and choose target representatives with
$\phi(P_j)=P'_j$ exactly. MPR58 realizes the specified isomorphism on
the entire Euclidean group; MPR59 supplies one positive collar width
and local finiteness. These are all the fields consumed by MPR60/MPR54.
The empty families are paired. Neither a proper sublattice nor an
unmarked cusp bijection suffices.
\item \textbf{The actual controlled maps.}
MPR49 aligns the inverse lift to have marking $\phi^{-1}$; MPR51--MPR53
prescribe mutually inverse cusp maps and bound the differential norms.
MPR50 bounds the composition errors over the compact cores upstairs.
MPR48 gives common constants $\lambda=L$ and $C=2\max(D,E)$ for both
maps. MPR61 therefore supplies continuity of $F$, both metric bounds,
and both composition bounds required later. No metric estimate on the
original lift $U$, injectivity of $F$, or core diffeomorphism is used.
\item \textbf{Canonical boundary before continuity.}
MPR62--MPR63 track complete lines in both directions with actual
point witnesses at distance at most $N$. MPR64 first constructs endpoint
maps and proves their inverse identities using the $D,E$ errors.
Only then does MPR65 prove continuity of those endpoint bijections.
Its reverse-tracking estimate is essential. The signed escape estimate
also supplies arbitrary interior sequences approaching the boundary,
so composition and exact equivariance concern this canonical boundary.
\item \textbf{Uniform disk control.}
Every normalization $aFb,b^{-1}Ga^{-1}$ has the same
$(\lambda,C,D,E)$. MPR69 chooses $B_*(\lambda,C)$ before the map,
and MPR71 chooses its disk ratio before the map and disk. MPR70 places
a closed inner disk inside the image of an open disk, with preimage
radius at most $1/\sqrt3<1$. This strict margin is the one needed by
MPR25's disjoint closed packings. No uniform Euclidean Holder bound
on the entire plane is assumed.
\item \textbf{The exact regularity conclusion.}
MPR25--MPR32 use disk control, planar homeomorphism, finite packing,
one-dimensional absolute continuity and Borel Fubini--Tonelli.
The result is a positive-measure Borel set where every rational
directional derivative exists and the horizontal derivative is nonzero.
MPR72 supplies this for every normalized controlled pair. Neither
Frechet differentiability nor a nonsingular full derivative is an input
to MPR76. The nonzero slopes there follow from injectivity of the limit.
\item \textbf{Returns for the current source group.}
MGL30 gives compact lift witnesses and countability. MGL32 gives the
full horoball family, also without orientability. Under $aFb$ the source
is $b^{-1}\Gamma b$, the centers are $\partial b^{-1}(P)$, and the core
is transported by the quotient isometry. Removing that countable set
preserves positive planar measure. MPR06 supplies arbitrarily late
real return times for each chosen endpoint, with no uniform return rate.
\item \textbf{An actual isometry-orbit limit.}
MPR66 bounds every fixed interior basepoint. MPR76/MPR77 use respectively
$(z,1)$ and $(w,1)$, whose homothety images lie exactly on the returning
rays. MPR74 uses compact lifts, continuity of $F$, proper isometry
evaluation and exact equivariance to obtain $aFb$ and its canonical
boundary $\partial a h\partial b$. Boundary compactness alone would
not identify this orbit. There is no use of an interior inverse of $F$
or convergence of the normalized coarse inverses.
\item \textbf{The four zooms and their coordinates.}
MPR75 retains the class and exact marking after each orbit extraction;
its reference to MPR78 names that theorem's hypotheses only. MPR76
chooses one subsequence before the line and real parameter. MPR77
chooses its center on an uncountable pencil line outside the current
countable exceptional set, without assuming that every line meets the
asterisk set. MPR15 preserves the first good direction in the actual
positive-homothety coordinates through the second two zooms.
MPR16 then gives an invertible real affine map.
\item \textbf{Recovering the marked isometry.}
For $h_4=\partial a h\partial b$, MPR41 is applied at
$\partial b(\infty)$ in the original source group. It supplies an
element of $b^{-1}\Gamma b$ moving infinity. MPR17 yields $A_4$,
and $a^{-1}A_4b^{-1}$ has boundary the original $h$. MPR10 uses
boundary faithfulness to obtain the actual $\phi$-equivariance and
descends the isometry. Geodesic interpolation from $U$ gives ordinary
homotopy; no bounded-track or proper-homotopy assertion is needed.
\item \textbf{Uniqueness in the same homotopy class.}
MPR41--MPR45 use positive finite volume with Riemannian densities,
finite-index covers and the full isometry group. The reversing case
is included. MPR46 lifts the given homotopy starting at the first
isometric lift, thereby giving the second lift the same marking.
Their quotient centralizes $\Gamma$, hence is the identity.
Arbitrary independently chosen lifts are not equated.
\item \textbf{Downstream scope.}
The result is qualitative rigidity of the complete hyperbolic models.
HG06/HG08 still need the actual core/end comparison, its induced marking,
and quantitative compactness with the stated parameter order.
Intrinsic cusp injection from MGL32 does not supply injection into an
ambient Ricci-flow slice. These remain separate obligations.
\end{enumerate}

\paragraph{Inherited interfaces that must be bound.}
The following twelve families organize the existing mathematical inputs;
they are not new axioms, new missing-theorem declarations, or claims
about an inspected PC release. Each exact declaration, elaborated type,
imports, normalization and axiom closure must be checked against the
accepted completed foundation. Empty binding fields are intentional.
General mathematics already written in the MGL/MPR proofs is not
reclassified as an imported theorem.

\begin{description}[leftmargin=2.3em,style=nextline]
\item[I00: space forms and scale]
Complete connected curvature-$-1$ smooth three-manifolds have universal
local-isometric covers by $X$ and free properly discontinuous full deck
actions. The standard metric rescaling preserves completeness and has
length/volume factors $\kappa,\kappa^3$. Consumers: MGL32, MPR61, MPR79.
\item[I01: covering and homotopy calculus]
Existence and uniqueness of lifts and homotopy lifts, the deck/fundamental
group identification, intermediate finite-index covers, quotient descent,
and local smoothness of lifts. Manifolds are second countable and have
relatively compact coordinate shrinks. Consumers: MGL30/MGL32,
MPR49--MPR52, MPR42, MPR46 and MPR10. Actual based choices must remain visible.
\item[I02: explicit hyperbolic models]
Ball, half-space and negative-last-coordinate hyperboloid identifications;
distance, tensor and density formulas; geodesics, endpoints and nearest
points; full Lorentz and Mobius/anti-Mobius actions; stabilizers and faithful
boundary action. These support the written model calculations in
MGL13--MGL29 and MPR17/MPR20/MPR34/MPR37/MPR41--MPR45.
No rigidity, general Morse theorem or quasiconformal regularity theorem
is imported under this heading.
\item[I03: properness and action topology]
Closed bounded subsets of $X$ and unit frames over compact sets are
compact. The full isometry group has its compact-open topology, with
proper evaluation $a\mapsto a(o)$ and compact-uniform action on the
closed-ball compactification. Thus convergent isometries converge
uniformly on the compact boundary, including inverses. Consumers:
MGL12, MPR62, MPR66 and MPR74. Properness of $X$ alone is not a statement
of properness of this evaluation map.
\item[I04: Riemannian metric and length calculus]
Hopf--Rinow, positivity/continuity of injectivity radius for these
quotients, differential bounds implying path-length bounds, and the
almost-everywhere chain rule and length formula for Lipschitz curves
in smooth charts. Consumers: the Margulis compact-core arguments,
MPR34/MPR62 and MPR50/MPR53/MPR55. No differentiability of an arbitrary
quasi-isometry is presumed.
\item[I05: smooth approximation and interpolation]
Whitney approximation with an ordinary homotopy for maps of smooth
manifolds without boundary, smooth cutoffs on the real line, and smooth
isometry-equivariant interpolation along the unique segment in $X$,
including coincident endpoints. Consumers: MPR49/MPR52 and MPR10.
Compact support, bounded homotopy tracks and relative core diffeomorphisms
are not part of these input contracts.
\item[I06: integration with densities]
Positive volume of a nonempty coordinate ball, local and global change
of variables for Riemannian densities, countable Borel partitions and
additivity, and integration of the displayed warped product densities.
Consumers: MGL02/MGL04/MGL08/MGL32 and MPR41--MPR44. The change-of-variables
interface must allow orientation reversals and finite-volume noncompact domains.
\item[I07: groups and finite-dimensional algebra]
Subgroups, centers, cosets, finite-index kernels, generated subgroups,
commutators and their behavior under homomorphisms; real operator norms,
matrix multiplication, complex $2\times2$ eigenlines and real orthogonal
eigenspaces. Consumers: MGL19--MGL29, MPR56/MPR58 and MPR45.
Zassenhaus and the finite-index Margulis step are written arguments,
not inherited rigidity assumptions.
\item[I08: lattices and affine actions]
The integer subgroup-basis theorem and finite-index rank preservation,
real extension of an isomorphism of lattice bases, finite sums and
finite-dimensional linear algebra. Consumers: MPR56 and MPR58.
The characteristic translation subgroup and exact finite-cocycle
averaging argument are proved there, so no unspecified affine realization
or Bieberbach theorem is needed for that step.
\item[I09: planar measure and one-dimensional analysis]
Lebesgue outer measure and area, disk area, completeness of null sets,
countable subadditivity, invariance of null sets under orthogonal maps,
Borel Fubini--Tonelli, and differentiability almost everywhere and
zero-derivative constancy for absolutely continuous real functions on
compact intervals. Consumers: MPR25--MPR32/MPR72. Borel derivative sets
and the interval-selection and packing arguments are written explicitly.
\item[I10: compact metric limits]
Sequential compactness, uniform continuity on compact sets, finite
packing, compact-domain Arzela--Ascoli and elementary limit calculus.
MPR67 uses Arzela--Ascoli for maps and inverses; MPR74 instead obtains
its orbit limit from I03. Countable selection of return times and the
uncountability of a real line are also used in MPR06/MPR77.
\item[I11: names of flat sections]
Classification of closed connected surfaces and multiplicativity of
Euler characteristic under a finite cover identify MGL32's nonorientable
flat section as a Klein bottle. Its geometric horoball, full stabilizer,
rank-two lattice and cusp-product conclusions are established before
this naming step. Those are the fields used in Mostow--Prasad.
The current packaged MGL32 theorem retains its naming conclusion.
\end{description}

\paragraph{Audit conclusion and next frontier.}
No new unfilled mathematical producer was found in these transfers,
relative to the interfaces above and the retained Margulis/source audits.
All 80 MPR and 33 MGL entries remain unformalized; the two register
corrections change no theorem statement or proof. The detailed record
\code{checks/evidence/revision118_mostow_dependency_audit.json} carries
the transfer/interface ledger, unchanged-text hashes and exact edge changes.
\code{reference_checks_revision118.md} distinguishes this turn's readings
from reused checks and records the bounded external correction check.

Revision119 develops the HG06 consumer in
Section~\ref{sec:hyp-quantitative-consumer}: a global homotopy equivalence
suffices for MPR79, without MPR09's product-collar conditions. The
quantitative compactness step is written there conditional on the still
qualified topological comparison. Revision120 completes the requested
final working audit in Section~\ref{sec:mostow-final-audit}, relative to
the stated inherited interfaces.
Accepted-foundation binding and independent mathematical review remain
necessary before claiming a formal theorem. No changing PC snapshot
inspection or new foundational Lean development is authorized by this audit.


\subsection{Final statement-and-proof audit of qualitative rigidity}
\label{sec:mostow-final-audit}

Revision120 completes the requested working review of the full written
Mostow--Prasad argument, starting from master119. This review reads the
78 MPR statement/argument entries, the two MPR section records, all
33 MGL statements and proofs, the HG04 target, and the surrounding
model, marking, cusp, regularity and inherited-interface conventions.
It is a mathematical review by the working assistant; the accompanying
hash and dependency checks only check record consistency. It is neither
independent certification nor a Lean proof.

\paragraph{The exact theorem reviewed.}
HG04 and MPR79 concern connected complete finite-volume hyperbolic
three-manifolds without boundary, with the same constant negative
sectional curvature, and a supplied homotopy equivalence between them.
The conclusion is a unique isometry in that homotopy class. Neither
compactness nor orientability is added. Common metric rescaling in HG01
reduces to curvature $-1$ and preserves the conclusion. After choosing
universal covers and a lift of the supplied map, the proof retains its
actual group isomorphism $\phi$ through the construction. The final
homotopy is an ordinary homotopy; it need not have bounded tracks.
There is no claim that an arbitrary supplied homotopy equivalence has
a quasi-isometric lift before the controlled representative is built.

\paragraph{Corrections made by this review.}
MPR25 and MPR29 now explicitly require a plane homeomorphism with
uniform disk control. Their earlier statements referred only to the
conclusion or disk control of MPR24, leaving the homeomorphism hypothesis
in the surrounding context. MPR25 uses injectivity for disjoint image
packings; MPR29 also uses it for disjoint image strips and continuity for
a compact finite-area bound. The MPR25 register now includes this
hypothesis. MPR29's register already included it. In MPR25 the proof
invokes the assumed disk property directly, so no bilipschitz producer
is introduced as an analytic prerequisite.

These are clarifications of the standalone contracts, not additional
hypotheses on HG04: MPR32 already assumes a plane homeomorphism, and
MPR68/MPR71/MPR72 provide precisely that map and the uniform open/closed
disk property for the general quasi-isometry route. Its finite constant
is greater than one. The actual producer-consumer chain remains valid.
The older MPR11 section record is also updated to point to MPR79:
QRG03 can use the homotopy equivalence itself, and MPR09 remains an
optional product-collar special case. The qualified radial producer
QRG00 is not promoted to a proved theorem.

\paragraph{Margulis and cusp audit.}
MGL26--MGL29 give the fixed matrix neighborhood and discrete-group
nilpotence, using the operator norm and a common depth of vanishing
commutators. MGL12's finite-index step uses the connected graph of
left cosets with edges $gA\mapsto sgA$; it does not assume $A$ normal.
A finite cover of the compact stabilizer is chosen before the uniform
small displacement. MGL19--MGL25 then supply the elementary-group
classification used in the constant-curvature three-dimensional case.
The source's norm convention and its right-coset proof are not silently
substituted for these written project arguments.

MGL13--MGL18 handle precise invariance and thin components. All powers
of a loxodromic generator enter the tube-radius maximum; a larger thin
threshold supplies the closure buffer. MGL02--MGL06 use finite volume
for compact thick parts, finitely many components, and exclusion of
rank-one cusps. MGL32 passes through the orientation cover and retains
full cusp stabilizers, flat torus or Klein-bottle cross-sections, and
actual peripheral inclusions. It does not identify the full group's
thin set with that of the orientation subgroup at one threshold.

MPR55--MPR60 recognize full cusp groups as maximal virtually
$\mathbb Z^2$ subgroups and identify their characteristic translation
lattices by finite conjugacy classes. Rank-two containment alone is
insufficient. Finite-holonomy averaging produces the affine cusp
conjugacy, including the orientation-reversing case. Disjoint enlarged
horoballs are locally finite, so the later prescriptions define maps
on the whole universal cover. MPR49--MPR54 align inverse markings by
homotopy lifting, prescribe inverse affine maps on ends, interpolate
smoothly across buffered collars, and bound derivatives and both
composition errors over compact quotient cores. MPR61 thus supplies
the controlled pair used in the theorem, including the empty-cusp case.

\paragraph{Boundary, regularity and returning-zoom audit.}
MPR62--MPR68 keep additive quasi-isometry errors, two-sided complete-line
tracking, and both coarse inverse bounds. Endpoint bijections are
constructed before continuity is proved; visual estimates yield
homeomorphisms and simultaneous compactness of map and inverse.
Three-point normalization yields bounded interior basepoints, but does
not itself imply membership in a marked isometry orbit.

MPR69--MPR71 give a single disk constant for the whole normalized class,
with a closed inner disk contained in the image of the open input disk.
MPR25--MPR32 use finite packings, approximate-supremum interval selection,
Borel directional-derivative loci, absolute continuity on almost every
line, and Fubini to obtain a positive-measure set of asterisks. Only
rational directional derivatives and one specified nonzero derivative
are required at this stage. The countable exceptional cusp set is
removed separately for each current conjugated group; no intersection
across all isometry normalizations is taken.

MPR74 uses compact returns and exact equivariance to extract an actual
isometry-orbit limit, requiring continuity of the chosen representative
but no interior inverse map. MPR76--MPR77 use real return times and the
basepoints $(z,1)$ or $(w,1)$ on the corresponding vertical rays.
Affine pencil limits are first obtained for one convergent map, then
extended by density; moving pencils converge to parallel affine lines
with their parametrizations. Four zooms preserve the first good
direction while producing the second. Auxiliary orbit isometries are
not used to change the plane coordinates for that argument. MPR41
provides an actual lattice element moving infinity; MPR17 then turns
the resulting affine map into a similarity. MPR78 concludes marked
boundary rigidity, with both orientation signs allowed.

\paragraph{Descent, uniqueness and source boundaries.}
MPR10 uses faithfulness of the boundary action to descend the isometry
with the original $\phi$, and equivariant geodesic interpolation to
recover the prescribed homotopy class. MPR41--MPR46 rule out a finite
boundary orbit and an invariant plane by finite-volume arguments with
densities, then handle the full isometry centralizer, including its
orientation-reversing case. Homotopy lifting therefore supplies the
exact uniqueness assertion, not just equality up to marking.

The source check reopens Schwartz's compact bilipschitz theorem,
extension/orbit and four-zoom arguments, analytic packing proof, and
finite-volume discussion, together with Ratcliffe's matrix and Margulis
passages. Schwartz's brief cusp-avoidance remark is not treated as a
proof of the finite-volume adapters above. Previously documented
source comparisons remain in force: the BP peripheral containment
qualification, the characteristic-translation correction in the flat
matching argument, the covering-selection and open/closed-disk margins,
and the distinction between oriented HG03 and MGL32. Exact versions,
locators, fresh versus reused readings, and bounded correction checks
are recorded in \code{reference_checks_revision120.md}.

\paragraph{Verdict and remaining work.}
After the two explicit hypothesis clarifications, this review found no
remaining mathematical gap in the written qualitative rigidity route
relative to the twelve inherited-interface families stated in
Section~\ref{sec:mostow-dependency-audit}. Those contracts cover the
model and covering geometry, elementary algebra, smooth approximation,
length and compactness arguments, Lebesgue analysis, and local volume
facts used here. Their accepted declarations, imports and axiom evidence
remain unbound. Thus the natural-language proof iteration and its
requested final working audit are complete at that boundary. This is
not a claim that HG04 has been formalized or that its foundations have
been independently certified.

No further Mostow--Prasad proof iteration is currently identified.
When an accepted PC release or explicitly stabilized interface is
available, bind these inherited contracts and implement the dependency
order already recorded. Independently, the next downstream blueprint
task is QRG00's radial core/end comparison for HG06; HG08 and ambient-flow
incompressibility also remain separate. New mathematical findings or
incompatible accepted bindings would reopen the affected proof segment.

\subsection{A controlled special case and the marking argument}
\label{sec:mostow-controlled-special-case}

\begin{lemma}[Product cusp extension of a supplied core map; MPR09]
\label{lem:mostow-cusp-product-extension}
Let $H,H'$ be connected hyperbolic manifolds given by compact smooth
cores $C,C'$ and finitely many standard cusp products with metrics
\[
 dt^2+e^{-2t}q_j,\qquad dt^2+e^{-2t}q'_j,\qquad t\geq0,
\]
where $q_j,q'_j$ are flat metrics on compact boundary sections.
Let $f_C:C\to C'$ be a smooth diffeomorphism matching the boundary
components. Assume the cusp coordinates extend to inner collars and,
on those collars, $f_C(x,t)=(b_j(x),t)$ for smooth boundary
diffeomorphisms $b_j$. Then $f_C$ extends by this same product formula
to a global smooth bilipschitz diffeomorphism $f:H\to H'$.
Its bilipschitz constant may depend on the supplied core map and metrics.
\end{lemma}
\begin{proof}
The collar agreement glues the formulas smoothly, including their
inverses. Compactness gives one $L\geq1$ bounding the differential
norms of all $b_j$ and $b_j^{-1}$ in their flat metrics. On a cusp,
\[
 \|d f(v,a)\|^2=a^2+e^{-2t}\|d b_j(v)\|_{q'_j}^2
 \leq L^2\bigl(a^2+e^{-2t}\|v\|_{q_j}^2\bigr),
\]
and the inverse has the same bound. Smoothness and compactness give
finite differential bounds for $f_C$ and its inverse on the cores.
Taking a maximum over this finite collection gives a global bound
$L_0\geq1$ for $df$ and $df^{-1}$. Integrate these bounds along every
piecewise smooth path and take the infimum of lengths. Applying the
same argument to the inverse gives
\[
 L_0^{-1}d_H(x,y)\leq d_{H'}(f(x),f(y))\leq L_0d_H(x,y).
\]
Connected Riemannian manifolds carry these path-length distances, so
this is the required global bilipschitz estimate.
\end{proof}
The flat sections need not be tori for this estimate; the oriented
three-dimensional HG02 application has tori. Collars, the actual core
diffeomorphism and its relation to a prescribed homotopy class are
hypotheses, not outputs of the norm calculation. A lift of this smooth
diffeomorphism is a diffeomorphism of the universal covers, and the
same differential/path-length argument makes it bilipschitz there.
Use the marking induced by that actual lift.

\begin{lemma}[Boundary marking, descent and homotopy; MPR10]
\label{lem:mostow-marked-descent}
Let $u:H\to H'$ be continuous, with continuous lift $U:X\to X$
equivariant for an isomorphism $\phi$. Suppose a $\phi$-equivariant
boundary homeomorphism $h$ is the boundary map of an isometry $A$.
Using faithfulness of the boundary action of $\operatorname{Isom}(X)$,
$A$ descends to an isometry $a:H\to H'$ homotopic to $u$.
\end{lemma}
\begin{proof}
The isometries $A\gamma$ and $\phi(\gamma)A$ have the same boundary
action, so faithfulness makes them equal. The inverse is equivariant
for $\phi^{-1}$, hence the two quotient maps are inverse local
isometries and preserve the Riemannian distances. For $s\in[0,1]$
join $U(x)$ to $A(x)$ by the unique hyperbolic geodesic segment and
take its parameter-$s$ point. This interpolation is continuous in
$(x,s)$ and commutes with all isometries. Consequently it is
$\phi$-equivariant and descends to a homotopy from $u$ to $a$.
There is no need for a uniform bound on $d(U(x),A(x))$.
\end{proof}
Boundary faithfulness follows, for example, because an isometry fixing
every boundary point preserves every geodesic; two intersecting
geodesics with different directions isolate each interior point.
Continuous dependence and uniqueness of hyperbolic segments are
inherited geometric APIs to bind. The conclusion is ordinary homotopy;
no proper-homotopy statement is being inferred.

\subsection{Source checks, iteration order and the quantitative consumer}
\label{sec:mostow-proof-frontier}

\paragraph{Consumer boundary; MPR11.}
KL Lemma 90.11 first constructs a diffeomorphism by its core/end
comparison, then applies rigidity, then proves quantitative closeness
by a separate compactness argument. The full MPR79 theorem now applies
directly to the resulting diffeomorphism, with its own induced marking;
MPR09's product-collar hypotheses are only needed for that optional
bilipschitz special case. Section~\ref{sec:hyp-quantitative-consumer}
writes the separate compactness step conditional on the topological input. HG06 retains its order
``fixed model, then accuracy, then approximation threshold'' and the
requirement that the target have at least as many cusps. Neither HG06
nor HG08 is closed by a proof of qualitative rigidity alone.

\paragraph{Exact evidence and version boundary.}
Schwartz's archived \code{mostow.tex} matches arXiv:2512.09774v5;
its labels \code{mmm}, \code{extend}, \code{MORSE}, \code{equ},
\code{one}, \code{star1}, \code{diff}, and \code{scope} identify the
main theorem, extension, normalization, orbit lemma, regularity,
zooming, analytic reduction and scope discussion. The pinned PDF has
29 pages, internally dated April 20, 2026. The main theorem is compact
and bilipschitz; Section 6, page 28, only remarks on the finite-volume
extension. BP D.3.12, printed151--152/PDF162--163, and D.3.14,
printed156/PDF167, supply the geometric cusp comparison. KL90.11 is
printed2813--2815/PDF227--229 in the archived version.
\code{reference_checks_revision105.md} records hashes, exact TeX lines,
reused correction checks and the limits of newly catalogued sources.

\paragraph{Next iterations.}
The subordinate register \code{mostow_prasad_proof_obligations.csv}
keeps MPR00--MPR79 attached to HG04. Its evidence status is now a
written proof relative to inherited interfaces; its formal status remains
unaccepted and unformalized.
Revisions106--107 supply the returning-zoom, boundary compactness,
disk and finite-packing arguments. Revision108 adds the project-specific
absolute-continuity and positive-measure regularity implication from
disk control, with inherited analysis inputs explicit. Revision109
supplies bilipschitz Morse tracking, compatible boundary extension
and triple-tube compactness. Revision110 supplies finite-volume
no-fixed-point/no-finite-orbit and no-invariant-plane arguments, the full
isometry centralizer, marked uniqueness, and the supplied-bilipschitz
comparison consumer. Do not restart these written arguments as open tasks.
Revision111 supplies the controlled quasi-isometry representative for
compact inputs and, conditional on matched geometric cusp data, for
cusped inputs. It fixes the inverse marking, prescribes inverse maps on
the ends, and bounds the two composition errors over compact quotient
cores. Revision114 uses MGL32 to supply the actual cusp matching and
full finite-volume controlled representative through MPR55--MPR61.
Revision115 supplies the general quasi-isometry boundary homeomorphism,
uniform normalization and boundary compactness with inverses through
MPR62--MPR68. Revision116 supplies uniform disk control and regularity
for this class through MPR69--MPR73, including cusp avoidance for every
normalized isometry composition. Revision117 supplies the returning-zoom
and isometry-orbit argument through MPR74--MPR78 and the full marked
existence and uniqueness conclusion in MPR79. The existing bilipschitz
statements retain their hypotheses. Revision118 audits the written
producer-to-consumer transfers and inherited interfaces, and corrects two
register edges without changing statements or proofs. Revision119
writes HG06's quantitative compactness step conditional on its
source-qualified topological comparison, preserving the parameter order.
Revision120 completes the final working statement-and-proof audit in
Section~\ref{sec:mostow-final-audit}; downstream radial and flow work
does not add a prerequisite to that theorem. The
individual core/cusp and compact-lift producers are now written in
MGL00/MGL32 and MGL30. Isometry properness, compactification-action
continuity and the inherited interfaces still require accepted bindings;
all new Lean implementations remain open.
\code{reference_checks_revision120.md} records the final working audit;
revision119 records the quantitative consumer;
revision118 records the dependency audit;
revision117 records returning zooms and the full marked assembly;
revision116 records disk control and regularity,
revision115 the general boundary comparison, revision114 peripheral and affine matching,
revision111 cusp prescription, and revisions106--110 the preceding evidence.

Francaviglia--Dunfield remains the checked alternative if this route's
remaining scope adapters prove more expensive. Selecting a route is an
engineering judgment, not a measured minimum of Lean effort. The
Magalh\~aes book is an available analytic supplement, not an adopted
theorem dependency; lack of Lehto--Virtanen is not a current blocker.
Every future revision should close named obligations or correct their
contracts, preserving source evidence and recording affected consumers.


\section{Quantitative rigidity and persistent thick models}
\label{sec:hyperbolic-persistence-contracts}

\begin{definition}[Smooth approximation data; HG05]
\label{def:hyp-smooth-approximation}
A quantitative approximation from a pointed model $(H,x)$ to $(H',x')$
specifies an open domain containing $B_h(x,R)$, a smooth embedding $f$
on that domain, $f(x)=x'$, an integer derivative order $k$, and
$\|f^*h'-h\|_{C^k(B_h(x,R),h)}<\delta$. The norm and its reference metric
are recorded. A comparison of maps on a smaller ball is a separate
distance bound in the target metric.
\end{definition}
For the source's $\delta$-approximation use $R=\delta^{-1}$ and a
specified integer order at least $\lceil\delta^{-1}\rceil$. This is a
sufficient strengthened finite-order convention for KL90.6--90.8; a
formal adapter must document the source's rounding and norm convention.
Nothing is required outside the recorded domain.

\begin{foundation}[Cusp-sensitive quantitative rigidity; HG06]
\label{found:hyp-quantitative-rigidity}
For each fixed pointed HG02 model $(H,x)$ and every $\eta>0$, there is
$\delta>0$ such that a HG05 $\delta$-approximation into a connected
complete finite-volume hyperbolic $H'$ with at least as many cusps as $H$
admits an isometry $I:H\to H'$ satisfying
\[
       \sup_{p\in B_h(x,\eta^{-1})}d_{h'}(I(p),f(p))<\eta.
\]
Choose $\delta$ small enough that this comparison ball is in the domain.
The isometry is not required to send $x$ exactly to $x'$.
\end{foundation}
KL90.11, printed2813--2815/PDF227--229, first compares compact cores and
cusp/tube ends, then uses HG04, then compactness of the maps. The target
component is connected throughout. Hyperbolic Dehn fillings explain why
the cusp-count condition cannot be removed. Neither constants uniform
over all reference manifolds nor prescribed boundary parametrizations
follow from this statement. MSM206's almost-isometry theorem gives a
stronger differentiable comparison on truncated cores; its parameter
dependence and nonsingular-flow hypotheses remain a separate comparison.

\subsection{The quantitative consumer after qualitative rigidity}
\label{sec:hyp-quantitative-consumer}

Revision119 separates two steps in KL90.11. The first produces a global
diffeomorphism using a comparison of compact cores and radial ends. The
second upgrades existence of some isometry to proximity to the actual
approximation. The following three written arguments supply the second
step, conditional on the first. They are downstream of MPR79 and add no
prerequisite to qualitative Mostow--Prasad rigidity. HG06 retains its
source-qualified status until the first step is expanded and checked.

\paragraph{The exact remaining topological input; QRG00.}
For the fixed pointed HG02 model $(H,x)$, supply a number $\delta_0>0$
such that, for every $0<\delta\leq\delta_0$, every HG05
$\delta$-approximation into a connected complete finite-volume
hyperbolic $H'$ with at least as many cusps as $H$ implies the existence
of a homotopy equivalence $u:H\to H'$. The curvature normalization is
the same on both manifolds. A global diffeomorphism, as constructed in
KL90.11, is sufficient. Neither a metric bound on $u$ nor agreement
between $u$ and the approximation on a specified collar is required
by this input. It is uniform over the target and the approximation,
after $(H,x)$ is fixed. QRG00 is a source-qualified producer requirement,
not a theorem proved in this revision.

\begin{theorem}[Quantitative rigidity from the topological input; QRG03]
\label{thm:hyp-quantitative-from-topology}
Fix a pointed HG02 model $(H,x)$ and assume QRG00. Then HG06 holds:
for every $\eta>0$ there is $\delta>0$, depending on this fixed model
and $\eta$, such that every approximation in HG06 admits an isometry
$I:H\to H'$ with
\[
 \sup_{p\in B_h(x,\eta^{-1})}d_{h'}(I(p),f(p))<\eta.
\]
The conclusion does not prescribe $I(x)=f(x)$ or the homotopy class
of an extension of $f$.
\end{theorem}
\begin{proof}
If the conclusion fails for an accuracy $\eta>0$, choose counterexamples
$f_i:H\supset U_i\to H_i'$ with approximation parameters
\[
 0<\delta_i\leq\min\{\delta_0/2,\eta/2,1/i\},
 \qquad \delta_i\longrightarrow0.
\]
Their domains exhaust every compact subset of $H$. QRG00 gives a homotopy
equivalence $u_i:H\to H_i'$. MPR79, with HG01 if needed for a common
curvature rescaling, gives a global isometry in this homotopy class.
Choose its inverse $J_i:H_i'\to H$ and put $F_i=J_i\circ f_i$.
These are smooth embeddings on $U_i$, and on every compact set
$F_i^*h=f_i^*h_i'\to h$ in $C^0$. The fixed-target compactness lemma
QRG02 below gives a subsequence converging uniformly on compact sets
to a global self-isometry $F$ of $H$. Then
\[
 I_i=J_i^{-1}\circ F,\qquad
 d_{h_i'}(I_i(p),f_i(p))=d_h(F(p),F_i(p)).
\]
The closed ball $\overline B_h(x,\eta^{-1})$ is compact. Eventually
the supremum on this closed ball is less than $\eta/2$, contradicting
the asserted failure on the open ball. Arbitrary choices of the
isometries $J_i$ are allowed: QRG02 prevents the corresponding
$F_i(x)$ from escaping into a cusp. The cusp-count hypothesis is used
in QRG00; compactness is applied only after all targets have been
identified isometrically with the fixed $H$.
\end{proof}

\begin{lemma}[Compactness of almost-isometries into a fixed model; QRG02]
\label{lem:hyp-fixed-target-compactness}
Let $(H,h)$ be a connected complete finite-volume hyperbolic
three-manifold. Let $U_i\subset H$ be open sets such that every compact
subset of $H$ is contained in all sufficiently late $U_i$, and let
$F_i:U_i\to H$ be smooth embeddings. Suppose
$F_i^*h\to h$ in $C^0$ on every compact set, where the tensor norm is
taken with respect to $h$. Then a subsequence converges uniformly on
compact subsets to a global isometry $F:H\to H$.
No basepoint condition or orientability assumption is needed.
\end{lemma}
\begin{proof}
Fix $x\in H$ and $r>0$ such that $\overline B_h(x,r)$ is contained
in a normal coordinate ball and $B_h(x,r)$ is contractible. For large
$i$ this closed ball lies in $U_i$, and on it
\[
             \tfrac14 h\leq F_i^*h\leq4h.
\]
QRG01 below, with $a=1/2$, gives
$\operatorname{inj}_H(F_i(x))\geq r/4$. These points lie in a compact
subset $K$ of this fixed finite-volume $H$. Indeed MGL32 supplies a
compact core and finitely many standard cusps. In a cusp, a nontrivial
translation loop on its compact flat section has length bounded above
by $C e^{-\kappa t}$ at depth $t$, uniformly over that section.
Here the common curvature is $-\kappa^2$. The full stabilizer has a
rank-two translation lattice, also in the nonorientable case. The
injectivity-radius displacement formula
\eqref{margulis:inj-displacement} therefore makes injectivity radius
tend uniformly to zero down each cusp. A fixed positive thick set is
closed and lies in a compact truncation. If there are no cusps, $H$
itself is compact.

For each fixed $R>0$ and sufficiently large $i$, the tensor bound
$F_i^*h\leq4h$ holds on $\overline B_h(x,3R)$. Minimizing geodesics
between points of $\overline B_h(x,R)$ stay in this larger ball, so
$F_i$ is $2$-Lipschitz on $\overline B_h(x,R)$. Its image lies in the
closed $2R$-neighborhood of $K$, which is compact by Hopf--Rinow.
Arzela--Ascoli on the closed integer-radius balls, followed by a
diagonal subsequence, gives a continuous map $F:H\to H$ and locally
uniform convergence. This uses only compact-domain metric compactness,
not a smooth compactness theorem for variable manifolds.

We check local preservation of distance, including the possible target
shortcuts. Fix $p\in H$ and $R>0$. On the compact
$\overline B_h(p,6R)$ choose bounds
\[
        a_i^2h\leq F_i^*h\leq b_i^2h,
        \qquad a_i,b_i\longrightarrow1,\quad a_i>0.
\]
For $q,z\in B_h(p,R)$, a minimizing source geodesic stays in
$B_h(p,3R)$, giving
$d_h(F_i(q),F_i(z))\leq b_i d_h(q,z)\leq2b_iR$.
For late $i$, $2b_i<3a_i$. Choose target paths from $F_i(q)$ to
$F_i(z)$ whose lengths decrease to this distance and are less than
$3a_iR$. QRG01, applied to $\overline B_h(q,4R)\subset B_h(p,6R)$,
lifts each path into that ball. Since $F_i$ is injective, its lifted
endpoint is $z$. The lower tensor bound gives
\[
 a_i d_h(q,z)\leq d_h(F_i(q),F_i(z))\leq b_i d_h(q,z).
\]
Passing to the limit proves that $F$ preserves distance on a
neighborhood of every $p$.

Such a map between equal-dimensional Riemannian manifolds is a smooth
local isometry. One may bind the inherited local distance-isometry
regularity statement here; in hyperbolic normal balls it also follows
from uniqueness of short geodesics and the hyperbolic cosine law,
which identifies their initial inner products and hence a linear
orthogonal map between tangent spaces. The expression in exponential
coordinates is then that orthogonal map. Completeness of the source
makes this local isometry a surjective covering: lift a finite-length
target geodesic from an image point, use completeness to continue the
lift at every finite endpoint, and use normal balls and uniqueness of
lifts to obtain evenly covered neighborhoods.

Finally this self-cover has one sheet. If two distinct preimages of
one point existed, covering transport along paths would give at least
two over every point. A countable partition of $H$ into Borel sets
subordinate to evenly covered coordinate balls, with the density
change-of-variables formula on two inverse branches, would give
$\Vol_h(H)\geq2\Vol_h(H)$, impossible because
$0<\Vol_h(H)<\infty$. Thus $F$ is a bijective local isometry and a
global isometry. This argument does not infer global injectivity merely
from the injectivity of the approximating maps.
\end{proof}

\begin{lemma}[Short inverse paths and a thick basepoint; QRG01]
\label{lem:hyp-short-inverse-paths}
Let $(M,g)$ and $(N,h)$ be connected Riemannian manifolds of the same
dimension. Let $f:U\to N$ be a smooth embedding of an open set, with
$\overline B_g(p,r)\subset U$ compact. Suppose $r>0$, $a>0$ and
$f^*h\geq a^2g$ on this closed ball. Every piecewise smooth path in
$N$ starting at $f(p)$ and having total length less than $ar$ lies in
$f(B_g(p,r))$ and lifts there by $f^{-1}$. If $B_g(p,r)$ is
contractible and $N$ is complete hyperbolic, then
\[
                   \operatorname{inj}_N(f(p))\geq ar/2.
\]
\end{lemma}
\begin{proof}
The equal-dimensional embedding is a diffeomorphism onto an open
subset. Start the inverse path at $p$. While it is in $B_g(p,r)$,
the tensor bound makes its length at most $a^{-1}$ times the target
length. If the total target length is $\ell<ar$, the inverse path
stays in the compact set $\overline B_g(p,\ell/a)\subset B_g(p,r)$.
Its endpoint at a putative first failure exists by the same length
bound on tails, and the local inverse there continues it. Thus no
failure occurs. A loop lifts to a loop because $f$ is injective.
When the ball is contractible, every target loop of length less than
$ar$ based at $f(p)$ is consequently null-homotopic in $N$.
For a complete hyperbolic quotient, a nonidentity deck transformation
with displacement less than $ar$ would project its geodesic segment
to a nontrivial loop of that length. The injectivity-radius formula
therefore gives the asserted lower bound. This proves only the
one-sided thick-basepoint estimate needed by QRG02, not the full
two-sided injectivity-radius comparison used in KL's radial argument.
\end{proof}

\paragraph{Source comparison and the remaining radial construction.}
KL90.6--90.9, printed2812--2813/PDF226--227, defines the partially
defined smooth approximations and the distance comparison between maps.
KL90.11, printed2813--2815/PDF227--229, treats orientable models at
curvature $-1/4$. Its first step chooses a source threshold $\mu_1$
with no thin tubes, then long compact cusp slabs and a truncation $W$.
The Busemann coordinate $\rho$ decreases down the cusp; the slabs are
$[-3L,-L]$, the boundary is $\rho=-2L$, and the stated principal
curvature is $-1/2$. These signs and units differ from the increasing
depth coordinate in MGL32; rescaling to curvature $-1$ changes lengths
by $1/2$ and injectivity-radius thresholds by the same factor.

To discharge QRG00 by this source route, still expand: the injectivity
comparison on buffered domains; exclusion of nearby target tube cores;
the long, low-curvature radial-curve estimate; transversality of each
image torus; its isotopy to the correct thin boundary; the simultaneous
core identification with distinct boundary components; and the final
cusp count. The last step uses a target thick core with exactly as many
boundary components as source cusps and at least that many target cusps,
so every complementary thin component is a cusp. An isotopy of one
torus alone is not the global core conclusion. Smooth extension across
the ends must supply an actual global diffeomorphism (a homotopy
equivalence suffices for QRG00), with no disconnected target component
silently ignored. These requirements retain the source proof; this
revision does not claim to have proved them.

MPR79 accepts the resulting arbitrary homotopy equivalence directly.
MPR09 remains a valid optional bilipschitz special case, but its exact
product collars are not an additional prerequisite for this application
of the full theorem. QRG03 uses no marking of the approximation outside
its given domain, and no harmonic-map or differentiable closeness
conclusion. Its compactness proof is an explicit project expansion of
KL's final paragraph, using the established Margulis cusp input. It is
not a verbatim theorem from the source. The inherited interfaces used
are path-length calculus, compactness, the injectivity-radius formula,
local distance-isometry regularity, complete local-isometry covering,
and integration with densities; accepted bindings and Lean proofs remain
pending. No new inspection of the changing PC snapshot is made.

\paragraph{Qualitative audit and separate downstream frontier.}
Revision120 completes the final working audit of the Mostow--Prasad
statement and proof in Section~\ref{sec:mostow-final-audit}. Accepted
inherited bindings and Lean implementations remain pending. The next
independent HG06 task is QRG00's radial construction and global core/end
comparison. It does not add a prerequisite to qualitative rigidity.
Successful document/static checks remain distinct from mathematical
review and from Lean verification.


% BEGIN REVISION160 QRG DEVELOPMENT
\subsection{The actual global core and end comparison}
\label{sec:hyp-actual-core-end-comparison}

Revision160 supplies QRG00 by expanding KL90.11's geometric route.
The earlier paragraph identifying it as open records revision119's state.
The new arguments use the established Margulis geometry and elementary
smooth and covering foundations; they do not use persistent flow maps,
ambient cusp incompressibility, or a theorem about embedded tori in an
arbitrary three-manifold. Fix curvature $-1/4$ for the construction;
HG01 transfers the result to any fixed common negative normalization.
Write $s$ for cusp depth, increasing toward infinity, so that
\[
 h=ds^2+e^{-s}q,\qquad
 \nabla^2s=-\tfrac12(h-ds\otimes ds).
\]
Thus $s=-\rho$ in KL's convention. All constants below are chosen
after the source $(H,x)$ and before the target or the approximation.

\begin{lemma}[Two short-curve obstructions; QRG04]
\label{lem:hyp-short-curve-obstructions}
A positive-length smooth regular curve in $\mathbb H^3_{-1/4}$ whose
unit-speed geodesic curvature is everywhere less than $1/2$ cannot
have equal endpoints. A smooth closed regular curve in a standard
hyperbolic cusp cannot have geodesic curvature everywhere less than
$1/2$. The first assertion permits a corner between the terminal
and initial velocities; the second concerns a smooth closed curve.
\end{lemma}
\begin{proof}
For the first assertion use distance $r$ from the common endpoint.
It has a positive interior maximum on the curve, away from $r=0$.
In hyperbolic space
\[
 \nabla^2r=\tfrac12\coth(r/2)(h-dr\otimes dr).
\]
At that maximum the second derivative along a unit-speed curve is
at least $\tfrac12\coth(r/2)-|\nabla_{\dot\gamma}\dot\gamma|>0$,
a contradiction. The possible corner is at the endpoints, not at
this interior maximum. For the second assertion take a minimum of
the periodic function $s\circ\gamma$. Its first derivative is zero,
and the cusp Hessian formula gives
\[
 (s\circ\gamma)''=-\tfrac12+
 ds(\nabla_{\dot\gamma}\dot\gamma)<0,
\]
contradicting the second derivative test at a minimum. The formulas
follow from the hyperbolic polar metric and the displayed warped
cusp metric, respectively, using the Koszul formula.
\end{proof}

\begin{lemma}[Buffered two-sided injectivity comparison; QRG05]
\label{lem:hyp-buffered-injectivity-comparison}
Let $f:U\to H'$ be a smooth embedding between complete hyperbolic
three-manifolds. Fix $p\in U$ and write $0<i=\operatorname{inj}_H(p)<\infty$.
Suppose $\overline B_h(p,R)\subset U$ is compact, $R>2i$, and
$a^2h\leq f^*h'\leq b^2h$ on this ball. Suppose also that the
image of an $h$-geodesic, after unit-speed reparametrization, has
$h'$-geodesic curvature less than $1/2$ there. Then
\[
 a\operatorname{inj}_H(p)\leq\operatorname{inj}_{H'}(f(p))
 \leq b\operatorname{inj}_H(p).
\]
On a fixed compact buffered domain these hypotheses hold with
$a,b\to1$ when $f^*h'\to h$ in $C^1$; the choices are uniform
over $H'$ and $f$.
\end{lemma}
\begin{proof}
If a nontrivial target geodesic loop based at $f(p)$ had length
$\ell<2ai$, QRG01 would lift the whole loop into $B_h(p,R)$.
Its endpoints coincide by injectivity of $f$. The lifted loop has
length at most $\ell/a<2i$. Each of its points is within half this
length of $p$, using the shorter of the two loop arcs. It therefore
lies in the normal ball $B_h(p,i)$ and contracts in that ball.
Its image would contract in $H'$, a contradiction. The deck
displacement formula gives the lower inequality, with equality
allowed by taking strict counterexamples.

A shortest nontrivial geodesic loop at $p$ has length $2i$ and lies
in $\overline B_h(p,i)$. Existence follows from discreteness of the
deck orbit and compactness of closed balls in the universal cover.
If its image contracted in $H'$, that image would lift to a regular
curve in $\mathbb H^3$ with equal endpoints and curvature less
than $1/2$, contrary to QRG04. The image loop is thus nontrivial
and has length at most $2bi$, proving the upper bound. No
injectivity of $f_*$ is assumed.

Put $\widetilde h=f^*h'$. The connection difference is
\[
 2\widetilde h((\nabla^{\widetilde h}-\nabla^h)_XY,Z)
 = (\nabla^h_X\widetilde h)(Y,Z)
  +(\nabla^h_Y\widetilde h)(X,Z)
  -(\nabla^h_Z\widetilde h)(X,Y).
\]
Small $C^1$ error and uniform positive definiteness make this tensor
small. Reparametrizing a geodesic changes its acceleration only
by the controlled tangential term. This proves the final assertion
and also controls inverse images of target unit-speed geodesics.
The argument uses the pulled-back metric on a fixed source domain;
it requires no target injectivity bound to obtain these estimates.
\end{proof}

\begin{lemma}[Deep cusp slabs exclude axes and force transversality; QRG06]
\label{lem:hyp-actual-radial-transversality}
Fix disjoint standard cusp charts $T_j\times[0,\infty)$ of $H$
and a compact core $C_0$ with these ends. Choose a target injectivity
threshold $\mu>0$ with $2\mu$ below the Margulis displacement
constant and $\mu<1/100$. For sufficiently large source depth $S$
and sufficiently small $C^1$ metric error on a fixed buffered
compact neighborhood, every image torus
\[
 Z_{j,u}'=f(T_j\times\{u\}),\qquad |u-S|\leq2,
\]
lies in a single target $\mu$-thin component, is at distance greater
than $2$ from its boundary and its tube axis if present, and is
transverse to the radial direction. These choices are uniform over
the connected oriented complete finite-volume target and $f$.
Each full band belongs to one thin component, though different
source bands have not yet been shown to belong to different ones.
\end{lemma}
\begin{proof}
Choose $S>50$ so that every source cusp slice at depths at least
$S-30$ has intrinsic diameter less than $1/100$ and ambient
injectivity radius less than $\mu/4$. MGL07 and a nontrivial
translation loop give these uniform bounds for the finite cusp
family. On the compact neighborhood choose the approximation
small enough that
\[
 a=0.99,\quad b=1.01,\quad
 |\nabla^h-\nabla^{\widetilde h}|_h<1/100,
\]
the QRG05 image-curve bound holds, and inverse images of target
geodesics have $h$-unit-speed curvature less than $1/10$.
Enlarge the fixed domain beforehand to contain the QRG05 buffers
at every point of the bands and their closed $20$-neighborhoods.
For example, first take that compact neighborhood, bound its
injectivity radius above by a finite $I$, and add its closed
$(2I+2)$-neighborhood. Hopf--Rinow makes all these choices compact.

Depth is $1$-Lipschitz until a curve exits its cusp through $s=0$.
Consequently every source ball of radius $20$ centered on these
bands lies in its one cusp at depth at least $S-22$.
QRG05 puts its image in the $\mu$-thin set. QRG01 lifts every
target path of length less than $10$ from a band point into this
source ball. In particular the entire target ball of radius $10$
lies in the same thin component and in this cusp image. This
proves the boundary-distance assertion with a substantial buffer.

If a target tube axis came within distance $2$ of a band point,
join to it by a path of length less than $3$. Its closed core
geodesic has length less than $2\mu$. Every point of this loop
is then reached by a path of length less than $4$ from the band
point, so the entire core loop lifts by $f^{-1}$ to this one cusp.
It is a smooth closed regular curve with curvature less than
$1/10$, contradicting QRG04. This excludes the nearby axis;
it does not yet exclude a tube as the ambient thin component.

At $z'=f(p)$ take the unit-speed radial geodesic
$\gamma':[ -\tau,\tau]\to H'$, with $\tau=1/8$ and
$\gamma'(0)=z'$. The preceding buffers keep it away from the
axis and the boundary. It is minimizing in the full target:
within the component the radial coordinate has gradient of norm
one and changes by $2\tau$; a competing path of length less than
$2\tau$ stays in the same component by the boundary-distance
buffer. Such a path therefore cannot be shorter.

Let $\gamma=f^{-1}\gamma'$ with the same parameter and let
$v=(s\circ\gamma)'(0)$. The metric and connection estimates give
$|\dot\gamma|_h\leq1/a$ and
$|(s\circ\gamma)''|<1$. If $|v|\leq1/2$, then
\[
 |s(\gamma(\tau))-s(\gamma(-\tau))|
 \leq\tau+\tau^2.
\]
Join these endpoints in the source by a vertical segment and a
path in one cusp slice, with length arbitrarily close to
$\tau+\tau^2+1/100$. The path stays in the buffered cusp domain.
Its image has length at most, in the limiting estimate,
\[
 1.01(1/8+1/64+1/100)<1/4=2\tau,
\]
contradicting minimization of $\gamma'$. Thus $|v|>1/2$.
In particular the inverse radial vector is not tangent to the
source torus, which is precisely the required transversality.
This short-path argument supplies the step compressed into the
long radial-curve sentence in KL90.11.
\end{proof}

\begin{lemma}[One radial sheet and the actual inward side; QRG07]
\label{lem:hyp-radial-one-sheet-side}
Under QRG06, each $Z_{j,u}'$ is a single radial graph over the
boundary torus of its thin component. With $S$ chosen still larger
if necessary and the same small-error bounds, the image of the
source side $s<S$ lies locally on the thickward side of $Z_{j,S}'$.
\end{lemma}
\begin{proof}
Away from the axis a target thin component is a product with
radial interval and compact connected torus cross-section.
Projection of $Z_{j,u}'$ to this torus is a local diffeomorphism
by QRG06. It is proper since the source torus is compact. Its
image is open and closed, hence the whole cross-section, and
it is a finite covering. The distinct points above any point
of the cross-section have distinct real radial heights because
the original torus is embedded. Sorting those heights gives
global smooth sheets: locally their order cannot change without
two points coinciding. A covering with more than one such sheet
would have disconnected total space. Thus the degree is one.
This is the one-sheet argument; transversality alone only gave
a covering.

Apply this to every $u\in[S-2,S+2]$. The radial graph heights
depend smoothly on $(\theta,u)$, by the inverse function theorem
for the projection together with the parameter. Their $u$
derivative is nonzero because $f$ is a local diffeomorphism.
Its sign is constant on this connected band. To determine that
sign use injectivity radius, without assuming that $f$ preserves
a radial coordinate. Sufficiently far down a source cusp,
MGL01 and the parabolic displacement formula give
\[
 \operatorname{inj}_H(\theta,u)
   =2\operatorname{arsinh}(c_j e^{-u/2}),\qquad c_j>0.
\]
The constant is determined by the shortest translation in the
fixed lattice, and the value is independent of $\theta$.
Increase $S$ so that
\[
 a\operatorname{inj}_H(T_j\times\{S-2\})
 > b\operatorname{inj}_H(T_j\times\{S+2\})
\]
for every $j$. QRG05 makes every target injectivity radius on
the shallower graph larger than every such radius on the deeper
graph. In a cusp, injectivity radius increases strictly in the
thickward radial direction. The same holds in a tube away from
its axis: the displacement formula for a loxodromic power is
\[
 \sinh^2(d_n(r)/4)
 =\sinh^2(n\ell/4)\cosh^2(r/2)
   +\sin^2(n\vartheta/2)\sinh^2(r/2),
\]
where $\ell>0$ is the generator's translation length. It is
strictly increasing in $r>0$, and the minimum over powers is
also strictly increasing: use a minimizing power at the larger
radius to compare with the minimum at the smaller one.
Hence increasing $u$ moves each graph strictly thinward.
The band fills the region between its outer graphs, by the
one-sheet coordinates and their nonzero $u$ derivative. The
side $s<S$ therefore corresponds to the thickward side, as
asserted. This fixes a side, not just an unoriented isotopy class.
\end{proof}

\begin{lemma}[All target boundary components and no hidden tubes; QRG08]
\label{lem:hyp-global-boundary-identification}
Let $W$ be the source core obtained by truncating every cusp at
depth $S$. Choose $\mu$ additionally so small that
$\mu<a\min_{C_0}\operatorname{inj}_H/4$.
With the fixed-domain choices of QRG05--QRG07 made also over $W$,
the target $\mu$-thick core is contained in $f(W)$.
Each target thin component contains exactly one component of
$f(\partial W)$, and $f(W)$ consists of the thick core and the
regions between these radial graphs and the thin boundaries.
In particular the number of target thin components equals the
number of source cusps. This conclusion does not yet use a
target cusp-count hypothesis.
\end{lemma}
\begin{proof}
The target thick core is a nonempty compact connected smooth
manifold with boundary. Compactness, finiteness of thin
components, and their product or solid-torus form are MGL01--MGL04.
For connectedness, join two points by a smooth path transverse
to the finitely many thin boundaries, and replace each excursion
through a thin component by a path in its connected boundary
torus. A small perturbation into the thick side gives a path in
the thick core. Nonemptiness is MGL03. The strict-threshold
convention produces no degenerate tube at equality.

The boundary of $f(W)$ lies strictly inside the thin set.
Membership in $f(W)$ is consequently constant on the connected
thick core. A point of $C_0$ has image injectivity radius at least
$a\min_{C_0}\operatorname{inj}_H>\mu$ by QRG05, so that core
meets $f(W)$ and is wholly contained in its interior.

There cannot be two image boundary tori in one thin component.
By QRG07 they are disjoint radial graphs and $f(W)$ lies locally
on the thickward side of each. Order all graphs in that component
and take two consecutive ones. The connected open region between
them has no boundary point of $f(W)$, so membership is constant
there. At its thinward boundary it must be inside $f(W)$; at its
thickward boundary it must be outside. This contradiction proves
uniqueness, including the possibility of nested source boundaries.

Suppose a target thin component contains no image boundary.
Since its boundary meets the thick core, the whole component is
inside $f(W)$. A cusp cannot be so contained, since $f(W)$ is
compact and the cusp is unbounded in the complete target.
If the component is a tube, its core geodesic $\sigma$ is wholly
inside $f(W)$. At every point of $\sigma$ the target injectivity
radius is less than $\mu$, so QRG05 gives source injectivity
radius less than $\mu/a<\min_{C_0}\operatorname{inj}_H$ along
$f^{-1}\sigma$. Thus the whole connected inverse loop lies
outside $C_0$, in a single source cusp. The inverse curve is
smooth and has geodesic curvature less than $1/10$. QRG04 gives
a contradiction. This excludes a tube hidden entirely inside
the image core; the local axis exclusion of QRG06 alone would
not have done so.

Every target thin component therefore has exactly one image
boundary graph. Its thickward region belongs to $f(W)$, since
it joins the thick core and has no further image boundary. Its
thinward region is outside $f(W)$ by QRG07 and connectedness.
These descriptions give the asserted equality of subsets, with
all boundary components and their labels. Radial collars show
that adjoining these graph-bounded product bands does not change
the diffeomorphism type of the thick core.
\end{proof}

\begin{theorem}[Actual QRG00 producer and quantitative rigidity; QRG09]
\label{thm:hyp-actual-qrg-producer}
For every fixed pointed HG02 source $(H,x)$ there is $\delta_0>0$
such that every HG05 $\delta$-approximation, $0<\delta\leq\delta_0$,
into a connected complete finite-volume hyperbolic target $H'$
with at least as many cusps as $H$ yields a diffeomorphism
$H\to H'$. The choice is uniform over $H'$ and the approximation.
No orientability hypothesis on the target needs to be added.
Consequently QRG00 and HG06 have written producers under the
retained foundations: for HG06 apply QRG03 to this conclusion.
\end{theorem}
\begin{proof}
If $H$ is closed, take the approximation radius large enough to
contain all of $H$. Its image is open by the inverse function
theorem and closed by compactness. Connectedness of $H'$ makes
the approximation a global diffeomorphism.

Otherwise first suppose $H'$ is oriented. Take $C_0$, its finitely
many cusp charts, $\mu$, $S$, and all compact buffers in the order
used in QRG05--QRG08. They involve only $H$. Enclose their union
in a fixed ball about $x$ and choose $\delta_0$ so that the HG05
domain contains a neighborhood of that closed ball, its derivative
order is at least two, and its metric error supplies all the stated
strict connection and tensor bounds. QRG08 applies to every such
target and approximation. If $k$ is the source cusp count, the
target has exactly $k$ thin components. Every target cusp is one
of these components, by the complete finite-volume Margulis
description. Having at least $k$ cusps forces all $k$ components
to be cusps: no target tube remains.

The actual image $f(W)$ is now a compact core of $H'$ with one
outgoing product cusp at each boundary graph. The single radial
graph description supplies these outgoing collars explicitly.
Compress each attached infinite collar smoothly into an inner
collar to identify $H$ with $\operatorname{int}W$ and $H'$ with
$\operatorname{int}f(W)$, as in MGL05. Composing those
diffeomorphisms with $f|_{\operatorname{int}W}$ gives the required
global diffeomorphism. It need not agree with $f$ on a prescribed
collar; QRG00 requires no such marking. This construction accounts
for the entire connected target, not just the component met by
one selected torus.

For a nonorientable target, pass temporarily to its connected
orientation double cover $\pi:H'^+\to H'$. An embedding of the
oriented connected source neighborhood has a lift $f^+$: push
forward its orientation to choose a section of the orientation
cover over its open image. This lift is an embedding with exactly
the same pulled-back metric. The cover is complete, hyperbolic
and of finite volume. Apply QRG08 to $f^+$; cusp counting is not
needed at this point. Its compact image contains the entire
$\mu$-thick core of $H'^+$. The nonidentity deck transformation
$\tau$ preserves that nonempty core, whereas
$f^+(W)$ and $\tau f^+(W)$ are disjoint: intersection would
contradict injectivity of $f$ and freeness of the orientation
cover. Both sets containing the core is impossible. Thus such
a sufficiently accurate embedding cannot have nonorientable
target. The already treated oriented case applies.

The finite choices above give one positive $\delta_0$, independent
of the target and $f$. HG01 deals with a different fixed common
curvature normalization, including every buffer and injectivity
threshold. A diffeomorphism is a homotopy equivalence, so it
supplies exactly QRG00. QRG03 then uses MPR79 and QRG02 to obtain
the required proximity to the given approximation.
\end{proof}

\paragraph{Source, scope, and the next analytic construction.}
The source route is KL90.11, printed2813--2815/PDF227--229 in
the archived February20,2013 version. QRG04--QRG09 are explicit
project expansions, including the short radial-path alternative,
the hidden-tube exclusion, the ordering proof of one sheet, and
the orientation-cover adapter. The source's fixed target-free
threshold and cusp count are retained. MGL01 uses displacement
$2\operatorname{inj}$; the present $\mu$ is an injectivity
threshold, which is why $2\mu$ must lie below its constant.
Hamilton1999 Sections8--10, printed707--716/PDF13--22, and
MSM206's \code{Nonsingular-ExistAlmostIsomNearIsom} supply the
comparison route, not an additional imported proof of these steps.
Exact readings, versions, and correction checks are recorded in
\code{reference_checks_revision160.md}.

No slow time dependence, simultaneous persistent family, or flow
derivative estimate follows from QRG09. HG07--HG08 and LP03/LP06
remain the next analytic producers. The larger analytic goal and
its requested final mathematical audit remain open at this stage.
No accepted-PC declaration binding or Lean proof is asserted.
% END REVISION160 QRG DEVELOPMENT

\begin{foundation}[Late smooth-limit and unscathed-window input; HG07]
\label{found:hyp-late-window}
For each fixed thickness threshold, supply the hypotheses and conclusions
of KL88.1 and 90.14 on the actual normalized surgery slices: complete
finite-volume hyperbolic pointed limits, a bounded total normalized volume,
and approximation maps which continue along unscathed windows with their
stated domain and error controls. The common scales and derivative orders
are explicit. These conclusions are not consequences of metric GH
convergence alone.
\end{foundation}
Chapter10 owns the inherited smooth compactness and estimates; Chapter12
binds their flow hypotheses. No new read of the moving PC snapshot is
needed to specify this interface.

% BEGIN REVISION163 THICK FLOW PRODUCTION
\subsection{Producing the late smooth limits and unscathed windows}
\label{sec:hyp-actual-longtime-flow-input}

This subsection supplies HG07 from the retained flow estimates. It
uses WBD03 only at order zero; that result uses KL92.13, KL70 and
KL81--86, not HG07 or persistent hyperbolic embeddings. Thus the
forward reference to Chapter13 is not a circular persistence argument.
The flow, its admissible surgery profile and its finite canonical
accuracy are fixed as in Section~\ref{sec:collapse-actual-ambient-whole-ball}.
Local Shi estimates on unscathed thick neighborhoods will provide
all higher orders here; they do not assert all-order thin estimates
for that fixed canonical accuracy.

\paragraph{Strict canonical-scale convention.}
KL69.1 takes its threshold scale in $(0,1)$. In WBD02's normalized
application of KL70.2, choose the fixed scale $1/2$ if this strict
range is retained literally, enlarge $D$ to exceed $4$, and use
$K(2\sqrt D,1/2)$ instead of $K(2\sqrt D,1)$. Since
$s/r_{\rm can}(t)\to0$, the stronger threshold assumption also
holds on the entire earlier rescaled flow. The proof and its
uniform conclusion are unchanged. This explicit endpoint adapter
avoids treating the numeral $1$ as an admissible value in the
source's open interval.

\paragraph{Branch convention.}
Use the nonextinct branch with $R_{\min}(t)<0$ at all sufficiently
late times, as in KL87. The elementary nonnegative-scalar alternative
is kept separate. If a regular slice is scalar-nonnegative, its
finitely many connected components are either flat or acquire strictly
positive scalar curvature on a common short smooth interval. Indeed,
$\partial_tR=\Delta R+2|\Ric|^2$ and the strong maximum principle give
this assertion, including the case where $R$ initially vanishes
identically but $\Ric$ does not. In dimension three $\Ric=0$ means
flatness. Flat components remain flat by uniqueness. The positive
components have a common positive scalar minimum $c$ after that
interval; scalar comparison forces their disappearance within
$3/(2c)$ further time units. Surgery preserves that lower bound.
Thus this alternative leaves only flat components, or an empty slice,
with the actual finite preceding history retained. This uses the
elementary positive-scalar argument of KL81.1, not a general finite-time
extinction theorem or a new default route through extinction. Flat
late components enter the separate nonnegative-curvature graph branch;
an empty slice enters LP10--LP11. If this alternative never occurs,
the negative-scalar branch below applies.

\begin{lemma}[Volume deficit forces hyperbolic limits; LTF01]
\label{lem:hyp-volume-deficit-hyperbolic-limit}
On the negative-scalar branch, the total normalized volume is bounded.
Any noncollapsed smooth parabolic limit of $t_j^{-1}g(t_js)$, with
$t_j\to\infty$, on a fixed positive $s$-interval has
\[
 \Ric(g_\infty(s))=-\frac1{2s}g_\infty(s),\qquad
 \sec(g_\infty(s))=-\frac1{4s}.
\]
The convergence is assumed on unscathed neighborhoods; the conclusion
is local and does not assume a complete limit.
\end{lemma}
\begin{proof}
Put $V(t)=\operatorname{vol}(M_t,g(t))$. The normalized scalar lower
bound is $R_{\min}(t)\geq-3/(2(t+1/4))$. Since $V'=-\int R$ on
smooth stages, and surgery decreases volume,
\[
 \overline V(t)=V(t)(t+1/4)^{-3/2}
\]
is nonincreasing and has a finite limit $V_\infty\geq0$.
On smooth stages the quantity $\mathcal R(t)=R_{\min}(t)V(t)^{2/3}$
satisfies
\begin{equation}\label{eq:hyp-normalized-scalar-deficit}
 \mathcal R'(t)\geq
 \frac23|R_{\min}(t)|V(t)^{-1/3}
       \int_{M_t}(R-R_{\min}(t))\,dV\geq0.
\end{equation}
The inequality holds in the lower-barrier, hence integrated, sense
if the scalar minimum is not differentiable. At surgery the volume
and negative scalar minimum change in the directions making
$\mathcal R$ nondecreasing. It is bounded above by zero and below
by its value at any fixed late initial time.

A noncollapsed smooth limit contains a fixed ball of positive volume
at some time, so $V_\infty>0$. The limits of $\mathcal R$ and
$\overline V$ show that $(t+1/4)R_{\min}(t)$ has a limit $c\geq-3/2$.
If $c>-3/2$, integrate $V'\leq-R_{\min}V$ for large $t$, including
its downward surgery jumps. For a sufficiently small $\eta>0$ this
gives $V(t)\leq C(t+1/4)^{-c+\eta}$, contradicting
$V_\infty>0$. Thus $c=-3/2$.

The limiting scalar curvature is consequently at least $-3/(2s)$.
If it is strictly greater at one point, smooth convergence gives a
fixed space-time neighborhood with a positive scalar gap. On its
physical realizations the integral in
\eqref{eq:hyp-normalized-scalar-deficit} is bounded below by
$c_1\sqrt{t_j}$ over a time interval of length $c_2t_j$.
The prefactor there is bounded below by $c_3t_j^{-3/2}$, because
$R_{\min}\sim-3/(2t)$ and $V(t)\sim V_\infty t^{3/2}$.
Each such interval therefore increases $\mathcal R$ by at least a
fixed positive number. Select a subsequence with disjoint time
intervals. This contradicts bounded monotonicity of $\mathcal R$.
There is no loss at intervening surgeries, which only increase it.
Hence $R_\infty(s)=-3/(2s)$ throughout the limit domain.

The scalar evolution equation now gives
$|\Ric|^2=3/(4s^2)=R^2/3$. Equality in the trace inequality implies
$\Ric=(R/3)g=-g/(2s)$. In dimension three the curvature tensor is
algebraically determined by Ricci, giving sectional curvature
$-1/(4s)$. This proves the claim and fixes the normalization used by
all later hyperbolic comparisons.
\end{proof}

\begin{lemma}[Actual thick points have a macroscopic seed; LTF02]
\label{lem:hyp-thick-point-macroscopic-seed}
Let $\rho(p,t)$ be the ambient curvature scale of KL89.1, with
$\rho=+\infty$ on a component of nonnegative sectional curvature.
Declare those infinite-scale points thin. For each fixed $w>0$,
there are $a(w),v(w)>0$ and $T(w)$ such that every $w$-thick point
at $t\geq T(w)$ satisfies
\[
 \rho(p,t)\geq a(w)\sqrt t,\qquad
 \sec\geq-(a(w)^2t)^{-1}\text{ on }B(p,a(w)\sqrt t),
\]
and this last ball has volume at least
$v(w)a(w)^3t^{3/2}$. All balls and thickness tests here are ambient.
\end{lemma}
\begin{proof}
For a component with a negative plane, compactness and continuity
of the minimum sectional curvature on distance balls give a finite
positive scale with minimum $-\rho^{-2}$. Suppose there are thick
points with $t_j\to\infty$ and $\rho_j/\sqrt{t_j}\to0$.
Their scale balls have volume at least $w\rho_j^3$. WBD03 at order
zero gives $|\Rm|\leq C(w)\rho_j^{-2}$ throughout each such ball
for large $j$. The bound extends to its closure by continuity.
At a point realizing its negative minimum, the least curvature-operator
eigenvalue in KL's convention is $-2\rho_j^{-2}$. Hamilton--Ivey
pinching and $t_j/\rho_j^2\to\infty$ therefore imply
$\rho_j^2R\to\infty$. This is a contradiction. One may use
arbitrarily close points in the open ball if the minimum is recorded
as an infimum; the same uniform bound applies.

This sequential contradiction supplies a uniform eventual positive
lower bound for $\rho/\sqrt t$ on the entire $w$-thick part.
Choose $a$ smaller than it. On the scale ball,
$\sec\geq-\rho^{-2}$ and $\operatorname{vol}B(p,\rho)\geq w\rho^3$.
Bishop--Gromov comparison down to $a\sqrt t\leq\rho$ gives
\[
 \frac{\operatorname{vol}B(p,a\sqrt t)}{(a\sqrt t)^3}
 \geq c_3w,
\]
where one may take a positive lower bound for the hyperbolic model
ratio on relative radii in $(0,1]$. Set $v=c_3w$.
The curvature lower bound weakens to $-(a^2t)^{-1}$ on the smaller
ball, proving all the assertions. Shrinking $a$ further preserves
the same lower volume ratio up to a universal factor.
\end{proof}

\begin{lemma}[Macroscopic seeds become hyperbolic on every fixed ball; LTF03]
\label{lem:hyp-macroscopic-seed-hyperbolic-balls}
Fix $a,v,L>0$. If at time $t$ an ambient ball $B(p,a\sqrt t)$ has
sectional curvature at least $-(a^2t)^{-1}$ and volume at least
$va^3t^{3/2}$, then, uniformly over all such points as $t\to\infty$,
\[
 \sup_{B(p,L\sqrt t)}|2t\Ric+g|_{g(t)}\longrightarrow0.
\]
In particular, on every fixed normalized radius there are uniform
curvature bounds and uniform positive lower volume bounds on
smaller balls. This is the actual spatial conclusion of KL88.1.
\end{lemma}
\begin{proof}
First prove the conclusion at the seed center. Shrink the seed by
Bishop--Gromov to meet KL81.3's upper-radius cutoff. Since
$h_*(t)/\sqrt t\to0$, its lower surgery-scale cutoff also holds at
late times. The resulting unscathed backward region has uniformly
bounded normalized curvature and a noncollapsed center. Local Shi,
the inherited curvature/volume-to-injectivity estimate, and smooth
flow compactness give a convergent parabolic subsequence from any
proposed counterexample. LTF01 makes its center Ricci tensor
hyperbolic, a contradiction. This argument applies to \emph{every}
fixed pair of seed constants, not just to the particular $a,v$ now
chosen.

Next retain a smaller parabolic seed satisfying KL84.1 with its
$1/(3r_0^2)$ curvature normalization and volume premise. For each
fixed $L$, part(b) of that proposition gives a number $Q=C(L)/t$
such that any point in a sufficiently large prescribed ball about
$p$ with $R\geq Q$ has a canonical neighborhood. Take $C(L)>1$.
The center scalar curvature is negative at late times by the first
paragraph. If some point in that ball has scalar curvature greater
than $Q$, a minimizing radial segment from $p$ has a point with
scalar curvature exactly $Q$. Its canonical neighborhood must be
a neck or cap: a closed positive-curvature component cannot contain
$p$. Its controlled geometry gives a smaller ball with radius a
fixed multiple of $\sqrt t$, a fixed lower sectional curvature
bound at that scale and a fixed positive volume ratio. Apply the
first paragraph to that ball. It says that its center scalar
curvature is negative at late times, contradicting $R=Q>0$.
Thus $R\leq C(L)/t$ on the prescribed ball. Hamilton--Ivey gives a
two-sided bound $|\Rm|\leq C'(L)/t$ there. Only parts(a),(b) of
KL84.1 are involved in this spatial step; no circular attempt is
made to meet part(c)'s radius bound while its enlargement factor
changes.

Perform this argument first on a ball large enough to contain all
$B(q,(L+a+1)\sqrt t)$ with $q\in B(p,L\sqrt t)$.
Each such larger ball contains the original noncollapsed seed.
Bishop--Gromov centered at $q$, now justified by the newly obtained
curvature bound on this whole larger domain, gives a uniform
positive volume ratio on $B(q,b\sqrt t)$ for every sufficiently
small fixed $b>0$. Choose $b$ so that the curvature lower bound is
at least $-(b^2t)^{-1}$. The center argument applied at every $q$
proves the asserted Ricci estimate, uniformly: a failure of
uniformity would give a sequence of these same fixed-constant seed
balls and contradict that argument. This derivation does not try
to deduce spatial volume noncollapse from a parabolic noncollapse
statement without checking its curvature hypotheses.
\end{proof}

\begin{proposition}[Unscathed windows and smooth static transport; LTF04]
\label{prop:hyp-actual-dyadic-windows}
Fix $w>0$, a normalized radius $L$ and an integer $k\geq0$.
For every sufficiently late $w$-thick point $(p,t)$, the ambient
ball $B_{g(t)}(p,L\sqrt t)$ is unscathed on $[t,2t]$.
Writing $Z_t=(M_t^+,t^{-1}g(t))$, the static-track maps $i_{t,u}$
satisfy, uniformly for $u\in[t,2t]$,
\[
 \|i_{t,u}^*g_{Z_u}-g_{Z_t}\|_{C^k(B_{Z_t}(p,L),g_{Z_t})}
 \longrightarrow0\quad(t\to\infty).
\]
Every fixed spatial buffer also has a backward unscathed interval
of length $c(w,L)t>0$ and bounds for all finite normalized curvature
derivative orders, with their constants allowed to depend on the
order. The maps, not merely their images, follow the specified
static flow tracks.
\end{proposition}
\begin{proof}
LTF02 gives a fixed macroscopic seed, and LTF03 applies on every
fixed enlarged ball. To establish the forward assertion first at
order zero, choose a small Ricci-defect tolerance $\eta$ and start
with a much larger spatial buffer. Up to a first failure of
$|2u\Ric+g|<\eta$, the normalized metric varies by at most a
controlled multiple of $\eta\,d\log u$. For $\eta$ small, sectional
curvatures are negative. Thus no surgery can affect the controlled
tracks. Distance and volume distortion preserve a definite
macroscopic seed at the transported center for $u\in[t,2t]$.
A first-exit length argument keeps its small ambient ball inside
the transported buffer; it is an ambient ball, not just an intrinsic
ball in the track domain.

At a first failure point, the transported radial paths have bounded
normalized length. Consequently that point lies in a fixed larger
ambient ball about the transported seed center. LTF03, applied at
that time with the preserved seed constants and tolerance
$\eta/2$, contradicts the first failure. A first surgery endpoint
is excluded by the same negative-curvature bound. Using a larger
initial buffer if necessary proves this on a neighborhood of the
whole requested ball. Since $\eta$ was arbitrary, the Ricci defect
tends to zero uniformly throughout these windows.

For backward buffers at time $t$, LTF03 gives curvature bounds on
a larger ambient ball and, by volume comparison from the central
seed, noncollapse on a fixed small macroscopic ball about every
point of the requested domain. Decrease their common radius if
necessary; the volume-ratio lower bound persists under this decrease.
Now WBD01 applies at every such center, since that radius is a fixed
multiple of $\sqrt t$ and $h_*(t)/\sqrt t\to0$.
Its constants are uniform over the domain, giving a common backward
interval and actual unscathed tracks. Combine with the forward
buffers just proved. Local Shi estimates on these enlarged space-time
regions give all finite normalized derivative bounds on the inner
windows, including their initial time $t$. This does not increase
the canonical-neighborhood derivative order of the chosen flow.

The volume lower bounds and curvature bounds also give uniform
injectivity bounds on each fixed inner normalized domain. In their
controlled charts, compactness or tensor interpolation upgrades the
vanishing Ricci defect to every fixed spatial derivative order:
a subsequence failing that conclusion would have a smooth tensor
limit, whose uniform zero-order limit is zero. Finally, along the
static tracks,
\[
 \frac{d}{du}(u^{-1}g(u))=-u^{-2}(g(u)+2u\Ric(g(u))).
\]
The normalized $C^k$ norm of this derivative is $o(1)/u$ uniformly
on the buffered dyadic window. For the connection comparison, use the variation formula
\[
 \partial_u\Gamma^k_{ij}=-g^{k\ell}
 (\nabla_i\Ric_{j\ell}+\nabla_j\Ric_{i\ell}-\nabla_\ell\Ric_{ij}).
\]
Since $\nabla g=0$,
$\nabla\Ric=(2u)^{-1}\nabla(g+2u\Ric)$. Integrating this identity
and its successive spatial derivatives compares the connections
with the fixed initial one, inductively in the derivative order.
The order-zero distortion compares the tensor norms. Integration over $[t,u]$,
whose logarithmic length is at most $\log2$, proves the displayed
$C^k$ estimate and its asserted uniformity.
\end{proof}

\begin{theorem}[Actual HG07 producer; LTF05]
\label{thm:hyp-actual-late-window-producer}
For the fixed admissible flow on the negative-scalar branch, each
sequence of $w$-thick points at times tending to infinity has a
subsequence converging smoothly, after normalization by $t^{-1}$,
to a complete connected oriented finite-volume hyperbolic three-manifold
of sectional curvature $-1/4$. The pointed limit family $\Lambda_w$
is compact. There is a nonincreasing $\beta_w(t)\to0$ such that
all late $w$-thick points have $\beta_w(t)$-approximations from
members of $\Lambda_w$, and static continuation on $[t,2t]$ has
error less than $\beta_w(t)$ on normalized balls of radius
$\beta_w(t)^{-1}$ through the requested increasing finite order.
These are HG07's actual limits, volume and unscathed-window inputs.
\end{theorem}
\begin{proof}
LTF02--LTF04 give noncollapse at the basepoint and uniform smooth
bounds on every fixed normalized spatial ball, with the buffers
required for smooth compactness. Take a diagonal subsequence over
integer radii and derivative orders. The limit is complete because
its pointed exhaustion contains every finite distance ball with
compact smaller closures. It is orientable: an orientation-reversing
loop would lie in one compact approximation domain and embed into an
orientable slice, which is impossible. Choose the orientation from
an oriented frame at the basepoint. LTF03, or LTF01 on its parabolic
neighborhoods, gives sectional curvature $-1/4$ at time one.
LTF01 bounds total normalized volume by some $V<\infty$; smooth
convergence on each compact ball gives limit volume at most $V$.
Thus the complete limit has finite volume, not just finite volume
on individual compact sets.

To check compactness of $\Lambda_w$, for its $j$th proposed member
choose an actual realizing $w$-thick slice at time greater than $j$
which is $1/j$-close on radius-$j$ domains through order $j$.
Apply the preceding extraction to these actual points. On each
fixed buffered domain the realizing errors vanish, so the same
limit is a subsequential limit of the proposed members and belongs
to $\Lambda_w$. This diagonal argument proves closedness as well
as precompactness; it does not presume that an arbitrary model is
an actual flow limit.

For every fixed finite accuracy request, all sufficiently late
$w$-thick points are approximated by some member of $\Lambda_w$.
Otherwise a sequence of failures has a pointed limit in that same
family, contradicting its failure of the request. If $\Lambda_w$
is empty, this argument says that the $w$-thick part is eventually
empty. Combine the approximation thresholds with LTF04's uniform
thresholds for radius $2n$, order $n$ and error $1/(2n)$, for each
integer $n$. Taking increasing thresholds and setting
$\beta_w=1/n$ on successive intervals gives the stated decreasing
function. Stronger finite requests can absorb the conventional
integer rounding and composition margins. Each threshold concerns
only finitely many conditions. The static maps are defined on an
open buffered domain containing the advertised balls, so shortening
to those balls changes none of the estimates.
\end{proof}

\paragraph{What this closes.}
LTF01--LTF05 supply HG07's flow producer, with its actual normalization,
uniformity and smooth time windows, under the retained basic analytic
and surgery foundations. They do not select a persistent model, impose
a marking, produce disjoint model images or prove thick coverage.
Those are HG08's geometric constructions using HG06 and HPS01--HPS05.
In particular, one cannot replace the remaining minimal-cusp induction
by the compactness of $\Lambda_w$ alone. Accepted PC bindings and Lean
proofs remain separate from this written producer.

% END REVISION163 THICK FLOW PRODUCTION

\begin{foundation}[Persistent hyperbolic cores; HG08]
\label{found:hyp-persistent-cores}
Using HG03, HG06 and HG07, construct the KL90.1 finite family of complete
finite-volume hyperbolic models and time-dependent approximation maps.
On each fixed compact truncation, for sufficiently large time these are
embeddings with pairwise disjoint images. Their normalized metric errors
tend to zero; their images cover the prescribed sufficiently thick region.
They have the source's slow time variation on the relevant unscathed
domains, not just independently chosen maps on a subsequence.
\end{foundation}
In the source notation, for a decreasing $\alpha(t)\to0$, the domains
exhaust balls of radius $\alpha(t)^{-1}$, the normalized metric error
is less than $\alpha(t)$ at the specified increasing derivative order,
and $|\partial_t f|_{g(t)}<\alpha(t)t^{-1/2}$. Fixed core/collar compact
sets must eventually lie in those domains. The proof chooses a
minimal-cusp limit, corrects successive approximations using HG06 and
smooths their time dependence. KL calls the last step a standard
smoothing argument: its full controlled construction and preservation of
embedding/disjointness remain named implementation obligations.

% BEGIN REVISION161 PERSISTENCE SMOOTHING
\subsection{Controlled transitions and smoothing in time}
\label{sec:hyp-controlled-time-smoothing}

KL's proof of90.1, printed2817/PDF231, passes from a family with
small jumps to a smooth slowly varying family. The following
arguments supply that passage on actual buffered compact domains.
They also specify what the discrete persistence construction must
produce. This subsection does not infer that construction from
the definition of a persistent family.

\begin{lemma}[Metric control upgrades close transitions; HPS01]
\label{lem:hyp-close-transition-regularity}
Let $(H,h)$ be a fixed smooth Riemannian manifold and let
$D\Subset\operatorname{int}D_1\Subset U$ be compact domains with
buffered neighborhoods. Suppose $a_i,b_i:U\to(N_i,g_i)$ are
smooth embeddings into targets of the same dimension as $H$,
and their pullback metrics tend to $h$ in
$C^\infty$ on compact subsets of $U$, and
\[
 \sup_{p\in D_1}d_{g_i}(a_i(p),b_i(p))\longrightarrow0.
\]
Then $a_i^{-1}b_i$ is defined on a neighborhood of $D$ for large
$i$ and tends to the identity there in every fixed finite $C^k$
norm. For a prescribed finite output order and error, only finitely
many metric derivatives and a sufficiently small distance error
are needed, with thresholds depending on the fixed source buffers.
\end{lemma}
\begin{proof}
Choose a positive radius $r$ such that closed $h$-balls of radius
$2r$ centered in a neighborhood of $D$ lie in $\operatorname{int}D_1$.
For large $i$ the pullback metrics have a uniform lower bound
$a_i^*g_i\geq h/4$. QRG01 lifts a minimizing short target path
from $a_i(p)$ to $b_i(p)$ into that source ball. If a target is
not complete, use a path whose length is arbitrarily close to
the distance; the strict smallness margin suffices. Thus
$b_i(p)\in a_i(U)$ and, for $\Phi_i=a_i^{-1}b_i$,
\[
 d_h(\Phi_i(p),p)\leq2d_{g_i}(a_i(p),b_i(p)).
\]
The same argument on the buffered neighborhood gives uniform
$C^0$ convergence there. The exact tensor identity is
\[
       b_i^*g_i=\Phi_i^*(a_i^*g_i).
\]
Choose finitely many source coordinate boxes covering $D$ whose
larger boxes lie in that neighborhood. Uniform $C^0$ closeness
puts $\Phi_i$ and the identity in the same larger box. Positive
metric bounds in the identity give a uniform first-derivative
bound. The isometry connection identity in those coordinates is
\[
 \partial_\alpha\partial_\beta\Phi_i^\ell
 =\Gamma(b_i^*g_i)^c_{\alpha\beta}\partial_c\Phi_i^\ell
  -\Gamma(a_i^*g_i)^\ell_{mn}(\Phi_i)
       \partial_\alpha\Phi_i^m\partial_\beta\Phi_i^n.
\]
It gives a uniform second-derivative bound. Differentiating it
inductively gives bounds for any specified finite number of map
derivatives from finitely many metric derivatives. Compact-domain
Arzela--Ascoli on smaller boxes now gives subsequential $C^k$
convergence. Every limit is the identity, since its $C^0$ limit
already is. A subsequence violating $C^k$ convergence would have
this same convergent subsubsequence, a contradiction.

For the last assertion fix the desired order, the boxes and the
buffers. Use enough metric derivatives to bound one more map
derivative than that order. If no positive joint smallness
threshold existed, a sequence with metric and distance errors
tending to zero would contradict the preceding finite-order
argument. This is a uniform qualitative modulus for the fixed
source data, not a computed or target-dependent constant.
\end{proof}

\begin{corollary}[The differentiable rigidity actually needed for transitions; HPS02]
\label{cor:hyp-differentiable-rigidity-on-compacts}
Fix a pointed HG02 model $H$, a compact domain $D\subset H$,
an integer $k\geq0$ and an error $\varepsilon>0$. A sufficiently
accurate HG05 approximation $f$ into a connected complete
finite-volume hyperbolic target with at least as many cusps as
$H$ admits an isometry $I:H\to H'$ such that
$I^{-1}f$ is $C^k$-close to the identity on $D$, with error less
than $\varepsilon$ in fixed buffered source charts.
\end{corollary}
\begin{proof}
Enclose $D$ and its buffers in a fixed source ball. If the claim
failed, choose counterexamples with HG05 parameters tending to
zero. QRG09 and QRG03 give isometries $I_i$ with distance error
tending to zero on increasingly large fixed balls: at stage $i$
request error less than $1/i$ on the ball of radius $i$, and
choose the approximation parameter below the corresponding
threshold. Apply HPS01 to $a_i=I_i$ and $b_i=f_i$ on the fixed
buffered neighborhood. Their pullback metrics are $h$ and a
sequence tending to $h$. The resulting $C^k$ convergence
contradicts the selected counterexamples. No homotopy marking
of $f$, canonical choice of $I$, or harmonic-map theorem is used.
\end{proof}

\begin{lemma}[A transition extends to a small compactly supported isotopy; HPS03]
\label{lem:hyp-compact-transition-isotopy}
Fix nested compact domains
$D\Subset\operatorname{int}D_1\Subset\operatorname{int}D_2$
in a complete smooth $(H,h)$. If a smooth map $\Phi$ is sufficiently
$C^1$-close to the identity on a neighborhood of $D_2$, there is
a smooth isotopy $E_\lambda:H\to H$, $0\leq\lambda\leq1$,
through diffeomorphisms, equal to the identity outside $D_2$,
with $E_0=\operatorname{id}$ and $E_1=\Phi$ on $D_1$.
For every fixed $k$, this construction tends to the identity
in $C^k$ if $\Phi$ does, uniformly in $\lambda$. Moreover
\[
 \sup_H|\partial_\lambda E_\lambda|_h
 \leq C\sup_{D_2}d_h(\Phi(p),p)
\]
for a constant determined by the fixed source buffers. The
paths $E_\lambda(D)$ remain in an arbitrarily prescribed fixed
neighborhood of $D$ when the transition is sufficiently small.
\end{lemma}
\begin{proof}
Choose a smooth cutoff $\chi$ supported in $\operatorname{int}D_2$
and equal to one on a neighborhood of $D_1$. The minimum normal
radius on this compact set is positive. For a sufficiently close
transition the vector
\[
 v(p)=\chi(p)\exp_p^{-1}(\Phi(p))
\]
is defined there and extends by zero to a smooth vector field
on $H$. Set $E_\lambda(p)=\exp_p(\lambda v(p))$.
On the fixed compact support, the derivative in $p$ is uniformly
close to the identity in local charts, for every $\lambda$.
Thus each $E_\lambda$ is a local diffeomorphism. It is proper
because it is the identity off a fixed compact set. On each
connected component its image is both open and closed, so it
is a surjective covering. Its homotopy to the identity is the
displayed family $E_\lambda$; the induced fundamental-group
homomorphism is consequently surjective, up to the basepoint
path. A connected covering with surjective induced homomorphism
has one sheet, by path and homotopy lifting. This proves global
injectivity, including pairs of points not in one coordinate box.

Differentiating the exponential expression on the fixed compact
normal neighborhood gives the asserted finite-order bounds.
Differentiation in $\lambda$ involves only the vertical derivative
of the exponential applied to $v$, giving the displayed speed
bound. The geodesic from $p$ to $E_\lambda(p)$ has length at most
$|v(p)|$, proving the neighborhood assertion. Cutoff derivatives
are fixed before the smallness threshold is chosen. No extension
of $\Phi$ itself to a global diffeomorphism is assumed.
\end{proof}

\begin{proposition}[Explicit smoothing of discrete persistent maps; HPS04]
\label{prop:hyp-discrete-map-smoothing}
Let $H$ be one fixed model, or a fixed finite disjoint union of
models, and let $t_j=2^jT$, $T>0$. Suppose a surgery flow supplies
the following actual data. At $t_j$ there is an embedding $f_j$
on an open source domain $U_j$, and it extends by static curves
to embeddings $f_j(t)$ on $[t_j,t_{j+1}]$ in an unscathed region.
The domains exhaust every source compact set. For every fixed
compact $D$ and finite order $k$,
\[
 \sup_{t\in[t_j,t_{j+1}]}
 \|t^{-1}f_j(t)^*g(t)-h\|_{C^k(D,h)}\longrightarrow0,
\]
and the endpoint jumps satisfy
\[
 \sup_{p\in D}d_{t_{j+1}^{-1}g(t_{j+1})}
       (f_j(t_{j+1})(p),f_{j+1}(p))\longrightarrow0.
\]
For a finite family assume the corresponding embeddings have
disjoint images on their supplied domains within each window.
Then on an open space-time domain containing $D\times[T_D,\infty)$
for every fixed compact $D$, there is a smooth family $F(t)$
of embeddings, with disjoint component images, such that on $D$
\[
 \|t^{-1}F(t)^*g(t)-h\|_{C^k(D,h)}\longrightarrow0,
 \qquad
 \sqrt t\sup_{p\in D}|\partial_tF(p,t)|_{g(t)}\longrightarrow0.
\]
The maps are obtained inside the old buffered images. No surgery
is crossed except within the supplied unscathed regions.
\end{proposition}
\begin{proof}
Fix a compact exhaustion $D_m\Subset\operatorname{int}D_{m+1}$.
At $t_{j+1}$, HPS01 makes
\[
 \Phi_j=f_j(t_{j+1})^{-1}f_{j+1}
\]
defined and close to the identity in every fixed finite order
on each fixed buffered compact set, for all large $j$.
This assertion uses the endpoint metric
$t_{j+1}^{-1}g(t_{j+1})$ for both maps. After discarding finitely many indices, choose integers
$m(j)\geq2$ tending to infinity sufficiently slowly that all domains and
smallness bounds needed on $D_{m(j)+2}$ hold, that HPS03 applies
there, and that its resulting isotopies $E_{j,\lambda}$ satisfy
\[
 \sup_{0\leq\lambda\leq1}
 \|E_{j,\lambda}-\operatorname{id}\|_{C^{m(j)+1}(D_{m(j)+1})}
   <1/m(j),\qquad
 \sup_H|\partial_\lambda E_{j,\lambda}|_h<1/m(j).
\]
Choose also the corresponding metric errors less than the
finite-order margins needed for the pullback conclusion below.
For every fixed $m$ these are finitely many eventual conditions;
take increasing thresholds and let $m(j)$ be the largest passed
stage, capped by $j$. This proves the diagonal choice without
assuming uniform constants over the full noncompact source.
Use the minimum of the finitely many component thresholds for
a finite family. The isotopy equals $\Phi_j$ near $D_{m(j)+1}$
at $\lambda=1$, and carries $D_{m(j)}$ inside $U_j$ throughout.

Fix one smooth nondecreasing function $\theta$ equal to zero
on $(-\infty,7/4]$ and one on $[15/8,\infty)$.
On the $j$-th interval define, on the indicated inner domain,
\[
 F(t)(p)=f_j(t)\big(E_{j,\theta(t/t_j)}(p)\big).
\]
Until $7t_j/4$ this is the old static map. From $15t_j/8$
onward it equals the new endpoint map transported backward
along the old unscathed window, since $E_{j,1}=\Phi_j$ on
the inner domain. That backward transport is defined: its
endpoint values lie in $f_j(t_{j+1})(U_j)$ by HPS01.
Near $t_{j+1}$ it and the new forward static map are therefore
the same static curves. They agree on a space-time neighborhood
of every point of a fixed sufficiently late compact domain.
All time and spatial derivatives match at the junction; continuity
alone is not the gluing criterion used here.

For each fixed $D_n$ choose $T_n$ beyond the finitely many
junctions where $m(j)<n+2$. On
\[
       \Omega=\bigcup_n\operatorname{int}D_n\times(T_n,\infty)
\]
the displayed interval formulas consequently define one smooth
map into the smooth unscathed part of the flow space-time.
At a fixed time they are restrictions of one common interval
formula. Each is an embedding, because $E_{j,\lambda}$ is a
diffeomorphism and $f_j(t)$ is an embedding on the image domain.
For several models the adjusted images lie in the disjoint old
images, proving disjointness without an inference from a limit
of injective maps. For the finitely many early intervals one
may shorten $\Omega$; no late conclusion changes.

On a fixed compact set, the pullback metric is
\[
 E_{j,\theta(t/t_j)}^*\big(t^{-1}f_j(t)^*g(t)\big).
\]
The uniform metric convergence and the finite-order isotopy
convergence prove the required $C^k$ assertion. Static maps
have zero spatial velocity along the flow time direction.
Thus the only velocity contribution is the isotopy term.
The normalized metric bound gives
$|df_j(t)|_{h,g(t)}\leq2\sqrt t$ for large $j$, whence
\[
 |\partial_tF|_{g(t)}
 \leq\frac{2\sqrt t\,\|\theta'\|_\infty}{t_jm(j)}
 \leq\frac{4\|\theta'\|_\infty}{m(j)\sqrt t}.
\]
This is a bound in the physical metric $g(t)$ and physical time
$t$. It is not the velocity bound in the normalized metric.
Since $m(j)\to\infty$, it proves the claimed slow variation.
The same construction makes the maps arbitrarily close to the
old static maps on each fixed compact set, uniformly in the
window, by the displacement bound in HPS03.
\end{proof}

% BEGIN REVISION166 SMOOTHING THRESHOLD NOTICE
\paragraph{Final-audit correction to the finite-family choice.}
In HPS04's preceding proof, the common starting threshold is the MAXIMUM
of the component thresholds. Only permitted error bounds are combined
by a minimum. LAU01 in Section~\ref{sec:collapse-longtime-final-analytic-audit}
proves the corrected simultaneous choice; the interpolation is unchanged.
% END REVISION166 SMOOTHING THRESHOLD NOTICE

\begin{lemma}[One decreasing accuracy function for the analytic bounds; HPS05]
\label{lem:hyp-common-accuracy-function}
For the family from HPS04 there are $T_0$ and a nonincreasing
$\alpha:[T_0,\infty)\to(0,1)$ tending to zero such that $F(t)$
is defined on the union of model balls of radius $\alpha(t)^{-1}$,
has normalized metric error less than $\alpha(t)$ there through
order $\lceil\alpha(t)^{-1}\rceil$, and has physical speed less
than $\alpha(t)t^{-1/2}$ there. Thick-part coverage is a separate
predicate; it is not a conclusion of this lemma.
\end{lemma}
\begin{proof}
For every integer $n\geq2$, choose a compact buffered domain
containing the union of closed model balls of radius $2n$.
HPS04 gives an eventual time after which the maps are defined
on a neighborhood of this domain, the $C^n$ metric error is
less than $1/(2n)$, and $\sqrt t$ times the speed is less than
$1/(2n)$. Choose strictly increasing thresholds $T_n\geq n$,
strictly later than those bounds, and set $\alpha(t)=1/n$ on
$[T_n,T_{n+1})$. This function has the stated properties.
The ambient map is already smooth on the open domain $\Omega$,
so jumps of the auxiliary accuracy function do not introduce
jumps of the maps. All integer requests have finite thresholds;
no single time is asserted to satisfy infinitely many derivative
orders. Later coverage predicates may enter this finite-stage
choice only after they have been proved on the stated balls.
\end{proof}

\paragraph{What the discrete construction still must supply.}
HPS04 replaces the source's smoothing sentence by an explicit
construction and proves embedding, disjoint-image preservation,
all-order local convergence and the correctly normalized speed.
Its actual inputs remain the successive maps, their common
unscathed windows, vanishing endpoint jumps and disjoint buffered
domains. HPS02 is available when choosing the rigidity corrections;
HPS01 also applies directly to their actual slice maps. The
minimal-cusp induction, finite family, coverage and HG07's actual
flow estimates are developed separately. This subsection alone
does not mark HG08, LP03, LP06 or LC90 complete.
% END REVISION161 PERSISTENCE SMOOTHING

\begin{lemma}[Finite-family volume budget; HG09]
\label{lem:hyp-finite-volume-budget}
Suppose every newly selected model contributes a disjoint core of
normalized volume at least $v/2>0$ at a sufficiently late common time,
and the total normalized slice volume is at most $V<\infty$. Then at
most $2V/v$ such selections are possible.
\end{lemma}
\begin{proof}
For any finite selection, choose cores capturing the required volume and
then one time meeting all its approximation and disjointness thresholds.
Finite additivity gives $kv/2\leq V$. An integer $k>2V/v$ contradicts
this inequality. The time may depend on the finite selection; no one
time works for a hypothetical infinite family by assumption. Exhaustion
of each finite-volume model, HG03's positive bound and HG08's controlled
embeddings are necessary inputs to the application.
\end{proof}

% BEGIN REVISION164 ACTUAL PERSISTENT FAMILY
\subsection{The actual persistent family, disjointness and thick coverage}
\label{sec:hyp-actual-persistent-family}

Throughout this subsection HG07 is supplied by LTF05 on the actual
flow. Distances in a slice are measured in $Z_t=(M_t^+,t^{-1}g(t))$;
between different connected components set that distance to $+\infty$.
Write
\[
 \Theta_t(p)=\frac{\operatorname{vol}_{Z_t}B(p,\rho_t(p))}{\rho_t(p)^3}
\]
when the normalized ambient curvature scale $\rho_t(p)$ is finite.
Points of infinite scale remain in the separate thin, nonnegative
branch. A $w$-thick point means a finite-scale point with
$\Theta_t(p)\geq w$. This is the same selector as LTF02, in normalized
units, not the injectivity-radius thick set in a hyperbolic model.

\begin{lemma}[One robust thick basepoint in every limit model; HPI01]
\label{lem:hyp-universal-thick-basepoint}
There are universal numbers $a,w_*>0$ such that every oriented
finite-volume hyperbolic model $H$ of curvature $-1/4$ has a point
$x$ with the following property. A sufficiently accurate $C^2$
embedding of $B_H(x,20)$ into a complete slice sends all of
$B_H(x,a)$ to points with finite curvature scale and
$\Theta_t\geq8w_*$. The accuracy threshold is independent of $H$.
For every $0<w\leq w_*$, the sets of underlying unpointed models
in $\Lambda_w$ and $\Lambda_{w_*}$ are therefore the same.
\end{lemma}
\begin{proof}
MGL08--MGL09 supply a universal $\iota>0$ and a point with
$\operatorname{inj}_H(x)\geq\iota$. Choose
$a<\min(\iota/10,1/100)$. The injectivity radius on $B_H(x,a)$
is at least $\iota-a$, so each radius-$a$ ball there has the same
positive volume $v_a$ as its ball in hyperbolic space.
For sufficiently small metric error, lengths change by factors
between $9/10$ and $11/10$, and sectional curvatures lie between
$-3/10$ and $-1/5$ on the image. These constants follow from the
intrinsic $C^2$ tensor estimate against the constant-curvature metric;
no coordinate injectivity bound away from the basepoint is needed.
QRG01, applied with the indicated source buffer, puts the ambient
radius-four ball about each image point in that controlled image.
Consequently its curvature scale lies between $1$ and $3$.
Its radius-one ball contains the image of the corresponding source
radius-$a$ ball and has volume at least $v_a/2$.
Thus $\Theta_t\geq v_a/54$. Choose $w_*=v_a/1000$.

If $H$ is a pointed limit from any fixed $w$-thick sequence, move
the basepoint in $H$ to such an $x$. This is a finite source distance.
The actual realizing embeddings include $B_H(x,20)$ at sufficiently
late members of that sequence. Their new basepoints are $w_*$-thick
by the first paragraph, and realize the same underlying $H$.
Thus every underlying model in any $\Lambda_w$ belongs to the
underlying-model set of $\Lambda_{w_*}$. Conversely
$\Lambda_{w_*}\subset\Lambda_w$ as pointed families when $w\leq w_*$.
The equality asserted is of underlying models; the allowed basepoints
need not be the same.
\end{proof}

\begin{lemma}[Actual inverse-comparison maps; HPI02]
\label{lem:hyp-inverse-comparison-estimates}
Suppose $f:(H,x)\to(Z,p)$ and $\phi:(H',x')\to(Z,p)$ are pointed
embeddings with small pullback metric errors on sufficiently large
buffered balls. On each fixed smaller source ball,
$A=\phi^{-1}f$ is defined and is a metric approximation from $H$ to
$H'$. Each prescribed finite $C^k$ error of $A^*h'-h$ tends to zero
with the corresponding sufficiently high finite errors of $f^*g-h$
and $\phi^*g-h'$. The bounds require no injectivity-radius lower bound
throughout the target core.
\end{lemma}
\begin{proof}
For a requested radius $R$, take the source domains to include
$B_H(x,2R+2)$ and $B_{H'}(x',8R+8)$, with length distortion less
than a factor two. The image of the latter contains the ambient
ball needed to contain $f(B_H(x,R))$, by QRG01 and the length bound
on images of radial paths. This defines the actual inverse map on
a neighborhood of the requested ball.

Put $g_0=f^*g$ and $k_0=\phi^*g$, so $A^*k_0=g_0$ exactly.
That identity bounds $dA$ and its inverse in the source and target
reference metrics. The isometry connection identity is
\[
 \nabla^{h,h'}dA
 =dA(\nabla^{g_0}-\nabla^h)
   +(\nabla^{h'}-\nabla^{k_0})(dA,dA).
\]
The connection differences and their finite derivatives are controlled
by the metric errors. Differentiating this identity inductively bounds
all required higher derivatives of $A$, intrinsically along $A$.
Finally use
\[
 A^*h'-h=(g_0-h)+A^*(h'-k_0).
\]
The tensor product and pullback differentiation rules give the claimed
finite-order estimates. For example, requesting two extra metric
derivatives is more than sufficient at each fixed target order.
All norms are measured on the buffered domains already specified.
This proves the composition assertion actually used in the correction
step, rather than presuming that an arbitrary inverse of close maps
has controlled derivatives.
\end{proof}

\begin{proposition}[Actual minimal-cusp correction and one persistent model; HPI03]
\label{prop:hyp-actual-minimal-cusp-persistence}
If the underlying-model set of $\Lambda_{w_*}$ is nonempty, choose a
model $H$ with the smallest number of cusps and a basepoint $x$ as in
HPI01. There is an actual smooth family of embeddings of exhausting
compact subsets of $H$ into the late slices, with normalized metric
error tending to zero in every finite order on each fixed compact
set, and physical velocity $o(t^{-1/2})$ there. Its basepoint and a
fixed small source neighborhood remain uniformly $w_*$-thick.
\end{proposition}
\begin{proof}
Cusp numbers are finite nonnegative integers by HG03, so a minimum
exists. Rebase as in HPI01, preserving membership in $\Lambda_{w_*}$.
Choose a fixed coarse approximation tolerance $\delta$ sufficiently
small for the following finite requests: HPI01 on its radius-twenty
buffer; QRG09 and HPS02 on the compact set needed for the first
correction; and a correction displacement less than $a/100$ on the
source basepoint neighborhood. Choose the larger source buffers and
finite derivative orders for HPI02 before choosing $\delta$.
LTF05's error function can then be made small enough for all these
requests by choosing the starting time late enough.

Take such a late realizing slice at time $T$ and an initial
$\delta$-approximation $f_0:H\to Z_T$. Put $t_j=2^jT$.
Assume a coarse approximation $f_j$ is given at $t_j$.
Its basepoint is uniformly thick by HPI01. LTF04--LTF05 transport
its fixed coarse source buffer along actual static tracks to all
of $[t_j,t_{j+1}]$, with errors still within the chosen coarse
margin. Write $f_j(t)$ for this static extension. Its endpoint
basepoint is still $w_*$-thick. LTF05 supplies a pointed
$\beta(t_{j+1})$-approximation
\[
 \phi_j:(H'_j,x'_j)\longrightarrow
       (Z_{t_{j+1}},f_j(t_{j+1})(x)),\qquad
 (H'_j,x'_j)\in\Lambda_{w_*}.
\]
HPI02 makes $A_j=\phi_j^{-1}f_j(t_{j+1})$ an actual approximation
from the fixed $H$ into $H'_j$. Minimality gives
$\#\operatorname{Cusps}(H'_j)\geq\#\operatorname{Cusps}(H)$.
QRG09 and HPS02 therefore give an isometry $I_j:H\to H'_j$ close
to $A_j$ on the prescribed compact set. Define
\[
 f_{j+1}=\phi_j I_j.
\]
Its pullback metric error is that of $\phi_j$, and its domain contains
balls about $x$ of radius tending to infinity: the displacement
$d_{H'_j}(I_jx,x'_j)$ is less than the fixed correction margin,
while $\beta(t_{j+1})^{-1}\to\infty$. Thus $f_{j+1}$ again meets
the coarse requirements, including the robust basepoint thickness.
This proves the actual induction at every dyadic time.

The corrections can and must improve with time. On any fixed source
compact set, the metric errors of $f_j$ tend to zero because the new
map is a vanishing-error approximation precomposed with an isometry
whose basepoint shift is bounded. LTF04 then gives the same conclusion
uniformly on its entire dyadic window. HPI02 makes the endpoint
comparisons $A_j$ increasingly accurate on each fixed compact set.
For each integer $n$, QRG09/HPS02 have a positive threshold for
correction error $1/n$ on a fixed radius-$n$ buffer through order $n$.
For sufficiently large $j$ that threshold is met. Choose $I_j$ with
the strongest of a slowly increasing finite list of these requests,
always retaining the initial coarse request. The finite-list request
is obtained by taking the largest compact set and strongest accuracy;
no intersection of unrelated choices of isometry is presumed.
The resulting endpoint jumps tend to zero on every fixed compact
set. HPS01 also gives their smooth transition control.

All inputs to HPS04 have now been produced for this model: actual
exhausting maps, actual unscathed windows, all-order compact metric
convergence and vanishing jumps. Apply its explicit interpolation
inside the old images. Choose its compact exhaustion slowly enough
that the basepoint correction stays inside $B_H(x,a/10)$.
Every point of that fixed neighborhood is uniformly thick by HPI01,
both during the static portions and during interpolation. The result
has the stated embedding, metric and physical speed conclusions.
A closed model is included: a sufficiently large source domain is
then all of $H$, whose compact codimension-zero image is an entire
slice component. HPS04 applies with compact support equal to $H$.
\end{proof}

\begin{lemma}[Thin cusp barriers and escape from old models; HPI04]
\label{lem:hyp-thin-cusp-barriers}
For a finite already constructed persistent family, choose sufficiently
deep fixed compact cusp truncations $C_i\subset H_i$. At all
sufficiently late times their boundary collars have
$\Theta_t<w_*/2$. A continuous admissible $w_*$-thick basepath
starting outside all their images cannot enter them. Moreover,
$w_*$-thick points outside these images escape normalized distance
from every old model basepoint as time tends to infinity.
\end{lemma}
\begin{proof}
In a source cusp write $h=ds^2+e^{-s}q$. A fixed-radius ball at
large depth $s$ is contained in a band of bounded depth width, whose
volume is at most $C e^{-s}\operatorname{area}(T,q)$.
On each fixed buffered compact band the persistent maps converge
smoothly to this metric. QRG01 puts the ambient radius-four balls
about the retained image band inside its larger controlled image.
Their curvatures are close to $-1/4$, so their curvature scales lie
between one and three. Their scale-ball volumes are consequently
bounded by a fixed multiple of the volume of the source band.
Choose the truncation levels so deep that the resulting thickness
ratio is less than $w_*/2$, with a strict margin, on the boundary
and a fixed collar on either side. There are only finitely many
cusps and models, so one common late time works. Closed old models
have no cusp boundary and their images are whole components.

The old boundary families are smooth in the unscathed flow space-time.
If an admissible continuous thick path changes membership in an old
core, it meets one of these boundaries, contradicting its thickness.
At a surgery time the core and its collar have unscathed neighborhoods,
so the same assertion uses the specified common identification there.
For a closed old model its whole-component track is both open and
closed locally in space-time, so membership cannot change either.

Finally suppose exterior thick points $p_j$ at times $t_j\to\infty$
have bounded normalized distance from one old basepoint. QRG01 for
that old model's increasingly large domains places each $p_j$ in
the image of a fixed source ball. Its inverse lies outside the fixed
core. For a closed old model this is already impossible. Otherwise
these inverse points lie in fixed compact portions of the cusp tails
beyond the chosen truncation. The same band estimate applies uniformly
on those portions, using a larger fixed source buffer when necessary.
It makes $p_j$ less than $w_*/2$-thick, a contradiction. This proves
the stated uniform escape conclusion; it is stronger than saying
that the points lie outside one arbitrarily chosen growing image.
\end{proof}

\begin{proposition}[Actual disjoint iteration and finite termination; HPI05]
\label{prop:hyp-actual-disjoint-finite-family}
The construction yields a finite, possibly empty, family of persistent
models $H_1,\ldots,H_k$ whose fixed compact images are pairwise
disjoint at sufficiently late common times. For every $w>0$, there
is no sequence of $w$-thick points escaping normalized distance from
all these model basepoints.
\end{proposition}
\begin{proof}
For a finite old family define $\Lambda_w^{\rm esc}$ using actual
$w$-thick realizing sequences whose distances from every old
basepoint tend to infinity. This family is compact: in LTF05's
closedness argument choose the $j$th realizing point also to have
all those distances greater than $j$. Rebasing a finite distance
in a limit preserves escape, because the realizing connecting path
has bounded normalized length. HPI01 therefore proves equality of
the underlying-model sets for all $0<w\leq w_*$, also for these
remaining families.

If a remaining family is nonempty, choose its minimal-cusp model
and start its realizing approximation at a point escaping the old
basepoints, hence outside their fixed deep cores. The old cores
are already defined at that late time. We verify the endpoint
premise needed to repeat HPI03 in this restricted family.
Every late $w_*$-thick point outside the old cores admits increasingly
accurate approximations from $\Lambda_{w_*}^{\rm esc}$, uniformly.
Otherwise a sequence of failures escapes by HPI04; LTF05 extracts
an actual limit in this same remaining family, a contradiction.
The unscathed-window estimates still apply to all these actual
thick points, independently of which limit family is used.

Run HPI03 with those endpoint approximations. Its static basepaths
are uniformly thick and hence stay outside the old cores by HPI04.
At a correction time the old and new basepoints are joined inside
the image of the fixed source basepoint ball by the small transition
isotopy. HPI01 makes that whole path thick, so the new endpoint
also stays outside. The same argument applies to the time smoothing:
at any time its basepoint lies in a connected, uniformly thick
source-ball image containing the old static basepoint. Such a ball
cannot meet an old thin boundary. Thus the actual smoothed basepath
stays exterior, and HPI04 shows that it escapes all old basepoints.
This verifies membership of all subsequent endpoint limits in the
remaining family throughout the induction; it is not assumed from
the initial location alone.

If two basepoints escape each other, images of any fixed source
compact sets are disjoint at late times: each image lies within a
fixed normalized distance of its own basepoint, by the length bound
on source paths. Finitely many such requirements have a common
late time. Thus the new model has the asserted disjointness from
all old models. Shrinking to a common slowly growing exhaustion
will give disjoint full advertised domains after the iteration ends.

Let $v>0$ be HG03's universal lower model volume at curvature $-1/4$,
and $V<\infty$ the total normalized slice-volume bound from LTF01.
For each selected model choose a compact core of volume at least
$3v/4$. At a common sufficiently late time these cores have pairwise
disjoint images, each of normalized volume at least $v/2$.
HG09 therefore bounds the number of selections by $2V/v$.
If a remaining escaping family were still nonempty after that many
selections, the construction just proved would add another model,
contradicting the bound. Hence the process terminates with all
escaping families empty. Equality of underlying models for small
$w$ handles every $w\leq w_*$; for larger $w$ use inclusion in the
$w_*$-thick part. If no family was initially nonempty the empty
selection has the same property, by LTF05.
\end{proof}

\begin{lemma}[Coverage first on fixed balls; HPI06]
\label{lem:hyp-fixed-thickness-coverage}
For the family from HPI05 and every fixed $w>0$, there is a finite
$R(w)$ such that every sufficiently late $w$-thick point is in one
of the images of $B_{H_i}(x_i,R(w))$. If the family is empty, the
$w$-thick part is eventually empty.
\end{lemma}
\begin{proof}
Otherwise choose times tending to infinity and $w$-thick points
whose distance from every old basepoint tends to infinity. LTF05
gives an escaping limit, contradicting HPI05. Thus all such points
are within a fixed normalized distance of at least one basepoint.
The persistent domains exhaust larger source balls with length
error tending to zero. QRG01 then puts all the required ambient
balls inside their images, with inverse points in a fixed larger
source ball. Increasing $R(w)$ gives the conclusion. This argument
uses ambient ball coverage, not a presumed inverse defined throughout
the complement of the old cores.
\end{proof}

\begin{lemma}[The matching radius for thick coverage; HPI07]
\label{lem:hyp-matching-radius-coverage}
For all sufficiently small fixed $w>0$, every sufficiently late
$w$-thick point belongs to an image of $B_{H_i}(x_i,1/w)$.
\end{lemma}
\begin{proof}
On any fixed buffered source compact set, the normalized ambient
curvature scales at its image points tend uniformly to $2$.
Indeed, QRG01 contains their ambient radius-four balls in a larger
controlled source image. Curvature convergence on those balls bounds
the scale below by $(1/4+\varepsilon)^{-1/2}$; the negative curvature
at the center bounds it above by $(1/4-\varepsilon)^{-1/2}$.
Both bounds tend to $2$. Bilipschitz ball inclusions under the same
inverse maps and convergence of the volume forms now give uniformly
on the fixed compact set
\[
 \Theta_t(f_i(t)(y))\longrightarrow
 \Theta_{H_i}(y):=\frac18\operatorname{vol}_{H_i}B_{H_i}(y,2).
\]
The needed continuity of ball volume is uniform on compact sets:
distance spheres have zero Riemannian volume, so dominated convergence
gives joint continuity in center and positive radius; compactness
then gives uniform continuity. No injectivity-radius identification
of the thickness selector is made.

For the finite family of models, the cusp band estimate in HPI04
shows that $\Theta_{H_i}(y)\leq C_i e^{-s}$ at depth $s$.
The source distance to the fixed basepoint is at most $C_i'+s$.
Thus all model points with $\Theta_{H_i}\geq w/2$ lie within
$C_0+C_1\log(1/w)$ of their basepoints, with common finite constants.
Closed models have bounded diameter and obey the same eventual bound.
For sufficiently small $w$, this radius is less than $1/(2w)$.

HPI06 first places the inverse points of all actual $w$-thick points
in fixed compact radius-$R(w)$ source sets. On those fixed sets,
take time large enough that the preceding convergence error is
less than $w/2$. Their inverse points then satisfy
$\Theta_{H_i}\geq w/2$, so lie in $B_{H_i}(x_i,1/(2w))$ and hence
in the advertised radius-$1/w$ balls. This extra cusp estimate is
why the conclusion has the required matching radius; a diagonal
argument applied only to unspecified radii $R(w)$ would not suffice.
\end{proof}

\begin{theorem}[Actual HG08 persistent-family producer; HPI08]
\label{thm:hyp-actual-persistent-family-producer}
There are $T_0$, a nonincreasing $\alpha(t)\to0$ and a finite,
possibly empty, family of complete connected oriented finite-volume
hyperbolic models $(H_i,x_i)$ of curvature $-1/4$, with maps
\[
 f(t):\coprod_iB_{H_i}(x_i,\alpha(t)^{-1})\longrightarrow M_t^+,
 \qquad t\geq T_0,
\]
which are embeddings with pairwise disjoint images, extend to a
smooth admissible family on open exhausting space-time domains,
and satisfy
\[
 \|t^{-1}f_i(t)^*g(t)-h_i\|_{C^{\lceil\alpha(t)^{-1}\rceil}}
 <\alpha(t),\qquad
 |\partial_t f_i|_{g(t)}<\alpha(t)t^{-1/2}.
\]
Their images contain the entire $\alpha(t)$-thick part.
In particular, every fixed compact core and buffered cusp collar
has actual disjoint time-dependent embeddings with these asymptotic
controls. No ambient cusp incompressibility is asserted here.
\end{theorem}
\begin{proof}
HPI03--HPI05 produce actual smooth families on fixed compact sets,
and HPI05 gives eventual pairwise separation. For each sufficiently
large integer $n$, impose the following finite list of late conditions:
all maps are defined on neighborhoods of radius-$2n$ balls; their
metric errors through order $n$ and their speeds multiplied by
$\sqrt t$ are less than $1/(2n)$ there; these larger images are
pairwise disjoint; and the $1/n$-thick part is covered by the
radius-$n$ images. HPS04, HPI05 and HPI07 supply these conditions
at every sufficiently late time. There are finitely many models
and pairs, so one threshold $T_n$ works. Choose the thresholds
strictly increasing and set $\alpha(t)=1/n$ on $[T_n,T_{n+1})$.
This is HPS05's diagonal with the coverage premise now actually
proved on its matching balls. Integer rounding gives exactly order
$n=\lceil\alpha^{-1}\rceil$ on each interval.

The underlying families are already smooth on open exhausting
domains by HPS04. Restricting the advertised radii according to
this step function does not create time discontinuities of those
families. The factor-two buffers give neighborhoods of all the
advertised balls and preserve their embedding and disjointness.
The velocity estimate is in physical $g(t)$ and physical time,
exactly as proved by HPS04. In the empty case, LTF05 and HPI05
make every fixed $1/n$-thick part eventually empty; choose the same
threshold diagonal and all map assertions are vacuous. This proves
the stated actual producer.
\end{proof}

\paragraph{Endpoint and remaining analytic work.}
HG08 is now supplied by the actual construction, using HG07 from
LTF05 and the completed written rigidity and smoothing developments.
The output gives the time-dependent cusp embeddings needed by the
separate ambient incompressibility arguments. The present theorem
does not replace HG10--HG13 or the Meeks--Yau interface. The remaining
in-scope analytic work is the actual intrinsic thin complement:
its long collars, equality or comparison of tested balls, the
large-radius exclusion, and the parameter order leading through
LC89 to the static graph theorem. The full analytic goal and its
ONE final mathematical audit remain open until that work is done.

% END REVISION164 ACTUAL PERSISTENT FAMILY

\section{The actual ambient cusp-incompressibility producer}
\label{sec:hyperbolic-ambient-producer}

Fix a persistent compact core $Y\subset H$, one cusp torus $Z=\partial V$,
and the actual maps $f_t$ into the late slice $M_t^+$. Write
$\widehat Y_t=f_t(Y)$, $\widehat Z_t=f_t(Z)$ and
$\widehat W_t=M_t^+\setminus\operatorname{int}(\widehat Y_t)$, with the
other hyperbolic cores and component conventions recorded when needed.
KL91.2 concerns the actual map
\[
        (f_t|_Z)_*: \pi_1(Z,z)\longrightarrow
                    \pi_1(M_t^+,f_t(z)).
\]
The following inputs expose the analytic part still missing from AT01.

\begin{foundation}[Least-area compressing disks; HG10]
\label{found:hyp-compressing-disks}
If the actual cusp map is compressible, supply an exterior compressing
disk and the source's least-area embedded representative with free boundary
on $\widehat Z_t$. The exterior is compact, the relevant boundary is
convex at the selected late times, and the primitive boundary class is
fixed up to inverse through the recorded kernel comparisons. Supply the
topological disk theorem and the precise Meeks--Yau existence hypotheses.
\end{foundation}
Model peripheral injection, van Kampen and the boundary-kernel argument
enter before minimization. They do not establish ambient injection. The
exact variational interfaces must be compared with inherited results;
this contract does not order a new minimal-surface foundation.

\begin{foundation}[Avoidance of surgery regions; HG11]
\label{found:hyp-disk-surgery-avoidance}
Supply KL91.10 for the minimizing disks: for sufficiently late time no
such disk meets the centre of the specified surgery necks. The proof's
stable minimal-surface estimates, smooth extraction, elliptic regularity
and product-cylinder rigidity keep their full hypotheses. Supply the
common compact neighborhood needed to move each disk across a time
interval, including a surgery time when applicable.
\end{foundation}

\begin{foundation}[Cusp boundary terms; HG12]
\label{found:hyp-disk-boundary-terms}
For every $D>0$ choose a cusp diameter bound $a_0>0$; then for
$0<a<a_0$ choose the required late-time threshold. In the original
metric the minimizing disk boundary satisfies the KL91.11 bounds
\[
       \int_{\partial N_t} k_g\,ds\leq D/2,
       \qquad L_{g(t)}(\partial N_t)\leq (D/2)\sqrt t.
\]
All signs use the source's free-boundary and exterior conventions.
\end{foundation}
KL only sketches this estimate and refers to Hamilton1999 Sections11--12.
The detailed primary proof and the source's parameter dependence remain
qualified. The order is $D$, then $a$, then late time. No time uniform
over all arbitrarily small cusp choices is asserted.

\begin{foundation}[Continuous least area and upper barriers; HG13]
\label{found:hyp-area-barrier-input}
Under HG08 and HG10--HG12, the hypothetical least-area function $A(t)$ is
continuous and nonnegative for all sufficiently late times. At each $t_0$
it has a smooth local upper barrier $\overline A$ with equality at $t_0$
and
\[
 \overline A'(t_0)<\frac{3A(t_0)}{4(t_0+1/4)}-2\pi+D.
\]
This is KL91.12 together with its continuity argument, using the scalar
lower bound, Gauss--Bonnet, moving-boundary variation and HG11. It includes
surgery times; a calculation only on smooth intervals is insufficient.
\end{foundation}

\begin{lemma}[Upper-barrier area contradiction; HG14]
\label{lem:hyp-area-contradiction}
For $0<D<2\pi$ no such function $A$ can exist on a whole late half-line.
\end{lemma}
\begin{proof}
Put $s=t+1/4$ and
\[
 F(t)=s^{-3/4}A(t)+4(2\pi-D)s^{1/4}.
\]
Each upper barrier gives an upper barrier for $F$ with strictly negative
derivative at contact: differentiating cancels the two area terms and
the constant $2\pi-D$. The elementary continuous upper-barrier comparison
principle implies that $F$ is nonincreasing. For completeness, if
$F(b)>F(a)$, choose $0<\epsilon<(F(b)-F(a))/(b-a)$ and put
$G(t)=F(t)-\epsilon(t-a)$. Choose a level strictly between $G(a)$ and
$G(b)$. The last crossing of that level by $G$
before $b$ has larger values immediately to its right; its smooth upper
barrier has negative derivative there, a contradiction. Thus
$F(t)\leq F(T)$, whereas nonnegativity of $A$ gives
$F(t)\geq4(2\pi-D)(t+1/4)^{1/4}\to\infty$.
\end{proof}
The barrier comparison is an elementary real-analysis interface for
inherited reuse; no differentiability of the minimizing value $A$ was
assumed.

\begin{theorem}[Conditional ambient cusp injection; HG15]
\label{thm:hyp-ambient-cusp-injection}
With HG08 and the qualified inputs HG10--HG13, and with the actual common
core maps of AT11, every selected cusp torus is injective on fundamental
groups into the late slice, after the stated truncation and time choices.
\end{theorem}
\begin{proof}
If the kernel is nonzero, AT11 and its smooth-stage analog preserve it
along the given core maps. HG10 supplies the minimizing disks and HG11--HG13
their global late area function and barriers. Choose $D<2\pi$ and apply
HG14. This rules out a nonzero kernel at sufficiently late times; the
same kernel comparisons yield the source's stated earlier interval.
Apply the argument to the finite cusp family and take a common threshold.
\end{proof}
This conditional theorem supplies the exact AT01 map when bound to a flow.
It does not identify an abstract group with an ambient subgroup by name,
nor does it use Mostow rigidity as an incompressibility argument.

% BEGIN REVISION167 AMBIENT DISK DEVELOPMENT
\subsection{A prescribed meridian and the actual global area function}
\label{sec:ambient-fixed-meridian-development}

\paragraph{Selected improvement of the HG10--HG15 route.}
The endpoint is injection into the actual slice. We use a prescribed
smooth meridian curve, rather than minimize its boundary over a torus
homotopy class. This changes the intermediate area function, but supplies
exactly HG14's global continuous function and upper barriers. Thus the
free-boundary orthogonality and Hamilton length estimate in the historical
HG10/HG12 contracts are not premises of the following proof. They remain
valid subjects for a separate comparison with that route. Meeks--Yau is
used in its \emph{fixed-boundary} form. Surgery avoidance is proved by a
local stability radius estimate. No limiting minimal surface, its
noncompactness, or cylinder rigidity is needed by this version.

\paragraph{Sources and retained foundation.}
The flow/surgery conventions are those of KL, February 20, 2013,
Definitions 68.1 and 73.1 and Section 91, especially printed
2760--2762 and 2818--2823/PDF 174--176 and 232--237. Hamilton 1999,
Sections 11--12, printed 716--728/PDF 22--34, is the area comparison
source, with a different, volume-normalized time variable.
The static disk input here is Meeks--Yau, \emph{Math. Z.} 179 (1982),
Theorems 1--2 and Assertion 2, printed 153--157/PDF 4--8 of the
author-hosted copy. This is a different paper from KL's 1980 reference.
Its printed Theorem 1 suppresses the necessary nullhomotopy hypothesis;
we explicitly provide a spanning disk before invoking it. Its proof
uses Morrey and the earlier embeddedness theorem, which remain named
classical minimal-surface inputs, not proofs supplied by this chapter.
The smooth Loop Theorem, elementary surface topology, first and second
variation, Gauss--Bonnet, compact-domain elliptic spectral theory and
smooth isotopy extension are also retained foundation. Accepted PC
bindings for these precise signatures are not asserted.

The radius calculation below specializes the proof of Meeks--P\'erez--Ros,
\emph{Stable constant mean curvature surfaces}, author PDF dated
June 16, 2008, Lemma 2.1 and Theorem 2.8, printed/PDF 4--5 and 8--9,
with the stability convention on printed/PDF 14--15. It is written
locally on a disk to avoid any tacit completeness or compact-sphere
exception. Bamler, arXiv:1411.6649v2, Proposition 2.2, PDF 5--8,
confirms the prescribed-loop area method. We do not import that
proposition: it assumes stationary loops and trivial surgeries, while
our loops move and the surgery comparison below allows the actual
finite sphere-cutting history. Exact hashes, source differences and
the search scope are in \code{reference_checks_revision167.md}.

\begin{lemma}[Actual exterior kernels through surgery; IMS01]
\label{lem:ambient-exterior-kernels}
Choose finitely many fixed compact cores of the HPI08 models, including
all their cusp boundary components, and choose a late half-line on
which the cores and buffered collars are embedded. Write
\[
 Y_t=\coprod_i f_i(t)(Y_i),\qquad
 W_t=M_t^+\setminus\operatorname{int}(Y_t).
\]
For each boundary label $e$, both kernels of
$\pi_1(T_e)\to\pi_1(M_t^+)$ and $\pi_1(T_e)\to\pi_1(W_t)$
are constant along the family, using the component containing that
torus and recording all basepoint tracks. The same conclusion holds
when a fixed truncation is moved a finite distance down a supplied
cusp collar.
\end{lemma}
\begin{proof}
On a compact interval without surgery, extend the velocity of the
finite compact core/collar family to a smooth ambient vector field of
compact support. Its flow gives diffeomorphisms of the pairs and of
their exteriors. Compact interval support supplies the needed complete
trajectories. AT09--AT10 then give the kernel comparison.

At a surgery time, all these compact hyperbolic collars lie in the
unchanged region: on them scalar curvature is negative at sufficiently
late times, whereas the cutting necks, caps and discarded components
have positive scalar curvature at the surgery scale. HPI08 supplies
their actual unscathed space-time neighborhoods. Cut at the unchanged
spherical cross-sections in KL73.1. Each connected common piece
containing a marked torus includes injectively into the nearby
pre-surgery slice and into the post-surgery slice. Indeed gluing along
a sphere gives a free product; a nonseparating sphere additionally
gives a free generator. Filling a spherical boundary with a ball
does not change the fundamental group. In every case the vertex group
of that common piece injects. These assertions hold equally after
removing the hyperbolic cores: the sphere collars are disjoint from
them, so the same sphere gluing describes the exteriors. Discarded
pieces may change other factors but cannot remove the marked piece.
This verifies the actual geometric premises of AT11; it does not
infer a reverse map of whole surgery slices.

The pre-surgery regular locus has the compact sphere-cut core used
in KL73.1; its inclusion into the earlier smooth slice has the same
free-product description, as in KL73.4. Thus the argument includes
the singular-time limit, rather than treating it as a diffeomorphic
closed pre-surgery slice. Local finiteness gives finitely many events
on each compact time interval and hence constancy on the entire
half-line. A finite change of cusp depth is an ambient collar isotopy
on any time slice containing that larger compact collar. It identifies
the exteriors and supplies the corresponding basepoint homotopy.
\end{proof}

\begin{lemma}[Choose an exterior meridian before shrinking it; IMS02]
\label{lem:ambient-prescribed-meridian}
If an actual seam map fails to be injective at one time, then some
boundary label of $W_t$ has a nonzero kernel. One can select, once and
for all under that supposition, a primitive nonzero class
$\mu\in\pi_1(T_e)$ which bounds a disk in the appropriate exterior
at every later time. In cusp coordinates of curvature $-1/4$,
\[
 h=ds^2+e^{-s}g_0,
 \qquad \gamma_b(u)=(b,\gamma_0(u)),
\]
one may take $\gamma_0$ to be a flat geodesic representing $\mu$.
The curve $\gamma_b$ is simple, has length $e^{-b/2}L_0$ and has
ambient curvature norm $1/2$. Its actual image $\gamma_t=f_e(t)\gamma_b$
bounds a smooth parametrized exterior disk for every sufficiently late time.
The depth $b$ is fixed after $\mu$, not uniformly over all slopes.
\end{lemma}
\begin{proof}
All ports into the hyperbolic cores are injective by the model peripheral
theorem. If all exterior ports were also injective, the finite collared
gluing theorem of Chapter 15 would make every seam injective into the
assembled slice. Take the contrapositive. This selects \emph{some}
compressible exterior port; it need not be the originally failing
label. The graph-of-groups normal form includes cycles and disconnected
exteriors, so no successive-separating-seam assumption enters.

Apply the smooth Loop Theorem to that boundary torus in its compact
exterior component. Its essential simple disk boundary represents a
primitive lattice vector: cutting a torus along an essential simple
curve gives an annulus, and an arc joining the two sides closes to a
curve of algebraic intersection one. A nontrivial multiple cannot
have this intersection number. IMS01 preserves this chosen class.
No claim that every projection of a Lagrangian has rank one, nor any
identification of the entire torus kernel, is needed.

The straight closed curve of a primitive lattice vector is embedded
and is isotopic on the torus to the original disk boundary. Torus
isotopy extends across an exterior collar and takes a compressing disk
to one with exactly this curve as boundary. Moving from $s=0$ to
$s=b$ is a further ambient collar isotopy once the larger collar is
embedded. At each later time IMS01 gives nullhomotopy of the same simple
curve in the exterior. Relative smooth approximation of a nullhomotopy
gives a smooth parametrized spanning disk, fixed on that curve.
Embeddedness at these later times will follow from IMS03.
The warped-product connection gives
$\nabla_{\dot\gamma_b}\dot\gamma_b=\tfrac12\partial_s$ for unit
tangent, and the induced flat metric is $e^{-b}g_0$. These prove
the two geometric formulas. The finite number $L_0$ can depend on
the hypothetical meridian; it is fixed before $b$ is chosen.
\end{proof}

\begin{lemma}[Actual fixed-boundary Meeks--Yau application; IMS03]
\label{lem:ambient-fixed-boundary-minimizer}
For a fixed $b$ as in IMS02 and every sufficiently late time, including
post-surgery times, let $\mathcal D_t$ be the smooth embedded disks in
the compact exterior with boundary exactly $\gamma_t$ and interior
in its interior. Then
\[
 A(t)=\inf_{D\in\mathcal D_t}\operatorname{Area}_{g(t)}(D)
\]
is finite, strictly positive and attained by a smooth embedded minimal
disk. Every attaining disk is two-sided and stable for compactly
supported interior variations. No orthogonal boundary condition is
claimed.
\end{lemma}
\begin{proof}
All boundary components of this exterior are actual cusp tori. Its
outward normal points toward the removed core and is asymptotic to
$-\partial_s$; with convention
$\mathrm{II}(v,v)=\langle\nabla_v\nu,v\rangle$, the limiting
second fundamental form is $\tfrac12 g_T$. The $C^2$ convergence
from HPI08 makes it positive definite on every component at one
common late time. The manifold is compact and smooth, and the
selected curve is a smooth regular Jordan curve. IMS02 gives a
nonempty spanning-disk class in this very exterior. Thus the smooth,
strictly mean-convex case of Meeks--Yau Theorems 1--2 applies.
It gives a smooth least-area embedded disk with interior disjoint
from the boundary and no boundary branch point. The disk supplied
by the theorem minimizes among parametrized spanning disks, so its
area equals the infimum over embedded disks defined here.

In particular, any attaining embedded disk is a global minimizer
among the same fixed-boundary competitors. Interior first variation
makes it minimal, and second variation is nonnegative. The oriented
disk in an oriented three-manifold has a global unit normal, so this
is two-sided stability with all compactly supported test functions,
not merely volume-preserving stability. Smooth immersion on a
nonempty open subdisk gives positive area; the spanning competitor
gives finiteness. The post-surgery exterior is again a compact smooth
manifold with the same convex torus boundaries, so this application
also covers those times. Existence here is the precise retained
Meeks--Yau input applied to verified geometric hypotheses; it is not
an asserted new proof of that classical theorem.
\end{proof}

\begin{lemma}[Boundary estimates for the prescribed curve; IMS04]
\label{lem:ambient-prescribed-boundary-estimates}
For every $D>0$, the hypothetical meridian of IMS02 admits a fixed
depth $b$ and then a time $T$ such that every IMS03 minimizer has
\[
 \int_{\partial D_t} k_{D_t}\,d\ell<D/4,
 \qquad
 \int_{\partial D_t}|\partial_t\gamma_t|\,d\ell<D/4
 \qquad(t\geq T).
\]
Here $k_{D_t}$ is the inward-conormal geodesic curvature used in
Gauss--Bonnet. More explicitly, for any fixed $\ell>0$ one can
choose $b$ and then $T$ so that
\[
 L_{g(t)}(\gamma_t)<\ell\sqrt t,
 \quad |\kappa_{\gamma_t}|_{g(t)}<t^{-1/2},
 \quad |\partial_t\gamma_t|_{g(t)}<\alpha(t)t^{-1/2}.
\]
\end{lemma}
\begin{proof}
First choose $b$ with $e^{-b/2}L_0<\ell/2$. On a fixed buffered
collar of this curve HPI08 gives $C^2$ convergence of
$t^{-1}f_t^*g(t)$ to $h$, and its physical velocity estimate.
The length and connection formulas then give the displayed bounds
after increasing $T$. In particular curvature norms scale by
$t^{-1/2}$ and length by $\sqrt t$. On any smooth immersed disk
with this boundary, $k_{D_t}$ is the component of the ambient
curvature vector in its inward unit conormal. Hence
$k_{D_t}\leq|\kappa_{\gamma_t}|$ without an angle or orthogonality
assumption. The two integrals are at most $\ell$ and
$\alpha(t)\ell$, respectively. Choose $\ell<D/4$ and finally
$T$ so that $\alpha(t)<1$. This proves the strict inequalities.
All parameters are chosen after the fixed meridian. A small torus
diameter by itself is not a uniform bound for every primitive slope.
\end{proof}

\begin{lemma}[A local stability radius bound; IMS05]
\label{lem:ambient-stability-radius}
Let $\Sigma$ be a smooth two-sided stable minimal disk in an oriented
three-manifold. Suppose the closed intrinsic ball
$\overline B_\Sigma(p,r)$ is contained in its interior and ambient
scalar curvature is at least $\sigma>0$ on a neighborhood of that
ball. Suppose also that its intrinsic radius-$r$ sphere is nonempty.
Then
\[
                 r\leq 2\pi\sqrt{\frac{2}{3\sigma}}.
\]
No ambient curvature-derivative or injectivity bound is required.
\end{lemma}
\begin{proof}
Put $q=(R+|\mathrm{II}_\Sigma|^2)/2$. The Gauss equation changes
the stability form into
\[
 \int_\Sigma\bigl(|\nabla v|^2+(K_\Sigma-q)v^2\bigr)\geq0.
\]
Choose a connected relatively compact smooth domain $\Omega$ in
$\operatorname{int}\Sigma$ containing the closed ball and a small
neighborhood of it. Its first Dirichlet eigenvalue for
$-\Delta+K_\Sigma-q$ is $\lambda\geq0$, by that form. The
positive first eigenfunction $u$ satisfies
$\Delta u=(K_\Sigma-q-\lambda)u$ and is smooth and strictly
positive on a neighborhood of the closed ball. This local spectral
construction avoids a completeness hypothesis for a global Jacobi
function. Only on the closed ball do we use $q\geq\sigma/2$.

For the metric $\widetilde g=u^2g_\Sigma$, minimize length from
$p$ to the intrinsic radius-$r$ sphere, staying in the closed ball.
Compactness and the positive minimum of $u$ there give a minimizer.
Stop at its first encounter with that sphere; its interior is a
smooth minimizing $\widetilde g$-geodesic. Parameterize by original
arc length $s\in[0,L]$. Then $L\geq r$. Write $v=\log u$.
The conformal curvature formula gives
\[
 \widetilde K=u^{-2}(q+\lambda+|\nabla v|^2).
\]
Second variation of this geodesic, with fixed endpoints and normal
test field $\phi=\sqrt u\,\psi$, gives
\begin{align*}
 0&\leq\int_0^L u^{-1}
       \bigl((\phi')^2-(q+\lambda+|\nabla v|^2)\phi^2\bigr)\,ds\\
  &\leq\int_0^L\bigl((\psi')^2+v'\psi\psi'
                   -\tfrac34(v')^2\psi^2-\tfrac\sigma2\psi^2\bigr)\,ds\\
  &\leq\int_0^L\bigl(\tfrac43(\psi')^2
                                  -\tfrac\sigma2\psi^2\bigr)\,ds.
\end{align*}
The last step completes the square
$-\tfrac34(v'\psi-\tfrac23\psi')^2\leq0$.
Use $\psi(s)=\sin(\pi s/L)$, or compactly supported approximations
with the same limiting integrals. We obtain
$4\pi^2/(3L^2)\geq\sigma/2$ and the asserted bound.
The nonempty sphere is explicit: applying this argument to an entire
closed stable sphere would not provide the required endpoint and
would be invalid. In the applications below the surface is a compact
disk whose boundary lies farther away than $r$.
\end{proof}

\begin{lemma}[Quantitative exclusion from long positive neck bands; IMS06]
\label{lem:ambient-neck-exclusion}
In a normalized neck band $S^2\times(-20,20)$ suppose $R\geq1/2$
and $|dz|\leq2$. A compact embedded stable minimal disk whose
boundary is outside the closed band cannot meet its middle sphere.
The same assertion holds on any translated band and after constant
metric rescaling. It applies uniformly to all IMS03 minimizers on
the unchanged long parts of sufficiently late surgery necks and on
nearby smooth slices.
\end{lemma}
\begin{proof}
If a disk contains a point $p$ on the middle sphere, any path on the
disk from $p$ to its boundary must first exit the neck band. Until
then $|dz|\leq2$, so that path has length at least $10$ in the
normalized metric. In fact the closed intrinsic radius-$8$ ball
about $p$ stays in the open band, is compact and is disjoint from
the disk boundary. Its radius-$8$ sphere is nonempty: take a path
to the disk boundary and use continuity of intrinsic distance.
IMS05 would give
$8\leq 4\pi/\sqrt3<8$, a contradiction. This proves exclusion
without assuming that intrinsic disk balls equal ambient balls.

On a sufficiently small $\delta$-neck, scale by $Q=R(x)$ so
that the cylindrical model has scalar curvature one. $C^2$
closeness gives the strict inequalities $R>1/2$ and $|dz|<2$
on every fixed buffered band. The actual torus boundaries have
negative scalar curvature, so they are outside these positive
bands. Constant scaling preserves minimality and stability and
changes both intrinsic radii and scalar curvature by the reciprocal
factors in IMS05. At a fixed surgery event the unchanged compact
bands persist smoothly for a short interval on both sides; their
strict estimates persist there as well. Taking a minimum interval
over the finite neck set gives a common interval. The argument
applies to every minimizing disk, independently of its area.
\end{proof}

\begin{lemma}[One unchanged compact carrier for all nearby minimizers; IMS07]
\label{lem:ambient-common-disk-carrier}
For every sufficiently late time $t_0$, including a surgery time,
there are a compact carrier $K$ in the unchanged region and a small
open interval $I$ about $t_0$ such that all IMS03 minimizing disks
at times in $I$ lie in the interior of the copies of $K$, apart
from their prescribed boundary on the moving torus. The metric on
a neighborhood of $K$ is smooth across $t_0$. There are smooth
diffeomorphisms of neighborhoods, $P_t$, with $P_{t_0}=\mathrm{id}$,
which take every cusp boundary and its prescribed curve at $t_0$
to the corresponding data at $t$, preserve the exterior side,
and are the identity near the artificial sphere boundaries of $K$.
\end{lemma}
\begin{proof}
At a regular time use the compact slice and a surgery-free interval.
At a surgery time use the coordinates in KL73.1: the unchanged
part of each neck is $[0,\delta^{-1})\times S^2$, and the new
cap and all metric changes lie on the other side of $z=0$.
Since $\delta(t)\to0$, increase the starting time so that these
coordinates include $[0,100]\times S^2$. For every $z_0\in[40,60]$
the band $(z_0-20,z_0+20)\times S^2$ satisfies IMS06, with a
small common margin. All minimizing disks therefore avoid the
entire closed slab $[40,60]\times S^2$. The same is true for
all nearby times, after shortening the interval to preserve the
strict metric estimates on the finite collection of buffered bands.

Cut all necks at $z=50$ and keep the common compact pieces on
the low-curvature side which contain the marked cusp boundaries.
This is a compact subset of the pre-surgery regular locus, and
it is unchanged by surgery. The torus boundary is on this side:
the discarded positive-curvature region and cap side do not meet
the persistent negative-curvature core collars. A connected disk
with that boundary cannot enter any removed horn, new cap or
other discarded piece without crossing a cutting sphere. IMS06
excludes each crossing. The excluded slab gives a uniform open
buffer from the artificial boundary, not merely avoidance of a
single two-sphere. Finitely many retained components may be used
for $K$; only the component containing a given disk matters.
KL's smooth unscathed identification on this compact subset gives
a smooth metric on its neighborhood across the event.

On the finite compact cusp collars use the actual family $f_t$:
locally $f_t f_{t_0}^{-1}$ moves the collars and the parametrized
curves. Extend its velocity using a cutoff to a vector field
supported in a slightly larger collar, still away from all
artificial spheres. Integrate on a small common interval. The
resulting $P_t$ takes the marked hypersurfaces and curves to their
actual positions, preserves their sides, and fixes the artificial
boundaries. Compact support ensures all trajectories exist.
The same construction applies to the finitely many other cusp
boundaries of the exterior. This proves the simultaneous statement.
\end{proof}

\begin{lemma}[Continuity of the actual minimum across surgery; IMS08]
\label{lem:ambient-area-continuity}
The IMS03 function $A$ is continuous on a sufficiently late half-line,
including every surgery time. No continuous choice of minimizing disk
is required.
\end{lemma}
\begin{proof}
Fix $t_0$ and use IMS07 to identify $K$ with its copies at nearby
times. On $K$, the metrics $P_t^*g(t)$ converge smoothly to $g(t_0)$.
For every $\epsilon>0$, on a sufficiently small interval their
quadratic forms lie between $e^{-2\epsilon}g(t_0)$ and
$e^{2\epsilon}g(t_0)$. Two-dimensional areas consequently differ
by factors at most $e^{2\epsilon}$.

Transport a minimizer at $t_0$ forward by $P_t$. It is an admissible
embedded exterior disk with the exact prescribed boundary at $t$;
hence $A(t)\leq e^{2\epsilon}A(t_0)$. Transport any minimizer
at $t$ back by $P_t^{-1}$, which is possible because \emph{every}
such minimizer lies in the common $K$. This gives
$A(t_0)\leq e^{2\epsilon}A(t)$. The two comparisons prove
continuity. They are valid on both sides of a surgery time because
all disks and transports lie in the same unchanged metric region.
No global diffeomorphism of the surgery slices, no assumption of
trivial surgery, and no convergence theorem for minimizers is used.
\end{proof}

\begin{lemma}[Actual smooth local upper barriers; IMS09]
\label{lem:ambient-area-upper-barriers}
For every $D>0$ choose the fixed meridian depth and late threshold
as in IMS04, enlarging that threshold for IMS07. At each $t_0$
there is a smooth local function $\overline A\geq A$, with equality
at $t_0$, such that
\[
 \overline A'(t_0)<\frac{3A(t_0)}{4(t_0+1/4)}-2\pi+D.
\]
This includes post-surgery contact times and a two-sided neighborhood
in time of each such contact.
\end{lemma}
\begin{proof}
Choose an IMS03 minimizer $\Sigma$ at $t_0$. Define
$\overline A(t)=\operatorname{Area}_{g(t)}(P_t\Sigma)$ using
IMS07. These are smooth embedded admissible disks on a common
unchanged compact neighborhood; thus the area is smooth, dominates
the infimum and agrees with it at contact. Let $V$ be the variation
field at $t_0$ and $\eta$ the disk's outward conormal. At its
boundary $V$ agrees with the velocity of the parametrized curve.
First variation and $\partial_tg=-2\Ric$ give
\[
 \overline A'(t_0)
 =-\int_\Sigma\operatorname{tr}_{T\Sigma}\Ric\,dA
                     +\int_{\partial\Sigma}\langle V,\eta\rangle\,d\ell.
\]
The interior motion term vanishes because $\Sigma$ is minimal.
For a minimal surface in dimension three the Gauss equation gives
\[
 -\operatorname{tr}_{T\Sigma}\Ric
       =-\tfrac12R-K_\Sigma-\tfrac12|\mathrm{II}_\Sigma|^2
       \leq-\tfrac12R-K_\Sigma.
\]
Use the actual scalar lower bound
$R\geq-3/[2(t_0+1/4)]$ and
$\int_\Sigma K_\Sigma+\int_{\partial\Sigma}k_\Sigma=2\pi$.
The resulting estimate is
\[
 \overline A'(t_0)\leq
 \frac{3A(t_0)}{4(t_0+1/4)}-2\pi
 +\int_{\partial\Sigma}k_\Sigma\,d\ell
 +\int_{\partial\Sigma}|V|\,d\ell.
\]
IMS04 bounds the last two terms by $D/2<D$, providing strict
slack. The unchanged region satisfies the Ricci equation through
the surgery time by smooth restriction of the flow there; no
derivative of a jump in the ambient topology is being taken.
This establishes the full HG13-type producer for the prescribed
boundary area function.
\end{proof}

\begin{theorem}[Actual ambient cusp injection; IMS10]
\label{thm:ambient-actual-cusp-injection}
Under the retained foundation stated above and the actual HPI08
family, every actual fixed cusp interface is injective on fundamental
groups into the late flow slice, throughout its sufficiently late
persistent interval. This supplies HG15's conclusion and AT01 for
the actual maps. The conclusion also holds for the actual slowly
moving truncations of TCF01 once their collars include the fixed
reference tori.
\end{theorem}
\begin{proof}
Choose one late half-line containing the finite family of fixed
core collars and satisfying IMS01. Suppose some seam map at some
time in it is not injective. IMS02 selects an exterior boundary
label and a primitive nonzero class which remains compressible at
every later time. Fix $D=\pi$. Choose its cusp depth only after
this class is fixed, then enlarge the time threshold for the
finitely many compact collars, IMS03--IMS04 and IMS06--IMS07.
All choices are finite and eventual; no supremum over possible
meridians or surgery times is taken.

IMS03 and IMS08 now supply a nonnegative continuous area function
on the entire resulting half-line. IMS09 supplies exactly the
strict upper barriers of HG14. HG14 gives a contradiction.
Consequently every exterior port, and hence by finite collared
gluing every actual seam, is injective. Equivalently, the initial
supposition that any seam failed is impossible. IMS01 transports
the conclusion back over the fixed persistent interval.

At a time with a TCF01 truncation, the segment of the actual
embedded cusp between the fixed reference torus and the selected
torus gives a homotopy of their inclusion maps, with its basepoint
track recorded. These maps therefore have conjugate images and
the same kernel after the product identification of the two tori.
The selected torus is injective. This is a pointwise collar
comparison at each time; a new time-dependent minimizing problem
for the diverging depth $S(t)$ is unnecessary. Take the maximum
of the finitely many eventual domain thresholds. Empty hyperbolic
families and closed hyperbolic components have no cusp ports and
require no disk argument. No graph-manifold property of the exterior
is used to establish its boundary injections.
\end{proof}

\paragraph{Precise completion scope before the final audit.}
IMS01--IMS10 provide the actual geometric producers and the global
area function through surgery, relative to the explicitly retained
classical foundation. The historical free-boundary versions of
HG10 and HG12 have been replaced for this endpoint, not asserted
proved verbatim. In particular the independent Meeks--Yau formalization
remains a prerequisite at implementation time. The final mathematical
audit of this new ambient-incompressibility goal has not yet occurred.
The earlier analytic audit166 and static audit159 are separate and
remain in force. Essential graph refinement and reconstruction to the
initial manifold remain subsequent work.
% END REVISION167 AMBIENT DISK DEVELOPMENT

% BEGIN REVISION168 FINAL AMBIENT AUDIT
\subsection{Final audit of actual ambient cusp incompressibility}
\label{sec:ambient-final-mathematical-audit}

This is the ONE final mathematical audit requested for the
ambient-incompressibility goal begun from revision166. It checks
IMS01--IMS11, their HG14 consumer and their actual HPI08/TCF01/AT01
maps. It is separate from the completed analytic audit166 and
static audit159. Two hypotheses are made explicit below: compactness
of the intrinsic ball in the radius proof, and the small neck
accuracy required of the single surgery profile. We also spell out
the carrier identifications used in both directions at surgery.
These clarifications supersede the corresponding shorthand in
revision167 without changing its fixed-meridian proof route.

\begin{lemma}[Compact-ball interpretation of the stability estimate; IAU01]
\label{lem:ambient-audit-compact-radius}
In IMS05 require that the closed intrinsic ball be compact, in
addition to being contained in the surface interior and having a
nonempty radius-$r$ sphere. This hypothesis is automatic for every
IMS03 disk used in IMS06. The radius estimate and the exclusion
$8\leq4\pi/\sqrt3<8$ then follow with no completeness assumption
on a larger minimal surface.
\end{lemma}
\begin{proof}
An IMS03 disk is the smooth embedded image of a compact closed
disk, with a smooth positive definite induced metric up to its
boundary. Its intrinsic metric induces its compact topology.
Thus a closed intrinsic ball is compact. In IMS06 the Lipschitz
bound for $z$ makes its distance to the disk boundary at least
$10$, so the closed radius-$8$ ball lies in the surface interior.
A path to the nonempty disk boundary supplies a point at distance
$8$. A small neighborhood of this compact ball still lies in the
interior; a finite coordinate cover and a regular level of a smooth
cutoff give the relatively compact smooth domain $\Omega$ used
in IMS05, taking its connected component containing the ball.

On that domain stability makes the first Dirichlet eigenvalue
nonnegative. The positive first eigenfunction has a positive
minimum on the compact ball, so its conformal metric is uniformly
comparable to the original metric there. A minimizing sequence
of curves from the center to the compact nonempty distance sphere
has uniformly bounded original length. Arc-length
reparametrization, compactness and lower semicontinuity give a
minimizer. Until its first contact with the sphere it is an
unconstrained smooth conformal geodesic. Compactly supported
normal variations on that open segment remain inside the ball.
The fixed-endpoint second variation therefore applies; the sine
test function follows by approximation in $H^1_0(0,L)$.
The IMS05 calculation gives
$\sigma/2\leq4\pi^2/(3L^2)$ with $L\geq r$.
This verifies every localization used there. Merely saying that a
ball in an arbitrary incomplete open surface is closed would not
supply the compactness step; that stronger interpretation is not
used here. Neither a noncompact pointed disk limit nor the
classification of complete stable surfaces is an input.
\end{proof}

\begin{lemma}[One surgery profile and exact common carriers; IAU02]
\label{lem:ambient-audit-profile-carrier}
There is a universal neck accuracy $\delta_*>0$ sufficient for the
IMS06--IMS09 constructions. Their flow hypothesis is that the
single admissible profile fixed before constructing the flow
satisfies $\delta(t)<\delta_*$ eventually. The standard choice
$\delta(t)\to0$ suffices. With that hypothesis the compact carrier
and its boundary-moving isotopy in IMS07 give both area comparisons
in IMS08 and the two-sided smooth barrier in IMS09, including at
singular times.
\end{lemma}
\begin{proof}
Use the scalar-one cylindrical normalization of KL72.24. On the
fixed model band $S^2\times[10,90]$, scalar curvature is $1$ and
$|dz|=1$. Continuity of the inverse metric and curvature in the
$C^2$ topology gives a number $\eta>0$ such that $\eta$-closeness
implies $R>1/2$ and $|dz|<2$ on that entire band. Choose
$\delta_*<\min\{1/200,\eta\}$, decreased if necessary for the
source's neck-chart norm convention. A $\delta$-neck of that
accuracy has the required derivatives and contains $[0,100]$.
Thus every translated length-$40$ band centered in $[40,60]$
satisfies IMS06 simultaneously. Only one fixed compact band's
$C^2$ norm is tested, not a separate time interval for infinitely
many centers.

This is an explicit condition on the surgery profile, not a
consequence of nonincrease alone. The retained existence theorem
permits a smaller positive profile before the flow is constructed.
For example its bound can be replaced, at that construction stage,
by the running minimum with $\delta_*/2$ and $(t+2)^{-1}$,
together with the already required KL85/86.9 bounds. The actual
HPI08 and thin-carrier producers then use that same fixed flow.
No profile is changed after selecting a disk, meridian or surgery
event. This makes the standard late-profile convention in IMS07
precise; an arbitrary surgery flow with only nonincreasing
$\delta$ would require the additional smallness hypothesis.

At a fixed event, cut the unchanged necks at $z=50$ and let $K$
be the union of the compact retained pieces containing the marked
cores and their cusp collars. Here take the whole pieces, before
removing the hyperbolic interiors. Every relevant minimizing disk
is contained in $K$ and avoids its artificial boundary by the
open slab buffer. The marked collars lie in $\operatorname{int}K$.
Use KL68's actual static identifications on a neighborhood of $K$;
its metric is smooth through the event. This follows from the
regular-spacetime assertion following Definition68.1, printed
2746--2748/PDF160--162, not from an identification of entire
pre- and post-surgery manifolds.

The cutoff vector field can be supported in
$\operatorname{int}K$, and its flow $P_t$ consequently maps $K$
onto itself, fixing its artificial boundary. It maps all marked
cores, exterior sides and the parametrized prescribed curve to
their actual data at $t$. Hence
\[
 P_t\bigl(K\cap W_{t_0}\bigr)=K\cap W_t.
\]
All minimizers at $t_0$ and at nearby times lie in these sets.
Both $P_t$ and $P_t^{-1}$ therefore preserve precisely the disk
competitor classes needed for the two inequalities in IMS08.
No minimizing disk is asserted to exist on the possibly noncompact
backward singular slice: at the event $A(t_0)$ uses the compact
\emph{post}-surgery exterior, while for $t<t_0$ it uses the
nearby compact smooth exterior. The common $K$ connects them.
The same transport of a single contact minimizer supplies the
smooth upper barrier on an open interval across the event.
There is no need for a uniform neighborhood size over all events;
local finiteness and eventwise neighborhoods suffice.
\end{proof}

\begin{theorem}[Audited actual cusp and carrier injections; IAU03]
\label{thm:ambient-audited-written-endpoint}
For the actual persistent family and late thin decomposition, on
the single admissible surgery flow with IAU02's profile convention,
every fixed reference cusp torus and every corresponding TCF01
interface is injective on fundamental groups into its actual
late slice component. AT01 is supplied for those maps; AT02 then
gives injection of the actual core and thin-carrier maps. This is
a written mathematical endpoint under the explicitly retained
classical foundations, not an accepted Lean theorem.
\end{theorem}
\begin{proof}
The finite collared gluing theorem gives the following useful
contrapositive: a failed ambient seam injection forces some
exterior boundary injection to fail. That boundary need not be
the original label. Conversely, one can start the IMS02 construction
from \emph{any} nonzero exterior kernel. The Loop Theorem then
selects one essential simple meridian. IMS01 preserves its kernel
in the actual exterior through smooth stages and sphere surgeries;
all common-piece inclusions are the actual ones and all basepoint
tracks are retained. Discarded components cannot contain a marked
negative-scalar collar. There is no rank-one projection assertion
and no assumed identification of a torus kernel with a model group.

Fix that primitive meridian first. Choose $D=\pi$, then its finite
cusp depth, then a late threshold for all finitely many core and
collar requirements. The flat representative has exponentially
small normalized length, and normalized ambient curvature $1/2$.
IMS03 applies the fixed-boundary Meeks--Yau theorem to the actual
compact mean-convex exterior and the already nullhomotopic regular
Jordan curve. Its minimizer is a smooth properly embedded disk;
therefore the compactness convention in IAU01 applies. IMS04
bounds both boundary terms independently of the disk's angle
with the torus. IMS06 and IAU02 exclude every minimizer from the
actual surgery bands and place it in the common unchanged carrier.

IMS08 gives continuity of the attained nonnegative minimum over
the entire late half-line, not just on smooth intervals. IMS09
gives at each contact a smooth upper barrier with
\[
 \overline A'<\frac{3A}{4(t+1/4)}-\pi.
\]
The scalar bound is KL79.11 in physical time (printed 2780,
PDF 194); the cusp normalization is $t^{-1}g(t)$ with
sectional curvature $-1/4$. Gauss--Bonnet uses the inward disk
conormal, while boundary motion uses the outward conormal in
first variation. The signs in IMS09 follow with these conventions.
HG14 now forbids the global area function. This rules out every
hypothetical nonzero exterior kernel, and the finite gluing theorem
therefore proves all ambient seam injections. IMS01 also carries
the conclusion back along each fixed persistent interval.

The actual embedded cusp product at each time identifies the
reference torus inclusion with the selected TCF01 inclusion,
up to its stated basepoint path. Thus the actual moving interface
has the same kernel. IMS11 binds precisely these maps at the
TCF07 common time, and simultaneous gluing supplies the carrier
injections. We do not cancel a carrier map from an injective
composite. Empty cusp families and closed hyperbolic components
retain their vacuous or identity cases. This completes the
requested written application through surgery.
\end{proof}

\paragraph{Audit result and retained implementation obligations.}
The single final audit is complete with the compact-ball and
profile hypotheses clarified and the common-carrier comparison
made explicit. The review and requirement matrix are
\code{AUDIT_AMBIENT_REVISION168.md} and
\code{ambient_goal_audit_revision168.csv}. The static fixed-boundary
Meeks--Yau theorem, including its Morrey/embeddedness/regularity
foundation, remains a named input awaiting its independent
formalization. The smooth Loop Theorem, variation, Gauss--Bonnet,
elliptic spectral and smooth/topological foundations also require
precise accepted-release bindings. This chapter proves their actual
geometric application; it does not declare them Lean-verified.
No theorem about free-boundary minimizers or a uniform bound over
all slopes is inferred. Essential graph refinement, model conversion,
finite-history reconstruction and the full GC endpoint remain
subsequent work. The superseded ``audit pending'' and earlier open
frontiers are historical development notes.
% END REVISION168 FINAL AMBIENT AUDIT

\section{Complete interiors and the Chapter 12 handoff}
\label{sec:hyperbolic-producer-handoff}

\begin{lemma}[Opening a truncated core to the full model; HG16]
\label{lem:hyp-whole-interior}
HG02's actual smooth cusp and collar data give a diffeomorphism
$\psi:C^\circ\to H$. Its pullback $\psi^*h$ is complete, hyperbolic and
has the same finite volume as $(H,h)$.
\end{lemma}
\begin{proof}
In each of the finitely many disjoint boundary collars replace a finite
open interval by a half-infinite end. More explicitly, an increasing
diffeomorphism $\theta:(-\epsilon,0)\to(-\epsilon,\infty)$ can equal the
identity near $-\epsilon$: integrate a positive smooth function equal to
$1$ there whose integral diverges as the argument tends to $0$.
Use the same torus coordinate, and use the identity off the collars.
The maps and inverses agree on their overlaps and so define $\psi$.
For a closed model take the identity. Pullback makes $\psi$ an isometry;
lengths, Cauchy convergence, local model charts and the volume measure
transport as in AT26. This is not restriction of $h$ to a finite core.
\end{proof}

\begin{theorem}[Conditional thick certificate producer; HG17]
\label{thm:hyp-thick-certificate}
For one admissible common late time, HG08, HG15 and HG16 give finite
hyperbolic core records, pairwise disjoint late embeddings, actual ambient
boundary injections and complete finite-volume interior witnesses.
The core embeddings, full-interior diffeomorphisms and boundary labels
are retained as distinct maps. Placement in chosen prime carriers and
transport to the initial manifold are later Chapter15 operations.
\end{theorem}
\begin{proof}
Choose the finite truncations first, then a time containing their required
collars and satisfying the finite family of thresholds. Apply HG16 to
each model and HG15 to each boundary map, retaining HG08's actual
embeddings and disjointness. Store these outputs on those same carriers.
No new ambient injection or full-interior metric is inferred from a tag.
\end{proof}

\begin{definition}[Selected proof route and remaining bindings; HG18]
\label{def:hyp-route-selection}
For the detailed stabilization dependency map use KL90's smooth
compactness and controlled smoothing route. Hamilton1999 and MSM206 remain
the geometric/analytic guides and comparison route. Harmonic-map stability
is a possible alternative implementation, not an additional compulsory
prerequisite of this selected proof. Both routes require finite-volume
rigidity. Morgan--Tian2014 remains the global spine; the inspected MT2.11
body is the archived2012 manuscript, not a verified published-book crosswalk.
\end{definition}
HG03 now has the written Margulis/cusp proof in
Section~\ref{margulis:section}, relative to its explicit inherited
interfaces. HG04, HG06--HG08 and HG10--HG13 remain unfinished
source-qualified producers. The static HG03 proof does not discharge
their rigidity, flow, quantitative or ambient-injection obligations.
Accepted inherited bindings and Lean formalization remain open for all. Chapter12 must bind HG07 and HG17 to one actual flow and
normalization; Chapter16 only consumes the supplied model witnesses.

\chapter{Ricci-flow PDE interfaces and surviving gaps}
\label{ch:ricci-flow-pde}

\paragraph{Scope.} Only the PDE contracts that survive completed-foundation comparison.

\paragraph{Planned sections.}
\begin{enumerate}[leftmargin=*]
\item Time-dependent metrics.
\item Equation.
\item Solution regularity.
\item Existence and uniqueness.
\item Evolution.
\item Convention adapters.
\end{enumerate}

\paragraph{Existing material and present task.} Existing RicciFlowPDE contract and historical T09-T16 discoveries in Appendix C. Write the mathematical input/output chain; retain short-time, scalar and Hamilton-Ivey results as inherited candidates.

\section{Inherited interfaces before gap classification}
The existing \code{RicciFlowPDE} inventory row is the contract anchor.
The detailed historical discoveries for short-time solutions, scalar
comparison, curvature evolution and tensor estimates remain in
Appendix~\ref{ch:m2a-evidence-appendix}. Write a mathematical definition
and input/output worksheet for each consumer without binding it to a
changing declaration name. A result becomes a genuine gap only after
the completed-release comparison; no absence claim follows from this
chapter's current outline.

\chapter{Perelman estimates and geometric compactness}
\label{ch:perelman-estimates}

\paragraph{Scope.} Estimate and compactness interfaces used by surgery and long-time analysis.

\paragraph{Planned sections.}
\begin{enumerate}[leftmargin=*]
\item Functionals.
\item Monotonicity.
\item Reduced geometry.
\item Noncollapsing.
\item Curvature control.
\item Smooth pointed-flow compactness.
\end{enumerate}

\paragraph{Existing material and present task.} Existing M4 inventory and inherited estimate candidates. Separate static metric compactness from smooth flow compactness and list exact downstream consumers.

\section{Estimate-to-consumer map}
For each estimate, record which surgery or long-time consumer needs it,
with its time domain, normalization, scale, regularity and surgery
compatibility. Reuse the existing PC candidates and source-role ledger.
Perelman's papers remain primary; Kleiner--Lott supplies detailed analytic
expansion and Morgan--Tian the global assembly comparison.

\section{Two kinds of compactness}
Chapter~\ref{ch:metric-geometry} concerns static metric approximations.
The present chapter owns the smooth pointed-flow compactness interface
and the estimate hypotheses supporting its use. The existing
\code{HasSmoothCheegerGromovLimit} row is not replaced by metric
Gromov--Hausdorff convergence. Exact bridges remain separate targets.

\chapter{Ricci flow with surgery}
\label{ch:surgery-topic}

\paragraph{Scope.} The typed flow/event input and its continuation and invariant interfaces.

\paragraph{Planned sections.}
\begin{enumerate}[leftmargin=*]
\item Canonical neighborhoods.
\item Event geometry.
\item Caps.
\item Continuation.
\item Smooth stages.
\item Bounded-time history.
\item Analytic producer boundary.
\end{enumerate}

\paragraph{Existing material and present task.} Existing producer boundary; T17-T21 records remain in Appendix C. Describe event interfaces and convention obligations without probing the unfinished PC snapshot.

% BEGIN migrated B023
\section{Analytic producer boundary and PC-reuse audit}
\label{sec:analytic-producer-boundary}

M1f fixes the endpoint consumed by analysis. The next contract is the reverse
direction: a completed Ricci-flow-with-surgery development must produce an
\code{LongTimeOutcome}, with its late branch yielding an M1e
\code{ThickThinCollapsePackage}. This revision separates the PC results that
are imported from the genuinely GC-specific long-time output.

\subsection{The producer theorem}

The analytic workstream has one public output theorem:
\begin{equation}
\begin{aligned}
 &\texttt{long\_time\_outcome} : \\
 &\quad \texttt{RicciFlowSurgeryPackage M} \\
 &\quad\longrightarrow \texttt{LongTimeOutcome M}.
\end{aligned}
\end{equation}

Together with M1f it yields the final conditional composition:
\begin{equation}
\begin{aligned}
 &\texttt{PoincareAdapter} \\
 &\quad\longrightarrow \texttt{RicciFlowSurgeryPackage M} \\
 &\quad\longrightarrow \texttt{Geometrizes M}.
\end{aligned}
\end{equation}

\code{long_time_outcome} is the global analytic frontier. Its late branch includes
the long-time thick-region hyperbolic limits, the local-collapse theorem for
the thin region, the graph-manifold conversion, and the topology transport
through surgery. It does not include the M1 assembly proof.

\subsection{Fields of \code{RicciFlowSurgeryPackage}}

The producer input is an interface with six groups of fields:

\setlength{\tablebodywidth}{\dimexpr\textwidth-4\tabcolsep\relax}
\begin{longtable}{@{}P{0.280899\tablebodywidth}P{0.471910\tablebodywidth}P{0.247191\tablebodywidth}@{}}
\textbf{Group} & \textbf{Required content} & \textbf{Origin}\\\hline
\endhead
Initial data & A closed connected orientable smooth 3-manifold $M$, an initial
Riemannian metric, and normalization conventions. & PC core plus GC adapter.\\
Smooth-flow data & Short-time existence, uniqueness, evolution identities,
continuation, compactness, derivative estimates, and rescaling interfaces. &
Primarily direct PC reuse; audit required.\\
Surgery data & Typed finite or locally finite surgery history, canonical
neighborhood alternatives, surgery scales, continuation maps, and invariant
preservation. & Completed PC release, adapted at the public boundary.\\
Perelman data & Reduced volume/entropy monotonicity, noncollapsing, curvature
scale control, and three-dimensional pinching. & Direct PC reuse or a named
GC adapter.\\
Long-time prerequisites & Time and rescaling conventions, admissible
parameters and already proved estimates on the supplied flow. No
long-time outcome or region-certificate producer is an input field.
& Audited inputs to M6; conclusions are proved separately.\\
Audit data & Exact public theorem names, imports, hypotheses, and transitive
axiom closures for every reused component. & PCReuseAudit ledger.\\
\end{longtable}

The structure may contain imported fields or references to imported packages.
It contains no \code{LongTimeOutcome}, \code{LateThickThinOutcome},
\code{ThickThinCollapsePackage}, or function returning those conclusions.
Pointed hyperbolic limits, graph-manifold conclusions and their global
transport are outputs of M6, not fields making its producer a projection.
Every bounded time interval has a finite surgery-event set, and a selected
slice has a finite preceding history; no globally finite event count is
required. Maximal continuation and the exceptional alternative must be
proved with their full hypotheses, not inferred merely from failure of
arbitrarily late nonempty slices.
The structure must not hide a theorem statement as an unreviewed axiom. A temporary
sandbox version is allowed during interface development; the release target
requires every field to be proved or explicitly approved under the project
trust policy.

\subsection{PC reuse classes}

Every candidate PC declaration is assigned one of three classes after the
completed release is pinned:

\begin{description}[leftmargin=2.9cm,style=nextline]
\item[Direct] The public declaration has compatible hypotheses and conclusion.
GC imports it without changing its mathematical content, records its exact
axiom closure, and writes only notation adapters if needed.
\item[Adapter] The PC declaration is mathematically sufficient but uses a
  different manifold, time, surgery-history, or normalization type. GC proves
  a translation theorem and keeps the imported implementation out of its
  public endpoint.
\item[GC-new] No completed PC declaration supplies the required conclusion.
  GC formalizes the theorem, using PC results only for its audited analytic
  inputs.
\end{description}

The class is assigned per theorem, not per directory. A large imported module
may contain both direct inputs and declarations that are unsuitable for GC.

\subsection{PCReuseAudit gate}

The GC project starts its analytic phase only after the following audit rows
are complete for the pinned PC release:

\begin{enumerate}[label=\textbf{R\arabic*:},leftmargin=*]
\item public import path and exact declaration name;
\item full theorem signature, including manifold category, regularity, time
  domain, compactness, and normalization hypotheses;
\item a minimal adapter theorem, when the class is \code{adapter};
\item build target and a clean root-build result;
\item \code{#print axioms} output and transitive provenance; and
\item final classification as \code{direct}, \code{adapter}, or \code{GC-new}, with the
  downstream consumer in the M2--M6 plan.
\end{enumerate}

The first priority rows are the global flow-with-surgery object,
canonical-neighborhood theorem, surgery-event and continuation interfaces,
invariant-preservation theorem, flow compactness, reduced-volume/noncollapsing
results, Hamilton--Ivey pinching, and the standard-cap/cut-cap topology. The
finite-extinction tree is not a required row and remains outside the default
GC route.

\subsection{Producer work packages after the audit}

Once PCReuseAudit passes, the analytic work splits as follows:

\begin{itemize}[leftmargin=*]
\item \textbf{A0: adapter layer.} Import the direct results and prove the
  notation, normalization, and history adapters classified as \code{adapter}.
\item \textbf{A1: long-time slice.} Define thick and thin selectors and the
  quantifier order for sufficiently late surgery slices.
\item \textbf{A2: hyperbolic output.} Prove pointed convergence, finite-volume
  hyperbolic limits, cusp/core extraction, and the M1e thick certificate.
\item \textbf{A3: local collapse.} Formalize the selected Morgan--Tian or
  Kleiner--Lott collapse theorem with its exact lower-curvature, derivative,
  boundary, and scale hypotheses.
\item \textbf{A4: graph-manifold output.} Convert local collapse to a finite
  graph-manifold certificate and then to the M1e thin-region format.
\item \textbf{A5: surgery transport.} Transport all region, boundary, and
  incompressibility data through the audited surgery history.
\item \textbf{A6: producer composition.} Package A1--A5 as
  the late-branch package conversion, and combine it with terminal data
  in \code{long_time_outcome}.
\end{itemize}

A2--A5 are GC consumer work packages. Their exact direct-reuse, adapter
or surviving-new classifications are determined per theorem after the
completed-release comparison. Reuse of inherited estimates and topology
is mandatory wherever their audited scope suffices; no present absence
claim follows from these package names.

\subsection{No finite-extinction dependency}

The producer contract follows the long-time route. It does not require a
finite-time extinction witness, a finite surgery-history endpoint, or a
non-aspherical extinction theorem. If a later proof attempt discovers a
specific missing implication, it must add a new theorem-inventory row and a
written dependency argument before reopening the extinction workstream.

\subsection{First analytic acceptance target}

Before proving the complete producer theorem, the analytic project should
compile a mock producer whose output is an M1e package for the four standard
examples. This test checks the exact handoff without pretending that the
long-time theorem has been proved. The first new mathematics remains the M3 Ricci-flow PDE gaps.
The hyperbolic and collapse outputs belong to the later M6 tranche.

The audit template for this boundary is recorded in
\code{pc_reuse_audit_template.csv}.
% END migrated B023

\chapter{Long-time behavior and thick/thin decomposition}
\label{ch:long-time-topic}

\paragraph{Scope.} Normalization, regular slices and the honest late-or-terminal outcome.

\paragraph{Planned sections.}
\begin{enumerate}[leftmargin=*]
\item Scale normalization.
\item Quantifiers.
\item Late slices.
\item Thick output.
\item Thin hypotheses.
\item Stabilization.
\item Terminal alternative.
\end{enumerate}

\paragraph{Existing material and present task.} M6 contract and \code{m6_longtime_api.csv}. Keep branch producers separate from output records and expose remaining H1-H4 application obligations.

% BEGIN migrated B026
\section{M6 long-time thick/thin analytic contract}
\label{sec:m6-long-time-analytic-contract}

After the inherited foundation and the Ricci-flow PDE, estimate, compactness,
and surgery interfaces have been audited, the first genuinely GC-specific
analytic task is to prove the long-time thick/thin output. This revision fixes its
quantifier discipline and separates four analytic conclusions that are often
compressed into one sentence: late-slice selection, hyperbolic thick limits,
local collapse of the thin part, and finite topological stabilization.

\subsection{Output theorem and witness form}

The global producer theorem has two explicit branches. In all earlier
contract displays, a long-time output interface is interpreted through this
disjunction; no unconditional existence of late slices is assumed:
\begin{equation}
\begin{aligned}
 &\texttt{long\_time\_outcome} : \\
 &\quad \texttt{RicciFlowSurgeryPackage M} \\
 &\quad\longrightarrow \texttt{LongTimeOutcome M}.
\end{aligned}
\end{equation}
The late-slice branch has a producer and a separate package conversion:
\begin{equation}
\begin{aligned}
 &\texttt{late\_thick\_thin\_outcome} : \\
 &\quad \texttt{LateTimeAvailable flow} \\
 &\quad\longrightarrow \texttt{LateThickThinOutcome M},\\[4pt]
 &
  \texttt{assemble\_late\_thick\_thin} : \\
 &\quad \texttt{LateThickThinOutcome M} \\
 &\quad\longrightarrow \texttt{ThickThinCollapsePackage M}.
\end{aligned}
\end{equation}

Here \code{flow : RicciFlowSurgeryPackage M} is an explicit parameter.
\code{LateTimeAvailable flow} means that every time bound admits a later
nonempty regular slice with valid preceding history. The global producer
proves the late alternative or provides terminal exceptional data; the
public geometrization theorem has no extra late-existence hypothesis.

Internally, the late-slice branch should first construct a
\code{LongTimeOutputWitness}. The witness
contains a valid late flow slice, a finite history segment from the initial
manifold to that slice, a normalization scale, thick and thin finite index
types, and all region/limit/collapse data needed to build the M1e package.
The conversion
\[
  \texttt{LongTimeOutputWitness}
  \longrightarrow
  \texttt{ThickThinCollapsePackage M}
\]
is a separate structural theorem, so the analytic proof is not allowed to
hide package fields inside an informal choice of ``sufficiently late time.''
The other branch is a \code{TerminalExceptionalOutcome} carrying explicit terminal
prime and model data. It is consumed by the exceptional-prime and model-
realization interfaces in C10. This branch is a logical interface, not a
request to formalize Perelman's finite-extinction paper; if the selected PC
adapter already exports the required terminal data, M2A records that reuse.

\subsection{Quantifier normal form}

The formal proof should use the following normal form. For every tolerance
$\varepsilon>0$, there is a late-slice threshold and a valid selected slice
whose rescaled flow data satisfy:

\begin{enumerate}[label=\textbf{L\arabic*:},leftmargin=*]
\item the selected time lies in a valid smooth interval of the surgery flow
  and its preceding history segment has the audited topology maps;
\item a positive normalization scale is fixed, with all curvature and volume
  comparisons stated using that same scale;
\item the thick selector has a uniform noncollapsing/volume lower bound at the
  chosen scale and admits pointed smooth limits;
\item the thin selector satisfies the lower-curvature, small-volume, boundary,
  and derivative hypotheses of the selected local-collapse theorem;
\item the thick and thin selectors, together with exceptional prime data,
  cover each prime carrier after the prescribed collars are removed; and
\item the region labels, boundary tori, and incompressibility proofs transport
  back through the selected history segment to the original $M$.
\end{enumerate}

The output package needs one witness, not a canonical decomposition for every
late time. \code{LateSliceExtraction} selects valid \code{LongTimeSliceData}
from a sequence with the required common hypotheses. The branch theorems
and \code{FiniteRegionStabilization} then construct the full
\code{LongTimeOutputWitness} satisfying L1--L6. No finished package or its
\code{PackageTransport} projection is an input to that construction.
A common index satisfying all branch hypotheses must be justified; an
individual subsequence for each condition does not suffice.

\subsection{Thick branch}

The thick theorem has the following contract:
\begin{equation}
\begin{aligned}
 &\texttt{HyperbolicThickOutput} : \\
 &\quad \texttt{LongTimeSliceData} \\
 &\quad\longrightarrow \texttt{FiniteHyperbolicRegionCertificates}.
\end{aligned}
\end{equation}

Its hypotheses explicitly include the rescaling convention, the lower volume
or injectivity-radius bound, curvature derivative bounds on the selected
balls, and the pointed convergence topology. Its output includes:

\begin{itemize}[leftmargin=*]
\item a finite family of thick region indices and embeddings into the slice;
\item pointed smooth convergence to complete hyperbolic limits;
\item finite-volume and cusp/core extraction data;
\item boundary-torus labels and incompressibility proofs; and
\item transport maps into the M1c cut carriers.
\end{itemize}

Hamilton's 1999 classification and MSM206 Chapters 31--34/Appendix K guide
the nonsingular and hyperbolic-limit decomposition. Perelman's estimates and
the audited PC compactness/noncollapsing results supply the flow hypotheses.
The theorem must state which part is imported and which part is newly proved.

\subsection{Thin branch}

The thin theorem has the contract
\begin{equation}
\begin{aligned}
 &\texttt{CollapsedThinOutput} : \\
 &\quad \texttt{LongTimeSliceData} \\
 &\quad\longrightarrow \texttt{FiniteGraphManifoldRegionCertificates}.
\end{aligned}
\end{equation}

The selected collapse source must be fixed before formalizing its proof. The
candidate route is Kleiner--Lott's boundary theorem with curvature-derivative
control, compared against Morgan--Tian's Alexandrov-space proof. The Lean
statement must expose:

\begin{itemize}[leftmargin=*]
\item the lower-curvature convention and rescaled metric;
\item the volume scale and all orders of choosing constants;
\item the derivative bounds and the balls on which they hold;
\item boundary collars, incompressibility, and regularity hypotheses;
\item the precise local fibration/graph-manifold output; and
\item the conversion from local models to one finite graph-manifold region
  with M1c boundary labels.
\end{itemize}

Shioya--Yamaguchi and Cheeger--Fukaya--Gromov remain comparison/support
sources. They are not silently substituted for the selected local-collapse
theorem.

\subsection{Finite stabilization and transport}

The analytic literature often produces subsequences, local charts, or regions
that vary with the time parameter. M6 must prove a finite stabilization lemma:
\begin{equation}
\begin{aligned}
 &\texttt{FiniteRegionStabilization} : \\
 &\quad \texttt{LongTimeSequenceData} \\
 &\quad\longrightarrow \texttt{LongTimeOutputWitness}.
\end{aligned}
\end{equation}

This lemma chooses one finite region index set and constructs the raw
transport and coverage witnesses from the branch outputs and audited
history. It does not use \code{PackageTransport}, which only projects a
witness from an already constructed M1e package. These raw witnesses
are then fields of \code{LongTimeOutputWitness} and \code{LateThickThinOutcome}.
The later \code{AssembleLateThickThin} theorem converts that supplied data
to the package without rerunning stabilization. This distinguishes data
types, their producers and their projections.
The maps compare the selected regions with the original manifold. It does not claim uniqueness of the
decomposition. Its obligations are finite coverage, consistent boundary
labels, disjointness, and preservation of incompressibility through the
selected surgery-history segment.

\subsection{Analytic work packages}

The quantifier-normalized output splits into these independent packages:

\begin{description}[leftmargin=2.8cm,style=nextline]
\item[M6A: slice and scale] Define valid late slices, rescaling, curvature,
volume, and noncollapsing predicates against the audited PC types.
\item[M6B: thick limits] Prove the pointed hyperbolic convergence and
finite-volume core extraction theorem.
\item[M6C: collapse hypotheses] Formalize the selected local-collapse
hypotheses, including boundary and derivative control.
\item[M6D: graph output] Prove the local-collapse-to-graph-manifold theorem
and convert it to the M1e thin-region certificate.
\item[M6E: stabilization] Choose finite region indices, prove coverage, and
transport labels and incompressibility through surgery.
\item[M6F: producer] Package M6A--M6E as
the late-branch package conversion; \code{long_time_outcome} also handles
terminal exceptional data.
\end{description}

M6B and M6D may proceed independently after M6A fixes the scale and
quantifier conventions. M6E depends on the PC surgery-history adapter and on
both region branches. None of M6A--M6F requires finite-time extinction.

\subsection{Acceptance tests}

Before the full theorem is attempted, the analytic interface should pass:

\begin{itemize}[leftmargin=*]
\item a rescaling test showing that all branch inequalities use one scale;
\item a quantifier test distinguishing a chosen slice from an infinite
  subsequence;
\item a mock thick limit that produces a complete finite-volume hyperbolic
  certificate;
\item a mock collapsed region satisfying the selected local-collapse
  hypotheses and producing a graph-manifold certificate;
\item a failure test when a derivative, boundary, or coverage hypothesis is
  omitted; and
\item an axiom/import audit showing that the long-time module consumes only
  the PC handles and source-specific new theorems listed in the inventory.
\end{itemize}

The field and theorem audit for this layer is recorded in
\code{m6_longtime_api.csv}. The earlier \code{m5_longtime_api.csv} is retained only as
historical revision material and is not a current dependency.
% END migrated B026

\section{Actual producer bindings and the common late witness}
\label{sec:longtime-producer-bindings}

The LP contracts in {\upshape\code{longtime_producer_contracts.csv}} refine M6.
They separate record definitions, supplied hypotheses and proof producers.
No output-producing function is added to the raw surgery-flow package.

\begin{foundation}[Global time domain and finite-prefix interface; LP01]
\label{found:longtime-global-domain}
Bind the accepted global surgery-flow theorem to actual regular slices,
including empty slices. For every time bound there is a later admissible
regular slice, with its covering finite preceding event history and all
smooth-stage comparisons. Empty continuation, if used after termination,
must be a defined part of the flow interface. Each selected bounded
history has finitely many events, not merely a discrete event set.
\end{foundation}
The archived MT Theorem1.2 and its application on printed12--13/PDF16--17
guide this binding. A globally defined flow with locally finite events
provides regular choices on later compact intervals. The actual finite
event enumeration and event/discard conventions must still be supplied
through Chapter11. The existing PC snapshot is not accepted as this
release interface. A maximal flow on an unspecified shorter time domain
does not satisfy LP01 automatically.

\begin{definition}[One normalized slice record; LP02]
\label{def:longtime-same-slice}
An admissible slice record fixes the actual carrier $N=M_t^+$, its
regular time $t>0$, the covering history, a single positive rescaling,
and all thick/thin selectors and comparison maps on that carrier.
Use HG01 for every scale translation. When a subregion carries the
ambient restricted metric, distinguish ambient balls from its intrinsic
balls and state every containment needed by the selected collapse theorem.
\end{definition}

\begin{foundation}[Ordered parameter and application inputs; LP03]
\label{found:longtime-application-inputs}
Choose the finite derivative order and control functions, then the
collapse tolerance, truncations and geometric auxiliary parameters in
the source's order, before taking a late-time threshold. Supply the
flow estimates on every required ball and collar, and the compatible
finite system of parameter constraints. Retain H1--H4 and FC44--FC46:
a mixed upper/lower constraint requires a nonempty admissible interval.
\end{foundation}
This is the qualified flow-to-static bridge. Center-only estimates are
not silently promoted to whole-ball estimates. The boundary radius and
physical collar-width requirements of KL16.1 remain those recorded in
Chapter13, not the superficially similar MT2012 hypotheses.

\begin{lemma}[A single common late choice; LP04]
\label{lem:longtime-common-choice}
Once LP03's parameters are fixed, suppose every member of a finite family
of required predicates holds at every sufficiently late admissible slice.
If {\upshape\code{LateTimeAvailable}} holds, one admissible nonempty slice satisfies
all these predicates and the fixed time bound simultaneously.
\end{lemma}
\begin{proof}
Take a bound above the finite maximum of their thresholds and the fixed
bound, using only the latter if the family is empty. Apply late
availability and retain the selected slice's LP01 history. This is AT29.
Separate unrelated subsequences do not satisfy the premise; neither does
pointwise eventuality over an infinite family of derivative orders.
\end{proof}

\begin{proposition}[Thick-output binding; LP05]
\label{prop:longtime-thick-binding}
On the LP04 slice, bind HG07 to the flow estimates and apply HG17 with
the fixed normalization, core family and cusp choices. The resulting
{\upshape\code{HyperbolicThickOutput}} retains the actual embeddings, complete
interior metrics and ambient cusp maps on that very slice.
\end{proposition}
\begin{proof}
The chosen predicates supply the hypotheses of HG08 and HG15 at the same
time and on the same domains. HG17 constructs the finite records. Store
the LP02 normalization and the actual carrier alongside them; no new
limit or time is chosen in this packaging step.
\end{proof}

\begin{foundation}[Thin-output application; LP06]
\label{found:longtime-thin-binding}
The complementary compact thin carriers of the selected truncation,
with their actual inherited boundary labels and collars, satisfy LP03's
chosen KL hypotheses. Apply LC90/FC45 on those same carriers, or along
specified marked diffeomorphisms, to produce their finite static graph
certificates. All source-qualified Chapter13/14 inputs remain premises.
\end{foundation}
This binds {\upshape\code{CollapsedThinOutput}}. It does not yet give essential
internal graph cuts or all model metrics. Those are Chapter5--7 producers.
Empty thick families, closed hyperbolic components, and components with
no boundary are handled separately in the same finite component ledger.

\begin{lemma}[Joint regions and actual coverage; LP07]
\label{lem:longtime-joint-regions}
Let the finitely many LP05 cores be disjoint smoothly embedded compact
codimension-zero submanifolds with their disjoint collars. Define the
complementary carriers by the closure of their complement. With LP06's
smooth boundary/component data, these cores and complements have exact
coverage, actual matched boundary maps and the induced orientation signs
on $N$. Their boundary tori inject into each complementary component
whenever their composites into $N$ inject.
\end{lemma}
\begin{proof}
Every point is either in a core interior or its closed complement;
intersections are precisely the listed boundary tori. The supplied
collars identify the two sides, giving inverse pairing maps and opposite
induced boundary orientations. Compact smooth manifolds have finitely
many components, since their open components cover a compact space;
retain those actual components. For a boundary map $a:T\to W$ and
inclusion $b:W\to N$, injectivity of $(ba)_*$ implies injectivity of
$a_*$ by AT09. It does not imply injectivity of $b_*$. The identity
on these collared pieces gives the gluing identification with $N$.
\end{proof}
Thin smallness and the collapse theorem did not follow from taking a
complement: they remain LP06. No union of independently selected regions
is substituted for this joint construction.

\begin{proposition}[Finite stabilization and raw reconstruction; LP08]
\label{prop:longtime-finite-reconstruction}
From LP07 and the selected finite history, supply AT13/AT16's relative
prime and good-block producers, AT21/AT22's event comparisons and AT25's
actual exceptional factors. Then AT31 constructs the chosen prime/cut
and reconstruction data AT30, on the initial carrier. Retain the finite
late-region ledger and all correspondence maps as well.
\end{proposition}
\begin{proof}
LP07 supplies AT12's exact selected late carrier and HG15 supplies AT01.
Use AT15--AT18 for chosen late primes, then AT24 to flatten the covering
history and AT25 for its classified exceptions. Retain the normalized
sphere-cycle ledger, counting each converted cycle only once. This is
the conditional AT31 construction; it reads no field from an already
built {\upshape\code{ThickThinCollapsePackage}}.
\end{proof}
This is the intended {\upshape\code{FiniteRegionStabilization}} producer. It
constructs a finite compatible witness; it does not assert eventual
absence of surgery, canonical decomposition, or a reverse-time map on
the entire late slice. Relative graph refinement and model realization
are not discharged by changing their names to ``stabilization.''

\begin{proposition}[Late outcome production; LP09]
\label{prop:longtime-late-producer}
Given late availability, the qualified input theorems above and the
required actual model-carrier data, LP04--LP08 produce a
{\upshape\code{LongTimeOutputWitness}} and hence {\upshape\code{LateThickThinOutcome}}.
The witness retains one time, one scale, the thick/thin data, the chosen
prime reconstruction and exact map/coverage evidence.
\end{proposition}
\begin{proof}
Choose LP04 once and perform the two branch constructions on that
carrier. LP07 proves their joint coverage, and LP08 constructs the raw
transport and refinement data. Insert those constructed values into
the witness record. Definitions of the record and proofs producing
its values have distinct dependency edges.
\end{proof}

\section{Nonlate selection and the explicit terminal branch}
\label{sec:longtime-terminal-producer}

\begin{lemma}[Nonlate implies a selected empty slice under LP01; LP10]
\label{lem:longtime-nonlate-empty}
Assume LP01. If late availability fails, there is an admissible regular
empty slice with a covering finite preceding history.
\end{lemma}
\begin{proof}
Negating ``for every bound there is a later admissible nonempty slice''
gives a bound $B$ above which no admissible slice is nonempty. LP01
provides an admissible regular slice at some $t>B$, including its finite
history. Its carrier must be empty. This argument uses classical logic
and LP01's unbounded regular-time supply; it is not valid for an
arbitrarily truncated or incompletely specified flow.
\end{proof}

\begin{proposition}[Explicit terminal production; LP11]
\label{prop:longtime-terminal-producer}
Given LP10's empty slice and AT21/AT22/AT25 data for its covering finite
history, AT28 produces the explicit terminal exceptional and
reconstruction data required by the terminal case of AT30. Model
witnesses come from the actual classified
factors and the separately proved Chapter7/PC interfaces.
\end{proposition}
\begin{proof}
Start with the empty selected-slice ledger and evaluate its finite
history by AT28. Retain every discarded factor and additional sphere
product required by the reconstruction. Apply AT25's actual recognition
and model witnesses, keeping $S^3$ identity components distinct from
an empty manifold. This constructs the terminal record; failure of late
availability alone supplied none of those exceptional fields.
\end{proof}
This is a conditional evaluation of actual history, not a finite-time
extinction theorem for a class of initial manifolds. It has no local
collapse, Mostow rigidity or ambient-cusp premise.

\begin{theorem}[Conditional long-time outcome; LP12]
\label{thm:longtime-outcome-production}
Once LP01 and the required branch input theorems are supplied for the
chosen admissible global flow, the constructions above produce
{\upshape\code{LongTimeOutcome}}: either LP09's late outcome or LP11's explicit
terminal outcome. No such outcome is a field of the input flow package.
\end{theorem}
\begin{proof}
Decide late availability by classical excluded middle. In its positive
case apply LP09. In its negative case apply LP10 and LP11. These are
alternative values of the outcome type, not simultaneous existence
requirements for both branches. Library theorems proving each conditional
branch remain proof dependencies of this case-split theorem.
\end{proof}

The new producer links are recorded in the production inventory and in
the typed integration ledger. They close missing \emph{recorded paths},
not the mathematical qualifications of HG, LP, LC, FC or AT inputs.
Chapter16's outcome eliminator uses LP12; its pure structural assembly
module does not import LP12 or any Ricci-flow implementation.

\chapter{Local collapsing theory}
\label{ch:local-collapse-topic}

\paragraph{Current revision 96 reading guide.}
The chronological progress paragraphs below retain their historical
context. The current completion/status table is in
Section~\ref{sec:collapse-completion-audit}. Chapter 13 now owns local
adapted-coordinate existence and local fiber topology; Chapter 14 owns
compatibility and assembly. Blueprint coverage is complete, with
source-qualified producer and formalization gates listed explicitly.

\paragraph{Scope.} The static local-collapse route and its separate flow application.

\paragraph{Planned sections.}
\begin{enumerate}[leftmargin=*]
\item Exact local hypotheses.
\item Volume/curvature scales.
\item Local models and strata.
\item Boundary conditions.
\item Comparison route.
\item Flow application.
\end{enumerate}

\paragraph{Existing material and present task.} KL14/MT2012 comparison and I04. LC01--LC15 write the volume-scale and local-volume arguments for KL 6.18. LC16--LC20 exclude the three-stratum; LC21--LC43 add cone scales, radial smoothing, sublevel topology and quantitative core transfer. Boundary adaptation, inherited formal bindings and smooth/fibration producers remain explicit.

% BEGIN migrated B037
\section{Collapse: the two principal proof sources}
\label{sec:collapse-sources}

\textbf{Revision decision.} Morgan--Tian and Kleiner--Lott (2014) have central
roles. The previous assignment of the collapse interface to
Shioya--Yamaguchi and Cheeger--Fukaya--Gromov was too broad. Local collapse
under a lower sectional-curvature bound must be distinguished from global
small-volume results and from two-sided bounded-curvature collapse.

\textbf{Morgan--Tian.} Use \emph{The Geometrization Conjecture}
\cite{MorganTian2014} both for its collapse proof and for the connection
between long-time flow and the topological conclusion. Its treatment develops
Gromov--Hausdorff limits and low-dimensional Alexandrov geometry, obtains
local topological descriptions, and assembles them. The archived file is
\code{MorganTianGeom.pdf}, but revision 45 identifies its body as the August
7, 2012 manuscript. Its body numbering is not the numbering in its prefixed
contents or a verified locator in the 2014 publication.

\textbf{Kleiner--Lott.} \emph{Locally collapsed 3-manifolds}
\cite{KleinerLottCollapse2014}, Ast\'erisque 365 (2014), pp.~7--99, supplies a
dedicated proof. It uses volume scales, stratification, and compatible local
fibrations. Sections~16--17 treat boundary and application to geometrization;
Section~18 discusses removing curvature-derivative assumptions. This paper
is separate from \emph{Notes on Perelman's papers} \cite{KleinerLott2008}.
The local copy is \code{KleinerLottAsterisqueLocalCollapse.pdf}.

\textbf{Provisional recommendation.} Start the comparison with
Kleiner--Lott's boundary theorem with curvature-derivative control as the
candidate to formalize, and use Morgan--Tian as the full-GC guide and an
independent collapse proof. This is a formalization judgment, not a claim that
one proof is mathematically superior. Derivative estimates may already be
available from the completed PC library and its long-time extensions. We must
prove that they hold on the required balls with the required uniformity.
The version without derivative assumptions may require additional stability
theory; it should not be selected merely because its statement is shorter.
Both routes require geometric and topological foundations beyond PC.

We should choose one proof of the collapse theorem for the first complete GC
formalization. The other remains a comparison source unless it offers a
specific lemma with a verified compatible interface. We do not assume that
the two proofs' local models, scales, or graph-manifold definitions coincide.
% END migrated B037

% BEGIN migrated B038
\section{The interfaces that the comparison must resolve}
\label{sec:collapse-comparison}

The comparison in Section~\ref{sec:independent-collapse-contracts} begins
the exact-statement audit and distinguishes the archived Morgan--Tian
manuscript from the published book. The following five checks remain the
acceptance criteria:
\begin{enumerate}[label=\textbf{C\arabic*:},leftmargin=*]
\item \textbf{Region and boundary.} Specify the compact complement of
  truncated hyperbolic pieces. Record the collar geometry, boundary size,
  regularity, and every incompressibility assumption. A closed-manifold
  collapse theorem alone does not cover this region.
\item \textbf{Scales and quantifiers.} Compare curvature and volume scales,
  metric rescalings, and the order of choosing constants and large times.
  Translate any sequential conclusion into the form actually needed by GC.
\item \textbf{Analytic inputs.} Record the precise derivative order, the balls
  on which estimates hold, and which bounds are uniform. Assign each input
  either to the completed PC library or to a new long-time theorem.
\item \textbf{Topological output.} Compare graph-manifold definitions and
  prove any required conversion. Auxiliary cutting tori in a graph-manifold
  presentation need not be the incompressible tori required at the GC
  endpoint. Torus bundles, product regions, and twisted bundles require
  explicit treatment when passing to geometric pieces.
\item \textbf{Application and reconstruction.} Prove that the actual thin
  regions meet the selected hypotheses, then use their graph-manifold
  structure and the hyperbolic regions to recover the initial manifold's
  decomposition through the surgery history. These are separate obligations
  from the collapse theorem itself.
\end{enumerate}

This yields independently assignable developments once definitions agree:
local collapse, the estimates and collars needed to apply it to flow, and the
conversion of its graph-manifold output to the GC endpoint. The collapse proof
does not depend on the construction of Ricci flow; its application does.

\textbf{Supporting references.} Keep Shioya--Yamaguchi (2005)
\cite{ShioyaYamaguchi2005} for related results and comparison.
Cheeger--Fukaya--Gromov \cite{CheegerFukayaGromov1992} becomes a required
dependency only if a selected proof step actually invokes an appropriate
bounded-curvature result. Neither entire theory is a mandatory extra project
merely because it belongs to the history of collapsing.

\textbf{Version control.} Record the local filename, source version, printed
locator, and PDF page for each selected theorem. Include applicable errata;
Lott provides a correction sheet for the Ast\'erisque collapse paper at
\url{https://math.berkeley.edu/~lott/corrections.pdf}. No source-level
comparison in this revision certifies a Lean proof or finalizes the route.
% END migrated B038

% BEGIN migrated B039
\section[Independent collapse contracts]{Independent collapse contracts while the PC release is pending}
\label{sec:independent-collapse-contracts}

\textbf{Revision 45 scheduling decision.} Pause further declaration-level
expansion of the provisional PC audit. Preserve the 111 discoveries, 333
declaration probes, twenty consumer drafts and their release gates. Resume
exact adapter work when a completed release can be pinned, or when a
specific upstream interface has been deliberately stabilized. The current
snapshot is not a compatibility target for production GC code.

The active work is mathematical contract and proof decomposition for local
collapse, followed by graph-manifold refinement and the hyperbolic boundary
bridge. These specifications do not require PC's eventual declaration names.
Their later Lean implementations will still use its audited geometric
foundation. This changes the work schedule, not the foundation or the M3
gate: the first genuinely new analytic layer is still the surviving
Ricci-flow PDE gaps after completed-release comparison.

\subsection{Exact copies and locators}

The source comparison is indexed in \code{collapse_source_contracts.csv}.
The archived Kleiner--Lott paper is the published Ast\'erisque article:
printed page $p$ is PDF page $p-5$. Definitions 1.1--1.2 and Theorem 1.3
are on printed page 8 (PDF 3); Theorem 16.1 is on printed pages 87--88
(PDF 82--83); Theorems 17.1 and 17.3 are on pages 90--91 (PDF 85--86).

There is a material version distinction in \code{MorganTianGeom.pdf}.
PDF pages 2--4 contain a contents listing with Chapter 6 and Theorem 6.2.1,
but the body begins on PDF page 5 with an August 7, 2012 title page.
The actual body uses Definition 5.2 and Proposition 5.4 on printed page 55
(PDF 59), and Theorem 5.5 on printed page 56 (PDF 60).
These are locators for the archived 2012 manuscript
\cite{MorganTian2012Archive}; they are not verified 2014-book locators.
The source manifest now records this distinction without changing the file.
Morgan--Tian 2014 remains the global spine; its published body and theorem
crosswalk must be verified before transferring these locators to it.

\subsection{Candidate local theorem: KL14 Theorem 16.1}

This is a theorem about a single compact connected orientable Riemannian
three-manifold with boundary. It contains no Ricci-flow or surgery-history
input. Let $c_3$ be the Euclidean unit three-ball volume. Its parameter order is
\[
 \forall K\in\N,\ K\geq10,\quad
 \forall A:(0,\infty)\longrightarrow(0,\infty),\quad
 \exists w_0\in(0,c_3),\quad \forall(M,g).
\]
The four geometric premises are as follows; all quantities use the same
metric normalization.
\begin{enumerate}[label=\textbf{KL\arabic*:},leftmargin=*]
\item Each boundary component has diameter at most $w_0$.
\item Each boundary component $X$ admits an embedding of pairs from the
  $100$-neighborhood of the boundary of a hyperbolic cusp into $(M,X)$.
  The embedding is $C^{K+1}$ and $w_0$-close to an isometry in the source's
  sense. Section 16 uses the cusp metric
  $dz^2+e^{-z}(dx^2+dy^2)$ modulo a lattice, with sectional curvature $-1/4$.
\item At every $p$ with $d(p,\partial M)\geq10$, collapse holds at the
  curvature scale $R_p$:
  \[
    \Vol_g B_g(p,R_p)\leq w_0 R_p^3.
  \]
\item For every $p\in M$, $w'\in[w_0,c_3)$, integer $0\leq k\leq K$,
  and positive $r\leq R_p$, a lower volume ratio triggers derivative
  control throughout the ball:
  \[
   \begin{aligned}
    \Vol_g B_g(p,r)&\geq w'r^3\quad\Longrightarrow\\
    |\nabla^k\Rm|_g(q)&\leq A(w')r^{-(k+2)}
       \quad\text{for every }q\in B_g(p,r).
   \end{aligned}
  \]
\end{enumerate}
The conclusion is graph-manifold structure in Definition 1.2: a finite
disjoint family of interior embedded tori such that the completed cut
components are circle bundles over surfaces. That definition does not
assert incompressibility of its cutting tori. Neither a final geometric
certificate nor transport to the initial manifold is a conclusion of
Theorem 16.1. Boundary incompressibility is not one of its four premises.

\textbf{Definition obligations retained.} Definition 1.1 assigns
$R_p=\infty$ on a nonnegatively curved connected component, and otherwise
characterizes $R_p$ by the infimum of sectional curvature on its ball.
The extended-radius case must not be silently replaced by an ordinary
positive real or discarded. Specify the boundary metric, distance and ball
conventions, and the precise $C^{K+1}$ closeness predicate before translating
the statement into Lean. A region's balls must not be silently identified
with ambient flow balls. An empty boundary needs a separate distance
convention or the closed theorem. A disconnected thin carrier is treated
componentwise with explicit finite component data and common parameters.
These are specification obligations, not missing PC theorems.

\clearpage
\subsection{Comparison with the archived Morgan--Tian statement}

\setlength{\tablebodywidth}{\dimexpr\textwidth-4\tabcolsep\relax}
\begin{longtable}{@{}P{0.19\tablebodywidth}P{0.39\tablebodywidth}P{0.42\tablebodywidth}@{}}
\textbf{Feature} & \textbf{KL14, Theorem 16.1} & \textbf{Archived MT12, Theorem 5.5}\\\hline
\endhead
Logical form & Fixed $K,A$ determine one $w_0$; apply to one connected compact manifold. & Sequence $(M_n,g_n)$ with $w_n>0$ tending to zero; conclusion for all sufficiently large $n$. Closed or convex boundary is allowed.\\
Collapse radius & Curvature scale $R_p$; the volume premise is required away from the boundary by distance at least $10$. & For every $x\in M_n$, some $\rho_n(x)>0$ has sectional curvature at least $-\rho_n(x)^{-2}$ and volume at most $w_n\rho_n(x)^3$.\\
Boundary & Small diameter and a $C^{K+1}$ nearly isometric cusp collar of depth $100$. & A uniform constant $C$ bounds diameter by $Cw_n$; locally convex incompressible tori; trivial collars containing the distance-$1$ neighborhood, with sectional curvature between $-5/16$ and $-3/16$.\\
Derivatives & Orders $0$ through $K$, on the entire triggered ball, for every positive $r\leq R_p$. & For every $w'>0$, a radius cutoff and constants for all derivative orders; eventually in $n$, positive radii below that cutoff with the stated volume and curvature lower bounds give estimates at the center.\\
Graph output & Definition 1.2: torus cuts into circle bundles, without an incompressibility clause. & Definition 5.2: prime factors split along incompressible tori into Seifert-fibered pieces.\\
\end{longtable}

There is no automatic implication between these recorded input packages.
In particular, center estimates do not immediately supply ball estimates,
and a sequential radius cutoff does not supply all radii up to $R_p$.
Choosing $A$ after choosing $w_0$ would reverse the local theorem's
quantifiers. Eventual assertions for each individual $w'$ do not, by
themselves, yield one tail uniform over the entire interval $[w_0,c_3)$.

The graph-manifold definitions also require a proved conversion. The
archived MT Proposition 5.4 additionally assumes the boundary of the
embedded truncated hyperbolic region is incompressible in the ambient
closed manifold. Its proof discusses prime factors, product pieces and
twisted bundles, and may replace certain tori by Klein bottles. It is a
comparison for reconstruction, not a direct proof of the project's
torus-based endpoint contract.

\subsection{The flow application stays a separate obligation}

KL14 Section 17 normalizes by $\widehat g(t)=t^{-1}g(t)$ and uses
hyperbolic limits with sectional curvature $-1/4$. Theorem 17.1 supplies
finite hyperbolic pieces, increasingly large almost-isometric images,
coverage of a shrinking thick part, and ambient incompressibility of
cuspidal tori. Theorem 17.3 applies the boundary collapse theorem to the
complements of truncated images along a sequence $t^\alpha\to\infty$.
It is not an assertion about an arbitrary region labelled thin.

The proof of Theorem 17.3 checks the diameter and cusp collar from the
hyperbolic embeddings, and interior collapse from thick-part coverage.
For the derivative premise it first uses the small-radius Ricci-flow
estimate cited there as KL08 Lemma 92.13. Its additional argument uses
Bishop--Gromov comparison and hyperbolic closeness (KL08 Lemma 88.1),
forcing a noncollapsed admissible radius to lie below the small-radius
cutoff. The resulting common $A$ and the choice of the eventual tail must
be justified explicitly. These KL08 locators are leads cited by KL14;
their exact archived statements have not been compared in this iteration.

The inventory therefore separates the geometric \code{CollapseHypotheses}
from the \code{CollapsedThinOutput} application on \code{LongTimeSliceData}.
The local \code{CollapseToGraphManifold} theorem consumes the geometric
hypotheses and topological foundation, not a prior long-time decomposition.
Its eventual foundation imports remain release-gated. No replacement
manifold, curvature, or metric library is introduced.

\subsection{Corrections, limits and next bounded work}

The author's correction sheet dated May 15, 2015 was checked at
\url{https://math.berkeley.edu/~lott/corrections.pdf}. It changes the final
interval in Corollary 6.5 to $(0,1+2\Lambda]$, adds $D^3$ to
Proposition 14.1(2), and corrects the basepoints and limiting spaces in
Lemma 20.2. Record these at the affected proof steps. The archived
\code{MorganTianCorrection19-2.tex} concerns curve shrinking in the
Poincar\'e book; it is not an errata audit of this collapse comparison.

KL14 Theorem 18.1 removes derivative assumptions for \emph{closed}
manifolds. This does not silently replace the boundary theorem. Retain the
derivative-controlled boundary route as provisional; H1--H4 remain open.
No source comparison is a Lean proof or a release acceptance.

Chapter~\ref{ch:seifert-graph} decomposes the graph-manifold output conversion into
source-located proof targets and exposes an endpoint compatibility gate.
The boundary-scale and uniformity obligations above remain visible for the
later flow application. This work proceeds without expanding provisional
PC declaration probes.
% END migrated B039

\section{Volume scales and the dimension of collapsed models}
\label{sec:collapse-volume-models}

This section expands KL Lemma 6.18, printed 44--45/PDF 39--40.
It separates the volume estimate, geometric compactness and the
full-dimensional volume input. All manifolds here have no boundary;
pointed manifolds are connected, or restricted to the basepoint component.
The application to the boundary theorem and to surgery slices remains
separate. The contracts LC01--LC09 are recorded in
\code{local_collapse_contracts.csv}. Their written arguments are not
Lean proofs. LC06 records the local volume implication that is now
supplied by LC15 in Section~\ref{sec:collapse-local-volume-proof}.

We write $\rho$ for KL's modified scale, and $w$ for
its parameter $\bar w$. Rescaling the metric distance by $\rho^{-1}$
means replacing the Riemannian tensor $g$ by $\rho^{-2}g$.
All volume factors below are therefore $\rho^{-3}$.

\clearpage
\subsection{First volume scales and the standing sequence}

\begin{lemma}[Attained volume scale on a closed manifold; LC01]
\label{lem:collapse-volume-scale}
Let $(M,g)$ be a closed connected smooth Riemannian three-manifold,
$p\in M$, and $0<w<c_3$, where $c_3$ is the Euclidean unit-ball
volume. Put
\[
 r_p(w)=\inf\{r>0:\Vol_g B_g(p,r)=wr^3\}.
\]
Then $0<r_p(w)<\infty$, equality holds at this first radius,
and $\Vol_g B_g(p,s)>ws^3$ for $0<s<r_p(w)$.
If $0<w'<w<c_3$, then $r_p(w)<r_p(w')$.
The proof uses the inherited continuity and small-ball asymptotics
of Riemannian volume; these analytic facts are not re-proved here.
\end{lemma}
\begin{proof}
The continuous ratio $\Vol B(p,r)/r^3$ tends to $c_3$ as
$r\downarrow0$ and to zero as $r\to\infty$, since total volume
is finite. The intermediate value theorem gives a root. No root
lies in a sufficiently small positive interval. A sequence of roots
tending to their positive infimum and continuity show attainment.
Before that first root the ratio must exceed $w$, since otherwise
continuity from a smaller radius would give an earlier root.
At $r_p(w')$ the ratio is $w'<w$, so the intermediate value theorem
gives an earlier $w$-root. This argument does not assume that the
Euclidean-normalized volume ratio is globally monotone.
\end{proof}

This is KL Definition 1.5, printed 10/PDF 5, specialized to closed
manifolds. On general complete manifolds a volume scale can be infinite;
LC01 does not discard that case from the source definition. The curvature
scale $R_p$ is also allowed to be infinite on a nonnegatively curved
component, as in KL Definition 1.1.

For the closed sequence in KL Standing Assumption 5.2 retain only the
following part for the present metric argument:
\[
 R_p\geq\alpha r_p(1/\alpha)
 \quad\text{for every }p\in M^\alpha,
 \qquad \alpha\longrightarrow\infty.
\]
We work on the tail where $1/\alpha<c_3$. The curvature-scale convention
means that $\sec_g\geq-R_p^{-2}$ on its defining ball when $R_p<\infty$;
$R_p=\infty$ means nonnegative sectional curvature on the component.
Thus for fixed $0<v<c_3$, once $1/\alpha\leq v$, LC01 gives
$R_p\geq\alpha r_p(v)$ uniformly in $p$. This implication also makes
sense when $R_p=\infty$, without dividing infinity by a real number.
The derivative bounds and non-graph conclusion in 5.2 are not premises
of this metric step. Their roles elsewhere in KL are retained.

\paragraph{Inherited comparison interface.}
For a complete Riemannian three-manifold, if
$\operatorname{Ric}\geq-2\kappa^2g$ on $B(q,L)$, the local
Bishop--Gromov input compares concentric balls of radii
$0<a\leq b<L$ by
\[
 \frac{\Vol B(q,b)}{V_\kappa(b)}
 \leq\frac{\Vol B(q,a)}{V_\kappa(a)},\qquad
 V_\kappa(t)=4\pi\int_0^t
       \left(\frac{\sinh(\kappa u)}{\kappa}\right)^2du.
\]
The integrand is $u^2$ at $\kappa=0$; endpoint radii follow by
continuity when the curvature bound holds on the open ball up to that
radius. A sectional bound $-\kappa^2$ supplies this Ricci bound in
dimension three. This is an inherited geometric comparison input,
with its local radial domain explicit, not a newly verified PC declaration.
The radial Jacobian proof in GSM77, labels
\code{Thm bish vol compr} and \code{Blade Runner PD}, and BBI 10.6.6
explain the localization: the minimizing radial segments of length below
$L$ stay in $B(q,L)$. Cut-locus and volume-identification bindings remain
part of the inherited interface.

\subsection{The modified scale, with the corrected interval exposed}

\begin{proposition}[Slowly varying modified volume scale; LC02]
\label{prop:collapse-modified-scale}
Fix $\Lambda>0$ and $0<w<c_3$, and put
\[
 S=1+2\Lambda^{-1},\qquad
 T=1+2\max\{\Lambda,\Lambda^{-1}\},\qquad
 w'=\frac{w}{2S^3}.
\]
For all sufficiently large $\alpha$, uniformly over $p\in M^\alpha$,
there is a smooth positive $\Lambda$-Lipschitz function $\rho$ with
\[
 \tfrac12 r_p(w)\leq\rho(p)\leq2r_p(w').
\]
This is conditional on the inherited local Bishop--Gromov interface
and compact-manifold Lipschitz smoothing in KL Corollary 3.15.
The tail can depend on $w,\Lambda$ and the supplied sequence.
\end{proposition}
\begin{proof}
Set $l(p)=r_p(w)$ and $u(p)=r_p(w')$. We prove
$l(p)-(\Lambda/2)d(p,q)\leq u(q)$ for all $p,q$.
Otherwise $d(p,q)<2l(p)/\Lambda$ and $u(q)<l(p)$, giving
\[
 B(p,l(p))\subset B(q,Sl(p))
       \subset B(p,(1+4\Lambda^{-1})l(p)).
\]
For any fixed small $c>0$, the standing curvature-scale bound puts
sectional curvature at least $-c^2l(p)^{-2}$ on the outer ball, on a
tail independent of $p,q$. This follows by taking
$\alpha>1+4\Lambda^{-1}$ and $\alpha\geq c^{-1}$, as well as
$1/\alpha\leq w$. Hence local Bishop--Gromov applies at $q$.
Write $x=u(q)/l(p)\in(0,1)$ and
\[
 J_c(t)=\int_0^t\left(\frac{\sinh(cv)}c\right)^2dv.
\]
LC01 and the first ball inclusion give
$\Vol B(q,Sl(p))\geq wl(p)^3=2w'S^3l(p)^3$.
The comparison and $\Vol B(q,u(q))=w'u(q)^3$ therefore imply
\[
 \frac{x^3}{J_c(x)}\geq\frac{2S^3}{J_c(S)}.
\]
For $0<t\leq T$, elementary monotonicity of $\sinh(v)/v$ gives
\[
 1\leq\frac{3J_c(t)}{t^3}
 \leq\left(\frac{\sinh(cT)}{cT}\right)^2.
\]
Choose $c$ so small that the last bound is less than $3/2$.
Then the left side of the preceding comparison is at most $3$,
while its right side is greater than $4$, a contradiction.
This establishes the envelope inequality without assuming $S\leq3$.

MC19 supplies a positive $(\Lambda/2)$-Lipschitz function $f$
with $l\leq f\leq u$. Its minimum $m$ is positive on the compact
manifold. Apply the cited smoothing input with
$0<\eta<\min\{m/2,\Lambda/2\}$ to obtain smooth $\rho$ with
$\|\rho-f\|_\infty<\eta$ and Lipschitz constant at most
$\Lambda/2+\eta<\Lambda$. Then
$\rho>f/2\geq l/2$ and $\rho<3f/2\leq2u$.
\end{proof}

\paragraph{What the correction sheet literally says.}
The May 15, 2015 author sheet replaces the proof's interval $(0,3]$
by $(0,1+2\Lambda]$. The archived equations (6.4), (6.7)--(6.9)
instead use $1+2\Lambda^{-1}$. Both were visually checked.
We retain this reciprocal-notation discrepancy explicitly; we do not
attribute $(0,S]$ to the sheet. The proved estimate on $(0,T]$ contains
both endpoints and is sufficient for every fixed $\Lambda>0$.
This is our stated domain repair, not a claim of a new author erratum.

\subsection{Small upper volume and fixed-parameter noncollapse}

\begin{lemma}[Volume upper bound at twice the modified scale; LC03]
\label{lem:collapse-small-volume-upper}
Let $(M,g)$ be complete Riemannian of dimension three, with $p\in M$
and $\rho>0$. Suppose $w>0$, a finite volume scale $r=r_p(w)$ is attained,
$r\leq2\rho$, and $\sec_g\geq-(2\rho)^{-2}$ on $B(p,2\rho)$.
With $I(t)=\int_0^t\sinh^2u\,du$ and $C_H=3I(1)$, local
Bishop--Gromov gives
\[
 \frac{\Vol_g B_g(p,2\rho)}{(2\rho)^3}\leq C_H w,
 \qquad
 \Vol_{\rho^{-2}g}B_{\rho^{-2}g}(p,2)\leq8C_Hw.
\]
\end{lemma}
\begin{proof}
Set $t=r/(2\rho)\in(0,1]$. Comparison at radii $r$ and $2\rho$
with curvature model $-(2\rho)^{-2}$ yields
\[
 \frac{\Vol B(p,2\rho)}{(2\rho)^3}
 \leq w I(1)\frac{t^3}{I(t)}\leq3wI(1),
\]
since $\sinh u\geq u$ gives $I(t)\geq t^3/3$.
The case $t=1$ also satisfies the inequality, directly or by continuity.
The second assertion is the Riemannian volume scaling identity.
No lower bound on $t$ is needed. This is KL (6.19)--(6.20).
\end{proof}

\begin{lemma}[Lower volume for a fixed modified parameter; LC04]
\label{lem:collapse-fixed-parameter-volume}
Let $(M,g)$ be a complete Riemannian three-manifold, $p\in M$,
$0<w'<c_3$, and $u=r_p(w')$ finite and attained. Suppose
$0<\rho\leq2u$ and
$\sec_g\geq-u^{-2}$ on $B(p,u)$. Then, under the same comparison input,
\[
 \frac{\Vol_g B_g(p,\rho)}{\rho^3}
 \geq\frac{w'}{24I(1)}>0.
\]
\end{lemma}
\begin{proof}
Put $t=\rho/(2u)\leq1$. Comparing the smaller ball of radius
$\rho/2$ with the $u$-ball gives
$\Vol B(p,\rho/2)\geq w'u^3I(t)/I(1)$.
Use $I(t)\geq t^3/3$ and $B(p,\rho/2)\subset B(p,\rho)$.
\end{proof}

LC04 is the volume part of KL 6.10, printed 42--43/PDF 37--38.
For fixed $w,\Lambda$, $w'>0$ is fixed. In the contradiction for
6.18, however, $w_i\to0$ and therefore $w_i'\to0$.
The lower bound in LC04 then tends to zero. One must not use it as a
uniform positive noncollapse bound for that diagonal sequence, nor use
KL 6.10's fixed-parameter smooth compactness without uniform derivative
data. LC03 supplies the volume that actually tends to zero.

\subsection{Geometric extraction and the separate volume input}

\begin{lemma}[Curvature-scale exhaustion at a chosen scale; LC05]
\label{lem:collapse-rescaled-exhaustion}
Let $(M_i,g_i,p_i)$ be complete pointed Riemannian three-manifolds,
$\rho_i>0$, and $H_i\to\infty$ positive, with
\[
 \sec_{g_i}\geq-(H_i\rho_i)^{-2}
      \quad\hbox{on }B_{g_i}(p_i,H_i\rho_i).
\]
Set $\widehat g_i=\rho_i^{-2}g_i$. Conditional on the inherited
Riemannian-to-local-Alexandrov comparison and dimension interfaces,
and AC41's retained comparison input, a subsequence of
$(M_i,\widehat g_i,p_i)$ converges to a complete proper geodesic
nonnegatively curved space of Hausdorff dimension at most three.
\end{lemma}
\begin{proof}
Distances and volumes scale by $\rho_i^{-1}$ and $\rho_i^{-3}$;
sectional curvature scales by $\rho_i^2$. Thus the $H_i$-ball in
$\widehat g_i$ has sectional curvature at least $-H_i^{-2}$.
Completeness and the length property are preserved. Smooth
three-dimensional metric neighborhoods have Hausdorff dimension three;
the stated local comparison interface supplies the Alexandrov bound.
Apply AC41 with $c_i=H_i^{-2}$, $r_i=H_i$, and $n=3$.
No lower volume or curvature-derivative bound enters this extraction.
\end{proof}

\begin{foundation}[Full-dimensional local volume input; LC06]
\label{found:collapse-local-volume-input}
For sequences as in LC05 that converge to a complete proper
nonnegative Alexandrov limit $(X,x)$, the needed input is:
\[
 \dim_H X>2\quad\Longrightarrow\quad
 \liminf_i\Vol_{\widehat g_i}B_{\widehat g_i}(p_i,2)>0.
\]
The positive bound may depend on the particular limiting space and
sequence; no uniform bound over all three-dimensional limits is asserted.
Revision 80 recorded this as an open mathematical input. LC15 below
now supplies the local argument with the inherited comparison and
volume interfaces retained; it is not an accepted axiom or a proved declaration.
\end{foundation}

KL 6.18 uses this implication in the sentence following its assertion
that the limit is three-dimensional. Two source facts are involved:
finite-dimensional Alexandrov dimension is integral (BGP 6.4--6.5,
printed 20--21/PDF 21--22), and full-dimensional convergence prevents
local volume collapse. BGP Theorem 10.8, printed 42/PDF 43, gives
Hausdorff-measure convergence for \emph{compact} $n$-dimensional
Alexandrov spaces with a common global lower curvature bound.
Its proof at 10.10, printed 43/PDF 44, uses the almost-isometry
Theorem 9.8 and the singular-set estimate 10.9. The theorem as printed
cannot be applied to arbitrary restricted closed balls as if they were
Alexandrov length spaces. Nor do the LC05 manifolds have a common
curvature bound on the whole manifold.

The direct proof of LC06 below localizes the full-dimensional argument to
fixed buffered balls, uses the same pointed approximation maps, and identifies normalized
three-dimensional Hausdorff measure with Riemannian volume. LC10--LC15
prove a fixed distance-coordinate neighborhood and a positive image-cube
bound near a lifted rank-three packet. AC40's upper-dimensional estimate
alone would not suffice. The full measure-convergence theorem remains
a separate result, and no theorem is classified as absent from PC.

\begin{proposition}[Vanishing normalized volume gives a low-dimensional limit; LC07]
\label{prop:collapse-low-dimensional-limit}
Assume LC06 and the inputs of LC05. Suppose additionally that
$w_i>0$, $w_i\to0$, each $r_{p_i}(w_i)$ is finite and attained, and
$r_{p_i}(w_i)\leq2\rho_i$. Then a subsequence of
$(M_i,\rho_i^{-2}g_i,p_i)$ converges to a complete proper
nonnegative Alexandrov space of Hausdorff dimension at most two.
\end{proposition}
\begin{proof}
LC05 gives a limit $X$ of dimension at most three. On a tail $H_i>2$,
the curvature hypothesis implies the bound in LC03 on $B(p_i,2\rho_i)$.
Thus $\Vol_{\widehat g_i}B(p_i,2)\leq8C_Hw_i\to0$.
If $\dim_H X>2$, LC06 contradicts this upper bound. Hence
$\dim_H X\leq2$. This does not invoke smooth compactness or assume
that dimension is automatically preserved under metric convergence.
\end{proof}

\subsection{Uniform models and the order of parameter choices}

\begin{theorem}[Uniform metric model contract; LC08]
\label{thm:collapse-uniform-model}
Conditional on LC06 and the retained comparison interfaces, for every
$0<\sigma<1$ there are $0<w_*<c_3$ and $H\geq4$ such that the
following holds for every complete pointed Riemannian three-manifold
$(M,g,p)$ and every $\rho>0$.
If $0<w<w_*$, $r_p(w)$ is finite and attained, $r_p(w)\leq2\rho$,
and
\[
 \sec_g\geq-(H\rho)^{-2}\quad\hbox{on }B_g(p,H\rho),
\]
then there is an actual pointed KL $\sigma$-approximation from
$(M,\rho^{-2}g,p)$ to some complete proper nonnegative Alexandrov
space of dimension at most two. The constants are independent of
$M,p,\rho,w$ subject to those conditions. No explicit modulus or
identity of pointed GH distance conventions is claimed.
\end{theorem}
\begin{proof}
If this fails for a fixed $\sigma$, choose counterexamples with
$H_i=i+4$ and $0<w_i<\min\{c_3/2,1/i\}$ for $i\geq1$.
LC07 gives a subsequential pointed limit $X$ of dimension at most two.
MC11 converts the actual bounded-radius convergence maps to pointed
KL $\sigma$-approximations to this fixed $X$ for all sufficiently late
indices. That contradicts their defining failure. The radius padding
and possible error enlargement are precisely MC11's convention bridge;
no metric on isometry classes is being numerically identified with it.
\end{proof}

\begin{corollary}[KL 6.18 with its suppressed tail restored; LC09]
\label{cor:collapse-kl618-tail}
Fix $\Lambda>0$ and $0<\sigma<1$. Subject to LC06 and the retained
comparison and smoothing inputs, there is $w_*(\sigma,\Lambda)>0$
such that for every fixed $0<w<\min\{w_*,c_3\}$ and every standing
closed sequence above, there is an index $\alpha_0$ for which every
$\alpha\geq\alpha_0$, every $p\in M^\alpha$, and every modified
scale satisfying LC02's bounds has LC08's low-dimensional metric model.
The model and its approximation may depend on $p,\alpha$ and the scale.
\end{corollary}
\begin{proof}
Choose $w_*$ and $H$ from LC08 and fix $w,\Lambda$ afterwards.
Set $w'=w/[2(1+2\Lambda^{-1})^3]$. On a tail,
$1/\alpha\leq w'$, and LC01 and the standing assumption give
\[
 R_p\geq\alpha r_p(w'),\qquad
 \rho(p)\leq2r_p(w'),\qquad
 r_p(w)\leq2\rho(p).
\]
For $\alpha>2H$ the curvature-scale convention therefore gives
$\sec_g\geq-(H\rho(p))^{-2}$ on $B(p,H\rho(p))$, uniformly in $p$.
Use nonnegative curvature directly if $R_p=\infty$. LC08 applies.
If existence of a smooth scale is needed, further enlarge the tail
using LC02. The threshold $\alpha_0$ can depend on $w$; it need not
work simultaneously for all $w\downarrow0$.
\end{proof}

The source's double sequence is now explicit. If one argues directly
by contradiction, first select $w_i\to0$ and then, for each fixed
$w_i'=w_i/[2(1+2\Lambda^{-1})^3]$, select a sufficiently late
$j(i)$ so that the curvature radius dominates $i\rho_i$.
Selecting one common tail before varying $w$ would reverse the
quantifiers. LC08 packages the same compactness contradiction by
separating small volume and a sufficiently large curvature buffer.

\paragraph{Consumers and next proof task.}
LC09 supplies only a metric model, conditional on LC06. It does not
produce a homeomorphism, smooth adapted coordinates, a circle bundle,
a graph-manifold conclusion or an application to boundary regions.
Section~\ref{sec:collapse-three-stratum-exclusion} supplies KL 7.2's
uniform separation argument using full approximate products.
The following section supplies the local full-dimensional volume
implication LC06 through distance coordinates and a fixed image cube.
The remaining comparison and formal volume bindings stay explicit. Chapter 14
retains smooth adapted coordinates and fibration compatibility;
mathlib pinning and PC-specific probes remain deferred.

\section{Local volume from distance coordinates}
\label{sec:collapse-local-volume-proof}

This section supplies the localized implication LC06, using the written
metric arguments of Chapter 4. The mechanism is a fixed cube in the
image of three distance coordinates. No convergence of measures on the
whole limit is needed. The comparison is with BGP 5.4/5.8 and 6.3--6.5,
and BBI 10.8.15 and 10.8.20--10.8.23. BGP 9.8/9.12 and 10.8/10.10
were read for the stronger almost-isometry and volume-convergence route;
those theorems are not prerequisites of the proof selected here.
The constants below are project estimates. Source locations, correction
checks and reading limits are in \code{reference_checks_revision81.md}.

\subsection{A nearby chart has the full dimension}

\begin{lemma}[Dimension transported from one bounded chart; LC10]
\label{lem:collapse-chart-global-dimension}
Let $X$ be a length space on which the curvature-$-1$ four-point
comparison holds globally. Suppose $q\in X$, $\rho>0$, and there is
an $L$-bilipschitz map from $B(q,\rho)$ onto an open subset of
$\mathbb R^m$, where $m\geq1$ and $L\geq1$.
Then $\dim_H X=m$.
\end{lemma}
\begin{proof}
Translate the chart to send $q$ to zero. Its image lies in
$[-L\rho,L\rho]^m$. For every $\sigma>0$, the finite grid argument
used in AC33 gives
\[
 P_{B(q,\rho)}(\sigma)
 \leq\left(2+\frac{4L^2\rho\sqrt m}{\sigma}\right)^m.
\]
For each integer $N\geq1$, put
\begin{align*}
 W_N&=\{y\in X:N^{-1}\leq d(q,y)\leq N\},\\
 t_N&=\min\{1/2,\rho/(2N)\},\qquad
 \lambda_N=\frac{t_NN}{\sinh N}>0.
\end{align*}
Apply AC32 with the common comparison set $\Omega=X$, $a=N^{-1}$,
$D=N$ and $t=t_N$. Its finite-set transport gives
\[
 P_{W_N}(\varepsilon)
 \leq\left(2+\frac{8L^2\rho\sqrt m}
                         {\lambda_N\varepsilon}\right)^m
 \leq C_N\varepsilon^{-m}\qquad(0<\varepsilon\leq1)
\]
for a finite $C_N$. A greedy $\varepsilon$-separated set must stop
within this bound and is then an internal $\varepsilon$-net.
AC39 therefore gives $\dim_H W_N\leq m$. No uniform bound on $C_N$
as $N\to\infty$ is needed. Since
$X=\{q\}\cup\bigcup_{N\geq1}W_N$, countable stability of Hausdorff
dimension gives $\dim_H X\leq m$. The chart's nonempty open image
has dimension $m$, so the Lipschitz coordinate map forces
$\dim_H X\geq m$. The annuli and chart balls carry restricted
ambient metrics; none is declared to be a length space.
\end{proof}

\begin{lemma}[A rank-three packet near the limit basepoint; LC11]
\label{lem:collapse-full-rank-packet}
Let $(X,x)$ be complete proper geodesic with global nonnegative
Alexandrov comparison and $2<\dim_H X\leq3$.
There are $q\in B(x,1/4)$ and six distinct anchors
$a_1,b_1,a_2,b_2,a_3,b_3\in X$, all different from $q$, whose
curvature-$-1$ comparison angles at $q$ satisfy the paired-packet
inequalities with quality
\[
 \tau_3=\frac1{320000}.
\]
There is no equal-radius or prescribed anchor-length requirement.
\end{lemma}
\begin{proof}
First the global zero-curvature comparison implies the global
curvature-$-1$ four-point inequality with the same distances.
For completeness, for a triangle with central sides $a,b>0$ and
Euclidean comparison angle $\theta_0$, put
$s=(1-\cos\theta_0)/2$. Its opposite side $c$ satisfies
$c^2=(1-s)(a-b)^2+s(a+b)^2$. Convexity of
$u\mapsto\cosh\sqrt u$ on $[0,\infty)$, seen from its power series,
gives
\[
 \cosh c\leq(1-s)\cosh(a-b)+s\cosh(a+b)
           =\cosh a\cosh b-\cos\theta_0\sinh a\sinh b.
\]
The hyperbolic cosine law therefore gives
$\theta_{-1}\leq\theta_0$, including degenerate angles by continuity.
Summing at a center transfers the four-point inequality.

Apply AC28 with $n=3$, $\Omega=X$ and $O=B(x,1/4)$.
Retain the maximal admissible rank $m$ and its witnessing anchors
from that proof, as well as its chart. LC10 shows that
$\dim_H X=m$. Since $2<\dim_H X\leq3$ and $m$ is an integer,
$m=3$. The retained packet has quality
$\tau_m=200^{m-4}/1600$, hence the stated value at $m=3$.
This uses the actual finite hierarchy, not an assertion that any
previously chosen lower-rank packet can be improved to arbitrary quality.
In particular no independent integral-dimension or regular-tangent
input has been added.
\end{proof}

\subsection{Lift one packet and keep one neighborhood}

\begin{lemma}[A uniform lifted comparison packet; LC12]
\label{lem:collapse-uniform-lifted-packet}
Let $(M_i,\widehat g_i,p_i)\to(X,x)$ be a convergent sequence
as in LC06, and suppose $2<\dim_H X\leq3$. Retain LC11's $q$ and anchors,
and put $\delta=1/600$.
There exist $0<a\leq A$, $0<r<1/4$ and a single tail on which
one can choose $q_i$ and six distinct anchor lifts so that
$d(p_i,q_i)<1/2$ and the following hold:
\begin{enumerate}[leftmargin=*]
\item all anchors and $B(q_i,2r)$ lie in one common
curvature-$-1$ comparison region for the ambient distance;
\item throughout $B(q_i,2r)$ the anchor distances belong to $[a,A]$
and the paired inequalities have quality $\delta$;
\item $\overline B(q_i,r)$ is complete.
\end{enumerate}
The numbers $a,A,r,\delta$ do not depend on the late index $i$.
No additional subsequence is taken.
\end{lemma}
\begin{proof}
Choose a fixed $R>2$ larger by more than $2$ than all distances
from $x$ to the seven chosen points. Pointed convergence supplies
actual fixed-radius approximation maps with errors tending to zero;
no prescribed $1/i$ rate is assumed. AC34 supplies simultaneous
lifts whose mutual distances converge to those of the chosen points.
All lifts eventually lie in $B(p_i,R)$ and are distinct, because the
finite limiting list has positive minimum separation.

On a tail $H_i>16R$ and $H_i>1$. The rescaled sectional curvature
is at least $-1$ on $B(p_i,16R)$. Apply the inherited local
Riemannian comparison interface and AC36 with curvature parameter $1$.
This gives ambient four-point comparison on $B(p_i,R)$.
The intrinsic open $16R$-ball has the required local dimension three:
smooth metric neighborhoods are locally bilipschitz to Euclidean
ones, and a countable cover supplies the intrinsic bound. Local
compactness and the complete buffers are supplied by completeness
of the Riemannian source. The open ball itself is not asserted complete.

Let $a_0$ and $A_0$ be the minimum and maximum distances from $q$
to its anchors. Set $a=a_0/2$ and $A=A_0+1$.
Since $\tau_3<\delta/4$, continuity of the finitely many
hyperbolic comparison-angle formulas gives $\xi>0$ such that
perturbing every relevant side length by less than $\xi$ preserves
quality $\delta$ and central distances in $[a,A]$.
Shrink $\xi$ if needed for the positive-distance bounds.
Choose $0<r<\min\{1/4,\xi/8\}$. On a single further tail,
all lifted side-length errors are less than $\xi/2$ and
$d(p_i,q_i)<1/2$.
For every $z\in B(q_i,2r)$ and every anchor lift $c_i$,
\[
 |d(z,c_i)-d(q,c)|
 \leq d(z,q_i)+|d(q_i,c_i)-d(q,c)|<2r+\xi/2<\xi.
\]
Anchor-to-anchor sides already have error less than $\xi/2$.
Thus the same packet and bounds work at every such $z$ and every
late index. Also $B(q_i,2r)\subset B(p_i,1)\subset B(p_i,R)$.
The closed $r$-ball is a closed subset of the complete source.
It is crucial that this $r$ was chosen from one limiting configuration,
rather than by invoking qualitative openness separately for each $i$.
\end{proof}

\begin{lemma}[A common cube in the distance-coordinate image; LC13]
\label{lem:collapse-uniform-image-cube}
In LC12, let
\begin{align*}
 F_i(z)&=(d(z,a_{1,i}),d(z,a_{2,i}),d(z,a_{3,i})),\\
 \mu&=1-3\delta^2-12\delta,\qquad
 t_0=\min\{1,a/2,\delta^2/K(a,A)\},\\
 e_*&=\min\{t_0,\mu r/2\}>0,
\end{align*}
where $K(a,A)$ is AC19's constant. Then $\mu\geq9/10$ and,
for every index in the common tail,
\[
 F_i(q_i)+[-e_*/4,e_*/4]^3
           \subset F_i\bigl(B(q_i,r/2)\bigr).
\]
Each $F_i$ is $\sqrt3$-Lipschitz for the Euclidean target norm.
The cube size is independent of $i$; no injectivity is required.
\end{lemma}
\begin{proof}
For a target $w$ in the displayed cube, its initial $\ell^1$ residual
is at most $3e_*/4<e_*$. Use AC22 successively with the
quality-$\delta$ packet on $B(q_i,2r)$.
The AC23 complete-buffer argument gives geometric decay of the
residual and a total displacement at most $e_0/\mu<r/2$.
Each step has length at most $e_0<t_0$ and stays in the common
packet region. Completeness of $\overline B(q_i,r)$ gives a limit
$z$ with $F_i(z)=w$ and $d(q_i,z)<r/2$.
The constants depend only on the fixed $a,A,r,\delta$, so the
same cube size works on the whole tail. Finally each distance
coordinate is $1$-Lipschitz, whence the Euclidean bound $\sqrt3$.
The same proof includes a target with zero initial residual.
\end{proof}

\subsection{The measure estimate and the LC06 implication}

We use normalized three-dimensional Hausdorff measure $\mathcal H^3$,
chosen to equal Lebesgue measure on Euclidean space, and its inherited
identification with Riemannian volume on a smooth three-manifold.
This is BBI Definition 5.5.1 and Theorem 5.5.5, printed 193--195/PDF
208--210, together with 1.7.7--1.7.12, printed 19--21/PDF 34--36.
The normalization is separate from the unnormalized contents used
in AC39, whose vanishing-measure conclusions are unaffected by a positive
constant. The comparison with the accepted PC volume remains a release
binding, not a new foundational library or an absence claim.

\begin{lemma}[A Lipschitz image cube bounds the source volume; LC14]
\label{lem:collapse-cube-volume-bound}
Let $U$ be a Borel subset of a smooth Riemannian three-manifold,
with its restricted ambient distance. Suppose
$F:U\to\mathbb R^3$ is $\sqrt3$-Lipschitz and its image contains
a Euclidean cube of side $e/2$, where $e>0$.
Under the stated volume identification,
\[
 \Vol(U)\geq\frac{(e/2)^3}{3\sqrt3}>0.
\]
Neither injectivity nor differentiability of $F$ is required.
\end{lemma}
\begin{proof}
For an $L$-Lipschitz map, the images of a countable cover have
diameters at most $L$ times the original diameters. Taking infima
over covers and then letting the cutoff tend to zero proves
$\mathcal H^3(F(U))\leq L^3\mathcal H^3(U)$.
This is BBI 1.7.8(4); it holds for Hausdorff outer measure, so no
measurability of the image is needed. Monotonicity, the normalized
Euclidean cube volume and the Riemannian identification give
\[
 (e/2)^3\leq\mathcal H^3(F(U))
        \leq3\sqrt3\,\mathcal H^3(U)=3\sqrt3\,\Vol(U).
\]
This argument does not invoke an area formula or a measure-convergence
theorem. It does not replace the ambient distance on $U$ by an
intrinsic path metric.
\end{proof}

\begin{theorem}[The localized full-dimensional volume implication; LC15]
\label{thm:collapse-local-volume-proved}
LC06 follows from the written Chapter 4 arguments and the retained
local Riemannian comparison and volume-identification interfaces.
More precisely, if the rescaled sequence in LC06 converges to $(X,x)$
and $\dim_H X>2$, then there are $v>0$ and an index $i_0$ such that
\[
 \Vol_{\widehat g_i}B_{\widehat g_i}(p_i,2)\geq v
                      \qquad(i\geq i_0).
\]
The bound may depend on the limit and its chosen finite configuration;
it is not uniform over all three-dimensional limits.
\end{theorem}
\begin{proof}
Apply AC40 with $n=3$ directly to the supplied pointed limit $X$,
using the rescaled curvature bounds of LC05. This gives
$\dim_H X\leq3$ without extracting a different limit. LC11 supplies
a rank-three packet near $x$; LC12
lifts it on one common neighborhood, and LC13 gives a common image
cube. Put $U_i=B(q_i,r/2)$ and
$v=(e_*/2)^3/(3\sqrt3)>0$.
LC14 gives $\Vol_{\widehat g_i}(U_i)\geq v$ for every late $i$.
Since $d(p_i,q_i)<1/2$ and $r<1/4$, this whole ball is contained
in $B(p_i,2)$. Monotonicity gives the desired inequality and hence
LC06's positive lower limit. All choices were made on the convergent
sequence itself; passing to a further subsequence would not by itself
prove the asserted lower limit.
\end{proof}

\paragraph{Effect on the model theorem.}
LC07--LC09 now use LC15 through LC06; they no longer have an unfilled
local volume proof as an additional premise. They still retain the
explicit source-qualified comparison, smoothing where used, and
inherited formal bindings. No theorem has been elaborated in Lean.
The argument proves only the positive-volume implication needed here,
not the full BGP volume-convergence theorem, an almost-isometry theorem,
or smooth fibration compatibility. The following section proves KL 7.2's
uniform exclusion of three-splittings; Section~\ref{sec:collapse-cone-scales}
then begins the cone-scale application.

\section{Excluding the three-stratum}
\label{sec:collapse-three-stratum-exclusion}

KL Definition 7.1 and Lemma 7.2, printed 46/PDF 41, use the
largest approximately split Euclidean rank at the modified volume
scale. The printed separation proof refers to a pointed GH distance
from $\mathbb R^3$. We write its uniform conclusion using the actual
maps of MC06 and AC49, retaining an arbitrary transverse factor in a
three-splitting. The numerical parameter in KL Definition 3.2 is not
itself a symmetric metric, so its errors are composed with the losses
proved in Chapter 3. The model compactness and product-limit steps
come from AC41 and AC50--AC51. Source comparisons and reading limits
are recorded in \mbox{\code{reference_checks_revision82.md}}.

\subsection{The finite rank convention}

\begin{definition}[Largest admitted rank at a chosen scale; LC16]
\label{def:collapse-maximal-rank}
Let $(M,d)$ be a metric space, let $\rho:M\to(0,\infty)$ be a
chosen scale function, and fix
\[
 0<\beta_1<\beta_2<\beta_3<1,\qquad \beta_0=0.
\]
For each $p\in M$, put $Z_p=(M,\rho(p)^{-1}d,p)$. For
$j\in\{1,2,3\}$, let $P_j(p)$ mean that $Z_p$ admits a normalized
$(j,\beta_j)$-splitting, and set
\[
 J(p)=\{0\}\cup\{j\in\{1,2,3\}:P_j(p)\},
 \qquad s(p)=\max J(p).
\]
The $k$-stratum is $\{p:s(p)=k\}$. Rank zero is supplied by the
exact identity $Z_p\cong\mathbb R^0\times Z_p$ of AC42; it is not
an invocation of AC49 at its excluded parameter zero. Equivalently,
rank $k$ is admitted and every larger rank in $\{0,1,2,3\}$ is
excluded at its own prescribed tolerance. The four sets are disjoint
and cover $M$, since every nonempty finite set $J(p)$ has a unique
maximum.
\end{definition}

This is a finite rank assignment. It does not assert that the sets
are smooth strata, that lower ranks are admitted at their smaller
tolerances, or that choices of splittings vary continuously with $p$.
For a Riemannian metric $g$, $Z_p$ uses the distance of
$\rho(p)^{-2}g$ through the inherited metric-scaling identification.
No continuity of $\rho$ is needed for this definition.

\subsection{A uniform separation with arbitrary product factors}

\begin{lemma}[A common limit through low-dimensional models; LC17]
\label{lem:collapse-model-limit-transfer}
Let $(C_i,c_i)$ be complete proper geodesic nonnegatively curved
Alexandrov spaces with $\dim_H C_i\leq2$. Suppose that arbitrary
pointed metric spaces $(Z_i,z_i)$ have whole-space pointed KL
$\sigma_i$-approximations
\[
 f_i:(Z_i,z_i)\longrightarrow(C_i,c_i),
 \qquad 0<\sigma_i<1,\qquad \sigma_i\longrightarrow0.
\]
There is a subsequence on which both $C_i$ and $Z_i$ converge to the
same complete proper geodesic nonnegative Alexandrov space $(C,c)$
with $\dim_H C\leq2$. No curvature, dimension, length, completeness
or properness assumption on $Z_i$ is required.
\end{lemma}
\begin{proof}
Apply AC41 to $C_i$ with $n=2$, curvature parameters
$a_i=1/(i+2)$ and controlled radii $r_i=i+2$. Global nonnegative
comparison supplies these weaker local lower bounds. Its local
dimension hypotheses hold for the intrinsic open balls: inside a
sufficiently small ball about any interior point, an ambient minimizing
segment between nearby points stays in the given open ball, by the
triangle inequality. The local intrinsic and ambient distances
therefore agree there. A countable such cover exists by properness
and separability. Hausdorff dimension on the intrinsic ball is thus
at most two by countable stability. The usual complete buffers come
from completeness of $C_i$; the open balls themselves need not be
complete. AC41 gives the stated subsequential limit $C$.

For any fixed $R>0$, choose $S=2R+4$. On this subsequence,
choose actual maps $g_i:C_i\to C$ satisfying
$\mathcal A_{S,e_i}$ on the closed $S$-ball with $e_i\to0$.
MC11 gives $\mathcal A_{S,3\sigma_i}(f_i)$ eventually, since
$3\sigma_i<S<\sigma_i^{-1}$. On a further tail,
\[
 2(3\sigma_i+e_i)<R,\qquad R+3\sigma_i\leq S.
\]
MC08 therefore gives
\[
 \mathcal A_{R,\,2(3\sigma_i+e_i)}(g_i\circ f_i).
\]
The intermediate images lie in the actual $S$-domain. This proves
$Z_i\to C$ in MC06's all-radii sense, with the same limit and
subsequence as the models. The tails and auxiliary maps may depend
on $R$; no prescribed convergence rate or fixed-error symmetry is
used. AC41's retained comparison qualification remains explicit.
\end{proof}

\begin{theorem}[Uniform exclusion of a three-splitting; LC18]
\label{thm:collapse-uniform-three-splitting-exclusion}
There exists $0<\eta<1/10$ such that the following data cannot
coexist, for any $0<\sigma,\beta\leq\eta$:
\begin{enumerate}[leftmargin=*]
\item a pointed metric space $(Z,z)$ and a whole-space pointed KL
$\sigma$-approximation $f:Z\to C$, where $(C,c)$ is complete,
proper, geodesic, nonnegatively curved and $\dim_H C\leq2$;
\item a normalized $(3,\beta)$-splitting
$\phi:Z\to\mathbb R^3\times Y$ with any pointed metric factor $Y$.
\end{enumerate}
The tolerance $\eta$ is independent of $Z,C,Y$ and of both maps.
The factor $Y$ need not be complete, proper, a length space or a
singleton. No numerical value of $\eta$ is claimed.
\end{theorem}
\begin{proof}
If the assertion fails, choose such data for every
$\eta_i=1/(i+20)$, $i\geq1$, with
$0<\sigma_i,\beta_i\leq\eta_i$. LC17 gives a subsequence on
which $Z_i$ converges to one complete proper $C$ with
$\dim_H C\leq2$. Apply AC50--AC51 to the supplied product maps
at the fixed rank $3$ and errors $\beta_i\to0$. Those lemmas
have no assumption that the rank be at most a source or limit
dimension. They give, after another subsequence, an onto pointed
isometry
\[
 \alpha:C\longrightarrow\mathbb R^3\times Y_\infty,
 \qquad \alpha(c)=(0,y_\infty).
\]
The map $u\mapsto\alpha^{-1}(u,y_\infty)$ is an isometric
embedding of $\mathbb R^3$ into $C$. Isometry invariance and
monotonicity of Hausdorff dimension give
\[
 3=\dim_H\mathbb R^3\leq\dim_H C\leq2,
\]
a contradiction. These dimension properties are BBI Proposition
1.7.19(1),(3),(4), printed 23/PDF 38. No product-dimension equality
for arbitrary factors is needed. AC51 constructs the product
metrically; the line-splitting theorem AC43 and long-strainer
producers are not used.
\end{proof}

Taking $Z=\mathbb R^3$ and its exact product with a singleton gives,
in particular, a uniform obstruction to KL approximations from
$\mathbb R^3$ to this class of models. The theorem above supplies
the stronger common-source conclusion actually needed by LC16.
It does not replace an arbitrary factor $Y$ by a point, or assert
that the infimum distance to the model class is attained. Nor is
$\eta$ identified with the printed $c/4$ in an unspecified
metrization of pointed GH topology.

\subsection{The local criterion and the standing sequence}

\begin{corollary}[A local curvature-and-volume exclusion; LC19]
\label{cor:collapse-local-no-three-splitting}
Fix $\eta$ from LC18. There are $0<w_0<c_3$ and $H\geq4$ such
that, if $(M,g,p)$ is a complete pointed Riemannian three-manifold,
$\rho>0$, $0<w<w_0$, $r_p(w)$ is finite and attained,
$r_p(w)\leq2\rho$, and
\[
 \sec_g\geq-(H\rho)^{-2}\quad\hbox{on }B_g(p,H\rho),
\]
then $(M,\rho^{-2}g,p)$ has no normalized $(3,\beta)$-splitting
for any $0<\beta\leq\eta$. Thus for any LC16 thresholds with
$\beta_3\leq\eta$, the rank at this point and scale is at most
two. The constants are independent of $M,p,\rho,w$ subject to
these hypotheses.
\end{corollary}
\begin{proof}
Use LC08 with the fixed approximation tolerance $\sigma=\eta/2$.
Its $w_0,H$ supply an actual KL $\sigma$-approximation to a
complete proper nonnegative Alexandrov model of dimension at most
two. LC18 excludes the proposed three-splitting, and LC16 gives
the rank conclusion. LC08 now uses the written LC15 producer
through LC06. Its inherited comparison and volume-identification
interfaces remain; this corollary supplies no new analytic input.
\end{proof}

\begin{theorem}[No three-stratum on the eventual collapsed tail; LC20]
\label{thm:collapse-eventual-no-three-stratum}
Fix $\Lambda>0$ and LC16 thresholds with $\beta_3\leq\eta$.
For the standing closed sequence and modified scales of LC09,
there is $w_0=w_0(\Lambda)>0$ such that, for every fixed
$0<w<\min\{w_0,c_3\}$, there is an index $\alpha_0$ for which
all $\alpha>\alpha_0$, all $p\in M^\alpha$ and every modified
scale satisfying LC02's bounds admit no normalized
$(3,\beta_3)$-splitting at that scale. Any choice of such a scale
function consequently has only ranks $0,1,2$ in LC16.
\end{theorem}
\begin{proof}
Apply LC09 with $\sigma=\eta/2$ and the fixed $\Lambda$ to
choose $w_0$, then fix $w$, and only then select the tail
$\alpha_0$. LC09 supplies a model and actual approximation for
every indicated point and scale on this one tail. LC18 rules out
a three-splitting at every one of them. LC16's finite maximum
therefore belongs to $\{0,1,2\}$. The tail need not work
simultaneously for every $w\downarrow0$. If a smooth modified
scale is required, its existence retains LC02's smoothing input;
the exclusion itself is pointwise and needs no derivative of the
scale function.
\end{proof}

This writes the conclusion of KL 7.2 in the standing-sequence
convention with its suppressed tail restored. A three-dimensional
smooth source can have only these coarse ranks at the chosen scale;
its manifold dimension has not changed. The result asserts no
circle fibration, regularity of strata, factor-diameter bound or
global graph-manifold conclusion. Existing PC release and source
compatibility gates remain. The following section begins the
Section 11 cone-scale analysis and its annular metric application;
smooth adapted coordinates and actual fibration compatibility
remain in Chapter 14.


\section{Cone scales and outward points on annuli}
\label{sec:collapse-cone-scales}

KL Section 11 begins with a cone approximation at a scale larger than
$\rho_p$, valid near every point, and later selects balls containing
zero-stratum points. We first write the metric scale-selection argument.
It has a separate input: a cone at infinity for each possible normalized
limit. Cone closure AC85 preserves radial maps already supplied on the
approximating spaces; it does not produce such maps on arbitrary
rescalings of a nonnegatively curved manifold.

\subsection{The limit model and the order of scales}

\begin{definition}[A supplied cone at infinity; LC21]
\label{def:collapse-cone-at-infinity}
For a complete proper pointed metric space $(N,n)$, a cone-at-infinity
package consists of a complete proper pointed metric space $(C,o)$,
radial maps satisfying AC82, and the following approximation data:
for every $0<\varepsilon<1$ there is $R_0>0$ such that for every
real $R\geq R_0$ there is an actual pointed KL
$\varepsilon$-approximation
\[
 (R^{-1}N,n)\longrightarrow(C,o).
\]
The notation $R^{-1}N$ divides distances by $R$; a Riemannian metric
tensor is divided by $R^2$. The target and its apex are fixed as $R$
varies. A one-point cone is allowed. Curvature, dimension and a smooth
structure on $N$ are additional properties when a consumer asks for them.
\end{definition}

KL Section 3.3, printed 23/PDF 18, asserts that a nonnegatively curved
Alexandrov space has a Tits cone, the pointed limit as $R\to\infty$,
and that this limit is again nonnegatively curved and a metric cone.
The passage states this fact without supplying its proof. BBI Definition
8.2.7, printed 276/PDF 291, defines a cone at infinity \emph{if the
limit exists}; that definition alone is not an existence theorem.
BBI uses a multiplier $\lambda\to0$, corresponding to $\lambda=R^{-1}$.
BBI Definition 8.2.1 uses the weaker term ``cone'' for pointed
self-similarity under all rescalings; that alone supplies no AC82
radial maps. The radial-map convention is translated by AC82--AC83, retaining the
angular cutoff correction to BBI 3.6.16. The Tits-cone existence and
cone-structure theorem remains a source-qualified input for the
application below; neither AC41 nor AC85 discharges it.

\begin{lemma}[Transfer at one fixed enlarged scale; LC22]
\label{lem:collapse-fixed-cone-scale-transfer}
Suppose $(X_i,p_i)\to(N,n)$ in the pointed sense of MC06--MC11.
Fix $R>0$ and $0<\delta<1$, and set
\[
 U=\delta^{-1}+\delta,\qquad S=2U+4,\qquad
 \varepsilon=\min\{\delta/100,\,1/(2S)\}.
\]
If there is an actual pointed KL $\varepsilon$-approximation
$g:(R^{-1}N,n)\to(C,o)$, then all sufficiently late
$(R^{-1}X_i,p_i)$ admit actual pointed KL
$\delta$-approximations to this same $(C,o)$.
\end{lemma}
\begin{proof}
Fixed rescaling preserves pointed convergence by MC10. Hence, on a
tail, choose
$f_i:\overline B_{R^{-1}X_i}(p_i,S)\to R^{-1}N$
satisfying $\mathcal A_{S,a_i}$ with $0<a_i<\delta/100$.
MC11 restricts $g$ to $\mathcal A_{S,3\varepsilon}$:
indeed $3\varepsilon<S<\varepsilon^{-1}$.
The buffer $S=2U+4$ ensures $U+a_i\leq S$, and
$2(a_i+3\varepsilon)<U$. MC08 therefore gives
\[
 \mathcal A_{U,\,2(a_i+3\varepsilon)}(g\circ f_i),\qquad
 2(a_i+3\varepsilon)<8\delta/100<\delta/4.
\]
Enlarging the error to $\delta/4$ only weakens distortion and shrinks
the required coverage ball. The converse conversion in MC11 now
restricts to the open $\delta^{-1}$-ball and extends by $o$ outside
it, giving the required whole-space map. Every composite is evaluated
inside its tested domain; the fixed $R$ is chosen before the tail.
\end{proof}

\begin{theorem}[A uniformly bounded interval containing a cone scale; LC23]
\label{thm:collapse-bounded-cone-scale}
Let $M^\alpha$ be pointed at any chosen $p\in M^\alpha$, with a
positive scale $\rho_\alpha(p)$. Assume the following sequential model
property: for every $\alpha_i\to\infty$ and every $p_i\in M^{\alpha_i}$,
the normalized spaces
$(\rho_{\alpha_i}(p_i)^{-1}M^{\alpha_i},p_i)$ have a subsequence
converging pointedly to a complete proper $(N,n)$ in a fixed class
$\mathcal N$, and every member of $\mathcal N$ has an LC21 package.
Then for every $0<\delta<1$ and every $T>1$, there are $V\geq T$
and $\alpha_0$ such that, for all $\alpha>\alpha_0$ and all
$p\in M^\alpha$, there are $s\in[T,V]$, a model
$(N_p,n_p)\in\mathcal N$ and its supplied cone $(C_p,o_p)$ with an
actual pointed KL $\delta$-approximation
\[
 ((s\rho_\alpha(p))^{-1}M^\alpha,p)\longrightarrow(C_p,o_p).
\]
Thus $r_p^0=s\rho_\alpha(p)$ lies in $[T\rho_\alpha(p),V\rho_\alpha(p)]$.
The constants are uniform in $p$ on this one tail.
\end{theorem}
\begin{proof}
If the assertion fails for fixed $\delta,T$, choose for every integer
$i\geq1$ an index $\alpha_i>i$ and a point $p_i$ at which no such
map exists for any $s\in[T,iT]$ and any model in $\mathcal N$.
The assumed property supplies a subsequence converging to one
$(N,n)\in\mathcal N$. Take $\varepsilon$ from LC22, and then use
LC21 to choose a single $R\geq T$ for which $R^{-1}N$ has a KL
$\varepsilon$-map to its supplied cone $C$. LC22 supplies KL
$\delta$-maps from the normalized sources further divided by this
fixed $R$, on a tail of the same subsequence. Eventually $R\leq iT$,
contradicting the chosen failure at $p_i$.
\end{proof}

Only one suitable scale in a bounded interval is asserted. There is no
uniform rate of convergence to infinity for every model in the class,
no claim that every $s\geq T$ works, and no continuous choice of $s$
or $C_p$. The limit, then $R$, then the source tail are chosen in that
order. Replacing $R$ by an arbitrary diverging $R_i$ would require
control of approximation domains growing at that rate, which pointed
convergence alone does not give. Compact limit models are included:
if $N$ is compact, $\operatorname{diam}(R^{-1}N)\to0$, so its LC21
package can use the point cone.

\subsection{The KL application and the remaining good-annulus inputs}

\begin{corollary}[Metric part of the good-annulus conclusion; LC24]
\label{cor:collapse-kl-metric-cone-scale}
Fix $\Lambda>0$ and a fixed $0<w<c_3$, with the positive
$w'$ and admissible modified scales $\rho_\alpha(p)$ of LC02.
In KL Standing Assumption 5.2, retain the full sequential smooth
compactness input used in Lemma 6.10(3): every pointed sequence at
these scales has a subsequence converging in the stated pointed
$C^K$ sense to a complete nonnegatively curved Riemannian
three-manifold, with the corresponding pointed GH convergence.
Also retain the Tits-cone input in LC21 for these limit manifolds.
Then for every $0<\delta<1$ and $T>1$ there are $V\geq T$ and
$\alpha_0$ such that every $\alpha>\alpha_0$ and every $p$ have
$r_p^0\in[T\rho_\alpha(p),V\rho_\alpha(p)]$ and an actual KL
$\delta$-map to the cone of such a limit manifold.
\end{corollary}
\begin{proof}
The retained smooth compactness statement and its GH component give
exactly LC23's sequential model property. Completeness and the
finite-dimensional Riemannian structure give properness of the limit;
this is part of the inherited geometric binding. Apply LC23 to this
class with its supplied Tits cones. The metric tensor at the final
source scale is $(r_p^0)^{-2}g^\alpha$.
\end{proof}

This is an explicit conditional use of smooth compactness, not a
proof of it from metric precompactness. KL 6.10(1)--(2), printed
42--43/PDF 37--38, supplies fixed-$w'$ noncollapse and bounds on
$\nabla^k\mathrm{Rm}$ for $0\leq k\leq K$, on controlled balls;
its proof of (3) invokes smooth compactness and passes the lower
curvature bound to the limit. LC04 provides the noncollapse component,
and LC02 the modified-scale bounds, but neither replaces derivative
control or the smooth compactness binding. The source has $K\geq10$.
Keep $\Lambda,w,w'$ and the standing derivative-control function fixed
before taking the subsequence or choosing $V$ and $\alpha_0$. These
constants may depend on that fixed data; no common bound or common
tail as $w'\downarrow0$ is asserted. PC-specific probing remains paused.

KL 11.1, printed 59--60/PDF 54--55, also asserts a diffeomorphism
$B(p,r_p^0)\simeq N_p$ and absence of critical points of the distance
function on $A(p,r_p^0/100,r_p^0)$. Its proof uses eventual topology
of large balls in the nonnegative limit and transfer through smooth
convergence. LC24 proves neither assertion: its model need not yet be
identified topologically with the source ball. The source proof's
opening negates conclusion (1), while its later contradiction carries
all three conclusions; a full formalization must negate the joint
failure and produce all witnesses at the same scale. Section 18's
replacement of smooth convergence by Alexandrov stability, printed
92--93/PDF 87--88, changes both hypotheses and outputs and is a
separate route. No stability, soul classification, ball diffeomorphism
or finite-time extinction is silently imported here. Remark 11.2's
low-dimensional \emph{conical} comparison also requires a separate
scale/limit argument; LC09 alone is not that argument.

\subsection{An outward point with an exact radial distance}

\begin{lemma}[Metric radial extension from a cone approximation; LC25]
\label{lem:collapse-cone-outward-point}
Let $(X,p)$ be a geodesic metric space, and let
$\phi:(X,p)\to(C,o)$ be an actual pointed KL
$\delta$-approximation to a cone with AC82 radial maps. Fix $0<a\leq b$
and assume
\[
 0<\delta<\min\{a/10,\,1/(4b+20)\}.
\]
For every $q$ with $a\leq r=d(p,q)\leq b$, there is $q'\in X$ with
\[
 d(p,q')=2r,\qquad r\leq d(q,q')<r+15\delta.
\]
In particular, choosing also $15\delta<\mu a$ gives
$d(q,q')<(1+\mu)r$. The same tolerance works for all these $q$.
\end{lemma}
\begin{proof}
Set $t=d_C(o,\phi(q))$. The tested basepoint pair gives
$|t-r|\leq\delta$, so $t>0$. By AC83 the point
\[
 z=H_{(2r+4\delta)/t}(\phi(q))
\]
has radius $2r+4\delta$ and
$d_C(\phi(q),z)=2r+4\delta-t$. This radius is less than
$\delta^{-1}-\delta$ by the assumed bound on $\delta$.
The KL coverage infimum therefore supplies $x\in B_X(p,\delta^{-1})$
with $d_C(\phi(x),z)<2\delta$; attainment at error $\delta$ is not
assumed. All pairs used below are inside the tested source ball.
Distortion and the triangle inequality give
\[
 2r+\delta<d_X(p,x)<2r+7\delta,\qquad
 d_X(q,x)<r+8\delta.
\]
Take the point $q'$ at distance $2r$ from $p$ on a minimizing
segment from $p$ to $x$. Then $d_X(x,q')<7\delta$, so
$d_X(q,q')<r+15\delta$. Finally the reverse triangle inequality
and $d_X(p,q')=2d_X(p,q)$ give $d_X(q,q')\geq r$.
If $C$ is a point, the nonempty annulus would force
$t\geq r-\delta>0$, a contradiction; thus that degenerate case
causes no missing witness. No curvature or differentiability was used.
\end{proof}

The first paragraph of KL 11.3's proof, printed 60/PDF 55, prints
$d(p,q')=2d(p,q)$ and the \emph{lower} estimate
$d(q,q')\geq(1-\mu)d(p,q)$ before asserting an almost-opposite angle
at $q$. That lower estimate follows already from the reverse triangle
inequality and is insufficient for the angle inference. For example,
in the Euclidean line take $p=0,q=1,q'=-2$: it satisfies the printed
distance requirements for $\mu>0$, but the two directions at $q$
coincide. LC25 independently supplies the needed \emph{upper} estimate
and an exact radial distance. This is a project proof repair, not an
author-issued correction; the inspected May 15, 2015 author sheet
lists no correction to 11.3.

The following section turns LC25 into uniformly almost-opposite
directions using lower-curvature comparison with complete buffers,
and records the exact distance-smoothing interface in KL 3.16/11.3. Keep the metric proof in Chapter 13; actual smooth adapted
coordinates and fibration compatibility remain in Chapter 14. The
large-ball topology needed for all of KL 11.1 stays explicit, with
user-confirmed smooth Schoenflies an inherited release candidate.


\section{Annular angles, radial smoothing and cutoffs}
\label{sec:collapse-radial-smoothing}

This section supplies the analytic part of KL 11.3, conditional on the
explicit smoothing input below. LC25 gives an outward point at exactly
twice the radius. We now bound the angle to \emph{every} minimizing
direction toward the center, verify the domains for smoothing, and
construct a smooth annular cutoff. The sublevel-set diffeomorphism in
KL 11.3(4) remains a separate topological assertion.

\subsection{A quantitative angle estimate with complete buffers}

\begin{lemma}[An almost-straight model triangle; LC26]
\label{lem:collapse-model-opposite-angle}
Let $0<r\leq b$, $0<\mu\leq1$, $r\leq\ell\leq(1+\mu)r$,
and $\kappa\geq0$ with $\kappa b\leq1/3$. The model angle
$\vartheta=\Theta_{-\kappa^2}(r,\ell,2r)$ satisfies
\[
 0\leq1+\cos\vartheta\leq6\mu.
\]
Degenerate triangles are included. In particular, if $0<\tau<\pi/2$
and $6\mu<1-\cos\tau$, then $\vartheta>\pi-\tau$.
\end{lemma}
\begin{proof}
Put $h=\ell-r\in[0,\mu r]$. The triangle inequalities hold because
$r\leq\ell\leq2r$. For $\kappa>0$, the hyperbolic cosine law,
with curvature parameter $-\kappa^2$, gives
\[
 1+\cos\vartheta=
 \frac{\cosh(\kappa(2r+h))-\cosh(2\kappa r)}
      {\sinh(\kappa r)\sinh(\kappa\ell)}.
\]
Since $\sinh t\geq t$, $\sinh t\leq t\cosh t$ for $t\geq0$,
and $\kappa(2r+h)\leq3\kappa b\leq1$, the numerator is at most
\[
 \kappa h\sinh(\kappa(2r+h))
 \leq \kappa^2 h(2r+h)\cosh1
 \leq 6\kappa^2\mu r^2.
\]
Here $h\leq\mu r\leq r$ and $\cosh1<2$. The denominator is at
least $\kappa^2r\ell\geq\kappa^2r^2$. For $\kappa=0$ the
Euclidean cosine law instead gives
\[
 1+\cos\vartheta=\frac{h(4r+h)}{2r(r+h)}
                 \leq\tfrac52\mu\leq6\mu.
\]
These formulas also give angle $\pi$ when $h=0$. Finally cosine
is strictly decreasing on $[0,\pi]$ and
$-1+6\mu<-\cos\tau=\cos(\pi-\tau)$.
\end{proof}

\begin{theorem}[One outward direction opposes all inward directions; LC27]
\label{thm:collapse-annular-directions}
Let $(M,g)$ be a complete connected smooth Riemannian manifold without
boundary. Suppose $\sec_g\geq-\kappa^2$ on $B_g(p,20b)$, where
$0<a\leq b$, $\kappa\geq0$ and $\kappa b\leq1/3$.
Let an actual pointed KL $\delta$-map from $(M,p)$ to an AC82 cone
be supplied. Fix $0<\tau<\pi/2$ and assume
\[
 0<\delta<\min\left\{\frac a{30},\frac1{4b+20},
                         \frac{a(1-\cos\tau)}{90}\right\}.
\]
For each $q$ with $a\leq d_g(p,q)\leq b$, let $V_q$ be the set
of initial \emph{unit} velocities of all minimizing segments from
$q$ to $p$. There is a unit vector $w_q\in T_qM$ such that
\[
 \angle(v,w_q)>\pi-\tau\quad(v\in V_q),\qquad
 \operatorname{diam}_{g_q}V_q\leq4\sin(\tau/2)<2\tau.
\]
In particular $q$ is noncritical for $d_p$ in KL Section 3.4's sense.
No continuous choice of $w_q$ is asserted. The local comparison
input recorded in AC02 remains source-qualified.
\end{theorem}
\begin{proof}
Complete Riemannian geometry supplies minimizing segments and properness.
LC25 gives $q'$ with $d(p,q')=2r$ and
$r\leq\ell=d(q,q')<r+15\delta$, where $r=d(p,q)$.
Choose one minimizing segment $q$ to $q'$ and let $w_q$ be its initial
unit velocity. Use LC26 with $\mu=15\delta/a<1/2$.
Its angle is greater than $\pi-\tau$ by the last bound on $\delta$.

Here are the domains needed for the local Toponogov input. Put
$\Omega=B(p,20b)$ with its intrinsic metric. Minimizing segments
$p$ to $q$, $p$ to $q'$ and $q$ to $q'$ have lengths at most $2b$
and lie in $B(p,3b)$. Distances between their vertices agree with
intrinsic distances in $\Omega$, and the chosen segments are also
intrinsically minimizing. For any vertex $v$, the intrinsic closed
$8b$-ball is complete: an intrinsic Cauchy sequence is ambient Cauchy,
and its ambient limit $z$ satisfies
\[
 d(p,z)\leq d(p,v)+8b\leq10b<20b.
\]
Local equality of intrinsic and ambient metrics inside $\Omega$
then gives intrinsic convergence. The same argument works for
vertices of radial subtriangles in $B(p,3b)$, with $11b<20b$.
Use $D=4b$, so all sides are strictly shorter than $D$ and these
complete balls have radius $2D$. This verifies the KL Section 3.3
safeguard; the open set $\Omega$ is not itself declared complete.
The curvature is nonpositive in the comparison model, so there is no
positive-curvature perimeter restriction.

Lower comparison implies that the actual angle at $q$ is at least
$\Theta_{-\kappa^2}(r,\ell,2r)$. This holds for every chosen
minimizing segment from $q$ to $p$, with the same outward segment.
The Riemannian identification of the segment angle with its tangent
angle is an inherited geometric input. Thus
$\|v+w_q\|<2\sin(\tau/2)$ for every $v\in V_q$.
The triangle inequality gives the displayed diameter bound; its
strict upper bound $2\tau$ follows from $\sin t<t$ for $t>0$.
Since $\tau<\pi/2$, $w_q$ makes angle greater than $\pi/2$ with
all inward directions, proving noncriticality directly from its
definition, without using a Morse or isotopy theorem.
\end{proof}

The small cone error and the curvature bound have separate roles.
Pointed GH closeness to a cone alone supplies neither an angle estimate
nor differentiability. The constants above are sufficient project
bounds, not values quoted from KL.

\subsection{What is supplied by the smoothing theorem}

For a Lipschitz function $f$ on a smooth Riemannian manifold, KL
Section 3.6 defines its generalized gradient at $q$ by transporting
nearby classical gradients radially to $T_qM$, taking their closed
convex hull, and intersecting as the neighborhood shrinks. Normal
neighborhoods are used for this transport. The resulting object is a
set of vectors, not an ordinary gradient at a nondifferentiable point.

\begin{foundation}[Localized distance smoothing; LC28]
\label{input:collapse-distance-smoothing}
For every $\varepsilon>0$ there is $\theta>0$ with the following
property. Let $M$ be complete connected and smooth Riemannian without
boundary, let $Y\subset M$ be nonempty and closed, and let
$U\subset M\setminus Y$ be open. At $q\in U$, let $V_q(Y)$ be
the initial unit velocities of all minimizing segments to nearest
points of $Y$. Suppose $\operatorname{diam}_{g_q}V_q(Y)<\theta$
for every $q\in U$. For every compact $C\subset U$ and every
$e>0$ there is a globally Lipschitz $F:M\to\mathbb R$ such that
\[
 \begin{gathered}
 F\text{ is smooth on an open neighborhood of }C,\qquad
 \|F-d_Y\|_\infty<e,\\
 F=d_Y\text{ on }M\setminus U,\qquad
 \operatorname{Lip}_g(F-d_Y)\leq\varepsilon.
 \end{gathered}
\]
This is the source-qualified input from KL Corollary 3.16(1)--(3),(5),
printed 26--27/PDF 21--22; it is not a proved Lean declaration.
\end{foundation}

The approximation error $e$ is chosen after the direction-diameter
threshold $\theta(\varepsilon)$; arbitrarily small value error does
not replace the global Lipschitz bound on the \emph{difference}.
The difference is real-valued and may have either sign. Our current
consumer does not need the additional angular clause 3.16(4).

KL's proof identifies the generalized gradient of $d_Y$ with the
negative convex hull of $V_q(Y)$, uses upper semicontinuity under
parallel transport, and enlarges that set to an open fiberwise-convex
set with small fibers. Lemma 3.14 mollifies inside this set, agrees
with the original function off $U$, and smooths near the compact set.
The supplied proof explicitly omits the mollification details and
refers to Grove--Shiohama. Its final global Lipschitz conclusion uses
an almost-everywhere gradient bound for the difference. These remain
parts of the smoothing interface to bind to inherited analysis or to
expand from its sources; no claim is made that generic value-only
smooth approximation suffices, or that this analysis is absent from PC.

\begin{lemma}[A smooth distance perturbation has nonzero gradient; LC29]
\label{lem:collapse-distance-perturbation-gradient}
Let $M$ be complete connected smooth Riemannian without boundary,
$p\in M$, and $0\leq\varepsilon<1$. Suppose $F:M\to\mathbb R$
is smooth on an open $W\subset M\setminus\{p\}$ and
$F-d_p$ is globally $\varepsilon$-Lipschitz. Then for all $q\in W$,
\[
 1-\varepsilon\leq\|\nabla F(q)\|_g\leq1+\varepsilon.
\]
The use of Rademacher's theorem and the usual Riemannian
measure/calculus facts retains their inherited formal bindings.
\end{lemma}
\begin{proof}
The distance is $1$-Lipschitz. At any point $x\ne p$ where it is
differentiable, its gradient has norm at most one. Along a unit-speed
minimizing segment from $x$ to $p$, its one-sided derivative at $x$
is $-1$. Differentiability and the chain rule therefore force
$\|\nabla d_p(x)\|=1$. Rademacher's theorem gives such points
almost everywhere. At each of them in $W$, differentiating the
Lipschitz difference gives
$\|\nabla F(x)-\nabla d_p(x)\|\leq\varepsilon$, hence both bounds.
Every nonempty open Riemannian set has positive volume, so this
full-measure set is dense in $W$. Continuity of $\nabla F$ and its
norm extends the inequalities to all of $W$. This argument requires
no first-variation formula through a cut point and no assertion that
$d_p$ is smooth on $W$. BBI 5.5.7, printed 196/PDF 211, states the
Rademacher input without proof and refers to Federer.
\end{proof}

\subsection{The radial function at a supplied cone scale}

\begin{theorem}[A radial function and its noncritical level band; LC30]
\label{thm:collapse-radial-function}
For every $0<\varepsilon<1/4$ there is $\delta_0>0$ with the
following property, conditional on LC28 and the retained local
comparison and analytic inputs. Let $(M,g,p)$ be complete connected
smooth Riemannian without boundary, with
\[
 \sec_g\geq-1/60^2\quad\hbox{on }B_g(p,400),
\]
and suppose an actual KL $\delta$-map to an AC82 cone is supplied,
where $0<\delta<\delta_0$. For every $0<e<1/40$ there is a
Lipschitz $\eta:M\to[0,\infty)$ satisfying
\[
 \begin{gathered}
 U=\{1/20<d_p<20\},\qquad C=\{1/10\leq d_p\leq10\},\\
 \eta\text{ is smooth near }C,\qquad
 \|\eta-d_p\|_\infty<e,\qquad
 \operatorname{Lip}_g(\eta-d_p)\leq\varepsilon,\\
 \eta=d_p\text{ on }M\setminus U,\qquad
 1-\varepsilon\leq\|\nabla\eta\|_g\leq1+\varepsilon
                     \text{ on }C.
 \end{gathered}
\]
Moreover $\eta^{-1}([1/5,2])$ is contained in $C$, so $\eta$ is
smooth and has no critical points on a neighborhood of this level band.
No diffeomorphism type of a sublevel set is asserted.
\end{theorem}
\begin{proof}
Choose $\theta(\varepsilon)$ from LC28 and put
$\tau=\min\{\pi/4,\theta/4\}$. Set $a=1/20$, $b=20$ and take
$\delta_0$ to be the positive minimum in LC27. With $\kappa=1/60$,
$\kappa b=1/3$. LC27 gives
$\operatorname{diam}V_q\!<2\tau\leq\theta/2<\theta$ on $U$.
Properness makes $C$ compact, and $C\subset U$ with positive margins.
Apply LC28 with $Y=\{p\}$, this $U,C$ and the prescribed $e$.
The resulting $F$ equals the nonnegative $d_p$ off $U$; on $U$,
$F>d_p-e>1/20-1/40>0$. Thus it defines the stated $\eta$, with
$\eta(p)=0$. LC29 gives its gradient bounds near $C$ after shrinking
the smooth neighborhood to avoid $p$.
Finally, $\eta(q)\in[1/5,2]$ implies
\[
 1/5-e<d_p(q)<2+e,
 \qquad [1/5-e,2+e]\subset(1/10,10).
\]
This gives a neighborhood of the level band inside the smooth region;
$1-\varepsilon>0$ rules out critical points there.
\end{proof}

KL 2.2 defines $A(p,r,R)=\overline{B(p,R)}\setminus B(p,r)$,
with the inner ball open and the outer ball closed. In this complete
Riemannian setting the annulus is compact, though generally not open.
Corollary 3.16 requires an open $U$ and a compact $C$. LC30 uses a
buffered open region $U$ around the closed radial shell $C$, so it
supplies all the points in KL 11.3's smaller annulus and verifies the
actual hypotheses of the smoothing input. A small value error alone
would not prevent critical points; LC29 uses the Lipschitz bound on
the difference. Revision 93 corrects the earlier transcription of the
source annulus; LC30's statement and proof are unchanged.

\subsection{The eventual scale and a cutoff with checked endpoints}

Fix a smooth nonincreasing $\phi:\mathbb R\to[0,1]$, equal to one
on $(-\infty,0]$ and zero on $[1,\infty)$. For $a<b$ define
\[
 \chi_{a,b}(t)=\phi\!\left(\frac{t-a}{b-a}\right),\qquad
 \Phi(t)=\chi_{-3/10,-1/5}(-t)\,\chi_{4/5,9/10}(t).
\]
Then $\Phi=0$ for $t\leq1/5$ or $t\geq9/10$, and $\Phi=1$
on $[3/10,4/5]$. It is smooth, valued in $[0,1]$, and has a finite
constant $L_\Phi=\sup_t|\Phi'(t)|$.
These endpoint properties follow by substituting into $\phi$.
The existence of a smooth scalar cutoff is an inherited calculus input.

KL equation (2.1), printed 19/PDF 14, instead prints
$\phi(a+(b-a)t)$ while claiming the endpoints $a,b$. Those statements
do not agree: for $a=2,b=3,t=2$ the printed expression is
$\phi(4)=0$ where the claimed value is one. We use the inverse affine
change above. This is a directly checked project repair, not a
correction listed in the inspected May 15, 2015 author sheet.

\begin{theorem}[Eventual radial function and compact annular cutoff; LC31]
\label{thm:collapse-eventual-radial-cutoff}
Retain LC24's fixed data, smooth compactness and Tits-cone inputs,
and the smoothing/comparison inputs of LC30. Fix
$0<\varepsilon<1/4$, $0<e<\min\{\varepsilon,1/40\}$ and $T>1$.
There are $V\geq T$ and $\alpha_0$ such that for every
$\alpha>\alpha_0$ and every $p\in M^\alpha$ there are a scale
$r_p^0\in[T\rho_\alpha(p),V\rho_\alpha(p)]$ and a radial function
$\eta_p$ satisfying LC30 in the metric $\widehat g=(r_p^0)^{-2}g^\alpha$.
The function $\zeta_p=\Phi\circ\eta_p$ is globally smooth and compactly
supported, with $0\leq\zeta_p\leq1$, equal to one wherever
$3/10\leq\eta_p\leq4/5$, and
\[
 \begin{gathered}
 \operatorname{supp}\zeta_p
 \subset\{1/5-e\leq d_{\widehat g}(p,\cdot)\leq9/10+e\},\\
 \|\nabla\zeta_p\|_{\widehat g}\leq L_\Phi(1+\varepsilon).
 \end{gathered}
\]
The scale, cone map, radial function and cutoff refer to the same
pointed source. No ball/sublevel topology or adapted-fibration
compatibility is asserted.
\end{theorem}
\begin{proof}
Choose $\delta=\delta_0(\varepsilon)/2$ in LC24, decreasing it
further below one if needed, and obtain $V$ and a first tail. The
curvature bound required by LC30 can be imposed on a later tail
\emph{after} $V$ is fixed. Indeed eventually $1/\alpha<w'$, and
volume-scale monotonicity, LC02 and Standing Assumption 5.2(1) give
\[
 \rho_\alpha(p)\leq2r_p(w')\leq2r_p(1/\alpha),\qquad
 R_p\geq\alpha r_p(1/\alpha)\geq(\alpha/2)\rho_\alpha(p).
\]
The case $R_p=\infty$ satisfies the same finite-radius conclusions.
Put $H=\alpha/2$. At scale $\rho_\alpha(p)$, the sectional curvature
is at least $-H^{-2}$ on the $H$-ball. Writing
$s=r_p^0/\rho_\alpha(p)\in[T,V]$, at the cone scale it is at least
$-(s/H)^2$ on the $(H/s)$-ball. Increase the tail until $H>400V$.
Then $H/s>400$ and $s/H<1/400<1/60$, so LC30 applies to the
\emph{same} cone-scale map selected by LC24. This is uniform over
points at the fixed $w,w'$; no common tail as $w'\downarrow0$ is claimed.

Take the function $\eta_p$ from LC30 and compose with $\Phi$. If
$\eta_p(q)\in[1/5,9/10]$, its distance from $p$ lies between
$1/5-e$ and $9/10+e$, strictly inside the inner region $C$ of LC30.
There $\eta_p$ is smooth and the chain rule bounds the gradient.
At every other point outside this closed preimage, continuity makes
$\Phi\circ\eta_p$ identically zero on a neighborhood. Thus the
composition is smooth everywhere, including transition endpoints,
and its support lies in the displayed compact annulus. Properness,
not compactness of the whole manifold, suffices for this support claim.
The gradient estimate follows from LC30 on that annulus and from zero
outside it.
\end{proof}

In original units, the same scalar function satisfies
\[
 \operatorname{Lip}_{g^\alpha}
   (\eta_p-d_{g^\alpha}(p,\cdot)/r_p^0)
       \leq\varepsilon/r_p^0,\qquad
 \|\nabla\zeta_p\|_{g^\alpha}
       \leq L_\Phi(1+\varepsilon)/r_p^0.
\]
LC30--LC31 give KL 11.3(1)--(3),(5) and the noncritical portion of (4),
with explicit smoothing and inherited bindings. They do not establish
that the sublevel sets are disk bundles over a soul, nor identify the
source ball with the model $N_p$. The following section proves sublevel comparisons and specifies the
remaining model-core transfer for 11.1/11.3; selection in 11.5 follows.
Chapter 13 owns these radial scalar functions; Chapter 14 owns their
use as adapted coordinates and the compatibility of actual fibrations.


\section{Radial sublevels and the model-core transfer}
\label{sec:collapse-sublevel-topology}

We now prove a sublevel comparison from the radial function already
constructed. The comparison with a distance ball is a homeomorphism;
the comparison between smooth transverse cores is a diffeomorphism.
The remaining geometric task is to produce such a core from the
\emph{same} nonnegative model and at the \emph{same} selected scale.
This last task is recorded as data to be supplied, not as a consequence
of metric cone approximation.

\subsection{A compact smooth band and its flow coordinates}

Throughout the next five statements, retain an LC30 radial function
on a complete connected smooth Riemannian manifold without boundary.
Write $d=d_g(p,\cdot)$, $0<\varepsilon<1/4$, $0<e<1/40$, and set
\[
 a=1/8,\qquad b=3,\qquad
 A_t=\{\eta\leq t\},\qquad K=\eta^{-1}([a,b]).
\]
The curvature and smoothing inputs used to produce $\eta$ remain
those of LC30. The arguments below only use its stated output and
inherited regular-level, cutoff and ordinary differential-equation theory.

\begin{lemma}[Proper sublevels and a smooth enlarged band; LC32]
\label{lem:collapse-proper-sublevels}
The function $\eta$ is proper, $K$ is compact, and $\eta$ is smooth
with $1-\varepsilon\leq\|\nabla\eta\|\leq1+\varepsilon$ on an open
neighborhood of $K$. For each $t\in[a,b]$, $A_t$ is a compact smooth
codimension-zero domain, with
\[
 \operatorname{int}A_t=\{\eta<t\},\qquad
 \partial A_t=\eta^{-1}(t).
\]
An empty level and an empty boundary are allowed. Smoothness of
$\eta$ in the interior near $p$ is not required.
\end{lemma}
\begin{proof}
A closed sublevel $A_t$ is contained in the closed distance ball of
radius $t+e$; properness of the Riemannian distance makes it compact.
The same argument applies to inverse images of compact subsets of
the real line, proving properness. On $K$,
\[
 1/10<a-e<d<b+e<10.
\]
Thus $K$ lies strictly inside the smooth region of LC30, and its
gradient bound follows there from LC29. At any point where $\eta=t$,
the regular-level theorem supplies a smooth local coordinate
$\eta-t$. In that chart $A_t$ is a closed half-space. At a point
where $\eta<t$, continuity puts an entire ambient neighborhood in
$A_t$; this neighborhood inherits the ambient smooth structure,
even if $\eta$ itself is nonsmooth there. These charts also prove
both assertions about interior and boundary. Compactness was already
proved, and the empty-level case needs no boundary chart.
\end{proof}

\begin{lemma}[Normalized gradient product on the band; LC33]
\label{lem:collapse-radial-flow-product}
Let $\Sigma=\eta^{-1}(1)$. There is a smooth product diffeomorphism
of manifolds with boundary
\[
 F:\Sigma\times[a,b]\longrightarrow K,\qquad
 \eta(F(x,u))=u,\qquad F(x,1)=x.
\]
On this band its coordinate curves have velocity
$X=\nabla\eta/\|\nabla\eta\|^2$ and speed at most
$(1-\varepsilon)^{-1}$. In particular, for $a\leq s<t\leq b$,
\[
 d(F(x,t))-d(F(x,s))
 \geq c_\varepsilon(t-s),\qquad
 c_\varepsilon=\frac{1-2\varepsilon}{1-\varepsilon}>0.
\]
If $\Sigma$ is empty, $K$ is empty and the product statement has
its usual empty interpretation.
\end{lemma}
\begin{proof}
Choose a relatively compact open neighborhood of $K$ on which
$\eta$ is smooth with nonzero gradient. A smooth bump equal to one
near $K$ extends $X$ to a compactly supported smooth vector field
$\widetilde X$ on $M$. Its flow $\varphi_t$ is complete by inherited
ODE theory. While a curve is in $K$, its $\eta$-value increases at
unit speed. The same identity holds on a neighborhood of $K$, so a
first-exit argument shows that a curve starting at level $u$ remains
in $K$ for all times between $a-u$ and $b-u$.
Consequently
\[
 F(x,u)=\varphi_{u-1}(x),\qquad
 F^{-1}(q)=\bigl(\varphi_{1-\eta(q)}(q),\eta(q)\bigr).
\]
Both formulas extend smoothly to neighborhoods of the respective
boundary levels; the flow law proves they are inverse. The inverse
formula also shows that a nonempty $K$ forces a nonempty $\Sigma$.
The norm bound on $\nabla\eta$ gives the speed bound. Put
$h=\eta-d$, which is globally $\varepsilon$-Lipschitz. Along a
coordinate curve,
\[
 \begin{split}
 d(F(x,t))-d(F(x,s))
 &=t-s-h(F(x,t))+h(F(x,s))\\
 &\geq (t-s)-\varepsilon\operatorname{Length}
                         (F(x,\cdot)|_{[s,t]})\\
 &\geq\left(1-\frac{\varepsilon}{1-\varepsilon}\right)(t-s).
 \end{split}
\]
This finite-difference argument requires no derivative of the distance
along the flow, including at cut points.
\end{proof}

The source for the inherited flow machinery is John M. Lee,
\emph{Introduction to Smooth Manifolds}, in the archived second-edition TeX: \code{c09.tex}, labels \code{theorem:flows},
\code{theorem:cpctly-supp-complete} and \code{thm:time-dept-flow}.
The regular-level proof is in \code{c05.tex},
\code{cor:reg-level-set}; the regular-sublevel proposition there leaves
its proof as an exercise. LC32 supplies the needed local argument.
The January 3, 2026 author errata corrects the flow proof's exit-time
formula on printed page 213; that correction is retained in our source
review. No new flow or manifold foundation is being proposed.

\subsection{Distance balls and smooth graphs have different outputs}

\begin{theorem}[Distance-sublevel homeomorphism with small tracks; LC34]
\label{thm:collapse-distance-sublevel-homeomorphism}
For every $\rho\in[1/5,2]$, there is a continuous ambient isotopy
$H_\lambda:M\to M$, $0\leq\lambda\leq1$, consisting of
homeomorphisms, with $H_0=\mathrm{id}$, equal to the identity off $K$,
such that
\[
 H_1(A_\rho)=\overline B_g(p,\rho),\qquad
 d_g(H_\lambda(q),q)<\frac e{1-\varepsilon}
                  \quad(q\in M).
\]
It also carries $\partial A_\rho$ to the distance sphere and the
interior to the open distance ball. The distance sphere is not
asserted smooth, and this isotopy is not asserted to be smooth.
\end{theorem}
\begin{proof}
For each $x\in\Sigma$, LC33 makes $u\mapsto d(F(x,u))$ continuous
and strictly increasing, with lower slope $c_\varepsilon$.
Its values at $a,b$ are respectively less than $a+e<1/5$ and greater
than $b-e>2$. There is a unique $h_\rho(x)\in(a,b)$ with
$d(F(x,h_\rho(x)))=\rho$, and
\[
 |h_\rho(x)-\rho|<e.
\]
The lower-slope bound proves continuity: for fixed $x_0$, it bounds
$|h_\rho(x)-h_\rho(x_0)|$ by
$c_\varepsilon^{-1}|d(F(x,h_\rho(x_0)))-\rho|$.

Set $h_\lambda(x)=(1-\lambda)\rho+\lambda h_\rho(x)$.
On each $[a,b]$ fiber use the increasing piecewise affine map
\[
 P_{\lambda,x}(u)=
 \begin{cases}
 a+\dfrac{h_\lambda(x)-a}{\rho-a}(u-a),&a\leq u\leq\rho,\\[4pt]
 h_\lambda(x)+\dfrac{b-h_\lambda(x)}{b-\rho}(u-\rho),&\rho\leq u\leq b.
 \end{cases}
\]
It fixes both endpoints, sends $\rho$ to $h_\lambda(x)$, and has
$|P_{\lambda,x}(u)-u|<e$. The formulas and their inverses are
continuous in all parameters. Conjugate by $F$ and extend by the
identity off $K$. Endpoint agreement proves the pasting is continuous,
also for the inverse and the parameter $\lambda$. At $\lambda=1$,
strict monotonicity along each fiber gives exactly the desired
sublevel, level and strict-sublevel sets in $K$. Off $K$, the inequalities
$a+e<\rho<b-e$ give the same classifications for $d$ and $\eta$.
The speed bound in LC33 proves the track estimate. If $K$ is empty,
these off-band classifications prove the result with the identity.
\end{proof}

\begin{lemma}[Smooth radial graph isotopy; LC35]
\label{lem:collapse-smooth-graph-isotopy}
Let $\rho\in(a,b)$ and let $h:\Sigma\to(a,b)$ be smooth. Define
\[
 D_h=A_a\ \cup\ \{F(x,u):x\in\Sigma,\ a\leq u\leq h(x)\}.
\]
There is a smooth ambient isotopy, supported in a compact subset of
$\eta^{-1}((a,b))$, carrying $A_\rho$ to $D_h$. In particular all
$A_\rho$, $\rho\in[1/5,2]$, are mutually diffeomorphic by smooth
ambient isotopies.
\end{lemma}
\begin{proof}
If $\Sigma$ is empty, the assertion is the identity. Otherwise
compactness of $\Sigma$ gives a positive margin from both endpoints
for $\rho$ and every value of $h$. Put
$h_\lambda=(1-\lambda)\rho+\lambda h$ and choose a smooth scalar
bump $\chi$, equal to one near zero, whose support has radius smaller
than half that margin. On the product band use the time-dependent
vertical vector field
\[
 W_\lambda(x,u)=
 (h(x)-\rho)\chi(u-h_\lambda(x))\,\partial_u.
\]
Push it forward by $F$ and extend by zero. This is smooth and has
one common compact spatial support strictly inside the band. It
extends in time to a neighborhood of $[0,1]$, so inherited
time-dependent ODE theory, with compactness preventing finite-time
escape, gives a smooth ambient isotopy on this interval. The curve
$u(\lambda)=h_\lambda(x)$ solves the ODE starting at $\rho$.
Uniqueness makes each fiber map increasing and fixes the regions
near its endpoints; hence it maps the portion below $\rho$ onto the
portion below $h(x)$. Outside the band the map is the identity.
Taking $h$ to be a constant gives the last assertion. No smoothness
of the function $h_\rho$ in LC34 is used here.
\end{proof}

\subsection{The precise core that a model must supply}

\begin{lemma}[Outward transversality recognizes a smooth core; LC36]
\label{lem:collapse-transverse-core}
Suppose $D\subset M$ is a compact smooth codimension-zero domain,
\[
 A_a\subset\operatorname{int}D,\qquad D\subset\{\eta<b\},
\]
and $X=\nabla\eta/\|\nabla\eta\|^2$ is strictly outward transverse
to $\partial D$. Precisely, every boundary point has a local smooth
defining function $f$ with $D=\{f\leq0\}$ and $df(X)>0$ there.
Then $D=D_h$ for a smooth $h:\Sigma\to(a,b)$, and $D$ is smoothly
ambient isotopic to every $A_\rho$ with $\rho\in[1/5,2]$.
\end{lemma}
\begin{proof}
The enclosure places $\partial D$ strictly inside the band. Each
fiber starts inside $D$ and ends outside it. At a boundary crossing,
$\partial_u(f\circ F)=df(X)>0$, so the crossing is from inside to
outside. The crossings on a fixed compact fiber form a closed set
and each is isolated by this derivative test, hence there are finitely
many. There is at least one crossing, and there cannot be a second:
that would require an intervening outside-to-inside crossing.
Thus the boundary meets each fiber exactly once. The implicit
function theorem makes its height $h(x)$ locally smooth; uniqueness
patches these heights to one smooth function. The enclosure also
excludes extra pieces below or above the band, so $D=D_h$. Apply
LC35. In the empty-band case the enclosure directly gives $D=A_a$
and all these sublevels agree.
\end{proof}

\begin{definition}[A model-core transfer packet; LC37]
\label{def:collapse-model-core-packet}
For a chosen LC31 scale, cone map and radial function, retain the
\emph{same} nonnegative model $N$ whose cone occurs in the LC24
witness. A model-core transfer packet consists of one of these
alternatives, with actual maps in supplied smooth structures:
\begin{enumerate}
\item $N$ is compact, there is a smooth embedding $j:N\to M$
between connected three-manifolds, and $\eta<1/5$ on all of $M$.
\item $N$ is noncompact. There are a supplied compact soul $S$,
a compact smooth domain $D_N\subset N$ with a specified
 diffeomorphism $D_N\simeq D(\nu S)$ to the closed unit normal disk
bundle, and a smooth embedding $j:O\to M$ from an open neighborhood
$O$ of $D_N$. The image $D=j(D_N)$ satisfies the enclosure and
strict outward transversality in LC36 for this very $\eta$ and $X$.
\end{enumerate}
This definition does not assert existence of either alternative.
In the noncompact alternative the normal bundle has positive rank;
its disk bundle is compact because its base is compact.
\end{definition}

\begin{theorem}[Conditional soul type of every radial sublevel; LC38]
\label{thm:collapse-soul-sublevel-type}
Given an LC37 packet, every $A_\rho$ for $\rho\in[1/5,2]$ is
 diffeomorphic to $N$ in its compact alternative and to $D(\nu S)$
in its noncompact alternative. In the latter case the comparison
with the embedded model core is a smooth ambient isotopy. In either
case LC34 gives a homeomorphism of the corresponding distance ball
with this compact domain, including its boundary and interior.
\end{theorem}
\begin{proof}
In the compact case an equal-dimensional smooth embedding is open
by the inverse function theorem and has closed image by compactness.
Its nonempty image is all of connected $M$, so it is a diffeomorphism.
The strict global bound on $\eta$ gives $A_\rho=M$. In the noncompact
case $j$ identifies $D_N$ smoothly with $D$; LC36 identifies $D$
with $A_\rho$, and composition with the supplied disk-bundle map
proves the claim. LC34 supplies the distance-ball comparison.
\end{proof}

\subsection{Source comparison and the remaining producer}

KL 11.3(4), printed 60--61/PDF 55--56, first compares the distance
and smoothed sublevels by a common increasing flow. It then identifies
interiors and invokes a compact three-manifold boundary-recognition
argument, citing Schoenflies and Stallings. A homeomorphism obtained
in this way must still be bound to the requested differentiable
conclusion; we do not silently promote a topological map to a smooth
one. Smooth Schoenflies remains a user-confirmed inherited PC candidate.

Remark 11.4 instead describes smoothing large noncritical distance
balls in the nonnegative model and relating their smooth isotopy class
to a tubular neighborhood of the soul. LC35--LC38 make a sufficient
\emph{direct smooth-transfer interface} for that route. They do not
claim that the remark already supplies our enclosure, outward
transversality or source embedding, nor that an arbitrary interior
homeomorphism supplies them. The soul theorem cited in KL 3.11 is
Cheeger--Gromoll (1972); its proof and the extra isotopy-class assertion
are not newly verified here. The normal disk bundle and its map remain
explicit packet data.

The next producer must transfer a whole neighborhood of a smooth model
core and check enclosure and outward transversality against the selected
radial function. A $C^K$ metric limit or convergence map does not by
itself supply LC37's smooth maps: the differentiable-category binding
and $C^1$ control remain obligations. The compact alternative needs
a global embedding and the small radial range. Noncriticality alone
supplies neither alternative.

Uniform existence must negate the \emph{joint} cone-map, radial-function
and model-core conclusion. One fixed limit, one large radius and one
later tail must supply all witnesses; separately chosen scales do not
suffice. LC37 is a sufficient route target, not a proved producer or a
reformulation of all KL 11.1. Producing it and then KL 11.5 ball selection
remain Chapter 13 work. Chapter 14 owns adapted coordinates and
compatible fibrations. All inherited and source-qualified gates remain
explicit.

\section{Transferring a core through a buffered smooth comparison}
\label{sec:collapse-core-stability}

We now give explicit tests for the source-side part of LC37. A smooth
comparison on a larger ball supplies enclosure. A separate estimate
on the radial differential supplies outward transversality. The metric
and the scalar function have different convergence obligations.
The soul type of a model core remains a supplied geometric input.

\subsection{A larger embedded ball rules out short exits}

\begin{lemma}[Buffered distance and ball coverage; LC39]
\label{lem:collapse-buffered-embedding}
Let $(N,g,n)$ and $(M,\widehat g,p)$ be complete connected smooth
Riemannian manifolds without boundary of the same dimension. Let
$L>0$, $0<\lambda<1$, and let $j:O\to M$ be a smooth embedding,
where $O\subset N$ is open and contains $\overline B_g(n,L)$,
with $j(n)=p$. Suppose on that closed ball
\[
 (1-\lambda)^2g\leq h=j^*\widehat g\leq(1+\lambda)^2g.
\]
Then
\[
 B_{\widehat g}(p,(1-\lambda)L)\subset j(B_g(n,L)).
\]
If $0<r<L$ and $(1+\lambda)r<(1-\lambda)L$, then, for
$x\in\overline B_g(n,r)$,
\[
 (1-\lambda)d_g(n,x)\leq d_{\widehat g}(p,j(x))
                   \leq(1+\lambda)d_g(n,x).
\]
No global distance preservation follows merely from local metric
closeness; the larger domain is part of the assertion.
\end{lemma}
\begin{proof}
The embedding is a diffeomorphism onto an open image. Its restriction
to the compact closed $L$-ball is a homeomorphism onto a compact set.
Consequently every boundary point of $j(B_g(n,L))$ lies in
$j(\{d_g(n,\cdot)=L\})$. This argument uses no regularity of the
distance sphere.

A curve from $p$ that leaves $j(B_g(n,L))$ has a first exit.
The preceding boundary description and the inverse on $j(O)$ lift
the curve up to that exit, starting at $n$ and ending at distance $L$.
Length comparison gives length at least $(1-\lambda)L$ before
exit. Every point of the displayed source ball can be reached by a
curve shorter than this number, so its curve cannot exit. This proves
coverage. The same reasoning applies to every piecewise smooth
competitor used to compute distance; the lift exists near the exit
because the whole closed ball lies in $O$.

For the upper distance bound, map a minimizing $g$-segment from $n$
to $x$; it stays in the closed $r$-ball. For the lower bound, any
competitor that stays in the image of the $L$-ball lifts and has length
at least $(1-\lambda)d_g(n,x)$. Any competitor that exits already
costs at least $(1-\lambda)L$, a larger number. Taking the infimum
proves the bound. Completeness supplies compact balls and minimizing
segments through the inherited Riemannian foundation.
\end{proof}

The local source for these metric conventions is Lee's archived
\emph{Introduction to Riemannian Manifolds}, second-edition TeX,
\code{c02.tex}: \code{lemma:immersion-pullback-metric},
\code{prop:prop-length} and \code{prop:isometry-vs-isometry}.
Those three propositions leave proofs as exercises; LC39 writes the
needed length and first-exit argument. The source explicitly warns
that a local isometry need not preserve global distances. Its
July 25, 2026 errata distinguishes boundaryless domains in the
isometry discussion; all ambient manifolds here are boundaryless.

\subsection{A quantitative outward-sign test}

\begin{lemma}[Stability of an outward conormal; LC40]
\label{lem:collapse-conormal-stability}
On a finite-dimensional real inner-product space with metric $g$,
let $h$ be another positive-definite metric with
$(1-\delta)g\leq h\leq(1+\delta)g$, $0\leq\delta<1$.
For covectors $\beta,\xi$ with
$\|\beta\|_{g^*}\geq m>0$ and
$\|\xi-\beta\|_{g^*}\leq\sigma$, one has
\[
 \beta(h^{-1}\xi)
 \geq\frac{\|\beta\|_{g^*}^2}{1+\delta}
       -\frac{\sigma\|\beta\|_{g^*}}{1-\delta}>0
 \quad\hbox{if}\quad
 \sigma<\frac{1-\delta}{1+\delta}m.
\]
Here $h^{-1}$ raises a covector to a vector; every covector norm is
measured using the same $g^*$.
\end{lemma}
\begin{proof}
Write $h(v,w)=g(Av,w)$ with $A$ positive and self-adjoint. Its
eigenvalues lie in $[1-\delta,1+\delta]$, so those of $A^{-1}$
lie in $[(1+\delta)^{-1},(1-\delta)^{-1}]$.
Split $\xi=\beta+(\xi-\beta)$. The first summand gives the
displayed positive quadratic lower bound; Cauchy--Schwarz and the
operator norm of $A^{-1}$ bound the magnitude of the second by
$\sigma\|\beta\|_{g^*}/(1-\delta)$. Factor out
$\|\beta\|_{g^*}$ and use its lower bound $m$ for strict positivity.
This is pointwise linear algebra, independent of the existence of
radial functions or smooth convergence.
\end{proof}

Lee's \code{c02.tex}, \code{eq:raise-index} and
\code{eq:define-gradient}, gives the inverse-metric convention.
It also makes the following naturality identity immediate. If
$u=\eta\circ j$ and $h=j^*\widehat g$, then
\[
 dj(\nabla_hu)=\nabla_{\widehat g}\eta.
\]
Indeed pair both sides with $dj(v)$, use the pullback identity and
the chain rule, and use that $dj$ is onto in equal dimensions.
Thus a positive conormal pairing computed on $N$ tests the actual
source radial gradient, including its direction after normalization.

\begin{theorem}[Finite comparison data produce a transverse core; LC41]
\label{thm:collapse-quantitative-core-transfer}
Use LC39 with $L=10$ and $0<\lambda<1/10$. Suppose $M$ carries
an LC30 radial function with $0<e<1/40$ and $0<\varepsilon<1/4$.
Let $D_N\subset N$ be a compact smooth codimension-zero domain with
\[
 \overline B_g(n,1/2)\subset\operatorname{int}D_N,
 \qquad D_N\subset B_g(n,2).
\]
Suppose a smooth function $v$ on a neighborhood of $\partial D_N$
defines $D_N$ there by $v\leq0$, with $\|dv\|_{g^*}\geq m>0$
on its boundary. Set $u=\eta\circ j$. Assume on $\partial D_N$
\[
 (1-\delta)g\leq h\leq(1+\delta)g,\qquad
 \|du-dv\|_{g^*}\leq\sigma
 <\frac{1-\delta}{1+\delta}m,\qquad 0\leq\delta<1.
\]
Then $D=j(D_N)$ satisfies every hypothesis of LC36. In particular
it is smoothly ambient isotopic to all $A_\rho$, $\rho\in[1/5,2]$.
If $N$ is the same model as the chosen LC31 witness and $D_N$ has
the supplied soul/normal-disk-bundle data of LC37, these actual data
produce its noncompact alternative and hence the LC38 conclusion.
\end{theorem}
\begin{proof}
If $q\in A_a$, then $d_{\widehat g}(p,q)<a+e<3/20$.
LC39 gives a preimage $x$ in the $10$-ball. A curve from $p$ to $q$
of length less than $a+e$ cannot exit this image. Lifting it gives
\[
 d_g(n,x)<\frac{a+e}{1-\lambda}<\frac16<\frac12.
\]
Thus $A_a\subset\operatorname{int}D$. The upper distance bound
of LC39 on the $2$-ball also gives
\[
 \eta(j(x))<2(1+\lambda)+e<89/40<3\qquad(x\in D_N).
\]
Hence $D\subset\{\eta<b\}$. Both bounds refer to ambient source
distance, not just an intrinsic distance inside $j(O)$.

Every boundary point of $D_N$ has $g$-distance strictly greater than
$1/2$ and less than $2$. LC39 therefore puts its image strictly
between distances $9/20$ and $11/5$ from $p$, inside the smooth
region of $\eta$. Thus $u$ is smooth near this compact boundary.
Take $\beta=dv$, $\xi=du$ in LC40. The source defining function
$f=v\circ j^{-1}$ satisfies, by gradient naturality,
\[
 df(\nabla_{\widehat g}\eta)=dv(\nabla_hu)>0.
\]
The nonzero source gradient has positive squared norm, so dividing
by it preserves this outward sign for $X$. Compactness and the
smooth-domain structure transfer through the embedding. Apply LC36,
and, when the additional soul data are supplied, LC37--LC38.
\end{proof}

\subsection{What a convergent sequence must supply}

\begin{corollary}[Eventual transfer of one fixed model core; LC42]
\label{cor:collapse-eventual-core-transfer}
Fix the model $N$, core $D_N$, defining function $v$ and $m>0$
of LC41. Suppose $j_i:O_i\to M_i$ are actual pointed smooth
embeddings with $\overline B_g(n,10)\subset O_i$ and
$h_i=j_i^*\widehat g_i\to g$ uniformly on this closed ball.
Suppose each $M_i$ has the stated LC30 radial output, and
\[
 \sup_{\partial D_N}\|d(\eta_i\circ j_i)-dv\|_{g^*}\leq m/4
\]
eventually. Then $j_i(D_N)$ satisfies LC36 on one common tail.
If the model and soul data match the selected cone witnesses, it
gives LC37 packets and LC38 sublevel types on that tail.
\end{corollary}
\begin{proof}
Uniform metric convergence on the compact ball implies, eventually,
$(1-1/10)g\leq h_i\leq(1+1/10)g$.
These bounds imply the LC39 bounds with $\lambda=1/16<1/10$,
since $(15/16)^2<9/10$ and $11/10<(17/16)^2$.
For the boundary sign take $\delta=1/10$ and $\sigma=m/4$;
$1/4<9/11=(1-\delta)/(1+\delta)$. Intersect this metric tail
with the assumed differential-control tail and apply LC41.
Neither $N$ nor $D_N$ changes with $i$.
\end{proof}

The boundary differential bound is an \emph{additional hypothesis}.
It is not inferred from uniform convergence of values or metrics.
For example, on the real line take $v(x)=x$ and
$u_k(x)=x-2\sin(kx)/k$. Then $\|u_k-v\|_\infty\leq2/k$,
whereas $u_k'(0)=-1$ and $v'(0)=1$, with the metric unchanged.
This example isolates the insufficiency of value control; it is not
a counterexample satisfying the stronger LC30 Lipschitz-difference
and annular hypotheses. Those hypotheses may help prove the needed
bound, but no comparison of independently chosen smoothings has yet
been established. Full $C^1$ convergence to the defining function is
sufficient; LC42 only asks for the displayed eventual boundary error.

\begin{lemma}[Compact model transfer at a fixed enlarged scale; LC43]
\label{lem:collapse-compact-core-transfer}
Let $(N,g,n)$ be a nonempty compact connected smooth Riemannian
three-manifold without boundary. Fix $0<e<1/40$, and choose $R>0$
such that
\[
 \frac{2\operatorname{diam}_gN}{R}+e<1/5.
\]
Suppose $j_i:N\to M_i$ are pointed smooth embeddings into connected
smooth three-manifolds, and $j_i^*g_i\leq4g$ on $N$ eventually.
If $\eta_i$ is the selected radial function at scale $R$ with
$|\eta_i-d_{R^{-2}g_i}(p_i,\cdot)|<e$, then
$\eta_i<1/5$ on all of $M_i$ on that tail. When $N$ is the
same model as the cone witness, this gives the compact alternative
of LC37. Such an $R$ can also exceed any prescribed finite lower
bound, and is chosen before the tail.
\end{lemma}
\begin{proof}
Each embedding has open image by equal dimensionality and closed
image by compactness. Its nonempty image is all of connected $M_i$.
Map minimizing curves in $N$ and use $j_i^*g_i\leq4g$ to get
$d_{g_i}(p_i,j_i(x))\leq2\operatorname{diam}_gN$.
Divide distances by $R$ and use the value error to obtain the strict
global bound. The model is compact, so the diameter is finite; this
also permits increasing $R$ to meet other fixed lower bounds.
The lemma assumes actual global embeddings, not merely maps on a
small coordinate ball.
\end{proof}

\subsection{Source scope and the next geometric obligation}

KL Definition 3.4, printed 22/PDF 17, uses $C^{K+1}$ pointed
maps, a $C^K$ metric-difference bound and a GH approximation bound;
the manifolds themselves have $C^{K+1}$ transitions and $C^K$
metrics. Its text does not insert an embedding hypothesis into that
definition. Lemmas 3.5 and 6.10 retain smooth compactness as an
input. Passing to embeddings of whole neighborhoods in the model-to-source
direction and binding them to the smooth category used in LC37
therefore remain explicit tasks. A local nonsingular differential
does not alone give global injectivity on the comparison ball.

The source uses that convergence in the last step of 11.1, printed
59--60/PDF 54--55. LC39--LC43 now write the enclosure and
perturbation calculations that can be used once suitable maps and
one model core are supplied. Remark 11.4, printed 61/PDF 56,
motivates the eventual soul-core isotopy class; the proof of that
assertion and the construction of a radially controlled representative
have not been newly verified here. In particular, a tubular
neighborhood chosen without radial control need not satisfy LC41.

The next geometric target is to choose a sufficiently large fixed
scale in a noncompact limit and a smooth soul-type core with the
displayed enclosure and a defining function whose differential is
compatible with the source radial functions. An outward conormal
estimate weaker than full differential closeness would also suffice
if proved with a positive margin. Compare minimizing directions and
the source-qualified smoothing construction before choosing that
route; metric convergence alone is not the missing theorem.

All comparisons must use the same rescaled metric, point and cone
model. A uniform statement still requires negating the joint
conclusion, choosing one limit, then one sufficiently large fixed
$R$, and then a common later tail for cone approximation, curvature,
radial smoothing and core transfer. LC42 is a fixed-model transfer
result, not that full uniform producer. The compact alternative can
use the point cone after taking $R$ sufficiently large. KL11.5
selection remains subsequent Chapter 13 work; adapted coordinates
and actual fibration compatibility remain in Chapter 14.


\clearpage
\section{Inherited PC geometry and the remaining radial compatibility}
\label{sec:collapse-pc-reuse-binding}

The September 22, 2026 audit compared the blueprint with the private
\code{qinz1yang/differential-geometry-dev} repository, branch
\code{codex/pc-sorry-free}, at commit \code{54ad4d8ec040}.
The remote branch and the clean local checkout agreed. The earlier
September 15 snapshot was not used as the completed release. The audit
and its exact declarations, source hashes, commands and coverage are in
\code{pc_audit/revision87/} and \code{PC_REUSE_AUDIT_REVISION87.md}.
This is a component reuse audit, not acceptance of \code{PCReleaseGate}.

\subsection{What is already available}

The following are proof-bearing PC declarations, not just definitions
of desired outputs. Their inspected compiled axiom closures contain
only \code{propext}, \code{Classical.choice} and \code{Quot.sound}.
The selected module builds succeeded; ten source modules were also
freshly elaborated in an isolated copy of PC's existing environment.
No GC mathlib environment was pinned.

\begin{itemize}[leftmargin=*]
\item \code{exists_soul_diffeomorphism} supplies the compact soul and
an actual diffeomorphism from its normal-bundle total space to the
complete noncompact nonnegatively curved manifold, fixing the zero
section. Its bundle instances, compatible metric norm and boundaryless
hypotheses remain part of the interface. The weaker name
\code{cheegerGromollSoulTheorem} is a proposition definition;
it must not be substituted for this proof-bearing declaration.
\item \code{exists_soul_set_with_radial_escape_flow} supplies a smooth
complete outward flow relative to distance from the soul. The proof of
\code{exists_soul_normal_diffeomorph} uses the same soul, its normal
tube and that flow. The lower-level
\code{exists_normalFlow_diffeomorph} retains equality to the explicit
normal-flow map; the high-level theorem forgets that equality.
Retain the lower-level witness when radial control is needed.
\item \code{minimizing_directions_close_at_infinity} and
\code{exists_smooth_outward_field_at_infinity} concern distance from a
\emph{point}. The latter supplies a smooth field $V$ with
$\|V\|<2$ and $g(V,u)\leq-1/4$ against every inward unit minimizing
direction $u$ outside a sufficiently large ball. These are inherited
inputs, not new soul-theorem tasks.
\item \code{distance_levelProduct_at_infinity} and
\code{distance_shell_and_exterior_homeomorph} supply large radial
products and exteriors in the topological category. Their conclusion
is a homeomorphism, not the smooth disk-core identification of LC37.
\item \code{PointedRiemannianConvergenceMaps} stores actual pointed
smooth partial diffeomorphisms with open source sets exhausting every
compact set. \code{compactLimit_eventually_globalizes} makes them
pointed global diffeomorphisms on one tail when the limit is compact
and the approximating manifolds are connected. The audit consumer
\code{GCAudit.compact_model_maps} was elaborated with precisely these
inputs. Thus LC43's map hypothesis has a checked conditional PC binding.
\item \code{pointed_metric_eventually_quadratic_bounds} supplies the
uniform pullback inequalities on compact sets from convergence of the
\emph{canonical} source metrics. The inspected compactness producer
returns this canonical identification. Arbitrary reference metrics
cannot be silently substituted.
\end{itemize}

The audit also confirms the inherited Hamilton--Ivey and smooth
Schoenflies declarations. It does not require their reproving for GC.
Of the 111 old prepared declaration names, 110 resolve unchanged;
P63's operator exists under \code{DifferentialGeometry.PDE.RicciFlow},
not its prepared \code{Geometry.Operator} namespace. All 111 resolved
candidates have standard-only inspected axiom closures. That does not
by itself verify their proposed GC consumer types.

\subsection{Three interfaces that must still be matched}

First, the inspected producer
\code{exists_canonical_metric_compact_limit_of_local_curvature_injectivity}
requires, for every radius and \emph{every} derivative order, an eventual
curvature-derivative bound, together with injectivity lower bounds on
bounded balls. LC24 retains KL's fixed finite $K\geq10$ input.
The quantifiers
\[
 \forall R>0\;\forall m\in\mathbb N\;\exists C_{R,m}
 \quad\hbox{and}\quad
 \forall R>0\;\forall m\leq K\;\exists C_{R,m}
\]
are different. Reuse of the smooth maps is conditional on obtaining
this PC convergence package, or on binding a suitable finite-regularity
variant. Noncollapse must also supply the required injectivity bounds.
Neither step follows just by naming the compactness theorem. The
same issue affects MC20 and the LC24--LC31 application, rather than
invalidating the independent metric compactness arguments of Chapter 3.

Second, distance to a soul $S$ and distance to the chosen point $p$
are different functions. Separate existential applications may choose
different souls, fields or diffeomorphisms. A disk core extracted from
one retained normal-flow construction must use that construction's
own witnesses. Compactness of such a core gives some outer radius;
it does not give LC41's half-ball/two-ball enclosure ratio or an outward
normal estimate for the selected $\eta$.

Third, the final transfer needs one model, one core, one enlarged scale
and one common tail. Pointwise distance convergence and existence of
some minimizing arms do not give uniform control of \emph{all}
minimizing directions along a moving boundary. The source theorem
\code{eventually_exists_minimizingArms_comparisonAngle_lt} provides
selected arms for fixed points, so that distinction remains explicit.

\subsection{A weaker sufficient condition for radial outwardness}

LC40--LC42 use a strong bound comparing two smooth differentials.
The following criterion uses the actual radial directions instead.
It reuses LC29's differentiability and full-support argument, with
Rademacher and Riemannian calculus still inherited inputs.

\begin{lemma}[Lipschitz radial perturbations preserve a direction margin; LC44]
\label{lem:collapse-direction-margin}
Let $(M,g)$ be complete connected smooth Riemannian without boundary,
$p\in M$, and $W\subset M\setminus\{p\}$ open. Let $F$ be smooth
on $W$, let $F-d_p$ be globally $\varepsilon$-Lipschitz with
$\varepsilon\geq0$, and let $Z$ be a continuous tangent field on $W$.
Suppose $a,B>0$ and, at every $q\in W$,
\[
 \|Z_q\|\leq B,\qquad
 g_q(Z_q,u)\leq-a
\]
for every inward unit minimizing direction $u:q\to p$.
Then
\[
 dF_q(Z_q)\geq a-\varepsilon B \qquad(q\in W).
\]
In particular this derivative is positive if $\varepsilon B<a$.
If, on the smooth boundary of a compact domain $D\subset M$, the
field $Z$ equals its outward unit normal and $\partial D\subset W$, then the same strict
inequality makes $\nabla F$ point outward along that boundary.
The hypotheses must hold on an open neighborhood of the boundary,
not merely on its measure-zero set of points.
\end{lemma}
\begin{proof}
At almost every $q\in W$, both the distance and its Lipschitz
difference from $F$ are differentiable. A minimizing unit-speed
segment to $p$ gives $d(d_p)_q(u)=-1$. Since $d_p$ is
$1$-Lipschitz, Cauchy--Schwarz and equality in that inequality imply
$\nabla d_p(q)=-u$ and $\|\nabla d_p(q)\|=1$. Differentiating the
Lipschitz difference gives
$\|\nabla F(q)-\nabla d_p(q)\|\leq\varepsilon$. Consequently
\[
 dF_q(Z_q)
 =-g_q(u,Z_q)+g_q(\nabla F(q)-\nabla d_p(q),Z_q)
 \geq a-\varepsilon B.
\]
Every nonempty open Riemannian set has positive volume, so these
points are dense in $W$. Continuity of $dF(Z)$ extends the inequality
to all of $W$. If $Z=\nu$ on $\partial D$, gradient duality gives
$g(\nabla F,\nu)=dF(\nu)>0$, which is the required outward sign.
\end{proof}

For the PC point-distance field, $a=1/4$ and $B=2$ give positivity
whenever $\varepsilon<1/8$, on the open exterior where its direction
bound holds. This is a statement about $dF(V)$. To use LC36 one still
needs the sign against the actual outward normal of the chosen core.
Two positive pairings with a common field do not imply that sign:
in $\mathbb R^2$, with $\nu=(1,0)$, $V=(1,2)$ and $w=(-1,2)$,
\[
 \langle V,\nu\rangle=1>0,\qquad
 \langle w,V\rangle=3>0,\qquad
 \langle w,\nu\rangle=-1<0.
\]
Thus simply flowing a disk core by a PC outward field and proving
$d\eta(V)>0$ does not complete the gradient-transversality argument.
LC44 supplies a sufficient normal-direction test; it does not assert
that the required core and uniform margin have been constructed.

\subsection{Revised implementation order and release evidence}

First retain one PC soul/normal-flow witness and identify a compact
smooth normal disk core. Next establish its radial enclosure and an
outward-normal margin for point-distance directions, or prove directly
the stronger conormal criterion of LC40. Match that margin to source
minimizing directions on one buffered comparison domain and one tail;
LC44 then applies to the \emph{selected} LC30 smoothing without requiring
independently chosen smoothings to converge in $C^1$. Resolve the
finite-order compactness binding in parallel with this mathematical
comparison. Only then assemble LC37 and the joint KL 11.1 producer.
These are compatibility tasks using inherited PC geometry.

The repository-wide lexical scan covered 12,610 Lean source files and
found six \code{sorry} sites. The diagnostic aggregate imported every
available root module, excluding one nonexistent root import, and
inspected axiom closures for 225,569 compiled declarations, including
generated declarations; 24 depend on \code{sorryAx}. This is not a
fresh proof-by-proof kernel replay or a statement-correctness review
of all those declarations. The root build fails on the missing
\code{DoubledTameConsumer} import. The inspected Poincare assembly
also retains controlled-extinction and actual-event reconstruction
premises. Accordingly the completed-release gate remains pending even
though the individual geometric components above are usable audit
candidates. GC retains its late-thick/thin-or-terminal-exceptional
outcome and does not acquire a default finite-extinction dependency.


\clearpage
\section[Normal-flow cores and common fields]{Normal-flow cores and a common-field comparison}
\label{sec:collapse-common-field-core}

The normal-direction criterion in LC44 is sufficient for the original
core itself to satisfy LC36. It is not necessary for a smooth comparison
of that core with a radial sublevel. A smooth field that points outward
from the core and increases $\eta$ can supply its own product coordinates.
The resulting ambient isotopy then moves the core to a regular sublevel,
where the radial gradient is automatically outward. This gives the
unchanged LC37 packet using a \emph{modified} embedding.
The counterexample after LC44 remains valid: it excludes an inference
about the original core's gradient, not this different flow construction.

\subsection{Compact disk cores from one inherited normal-flow map}

\begin{lemma}[Disk cores and the normal-flow coordinate; LC45]
\label{lem:collapse-normal-flow-cores}
Let $S$ be a compact smooth manifold without boundary, let
$E\to S$ be a smooth real vector bundle of finite positive rank with
a smooth fiber inner product, and let $e:E\to N$ be a supplied
smooth diffeomorphism onto a boundaryless smooth manifold. Write
\[
 Q(q,v)=\|v\|^2,\qquad
 u(e(q,v))=\|v\|,\qquad D_T=e(\{Q\leq T^2\})\quad(T>0).
\]
Then $u$ is continuous and proper and is smooth off $e(S)$;
$u^2$ is smooth everywhere. Each $D_T$ is a compact smooth domain,
with boundary $\{u=T\}$, and
\[
 D_T\simeq D(E),\qquad D_T\subset\operatorname{int}D_{T'}
 \quad(0<T<T'),\qquad \bigcup_{T>0}\operatorname{int}D_T=N.
\]
The disk-bundle diffeomorphism is the restriction of
$(q,v)\mapsto e(q,Tv)$ from the unit disk bundle.
Every compact subset of $N$ lies in $\operatorname{int}D_T$ for
all sufficiently large $T$.

Suppose additionally that one retained smooth flow $\varphi$ with
smooth generator $V$ and one $\ell>0$ satisfy
\[
 e(q,rv)=\varphi_{r-\ell}(e(q,\ell v))
 \quad(\|v\|=1,\ r>\ell).
\]
Then $du(V)=1$ on $\{u>\ell\}$. In particular $V$ is strictly
outward transverse to every $\partial D_T$ with $T>\ell$.
The bundle, map and flow in these assertions are the \emph{same}
chosen witnesses.
\end{lemma}
\begin{proof}
A finite trivializing cover of the compact base has compact subsets
$K_j$, each contained in a trivialization, whose interiors cover $S$.
In each such chart the least eigenvalue of the fiber metric has a
positive lower bound on $K_j$. Thus $\{Q\leq T^2\}$ over $K_j$
is a closed subset of $K_j$ times a bounded closed Euclidean ball,
and is compact. The finitely many chart pieces prove compactness of
the entire disk bundle. This also proves properness of the fiber
norm, and hence of $u$, since $e$ is a homeomorphism.

Smoothness of $Q$ follows in bundle coordinates. At a nonzero vector,
the fiber curve $s\mapsto(q,(1+s)v)$ has
\[
 \left.\frac{d}{ds}\right|_{s=0}Q(q,(1+s)v)=2Q(q,v)>0.
\]
Thus every positive value of $Q$ is regular. The inherited regular-level
and local half-space arguments give the smooth boundary and domain;
fiber scaling and $e$ give the stated diffeomorphism. Continuity of $u$
and its maximum on any compact set give nesting and exhaustion.

For $r>\ell$ and $r+t>\ell$, the retained flow identity and its
composition law give
\[
 \varphi_t(e(q,rv))=e(q,(r+t)v),\qquad
 u(\varphi_t(e(q,rv)))=r+t.
\]
Differentiate at $t=0$. Near $\partial D_T$ the defining function
$u-T$ has derivative $1$ on $V$, giving the outward sign.
\end{proof}

For the PC soul, the inputs come from the \emph{same} construction in
\code{SoulDiffeomorph.lean}, retaining the map equality exported by
\code{exists_normalFlow_diffeomorph} in
\code{NormalFlowGluing.lean:346--425}. Its
\code{normalFlowMap} is the normal exponential below $\ell$ and
$\varphi_{r-\ell}$ applied to the tube sphere above $\ell$.
The exact source is the revision-87 pinned PC commit
\code{54ad4d8ec0408e71da2247daed5ecee2a574f7ec}; the high-level
existential soul theorem alone forgets the needed equality.
The two small consumer lemmas in
\code{pc_audit/revision88/RadialFlowConsumer.lean} check the ray formula
and radial translation using this actual definition. They do not
formalize all of LC45, compactness of disk bundles, or core transfer.

LC45 produces compact smooth disk cores conditional on the inherited
actual diffeomorphism, and the retained construction makes their field
transversality explicit. It does \emph{not} compare $u$ with distance
to the chosen basepoint. Compact exhaustion alone gives neither the
half-ball/two-ball enclosure in LC41 nor a uniform direction margin.
Distance to the soul, bundle radius and point-distance remain distinct.

\subsection{Product coordinates for an increasing smooth field}

Use LC32's notation
$a=1/8$, $b=3$, $K=\eta^{-1}([a,b])$ and
$\Sigma=\eta^{-1}(1)$ for the selected LC30 function.

\begin{lemma}[A common field supplies a smooth product band; LC46]
\label{lem:collapse-common-field-band}
Let $Y$ be smooth on an open neighborhood of $K$, with
$d\eta(Y)>0$ at every point of $K$. There is a diffeomorphism
\[
 \Psi_Y:\Sigma\times[a,b]\longrightarrow K,\qquad
 \eta(\Psi_Y(x,s))=s,
\]
whose coordinate curves are integral curves of $Y/d\eta(Y)$.
If $\Sigma$ is empty then $K$ is empty and the statement has the
empty interpretation. No equality of $Y$ with $\nabla\eta$ is required.
\end{lemma}
\begin{proof}
For nonempty $K$, positivity and compactness give a positive minimum
for $d\eta(Y)$ on $K$. Shrink to a relatively compact open neighborhood
where this pairing is positive and $\eta$ is smooth. Extend
$Y/d\eta(Y)$ by a smooth bump, equal to one near $K$, to a compactly
supported vector field on $M$. Its complete flow $\psi$ satisfies
$d\eta(\dot\psi)=1$ near $K$. A first-exit argument confines each
trajectory starting at level $s$ to $K$ for all times from $a-s$
to $b-s$. The formulas
\[
 \Psi_Y(x,s)=\psi_{s-1}(x),\qquad
 \Psi_Y^{-1}(q)=(\psi_{1-\eta(q)}(q),\eta(q))
\]
are inverse by the flow law and extend smoothly past both end levels.
They also show that a nonempty $K$ meets level $1$.
The empty case is immediate.
\end{proof}

\begin{lemma}[A common outward field recognizes the core; LC47]
\label{lem:collapse-common-field-isotopy}
Let $D\subset M$ be a compact smooth codimension-zero domain with
$A_a\subset\operatorname{int}D$ and $D\subset\{\eta<b\}$.
Suppose $Y$ satisfies LC46 and is strictly outward transverse to
$\partial D$: for each local defining function $f$ oriented by
$D=\{f\leq0\}$, require $df(Y)>0$ there. Then for each
$\rho\in[1/5,2]$ a smooth ambient isotopy $H_t$, $0\leq t\leq1$,
satisfies
\[
 H_0=\mathrm{id},\qquad H_1(D)=A_\rho.
\]
The isotopy has one compact support contained in $\eta^{-1}((a,b))$,
so it fixes a neighborhood of $A_a$. Its construction uses
$\Psi_Y$; it makes no assertion that the original $D$ is a graph
in LC33's gradient coordinates.
\end{lemma}
\begin{proof}
Every $\Psi_Y$ fiber begins inside $D$ and ends outside it. At a
boundary crossing its defining function has derivative
$df(Y)/d\eta(Y)>0$. Each crossing is isolated, and the closed set
of crossings on the compact fiber is finite. A second outward
crossing would require an intervening inward crossing; hence there
is exactly one. The implicit function theorem gives a unique smooth
height $h:\Sigma\to(a,b)$ with
\[
 D=A_a\cup\{\Psi_Y(x,s):a\leq s\leq h(x)\}.
\]
Compactness gives a uniform margin from $h(\Sigma)$ and $\rho$ to
both endpoints. Put $h_t=(1-t)h+t\rho$ and take a smooth scalar
bump $\chi$, equal to one near zero, supported in less than half
this margin. The time-dependent vertical field
\[
 (\rho-h(x))\chi(s-h_t(x))\,\partial_s
\]
has one common compact support strictly inside the product band.
Push it forward by $\Psi_Y$ and extend by zero. Its smooth flow
exists for $0\leq t\leq1$, with a smooth inverse, by inherited
ODE theory and the common compact support. The curve $s=h_t(x)$
solves its initial-value problem. Uniqueness makes the fiber maps
increasing, so they carry the whole subgraph below $h$ to the one
below $\rho$. This proves the asserted equality, including boundary
and interior. In the empty-band case the enclosure forces
$D=A_a=A_\rho$ and the identity suffices.
\end{proof}

The flow and bump inputs are inherited differential topology. Lee's
archived second-edition \code{c02.tex:1180--1208} supplies the bump
construction; \code{c09.tex}, labels \code{theorem:flows},
\code{theorem:cpctly-supp-complete} and \code{thm:time-dept-flow},
supply the flow statements and proofs. The recorded January 3, 2026
exit-time correction remains in force. KL 11.3(4), printed 60/PDF55,
also uses a common increasing field, while Remark 11.4, printed61/PDF56,
compares the model's smoothed balls with the soul tube. These passages
motivate this sufficient route; they do not provide its quantitative
source hypotheses. LC47 uses a smooth core and smooth field throughout,
so no distance sphere is promoted to a smooth hypersurface.

\subsection{A collar margin yields the original packet after isotopy}

\begin{theorem}[Collar-field transfer and packet repair; LC48]
\label{thm:collapse-collar-field-transfer}
Use LC39 with $L=10$ and $0<\lambda<1/10$, writing its pointed
embedding as $j_0:O\to M$. Let the complete source $M$ carry the
selected LC30 function, with $0<e<1/40$ and $0<\varepsilon<1/4$.
Let a compact smooth model core $D_N\subset N$ satisfy
\[
 \overline B(n,1/2)\subset\operatorname{int}D_N,
 \qquad D_N\subset B(n,2).
\]
Suppose a smooth model field $V$ is defined near $\partial D_N$ and
is strictly outward transverse there. On an OPEN source neighborhood
$W\subset M\setminus\{p\}$ of $\partial j_0(D_N)$ contained in the smooth region of $\eta$,
let $Z=(j_0)_*V$ be defined and satisfy, for fixed $\alpha,B>0$,
\[
 \|Z\|\leq B,\qquad \widehat g(Z,w)\leq-\alpha,
 \qquad \varepsilon B<\alpha,
\]
for EVERY inward UNIT minimizing direction $w$ to $p$, at every
point of $W$. Then $j_0(D_N)$ is smoothly ambient isotopic to every
$A_\rho$, $\rho\in[1/5,2]$, by isotopies supported in a compact
subset of $\eta^{-1}((a,b))$.

If this is the SAME noncompact model, scale, cone witness and $\eta$
as in LC37, and $D_N$ has its specified normal-disk-bundle
identification (for example from LC45), there is a modified pointed
embedding $j_1:O\to M$ producing the noncompact LC37 packet.
One may take $j_1(D_N)=A_1$. The original metric comparison belongs
to $j_0$; no small metric distortion or small tracks are asserted
for $j_1$ or the isotopy.
\end{theorem}
\begin{proof}
LC39's first-exit estimate, exactly as in LC41, gives
$A_a\subset\operatorname{int}j_0(D_N)$ and
$j_0(D_N)\subset\{\eta<b\}$: the lifted inner radius is less than
$1/6$, and $\eta<89/40<3$ on the image core. The image boundary
has source distances strictly between $9/20$ and $11/5$, in the
smooth region. Pushforward preserves outward transversality, since
for a model defining function $v$,
\[
 d(v\circ j_0^{-1})(Z)=dv(V)>0.
\]
LC44 on $W$ gives $d\eta(Z)\geq\alpha-\varepsilon B>0$.
This uses its directional inequality, not its optional normal clause.

Shrink $W$ within $\eta^{-1}((a,b))$ around the compact boundary.
Choose a smooth bump $\theta$ equal to one near that boundary,
with compact support in $W$ and $0\leq\theta\leq1$. Such a bump
is obtained by first choosing a smaller neighborhood with compact
closure in $W$. On a neighborhood of $K$ define
\[
 Y=\theta Z+(1-\theta)\nabla\eta,
\]
where $\theta Z$ extends by zero. The extension is smooth because
its support is contained in the domain of $Z$. On $K$,
\[
 d\eta(Y)=\theta\,d\eta(Z)+(1-\theta)\|\nabla\eta\|^2>0;
\]
outside the bump support use the second term alone. Along the core
boundary $Y=Z$, so $Y$ is strictly outward. LC47 gives the isotopy.

Take its $\rho=1$ instance and set $j_1=H_1\circ j_0$.
It is a smooth embedding on the SAME open $O$ and maps $D_N$ to
$A_1$. The isotopy fixes $A_a$, which contains $p$ because
$|\eta(p)-d(p,p)|<e<a$, so $j_1(n)=p$. The domain $A_1$
contains $A_a$ in its interior and lies in $\{\eta<b\}$. On its
boundary $f=\eta-1$ has $df(X)=1$ for LC36's normalized radial
gradient $X$. Thus all original LC37 clauses hold for $j_1$, without
weakening that packet or asserting them for $j_0$. LC38 then gives
the stated soul type and LC34 the distance-ball homeomorphism.
\end{proof}

The repair changes the embedding, not the chosen source radial function,
scale, cone map or nonnegative model. Retain $j_0$ separately whenever
metric estimates are needed. No derivative convergence of independently
chosen radial smoothings is assumed, and the margin is needed only
on an open collar, not on the whole radial band.

Still needed are one fixed scale and an LC45 core with the required
enclosure, plus a uniform source collar margin for its OWN field against
ALL point-minimizing directions, on one common tail with the cone and
LC30 witnesses. PC supplies strict soul-direction pairings; its separate
point-distance field requires a compatibility proof before substitution.
The finite-$K$ versus all-order compactness and smooth-category bindings
also remain open. LC48 does not prove the full joint good-annulus
producer. Ball selection remains in Chapter13 and compatible fibrations
in Chapter14.

\clearpage
\section[Uniform collar margins]{Uniform collar margins under metric convergence}
\label{sec:collapse-uniform-collar}

The source-collar hypothesis of LC48 can be reduced to a strict
point-distance condition on the boundary of ONE fixed model core.
The reduction has two parts: compactness gives a model neighborhood
and a positive margin, and buffered metric convergence transfers them
to all source minimizing directions on one common tail. Neither part
constructs the required enclosed model core or identifies its own
outward field with PC's separately chosen point-distance field.

All manifolds below are finite-dimensional, connected, boundaryless
and smooth, with complete smooth Riemannian metrics. For $q\ne n$ put
\[
 \mathcal U_n(q)=\{v\in T_qN:\ |v|_g=1,
              \ \exp_q(d_g(n,q)v)=n\}.
\]
These are ALL inward unit minimizing directions. Completeness gives
nonemptiness; the equality uses the distance as travel time, so the
unit geodesic is minimizing. No choice of one preferred geodesic is
part of the definition.

\subsection{A strict boundary condition gives a compact collar}

\begin{lemma}[Uniform model collar from strict point-directions; LC49]
\label{lem:collapse-model-collar-margin}
Let $C_0$ be a nonempty compact subset of $B_g(n,3)\setminus\{n\}$,
and let $V$ be smooth on an open neighborhood of $C_0$. Suppose
\[
 g_q(V_q,v)<0\quad(q\in C_0,\ v\in\mathcal U_n(q)).
\]
Then there are $\alpha,B>0$ and an open neighborhood $U$ of $C_0$
whose closure $C$ is compact, lies in $B_g(n,3)\setminus\{n\}$ and
in the domain of $V$, and satisfies
\[
 |V_q|_g\leq B,\qquad g_q(V_q,v)\leq-2\alpha
       \quad(q\in C,\ v\in\mathcal U_n(q)).
\]
All constants and the collar are chosen before a source sequence index.
\end{lemma}
\begin{proof}
The unit tangent bundle over a compact base set is compact. In that
bundle the endpoint condition defining $\mathcal U_n$ is closed,
by continuity of distance and the exponential map. Over $C_0$ the
continuous function $-g(V,v)$ therefore has a positive minimum $m$.
Choose $\alpha=m/4$. First place $C_0$ in a relatively compact
neighborhood inside the specified open domain, away from $n$ and
from the radius-$3$ boundary.

If the inequality $g(V,v)<-3\alpha$ did not hold in a neighborhood
of $C_0$, there would be points tending to $C_0$ and bad unit
minimizing directions above them. Compactness of the unit tangent
bundle over the preliminary compact neighborhood gives a convergent
subsequence. Its limit satisfies the endpoint condition and has
$g(V,v)\geq-3\alpha$, contradicting $g(V,v)\leq-4\alpha$ on
$C_0$. Shrink once more so that the closure remains inside the good
neighborhood. Continuity bounds $|V|$ there; increase the bound to a
positive $B$. This gives the stated weaker margin $2\alpha$.
\end{proof}

PC already proves the fixed-metric openness mechanism in
\code{SoulAngles.lean:343--386},
\code{eventually_minimizing_inner_lt}. Its proof uses the proper
unit-tangent projection and the closed bad-direction set.
\code{NoncriticalDistance.lean:31--63} supplies the corresponding
fixed-fiber definition, nonemptiness and compactness for distance to
a compact set; specialize to $\{n\}$. A locally defined $V$ can
be extended near $C_0$ by a bump before using the global-field PC API.
LC49 is the compact-neighborhood packaging of this inherited argument.
It does not replace point-directions by directions to the soul.

\subsection{Every source minimizing direction has a model limit}

\begin{lemma}[Buffered compactness of minimizing directions; LC50]
\label{lem:collapse-minimizing-direction-limit}
Let $O\subset N$ be open and contain $\overline B_g(n,10)$.
For each $i$, let $j_i:O\to M_i$ be a smooth embedding between
manifolds of the SAME dimension, with $j_i(n)=p_i$. Suppose the
pullback metrics $h_i=j_i^*g_i$ converge to $g$ in $C^1$ on compact
subsets of $O$. Fix a compact
$C\subset B_g(n,3)\setminus\{n\}$.
For ANY sequence $q_i\in C$ and ANY inward $g_i$-unit minimizing
directions $w_i$ from $j_i(q_i)$ to $p_i$, a subsequence satisfies
\[
 q_i\longrightarrow q\in C,\qquad
 (d j_i|_{q_i})^{-1}w_i\longrightarrow v\in\mathcal U_n(q)
\]
in the tangent bundle of $N$. This is a subsequential assertion for
arbitrary choices, not uniqueness or continuity of a selected direction.
\end{lemma}
\begin{proof}
Compactness of the closed $10$-ball and positive definiteness give
$\lambda_i\to0$, eventually $0<\lambda_i<1/10$, with
\[
 (1-\lambda_i)^2g\leq h_i\leq(1+\lambda_i)^2g
                     \quad\hbox{on }\overline B_g(n,10).
\]
LC39 gives source-ball coverage and the radial distance bounds.
Write $r_i=d_g(n,q_i)<3$ and
$l_i=d_{g_i}(p_i,j_i(q_i))$. Every source minimizing geodesic
from $j_i(q_i)$ to $p_i$ has length $l_i\leq(1+\lambda_i)r_i<33/10$.
Every point on it has source distance at most $l_i$ from $p_i$,
strictly below $10(1-\lambda_i)>9$. Thus the ENTIRE geodesic lies
in $j_i(B_g(n,10))$ and has a unique lift. The inverse length
bound, applied from its endpoint $n$, gives
\[
 d_g(n,x)\leq\frac{l_i}{1-\lambda_i}
       \leq\frac{1+\lambda_i}{1-\lambda_i}r_i
       <\frac{11}{3}<4
\]
at every lifted point. This is a global-minimizer argument; local
isometry alone would not provide this confinement.

Reparametrize each lifted geodesic as $\gamma_i:[0,1]\to N$,
from $q_i$ to $n$, with constant $h_i$-speed $l_i$.
Its $g$-speed is bounded by $11/3$, and its image lies in the
compact closed $4$-ball. Arzela--Ascoli gives a uniformly convergent
subsequence and $q_i\to q\in C$. Cover the limiting path by
finitely many coordinate neighborhoods, with compact buffers, and
partition $[0,1]$ into finitely many overlapping chart intervals.
Uniform convergence puts the tails in those same buffers.

On each buffer, $C^1$ metric convergence implies uniform convergence
of the inverse metrics and of
\[
 \Gamma(h_i)^k_{ab}
 =\tfrac12(h_i)^{kc}
   (\partial_a(h_i)_{bc}+\partial_b(h_i)_{ac}
                       -\partial_c(h_i)_{ab}).
\]
Coordinate velocities are bounded by the metric speed bound and
uniform chart equivalence. The geodesic equation
$\ddot\gamma_i^k=-\Gamma(h_i)^k_{ab}\dot\gamma_i^a\dot\gamma_i^b$
therefore bounds accelerations uniformly. A finite further
Arzela--Ascoli extraction gives $C^1$ convergence on all chart
intervals, including one-sided derivatives at $0$ and $1$.
Passing to the integral form of the equation proves that the limit
is a $g$-geodesic. The overlapping limits agree by the coordinate
change rule; smoothness then follows from the smooth limiting equation.
No bound on derivatives of the Christoffel symbols is used.

The two-sided LC39 estimates give $l_i\to l=d_g(n,q)>0$.
The limit has constant $g$-speed $l$, joins $q$ to $n$, and hence
has length equal to their distance. Since
\[
 \dot\gamma_i(0)=l_i(d j_i|_{q_i})^{-1}w_i,
\]
the limiting initial unit vector belongs to $\mathcal U_n(q)$.
This proves the assertion for arbitrary directions, including
directions based at cut-locus points.
\end{proof}

Lee's archived IRM second-edition \code{c04.tex:1495--1614}
gives the geodesic equation and its first-order form, with the
existence/uniqueness proof; \code{c05.tex:600--815} gives the
metric-to-Christoffel calculation. The July 25, 2026 author errata
retain the boundary and endpoint qualifications. Here the ambient
manifolds are boundaryless and endpoint derivatives are obtained
from the closed-interval integral equation. Arzela--Ascoli and
local calculus are inherited foundation inputs. The argument above
is the written moving-metric adapter, not a theorem quoted verbatim
from Lee and not a new compactness theorem for producing the $j_i$.
PC's \code{tendsto_pointed_map_distance} and
\code{eventually_exists_minimizingArms_comparisonAngle_lt} remain
reuse candidates; their fixed-point and existential-arm signatures
alone do not assert this uniform ALL-direction conclusion.

\subsection{The common tail and the remaining model problem}

\begin{theorem}[Uniform source collar and conditional packet; LC51]
\label{thm:collapse-uniform-collar-packet}
Use the maps and $C^1$ convergence of LC50, with $N$ noncompact.
Let a fixed compact
smooth model core $D\subset N$ contain $\overline B_g(n,1/2)$
in its interior and lie in $B_g(n,2)$. Let its OWN smooth field $V$
be defined near $\partial D$, strictly outward transverse there,
and satisfy
\[
 g(V,v)<0\quad(q\in\partial D,\ v\in\mathcal U_n(q)).
\]
Apply LC49 to $C_0=\partial D$, fixing $U,C,\alpha,B$.
Then for all sufficiently large $i$, on the OPEN collar
$W_i=j_i(U)$ the field $Z_i=(j_i)_*V$ satisfies
\[
 |Z_i|_{g_i}\leq2B,\qquad
 g_i(Z_i,w)\leq-\alpha
 \quad\hbox{for EVERY inward UNIT direction to }p_i.
\]
Suppose also that the selected LC30 functions $\eta_i$ occur on
this SAME scale and tail, with errors $e_i<1/40$ and Lipschitz
errors bounded by a fixed
\[
 0<\varepsilon<\min\{1/4,\alpha/(2B)\}.
\]
Then LC48 applies on a common tail. If $D$ has the supplied
normal-disk-bundle identification and the same model/cone data
required by LC37, the modified embeddings give the original
noncompact LC37 packets, with image core $A_{i,1}$.
\end{theorem}
\begin{proof}
The boundary is nonempty: otherwise the nonempty compact domain
would be both open and closed in the connected noncompact model.
It is compact and lies strictly between model radii $1/2$ and $2$,
so LC49 applies. Uniform $C^0$ metric convergence gives
$|V|_{h_i}\leq2B$ on $C$ eventually. If the directional bound
failed arbitrarily far out, select an increasing sequence of bad
indices, points $q_i\in C$ and source minimizing directions with
$g_i(Z_i,w_i)>-\alpha$. LC50 gives a subsequential model direction.
Continuity of $V$ and metric convergence imply
\[
 g_i(Z_i,w_i)
 =h_i\bigl(V,(d j_i)^{-1}w_i\bigr)
 \longrightarrow g(V,v)\leq-2\alpha,
\]
contradicting the lower bound $-\alpha$ on the limit. Thus the
margin holds uniformly on $C$, and hence on $U$, with one tail
for ALL points and directions.

The core enclosure and LC39 put its image boundary in the smooth
region of every $\eta_i$. Intersect $j_i(U)$ with that region
and, if necessary, the open band used in LC48; this remains an
open neighborhood of the boundary and preserves the estimates.
Pushforward preserves outward transversality. Intersect the tails
for these estimates, $\lambda_i<1/10$, the given smoothing bounds
and the same-model/cone witnesses. Since $2B\varepsilon<\alpha$,
LC48 gives the required isotopies and modified pointed embeddings.
The original $j_i$ retain the metric estimates; those estimates
are not asserted for the modified embeddings.
\end{proof}

The quantifier order matters: first fix the model, scale, core and
its own field, then the collar and constants, then one tail.
The final smoothing inequality is an explicit hypothesis. This
argument does not justify choosing a globally prescribed smoothing
tolerance after seeing the limiting core. Multiplying $V$ by a
positive constant scales both $\alpha$ and $B$, so it cannot
improve the ratio $\alpha/B$.

There is also a concrete limitation to fixed-time cores. In the
Euclidean plane over a one-point base, the diffeomorphism
$e(x,y)=(x,y/8)$ gives $u(X,Y)=\sqrt{X^2+64Y^2}$ and disk cores
with semiaxes $T$ and $T/8$. On $u\geq\ell>0$ let
$V=(X,Y)/u$, extending smoothly by a cutoff that vanishes near
the origin. Then $du(V)=1$, $|V|\leq1$, and its pairing with
the inward point-direction is at most $-1/8$ on this exterior.
The LC45 flow identity holds there. Yet no scale $R>0$ has
\[
 \overline B(0,R/2)\subset\operatorname{int}D_T
                \subset D_T\subset B(0,2R),
\]
because the first inclusion requires $R<T/4$ and the last
requires $R>T/2$. This tests the abstract LC45 data; the map
is not the normal exponential on a Euclidean tube, so it is not
a counterexample to additional properties of PC's specific map.
It shows why exhaustion and a bounded outward field alone do not
establish the required radius ratio.

The next model task is therefore a core adapted to point-distance,
potentially a variable-height section of the inherited normal flow,
with its disk-bundle identification and a compatible outward field.
LC49--LC51 supply the subsequent uniform transfer once those inputs
and a compatible smoothing budget are obtained. Fixed finite-$K$
compactness still must produce the actual maps and the required
smooth category; using only $C^1$ convergence in LC50 does not
solve that producer problem. Full LC37/KL11.1 production remains
open. This work stays in Chapter 13.

\clearpage
\section[Point-distance model cores]{Point-distance cores from a compatible normal flow}
\label{sec:collapse-point-distance-core}

We now construct the noncompact smooth model core needed in the previous
section. The construction retains the inherited soul and its normal tube,
but chooses the field \emph{before} choosing the global normal-flow map.
A smoothed point-distance level will be a variable-height section of
those flow coordinates. This supplies the radius ratio that constant
fiber-radius cores did not provide. Cone existence, the LC28 smoothing
input and the inherited smooth foundation retain their recorded
qualifications; no new Lean implementation is asserted.

Throughout this section, $(N,g,p)$ is a complete connected noncompact
finite-dimensional smooth Riemannian manifold without boundary, with
$\sec_g\geq0$, and with a supplied LC21 cone-at-infinity package. All
metric and tangent-bundle conventions are those of the inherited PC
foundation. Write $\mathcal U_p(q)$ for ALL inward UNIT minimizing
directions to $p$, and $\mathcal U_S(q)$ for ALL inward UNIT directions
realizing $d(q,S)$ for a nonempty compact set $S$. These are two families:
minimizing distance to $S\cup\{p\}$ would in general discard one of them.

\subsection{A common outward direction far from a compact set}

\begin{lemma}[One direction for both targets; LC52]
\label{lem:collapse-common-outward-direction}
For each nonempty compact $S\subset N$, there is $A>0$ such that,
for every $q$ with $d(p,q)\geq A$, there is a unit vector $w\in T_qN$
with
\[
 g_q(w,v)<-1/2
 \quad\text{for EVERY }v\in\mathcal U_p(q)\cup\mathcal U_S(q).
\]
The same $w$ works for both families. No continuous choice of $w$ is
asserted. The radius $A$ may depend on the model and $S$; the margin does not.
\end{lemma}
\begin{proof}
Choose $D>0$ with $d(p,s)\leq D$ for all $s\in S$, and put
$\mu=1/100$. Fix $\delta>0$ with
$\delta<\min\{1/24,\mu/15\}$. LC21 supplies an actual cone
approximation at every sufficiently large scale $r$. Enlarge the
threshold $A$ so that $D/r\leq\mu$ whenever $r\geq A$.
For $r=d(p,q)\geq A$, apply LC25 to $r^{-1}N$, at normalized
radius $1$ (its parameters $a=b=1$). Rescaling back gives $q'$ with
\[
 d(p,q')=2r,\qquad r\leq b:=d(q,q')<(1+\mu)r.
\]
Fix one minimizing segment $q\to q'$ and let $w$ be its unit initial
velocity. For any $s\in S\cup\{p\}$ and any minimizing segment
$q\to s$, put $a=d(q,s)$ and $c=d(s,q')$. The triangle inequality gives
\[
 (1-\mu)r\leq a\leq(1+\mu)r,\qquad c\geq(2-\mu)r.
\]
In particular both initial directions are defined. For the Euclidean
comparison angle $\bar\theta$ at $q$, the cosine rule gives
\[
 1+\cos\bar\theta
 =\frac{(a+b)^2-c^2}{2ab}
 \leq\frac{12\mu+3\mu^2}{2(1-\mu)}
 =\frac{401}{6600}<\frac12.
\]
The complete nonnegative-curvature angle comparison inherited in AC02
implies $\theta\geq\bar\theta$ for the \emph{chosen} two minimizing
segments, so $g(w,v)=\cos\theta<-1/2$. If the triangle is degenerate,
the same comparison has its limiting meaning; in the equality
$c=a+b$ case the concatenation is minimizing and the angle is $\pi$.
The other collinear equality cases cannot violate the displayed bound.
This works for every $s$ and every minimizing segment to it, hence for
the stated two direction families. The far threshold precedes $q$;
no rate uniform over all models has been used.
\end{proof}

The new estimate is a compact-target extension of the point argument,
using LC25's repaired \emph{upper} bound. It is not inferred from the
insufficient lower bound printed in KL 11.3. The complete Riemannian
comparison binding is inherited; the LC21 Tits-cone input is still explicit.

\subsection{Smooth patching with both direction constraints}

\begin{lemma}[A bounded field for both distances; LC53]
\label{lem:collapse-common-outward-field}
There are $0<A_0<A_1$ and a global smooth field $X$ such that
\[
 |X|_g<2,\qquad X=0\text{ on }\overline B(p,A_0),
\]
\[
 g(X,v)\leq-1/4\quad
 \bigl(d(p,q)\geq A_1,\quad
 v\in\mathcal U_p(q)\cup\mathcal U_S(q)\bigr).
\]
These inequalities test EVERY direction in each family, including
points in either cut locus. There is no smooth choice of minimizing arms.
\end{lemma}
\begin{proof}
Take $A_0$ larger than the LC52 threshold and a ball containing $S$,
and put $A_1=A_0+1$. At a point with $d(p,q)\geq A_1$, extend the
LC52 unit vector to a smooth local field $W$. The inherited fixed-metric
openness arguments for point-distance and distance to a closed set give,
after intersecting neighborhoods, $g(W,v)<-1/4$ for BOTH families.
Shrink again so that $|W|<2$ and $d(p,\cdot)>A_0$ there. The openness
uses compact unit tangent fibers and closed endpoint conditions, not
uniqueness of geodesics. Below radius $A_1$, the zero field on the open
ball is a valid local choice.

For each fiber impose simultaneously $|z|<2$, the condition $z=0$
if $d(p,q)\leq A_0$, and all the linear inequalities
$g(z,v)\leq-1/4$ from BOTH families if $d(p,q)\geq A_1$.
This is a convex set. The preceding local fields take values in these
sets on neighborhoods covering $N$. A subordinate locally finite smooth
partition of unity gives a global smooth convex combination satisfying
every constraint. Local finiteness makes the sum smooth and finite at
each point; convexity preserves even the strict norm bound. This is the
inherited convex-section patching argument, with two families of
half-spaces instead of one.
\end{proof}

PC's \code{SoulAngles.lean:343--468} supplies point-direction openness
and the single-target field proof. The parallel closed-set statement
\code{eventually_infDist_minimizing_inner_lt}, in
\code{DistanceField.lean:67--112}, supplies the second openness input.
The actual patcher is
\code{exists_contMDiffSection_forall_mem_convex_of_local}, in the PC
checkout's pinned mathlib \code{Geometry/Manifold/PartitionOfUnity.lean:591--631}.
These are source-inspected inherited ingredients, not a claim that the
new two-target theorem was already elaborated. That mathlib belongs to
the PC audit environment; production GC remains unpinned.

\subsection{Retain the soul tube and choose a compatible flow}

\begin{lemma}[Point-outward normal-flow data; LC54]
\label{lem:collapse-compatible-normal-flow}
Use one inherited PC soul $S$, its positive-rank normal bundle $E=\nu S$,
and the bounded radial-field and normal-tube data of
\textup{\code{exists_soul_set_with_radial_outward_field}}.
There is a choice of global smooth field $V$, complete smooth flow
$\varphi$, $\ell>0$ and actual smooth diffeomorphism $e:E\to N$ with
\[
 |V|_g\leq2,\qquad e(q,0)=q,
 \qquad e(q,tv)=\varphi_{t-\ell}(e(q,\ell v))
 \quad(|v|_g=1,\ t>\ell).
\]
For some $A_2>0$, independent of $q$,
\[
 g(V,w)\leq-1/4
 \quad(d(p,q)\geq A_2,\ w\in\mathcal U_p(q)).
\]
The chosen $e$ is the actual normal-flow map of this SAME $V$ and
$\varphi$, with the retained soul and tube. It need not equal any
previously chosen existential soul diffeomorphism.
\end{lemma}
\begin{proof}
Choose the inherited tube radius $\epsilon>0$ and radial-data radius
$r_0>0$ for this same $S$, and set $r=\min\{r_0,\epsilon\}>0$.
The lower-level bounded-field theorem supplies $W_0$ with $|W_0|\leq2$,
agreement with $\nabla d_S$ on a neighborhood of
$\{r/4\leq d_S\leq r/2\}$, and strictly negative pairing with every
member of $\mathcal U_S(q)$ when $d_S(q)\geq r/8$.
Obtain $X$ from LC53 for this $S$.

Choose a smooth function $0\leq\chi\leq1$, zero on a neighborhood of
both that compact tube annulus and $\overline B(p,A_1)$, and equal to
one outside a larger compact set. It is one minus a compactly supported
smooth bump equal to one near the first compact set; it is not a
composition with the possibly nonsmooth function $d_p$. Set
\[
 V=(1-\chi)W_0+\chi X.
\]
The norm bound survives convex combination, and the old gradient
agreement is unchanged near the tube annulus. Where $\chi>0$, the
LC53 soul-direction bound is valid. Thus $V$ remains strictly outward
to $S$ on $\{d_S\geq r/8\}$, including points where $\chi=1$.
Outside a sufficiently large ball it equals $X$, giving the point
margin with one $A_2$.

Apply the inherited bounded-field completeness and escape-flow theorem
to this $V$, with bound $2$ and inner radius $r/8$.
Set $\ell=3r/8$ and $\delta=r/16$. Then
$[\ell-\delta,\ell+\delta]\subset(r/4,r/2)$ and
$\ell+\delta<\epsilon$. Tube calibration and the retained equality
$V=\nabla d_S$ show that normal geodesics and this flow coincide on
the overlap, by integral-curve uniqueness. The actual inherited
\code{exists_normalFlow_diffeomorph} therefore applies. Keep its
map equality, rather than projecting to the weaker headline theorem.
The definition of \code{normalFlowMap} gives the stated ray identity.
LC45 can now use this same $e,V,\varphi$ to obtain $u,D_T$ and
$du(V)=1$ on $\{u>\ell\}$.
\end{proof}

This is a new written adapter of existing PC machinery. The source
proof in \code{SoulDiffeomorph.lean:31--136} makes the same choices of
$r,\ell,\delta$; the field modification above is inserted before its
flow/gluing step. The bounded input is retained from
\code{SoulRadialField.lean:31--67}, since its escape-flow headline
forgets that bound. The reused consumers are
\code{DistanceEscape.lean:471--496},
\code{RadialNormalFlow.lean:105--126} and
\code{NormalFlowGluing.lean:346--425}. No PC source was edited.

\subsection{Cut the flow coordinates at a smoothed distance level}

\begin{theorem}[An enclosed model disk core with a fixed margin; LC55]
\label{thm:collapse-point-distance-model-core}
Retain LC54's witnesses and the LC28 smoothing input used by LC30.
There is $R_*>0$ such that for EVERY real $R\geq R_*$ the scaled
model $g_R=R^{-2}g$ admits a selected LC30 function $\zeta_R$, with
\[
 \|\zeta_R-d_{g_R}(p,\cdot)\|_\infty<1/80,
 \qquad \operatorname{Lip}_{g_R}(\zeta_R-d_{g_R})\leq1/64,
\]
whose sublevel $D=\{\zeta_R\leq1\}$ is a compact smooth domain and
\[
 \overline B_{g_R}(p,1/2)\subset\operatorname{int}D,
 \qquad D\subset B_{g_R}(p,2),\qquad D\simeq D(\nu S).
\]
The disk-bundle identification is an actual smooth diffeomorphism of
manifolds with boundary. On the OPEN collar
$U_R=\{3/4<d_{g_R}(p,\cdot)<5/4\}$, which contains $\partial D$,
the SAME field $V_R=RV$ satisfies
\[
 |V_R|_{g_R}\leq2,\qquad
 g_R(V_R,w)\leq-1/4\quad(w\in\mathcal U^{g_R}_p(q)),
 \qquad d\zeta_R(V_R)\geq7/32>0.
\]
In particular it is strictly outward transverse to this core. The
numerical errors and margins precede the model; only $R_*$ and the
chosen core depend on it.
\end{theorem}
\begin{proof}
Fix $T_0>\ell$ and the compact LC45 core $D_{T_0}$. Choose $R_*$
large enough that for every $R\geq R_*$,
\[
 R>2A_2,\qquad D_{T_0}\subset B_g(p,R/2),
\]
and the LC21 approximation at scale $R$ is sufficiently accurate for
LC30 with Lipschitz error $1/64$. Nonnegative sectional curvature
satisfies LC30's lower bound at every scale. Choose its uniform
function error $e=1/80$. The same fixed cone is used for every $R$.
LC32 supplies compactness and smoothness of the sublevel domain.
The error bound gives the asserted strict enclosure and
\[
 \partial D\subset\{79/80<d_{g_R}(p,\cdot)<81/80\}
 \subset U_R.
\]
On $U_R$, $d_g(p,\cdot)>3R/4>A_2$. A $g_R$-unit direction is
$R$ times its $g$-unit counterpart, so both the norm and pairing
bounds in the statement follow from LC54. The function is smooth
on a neighborhood of this collar. LC44 gives
\[
 d\zeta_R(V_R)\geq\frac14-\frac2{64}=\frac7{32}.
\]

It remains to identify the core; enclosure alone would not do this.
We have $D_{T_0}\subset\operatorname{int}D$. By compactness of $D$
and properness of $u$, choose $T_1>T_0$ with
$D\subset\operatorname{int}D_{T_1}$. Let $\Sigma_E$ be the compact
unit sphere bundle of $E$ in the original metric. On
$\Sigma_E\times[T_0,T_1]$, the actual coordinates
\[
 \Theta((q,v),t)=e(q,tv)
\]
identify $V$ with $\partial_t$. Every fiber starts strictly inside
$D$ and ends strictly outside. Its intersections with $\partial D$
are zeros of $\zeta_R\circ\Theta-1$. This function is continuous
on the whole interval and smooth near every zero; at a zero its
derivative is at least $7/(32R)>0$. The zero set is compact and
has no accumulation point, hence is finite. All crossings are from
inside to outside, so there is exactly one. This does not assert
monotonicity along the entire fiber away from the level.

The implicit function theorem now gives a smooth height
$h:\Sigma_E\to(T_0,T_1)$, and
\[
 D=D_{T_0}\cup
 \{\Theta(z,t):z\in\Sigma_E,\ T_0\leq t\leq h(z)\}.
\]
Choose $T_*\in(T_0,\min h)$. The same compact-support vertical-flow
construction as in the proof of LC47, now in these normal-flow
coordinates, moves the constant height $T_*$ to $h$: put
$h_s=(1-s)T_*+sh$, and use
\[
 (h(z)-T_*)\beta(t-h_s(z))\,\partial_t\qquad(0\leq s\leq1),
\]
where $\beta=1$ near zero and has support smaller than half the
uniform distance of all $h_s(z)$ to $T_0,T_1$. The field has one
common compact support inside the open band and extends by zero
to $N$. Its time-one diffeomorphism $H_1$ fixes $D_{T_0}$ and maps
$D_{T_*}$ onto $D$, including their interiors and boundaries.
Inherited compact-support ODE theory gives the isotopy and its inverse.
Thus $H_1$ composed with LC45's fiber-scaling disk map is the desired
diffeomorphism. For the unit disk bundle measured in $g_R$, the formula
is explicitly
\[
 (q,v)\longmapsto H_1(e(q,(T_*/R)v));
\]
normal subspaces are unchanged by constant metric scaling, while their
unit disk radii are not. This accounts for that convention as well.
\end{proof}

\paragraph{Consequences and next binding.}
The conditional smooth noncompact \emph{model-core} construction is
now written. It uses a variable height, so it does not contradict the
constant-height ellipsoidal obstruction in the previous section.
The proof follows the soul-tube comparison suggested by KL Remark 11.4,
printed 61/PDF56, while explicitly retaining the common flow and smooth
boundary category. No assertion that an unsmoothed distance sphere is
smooth, and no interior-homeomorphism-to-diffeomorphism upgrade, is used.

The collar has compact closure away from $p$ inside the radius3 ball,
and its model bounds allow the explicit LC51 choices
$\alpha=1/8$, $B=2$. Thus that proof's source bounds are $1/8$ and $4$.
A source smoothing tolerance fixed \emph{before} selecting the model,
with $0<\varepsilon<1/64$, satisfies $4\varepsilon<1/8$.
This removes the model-dependent smoothing-budget obstruction for
these constructed cores. It does not produce the LC50 embeddings:
actual buffered $C^1$ maps in the required smooth category and their
finite-$K$ compactness producer remain to be bound. Metric estimates
stay on the original embeddings when LC48 modifies them.

Next combine this model construction with the compact-model branch and
one joint contradiction sequence at a fixed enlarged scale, recording
exactly what the finite-regularity map producer must supply. Keep the
same model, cone, scale, radial-function witnesses and one common tail.
The LC21/LC28 source-qualified inputs and PC release gate remain open;
full LC37/KL11.1 production and a production GC Lean build are not claimed.
Ball selection remains subsequent Chapter13 work; compatible fibrations
remain in Chapter14.

\clearpage
\section[Joint model assembly]{Joint cone, radial-function and model-core assembly}
\label{sec:collapse-joint-model-assembly}

LC43 and LC55 supply the two model branches for one joint contradiction.
The sufficient input below retains actual maps to the same smooth
limit. Producing it from KL's finite-$K$ hypotheses remains open.

\subsection{The convergence data consumed by core transfer}

\begin{definition}[Buffered smooth-model convergence data; LC56]
\label{def:collapse-buffered-model-data}
Let $(M_i,g_i,p_i)$ be complete connected smooth Riemannian
three-manifolds without boundary. A buffered smooth-model package
for this sequence consists of the following data.
\begin{enumerate}[leftmargin=*]
\item A complete connected smooth Riemannian three-manifold
$(N,g,n)$ without boundary, with $\sec_g\geq0$, and pointed GH
convergence $(M_i,p_i)\to(N,n)$ for these very metrics and points,
in the MC06--MC11 conventions.
\item Actual pointed smooth embeddings $j_i:U_i\to M_i$, with
$U_i\subset N$ open and $n\in U_i$. EVERY compact subset of $N$
is contained in $U_i$ for ALL sufficiently large $i$. In fixed
model charts, the actual pullbacks satisfy
\[
 j_i^*g_i\longrightarrow g\quad\text{in }C^1
 \text{ on every compact subset of }N.
\]
This is convergence to the SAME model metric, not to auxiliary
reference metrics. No monotonicity of $U_i$ is required.
\item A supplied LC21 cone package $(C,o)$ for $(N,n)$.
When $N$ is compact, choose the point cone at this stage; the
constant maps give its LC21 package since $\operatorname{diam}(N)/R\to0$.
When $N$ is noncompact, keep the existing Tits-cone qualification.
\item Positive numbers $H_i\to\infty$ and the SEPARATE curvature
bounds
\[
 \sec_{g_i}\geq-H_i^{-2}\quad\text{on }B_{g_i}(p_i,H_i).
\]
These bounds are supplied independently of the $C^1$ convergence.
\end{enumerate}
The smooth soul, normal-tube and flow foundation used for a noncompact
$N$ remains inherited. The package is a sufficient specification, not
an existence theorem or an accepted PC release binding.
\end{definition}

The exhaustion quantifier is essential: a union of domains equal to
$N$ would not guarantee that one late comparison map contains a
prescribed compact core. For every fixed $R>0$, properness and item(2)
give a common restriction
\[
 O_R=B_g(n,11R)\subset U_i
\]
on a tail, by first putting the compact closed $11R$-ball in $U_i$.
After the FIXED rescaling $g_R=R^{-2}g$, $g_{i,R}=R^{-2}g_i$,
these restricted embeddings meet LC50's open-neighborhood requirement
for the CLOSED radius10 ball, with the same $C^1$ convergence.
If $N$ is compact, apply exhaustion to $N$ itself: eventually
$U_i=N$, and the embeddings are global diffeomorphisms onto connected
$M_i$ by equal-dimensional openness and compactness.

PC already stores stronger map objects: the pointed smooth partial
diffeomorphisms in \code{PointedRiemannianConvergenceMaps}, with
\code{source_exhausts} and \code{source_subset}. Its
\code{compactLimit_eventually_globalizes} proves the compact conclusion.
The revision87 compact-map consumer remains checked evidence. The
package's coordinate $C^1$ statement must still be matched to PC's
canonical tensor-derivative convergence conventions; the checked
quadratic-bounds theorem supplies its $C^0$ metric inequalities.
Neither that theorem nor item(2) alone asserts item(4) or produces a
smooth nonnegative limit from finite-order curvature hypotheses.

\subsection{One fixed scale handles either model branch}

\begin{lemma}[Eventual joint witnesses from one model; LC57]
\label{lem:collapse-fixed-joint-witnesses}
Assume LC56 data and retain LC28, the comparison/analytic bindings of
LC30 and the inherited smooth foundation. Fix, BEFORE choosing the
model scale,
\[
 0<\delta<1,\qquad 0<\varepsilon<1/64,\qquad
 0<e<\min\{\varepsilon,1/40\},\qquad T>1.
\]
There are ONE fixed $R\geq T$ and $i_0$ such that for every
$i\geq i_0$, at the metric scale $R^{-2}g_i$, there are jointly:
\begin{enumerate}[leftmargin=*]
\item an actual pointed KL $\delta$-approximation to the SAME
supplied $(C,o)$ of $N$;
\item a selected LC30 function $\eta_i$ with value error $e$ and
Lipschitz-difference error $\varepsilon$, together with LC31's smooth
compactly supported cutoff and its stated bounds;
\item an LC37 packet using this SAME model $N$, scale, cone and
selected $\eta_i$.
\end{enumerate}
Consequently LC38 gives the smooth type of every radial sublevel
$A_{i,\rho}$, $\rho\in[1/5,2]$, on that common tail. This is the
compact type $N$ or the noncompact model's closed normal disk bundle.
The original metric-comparison embeddings are retained separately
from any modified packet embeddings.
\end{lemma}
\begin{proof}
Let $\delta_0(\varepsilon)>0$ be the LC30 threshold, and choose
\[
 0<\tau<\min\{\delta/4,\delta_0(\varepsilon),1\}.
\]
Use LC22's tolerance for target error $\tau$. By LC21, the model
has an actual approximation with this tolerance at EVERY sufficiently
large real scale $R$. If $N$ is noncompact, increase the lower
threshold for $R$ to satisfy LC55 as well. If $N$ is compact,
instead require
\[
 2\operatorname{diam}_g(N)/R+e<1/5.
\]
Choose one such $R\geq T$, once and for all. LC22 transfers the
model cone approximation to actual source KL $\tau$-maps on a tail.
The same maps are KL $\delta$-maps. For a target of radius less
than $\delta^{-1}-\delta$, choose a lift with target error less
than $2\tau$. Its source radius is less than
$\delta^{-1}-\delta+3\tau<\delta^{-1}$, so it survives
the restriction. Distortion and coverage errors are below $\delta$.

Next use LC56 item(4). Eventually $H_i>400R$. At scale $R$ the
curvature bound is $-(R/H_i)^2$ on the ball of radius $H_i/R$;
therefore the radius400 source ball has sectional curvature at least
$-1/60^2$. Apply LC30 to the SAME $\tau$-maps and choose $\eta_i$.
LC31's cutoff construction and bounds use these selected functions;
its independent scale-selection theorem is not invoked to choose a
second scale. Here ``LC31 data'' means those output clauses at the
jointly chosen scale.

If $N$ is compact, exhaustion supplies global pointed embeddings on
a tail. Uniform $C^0$ metric convergence gives $j_i^*g_i\leq4g$.
LC43 gives $\eta_i<1/5$ on all of $M_i$, and hence precisely the
compact LC37 alternative for these same witnesses.

If $N$ is noncompact, LC55 supplies at this fixed $R$ a compact
smooth core $D\subset(N,g_R)$, its normal-disk-bundle identification,
and its OWN outward field $V_R$ with norm at most2 and point-direction
pairing at most$-1/4$ on the open collar. Its compact closure lies
away from $n$ and inside the radius3 ball. The common restriction
$O_R$ above supplies the actual LC50 maps. The LC51 proof applies
with the explicit constants
\[
 \alpha=1/8,\qquad B=2,\qquad
 2B\varepsilon=4\varepsilon<1/8.
\]
One tail works for ALL collar points and ALL minimizing directions.
LC48 then modifies the embedding by an ambient isotopy, sending the
image of $D$ to $A_{i,1}$ and giving the ORIGINAL LC37 gradient
transversality requirement. It does not change $N,C,R$ or $\eta_i$.
The metric bounds still belong to the original $j_i$, with no
small-distortion assertion for the modified embedding.

Intersect the finitely many tails after the model and $R$ are fixed.
In each branch all three conclusions now hold together. LC38 gives
the stated sublevel types. The selected $R$ is independent of $i$;
no convergence estimate on balls of an arbitrarily growing radius
$R_i$ has been used.
\end{proof}

\subsection{Uniform bounded scales by negating the joint conclusion}

\begin{theorem}[Conditional uniform joint scale interval; LC58]
\label{thm:collapse-uniform-joint-scale}
Let $M^\alpha$ be complete connected smooth Riemannian
three-manifolds without boundary, indexed by positive integers, and
let $\rho_\alpha(p)>0$ be prescribed scales. Suppose that for EVERY
sequence $\alpha_i\to\infty$ and EVERY choice $p_i\in M^{\alpha_i}$,
the normalized pointed sequence
\[
 (M^{\alpha_i},\rho_{\alpha_i}(p_i)^{-2}g^{\alpha_i},p_i)
\]
has a subsequence carrying ONE LC56 package. Retain the LC57
smoothing, comparison and inherited foundation inputs.
Fix $\delta,\varepsilon,e,T$ in LC57's ranges, before this extraction.
Then there are $V\geq T$ and $\alpha_0$ such that, for every
$\alpha>\alpha_0$ and every $p\in M^\alpha$, some
\[
 s\in[T,V],\qquad r_p^0=s\rho_\alpha(p)
\]
gives ALL THREE LC57 conclusions at scale $(r_p^0)^{-2}g^\alpha$,
with one model, its cone, one selected radial function and one
matching core packet. All LC38 sublevel conclusions and LC31 cutoff
bounds hold for these same witnesses.
\end{theorem}
\begin{proof}
If the conclusion fails for the fixed parameters, for each integer
$i\geq1$ choose $\alpha_i>i$ and $p_i$ such that NO scale
$s\in[T,iT]$ and NO allowed collection of witnesses satisfies the
joint conclusion. Extract a subsequence with indices $i_j\to\infty$
and one LC56 model package, by hypothesis. LC57 gives one fixed
$R\geq T$ and a common tail of that subsequence with all three
conclusions. Eventually $R\leq i_jT$. At that index, $s=R$ and
$r_{p_{i_j}}^0=R\rho_{\alpha_{i_j}}(p_{i_j})$ are permitted choices
contradicting the joint failure. This proves existence of $V$ and
$\alpha_0$, without producing a numerical rate or a continuous
choice of scales or models.
\end{proof}

The quantifier order is: fix errors and standing data; extract a
model if uniformity fails; choose one $R$; then choose one source
tail. It is not enough to obtain separate cone, smoothing and topology
scales, even if each lies in a bounded interval. The constants may
depend on the fixed noncollapse and curvature-derivative data. There
is no uniform assertion as the noncollapse parameter tends to zero.

\subsection{The finite-regularity producer still required by the application}

For the KL application at the modified volume scale, the curvature
component of LC56 is already accounted for by the scale calculation
in LC31. At fixed $w'>0$, eventually $1/\alpha<w'$, and
\[
 R_p\geq(\alpha/2)\rho_\alpha(p).
\]
Thus on a selected normalized sequence one may take
$H_i=\alpha_i/2$, including the infinite-curvature-scale case.
This argument uses the standing volume/curvature hypotheses; it does
not derive a second-derivative bound from $C^1$ convergence.

KL Definition3.4 and Lemma3.5, printed22/PDF17, use $C^{K+1}$
transition/comparison maps and $C^K$ metrics. The compactness lemma
assumes $C^{K+2}$-smooth Riemannian inputs in that convention,
noncollapse and bounds on $\nabla^k\mathrm{Rm}$ for $0\leq k\leq K$.
Its harmonic-coordinate paragraph cites Petersen Chapter10, whose
proof has not been newly audited here. Lemma6.10, printed42--43/PDF37--38,
supplies fixed-$w'$ noncollapse and these FINITELY many derivative
orders, then invokes pointed $C^K$ compactness. Definition3.4 itself
does not assert embeddings of whole model-ball neighborhoods.

The inspected PC producer
\code{exists_canonical_metric_compact_limit_of_local_curvature_injectivity}
requires a local bound for EVERY derivative order and an eventual
injectivity lower bound at EVERY point of EACH bounded ball. It
returns actual smooth maps and canonical metric data. Its quantifiers
are stronger than the presently supplied finite-$K$ hypotheses; smooth
individual inputs do not give uniform bounds at missing orders.

\begin{center}
\begin{tabular}{p{.27\textwidth}p{.65\textwidth}}
\textbf{Binding} & \textbf{What remains to be supplied}\\
Actual maps & Model-to-source embeddings on neighborhoods of whole
closed balls, preserving the marked points; local nonsingularity alone
is insufficient.\\
Metric conventions & Coordinate $C^1$ convergence to the SAME limit
metric from the canonical PC tensor data, with explicit compact
restrictions.\\
Finite regularity & A finite-order route with justified differentiable
category, or independently proved stronger bounds. Merely changing
$K$ to infinity is not a proof.\\
Injectivity & The bounded-ball lower bounds needed by the selected
producer, derived from the available noncollapse and curvature data,
or a producer whose hypotheses already match those data.\\
Smooth model & The limit metric and the soul/flow/core machinery in a
compatible category. Changing the atlas alone does not make a finite-
regularity metric smooth or preserve its curvature under arbitrary
smoothing.\\
\end{tabular}
\end{center}

These are inherited-foundation bindings and adapters. No second soul
or compactness foundation is proposed, and no production environment
is pinned. The original PC release gate remains pending at the
recorded audit commit. The cone and distance-smoothing inputs remain
source-qualified. LC58 is a conditional joint producer; it is not
acceptance of its sequential LC56 premise under KL's finite-$K$ data.

KL11.1's proof, printed59--60/PDF54--55, initially negates its cone
clause before using topology. LC58 negates its entire output conjunction.
Its smooth diffeomorphisms concern SMOOTHED sublevels; LC34 gives
distance-ball homeomorphisms. The original open-ball diffeomorphism
still needs a smooth category/flow argument. Next bind the finite-
regularity model data and that comparison. KL11.5 ball selection remains
later Chapter13 work; compatible fibrations remain in Chapter14.

\clearpage
\section[Open distance balls]{Smooth interiors of distance sublevels}
\label{sec:collapse-open-ball-smooth-type}

The boundary height in LC34 is continuous, with no assumed
smoothness. Nevertheless the open subgraph inherits a smooth structure
from the product. A smooth coordinate tending to infinity at its
upper end compares its interior with a constant-height subgraph.
This resolves the open-ball comparison left in Section~
\ref{sec:collapse-joint-model-assembly}; the finite-$K$ production
of the smooth model package remains open.

\subsection{Straightening an open subgraph without smoothing its boundary}

\begin{lemma}[Smooth interior straightening of a continuous graph; LC59]
\label{lem:collapse-open-graph-straightening}
Let $\Sigma$ be a compact smooth manifold without boundary, let
$a<c<b$, and let $h:\Sigma\to(c,b)$ be continuous. Choose
$r\in(c,b)$ and assume
\[
 c<\min\{r,\min_{x\in\Sigma}h(x)\}
\]
when $\Sigma$ is nonempty. With the smooth structures inherited from
$\Sigma\times\mathbb R$, there is a fiber-preserving diffeomorphism
\[
 P:\Sigma\times(a,r)\longrightarrow
 \Omega_h:=\{(x,u):a<u<h(x)\}
\]
that is the identity for $u\leq c$. It is strictly increasing on each
fiber. No smoothness of $h$, boundary extension or ambient smooth
isotopy is asserted. If $\Sigma$ is empty, use the empty map.
\end{lemma}
\begin{proof}
Assume $\Sigma$ is nonempty. First construct a smooth positive
majorant $m$ of the continuous function
$q(x,u)=(h(x)-u)^{-1}$ on the OPEN manifold $\Omega_h$.
For each point choose an open neighborhood on which $q$ is bounded
above by a positive constant $C_j$. An inherited subordinate smooth
locally finite partition of unity $\{\psi_j\}$ gives
\[
 m=\sum_j C_j\psi_j>q>0.
\]
Only locally finitely many summands occur; every nonzero summand
uses a constant exceeding $q$ at that point. This proves both
smoothness and the strict bound without differentiating $h$.

Choose $\beta>0$ with
$c+\beta<\min\{r,\min h\}$ and a smooth scalar cutoff $\chi$
with $0\leq\chi\leq1$, equal to zero for $u\leq c$ and to one
for $u\geq c+\beta$. Define on $\Omega_h$
\[
 w_h(x,u)=1+\chi(u)m(x,u),\qquad
 Q_h(x,u)=c+\int_c^u w_h(x,s)\,ds.
\]
The integral is defined because every interval between $c$ and $u$
lies below $h(x)$. It is smooth in $(x,u)$: near any fixed point,
continuity of $h$ puts all these compact integration segments inside
one open product neighborhood in $\Omega_h$. Equivalently rewrite
the integral as
\[
 (u-c)\int_0^1 w_h(x,c+t(u-c))\,dt
\]
and differentiate on that fixed compact interval. Continuous
derivatives are uniformly bounded on a smaller compact coordinate
neighborhood; the usual difference-quotient argument and iteration
justify every order. Also $\partial_u Q_h=w_h>0$, and $Q_h=u$
for $u\leq c$. For each FIXED $x$, when $u>c+\beta$,
\[
 Q_h(x,u)\geq Q_h(x,c+\beta)
       +\log\frac{h(x)-c-\beta}{h(x)-u}
       \longrightarrow+\infty\quad\text{as }u\uparrow h(x).
\]
Hence $\widehat Q_h(x,u)=(x,Q_h(x,u))$ is a bijection from
$\Omega_h$ onto $\Sigma\times(a,\infty)$. Its derivative in
product charts is triangular with positive last diagonal entry.
The inherited inverse function theorem gives local smooth inverses;
bijectivity patches them to one smooth inverse. Construct
$\widehat Q_r$ in exactly the same way for the constant height $r$,
using the same $c$. Then
\[
 P=\widehat Q_h^{-1}\circ\widehat Q_r
\]
has all the required properties. In particular both coordinate
maps are the identity below $c$, and no regularity of a hitting-time
function has been used.
\end{proof}

The partition-of-unity foundation is already available in PC's pinned
mathlib dependency: \code{exists_contMDiffMap_forall_mem_convex_of_local_const}
applied to the intervals $(q(x,u),\infty)$ gives precisely the smooth
majorant. The archived Lee second-edition sources supply the explanatory
comparison: \code{c02.tex}, \code{thm:pou}, and \code{c04.tex},
\code{thm:ift-manifolds}. Both supplied proofs were read for this
application. The integral theorem \code{appC.tex},
\code{thm:diff-under-int}, states the needed finite-order rule but
refers elsewhere for its proof; the local compact-interval argument
above states the required justification. This is a consumer of
inherited differential topology and calculus, not a new foundation
or a newly checked Lean adapter.

\subsection{A smooth comparison of the actual open sets}

\begin{theorem}[Smooth open-distance-ball comparison with small displacement; LC60]
\label{thm:collapse-open-ball-diffeomorphism}
Retain the LC32--LC34 setting, with $0<\varepsilon<1/4$ and
$0<e<1/40$. For every $\rho\in[1/5,2]$ there is a diffeomorphism
in the inherited ambient smooth structures
\[
 J_\rho:\operatorname{int}A_\rho\longrightarrow B_g(p,\rho),
 \qquad
 J_\rho(q)=q\quad\text{if }\eta(q)\leq\rho-2e,
\]
with
\[
 d_g(J_\rho(q),q)<\frac{3e}{1-\varepsilon}.
\]
In particular it fixes $p$. This is a diffeomorphism between OPEN
sets, with no claim of smoothness on the distance sphere, smooth
extension to their closures or smooth dependence on $\rho$.
\end{theorem}
\begin{proof}
Use the SAME product $F$ and height $h_\rho$ constructed in LC33--LC34.
Put $c=\rho-2e$. Since $a=1/8$, $b=3$, and
$|h_\rho-\rho|<e$, we have
\[
 a<c<\rho-e<h_\rho(x)<\rho+e<b.
\]
LC59 applies to the constant height $\rho$ and continuous height
$h_\rho$. Conjugate its map by $F$ on the open product band and
extend by the identity below the band. The extension is smooth with
smooth inverse because it is already the identity on the WHOLE
overlap $a<u<c$, not merely at the level $a$.
Strict increase of $d\circ F$ identifies the target exactly with
$B_g(p,\rho)$ in the band. Outside the band, the inequalities
$a+e<\rho<b-e$ give the same classification by $d$ and $\eta$.
Thus the extension is a bijection of the indicated open sets.

If a point moves, its original coordinate satisfies $c<u<\rho$,
and its image coordinate satisfies $c<v<h_\rho(x)<\rho+e$.
Consequently $|v-u|<3e$; the speed bound of LC33 bounds the length
of the fiber segment by $3e/(1-\varepsilon)$. If it does not move,
the estimate is immediate. Finally $\eta(p)<e<c$, so $p$ is fixed.
When $\Sigma$ is empty, LC33 gives $K=\varnothing$ and the
off-band classification makes the identity the required map.
\end{proof}

LC34 remains the separate ambient HOMEOMORPHISM with the sharper
track bound $e/(1-\varepsilon)$. LC60 constructs a different smooth
map of the interiors; it does not promote LC34's piecewise affine
map to a diffeomorphism. No boundary-recognition or Schoenflies
argument is needed for this interior comparison.

\subsection{The model type at the already selected joint scale}

\begin{corollary}[Conditional smooth model type of open balls; LC61]
\label{cor:collapse-joint-open-ball-type}
In LC57 or LC58, retain all its hypotheses and the SAME selected
model, scale, cone and radial function. In the noncompact branch
retain an inherited actual diffeomorphism
$\Theta:\nu S\to N$ for the SAME soul and normal bundle used by
the LC37 packet; the LC54--LC55 construction retains such a witness.
Then every open distance ball of normalized radius
$\rho\in[1/5,2]$ is diffeomorphic to that same $N$.
In LC58 this holds at the SAME scale $r_p^0=s\rho_\alpha(p)$,
with the SAME uniform interval $s\in[T,V]$. No new scale, tail or
compactness extraction is needed.
\end{corollary}
\begin{proof}
LC60 identifies the open ball with $\operatorname{int}A_\rho$.
In the compact branch LC38 identifies this with $N$. In the
noncompact branch LC38 restricts to a diffeomorphism with the
interior of the closed unit normal disk bundle. For its smooth
fiber norm, the mutually inverse bundle maps
\[
 v\longmapsto\frac{v}{\sqrt{1-\|v\|^2}},\qquad
 z\longmapsto\frac{z}{\sqrt{1+\|z\|^2}}
\]
identify that open disk bundle with the entire $\nu S$. They are
smooth also at the zero section, since the squared norm is smooth
and the denominators are strictly positive on their respective
domains. Compose with the retained $\Theta$. All inputs concern
the already selected LC57--LC58 witnesses. The construction is made
for each $\rho$ on every index in the existing tail; it introduces
no limiting estimate and hence no intersection of uncountably many
new tails.
\end{proof}

KL Lemma11.1(2), printed59--60/PDF54--55, requests the smooth
type of an OPEN ball. Lemma11.3(4) and Remark11.4,
printed60--61/PDF55--56, distinguish its interior comparison from
the compact-domain argument. LC59--LC61 now supply the open-ball
smooth comparison directly, conditional on the existing smooth
package and inherited normal-bundle map. The finite-$K$ hypotheses
have still not been bound to LC56; the cone and smoothing inputs
and the PC release gate retain their qualifications. There is no
new production Lean proof or full KL11.1 acceptance.

Next resolve the finite-order actual-map and model-regularity
binding, or develop the ensuing ball-selection argument with that
producer retained as an explicit hypothesis. Ball selection remains
Chapter13; compatible fibrations remain Chapter14.

\clearpage
\section[Selecting zero-stratum balls]{Selecting zero-stratum balls from fixed pointwise data}
\label{sec:collapse-zero-ball-selection}

This section separates the metric selection in KL Lemma 11.5 from
its geometric producers. Its supplied proof and that of Lemma 19.1
were checked in the archived Ast\'erisque text, printed 61--62 and
93--94/PDF 56--57 and 88--89. The radius assignment need not be
continuous, and the set of candidate centers need not be closed.
The smooth-model package, annular splitting and end-count conclusions
are not consequences of a covering lemma. LC62--LC64 below are written
metric arguments, using an inherited Vitali theorem; no new Lean
elaboration or production environment pin is claimed.

\subsection{Maximal balls and the neighboring scale bound}

\begin{lemma}[LC62: maximal doubling-containing balls]
\label{lem:collapse-maximal-candidate-balls}
Let $X$ be a metric space, $Z\subset X$, and let $r:X\to(0,\infty)$
satisfy $r(p)\leq R_{\max}<\infty$ for every $p$. Write $B_p=B(p,r(p))$
for the OPEN ball and put
\[
 P_0=\{p\in X:B_p\cap Z\ne\varnothing\},\qquad
 p\prec q\quad\Longleftrightarrow\quad
       2r(p)<r(q)\ \text{ and }\ B_p\subset B_q.
\]
Let $V\subset P_0$ be the set of maximal elements for this strict
order restricted to $P_0$. Every $p\in P_0$ has $v\in V$ with
$p=v$ or $p\prec v$. Moreover, for $v\in V$ and $q$ with $d(v,q)\leq10r(v)$,
\[
 r(q)\leq20r(v).
\]
If $T>0$ and $\rho:X\to(0,\infty)$ also satisfy
$T\rho(q)\leq r(q)$ for every $q$, then
$r(v)/\rho(q)\geq T/20$ on that same CLOSED ball.
\end{lemma}
\begin{proof}
The relation is irreflexive and transitive. Fix $p\in P_0$ and
choose an integer $n\geq1$ with $2^n r(p)\geq R_{\max}$.
A chain of $n$ strict steps starting at $p$ would end with a radius
strictly larger than $2^n r(p)$, contradicting the upper bound.
If there were no maximal element above $p$ (allowing equality),
every element reached could be extended by another strict step.
Finite induction would produce that forbidden chain. Thus such a
$v$ exists. This argument uses neither continuity nor attainment of
a supremum, nor a positive lower bound uniform over all of $P_0$.

For the local estimate, suppose instead $r(q)>20r(v)$.
For $x\in B_v$, even when $d(v,q)=10r(v)$, the triangle
inequality gives $d(q,x)<11r(v)<r(q)$, so $B_v\subset B_q$. The point of $Z$ in
$B_v$ puts $q\in P_0$, and $2r(v)<r(q)$ gives $v\prec q$, a
contradiction. Finally
$T\rho(q)\leq r(q)\leq20r(v)$, and division by the positive
$20\rho(q)$ gives the asserted ratio. If $P_0$ is empty, all
statements about its elements are vacuous.
\end{proof}

\subsection{The inherited selection theorem and finiteness}

The inspected PC audit dependency already proves the abstract
Vitali intersection-and-radius selection theorem in
\code{Mathlib/MeasureTheory/Covering/Vitali.lean},
\code{Vitali.exists_disjoint_subfamily_covering_enlargement}.
For a family of nonempty sets $B_p$, nonnegative sizes $r(p)$ bounded
above, and any $\tau>1$, it selects a pairwise disjoint subfamily
such that every original set intersects a selected $B_i$ with
$r(p)\leq\tau r(i)$. Its full supplied proof was read at the recorded
mathlib commit. The open-ball enlargement wrapper alone hides the
intersection witness needed here, so the intended consumer uses the
abstract theorem with $\tau=2$. This is a reuse target, not a new
Vitali foundation or an already elaborated GC adapter.

\begin{lemma}[LC63: finite disjoint selection with radius control]
\label{lem:collapse-finite-disjoint-radius-selection}
Let $P$ be a totally bounded subset of a metric space $X$. Assign
radii satisfying $0<r_{\min}\leq r(p)\leq R_{\max}<\infty$ for
$p\in P$. There is a FINITE subset $I\subset P$ such that the
OPEN balls $B_i=B(i,r(i))$, $i\in I$, are pairwise disjoint and
\[
 \forall p\in P\quad\exists i\in I:\quad
 B_p\cap B_i\ne\varnothing,\qquad r(p)\leq2r(i).
\]
For each such pair, $d(p,i)<3r(i)$ and $B_p\subset B(i,5r(i))$.
No closedness of $P$ or continuity of its radius assignment is assumed.
\end{lemma}
\begin{proof}
Apply the inherited abstract Vitali theorem with $\tau=2$.
Every $B_p$ is nonempty since it contains $p$; positivity and the
upper bound verify its size hypotheses. This supplies a possibly
infinite disjoint family with the stated intersection and radius
properties.

Distinct selected centers have distance at least $r_{\min}$:
if $d(i,j)<r_{\min}\leq r(i)$, then $j\in B_i\cap B_j$.
Total boundedness supplies a finite cover of $P$ by open balls of
radius $r_{\min}/3$. Two selected centers in one covering ball
would have distance less than $2r_{\min}/3$, which is impossible.
Each covering ball therefore contains at most one selected center,
and the selected family is finite. This is the finite-cover/packing
argument of MC02; it does not claim a cardinality bound uniform in
other spaces or in a sequence with shrinking $r_{\min}$.

Choose $y\in B_p\cap B_i$. Then
\[
 d(p,i)\leq d(p,y)+d(y,i)<r(p)+r(i)\leq3r(i).
\]
For $x\in B_p$, the same inequality gives
$d(x,i)<2r(p)+r(i)\leq5r(i)$.
If $P=\varnothing$, take $I=\varnothing$.
\end{proof}

\subsection{A conditional small-core cover with unchanged witnesses}

\begin{theorem}[LC64: zero-set cover from annular exclusion]
\label{thm:collapse-zero-small-core-cover}
Let $X$ be a compact metric space, $Z\subset X$ an arbitrary subset,
$\rho:X\to(0,\infty)$ continuous, and $0<T\leq U<\infty$.
Fix ONCE a radius $r(p)$ for every $p\in X$ with
\[
 T\rho(p)\leq r(p)\leq U\rho(p).
\]
Use the $P_0$ and $V$ of LC62. There is a finite $I\subset V$ with
pairwise disjoint OPEN balls $B(i,r(i))$, each meeting $Z$, such that
\[
 \begin{gathered}
 q\in B(i,10r(i))\quad\Longrightarrow\quad
 r(q)\leq20r(i),\qquad r(i)/\rho(q)\geq T/20,\\
 Z\subset\bigcup_{i\in I}B(i,5r(i)).
 \end{gathered}
\]
If, in addition, EACH selected center satisfies the annular exclusion
\[
 Z\cap\{x\in X:r(i)/10\leq d(i,x)\leq10r(i)\}
       =\varnothing,\qquad i\in I,
\]
then
\[
 Z\subset\bigcup_{i\in I}B(i,r(i)/10).
\]
All radii and any attached pointwise witnesses are restricted from
the ORIGINAL fixed assignment; none is selected again after taking $I$.
The selected centers need not themselves lie in $Z$.
\end{theorem}
\begin{proof}
If $X$ or $Z$ is empty, use $I=\varnothing$; the conclusions hold.
Otherwise continuity, positivity and compactness give numbers
$0<\rho_{\min}\leq\rho(p)\leq\rho_{\max}<\infty$.
The original radius assignment therefore has the uniform positive
lower bound $T\rho_{\min}$ and finite upper bound $U\rho_{\max}$.
Apply LC62 and then LC63 to the totally bounded subset $V$ of $X$.
The resulting $I\subset V$ is finite and disjoint. Membership in
$P_0$ says each selected ball meets $Z$, and LC62 supplies the two
neighboring-scale estimates without changing any radius.

For $z\in Z$, positivity gives $z\in B_z$, so $z\in P_0$.
Choose $v\in V$ above $z$ using LC62. Either $v=z$ or
$B_z\subset B_v$; in both cases $z\in B_v$.
LC63 supplies $i\in I$ with $B_v\cap B_i\ne\varnothing$ and
$r(v)\leq2r(i)$, whence $z\in B(i,5r(i))$.
If the additional exclusion holds, $d(i,z)<5r(i)<10r(i)$
forces $d(i,z)<r(i)/10$. Equality at the INNER boundary is excluded
by the non-strict lower inequality in the hypothesis. This proves
the cover by open small cores. Restricting a fixed assignment to
$I$ changes no radius or attached witness and requires no new limit,
subsequence, or uniform tail.
\end{proof}

\subsection{Application boundary and the next geometric obligations}

For KL 11.5, take $Z$ to be the zero-stratum set defined in LC16,
$\rho$ to be the positive continuous volume-scale function, and
$r(p)=r_p^0$. Conditional on the pointwise package, LC58 supplies
one uniformly bounded scale interval and LC61 supplies the smooth
open-ball comparison using the same witnesses. Fix that entire
assignment before LC62--LC64. The inherited scale regularity yields
the extrema used here even though $p\mapsto r_p^0$ need not be
continuous. The metric results establish KL 11.5(1),(2) and, ONLY
when annular exclusion is supplied, (4). They do not accept the
finite-$K$ producer required to obtain the original pointwise package.

KL's notation in Section 2.2 is
$A(p,a,b)=\overline{B(p,b)}\setminus B(p,a)$.
On the complete Riemannian manifold in the application this is the
closed radial shell, including BOTH boundary spheres. LC62 also proves the neighboring
scale bound on the outer sphere, so the annular consumer need not
pass from an open-ball estimate by continuity of the chosen radii.
The earlier
explanatory paragraph following LC30 is corrected in this revision:
the source annulus has a closed outer ball. LC30's own buffered open
$U$ and compact closed $C$, and its proof, were already correct.
This correction is to our transcription, not an author erratum.
The strict small-core conclusion above specifically needs exclusion
at $d(i,z)=r(i)/10$.

Two geometric tasks remain before all of KL 11.5 can be accepted.
First, its part (3) must produce a $(1,\beta_1)$-splitting at every
point of the CLOSED annulus from the cone comparison, the neighboring
scale bound, and KL 4.15. Its radial coordinate in the metric
rescaled by $\rho(q)^{-1}$ is
\[
 \frac{r_i^0}{\rho(q)}\,\eta_i,
\]
with the quality asserted through KL 4.28 and the Lipschitz comparison
of KL 11.3. The additive centering convention must be bound to the
actual adapted-coordinate definition. Admitting this one-dimensional
splitting excludes rank zero; with LC20 it leaves ranks one and two.
The required parameter order and quantitative transfer remain to be
written, not inferred from LC63's metric cover.

Second, part (5) must show that the SAME model $N_i$ has at most one
end. The source argues that multiple ends force its Tits cone to be
$\mathbb R$, which gives a one-splitting at every point of
$B(i,r_i^0)$; this contradicts that ball's intersection with $Z$.
The splitting/end theorem and uniform transfer through the changing
$\rho(q)$ scale still require their precise bindings. This paragraph
records the supplied source route, not a fresh proof audit of the
end theorem. Retain both tasks here in Chapter 13; compatible
fibrations and their gluing remain Chapter 14 work.

\clearpage
\section[Annular splitting and radial normalization]{Annular splitting, buffered smoothing and radial normalization}
\label{sec:collapse-annular-splitting-normalization}

KL 11.5(3) has two conclusions: a one-splitting on the closed
annulus and an adapted-coordinate assertion for the prescribed radial
function. We supply the first from the existing Chapter 4 results,
then specify the normalization and a stability argument for the second.
The final comparison to the original radial function remains explicit.
The source checks cover KL 3.6, 4.15, 4.21--4.28 and 11.3--11.5;
precise locators and reading limits are recorded in
\code{reference_checks_revision94.md}. No new Lean result is claimed.

\subsection{Shorten the annular strainer to the required exact scale}

Set $a=1/10$, $b=10$ and $\kappa=1/60$. For $0<\sigma<1$ put
\[
 \begin{split}
 L&=\sigma^{-1},\\
 c_\sigma&=\frac2{\sqrt\sigma}
    \operatorname{arsinh}\!\left(\sinh(\sqrt\sigma L)
                                      \cos(\sigma/2)\right)<2L.
 \end{split}
\]
Thus $c_\sigma$ is the opposite side of the equal-leg model triangle
with legs $L$, angle $\pi-\sigma$ and curvature $-\sigma$.
Choose $0<\theta<\pi/2$ with
$2L\cos(\theta/2)>c_\sigma$, and define
\[
 \begin{split}
 \delta_\sigma&=\min\{a/60,\,1/(4b+20),\,
                              a(1-\cos\theta)/60\},\\
 \Lambda_\sigma&=\max\{2L/a,\,\kappa/\sqrt\sigma,\,2\}.
 \end{split}
\]
These are sufficient project constants, not constants quoted from KL.
The potentially very small $\theta$ is chosen AFTER $\sigma$.

\begin{lemma}[LC65: annular strainers at an exact prescribed scale]
\label{lem:collapse-annular-exact-strainer}
Let $(M,g_0,p)$ be a complete connected smooth Riemannian
$n$-manifold without boundary, with $\sec_{g_0}\geq-\kappa^2$
on $B_{g_0}(p,400)$. Supply an actual pointed KL $\delta$-map
from $(M,d_{g_0},p)$ to an AC82 cone, with
$0<\delta<\delta_\sigma$.
For every $q$ in the CLOSED shell $a\leq d_{g_0}(p,q)\leq b$
and every $\lambda\geq\Lambda_\sigma$, the metric
$g_q=\lambda^2g_0$ has curvature at least $-\sigma$ on
$B_{g_q}(q,L)$, and $q$ has a KL one-strainer of quality $\sigma$
at EXACT scale $L$ in that metric, using comparison curvature
$-\sigma$. The inward anchor lies on a minimizing segment from
$q$ to $p$; it need not equal $p$. The local comparison and
Riemannian-to-Alexandrov bindings remain those of AC02/AC64.
\end{lemma}
\begin{proof}
Write $D=d_{g_0}(p,q)\in[a,b]$.
LC25 gives $q'$ with $d_{g_0}(p,q')=2D$ and
$D\leq\ell=d_{g_0}(q,q')<D+15\delta$.
On a minimizing segment from $q$ to $q'$, choose $z$ at distance
$D$ from $q$. Then
\[
 2D-15\delta<d_{g_0}(p,z)\leq2D.
\]
Indeed $d(p,z)\geq d(p,q')-d(z,q')=3D-\ell$.
Put $E=2D-d(p,z)\in[0,15\delta)$.
For the comparison angle $\omega$ of the equal-leg triangle
$(q;p,z)$ in curvature $-\kappa^2$, the model cosine law gives
\[
 \begin{aligned}
 1+\cos\omega
 &=\frac{\cosh(2\kappa D)-\cosh(\kappa(2D-E))}
         {\sinh^2(\kappa D)}\\
 &\leq2\kappa E\coth(\kappa D)
 \leq\frac{4E}{D}<\frac{60\delta}{a}<1-\cos\theta.
 \end{aligned}
\]
Here $\kappa D\leq1/6$, and $x\coth x\leq2$ for
$0<x\leq1$ follows from $\sinh x\geq x$ and $\cosh1<2$.
The degenerate case $E=0$ has angle $\pi$.
Consequently $\omega>\pi-\theta$.

Put $s=L/\lambda\leq a/2<D$. Shorten the chosen segments
$q$ to $p$ and $q$ to $z$ to length $s$, obtaining $a_q,b_q$.
Apply AC64 at center $q$, with its radial length $D$ and outer
radius $100$: $8D\leq80<100$, and
$B_{g_0}(q,100)\subset B_{g_0}(p,110)\subset B_{g_0}(p,400)$.
Completeness, local compactness and local Hausdorff dimension $n$
come from the inherited smooth Riemannian foundation; curvature
uses the retained local comparison convention. AC64's proof gives
the complete intrinsic buffers for every side replacement. It follows
that the shortened comparison angle in curvature $-\kappa^2$
is still greater than $\pi-\theta$.

If $c=d_{g_0}(a_q,b_q)$, the equal-leg model identity implies
\[
 \sinh(\kappa c/2)>\cos(\theta/2)\sinh(\kappa s),
 \qquad c>2s\cos(\theta/2).
\]
For the second implication use
$\sinh(tu)\leq u\sinh t$ for $t\geq0$, $0\leq u\leq1$.
In the metric $g_q$, both legs have length $L$ and the opposite
side is $\lambda c>2L\cos(\theta/2)>c_\sigma$.
Strict monotonicity of the model opposite side with its angle gives
a comparison angle greater than $\pi-\sigma$ in curvature
$-\sigma$. This proves the one-strainer claim without assuming
that the comparison angle is unchanged when its model curvature changes.
Finally $\sec_{g_q}\geq-\kappa^2/\lambda^2\geq-\sigma$;
the required $L$-ball corresponds to the $s$-ball about $q$ and
lies in the verified curvature region. All choices are uniform in
$q$ and in $\lambda\geq\Lambda_\sigma$.
\end{proof}

\subsection{Exclude the zero stratum and apply the selected cover}

\begin{theorem}[LC66: annular one-splitting and the selected small-core cover]
\label{thm:collapse-annular-splitting-cover}
Fix $0<\beta_1<\beta_2<\beta_3<1$ and use the rank convention
LC16 on a closed connected smooth Riemannian three-manifold $(M,g)$.
Retain a positive continuous $\rho$ and a fixed pointwise assignment
$T\rho(p)\leq r(p)\leq U\rho(p)$ as in LC64.
Choose $\sigma=\sigma(3,\beta_1)$ from AC55, then the constants
of LC65, and require
\[
 T\geq20\Lambda_\sigma.
\]
Take the finite selection $I$ from LC64, before imposing its extra
annular exclusion. Suppose that for every $i\in I$, the metric
$g_i=r(i)^{-2}g$ on the SAME manifold has the LC65 curvature bound
and cone approximation with error less than $\delta_\sigma$.
Then every $q$ with
\[
 r(i)/10\leq d_g(i,q)\leq10r(i)
\]
admits a normalized $(1,\beta_1)$-splitting of
$(M,\rho(q)^{-2}g,q)$ and is not in the zero stratum. If LC20's
no-three-stratum conclusion also holds, its rank is one or two.
Consequently the selected OPEN tenth-radius balls cover the entire
zero stratum. These assertions are conditional on AC55's retained
comparison/splitting inputs and the explicitly supplied cone packages.
\end{theorem}
\begin{proof}
Since $i$ is a maximal candidate, LC62, including its outer-boundary
case, gives $\lambda=r(i)/\rho(q)\geq T/20\geq\Lambda_\sigma$.
LC65 supplies an exact-scale one-strainer and the curvature and
dimension hypotheses for AC55, all in
$\lambda^2g_i=\rho(q)^{-2}g$.
AC55 gives the required splitting. Its coordinate for the inward
anchor $a_q$ is $L-d_{g_q}(a_q,\cdot)$; reflection of the
Euclidean factor gives the outward-oriented coordinate
$d_{g_q}(a_q,\cdot)-L$ without changing any approximation error.
This coordinate is based on the SHORTENED anchor, so this step does
not assert that it is the original centered distance to $i$.

Admitting rank one implies $s(q)\geq1$ directly from LC16's finite
maximum. Exclusion of rank three, when supplied, leaves ranks one
and two. Thus the zero stratum misses the whole closed shell,
including both boundary spheres. LC64 now gives its open small-core
cover. Neither this exclusion nor the cover requires smooth adapted
coordinates or changes any of the selected radii or model witnesses.
\end{proof}

This proves the annular stratum and covering portions of KL 11.5,
conditional on the stated geometric inputs. Choosing $\beta_1$ first,
then $\sigma,\theta,\delta_\sigma,\Lambda_\sigma$, then taking
$T\geq20\Lambda_\sigma$, makes the order of constants explicit.
If the pointwise cone package is produced through LC58, these
parameters must be fixed before its uniform interval and tail are
chosen. The finite-$K$ producer remains open. No dependence on an
unproved adapted-coordinate assertion has entered the covering proof.

\subsection{A buffered radial choice and its exact change of units}

\begin{lemma}[LC67: buffered smoothing and centered radial normalization]
\label{lem:collapse-buffered-radial-normalization}
For every $0<\varepsilon<1/4$ there is $\delta_0>0$ such that,
under LC30's geometric hypotheses with $0<\delta<\delta_0$, its
radial function $\eta$, for every $0<e<1/40$, can be chosen
smooth on a neighborhood of
\[
 C^+=\{x:3/40\leq d_{g_0}(p,x)\leq11\},
\]
while retaining its value error $e$, global difference-Lipschitz,
equality off $U$ and
gradient conclusions on the original $C$.
This choice is made BEFORE joint witness construction and ball selection.
For such a fixed $\eta$, every $q$ with
$1/10\leq d_{g_0}(p,q)\leq10$ and every $\lambda\geq80$
have the following property. Put $g_q=\lambda^2g_0$ and
\[
 \begin{split}
 u_q(x)&=d_{g_q}(p,x)-d_{g_q}(p,q),\\
 \psi_q(x)&=\lambda\bigl(\eta(x)-\eta(q)\bigr).
 \end{split}
\]
Then $\psi_q(q)=u_q(q)=0$, $\psi_q$ is smooth on
$B_{g_q}(q,1)$, and
\[
 \operatorname{Lip}_{g_q}(\psi_q-u_q)\leq\varepsilon,
 \qquad |\psi_q(x)-u_q(x)|\leq\varepsilon d_{g_q}(q,x).
\]
In particular the value error on that unit ball is less than
$\varepsilon$ and $\psi_q$ is $(1+\varepsilon)$-Lipschitz.
The uncentered source coordinate is $\lambda\eta$ and has the
same derivative; an arbitrary already selected LC30 witness is not
silently assumed to have this larger smooth neighborhood.
\end{lemma}
\begin{proof}
Repeat LC30's use of LC27 on
$U=\{1/20<d_{g_0}(p,\cdot)<20\}$, with the direction tolerance
prescribed by LC28. Use $C^+$ in LC28 in place of $C$.
It is compact by properness and has positive margins inside the
same open $U$. The same direction estimate applies throughout $U$.
Thus LC28 supplies smoothness near $C^+$, the arbitrarily prescribed
value error, equality to the distance off $U$, and
$\operatorname{Lip}_{g_0}(\eta-d_{g_0,p})\leq\varepsilon$.
The nonnegativity argument in LC30 is unchanged, and LC29 gives
the gradient bounds, in particular on the original $C\subset C^+$.
The cutoff conclusions referring to the old levels use exactly the
same original estimates. This is an allowed stronger choice at the
smoothing step, not a later replacement of selected data.

The $g_q$ unit ball is the $g_0$ ball of radius $1/\lambda$.
For its points the distance to $p$ lies strictly between
$1/10-1/80=7/80>3/40$ and $10+1/80<11$.
It therefore lies in the smooth region. Write
$f=\eta-d_{g_0,p}$. Then
\[
 \psi_q-u_q=\lambda(f-f(q)),\qquad
 d_{g_q}=\lambda d_{g_0}.
\]
Multiplying both the function and the metric by the SAME $\lambda$
preserves its Lipschitz constant. Evaluating at $q$ gives the
pointwise estimate. Adding the one-Lipschitz $u_q$ proves the last
claim. Centering avoids the potentially large bound
$\lambda\|\eta-d_{g_0,p}\|_\infty$; no upper bound on $\lambda$
is needed. In the application, $\lambda=r(i)/\rho(q)$, so
$g_q=\rho(q)^{-2}g$ as required.
\end{proof}

\subsection{What the adapted-coordinate conclusion still requires}

For clarity, in rank one KL Definition 4.21 says the following.
Fix a normalized $(1,\beta)$-splitting $\alpha=(\Phi,v)$ at $q$,
and $0<\beta\leq\gamma<1$. A function $\phi$ on $B(q,1)$ is
$\alpha$-adapted of quality $\gamma$ if it is smooth and
$(1+\gamma)$-Lipschitz, its image has Hausdorff distance at most
$\gamma$ from $(\phi(q)-1,\phi(q)+1)$, and for EVERY
$x\in B(q,1)$, $y\in B(q,\gamma^{-1})$ with $d(x,y)>1$,
and EVERY unit initial velocity $w$ of a minimizing segment from
$x$ to $y$, one has
\[
 \left|D\phi_x(w)-\frac{\Phi(y)-\Phi(x)}{d(x,y)}\right|<\gamma.
\]
The far-point test domain uses the ADAPTED QUALITY $\gamma$, not
$\beta$. Adding a constant to $\phi$ preserves all three conditions;
its target interval is centered at $\phi(q)$, not necessarily zero.
Reflecting both $\phi$ and $\Phi$ also preserves them. These are
immediate from the displayed conditions. This local comparison
belongs here in Chapter 13; general compatible fibrations remain
Chapter 14 work.

\begin{lemma}[LC68: stability of a supplied adapted radial comparison]
\label{lem:collapse-adapted-radial-stability}
Use the SAME $(1,\beta)$-splitting $\alpha$ and suppose $\phi(q)=0$
is an $\alpha$-adapted coordinate of quality $\gamma$, with
$0<\beta\leq\gamma<\zeta<1$.
Suppose a smooth $\psi$ on $B(q,1)$ has $\psi(q)=0$ and
$\psi-\phi$ is $h$-Lipschitz there for the AMBIENT distance,
where $h\geq0$ and $\gamma+h\leq\zeta$.
Then $\psi$ is $\alpha$-adapted of quality $\zeta$.
In LC67's setting it suffices to supply, for this same splitting,
such a $\phi$ with
\[
 \operatorname{Lip}_{B(q,1)}(u_q-\phi)\leq t,
 \qquad \gamma+t+\varepsilon\leq\zeta.
\]
No existence of that reference $\phi$ or its comparison bound is
asserted by this lemma.
\end{lemma}
\begin{proof}
The Lipschitz constants add to at most $1+\gamma+h\leq1+\zeta$.
Since both functions vanish at $q$, their values differ by at most
$h d(q,x)\leq h$ on the unit ball. Their images have Hausdorff
distance at most $h$, using the same source points; the triangle
inequality for the two directed set distances then gives image error
at most $\gamma+h\leq\zeta$ from the SAME interval $(-1,1)$.
This uses infimum distances and does not assume nearest image points
or closed images.

For a unit tangent vector $w$ at a point of the open unit ball,
differentiate the ambient $h$-Lipschitz difference along a short
smooth curve with initial velocity $w$. Its derivative has absolute
value at most $h$. The far-point domain
$B(q,\zeta^{-1})$ is contained in $B(q,\gamma^{-1})$;
hence the adapted derivative test for $\phi$ gives a strict error
less than $\gamma+h\leq\zeta$ for $\psi$, for every required
minimizing direction. This proves all three conditions without
assuming that $B(q,1)$ is geodesically convex.
For the radial specialization, LC67 bounds
$\operatorname{Lip}(\psi_q-u_q)$ by $\varepsilon$; add the
explicit $t$ bound and take $h=t+\varepsilon$.
\end{proof}

KL 4.23 supplies some adapted coordinates from a sufficiently good
splitting and the stated local curvature bound, using smoothing and
the comparison estimate in 4.25. KL 4.28 compares TWO functions for
the SAME splitting: one is already adapted; the other is either
adapted or satisfies its specified directional tests. Its proof uses
derivatives and an integration argument with common centering.
The Lipschitz difference between $\eta$ and the radial distance,
by itself, does not verify those directional tests or identify the
coordinate of the shortened-anchor splitting in LC66.

The remaining radial task is therefore concrete: construct a splitting
with the original centered distance $u_q$ (or compare the LC66
coordinate to it with sufficient control), and produce the reference
adapted coordinate and the difference bound required by LC68, or
verify KL 4.28's alternative hypotheses directly. Preserve the
outward sign, the factor $r(i)/\rho(q)$, the buffered smooth domain,
the same splitting, and all tested minimizing directions.
Use a cone error below both LC65's $\delta_\sigma$ and LC67's
$\delta_0$. When setting up new joint witnesses, include LC67's stronger smooth
choice at their original construction and take $T$ large enough for
both LC66 and $T/20\geq80$. Previously fixed witnesses need the
explicit domain check; none is replaced after selection.
The at-most-one-end model argument and the finite-$K$ producer
remain open. No full KL 11.5 or compatible-fibration conclusion is
claimed by this iteration.

\clearpage
\section[The prescribed radial coordinate]{Recovering the prescribed radial coordinate and its adaptedness}
\label{sec:collapse-prescribed-radial-coordinate}

The remaining coordinate step in KL 11.5(3) uses the ORIGINAL
selected center, not the shortened inward anchor of LC66. We first
retain that original distance in a compactness argument. We then
compare gradients by testing along actual minimizing segments to
that center. The latter argument writes the rank-one mechanism
behind KL 4.28 without presupposing an ambient Lipschitz bound
between two smooth coordinates on a possibly nonconvex ball.
KL 4.23's existence theorem is isolated as a source-qualified input;
its full implementation is not proved by the following application.
See \code{reference_checks_revision95.md} for exact reading limits.

\subsection{Calibration of the original distance on an outward segment}

\begin{lemma}[LC69: the original center calibrates the outward prefixes]
\label{lem:collapse-original-radial-calibration}
Use LC65's constants and construction, and put
$D=d_{g_0}(p,q)$. Let $\xi:[0,D]\to M$ be the chosen unit-speed
minimizing segment from $q$ to $z$. For every $0<t\leq D$,
\[
 0\leq D+t-d_{g_0}(p,\xi(t))
       <2t(1-\cos\theta).
\]
In the metric $g_q=\lambda^2g_0$, writing
$u(x)=d_{g_q}(p,x)-d_{g_q}(p,q)$, the corresponding prefix at
normalized distance $T=\lambda t$ satisfies
\[
 T-2T(1-\cos\theta)<u(\xi(T/\lambda))\leq T.
\]
This estimate has no factor $\lambda\delta$ and no upper bound on
$\lambda$. On every inward prefix of a minimizing segment to $p$,
$u$ equals minus the normalized prefix length exactly.
\end{lemma}
\begin{proof}
LC65 gives an angle greater than $\pi-\theta$ for the equal-leg
triangle $(q;p,z)$ in comparison curvature $-\kappa^2$.
Apply AC64 with one radial length still $D$ and the other shortened
to $t$, in the same buffer $8D<100$ used by LC65. The opposite
side $d(p,\xi(t))$ is therefore at least the model side $c$ with
legs $D,t$ and angle $\pi-\theta$.
Since $\theta<\pi/2$, the hyperbolic cosine law implies $c\geq D$.
It also gives
\[
 \cosh(\kappa(D+t))-\cosh(\kappa c)
    =(1-\cos\theta)\sinh(\kappa D)\sinh(\kappa t).
\]
On the interval $[c,D+t]$ the derivative of the left model
function is at least $\kappa\sinh(\kappa D)$. Consequently
\[
 D+t-c\leq (1-\cos\theta)\frac{\sinh(\kappa t)}{\kappa}
          <2t(1-\cos\theta),
\]
because $0<\kappa t\leq\kappa D\leq1/6$ and
$\sinh x/x\leq\cosh x<2$ there. The triangle inequality supplies
the nonnegative lower bound. Multiply by $\lambda$ for the
normalized assertion. Along a minimizing segment to $p$, the
remaining distance decreases by the traversed length, giving the
exact inward assertion without differentiability of the distance.
\end{proof}

\begin{theorem}[LC70: a splitting with the original radial coordinate]
\label{thm:collapse-original-radial-splitting}
For every integer $n\geq1$ and $0<\beta<1$ there are
$\delta_*>0$ and $\Lambda_*>0$ with the following property.
Let $(M,g_0,p)$ be a complete connected smooth Riemannian
$n$-manifold without boundary, with
$\sec_{g_0}\geq-1/60^2$ on $B_{g_0}(p,400)$, and supply an
actual pointed KL $\delta$-map to an AC82 cone, where
$0<\delta<\delta_*$. Then for EVERY
\[
 1/10\leq d_{g_0}(p,q)\leq10,\qquad\lambda\geq\Lambda_*,
\]
there is a normalized $(1,\beta)$-splitting of
$(M,\lambda^2g_0,q)$ whose real coordinate is EXACTLY
\[
 u_q(x)=d_{\lambda^2g_0}(p,x)-d_{\lambda^2g_0}(p,q).
\]
The compactness, local comparison and global splitting inputs are
those retained in AC61. No explicit modulus, upper bound on
$\lambda$, or monotonicity of a previously chosen modulus is claimed.
\end{theorem}
\begin{proof}
Suppose the assertion fails for fixed $n,\beta$. For each integer
$i\geq2$, take $\sigma_i=1/i$ in LC65 and choose its angle
$0<\theta_i<1/i$ while still satisfying
$2\sigma_i^{-1}\cos(\theta_i/2)>c_{\sigma_i}$.
Failure supplies a counterexample with
\[
 0<\delta_i<\min\{\delta_{\sigma_i},1/i\},\qquad
 \lambda_i\geq\max\{\Lambda_{\sigma_i},i\}.
\]
LC65 gives an exact-scale one-strainer in
$h_i=\lambda_i^2g_{0,i}$, with length $L_i=\sigma_i^{-1}$,
quality $\sigma_i$, and curvature at least $-\sigma_i$ on
$B_{h_i}(q_i,L_i)$. Use its actual inward and outward minimizing
segments for the prefixes. AC61's proof applies to these segments:
they are admissible almost-shortest paths with any positive error
tending to zero. Thus, after one subsequence, the spaces have a
complete proper nonnegative limit $(X,o)$, the prefixes converge to a
line, and an onto splitting $X=\mathbb R\times Y$ aligns its
OUTWARD ray with the positive axis. Write $U$ for its real
coordinate and $f_i$ for the same pointed approximation maps.

Put $u_i=d_{h_i}(p_i,\cdot)-d_{h_i}(p_i,q_i)$. These functions
are one-Lipschitz and vanish at $q_i$. For fixed $T>0$, and all
large $i$, denote the outward and inward time-$T$ prefixes by
$x_i^+(T)$ and $x_i^-(T)$. LC69 gives
\[
 \begin{split}
 T-2T(1-\cos\theta_i)&<u_i(x_i^+(T))\leq T,\\
 u_i(x_i^-(T))&=-T.
 \end{split}
\]
Fix $R>0$ and put $B=R+1$. For $T>B$, let $e_i$ bound the
approximation distortion on a fixed ball containing the source
$R$-ball and both prefixes, and let $z_i(T)$ bound the distances
of their images from the limiting points $(\pm T,y)$.
For this fixed $T$, both errors tend to zero. The one-Lipschitz
bound and the two calibrations give, for $w=f_i(x)$,
\[
 \begin{split}
 T-d_X(w,(T,y))-E_i(T)&\leq u_i(x)\\
                       &\leq d_X(w,(-T,y))-T+E_i(T),\\
 E_i(T)&=2T(1-\cos\theta_i)+e_i+z_i(T).
 \end{split}
\]
Eventually $d_X(o,w)\leq B$. The product computation in AC60,
namely rationalization of $\sqrt{(T\pm U(w))^2+r^2}$, yields
\[
 \sup_{x\in\overline B_{h_i}(q_i,R)}
       |u_i(x)-U(f_i(x))|
 \leq E_i(T)+\frac{B^2}{2(T-B)}.
\]
First choose $T$ large, then $i$ large. This proves uniform
convergence of the ORIGINAL distance coordinates on every fixed
ball. In particular no estimate $\lambda_i\delta_i\to0$ was
needed or asserted. AC52 now restores $u_i$ as the EXACT real
coordinate of a $(1,\beta)$-splitting for all sufficiently large
$i$, contradicting the chosen counterexamples.
\end{proof}

\subsection{The source input for an adapted reference coordinate}

\begin{foundation}[LC71: rank-one adapted-coordinate existence]
\label{input:collapse-adapted-reference}
For every $n\geq1$ and $0<\gamma<1$ there is
$0<\nu_0<\min\{\gamma,1/4\}$ such that the following holds
for every $0<\nu\leq\nu_0$.
Let $(M,h,q)$ be complete connected smooth Riemannian of
dimension $n$ and without boundary. Suppose
\[
 \sec_h\geq-\nu^2\quad\hbox{on }B_h(q,\nu^{-1}),
\]
and let $\alpha$ be a supplied normalized $(1,\nu)$-splitting.
There is a smooth $\alpha$-adapted coordinate $\phi$ of quality
$\gamma$ on $B_h(q,1)$, and it can be centered so $\phi(q)=0$.
This is the rank-one specialization of KL 4.23, with its curvature
SQUARE retained. It is a source-qualified existence input, not a
newly proved theorem or a declaration validated against PC.
\end{foundation}

The source statement supplies a sufficiently small parameter.
For every smaller $\nu$, the direct proof in LC79 below supplies
the SAME-map conclusion with all source-domain margins checked.
Bare numerical weakening of a KL error is insufficient without
checking that coverage witnesses survive the restricted domain.
Translation of $\phi$ preserves Definition 4.21.

The proof packet has three parts. AC56 supplies the product-limit
route used by KL 4.16. Local comparison controls the long anchors
and the directional estimates of Sublemma 4.25. LC28 supplies the
localized smoothing. The source also compresses the Lipschitz and
image-coverage estimates into a similar compactness argument.
These must be retained in implementing LC71; the three dependency
links do not by themselves prove the input. In particular,
Sublemma 4.25 passes to limits of COMPARISON angles, denoted with
a tilde in the source. It is not an assertion of automatic convergence
of arbitrary actual angles under GH convergence. Complete intrinsic
buffers and the smoothing interface retain their existing qualifications.
The following consumer does not use KL 4.28 as a black box.

\subsection{A calibrated direction controls the gradient}

\begin{lemma}[LC72: transfer from a calibrated radial direction]
\label{lem:collapse-calibrated-gradient-transfer}
Let $(M,h,q)$ be complete connected smooth Riemannian without
boundary, and let $d_h(p,q)>3$. Suppose a normalized
$(1,\nu)$-splitting $\alpha$ has real coordinate
$u=d_h(p,\cdot)-d_h(p,q)$, and a smooth $\alpha$-adapted
coordinate $\phi$ of quality $\gamma$ is supplied, where
$0<\nu<\gamma<1/3$ and $\phi(q)=0$.
Let $\psi$ be smooth on $B_h(q,1)$ with $\psi(q)=0$ and
\[
 \operatorname{Lip}_{B_h(q,1)}(\psi-u)\leq\varepsilon,
 \qquad 0<\varepsilon<1/4,
\]
using the AMBIENT distance. Set
\[
 A(\gamma,\varepsilon)
   =\sqrt{4\gamma+\gamma^2}+\sqrt{4\varepsilon+\varepsilon^2}.
\]
Then, on that unit ball,
\[
 \|D\psi-D\phi\|\leq A(\gamma,\varepsilon),\qquad
 |\psi(x)-\phi(x)|\leq A(\gamma,\varepsilon)d_h(q,x).
\]
For any $\gamma<\zeta<1$ with
$\gamma+A(\gamma,\varepsilon)\leq\zeta$, the function $\psi$
is $\alpha$-adapted of quality $\zeta$. No existence of $\phi$
is asserted here, and no global ambient Lipschitz bound on
$\psi-\phi$ is inferred from the derivative estimate.
\end{lemma}
\begin{proof}
Fix $x\in B_h(q,1)$ and choose ANY minimizing segment from $x$
to $p$, with initial unit vector $v$. Its length exceeds two.
Let $y$ be its point at distance two from $x$. Then
$d_h(q,y)<3<\gamma^{-1}$, and
\[
 d_h(x,y)=2,\qquad u(y)-u(x)=-2.
\]
The adapted derivative test for $\phi$, applied to this segment,
gives $|D\phi_x(v)+1|<\gamma$. Also
$\|\nabla\phi(x)\|\leq1+\gamma$.
On a sufficiently short initial portion of the SAME segment,
$u$ decreases at unit speed exactly, even when it is not
differentiable at $x$. The difference-Lipschitz hypothesis and
smoothness of $\psi$ therefore give
$|D\psi_x(v)+1|\leq\varepsilon$.
Moreover $\psi$ is $(1+\varepsilon)$-Lipschitz, so its gradient
norm is at most $1+\varepsilon$.

For $F=\phi$ or $\psi$, with the corresponding error
$e=\gamma$ or $\varepsilon$, expand the squared norm:
\[
 \|\nabla F(x)+v\|^2
 =\|\nabla F(x)\|^2+1+2DF_x(v)
 \leq(1+e)^2+1+2(-1+e)=4e+e^2.
\]
The triangle inequality gives the claimed derivative bound.
A minimizing segment from $q$ to $x$ stays inside $B_h(q,1)$,
so integration from their common zero value proves the pointwise
bound. This uses only basepoint-to-point segments, not convexity
of the entire unit ball.

The images of $\phi$ and $\psi$ have Hausdorff distance at most
$A(\gamma,\varepsilon)$, using the same source points. The
image condition for $\phi$ therefore gives error at most
$\gamma+A(\gamma,\varepsilon)\leq\zeta$ from $(-1,1)$.
Since $\sqrt{4\varepsilon+\varepsilon^2}\geq\varepsilon$,
$\psi$ is $(1+\zeta)$-Lipschitz. Finally, every far-point test
for quality $\zeta$ lies inside the test ball for quality $\gamma$.
The derivative difference bound transfers the STRICT error
$<\gamma$ to error $<\gamma+A(\gamma,\varepsilon)\leq\zeta$,
for EVERY required minimizing direction. This verifies all of
Definition 4.21, without a first-variation formula at a cut point.
\end{proof}

\subsection{Binding the selected radial function}

\begin{theorem}[LC73: the prescribed annular adapted coordinate]
\label{thm:collapse-prescribed-annular-adapted}
Fix $0<\beta_1<\zeta<1$. Conditional on LC71 and the retained
comparison, splitting and smoothing inputs, choose
$\varepsilon,\delta_*,\Lambda_*>0$ as in the proof below.
Use LC64's fixed scale assignment and selection on a closed
connected smooth Riemannian three-manifold, with
$T\geq20\Lambda_*$. At each selected center $i$, supply the
LC65 curvature bound and an actual cone map of error less than
$\delta_*$ in $g_i=r(i)^{-2}g$. Supply the LC67 buffered radial
function $\eta_i$, with difference-Lipschitz error at most
$\varepsilon$, chosen at the ORIGINAL joint construction.
For every point in the CLOSED shell
$r(i)/10\leq d_g(i,q)\leq10r(i)$, there is a normalized
$(1,\beta_1)$-splitting $\alpha_q$ of $(M,\rho(q)^{-2}g,q)$
with original real coordinate
\[
 u_q(x)=d_{\rho(q)^{-2}g}(i,x)-d_{\rho(q)^{-2}g}(i,q),
\]
for which the prescribed function
\[
 \psi_q=\frac{r(i)}{\rho(q)}\bigl(\eta_i-\eta_i(q)\bigr)
\]
is adapted of quality $\zeta$. The uncentered source function
$(r(i)/\rho(q))\eta_i$ has the same adaptedness, with its
image interval centered at its own value at $q$.
No selected radius, model or radial function is replaced.
\end{theorem}
\begin{proof}
Choose $0<\gamma<\min\{1/3,\zeta\}$ and
$0<\varepsilon<1/4$ so small that
$\gamma+A(\gamma,\varepsilon)<\zeta$.
Take $0<\nu<\min\{\nu_0(3,\gamma),\beta_1/4,\gamma\}$
from LC71, then take LC70's constants for $(n,\beta)=(3,\nu)$.
Decrease $\delta_*$ to satisfy also LC67's smoothing threshold
for this $\varepsilon$, and enlarge $\Lambda_*$ to be at least
\[
 \max\{80,\nu^{-1}\}.
\]
Fix all these constants BEFORE choosing the joint witnesses,
their uniform scale interval and the selected balls.

For $q$ in the closed shell, LC62 gives
$\lambda=r(i)/\rho(q)\geq T/20\geq\Lambda_*$.
LC70 supplies a $(1,\nu)$-splitting with real coordinate exactly
$u_q$ in $h_q=\lambda^2g_i=\rho(q)^{-2}g$.
It is also a $(1,\beta_1)$-splitting: a coverage lift has error
less than $2\nu$ and radius less than
$\beta_1^{-1}-\beta_1+3\nu<\beta_1^{-1}$.
Thus the SAME map survives the restriction, as in the corrected LC57.
For LC71 the required ball has $g_i$-radius
$(\lambda\nu)^{-1}\leq1$ about $q$, so it lies in
$B_{g_i}(i,11)\subset B_{g_i}(i,400)$.
Furthermore
\[
 \sec_{h_q}\geq-\frac{1}{60^2\lambda^2}\geq-\nu^2.
\]
Thus LC71 gives a centered adapted reference $\phi_q$ of quality
$\gamma$ for this SAME splitting. LC67 gives smoothness of
$\psi_q$ on the normalized unit ball and
$\operatorname{Lip}(\psi_q-u_q)\leq\varepsilon$ there.
Finally $d_{h_q}(i,q)\geq\lambda/10\geq8>3$.
All hypotheses of LC72 hold, proving adaptedness of quality
$\zeta$. Translation gives the uncentered assertion directly
from Definition 4.21.
\end{proof}

LC69--LC73 provide a written conditional route for the prescribed
radial part of KL 11.5(3). Together with LC66 they retain the
annular rank and small-core coverage assertions. The earlier open
coordinate comparison is superseded by this route; no claim is made
that LC68's stronger ambient difference hypothesis has been proved.
LC71 remains a named source-qualified input with its proof packet
still to implement. This is separate from the finite-$K$ smooth-model
producer, which also remains open, and from the PC release gate.

The next geometric application is KL 11.5(5): prove that the SAME
selected model has at most one end, with its multiple-end splitting,
Tits-cone identification and transfer to every point of the selected
unit-radius ball checked explicitly. That remains Chapter 13 work.
General compatible fibrations and their gluing remain Chapter 14.
No new Lean theorem, full KL 11.5 acceptance or production environment
pin is claimed here.

\clearpage
\section[The selected model has at most one end]{Ends, line cones and the selected model}
\label{sec:collapse-selected-model-ends}

We now supply KL 11.5(5) for the SAME model attached to each
selected ball. The argument needs neither the classification list
in KL 3.11 nor the adapted radial coordinate. The meaning of
``at most one end'' is that, for every compact set, its complement
has at most one unbounded connected component. In the present
complete proper connected manifold setting this is the usual
zero-or-one-end alternative; compact models have zero ends.

\begin{lemma}[LC74: two ends produce a line]
\label{lem:collapse-ends-produce-line}
Let $X$ be a proper geodesic space. If $X\setminus K$ has two
unbounded connected components for some compact $K$, then $X$
contains an isometric copy of $\mathbb R$.
\end{lemma}
\begin{proof}
Fix $o\in X$ and choose points $x_j,y_j$ in those two components
whose distances to $o$ tend to infinity. A minimizing segment
joining them meets $K$: otherwise its connected image would lie
in a single component of $X\setminus K$.
Choose a meeting point $z_j\in K$ and parametrize the segment
by arc length with $z_j$ at time zero. Both parameter lengths tend
to infinity, since $d(o,z_j)$ is bounded. On every bounded time
interval the segments are one-Lipschitz and lie in one compact
ball. A diagonal compactness argument on rational times, followed
by the one-Lipschitz extension to real times, gives a limit
$\gamma:\mathbb R\to X$. Along each original segment,
$d(\gamma_j(s),\gamma_j(t))=|s-t|$; continuity of distance gives
the same equality in the limit. Thus $\gamma$ is a line.
\end{proof}

\begin{lemma}[LC75: the supplied cone of a multiple-ended model is a line]
\label{lem:collapse-model-line-cone}
Let $(N,n)$ be a complete connected proper nonnegatively curved
Alexandrov space, with the retained AC43 splitting input and a
supplied LC21 cone $(C,o)$. If $N$ has more than one end, then
$N$ is isometric to $\mathbb R\times Y$ with $Y$ compact, and
there is a pointed onto isometry $(C,o)\cong(\mathbb R,0)$.
This applies in particular to the SAME nonnegative Riemannian
model in LC56, with its inherited Riemannian-to-Alexandrov binding.
\end{lemma}
\begin{proof}
LC74 gives a line and AC43 gives a product splitting. AC44
ensures that $Y$ is proper and geodesic. We show that unbounded
$Y$ would force $N$ to have at most one end.
For any compact $K\subset\mathbb R\times Y$, choose $R>0$ with
$K\subset[-R,R]\times\overline B_Y(y_0,R)$ and choose
$z_*\in Y$ with $d_Y(y_0,z_*)>R+1$.
Every point outside the compact box
$[-R-1,R+1]\times\overline B_Y(y_0,R+1)$ can be joined to
$(R+2,z_*)$ in the complement of $K$. If its $Y$ coordinate
lies outside the latter ball, first move vertically to height
$R+2$, then move within $Y$ at that height. Otherwise its real
coordinate has absolute value greater than $R+1$; first move
in $Y$ to $z_*$ at that height, then move vertically.
All these paths avoid $K$. Every unbounded component of the
complement meets the exterior of the compact box, so there is
only one. This contradiction shows $Y$ bounded, hence compact.

Translate the real coordinate so the marked point is $(0,y_0)$,
and set $D=\operatorname{diam}Y$. Projection
\[
 F_R:(R^{-1}N,n)\longrightarrow(\mathbb R,0),\qquad
 F_R(t,y)=t/R
\]
has distortion at most $D/R$, by the Pythagorean formula.
Every target point $s$ has preimage $(Rs,y_0)$ at normalized
distance $|s|$ from the basepoint. Thus it is a KL approximation
of arbitrarily small error for every sufficiently large real $R$.
The LC21 maps also make these SAME rescaled models converge to
$C$. Apply MC24 uniqueness of complete proper pointed limits
along any sequence $R\to\infty$ to obtain the claimed actual
pointed isometry. No cone witness or model is replaced.
\end{proof}

\begin{lemma}[LC76: a line approximation excludes rank zero on the unit ball]
\label{lem:collapse-line-cone-unit-splitting}
For every $0<\beta<1$ there are $\delta_\ell>0$ and
$\Lambda_\ell>0$ such that the following holds.
Let $(M,g_0,p)$ be complete connected smooth Riemannian of
dimension three, with $\sec_{g_0}\geq-1/60^2$ on $B(p,400)$,
and supply an actual pointed KL $\delta$-map to $\mathbb R$,
where $0<\delta<\delta_\ell$.
For EVERY $q\in B_{g_0}(p,1)$ and EVERY
$\lambda\geq\Lambda_\ell$, the pointed manifold
$(M,\lambda^2g_0,q)$ admits a normalized $(1,\beta)$-splitting.
The AC55 and AC64 qualifications are retained.
\end{lemma}
\begin{proof}
Choose $\sigma=\sigma(3,\beta)$ from AC55, set $L=1/\sigma$,
and choose $\theta$ as in LC65, so
$2L\cos(\theta/2)>c_\sigma$. Require
\[
 \delta_\ell\leq\min\{1/100,3(1-\cos\theta)/44\},\qquad
 \Lambda_\ell\geq\max\{L,(1/60)/\sqrt\sigma,2\}.
\]
If $F$ is the line approximation, the points $F(q)\pm4$ lie
strictly in its coverage region: $|F(q)|<1+\delta$.
Choose lifts $a,b$ with target error less than $2\delta$;
coverage is an infimum and need not be attained. The distortion and coverage bounds give
\[
 4-3\delta<d(q,a),d(q,b)<4+3\delta,\qquad
 d(a,b)>8-5\delta.
\]
Let $D$ be the smaller of the two radial lengths, and shorten
only the longer segment to length $D$. The two resulting endpoints
have mutual distance $c$ with
\[
 3<D<5,\qquad 0\leq E=2D-c<11\delta.
\]
Indeed the loss from shortening is at most the difference of the
original lengths, so the new excess is at most the old excess.
The equal-leg computation of LC65, with $\kappa=1/60$, gives
$1+\cos\omega\leq4E/D<44\delta/3<1-\cos\theta$.

Use AC64 inside $B(q,100)\subset B(p,101)$, with $8D<40<100$,
to shorten both legs to $L/\lambda\leq1<D$.
Exactly as in LC65, the scaled opposite side exceeds
$2L\cos(\theta/2)>c_\sigma$, so the new comparison angle in
curvature $-\sigma$ exceeds $\pi-\sigma$.
The curvature on the normalized $L$-ball is at least $-\sigma$.
AC55 applies. All constants were chosen independently of $q$
and of every larger $\lambda$; in particular the argument covers
points at the original center, where an annular argument would not.
\end{proof}

\begin{theorem}[LC77: the same selected model has at most one end]
\label{thm:collapse-selected-model-one-end}
Use LC64's fixed radii and selected centers, and LC16's zero
stratum $Z$. Attach to each selected center $i$ the SAME
nonnegative model $N_i$, its LC21 cone $C_i$, an actual
$\delta$-cone map from $(M,r(i)^{-2}g,i)$ to $C_i$, and
LC76's curvature bound at that scale.
Choose $\delta<\delta_\ell(\beta_1)$ and
$T\geq20\Lambda_\ell(\beta_1)$ BEFORE the joint witnesses and
selection. Then every selected $N_i$ has at most one end.
\end{theorem}
\begin{proof}
If $N_i$ had more than one end, LC75 would identify its SAME
supplied cone with the line. Compose the actual cone map with
that pointed isometry. For every $q\in B_g(i,r(i))$, LC62
supplies $\lambda=r(i)/\rho(q)\geq T/20$.
LC76 therefore gives a $(1,\beta_1)$-splitting at the
$\rho(q)$ scale. LC16 implies $q\notin Z$.
But LC64 selects only candidate balls that MEET $Z$, a
contradiction. The selected center itself need not lie in $Z$.
This proof does not require LC20's exclusion of rank three,
adapted-coordinate existence, or a classification of compact souls.
\end{proof}

The parameters for LC73 and LC77 must be imposed simultaneously
when the pointwise packages are first produced: take the smaller
cone tolerance and the larger lower scale threshold. The selected
model, cone, radial function and radius then remain fixed.
This completes the end-count application from its supplied geometric
packages; producing those packages with the correct finite regularity
is a separate obligation addressed below.

\clearpage
\section{A written rank-one adapted-coordinate construction}
\label{sec:collapse-rank-one-construction}

The source input LC71 can be supplied by the following argument,
conditional on LC28 and inherited local comparison and first variation.
This expands KL 4.23--4.25 without assuming convergence of actual
angles or using a compactness theorem for smooth metrics.
The first-variation formula used here is KL (3.9), printed 24/PDF 19,
citing BBI Corollary 4.5.7: the right derivative of distance along a
minimizing segment is minus the largest cosine of an angle to a
minimizing direction toward the distance center. It applies at cut
points as a one-sided formula; it does not assert differentiability.

\begin{lemma}[LC78: long product anchors control every tested direction]
\label{lem:collapse-product-anchor-directions}
Fix $R>2$ and $\tau>0$. There are $s>2R+10$ and $\nu_*>0$
such that the following holds. Let $(M,h,q)$ be complete smooth,
with $\sec_h\geq-\nu^2$ on $B(q,\nu^{-1})$, and supply a
normalized $(1,\nu)$-splitting
$\alpha=(\Phi,v):M\to\mathbb R\times Y$, where
$0<\nu<\nu_*$. Choose lifts $a_\pm$ whose images are within
$2\nu$ of $(\pm s,y_0)$. Set
$h_+(x)=d(q,a_+)-d(x,a_+)$. Then
\begin{enumerate}[leftmargin=*]
\item $|h_+(x)-\Phi(x)|<\tau$ for $x\in B(q,R)$;
\item for $x\in B(q,2)$, all minimizing unit directions from $x$
  to $a_+$ have angular diameter less than $\tau$;
\item for $x\in B(q,2)$, $z\in B(q,R)$ with $d(x,z)>1$,
  every minimizing initial unit direction $w$ from $x$ to $z$,
  and every such direction $u_+$ toward $a_+$,
\[
 \left|\langle u_+,w\rangle-
           \frac{\Phi(z)-\Phi(x)}{d(x,z)}\right|<\tau.
\]
\end{enumerate}
The anchors and every comparison triangle lie inside the specified
curvature ball, with the margin required by AC02's local comparison.
\end{lemma}
\begin{proof}
First do the distance calculation in the product. Write
$x=(t,y)$, $z=(t',y')$, and $r=d_Y(y,y_0)$.
On the fixed $R$-ball,
\[
 d(x,(s,y_0))=\sqrt{(s-t)^2+r^2},\qquad
 0\leq d(x,(s,y_0))-(s-t)\leq\frac{r^2}{2(s-R)}.
\]
Thus $s-d(x,(s,y_0))\to t$ uniformly as $s\to\infty$.
For $x\in B(q,2)$ and $z\in B(q,R)$, put
$\ell=d(x,z)>1$. The Euclidean comparison cosine at $x$ for
$(a_+,x,z)$ is
\[
 \frac{d(x,a_+)^2+\ell^2-d(z,a_+)^2}
      {2d(x,a_+)\ell}
 =\frac{t'-t}{\ell}+O_R(s^{-1}).
\]
This follows by expanding the two squares: the leading numerator
is $2s(t'-t)$ and the remaining terms are bounded in terms of $R$.
The lower bounds $\ell\geq1$ and $d(x,a_+)\geq s-2$ make the
estimate uniform. For the negative anchor the leading term changes
sign. The comparison cosine between the two anchors tends uniformly
to $-1$ on $B(q,2)$.

Choose $s$ so these three errors are smaller than a prescribed
$e>0$. Then choose $\nu$ small. Distortion and the lift errors
change every relevant side length by a quantity tending uniformly
to zero. Comparison cosines in curvature $-\nu^2$ converge uniformly
to their Euclidean counterparts for these bounded side lengths;
the two adjacent lengths in every denominator are bounded away
from zero. This also covers limiting degenerate triangles by the
continuous cosine formula. Require additionally
$\nu^{-1}>100(s+R+1)$. The buffered local comparison of AC02
then applies to every choice of the minimizing segments.

Actual angles are at least their comparison angles. Consequently,
for any anchor directions $u_+,u_-$ and tested direction $w$,
\[
 \begin{split}
 \langle u_+,w\rangle&\leq
       \frac{\Phi(z)-\Phi(x)}{d(x,z)}+2e,\\
 \langle u_-,w\rangle&\leq
       -\frac{\Phi(z)-\Phi(x)}{d(x,z)}+2e,\qquad
 \|u_++u_-\|<2\sqrt e.
 \end{split}
\]
The second inequality and the last bound give the reverse bound
for $u_+$. Fixing one $u_-$ also bounds the diameter of all
$u_+$ by $4\sqrt e$. Choose $e$ with
$2e+2\sqrt e<\tau$ and
$4\sqrt e<2\sin(\min\{\tau,1\}/2)$. The latter norm bound
implies angular diameter less than $\tau$. Make the
value error smaller than $\tau$. This proves all three claims.
The constants are independent of the manifold, the residual metric
space and the points in the stated balls.
\end{proof}

\begin{theorem}[LC79: supplying the LC71 contract]
\label{thm:collapse-rank-one-adapted-existence}
Conditional on LC28, inherited first variation and AC02's buffered
Riemannian comparison binding, the conclusion of LC71 holds for
all sufficiently small splitting parameters. No finite-$K$ model
producer is used.
\end{theorem}
\begin{proof}
Fix $0<\gamma<1$. Choose
$0<\varepsilon<\gamma/16$ and obtain LC28's
$\theta(\varepsilon)>0$. Apply LC78 with
$R>\max\{2,\gamma^{-1}\}$ and
$0<\tau<\min\{\theta(\varepsilon),\gamma/16\}$.
Choose its parameters small enough also that
$\nu<\min\{\gamma/64,1/4\}$.
Take $U=B(q,2)$ and $C=\overline B(q,1)$;
$C$ is compact and $U$ misses $a_+$.
LC28 supplies $F$ smooth near $C$ with
$\operatorname{Lip}(F-d_{a_+})\leq\varepsilon$.
Put $\phi(x)=F(q)-F(x)$, so $\phi(q)=0$ and
\[
 \operatorname{Lip}\phi\leq1+\varepsilon,\qquad
 |\phi(x)-h_+(x)|\leq\varepsilon d(q,x).
\]
For a tested minimizing segment from $x$ to $z$ in Definition 4.21,
first variation expresses the right derivative of $h_+$ as the
largest $\langle u_+,w\rangle$. LC78 bounds EVERY such cosine.
The one-sided derivative of the Lipschitz difference, whenever
formed as the difference of these two existing derivatives, has
absolute value at most $\varepsilon$. Hence
\[
 \left|D\phi_x(w)-
       \frac{\Phi(z)-\Phi(x)}{d(x,z)}\right|
       <\tau+\varepsilon<\gamma.
\]
This argument uses no distance gradient at a cut point.

For image coverage lift $(\pm(1-\gamma/4),y_0)$ through
$\alpha$. Decrease $\nu$ further so the two lifts lie in
$B(q,1)$ and their $\Phi$-values differ from the indicated real
numbers by less than $\gamma/16$. The value estimates above
then give one $\phi$-value greater than $1-\gamma/2$ and
one less than $-1+\gamma/2$. The unit ball is path connected
by minimizing paths through $q$, so continuity and the intermediate
value theorem give every intermediate value. Conversely
$|\phi(x)|<1+\varepsilon$ on that ball. Its image therefore has
Hausdorff distance at most $\gamma$ from $(-1,1)$.
These are all three clauses of the SAME-splitting adapted-coordinate
contract. All tolerances were chosen before the manifold and map. Halving
the final strict threshold gives the form $0<\nu\leq\nu_0$.
\end{proof}

\paragraph{Current disposition of LC71.}
Its original source-input record is retained as historical evidence.
LC79 now supplies that record by a written conditional proof, with
LC28, first variation and local comparison still explicit formal
bindings. LC73 may use LC79 as its existence witness. There is no
new Lean declaration and no claim that these inherited ingredients
are absent from PC.

\begin{proposition}[LC80: the zero-stratum local export]
\label{prop:collapse-zero-export}
Supply the joint model/core packages of LC56--LC61, the zero-stratum
selection hypotheses of LC64, and the inherited inputs used above.
Require $0<\varepsilon<1/64$. Choose the smoothing buffer, cone tolerance and lower scale ratio
simultaneously to satisfy LC66, LC73 and LC77 BEFORE constructing
and selecting the witnesses. Then the selected family exports:
\begin{enumerate}[leftmargin=*]
\item finitely many pairwise disjoint open balls, each meeting the
  zero stratum; the open tenth-radius balls cover that stratum;
\item the SAME radius, nonnegative model, cone and radial function
  at each index, with the prescribed core/interior identification;
\item a normalized one-splitting at every point of the CLOSED
  shell $[1/10,10]$, with the original radial coordinate and its
  prescribed smooth adapted version; and
\item at most one end for every attached model.
\end{enumerate}
\end{proposition}
\begin{proof}
LC64 gives finiteness and disjointness. LC66 supplies its annular
exclusion premise. LC61 supplies the interior identification.
LC79 supplies LC71, so LC73 applies to the original radial function.
LC77 gives the end count. Taking the minimum of finitely many
positive error thresholds and the maximum of their scale thresholds
before choosing the joint data preserves each assertion.
\end{proof}
This is a conditional local export, not an acceptance of the
finite-$K$ producer or the full static graph-manifold theorem.
\clearpage
\section{Finite regularity and the other local models}
\label{sec:collapse-local-model-completion}

The remaining local types are circle fibers over two-dimensional
regions, disk fibers near an edge, and sphere or torus fibers in the
slim one-stratum. Their existence and topology belong here.
Making different local maps compatible, changing them on overlaps,
and extracting the global graph presentation belong to
Chapter~\ref{ch:compatible-fibrations}. The following contracts state
what that chapter may consume. A source-qualified producer is an
implementation target, not a temporary axiom or a completed Lean proof.

\begin{foundation}[LC81: finite-order model production and its category gate]
\label{input:collapse-finite-order-models}
Fix $K\geq10$, positive $v,r$, and a bound function $A(R)>0$.
The finite-order compactness input from KL 3.5 concerns complete
pointed $C^{K+2}$ Riemannian $n$-manifolds with
\[
 \Vol B(p,r)\geq v,\qquad
 |\nabla^k\Rm|\leq A(R)\quad\hbox{on }B(p,R),
 \quad 0\leq k\leq K,\ R>0.
\]
Its output has a $C^{K+1}$ atlas and a $C^K$ metric, with pointed
$C^K$ convergence in Definition 3.4. If the same sequence has
expanding-ball lower curvature bounds tending to zero, the model
has nonnegative sectional curvature. KL 6.10 applies this input at
fixed modified volume scale; KL 6.16 additionally tracks a supplied
approximate $j$-splitting and produces a complete nonnegative
$C^K$ factor of dimension $3-j$.

The GC implementation packet must additionally record:
\begin{enumerate}[leftmargin=*]
\item model-to-source embeddings on each required WHOLE buffered
  neighborhood, their image containment and eventual common tail;
\item the SAME metric, basepoint and approximate splitting, including
  the map-comparison estimates needed by the local coordinates;
\item explicit coordinate $C^1$ convergence for gradients and normals,
  and $C^2$ control where curvature is passed to the limit; and
\item a finite-regularity version of the core, normal-flow, isotopy
  and bundle interfaces, or an independently proved upgrade supplying
  LC56's stronger smooth metric package.
\end{enumerate}
These last obligations are not extra conclusions asserted by the bare
text of Definition 3.4. In particular finite $K$ does NOT supply the
all-orders hypotheses of the inspected PC compactness producer.
Smoothing an atlas does NOT smooth its metric or preserve curvature.
The finite-regularity binding remains open; LC56 is a sufficient
smooth package, not a consequence already proved from this input.
\end{foundation}

\paragraph{Regularity budget and implementation order.}
The selected route keeps the finite-order model. With $K\geq10$,
its metric has more derivatives than required for sectional curvature,
first variation, $C^1$ vector fields and their flows. This numerical
observation does not itself port a theorem stated for smooth metrics.
First bind finite-order compactness and actual embeddings; then state
and prove the finite-$C^r$ versions of the normal-flow and perturbation
arguments; finally use smooth approximations of maps on the original
smooth manifold, with rank and isolating margins preserved, for the
smooth bundle output. Each approximation must include a $C^1$ bound
and compact support, not just a value bound. The soul and classification
input must also be in the finite category: KL 3.11 explicitly retains
$C^K$ metrics. No change to the PC foundation or GC environment is made.

\subsection{Finite-category consumers without smoothing the metric}
\label{sec:collapse-finite-consumer-bridge}

This revision starts from master125. The following LFR records are
subordinate developments of LC28 and LC81--LC85; their proofs are
written mathematical arguments, not new Lean declarations. They keep
the original metric, distance and curvature. The inherited analytic
bindings are Euclidean mollification and Rademacher, partitions of
unity, the inverse function theorem and finite-regularity ODEs.
Finite-order compactness, geometric splitting and the required model
classification still have to produce the actual data to which these
arguments apply. In particular an unqualified smooth soul theorem is
not replaced by a statement that $K\geq10$ is large.

\begin{lemma}[LFR01: compatible structures and a derivative ledger]
\label{lem:collapse-finite-derivative-ledger}
Let $N$ be a Hausdorff second-countable $C^{K+1}$ manifold with a
positive definite $C^K$ metric $g$, $K\geq3$. A compatible smooth
structure may be chosen so that the identity in both directions is
$C^{K+1}$. In that structure the SAME tensor is $C^K$, with the same
lengths, distance, completeness and sectional curvature. In smooth
coordinates, $\Gamma$ is $C^{K-1}$, curvature is $C^{K-2}$, and the
geodesic flow and exponential map are at least $C^{K-1}$ on their
open domains. A change through a merely $C^s$ diffeomorphism generally
guarantees only $C^{\min\{K,s-1\}}$ coefficients for the pulled-back
metric. None of these assertions makes $g$ smooth.
\end{lemma}
\begin{proof}
Hirsch, Chapter 2, Theorems 2.9--2.10, supplies the compatible
structure; transporting it by the stated diffeomorphism gives the
identity formulation. The tensor transformation formula uses one
derivative of each coordinate change, so it preserves $C^K$ here.
Length integrals are unchanged by the chain rule, hence so are their
infimum distance and completeness. The Christoffel formula
\[
 \Gamma^k_{ij}=\tfrac12 g^{k\ell}
 (\partial_i g_{j\ell}+\partial_j g_{i\ell}-\partial_\ell g_{ij})
\]
and the curvature formula, involving first derivatives of $\Gamma$
and quadratic contractions of $\Gamma$, give the stated orders and
coordinate-invariant curvature. The geodesic equation is a first-order
system on the tangent bundle with $C^{K-1}$ coefficients, so its local
flow and its restriction defining the exponential map have that order.
The same tensor transformation formula proves the last assertion.
For example normal coordinates need not retain $C^K$ metric
coefficients; DeTurck--Kazdan, Lemma 1.2 and Theorem 2.1, explicitly
warn against that conclusion. Use the chosen compatible charts for
the ledger, not repeated unbudgeted normal-coordinate changes.
\end{proof}

\begin{theorem}[LFR02: localized distance smoothing in the finite category]
\label{thm:collapse-finite-distance-smoothing}
LC28 holds on a smooth carrier with a complete $C^2$ Riemannian
metric. Smoothness of its output function refers to that carrier;
the metric is unchanged. One may take
$\theta(\varepsilon)=\min\{1,\varepsilon\}/100$.
In addition, $F-d_Y$ has compact support in $U$. On its smooth
neighborhood of $C$, for every $v\in V_x(Y)$,
\[
             \|\nabla_gF(x)+v\|_g<\varepsilon.
\]
Here minimizing-geodesic existence and continuity of the geodesic
flow are used for a complete $C^2$ metric; no first-variation formula
at a cut point is needed for this smoothing lemma.
\end{theorem}
\begin{proof}
Put $f=d_Y$ and $a=\min\{1,\varepsilon\}/100$. The sets $V_x(Y)$
are nonempty compact. They are upper semicontinuous in any continuous
tangent trivialization on $M\setminus Y$: if $x_i\to x$ and
$v_i\in V_{x_i}(Y)$, the nearest endpoints stay in one compact
metric ball. Pass to convergent endpoints and directions; continuous
dependence of the geodesic equation and continuity of $d_Y$ show
that the limit segment is again minimizing to $Y$.
At a differentiability point of $f$, its gradient is $-v$ for every
$v\in V_x(Y)$. Indeed its norm is at most one, while its derivative
along each such unit-speed segment is $-1$: the triangle inequality
gives $d_Y(\sigma(t))=d_Y(x)-t$ before its nearest endpoint.
Equality in Cauchy--Schwarz
then gives the assertion. Rademacher applies in the smooth charts.

Use covectors in those charts to avoid any regularity assumption on
parallel transport. At each point $b\in U$, shrink a relatively
compact chart so that the following holds on its closure, with a
slightly larger chart still in $U$: all coordinate covectors
$-g_y(v,\cdot)$, $v\in V_y(Y)$, lie in a single convex bounded
set $A_b$ of diameter less than $4a$, measured in the dual metric
at EVERY point of the smaller chart. To see this, at $b$ their
diameter is less than $a$; upper semicontinuity puts nearby sets
in an arbitrarily small neighborhood of that one, and continuity
of $g^{-1}$ compares the norms. Take its convex hull and a small
open thickening. This construction controls all directions, not a
chosen continuous minimizing arm.

Choose finitely many such charts covering a neighborhood of $C$,
and smooth $\psi_i\geq0$ supported compactly in their smaller
charts, such that $\sum_i\psi_i=1$ near $C$ and
$0\leq\sum_i\psi_i\leq1$ everywhere. The supports and the
neighborhood on which their sum is one are chosen before mollifying.
Convolve $f$ in chart $i$ with a nonnegative smooth Euclidean
kernel of sufficiently small radius, obtaining $f_i$ on a
neighborhood of $\operatorname{supp}\psi_i$. Its differential is
the average of the weak differentials of $f$ in that chart. All
these covectors lie in $A_b$, so, at every such $x$ and every
$v\in V_x(Y)$,
\[
             \|df_i(x)+g_x(v,\cdot)\|_{g_x^*}<4a.
\]
No derivative of the metric occurs in convolution. Choose its radii
also so that $\|f_i-f\|_\infty<e/2$ on the respective supports
and
\[
       \sum_i\|d\psi_i\|_{g^*,\infty}
                      \|f_i-f\|_\infty<a.
\]
These are finitely many choices on compact sets. Define
\[
 F=f+\sum_i\psi_i(f_i-f).
\]
The summands extend by zero, have compact support in $U$, and are
Lipschitz. The function $F$ is smooth where $\sum_i\psi_i=1$;
it equals $f$ off their union and satisfies $\|F-f\|_\infty<e$.
At almost every point, differentiation gives
\[
 d(F-f)=\sum_i\psi_i(df_i-df)+\sum_i(f_i-f)d\psi_i,
 \qquad \|d(F-f)\|_{g^*}<5a<\varepsilon.
\]
Outside the supports it vanishes. A locally Lipschitz function with
this almost-everywhere bound is globally $\varepsilon$-Lipschitz
for the length metric. For completeness, in a sufficiently small
coordinate ball compare $g$ with the constant metric at its center,
apply the Euclidean weak-gradient Lipschitz estimate, and obtain
a $(1+\delta)\varepsilon$ bound along curves in that ball. Subdivide
a piecewise smooth curve into such balls, add the estimates, and
let $\delta\downarrow0$ and its length decrease to the distance.
This avoids differentiating along a curve that might lie in the
ambient exceptional null set. The same computation where
$\sum_i\psi_i=1$, now comparing every $df_i$ with any fixed
$-g_x(v,\cdot)$, gives the final $5a$ bound at every smooth point.
\end{proof}

This fills the omitted localization and partition-derivative argument
behind KL 3.14--3.16 (printed 26--27/PDF 21--22).
Grove--Shiohama, Section 2, Proposition 2.1 and Theorem 2.2
(printed 207--209/PDF 8--10 in the archived copy), obtains controlled
gradients by exponential-kernel averaging on a smooth metric. LFR02
instead uses coordinate covectors and only smooths the function.
The finite-category proof is the argument above, not a claim that
their exponential kernel is $C^\infty$ for a finite-$C^K$ metric.
LC29 and the LC78--LC79 use of the Lipschitz DIFFERENCE now have a
written smoothing producer, with their inherited formal bindings
still pending.

\begin{theorem}[LFR03: finite-regularity buffered fiber stability]
\label{thm:collapse-finite-fiber-isotopy}
In LC82 replace the smooth family by a jointly $C^r$ map
$F:U\times[0,1]\to\mathbb R^k$, $3\leq r<\infty$, defined
$C^r$ up to the time endpoints. Take $U$ to have a chosen smooth
structure and $X$ to be smooth. Retain transversality to BOTH $X$
and $\partial X$ at every time, and the SAME compact set
$Q\Subset U$ containing ALL $f_t^{-1}(X)$ in its interior.
Then a compactly supported ambient $C^{r-1}$ isotopy of $U$
identifies these fibers and their boundaries. Its support can be
chosen in an arbitrarily prescribed relatively compact neighborhood
of their union inside $U$. No metric regularity is used.
\end{theorem}
\begin{proof}
Near an interior point of $X$, write it as $z=0$ in target
coordinates. Near a boundary point write it as
$z=0$, $b\geq0$, with boundary $z=b=0$.
On the inverse image, transversality gives a right inverse for
$D_x(z\circ F)$ in the first case and for
$D_x((z,b)\circ F)$ in the boundary case. On a small neighborhood
choose one nonsingular minor and solve respectively
\[
 D_x(z\circ F)v=-\partial_t(z\circ F),\qquad
 D_x((z,b)\circ F)v=-\partial_t((z,b)\circ F).
\]
The solutions are $C^{r-1}$ ambient time-dependent vector fields.
The second equation is imposed throughout a boundary neighborhood,
not just at the boundary. It both lifts time on $F^{-1}(X)$ and
is tangent to $F^{-1}(\partial X)$.

The trace $W=F^{-1}(X)$ is closed in $Q\times[0,1]$, hence
compact. Choose finitely many of these neighborhoods covering $W$,
with interior-type neighborhoods disjoint from its boundary stratum.
A partition of unity on an ambient neighborhood of $W$ combines
the local fields. On $W$ the time-lifting equation is affine in
$v$, so it is preserved by a convex combination; on the boundary
only the boundary-type fields contribute. Thus $(v,1)$ is tangent
to $W$ and to its side boundary. Multiply by a cutoff equal to one
near $W$ and supported in the prescribed neighborhood, then extend
by zero. Up to time endpoints the construction uses local $C^r$
extensions and a one-sided time interval, without a corner-doubling
theorem. The finite implicit function theorem identifies the local
trace with a half-space times time, so tangency and ODE uniqueness
preserve both strata.

The resulting $C^{r-1}$ field has one compact spatial support.
Its nonautonomous flow exists from every time in $[0,1]$ to every
other time and is jointly $C^{r-1}$; solve backwards for its inverse.
An orbit starting in a fiber stays in $W$, and the reverse flow
proves equality with the next fiber, including the boundary. Outside
the support the flow is the identity. This is the claimed isotopy.
\end{proof}

\begin{lemma}[LFR04: smooth fiber type from a finite diffeomorphism]
\label{lem:collapse-smooth-fiber-type}
Let $A,B$ be compact smooth manifolds, with or without boundary.
If a $C^2$ diffeomorphism $h:A\to B$ of manifolds with boundary
is supplied, then they are smoothly diffeomorphic. This assertion
does not identify $h$ itself with the smooth diffeomorphism, and
provides no curvature or comparison estimate for a replacement map.
\end{lemma}
\begin{proof}
For empty boundary use $C^1$ approximation and openness of the
diffeomorphisms (Hirsch, Chapter 2, Theorems 2.6--2.7).
Here is the additional boundary argument needed in this application.
Choose smooth collars of $A$ and $B$. On a sufficiently short
compact collar of $\partial A$ write
$h(x,t)=(b(x,t),c(x,t))$, where $c(x,0)=0$ and
$\partial_tc(x,0)>0$. Factor
\[
 c(x,t)=t a(x,t),\qquad
 a(x,t)=\int_0^1\partial_tc(x,ut)\,du.
\]
Then $a$ is $C^1$ and has a positive lower bound after shortening
the collar. Smoothly approximate $a$ in $C^1$, retaining positivity.
Approximate $b$ in $C^1$ by a smooth map to $\partial B$:
extend across the end of the parameter interval in local charts,
embed $\partial B$ in Euclidean space, mollify and project through
its smooth tubular neighborhood. The map
$(\widetilde b(x,t),t\widetilde a(x,t))$ is smooth near the boundary,
preserves it, and is arbitrarily $C^1$ close to $h$ there.
On a fixed smaller collar keep this map; across an outer collar
annulus interpolate to $h$ in a tubular neighborhood of a smooth
embedding of $B$. That annulus lies in the interior of $B$, so
sufficiently close interpolation stays there. The new map is $C^1$
close to $h$ and smooth near the boundary. Relative approximation
on the remaining interior, fixed near that smaller collar, makes
it smooth everywhere.

Sufficient $C^1$ closeness gives invertible differentials and an
inward positive normal derivative, hence a local diffeomorphism of
half-spaces at the boundary. Injectivity follows on close pairs
from finitely many inverse-function charts and on separated pairs
from compactness and injectivity of $h$. The image is both open
and closed; closeness makes it meet every target component. Thus
the smooth map is bijective and its local smooth inverses agree.
This proves the assertion including the boundary, without treating
unrestricted approximation as boundary-preserving. It implements
the needed case of Hirsch, Chapter 2, Section 3, rather than using
the problematic intermediate convolution argument noted by Granja.
\end{proof}

\begin{proposition}[LFR05: what the finite model transfers to a smooth packet]
\label{prop:collapse-finite-packet-consumer}
Supply an actual $C^r$ model-to-source diffeomorphism onto its open
image, $r\geq3$, on the WHOLE buffered model domain. Supply a
jointly $C^r$ interpolation between the model map and the pullback
of the smooth source map, satisfying LFR03 on that domain. If its
model fiber is $C^{r-1}$ diffeomorphic to a specified compact smooth
fiber $H$ (allowing boundary), and the ENTIRE relevant source fiber
is contained in the image of this embedding, then the source fiber is smoothly
diffeomorphic to $H$. If the smooth source map has a proper
submersion restriction, also on its boundary when present, that
restriction is a smooth bundle with fiber $H$. It is trivial over
a ball or interval. The embeddings used to compare metrics and
gradients remain the original supplied embeddings.
\end{proposition}
\begin{proof}
Apply LFR03 on the model domain and then the actual embedding to
the terminal fiber. The source enclosure ensures that this image is
the entire source fiber, not merely its part in the comparison chart.
Composition gives a $C^{r-1}$ diffeomorphism
from $H$ to the source fiber, including its boundary. Since
$r-1\geq2$, LFR04 supplies its smooth type. The proper-submersion
theorem, with a boundary-tangent lift when needed, gives the bundle;
its trivialization over a ball or interval is an inherited binding.
The isotope and the final type-identifying diffeomorphism are not
used as metric convergence maps.
\end{proof}

\paragraph{Core and flow consumers.}
LC45--LC47 likewise need no smooth metric once their actual data
are supplied in finite regularity. For a $C^r$ bundle quadratic
function and $C^r$ normal-flow diffeomorphism $e$, $r\geq3$, the
proof of LC45 gives $C^r$ disk cores and the same radial identities.
Compactness is topological and regularity of a nonzero level follows
from the nonvanishing fiber derivative $2\|v\|^2$. In LC46, for
$\eta,Y$ of class $C^r$, the normalized field $Y/d\eta(Y)$ is
$C^{r-1}$; the displayed flow and inverse therefore give a
$C^{r-1}$ band product. In that product LC47's implicit crossing
height is $C^{r-1}$. Its vertical field and its flow are
$C^{r-1}$, and conjugating the resulting maps by the band product
gives the same $C^{r-1}$ ambient isotopy. One need not differentiate
the pushed-forward field and silently spend another derivative.
These are finite versions of the written consumer proofs, conditional
on the supplied soul, normal tube, same flow and outward inequalities.
They do not construct that geometric input. LFR04 applies to compact
smooth end fibers after the resulting finite diffeomorphism is known.

\paragraph{Remaining LC81 producer obligations.}
The bridge now has proved consumer steps, not a smooth-metric upgrade.
The producer must still deliver $C^r$ maps and fibers with $r\geq3$
on the actual common buffers, coordinate convergence, splitting
factors and the required finite-category core/classification data.
Record derivative losses before invoking LFR03 or LFR05; merely
asserting a $C^K$ metric is not that record. In particular KL 3.11
(printed 24--25/PDF 19--20) explicitly states finite-$C^K$
classification, but its noncompact proof refers to Cheeger--Gromoll.
The latter's Theorem 2.2 (printed 423--428/PDF 12--17) announces
a diffeomorphism and then gives a homeomorphism construction, leaving
the smoothing work implicit at printed 424/PDF 13. This is not
an inspected finite-regularity normal-flow construction. The existing
LC54 route, with its actual soul/tube and outward-field inputs, is
the appropriate place to finish that obligation. Compact
three-model classification also remains separate from the surface
fiber classification used by LC84--LC85. Reuse the frozen PC CGT
injectivity estimate for the original smooth manifolds; it does not
by itself supply finite-order compactness or these limit-model data.

\begin{lemma}[LC82: buffered stability of a compact fiber with boundary]
\label{lem:collapse-buffered-fiber-stability}
Let $f_t:U\to\mathbb R^k$, $0\leq t\leq1$, be a smooth family
on an open manifold domain, and let $X\subset\mathbb R^k$ be a
closed embedded submanifold with boundary. Suppose every $f_t$ is
transverse to both $X$ and $\partial X$, and that one compact
$Q\subset U$ contains ALL $f_t^{-1}(X)$ in its interior.
Then $f_0^{-1}(X)$ and $f_1^{-1}(X)$ are ambiently isotopic in $U$,
including their boundaries, conditional on the inherited submersion
and compactly supported flow bindings.
\end{lemma}
\begin{proof}
For $F(x,t)=f_t(x)$, the set $W=F^{-1}(X)$ is closed in
$Q\times[0,1]$, hence compact. Transversality on the interior
and boundary gives a manifold over $[0,1]$, with corners where
the fiber boundary meets a time endpoint.
Its projection to time is a submersion on both strata: at any point,
solve $Df_t(v)+\partial_t f_t\in T X$, and on the boundary solve
it in $T\partial X$. The respective transversality hypotheses are
exactly the solvability conditions. Local solutions can be combined
by a partition of unity; first prescribe a solution tangent to the
boundary and extend it over a collar. The resulting time-lifting
field is tangent to that boundary. Compactness gives its flow for
all $t\in[0,1]$, and a compactly supported extension in $U$
gives an ambient isotopy. This is the proof of KL 21.1 with the
compact isolating set made explicit.
\end{proof}
For a straight-line perturbation family, quantitative $C^1$ bounds
on $Q$ preserve transversality, but a SEPARATE image-separation
condition must keep every inverse image inside $\operatorname{int}Q$.
Control on $Q$ alone cannot prevent a new component outside $Q$.
This is the safe contract for using KL 21.3 and is the same enclosure
principle already used in LC37--LC43.

\begin{proposition}[LC83: the circle-fiber local model]
\label{prop:collapse-two-stratum-local}
The two-stratum producer must supply, at scale $\rho(p)$, a
smooth map $\eta:B(p,200)\to\mathbb R^2$, $\eta(p)=0$, with
rank two, Lipschitz constant at most two, and
\[
 \eta^{-1}(B(0,100))\subset B(p,102),\qquad
 \eta^{-1}(0)\subset B(p,2).
\]
Distances here use $\rho(p)^{-2}g$, on a complete connected smooth
Riemannian three-manifold.
Then $\eta:\eta^{-1}(B(0,100))\to B(0,100)$ is a proper
surjective submersion with connected circle fibers and is a trivial
bundle. In particular the inverse image of $B(0,R)$ is
$S^1\times B(0,R)$ for $0<R<100$, conditional on the inherited
proper-submersion and one-manifold classification bindings.
\end{proposition}
\begin{proof}
The inverse image of a compact subset of $B(0,100)$ is closed
in the compact $\overline B(p,102)$, strictly inside the map
domain; thus it is compact. A submersion has open image, and a
proper map has closed image. The image contains zero, so it is the
whole connected base. The proper-submersion theorem makes this a
bundle, trivial over the contractible ball, with compact fibers.
Any two points of the zero fiber can be joined through $p$ by
minimizing paths in $B(p,2)$. The Lipschitz bound keeps their
images in $B(0,4)$, so this path lies in the bundle domain.
In a trivial bundle such a path connects points in the same fiber
component. Thus the zero fiber, and every fiber, is a connected
compact one-manifold without boundary, hence a circle.
\end{proof}

\paragraph{Producing the two enclosures.}
KL 8.1--8.4, printed 47/PDF 42, use a small residual factor and
adapted coordinates. The assertion that EVERY arbitrary target
factor has small GLOBAL diameter is too strong: unused remote
points can be added beyond the approximation's tested ball, as
recorded in revision82. The needed replacement is only this:
for the ACTUAL map on $B(p,200)$, the residual coordinate lies
within distance one of its marked point, and
$|\eta-\Phi|<1/10$, with distortion less than $1/10$.
Pythagoras then gives both enclosures in LC83.
This local residual bound follows by the AC41/AC51 contradiction
argument: if it failed as splitting error and collapsed-model error
tend to zero, the actual factor maps on the fixed 201-ball would
have a nontrivial limiting factor. AC48 forces the factor of a
rank-two splitting in a limit of dimension at most two to be a point.
A violating sequence in that fixed ball has convergent images in
AC51's construction, a contradiction. No conclusion about remote
parts of the original factor is drawn. Revision126 supplies rank-two adapted-coordinate existence, its value
estimate and rank test in LFR06 below. LC79 is retained as the
rank-one proof; LFR07 gives the circle-packet application.
The cutoff $\Phi_{8,9}(|\eta|)$ has support strictly inside the
map domain, so its zero extension is smooth.

\subsection{Rank-two coordinates and the circle packet}
\label{sec:collapse-rank-two-construction}

\begin{theorem}[LFR06: actual rank-two adapted coordinates]
\label{thm:collapse-rank-two-adapted-existence}
Fix $0<\gamma<1/10$ and $R>\max\{2,\gamma^{-1}\}$.
For sufficiently small $\nu>0$, a normalized $(2,\nu)$-splitting
$\alpha=(\Phi,v):M\to\mathbb R^2\times Y$ of a complete
smooth pointed Riemannian manifold $(M,g,q)$, with
$\sec_g\geq-\nu^2$ on $B(q,\nu^{-1})$, has smooth
coordinates $\eta:B(q,1)\to\mathbb R^2$ with $\eta(q)=0$,
\[
 \operatorname{rank}D\eta=2,\qquad
 \operatorname{Lip}\eta\leq1+\gamma,\qquad
 |\eta-\Phi|<\gamma\quad\hbox{on }B(q,1).
\]
Their image has Hausdorff distance less than $\gamma$ from
$B(0,1)$. For $x\in B(q,1)$, $z\in B(q,R)$,
$d(x,z)>1$, and EVERY minimizing initial unit direction $w$ to $z$,
\[
 \left|D\eta_x(w)-
            \frac{\Phi(z)-\Phi(x)}{d(x,z)}\right|<\gamma.
\]
The parameters depend on $R,\gamma$, not on $M$ or the residual
factor. The geometric inputs are LC78's buffered local comparison
and inherited first variation. LFR02 supplies the smoothing.
\end{theorem}
\begin{proof}
Choose long product anchors $a_{j\pm}$ lifting
$(\pm s e_j,y_0)$, $j=1,2$, with error $2\nu$.
Let $h_j(x)=d(q,a_{j+})-d(x,a_{j+})$. Each real coordinate of
the splitting is a rank-one product projection, with the other
real coordinate included in its residual factor. The calculation
in LC78 thus gives, on fixed buffered balls, arbitrarily small
value error, arbitrarily small diameter of all directions to each
$a_{j+}$, and its scalar tested-direction estimate, by first
increasing $s$ and then decreasing $\nu$. These statements can
be imposed simultaneously for both coordinates. Use a larger fixed
ball, for example radius $4R$, in that calculation and arrange
smoothing on a neighborhood of $\overline B(q,3)$.

There is one additional estimate: independently smoothing two
one-dimensional coordinates would only give a crude $\sqrt2$
Lipschitz bound. We must control the Gram matrix. The same product
distance calculation gives comparison cosines tending uniformly to
zero between anchors with different indices, and to $-1$ between
opposite anchors of the same index, at every $x\in B(q,3)$.
All adjacent lengths are $s+O_R(1)$ and the mixed opposite side
is $\sqrt2s+O(\nu)$. Apply the buffered lower-curvature comparison
to all four sign choices. For every choice of minimizing directions
$u_{j\pm}$ it gives
\[
 \langle u_{1+},u_{2+}\rangle\leq o(1),\qquad
 \langle u_{1+},u_{2-}\rangle\leq o(1),\qquad
 \|u_{2+}+u_{2-}\|=o(1).
\]
The last two inequalities give the reverse bound in the first.
Thus $|\langle u_{1+},u_{2+}\rangle|=o(1)$ uniformly, in the
ordered limit $s\to\infty$ followed by $\nu\to0$.

Apply LFR02 to each distance function, taking a Lipschitz-difference
tolerance $\epsilon$ and direction threshold sufficiently small
before selecting $s,\nu$. Set
$\eta_j(x)=F_j(q)-F_j(x)$. At common differentiability points
of the two distances,
\[
 \|\nabla\eta_j-u_{j+}\|\leq\epsilon.
\]
Consequently the Gram matrix
$G=(\langle\nabla\eta_i,\nabla\eta_j\rangle)_{i,j=1}^2$
is uniformly as close to the identity as desired. The distances
are simultaneously differentiable on a full-measure, hence dense,
set; continuity extends the estimate throughout $B(q,3)$.
Choose the errors so that
\[
          \|G-I\|_{\mathrm{op}}<\gamma/4.
\]
Both eigenvalues are positive, proving rank two, and
$\|D\eta\|\leq\sqrt{1+\gamma/4}<1+\gamma/4$.
The asserted Lipschitz bound on $B(q,1)$ uses this larger smooth
domain: every minimizing segment between two of its points stays
in $B(q,3)$, where the differential estimate holds. It is not an
inference from a derivative estimate on a possibly nonconvex unit
ball alone.

The Lipschitz DIFFERENCE bounds give
$|\eta_j-h_j|\leq\epsilon d(q,\cdot)$. Combine them with the
product value estimate and choose errors to obtain
$|\eta-\Phi|<\gamma/8$ on $B(q,1)$. The componentwise
LC78 and first-variation argument used in LC79 gives the tested
vector estimate after multiplying the scalar error by $\sqrt2$.
For image coverage take any $y\in B(0,1)$ and lift
$((1-\gamma/4)y,y_0)$ through the SAME splitting. The lift lies
in $B(q,1)$ for $3\nu<\gamma/8$, and its $\eta$ image is
within $\gamma/4+2\nu+\gamma/8<\gamma$ of $y$.
Conversely $\eta(q)=0$ and the Lipschitz estimate put every image
point within $\gamma$ of the unit ball. This is Hausdorff coverage;
no surjectivity or degree theorem has been smuggled into it.
\end{proof}

\begin{corollary}[LFR07: a supplied circle packet at a two-stratum point]
\label{cor:collapse-rank-two-circle-packet}
Retain LC83's actual two-stratum residual-factor estimate on the
tested $201$-ball, obtained from AC41/AC48/AC51 and the collapsed
model comparison. Choose the SAME splitting sufficiently accurate
for LFR06 after fixed rescaling. Then the rank-two coordinate
producer, the two LC83 enclosures, the proper trivial circle
bundle over $B(0,100)$ and its smooth cutoff all exist. This
conditional local packet no longer requires an unexpanded
rank-two adapted-coordinate input.
\end{corollary}
\begin{proof}
Apply LFR06 after rescaling the normalized source by $201^{-1}$,
and multiply its output coordinates by $201$ on returning to the
original scale. Fixed-rescaling and restriction of the actual
approximation use the buffered coverage bookkeeping of MC10--MC11;
choose the original error after this fixed scale and the desired
coordinate error have been fixed. Arrange
$|\eta-\Phi|<1/10$ on $B(q,200)$, distortion less than $1/10$,
rank two and Lipschitz constant at most two. The actual residual
coordinate there has distance less than one from $y_0$.
For $|\eta(x)|<100$, the pointed distortion and product metric give
\[
 d(q,x)<\sqrt{(100+1/10)^2+1}+1/10<102.
\]
For $\eta(x)=0$ the same estimate is
$d(q,x)<\sqrt{(1/10)^2+1}+1/10<2$.
LC83 now proves properness, surjectivity, connected circle fibers
and triviality. The cutoff $\Phi_{8,9}(|\eta|)$ equals one near
$\eta=0$ and vanishes on a neighborhood of the boundary of the
map domain: its closed support lies in the compact enclosure of
$|\eta|\leq9$. Its extension by zero is smooth, including at
$\eta=0$ since the profile is constant there.
\end{proof}

The remaining geometric input in LFR07 is precisely the tested
residual-factor estimate and its collapsed-model hypotheses, not
small diameter of an arbitrary entire target factor. The circle
construction takes place on the original smooth manifold and needs
no finite-order nonnegative limit metric. Edge and slim packets
still need their endpoint/bounded-factor producers, the same-metric
finite model embeddings, model disk or compact-surface identification,
and interpolation estimates on one isolating domain. LFR02--LFR05
remove the corresponding smoothing and finite-isotopy consumer
obstructions when those data are supplied; they do not assert the
missing geometric estimates. No global compatibility or graph
assembly is claimed here.

\subsection{Actual embeddings and regularity of an exact product}
\label{sec:collapse-finite-embeddings-products}

Revision127 fills two producer clauses that were left separate in LFR05:
injectivity and image containment of the ACTUAL comparison maps, and
finite regularity of an exact metric splitting. It does not assume a
differentiable structure on the metric factor. The compactness theorem
must still produce a limit and the pointed comparison maps. All metric
and curvature estimates below concern the unchanged metrics.

\begin{lemma}[LFR08: propagating the basepoint volume lower bound]
\label{lem:collapse-local-injectivity-from-basepoint}
Fix $n\geq2$, $r,v,S>0$ and $A>0$. Let $(M,g,p)$ be a complete
smooth pointed $n$-manifold with
\[
 \operatorname{Vol}_g B(p,r)\geq v,\qquad
 |\operatorname{Rm}|_g\leq A
       \quad\hbox{on }B(p,2S+r+2).
\]
There is $\iota=\iota(n,r,v,S,A)>0$ such that
$\operatorname{inj}_g(x)\geq\iota$ for every $x\in B(p,S)$.
In particular the same constant works for a family with these data;
no global curvature or global injectivity bound is required.
This uses the inherited LOCAL Bishop--Gromov and CGT estimates.
\end{lemma}
\begin{proof}
Put $D=S+r+1$, $\kappa=\max\{1,A\}$ and let
$V_{-\kappa}(t)$ denote the volume of a radius-$t$ ball in the
simply connected constant-curvature $-\kappa$ $n$-space. For
$x\in B(p,S)$, the inclusions
\[
 B(p,r)\subset B(x,D)\subset B(p,2S+r+1)
\]
and $\operatorname{Ric}\geq-(n-1)\kappa g$ on the latter ball
give, for $0<b\leq1$,
\[
 \operatorname{Vol} B(x,b)
       \geq v\,\frac{V_{-\kappa}(b)}{V_{-\kappa}(D)}=:v_b>0.
\]
Only the centered comparison up to radius $D$ is used; all its
radial minimizing segments stay in the displayed curvature buffer.
Choose a common Jacobi radius $a\leq1$ with
$a\leq\pi/\sqrt\kappa$, small enough that the exponential map
is a local diffeomorphism on its tangent $a$-ball under
$|\operatorname{Rm}|\leq\kappa$ on $B(x,1)$. The local Jacobi
estimate makes this choice depend only on $n,\kappa$, not $x$ or
$M$. Apply the retained local-volume CGT estimate with its parameters
$R=a$, $r_0=s=a/8$ and $q=\sqrt\kappa$. Its strict inequalities
$r_0+2s<R$ and $r_0<R/4$ hold. Completeness supplies compact
closed $s$-balls, and all curvature and Ricci hypotheses lie in
$B(x,1)$. It gives
\[
 \operatorname{inj}_g(x)\geq
 \frac{a}{16}\,
 \frac{v_{a/8}}{V_{-\kappa}(a/8)+V_{-\kappa}(a/4)}
 =:\iota>0.
\]
The denominator is the ordinary and exponential-pullback volume
comparison denominator; the latter does not presume injectivity.
This is exactly the role of CGT here, on the ORIGINAL smooth metric.
\end{proof}

The retained PC source names are
\code{exists_uniform_local_cgt_radius} and
\code{intrInj_ge_vol_of_local_ball}, in
\code{LocalCGTInjectivity.lean}, lines 89 and 123 of the frozen
revision87 copy. The first proof chooses its radius by the
dimension-and-curvature Jacobi estimate; its bare existential
statement alone would not establish uniformity across varying metrics.
The exact source commit and earlier axiom census are retained in the
revision127 source record. This is reuse of that discovery, not a new
PC probe, an assertion of release acceptance, or a new Lean run.

\begin{lemma}[LFR09: uniform short joining geodesics]
\label{lem:collapse-uniform-short-geodesics}
Let $C\subset W$ be compact inside an open domain of a smooth
manifold, and let positive definite $C^2$ metrics $h_i$ converge
to a $C^2$ metric $g$ in $C^2$ on compact subsets of $W$.
Distance $d_g$ is taken in a fixed ambient manifold carrying $g$.
There are $\tau,L>0$ and a tail of the sequence such that, for
$x,y\in C$ with $d_g(x,y)<\tau$, an $h_i$-geodesic from $x$
to $y$ stays in $W$ and has $h_i$-length at most $L d_g(x,y)$.
Completeness of $h_i$ is not assumed.
\end{lemma}
\begin{proof}
Cover $C$ by finitely many nested coordinate balls with their larger
closures in $W$. Shrink the smaller balls so that every sufficiently
$d_g$-close pair with first point there lies in the corresponding
larger ball, and coordinate difference is at most a fixed constant
times $d_g(x,y)$. For the last assertion, a curve of sufficiently
small $g$-length cannot exit the larger ball: on its compact closure
the metric bounds coordinate speed below, and the smaller ball has a
positive coordinate margin. Apply this to curves with lengths tending
to the ambient distance. Thus no convexity of $C$ is needed.

In each chart the Christoffel symbols of $h_i$ converge in $C^1$.
The geodesic ODE and its first variations therefore converge uniformly
on a common small set of initial points and velocities, up to time
one, to those for $g$. There is a common velocity radius $b>0$
such that these geodesics stay in the larger chart and, in coordinate
velocities,
\[
 \left\|D_w\bigl(\exp_x^{h_i}(w)-x\bigr)-I\right\|<1/2
                         \qquad (|w|\leq b).
\]
This follows by continuity at $w=0$, where the derivative is $I$,
compactness of the initial-point set, and $C^1$ dependence of the
flows. If $|y-x|<b/2$, the map
\[
 w\longmapsto y-x-\bigl(\exp_x^{h_i}(w)-x-w\bigr)
\]
maps the closed $b$-ball to itself and is a contraction. Its fixed
point satisfies $\exp_x^{h_i}(w)=y$ and $|w|\leq2|y-x|$.
Uniform metric upper bounds at $x$ bound the length of that geodesic
by a fixed multiple of $|y-x|$. Take the minimum admissible distance
threshold and maximum length constant over the finite chart family.
\end{proof}

\begin{proposition}[LFR10: the comparison maps are buffered embeddings]
\label{prop:collapse-actual-buffered-embeddings}
Let $(N,g,p)$ be a complete connected $C^K$ Riemannian $n$-manifold
on a compatible smooth carrier, $K\geq3$. Let $(M_i,g_i,p_i)$ be
complete smooth pointed $n$-manifolds with a common basepoint volume
lower bound and uniform full-curvature bounds on every fixed pointed
ball, as in LFR08. Supply actual $C^{K+1}$ maps
$f_i:U_i\to M_i$, where every compact subset of $N$ is eventually
contained in $U_i$, such that $f_i(p)=p_i$ and, on every fixed ball,
\[
 f_i^*g_i\longrightarrow g\text{ in }C^2,\qquad
 \sup_{x,y}\left|d_{g_i}(f_i(x),f_i(y))-d_g(x,y)\right|
                                      \longrightarrow0.
\]
Convergence is on neighborhoods of the closed balls; derivatives
are measured in the fixed compatible charts on $N$. For every
$0<r<R<\infty$, eventually the SAME map restricts to a
$C^{K+1}$ embedding of $B_N(p,R)$ and
\[
                   B_{M_i}(p_i,r)\subset f_i(B_N(p,R)).
\]
It is an embedding on a neighborhood of the closed $R$-ball if one
starts with a larger buffer. A finite collection of radii has one
common tail. No dense-image premise is needed for this conclusion.
\end{proposition}
\begin{proof}
Work first on the closed $(R+1)$-ball and its open neighborhood.
Completeness makes this a compact subset of $N$.
The tensors $h_i=f_i^*g_i$ are eventually positive definite there.
Consequently $Df_i$ has rank $n$, so $f_i$ is a local
$C^{K+1}$ diffeomorphism and a local isometry for $h_i$.
By LFR09 there is a uniform short joining geodesic for any two
sufficiently close points of the closed $R$-ball, in this buffer.
LFR08 gives a common injectivity lower bound $\iota$ on
$B_{M_i}(p_i,R+1)$.

Write $\epsilon_i\to0$ for the distortion bound on the buffer.
If $x,y\in B_N(p,R)$ have $f_i(x)=f_i(y)$, then
$d_g(x,y)\leq\epsilon_i$. Their LFR09 geodesic has length at
most $L\epsilon_i<\iota$ for large $i$. Its image under the
local isometry is a geodesic from $f_i(x)$ to itself, with that
same length. The starting image lies in $B(p_i,R+1)$ by the
pointed distortion bound. Injectivity of its exponential map on
the tangent $\iota$-ball makes the initial velocity zero. Hence
$x=y$. An injective local diffeomorphism is a diffeomorphism onto
its open image, proving the embedding assertion.

Let $q\in B(p_i,r)$ and lift a minimizing geodesic from $p_i$
to $q$, beginning at $p$, using the local inverses on $B_N(p,R)$.
At every time when this lift $x(t)$ is defined,
\[
 d_g(p,x(t))\leq d_{g_i}(p_i,f_i(x(t)))+\epsilon_i
                                       \leq r+\epsilon_i<R.
\]
For a fixed tail choose $r'$ with $r+\epsilon_i<r'<R$.
The lift stays in the compact closed $r'$-ball. At a putative last
time, choose a convergent subsequence of lifted points. Its limit
is in the domain and has a local inverse neighborhood; on a terminal
time interval this inverse extends the lift. Uniqueness follows
from the already proved injectivity. Thus lifting reaches $q$,
which proves image containment. Repeat the argument with a larger
radius for the closed-ball assertion, and take a maximum of finitely
many sequence thresholds for finitely many buffers.
\end{proof}

This supplies the two missing domain properties of a KL3.4-type
comparison map, provided that compactness has produced that map in
the LIMIT-TO-SOURCE direction. It does not infer $C^1$ convergence
from GH approximation, invert a map without injectivity, or replace
the original comparison map by a new fiber-type map. Petersen's
author manuscript, Section11.3.2, already builds embeddings and
source-ball containment into its convergence definition. Passing
from KL's definition to that stronger definition needs precisely
the additional argument above (or compactness directly in the
stronger convention). A fiber enclosed in $B(p_i,r)$ is now wholly
in the image; model-side compactness alone still does not establish
that source enclosure.

\begin{theorem}[LFR11: finite regularity of an exact metric splitting]
\label{thm:collapse-exact-splitting-regularity}
Let $(N,g)$ be a complete connected Riemannian $n$-manifold with
$C^K$ metric on a compatible smooth carrier, $K\geq3$. Supply
a surjective distance isometry
\[
 I:(N,d_g)\longrightarrow\mathbb R^j\times Y
\]
with the Euclidean $\ell^2$ product metric, $0\leq j\leq n$;
no manifold structure on $Y$ is assumed. Put $t=\pi_{\mathbb R^j}I$.
Then $t$ is $C^{K+1}$ and $Z=t^{-1}(0)$ is a $C^{K+1}$
embedded submanifold with a complete induced $C^K$ metric $h$.
The map $y\mapsto I^{-1}(0,y)$ is an isometry from $Y$ onto
$(Z,d_h)$, and the actual product map
\[
 \Psi:\mathbb R^j\times Z\longrightarrow N,\qquad
 \Psi(u,z)=I^{-1}(u,\pi_Y I(z))
\]
is a $C^{K+1}$ diffeomorphism with
$\Psi^*g=du^2+h$. If $\sec_g\geq0$, then $\sec_h\geq0$.
An orientation of $N$, together with the ordered Euclidean factor,
orients $Z$. Compatible smooth structures may be chosen as in
LFR01; neither metric is thereby made smooth.
\end{theorem}
\begin{proof}
There is nothing to prove for $j=0$, so take $j>0$.
First obtain differentiability without any coordinates on $Y$.
For $p\in N$, write $I(p)=(u_0,y_0)$ and choose $s>0$ small.
The points $a_\pm=I^{-1}(u_0\pm s e_a,y_0)$ lie in a common
convex normal neighborhood of $p$. For $x$ near $p$, the product
distance identity gives
\[
 t_a(x)-t_a(p)
             =\frac{d_g(x,a_-)^2-d_g(x,a_+)^2}{4s}.
\]
The local exponential maps and their inverses are $C^{K-1}$ by
the geodesic ODE and inverse function theorem. Squared distance
from each fixed $a_\pm$ is the metric norm squared of that local
logarithm, hence is $C^{K-1}$. Thus $t_a$ is $C^{K-1}$,
in particular $C^2$. Each $t_a$ is also $1$-Lipschitz.

The lift under $I^{-1}$ of each coordinate line through $I(p)$
is a unit-speed distance-minimizing Riemannian geodesic. Along it
$dt_b(c_a')=\delta_{ab}$. The Lipschitz bound and equality in
Cauchy--Schwarz imply $X_a:=\nabla t_a=c_a'$ at $p$.
Consequently the $X_a$ are orthonormal everywhere and $t$ has
rank $j$. Moreover $t_a$ is affine along every locally minimizing
geodesic in $N$. Indeed equality in the triangle inequality for
a product geodesic, followed by equality in the Euclidean norm
triangle inequality for its two component distances, makes its
Euclidean projection affine with the same parameter. Applying this
on short geodesic intervals proves the assertion for every geodesic.
It follows that $\operatorname{Hess}t_a(v,v)=0$ for every $v$.
The symmetric Hessian and polarization give $\nabla X_a=0$.

Initially $X_a$ is $C^{K-2}$. In the chosen smooth coordinates
the parallel equation reads
\[
                 \partial_b X_a^c=-\Gamma^c_{bd}X_a^d.
\]
Since $\Gamma$ is $C^{K-1}$, this equation successively raises
$X_a$ to $C^{K-1}$ and then to $C^K$; here $K\geq3$ ensures
the initial equation is classical. Now $dt_a=g(X_a,\cdot)$ is
$C^K$, so $t_a$ is $C^{K+1}$. The implicit function theorem
gives the asserted structure on $Z$ and $C^K$ induced metric.
Its second fundamental form vanishes because all normal fields
$X_a$ are parallel. Thus it is totally geodesic and its sectional
curvature equals the ambient curvature on its tangent two-planes.

For two points of $Z$, a minimizing ambient geodesic exists by
completeness. Its $t$ coordinate is affine with both endpoints
zero, so the geodesic stays in $Z$. Therefore its intrinsic
distance equals ambient distance, which by the product formula
is the distance of the corresponding points of $Y$. This proves
the stated factor isometry, connectedness of $Z$, and completeness
of $h$ since $Z$ is closed in complete $N$. In zero dimension
these assertions say that $Z$ is a single point. In positive
dimension the ordered normal frame $X_1,\ldots,X_j$ and the
orientation of $N$ give an orientation of $Z$.

The unit parallel fields are complete and commute. Their $C^K$
flows act by the exact coordinate translations, as their integral
curves are the lifted coordinate lines. Composing these flows
from $z\in Z$ gives the displayed $\Psi$ and its inverse
$x\mapsto(t(x),\operatorname{Fl}_{-t(x)}(x))$, both initially
$C^K$. The flows preserve $g$ since the fields are parallel,
and their coordinate directions are orthonormal and perpendicular
to $Z$. Hence $\Psi^*g=du^2+h$ pointwise.

Finally use the isometry identity to recover the remaining derivative.
Choose compatible smooth charts on $Z$. For the already $C^2$ map
$\Psi$, the Levi--Civita connection transformation formula is
\[
 \partial_i\partial_j\Psi^a
   =\Gamma_{\mathrm{prod},ij}^{c}\partial_c\Psi^a
    -\Gamma_{N,bc}^{a}(\Psi)
                         \partial_i\Psi^b\partial_j\Psi^c.
\]
Both metrics are $C^K$, and $\Psi$ is $C^K$, so the right-hand
side is $C^{K-1}$. Thus $\Psi$ is $C^{K+1}$; its inverse has
the same class by the inverse function theorem. This is the
connection bootstrap in Matveev--Troyanov2.2/(2.4), with the
initial $C^2$ regularity proved here by squared distances and
parallel fields. No general low-regularity isometry theorem or
preexisting differentiable structure on $Y$ has been assumed.
\end{proof}

\paragraph{The budget for the LC84--LC85 fiber consumers.}
The exact splitting is still an input; LFR11 is not an Alexandrov
or Riemannian splitting existence theorem. Once the finite-$C^K$
limit, the actual $C^{K+1}$ convergence maps and that metric
splitting have been produced, LFR10--LFR11 give a complete
$C^K$ surface factor in the rank-one three-dimensional case,
with actual $C^{K+1}$ product and buffered comparison maps.
Composing them spends no derivatives. A model function supplied
in $C^{K-1}$, together with the pullback of a smooth source
function, therefore has a jointly $C^{K-1}$ straight interpolation.
If the geometric estimates supply LFR03's transversality to both
strata and its SAME compact isolating set, the isotope is
$C^{K-2}$. For $K\geq4$ this is at least $C^2$; the retained
$K\geq10$ allowance is ample for this particular chain. A model
fiber identification must likewise be at least $C^{K-2}$ in this
chain, or the common lower order must be recorded explicitly.
After the geometric source-fiber enclosure is supplied, LFR10
and LFR05 then transfer its smooth type and bundle structure.
No regularity for an unspecified soul, tube, model function or
fiber identification follows from this arithmetic alone.

\paragraph{What remains to produce the local models.}
KL3.5 supplies the finite-compactness route via harmonic coordinates;
its proof is a short invocation, not a written localization of all
chart and transition estimates. The source-qualified obligation is
now to produce those charts and the actual pointed maps under
BALLWISE bounds for derivatives $0,\ldots,K$, with the exact
integer $C^K$ output. Petersen's author-manuscript11.3.6 uses
uniform chart norms and H\"older exponents; its local exhaustion
and transition regularity must be justified before binding it to
that contract. A theorem with global bounded geometry cannot be
applied merely because each pointed ball has its own bound.
LFR08--LFR10 discharge the volume/injectivity and buffered-map
clauses, and LFR11 discharges the differentiable-factor clause.
They do not construct the remaining compactness data, exact metric
splitting, finite soul/core, or edge/slim interpolation estimates.

\subsection{Finite compactness sufficient for the local consumers}
\label{sec:collapse-finite-compactness-construction}

The actual GC source manifolds are smooth, but their available
curvature bounds stop at a finite order $K$. The following construction
uses those bounds to produce a $C^{K-1}$ metric and actual $C^K$
comparison maps. This deliberately records one derivative of loss.
It does not claim KL3.5's stronger integer-$C^K$ output, nor change
the historical LC81 statement. For the specified fiber consumers,
the loss leaves sufficient regularity: at the standing $K\geq10$
we retain a $C^9$ metric, $C^{10}$ product/comparison maps, and
at least $C^7$ isotopies from explicitly supplied $C^8$ model maps.
The geometric fibers and their interpolation estimates remain to be
constructed; regularity alone does not provide them.

\begin{lemma}[LFR12: local normal-coordinate estimates to a finite order]
\label{lem:collapse-finite-normal-coordinate-bounds}
For $n,K\geq1$, $A>0$ and $\iota>0$, there are $a,C>0$,
depending only on these data, with $2a<\min\{1,\iota\}$, as
follows. Suppose $g$ is smooth, $\operatorname{inj}(x)\geq\iota$,
and $|\nabla^q\operatorname{Rm}|\leq A$, $0\leq q\leq K$,
on $B(x,1)$. In any orthonormal normal chart at $x$, on the
closed Euclidean $a$-ball its metric matrix $H$ satisfies
\[
 \tfrac12 I\leq H\leq2I,\qquad
 \max_{|\alpha|\leq K}
       \bigl(|\partial^\alpha H|+|\partial^\alpha H^{-1}|\bigr)
                                                       \leq C.
\]
The estimates use only this ball and these finitely many curvature
derivatives. In particular the constants may be chosen separately
on each larger pointed ball in a sequence.
\end{lemma}
\begin{proof}
Here is a radial-frame proof with explicit curvature conventions.
The local Jacobi estimate gives the displayed ellipticity after
choosing a sufficiently small radius. Work on a normal ball with
twice the final radius, so every supremum used below is over a
compact set where the chart is defined. Parallel transport an
orthonormal frame from the center along radial geodesics. Write
\[
 \partial_i=E_i^a e_a,\qquad
 \nabla_{\partial_i}e_b=(A_i)^a{}_b e_a,\qquad
 F_{ji}=\partial_jA_i-\partial_iA_j+[A_j,A_i].
\]
Thus $F_{ji}$ is the matrix of
$R(\partial_j,\partial_i)$ in this frame. This fixes the convention;
we do not absorb a factor $1/2$ into a curvature two-form coefficient.
The radial parallel identities are $x^iA_i=0$ and $x^iE_i=x$.
Differentiating the first and using the formula for $F$ gives
$(x\cdot\partial+1)A_i=x^jF_{ji}$. Torsion freeness gives
$(x\cdot\partial+1)E_i=e_i+A_i x$. Integrating along a radius,
\[
 A_i(x)=\int_0^1 t x^jF_{ji}(tx)\,dt,\qquad
 E_i(x)=e_i+\int_0^1t A_i(tx)x\,dt.                 \tag{LFR12.1}
\]
Here the first $e_i$ on the right is the constant coordinate column.
The zero-order bounds on $E,E^{-1}$ follow from ellipticity;
the first identity then bounds $A_i$ by $c(n)A|x|$.

Let $T_q$ denote ALL components of $\nabla^q\operatorname{Rm}$
in the transported orthonormal frame, including its derivative
slots. Covariant differentiation of these components gives
\[
       \partial_i T_q=E_i^a(T_{q+1})_a+A_i*T_q,
\]
where $*$ is the finite signed sum of the connection action on
each tensor slot. Consequently order-$b$ ordinary derivatives
of $T_0$ are bounded using $T_0,\ldots,T_b$ and derivatives
of $E,A$ only through order $b-1$. This follows by induction
from the displayed identity, and avoids confusing coordinate
slots with transported-frame slots.

Suppose all lower orders of $E,A$ have been bounded and $1\leq b\leq K$.
The preceding identity bounds $\partial^b T_0$. Since
$F_{ji}$ is a contraction of $T_0$ with $E_j,E_i$,
differentiating (LFR12.1) on a ball of radius $a$ yields
\[
 A_b\leq c(n)A a E_b+B_b,
 \qquad E_b\leq a A_b+B'_b,
\]
where $E_b,A_b$ are the maximum absolute order-$b$ component
derivatives on that closed ball. The constants $B_b,B'_b$ use
only the already bounded lower orders and the given covariant
derivative bounds. The highest derivatives occur linearly: in
$F$ they differentiate one of its two $E$ factors, and in the
second integral they differentiate $A$. Derivatives hitting $x$
or the other factors belong to the lower-order terms. Factors
$t^b$ in the integrals are at most one. Choose $a$ so that
$c(n)Aa^2<1/2$ as well as the earlier radius restrictions.
Combining the two inequalities and absorbing that term bounds
$E_b$, then $A_b$. Each individual smooth chart has finite
suprema before absorption, so this does not assume the desired
uniform bound. Induction stops at $b=K$.

Finally $H=E^{\mathrm T}E$ gives its derivative bounds.
Differentiating $H^{-1}H=I$ bounds the inverse derivatives
recursively, using the ellipticity bound for its zero order.
All radii and constants were chosen from $n,K,A,\iota$.
\end{proof}

This is the interior estimate of Schick2.5(a1), arXiv
\code{math/0001108v1}, printed/PDF4, with its synchronous-frame
proof on printed/PDF5--10. The argument above uses its own
consistent frame and curvature notation and the short-radius
absorption argument, rather than copying the component identities
verbatim. The retained arXiv snapshot of Schick's paper, the published
2001 article and KL's harmonic-coordinate route are distinct sources.
No boundary-coordinate estimate is needed in this construction.

\begin{lemma}[LFR13: uniform regularity of actual transition isometries]
\label{lem:collapse-finite-transition-bounds}
Let $H_i,G_i$ be smooth positive definite metric matrices on
fixed Euclidean chart domains. Suppose they and their inverses
have uniform $C^K$ bounds and uniform ellipticity, $K\geq1$.
For actual smooth transition diffeomorphisms $\tau_i$ with
$\tau_i^*G_i=H_i$, their derivatives of orders $1,\ldots,K+1$
are uniformly bounded on every common domain where the stated
metric bounds apply. The zero order is bounded whenever their
images stay in a fixed bounded chart. Under this additional
image bound, on compact subsets of a fixed overlap a subsequence
converges in $C^K$; if a
pointwise transition limit has already been identified, every
such subsequential limit is that same map. Apply the assertion
also to the inverse transitions.
\end{lemma}
\begin{proof}
Ellipticity and $(D\tau_i)^{\mathrm T}G_i(\tau_i)D\tau_i=H_i$
bound $D\tau_i$ and its inverse. Naturality of the Levi--Civita
connection gives
\[
 \partial_a\partial_b\tau_i^c
   =\Gamma(H_i)^d_{ab}\partial_d\tau_i^c
    -\Gamma(G_i)^c_{ef}(\tau_i)
                          \partial_a\tau_i^e\partial_b\tau_i^f.
\]
The Christoffel symbols have uniform bounds through order $K-1$.
First the displayed identity bounds second derivatives. Its
successive derivatives of order $b-1$, $1\leq b\leq K$,
bound derivatives of $\tau_i$ of order $b+1$ using only
already bounded derivatives through order $b$ and the metric
coefficients through order $b$. This is ordinary finite chain
and product differentiation; no derivative of order $K+1$ of
either metric is taken. On a smaller compact overlap, the extra
derivative bounds imply equicontinuity of order-$K$ derivatives.
Arzela--Ascoli in finitely many Euclidean balls gives $C^K$
subconvergence. Pointwise identification fixes its limit.
For inverse transitions exchange the two metrics.
\end{proof}

Schick3.3--3.9, printed/PDF17--18, gives an ODE-based coordinate
change estimate with a different derivative budget. LFR13 uses
the additional EXACT isometry identity of a transition, as in
Matveev--Troyanov2.2/(2.4), printed2703/PDF5. It is not a
generic assertion about arbitrary coordinate changes. The two
metrics in this lemma are already smooth coordinate pullbacks
on the original manifold; no low-regularity isometry theorem
is invoked to begin the argument.

\begin{theorem}[LFR14: construction of the finite models and comparison maps]
\label{thm:collapse-finite-consumer-compactness}
Fix $n\geq2$, $K\geq5$, $r,v>0$ and a positive bound function
$A(R)$. Let $(M_i,g_i,p_i)$ be complete connected SMOOTH
pointed $n$-manifolds with
\[
 \operatorname{Vol}B(p_i,r)\geq v,\qquad
 |\nabla^q\operatorname{Rm}|\leq A(R)
       \text{ on }B(p_i,R),\quad 0\leq q\leq K, R>0.
\]
A subsequence has a complete connected pointed limit $(N,g,p)$
with a $C^K$ atlas and a $C^{K-1}$ metric. On a compatible
smooth carrier there are open sets $U_i$ eventually containing
each compact subset of $N$, and actual pointed $C^K$ maps
$j_i:U_i\to M_i$, such that on every compact buffered domain
\[
             j_i^*g_i\longrightarrow g\quad\hbox{in }C^{K-1}.
\]
Their pointed metric distortions tend to zero on every fixed
ball. They can be chosen to be embeddings on their whole domains,
with $j_i(p)=p_i$, and for every fixed $0<a<b$ eventually
\[
                B_{M_i}(p_i,a)\subset j_i(B_N(p,b)).
\]
All required finitely many buffers share a common tail. If the
$M_i$ are oriented, $N$ can be oriented and the maps preserve
orientation. If in addition $\sec_{g_i}\geq-\epsilon_i$ on
$B(p_i,L_i)$ with $\epsilon_i\to0$ and $L_i\to\infty$,
then $\sec_g\geq0$. No bounds above order $K$ are used.
\end{theorem}
\begin{proof}
\emph{1. Uniform data on each fixed ball and one metric limit.}
For an integer $s>0$, use LFR08 with curvature bound
$A(2s+r+4)$ to obtain a common injectivity lower bound on
$B(p_i,s+1)$. This uses that particular larger ball; $A$
need not be monotone. Curvature derivatives on the same buffer
are bounded by the same value. LFR12 gives normal charts of
one positive radius $a_s$ at all points of $B(p_i,s)$, with
uniform metric $C^K$ bounds and ellipticity on a larger chart
than will be retained. The radius may decrease and the bounds
may increase with $s$.

For each fixed small $\eta>0$, the volume propagation in LFR08
gives a positive lower bound for every $\eta$-ball centered
in $B(p_i,s)$, using one larger buffer. Bishop--Gromov at
$p_i$ gives a finite upper bound for the containing $(s+1)$-ball.
Disjoint $\eta$-balls therefore have a uniformly bounded number.
Maximal separated sets give the internal finite nets of MC13;
larger $\eta$ follow by using a smaller one. MC13 and MC23,
with the inherited preservation of approximate midpoints, give
one complete proper geodesic pointed metric limit $(X,d,p)$.
All subsequent extractions refine this subsequence and use this
SAME metric limit, rather than selecting unrelated limits of balls.

\emph{2. Local charts and local metric identity.}
Choose a countable dense set of centers in $X$, including $p$,
with approximating centers in $M_i$. At a center in the $s$-ball
use the preceding normal charts on a fixed larger Euclidean ball.
On sufficiently small inner balls these charts are uniformly
bilipschitz for the AMBIENT distances. For the lower bound, a
curve that exits the larger chart must cross a fixed radial
margin and thus cannot shorten the distance of two sufficiently
near inner points; inside the chart use ellipticity. For the
upper bound use the Euclidean segment and the metric upper bound.
The centered radial identity is exact:
$d(p_{i,a},\phi_{i,a}(w))=|w|$ within the normal chart.

Finite nets in the compact Euclidean parameter balls and a diagonal
extraction give uniformly convergent chart maps into the GH limit,
denoted $\phi_a$. Their local lower bound makes them injective.
Their images contain an open ball about the limiting center:
if $z$ is strictly inside that ball, choose source approximants
$z_i$. They are eventually in the normal chart, and their inverse
coordinates lie in a fixed compact inner Euclidean ball. A
subsequence of those coordinates converges to a preimage of $z$.
Shrinking the retained charts gives homeomorphisms onto open
neighborhoods. More explicitly, the limiting radial identity
$d(p_a,\phi_a(w))=|w|$ and this inverse-coordinate argument
identify each sufficiently small parameter ball with the whole
open metric ball of the same radius about $p_a$. No invariance
of domain theorem is being substituted for coverage. The dense
centers and positive radius on each
fixed pointed ball make these neighborhoods cover $X$.

The metric matrices in these charts have uniform $C^K$ bounds.
On smaller compact parameter balls Arzela--Ascoli gives convergence
in $C^{K-1}$ to positive definite matrices $H_a$. The convergence
of their induced LOCAL distances follows directly from uniform
ellipticity and $C^0$ convergence: compare lengths of each curve
in a fixed chart, using $(1-\delta_i)^2H_a\leq H_{i,a}
\leq(1+\delta_i)^2H_a$. Curves relevant to sufficiently close
inner endpoints remain in that chart by the radial exit margin.
Thus the metric of $H_a$ agrees locally with $d$ under $\phi_a$.
The single countable diagonal includes all chart centers, compact
subdomains and derivative orders through $K-1$.

\emph{3. Transition maps and the global finite metric.}
On each compact subset of a limiting overlap, the source
transition $\phi_{i,b}^{-1}\phi_{i,a}$ is eventually defined:
the image point is a positive distance inside the larger $b$-chart,
and its GH approximants retain that margin. The preceding
coordinate compactness identifies its pointwise limit as
$\phi_b^{-1}\phi_a$. LFR13 supplies uniform $C^{K+1}$ bounds,
so on a further countable diagonal the transitions and inverse
transitions converge in $C^K$. Their cocycle identities pass
to the limit; inverse identities and nonsingular differentials
pass as well. Thus the $\phi_a$ form a $C^K$ atlas. Passing
to the pullback identity in $C^{K-1}$ makes the $H_a$ one
$C^{K-1}$ tensor $g$ in that atlas.

The space $X$ is Hausdorff and second countable because it is
proper and exhausted by compact metric balls with finite nets.
Its locally identified length metric is globally $d$: subdivide
any $d$-geodesic into finitely many chart pieces, and conversely
subdivide any piecewise differentiable curve and compare lengths.
Thus $(X,g)$ is complete and connected. Denote it by $N$ and
choose its compatible smooth carrier by LFR01, with $K-1$ in
place of that lemma's metric order. The tensor remains $C^{K-1}$.

\emph{4. Constructing actual maps on each compact buffer.}
On a limiting chart put $f_{i,a}=\phi_{i,a}\phi_a^{-1}$.
These are $C^K$ maps to the ORIGINAL $M_i$, and they converge
to the identity in the GH identification. On a compact overlap,
their coordinate difference tends to zero in $C^K$, by the
transition convergence just proved. Also $f_{i,a}^*g_i\to g$
in $C^{K-1}$. We give the finite patching step so that actual
maps, rather than a compatible list of local germs, are produced.

Cover a compact buffer by finitely many smaller chart domains
whose closures lie inside larger charts. Start with the chart
at $p$, choosing $\phi_{i,0}(0)=p_i$. Inductively a map $F_i$
is defined on a neighborhood of the compact pieces already
covered and is $C^K$ close to every local $f_{i,a}$ there.
Add the next compact chart piece. Shrink the combined working
neighborhood inside the union of the old domain and the new
chart, retaining all compact pieces. Choose a fixed partition
$\lambda+\mu=1$, subordinate to those two domains, with
supports compactly inside them on that working neighborhood.
It may be chosen with $\mu=0$ on a fixed smaller neighborhood
of $p$ contained in the first chart.

Where $\mu\ne0$, use the new target chart to define
\[
 \widetilde F_i(x)=\phi_{i,a}\bigl(
     \lambda(x)\phi_{i,a}^{-1}(F_i(x))
                +\mu(x)\phi_a^{-1}(x)\bigr).       \tag{LFR14.1}
\]
The first summand is used only on the support of $\lambda$;
elsewhere the formula is $f_{i,a}$. Where $\mu=0$ keep $F_i$.
On the overlap both coordinate values are close to
$\phi_a^{-1}(x)$ in a compact inner ball of the larger chart,
so their segment stays in its domain. The support margins make
these prescriptions agree smoothly up to class $C^K$; no map
is evaluated outside its domain. The cutoffs and their derivative
bounds are fixed BEFORE taking $i$ large. Finite product and
chain rules show that (LFR14.1) remains $C^K$ close to all the
local maps, since it changes $F_i$ by a fixed cutoff times a
coordinate difference tending to zero. This completes the finite
induction and preserves the original chart map near $p$.

The resulting actual $F_i$ are pointed and have
$F_i^*g_i\to g$ in $C^{K-1}$ on the compact buffer. Let
$q_i:M_i\to X$ be the pointed metric approximations used to
construct the limiting charts, on suitably enlarged balls.
Uniform convergence of $q_i\phi_{i,a}$ to $\phi_a$, and the
uniformly small coordinate changes in patching, give
$q_iF_i\to\operatorname{id}$ uniformly there. Therefore
\[
 \begin{split}
 |d_i(F_i(x),F_i(y))-d(x,y)|
 &\leq\operatorname{dis}(q_i)
       +d(q_iF_i(x),x)+d(q_iF_i(y),y)\longrightarrow0.
 \end{split}
\]
All points are in the chosen common approximation domain.
This proves DISTANCE distortion control without requiring
the $q_i$ to be continuous. This conclusion
does not follow from differential convergence on that buffer alone.

\emph{5. Exhaustion, embeddings and retained geometric properties.}
Repeat step4 on compact balls with integer radii and larger
buffers. For the first $s$ chart subdomains and the $s$-ball,
choose a tail making the coordinate, metric and distance errors
less than $1/s$. Apply LFR10 with its metric order $K-1$ to
these actual maps; its proof on any fixed larger buffer uses
only the maps and estimates on that buffer. Increase this tail
so the maps embed the prescribed open $s$-ball and cover the
smaller source $(s-1)$-ball. Choose these thresholds increasing.
For index $i$, use the map from a radius $s(i)\to\infty$ whose
threshold has been passed, restricting its domain to that open
$s(i)$-ball. This constructs one sequence $j_i$ with the
asserted convergence on EVERY fixed compact subset. No agreement
between maps from two different radii is required. LFR10 then
gives the stated containment for every fixed $a<b$, as well as
a common tail for any finite list of buffers.

For orientations choose oriented frames in step2. The transition
Jacobians are positive, and their nonsingular limits stay positive.
Patching maps are $C^1$ close to these oriented local maps, so
are orientation preserving. Finally curvature is a continuous
expression in the metric, its inverse, and its first two
derivatives. Since $K-1\geq4$, $C^2$ convergence is available.
Images of each fixed compact buffer remain in a fixed source
ball by the pointed distortion estimate. The additional
expanding-ball lower sectional bound therefore passes to $g$.
This proves all assertions, using no derivative above order $K$.
\end{proof}

The metric extraction is MC13/MC23. The Riemannian chart and
patching construction follows the structure of Petersen's
author-manuscript11.3.6, printed358--361/PDF374--377, with
ballwise constants, compact overlap supports and basepoint
preservation made explicit. LFR13 supplies its transition
regularity here from smooth source isometries. The proof does
not apply a global bounded-geometry theorem to local bounds,
or replace a finite curvature budget by all-order compactness.

\begin{corollary}[LFR15: the same finite model tracks the actual splitting]
\label{cor:collapse-finite-model-splitting-producer}
In LFR14 assume its expanding-ball lower sectional bounds and
supply normalized $(j,\delta_i)$-splittings
$\alpha_i=(u_i,v_i):M_i\to\mathbb R^j\times Y_i$ with fixed
$0<j<n$ and $\delta_i\to0$. Along a common subsequence the
SAME complete $C^{K-1}$ limit admits an exact product
$\mathbb R^j\times Z$, with $Z$ a complete nonnegative
$C^{K-1}$ Riemannian $(n-j)$-manifold, and the actual product
diffeomorphism is $C^K$. The maps $j_i$ may be retained so that
on every fixed compact subset of $N$,
\[
                 u_i\circ j_i\longrightarrow t
                                \quad\hbox{uniformly},
\]
where $t$ is the exact Euclidean coordinate of that product.
The convergence of pulled-back metrics remains $C^{K-1}$.
When the source manifolds are oriented, so is $Z$ with the
ordered Euclidean normal factor. No differentiability of the
original maps $\alpha_i$ is asserted.
\end{corollary}
\begin{proof}
Use MC09 to take metric quasi-inverses $h_i$ of the actual $j_i$
on growing balls, with errors tending to zero. They map source
to limit and satisfy $h_i j_i\to\operatorname{id}$ uniformly
on fixed balls. AC50--AC51 apply to the supplied approximate
products and this same proper limit. In the proof of AC51
choose these $h_i$ as its approximation maps. It produces an
onto pointed isometry $I:N\to\mathbb R^j\times Y$ and
\[
       \sup_{x\in B(p_i,R)}|u_i(x)-\pi_E I(h_i(x))|
                                              \longrightarrow0
\]
for every fixed $R$. Evaluating at $x=j_i(z)$, whose pointed
radius is uniformly bounded on a fixed compact set, and using
the $1$-Lipschitz map $\pi_E I$ proves the stated convergence.
Apply LFR11 with metric order $K-1$ to this ACTUAL metric
isometry. It constructs the factor $Z$, its complete induced
metric, its orientation and the $C^K$ product map. Neither
the metric nor the comparison maps have been replaced.
\end{proof}

\paragraph{Disposition of the LC81 bridge after this construction.}
For the original smooth manifolds used by local collapse,
LFR14--LFR15 now produce the finite models, actual buffered
embeddings, and complete differentiable splitting factors with
the original splitting coordinates retained. Thus those clauses
are no longer merely requested data. The chosen output order is
$m=K-1$, not $K$. Apply LFR02--LFR05 and LFR11 with that
order: explicitly supplied $C^{m-1}$ model maps give
$C^{m-2}=C^{K-3}$ isotopies. Since $K\geq5$, the latter are
at least $C^2$, sufficient for smooth fiber-type recovery.
The standing $K\geq10$ supplies further margin. This completes
the compactness-to-category part needed by those consumers;
the stronger historical LC81 regularity, LC56's smooth-metric
package, model soul/core and classification constructions, and
the edge/slim geometric estimates are separate. In particular
uniform value control of $u_i j_i-t$ is not a $C^1$ estimate
for an adapted source map. That derivative estimate must still
be proved using its actual minimizing directions and the
unchanged comparison metrics.

\subsection{A complete slim packet in the finite category}
\label{sec:collapse-finite-slim-construction}

This construction supplies LC85's compact factor, its smooth fiber
type, and the actual source-fiber enclosure. It uses the original
smooth source metric and the finite product from LFR15. The metric
on that product is not upgraded to a smooth nonnegative metric.
The surface calculation below needs only the inherited smooth
Gauss--Bonnet theorem and smooth classification of closed orientable
surfaces, not the full three-dimensional nonnegative classification.

\begin{lemma}[LFR16: eventual ballwise bounds and a compact factor]
\label{lem:collapse-finite-compact-factor}
In the three-dimensional oriented case of LFR14--LFR15, replace
the curvature-derivative hypothesis by the following eventual version:
for every $R>0$ there is $i_R$ such that the bounds $0,\ldots,K$
by $A(R)$ hold on $B(p_i,R)$ for all $i\geq i_R$.
Keep the common basepoint volume lower bound, the expanding-ball
sectional lower bounds tending to zero, and the supplied normalized
$(1,\delta_i)$-splittings, $\delta_i\to0$.
The conclusions of LFR14--LFR15 still hold. If in addition every
supplied residual factor $Y_i$ has diameter at most a fixed $D$,
the SAME factor $Z$ is a compact connected orientable surface
without boundary, with $\operatorname{diam}Z\leq D$.
Its metric is $C^{K-1}$ and nonnegative; its compatible carrier
and actual product map have the regularity recorded in LFR15.
\end{lemma}
\begin{proof}
Every estimate in step1 of LFR14 uses finitely many fixed pointed
balls. Start that estimate after the maximum of their thresholds
$i_R$. The normal charts, their inner buffers and the packing bounds
then hold on that tail. MC13 explicitly permits eventual net bounds.
The later chart, overlap and map constructions likewise use finitely
many buffers at each stage. In the countable diagonal, increase its
stage-$s$ threshold to exceed all the finitely many $i_R$ used through
that stage. The nested subsequences and the final growing-domain
selection still tend to infinity. Thus each fixed-ball conclusion
holds eventually, exactly as in LFR14. No member of the sequence
is asserted to satisfy uniform derivative bounds on every ball.
The expanding curvature and splitting arguments of LFR15 are unchanged.

AC50 constructs a complete proper limit $Y$ of the actual $Y_i$.
For $y,y'\in Y$, choose a fixed pointed ball containing both with
strict buffer. Approximation coverage supplies points $y_i,y_i'$
in the corresponding factor balls whose images tend to $y,y'$.
Their distances are at most $D$, and distortion tends to zero, so
$d_Y(y,y')\leq D$. Hence the WHOLE space $Y$ is bounded and proper,
and is compact. This uses the global diameter hypothesis on $Y_i$,
not a bound on just one small ball. LFR11 constructs $Z$ isometric
to $Y$, so $Z$ is compact. The limit $N$ is connected, and projection
$\mathbb R\times Z\to Z$ is continuous and onto, so $Z$ is connected.
Its absence of boundary and its induced orientation follow from
LFR11's regular-level construction in the oriented boundaryless
three-manifold, with the ordered real normal factor.
\end{proof}

\begin{lemma}[LFR17: closed nonnegative finite surfaces have smooth type]
\label{lem:collapse-finite-surface-classification}
Let $Z$ be a nonempty compact connected orientable surface without
boundary, on a compatible smooth carrier, with a $C^2$ Riemannian
metric $k$ and Gaussian curvature $\kappa_k\geq0$.
Then its smooth carrier is diffeomorphic to $S^2$ or $T^2$.
In the torus case $\kappa_k$ vanishes identically. No smooth
nonnegative approximation of $k$ is assumed or produced.
\end{lemma}
\begin{proof}
Choose a finite smooth coordinate cover and a subordinate smooth
partition of unity with compact supports in its charts. On slightly
larger coordinate domains, convolve the symmetric coefficient
matrices of $k$ with Euclidean mollifiers. Express the resulting
local smooth symmetric tensors back in those same charts and sum
with the partition. Finite product differentiation and the usual
$C^2$ convergence of mollification give smooth symmetric tensors
$k_j\to k$ in $C^2$. Positivity is open: relative to any fixed
smooth reference metric the least eigenvalue of $k$ has a positive
minimum on the compact surface. For large $j$, $k_j$ are therefore
Riemannian metrics. Their curvatures need NOT be nonnegative.

In the fixed finite chart cover, curvature is a continuous expression
in the metric, its inverse and its first two derivatives. Thus
$\kappa_{k_j}\to\kappa_k$ uniformly, and the volume densities also
converge uniformly relative to the fixed smooth density. Smooth
Gauss--Bonnet on the SAME smooth surface gives
\[
 2\pi\chi(Z)=\int_Z\kappa_{k_j}\,dA_{k_j}
       \longrightarrow\int_Z\kappa_k\,dA_k\geq0.
\]
This is the only use of the auxiliary smooth metrics. All subsequent
geometric comparisons keep $k$.

The inherited smooth classification of closed connected orientable
surfaces gives an integer $g\geq0$, a smooth diffeomorphism to the
genus-$g$ model, and $\chi(Z)=2-2g$. Hence $g=0$ or $1$.
This is classification up to DIFFEOMORPHISM, not an unproved upgrade
of a homeomorphism. If $g=1$, the last integral is zero. A continuous
nonnegative function with zero integral against a positive Riemannian
density vanishes everywhere: a positive value would stay positive on
a coordinate neighborhood of positive area. This proves the final claim.
\end{proof}

For the smooth inputs, Lee's archived IRM \code{c09.tex:792--975}
contains the Gauss--Bonnet formula and its proof, and
\code{c09.tex:997--1211} contains triangulations, the global theorem
and proof, and the surface-sign corollary. Its classification discussion
is topological. Hirsch's corrected1994 scan supplies the SMOOTH
classification in Theorem9.3.5, printed204--205/PDF109, with the
preceding Morse decomposition on printed200--204/PDF107--109.
The genus definition and $2-2g$ formula on printed190--191 are on
PDF103: that scan has PDF102 and PDF103 out of printed order.
These foundational smooth results are to be reused from the accepted
PC foundation. The argument above is their finite-metric adapter;
no new PC declaration has been inspected or accepted.

\begin{lemma}[LFR18: every vertical source direction converges]
\label{lem:collapse-finite-vertical-directions}
Let $N=\mathbb R\times Z$ have a complete product metric $g$ of
class $C^m$, $m\geq2$, on its compatible differentiable carrier.
Supply actual $C^3$ growing-domain embeddings $j_i:N\supset U_i\to M_i$
with the pointed distance, source-ball coverage and $C^1$ pullback
convergence of LFR14. Fix $\ell>0$ and a compact set $C\subset N$.
For $x=(t,z)\in C$, put $x^+=(t+\ell,z)$.
For EVERY source minimizing unit direction $w$ from $j_i(x)$
to $j_i(x^+)$,
\[
 \sup_{x\in C,w}
       \bigl\|(d j_i|_x)^{-1}w-\partial_t\bigr\|_g
                                                     \longrightarrow0.
\]
No selected family of minimizing geodesics is required to be continuous.
\end{lemma}
\begin{proof}
If uniform convergence fails, choose bad points $x_i\in C$ and
bad minimizing segments. Their endpoints converge after extraction,
and their lengths tend to $\ell$ by the distance distortion estimate.
The starting points have bounded source radius; a minimizing segment
of bounded length stays in a fixed larger source ball. LFR14's
coverage with strict radii puts that entire ball in the image of a
fixed still larger model ball. Hence each segment has a unique lift
whose image is in one compact model ball. On that ball the pulled-back
metrics are uniformly comparable with $g$.

Parametrize the lifts on $[0,1]$ at constant source speed. Their
$g$-speeds are uniformly bounded. Arzela--Ascoli gives uniform path
convergence on a subsequence. In finitely many buffered charts along
the limit path, $C^1$ metric convergence gives uniform convergence
of the Christoffel symbols. The coordinate geodesic equations bound
accelerations. A further finite extraction gives $C^1$ convergence
on chart intervals, including the endpoint velocities. Passing to
the integral equation gives a $g$-geodesic of length $\ell$ between
$x$ and $x^+$. This is LC50's argument with both endpoints moving;
it uses only a $C^2$ limiting metric, not a smooth metric upgrade.

The unique minimizing segment between these product endpoints is
vertical. Indeed length is at least the variation of the real
coordinate, which is at least $\ell$. Equality with length $\ell$
forces the transverse speed to vanish and the real speed to be
nonnegative; the constant-speed parametrization fixes that speed.
Thus its initial unit vector is $\partial_t$. Endpoint velocity
convergence contradicts the selected bad directions. This proves
the assertion for all choices simultaneously.
\end{proof}

\begin{lemma}[LFR19: rank-one coordinates with a separate value tolerance]
\label{lem:collapse-buffered-scalar-coordinate}
Fix $L>0$, $T>2L$, $e>0$ and $0<\sigma<1$.
For sufficiently small $\beta>0$, a supplied normalized
$(1,\beta)$-splitting $\alpha=(u,v)$ of a complete smooth pointed
manifold, with $\sec\geq-\beta^2$ on $B(p,\beta^{-1})$,
has a scalar coordinate $\eta$, smooth on a neighborhood of
$\overline B(p,L)$, such that
\[
 \eta(p)=0,\quad\operatorname{Lip}(\eta|_{B(p,L)})\leq1+\sigma,
 \qquad |\eta-u|<e\quad\hbox{on }B(p,L).
\]
Its image is within Hausdorff distance $\sigma L$ of $(-L,L)$.
For $x\in B(p,L)$, $x'\in B(p,T)$ with $d(x,x')>L$, and
every minimizing initial unit direction $w$ from $x$ to $x'$,
\[
 \left|D\eta_x(w)-\frac{u(x')-u(x)}{d(x,x')}\right|<\sigma.
                                                        \tag{LFR19.1}
\]
The threshold depends only on $L,T,e,\sigma$ and the dimension.
\end{lemma}
\begin{proof}
Rescale distances by $L^{-1}$ and use the SAME splitting, with
Euclidean coordinate $u/L$ and rescaled residual metric. The buffered
rescaling conversion MC10--MC11 gives its required normalized small
parameter once $\beta$ is small; the rescaled sectional bound is
$-L^2\beta^2$, which also tends to zero on expanding balls.
Choose an auxiliary quality
\[
 0<\gamma<\min\{\sigma/4,e/(4L),L/(2T),1/10\}.
\]
Use the construction in LC79, with its anchor calculation LC78 and
smoothing now supplied by LFR02. In that construction choose the
anchor value error and the Lipschitz-difference error each below
$\gamma/4$ on the rescaled unit ball. It then gives
$|\phi-u/L|<\gamma$, in addition to its stated Lipschitz, image
and ALL-direction estimates. The tested outer ball contains the
rescaled $T$-ball since $\gamma^{-1}>T/L$; the length test is
$d/L>1$. The smoothing compact set is the closed unit ball,
so smoothness holds on a neighborhood of it. The Lipschitz bound
comes from the globally Lipschitz DIFFERENCE and the distance
function, and does not assume the unit ball is geodesically convex.
Put $\eta=L\phi$. Values, image distances and tested lengths
rescale by $L$, while unit-direction derivatives and Lipschitz
constants retain their bounds. The asserted tolerances follow.
All choices precede the manifold and its actual approximation map.
\end{proof}

\begin{theorem}[LFR20: the complete local slim surface packet]
\label{thm:collapse-finite-slim-packet}
Fix $\Delta\geq1$, $K\geq5$, $r,v>0$ and $A(R)>0$.
Set
\[
 L=10^6\Delta,\quad D=10^3\Delta,\quad e=\Delta/100,
 \quad a=9L/10,\quad b=19L/20.
\]
For every $0<\sigma\leq1/100$, there is $\beta_0>0$ as follows.
Let $0<\beta<\beta_0$ and let $(M,g_M,p)$ be a complete connected
oriented smooth three-manifold satisfying
\[
 \Vol B(p,r)\geq v,\qquad
 |\nabla^q\Rm|\leq A(R)\text{ on }B(p,R),
 \quad 0\leq q\leq K,\quad 0<R<\beta^{-1},
\]
and $\sec\geq-\beta^2$ on $B(p,\beta^{-1})$.
Supply an actual normalized $(1,\beta)$-splitting
$\alpha=(u,v_Y):M\to\mathbb R\times Y$ with
$\operatorname{diam}Y\leq D$. Then:
\begin{enumerate}[leftmargin=*]
\item LFR19 supplies $\eta$ with $T=L/\sigma$, value error $e$
  and quality $\sigma$ on $B(p,L)$.
\item For EVERY such $\eta$, the open set
  $\Omega=\{x\in B(p,L):|\eta(x)|<a\}$ is a smooth proper
  trivial bundle over $(-a,a)$ under $\eta$. Its connected fiber
  is smoothly $S^2$ or $T^2$. In particular $\eta^{-1}(0)$ means
  the ENTIRE zero fiber in its stated domain.
\item There is a complete finite product $N=\mathbb R\times Z$
  with $\operatorname{diam}Z\leq D$, a compact connected oriented
  nonnegative $C^{K-1}$ surface $Z$, and an actual pointed $C^K$
  embedding $j$ of a neighborhood of $[-b,b]\times Z$ into $M$.
  With $f=\eta\circ j$ there,
\[
 |f-t|<2e,\qquad \partial_t f>3/4.
\]
  Its image contains EVERY source fiber over $[-a,a]$ in $B(p,L)$.
  The straight interpolation from $t$ to $f$ is transverse, with
  all inverse images of $[-a,a]$ enclosed in the same compact
  $Q=[-a-3e,a+3e]\times Z\subset(-b,b)\times Z$.
\item A smooth profile of $\eta$ equal to one on
  $|\eta|\leq8\cdot10^5\Delta$ and zero for
  $|\eta|\geq9\cdot10^5\Delta$, with closed support strictly
  inside $|\eta|<9\cdot10^5\Delta$, extends by zero to a smooth
  compactly supported function on $M$.
\end{enumerate}
The model embedding can additionally meet any prescribed positive
finite-buffer $C^{K-1}$ metric and pointed-distance comparison
tolerance, by decreasing $\beta_0$ after that tolerance is fixed.
This does not claim a smooth model metric. The derivative bounds
are required only on the displayed expanding range of radii.
\end{theorem}
\begin{proof}
LFR19 proves existence of the coordinates after a uniform reduction
of $\beta_0$. We prove the remaining assertions uniformly for ALL
coordinates satisfying its displayed estimates. If this fails for
arbitrarily small $\beta$, choose a sequence $\beta_i\to0$ of
counterexamples and their actual maps and coordinates. For each fixed
$R$, eventually $R<\beta_i^{-1}$, so LFR16 applies with the fixed
bound function $A$. It supplies, along one subsequence, a complete
product $N=\mathbb R\times Z$, the SAME actual pointed embeddings
$j_i$, and
\[
 h_i=j_i^*g_i\to g_N\text{ in }C^{K-1},\qquad
 u_i j_i\to t\text{ uniformly on every fixed compact buffer}.
\]
The factor is compact with diameter at most $D$ by LFR16 and has
smooth type $S^2$ or $T^2$ by LFR17. All of the following finitely
many comparisons share one tail. The basepoint is $(0,z_0)$.

\emph{1. A whole cylinder in the controlled source domain.}
The compact cylinder $C=[-b,b]\times Z$ lies in the model ball
of radius $\sqrt{b^2+D^2}<0.96L$. Pointed distance convergence
therefore puts $j_i(C)$ strictly inside $B(p_i,L)$ eventually.
Set $f_i=\eta_i j_i$ on a neighborhood of $C$; it is $C^K$.
The retained value estimate gives $|f_i-t|<2e$ on $C$ for large $i$.
No derivative conclusion has yet been inferred from this value bound.

\emph{2. The actual derivative estimate.}
For $x=(t,z)\in C$, test against $x^+=(t+2L,z)$.
These points lie in a fixed compact model ball of radius less than
$3L$. Their source distances from each other tend uniformly to $2L$,
and their pointed radii are eventually less than $4L<T=L/\sigma$.
Thus (LFR19.1) applies to EVERY minimizing source direction $w_i$.
LFR18 says that its pulled-back vector tends uniformly to
$E=\partial_t$. The quotient in (LFR19.1) tends uniformly to one,
because the original coordinates $u_i j_i$ converge uniformly to
$t$ on this fixed larger buffer. Moreover
$|D\eta_i|\leq1+\sigma$ on $B(p_i,L)$ and $h_i\to g_N$.
Consequently
\[
       D f_i(E)\geq1-\sigma-o(1)>3/4\quad\hbox{on }C.
                                                        \tag{LFR20.1}
\]
For clarity the stronger quantitative conclusion is
$\|df_i-dt\|_{g_N}^2\leq4\sigma+\sigma^2+o(1)$:
use $|df_i|_{g_N}\leq1+\sigma+o(1)$, $|dt|=1$, and (LFR20.1)
in the polarization identity. For fixed $\sigma$ this is a bound,
not convergence to zero. Only the positive derivative is needed below.

\emph{3. Enclosing the ENTIRE source fibers.}
For any $x\in B(p_i,L)$ with $|\eta_i(x)|\leq a$, the SAME
original approximation, its pointed distortion and the global factor
diameter bound give
\[
 d(p_i,x)\leq
       \sqrt{(a+e)^2+D^2}+\beta_i<0.91L.             \tag{LFR20.2}
\]
Here $B(p_i,L)$ is inside its tested domain once $\beta_i^{-1}>L$.
LFR14's strict-radius coverage gives
$B(p_i,0.92L)\subset j_i(B_N(p,0.93L))$ eventually.
Thus every such $x$ has a preimage with $|t|<0.93L<b$,
and hence belongs to $j_i((-b,b)\times Z)$.
There cannot be an additional source-fiber component outside that
image. This conclusion uses (LFR20.2), not model-side properness alone.
The same inequality for the closed slab has a strict margin inside
the smooth domain, so that slab is compact: it is closed inside a
fixed compact source ball of radius between $0.91L$ and $L$.

\emph{4. Graphs, interpolation and smooth bundles.}
At $t=\pm b$, the value bound implies $f_i(b,z)>a$ and
$f_i(-b,z)<-a$, since $b-a>2e$. For each $s\in[-a,a]$
and $z\in Z$, (LFR20.1) and the intermediate value theorem give
one and only one root $t=t_i(s,z)\in(-b,b)$ of $f_i(t,z)=s$.
The finite implicit function theorem makes $t_i$ jointly $C^K$.
The whole fiber is therefore the $C^K$ graph over $Z$, and its
image under the actual $j_i$ is precisely the entire source fiber
by step3. This gives connectedness as well as its finite differentiable
type. The source fiber is a smooth regular level of the smooth
function $\eta_i$. LFR04 upgrades its $C^K$ diffeomorphism with
the compatible smooth $Z$ to a smooth diffeomorphism, since $K\geq5$.

For $F_{i,u}=(1-u)t+u f_i$, $0\leq u\leq1$, the same derivative
is greater than $3/4$. If $|F_{i,u}|\leq a$, then
$|t|\leq a+2e$, so all these inverse images lie in the interior
of the SAME compact $Q$ in the statement. The jointly $C^K$
family satisfies LFR03, giving a compactly supported $C^{K-1}$
isotopy for its zero fibers. These maps are type comparisons;
they do not replace the metric embeddings $j_i$.

On the original smooth source, $\eta_i:\Omega_i\to(-a,a)$ is a
surjective smooth submersion. It is proper: the inverse image of
any compact target set is a closed subset of the compact closed
slab from step3. To see smooth triviality directly, use the smooth
vector field $X_i=\nabla\eta_i/|\nabla\eta_i|^2$ on $\Omega_i$.
Along its flow $d\eta_i(X_i)=1$. For any compact subinterval of
$(-a,a)$, properness traps these trajectories in a compact set,
so they exist for all the finite times needed to reach its levels.
The flow and its reverse give a smooth product with the zero fiber.
No assertion of global completeness of $X_i$ is needed.

\emph{5. Cutoff, comparison tolerance and uniform threshold.}
Choose a smooth real profile supported in $[-0.89L,0.89L]$ and
equal to one on $[-0.8L,0.8L]$. Thus its CLOSED support is
strictly inside $(-a,a)$, not merely zero at its endpoints.
Its composition with $\eta_i$ has support
in the compact closed slab of step3, strictly inside $B(p_i,L)$.
It therefore extends by zero smoothly and with compact support.
All requested comparison estimates follow on the same finite buffers
from LFR14--LFR15; for example one may measure the metric error by
$\sum_{q=0}^{K-1}\sup|\nabla_{g_N}^q(h_i-g_N)|$ there.
Coordinate convergence implies this fixed finite norm tends to zero.
Choose the tail also for any prescribed positive comparison tolerance.
We have produced every asserted object for all sufficiently late
counterexamples, a contradiction. This proves existence of one
threshold depending on the fixed data and, when requested, the
comparison tolerance. No tolerance has been reselected after choosing
an individual bad manifold.
\end{proof}

\paragraph{Scope of the slim completion.}
This proves LC85's local geometric producer under its explicit
finite-derivative, noncollapse, small lower-curvature defect and
bounded-factor splitting inputs. The constant controlling the curvature
defect and the error of the supplied splitting have been combined in
$\beta$ for the statement; small splitting error alone does not imply
the curvature bound. In the collapse application their separate
parameters must both meet the chosen threshold, with buffered
weakening of approximations as in MC11. The normalized analytic
hypotheses still come from the inherited scale estimates and the
static/flow hypotheses, not from the definition of a slim point.

KL10.1--10.3, printed58/PDF53, is the source route. Its compactness,
vertical-direction and perturbation steps are expanded here in the
finite category. The rescaled adapted test uses length $>L$ and
outer radius $L/\sigma$; the test shift $2L$ satisfies it explicitly.
The same product supplies the model, full fiber, compact interpolation
trace and cutoff. Smooth Gauss--Bonnet and smooth surface classification
are inherited foundations, not new PC absence claims. The edge disk
producer, its endpoint/core geometry, general line splitting, Tits-cone
existence and full three-dimensional model classification remain
separate. Global overlap compatibility and assembly stay in Chapter14.


\subsection{Finite surface topology and the edge model disk}
\label{sec:collapse-finite-edge-core}

The slim consumer used a compact surface. The edge consumer needs a
disk in a possibly noncompact surface, with a controlled collar on the
SAME finite metric. The following adapter uses the topological part of
the inherited soul construction. It does not infer a smooth metric or
a smooth normal-flow map from a homeomorphism. Riemannian comparison
and the elementary manifold, ODE and surface-topology foundations are
inherited inputs; the general Alexandrov comparison/splitting bindings
remain separate.

\begin{lemma}[LFR21: the finite surface soul-topology adapter]
\label{lem:collapse-finite-surface-soul-topology}
Let $Z$ be a complete connected noncompact surface, on a compatible
smooth carrier, with a $C^m$ metric $k$, $m\geq4$, and $\sec_k\geq0$.
The topological Cheeger--Gromoll construction gives a compact connected
boundaryless totally geodesic submanifold $S$ of dimension zero or one,
and a HOMEOMORPHISM $\nu S\to Z$ fixing $S$.
Here $S$ can be taken $C^{m-1}$ and its normal bundle $C^{m-2}$;
the conclusion asserts no differentiability of that homeomorphism.
In particular, if $Z$ is orientable, it is homeomorphic to
$\mathbb R^2$ or $S^1\times\mathbb R$.
\end{lemma}
\begin{proof}
We port the topological construction, including its maps, rather than
the unexpanded smoothing assertion in Cheeger--Gromoll 2.2.
The required finite calculus is as follows. In the fixed smooth charts,
$\Gamma$ is $C^{m-1}$, curvature is $C^{m-2}$, and geodesic flow and
the exponential map are $C^{m-1}$. The geodesic variations used below
are therefore at least $C^3$. First and second variation and the index
form are obtained by differentiating these variations and integrating
by parts; no infinite derivative is used. In a parallel frame the
Jacobi and Riccati equations have continuous curvature coefficients.
The matrix comparison proof uses a maximum eigenvector and a one-sided
derivative at a positive maximum, so it applies to these coefficients.
At the initial endpoint it requires $f(a,x)\leq0$, not equality;
this incorporates Lee's May31,2022 correction to printed page326.
It gives both the point version and the short transverse version of
Rauch comparison. The latter is used only before focal points, on a
uniformly short compact tube. The inherited Riemannian triangle
comparison, including the two-minimal-side version, has the same
first/second-variation requirements. This is a finite-category use of
those comparison bindings, not a new general Alexandrov theorem.

For clarity we track every stage of the topological construction.
Fix $o\in Z$. For every ray $c$ from $o$, let
$B_c=\bigcup_{s>0}B(c(s),s)$ and write $c_t(s)=c(t+s)$.
The comparison argument of Cheeger--Gromoll 1.2 shows that
$Z\setminus B_c$ is totally convex, with ALL geodesic segments
between its endpoints required to stay in it. Define
\[
 C_t=\bigcap_{c}(Z\setminus B_{c_t}),\qquad t\geq0.
\]
These sets are closed and totally convex, contain $\overline B(o,t)$,
and exhaust $Z$. They are compact: otherwise minimizing segments from
$o$ to escaping points in one $C_t$ have a subsequence converging to a
ray contained in $C_t$, while that ray's own half-space excludes its
points beyond parameter $t$. Properness and compactness of the unit
tangent circle justify this extraction. The triangle-inequality
calculation in 1.3 also gives, for $T>t$,
\[
 C_t=\{x\in C_T:d(x,\partial C_T)\geq T-t\}.
\tag{LFR21.1}
\]
The distance here is the ambient distance; minimizing connections in
a totally convex set realize its intrinsic distance as well.

The convex-set structure argument in 1.4--1.7 is valid with
$C^{m-1}$ submanifolds in place of smooth submanifolds. In a small
convex ball, maximize the dimension of a submanifold contained in the
set. If a point leaves its local tangent sheet, coning a transverse
patch to that point with short geodesics produces a submanifold of
one higher dimension, a contradiction. These cones are $C^{m-1}$.
The resulting interior is totally geodesic. Radial parametrization
from an interior point gives continuous boundary charts, since each
ray hits that boundary at most once and the hitting radius is
continuous. Thus the closed set is a topological manifold with
boundary, with $C^{m-1}$ totally geodesic interior; its boundary is
not asserted differentiable.

For a compact totally convex set $C$ with nonempty boundary, put
$C^a=\{x\in C:d(x,\partial C)\geq a\}$. The three-angle
comparison in 1.10 gives concavity of distance to this boundary along
each geodesic in $C$. Its short variations, parallel fields and
supporting half-spaces all have the finite regularity just verified.
Hence every $C^a$ is totally convex. The maximum set has strictly
smaller dimension: along a minimizing segment to the boundary the
distance immediately decreases, so the maximum set has no relative
interior of full dimension. Repeating this operation terminates in
at most two dimension drops, at a compact boundaryless totally
geodesic $S$. Connectedness follows at each step from convexity.
Since the initial $C_t$ is a proper compact subset of the noncompact
surface, $\dim S<2$. The geodesic charts give $S$ order $C^{m-1}$;
orthogonal complementation gives its normal bundle order $C^{m-2}$.
Its short normal exponential tube is at least $C^2$ and is a genuine
tube by the finite inverse-function theorem and compactness of $S$.

Here is the map construction that carries the topology through the
possibly nonsmooth boundaries. On a compact neighborhood choose the
positive short radius from 2.3: balls are strictly convex and distance
to their center strictly increases along a geodesic starting with
nonnegative outward pairing. Uniform continuity of the compact family
$C^a$ makes successive levels close enough to lie within that radius.
For a point $x$ in such a layer the nearest point $h(x)$ on a slightly
deeper $C^{a'}$ is unique. If two minimizers existed, the geodesic in
$C^{a'}$ joining them and first variation would contradict that
strict increase. Compactness then makes $h$ continuous. The vectors
\[
             v_x=(\exp_x|_{B(0,s)})^{-1}(h(x))
\]
depend continuously on $x$. The arcs $\exp_x(uv_x)$ have unique
successive level crossings. Distinct arcs cannot merge outside the
deeper set, by nearest-point uniqueness and uniqueness of geodesics.
Extending each arc backwards within its convex ball proves coverage
of the whole layer. Their crossing times are continuous, so they
give the collar homeomorphisms of 2.4(2).

The same arcs give 2.4(4): reparametrize each arc by its distance
from the deeper set to identify a near-maximum level with a small
metric neighborhood of that set, fixing the deeper set. For a
subsequent lower-dimensional set, 2.4(5) applies the same construction
to the intersections of the original convex set with its small
neighborhoods. Its uniqueness argument uses the same short-radius
distance monotonicity; continuity at the fixed deeper set follows
because the reparametrized arcs have length tending to zero there.
A finite subdivision of each level interval and the at most two
dimension drops give the homeomorphism from a closed normal disk
bundle to $C_t$ in 2.5(3). Only its innermost identification is the
finite normal tube; the outer maps remain topological.

The same $S$ works for every expanding stage: (LFR21.1) erodes
each $C_j$ to the common $C_0$, after which the remaining finite
flag is chosen once. Finally (LFR21.1) and the collar construction identify each shell
$C_{j+1}\setminus\operatorname{int}C_j$ with the corresponding
normal-bundle shell. Choose each collar's boundary parametrization
to agree with the preceding one. Gluing gives an actual bijection
$\nu S\to Z$, continuous with continuous inverse: each point lies
in the interior of one compact stage on both exhaustions. It fixes
$S$. No Hilbert-manifold Morse argument and no smoothing of the
resulting homeomorphism is needed.

If $\dim S=0$, connectedness makes $S$ a point, and $\nu S$ is a
plane. Otherwise the inherited compact one-manifold classification
makes $S$ a circle. Orient $S$. Together with the orientation of
$Z$, this orients its normal line bundle; an oriented real line
bundle is trivial, by choosing positive local sections and patching
them. Its total space is $S^1\times\mathbb R$.
\end{proof}

The checked source is Cheeger--Gromoll, Section1 and the topological
part of Section2, printed415--428/PDF4--17. Page424 explicitly
separates the displayed homeomorphism construction from the omitted
differentiable refinement. The finite-coordinate comparison check
also uses Lee's archived \code{c11.tex:330--982}, especially the
continuous-coefficient matrix theorem. This adapter does not supply
LC54's actual smooth normal-flow data for a finite metric.

\begin{lemma}[LFR22: the finite surface types and the flat cover]
\label{lem:collapse-finite-complete-surface-types}
A complete connected orientable $C^m$ surface, $m\geq4$, with
nonnegative Gaussian curvature is homeomorphic to one of
$S^2,T^2,\mathbb R^2,S^1\times\mathbb R$. In the compact case
these are smooth types on its compatible carrier, and the torus
metric is flat. A complete flat surface has Euclidean universal
Riemannian cover even in this finite category.
\end{lemma}
\begin{proof}
For a noncompact surface use LFR21. For a compact surface use
LFR17, which passes Gauss--Bonnet on the original finite metric and
uses inherited smooth surface classification. In the torus case
the nonnegative continuous curvature has zero integral and vanishes
everywhere.

To verify the last assertion, give the universal cover its lifted
metric. It is complete: a finite-length Cauchy tail projects into
a sufficiently small evenly covered metric ball and stays in one
sheet. At a point of this cover the exponential map is defined
on the entire Euclidean tangent plane. Flatness makes every radial
Jacobi field with initial value zero the parallel translate of $tv$.
Consequently the differential of the exponential preserves the
inner product at every tangent vector; it is a local isometry.
A local isometry with complete domain covers a connected target:
lift a piecewise smooth path successively in local inverse charts;
a finite endpoint is reached because the lift has the same finite
length and is Cauchy. The same argument on short normal balls gives
the evenly covered neighborhoods. Since the target universal cover
is simply connected, this covering from $\mathbb R^2$ is one-to-one.
Lengths, and hence distances, agree in both directions. The argument
uses the finite Jacobi equation and at least $C^1$ exponential maps,
not a smoothness upgrade of the metric.
\end{proof}

\begin{theorem}[LFR23: an endpoint interval gives an actual disk band]
\label{thm:collapse-finite-endpoint-disk-band}
Let $(Z,k,z_0)$ be a complete connected orientable $C^m$ surface,
$m\geq4$, with nonnegative curvature. Write $r=d(z_0,\cdot)$.
Suppose there is a map
\[
 q:\overline B(z_0,10)\longrightarrow[0,10],\qquad q(z_0)=0,
\]
of distortion at most $\delta$, with $\delta$-dense image.
For sufficiently small universal $\delta>0$ the following hold.
For $1/2\leq r(x)\leq37/4$, the inward unit minimizing-direction
set $V_x(z_0)$ has diameter at most $2\sqrt{600\delta}$.
There is a smooth outward field on a neighborhood of this closed
band, with norm less than two and pairing less than $-3/4$ with
EVERY inward direction. For every $a\in[1,9]$,
$\overline B(z_0,a)$ is a topological closed disk. Its boundary is
a connected Lipschitz circle, and the distance levels in this
range are related by the hitting-time homeomorphisms of the SAME
outward field.
\end{theorem}
\begin{proof}
Assume initially $\delta<10^{-3}$. Coverage of the endpoint $10$
gives a point of distance greater than $9.5$ from $z_0$. Trim a
minimizing segment to obtain $x^+$ with $r(x^+)=9.5$. For a point
$x$ in the stated band put $a=r(x)$ and $b=d(x,x^+)$. Distortion
and the based condition give
\[
 |b-(9.5-a)|\leq3\delta,\qquad
 0\leq a+b-9.5\leq3\delta,\qquad b>0.24.
\]
Fix ANY unit minimizing direction $w$ from $x$ to $x^+$. For
EVERY inward unit direction $v$, triangle comparison and the
Euclidean cosine formula give
\[
 1+\langle v,w\rangle
 \leq\frac{(a+b)^2-9.5^2}{2ab}<300\delta,
 \qquad \|v+w\|^2<600\delta.
\tag{LFR23.1}
\]
Thus any two inward directions are within $2\sqrt{600\delta}$.
This is an all-choice comparison, not convergence of a selected
angle. Degenerate triangles have the same limiting inequality.
After decreasing $\delta$, the negative of any inward direction
has pairing less than $-7/8$ with all the others. Upper
semicontinuity of the direction sets extends each such vector to
a smooth local field with norm less than two and pairings less
than $-3/4$. A finite partition of unity patches these fields
near the compact band. Convexity preserves both inequalities.
Slightly larger and smaller radii provide an open neighborhood of
the closed band throughout this construction.

First variation implies that distance strictly increases along
the field, with lower rate $3/4$. Flow segments between fixed
levels stay in a compact annulus; they exist until reaching the
target level, in time bounded by the level difference divided by
$3/4$. The hitting time is unique and continuous. In each smooth
flow box it is Lipschitz, because distance is locally Lipschitz
and the increase in the flow direction has a fixed positive
lower bound. The distance level is therefore a Lipschitz graph
in that flow box. This proves that each level is a compact
one-manifold and that the entire band is its product with an
interval. No smoothness of the distance function at its cut locus
has been used.

The level $r=5$ is connected. Indeed two of its points have
distance at most $3\delta$, since each of their $q$-values is
within $\delta$ of $5$. A minimizing segment joining them stays
in $4<r<6$. Project this path to the $r=5$ level by the
continuous hitting-time map. This is a path joining the two
points in the level. Thus this level, and every level in $[1,9]$,
is one circle, not a disjoint union of nearby circles.

Each closed ball in that range is a compact connected surface
with this single boundary: radial minimizing paths connect its
interior, and the flow boxes supply boundary half-charts. It is
proper, since $x^+$ lies outside. Apply LFR22. In a plane or a
sphere such a domain is a disk, by the inherited Jordan separation
theorem. In a cylinder its boundary cannot be an essential circle,
as both sides of an essential circle are noncompact. A nullhomotopic
circle has a disk side and a noncompact side, so compactness again
selects the disk. In a torus a separating circle bounds a disk;
the only nondisk possibility for the closed ball is its complementary
once-punctured torus. These are statements about embedded compact
surface domains with a SINGLE locally flat boundary, not arbitrary
subsets of the four surfaces.

Exclude that last possibility for the radius-one ball. If it is
a once-punctured torus, its complement $U$ is an open disk and
contains the band $1<r<9$. The torus metric is flat by LFR22.
Lift $U$ to its Euclidean universal cover; the lift is one-to-one
since $U$ is simply connected. Choose $x$ with $r(x)=5$. Every
path of length less than four from $x$ stays in $U$. Thus the
lift contains the Euclidean unit ball about its lift of $x$:
otherwise a radial segment of length at most one would first
exit through a lift of the boundary of $U$, contradicting that
distance bound. Three vertices of an equilateral triangle of
side $1/2$ centered there project into $B(z_0,6)$ with pairwise
AMBIENT distances $1/2$. To check the last word, any shorter
minimizer stays within distance two of $x$, hence in $U$; its
lift joins the same endpoints and cannot shorten a Euclidean
segment. An interval cannot contain approximations to these
three distances with distortion $\delta<1/6$: the largest
interval distance is the sum of the other two, forcing
$1/2\leq3\delta$. This contradiction excludes the torus case.
Finally the common band collars identify the radius-one ball
with every ball of radius in $[1,9]$, by adding or reparametrizing
its collar. All are disks.
\end{proof}

This expands KL3.12, printed25--26/PDF20--21. Its connectedness
step now uses short connecting paths and flow projection. Its flat
torus exclusion uses an embedded lifted disk and checks ambient
distances, so it does not infer metric noncollapse from a local
isometry alone. A pointed GH bound gives the displayed based map
after the usual controlled increase of the error; its constants
are chosen before any noncollapse or derivative-bound data.

\begin{theorem}[LFR24: a smooth edge-model core on the finite metric]
\label{thm:collapse-finite-edge-model-core}
Fix $0<\varepsilon<1/100$. For all sufficiently small endpoint
errors $\delta=\delta(\varepsilon)$ in LFR23 and EVERY later
$0<\mu<1/100$, there is a function $h$ smooth on a neighborhood of $\overline B(z_0,9)$
in the compatible carrier, with the following properties:
\[
 h=0\text{ near }\overline B(z_0,1/2),\qquad
 0\leq h\leq2\text{ where }r\leq1,
\]
\[
 |h-r|<\mu,\quad
 \|\nabla_k h+v\|<\varepsilon\quad
 (2.1\leq r\leq9,\ v\in V_x(z_0)).
\tag{LFR24.1}
\]
For every $s\in[3,6]$ the WHOLE sublevel in this domain,
\[
 D_s=\{x\in B(z_0,9):h(x)\leq s\},
\]
is a smooth compact disk, containing $\overline B(z_0,s-\mu)$
and contained in $B(z_0,s+\mu)$. The SAME smooth outward field
can be used with $dh(V)>1/2$ on the collar $2.1\leq r\leq8$.
The metric is still $k$.

At any physical scale $\Delta>0$, apply this to $\Delta^{-1}Z$.
On the product domain
$(-6\Delta,6\Delta)\times B_Z(z_0,9\Delta)$ the map
$(t,z)\mapsto(t,\Delta h(z))$ is smooth in the compatible product
carrier, is transverse to the interior and boundary of
$\{0\}\times(-\infty,4\Delta]$, and has disk preimage there.
Its ENTIRE inverse image of
$[-4\Delta,4\Delta]\times(-\infty,4\Delta]$ is enclosed in
the interior of the compact set
\[
 Q=[-4.5\Delta,4.5\Delta]\times\overline B_Z(z_0,5\Delta).
\tag{LFR24.2}
\]
\end{theorem}
\begin{proof}
Use LFR02 with target $\{z_0\}$, the open annulus
$1/2<r<37/4$, and compact smoothing set
$3/4\leq r\leq9.1$. Choose $\delta$ so that LFR23's direction
diameter is below LFR02's $\theta(\varepsilon)$ throughout a
slightly larger annulus. This choice is independent of $\mu$.
It gives $F$ with $|F-r|<\mu$, with
$\operatorname{Lip}(F-r)<\varepsilon$, smooth near that compact
annulus, and $\|\nabla_kF+v\|<\varepsilon$ there for every
inward direction. Outside the smoothing annulus $F=r$.

Choose a smooth nondecreasing real function $\psi$ with
$\psi(u)=0$ for $u\leq1$, $\psi(u)=u$ for $u\geq2$, and
$0\leq\psi(u)\leq2$ for $1\leq u\leq2$. Put $h=\psi\circ F$. Where $F$ need not be smooth near the core it is
less than one, so $h$ is locally constant. Every transition point
lies in the smoothing annulus. Thus $h$ is genuinely smooth
on a neighborhood of the closed radius-nine ball, while it agrees with $F$ for
$r\geq2.1$. This retains (LFR24.1) on the entire boundary collar;
it does not claim a small value error for $h-r$ in the constant
core. Also $h\leq2$ wherever $F\leq2$, and hence all levels
$s\in[3,6]$ satisfy $h=s$ if and only if $F=s$.

The value bound gives the two closed-ball inclusions and places
the entire sublevel strictly inside radius $6.01<9$. It is closed
inside a compact ball of slightly larger radius, hence compact.
On its boundary $\|\nabla_kF\|>1-\varepsilon$, so that boundary
is a smooth regular level. For LFR23's same field, $|V|<2$ and
$\langle V,v\rangle<-3/4$, and consequently
\[
 dF(V)>3/4-2\varepsilon>1/2.
\]
The outward flow identifies the region between the topological
circle $r=2$ and $F=s$ with a collar: it starts with $F<2.01<s$,
reaches $F>s$ before radius seven, and crosses each intermediate
level exactly once. Thus $D_s$ is the radius-two topological disk
of LFR23 with a collar attached, and is itself a topological disk.
It is already a smooth compact surface with boundary in the
compatible carrier. The inherited smooth classification with
boundary, specifically Hirsch9.3.7, now makes it diffeomorphic
to the standard closed disk ($\chi=1$, one boundary component).
No homeomorphism is silently declared to be a diffeomorphism.

For the scaled product, its first coordinate has differential
$dt$ everywhere. On the level $h=4$, the second differential is
nonzero and annihilates the real direction; the two are independent.
These facts give transversality to both strata of the stated
half-line. Its preimage is $\{0\}\times D_4$. The whole slab
preimage is $[-4\Delta,4\Delta]\times D_4$, which lies strictly
inside (LFR24.2) by $4+\mu<5$ and $4<4.5$.
\end{proof}

\paragraph{What this supplies to LC84.}
The complete finite surface, its disk topology, a globally smooth
function on the required model neighborhood, all-direction collar
estimates, both transversality conditions and a compact model
enclosure are now constructed. The physical normalization is
\emph{first} $\Delta^{-1}Z$, as required to use KL3.12 in KL9.21.
The endpoint threshold is independent of the volume lower bound
$w'$ and of the derivative-bound function. The sharper source
comparison tolerance may depend on them, preserving KL9.23's order.

This does not yet complete the edge packet. One must produce the
strong/weak edge-set comparisons and the source smoothed distance,
retain the derivative of the variable scale, compare ALL nearest-set
directions in the source collar, and enclose the ENTIRE source
fiber before the finite interpolation. Model-side enclosure alone
does not exclude additional source components. The source distance
is smooth only on its collar in KL9.12; it needs the same constant
core modification before applying a globally smooth-map isotopy
theorem. That modification must preserve the tested boundary
level and actual domain. The zero-model normal-flow/classification
consumer and general Alexandrov/Tits inputs remain distinct.


\subsection{Transferring the finite edge disk to the source}
\label{sec:collapse-finite-edge-transfer}

This subsection proves the source-manifold part of the disk argument
from an explicit coarse-border hypothesis. It does not declare that
hypothesis to follow merely from the definition of an edge point.
In particular, the strong/weak edge-set producer in KL9.7, its changing
scale and its buffers remain a separate obligation below. The metric
model, original comparison embedding, smoothed function and entire
source fiber are kept distinct throughout.

\paragraph{A buffered coarse-border hypothesis.}
Fix $\Delta\geq1$ and $0<\tau<10^{-4}$. In a complete connected pointed smooth
manifold $(M,p)$, let $A$ be closed and contain $p$. A coarse-border
chart for $(M,p,A)$ means an ACTUAL map
\[
 Q=(u,b):B(p,200\Delta)\longrightarrow\mathbb R\times[0,\infty),
 \qquad Q(p)=(0,0),
\]
with distance distortion at most $\tau\Delta$, whose image is
$\tau\Delta$-dense in
$[-100\Delta,100\Delta]\times[0,100\Delta]$, and satisfying
\[
 b(a)\leq\tau\Delta\qquad(a\in A\cap B(p,190\Delta)).
\]
For every $|s|\leq100\Delta$ there must also be an actual
$a\in A\cap B(p,190\Delta)$ such that
\[
                       |Q(a)-(s,0)|\leq\tau\Delta.             \tag{LFR25.1}
\]
The second clause is coverage of the border, not just containment
near it. The conditions are imposed on $A=\overline{E'}$ in the
application, with positive slack inherited from larger buffers.
They do not assert that $\overline{E'}$ satisfies a strict edge
predicate. All distances in this hypothesis are ambient distances.

\begin{lemma}[LFR25: a coarse border controls every nearest direction]
\label{lem:collapse-edge-closed-set-directions}
For a coarse-border chart, write $a(x)=d_A(x)$. Then
\[
 |a(x)-b(x)|\leq2\tau\Delta\qquad(x\in B(p,70\Delta)).       \tag{LFR25.2}
\]
Suppose also $\sec\geq-\kappa^2$ on $B(p,1000\Delta)$ and
$\kappa\Delta\leq1/100$. If $x\in B(p,30\Delta)$ and
$\Delta/2\leq a(x)\leq12\Delta$, there is $y\in B(p,70\Delta)$
such that
\[
 |d(x,y)-a(x)|\leq4\tau\Delta,
 \qquad |a(y)-2a(x)|\leq7\tau\Delta.                         \tag{LFR25.3}
\]
Every minimizing unit direction $w$ from $x$ to this $y$ obeys
\[
 \|v+w\|^2\leq200\tau\quad(v\in V_x(A)),
 \qquad \operatorname{diam}V_x(A)\leq2\sqrt{200\tau}<30\sqrt\tau.
                                                                    \tag{LFR25.4}
\]
The local Riemannian comparison input retains AC02's qualification.
No noncollapse or upper curvature derivative bound enters these constants.
\end{lemma}
\begin{proof}
Properness gives a nearest point of the nonempty closed set $A$.
For $x\in B(p,70\Delta)$ it lies in $B(p,140\Delta)$, since
$p\in A$ and $a(x)\leq d(p,x)$. Its height is at most $\tau\Delta$,
so distortion gives $a(x)\geq b(x)-2\tau\Delta$. Also
$|u(x)|<71\Delta$. The border coverage clause at $s=u(x)$ gives
$a(x)\leq b(x)+2\tau\Delta$. This proves (LFR25.2); no nearest
point outside the tested border buffer has been ignored.

Use rectangle coverage to choose $y$ with
$|Q(y)-(u(x),2b(x))|\leq\tau\Delta$. Here
$|u(x)|<31\Delta$ and $b(x)<13\Delta$, so its pointed distance
is less than $70\Delta$. Distortion gives
$|d(x,y)-b(x)|\leq2\tau\Delta$. Applying (LFR25.2) at both
points gives (LFR25.3).

Let $z\in A$ be any nearest point of $x$, put
$a=a(x)$, $\ell=d(x,y)$ and $c=d(z,y)$. Then
\[
 0\leq a+\ell-c\leq11\tau\Delta,\quad
 a\geq\Delta/2,\quad \ell\geq(1/2-4\tau)\Delta,
 \quad a+\ell\leq(24+4\tau)\Delta.
\]
The same hyperbolic calculation as LC26, now without requiring equal
radial sides, gives for the comparison angle $\vartheta$ at $x$
\[
 1+\cos\vartheta
 =\frac{\cosh(\kappa(a+\ell))-\cosh(\kappa c)}
        {\sinh(\kappa a)\sinh(\kappa\ell)}
 \leq\frac{2(a+\ell-c)(a+\ell)}{a\ell}<100\tau.
\]
The bound uses $\kappa(a+\ell)<1$ and $\cosh1<2$.
At $\kappa=0$ the Euclidean cosine formula gives the same bound.
All chosen triangle vertices and sides lie in $B(p,100\Delta)$.
For the local comparison domain $B(p,1000\Delta)$ use $D=100\Delta$:
all side lengths are less than $D$, and the intrinsic closed
$2D$-balls at triangle and radial-subtriangle vertices are complete.
Indeed an intrinsic Cauchy sequence has an ambient limit of radius
at most $300\Delta<1000\Delta$, and local intrinsic/ambient metric
equality gives intrinsic convergence, as in LC27. Lower comparison
therefore applies to EVERY chosen pair of minimizing arms. It implies
$\langle v,w\rangle\leq-1+100\tau$, which proves (LFR25.4).
One fixed choice of $w$ bounds the diameter of all inward directions.
\end{proof}

\begin{lemma}[LFR26: comparison of inward directions without an exact set limit]
\label{lem:collapse-edge-model-direction-transfer}
Suppose the original finite-model embeddings
$j_i:N=\mathbb R\times Z\supset U_i\to M_i$ have the growing-domain
coverage, pointed distance convergence and $C^1$ pullback metric
convergence of LFR14--LFR15. The complete product metric is $C^m$,
$m\geq2$, nonnegative, with basepoint $(0,z_0)$; put $r(t,z)=d_Z(z,z_0)$.
Each $M_i$ has a coarse-border chart $Q_i=(u_i,b_i)$ with the SAME
$\Delta,\tau$, a closed set $A_i$, and the curvature bound of LFR25.
Assume $u_i j_i\to t$ uniformly on fixed compact sets, and, for
$0<h<1/100$,
\[
 \limsup_i\sup_{B_N((0,z_0),100\Delta)}|b_i j_i-r|\leq h\Delta.
                                                               \tag{LFR26.1}
\]
On the compact collar $|t|\leq10\Delta$, $2.5\Delta\leq r\leq6.5\Delta$,
EVERY source nearest-set direction $v_i\in V_{j_i(x)}(A_i)$ and
EVERY product radial inward direction $v\in V_x(\mathbb R\times\{z_0\})$
satisfy the uniform bound
\[
 \limsup_i\sup_{x,v_i,v}\|(d j_i|_x)^{-1}v_i-v\|\leq40\sqrt{h+\tau}.
                                                               \tag{LFR26.2}
\]
The constants are independent of noncollapse. No Hausdorff convergence
of $A_i$ to the model axis, and no continuous choice of an arm, is assumed.
\end{lemma}
\begin{proof}
Take any sequence of tested points and directions. Their source points
are in $B(p_i,30\Delta)$; their distances to $A_i$ lie in
$[\Delta/2,12\Delta]$ eventually, by (LFR25.2) and (LFR26.1).
Use LFR25's outward points $y_i$ and choose any minimizing outward
arms. The actual source-ball coverage lifts $y_i$ and those entire
segments to one compact model ball, with radius less than $100\Delta$.
After extraction the tested points converge to $x=(t,z)$ and the
lifted endpoints to $y$. Uniform convergence of $u_i j_i$ and the
choice of $y_i$ give $|t(y)-t(x)|\leq\tau\Delta$.
Writing $\ell=d_N(x,y)$, the value estimates give
\[
 |\ell-r(x)|\leq(h+6\tau)\Delta,\qquad
 r(y)\geq2r(x)-(3h+\tau)\Delta.
\]
For $o_x=(t(x),z_0)$, therefore,
\[
 0\leq r(x)+\ell-d_N(o_x,y)\leq(4h+7\tau)\Delta,
 \qquad r(x)\geq2.5\Delta,\quad\ell>2\Delta.
\]
Nonnegative Riemannian comparison in the SAME finite metric gives,
for every radial inward direction $v$ and every limiting outward
minimizing direction $w$,
$\|v+w\|^2\leq20(h+\tau)$. Here one uses the Euclidean
cosine formula with the displayed defect; the coefficient follows
from $1/r(x)+1/\ell<0.9/\Delta$.

For clarity, the outward source directions converge on a subsequence
in the required tangent trivialization. The proof of LFR18 applies
up to its final uniqueness step: lifted speeds are bounded, the
Christoffel symbols converge uniformly, and the coordinate geodesic
integral equations give $C^1$ convergence including the starting
velocities. The limit is minimizing by distance convergence. It need
not be unique, and the preceding comparison holds for every possible
limit. Radial model directions likewise have convergent subsequences,
by compactness of their endpoints and continuity of the finite geodesic
flow. LFR25 bounds every pulled-back inward source direction within
$\sqrt{200\tau}+o(1)$ of the negative outward direction. Thus the
limsup of the tested difference is at most
$\sqrt{200\tau}+\sqrt{20(h+\tau)}<40\sqrt{h+\tau}$.
A sequence violating the uniform bound would admit precisely such an
extraction, a contradiction. Upper curvature/volume bounds are used
only in producing the finite model and its convergent embeddings,
not in the displayed angular tolerance.
\end{proof}

\begin{lemma}[LFR27: scaled distance smoothing with a smooth constant core]
\label{lem:collapse-edge-scaled-core-smoothing}
In LFR25 choose $0<\varepsilon,\mu<1/100$ and require
$30\sqrt\tau<\theta(\varepsilon)$ from LFR02. Let $\rho$ be smooth
and positive on $B(p,100\Delta)$, with $\rho(p)=1$ and
$|d\rho|\leq\Lambda$, where $\lambda=100\Delta\Lambda<1/100$.
There exists a global nonnegative Lipschitz function $F$ with
\[
 \|F-d_A\|_{\infty,M}<\mu\Delta,
 \qquad \operatorname{Lip}_M(F-d_A)\leq\varepsilon.             \tag{LFR27.1}
\]
It is smooth on a neighborhood of
$C=\overline B(p,20\Delta)\cap\{.75\Delta\leq d_A\leq10.5\Delta\}$,
and there $\|\nabla F+v\|<\varepsilon$ for EVERY $v\in V_x(A)$.
Put $\eta=F/\rho$. On this collar,
\[
 d\eta=\rho^{-1}dF-F\rho^{-2}d\rho,
 \qquad \|d\eta-dF\|<150\Delta\Lambda.                        \tag{LFR27.2}
\]
Fix a smooth nondecreasing profile $\psi$, constant zero on
$(-\infty,1]$, equal to the identity on $[2,\infty)$, with
$0\leq\psi\leq2$ on $[1,2]$ and $\operatorname{Lip}\psi\leq4$.
Then $H=\Delta\psi(\eta/\Delta)$ is smooth on
\[
              B(p,20\Delta)\cap\{d_A<10.25\Delta\}.          \tag{LFR27.3}
\]
For EVERY $s\geq2\Delta$, $H\leq s$ if and only if $\eta\leq s$;
for $s>2\Delta$ their level sets and differentials agree there.
Thus this modification preserves the ACTUAL tested sublevel and
boundary, including all components.
\end{lemma}
\begin{proof}
Apply LFR02 on
$U=B(p,30\Delta)\cap\{\Delta/2<d_A<12\Delta\}$ with the displayed
compact $C$. LFR25 verifies its all-choice direction hypothesis.
The output equals $d_A$ off a compact subset of $U$; since
$d_A>\Delta/2$ there and $\mu<1/100$, it is nonnegative everywhere.
The value tolerance is chosen after the direction tolerance.
On $B(p,100\Delta)$, $|\rho-1|\leq\lambda$, by integrating along
minimizing radial segments. On $C$, $|dF|<1+\varepsilon$ and
$|F|<(10.5+\mu)\Delta$. The quotient rule gives
\[
 \|d\eta-dF\|\leq
 \frac{\lambda(1+\varepsilon)}{1-\lambda}
 +\frac{(10.5+\mu)\Delta\Lambda}{(1-\lambda)^2}
 <150\Delta\Lambda.
\]
This is a derivative estimate, not a consequence of value closeness.
In particular the term with $d\rho$ has not been dropped.

A profile as stated is obtained by smoothing a nondecreasing transition
from zero to the identity inside $(1,2)$, matching it near both ends;
its derivative can be kept below four. Where $F/\rho$ is not known
smooth near the inner core in (LFR27.3), it is less than $\Delta$,
so $H$ is locally constant. Every transition point lies in the smooth
collar $C$ with strict margins. Points farther out in (LFR27.3) also
lie in that collar. This proves smoothness on the whole domain.
The scalar profile itself has
$\psi(q)\leq s/\Delta\Longleftrightarrow q\leq s/\Delta$
for $s\geq2\Delta$; for $s>2\Delta$ it is the identity near that
level. This proves both set equality and equality of differentials.
For a finite collection of such local collars in the SAME manifold,
one can apply LFR02 once to their compact union in the union of their
open smoothing regions, choosing the value tolerance below all the
finitely many prescribed physical tolerances. The resulting $F$ is
shared by those packets; separate local mollifications are not forced.
For charts with different normalizing scales, perform this one smoothing
in the physical metric. The local value tolerances become
$\mu\Delta\rho(p)$; the norm of the gradient error is unchanged
when both lengths and distance functions are rescaled.
\end{proof}

\begin{theorem}[LFR28: the whole source edge disk bundle]
\label{thm:collapse-finite-source-edge-packet}
There are universal positive constants $\tau_0,\sigma_0,\varepsilon_0,
\mu_0,\lambda_0$ with the following parameter order. Choose
$\Delta\geq1$, $0<\tau\leq\tau_0$, $0<\sigma\leq\sigma_0$,
$0<\varepsilon\leq\varepsilon_0$, $0<\mu\leq\mu_0$ and
$100\Delta\Lambda\leq\lambda_0$, meeting LFR27's direction condition.
These choices precede $K\geq5$, noncollapse data $r_0,v>0$ and
finite derivative bounds $A(R)>0$. For all sufficiently small
$\beta>0$, depending also on these later data, suppose:
\begin{enumerate}[leftmargin=*]
\item $(M,g,p)$ is a complete connected oriented smooth three-manifold
with $\Vol B(p,r_0)\geq v$, $|\nabla^q\Rm|\leq A(R)$ on
$B(p,R)$ for $q\leq K$, $R<\beta^{-1}$, and
$\sec\geq-\beta^2$ on $B(p,\beta^{-1})$;
\item an ACTUAL $(1,\beta)$-splitting $\alpha=(u,v_Y)$ is supplied;
its first coordinate is the first coordinate of a coarse-border
chart $Q=(u,b)$ for a closed set $A\ni p$;
\item $\rho$ satisfies LFR27, and $F$ is any one of its outputs
with the displayed tolerances; put $\eta=F/\rho$.
\end{enumerate}
LFR19 constructs a tangential coordinate $f$ on $B(p,100\Delta)$,
with value error $\mu\Delta$, quality $\sigma$ and outer test
radius $1000\Delta$. For EVERY such $f$ the map
\[
 f:\Omega=\{x\in B(p,100\Delta):|f(x)|<4\Delta,
                         \ \eta(x)\leq4\Delta\}
                      \longrightarrow(-4\Delta,4\Delta)       \tag{LFR28.1}
\]
is a proper smooth trivial bundle with connected fiber the closed
smooth disk. Its boundary is exactly $\eta=4\Delta$, where $df,d\eta$
have rank two. There is one complete nonnegative $C^{K-1}$ product
model $\mathbb R\times Z$ and an actual $C^K$ embedding of the
whole buffered model domain from LFR24, containing the ENTIRE source
closed slab $|f|\leq4\Delta$, $\eta\leq4\Delta$. The interpolation
of the model map with $(f,H)$ uses that same embedding and the same
compact set (LFR24.2), at all times and all levels $|s|\leq4\Delta$.
Any prescribed positive finite-buffer metric/distance comparison
tolerance may be imposed by further reducing $\beta$.

There is also a smooth compactly supported source cutoff equal to
one on $|f|\leq8\Delta$, $\eta\leq8\Delta$, with closed support
strictly inside $|f|<9\Delta$, $\eta<9\Delta$ in its stated domain.
The endpoint and coarse-border tolerances are independent of noncollapse.
The theorem constructs the transfer from the listed geometric data;
it does not yet produce (LFR25.1) from the edge predicates.
\end{theorem}
\begin{proof}
Fix the profile $\psi$ of LFR27 once. One sufficient choice is
$\sigma_0=10^{-10}$ and
$\varepsilon_0=\mu_0=\lambda_0=10^{-8}$. Decrease $\tau_0>0$
so that $\tau_0<10^{-8}$, $h=40\sqrt{\tau_0}<1/100$,
$40\sqrt{h+\tau_0}<10^{-5}$, and $20\tau_0$ meets LFR24's
endpoint threshold with gradient error $10^{-8}$.
LFR27's smaller direction condition is explicitly retained for
any subsequently smaller $\varepsilon$. These are choices before
$v,r_0,A$; none uses an injectivity radius.

If the claimed uniform $\beta$ threshold fails, take counterexamples
$\beta_i\to0$, with all these other data fixed, and their supplied
charts, functions and closed sets. LFR14--LFR16 give a complete
nonnegative product $N=\mathbb R\times Z$, with $Z$ an oriented
$C^{K-1}$ surface, and the SAME actual embeddings $j_i$, metric
convergence, source-ball coverage, and $u_i j_i\to t$ on every fixed
buffer. All following assertions hold on one common tail.

\emph{1. The endpoint model and height comparison.}
For $x=(t,z)$ in a fixed model ball of radius $100\Delta$, compare
the pointed-distance identities
\[
 d_N(p,x)^2=t^2+r(x)^2,\qquad
 |Q_i(j_i(x))|^2=(u_i j_i(x))^2+(b_i j_i(x))^2.
\]
Their distance error is at most $\tau\Delta+o(1)$.
The difference of squares, $u_i j_i\to t$, and nonnegativity of
both heights give
\[
 \limsup_i\sup |b_i j_i-r|\leq20\sqrt\tau\,\Delta.
                                                               \tag{LFR28.2}
\]
For example the squared-height error is at most
$201\tau\Delta^2+o(1)$, and $|a-b|^2\leq|a^2-b^2|$ for
$a,b\geq0$. Hence LFR26 applies with $h=40\sqrt\tau$.

Here is also an actual endpoint approximation, without assuming
continuity of any $Q_i$. Take a finite $\tau\Delta$-net of
$\overline B_Z(z_0,10\Delta)$, including $z_0$, and extract limits
of the finitely many heights $b_i j_i(0,z)$ there. On that slice,
pairwise distances differ from differences of heights by at most
$\tau\Delta$ in the limit, and the basepoint height is zero.
Assign every surface point the limiting height of a nearest net point,
then truncate to $[0,10\Delta]$. Before truncation the distortion is
at most $3\tau\Delta$; the heights exceed $10\Delta$ by at most
$\tau\Delta$. Thus the truncated map has distortion less than
$6\tau\Delta$ and sends $z_0$ to zero.
For coverage, choose rectangle lifts of $(0,s)$ at a finite
$\tau\Delta$-mesh with $0\leq s\leq(10-4\tau)\Delta$.
Their pointed radii are at most $s+2\tau\Delta<10\Delta$.
Coverage by $j_i$ and compactness give limiting product points
$(t,z)$ with $|t|\leq\tau\Delta$ and $r(z)<10\Delta$.
Compare them with a net point within $\tau\Delta$ of $z$.
Distortion and the rectangle error put that net point's limiting
height within $4\tau\Delta$ of $s$; truncation and the final
endpoint interval cost at most another $6\tau\Delta$.
The map is therefore $20\tau\Delta$-dense. After scaling distances
by $\Delta^{-1}$ it meets LFR23 with error $20\tau$.
This argument controls the marked endpoint as well as the interval
shape; no limit of discontinuous maps is asserted.

Use the construction in LFR24's proof: LFR02 gives a model
smoothing $G$ of $r$ with GLOBAL value error $\mu\Delta$
and gradient error $10^{-8}$ on its collar. Use the
SAME profile to put $H_N=\Delta\psi(G/\Delta)$.
Its whole $4\Delta$ sublevel is the compact smooth disk there.

\emph{2. Source derivatives on the same finite model.}
Apply LFR19 with $L=100\Delta$, $T=1000\Delta$ to construct $f_i$.
On the compact cylinder $[-6\Delta,6\Delta]\times\overline B_Z(z_0,9\Delta)$,
test its differential against the actual source endpoints corresponding
to $(t+200\Delta,z)$. Their separation tends to $200\Delta>L$,
and both pointed radii lie in the tested outer ball. LFR18 and the
LFR19 all-direction estimate give, uniformly,
\[
 |f_i j_i-t|<2\mu\Delta,\qquad
 \|d(f_i j_i)-dt\|^2\leq4\sigma+\sigma^2+o(1)<10^{-6}.        \tag{LFR28.3}
\]
As in LFR20, a fixed adapted quality gives a bound, not convergence
of that error to zero.

On $|t|\leq6\Delta$, $2.5\Delta\leq r\leq6.5\Delta$, combine
LFR26 with (LFR27.1). Pullback covectors are compared in the SAME
metric $g_N$; $j_i^*g_i\to g_N$ accounts for the norm change.
Every inward model direction is within $40\sqrt{h+\tau}+o(1)$
of every pulled-back source nearest direction. The quotient-rule
error is at most $150\Delta\Lambda$, by LFR27. The model smoothing
has its own $10^{-8}$ error against those same radial directions.
It follows that
\[
 \|d(\eta_i j_i)-dG\|<10^{-3}                                \tag{LFR28.4}
\]
on that collar. The estimates were obtained from actual minimizing
arms, not by differentiating a GH approximation or a value bound.

Throughout the larger cylinder, (LFR25.2), (LFR28.2), the smoothing
value errors, and $|\rho_i-1|\leq\lambda_0$ imply
\[
 |\eta_i j_i-G|<\Delta/1000,
 \qquad |H_i j_i-H_N|<\Delta/100.                             \tag{LFR28.5}
\]
The latter uses $\operatorname{Lip}\psi\leq4$, including in the
constant core. The image cylinder lies in $B(p_i,20\Delta)$ and
has $d_{A_i}<10.25\Delta$, so $H_i$ is smooth on a neighborhood
of its WHOLE image. Equations (LFR28.4) compare $H_i j_i$ with
$H_N$ wherever their values are near $4\Delta$, since both profiles
are the identity there. No smoothness of $F_i$ at $A_i$ is required.

\emph{3. Enclosure of every source component.}
If $x\in B(p_i,100\Delta)$, $|f_i(x)|\leq4\Delta$ and
$\eta_i(x)\leq4\Delta$, then
\[
 |u_i(x)|\leq(4+\mu)\Delta,\quad
 a_i(x)\leq\{4(1+\lambda_0)+\mu\}\Delta.
\]
For (LFR25.2) one must first put $x$ in its radius-$70\Delta$
domain. Use the border condition directly: a nearest point has
radius at most $100\Delta+a_i(x)<105\Delta<190\Delta$.
Distortion thus gives $b_i(x)\leq a_i(x)+2\tau\Delta$.
Consequently
\[
 d(p_i,x)\leq\Delta\sqrt{(4+\mu)^2+
                \{4(1+\lambda_0)+\mu+2\tau\}^2}+\tau\Delta
                         <6\Delta.                           \tag{LFR28.6}
\]
This uses a tested point of the ORIGINAL chart, before assuming it
lies in the model image. Strict-radius coverage now puts the entire
source slab in $j_i(B_N(p,6.1\Delta))$. Its preimage has
$|t|<4.01\Delta$ by the original value error and uniform
$u_i j_i\to t$ on the radius-seven ball, and $r<4.02\Delta$
by the height comparison. This step does not assume in advance
that the preimage lies in the cylinder used in (LFR28.3). Thus it is in
$U=(-6\Delta,6\Delta)\times B_Z(z_0,9\Delta)$.
The same bound makes the source closed slab compact, since it is
closed inside a compact ball of radius strictly between $6\Delta$
and $100\Delta$. In particular there are no unexamined extra fibers.

\emph{4. The compact interpolation and smooth disk bundle.}
On $U$ set
\[
 \mathcal F_s(x)=\big((1-s)t+s f_i j_i,
                       (1-s)H_N+s H_i j_i\big),\qquad0\leq s\leq1.
\]
For every target first-coordinate value $|c|\leq4\Delta$, all
inverse images of $X_c=\{c\}\times(-\infty,4\Delta]$ lie
in the interior of the SAME
$Q=[-4.5\Delta,4.5\Delta]\times\overline B_Z(z_0,5\Delta)$.
Indeed (LFR28.3) bounds $|t|$ by $4.01\Delta$; (LFR28.5)
bounds $H_N$ by $4.01\Delta$, and hence $r<4.02\Delta$.
On the inverse image of $\partial X_c$, instead
$3.98\Delta<r<4.02\Delta$. There (LFR28.4) applies.
The first covector is within $10^{-3}$ of $dt$ and the second
within $10^{-3}$ of $dH_N$, while
$dH_N(\partial_t)=0$ and $\|dH_N\|>1-10^{-8}$.
Evaluating on $\partial_t$ and the unit model gradient direction
of $H_N$ gives a two-by-two determinant greater than $0.99$.
Thus the maps are transverse to BOTH strata. In the interior the
first covector alone is nonzero. The jointly $C^K$ family satisfies
LFR03; its isotopy is $C^{K-1}$ on this same compact buffer.

The model fiber is a smooth disk by LFR24. At the source end,
$H_i\leq4\Delta$ is EXACTLY $\eta_i\leq4\Delta$, and both
boundary functions agree near $4\Delta$, by LFR27. Step3 identifies
the image with the entire source fiber. LFR04 therefore gives its
smooth disk type, since $K-1\geq2$. This also proves connectedness.
On the source, $f_i$ is a smooth submersion on (LFR28.1), including
its boundary, and is proper by the compact closed slab. A smooth
field $X$ with $df_i(X)=1$, tangent to $\eta_i=4\Delta$, is
constructed by solving the two differential equations near the
boundary and the single one in the interior, and combining solutions
by a partition of unity. Properness traps its trajectories over each
compact subinterval. Its forward and reverse flows give a smooth
trivialization over $(-4\Delta,4\Delta)$.

\emph{5. A cutoff and the uniform conclusion.}
Choose smooth profiles with plateau $[-8\Delta,8\Delta]$ in the
first variable and $(-\infty,8\Delta]$ in the second, supported
respectively in $[-8.9\Delta,8.9\Delta]$ and $(-\infty,8.9\Delta]$.
Their product evaluated at $(f_i,\eta_i)$ is smooth: the second
profile is constant near the nonsmooth core, and its transition lies
in LFR27's smooth collar. Repeating (LFR28.6) with $9$ in place of
$4$ puts its closed support inside $B(p_i,13\Delta)$, strictly
inside the chart. Extension by zero is therefore smooth and compactly
supported, with the asserted strict support margins.

Any additional prescribed positive comparison tolerance holds on
the same finite buffers on a sufficiently late tail. All conclusions
have now been produced in every sufficiently late counterexample,
contradicting its defining failure. This proves the uniform threshold.
The separate lower-curvature and splitting smallness requirements
were combined into $\beta$ only for the statement: small splitting
error alone supplies neither curvature nor noncollapse. In KL6.10--6.16
the normalized curvature defect tends to zero with the ambient sequence
index; the sufficiently late index may depend on the already fixed
parameters. It is not the endpoint approximation parameter of KL9.23.
\end{proof}

\paragraph{What is now closed, and the retained producer boundary.}
LFR25--LFR28 supply distance regularization, variable-scale derivatives,
comparison of ALL inward directions, the smooth extension across the
core, ENTIRE source-fiber enclosure, both transversality conditions,
compact interpolation and the actual proper smooth disk bundle.
The finite-$C^{K-1}$ model metric is unchanged. Angular tolerances are
fixed before noncollapse; only the finite-model comparison threshold
uses the later volume and derivative data. Exact nearest-set convergence
is unnecessary for this argument.

To obtain the full LC84 packet from the original edge predicates one
must still construct the buffered coarse-border chart (LFR25.1), with
closure, domain restrictions, scale recentering and parameter order,
and the full rank-two adapted collar of KL9.20. The latter is stronger
than rank two at the tested disk boundary. One must also arrange the
common smoothing on the selected edge family; LFR27 explains how a
finite family shares the same function after its uniform geometric
hypotheses are verified. Strong-edge density/coverage from KL9.10
remains separate. No missing edge-set hypothesis is inferred from
this transfer theorem. The zero-model normal-flow/3D classification,
general Alexandrov comparison/splitting and Tits-cone obligations
remain as previously recorded; accepted PC foundation bindings remain
pending, with no fresh probe of the moving checkout.

KL9.7,9.12,9.17--9.23, printed50,52--56/PDF45,47--51, supplies the
source route; KL6.10--6.16, printed42--44/PDF37--39, supplies the
sequence/finite-compactness context. The construction above expands
the source's compactness argument with explicit buffers and retains
its distinction between endpoint tolerance and finite-model comparison.

\subsection{Producing the closed edge border and one shared distance function}
\label{sec:collapse-edge-border-producer}

The input left explicit in LFR25 is now produced from actual edge maps.
The argument is metric: it needs neither Alexandrov compactness nor
noncollapse. Curvature and finite differentiability enter only when
LFR25--LFR28 consume its output. In particular, the numerical quality
of this producer is chosen BEFORE the finite-model analytic data.

\paragraph{Source convention and parameter order.}
KL Definition9.2, printed49/PDF44, has $C\geq200\Delta$ (the page image,
not its OCR, verifies the non-strict sign). The retained LC84 and FC31
contracts use $C>200\Delta$. We keep those historical contracts intact
and give an explicit interval-extension adapter below. The proofs of
this subsection work for either edge convention; the witnesses produced
along the border satisfy the STRICT convention. No extra positive
margin in the given $C$ is assumed.

All approximations here use the actual whole-space pointed maps and
open-ball KL convention MC06, with coverage expressed as a distance
to the image. Actual witnesses are chosen with error less than twice
the stated error; an infimum is not assumed attained. Product metrics
are Euclidean product metrics. No continuity, properness or length
structure of a residual factor is assumed.

Fix $\Delta\geq1$ and $0<\tau<10^{-4}$. Write
$b=\beta_E$, $s=\sigma_E$, $b'=\beta_{E'}$, $s'=\sigma_{E'}$.
The following deliberately generous sufficient constraints are used:
\begin{align*}
 W&=\min\{\tau\Delta/10^9,\ (10^6\Delta)^{-1},\ 10^{-4}\},\\
 0<b',s'&<W,\qquad 0<s<10^{-5}\min\{b',s'\},\\
 0<\Lambda&<\min\{(10^6\Delta)^{-1},\ s'/(10^8\Delta^2)\},\\
 0<b&<10^{-5}\min\{s,b',s'\}.                  \tag{LFR29.1}
\end{align*}
Here $\rho:M\to(0,\infty)$ is $\Lambda$-Lipschitz in the physical
metric $d$, and $d_p=d/\rho(p)$. The normalized function
$\rho/\rho(p)$ is $\Lambda$-Lipschitz for $d_p$ as well. At an edge
point the supplied maps are
\[
 F_p=(u,v):(M,d_p,p)\longrightarrow\mathbb R\times(Y,y_0),
 \qquad G_p:(Y,y_0)\longrightarrow([0,C],0),
\]
with KL qualities $b,s$ for a strong point and $b',s'$ for a weak point.
Let $E,E'$ be the corresponding sets, using one fixed choice of the
strict or non-strict interval convention. Put $A=\overline{E'}$ in the
physical metric (the closure is unchanged by constant normalization).
The extra bound on $\Lambda$ depends on the EARLIER $s',\Delta$; it is
acyclic in KL15.4/FC44. It is a project sufficient constraint, not a
claim that KL9.7's displayed $\overline\Lambda(\epsilon,\Delta)$ already
states this stronger dependency. The strong splitting quality $b$ can
subsequently be decreased using noncollapse and derivative data; $s$,
$b',s'$ and $\Lambda$ need not be decreased with those later data.

\begin{lemma}[LFR29: four directions cannot be based on the half-plane border]
\label{lem:collapse-half-plane-four-point-obstruction}
There are no four points $q_1^+,q_1^-,q_2^+,q_2^-$ in
$\mathbb R\times[0,\infty)$ for which, with $0<e\leq10^{-3}$,
\[
 \big||q_i^\pm|-1\big|\leq e,\qquad
 \big||q_i^+-q_i^-|-2\big|\leq e,\qquad
 \big||q_1^\pm-q_2^\pm|-\sqrt2\big|\leq e.
\]
In the last condition all four independent choices of signs are tested.
\end{lemma}
\begin{proof}
For each opposite pair the parallelogram identity gives
\[
 |q_i^++q_i^-|^2
 \leq4(1+e)^2-(2-e)^2=12e+3e^2\leq15e.
\]
Both vertical coordinates are nonnegative, so each is at most
$\sqrt{15e}$. Each horizontal coordinate consequently has absolute
value between $1-17e$ and $1+e$, since its square is at least
$(1-e)^2-15e\geq1-17e$. Two of the four nonzero horizontal
coordinates have the same sign. Their horizontal difference is at
most $18e<20e$ and their vertical difference at most $\sqrt{15e}$.
Their squared distance is at most $400e^2+15e<1$. Every distinct
pair was required to have distance at least $\sqrt2-e>1$, a
contradiction. This is an elementary finite metric obstruction.
\end{proof}

\begin{lemma}[LFR30: an actual composite edge chart with buffered lifts]
\label{lem:collapse-composite-edge-chart}
Suppose $F=(u,v):(X,p)\to\mathbb R\times(Y,y_0)$ is a pointed KL
$b$-approximation and $G:(Y,y_0)\to([0,C],0)$ is a pointed KL
$s$-approximation. Set $\delta=b+s$ and $Q=(u,Gv)$. If
$S+10\delta<\min\{b^{-1},s^{-1}\}$, then $Q$ has distortion at
most $\delta$ on $B(p,S)$. Every $z=(t,h)\in\mathbb R\times[0,C]$
with $|z|\leq S-4\delta$ has an ACTUAL witness $x\in B(p,S)$ with
\[
 |Q(x)-z|<3\delta,\qquad d(p,x)<|z|+3\delta.       \tag{LFR30.1}
\]
For $z=(t,0)$ one may instead lift $(t,y_0)$ directly through $F$:
\[
 |F(a)-(t,y_0)|<2b,\qquad
 d(p,a)<|t|+3b,\qquad d_Y(v(a),y_0)<2b.            \tag{LFR30.2}
\]
Finally, if $C\geq200\Delta$ and $0<s<1/8$, the SAME point function
$G$, regarded as taking values in
$[0,\max\{C,200\Delta+s\}]$, is a KL $4s$-approximation.
Its new target has length strictly greater than $200\Delta$.
\end{lemma}
\begin{proof}
For $x\in B(p,S)$, $d_Y(v(x),y_0)\leq d(p,x)+b<s^{-1}$.
The product norm changes by at most the change in its second
component, so the two distortions add. For coverage choose $y$
with $|G(y)-h|<2s$ in the original tested source ball of $G$.
Its radius is less than $h+3s$. Thus $(t,y)$ has radius less than
$|z|+3s$ and is in the coverage range of $F$. Choose $x$ with
$|F(x)-(t,y)|<2b$. It has radius less than $|z|+3s+3b<S$.
Both $v(x)$ and $y$ are in the tested ball of $G$. First compare
$Q(x)$ with $(t,G(y))$, at cost at most $2b+s$, and then with
$(t,h)$, at cost less than $2s$. This proves (LFR30.1).
The margins $S+10\delta$ ensure all original coverage radii,
including their subtracted errors, rather than just distortion
radii, contain these targets. The same argument with the direct
$F$-target $(t,y_0)$ gives (LFR30.2).

For the last assertion the target extension has length at most $s$.
For a tested new target $h<(4s)^{-1}-4s$, replace it by
$h_0=\min\{h,C\}$ and use original coverage with error less than
$2s$. Its source witness has radius less than $h_0+3s<(4s)^{-1}$.
Its error from $h$ is less than $3s<4s$. Distortion on the smaller
source ball is at most $s$ and the basepoint is unchanged. This
adapter weakens endpoint quality by a fixed factor, with no new
strong-splitting dependency. It does not identify $C\geq200\Delta$
with $C>200\Delta$ without changing the approximation contract.
\end{proof}

\begin{lemma}[LFR31: changing scale and repairing the endpoint length]
\label{lem:collapse-edge-rescale-recenter}
Assume (LFR29.1) and strong maps at $p$ with $C\geq200\Delta$.
Let $a$ satisfy $d_p(p,a)\leq101\Delta$ and
$d_Y(v(a),y_0)<2b$. Then $a$ has actual weak edge maps in its OWN
normalization $d_a$, with endpoint length strictly greater than
$200\Delta$. In particular $a\in E'$ for either convention above.
\end{lemma}
\begin{proof}
Put $q=\rho(a)/\rho(p)$ and $\lambda=q^{-1}$. Lipschitz control gives
\[
 |q-1|\leq101\Delta\Lambda,
 \qquad |\lambda-1|<102\Delta\Lambda,
 \qquad \tfrac12<q,\lambda<2.
\]
On the residual factor use the metric $\lambda d_Y$ and basepoint
$v_a=v(a)$. Keep the actual point function $v$ and set
\[
 F_a(x)=(\lambda(u(x)-u(a)),v(x)),\qquad d_a=\lambda d_p.
\]
Its distortion is at most $\lambda b<2b<b'$. The original $F$-ball
contains the ball
\[
 B_{d_p}(a,\lambda^{-1}(b')^{-1}).
\]
Its target coverage contains the corresponding ball about $F(a)$.
The sufficient margin $101\Delta+2/b'+2<1/b$ follows from (LFR29.1).
For a target of new radius $r<(b')^{-1}-b'$, original coverage gives
error less than $2\lambda b$ and a witness of new distance from $a$
less than $r+3\lambda b<(b')^{-1}$. Thus this SAME recentered map
is a KL $b'$-approximation, with its exact new basepoint.

It is not legitimate simply to keep endpoint length $\lambda C$:
that length may be below $200\Delta$. Define instead
\[
 L=\max\{\lambda C,\ 200\Delta+s'/100\},\qquad
 G_a(z)=\begin{cases}0,&z=v_a,\\ \lambda G(z),&z\ne v_a.
                    \end{cases}                              \tag{LFR31.1}
\]
This is a whole-space map into $[0,L]$, based at zero. Its endpoint
extension and its one-point change obey
\begin{align*}
 0\leq L-\lambda C
 &\leq200\Delta(1-\lambda)_++s'/100
   <20400\Delta^2\Lambda+s'/100<s'/50,\\
 \theta:=\lambda G(v_a)&\leq\lambda(2b+s)<2(2b+s).
                                                               \tag{LFR31.2}
\end{align*}
There is no upper bound on $C$ in this calculation. On the new
$(s')^{-1}$-ball the original points have radius at most
$2b+2/s'<s^{-1}$. Distortion is bounded by
$\lambda s+2\theta\leq6s+8b<s'/10$.

For a tested $c\in[0,L]$, $c<(s')^{-1}-s'$, put
$c_0=\min\{c,\lambda C\}$. Original $G$-coverage at $c_0/\lambda$
is available: its radius is at most $2/s'$, far inside $s^{-1}-s$.
Choose a witness $z$ with original image error less than $2s$.
Its new source radius about $v_a$ is less than
\[
 c_0+\lambda(2b+3s)<c+s'/10<(s')^{-1},
\]
and its image error after the one-point repair is less than
$(L-\lambda C)+2\lambda s+\theta<s'/10$. This proves the required
KL $s'$-approximation. Formula (LFR31.1) supplies the STRICT endpoint
length without strengthening the original hypothesis on $C$.
No continuity of either approximation has been used.
\end{proof}

\begin{theorem}[LFR32: the closed coarse border is supplied by the edge predicates]
\label{thm:collapse-edge-coarse-border-producer}
Let $(M,d)$ be a complete connected smooth Riemannian manifold, with
positive $\Lambda$-Lipschitz $\rho$ and the edge predicates above,
satisfying (LFR29.1). For every $p\in E$, the actual composite
$Q_p=(u,b_p)=(u,G_pv)$ on $B_{d_p}(p,200\Delta)$, and the SAME
closed set $A=\overline{E'}$, satisfy every part of LFR25.1 with
parameters $\Delta,\tau$. In particular, $p\in A$,
\[
 b_p(a)\leq\tau\Delta\quad
 (a\in A\cap B_{d_p}(p,190\Delta)),                         \tag{LFR32.1}
\]
and every $(t,0)$ with $|t|\leq100\Delta$ has a witness in
$E'\cap B_{d_p}(p,190\Delta)$ of image error less than $\tau\Delta$.
All constants are independent of noncollapse and curvature derivatives.
\end{theorem}
\begin{proof}
Write $\delta=b+s$. LFR30 with $S=200\Delta$ gives distortion
$\delta<\tau\Delta$ and actual $\tau\Delta$-coverage of the required
rectangle $[-100\Delta,100\Delta]\times[0,100\Delta]$, whose points
have norm at most $100\sqrt2\Delta<142\Delta$. The lower bound on
$C$ contains this rectangle. LFR31 with $a=p$ proves $p\in E'$.

\emph{Exclude weak edge points high above the border.}
Suppose $x\in E'\cap B_{d_p}(p,191\Delta)$ and
$b_p(x)\geq\tau\Delta/2$. Put $R=\tau\Delta/100$.
The four points $Q_p(x)\pm Re_1,Q_p(x)\pm Re_2$ lie in
$\mathbb R\times[0,C]$: their lower height is positive, and their
upper height is less than $192\Delta<C$. Their norms are below
$192\Delta$. LFR30 supplies four witnesses $x_i^\pm$ of image
error less than $3\delta$, each in $B_{d_p}(p,193\Delta)$.
Using distortion on $B(p,200\Delta)$, their distances from $x$
differ from $R$ by at most $4\delta$, their opposite-pair distances
from $2R$ by at most $7\delta$, and all cross-pair distances from
$\sqrt2R$ by at most $7\delta$.

Let $q_x=\rho(x)/\rho(p)$. Here $|q_x-1|<191\Delta\Lambda$, so
$1/2<q_x<2$. Use the weak composite $Q_x$ in its OWN scale,
based at $x$. Its distortion on the ball of radius $4R/q_x$
is at most $b'+s'$ by LFR30; all four witnesses lie in that ball.
The large weak domains required by LFR30 follow from $b',s'<W$.
The four vectors
\[
 z_i^\pm=\frac{q_x}{R}Q_x(x_i^\pm)
          \in\mathbb R\times[0,\infty)
\]
therefore satisfy the forbidden estimates of LFR29 with
\[
 e=\frac{7\delta+2(b'+s')}{R}<10^{-3}.
\]
This contradiction proves $b_p(x)<\tau\Delta/2$ for every such
weak point. The proof uses both pointed charts on actual points;
it does not take a purported continuous limit of their compositions.

\emph{Pass to closure without assuming continuity of the chart.}
For $a\in A\cap B(p,190\Delta)$ choose $x_n\in E'$ tending to $a$.
Eventually $x_n\in B(p,191\Delta)$. The coarse Lipschitz estimate
from the original chart gives
\[
 |b_p(a)-b_p(x_n)|\leq d_p(a,x_n)+\delta.
\]
Consequently $b_p(a)\leq\tau\Delta/2+\delta<\tau\Delta$,
which is (LFR32.1). The extra radius-one buffer is used here.

\emph{Supply border witnesses which are themselves weak edge points.}
For $|t|\leq100\Delta$, use the DIRECT $F_p$-lift in (LFR30.2).
It has $d_p(p,a)<101\Delta$ and $d_Y(v(a),y_0)<2b$.
LFR31 proves $a\in E'$ at its own scale. Moreover,
$|Q_p(a)-(t,0)|\leq2b+s<\tau\Delta$: compare the product norm
of $F_p(a)-(t,y_0)$ and use $G_p$-distortion from its basepoint.
This is the other direction of the coarse-border condition, with
witnesses in the required larger ball. Every clause of LFR25.1,
including the closed-set and exact first-coordinate requirements,
has now been obtained from the edge maps themselves.
\end{proof}

\begin{corollary}[LFR33: a finite verified edge family shares one distance function]
\label{cor:collapse-finite-edge-family-shared-smoothing}
Let $p_1,\ldots,p_N\in E$, $N\geq1$, satisfy LFR32 with common
$\Delta,\tau$ meeting LFR28's universal tolerances. Assume $M$ is oriented and
three-dimensional, $\rho$
is smooth and that every normalized pointed manifold at $p_i$ meets
LFR28's analytic hypotheses with common $K,r_0,v,A(R)$ and sufficiently
small lower-curvature and strong-splitting errors. Here $b$ is chosen
after the common analytic data and BEFORE fixing the family, still
subject to (LFR29.1); impose also $100\Delta\Lambda\leq\lambda_0$.
Then there is ONE nonnegative global
Lipschitz function $F$ approximating the physical distance $d_A$ such
that $\eta=F/\rho$ serves simultaneously in all these edge disk
packets. Each has the actual proper trivial smooth $D^2$-bundle and
supported cutoff of LFR28. The SAME constant-core function
$H=\Delta\psi(\eta/\Delta)$ supplies their smooth interpolation maps.
The assertion is for a finite verified family, not a proof that these
points cover the edge region or satisfy the full adapted-collar contract.
\end{corollary}
\begin{proof}
LFR32 supplies the common CLOSED set $A=\overline{E'}$ and all
normalized coarse-border charts with the original tangential
coordinates. Let $\rho_i=\rho(p_i)$. In physical units set
\begin{align*}
 U_i&=B(p_i,30\Delta\rho_i)
             \cap\{\tfrac12\Delta\rho_i<d_A<12\Delta\rho_i\},\\
 C_i&=\overline B(p_i,20\Delta\rho_i)
             \cap\{.75\Delta\rho_i\leq d_A\leq10.5\Delta\rho_i\}.
\end{align*}
Each $C_i$ is compact by completeness and $C_i\subset U_i$.
LFR25 gives the all-nearest-direction diameter bound on each $U_i$;
the same assertion at a point does not depend on which scale verifies
it. Thus it holds on $U=\bigcup_i U_i$. The compact union
$C=\bigcup_i C_i$ lies in $U$, which is separated from $A$ by at least
$\tfrac12\Delta\min_i\rho_i>0$.

Apply LFR02 ONCE in the physical metric, with direction error
$\varepsilon_0$ and global value error
$e<\mu_0\Delta\min_i\rho_i$, decreasing $e$ below half that last
separation if needed. It gives $F=d_A$ off $U$, smooth near $C$,
with global Lipschitz-difference bound $\varepsilon_0$ and the
all-choice derivative estimate there. It is nonnegative: off $U$
it is a distance, and on $U$ its error is smaller than $d_A$.
In the normalization at $p_i$, put $F_i=F/\rho_i$ and
$\widehat\rho_i=\rho/\rho_i$. Then
\[
 \|F_i-d_{p_i,A}\|_{\infty,M}<\mu_0\Delta,
 \quad\operatorname{Lip}_{d_{p_i}}(F_i-d_{p_i,A})\leq\varepsilon_0,
 \quad \frac{F_i}{\widehat\rho_i}=\frac{F}{\rho}=\eta.
\]
Simultaneous constant rescaling of the metric and distance function
preserves the covector error norm. Hence each normalization satisfies
LFR27, including the full quotient derivative and smooth constant core.
LFR28 now applies at every $p_i$, with this SAME $\eta$ and $H$ and
its own tangential adapted function. Its finite-model threshold is
uniform because the analytic input data were common. The potentially
smaller physical value error depending on the finite family is chosen
only when smoothing; it does not alter any earlier geometric tolerance.
No compatibility of separately chosen smoothings is asserted or needed.
\end{proof}

\paragraph{Updated boundary of the edge packet.}
The coarse-border producer previously left open in LFR25--LFR28 is
now written, including closure, actual witnesses, changing scale and
endpoint repair. LFR33 supplies a common smoothing for a finite family
whose hypotheses are verified. The full rank-two adapted collar of
KL9.20 and the strong-edge density/coverage statements of KL9.10 are
still distinct obligations. Boundary rank two and the disk topology
alone do not imply them. The zero-model actual normal-flow map and
three-dimensional classification, and general Alexandrov comparison,
splitting and Tits-cone inputs, remain as recorded. All results here
are written project arguments; inherited analytic/geometric operations
and accepted-release PC/Lean bindings remain explicit, with no new
probe of the moving checkout and no final-goal completion claim.

The source comparison is KL3.2 (printed21/PDF16),4.10 and its proof
(printed30--31/PDF25--26),9.2--9.7 (printed49--50/PDF44--45), and
15.4 with its ordering explanation (printed87/PDF82). The source
9.7 limit argument is replaced here by a finite four-point obstruction
and explicit maps. This expands its border assertion for the buffered
closed-set consumer; it is not an additional author erratum.

\subsection{The full adapted collar for the original edge coordinates}
\label{sec:collapse-full-edge-adapted-collar}

The disk boundary at height $4\Delta$ is only one part of the
edge consumer. KL9.20 requires the ORIGINAL pair $(f,\eta)$ to be
adapted throughout the band with height in $[\Delta/10,10\Delta]$.
In particular LFR27's smoothing region, starting at $3\Delta/4$,
does not by itself supply that band. We first enlarge the smoothing
region, and then prove all three adapted-coordinate conditions.
We never replace the original pair by the reference coordinates
used to estimate it.

In this subsection the metric is normalized at the strong edge center
$p$, so $\rho(p)=1$. A coarse-border chart means exactly LFR25.1,
with its actual map $Q=(u,b_Q)$, closed set $A$ and buffers. In the
application LFR32 supplies it with $A=\overline{E'}$ and with $u$
the first coordinate of the original rank-one splitting. Distances
and unit vectors below refer to the metric explicitly indicated.

\begin{lemma}[LFR34: distance smoothing down to the full collar]
\label{lem:collapse-edge-low-collar-smoothing}
In LFR25 the outward-point and triangle argument also applies when
$\Delta/50\leq d_A(x)\leq12\Delta$. It gives
\[
 \|v+w\|^2\leq4600\tau,
 \qquad \operatorname{diam}V_x(A)<140\sqrt\tau.       \tag{LFR34.1}
\]
Retain its curvature assumptions. Given $0<\varepsilon<1/100$ and
$0<\mu<10^{-3}$, choose $140\sqrt\tau<\theta(\varepsilon)$.
Then there is a global nonnegative Lipschitz $F$ satisfying LFR27.1,
smooth on a neighborhood of
\[
 C^+=\overline B(p,20\Delta)\cap
           \{\Delta/25\leq d_A\leq10.5\Delta\},
\]
with $\|\nabla F+v\|<\varepsilon$ for every nearest direction there.
Its difference from $d_A$ has compact support in
\[
 U^+=B(p,30\Delta)\cap\{\Delta/50<d_A<12\Delta\}.
\]
If $\rho$ meets LFR27's scale assumptions, the full quotient estimate
$\|d(F/\rho)-dF\|<150\Delta\Lambda$ holds on $C^+$ as well.
The SAME profile $H=\Delta\psi(F/(\Delta\rho))$ still has LFR27's
smooth-core and exact-sublevel conclusions. Thus this output meets
the displayed smoothing contract consumed by LFR28.
\end{lemma}
\begin{proof}
The value estimate $|d_A-b_Q|\leq2\tau\Delta$ and the outward
lift $(u(x),2b_Q(x))$ in LFR25 used no lower bound $\Delta/2$.
For the resulting triangle use instead
\[
 a\geq\Delta/50,\quad
 \ell\geq(1/50-4\tau)\Delta,\quad a+\ell-c\leq11\tau\Delta.
\]
The same buffered comparison gives
\[
 1+\cos\vartheta\leq
 22\tau\left(50+\frac1{1/50-4\tau}\right)<2300\tau
       \quad(\tau<10^{-4}).
\]
Hence (LFR34.1) holds for ALL minimizing arms. The comparison
buffers and the bound on $\kappa(a+\ell)$ are unchanged.
Apply LFR02 to $C^+\subset U^+$ with global value tolerance
$\mu\Delta$ and Lipschitz-difference tolerance $\varepsilon$.
Nonnegativity follows from $\mu<10^{-3}<1/50$ on $U^+$ and
from equality with the distance elsewhere. The quotient calculation
in LFR27 uses only the upper bound $10.5\Delta$ on $d_A$ and so
applies unchanged. Its profile transitions occur in $C^+$; below
that region the profile is constant. The profile identities, and
hence the ACTUAL sublevel and boundary, are unchanged. We use the
displayed output contract of LFR27, not its particular narrower
choice of smoothing support.
\end{proof}

\paragraph{The scale-$100$ adapted contract.}
For a metric $d_x$ and an actual pointed two-splitting
$\Phi_x:(M,d_x,x)\to(\mathbb R^2,0)$, a smooth map $J$ on
$B_{d_x}(x,100)$ is adapted of quality $\gamma$ when its ambient
Lipschitz constant is at most $1+\gamma$, its image is within
Hausdorff distance $100\gamma$ of $B(J(x),100)$, and
\[
 \left|DJ_y(w)-\frac{\Phi_x(z)-\Phi_x(y)}{d_x(y,z)}\right|
 <\gamma                                                        \tag{LFR35.1}
\]
for EVERY $y\in B_{d_x}(x,100)$, $z\in B_{d_x}(x,100/\gamma)$
with $d_x(y,z)>100$, and every initial $d_x$-unit minimizing
direction $w$. This is exactly Definition4.21 rescaled as in
Remark4.35: the image error, tested length and outer test radius
all rescale. Translation by $J(x)$ is allowed, and $\Phi_x$ may
likewise be translated to mark its target at $J(x)$.

\begin{lemma}[LFR35: a plane comparison and reference coordinates on a buffer]
\label{lem:collapse-edge-plane-reference}
Fix $0<\gamma<1/100$, $0<\beta_2<\gamma/1000$ and $\iota>0$.
There is $\Delta_0(\gamma,\beta_2,\iota)\geq10^6$ with the
following property for every $\Delta\geq\Delta_0$. Choose $\tau$
sufficiently small in terms of these data, and then $\kappa$ sufficiently
small, with $\sec\geq-\kappa^2$ on $B(p,10000\Delta)$.
For any coarse-border chart and any
\[
 x\in B(p,15\Delta),\qquad
 .09\Delta\leq d_A(x)\leq10.1\Delta,
 \qquad .99\leq q\leq1.01,
\]
put $d_x=d/q$. The ACTUAL map
\[
 \Phi_x(y)=\frac{Q(y)-Q(x)}q\quad(y\in B(p,200\Delta)),
                                                               \tag{LFR35.2}
\]
extended by zero elsewhere, is a pointed KL $\beta_2$-approximation
to $\mathbb R^2$, hence a two-splitting with singleton residual factor.
Choose actual lifts $a_{j\pm}$ of $Q(x)\pm s e_j$, where
$s=\Delta/40$, with $Q$-error at most $\tau\Delta$.
There is a smooth reference map $\chi_x$, defined near
$\overline B_{d_x}(x,300)$, with $\chi_x(x)=0$, such that:
\begin{enumerate}[leftmargin=*]
\item on $B_{d_x}(x,100)$ it is adapted to $\Phi_x$ of quality
  $\gamma/4$, with $|\chi_x-\Phi_x|<100\gamma/32$;
\item on $B_{d_x}(x,300)$ its differential Gram matrix satisfies
  $\|D\chi_x(D\chi_x)^*-I\|_{\rm op}<\gamma/4$; and
\item on that same larger ball, for every minimizing $d_x$-unit
  direction $v_{j+}$ toward $a_{j+}$,
  $\|\nabla_{g_x}(\chi_x)_j-v_{j+}\|_{g_x}<\iota$.
\end{enumerate}
All choices are uniform over the indicated $x,q$, and require no
noncollapse, upper curvature bounds or finite-model theorem.
\end{lemma}
\begin{proof}
First require $\Delta>10^5/\beta_2$. From LFR25.2 the height
$b_Q(x)$ lies in $[.08\Delta,10.2\Delta]$ if $\tau$ is small.
Both the four anchor targets and all targets
$Q(x)+qz$, $|z|\leq\beta_2^{-1}$, lie strictly inside
$[-100\Delta,100\Delta]\times[0,100\Delta]$. Every tested
source ball also lies inside $B(p,200\Delta)$.
Distortion of (LFR35.2) is at most $\tau\Delta/q$.
For $|z|<\beta_2^{-1}-\beta_2$, rectangle coverage gives a witness
with image error at most $\tau\Delta/q$ and source radius at most
$|z|+2\tau\Delta/q<\beta_2^{-1}$, once
$4\tau\Delta<\beta_2$. Thus BOTH coverage and distortion have the
KL quality claimed; the extension off the buffer is irrelevant to
these tests. Rescaling lengths and coordinates by $1/100$ gives
KL quality $100\beta_2<\gamma/4$: its coverage targets have the
stronger original radius margin $10000\beta_2$, so the original
witnesses still fit. Thus the SAME map also meets the approximation
quality required by the rescaled adapted definition. Anchor lifts lie in $B(p,20\Delta)$ by their pointed
norms and distortion.

We spell out the reference-coordinate construction, since the anchor
scale is chosen before the later splitting error. Work in $d_x$,
put $s_x=s/q$, and write $z_y=\Phi_x(y)$. On every fixed ball
about $x$, all distances among tested points and anchors differ,
by at most a fixed multiple of $\tau\Delta/q$, from the corresponding
Euclidean distances among $z_y$ and $\pm s_x e_j$. Set
\[
 h_j(y)=d_x(x,a_{j+})-d_x(y,a_{j+}).
\]
For $|z_y|\leq R+1$ the Euclidean calculation is
\[
 0\leq |s_xe_j-z_y|-(s_x-(z_y)_j)
       \leq\frac{(R+1)^2}{2(s_x-R-1)}.                 \tag{LFR35.3}
\]
Take $R>400/\gamma+1000$, fixed BEFORE increasing $\Delta$.
For $y\in B_{d_x}(x,300)$ and $z\in B_{d_x}(x,400/\gamma)$,
$d_x(y,z)>100$, expansion of the Euclidean cosine formula yields
\[
 \cos\widetilde\angle_y(a_{j+},z)
   =\frac{(z_z)_j-(z_y)_j}{d_x(y,z)}+O_R(s_x^{-1})+o_\tau(1).
                                                               \tag{LFR35.4}
\]
Here, with $\Delta$ fixed, $o_\tau(1)$ tends to zero as $\tau\to0$;
the lower tested length $100$ and anchor length $s_x-R-1$ bound
all denominators away from zero. The analogous negative-anchor
formula has the opposite leading term. Opposite-anchor comparison
cosines tend to $-1$, and mixed-index comparison cosines to zero,
uniformly on the radius-$400$ ball, first as $\Delta\to\infty$
and then as $\tau\to0$. These facts follow from the same squared
Euclidean distances; no limit of actual angles is asserted.
For fixed $\Delta$ these Euclidean formulas are uniform limits of
the curvature-$-q^2\kappa^2$ cosine formulas as $\kappa\to0$.

Apply buffered lower comparison to every choice of arms. Opposite
arms have sum tending to zero, so the upper cosine estimates for
both signs give the lower estimates as well. This gives small
diameter of ALL directions toward each positive anchor, the
all-direction version of (LFR35.4), and an almost identity Gram
matrix for the two positive unit directions. Triangles lie inside
$B(p,100\Delta)$; intrinsic comparison buffers are contained in
$B(p,10000\Delta)$ as in LFR25. No comparison at a cut point is
replaced by differentiability of distance there.

Apply LFR02 to each anchor distance, with smoothing compact set
$\overline B_{d_x}(x,300)$ in the radius-$400$ open ball, and
Lipschitz-difference error $\epsilon$ chosen smaller than
$\min\{\iota,\gamma/1000\}$. Write the smooth distances as $D_j$
and set $(\chi_x)_j=D_j(x)-D_j$. LFR02 gives the asserted
comparison with EVERY positive-anchor direction. It also gives
$|(\chi_x)_j-h_j|\leq\epsilon d_x(x,\cdot)$.
Choose the preceding geometric tolerances smaller if necessary.
The positive-direction Gram estimate then gives item2, and
(LFR35.3) gives the stated value bound on the radius-$100$ ball.
First variation, or the LFR02 gradient estimate together with the
all-direction comparison above, gives (LFR35.1) with error $\gamma/4$
for the full outer radius $400/\gamma$. Components are combined in
the Euclidean norm, so their individual errors are chosen below
$\gamma/(16\sqrt2)$.

The Gram estimate on the radius-$300$ ball bounds the operator norm
by $\sqrt{1+\gamma/4}<1+\gamma/8$. Every minimizing segment
between two points of $B_{d_x}(x,100)$ stays in that larger ball;
integration therefore gives the AMBIENT Lipschitz bound. For image
coverage lift $(1-\gamma/16)z$ through $\Phi_x$, $|z|<100$.
The original rectangle witnesses just constructed lie inside the
radius-$100$ source ball after decreasing $\tau$; their images,
combined with the value bound, are within $100\gamma/4$ of $z$.
Conversely the Lipschitz bound and $\chi_x(x)=0$ put every image
point within $100\gamma/4$ of that target ball. This proves every
adapted condition, without requiring the ball to be convex or the
approximation map to be continuous.
\end{proof}

\begin{lemma}[LFR36: the prescribed pair is pinned to the same two directions]
\label{lem:collapse-edge-prescribed-gradient-pinning}
Retain a coarse-border chart and an actual $(1,b)$-splitting
$\alpha=(u,v)$ with the SAME $u$. Let $\rho(p)=1$, let $\rho$
be smooth positive and $\Lambda$-Lipschitz, and let $F$ be an LFR34
output with errors $\varepsilon,\mu$. Put $\eta=F/\rho$.
Let $f$ be any LFR19 output with $L=100\Delta$, $T=1000\Delta$,
value error $\mu\Delta$ and quality $\sigma$.
For any prescribed $\iota>0$, the following holds by first taking
$\Delta$ large, then $\tau,\sigma,\varepsilon,\mu,\Delta\Lambda$
small, and then $b,\kappa$ small, with
$\sec\geq-\kappa^2$ on $B(p,10000\Delta)$ and the original
splitting domain large enough. Every point of the ENTIRE tested band
\[
 x\in B(p,100\Delta),\qquad |f(x)|\leq10\Delta,
             \Delta/10\leq\eta(x)\leq10\Delta            \tag{LFR36.1}
\]
lies in $B(p,15\Delta)$ and satisfies
$.09\Delta\leq d_A(x)\leq10.1\Delta$.
Set $q=\rho(x)$, $d_x=d/q$. On $B_{d_x}(x,300)$ the same $F$
and $\eta$ are smooth, and, for the anchors of LFR35 and EVERY
choice of their minimizing unit directions,
\[
 \|\nabla_{g_x}f-v_{1+}\|_{g_x}<\iota,
 \qquad \|\nabla_{g_x}\eta-v_{2+}\|_{g_x}<\iota.       \tag{LFR36.2}
\]
The order ``large $\Delta$, then small errors'' can be imposed
simultaneously with LFR35 and LFR29.1. These estimates are independent
of noncollapse and upper curvature derivative bounds.
\end{lemma}
\begin{proof}
Require $100\Delta\Lambda,\mu<10^{-6}$. On $B(p,100\Delta)$
we have $|\rho-1|\leq100\Delta\Lambda$. The band inequalities and
global value estimate give $.09\Delta<d_A(x)<10.1\Delta$.
For a nearest $z\in A$, initially $d(p,z)<111\Delta<190\Delta$;
we have NOT yet used the smaller-ball estimate LFR25.2. The ORIGINAL
chart gives $b_Q(x)\leq d_A(x)+2\tau\Delta$ and
$|u(x)|\leq(10+\mu)\Delta$, hence
\[
 d(p,x)\leq\Delta
 \sqrt{(10+\mu)^2+(10.1+2\tau)^2}+\tau\Delta<15\Delta.
                                                               \tag{LFR36.3}
\]
Now $.99<q<1.01$. For $y\in B_{d_x}(x,300)$ we have
$d(x,y)<600$; taking $\Delta\geq10^6$ puts the WHOLE such ball
strictly inside $C^+$ and $B(p,16\Delta)$. Thus the shared smoothing
is smooth on this larger buffer, and its all-choice gradient estimate
applies there.

\emph{The vertical coordinate.}
Work temporarily in the $p$-normalized metric. Let $z\in A$ be any
nearest point to $y$, let $a=d_A(y)$, let $\ell=d(y,a_{2+})$ and
$c=d(z,a_{2+})$. The original chart, its border containment, and the
anchor lift give
\[
 a\geq\Delta/25,\qquad \ell\geq\Delta/80,\qquad
 0\leq a+\ell-c\leq1200+9\tau\Delta.                  \tag{LFR36.4}
\]
For the last bound use
$a\leq b_Q(x)+600+3\tau\Delta$,
$\ell\leq s+600+3\tau\Delta$, and
$c\geq b_Q(x)+s-3\tau\Delta$.
The lower bound on $\ell$ follows from
$s-600-3\tau\Delta\geq\Delta/80$ for the chosen large $\Delta$
and small $\tau$. The same comparison calculation as LFR34 yields
\[
 \|v_A+v_{2+}\|_g^2
 \leq504000/\Delta+3780\tau,                         \tag{LFR36.5}
\]
where both vectors are unit for the $p$-normalized metric and
$v_A$ is any inward nearest direction. Here $\kappa\Delta$ is
small enough that $\kappa(a+\ell)<1$. Thus every inward direction
is nearly opposite to every positive vertical-anchor direction.

LFR34 controls $\nabla F+v_A$ and retains the FULL derivative
\[
 d\eta=\rho^{-1}dF-F\rho^{-2}d\rho.
\]
Consequently $\nabla\eta$ differs from $-v_A$ by at most
$\varepsilon+150\Delta\Lambda$ in this metric. Passing to $g_x$
changes a covector norm by the factor $q$, whereas a unit tangent
vector changes by the factor $q$. The resulting additional difference
is bounded by $|q-1|$. The second inequality in (LFR36.2) follows
from (LFR36.5), after the indicated ordered choices.

\emph{The tangential coordinate; retain the long tested length.}
The local anchor $a_{1+}$ is too close for the length test $>100\Delta$
in LFR19. For each $y$ instead use the ACTUAL original splitting
$\alpha$ to lift
\[
 (u(y)\pm T_0,v(y)),\qquad T_0=200\Delta,
\]
with error less than $2b$, to points $l_\pm$. They lie in
$B(p,250\Delta)\subset B(p,1000\Delta)$ and satisfy
$|d(y,l_\pm)-T_0|<3b$. For every minimizing direction $w_+$
toward $l_+$, LFR19 therefore gives
\[
 df_y(w_+)\geq1-\sigma-O(b/\Delta),\qquad
 \|df_y\|\leq1+\sigma.
\]
The elementary identity
$\|\nabla f-w_+\|^2=\|\nabla f\|^2+1-2df(w_+)$ bounds this
squared norm by $4\sigma+\sigma^2+O(b/\Delta)$.
It remains to compare that long direction with $v_{1+}$.

Write $D=u(a_{1+})-u(y)$ and $r=d(y,a_{1+})$. The coarse chart
implies
\[
 D\geq s-600-2\tau\Delta,\quad
 r\leq s+600+3\tau\Delta,\quad r\geq\Delta/80.
\]
Thus $D/r\to1$ uniformly as $\Delta\to\infty$ and then
$\tau\to0$. In the EXACT target product of $\alpha$, the squared
distance from $(u(y)-T_0,v(y))$ to $\alpha(a_{1+})$ is
\[
 T_0^2+2T_0D+d(\alpha(y),\alpha(a_{1+}))^2.
\]
Source distortion and the lift error change the squared sides by
$O(b\Delta)$, so the Euclidean comparison cosine between $l_-$
and $a_{1+}$ at $y$ equals $-D/r+O(b/\Delta)$. The comparison
cosine between $l_+$ and $l_-$ tends to $-1$ as well. For fixed
$\Delta$, their hyperbolic comparison cosines converge uniformly to
these quantities as $\kappa\to0$; adjacent lengths are bounded
below by $\Delta/80$. Lower comparison gives
$\|w_-+v_{1+}\|\to0$ and $\|w_-+w_+\|\to0$ for ALL choices.
This proves the tangential estimate in the original metric, and
$|q-1|\to0$ gives it in $g_x$.

All of these long triangles have vertices in $B(p,250\Delta)$,
sides of length less than $1000\Delta$, and all comparison buffers
inside $B(p,10000\Delta)$. The intrinsic-completeness check is the
same buffered argument as LFR25, now with $D=1000\Delta$.
This records the larger curvature domain required by the long arms.
Only distance formulas and lower comparison were used; convergence
of actual angles, or differentiation of $Q$, was not assumed.
\end{proof}

\begin{lemma}[LFR37: buffered stability of all adapted conditions and rank]
\label{lem:collapse-buffered-vector-adapted-stability}
Let $\Phi$ be a fixed actual pointed two-splitting on a complete
pointed Riemannian manifold $(M,x)$ and let $0<\gamma<1/100$.
Suppose $\chi,J$ are smooth near $\overline B(x,300)$,
$\chi(x)=0$, and $\chi|_{B(x,100)}$ is adapted to $\Phi$ of
quality $\gamma/4$. Assume on $B(x,300)$ that
\[
 \|D\chi(D\chi)^*-I\|<\gamma/4,
 \qquad \|DJ-D\chi\|_{\rm op}<\gamma/100.
\]
Then $J|_{B(x,100)}$ has rank two and is adapted to the SAME
$\Phi$ of quality $\gamma$, centered at $J(x)$.
\end{lemma}
\begin{proof}
Integrate $D(J-J(x)-\chi)$ along minimizing segments. For endpoints
in $B(x,100)$, those segments lie in $B(x,300)$, so the difference
is ambient $\gamma/100$-Lipschitz and has norm less than $\gamma$
on the radius-$100$ ball. This gives an ambient Lipschitz bound
$1+\gamma/4+\gamma/100<1+\gamma$ for $J$, and image Hausdorff
error at most $25\gamma+\gamma<100\gamma$ from $B(J(x),100)$.
For every tested minimizing direction, the derivative errors add to
less than $\gamma/4+\gamma/100<\gamma$; the required outer radius
$100/\gamma$ is contained in the reference outer radius $400/\gamma$.
The smallest singular value of $D\chi$ is at least
$\sqrt{1-\gamma/4}$. For every unit target covector $a$,
\[
 \|(DJ)^*a\|\geq\sqrt{1-\gamma/4}-\gamma/100>0.
\]
Thus $DJ$ is onto and has rank two. These are ALL three adapted
conditions plus rank, not just a pointwise submersion estimate.
\end{proof}

\begin{theorem}[LFR38: the original edge pair has the full adapted collar]
\label{thm:collapse-full-edge-collar-packet}
Fix $0<\gamma<1/100$ and $0<\beta_2<\gamma/1000$.
Choose $\Delta$ sufficiently large depending on these two earlier
parameters. Choose the edge smoothing and tangential qualities,
weak-edge qualities, endpoint quality and scale Lipschitz bound
sufficiently small in the order FC44, meeting LFR29.1. Finally
choose the strong splitting error sufficiently small, and supply
the separate sufficiently small normalized lower-curvature bound on
all required balls. Then at every strong edge point the distance
smoothing can be chosen as in LFR34 so that the ORIGINAL pair
\[
 J=(f,\eta),\qquad\eta=F/\rho,
\]
with any tangential coordinate $f$ meeting the displayed LFR19
bounds, has the following property. At EVERY point of the full band
(LFR36.1), in the point's OWN scale $d_x=d/\rho(x)$, there is an
actual $(2,\beta_2)$-splitting $\Phi_x$ such that $J$ on
$B_{d_x}(x,100)$ has rank two and satisfies ALL the scale-$100$
adapted conditions of quality $\gamma$. If LC20's three-stratum
exclusion is supplied, these are two-stratum points in the sense
of LC16/KL7.1.

For a finite verified strong-edge family with common normalized
analytic data, the ONE smoothing may simultaneously have this full
collar property and supply all disk bundles and cutoffs of LFR33.
The strong tolerance is fixed from those data BEFORE choosing the
family. Thus the local edge disk and its full adapted collar use
exactly the same $F/\rho$, not independently chosen functions.
Strong-edge density and coverage remain a separate assertion.
\end{theorem}
\begin{proof}
Set $\iota=\gamma/1000$. First make $\Delta$ large enough for
LFR35 and LFR36. Choose small tangential and smoothing errors
$\sigma,\varepsilon$ and value error $\mu<10^{-6}$; the direction
threshold $\theta(\varepsilon)$ is now fixed. Choose $\tau$ small
enough for LFR34--LFR36 and for $4\tau\Delta<\beta_2$.
LFR29.1 then chooses the weak-edge qualities and strong endpoint
quality, followed by $\Lambda$, also making $\Delta\Lambda$ as
small as LFR36 needs. All these choices precede noncollapse and
curvature-derivative data. Decrease the strong splitting tolerance
later to meet LFR19, LFR36 and, when needed, LFR28. The sufficiently
late curvature bound is imposed separately; a small splitting error
alone does not imply it. These are additional acyclic inequalities
in the existing FC44 order, not a replacement parameter architecture.

LFR32 produces the coarse border. LFR34 gives the enlarged smoothing;
LFR36 encloses every tested band point BEFORE invoking a small-ball
comparison, and gives the larger smooth buffer in its own metric.
Use LFR35's actual comparison map and its reference coordinates.
For every point of the radius-$300$ ball, choose any positive-anchor
minimizing directions. LFR35 and LFR36 compare each component
covector of $DJ-D\chi_x$ with the SAME unit direction, so
\[
 \|DJ-D\chi_x\|_{\rm op}\leq2\sqrt2\,\iota<\gamma/100.
\]
LFR37 proves the full desired conclusion. Both the actual comparison
map and the two components of $J$ have been retained. To identify
the stratum, use the exhibited two-splitting and the SUPPLIED absence
of three-splittings; rank two alone would not supply that absence.

For the finite family, replace each $U_i,C_i$ in LFR33 by the
physical rescalings of $U^+,C^+$. Apply LFR02 ONCE to the compact
union, with value tolerance smaller than $\mu\Delta\min_i\rho(p_i)$
and the common Lipschitz-difference tolerance. The same rescaling
calculation gives LFR34 at every center, hence the full collar just
proved. These stronger smoothness conclusions also include LFR27's
entire displayed output, so LFR28 still gives the same proper smooth
disk bundles, sublevels, boundary functions and supported cutoffs.
Its finite model metric and actual embeddings are unchanged. The
constant-core modification is used for disk interpolation; it is
NOT substituted for $\eta$ on the lower adapted collar, where its
being constant would destroy rank. This explicitly separates the two
uses of the SAME distance smoothing without changing either function.
\end{proof}

\paragraph{Source and remaining scope.}
KL Definition4.21,4.23--4.28 and Remark4.35, printed34--38/PDF29--33,
provide the adapted-coordinate route; KL9.17--9.20, printed54--55/
PDF49--50, specify the full band. The reference construction above
uses explicit distance and comparison estimates in place of an
asserted convergence of actual angles. KL7.1--7.2, printed46/PDF41,
is why the three-stratum exclusion is retained, and KL15.4,
printed87/PDF82, supplies the order used for the additional constraints.
The author's current correction sheet adds no entry for these passages.
These are written project proofs using the inherited local comparison,
first-variation and smoothing bindings, not Lean declarations.
The local disk plus full collar is now written under its normalized
inputs. KL9.10's strong-edge density/coverage and the zero-model
normal-flow/three-dimensional classification remain. General
Alexandrov splitting/comparison and Tits-cone existence are not proved
by this source-manifold argument, and the accepted-PC gate is unchanged.

\subsection{Strong-edge density and an actual finite disk cover}
\label{sec:collapse-strong-edge-coverage}

The local packets just constructed must occur near EVERY nonslim
rank-one point. We expand KL9.1 and9.10--9.11, keeping two different
accuracies: the one-dimensional model is chosen before noncollapse,
whereas the original rank-one splitting can be sharpened later.
Retaining its EXACT real coordinate makes the residual-factor
comparison elementary. In particular this step does not require a
new application of AC43 or AC76. The existing Alexandrov compactness
comparison input and AC46's one-dimensional recognition remain explicit.

All metrics in LFR39--LFR42 are normalized at the indicated basepoint.
KL approximations have the open-ball convention of AC49. The slim
predicate means existence of an ACTUAL $(1,\beta_1)$-splitting with
factor diameter strictly less than $1000\Delta$, at the ORIGINAL
threshold $\beta_1$; weakening that threshold would not suffice.

\begin{lemma}[LFR39: a one-dimensional model with the prescribed real coordinate]
\label{lem:collapse-prescribed-low-dimensional-model}
For every $0<e<1$ there exists $0<a_0<e/10$ with the following
property, conditional on the comparison input retained in AC41.
Suppose $(Z,p)$ admits both a pointed KL $a_0$-approximation to a
complete proper geodesic nonnegative Alexandrov space of Hausdorff
dimension at most two and an ACTUAL $(1,a_0)$-splitting
$\phi=(u,v):(Z,p)\to\mathbb R\times(X,x_0)$.
Then there is an ACTUAL $(1,e)$-splitting
\[
 F=(u,w):(Z,p)\longrightarrow\mathbb R\times(Y,y_0),
                                                               \tag{LFR39.1}
\]
with the SAME whole-space real coordinate $u$, where $Y$ is complete,
proper, geodesic, nonnegative Alexandrov and $\dim_HY\leq1$.
Sufficiently smaller input qualities are allowed with the same maps.
\end{lemma}
\begin{proof}
Otherwise choose counterexamples with both input qualities tending to
zero. LC17 gives a common complete proper geodesic nonnegative limit
$C$ with $\dim_H C\leq2$. AC50--AC51 applied to the ORIGINAL product
maps give an onto exact splitting $C\to\mathbb R\times Y$, together
with convergence of their real coordinates on every fixed ball.
AC44--AC45 supply the asserted geometry and dimension of $Y$.
AC52 now transfers this splitting back with real coordinate exactly
$u$, giving (LFR39.1) on a sufficiently late tail, a contradiction.
AC52 explicitly extends the residual coordinate off the tested buffer;
no claim of global control there is needed. Neither local recognition
AC46 nor line splitting AC43 is used in this argument: the exact
product is supplied by the converging approximate products themselves.

For completeness, a pointed KL $b$-map with $0<b<a_0/4$ is also
an $a_0$-map. A target of radius $r<a_0^{-1}-a_0$ has an original
coverage witness of error $<2b$ and source radius $<r+3b<a_0^{-1}$.
The original target and source domains contain these tests, and
$b<a_0$ bounds distortion. Thus thresholds can be decreased without
silently discarding the target-coverage margin.
\end{proof}

\begin{lemma}[LFR40: elementary cancellation of the common real coordinate]
\label{lem:collapse-common-coordinate-factor-map}
Fix $L\geq1$, $0<\eta<1/100$ and $\theta>0$. There is $e_0>0$
such that if
\[
 \phi=(u,v):Z\to\mathbb R\times(X,x_0),\qquad
 F=(u,w):Z\to\mathbb R\times(Y,y_0)
\]
are pointed KL $b$- and $e$-approximations with $0<b,e<e_0$, then
there is a pointed KL $\eta$-approximation $G:X\to Y$ and
\[
 d_Y\bigl(G(v(z)),w(z)\bigr)<\theta
                    \quad(z\in B_Z(p,L)).                    \tag{LFR40.1}
\]
No curvature, length structure or continuity of the factors or maps
is assumed. In particular $e_0$ depends on $L,\eta,\theta$, not on a
subsequently prescribed strong-splitting quality.
\end{lemma}
\begin{proof}
Put $R=4(L+\eta^{-1}+1)$ and require
\[
 e_0<\min\left\{\frac\eta{100},\frac1{100R},
                 \frac{\theta^2}{100(L+1)},\frac1{100}\right\}.
                                                               \tag{LFR40.2}
\]
For each $x\in B_X(x_0,R)$ choose an actual $\phi$-lift $a_x$ of
$(0,x)$, with error $<2b$ and $d(p,a_x)<d_X(x_0,x)+3b$.
Choose $a_{x_0}=p$. Set $G(x)=w(a_x)$ on this ball and $G(x)=y_0$
elsewhere. All displayed lifts lie in both original tested source
balls, and all target lifts are inside their coverage radii.
For two such factor points,
\[
 \big|d(a_x,a_{x'})-d_X(x,x')\big|<5b,
 \qquad |u(a_x)-u(a_{x'})|<4b.
\]
The product formula for $F$ gives distortion of $G$ at most $e+9b$.

For $y\in B_Y(y_0,\eta^{-1}-\eta)$, write $r=d_Y(y_0,y)$ and
lift $(0,y)$ through $F$ to $z$, with error $<2e$ and
$d(p,z)<r+3e$. Put $x=v(z)$. Then
\[
 d_X(x_0,x)<r+3e+b<\eta^{-1},\qquad
 d(z,a_x)<2e+3b.
\]
Consequently $d_Y(G(x),y)<5e+3b<\eta$, proving coverage with a
witness in the REQUIRED source ball, not merely in $B_X(x_0,R)$.
The same estimates cover $x=x_0$, whose lift was fixed separately.
Thus $G$ is a pointed KL $\eta$-map.

Finally take $z\in B_Z(p,L)$ and $x=v(z)$. Put
$U=|u(z)-u(a_x)|\leq L+3b$ and $D=e+3b$.
The ORIGINAL product map gives $d(z,a_x)\leq U+3b$; the second
product map therefore gives
\[
 d_Y(w(z),G(x))^2\leq(U+D)^2-U^2
                  \leq(2L+8)(e+3b)<\theta^2.                \tag{LFR40.3}
\]
Here $b,e<1/100$, and (LFR40.2) supplies the last strict bound.
The estimate uses the identity of the two real coordinates, not
closeness of unrelated coordinate maps.
\end{proof}

\begin{lemma}[LFR41: a bounded comparison factor gives the original slim predicate]
\label{lem:collapse-empty-annulus-slim-producer}
Let $\Delta\geq1$, $0<b<(1000\Delta)^{-1}$, and let $(Z,p)$ be
complete and geodesic. Suppose an ACTUAL $(1,b)$-splitting
$\phi=(u,v)$ has residual factor $X$, and there is a pointed KL
$\eta$-map $G:X\to Y$ with $0<\eta<(10^4\Delta)^{-1}$.
If $\operatorname{diam}Y\leq500\Delta$, then $(Z,p)$ has an ACTUAL
$(1,b)$-splitting whose residual factor has diameter less than
$1000\Delta$. The quality is still exactly $b$.
\end{lemma}
\begin{proof}
There is no $x\in X$ with
$600\Delta\leq d_X(x_0,x)\leq900\Delta$: these points would be
inside the tested $G$-ball, whereas pointed distortion would give
$d_Y(y_0,G(x))\geq600\Delta-\eta>500\Delta$.
For any $z\in B_Z(p,b^{-1})$, take a minimizing path from $p$ to
$z$, entirely inside that ball. Subdivide it into steps shorter than
$\Delta$. The radial function $r(t)=d_X(x_0,v(t))$ changes by less
than $\Delta+b$ at each step, by distortion of $\phi$ and the product
projection bound. It starts at zero and cannot jump across the empty
closed annulus $[600\Delta,900\Delta]$. Hence
\[
 v\bigl(B_Z(p,b^{-1})\bigr)\subset X_0:=B_X(x_0,600\Delta).
                                                               \tag{LFR41.1}
\]
No inverse image of a factor ball has been asserted open or closed.
Since $G$ controls all pairs in $X_0$,
\[
 \operatorname{diam}X_0\leq500\Delta+\eta<1000\Delta;
\]
a radius bound alone would give the insufficient $1200\Delta$.
Keep $\phi$ unchanged on its ENTIRE tested ball, and outside it
replace the residual coordinate by $x_0$, keeping $u$.
This defines a whole-space map to $\mathbb R\times X_0$ with the
original distortion bound $b$. Every tested target in this smaller
product has the original coverage witnesses, all in that same
source ball and hence with residual coordinates in $X_0$ by
(LFR41.1). Their distances in the restricted factor are unchanged.
This proves the ORIGINAL $(1,b)$-splitting, without a loss of quality.
\end{proof}

\begin{lemma}[LFR42: a nonslim rank-one model has a nearby endpoint]
\label{lem:collapse-nonslim-endpoint-model}
Fix $0<\beta_2<1/100$ and $\Delta>100/\beta_2$.
Let $(Z,p)$ have no $(2,\beta_2)$-splitting. Suppose $F=(u,w)$ is a
pointed $(1,e)$-splitting with $e<\min\{\beta_2/100,(1000\Delta)^{-1}\}$,
where $Y$ satisfies LFR39 and $\operatorname{diam}Y>500\Delta$.
Conditional on AC46, $Y$ can be identified with $[0,C]$, $C>500\Delta$,
or $[0,\infty)$, with its original basepoint at a height
\[
                           0\leq a<\Delta/2.                 \tag{LFR42.1}
\]
In the finite case a reflection may be needed. The conclusion does
not replace a positive basepoint height by zero.
\end{lemma}
\begin{proof}
AC47 gives singleton, line, ray, segment or circle, with their actual
length metrics. The singleton is excluded by the diameter assumption.
A circle here has circumference greater than $1000\Delta$.
For a line or such a circle, or for a ray/segment whose basepoint is
at distance at least $\Delta/2$ from every endpoint, the product
has an EXACT pointed plane chart on its radius-$\Delta/4$ ball,
with its restricted pairwise distances preserved. In the circle case,
two factor coordinates in this chart differ by less than $\Delta/2$,
well below half the circumference.

Write $T=\beta_2^{-1}$. On $B_Z(p,2T)$ compose $F$ with this chart
and extend by zero elsewhere. Its tested radius-$T$ ball has distortion
at most $e$. For a plane target of radius $r<T-\beta_2$, regard it
as the same point of the product chart. Original $F$-coverage gives
error $<2e$ and a source witness of radius $<r+3e<T$.
All these images lie strictly inside the plane chart because
$2T+e<\Delta/4$. This is an ACTUAL pointed $(2,\beta_2)$-splitting
with singleton residual factor, contrary to the hypothesis.
The only remaining models have an endpoint within $\Delta/2$.
A finite segment has length greater than $500\Delta$; reflect it
if necessary. This proof uses explicit charts of the classified
one-dimensional models, not a strainer-to-splitting implication.
\end{proof}

\begin{lemma}[LFR43: border lifts are strong edges in their own normalization]
\label{lem:collapse-strong-border-lift}
Let $\Delta\geq1$, $0<s<1/100$, $0<b_E<1/100$ and
$0<\Lambda\Delta<10^{-6}$. In the normalization $d_p=d/\rho(p)$
of a complete connected smooth manifold, let $\rho>0$ be
$\Lambda$-Lipschitz for $d$. Suppose
\[
 \phi=(u,v):M\to\mathbb R\times(X,x_0),\qquad
 F=(u,w):M\to\mathbb R\times(Y,a)
\]
have qualities $b_1,e$, where $Y=[0,C]$, $C>500\Delta$, or
$Y=[0,\infty)$, and $0\leq a<\Delta/2$.
Suppose $G:X\to(Y,a)$ is a pointed KL $\eta$-map and
\[
 |G(v(z))-w(z)|<\theta\quad(z\in B_{d_p}(p,1000\Delta)),
\]
where
\begin{gather*}
 \eta<\min\{s/1000,(10^5\Delta)^{-1}\},\qquad
 \theta<s/1000,\qquad e<\min\{s/1000,(1000\Delta)^{-1}\},\\
 b_1<\min\{b_E/100,(100\Delta+100/b_E)^{-1}\}.       \tag{LFR43.1}
\end{gather*}
For EVERY $|t|\leq12\Delta$, choose an actual $F$-lift $q$ of $(t,0)$
with error $<2e$. Then $d_p(p,q)<13\Delta$, and $q$ has strong
$(b_E,s)$ edge maps at its OWN scale $d_q=d/\rho(q)$, with a FINITE
endpoint interval of length strictly greater than $200\Delta$.
The small-model tolerances $e,\eta,\theta$ need not depend on $b_E$.
\end{lemma}
\begin{proof}
Original coverage supplies the lift with
\[
 d_p(p,q)<\sqrt{t^2+a^2}+3e<13\Delta.
\]
Put $r=\rho(q)/\rho(p)$. Then $.99<r<1.01$, and
$h=G(v(q))/r<(2e+\theta)/r<s/100$.
Keep the ORIGINAL factor $X$, scale its metric by $r^{-1}$,
base it at $v(q)$, and use
\[
 \phi_q(z)=\bigl((u(z)-u(q))/r,v(z)\bigr).
\]
Its distortion is at most $b_1/r<2b_1$.
The old source and coverage domains contain the new tests, since
$13\Delta+2/b_E+2<b_1^{-1}$ by (LFR43.1).
For a new target of radius $R<b_E^{-1}-b_E$, the old coverage
witness has new image error $<2b_1/r$ and new distance from $q$
less than $R+3b_1/r<b_E^{-1}$. Thus $\phi_q$ is a pointed KL
$b_E$-splitting. Only this step asks $b_1$ to depend on $b_E$.

If $Y=[0,C]$, put $D=C/r>200\Delta$ and initially set
$H(z)=G(z)/r$. If $Y$ is a ray, put
\[
 D=\max\{201\Delta,4/s\},\qquad H(z)=\min\{G(z)/r,D\}.
\]
In both cases replace $H(v(q))$ by zero, leaving all other values
unchanged. This is a whole-space pointed map to $([0,D],0)$.
For a point in the new $s^{-1}$ source ball about $v(q)$, its old
radius is less than $13\Delta+b_1+2/s<\eta^{-1}$.
In the ray case, its untruncated image height is less than
$14\Delta+3/s<D$, so the truncation does NOTHING on the tested
ball. The old distortion, scaling and one-point repair therefore
give new distortion at most $\eta/r+2h<s$.

For a target height $c<s^{-1}-s$ in $[0,D]$, the original target
height $rc$ lies in $Y$ and has distance at most $2/s+\Delta/2$
from $a$, well inside the original $G$-coverage domain.
Choose $z$ with $|G(z)-rc|<2\eta$.
Both $z$ and $v(q)$ lie in the tested old $G$-ball. Hence its new
source radius is less than
\[
 c+h+3\eta/r<c+s/10<s^{-1}.
\]
Its new image error is less than $2\eta/r+h<s$, including when
$z=v(q)$ was the repaired point. In the ray case the preceding
source-radius estimate also ensures that clipping is inactive at
this witness. This proves the KL $s$-map and the strict finite
endpoint convention. No continuity, and no positive margin in an
unrelated edge interval, was used.
\end{proof}

\begin{theorem}[LFR44: strong-edge density and coverage by the actual disk packets]
\label{thm:collapse-strong-edge-density-cover}
Let $(M,g)$ be a complete connected smooth three-manifold and
$\rho>0$ a $\Lambda$-Lipschitz scale. Retain the edge predicates and ordered constraints of LFR29.1,
with $\Delta>100/\beta_2$, and decrease
$b'=\beta_{E'},s'=\sigma_{E'}$ below $10^{-8}$ and
$\Lambda\Delta$ below $10^{-6}$ if necessary.
There is a collapsed-model tolerance $a_*>0$, chosen using only
$\Delta,\beta_2,s=\sigma_E$, and a rank-one threshold
$b_{1,*}>0$, allowed to depend also on $b_E=\beta_E$, such that:
assume every point under consideration admits a pointed KL
$a_*$-approximation to a complete proper geodesic nonnegative
Alexandrov space of dimension at most two. Use
$0<\beta_1<b_{1,*}$ in LC16 and the ORIGINAL slim predicate.
Conditional on the retained comparison input and AC46,
\begin{enumerate}[leftmargin=*]
\item every nonslim one-stratum point $p$ has some $q\in E$ with
  $d(p,q)<\Delta\rho(q)$;
\item for every such $p$ and every
  $p'\in E'\cap B(p,10\Delta\rho(p))$, there is $q\in E$ with
  $d(p',q)<\rho(q)$.
\end{enumerate}
If $M$ is compact, one can select a finite family $q_i\in E$ with
pairwise disjoint balls $B(q_i,\Delta\rho(q_i)/3)$ such that the
$B(q_i,3\Delta\rho(q_i))$ cover both the nonslim one-stratum and
all the weak points in item2. If the common analytic/curvature
hypotheses and parameter bounds of LFR33 and LFR38 are supplied
at these centers, the SAME family has ONE distance smoothing,
actual proper smooth disk bundles and full adapted collars, and
those actual disk domains cover these sets. Their cutoffs equal
one there. No choice of a new strong tolerance is made after the
family is selected.
\end{theorem}
\begin{proof}
\emph{Choose the model accuracy before the strong-splitting accuracy.}
Set $L=1000\Delta$ and choose
$\eta<\min\{s/1000,(10^5\Delta)^{-1}\}$ and $\theta<s/1000$.
Choose $e>0$ satisfying LFR40.2 and
\[
 e<\min\{10^{-8},s/1000,\beta_2/100,(1000\Delta)^{-1}\},
\qquad a_*<a_0(e)/4,
\]
where $a_0(e)$ is from LFR39. Finally choose $b_{1,*}$ below
$a_0(e)/4$, the $b$ bound of LFR40.2, and all bounds of LFR43.1,
also below $(1000\Delta)^{-1}$ and $\beta_2$.
These choices allow the strong $b_E$ to be fixed using later
noncollapse/derivative data, then $\beta_1$ chosen still smaller.
The collapsed-model tolerance $a_*$ has NOT been chosen using $b_E$.
Its source application is LC08--LC09, using LC15, and the sufficiently late
uniform-in-center curvature tail. The different local curvature
hypotheses for disk/collar production remain separately supplied.

Fix a nonslim one-stratum point and normalize at $p$.
Its original $\beta_1$-splitting $\phi=(u,v)$ and LFR39 give the
model $F=(u,w)$, with EXACTLY the same real coordinate.
LFR40 supplies $G$ and (LFR40.1). LFR41 excludes
$\operatorname{diam}Y\leq500\Delta$, since that would give the
original slim predicate. By LC16 there is no $(2,\beta_2)$-splitting;
LFR42 therefore puts the model at height $a<\Delta/2$ above an
endpoint, retaining the possible ray. LFR43 applies to every
border lift with $|t|\leq12\Delta$.
At $t=0$ it gives a strong edge $q$ with
\[
 d_p(p,q)<a+3e<.51\Delta<\Delta\rho(q)/\rho(p),
\]
which is item1.

\emph{Weak edge points cannot be high above this model border.}
Let $p'$ be as in item2 and put $F(p')=(t,h)$.
Then $|t|<11\Delta$ and $0\leq h<11\Delta$.
Suppose $h\geq1/20$, and put $R=1/100$.
The four targets $F(p')\pm Re_1,F(p')\pm Re_2$ lie in the same
product interval (or half-plane), well inside the original coverage
range. Choose actual $F$-lifts of error $<2e$. Distances from
$p'$ to these four points differ from $R$ by at most $3e$,
and their pairwise distances differ from the four-axis Euclidean
configuration by at most $5e$ (the weaker bounds $4e,7e$ also hold).
All points lie in $B_{d_p}(p,12\Delta)$.

Put $r'=\rho(p')/\rho(p)$; then $.99<r'<1.01$.
The weak composite map based at $p'$ has distortion at most
$b'+s'$ on its radius-one ball by LFR30. The four lifts lie in
that ball in $d_{p'}$. Multiply their weak images by $r'/R$.
They lie in the half-plane based at its border and would satisfy
LFR29's forbidden four-point estimates with error
\[
 \frac{7e+2(b'+s')}{R}<10^{-3}.
\]
Thus $h<1/20$. Choose an actual $F$-lift $q$ of $(t,0)$.
LFR43 says $q\in E$ in its own scale, while
\[
 d_p(p',q)<h+3e<1/10<\rho(q)/\rho(p).
\]
This proves item2 directly. It does not apply a strong-centered
border lemma at $p$ before knowing that $p$ is strong.

\emph{Finite selection and the ACTUAL disk domains.}
If $E$ is empty, item1 says the nonslim stratum is empty and there
is no item2 region; choose the empty family. Otherwise compactness
and the positive minimum of $\rho$ give a finite upper bound for
cardinalities of disjoint $\Delta\rho/3$ balls. Choose a family of
maximal cardinality. Every omitted $q\in E$ has its small ball
meeting one selected small ball, centered at $q_i$. Put
$r_q=\rho(q)/\rho(q_i)$ and $k=\Lambda\Delta/3$.
The intersection and the Lipschitz bound imply
\[
 \frac{1-k}{1+k}\leq r_q\leq\frac{1+k}{1-k}<1.01,
 \qquad
 d(q,q_i)<\frac\Delta3(r_q+1)\rho(q_i)<.7\Delta\rho(q_i).
                                                               \tag{LFR44.1}
\]
The same conclusion is immediate if $q=q_i$.
Items1--2 and $\Delta\geq1$ now put each point in question in
some $B(q_i,3\Delta\rho(q_i))$; in fact the preceding estimates
give the stronger constant $2$. This also proves the usual
$\Delta\rho(q_i)$-ball cover of $E$.

The selected family is now fixed, at the already chosen strong
quality. Apply LFR33 and LFR38 using the common supplied analytic
data, choosing the common value error $\mu<10^{-6}$.
Let $\eta=F_A/\rho$ denote their ONE smoothed distance to
$A=\overline{E'}$, and let $f_i$ be the original tangential
coordinate at $q_i$. In its normalization, any $x$ with
$d(q_i,x)<3\Delta\rho(q_i)$ satisfies
\[
 |f_i(x)|<3\Delta+b_E+\mu\Delta<3.1\Delta,
 \qquad
 \eta(x)<\frac{(3+\mu)\Delta}{1-3\Delta\Lambda}<3.1\Delta.
                                                               \tag{LFR44.2}
\]
Here $q_i\in E\subset E'$, so its distance to $A$ is zero;
the value bounds and the derivative-independent scale estimate
are sufficient. These strict inequalities put the entire ball
inside the ACTUAL LFR28 disk domain $|f_i|<4\Delta$, $\eta\leq4\Delta$,
and inside its cutoff plateau. The same smoothing supplies its
full lower adapted collar by LFR38. Thus the output is a cover by
the constructed packets, not just a cover by unrelated metric balls.
Multiplicities retain LC86/FC08's stated whole-ball curvature inputs;
no new uniform bound on the total number of centers is asserted.
\end{proof}

\paragraph{Source comparison and remaining scope.}
KL9.1 and9.10--9.11, printed49--52/PDF44--47, supply the route;
9.23--9.24, printed56/PDF51, separate the late splitting tolerance
and finite selection. LFR39 strengthens the model statement by
preserving the original real coordinate, using AC51--AC52's written
map-level convergence. LFR40 replaces a second compatibility theorem
by squared product-distance estimates. LFR41 supplies the coarse-chain
repair to the source's unjustified open-and-closed preimage assertion,
including the global factor diameter and unchanged slim quality.
LFR42 uses the weaker endpoint height $\Delta/2$, sufficient for
9.10's conclusion, instead of the source's $\Delta/10$; the ray is
retained and truncated only in LFR43 with explicit coverage margins.
Both strong edge maps are on the SAME recentered original factor.
The parameter choices respect FC44/15.4; the additional explicit
bounds are project sufficient conditions, not a transcription of
every displayed source dependency. The current author correction
sheet adds no entry for these passages.

This completes the written strong-edge density and finite edge-packet
cover under its stated collapsed-model, analytic and recognition inputs.
General Alexandrov local comparison and AC46 recognition are NOT proved
by this argument. The zero-model actual normal-flow/classification and
Tits-cone inputs remain. The inherited PC bindings, full Chapter14
compatibility and accepted Lean evidence are distinct. This is neither
formal completion nor the final audit of the user's entire goal.

\subsection{Finite soul flows and a smooth carrier for the zero packet}
\label{sec:collapse-finite-zero-core}

The zero-model consumer needs an actual normal-flow map, not merely
an abstract topological soul theorem. We construct that map in finite
regularity, and then use it to choose a suitable smooth carrier. The
metric is NEVER smoothed. The cost of this choice is recorded below:
starting with the LFR14 metric order $m=K-1$, the zero-model carrier
has metric order $K-5\geq5$. This suffices for the geometric operations
used here. The earlier surface and edge constructions keep their own
carriers; no change to their witnesses is required.

\begin{lemma}[LFR45: a finite soul with strict outward directions]
\label{lem:collapse-finite-soul-separators}
Let $(N,g)$ be a complete connected noncompact three-manifold without
boundary, on a smooth carrier, with $g$ of class $C^m$, $m\geq8$, and $\sec_g\geq0$.
There is one compact connected boundaryless totally convex, totally
geodesic $C^{m-1}$ submanifold $S$, of dimension at most two, with
$C^{m-2}$ normal bundle $E=\nu_g S$. For every $q\notin S$ there is
a unit vector $v\in T_qN$ such that
\[
       g(v,u)<0\qquad(u\in\mathcal U_S(q)),             \tag{LFR45.1}
\]
where $\mathcal U_S(q)$ contains ALL unit minimizing directions to
$S$. There is a positive normal-tube radius $\epsilon$ on which the
normal exponential $\Phi:E_{<\epsilon}\to\{d_S<\epsilon\}$ is a
$C^{m-2}$ diffeomorphism, fixes $S$, and satisfies
$d_S(\Phi(s,w))=|w|$. Here $d_S=d_g(\cdot,S)$.
\end{lemma}
\begin{proof}
The convex exhaustion, boundary-distance concavity and convex-set
structure proof in LFR21 are dimension independent. Apply their
construction in dimension three, allowing at most three dimension
drops. For precision, retain the compact totally convex exhaustion
$C_t$ from Cheeger--Gromoll1.3 and its erosion identity, then the
finite flag of maximum boundary-distance sets from1.9--1.11. Each
nonterminal level has totally geodesic relative interior; a terminal
set with empty relative boundary is the required $S$. Its dimension
is less than three: a boundaryless full-dimensional compact submanifold
would be open and closed in the connected noncompact $N$.

The finite calculus in LFR21 applies unchanged: exponential charts and
geodesic variations have class $C^{m-1}$, and the curvature coefficients
in Jacobi/Riccati comparison are $C^{m-2}$. The short convex-ball,
Rauch and complete two-minimal-side triangle comparisons use the
inherited Riemannian bindings with these finite operations. This does
not invoke general Alexandrov splitting. The relative-interior charts
give $S$ class $C^{m-1}$ and its orthogonal normal bundle class
$C^{m-2}$.

We need additional information from this SAME flag. For every
$q\notin S$ there is a compact totally convex set $C$ with
\[
                 q\in\partial_{\rm rel}C,
                 \qquad S\subset\operatorname{int}_{\rm rel}C.
                                                               \tag{LFR45.2}
\]
Outside the innermost exhaustion member, use the first $C_t$ containing
$q$. The erosion identity and continuity of distance give membership
at this first parameter, boundary membership there, and strict interior
containment of the innermost member when the parameter is positive.
Inside that member use the first flag stage $B$ for which
$q\in B\setminus B_{\max}$. If $a=d(q,\partial_{\rm rel}B)$,
take $C=\{x\in B:d(x,\partial_{\rm rel}B)\geq a\}$.
Here $a<\max d(\cdot,\partial_{\rm rel}B)$, so this erosion has the
same relative dimension as $B$ and its maximum set, hence $S$, lies
in its relative interior. The point $q$ is on its boundary: a short
initial part of a minimizing connection to the boundary decreases the
distance. This includes $a=0$. The flag terminates, so these cases
exhaust $N\setminus S$.

At a relative boundary point of a convex set, the cone of directions
entering its relative interior is a proper open convex cone in its
linear span. Cheeger--Gromoll1.4 and the discussion after1.6 prove
openness by coning a transverse interior patch with short geodesics;
1.7 and the supporting-half-space argument on printed420 give a
supporting linear functional. Their construction uses only finite
exponential charts, first variation, and compact unit tangent fibers.
Equivalently, separate this open convex cone from zero in its span.
A nonzero supporting functional is strictly positive on the open
cone: equality at an interior vector would, by a small perturbation
in the negative functional direction, contradict nonnegativity.

Every geodesic from $q$ to a point of the relative interior has initial
direction in this cone. To check the assertion even for a long segment,
apply1.4 backwards in short convex balls along it. If its relative
interior portion had a first endpoint at positive time,1.4 with a
nearby interior point and a preceding point would extend that portion
backwards, a contradiction. Total convexity keeps the entire segment
in $C$. Thus all directions from $q$ to $S$ in (LFR45.2) have strictly
positive pairing with one inward supporting unit vector. Its negative
gives (LFR45.1), simultaneously for every minimizing choice.

Finally the differential of the normal exponential at the zero section
identifies $T_sS\oplus E_s$ with $T_sN$. The finite inverse-function
theorem and compactness give an injective tube of one positive radius.
A minimizing segment to compact $S$ is normal by first variation.
After decreasing the radius, every point with $d_S<\epsilon$ therefore
has its unique such tube representation, and no shorter representation
is possible. This proves both the target equality and calibration.
It also proves that $d_S$ is $C^{m-2}$ away from $S$ in this tube.
\end{proof}

\begin{lemma}[LFR46: the actual finite normal-flow map]
\label{lem:collapse-finite-normal-flow-map}
For the SAME $S,E,\Phi$ of LFR45, there are $\ell>0$, a bounded
$C^{m-3}$ field $V$, its complete $C^{m-3}$ flow $\varphi$, and a
$C^{m-3}$ diffeomorphism $e:E\to N$, fixing the zero section, with
\[
 e(s,w)=
 \begin{cases}
 \Phi(s,w),&|w|\leq\ell,\\
 \varphi_{|w|-\ell}\bigl(\Phi(s,\ell w/|w|)\bigr),&|w|>\ell.
 \end{cases}                                                   \tag{LFR46.1}
\]
One can require $|V|\leq2$ and $V=\nabla d_S$ on a tube annulus
containing $\{d_S=\ell\}$. Inside this annulus, toward $S$, $V$
can be chosen as $a(d_S)\nabla d_S$, where $a$ is a smooth scalar
profile, zero near zero and one near $\ell$.
If the model has its supplied LC21 cone at a marked point $p$, this
SAME choice may also
satisfy, for some $A_2>0$,
\[
       g(V,u)\leq-1/4
       \quad(d(p,q)\geq A_2,\ u\in\mathcal U_p(q)).       \tag{LFR46.2}
\]
The map in (LFR46.1) is chosen after this field, not before it.
\end{lemma}
\begin{proof}
The direction set $\mathcal U_S(q)$ is nonempty and compact. Its graph
in the unit tangent bundle is closed: a limit of its endpoints satisfies
$\exp_q(d_S(q)u)\in S$, using continuity of the exponential and of
$d_S$. Unit tangent fibers over compact sets are compact. Therefore a
strict inequality against all these directions persists locally, with
a smaller margin. Extend the vector in (LFR45.1) to a smooth local
field in the present carrier, preserving norm less than two and strict
negativity on a smaller neighborhood. Convex patching by a locally
finite partition gives a smooth bounded outward field outside any
prescribed smaller tube, with zero extension near $S$. No smooth
selection of minimizing segments is used.

In the calibrated tube $\nabla d_S$ has norm one and pairs to $-1$
with the inward minimizing direction. Its regularity is at least
$C^{m-3}$. Blend the preceding field with $a(d_S)\nabla d_S$ on an
annulus strictly inside the tube, keeping the latter on a neighborhood
of $\{d_S\leq\ell+\delta\}$, where
$0<\delta<\ell$ and $\ell+\delta<\epsilon$. Choose $a=1$ from
some positive level $b<\ell-\delta$ through the gluing annulus, and
$a=0$ near zero. All blending supports avoid $S$. The result is
$C^{m-3}$, bounded by two, and strictly outward to every minimizing
soul direction on $\{d_S\geq b\}$.

For the optional point margin, LC52's triangle estimate and LC53's
convex patching apply in this finite category: they use the supplied
cone, finite first variation and triangle comparison, and continuity
of the closed direction graph. Their common field tests BOTH targets.
Blend it with the preceding $V$ using a smooth cutoff equal to zero
near the tube and the common field's inner forbidden ball, and one
outside a larger compact set, exactly as in LC54. The blend is used
only where that common field is also soul-outward. Both the bound and
strict soul inequalities survive; (LFR46.2) holds at infinity. The
prescribed inner radial profile is unchanged.

A bounded field on a complete proper Riemannian manifold has a complete
flow: a trajectory in finite time has bounded length and remains in a
compact ball, where the finite ODE continuation theorem applies. The
joint flow is $C^{m-3}$. Along its exterior trajectories, first variation
for distance to a compact set gives positive distance growth. More
explicitly, on each compact annulus $b\leq d_S\leq B$, compactness of
the direction graph gives $c_B>0$ with $-g(V,u)\geq c_B$ for every
minimizing direction. Distance composed with the flow is locally
Lipschitz, and its almost-everywhere derivative is at least $c_B$
while in that annulus. Integration proves strict increase, finite-time
crossing of each higher level, and a unique backward crossing of
$\{d_S=\ell\}$ from any point with $d_S\geq\ell$. A trajectory
cannot return across a level where this positive-growth condition holds.

Near $\ell$, the radial normal geodesics and the flow solve the SAME
ODE, since $V=\nabla d_S$ there. Thus the two formulas (LFR46.1)
agree on an open overlap, not merely on their common boundary.
They define a $C^{m-3}$ map. For an exterior point $x$, let $\tau(x)\leq0$
be its unique hitting time at level $\ell$. Its inverse image is
\[
 \frac{\ell-\tau(x)}{\ell}\,
                  \Phi^{-1}(\varphi_{\tau(x)}x).                 \tag{LFR46.3}
\]
The implicit function theorem at the hitting point applies to
$d_S(\varphi_t x)-\ell$: it is $C^{m-3}$ near that point and has
time derivative one. Hence $\tau$ is locally $C^{m-3}$. Formula
(LFR46.3) and the inner inverse $\Phi^{-1}$ agree on the overlap.
The flow law verifies both inverse identities. This supplies the
actual diffeomorphism and its retained ray formula in finite regularity.
\end{proof}

\begin{lemma}[LFR47: a smooth soul-flow carrier with a finite metric]
\label{lem:collapse-soul-flow-smooth-carrier}
For LFR46's data put $r=m-3$. There is a smooth vector bundle
$\widehat E\to\widehat S$ with smooth fiber norm, and a fiberwise
isometric $C^{m-2}$ bundle isomorphism $B:E\to\widehat E$ over a
compatible smooth carrier of $S$. Transport the smooth structure of
$\widehat E$ to the SAME underlying $N$ by
\[
                       F=e\circ B^{-1}.                         \tag{LFR47.1}
\]
Then $F$ is a smooth diffeomorphism in this new carrier; its identity
transition with the old carrier is $C^r$ in both directions. The
UNCHANGED metric has class at least $C^{r-1}=C^{m-4}$ in the new
carrier, with the same distances, completeness and nonnegative
sectional curvature. The soul, its normal bundle with induced fiber
norm, the actual normal-flow map and its bounded field $V$ are smooth
in this carrier. No smoothness of the Riemannian metric is asserted.
\end{lemma}
\begin{proof}
Choose a compatible smooth structure on the compact $C^{m-1}$ base.
Using finitely many orthonormal bundle frames of class $C^{m-2}$ and
scalar cutoffs $\rho_j$ with $\sum\rho_j^2=1$, embed $E$ isometrically
as a $C^{m-2}$ subbundle of a trivial Euclidean bundle. Its orthogonal
projection $P(s)$ is $C^{m-2}$. Approximate $P$ uniformly by a smooth
symmetric matrix field $A$. If the approximation is sufficiently close,
its spectrum lies in two disjoint neighborhoods of zero and one. Its
spectral projection $\widehat P$ onto the latter is smooth and uniformly
close to $P$. For example apply the smooth matrix function
\[
 \widehat P=\tfrac12\left(I+(2A-I)\bigl((2A-I)^2\bigr)^{-1/2}\right).
\]
Positive-definite inverse square roots are smooth, as follows locally
from their convergent matrix power series after scalar normalization.
The range $\widehat E$ is a smooth subbundle of the same rank. The
restriction $L=\widehat P|_E$ is invertible onto that range when
$\|\widehat P-P\|<1/2$. Its polar correction
$B=L(L^*L)^{-1/2}$ is a fiberwise isometric $C^{m-2}$ isomorphism.
Thus its map on closed disk bundles preserves the ACTUAL radius, not
just the total-space topology.

The global finite diffeomorphism $F$ in (LFR47.1) defines the new
smooth atlas. Its $C^r$ transition to the old atlas gives, by the tensor
pullback formula, metric regularity $C^{r-1}$. This is the derivative
loss: metric components contain first derivatives of $F$. As $r\geq5$,
curvature transformation, first variation and the comparison calculus
used here remain legitimate. All curves and their lengths correspond
under this same map; the metric space and curvature tensor have not
been replaced by an approximation.

Near the zero section $F$ is the original normal exponential composed
with the fiber isometry $B^{-1}$. Its vertical differential identifies
$\widehat E$ with the geometric normal bundle of $S$, and the induced
norm is precisely its smooth Euclidean fiber norm. The zero section
is smooth in this atlas. On the exterior, the retained ray formula
identifies $V$ with unit radial differentiation in $\widehat E$.
Inside, it identifies $V$ with $a(|w|)\partial_{|w|}$, zero near zero
and agreeing with unit radial differentiation near the seam. This is
a global smooth field on the smooth bundle. Its complete flow, norm
bound and point-direction inequalities are the same ones as before.
This proves every claimed smooth object without altering the metric.
\end{proof}

\begin{lemma}[LFR48: smooth buffered comparison maps for this carrier]
\label{lem:collapse-smooth-buffered-zero-maps}
Let $(N,g)$ be complete and connected without boundary, on a smooth
carrier with a $C^2$ Riemannian metric. Suppose actual pointed $C^s$
embeddings $j_i:U_i\to M_i$, $s\geq3$, into complete connected smooth
Riemannian manifolds of the SAME dimension, without boundary, eventually
contain each model compact set in their domains, converge in pointed
distance on compact buffers, and satisfy $j_i^*g_i\to g$ in coordinate
$C^1$ there.
There are actual pointed SMOOTH embeddings $\widetilde j_i$ on an
exhausting sequence of open buffered domains, with the same metric
$C^1$ convergence and pointed distance control. Each fixed smaller
source ball is eventually contained in the image of any fixed larger
model ball, in the LFR14 conventions. After LFR47 this applies to LFR14's maps at order $K-4$.
\end{lemma}
\begin{proof}
Fix compact buffers $Q\Subset O\Subset O^+\Subset U_i$, with the
marked point in their interiors. On a neighborhood of $\overline O$
approximate $j_i$ in $C^2$ by a smooth map into $M_i$, using finitely
many source charts, a smooth target embedding and its tubular retraction.
The target can be restricted to a relatively compact neighborhood of
$j_i(\overline {O^+})$. Local convolution, finite cutoffs and the
retraction give arbitrarily small $C^2$ error. The required tolerance
may depend on $i$ and its target metric; there is no uniform derivative
bound on these replacements.

Embedding is stable on the smaller compact buffer: finitely many
inverse-function charts handle close pairs, while compactness and
injectivity of $j_i$ give a positive target separation for the remaining
pairs. Sufficient $C^1$ closeness preserves both. In a target chart near
$p_i$, compose with a smooth cutoff translation taking the new image
of the marked point exactly to $p_i$. Its $C^2$ change tends to zero
with that displacement, and it is a diffeomorphism when small. This
pointed correction preserves embedding on $O$.

For each fixed $i$, pullback of its fixed smooth $g_i$ is continuous
from $C^2$ maps to $C^1$ tensors on these compact sets. Choose the
approximation to make this tensor difference less than $1/i$, and
the target metric displacement less than $1/i$. Therefore convergence
to the SAME $g$ and the original pairwise distance estimates survive.
Use an increasing compact exhaustion and a diagonal choice of tails
to obtain one such smooth map per index on expanding open domains.
Every finite collection of buffers is eventually included.

For coverage fix $0<a<b$ and first use a larger compact model ball
than the closed $b$-ball. On it, pullback metric convergence gives
$(1-\lambda_i)^2g\leq\widetilde j_i^*g_i\leq(1+\lambda_i)^2g$,
with $\lambda_i\to0$. Lift a minimizing source segment of length
less than $a$ from its pointed preimage, using local inverses. A first
exit from the model $b$-ball would require length at least
$(1-\lambda_i)b>a$, a contradiction. Compactness of the buffered
closed ball ensures continuation of the lift at any earlier endpoint.
Thus the entire smaller source ball is covered. The same first-exit
argument retains the distance comparisons on any prescribed buffer.
After LFR47 the fixed coordinate transition is $C^{K-4}$, so the
old $C^{K-1}$ tensor convergence still implies the required $C^1$
convergence. No model metric has been smoothed in this construction.
\end{proof}

\begin{theorem}[LFR49: the finite model supplies the joint zero packet]
\label{thm:collapse-finite-joint-zero-packet}
Use LFR14 for smooth complete connected oriented three-manifolds,
with $K\geq10$, including its expanding-ball lower sectional bounds.
Supply LC21's cone package for the SAME limit, using the point cone
in the compact case, and retain the inherited Riemannian comparison
and analytic bindings in LC30. Fix $\delta,\varepsilon,e,T$ in
LC57's ranges BEFORE choosing the scale. Then the conclusions of LC57 and
LC61 hold on one common tail, at one fixed enlarged scale, for the
same model, cone, selected source radial functions and actual source
balls. In particular the LC37 packet has actual smooth embeddings
and smooth normal-disk-bundle identifications in supplied carriers;
every $A_\rho$, $\rho\in[1/5,2]$, has the stated smooth type, and
the open distance balls in LC61 have the smooth type of that model.
The scale interval contradiction in LC58 may use this finite package
instead of its stronger smooth-metric LC56 premise, provided these
finite-package hypotheses hold for EVERY sequence of indices and
basepoints required there. A single sequence does not supply that premise.
\end{theorem}
\begin{proof}
LFR14 gives one limit of metric order $m=K-1\geq9$ and its actual maps.
The supplied bounds $\sec_{g_i}\geq-\epsilon_i$ on the
$L_i$-balls can be put in LC57's form by taking
$H_i=\min\{L_i,\epsilon_i^{-1/2}\}$, interpreting the second entry
as infinity when $\epsilon_i=0$. Then $H_i\to\infty$ and
$\sec_{g_i}\geq-H_i^{-2}$ on the $H_i$-ball. This uses the
SEPARATE supplied curvature bounds, not the later $C^1$ convergence.
If it is noncompact,
apply LFR45--LFR47, choosing the point-outward field before its flow
map. The zero carrier now has metric order $K-5\geq5$, its actual
soul-flow and normal bundle are smooth, and all distances and the
supplied cone remain unchanged. LFR48 supplies actual pointed smooth
comparison embeddings with compact-set eventual exhaustion and $C^1$
pullback convergence. If the limit is compact, keep its earlier smooth
carrier and apply LFR48 directly.

Here are the consumer checks that replace the unjustified smooth-metric
upgrade. LFR02 smooths scalar distance on this chosen carrier. The
metric enters through its continuous covector norm and the finite
first-variation statements. In the noncompact case use LC55's exact
numerical choices: function error $1/80$, Lipschitz difference $1/64$,
and an enlarged fixed $R$ with $R>2A_2$, an inner constant normal-flow
core inside $B(p,R/2)$, and sufficiently accurate approximation to the
retained cone. Its $D=\{\zeta_R\leq1\}$ is a smooth compact domain,
strictly between the closed half-ball and the open two-ball. Smoothness
is needed only near its boundary. The SAME field $RV$ obeys the
all-direction margin $-1/4$ and gives
$d\zeta_R(RV)\geq1/4-2/64=7/32$ there.

In the smooth coordinates of the actual normal-flow map, every radial
fiber crosses this boundary exactly once. The proof in LC55 uses
positive derivative at crossings, compactness and the smooth implicit
function theorem, not a smooth metric. Its compactly supported vertical
isotopy therefore gives the ACTUAL smooth disk-bundle identification,
with its stated factor $T_*/R$ for the rescaled unit-disk convention.
This is the constructed soul and field, not an unrelated existential
normal-bundle homeomorphism.

At this FIXED $R$, restrict LFR48's maps to a neighborhood of the
closed radius-$10R$ model ball. LC50's minimizing-direction argument
uses only $C^1$ metric convergence (continuous connections), compact
buffers, geodesic ODE uniqueness and the same source-ball coverage.
It applies to this $C^{K-5}$ metric. LC49--LC51 consequently give the
uniform source collar bounds $\alpha=1/8$ and $B=4$. Source smoothing
with $\varepsilon<1/64$, fixed before the model, gives
$\alpha-\varepsilon B>1/16$. LC48's smooth common-field isotopy
then modifies the actual smooth comparison embedding to send this
core to $A_1$. The comparison estimates remain attached to the
unmodified maps. All original LC37 enclosure and gradient conditions
now hold, and LC38 supplies all indicated sublevel types at once.

If the limit is compact, eventual exhaustion includes the WHOLE
model. Its smooth embedding is open and closed in connected $M_i$,
hence a diffeomorphism. Choose the fixed scale large enough for LC43's
diameter bound; that argument and the point cone give the compact
LC37 alternative, with an empty radial band allowed.

Finally LC57's scale and tail bookkeeping applies verbatim: choose
$R$ after the model but before the index, transfer the SAME cone at
that scale with its buffered witness correction, impose the separate
source lower curvature bound on one later tail, and retain the chosen
source radial functions. LC59--LC60 use the smooth source metric and
source flows; their open-ball comparison is unchanged. The noncompact
model has the actual smooth total-normal-bundle map just constructed,
so the unit-open-disk/full-bundle map in LC61 finishes its smooth
open-ball type. The compact branch is already global. The contradiction
in LC58 negates the whole joint conclusion and uses this one fixed
$R$ and one common tail, exactly as before. It requires no all-order
curvature estimate and no smooth nonnegative limit metric.
\end{proof}

The source proof inspected here is Cheeger--Gromoll1.1--1.11,
printed415--422/PDF4--11, especially the relative tangent cone and
supporting half-space on419--420/PDF8--9. The source's omitted smooth
refinement of2.2 is not being assumed. LFR46 expands the same
outward-field, escape and radial-gluing mechanism already retained
from the frozen PC soul files; it is a mathematical finite-category
adapter, not a new declaration check or a probe of the moving checkout.
The carrier construction and its four-derivative metric cost are
project arguments. KL3.11 and11.1--11.4, printed24--25 and59--61,
remain the consumers being compared.

For these zero-model consumers LC81's category gate now has a written
bridge. The remaining compact/noncompact three-model classification
and the general Alexandrov comparison, recognition, splitting and
Tits-cone existence are separate obligations. In particular the
supplied cone in LFR49 is still supplied. Accepted PC release binding
and eventual Lean implementation remain pending; no final full-goal
completion or final full-goal audit is asserted here.


\subsection{Finite nonnegative three-models and their actual core types}
\label{sec:collapse-finite-three-model-types}

LFR45--LFR49 leave a classification question, not another atlas
smoothing question. For a noncompact model we classify the SAME
normal bundle. For a compact model, only its smooth type is needed:
an auxiliary smooth metric may be used to determine that type, while
the finite metric used in all comparisons remains unchanged. The
following arguments use inherited smooth Ricci-flow existence,
extension, curvature evolution, maximum principles, derivative
estimates and smooth compactness, together with the inherited smooth
surface and covering foundations. Their accepted-release bindings
remain to be audited; no new PC declaration or axiom is introduced.

\begin{lemma}[LFR50: compact finite metrics admit an auxiliary smooth model]
\label{lem:collapse-compact-finite-smooth-type-model}
Let $P$ be a nonempty closed connected smooth three-manifold with a
$C^2$ Riemannian metric $g$ and $\sec_g\geq0$. There are a closed
connected smooth three-manifold $Q$, a smooth metric $h$ on $Q$ with
$\sec_h\geq0$, and a smooth diffeomorphism $e:Q\to P$. Consequently
$P$ admits the auxiliary smooth metric $(e^{-1})^*h$.
This conclusion does NOT identify that metric with $g$, or assert
that a Ricci flow of $g$ is smooth in its original coordinates.
\end{lemma}
\begin{proof}
The finite chart/partition construction in LFR17 works for symmetric
two-tensors in every dimension. It gives smooth Riemannian metrics
$g_i\to g$ in $C^2$ on the SAME smooth $P$. Positivity follows from
the positive minimum eigenvalue on a finite compact chart cover.
Continuity of the inverse metric and of the curvature formula gives
constants $B\geq1$, $d_0<\infty$, $v_0>0$, independent of $i$, and
$\epsilon_i\downarrow0$, after passing to a subsequence, such that
\[
 |\Rm_{g_i}|\leq B,\quad \sec_{g_i}\geq-\epsilon_i,
 \quad\operatorname{diam}_{g_i}P\leq d_0,
 \quad\Vol_{g_i}P\geq v_0.
 \tag{LFR50.1}
\]
Increasing finitely many bounds or discarding initial terms is harmless.
No uniform bound on the INITIAL higher curvature derivatives is assumed.

Run the inherited smooth Ricci flow $g_i(t)$ from each $g_i$.
Curvature evolution and the scalar maximum principle give
$|\Rm_{g_i(t)}|\leq B/(1-c_3Bt)$ on its interval of existence.
Choose $T>0$ with $c_3BT<1/2$. The bounded-curvature extension
theorem implies that EVERY flow exists beyond $T$, and
$|\Rm_{g_i(t)}|\leq2B$ for $0\leq t\leq T$. Thus $T$ is fixed
before taking any limit. Integrating $\partial_tg_i=-2\Ric_{g_i}$
gives uniform bilipschitz bounds between $g_i(t)$ and $g_i(0)$,
and hence a common diameter bound $D$ and volume lower bound $v>0$.

We also retain a lower-curvature error tending to zero on this SAME
time interval. In dimension three every two-form is decomposable;
choose its fixed positive normalization so the sectional lower bound
is exactly the curvature-operator lower bound. Work in
evolving orthonormal frames. Its reaction is a fixed quadratic map
$\mathcal Q$, locally Lipschitz with constant $C B$ on the bounded
operator set. For a positive semidefinite operator $A$ and a null
unit vector $w$, $\langle\mathcal Q(A)w,w\rangle\geq0$:
in a diagonal frame the null eigenvalue reaction is a positive
constant times the product of the other two nonnegative eigenvalues.
At a first null vector of $\mathcal R_i+f_i I$, with
$f_i(t)=(\epsilon_i+\delta)e^{C'Bt}$ and $C'$ sufficiently large,
replace $\mathcal R_i$ by the positive semidefinite
$\mathcal R_i+f_i I$. The change in reaction has norm at most
$CBf_i$, whereas $f_i'>CBf_i$. Here we discard finitely many
$i$ and take small $\delta>0$, so $f_i\leq1$ on $[0,T]$ and
the bounded operator set used for $C$ is fixed. The tensor maximum
principle, followed by $\delta\downarrow0$, proves
\[
 \sec_{g_i(t)}\geq-C\epsilon_i e^{C'Bt},\qquad 0\leq t\leq T.
 \tag{LFR50.2}
\]
This is an explicit error estimate, not an application of exact
nonnegativity preservation to the slightly negative $g_i$.

Fix $t_*=T/2$. The inherited Bernstein--Bando--Shi estimates give
ALL curvature derivative bounds for $g_i(t_*)$, with constants
depending on the derivative order, $B$ and $t_*$, not on $i$.
Bishop--Gromov and (LFR50.1)'s evolved volume/diameter bounds give
a uniform positive lower bound for each fixed small ball at every
center: compare its volume ratio with radius $D+1$, which contains
all of $P$, using the uniform Ricci lower bound. The inherited smooth CGT estimate then gives a common
positive injectivity radius. Smooth pointed compactness is now
applicable with its ACTUAL all-order hypotheses. Let $(Q,h,q)$ be
its complete connected smooth limit, retaining the pointed
whole-buffer embeddings and distance comparisons.

Distances between any two limit points are at most $D$, by applying
those comparisons on a compact buffer containing both points and
then passing to the limit. Thus $Q$ is bounded. Completeness and
finite-dimensional Riemannian properness make it compact.
The convergence domains eventually contain ALL of $Q$. Each
resulting smooth equal-dimensional embedding $e_i:Q\to P$ is open
by the inverse function theorem and closed by compactness; its image
is nonempty, so connectedness gives a diffeomorphism. Smooth metric
convergence and (LFR50.2) give $\sec_h\geq0$. Choose one such $e_i$.
The construction never takes a Ricci flow limit at time zero, never
identifies its positive-time metric with $g$, and never changes a
metric in an existing local packet.
\end{proof}

\begin{lemma}[LFR51: orientation determines the finite soul bundle types]
\label{lem:collapse-oriented-finite-soul-types}
Let $N$ be an orientable model in LFR45--LFR47, retaining the
original order $m\geq8$ and its actual smooth bundle model
$\widehat E\to\widehat S$. Its smooth total-space and disk-bundle
types are precisely among the following possibilities:
\[
\begin{array}{c|c|c}
 \widehat S&\widehat E&D(\widehat E)\\\hline
 \{*\}&\mathbb R^3&D^3\\
 S^1&S^1\times\mathbb R^2&S^1\times D^2\\
 S^2&S^2\times\mathbb R&S^2\times[-1,1]\\
 T^2&T^2\times\mathbb R&T^2\times[-1,1]\\
 \mathbb {RP}^2&o(\mathbb {RP}^2)&D(o(\mathbb {RP}^2))\\
 \mathbb K&o(\mathbb K)&D(o(\mathbb K))
\end{array}
\tag{LFR51.1}
\]
Here $\mathbb K$ is the Klein bottle and $o(S)$ is the real
orientation line bundle. These are smooth bundle identifications,
including closed disk bundles after matching the fiber norms.
The metric $g$ on $N$ is not replaced by a standard metric from
this table.
\end{lemma}
\begin{proof}
The zero-dimensional connected soul is a point. A closed connected
one-dimensional soul is a circle. If its dimension is two, total
geodesy and the Gauss equation give nonnegative Gaussian curvature
for the induced finite metric. Pulling back through the compatible
smooth carrier preserves at least $C^2$ regularity. LFR17's
mollification and Gauss--Bonnet argument gives $\chi(S)\geq0$.
For a nonorientable $S$, apply the same argument on its orientation
double cover and use $\chi(\widetilde S)=2\chi(S)$. The inherited
smooth classification of closed surfaces therefore leaves exactly
$S^2,T^2,\mathbb {RP}^2,\mathbb K$. In the last Euler-characteristic
zero cases the continuous nonnegative Gaussian curvature vanishes,
since its integral is zero. This uses the retained finite soul data;
we do not apply LFR47 again to the already downgraded metric.

Along the zero section, $TN|_S=TS\oplus\nu S$ as finite vector
bundles. The ambient orientation is transported by LFR47's actual
map to the smooth total space of $\widehat E$, where the same
orientation relation holds. Over the circle the rank-two bundle is
therefore oriented. Cut the circle at one point, trivialize over the
interval with orthonormal frames, and let $A\in SO(2)$ be the
endpoint transition. A smooth path in $SO(2)$ from $I$ to $A$,
constant near the endpoints, changes the trivialization to a
periodic one. This proves the actual norm-preserving trivialization.

For a surface soul the bundle has rank one. A choice of unit normal
and the ambient orientation determine an orientation of $TS$, and
changing the normal reverses that orientation. Consequently its
unit sphere bundle is canonically the orientation double cover of
$S$. In smooth local frames this identification is locally constant
in the sign and hence smooth. Extending radially gives
\[
 \widehat E\simeq
 (\widetilde S_{\rm or}\times\mathbb R)/
       ((x,t)\sim(\tau x,-t)),
 \tag{LFR51.2}
\]
with the absolute-value fiber norm. For orientable $S$ its orientation
cover has two components, so this is the trivial line bundle; for
nonorientable $S$ it is $o(S)$. The smooth diffeomorphism of the base
to its standard surface lifts to its orientation cover. These
identifications thus preserve unit disks and sphere boundaries.
For any preliminary bundle isomorphism with a different smooth norm,
fiberwise polar normalization makes it isometric without changing
its base map. Composing with LFR47 keeps the ACTUAL soul/core witness.
\end{proof}

\begin{lemma}[LFR52: explicit twisted cores and their ends]
\label{lem:collapse-twisted-core-identifications}
In (LFR51.1), $D(o(\mathbb {RP}^2))$ is diffeomorphic, including
boundary, to $\mathbb {RP}^3\setminus\operatorname{int}D^3$.
The bundle $D(o(\mathbb K))$ is the orientable twisted interval
bundle over the Klein bottle and has one torus boundary component.
The first, second, fifth and sixth total spaces in (LFR51.1) have
one end; the third and fourth have two ends. Every assertion uses
the indicated ACTUAL disk-bundle type, rather than a recognition
of an arbitrary core from its interior alone.
\end{lemma}
\begin{proof}
Write the first twisted disk bundle as
$(S^2\times[-1,1])/((x,t)\sim(-x,-t))$ and fix $0<a<1$.
The smooth map
\[
 [x,t]\longmapsto[\sqrt{1-a^2t^2}\,x,at]\in S^3/\{\pm1\}
       =\mathbb {RP}^3
 \tag{LFR52.1}
\]
identifies it with the closed equatorial band $|u|\leq a$ in
projective three-space. Its inverse is obtained by choosing either
unit representative, reading $t=u/a$, and normalizing the first
three coordinates. The two choices differ by the bundle involution.
The two complementary open polar caps of $S^3$ are interchanged by
the antipodal map; their quotient is one standard open three-ball.
The displayed map is a diffeomorphism up to the boundary, where
the first three coordinates remain nonzero. Thus no interior-only
homeomorphism has been promoted to a smooth core identification.

For the Klein bottle use its standard orientation cover
$T^2=\mathbb R^2/\mathbb Z^2$, with free deck involution
\[
 \tau(x,y)=(x+1/2,-y),\qquad
 D(o(\mathbb K))=(T^2\times[-1,1])/
       ((x,y,t)\sim(x+1/2,-y,-t)).
 \tag{LFR52.2}
\]
The involution reverses both the surface orientation and the normal
coordinate, so the total space is orientable. Its two boundary
tori are interchanged, giving exactly one torus boundary. The same
formula with $\mathbb R$ in place of $[-1,1]$ gives a flat model
for this smooth bundle TYPE. It does not prove that the original
metric on $N$ is that flat metric.

For any positive-rank normed bundle over a compact base, its closed
radius-$R$ disk bundles are a cofinal compact exhaustion. Outside
one such disk, radial normalization is a smooth product
$S(\widehat E)\times(R,\infty)$. Hence its number of ends is the
number of components of the unit sphere bundle. Rank at least two
gives a connected sphere bundle over the connected base. For rank
one, this sphere bundle is the orientation cover from (LFR51.2),
which has two components exactly when $S$ is orientable. This proves
all the end assertions. LFR47's global diffeomorphism transports the
cofinal exhaustion to $N$. This is an end-count proof, not a proof
of the stronger same-metric isometric splitting clause of KL3.11;
that clause stays separate from these topological consumers.
\end{proof}

\begin{theorem}[LFR53: compact finite three-model classification]
\label{thm:collapse-compact-finite-classification}
Let $P$ be closed, connected and orientable with a $C^2$ metric of
nonnegative sectional curvature on a smooth carrier. Then its
smooth type is $S^1\times S^2$, $\mathbb {RP}^3\#\mathbb {RP}^3$,
a spherical space form, or a compact orientable flat quotient.
The last type can equivalently be written $T^3/\Gamma$ for a flat
torus and a finite group of orientation-preserving free isometries.
The classification determines a diffeomorphism type; it supplies
no comparison estimate for an auxiliary metric.
\end{theorem}
\begin{proof}
LFR50 supplies a smooth nonnegative metric on a manifold smoothly
diffeomorphic to $P$. Run a short smooth flow from this new metric.
The inherited tensor strong maximum principle makes the image of
the nonnegative curvature operator a parallel Lie subalgebra of
$\operatorname{so}(3)$ for a positive time interval. Such a subalgebra
has dimension zero, one or three: a two-dimensional subspace cannot
be a subalgebra since the cross product of two independent vectors
is nonzero and perpendicular to it.

In dimension three the metric has positive sectional, hence positive
Ricci, curvature. The inherited smooth Hamilton positive-Ricci
theorem gives a constant positive curvature metric on the SAME
closed manifold. Its universal Riemannian cover is the round
$S^3$ after a constant rescaling. The deck group is finite, acts
freely, and preserves orientation, giving the stated spherical type.
In dimension zero the metric is flat. Its complete simply connected
cover is Euclidean three-space, by parallel frames or the flat
exponential/Jacobi argument of LFR22 in dimension three. The deck
group is discrete and cocompact and acts freely by Euclidean isometries.
The retained Bieberbach lattice theorem supplies a normal full
translation lattice of finite index, giving $T^3/\Gamma$ with its
possibly nonrectangular flat torus metric. The finite group acts
freely and preserves orientation because the deck group does.

In dimension one, take the universal cover BEFORE choosing a unit
parallel form; the parallel line subbundle on $P$ need not be
orientable. On the simply connected cover, its Hodge dual gives
a global parallel unit vector, and inherited de Rham splitting
gives $(\Sigma,h)\times\mathbb R$. The Gaussian curvature of
$\Sigma$ is positive. It has a positive uniform lower bound because
it descends from the nonzero curvature eigenvalue on the compact
manifold at this fixed positive time. Bonnet--Myers therefore makes
$\Sigma$ compact. It is simply connected as a factor of the
universal cover, so smooth surface classification makes it $S^2$.
The parallel line is unique, so every deck transformation has form
\[
 (x,t)\longmapsto(Ax,\varepsilon t+b),
 \quad A\in\operatorname{Isom}(\Sigma,h),\quad\varepsilon=\pm1.
 \tag{LFR53.1}
\]
The projection of the deck group to $\operatorname{Isom}(\mathbb R)$
is injective. Indeed an element of its kernel is an orientation-
preserving diffeomorphism $A$ of $S^2$; such a map has a fixed point.
Indeed, if $A(x)\neq x$ everywhere, normalizing
$(1-s)A(x)-sx$ gives a homotopy from $A$ to the antipodal map,
contradicting their degrees $1$ and $-1$. Freeness therefore forces
the deck element to be the identity. Its image is discrete: a sequence of
distinct projected isometries tending to the identity has, by
compactness of $\operatorname{Isom}(\Sigma,h)$, a subsequence of
the full isometries whose pairwise differences tend to the identity.
This contradicts proper discontinuity of the deck action. Compactness
of $P$ makes the projected action cocompact. A discrete cocompact
subgroup of line isometries is either a translation lattice or an
infinite dihedral group, as follows by taking its least positive
translation and, if present, one reflection.

For the translation lattice the quotient is the mapping torus of
an orientation-preserving smooth diffeomorphism of $S^2$. The
inherited smooth surface mapping-class theorem makes that map
smoothly isotopic to the identity; its mapping torus is smoothly
$S^2\times S^1$. For the dihedral group choose two adjacent
reflections. Their lifts in the deck group square to the identity
by injectivity of the projection. Their actions on the sphere are
free orientation-reversing involutions. The quotient of a sphere
by either involution is a smooth $\mathbb {RP}^2$ by surface
classification; lifting this identification to universal covers
conjugates its action smoothly to the antipodal action. Each end
of a fundamental interval thus has a neighborhood of type
$D(o(\mathbb {RP}^2))$. The two neighborhoods are joined by a
sphere cylinder. LFR52 identifies this union with the connected
sum $\mathbb {RP}^3\#\mathbb {RP}^3$, using the usual smooth
sphere gluing convention. All cases are smooth identifications.
Only the auxiliary compact metric was used for this classification.
\end{proof}

\begin{theorem}[LFR54: classify the same finite local packet]
\label{thm:collapse-classified-finite-zero-packet}
Supply LFR49's finite joint zero packet for the original oriented
three-manifolds, with $K\geq10$, and the same-model at-most-one-end
conclusion of LC77. Then every indicated smooth radial sublevel has
one of the following smooth types:
\[
 D^3,\quad S^1\times D^2,\quad
 \mathbb {RP}^3\setminus\operatorname{int}D^3,\quad
 D(o(\mathbb K)),
 \tag{LFR54.1}
\]
or one of the compact types in LFR53. These identifications apply
to the ACTUAL sublevels and include their boundary; the distance
balls retain LC34's homeomorphisms with those sublevels. All metric
comparisons, radial functions, cone maps, fixed scales, common
tails, embeddings and transversality margins in the original
packet are retained.
\end{theorem}
\begin{proof}
For a noncompact limit, LFR49 and LC38 identify every indicated
sublevel with the SAME $D(\widehat E)$ from LFR47. LFR51 classifies
that bundle. LFR52 and LC77 exclude its two trivial surface-line
bundles and identify the remaining four disk types. Compose the
existing smooth core identification with these fixed bundle maps;
do not move its embedded domain or its boundary. This gives all
sublevel types simultaneously, without any new tolerance or tail.
The open-ball comparison of LC61 and the boundary/interior
homeomorphisms of LC34 remain the existing ones.

For a compact limit keep its original smooth carrier and finite
metric. LFR48 supplies the actual global smooth comparison
diffeomorphisms used in LFR49's compact case. The original metric
has at least $C^2$ regularity, so LFR53 classifies that carrier.
Compose the topological-type map with the global diffeomorphism.
The auxiliary metric produced inside LFR50 is discarded after
determining the smooth type: it is NEVER substituted into a cone
approximation, a source comparison map, or a radial estimate.
Hence there is no curvature or distance change in the local packet.
\end{proof}

\paragraph{Source comparison and remaining boundary.}
KL3.11, printed24--25/PDF19--20, asserts the finite-category
classification and cites Cheeger--Gromoll for the noncompact case
and Hamilton or Simon for the compact case. Cheeger--Gromoll8.1,
printed438--439/PDF27--28, classifies more strongly up to isometry;
LFR51--LFR52 prove the smooth bundle/core information required here.
The separate two-ended isometric splitting assertion has not been
silently inferred from a diffeomorphism. LC75 still retains its
general splitting input AC43. LFR53 expands the smooth three-case
argument in GSM77, section \emph{Rigidity of 3-manifolds with
nonnegative curvature}, with orientation and deck-action details;
the smooth Ricci-flow and surface foundations remain inherited.

The research supplement's Simon title was incorrect: the retained
arXiv 0612095v1 is \emph{Ricci flow of almost non-negatively curved
three manifolds}; the published2009 version has different numbering.
\emph{Deformation of $C^0$ Riemannian metrics in the direction of
their Ricci curvature} is a distinct2002 paper with author corrections.
LFR50 uses the simpler bounded-curvature compact argument and
inherited SMOOTH flow estimates. It does not import the2002
noncompact weak-metric assertions or the stronger2009 unbounded-
curvature existence theorem. Exact source versions, proof reading
and correction scope are recorded in reference\_checks\_revision136.

Thus the finite category and classification needed by the specified
zero-core consumer now have written bridges under the retained
foundations. General Alexandrov comparison, recognition AC46,
splitting AC43 and the LC21 Tits-cone producer remain distinct.
The source-scale simultaneous parameter assignment and final local
export must still be checked. This is not the final full-goal audit,
an unconditional local-collapse theorem, or Lean completion.


\subsection{Cones at infinity for the actual finite models}
\label{sec:collapse-finite-model-cone}

This subsection supplies LC21 for the finite Riemannian models used
by LC81 and LFR49. It does not assert the corresponding theorem for
an arbitrary Alexandrov space. Throughout, $N$ is a connected smooth
$n$-manifold without boundary, with a complete $C^m$ Riemannian metric
$g$, $m\geq4$, and $\sec_g\geq0$. Its carrier may be the consumer
carrier of LFR47; the metric is the SAME finite metric. The retained
Riemannian foundation supplies geodesic ODEs, parallel transport,
first/second variation, normal neighborhoods and Hopf--Rinow in their
finite differentiability range. The proofs below specify the finite
adapter and subsequent cone construction; no new PC declaration is
asserted or probed.

\paragraph{Source comparison and a qualification to the earlier research.}
GSM77, archived TeX lines4334--4515 and4582--4703, gives the variation
formulas and squared-distance support calculation; lines7476--7590
contain the lemma labelled \code{Eagles tour} and the nonnegative
Toponogov proof. The argument below expands its support-function step
and works at finite $m$. Mashiko--Nagano--Otsuka, \emph{The asymptotic
cones of manifolds of roughly non-negative radial curvature},
J.\ Math.\ Soc.\ Japan57(2005),55--68, gives the ray-angle route:
Theorem1.1 p58/PDF4, Proposition2.1 and2.3 p61/PDF7, and the proof
of Claim2.4 and Theorem0.1 pp62--63/PDF8--9. We specialize to
nonnegative sectional curvature and write the actual maps and radial
identity needed here; the broader radial-curvature theorem is not
imported. Their introduction p55/PDF1 reports a gap in Kasue1988's
compactification argument, referring to Drees1994. Kasue's
Section3.1--3.4 pp607--614/PDF16--23 was read for comparison, but
its compactification and sphere-flow maps are NOT dependencies here.
This qualifies the earlier research supplement's proposed use of
Kasue. Drees's original discussion has not been inspected, so we do
not claim an independent diagnosis of its precise error. KL3.3
p23/PDF18 and11.1 pp59--60/PDF54--55 specify the consumers; the
current author correction sheet has no new entry for them. Exact
source identities and the limited errata search are recorded in
\code{reference_checks_revision137.md}.

\begin{lemma}[Finite squared-distance comparison; LFR55]
\label{lem:collapse-finite-ray-comparison}
On $(N,g)$ as above, for every unit speed geodesic $\sigma$ and
$q\in N$, the function $t^2-d(q,\sigma(t))^2$ is convex on its
parameter interval. Consequently every minimizing hinge based at $p$
with initial unit vectors $v,w$ and arm lengths $a,b$ satisfies
\[
 d(\gamma_v(a),\gamma_w(b))\leq\lVert av-bw\rVert_{g_p}.
\]
Its Euclidean comparison angle $\theta_{v,w}(a,b)$ is nonincreasing
in each positive arm length, as long as both arms remain minimizing.
The metric space $(N,d_g)$ satisfies global nonnegative four-point
comparison. These conclusions use finite variation calculus, not a
smooth replacement of $g$ or general Alexandrov globalization.
\end{lemma}
\begin{proof}
Fix $t_0$ with $q\ne\sigma(t_0)$ and a unit speed minimizing segment
$c:[0,L]\to N$ from $q$ to $\sigma(t_0)$. Parallel translate
$\dot\sigma(t_0)$ along $c$ to obtain a unit field $W$ and put
$V(s)=(s/L)W(s)$. For small $h$ the variation
$c_h(s)=\exp_{c(s)}(hV(s))$ starts at $q$ and ends EXACTLY at
$\sigma(t_0+h)$. Compactness of this one segment gives one variation
interval; no global injectivity bound is used. The finite geodesic
flow and $m\geq4$ make this variation and its length $\ell(h)$ at
least $C^2$. Write $u=\langle W,\dot c\rangle$, constant along $c$.
First and second variation, with geodesic endpoint paths, give
\[
 \ell'(0)=u,\qquad
 \ell''(0)=\frac{1-u^2}{L}
 -\int_0^L\langle\operatorname{Rm}(V,\dot c)\dot c,V\rangle\,ds
 \leq\frac{1-u^2}{L}.
\]
Thus $(\ell^2)''(0)\leq2$. The function
$\psi(t) = t^2-\ell(t-t_0)^2$ is a local $C^2$ LOWER support for
$f(t)=t^2-d(q,\sigma(t))^2$, with $\psi''(t_0)\geq0$.
This holds for any minimizing choice, including a cut point. If
$q=\sigma(t_0)$, a normal neighborhood gives
$d(q,\sigma(t))^2=(t-t_0)^2$ locally, so $f$ itself is affine there.

For completeness, such lower supports imply convexity. If $f$ exceeds
one of its chords on $[a,b]$, subtract that chord and then
$\delta(t-a)(b-t)$ for a sufficiently small $\delta>0$. The result
has a positive interior maximum. At a maximizing point, the resulting
concave quadratic plus this maximum is an upper support $Q$ for $f$,
with $Q''=-2\delta$. Its difference from the local lower support
has a local minimum zero, forcing $\psi''\leq Q''<0$, a
contradiction. This proves the stated convexity without differentiating
distance at the cut locus.

For the hinge inequality, set $q=\gamma_v(a)$ and
$\sigma(t)=\gamma_w(t)$. Reverse the chosen segment from $p$ to $q$
and vary its endpoint along $\sigma$. First variation supplies an
upper support for squared distance at zero with derivative
$-2a\langle v,w\rangle$. Hence the right derivative of the convex
$f$ is at least $2a\langle v,w\rangle$. The supporting-line inequality
$f(b)\geq f(0)+2ab\langle v,w\rangle$ is exactly the asserted hinge
bound. A zero arm follows directly.

For angle monotonicity, keep $q=\gamma_w(b)$ fixed and apply convexity
along $\gamma_v|_{[0,a]}$. If $0<u\leq1$ and
$D=d(\gamma_v(a),q)$, then
\[
 d(\gamma_v(ua),q)^2
 \geq uD^2+(1-u)b^2-u(1-u)a^2.
\]
Substitution into the Euclidean cosine law gives
$\cos\theta_{v,w}(ua,b)\leq\cos\theta_{v,w}(a,b)$.
Interchange the two arms for the other monotonicity. Degenerate
triangles cause no problem since $\arccos$ is continuous on $[-1,1]$.
Finally, at any center choose minimizing segments to three other
points. The hinge bound makes each comparison angle at most the
corresponding angle of initial unit vectors. The three pairwise
angles of three unit vectors sum to at most $2\pi$: use the spherical
triangle inequality through the antipode of the middle vector.
This is the global four-point inequality, with coincident-point
conventions handled as in AC35. All derivatives above have order at
most two in the variation and curvature uses two metric derivatives;
$m\geq4$ suffices for the stated finite constructions.
\end{proof}

\begin{lemma}[Compact ray parameters and uniform distance limits; LFR56]
\label{lem:collapse-finite-ray-distance-limit}
Suppose $N$ is noncompact and fix $p\in N$. Let $A\subset S_pN$
be the initial directions of minimizing rays and put
$E=\{rv:r\geq0,\ v\in A\}\subset T_pN$, $E_S=E\cap\overline B(0,S)$.
Then $A$ is nonempty compact and $E_S$ is compact. For $R>0$ define
\[
 \Gamma_R(rv)=\gamma_v(Rr),\qquad \Gamma_R(0)=p,
 \qquad D_R(z,w)=R^{-1}d(\Gamma_Rz,\Gamma_Rw).
\]
There is a pseudometric $D$ on $E$ such that $D_R\downarrow D$
pointwise as real $R\to\infty$, uniformly on each $E_S\times E_S$.
For $r,s>0$ its formula is
\[
 D(rv,sw)^2=r^2+s^2-2rs\cos\theta_\infty(v,w),\qquad
 \theta_\infty(v,w)=\inf_{a,b>0}\theta_{v,w}(a,b).
\]
At a zero radius the distance is the other radius. Both $D_R$ and
$D$ are bounded above by the ambient Euclidean distance on $E$.
\end{lemma}
\begin{proof}
Properness and noncompactness give minimizing segments from $p$
whose lengths tend to infinity. Compactness of $S_pN$ gives a
subsequence of initial directions converging to $v$. Continuous
dependence for the complete finite geodesic ODE on every fixed time
interval shows $d(p,\gamma_v(t))=t$ for every $t\geq0$. Thus $v\in A$.
The same argument for a convergent sequence of ray directions proves
$A$ closed. Consequently $A$ and all $E_S$ are compact; continuity
of $\Gamma_R$ also holds at the origin by its exact radial distance.

LFR55 makes $\theta(a,b)$ decrease in each variable. It follows that
$\theta(Rr,Rs)\to\inf_{a,b>0}\theta(a,b)$ for any FIXED positive
$r,s$: once both $Rr$ and $Rs$ exceed prescribed arms, monotonicity
gives the upper bound, while the infimum gives the lower bound.
The cosine law proves the formula and convergence; it also shows
$D_R$ decreases with $R$. Each $D_R$ is a pseudometric since it is
a pullback of scaled distance. The triangle inequality passes to the
limit, so the displayed formula does not assume a cone triangle
inequality or an angular quotient theorem.

The hinge bound gives $D_R(z,w)\leq\lVert z-w\rVert$. In particular
\[
 |D_R(z,w)-D_R(z',w')|
 \leq\lVert z-z'\rVert+\lVert w-w'\rVert,
\]
and the same holds for $D$. A finite Euclidean $\delta$-net in $E_S$
therefore upgrades convergence on its finitely many pairs to uniform
convergence on $E_S^2$. More explicitly, the net errors plus $4\delta$
bound all errors. Hence
\[
 \omega_S(R):=\sup_{z,w\in E_S}(D_R(z,w)-D(z,w))\longrightarrow0
\]
for ALL sufficiently large real $R$, not merely a selected sequence.
No uniform rate over different model manifolds is asserted.
\end{proof}

\begin{lemma}[The proper quotient cone and its actual radial maps; LFR57]
\label{lem:collapse-finite-ray-quotient-cone}
The metric quotient $C=E/\{D=0\}$, with $o=[0]$, is complete and
proper, has Hausdorff dimension at most $n$, and has actual AC82 maps
$H_t[z]=[tz]$ for all real $t\geq0$.
\end{lemma}
\begin{proof}
The pseudometric triangle inequality makes the quotient distance
$d_C([z],[w])=D(z,w)$ well defined and separates classes. Its radial
norm is $d_C(o,[z])=\lVert z\rVert$. The quotient map from $E$ with
its ambient Euclidean metric is $1$-Lipschitz and maps $E_S$ ONTO
$\overline B_C(o,S)$. Thus these balls are compact. Every closed
bounded ball about any other center is closed in a larger apex ball,
and every Cauchy sequence lies in such a compact ball. This proves
properness and completeness. The $1$-Lipschitz image bound for
Hausdorff dimension, applied to $E_S\subset\mathbb R^n$ and then
to the countable union over integer $S$, gives $\dim_H C\leq n$.
There is no claim that the intrinsic angular boundary has dimension
$n-1$ or that it is connected.

The formula in LFR56 implies $D(tz,tw)=tD(z,w)$, including $t=0$,
so $H_t$ descends to the quotient. For $x=[rv]$, $y=[sw]$ with
positive radii, its SAME formula with independent radii gives
\[
 d_C(H_a x,H_b y)^2
 =a^2r^2+b^2s^2-2abrs\cos\theta_\infty(v,w).
\]
Here $rs\cos\theta_\infty(v,w)=B_o(x,y)$ by the formula at $a=b=1$.
Zero radii and zero parameters follow from the radial norm. Also
$H_0x=o$ and $H_1x=x$. These are precisely AC82, not just
self-similarity under simultaneous rescaling. No identification with
an intrinsic Tits boundary is required for this package.
\end{proof}

\begin{lemma}[Full-ball approximation at every large scale; LFR58]
\label{lem:collapse-finite-cone-actual-maps}
For every fixed $S>0$, define
\[
 \beta_S(R)=\sup_{x\in\overline B_N(p,RS)}
 R^{-1}d(x,\Gamma_R(E_S)).
\]
Then $\beta_S(R)\to0$ as real $R\to\infty$. There are actual
pointed maps from the closed radius-$S$ ball in $R^{-1}N$ ONTO
$\overline B_C(o,S)$ with distortion at most
$2\beta_S(R)+\omega_S(R)$. Consequently for every $0<\varepsilon<1$
there is $R_0$ such that EVERY real $R\geq R_0$ admits a whole-space
pointed KL $\varepsilon$-approximation $R^{-1}N\to C$.
\end{lemma}
\begin{proof}
All distances to $\Gamma_R(E_S)$ attain their minima; that set and
the source closed ball are compact. If $\beta_S$ fails to tend to
zero, choose $R_i\to\infty$ and $x_i$ for which the displayed
scaled distance is bounded below by some $\eta>0$. Put
$r_i=d(p,x_i)/R_i\leq S$ and pass to a subsequence with $r_i\to r$.
The case $r=0$ contradicts this lower bound because $p\in\Gamma_R(E_S)$.
If $r>0$, choose minimizing directions $v_i$ to $x_i$ and pass to
$v_i\to v$. Since their segment lengths $r_iR_i$ tend to infinity,
fixed-time ODE convergence, as in LFR56, shows $v\in A$. LFR55
applied at the ACTUAL long length $r_iR_i$ gives
\[
 R_i^{-1}d\bigl(x_i,\Gamma_{R_i}(r_iv)\bigr)
 \leq r_i\lVert v_i-v\rVert\longrightarrow0,
\]
a contradiction. Local convergence of geodesics alone would not
justify an estimate at these growing times; the hinge estimate does.

For each $x$ in the source ball choose $z_R(x)\in E_S$ minimizing
$d(x,\Gamma_Rz)$ and define $f_R(x)=[z_R(x)]$, taking $z_R(p)=0$.
The distance from $x$ to its chosen ray point is at most $R\beta_S$.
Two triangle inequalities followed by LFR56 give
\[
 |d_C(f_Rx,f_Ry)-R^{-1}d_N(x,y)|
 \leq2\beta_S(R)+\omega_S(R).
\]
These choices need not be continuous. They nevertheless give EXACT
surjectivity: for $[z]\in\overline B_C(o,S)$ take $x=\Gamma_Rz$.
The minimum is zero, so $D_R(z_R(x),z)=0$; since $D\leq D_R$,
we have $[z_R(x)]=[z]$. This also handles distinct rays that meet.
Moreover this witness has source radius exactly $\lVert z\rVert$.

Take $S=\varepsilon^{-1}$ and choose $R_0$ so that the distortion
bound is less than $\varepsilon/4$ for all $R\geq R_0$. Extend the
chosen map by $o$ outside the closed source ball. It is pointed and
has the required distortion on KL's tested open ball. Every target
point of radius less than $\varepsilon^{-1}-\varepsilon$ has the
exact preimage just constructed at the SAME radius, hence strictly
inside that tested source ball. Thus both clauses of KL3.2 and MC06
hold, including the coverage convention and exact basepoint. The
metric tensor normalization is $g/R^2$, with distance $d_g/R$.
\end{proof}

\begin{theorem}[LC21 for finite Riemannian consumers; LFR59]
\label{thm:collapse-finite-model-cone-package}
Every complete connected finite model $(N,g,p)$ above has an LC21
package. Its cone is complete, proper, geodesic, globally
nonnegatively curved in the four-point sense, and has Hausdorff
dimension at most $n$. For noncompact $N$ it is the SAME cone of
LFR56--LFR58; for compact $N$ it is the one-point cone. This supplies
the cone premise for LC23--LC24 and LC55--LC61 when their limits
are the finite Riemannian models of LFR14--LFR15 or LFR49. It does
not discharge their other hypotheses or the general Alexandrov
splitting and recognition inputs.
\end{theorem}
\begin{proof}
For noncompact $N$, LFR57 and LFR58 give LC21 with one target fixed
BEFORE $\varepsilon$ and $R$. Apply MC23 to any sequence of scales
tending to infinity, using the proper complete target and geodesic
source spaces, to obtain the geodesic property. For four points of
$C$, choose representatives $z_0,z_1,z_2,z_3\in E_S$ in one bounded
set. Their actual lifts $\Gamma_Rz_j$ have scaled pairwise distances
$D_R(z_i,z_j)\to D(z_i,z_j)$. LFR55's four-point comparison holds
at $\Gamma_Rz_0$, not merely at the original basepoint $p$.
AC35's cosine-law continuity passes it to the target; separated
limit points stay separated for large $R$, and its degenerate
conventions handle the remaining cases. Dimension and AC82 were
proved in LFR57. No general Alexandrov compactness or cone-existence
producer is used in this argument.

For compact $N$, use the constant whole-space map to a point;
its distortion is at most $\operatorname{diam}(N)/R$, and coverage
and radial identities are exact. This gives the same all-large-real
quantifier. The point has all the additional properties stated.

LFR14 produces metric order $K-1$, and LFR47's one-time change of
consumer atlas leaves the SAME metric space with metric order $K-5$.
For the standing $K\geq10$ both orders exceed four. Choose the cone
on the original model and retain it after that atlas change by the
identity isometry. The auxiliary compact classification metric of
LFR50 is NEVER substituted here. Thus LFR49 can retain its radial
function, normal-flow map, actual smooth cores and packet embeddings
while obtaining its previously supplied cone from this theorem.

For each fixed model and tolerance, the threshold $R_0$ may depend
on that model. A uniform threshold over a family has not been proved
or assumed. LC23's compactness contradiction first selects one limit
model, then ONE sufficiently large fixed scale for its cone, then
one common late index tail; LFR49's simultaneous annular/core
construction keeps this order. LC58 still requires the compactness
premise for EVERY index/basepoint sequence. In this way the existing
finite-window proof obtains its uniform conclusion without a new
model-uniform convergence rate. The radial identity, actual pointed
maps, scale and common tail now have their required finite-category
producer. General AC43 splitting, AC46 recognition and the final
simultaneous parameter/export binding remain separate obligations.
\end{proof}

\paragraph{Scope after this construction.}
The supplied-cone qualifications in the earlier LC21, LC24 and
LFR45--LFR54 records are now discharged for these finite Riemannian
consumers by LFR55--LFR59. Their historical text remains unchanged.
This is a written proof under retained finite Riemannian and metric
foundations, not an elaborated Lean declaration. It does not prove
a cone theorem for every Alexandrov space, smooth the finite metric,
or bind the moving PC snapshot. The source-checked ray construction
replaces the proposed Kasue route only for the specified consumers.


\subsection{Closing the remaining local comparison consumers}
\label{sec:collapse-written-comparison-consumers}

The larger comparison buffer of ALG05 does not justify replacing
AC02's old fixed margins without checking each use. The following
arguments supply the remaining comparison uses in the interior local
packets. Shortening arms by the triangle inequality BEFORE applying
comparison is the useful additional step. The original LC27 and LC65
hypotheses and constants are retained. LC69 keeps the ORIGINAL
radial center and its error proportional to the prefix length.
LFR25's original edge buffer already suffices.

The source applications are KL 11.3--11.5, printed 60--62/PDF
55--57, and 9.7/9.12, printed 50,52--53/PDF45,47--48.
The exact-radius strainer and prescribed-coordinate conventions
are KL4.15(2), printed31--32/PDF26--27, and4.28,
printed36--37/PDF31--32. These bodies were reopened; the local
comparison proof is the written ALG01--ALG08, not the sharp
localization statement quoted by the sources. As already recorded
at LC25, KL11.3 prints only a lower bound for the outward point's
distance; our existing LC25 supplies the necessary upper bound.
No new source correction is attributed to the authors.

In the arguments below the retained smooth Riemannian foundation
supplies completeness/properness, minimizing geodesics, the local
sectional-curvature-to-four-point implication, and identification of
ALG01's germ limits with tangent angles. The latter follows from
equal-small-arm distance expansion in a metric chart. Global or
buffered Alexandrov comparison is supplied by ALG05/ALG06 here;
it is not an additional inherited hypothesis. The finite nonnegative
models continue to use LFR55 on their unchanged finite metrics.

\begin{theorem}[LCP01: the original annular direction estimate with its original buffer]
\label{thm:collapse-annular-directions-written-comparison}
The full conclusion of LC27 holds under its ORIGINAL hypotheses,
including curvature control only on $B(p,20b)$ and its stated
$\delta$ bound, without the sharp local comparison input in AC02.
The same one outward direction works against EVERY inward
minimizing direction at each point.
\end{theorem}
\begin{proof}
For $r=d(p,q)\in[a,b]$, keep LC25's point $q'$ and one minimizing
segment from $q$ to $q'$, of length $\ell\in[r,r+15\delta)$.
Put $E=\ell-r<15\delta$ and $h=2r/5$. For ANY minimizing segment
from $q$ to $p$, let $u$ be its point at distance $h$ from $q$;
let $v$ be the point at distance $h$ on the fixed outward segment.
The triangle inequality, using both discarded tails, gives
\[
 2h-E\leq c:=d(u,v)\leq2h.
\]
Here $E<r/2<2h$, since $\delta<a/30$. For curvature $-\kappa^2$
the comparison angle $\omega$ of $(q;u,v)$ obeys, when $\kappa>0$,
\[
 \begin{split}
 1+\cos\omega
 &=\frac{\cosh(2\kappa h)-\cosh(\kappa c)}
          {\sinh^2(\kappa h)}\\
 &\leq2\kappa E\coth(\kappa h)
 \leq\frac{6E}{r}<\frac{90\delta}{a}<1-\cos\tau.
 \end{split}
\]
Indeed $\kappa h\leq2/15$, and for $0\leq t\leq2/15$,
\[
 t\coth t\leq\cosh t\leq(1-t^2)^{-1}
 \leq225/221<6/5,
\]
with the limiting value at zero. The power series for $\cosh$
proves the middle inequality, and $\sinh t\geq t$ proves the
first. Thus $2E/h$ times $6/5$ is $6E/r$.
For $\kappa=0$, the cosine formula gives
$1+\cos\omega\leq2E/h=5E/r<90\delta/a$.
These bounds include $E=0$ and collinear triangles.

Use $U=B(p,20b)$ with its intrinsic metric. The points and chosen
arms lie in $B(p,2b)$; all joining distances used above agree
with intrinsic distances, since ambient minimizing joins remain
in $B(p,4b)$. For each possible $u$, the intrinsic closed
$18b$-ball at $u$ is complete: its ambient Cauchy limit has
radius at most $d(p,u)+18b\leq(3/5+18)b<20b$, and local
metric equality gives intrinsic convergence. ALG05 applies to
our hinge with endpoint $u$, since
\[
              2h=4r/5\leq4b/5<18b/20.
\]
It makes the angle between the ORIGINAL initial directions at $q$
at least $\omega>\pi-\tau$. The outward segment was fixed before
the arbitrary inward choice. Thus every inward unit vector $z$
and the fixed outward unit vector $w$ satisfy
$\|z+w\|<2\sin(\tau/2)$. This proves the same diameter and
noncriticality conclusions as LC27, with its original tolerances.
\end{proof}

\begin{theorem}[LCP02: exact annular strainers without sharp localization]
\label{thm:collapse-annular-strainers-written-comparison}
LC65's original conclusion holds with its ORIGINAL
$\delta_\sigma,\Lambda_\sigma$, shell $[a,b]=[1/10,10]$
and curvature ball $B(p,400)$. Its inward and outward anchors
are still on the same chosen segments, at the exact prescribed
scale. No additional constraint on $\lambda$ is imposed.
\end{theorem}
\begin{proof}
Keep LC65's construction of equal arms $D=d(p,q)$ from $q$
to $p,z$ with $E=2D-d(p,z)<15\delta$. FIRST shorten both
arms to $r=\min\{D,1\}$. Their new endpoint distance $c_r$
satisfies $2r-c_r\in[0,E]$ by the triangle inequality; also
$a\leq r\leq1$. With $\kappa=1/60$, the equal-arm cosine
calculation in LC65 now gives
\[
 1+\cos\omega_r\leq2\kappa E\coth(\kappa r)
 \leq4E/r\leq4E/a<60\delta/a<1-\cos\theta.
\]
The degenerate case $E=0$ gives $\omega_r=\pi$ directly.
No comparison theorem was used for this first shortening.

The open $256$-ball at $q$ is contained in $B(p,266)\subset
B(p,400)$. ALG06 at $q$, radius $1$ and curvature $-\kappa^2$
therefore gives independent-arm comparison monotonicity on the
shortened radial segments, using the local Riemannian input and
complete ambient manifold. The FINAL radius
$s=L/\lambda\leq a/2<r$ is allowed by the original
$\Lambda_\sigma$. Both final endpoints are exactly the ones
obtained by shortening the original segments to $s$.
Their comparison angle is greater than $\pi-\theta$.
The remaining cosine-law and scaling calculation in LC65 is
unchanged: their opposite distance in the scaled metric exceeds
$2L\cos(\theta/2)>c_\sigma$, hence their curvature-$-\sigma$
comparison angle exceeds $\pi-\sigma$. The same scale bound
supplies curvature at least $-\sigma$ on the required $L$-ball.
This proves LC65 exactly, not a version with a larger curvature ball.
\end{proof}

\begin{theorem}[LCP03: keep the original radial center and the prefix error]
\label{thm:collapse-radial-prefix-written-comparison}
LC69 holds with its original constants, ORIGINAL center $p$ and
chosen outward segment $\xi:[0,D]\to M$. In particular its
error is still less than $2t(1-\cos\theta)$ for EVERY $0<t\leq D$,
and scaling introduces no $\lambda\delta$ error term.
\end{theorem}
\begin{proof}
Put $r=\min\{D,1\}$ and $z_r=\xi(r)$. Keep the inward arm
all the way to the ORIGINAL $p$, of length $D$; shorten only
the outward arm to $r$. The same discarded-tail inequality gives
\[
 E_r:=D+r-d(p,z_r)\in[0,E],\qquad E<15\delta.
\]
The comparison angle $\omega$ at $q$ of this mixed-arm triangle
satisfies
\[
 \begin{split}
 1+\cos\omega
 &\leq \frac{\kappa E_r\sinh(\kappa(D+r))}
                {\sinh(\kappa D)\sinh(\kappa r)}\\
 &\leq 2E_r(1/D+1/r)\leq4E/a
       <60\delta/a<1-\cos\theta.
 \end{split}
\]
Here $a\leq r\leq D\leq10$ and $\kappa(D+r)<1$; use
$\sinh u\leq u\cosh u$, $\sinh u\geq u$ and $\cosh1<2$.
The zero-excess case has angle $\pi$ and needs no division by an
excess. Thus the mixed-arm angle too is greater than $\pi-\theta$.

For $0<t<r$, let $x=\xi(t)$. Apply ALG05 in the intrinsic
$U=B(p,400)$ with endpoint $p$ and its complete closed
$300$-ball. The two adjacent hinges at $x$, with other endpoints
$q,z_r$, have arm sums at most
\[
 D+2r\leq12<300/20,\qquad D+r\leq11<300/20.
\]
All their distances are ambient: the required joins are inside
$B(p,22)$. Their comparison angles have sum at most their germ
limits, hence at most $\pi$ by ALG01. Its model-gluing calculation
therefore gives point-on-side lower comparison in $(p,q,z_r)$.
It follows that $d(p,\xi(t))$ is at least the model side with
arms $D,t$ and angle $\pi-\theta$. LC69's displayed cosine-law
calculation now gives
\[
 0\leq D+t-d(p,\xi(t))<2t(1-\cos\theta).
\]
The endpoint $t=r$ satisfies the same model comparison directly.
For $r<t\leq D$, no comparison theorem is needed: the remaining
tail to $z=\xi(D)$ gives
\[
 0\leq D+t-d(p,\xi(t))\leq E<15\delta
 <\tfrac a4(1-\cos\theta)<2t(1-\cos\theta).
\]
The last bound uses the ORIGINAL $\delta_\sigma$ and $t>r\geq a$.
Multiplication by $\lambda$ gives the original normalized error
$2T(1-\cos\theta)$, with $T=\lambda t$. Inward prefixes still
calibrate the original distance exactly. Nothing is relabelled as $p$.
\end{proof}

\begin{proposition}[LCP04: original radial coordinates, line exclusion and zero consumers]
\label{prop:collapse-zero-comparison-consumer-binding}
The comparison qualifications in LC66, LC70, LC73 and LC76--LC80
are supplied by the written arguments, with their original quantified
conclusions and other hypotheses retained. In particular the selected
zero-model packet keeps its original radial function, scale, model
and cone; no selected data are replaced.
\end{proposition}
\begin{proof}
LC66 uses LCP02, together with AC55. ALG07 supplies the compactness
comparison in AC55's route, ALG08 supplies long-strainer orthogonality,
and ALS01--ALS05 supply its global line splitting. The remaining
coordinate and product arguments are those already written in
AC52--AC61. LC66's shell and covering conclusions therefore retain
their exact original source hypotheses.

For LC70 retain its counterexample sequence, actual segments and
one subsequence. LCP03 supplies LC69's calibration at every fixed
normalized $T$, without requiring $\lambda_i\delta_i\to0$.
The original two-sided product squeeze in LC70 then gives its
EXACT original-distance coordinate; no shortened anchor is substituted.
LC73 uses that same coordinate and the original selected smoothed
radial function. LCP01 now supplies the direction comparison for
LC67's buffered smoothing; LFR02 supplies LC28. LC79 and ALG08
supply the adapted reference on the same splitting. The explicit
LC72 gradient transfer and LC73 threshold choices are unchanged.

For LC76, its thresholds are existential. After choosing $\sigma$
and $\theta$ as in its proof, require additionally
\[
 \delta_\ell\leq\min\{1/100,(1-\cos\theta)/44\}.
\]
This is an allowed choice independent of $q$ and every
$\lambda\geq\Lambda_\ell$. Its first equalization gives arms
$D\in(3,5)$ with excess $E<11\delta$. Shorten BOTH to length $1$
using only the triangle inequality. The new excess is at most $E$,
so the same cosine calculation gives $1+\cos\omega<44\delta
<1-\cos\theta$. ALG06 at $q$, radius $1$, applies inside
$B(q,256)\subset B(p,257)\subset B(p,400)$.
Shorten finally to $L/\lambda\leq1$ and finish LC76's original
model-curvature conversion and AC55 application. Thus its uniform
line-splitting exclusion holds at EVERY point of the unit ball.

LC77 consequently applies to the SAME cone and model at each
selected index. LC74's line from two ends and LC75's bounded-factor
cone identification use the written global splitting and LFR59's
actual cone maps. They introduce no sharp local comparison demand.
LC80 composes these results with the retained joint-zero data and
finite selection. No new restriction on a selected radius or
post-selection replacement is made. This closes their comparison
premises, while existence of the joint data and its parameter order
are still separate assertions to be assembled from LFR49/LFR54/LFR59.
\end{proof}

\begin{theorem}[LCP05: every edge nearest direction with the original buffer]
\label{thm:collapse-edge-directions-written-comparison}
All conclusions of LFR25 hold with its ORIGINAL coarse-border
hypotheses, $0<\tau<10^{-4}$ and curvature ball $B(p,1000\Delta)$,
without the sharp comparison input in AC02. The same applies
simultaneously to all nearest points and minimizing directions.
\end{theorem}
\begin{proof}
Retain LFR25's point $x$, the constructed outward point $y$, and
ANY nearest point $z\in A$. Its notation gives
\[
 \begin{gathered}
 d(p,z)<42\Delta,\qquad a=d(x,z)\in[\Delta/2,12\Delta],\\
 |\ell-a|\leq4\tau\Delta,\qquad \ell=d(x,y),\qquad
 a+\ell<(24+4\tau)\Delta<25\Delta.
 \end{gathered}
\]
All endpoints and chosen sides lie in $B(p,100\Delta)$, so their
intrinsic distances in $U=B(p,1000\Delta)$ equal their ambient
distances. The intrinsic closed $600\Delta$-ball at $z$ is
complete: an intrinsic Cauchy sequence has an ambient limit of
radius less than $(42+600)\Delta<1000\Delta$, and local metric
equality yields intrinsic convergence. ALG05 now applies with
endpoint $z$, since $a+\ell<25\Delta<600\Delta/20$.
Thus the angle between EVERY selected inward and outward arm is
at least its comparison angle.

The already written model estimate in LFR25 is uniform: writing
$e=a+\ell-d(z,y)\leq11\tau\Delta$, its bound is
\[
 1+\cos\omega\leq2e(1/a+1/\ell)
 \leq22\tau\left(2+(1/2-4\tau)^{-1}\right)<100\tau.
\]
The same Euclidean bound applies when $\kappa=0$.
Consequently $\|v+w\|^2\leq200\tau$ for every inward $v$
and outward $w$, exactly as in LFR25. All its preceding metric
estimates and subsequent diameter conclusion are unchanged.
There is no extra endpoint tolerance, scale-Lipschitz constraint,
noncollapse assumption or lost uniformity here.
\end{proof}

\begin{proposition}[LCP06: comparison route for the interior local packets]
\label{prop:collapse-local-comparison-route}
The local comparison uses in the written LC81 consumer constructions
for the zero, circle, edge and slim packets are supplied without
importing the sharp-buffer versions of AC02/AC36/AC64. This removes
that particular source-qualified mathematical input from these routes;
it does not assert the simultaneous existence of their input data.
\end{proposition}
\begin{proof}
There are four kinds of comparison use. First, growing-region
limits and actual approximate product maps use ALG07; their
long-strainer passages use ALG08, and their global splitting and
one-dimensional factors use ALS01--ALS05 and ALR01--ALR05.
This supplies the comparison/recognition premises in the collapsed
model and tested residual-factor arguments, including LC12's
explicit later tail in ALG08.

Second, radial zero-model geometry uses LCP01--LCP04. They bind
the annular direction, exact-scale strainer, original-distance
calibration, adapted radial coordinate and same-model end-count
steps, with the original buffers where just proved. Third, product
anchors in LC78, LFR06 and LFR19 use ALG08's choice of the
splitting threshold AFTER the long anchor radius. LFR02 supplies
their controlled smoothing. Thus circle, slim and tangential edge
coordinates use a written buffered comparison theorem for all
minimizing choices. Fourth, nearest-edge-set directions use LCP05;
LFR26--LFR28 and LFR34--LFR38 then retain their finite-model,
interpolation and full adapted-collar arguments. LFR39--LFR44 use
the same model-map and one-dimensional recognition route.

Comparison on the COMPLETE nonnegative finite models, including
the soul, normal-flow, surface and cone constructions, is LFR55
in the available differentiability range $m\geq4$. The metrics
have order at least $K-1\geq9$, or $K-5\geq5$ on the zero
consumer carrier, so that range is respected. ALG04 also supplies
global comparison whenever its local metric hypothesis is supplied.
None of these finite uses needs a smooth replacement metric.

The only retained geometric foundations at this comparison boundary
are the stated smooth/finite Riemannian variation and local comparison
interfaces, metric compactness/curve facts, and their eventual accepted
PC bindings. The stronger original sharp-margin AC contracts remain
separate and unproved here; no remaining interior-packet comparison
use is discharged just by declaring those parents complete. A single
simultaneous parameter assignment and packet export still require
their own composition proof. Boundary-collar extension and global
map adjustment remain distinct producers.
\end{proof}


\subsection{One parameter assignment produces the interior local families}
\label{sec:collapse-simultaneous-local-production}

We now bind the normalized finite-model constructions to the modified
volume scale. The setting is the CLOSED part of the collapse argument:
$M^\alpha$ are closed connected oriented smooth three-manifolds,
$K\geq10$ is fixed, and the first two premises of KL Standing
Assumption5.2 hold with ONE function $A'(C,v)>0$. Explicitly, for
$1/\alpha\leq v<c_3$ and $0<C<\alpha$,
\[
 \begin{gathered}
 R_p\geq\alpha r_p(1/\alpha),\\
 |\nabla^k\Rm|\leq A'(C,v)r_p(v)^{-(k+2)}
 \quad\hbox{on }B(p,Cr_p(v)),\quad0\leq k\leq K.
 \end{gathered}\tag{LPA.1}
\]
The nongraph premise is NOT used. Retain the inherited smooth/finite
Riemannian, analytic and topological foundations already identified in
the preceding constructions. General metric comparison, splitting and
one-dimensional recognition use ALG, ALS, ALR and LCP, as specified
above. These are written mathematical conclusions under those retained
foundations, not accepted-release or Lean bindings.

The source bodies are KL5.2/5.11 and6.10/6.16, printed39--44/PDF34--39;
7.1--8.5, printed46--48/PDF41--43;9.23--10.4, printed56--58/PDF51--53;
11.1--11.5, printed59--62/PDF54--57; and the ordering(15.4),
printed87/PDF82. The source's uniform-in-center tail convention is
made explicit below. Its finite-$C^K$ claims are implemented by the
consumer-sufficient $C^{K-1}$ metric route, with no smooth-metric upgrade.

\begin{lemma}[LPA01: common normalized analytic data at every center]
\label{lem:collapse-modified-scale-common-analytic-data}
Fix $\Lambda>0$ and $0<w<c_3$, and put
$w'=w/[2(1+2\Lambda^{-1})^3]$. On the LC02 tail choose any smooth
positive scale $\rho_\alpha$ satisfying its two volume-scale bounds.
Define
\[
 v_*={w'\over24I(1)},\qquad
 \mathcal A(R)=2^{K+2}A'(2R+2,w'),\qquad H_\alpha=\alpha/4.
\]
At EVERY $p\in M^\alpha$, in $h_{\alpha,p}=\rho_\alpha(p)^{-2}g^\alpha$,
once $1/\alpha\leq w'$ and $\alpha>4$,
\[
 \Vol_{h_{\alpha,p}}B(p,1)\geq v_*,\qquad
 \sec_{h_{\alpha,p}}\geq-H_\alpha^{-2}
                  \quad\hbox{on }B(p,H_\alpha).
 \tag{LPA01.1}
\]
For every $R>0$ with $2R+2<\alpha$, all orders $0\leq k\leq K$
satisfy $|\nabla^k\Rm|_{h_{\alpha,p}}\leq\mathcal A(R)$ on
$B_{h_{\alpha,p}}(p,R)$. These statements are uniform in $p$ and in
all such choices of the modified scale. In particular, for any fixed
$B>0$, all derivative tests with $0<R<B$ hold on ONE tail
$\alpha>2B+2$; no monotonicity of $A'$ is assumed.
\end{lemma}
\begin{proof}
Write $u=r_p(w')$ and $r=\rho_\alpha(p)$. LC01 and(LPA.1) give
$R_p\geq\alpha u$, while LC02 gives $r\leq2u$.
The curvature-scale convention supplies $\sec\geq-u^{-2}$ on
$B(p,u)$, so LC04 and volume scaling give the first assertion.
Also $H_\alpha r\leq\alpha u/2\leq R_p/2$.
On this ball the defining curvature bound is at least
$-R_p^{-2}\geq-(H_\alpha r)^{-2}$, giving the second assertion
by tensor rescaling. If $R_p=\infty$, use nonnegative curvature
directly; infinity is never a real radius in a volume formula.

For a tested $R$, take $C=2R+2$ in(LPA.1). Since $Rr\leq2Ru<Cu$,
its entire controlled ball contains the requested source ball.
The constant rescaling identity for the tensor norm gives
\[
 |\nabla^k\Rm|_{h_{\alpha,p}}
 =r^{k+2}|\nabla^k\Rm|_{g^\alpha}
 \leq(r/u)^{k+2}A'(2R+2,w')\leq\mathcal A(R).
\]
Every order is bounded by the same function because $k\leq K$.
The standing hypothesis already quantifies over ALL $C<\alpha$;
thus $\alpha>2B+2$ handles every real $R<B$ simultaneously.
No supremum of an arbitrary, possibly unbounded-on-compacts function
$A'$ has been taken. Positivity and finiteness hold at each argument.
\end{proof}

\begin{theorem}[LPA02: uniform joint zero witnesses from the finite category]
\label{thm:collapse-finite-uniform-joint-witnesses}
In LPA01, EVERY sequence $\alpha_i\to\infty$, $p_i\in M^{\alpha_i}$
has a subsequence with the finite package of LFR14, LFR49 and LFR59.
Consequently the joint conclusion of LC58 and LC61 holds at the
modified scale with the finite consumer category replacing LC56:
for each fixed $\delta,\varepsilon,e,T$ in LC57's ranges there are
$V\geq T$ and $\alpha_0$ such that EVERY $\alpha>\alpha_0$ and
EVERY center has ONE scale $s\in[T,V]$, model, its cone, selected
radial function and matching core. All indicated radial sublevels
and open-ball smooth types belong to those SAME witnesses.
The radial function may be required, at its original construction,
to have LC67's enlarged smooth buffer. No uniform cone-convergence
rate over all models is required.
\end{theorem}
\begin{proof}
LPA01 gives noncollapse at radius one by $v_*>0$, eventual bounds
by the FIXED $\mathcal A(R)$, and expanding lower-curvature balls
$H_{\alpha_i}\to\infty$. Use the eventual-tail argument in the
first paragraph of LFR16's proof for LFR14 itself: every chart,
packing and patching stage uses finitely many fixed radii, so start
it after their finite maximum of thresholds. The countable diagonal
still has increasing thresholds. This does not require a splitting
or bounded factor; those are used only in LFR16's later conclusions.
Thus one complete $C^{K-1}$ nonnegative model and actual $C^K$
pointed embeddings are obtained for this ARBITRARY sequence.

LFR59 constructs the cone of this same finite model. LFR49 applies
with that cone and the separate expanding sectional bounds just
proved. In its noncompact branch the consumer carrier has metric
order $K-5\geq5$, actual smooth normal-flow and core maps, and
smooth comparison embeddings with $C^1$ metric convergence.
In the compact branch it gives the actual global comparison
maps. In either case its full joint conclusion holds at one
fixed enlarged scale and on one common subsequence tail.

If the stronger LC67 buffer is requested, first fix its compact
shell and open smoothing neighborhood and decrease the internal
cone tolerance below its threshold. Use its LFR02 smoothing at
the ORIGINAL construction of each source radial function. It
has all LC30 bounds used by LFR49 and greater smoothness domain;
that proof uses only those bounds on its fixed core collar.
No function is changed after the core is attached. This extra
fixed demand is included before choosing the enlarged model scale.

To prove the uniform conclusion, negate the WHOLE conjunction.
There would be $\alpha_i>i$ and $p_i$ with no joint witness at
any $s\in[T,iT]$. The arbitrary-sequence extraction above gives
one limit; LFR49/59 give one fixed $R\geq T$ and a common tail
with all desired clauses. Eventually $R\leq iT$, a contradiction.
LC59--LC61 retain the original smooth source metric and flows and
therefore give every indicated open-ball type without another
scale or extraction. LFR49's finite normal-bundle map supplies
exactly their noncompact identification. This is the LC58 proof
with its previously conditional sequential premise now supplied.
The model, $R$ and cone threshold may depend on that extracted
limit; none is required to be uniform before the contradiction.
\end{proof}

\begin{lemma}[LPA03: the circle tolerance precedes noncollapse]
\label{lem:collapse-uniform-tested-circle-factor}
There is $a_2>0$ such that a pointed metric space with BOTH an
actual KL $\sigma$-map to a complete proper nonnegative model of
dimension at most two and an actual $(2,\beta)$-splitting, where
$0<\sigma,\beta<a_2$, has the following tested residual bound:
for every source point $x\in B(p,201)$ its actual residual
coordinate has distance less than one from the marked factor point.
For complete smooth three-manifolds with the sufficiently late local
curvature bounds, choosing $\beta$ smaller for a fixed requested
coordinate quality supplies LFR07 and LC83. Neither choice depends
on $w'$, $v_*$ or $\mathcal A$.
\end{lemma}
\begin{proof}
If no $a_2$ works, take both errors tending to zero and violating
points $x_i$ in the fixed201-ball. LC17 supplies one complete
proper nonnegative limit $X$ of dimension at most two. AC50--AC51
extract the ACTUAL two-product maps on this same subsequence,
with convergence on bounded balls, to an onto isometry
$X\simeq\mathbb R^2\times Y$. By AC48 its factor $Y$ is a point.
The sequence $x_i$ has a convergent subsequence under the actual
bounded-ball approximations: properness supplies compactness in
a slightly larger fixed ball. The residual images therefore
converge to that single factor point, contradicting their distances
at least one. This uses only the tested ball, and makes no claim
about unused remote points of an arbitrary target factor.

On the smooth source, LFR06 after the fixed rescaling by201 supplies
the rank-two coordinate and value error $<1/10$ by a tolerance
chosen from the requested quality and this numerical radius.
The curvature domain needed by that construction is a fixed finite
ball once the tolerance is fixed; LPA01 supplies it on a uniform
later tail. LFR07's two enclosures now give the whole circle bundle,
including connected fibers and its supported cutoff. The proof has
used metric compactness and smooth-source comparison/smoothing,
not finite-model noncollapse. It is therefore compatible with
choosing $\beta_2$ before the later parameter $w'$.
\end{proof}

\begin{theorem}[LPA04: an acyclic simultaneous choice for all four local types]
\label{thm:collapse-simultaneous-interior-parameters}
Fix the retained standing data $K,A'$ and positive local coordinate
and comparison qualities for the circle, edge collar and slim packets.
There is ONE parameter assignment, in the local part of FC44/(15.4),
and one uniform-in-center late tail, with the following properties:
LC83/LFR07 hold at every two-stratum point; LC85/LFR20 hold at every
slim one-stratum point; every strong edge point has LFR28's disk
packet and LFR38's full collar; and LPA02 holds at every center
with the buffered radial choices needed by LC80. LFR44's strong-edge
density applies to the SAME strata and edge predicates. All finite
metric models, actual embeddings and source functions are retained.
\end{theorem}
\begin{proof}
Here is an explicit sequence of choices. Each bound below is a
positive finite threshold supplied by the cited written result;
only bounds with already fixed arguments are combined.
\begin{enumerate}[leftmargin=*]
\item Choose $\beta_3$ below LC18's universal obstruction. Fix a
circle coordinate quality, also small enough for the requested
edge-collar quality $\gamma<1/100$. Choose
$\beta_2<\min\{\beta_3,\gamma/1000,a_2\}$ small enough for
LPA03's circle coordinate at its fixed scale. These choices use no
noncollapse data. Choose $\Delta>\max\{1,100/\beta_2\}$ large
enough for LFR35--LFR38. Fix the numerical domain/fiber margins
of LFR20, LFR28 and LFR38 at this $\Delta$.
\item Choose the slim, tangential edge, distance-smoothing and
value tolerances required in LFR20/LFR27/LFR28/LFR34--LFR38.
Decrease the coarse-border error $\tau$ to meet their all-direction
conditions and $4\tau\Delta<\beta_2$. Fix the zero radial
Lipschitz tolerance $\varepsilon<1/64$ as follows: choose
$\gamma_0>0$ and $\varepsilon$ with
$\gamma_0+A(\gamma_0,\varepsilon)<\zeta_{\rm slim}$, as in
LC73. This depends on the EARLIER slim quality, not $\beta_1$.
Fix $e<\min\{\varepsilon,1/40\}$ and the LC67 smoothing buffers.
\item Choose the weak qualities $b',s'$ and then the strong
ENDPOINT quality $s$ satisfying LFR29.1 and LFR38/LFR44's
additional inequalities. In particular the endpoint and coarse-border
errors are independent of later noncollapse. LFR44 now supplies
its model tolerance $a_*>0$, depending only on
$\Delta,\beta_2,s$. Choose the collapsed-model error
$\sigma_{\rm col}<\min\{a_*/4,a_2/4,\eta_{18}/4\}$, where
$\eta_{18}$ is LC18's error, with any fixed coverage padding
needed to use the same maps at these weaker accuracies.
Choose $\Lambda$ small enough for LFR29.1, LFR38, LFR44 and
all the fixed support radii, in particular
$100(2\cdot10^6)\Delta\Lambda<1$.
\item Apply LC09 with $\sigma_{\rm col},\Lambda$ and then choose
ONE $0<w<w_*$; define $w'$ as in LPA01. These choices fix
$v_*$ and $\mathcal A$ once and for all. There is no later
reduction of $w$ or $w'$ when a splitting error is made smaller.
\item Now choose the strong SPLITTING quality $b=\beta_E$ small
enough for LFR29.1, LFR19, LFR28 and LFR38, using the just fixed
analytic data. Then choose $\beta_1$ smaller than $\beta_2$ and $\zeta_{\rm slim}$,
LFR44's $b_{1,*}(b)$ and LFR20's slim threshold for these same
data. Any requested finite-buffer comparison tolerances are fixed
before this step and are included in its thresholds. An actual
slim factor has diameter $<1000\Delta$, so LFR20 applies with
its fixed weak inequality. The separate curvature and derivative
hypotheses for BOTH $b$ and $\beta_1$ will be supplied on a late
tail by LPA01; a small splitting error is not used to infer them.
\item Choose LC73's reference quality $\nu$ after $\beta_1$,
keeping the already fixed $\gamma_0,\varepsilon$. Choose the
cone error $\delta_0$ below its threshold, LC66's threshold and
LC77's line-exclusion threshold, and below the smoothing thresholds.
Choose $T>1$ so that $T/20$ exceeds all their lower scale thresholds,
including80. LCP01--LCP04 supply their comparison inputs.
Finally LPA02 supplies the upper scale bound $V$ for these fixed
$\delta_0,\varepsilon,e,T$. Thus $V$ is chosen LAST.
\end{enumerate}

We verify the simultaneous tail. LC02 supplies one smooth scale
at every late index; LC09 gives the collapsed model at EVERY center.
LC18/LC20 exclude three-strata there. For all fixed radii required
by the circle, edge and slim constructions, and for the expanding
ranges $R<b^{-1}$ and $R<\beta_1^{-1}$ used in LFR28/LFR20,
LPA01 provides the exact common analytic data and sectional
bounds after one further index. Its all-$R<B$ clause is why this
step does not require a maximum over infinitely many pointwise tails.
Take the maximum with LPA02's tail, and increase it so that
$H_\alpha>400V$ and all fixed support comparison buffers are
controlled. This makes every permitted zero scale in $[T,V]$
have its required radius400 lower curvature bound as well.

The circle conclusion is LPA03. The slim conclusion is LFR20 with
$r=1$, $v=v_*$ and $A=\mathcal A$. LFR32 supplies the strong
edge's coarse border, and LFR28/LFR38 supply the disk and collar
with those same analytic data. LFR44 supplies the density statements
because its model threshold was chosen BEFORE $w'$, and its
rank-one threshold AFTER $b$. For any later finite strong-edge
family, LFR33/LFR38 use ONE smoothing of $d_{\overline{E'}}$;
its physical value error is chosen using the positive minimum of
that finite family's scales. This changes no quality threshold,
edge predicate or model. The same smoothing gives all its disk
bundles, cutoffs and full lower collars. The constant-core function
used in disk interpolation is not substituted on that collar.
Finally LPA02 gives the joint zero data with the original functions.
This proves joint existence and the stated choice order, rather
than assuming all the local inputs simultaneously satisfiable.
\end{proof}

\begin{theorem}[LPA05: the selected zero packets keep all their witnesses]
\label{thm:collapse-simultaneous-selected-zero-packets}
Use LPA04's single assignment and tail. At every center fix ONCE
one LPA02 witness with $r(p)=s(p)\rho(p)$, $s(p)\in[T,V]$.
LC80 then produces its finite disjoint zero-ball family, tenth-radius
cover, prescribed annular adapted coordinates and same-model
at-most-one-end conclusions. Its actual smooth radial sublevels
have exactly the types in LFR54, including their boundaries, and
its open distance balls retain LC61's smooth identifications.
No continuity of the chosen pointwise witnesses is required.
\end{theorem}
\begin{proof}
The scale is positive and smooth on each compact manifold, so
$T\min\rho>0$ and $V\max\rho<\infty$. Fix the arbitrary
pointwise choices before applying LC62--LC64. The selected centers
lie in the set of balls MEETING the zero stratum; they need not
themselves be zero-stratum points. LC62 gives
$r(i)/\rho(q)\geq T/20$ on the entire closed ten-radius ball.
By the choices in LPA04, LC66 excludes the zero stratum on the
closed shell, so LC64 improves its cover to open tenth-radius balls.
LC73 applies to the originally attached LC67 radial function and
uses the SAME cone and scale, giving its adapted coordinate at
every point of that closed shell. LCP supplies these comparison
steps without the old sharp localization premise.

LC77 uses the same model cone: if that model had two ends, its
cone would be a line, and LC76 would produce a one-splitting at
EVERY point of the selected ball. This contradicts the ball's
meeting the zero stratum, including when its center has nonzero
rank. Thus the model has at most one end. LFR54 now classifies
the actual cores from LFR49, keeping all their metric comparison
maps and smooth embeddings. Composing fixed type identifications
changes neither the radial sublevel nor its boundary. LC61 uses
the retained finite normal-flow identification on the consumer
carrier. All assertions hold on the original common tail; the
continuum of sublevel radii needs no new tail intersection.
\end{proof}

\begin{theorem}[LPA06: simultaneous finite local cover with an early multiplicity bound]
\label{thm:collapse-simultaneous-finite-local-cover}
The SAME late manifolds and assignment of LPA04 admit finite
circle, strong-edge and slim families, together with LPA05's zero
family, whose actual domains cover the entire manifold. The circle,
edge and slim cutoffs equal one on their respective covering balls;
all cutoffs have closed supports strictly within their smooth domains.
The edge family uses ONE physical distance smoothing for both disks
and full adapted collars. The nonzero-family support multiplicity
is bounded by one numerical constant chosen BEFORE $\Delta,w'$.
The zero cutoffs have disjoint supports. This supplies the local
existence, topology, enclosure, cutoff and finite-cover fields in
items(1)--(3) of LC87 for the closed carrier.
\end{theorem}
\begin{proof}
Choose maximal-cardinality disjoint families with radii
$\rho/3$ on the two-stratum and $\Delta\rho/3$ on the slim
stratum. Compactness and $\min\rho>0$ prove finiteness.
For either family write $s=1$ or $s=\Delta$ and $r_p=s\rho(p)$.
If an omitted small ball meets a selected one, the Lipschitz bound
$\operatorname{Lip}r=s\Lambda$ gives comparable radii, with ratio
less than1.01, and center distance less than $.7r_i<r_i$.
Thus the $r_i$-balls cover the respective stratum. Construct their
packets using LPA04. On the circle covering ball $|\eta|<2<8$;
on a slim covering ball LFR20's value estimate gives
$|\eta|<2\Delta<8\cdot10^5\Delta$. Hence the cutoffs equal one
there, and both balls lie inside the respective proper bundle domain.
Use LFR44 for the finite strong-edge family. It covers every nonslim
one-stratum point by a $3\Delta\rho_i$-ball inside the ACTUAL
closed-disk domain and cutoff plateau. LFR33/LFR38 supply its one
shared smoothing after that finite family is fixed, as in LPA04.
The zero tenth-radius balls lie inside the actual smooth cores, since
$\eta<1/10+e<1/5$ there. Together with that cover and the absence of three-strata,
these actual local domains exhaust $M$.

Here is the uniform multiplicity argument, including its dependence.
Set the NUMERICAL constant $C=2\cdot10^6$. Each nonzero-family
cutoff is supported in $B(p_i,Cr_i)$: for circles use the compact
value enclosure, for edges the radius-$100\Delta$ map domain,
and for slim packets the radius-$10^6\Delta$ map domain.
For a point in several such support balls, within one family,
\[
 |r_i-r_j|\leq s\Lambda d(p_i,p_j)
 <Cs\Lambda(r_i+r_j).
\]
Since $Cs\Lambda<1/100$, all radii have ratios between $1/2$ and2.
Choose among the finitely many disjoint small balls one of minimum
VOLUME, say $B(p_0,r_0/3)$. Every other small ball lies in
$B(p_0,(3C+1)r_0)\subset B(p_0,4Cr_0)$.
Increase the common tail, if necessary, so that $H_\alpha>4C\Delta$.
LPA01 then supplies sectional curvature at least $-(Cr_0)^{-2}$
on that entire latter ball: its physical radius is at most
$4C\Delta\rho(p_0)$, and its requested lower bound is weaker
than the supplied one. Local Bishop--Gromov at $p_0$ gives
\[
 N\,\Vol B(p_0,r_0/3)
 \leq\Vol B(p_0,4Cr_0),\qquad
 N\leq{V_1(4)\over V_1(1/(3C))}
       \leq {27C^3V_1(4)\over c_3}.
\]
Minimum volume, not minimum radius, justifies the first inequality.
Choose the integer bound strictly larger than the last number;
it depends on the fixed numerical $C$ only, NOT on $\Delta,w'$
or the late manifold. Sum three such bounds for the three families.
It could therefore have been selected as the early multiplicity
parameter in FC44. The additional late curvature demand causes
no dependency cycle. This proof supplies all whole-ball hypotheses
that were conditional in LC86; it does not assert a uniform bound
on the total number of centers.

Zero cutoffs from LC31 are supported strictly within their selected
radius balls: their outer profile level is $9/10$ and their value
error is $e<1/40$. Those balls are pairwise disjoint. Their core
cover is supplied by the actual zero balls and smooth core packets;
the annular cutoff is not asserted to equal one at the center.
All compact supports, source-fiber enclosures, interpolation domains
and model embeddings are those already constructed. No map adjustment
has been performed to obtain this simultaneous local cover.
\end{proof}

\paragraph{Exact scope of the composition.}
LPA01--LPA06 remove the unbound source-scale and simultaneous-existence
premises for these four INTERIOR local families in the closed standing
sequence. In particular LC81 now has an actual consumer-sufficient
finite model and the three named smooth-fiber/core consumers, using
one admissible parameter order. The general sharp-margin Alexandrov
contracts and stronger smooth-metric LC56 remain separate.
Quantitative comparison and adjustment of DIFFERENT charts on overlaps,
LC87 item(4), and the boundary-collar extension in item(5) remain
separate Chapter4/14 and LC88 producers; this theorem does not assert
that the local maps already agree. Full local-goal review must check
the original requested consumers and all retained foundations against
these written constructions. This is not that final audit, a complete
static graph-manifold theorem, a Ricci-flow application or Lean completion.


\subsection{Final consumer audit of the finite-regularity development}
\label{sec:collapse-finite-final-consumer-audit}

\paragraph{Question and verdict.}
The task begun at revision125 was to supply LC81's finite-category
bridge for its actual local-model consumers, then construct the
rank-two, edge and slim packets. It also identified AC43, LC21,
LC28 and nonnegative-model classification as substantial dependencies.
The audit below follows those ORIGINAL requirements, rather than
using the number of new statements or passing document checks as a
completion criterion. The requested WRITTEN mathematical development
is complete under the explicitly retained foundations below. The
initial finite model has metric order $K-1$, not the stronger order
$K$ in the historical source contract; every named consumer has been
checked at that order. This is not a claim of Lean implementation or
acceptance of a completed PC release.

The original LC81 offered two alternatives: construct its finite
core/normal-flow/isotopy/bundle interfaces, or independently produce
the stronger smooth package LC56. The present proof implements the
FIRST alternative. The historical wording and all earlier proofs are
preserved. The following disposition supersedes their chronological
``still open'' statements only for the consumers listed here.

\setlength{\tablebodywidth}{\dimexpr\textwidth-4\tabcolsep\relax}
\begin{longtable}{@{}P{0.23\tablebodywidth}P{0.35\tablebodywidth}P{0.42\tablebodywidth}@{}}
\textbf{Original requirement} & \textbf{Actual producer} & \textbf{Audit of the interface}\\\hline
\endhead
Finite model and curvature & LFR08--LFR15; LPA01--LPA02 &
Finite normal-coordinate estimates, exact transition identities and
compact-buffer patching construct a $C^K$ atlas and $C^{K-1}$ tensor.
$C^2$ convergence passes curvature; no all-order bound is assumed.\\
Whole buffered embeddings & LFR10, LFR14; LFR48 for zero cores &
Actual pointed embeddings cover every strictly smaller source ball.
Distance distortion is proved separately from differential convergence.
Finite lists of buffers share a tail.\\
Same approximate splitting & LFR11, LFR15--LFR16 &
Quasi-inverses of the actual embeddings feed AC51. The original real
coordinate converges to the exact product coordinate. The factor's
finite manifold structure, completeness and orientation are produced.\\
Gradients and normals & LFR18--LFR20, LFR25--LFR28; LC50/LFR49 &
Every minimizing direction is controlled on compact buffers. Fixed
adapted quality gives a quantitative derivative error, not convergence
to zero by value estimates alone.\\
Finite isotopy with boundary & LFR03--LFR05; LFR20 and LFR28 &
Both target strata remain transverse and ALL interpolated inverse
images lie in one compact interior enclosure. Entire source fibers,
not only chart pieces, are identified before recovering smooth type.\\
LC28 distance smoothing & LFR02, LFR27, LFR33--LFR34 &
Coordinate covectors give arbitrary value error, a small GLOBAL
Lipschitz difference and estimates against ALL minimizing directions.
The finite edge family shares one physical smoothing.\\
Rank-two coordinates and LC83 & LFR06--LFR07; LPA03--LPA04 &
The Gram estimate gives rank two and the required Lipschitz bound.
Only the tested residual factor is bounded. Actual enclosures give
properness, connected circle fibers and the supported cutoff.\\
LC84 model and disk & LFR21--LFR24, LFR28--LFR33 &
Finite nonnegative surface topology produces a model disk. Actual
strong/weak predicates give the CLOSED coarse border, including scale
recentering. The source disk has its original boundary and rank there.\\
LC84 full collar and coverage & LFR34--LFR44; LPA04, LPA06 &
The ORIGINAL pair $(f,F/\rho)$ is adapted throughout the lower collar
in each point's scale. Density and finite selection cover the nonslim
stratum by the actual disk domains, with the same $F$.\\
LC85 slim packet & LFR16--LFR20; LPA04, LPA06 &
The bounded actual factor yields a compact complete oriented surface.
Its smooth type is $S^2$ or $T^2$; every whole source fiber and the
proper interval bundle are constructed with compactly supported cutoffs.\\
Finite soul and zero cores & LFR45--LFR49 &
One soul, calibrated normal tube and common outward field give an
ACTUAL finite normal-flow diffeomorphism. Its smooth bundle carrier
retains the metric, with the four-derivative cost recorded below.\\
Nonnegative model classification & LFR17, LFR21--LFR24, LFR50--LFR54 &
The same noncompact soul bundle gives the core types. Compact smooth
type is determined using an auxiliary metric, discarded before all
packet comparisons. Orientation and the two-ended alternatives are checked.\\
LC21 cone existence & LFR55--LFR59 &
Compact ray parameters give a proper quotient with the actual radial
cone identity and pointed maps at EVERY sufficiently large real scale.
The cone is fixed before the error and survives the carrier change.\\
AC43 and one-dimensional factors & ALS01--ALS05; ALR01--ALR05 &
Global four-point comparison gives an aligned ONTO splitting along
the specified line. Paired distance charts give exact interval and
half-interval recognition under the original Hausdorff dimension bound.\\
Alexandrov comparison at its uses & ALG01--ALG08; LCP01--LCP06 &
Written coarse localization suffices for growing-region limits and
long anchors. Shortened-arm arguments retain the ORIGINAL annular,
radial and edge buffers. Sharp AC02/AC36/AC64 are not imported.\\
One assignment and local families & LPA01--LPA06 &
All-center analytic data supply every required extraction; a joint
contradiction supplies one zero-scale window. The acyclic parameter
choice gives all four finite families on one tail, with actual-domain
coverage and an early support-multiplicity bound.\\
\end{longtable}

\paragraph{Derivative audit.}
Starting with bounds through curvature-derivative order $K\geq10$,
LFR12 gives uniform $C^K$ normal-coordinate coefficients; compactness
loses one derivative. LFR13 gains the extra transition derivative from
an EXACT isometry identity, so LFR14 gives metric $C^{K-1}$ and maps
$C^K$. No harmonic-coordinate endpoint regularity is asserted.
The product adapter LFR11, applied at metric order $K-1$, gives the
same order on the factor and a $C^K$ product map. Edge and slim
interpolations are jointly $C^K$; their isotopies are $C^{K-1}$,
more than the $C^2$ needed by LFR04 for compact smooth fiber type.

For the noncompact zero model, write $m=K-1\geq9$. LFR45 gives
soul order $m-1$ and normal-bundle order $m-2$; LFR46 gives field,
flow and global map order $m-3$. LFR47 changes the carrier through
that actual $C^{m-3}$ diffeomorphism, so the SAME tensor has order
$m-4=K-5\geq5$. The chosen soul-flow map and normed bundle are
smooth in this carrier. LFR48 supplies smooth comparison maps with
$C^1$ tensor convergence; separate expanding source curvature bounds
are retained. LFR55 needs only $m\geq4$ for its variation argument,
and LFR02 only a $C^2$ metric. These thresholds therefore hold on
both carriers. The cone is unchanged by the identity isometry.

\paragraph{Uniformity and witness audit.}
The parameters $K,A'$ are fixed first. LPA01 bounds ALL real radii
$R<B$ on one tail, without taking a supremum of an arbitrary $A'$.
LPA02 applies to every sequence of centers and indices, and negates
the entire joint zero packet when proving its uniform interval
$[T,V]$. Model, cone, scale, radial function and core are never chosen
independently. The cone convergence rate may depend on the extracted
model; the fixed enlarged scale precedes the late index.

LPA04 fixes endpoint errors before $w'$, then chooses analytic-dependent
strong splitting and the smaller one-stratum threshold. The radial
Lipschitz budget depends on earlier slim quality, not on that later
threshold. LPA05 keeps the pointwise witnesses chosen before selection,
including balls whose centers merely have balls MEETING the zero
stratum. LPA06 uses minimum VOLUME in its packing argument and a
whole controlled ball. Its multiplicity bound is independent of
$\Delta$ and $w'$. Acyclicity alone would not prove these inequalities;
the cited proofs provide them.

\paragraph{Retained foundation boundary.}
This is a relative mathematical development over the inherited
geometric, analytic and topological foundation required by the project.
Its leaves include finite-dimensional tensor/ODE and inverse-function
calculus, partitions and approximation, metric curve compactness,
Riemannian local variation and volume comparison, smooth CGT injectivity,
compatible structures, surface and bundle topology, and smooth Ricci-flow
existence, extension, derivative estimates, maximum principles and
positive-Ricci classification. The compact classification also names
de Rham splitting, smooth sphere mapping classes and Bieberbach's
lattice theorem. They are explicit mathematical imports; calling them
retained does NOT certify that their exact Lean declarations have been
found or validated in the accepted PC release.

The frozen CGT, soul-flow and smooth-compactness discoveries keep their
original evidence scope. In particular the frozen all-order compactness
producer was NOT used with finite-order hypotheses. No missing-name
search is evidence that PC lacks a theorem. Exact accepted-release
imports, category adapters, axiom checks and Lean implementations remain
the established later foundation/formalization task; this audit does
not bypass that gate or start a second foundation.

\paragraph{Scope at the handoff to global assembly.}
The named LC81/LC83--LC85 goal concerns existence and finite category of
the local models and their packets. LPA supplies these simultaneously
for the CLOSED standing sequence, including the local existence fields
LC87(1)--(3). The original chapter opening separates this work from
making different charts compatible in Chapter14. LC87(4)'s quantitative
comparisons between different charts and their subsequent adjustments
are therefore NOT claimed by the present completion. Neither is the
boundary augmentation LC88/LC87(5), the full static graph theorem, the
flow application or reconstruction. KL16.3--16.6 still requires its
boundary-specific scale and collar argument; a closed local theorem
alone does not prove the chosen boundary collapse theorem.

The sharper historical metric-order claim in LC81, smooth-metric LC56,
general positive-curvature globalization and an intrinsic identification
of the Tits boundary are stronger statements than these consumers need.
Their omission is justified by the actual producers and derivative
checks above, not by deleting a requested packet field. This audit
completes the requested local-model/finite-category blueprint development;
it does not declare every contract in Chapters3,4,13 or14 complete.

The requirement-by-requirement evidence, source-version and retained-
foundation records are in \code{LOCAL\_MODEL\_AUDIT\_REVISION143.md}
and \code{checks/evidence/revision143\_consumer\_audit.json}.
They distinguish written mathematical review from hash consistency,
PDF verification and Lean evidence. The static checker verifies the
records and preservation; it is not a proof checker. All earlier
statement/proof environments and historical revisions remain intact.


\subsection{Edge and slim geometric packet obligations}
\label{sec:collapse-edge-slim-obligations}


\begin{definition}[LC84: the edge disk-fiber packet]
\label{def:collapse-edge-packet}
An edge packet at $p$, at scale $\rho(p)$, records:
\begin{enumerate}[leftmargin=*]
\item the ACTUAL rank-one map to $\mathbb R\times Y$, and a
  second pointed approximation $(Y,y_0)\to([0,C],0)$ with
  $C>200\Delta$; the marked point is near an ENDPOINT;
\item strong and weak edge sets $E\subset E'$, their quantitative
  coarse-border comparison with $\mathbb R\times\{0\}$, and
  the distance function $d_{\overline {E'}}=d_{E'}$;
\item a tangential adapted function $\eta_p$ on $B(p,100\Delta)$,
  a smoothed distance $\rho_{E'}$, and
  $\eta_{E'}=\rho_{E'}/\rho$, including their value, derivative,
  smooth-domain and compact-support margins;
\item a buffered comparison with $\mathbb R\times Z$ from LC81,
  where $Z$ is a complete orientable nonnegative surface;
\item a model disk bounded by a regular smoothed distance level,
  and LC82's transversality and common compact enclosure for the
  interpolation to $(\eta_p,\eta_{E'})$.
\end{enumerate}
The local output is the proper disk bundle
\[
 \eta_p:\{ |\eta_p|<4\Delta,\ \eta_{E'}\leq4\Delta\}
       \longrightarrow(-4\Delta,4\Delta),
 \qquad\text{fiber }D^2.
\]
The boundary is the level $\eta_{E'}=4\Delta$; there the two
coordinate differentials must have rank two. This boundary condition
is necessary for a bundle of manifolds with boundary.
\end{definition}

\paragraph{Edge proof and its precise implementation obligations.}
KL 9.1--9.21, printed 49--56/PDF 44--51, give the source route.
The one-dimensional Alexandrov classification from Chapter 4 gives
point, line, ray, circle or interval factors. Interior line-like
points give rank two after quantitative rescaling; a nonslim
rank-one point is therefore near an endpoint model. KL 9.10 also
makes strong edge points coarsely dense in the relevant weak-edge
region. The proof uses recentering and the Lipschitz bound on $\rho$;
it must retain both map domains and the changing scale.

Two details require explicit implementation. First, the inverse
image of a factor ball under an arbitrary GH approximation is not
known to be open. For the empty-annulus step in KL 9.11, use this
coarse-chain argument instead. If the radial target distance changes by at most
$d(x,y)+\delta$ between source points (as for a factor projection
of a GH product map), subdivide a path from its basepoint
into steps of length $<h$. The radial target distance changes by
less than $h+\delta$ at each step. It cannot pass from below $a$
to above $b$ when the target annulus $[a,b]$ is empty and
$h+\delta<b-a$. Thus its tested image stays in the inner factor
ball. Restrict the factor and keep the original tested map and
coverage estimates. The comparison map from the restricted factor
to the small model bounds its diameter by that model diameter plus
the distortion; containment in a radius-$600\Delta$ ball alone
would give only $1200\Delta$, exceeding the slim threshold.
No continuity assumption is inserted.

Second, use the closed set $\overline {E'}$ for nearest-point and
smoothing arguments. Coarse comparisons stated with positive slack
pass to its closure by continuity of distance; outward-point and
angle estimates must be uniform on the larger buffered region.
This does not assert that a limit point satisfies a strict edge
predicate. If $E$ is empty no edge packet is requested. Otherwise
$E'$ is nonempty and closure does not change its distance function.
On a normalized band $\rho_{E'}/\rho=O(\Delta)$ the derivative is
\[
 d(\rho_{E'}/\rho)=\rho^{-1}d\rho_{E'}
             -\rho_{E'}\rho^{-2}d\rho.
\]
The second term contributes $O(\Delta\Lambda)$ at the local scale
when $|d\rho|\leq\Lambda$; it cannot be discarded as if $\rho$
were constant. Choose its tolerance before fixing $\Lambda$.

For the surface model, apply KL 3.12 to
$(\Delta^{-1}Z,z_0)$: its radius-10 ball is close to $[0,10]$.
This is the rescaling needed in the sentence preceding KL 9.21;
the unrescaled radius-$10\Delta$ ball is not compared directly to
that unit interval. Smooth a distance level in the noncritical band,
compare its differential with the source functions, and apply LC82
with $X=\{0\}\times(-\infty,4\Delta]$.
Properness and boundary rank then give the displayed bundle.
The endpoint-model, closure-uniform angle, model-embedding and
finite-regularity estimates are source-qualified producer targets,
not consequences of the definition of an edge packet.
KL 9.23 requires the endpoint error to be chosen independently of
$w'$; the sharper splitting tolerance may depend on $w'$.

\begin{definition}[LC85: the slim sphere-or-torus packet]
\label{def:collapse-slim-packet}
A slim packet records an ACTUAL $(1,\beta_1)$-splitting with a
chosen factor of diameter less than $10^3\Delta$, an adapted
coordinate $\eta$ on $B(p,10^6\Delta)$, and a comparison from
LC81 with $\mathbb R\times Z$, where $Z$ is a compact connected
orientable nonnegative surface. It includes LC82's compact enclosure
and transverse interpolation identifying $\eta^{-1}(0)$ with $Z$.
The fiber is $S^2$ or $T^2$, by inherited surface classification
and Gauss--Bonnet. A buffered proper restriction is a bundle over
an interval. The cutoff is supported in
$|\eta|<9\cdot10^5\Delta$, with a plateau on
$|\eta|\leq8\cdot10^5\Delta$.
\end{definition}
KL 10.1--10.4, printed 58/PDF 53, are the source.
Compactness of the model factor must be proved from the bounded
factor approximation and completeness; it is not inferred from a
single small metric ball. Orientation of $Z$ comes from the oriented
three-manifold and the oriented real factor. A torus is not discarded
because a sphere also occurs, and no global choice of fiber orientation
on future overlaps is asserted. The compact-enclosure clause replaces
any unsafe invocation of a perturbation theorem on only a small domain.

\begin{lemma}[LC86: finite local selections and overlap bounds]
\label{lem:collapse-local-selection-packet}
Let $M$ be compact with positive Lipschitz scale $\rho$, and fix
constants $a,b,C>0$. Suppose $\operatorname{Lip}\rho$ is small
enough that intersecting $a\rho$-balls have comparable radii and
centers in each other's $b\rho$-balls. A maximal pairwise disjoint
family of $a\rho$-balls with centers in a prescribed stratum is
finite, and its $b\rho$-balls cover that stratum. If, in addition, $2C\operatorname{Lip}\rho<1/2$ and
all enlarged comparison balls satisfy a common normalized lower
Ricci bound, the multiplicity of the $C\rho$-balls is bounded by
a constant depending only on $a,C$, dimension and that bound.
\end{lemma}
\begin{proof}
The positive minimum of $\rho$ and compactness bound the cardinality
of any disjoint family, so a family of maximal cardinality exists.
Every omitted center has its $a\rho$-ball meeting a chosen one,
which gives the asserted cover. At a point in several $C\rho$-balls,
small Lipschitz constant makes their radii uniformly comparable.
Their disjoint $a\rho$-balls lie in a fixed multiple of any one
of the radii. Choose among them a ball of smallest volume.
Bishop--Gromov bounds the volume of the containing larger ball
by a fixed multiple of that smallest ball. Disjointness bounds the
number of balls by the same multiple. Apply the argument first to
any finite subfamily, which suffices also for a multiplicity bound.
All curvature comparison balls, not just the small disjoint balls,
are premises of this estimate.
\end{proof}
KL 8.5, 9.24 and 10.4 use this with $a=1/3$ or
$a=\Delta/3$ and fixed larger support radii. The constants may
depend on the already fixed $\Delta$ and cutoff radii. Finiteness
is for each manifold; there is no uniform bound on the total number
of centers in a sequence. The additional exhaustion assertion
``rank one is slim or lies within $3\Delta\rho$ of a selected
edge center'' uses KL 9.10 and this selection, not compactness alone.
KL 9.26's overlap alternative and 9.28--9.29's auxiliary cutoffs
are exported with their margins to Chapter 14, where compatibility
of the coordinates is used; the selection lemma does not prove them.
\clearpage
\section{Local outputs, boundary collars and the flow application}
\label{sec:collapse-chapter-exports}

\paragraph{The model topology passed to Chapter 14.}
KL 3.11, printed 24--25/PDF 19--20, classifies complete connected
orientable nonnegative three-models, including its stated $C^K$
category. This is an inherited/source-qualified classification input,
not a new proof of Hamilton's compact classification or the soul theorem.
After LC77 excludes multiple ends for the selected zero models, the
remaining possibilities and their compact representatives are:
\begin{center}
\setlength{\tablebodywidth}{\dimexpr\textwidth-2\tabcolsep\relax}
\begin{tabular}{@{}P{.48\tablebodywidth}P{.52\tablebodywidth}@{}}
\textbf{Model} & \textbf{Core or compact alternative}\\\hline
$\mathbb R^3$ & $D^3$\\
$S^1\times\mathbb R^2$ & $S^1\times D^2$\\
Twisted real line bundle over $\mathbb {RP}^2$
  & $\mathbb {RP}^3\setminus\operatorname{int}D^3$\\
Orientable twisted real line bundle over the Klein bottle
  & Orientable twisted interval bundle\\
Compact model & $S^1\times S^2$, $\mathbb {RP}^3\#\mathbb {RP}^3$,\\
 & spherical space form, or compact flat quotient
\end{tabular}
\end{center}
This table supplies types only after the same-model core identification
has been provided; an arbitrary topologically equivalent core need not
have the required radial margins. Compact models are not forced to
have torus boundary. The author's correction adding $D^3$ to KL
14.1(2) must be retained in the Chapter 14 consumer.

\begin{definition}[LC87: the local-collapse export certificate]
\label{def:collapse-local-export-certificate}
A local-collapse export consists of a fixed manifold and metric,
positive modified scale, a single simultaneous parameter assignment,
and the following finite families with their actual maps:
\begin{enumerate}[leftmargin=*]
\item LC80 zero-model balls and cores, with original radial functions;
\item LC83 circle charts, LC84 edge disk packets and LC85 slim packets,
  with their buffered domains, proper restrictions and fiber types;
\item the LC86 covers and overlap bound, cutoffs with supports strictly
  inside their smooth domains, and the stratum exhaustion statement;
\item quantitative chart-overlap comparisons from Chapter 4, still
  distinguished from any subsequent adjustment of the charts; and
\item for a boundary carrier, LC88's boundary collar packets and the
  finite list of its connected components.
\end{enumerate}
Every item includes the witness metric and scale, map domain and image
buffer, regularity, error tolerance, core/fiber identification, and the
consumer that must preserve it. A certificate is DATA with proofs of
these fields. Its definition asserts neither existence from collapse
hypotheses nor compatibility after changing any map.
\end{definition}

\paragraph{Parameter order.}
All local quality requests from Chapter 14 must be fixed before
producing LC87. The finite list of resulting inequalities is solved
in the dependency order of KL 15.2/(15.4), printed 87/PDF 82;
Chapter 14 must expand that global order for its adjustment parameters.
Within the present chapter, choose fixed domain and fiber margins
first; then coordinate and smoothing tolerances; then the cone and
splitting errors and the lower scale thresholds; then the joint
witnesses and a common sequence tail; finally select the finite covers.
This sentence is a LOCAL scheduling rule, not a replacement for
(15.4). In particular $\Delta$ precedes the edge endpoint error,
that error precedes the volume choice $w'$, and the stronger edge
splitting error may be chosen after $w'$ (KL 9.23).
The modified-scale Lipschitz constant must meet all previously chosen
$\Delta$- and support-size constraints. The cusp parameter $r_\partial$
is chosen after the upper scale bound $V$ (KL's $\Gamma_0' $)
as required in Section 16. Taking minima or
maxima combines finitely many thresholds only after their dependency
order has been checked; it cannot resolve a circular inequality.

\begin{definition}[LC88: the boundary-collapse packet]
\label{def:collapse-boundary-packet}
For a compact connected orientable three-manifold with nonempty
boundary, use its INTRINSIC length metric. Record the four premises
of KL 16.1 in Section~\ref{sec:independent-collapse-contracts}, with
$K,A$ fixed BEFORE $w_0$. An embedding from the depth-100 cusp
collar is an embedding of PAIRS, with its boundary mapped onto the
labelled boundary component. Its reference curvature is $-1/4$ and
its metric is $dz^2+e^{-z}g_{T^2}$. The packet also records the
interior cover, the smoothed depth coordinate on the buffered collar,
and the boundary bundle restriction with fiber $T^2$.

The boundary variant of LC87 is to be produced by KL 16.4--16.6:
centers of the interior families have distance greater than ten from
the boundary; the collar cutoff is supported between depths 20 and
90 and equals one between 30 and 80; its retained piece is
$I\times T^2$. Either the whole connected carrier is already such
a product, or the retained boundary and zero-model regions are
pairwise disjoint. The comparison buffers and the endpoint levels
must be part of the packet, not inferred from the names of the sets.
\end{definition}

\paragraph{Boundary implementation and selected-center audit.}
KL 16.4 gives $1<R_p<3$, small modified scale and a slim adapted
coordinate on the depth band $(5,95)$. Its proof reduces to the
explicit cusp metric. Interior closed-manifold lemmas may be used
only on balls contained in the controlled interior or after a
specified extension whose metric agrees there. No complete ambient
manifold is silently substituted for the intrinsic carrier.
In applying KL 16.5, our selected zero-ball centers need not themselves
be zero-stratum points; their balls MEET that stratum. Thus use the
uniform line-splitting exclusion of LC76 throughout such a ball,
not an unsupported assertion about its center.
Here is a noncircular sufficient localization argument. On the
buffered retained cusp collar, take $w_0$ small enough that a sectional
curvature is at most $-1/8$. At an interior intersection point of
$B(i,r_i)$ with that collar, the zero packet's normalized lower bound
$-1/60^2$ implies
\[
 -\frac1{3600r_i^2}\leq-\frac18,
 \qquad r_i\leq\frac{\sqrt8}{60}<\frac1{20}.
\]
The pair embedding must include the buffer assertion that the
intrinsic one-neighborhood of the retained depth-$90$ collar,
away from boundary distance nine, lies in the smoothed depth band
$(5,95)$. It follows from the cusp length comparison and a first-exit
argument on its depth-$99$ domain, and is an explicit field of the
boundary packet. Since $d(i,\partial M)>10$, the entire selected
ball is away from boundary distance nine and lies within this
one-neighborhood. KL 16.4 therefore excludes zero rank at EVERY
point of the ball, contradicting that it meets the zero stratum.
There is no requirement to choose $\Lambda$ after the model-dependent
upper scale bound $V$.
This supplies the corrected logic once the collar map and buffer
estimates are bound. It does not prove disjointness of two different
boundary collars; that remains the separate product-or-disjoint
argument of KL 16.5 in the boundary producer.

Empty boundary uses the CLOSED theorem, with no invented distance to
an empty set. Disconnected carriers use finitely many connected
components with COMMON $K,A,w_0$; thresholds are combined only for
that finite list. A component with curvature scale $+\infty$ needs
its explicit nonnegative-curvature/classification branch. Never write
$\infty^3$ in a real-valued volume inequality and treat it as an
ordinary admissible radius. The static theorem supplies a raw graph
presentation, not ambient incompressibility or the final geometric
endpoint. The boundary producer and Chapter 14 assembly remain
source-qualified implementation targets.

\begin{lemma}[LC89: absorbing pointwise eventual estimates into one function]
\label{lem:collapse-common-derivative-function}
Fix $K$ and a sequence of smooth compact carriers $M_j$.
Suppose each carrier has a finite upper bound $R_j$ for every
finite radius to be tested, and finite bounds for
$|\nabla^k\Rm|$, $0\leq k\leq K$, on the entire carrier.
For each $w\in(0,c_3)$ suppose there are $N(w)$ and $B(w)<\infty$
such that for every $j\geq N(w)$ ALL the derivative tests with
volume ratio at least $w$ satisfy
\[
 r^{k+2}|\nabla^k\Rm|(q)\leq B(w)
 \quad(0\leq k\leq K,\ q\in B(p,r)).
\]
Then there is one positive function $A:(0,\infty)\to(0,\infty)$
for which these tests hold for EVERY $j$ and EVERY $w\in(0,c_3)$.
It can be fixed before invoking KL 16.1 to choose $w_0$.
\end{lemma}
\begin{proof}
For each $j$ choose a finite
\[
 D_j\geq\max_{0\leq k\leq K}
       R_j^{k+2}\sup_{M_j}|\nabla^k\Rm|.
\]
For $0<w<c_3$, set
\[
 A(w)=1+\max\bigl(\{B(w)\}\cup
                    \{D_j: j<N(w)\}\bigr),
\]
and put $A(w)=1$ elsewhere. The finite initial segment is controlled
by $D_j$, the tail by $B(w)$. Thus $A$ works simultaneously, even
when $N(w)$ is unbounded as $w$ varies. No common tail over all
$w\in[w_0,c_3)$ was assumed or extracted.
\end{proof}
The radius bound is real content. On a fixed compact connected
component containing a negatively curved tangent plane at $x$, a
finite number exceeding both its diameter and the inverse square
root of that negative curvature magnitude bounds every curvature
scale. Finitely many components give $R_j$. Nonnegative components
must first enter their separate branch. This observation does not
identify a carrier's intrinsic balls with the ambient flow balls.

\begin{proposition}[LC90: the static-to-flow interface and chapter exit]
\label{prop:collapse-static-flow-interface}
Supply the following inputs separately:
\begin{enumerate}[leftmargin=*]
\item the static boundary/closed collapse theorem assembled from
  LC87 in Chapter 14, with output in KL Definition 1.2;
\item actual normalized late carriers from Chapter 12, with
  $\widehat g(t)=t^{-1}g(t)$ and the same hyperbolic normalization
  $\sec=-1/4$ on their reference cusps;
\item proofs of the boundary size, collar, interior-volume and
  WHOLE-BALL derivative premises, with common parameters in the
  required order, including LC89 or an equally uniform argument;
\item the separate nonnegative-component and terminal-exceptional
  data, and the graph-output conversion and ambient transport maps.
\end{enumerate}
Then applying the static theorem componentwise gives the graph
part of the late thin output. Chapter 6 refines that raw presentation;
Chapters 15--16 transport and assemble it with the thick and
exceptional pieces. No Ricci-flow construction is an input to the
static local model arguments of this chapter.
\end{proposition}
\begin{proof}
Each connected carrier satisfies the SAME static theorem's hypotheses,
so instantiate it and retain the component labels. The finite union
is precisely the required raw graph output. Apply the separately
supplied conversion and transport maps; do not replace them by the
name ``graph manifold''. If the long-time outcome is terminal, use
its explicit exceptional data instead of requesting a nonexistent
late slice. This is conditional composition, not a proof of the
listed analytic or topological producers.
\end{proof}

\paragraph{Fresh flow-source audit.}
KL14 17.1--17.3, printed 90--91/PDF 85--86, motivate this composition.
The cited KL08 source is the archived February 20, 2013 version.
Its Theorem 92.3(3), printed 2825/PDF 239, and Lemma 92.13,
printed 2827--2828/PDF 241--242, give derivative bounds at the
CENTER, up to the finite order $\lfloor\epsilon^{-1}\rfloor$.
They do not by their literal statements give KL14 16.1's estimates
at every point of the triggered ball. Choose the surgery parameter
so this order is at least $K$, and prove a buffered/recentered
whole-ball adapter before using that premise. A center estimate
cannot merely be relabelled a ball estimate. LC89 addresses the
function/tail quantifiers only AFTER this analytic adapter is supplied.

KL08 88.1, printed 2807--2809/PDF 221--223, controls the normalized
Ricci tensor near the hyperbolic value on large neighborhoods under
its long-time noncollapse and curvature premises. The inherited
Section 87 setting is a normalized compact initial slice, the
nonextinct negative-scalar branch and the surgery-scale separation
$h_{\max}(t)/r(t)\to0$ imposed in Remark 86.9. These are branch
hypotheses, not conclusions of static collapse. The three-dimensional
Ricci-to-sectional translation and the actual neighborhood domains
must be included in the adapter. KL14 17.3 uses it and volume
comparison to rule out triggered radii above the small-radius cutoff.
This does not supply the missing small-radius whole-ball derivative
statement. H1--H4 remain explicit application/compatibility gates;
finite-time extinction is not added to the default route.
% BEGIN REVISION162 AMBIENT WHOLE BALL
\subsection{Actual whole-ball estimates on late ambient slices}
\label{sec:collapse-actual-ambient-whole-ball}

The mismatch identified above is resolved here for the ambient flow
balls. Restriction to the intrinsic complementary carriers is a
separate step. In particular, no equality of intrinsic and ambient
balls is used in this subsection.

\paragraph{Fixed flow and retained analytic estimates.}
Fix the normalized compact surgery flow before choosing any test ball.
Write $r_{\rm can}(t)$ for its nonincreasing canonical-neighborhood
scale, and $h_*(t)$ for the largest surgery radius on $[t/2,t]$.
Use the admissible profile of KL85's last paragraph and KL86.9:
$r_{\rm can}(t)\to0$, $h_*(t)/r_{\rm can}(t)\to0$, and the
KL84.1 estimates for each fixed spatial enlargement factor at all
sufficiently late times. This is one choice of surgery parameters for
one flow, not a change of the flow after a ball or sequence is chosen.
For example, first replace the finitely requested $\delta_A$ by their
running minima for integer $A$, then choose the surgery profile below
half that minimum for $A\leq\lceil2t\rceil$ at time $t$; impose the
scale-separation requirement at the same construction stage. The
inherited surgery-existence interface must permit this smaller profile.

The analytic inputs here are the scalar lower bound, Hamilton--Ivey
pinching, canonical-neighborhood control and its model derivative
estimates, KL70.1--70.2, the noncollapsed parabolic seed and bounded-distance
estimates KL81.3/84.1, and local Shi estimates. These are estimates on
an actual surgery flow, not static metric compactness. The center
estimate KL92.13 is also used, with its literal center conclusion:
for each $w>0$ there are $b(w)>0$, $C_0(w)<\infty$ and $T_0(w)$ such
that, for $t\geq T_0(w)$,
\begin{equation}\label{eq:collapse-center-estimate-input}
 \begin{gathered}
 0<s\leq b(w)\sqrt t,\quad
 \sec_{g(t)}\geq-s^{-2}\text{ on }B_{g(t)}(p,s),\\
 \operatorname{vol}_{g(t)}B_{g(t)}(p,s)\geq ws^3,\\
 \quad\Longrightarrow\quad
 s^2|\Rm|_{g(t)}(p)\leq C_0(w).
 \end{gathered}
\end{equation}
Its proof, not just its statement, was checked: KL83.1 supplies the
almost-Euclidean subball; above the surgery scale KL81.3 and Shi apply;
below it canonical neighborhoods and bounded-distance curvature control
exclude an arbitrarily large center scalar curvature. The following
arguments add the missing spatial uniformity. They do not replace
\eqref{eq:collapse-center-estimate-input} by an assumed whole-ball bound.
Accepted PC bindings for these inherited analytic interfaces remain pending.

\begin{lemma}[Buffered estimate above the surgery scale; WBD01]
\label{lem:collapse-macroscopic-whole-ball}
For each $w>0$ there are $\Lambda(w),C(w)<\infty$ and
$b(w),\tau(w)>0$ such that at sufficiently late times a ball as in
\eqref{eq:collapse-center-estimate-input} with
$\Lambda(w)h_*(t)\leq s\leq b(w)\sqrt t$ has an unscathed
backward neighborhood on the static tracks from $B_{g(t)}(p,2s)$,
of duration $\tau(w)s^2$, with
\[
 |\Rm|\leq C(w)s^{-2}.
\]
Consequently, for each fixed integer $m\geq0$,
\[
 \sup_{q\in B_{g(t)}(p,s)}s^{m+2}|\nabla^m\Rm|_{g(t)}(q)
 \leq C_m(w).
\]
No lower curvature bound on $B_{g(t)}(p,2s)$ is assumed.
\end{lemma}
\begin{proof}
Use KL81.3's proof on printed page~2794 (PDF page~208).
Retain its enlarged region instead of only the final radius-$s/4$ neighborhood.
Here are the buffers and the order of its constants. KL83.1 selects
$B(y,\theta s)\subset B(p,s)$ on which every subball has almost
Euclidean volume, with $\theta>0$ depending only on $w$ and the fixed
volume tolerance. KL84.2 gives a backward unscathed seed at $y$, after
$\Lambda$ is made large enough. Shrink its radius to $a s$ and its
backward time interval so that, at each time $u$ in a terminal interval
of length $c s^2$, the hypotheses of KL84.1 hold on
$P(y,u;as,-a^2s^2)$, including the bound
$|\Rm|\leq(3a^2s^2)^{-1}$ and a lower volume bound
$c_1(as)^3$. This shrink uses curvature, distance and volume distortion
inside the seed; take $2a^2<c$ before shortening the terminal interval.
The constants $a,c,c_1$ depend only on $w$ and the retained analytic
constants.

Choose the KL84.1 parameter larger than both $20/a$ and $1/c_1$.
Choose $b$ small enough for its part(c), for the seed estimates and
for Hamilton--Ivey's conversion of the resulting scalar bound to
$|\Rm|\leq C s^{-2}$. On the terminal interval this gives that bound
on the time-dependent balls $B_{g(u)}(y,20s)$. This conclusion uses
bounded-distance propagation from the seed; it does not apply
Bishop--Gromov outside the original lower-curvature ball. Enlarge
$\Lambda$ once more so that $Cs^{-2}$ is strictly below the
curvature threshold of every surgery with radius at most $h_*(t)$.
Also shorten the interval to lie in $[t/2,t]$. Thus these balls cannot
meet an inserted, excised or discarded surgery region. The usual
first-exit continuation defines their unscathed time-dependent family,
as in KL84's proof; the seed is already unscathed throughout.

Every point of $B_{g(t)}(p,2s)$ is joined to $y$ by a path of length
less than $3s$ at time $t$. Along its backward static transport,
$|\Ric|\leq C' s^{-2}$ bounds its length by
$3s\exp(C'\tau)$. Choose $\tau<c$ with this expression less than
$6s$. A first exit from $B_{g(u)}(y,20s)$ is then impossible.
The same reasoning applies to a fixed additional spatial buffer
around these tracks. It supplies the asserted static-track domain,
not just a collection of unrelated balls at different times.

For any $q\in B_{g(t)}(p,s)$, use a radius proportional to $s$ on
the earlier slice, contained in this buffered domain. Distance
comparison ensures that its inner half still contains the track
ending at $q$. Local Shi estimates at the positive elapsed time
$\tau s^2$ give $C_m(w)s^{-m-2}$. All constants are independent of
$q$, of its distance to the boundary of $B(p,s)$, and of $p,t,s$.
\end{proof}

\begin{lemma}[Uniform scalar bound below the surgery scale; WBD02]
\label{lem:collapse-microscopic-whole-ball}
Fix $w>0$ and $\Lambda<\infty$. For all sufficiently late times,
balls satisfying \eqref{eq:collapse-center-estimate-input} with
$0<s<\Lambda h_*(t)$ obey
\[
 \sup_{q\in B_{g(t)}(p,s)}s^2|\Rm|_{g(t)}(q)\leq C(w,\Lambda).
\]
\end{lemma}
\begin{proof}
Suppose otherwise and take counterexamples with $t_j\to\infty$.
Scale the metric by $s_j^{-2}$; put $R_j=s_j^2R$ on the final slice.
Its original ball now has radius one and sectional curvature at least
$-1$. By \eqref{eq:collapse-center-estimate-input}, $R_j(p_j)$ has a
fixed upper bound. Since $s_j\leq\Lambda h_*(t_j)\to0$ and
$s_j/r_{\rm can}(t_j)\to0$, the rescaled flows satisfy the
canonical-neighborhood assumption with the \emph{constant} threshold
scale $1$ on their entire preceding time intervals: an original
scalar curvature at least $s_j^{-2}$ is above the original threshold
at every earlier time. Their rescaled pinching functions are
\[
 \Phi_j(v)=s_j^2\Phi(v/s_j^2),\qquad
 \frac{\Phi_j(v)}v\longrightarrow0\quad(v>0\text{ fixed}).
\]
Here $\Phi$ is the fixed pinching function of the chosen flow; the
Hamilton--Ivey version at late times gives the same conclusion.

In dimension three, the lower sectional bound $-1$ implies that
an unbounded curvature norm on the unit ball forces an unbounded
scalar curvature there. Choose $q_j$ with $R_j(q_j)\to\infty$ and
fix $D>1$ larger than the center scalar upper bound. A minimizing
geodesic from $p_j$ to $q_j$ has length less than one. It has a point
$z_j$ with $R_j(z_j)=D$, and $d_j(z_j,q_j)<1$.
Apply KL70.2 to the rescaled flow with
\[
 Q=D,\qquad A=2\sqrt D,\qquad \widehat r=1.
\]
For large $j$ its pinching test holds because
$\Phi_j(D)/D\to0$. Its conclusion gives
$R_j(q_j)\leq K(2\sqrt D,1)D$, a contradiction.
The scale $\widehat r=1$ is essential: using the unscaled
$K(A,r_{\rm can}(t_j))$ would not give a uniform constant.
KL70.2 assumes the global a priori flow conditions, not a lower
sectional bound on a larger ball. Thus no boundary buffer has been
silently added. The contradiction proves a uniform eventual scalar
bound, and the original sectional lower bound converts it to the
stated curvature-norm bound.
\end{proof}

\begin{proposition}[Finite derivatives throughout the actual ball; WBD03]
\label{prop:collapse-ambient-whole-ball-derivatives}
Fix a finite integer $K\geq0$. Choose the canonical-neighborhood
accuracy before constructing the flow so that it is admissible for
KL70.2 and provides at least $K+2$ metric derivatives on its neck
and cap charts; the stronger request
$\lfloor\epsilon^{-1}\rfloor\geq K+4$ suffices with the stated
conventions. For each $w>0$ there are $b(w)>0$, $T(w)$ and finite
$C_m(w)$, $0\leq m\leq K$, such that for $t\geq T(w)$,
\begin{equation}\label{eq:collapse-actual-whole-ball-bound}
 \begin{gathered}
 0<s\leq b(w)\sqrt t,\quad
 \sec\geq-s^{-2}\text{ on }B_{g(t)}(p,s),\\
 \operatorname{vol} B_{g(t)}(p,s)\geq ws^3,\\
 \Longrightarrow\quad
 \sup_{q\in B_{g(t)}(p,s)}s^{m+2}|\nabla^m\Rm|_{g(t)}(q)
 \leq C_m(w)\quad(0\leq m\leq K),
 \end{gathered}
\end{equation}
provided the connected ambient slice component containing $p$ has
some negatively curved tangent plane. Nonnegative components retain
their separate classification branch.
\end{proposition}
\begin{proof}
Use WBD01 if $s\geq\Lambda(w)h_*(t)$, increasing its constants to
cover the finite list of derivative orders. Otherwise WBD02 bounds
$R(q)\leq Cs^{-2}$ simultaneously at every $q\in B(p,s)$.
For large $t$, $s<r_{\rm can}(t)$ uniformly in this second case.

If $R(q)>r_{\rm can}(t)^{-2}$, the point has a canonical
neighborhood. Types(c) and(d) in KL69.1 are closed positive-curvature
components, so cannot occur in this ambient component. For a neck
or cap the retained model estimates and the finite differentiable
closeness give
\[
 |\nabla^m\Rm|(q)\leq c_m R(q)^{(m+2)/2}
 \leq c_m C^{(m+2)/2}s^{-m-2},\qquad m\leq K.
\]
This also covers a cap point just inserted by surgery: it uses the
spatial model derivatives, not an unavailable backward parabolic
neighborhood through the surgery. The extra two metric derivatives
are retained when estimating curvature derivatives.

If $R(q)\leq r_{\rm can}(t)^{-2}$, the scalar lower bound at large
$t$ gives $|R(q)|\leq r_{\rm can}(t)^{-2}$. KL70.1 therefore bounds
$R$ by a fixed multiple of $r_{\rm can}(t)^{-2}$ on the backward
static tracks from a ball of radius $c r_{\rm can}(t)$ for a time
$c' r_{\rm can}(t)^2$, wherever the tracks are defined. Pinching
converts this to a two-sided curvature bound. Since
$h_*(t)/r_{\rm can}(t)\to0$, that bound excludes any surgery endpoint
of these tracks. The region is consequently unscathed, as in the
low-curvature case of KL92.13. Local Shi estimates, with fixed
smaller buffers, give
\[
 |\nabla^m\Rm|(q)\leq c_m' r_{\rm can}(t)^{-m-2}
 \leq c_m' s^{-m-2}.
\]
These ambient neighborhoods need not lie in $B(p,s)$; their control
comes from the flow and canonical estimates. No lower curvature
bound there is inferred from the trigger ball. All bounds are
uniform in $q$, including points arbitrarily close to its boundary.
Taking the maximum over the two scale cases and the two curvature
cases proves \eqref{eq:collapse-actual-whole-ball-bound}.
\end{proof}

\begin{corollary}[Normalization and exact remaining carrier adapter; WBD04]
\label{cor:collapse-normalized-whole-ball}
For $\widehat g(t)=t^{-1}g(t)$, WBD03 gives for every $w>0$
\[
 \sup_{q\in B_{\widehat g(t)}(p,r)}
 r^{m+2}|\nabla^m\Rm|_{\widehat g(t)}(q)\leq C_m(w),
 \qquad 0<r\leq b(w),\quad m\leq K,
\]
whenever this ambient ball has sectional curvature at least $-r^{-2}$
and volume at least $wr^3$, at a sufficiently late time depending
on $w$. The constants use one fixed flow and one fixed finite $K$.
\end{corollary}
\begin{proof}
Set $s=\sqrt t\,r$. Distances, volumes and sectional curvatures
scale respectively by $\sqrt t$, $t^{3/2}$ and $t^{-1}$ on returning
to the physical metric. Curvature derivative norms satisfy
\[
 |\nabla^m\Rm|_{\widehat g(t)}
 =t^{1+m/2}|\nabla^m\Rm|_{g(t)}.
\]
Hence $r^{m+2}|\nabla^m\Rm|_{\widehat g(t)}
=s^{m+2}|\nabla^m\Rm|_{g(t)}$, proving the claim.
\end{proof}

% BEGIN REVISION166 COMPONENT GUARD NOTICE
\paragraph{Final-audit scope clarification.}
WBD04 retains WBD03's negative-plane ambient-component hypothesis.
LAU02 checks it in every LTF02/TCF05 application. Closed nonnegative
components keep their separate classification branch.
% END REVISION166 COMPONENT GUARD NOTICE

\paragraph{Application boundary.}
WBD03--WBD04 close the CENTER versus WHOLE AMBIENT BALL mismatch,
including radii arbitrarily near an ambient curvature scale, when
they also meet the displayed small-radius cutoff. They do not
identify a complementary carrier's intrinsic curvature scale or
balls with the ambient ones. The remaining application must supply
the actual nearly cuspidal collars, compare a triggered intrinsic
ball with a controlled ambient domain, exclude triggered radii above
$b(w)$ using the thick-region estimates, and then apply LC89 with
one $A$ fixed before $w_0$. Those tasks, persistent coverage and the
common-time choices remain part of the active analytic goal.

% END REVISION162 AMBIENT WHOLE BALL

% BEGIN REVISION165 ACTUAL THIN APPLICATION
\subsection{The actual thin carriers and the static collapse application}
\label{sec:collapse-actual-thin-application}

We now supply the analytic premises of LC90 on the actual late flow.
The static inputs are PBR03 and BBR03, with their retained foundations;
HG08 is supplied by HPI08. Fix once and for all a finite integer
$K\geq20$ and the admissible surgery accuracy used in WBD03, with
$\lfloor\epsilon^{-1}\rfloor\geq K+4$. All metrics in this subsection
are the restrictions of $\widehat g(t)=t^{-1}g(t)$ on the SAME slice
$Z_t=(M_t^+,\widehat g(t))$. No doubling or replacement metric is used.
Times are admissible slice times, with the post-surgery convention.
The scalar-nonnegative and empty alternatives retain LTF01's preceding
branch convention and LP10--LP11's actual history data.

The source route is KL14, Theorems16.1 and17.3, printed87--88 and90--91,
PDF82--83 and85--86. Its brief application is expanded below at the
intrinsic-ball, boundary-near and parameter interfaces. The literal
center statement of KL92.13 remains a center statement; WBD03--WBD04
are the separate whole-ambient-ball adapter. The same-day source and
errata checks are recorded in \code{reference_checks_revision165.md}.

\begin{lemma}[Actual slow truncations and long cusp collars; TCF01]
\label{lem:collapse-actual-slow-cusp-truncations}
The finite persistent family admits smooth slowly escaping cusp
truncations whose complementary carriers $N_t\subset Z_t$ have the
following properties at all sufficiently late times.
They are compact orientable smooth manifolds, with finitely many
components and labeled torus boundary. Every boundary torus has an
actual cusp embedding of pairs of depth100, with a larger depth200
buffer in the same carrier. In the normalization
$h=dz^2+e^{-z}q_t$, their relative pulled-back metric errors through
order $K+4$, and their boundary diameters, tend uniformly to zero.
Every point of $N_t$ with finite ambient curvature scale satisfies
$\Theta_t(p)<\eta(t)$ for a function $\eta(t)\to0$.
The boundary tracks are smooth and their physical velocities are
$o(t^{-1/2})$. Closed persistent models are removed as whole components;
an empty family is allowed.
\end{lemma}
\begin{proof}
Choose disjoint standard cusps in each of the finitely many models,
with coordinates $s\geq0$ and metric $ds^2+e^{-s}q$. A common level
$S$ defines a compact truncation $C_i(S)$; for a closed model put
$C_i(S)=H_i$. Choose increasing levels $S_n\to\infty$ so that all
closed source balls $\overline B(x_i,n)$ lie in the INTERIORS of
$C_i(S_n)$, with a margin, and all level-$S_n$ tori have diameter
less than $1/n$. The finitely many closed models cause no difficulty.

For each $n$, the fixed compact sets through level $S_{n+1}+1000$
lie in the persistent domains at all sufficiently late times. HPI08
and HPI05 give disjointness of their images and relative metric errors
less than $1/n$ through order $K+4$. Require also
$\sqrt t\,|\partial_t f_i|<1/n$ on these sets, and that the entire
$1/n$-thick part lies in the images of $B(x_i,n)$, using HPI07.
These are finitely many eventual conditions. Choose thresholds $T_n$
meeting them, increasing so that
\[
 \log(T_{n+1}/T_n)\geq
 n(1+S_{n+1}-S_n)\|\theta'\|_\infty,
\]
where $\theta:[0,1]\to[0,1]$ is smooth, nondecreasing and constant
near its endpoints. This is possible recursively: the conditions
for the interval starting at $T_n$ involve the already specified
finite compact set through $S_{n+1}+1000$, not the future time $T_{n+1}$.
Interpolate $S_n$ to $S_{n+1}$ by $\theta$ in the variable $\log t$ on
$[T_n,T_{n+1}]$. The resulting $S(t)$ is smooth, tends to infinity,
and satisfies $t|S'(t)|\leq1/n$ on that interval.

Remove the interiors of the disjoint images $f_i(t)(C_i(S(t)))$.
Their closures are embedded compact submanifolds, so the complement
$N_t$ is compact and smooth, with exactly their boundary tori; a
closed model has open-and-closed image and removes a whole component.
A compact manifold has finitely many components, since its components
are open and admit a finite subcover. The actual collar maps are
\[
 e_{i,t}(y,z)=f_i(t)(y,S(t)+z),\qquad 0\leq z\leq200.
\]
They lie outside the removed interiors by injectivity and the
pairwise disjoint larger images. They cover the ENTIRE corresponding
boundary components. Their reference metric is
$dz^2+e^{-z}q_t$ with $q_t=e^{-S(t)}q$; translation in cusp depth
preserves the relative tensor norms. Thus the stated metric errors
and diameters tend to zero, independently of the changing torus size.

On $[T_n,T_{n+1}]$, the removed interiors contain all the
$1/n$-thick points by their source-ball margin. Hence every remaining
finite-scale point has $\Theta_t<1/n$; take $\eta(t)=1/n$ there.
For a moving boundary point, its velocity is the sum of the supplied
$\partial_t f_i$ and $df_i(\partial_s)S'(t)$. The second term has
physical norm at most $2\sqrt t\,|S'(t)|\leq2/(n\sqrt t)$.
This proves the physical speed assertion, including smooth passage
through the flat interpolation endpoints. The maps lie in the
persistent unscathed space-time domains, also across admissible
surgery identifications. In the empty-family case take $N_t=Z_t$;
HPI08 says every fixed thick part is eventually empty and gives the
same diagonal thinness assertion. No cusp or boundary claim is then
needed.
\end{proof}

\begin{lemma}[Interior curvature scales and balls are ambient; TCF02]
\label{lem:collapse-intrinsic-ambient-scale-equality}
For a component of $N_t$ with nonempty boundary, let
$D=d_{N_t}(p,\partial N_t)>10$. At sufficiently late times its
intrinsic curvature scale and the ambient curvature scale satisfy
\[
 R_{N_t}(p)=\rho_{Z_t}(p)\leq D-1<D.
\]
For every $0<r\leq R_{N_t}(p)$ the intrinsic and ambient balls,
their restricted metrics and their volumes are identical.
Consequently the interior-volume premise of BBR03 holds with
$\eta(t)$ in place of $w_0$.
\end{lemma}
\begin{proof}
Distances to the full boundary set agree intrinsically and ambiently.
Indeed, an ambient curve leaving the carrier must first hit that
boundary, and its initial segment lies in the carrier; approximation
by such curves proves both inequalities. In particular, whenever
$r<D$, every ambient curve from $p$ of length less than $r$ stays
inside, so $B_{Z_t}(p,r)=B_{N_t}(p,r)$. This also proves equality of
the restricted metric tensors and volume measures, not just a set
inclusion.

Choose an intrinsic minimizing curve from $p$ to a closest boundary
point, and let $q$ be its point one unit before the endpoint. Then
$d_{N_t}(p,q)=D-1$ and $d_{N_t}(q,\partial N_t)=1$, by minimality
and the distance function's Lipschitz inequality. BSA01 puts $q$
in its actual cusp collar and gives a plane with sectional curvature
at most $-1/8$. If an intrinsic admissible curvature radius exceeded
$D-1>9$, its ball would contain $q$ while requiring curvature greater
than $-1/81$, impossible. Thus $R_{N_t}(p)\leq D-1$.
The same point, at ambient distance at most $D-1$, proves
$\rho_{Z_t}(p)\leq D-1$ as well. The admissible radii form intervals
and their balls coincide for all radii below $D$, so their suprema
are equal. These scales are finite and strictly positive.
The ball at their common radius is still below $D$ and also coincides.
TCF01's ambient thinness now gives
$\Vol B_{N_t}(p,R_{N_t}(p))<\eta(t)R_{N_t}(p)^3$.
This argument is on the original restricted metric throughout.
\end{proof}

\begin{lemma}[All boundary-near derivative tests; TCF03]
\label{lem:collapse-boundary-near-all-radii}
There is a finite constant $E_K$ such that at every sufficiently
late time, for $d_{N_t}(p,\partial N_t)\leq10$, all
$0<r<R_{N_t}(p)$ and all $0\leq m\leq K$ satisfy
\[
 \sup_{q\in B_{N_t}(p,r)}
 r^{m+2}|\nabla^m\Rm|(q)\leq E_K.
\]
There is no volume assumption in this boundary-near assertion.
\end{lemma}
\begin{proof}
BSA01 gives $R_{N_t}(p)\leq D+3\leq13$ and puts $p$ in an
actual collar with height less than11. Every curve from $p$ of
length less than13 remains in this same collar, below height25:
before a putative first exit, $|dz|\leq1.01$ bounds its height
variation, whereas the only internal exit is beyond height98.
Thus the ENTIRE tested ball is in a uniformly controlled collar.

For completeness, these controls give derivative bounds without any
lower bound for the torus injectivity radius. On a local cover use
the reference hyperbolic metric $h$ and write $g=h+u$. The inverse
metric is uniformly bounded relative to $h$; the connection difference
is $\tfrac12g^{-1}*(\nabla^h u)$ with the usual three symmetrized
terms. Curvature is $\Rm(h)$ plus its first covariant derivative
and quadratic connection-difference terms, with contractions by
$g,h$. Repeated differentiation expresses $\nabla_g^m\Rm(g)$ in
finite tensor products of $g^{-1}$ and derivatives of $u$ through
order $m+2$, together with the parallel constant-curvature tensor of
$h$. TCF01 bounds all these factors uniformly. The estimates descend
from local covers and are independent of the lattice and its aspect
ratio. Hence $|\nabla^m\Rm|\leq C_K$ on the whole height25 region.
Take $E_K=C_K13^{K+2}$. This includes boundary centers and points
arbitrarily close to the boundary of a tested ball; it invokes no
backward parabolic neighborhood through an external boundary.
\end{proof}

\begin{lemma}[Large-radius derivative triggers disappear; TCF04]
\label{lem:collapse-actual-large-radius-exclusion}
Fix $w\in(0,\omega_3)$ and let $b=b(w)>0$ be WBD04's cutoff,
shrunk to at most one. At all sufficiently late times, an interior
point with $d_{N_t}(p,\partial N_t)>10$ cannot satisfy
\[
 b\leq r<R_{N_t}(p),\qquad
 \Vol B_{N_t}(p,r)\geq wr^3.
\]
The same assertion holds on a closed carrier component containing
a negatively curved tangent plane.
\end{lemma}
\begin{proof}
For the boundary case TCF02 makes all tested balls ambient, with
$r<\rho_{Z_t}(p)$; for the closed case the carrier component is an
entire ambient component, so the same statement is immediate.
Such a trigger has sectional curvature at least $-r^{-2}$ on its
whole ball. Bishop--Gromov there gives
\[
 \Vol B_{Z_t}(p,b)\geq c\,w b^3,\qquad
 \sec|_{B_{Z_t}(p,b)}\geq-b^{-2},
\]
with a numerical $c>0$: compare the model volume at radius $r$
in curvature $-r^{-2}$, which is a fixed constant times $r^3$,
with its radius-$b$ volume, which is at least $\omega_3b^3$.
Use endpoint volume continuity if necessary. Thus $b,cw$ are fixed
macroscopic seed constants in normalized units. LTF03 implies that
the sectional curvature at $p$ tends to $-1/4$, uniformly over any
sequence of such triggers. In particular $\rho_{Z_t}(p)<5$ at late
times. Positivity of that finite scale and monotonicity of ball
volume give
\[
 \Theta_t(p)
 =\frac{\Vol B_{Z_t}(p,\rho_{Z_t}(p))}{\rho_{Z_t}(p)^3}
 \geq w\left(\frac b5\right)^3>0.
\]
This contradicts TCF01's uniform ambient thinness. A purported
failure at arbitrarily late times would give exactly such a
sequence, proving the asserted eventual exclusion. No estimate at
an intrinsic ball touching a removed core has been used.
\end{proof}

\begin{proposition}[One fixed-thickness bound for every tested ball; TCF05]
\label{prop:collapse-actual-carrier-whole-ball-estimates}
After separating closed globally nonnegative-curvature components,
for every fixed $w\in(0,\omega_3)$ there are finite $B_K(w)$ and
$T_K(w)$ such that for $t\geq T_K(w)$, EVERY remaining carrier point,
$0<r<R_{N_t}(p)$ and $0\leq m\leq K$ with
$\Vol B_{N_t}(p,r)\geq wr^3$ satisfy
\[
 \sup_{q\in B_{N_t}(p,r)}
 r^{m+2}|\nabla^m\Rm|(q)\leq B_K(w).
\]
On each individual such compact carrier all curvature scales have
a finite common upper bound, and all indicated derivatives have
finite suprema.
\end{proposition}
\begin{proof}
TCF03 handles every boundary-near center, regardless of radius or
volume. For the other centers TCF02 identifies scales and balls
with the ambient ones; closed components require no adapter.
TCF04 excludes $r\geq b(w)$. WBD04 now controls the entire remaining
ambient ball, with constants $C_m(w)$. Its negative-plane component
hypothesis holds: a boundary component's collar supplies such a
plane in its ambient component, and a retained closed component
has one by the stated separation. Set
$B_K(w)=\max\{E_K,C_0(w),\ldots,C_K(w)\}$ and take the maximum of
the finitely many eventual time thresholds. These constants are
uniform over points, balls and carrier components for the fixed flow.

A closed carrier component is an open-and-closed subset of its
ambient component, since it is a codimension-zero embedded manifold
without boundary; hence it is that entire component. Closed components
with nonnegative sectional curvature enter the retained compact
classification and FC41 recognition branch, with actual graph
presentations and component labels. They are not assigned a finite
volume radius. For every remaining connected component choose one
negative plane. A number exceeding its intrinsic diameter and the
inverse square root of that plane's curvature magnitude bounds every
curvature scale, as in LC89's following observation. Compactness
bounds the curvature derivatives, including up to the smooth external
boundary. There are finitely many components, so maxima give the
individual-carrier bounds needed by LC89. These last bounds need not
be uniform in time.
\end{proof}

\begin{theorem}[Actual late thin graph carriers; TCF06]
\label{thm:collapse-actual-late-thin-graph-carriers}
Under the retained foundations of PBR03/BBR03 and the fixed admissible
flow hypotheses, every sufficiently late carrier $N_t$ from TCF01
has a finite smooth KL graph presentation, retaining its actual
external boundary labels. The closed and empty-boundary cases use
their stated separate conventions. This supplies LC90's analytic
static-application premises on the actual flow, rather than on
hypothetical carriers satisfying them.
\end{theorem}
\begin{proof}
Suppose that nongraph carriers existed at arbitrarily late admissible
times, and select such a sequence $t_j\to\infty$. Separate each one's
closed nonnegative components as in TCF05, retaining their supplied
graph presentations. On the remaining finite unions, TCF05 gives
for each fixed $w$ an eventual uniform $B_K(w)$ and an index $N(w)$.
Its last paragraph gives the finite radius and derivative bounds
on every individual member. LC89 therefore constructs ONE positive
finite function $A(w)$ for ALL members and ALL tests, absorbing the
finite exceptional initial segment separately for each $w$.
Extend $A$ positively outside $(0,\omega_3)$ if using the literal
KL16.1 domain. Fix this function BEFORE invoking the static theorem.

PBR03 and BBR03 give a common positive $w_0=w_0(K,A)$. At sufficiently
large indices TCF01 gives boundary diameters and relative $C^K$
cusp errors below $w_0$, on actual smooth embeddings of pairs of
depth100. TCF02 and $\eta(t_j)\to0$ give the interior-volume premise
at every point whose boundary distance is greater than10. On closed
remaining components it follows directly from ambient thinness and
the equality of ambient and intrinsic geometry. All whole-ball
derivative tests, for every $w\in[w_0,\omega_3)$, already satisfy
the SAME function $A$ by its construction. Thus every component
satisfies the precise static theorem to which it is submitted.
Apply BBR03 to nonempty-boundary components and PBR03 to closed ones;
restore the separately recognized nonnegative components. Their
finite union is a labeled graph presentation for $N_{t_j}$, contrary
to its choice. Empty unions have the empty presentation.

This contradiction proves the all-sufficiently-late assertion.
The order is $K$, fixed flow accuracy, proposed bad sequence, ONE
$A$, static $w_0$, then a sufficiently late index. There is no common
time extracted from infinitely many $T_K(w)$, no dependence of $A$
on an already selected $w_0$, and no claim that an individual
componentwise threshold is uniform over an unspecified family.
\end{proof}

\begin{corollary}[Actual common-time analytic output; TCF07]
\label{cor:collapse-actual-common-time-analytic-output}
On the late-slice branch, the analytic inputs of LP03--LP09 now admit
an actual common time: the finite persistent hyperbolic family, its
controlled disjoint compact cores and cusp collars, and the actual
complementary thin graph carriers in the same normalized slice.
The late/terminal-exceptional alternative is preserved.
\end{corollary}
\begin{proof}
HPI08 and TCF01 give the actual smooth families and eventual finite
lists of compact-domain, metric, velocity, boundary and coverage
conditions. TCF06 gives an eventual graph conclusion on the actual
complements. Take a time beyond their finitely many thresholds,
as in LP04; if another independently supplied finite list is needed
by a later consumer, enlarge that time as LP04 prescribes. Inclusion
maps and all external boundary labels are those of the constructed
carriers. In a terminal alternative retain LP10--LP11's recorded
exceptional history instead of choosing a nonexistent late slice.
\end{proof}

\paragraph{Exact completion boundary before the final audit.}
QRG09, LTF05, HPI08 and TCF06 supply the requested global hyperbolic
comparison, actual thick limits, persistent disjoint embeddings and
static collapse on the actual thin complement. WBD03--WBD04 and
TCF02--TCF05 provide the previously missing whole-ball and intrinsic
adapters. LC90's further graph-definition conversion, ambient
incompressibility (HG10--HG13/Meeks--Yau), essential-torus refinement
and final reconstruction remain distinct consumers and retained
work, not conclusions of this analytic iteration. In particular no
full LP12 geometrization result or accepted-PC/Lean proof is asserted.
All in-scope written producers are now present; the requested ONE
final mathematical audit is the next step, not a completed check here.
% END REVISION165 ACTUAL THIN APPLICATION

% BEGIN REVISION166 FINAL ANALYTIC AUDIT
\subsection{Final mathematical audit of the long-time analytic goal}
\label{sec:collapse-longtime-final-analytic-audit}

This is the ONE final mathematical audit requested for the analytic
goal begun from revision159. It reviews QRG04--QRG09, HPS01--HPS05,
WBD01--WBD04, LTF01--LTF05, HPI01--HPI08 and TCF01--TCF07, together
with the exact retained interfaces they consume. The separate completed
Chapter14 audit159 is not reopened. The review record is
\code{AUDIT_LONGTIME_REVISION166.md}, with its requirement-to-argument
matrix in \code{longtime_goal_audit_revision166.csv}. Static consistency,
source hashes and document builds are evidence about the records and
layout; they are not mathematical or Lean proof certification.

\begin{lemma}[Finite-family threshold correction; LAU01]
\label{lem:longtime-audit-common-threshold-repair}
In HPS04's finite-family diagonal, take the MAXIMUM of the component
starting-time or starting-index thresholds, and the MINIMUM of any
positive permitted error bounds. With this correction, all its
finite-family conclusions and the HPI08/TCF01 applications hold with
one common exhausting domain and accuracy function.
\end{lemma}
\begin{proof}
For a fixed integer stage $n$, there are finitely many components,
buffered domains, derivative orders and error requirements. Write
their eventual thresholds as $J_{n,1},\ldots,J_{n,k_n}$. Put $J_0=0$ and define
\[
 J_n=\max\{n,J_{n-1}+1,J_{n,1},\ldots,J_{n,k_n}\}.
\]
For an empty list use the first two entries. Then every $j\geq J_n$
satisfies ALL stage-$n$ requirements. Choose $m(j)$ as the largest
passed stage, capped by $j$, and impose the finite-order isotopy,
metric and speed requirements at that stage as in HPS04. Every fixed
stage eventually passes, so $m(j)\to\infty$. When a joint positive
smallness bound is needed, the minimum of finitely many positive
bounds is positive. Time thresholds and smallness bounds have opposite
inequality directions and must not be interchanged.

HPS04's sentence asking for the minimum of component thresholds is
therefore corrected to their maximum. Its formulas for the isotopy,
static transport, matching near junctions and physical speed are
unchanged. All adjusted component images still lie in the originally
disjoint buffered images. HPI08 and TCF01 already combine finite
lists by a common sufficiently late time; use the same maximum rule
there. This is a repair of the choice in the proof, not an added
hypothesis on the flow or a change of the geometric output.
\end{proof}

\begin{corollary}[Explicit scope of the normalized whole-ball estimate; LAU02]
\label{cor:longtime-audit-normalized-component-guard}
WBD04 carries the hypotheses of WBD03, including that the tested
ambient component contains a negatively curved tangent plane. On
such components its normalized estimate holds with one fixed flow,
one finite order $K$, and constants and eventual times depending
on $w$. Every invocation of WBD03--WBD04 in LTF02 and TCF05 meets
this component hypothesis. Closed nonnegative components use their
separate classification branch.
\end{corollary}
\begin{proof}
The rescaling identity in WBD04 changes norms and radii only; it
does not remove a hypothesis of WBD03. In LTF02 the ambient scale
is finite, so its component has a negative plane by the definition
and compactness argument there. In TCF05 a component with boundary
has a negative plane in its actual cusp collar, hence in its ambient
component. A closed carrier component is a whole ambient component;
the nonnegative ones were explicitly separated before that invocation.
Thus all actual uses satisfy the restriction. This clarification
prevents the shorter corollary statement from being read as an
unproved extension to arbitrary closed canonical alternatives.
Finite $K$ is fixed before canonical accuracy; all-order bounds on
thick windows arise separately from Shi on unscathed neighborhoods.
\end{proof}

\paragraph{Source hypotheses retained, with their status.}
The analytic inputs are the actual compact surgery flow with normalized
initial data and locally finite event history; its adjustable admissible
profile, canonical neighborhoods and finite model estimates; scalar
evolution and maximum principles; Hamilton--Ivey; local Shi and smooth
compactness with the curvature/volume-to-injectivity bridge; and the
source estimates KL70.1--70.2, KL83.1, KL84.1--84.2 (hence KL81.3),
and the literal center estimate KL92.13. The relevant statements and
proof passages were source-checked, including the time-dependent-ball
correction in KL86.8 and the separation in86.9. The project expansions
prove their actual domain, scale, regularity and parameter transfers;
they do not claim new Lean proofs of those input estimates or that
all of them have been located in the completed PC release. Their
eventual direct-reuse, adapter or GC-new classification remains the
Chapter10--11 release audit's task, with no fresh moving-snapshot probe.

The geometric inputs are the written Margulis and Mostow developments,
ordinary smooth covering and metric foundations, and PBR03/BBR03 with
their retained local-collapse and FC41 topological library. In particular
no persistent-family theorem or preexisting graph presentation of the
actual thin complement is an analytic input. HPI08 and TCF06 produce
those objects. Ambient cusp incompressibility and final reconstruction
are not smuggled into the input by a record field.

\paragraph{Normalization and dependency checks.}
The flow metric is normalized by $t^{-1}$, the hyperbolic sectional
curvature is $-1/4$, the cusp metric is $ds^2+e^{-s}q$, and the
physical velocity is $o(t^{-1/2})$. KL's curvature-operator eigenvalue
is twice sectional curvature; LTF02 uses $-2\rho^{-2}$ when the
sectional minimum is $-\rho^{-2}$. The strict canonical-scale adapter
in revision163 uses $1/2$ and $D>4$ in WBD02, not the excluded endpoint
of KL69.1's interval. The surgery accuracy is chosen once, with enough
finite metric derivatives for $K$ curvature derivatives. It is not
changed after the future collapse tolerance or test ball is selected.

The order of production is acyclic. WBD uses the listed local source
estimates and literal center bound, without hyperbolic persistence.
LTF uses WBD at order zero and produces actual limits/windows. QRG
and HPS are independent geometric constructions. HPI uses their
outputs and LTF to produce a finite disjoint persistent family with
matching-ball coverage. TCF constructs the actual carriers, identifies
interior balls, handles all boundary-near tests, excludes large-radius
triggers and applies LC89. The static tolerance is chosen after ONE
function $A$ has been fixed on a proposed bad sequence. The slowly
truncated family may be prepared in advance because its errors tend
to zero independently of that tolerance; the actual slice and cut
used by LP03 are selected afterward. This is not a reversal of the
consumer's parameter order.

\begin{theorem}[Audited written analytic endpoint; LAU03]
\label{thm:longtime-audited-written-analytic-endpoint}
Relative to the explicit retained foundations just listed, the written
analytic goal is complete: QRG00/HG06 have the global core/end producer;
HG07 has actual smooth hyperbolic limits and unscathed time windows;
HG08 has a finite disjoint controlled persistent family covering the
thick part; and the actual complementary late thin carriers satisfy
the static collapse hypotheses and have labeled finite graph
presentations. The analytic premises needed for a common LP04 slice
are supplied. The late/terminal-exceptional alternative is preserved.
\end{theorem}
\begin{proof}
QRG05 proves both injectivity comparisons; QRG06 proves radial
transversality and excludes nearby axes; QRG07 fixes one sheet and
the thickward side; QRG08 identifies all target boundaries and excludes
hidden tubes. QRG09's cusp count and orientation-cover argument finish
QRG00, and QRG03 gives HG06. HPS01--HPS04, with LAU01, supply the
finite-order transition and explicit time smoothing consumers.

WBD01 retains a genuinely larger controlled region around a source
seed and converts its time-dependent balls into buffered static tracks.
WBD02 uses a fixed rescaled canonical threshold, not an uncontrolled
$K(A,r_{\rm can}(t))$. WBD03 handles high canonical points, including
new caps, spatially and low points by unscathed Shi estimates. LAU02
retains their precise component scope. LTF01--LTF05 then give the
actual normalized hyperbolic limits, compact limit families and smooth
dyadic windows, with finite volume and uniform approximation thresholds.

HPI03 constructs actual minimal-cusp corrections whose jumps tend to
zero. HPI04--HPI05 use thin cusp barriers to keep new thick basepaths
outside the old cores and prove escape from every old basepoint.
Their fixed compact images are therefore disjoint, and HG09 terminates
the selection by volume. HPI06--HPI07 prove the additional matching
$1/w$ coverage before the final accuracy diagonal. HPI08 consequently
has every actual input required for its claimed smooth family.

TCF01 gives actual long collars and uniform ambient thinness. TCF02's
negative collar plane forces the interior curvature scale below the
boundary distance, giving equality of the tested intrinsic and ambient
balls. TCF03 covers all boundary-near radii. TCF04--TCF05 supply all
remaining whole-ball tests. LC89 and TCF06 choose one function before
the static tolerance and contradict any arbitrarily late nongraph
sequence, using PBR03/BBR03 on the same metrics. TCF07 and LP04 retain
one common admissible slice and the actual maps. Empty families, whole
closed hyperbolic components and nonnegative closed components use the
explicit separate cases. A terminal history is handled by LP10--LP11,
without substituting a general extinction theorem for the late route.
This proves the stated endpoint at the level of written mathematics.
\end{proof}

\paragraph{Audit result and completion boundary.}
The review found the finite-family threshold error repaired by LAU01
and the component-scope ambiguity clarified by LAU02. No further
unresolved in-scope implication was found relative to the explicitly
retained inputs. The audit is complete after these repairs. Earlier
paragraphs describing an open frontier record their historical revisions;
this endpoint and the current-view record govern the present status.
Ambient incompressibility, the Meeks--Yau interface, essential graph
refinement, model conversion, surgery reconstruction and accepted-PC/Lean
verification retain their separate obligations. In particular the whole
Geometrization formalization and LP12 are not declared complete.
% END REVISION166 FINAL ANALYTIC AUDIT

\clearpage
\section{Chapter 13 audit and completion boundary}
\label{sec:collapse-completion-audit}

Revision 96 completes the chapter's BLUEPRINT COVERAGE: every planned
topic has definitions, a theorem or producer contract, a source route,
explicit limitations and an identified consumer. This is not mathematical
closure of every source-qualified input and is not Lean completion.
The current disposition below supersedes earlier chronological paragraphs
saying ``next'', ``still needed'' or ``remains open'' when the indicated
written argument has since been supplied. Historical source and PC
records are retained with their original evidence scope.

\setlength{\tablebodywidth}{\dimexpr\textwidth-4\tabcolsep\relax}
\begin{longtable}{@{}P{.22\tablebodywidth}P{.43\tablebodywidth}P{.35\tablebodywidth}@{}}
\textbf{Material} & \textbf{Audit disposition} & \textbf{Unresolved binding or consumer}\\\hline
\endhead
LC01--LC15 & Scale estimates, fixed-parameter noncollapse and local
volume proof checked; LC15 supplies LC06. & Local comparison,
volume identification and smoothing remain inherited bindings.\\
LC16--LC25 & Finite rank, arbitrary-factor exclusion, same-limit cone
transfer and exact outward radius checked. & LC21 cone existence and
LC24 finite-order compactness remain source-qualified.\\
LC26--LC44 & Closed annuli, Lipschitz DIFFERENCE, smooth sublevels,
whole-ball embeddings and outward-sign tests checked. & LC28 smoothing;
actual core data and inherited calculus.\\
LC45--LC61 & Same soul/field, variable-height core, all-direction collar
transfer and joint witnesses checked; open-ball smooth type separated
from boundary topology. & LC56 smooth package is sufficient only;
LC81 finite-category producer/adapter remains open.\\
LC62--LC73 & Original radius assignment and closed shell retained;
LC79 supplies LC71 conditionally. & Buffered approximation weakening
corrected below; no new Lean proofs.\\
LC74--LC80 & Same-cone end count, product-anchor estimates and zero
export have written conditional proofs. & Splitting, first variation,
LC28 and the model package remain explicit inputs.\\
LC81--LC86 & Finite category and local circle/disk/slim interfaces
specified; properness and compact isolation distinguished from rank. &
Full-rank adapted and edge/slim geometric producers are source-qualified
implementation targets.\\
LC87--LC90 & Parameter, boundary and static/flow interfaces specified;
common derivative-function argument written. & Chapter 14 assembly;
whole-ball flow estimates; H1--H4; graph refinement and transport.\\
\end{longtable}

\paragraph{Revision126 update to this historical audit.}
LFR02 now supplies LC28 by a written finite-category smoothing proof;
LFR03--LFR05 supply the finite-isotopy and smooth-fiber consumers.
LFR06--LFR07 supply rank-two coordinates and the conditional circle
packet. Thus the corresponding open-producer wording in the historical
table is superseded. The finite-order model producer and its actual
buffered embeddings, finite soul/core and classification data, and
edge/slim geometric estimates remain open. LFR01's derivative ledger
is required at each use; no smooth metric or accepted Lean theorem
has been manufactured by choosing a compatible atlas. The final audit
of the user's full local-model objective has not yet been reached.

\paragraph{Revision127 update to this historical audit.}
LFR08--LFR10 now supply the basepoint-to-local injectivity adapter
and the injectivity and source-ball coverage of actual comparison
maps. LFR11 supplies a complete finite-$C^K$ factor and an actual
$C^{K+1}$ product diffeomorphism from an exact metric splitting,
without assuming a manifold structure on the factor. These clauses
supersede their open status above. Producing the compactness limit
and initial pointed maps with the required integer regularity,
producing exact splittings, and finishing the edge/slim core and
interpolation geometry remain. The stated $K\geq4$ fiber-consumer
budget applies only to the explicitly supplied $C^{K-1}$ model
functions and $C^{K-2}$ fiber identifications. No unspecified model
construction receives regularity merely from that count.

\paragraph{Revision128 update to this historical audit.}
LFR12--LFR15 now construct a consumer-sufficient finite limit,
its actual buffered pointed embeddings, and complete nonnegative
product factors from the original approximate splittings. This
supersedes the request for initial compactness maps in revision127
FOR THESE CONSUMERS: their model metric order is $m=K-1$ and
the actual map order is $K$. The proof uses only ballwise bounds
through curvature-derivative order $K$ on the original smooth
manifolds. It does not prove the stronger historical $C^K$
metric conclusion or the smooth-metric package LC56. The
geometric source-fiber enclosure, model core/type identification,
and derivative/transversality estimates for edge/slim packets
remain the next work. General line-to-splitting and Tits-cone
existence are not proved by this approximate-product construction.

\paragraph{Revision129 update to this historical audit.}
LFR16--LFR20 supply LC85's local slim geometry from its stated
normalized analytic and bounded-factor splitting inputs. The compact
factor, finite-surface type, derivative estimate, ENTIRE source-fiber
enclosure, compact interpolation trace, smooth bundle and cutoff are
written producers. The same metric comparisons and original coordinates
are retained. The surviving category work concerns the edge disk and
zero-model core/classification consumers; it is not an unspecified
smoothness upgrade for the slim surface. Source scale hypotheses,
accepted PC bindings, edge geometry, Tits-cone and other qualified
Alexandrov inputs, and global compatibility remain distinct. No new
Lean declaration or final-goal completion is asserted.

\paragraph{Revision130 update to this historical audit.}
LFR21--LFR24 now supply the edge consumer's finite surface topology,
endpoint disk band, smooth model sublevel, all-direction collar
estimate, both transversality conditions and compact model enclosure.
The topological soul adapter retains the inherited Riemannian comparison
bindings and asserts a homeomorphism only. Smooth disk type is obtained
separately from smooth surface classification with boundary. This does
not supply the zero-model actual normal-flow map, or promote a
finite metric to a smooth metric. Source edge-set geometry, normalized
distance derivatives and ENTIRE source-fiber enclosure remain before
the edge packet is complete; general Alexandrov and Tits inputs remain.
No new Lean or accepted-release binding is asserted.

\paragraph{Revision131 update to this historical audit.}
LFR25--LFR28 now prove the source edge-disk transfer from the stated
coarse-border chart: all-choice nearest directions, scaled smoothing,
the quotient derivative, an extension smooth through the core which
preserves the tested sublevel, entire-source-fiber enclosure, compact
interpolation, the proper smooth disk bundle and cutoff. The endpoint
and angular tolerances precede noncollapse; finite-model comparisons
use the later analytic data. A quantitative angular error suffices;
no exact nearest-set limit or smooth model metric is assumed.
The buffered coarse-border condition must still be produced from the
strong/weak edge predicates, with closure and changing-scale control.
The full rank-two adapted collar and strong-edge coverage remain;
boundary rank two alone does not supply them. The zero-model and
general Alexandrov/Tits obligations and accepted-PC gate are unchanged.

\paragraph{Revision132 update to this historical audit.}
LFR29--LFR32 supply the previously explicit coarse-border hypothesis
from strong/weak edge predicates, including actual coverage witnesses,
closure without continuity, changing normalization and endpoint-length
repair. Source9.2 has $C\geq200\Delta$, whereas retained LC84/FC31 use
$C>200\Delta$; LFR30 gives a fixed-quality-loss interval extension,
and LFR31 handles the loss at recentered scales. The additional scale
constraint depends on earlier weak-edge data, consistently with FC44.
LFR33 supplies ONE distance smoothing and actual disk bundles for a
finite verified edge family. Full rank-two adapted collars and strong
edge density/coverage are still open. Zero-model geometry, general
Alexandrov/Tits inputs and accepted-PC bindings remain separate; no
Lean declaration or final full-goal completion is asserted.

\paragraph{Revision133 update to this historical audit.}
LFR34--LFR38 now supply the full rank-two adapted edge collar of
KL9.20 for the ORIGINAL pair $(f,F/\rho)$. The distance smoothing
is extended below the earlier $3\Delta/4$ cutoff; the reference
plane map and smooth coordinates use explicit distance comparisons.
Long tangential anchors respect LFR19's tested length, and the full
quotient derivative is retained. A radius-$300$ differential buffer
supplies ambient Lipschitz control on the radius-$100$ ball; image
coverage and all minimizing-direction tests are proved in that point's
own scale. The SAME finite-family smoothing supplies the disk bundles
and lower collar. The constant-core function is not used as a collar
coordinate. Two-stratum membership additionally uses LC20; strong-edge
density/coverage still remains. Zero-model geometry, general Alexandrov/
Tits inputs and accepted-PC bindings are unchanged; no new Lean or
final full-goal completion is asserted.

\paragraph{Revision134 update to this historical audit.}
LFR39--LFR44 now supply the strong-edge density and coverage left
open in revision133. The low-dimensional model keeps the original
real coordinate; explicit product-distance estimates compare the
residual factors. The empty-annulus argument uses coarse chains,
retains the original slim quality, and proves the required factor
DIAMETER. Interior model charts give actual rank-two splittings;
near-endpoint lifts use the original sharper splitting and a separate
endpoint map on the SAME recentered factor. Rays have an explicit
finite-interval adapter. The collapsed-model tolerance is fixed before
noncollapse and the strong tolerance; only the final rank-one tolerance
uses the latter. A finite selected family now covers the nonslim
one-stratum and nearby weak-edge region by its ACTUAL disk domains,
with the shared smoothing and full adapted collars of LFR38.
The comparison and AC46 recognition inputs remain explicit; this
step needs no new invocation of general line splitting AC43 or
compatibility AC76. Zero-model geometry, general Alexandrov/Tits
inputs and accepted-PC bindings remain. No final full-goal audit or
Lean completion is claimed.

\paragraph{Revision135 update to this historical audit.}
LFR45--LFR49 now supply the finite soul/normal-flow and zero-core
category bridge left open above. The convex exhaustion and finite
flag give one soul with strict ALL-direction separation. A bounded
finite field, calibrated normal tube and unique level crossings give
the ACTUAL finite normal-flow diffeomorphism. A fiberwise isometric
smooth bundle model and that map then define the zero consumer's
smooth carrier. Its metric has order $K-5\geq5$, not infinite order;
its distances and curvature are unchanged. The same field, flow and
normal-disk bundle are smooth in this carrier. Buffered $C^2$
approximation supplies actual pointed smooth convergence embeddings
with $C^1$ metric convergence and source-ball coverage. LC55 and
LC57--LC61 then give the joint zero packet, on one fixed scale and
common tail, conditional on the SAME supplied cone and inherited
comparison inputs. This does not prove the stronger smooth-metric
LC56 package. Three-model classification, general Alexandrov/Tits
inputs and accepted-PC binding remain. No new Lean invocation or
final full-goal completion is asserted.

\paragraph{Revision136 update to this historical audit.}
LFR50--LFR54 supply the finite three-model classification required by
the zero-core consumer. For compact models, smooth approximating
metrics have one bounded-curvature flow interval; positive-time
all-order estimates and smooth compactness yield an AUXILIARY smooth
nonnegative metric on the same smooth type. This metric is used only
for classification and is discarded before any packet comparison.
The compact classification keeps orientation in the deck action;
parallel forms are chosen on the universal cover. For noncompact
models, orientation determines the SAME soul bundle; explicit smooth
maps identify its twisted cores and a cofinal disk exhaustion counts
its ends. Composing these maps with the existing packet maps gives
the actual radial sublevel types without moving their domains or
altering the finite metric, scale, cone, margins or common tail.
Smooth Ricci-flow, surface and covering foundations remain retained
inputs pending accepted-release binding. The stronger two-ended
same-metric isometry is separate; AC43, AC46 and LC21 and the final
simultaneous parameter/export check remain. No full-goal or Lean
completion is asserted.

\paragraph{Revision137 update to this historical audit.}
LFR55--LFR59 supply LC21 for the actual complete finite Riemannian
models, with no change of metric. Finite first/second variation gives
squared-distance supports, global hinge/four-point comparison and
monotone mixed-radius angles. Compact ray directions give uniformly
convergent pseudodistances on bounded Euclidean cones; their metric
quotient is proper and complete, has dimension at most that of the
model, and carries the full AC82 identity. Long hinges give uniform
whole-ball coverage. Nearest-ray choices produce actual pointed maps
with exact target-ball surjectivity at EVERY sufficiently large real
scale. The compact case uses a point cone. The limit is geodesic and
nonnegatively curved by the written metric limit and four-point
arguments. The threshold is model-dependent; LC23/LC58 still supply
the family-uniform finite window by their existing compactness order.
The original K-1 and consumer K-5 metrics both fit the finite ledger;
the auxiliary compact classification metric is never used here.
This discharges the supplied-cone qualification for these consumers,
not the general Alexandrov cone theorem or intrinsic Tits-boundary
identification. A reported gap in Kasue's compactification is recorded;
the adopted ray proof is independent of those compactification maps.
General local Alexandrov comparison, AC43 splitting, AC46 recognition
and final simultaneous parameter/export binding remain, with retained
foundation/accepted-PC bindings pending. No Lean or full-goal
completion is asserted.

\paragraph{Revision138 update to this historical audit.}
ALS01--ALS05 in Section~\ref{sec:alexandrov-written-line-splitting}
now prove AC43 from GLOBAL nonnegative four-point comparison.
Side comparison and armwise monotonicity follow by direct algebra.
Distances to every line are exactly quadratic; their signed Busemann
coordinate is Lipschitz and affine on every minimizing segment.
Properness gives two rays through each point, and saturation of that
coordinate joins them into an actual line. Quadratic comparison for
all positive and negative times proves uniqueness and the full
cross-distance identity. The zero slice and line translations give
actual inverse onto pointed product isometries aligned with the
specified line; factor geometry and dimension are retained.
This supplies the splitting premise for finite models and collapsed
limits once their written global-comparison hypotheses hold.
No actual-angle convergence, gradient flow, metric smoothing or
one-dimensional recognition is assumed. AC02/AC36's local-to-global
comparison, AC46 recognition and simultaneous local parameter/export
binding remain. Retained foundation and Lean bindings are pending;
this is not the final full-goal audit.

\paragraph{Revision139 update to this historical audit.}
ALR01--ALR05 in Section~\ref{sec:alexandrov-written-one-dimensional-recognition}
now supply AC46's exact metric recognition. A convex cosine-law
calculation weakens global zero-curvature comparison to the minus-one
normalization of the written paired charts. AC23/AC15 exclude rank two
using the original Hausdorff bound. AC27 gives injectivity of one
distance coordinate at EVERY segment interior; the original segment
then identifies the ACTUAL ambient ball isometrically with an interval.
Compact last-intersection arguments prove no exit before an endpoint
and produce the remaining half-interval charts, retaining boundary
points. No tangent-space dimension theorem, actual-angle convergence,
gradient flow or uniform recognition radius is assumed. AC47/AC48
and LFR43--LFR44 can use this written producer with their other inputs
unchanged. General local-to-global comparison and final simultaneous
local parameter/export binding remain. Accepted foundation/PC/Lean
bindings remain pending; this is not the final full-goal audit.

\paragraph{Revision140 update to this historical audit.}
ALG01--ALG08 in Section~\ref{sec:alexandrov-written-buffered-globalization}
now prove complete geodesic globalization for nonpositive lower bounds
and a local theorem with an explicit complete-ball factor of $20$.
They construct local germ angles directly, expand the cradle's
nondegeneracy estimate and localize almost-minimum selection and
geodesic existence to the supplied complete buffer. ALG06 supplies
ambient comparison with a $256R$ controlled window. ALG07 uses
$r_i\to\infty$ to obtain this larger window and the ORIGINAL
MC18/AC41 conclusions and net constants; no global source properness
or new input hypothesis is added. ALG08 separately checks AC65 with
shortening $L_i/1024$, LC12 with a later controlled-radius tail, and
LC78 with a smaller existential splitting threshold. The old sharp
AC02/AC36/AC64 margins remain source-qualified; the general complete
space theorem is not applied to an incomplete open set.
Remaining local-consumer buffers and simultaneous local parameter/export
binding are next. No metric replacement, positive-curvature theorem,
ultrapower or new PC foundation is used. Accepted foundation/PC/Lean
bindings remain pending; this is not the final full-goal audit.

\paragraph{Revision141 update to this historical audit.}
LCP01--LCP06 in Section~\ref{sec:collapse-written-comparison-consumers}
now supply the remaining interior-packet comparison uses. LCP01 first
shortens both arms to $2r/5$ and keeps LC27's ORIGINAL $20b$ buffer
and tolerances. LCP02 first shortens to $\min\{D,1\}$ and retains
LC65's exact scale, thresholds and $400$-ball. LCP03 keeps the ORIGINAL
radial center and all-prefix error proportional to prefix length, with
no rescaled $\lambda\delta$ term. LCP04 binds the zero consumers and
makes an allowed smaller existential LC76 threshold. LCP05 checks
LFR25's original edge buffer for EVERY nearest point and segment.
LCP06 records the resulting interior comparison routes. Original sharp
AC02/AC36/AC64 contracts remain separate; none is declared proved.
Simultaneous local parameter assignment and packet export still need
their composition proof. Retained foundation/PC/Lean bindings are
pending; this is not the final full-goal audit or a metric smoothing.

\paragraph{Revision142 update to this historical audit.}
LPA01--LPA06 in Section~\ref{sec:collapse-simultaneous-local-production}
now supply common modified-scale noncollapse and finite derivative data
at EVERY center, the arbitrary-sequence finite compactness premise,
and the joint bounded zero-scale producer. The tested circle-factor
estimate is uniform before noncollapse. An explicit six-stage parameter
choice produces circle, edge, slim and zero packets on one late tail;
endpoint tolerances precede $w'$, while strong splitting follows it.
Finite selections retain one edge distance smoothing, all actual source
fibers, the original zero functions and the same model/cone witnesses.
A whole-ball packing calculation gives support multiplicity independent
of $\Delta$, consistent with its earlier choice in FC44. These are the
local production fields of LC87(1)--(3) in the CLOSED standing sequence.
Chart-overlap adjustment, boundary extension and flow application remain
separate. The original requested consumers and retained foundations must
now undergo full-goal review; this update is not that final audit or Lean
completion. No model metric or already selected function is replaced.

\paragraph{Revision143 final audit disposition.}
Section~\ref{sec:collapse-finite-final-consumer-audit} records the final
consumer review of the task begun at revision125.
LFR01--LFR59, ALS01--ALS05, ALR01--ALR05, ALG01--ALG08,
LCP01--LCP06 and LPA01--LPA06 supply its named local consumers in
written mathematics under the explicitly retained foundations. The
finite-category route retains the metric and every needed derivative,
whole-source-fiber, same-witness and uniformity condition. Stronger
historical contracts are not promoted. This supersedes the earlier
``final local-goal audit remains'' notices, but does not close the
accepted-PC/Lean gate, LC87's other-chart comparison field, LC88's
boundary augmentation, global assembly or flow/reconstruction.

\paragraph{Corrections made during this audit.}
A smaller numerical KL error does not justify bare restriction of
coverage witnesses to a smaller tested source ball. In LC57 choose
$\tau<\delta/4$; in LC73 choose $\nu<\beta_1/4$.
For a target of radius less than $\delta^{-1}-\delta$, the
$\tau$-map supplies a witness of error less than $2\tau$ and
source radius less than $\delta^{-1}-\delta+3\tau<\delta^{-1}$.
Thus the SAME map has the required weaker quality. LC71's
all-smaller-parameter assertion is now supplied directly by LC79.
These repair proof bookkeeping, with no change to the theorem outputs.
LC76 and LC78 likewise use error $2\delta$ or $2\nu$ for lifts,
since the coverage infimum need not be attained.

The new local packets also retain the following source repairs:
LC83 controls only the tested part of a residual factor;
LC84 replaces a continuity argument for a GH map by a coarse chain,
uses closure for edge distance, retains the derivative of the scale
and rescales the surface model before applying KL 3.12;
LC82 requires one common compact isolating region for a perturbation;
and LC88 excludes zero points throughout the selected ball without
assuming its center is zero rank. These are project explanations and
repairs, not additional items in the author's correction sheet.

\paragraph{Relations with other chapters.}
\begin{itemize}[leftmargin=*]
\item Chapters 2--4 supply the inherited Riemannian/analytic layer,
  pointed maps, packing, comparison, splitting, strainers and cones.
  The metric compactness route does not depend on LC81's finite-order
  smooth producer. Nothing in Chapters 3--4 is replaced by a new library.
\item Chapter 14 consumes LC87's local data and proves compatibility,
  overlap adjustment, corners and global graph assembly. Its static
  theorem is then an input ONLY to LC90's flow application. At the
  level of proof stages this is acyclic: local data, global static
  assembly, flow application. No static lemma depends on that application.
\item Chapter 12 supplies actual late normalized carriers and analytic
  hypotheses. LC90 returns their raw thin graph data; it does not prove
  the Chapter 12 branch producer or ambient cusp incompressibility.
\item Chapter 6 converts raw torus cuts into the required relative graph
  output; Chapters 15--16 retain actual transport maps and assemble the
  late-or-terminal alternative. Neither compressible auxiliary cuts nor
  a static graph certificate are silently upgraded to the endpoint.
\end{itemize}

The complete contract-by-contract record is
\code{chapter13_audit_revision96.csv}, with the mathematical/source
findings and stage dependencies in \code{AUDIT_CHAPTER13_REVISION96.md}.
The independent dependency and specification checks test consistency,
not geometric truth. Historical PC evidence remains pinned to its
recorded audit commit; no fresh release acceptance, PC mutation,
Lean invocation or production mathlib pin accompanies this revision.
